\documentclass[11pt,a4paper]{article}
\usepackage{url}
\usepackage[margin=2.35cm]{geometry}
\usepackage{fontspec}
\usepackage{amsmath,amssymb,mathrsfs,mathtools,array,booktabs,pdflscape,adjustbox,diagbox}
\setmainfont{Malgun Gothic}[
  ItalicFont={Malgun Gothic},
  ItalicFeatures={FakeSlant=0.2},
  BoldItalicFont={Malgun Gothic Bold},
  BoldItalicFeatures={FakeSlant=0.2}
]
\setsansfont{Malgun Gothic}
\newfontfamily\latinfont{Times New Roman}
\XeTeXlinebreaklocale "ko"
\XeTeXlinebreakskip=0pt plus 1pt
\setlength{\parindent}{1.2em}
\setlength{\parskip}{0.45em}
\setlength{\emergencystretch}{3em}
\providecommand{\srcspaced}[1]{\textbf{#1}}
\providecommand{\srcfn}[2]{\begingroup\renewcommand{\thefootnote}{#1}\footnote{#2}\addtocounter{footnote}{-1}\endgroup}
\providecommand{\srcfnmark}[1]{\begingroup\renewcommand{\thefootnote}{#1}\footnotemark\addtocounter{footnote}{-1}\endgroup}
\providecommand{\srcfntext}[2]{\begingroup\renewcommand{\thefootnote}{#1}\footnotetext{#2}\endgroup}
\providecommand{\srcnumdisplay}[3][3.4em]{\par\smallskip\noindent\begin{minipage}[t]{#1}\vspace*{0pt}#2\end{minipage}\begin{minipage}[t]{\dimexpr\linewidth-#1\relax}\vspace*{0pt}\centering\(\displaystyle #3\)\end{minipage}\par\smallskip}
\providecommand{\srcdisplay}[2]{\par\smallskip\noindent\makebox[3.2em][l]{#1}\makebox[\dimexpr\linewidth-3.2em\relax][c]{$\displaystyle #2$}\par\smallskip}
\providecommand{\NoetherSrcNote}[2]{\textsuperscript{#1}\begingroup\renewcommand{\thefootnote}{}\footnotetext{\textsuperscript{#1}~#2}\endgroup}
\providecommand{\pair}[2]{(#1\,|\,#2)}
\providecommand{\eps}{\varepsilon}
\providecommand{\Hgrp}{\mathfrak H}
\providecommand{\Gg}{\mathfrak G}
\providecommand{\Hh}{\mathfrak H}
\providecommand{\oo}{\mathfrak o}
\providecommand{\Sp}{\operatorname{Sp}}
\providecommand{\glqq}{“}
\providecommand{\grqq}{”}
\providecommand{\mat}[1]{\begin{pmatrix}#1\end{pmatrix}}
\providecommand{\mA}{\mathfrak A}
\providecommand{\mB}{\mathfrak B}
\providecommand{\mK}{\mathfrak K}
\providecommand{\mL}{\mathfrak L}
\providecommand{\mR}{\mathfrak R}
\providecommand{\mS}{\mathfrak S}
\providecommand{\mT}{\mathfrak T}
\providecommand{\ma}{\mathfrak a}
\providecommand{\mb}{\mathfrak b}
\providecommand{\mo}{\mathfrak o}
\providecommand{\mJ}{\mathfrak J}
\providecommand{\D}{\Delta}
\providecommand{\A}{\Delta}
\providecommand{\R}{\varrho}
\providecommand{\mcell}[1]{\ensuremath{\begin{gathered}#1\end{gathered}}}
\providecommand{\bdelta}{\bar{\delta}}
\providecommand{\frakO}{\mathfrak O}
\providecommand{\frako}{\mathfrak o}
\providecommand{\frakp}{\mathfrak p}
\providecommand{\Afield}{\mathfrak A}
\providecommand{\Bfield}{\mathfrak B}
\providecommand{\Cfield}{\mathfrak C}
\providecommand{\Jdom}{\mathfrak J}
\providecommand{\Kfield}{\mathfrak K}
\providecommand{\Lsys}{\mathfrak L}
\providecommand{\Gdom}{\mathfrak G}
\providecommand{\Ssys}{\mathfrak S}
\providecommand{\Mfield}{\mathfrak M}
\providecommand{\Rfield}{\mathfrak R}
\providecommand{\Xreal}{\mathfrak X}
\providecommand{\Yreal}{\mathfrak Y}
\providecommand{\rcdots}{\,\cdots\,}
\providecommand{\qtext}[1]{\glqq #1\grqq{}}
\providecommand{\p}{\partial}
\providecommand{\dd}{\mathop{}\!\mathrm{d}}
\providecommand{\modu}[1]{\;\operatorname{mod} #1}
\providecommand{\ideal}[1]{\mathfrak{#1}}
\providecommand{\eqzero}[1]{\equiv 0\,(#1)}

\RequirePackage{tikz}
\newcommand{\NoetherRosetteGlyph}{%
  \begin{tikzpicture}[x=0.1363584905660377pt,y=-0.1363584905660377pt,baseline=(current bounding box.south)]
    \path[use as bounding box] (0,0) rectangle (108,108);
    \fill[even odd rule]
    (33.5,84) -- (32.5,83) -- (30.5,83) -- (29,81.5) -- (29,76.5) -- (30.5,75) -- (31.5,75) -- (32.5,74) -- (34.5,74) -- (36,72.5) -- (36,71.5) -- (35.5,71) -- (33.5,71) -- (33,70.5) -- (33,69.5) -- (30.5,67) -- (28.5,69) -- (26.5,69) -- (24.5,71) -- (23.5,71) -- (22.5,72) -- (21.5,72) -- (20.5,73) -- (18.5,73) -- (17.5,72) -- (16.5,73) -- (16,72.5) -- (16,68.5) -- (18.5,66) -- (19.5,66) -- (20,65.5) -- (20,64.5) -- (26,58.5) -- (26,56.5) -- (24,54.5) -- (24,53.5) -- (17,46.5) -- (17,45.5) -- (16.5,45) -- (15.5,45) -- (15,44.5) -- (15,42.5) -- (15.5,42) -- (17.5,42) -- (19.5,44) -- (20.5,44) -- (21,44.5) -- (20,45.5) -- (21,46.5) -- (21,47.5) -- (22,48.5) -- (22,49.5) -- (29,56.5) -- (29,57.5) -- (27,59.5) -- (27,60.5) -- (26,61.5) -- (26,62.5) -- (24,64.5) -- (24,65.5) -- (22,67.5) -- (22.5,68) -- (24.5,68) -- (25.5,67) -- (28.5,67) -- (29.5,66) -- (31.5,66) -- (32.5,65) -- (33.5,65) -- (34.5,64) -- (36.5,64) -- (37.5,63) -- (38.5,63) -- (39.5,62) -- (41.5,62) -- (42.5,61) -- (44.5,61) -- (45.5,60) -- (46.5,60) -- (47.5,59) -- (48.5,59) -- (49,58.5) -- (49,57.5) -- (49.5,57) -- (51.5,57) -- (52,56.5) -- (50.5,55) -- (49.5,55) -- (48.5,54) -- (47.5,54) -- (46.5,53) -- (45.5,53) -- (43,50.5) -- (43,49.5) -- (42,48.5) -- (42,47.5) -- (41.5,47) -- (40.5,47) -- (39.5,48) -- (33.5,48) -- (32,46.5) -- (32,44.5) -- (34,42.5) -- (34,40.5) -- (33.5,40) -- (32.5,40) -- (31.5,39) -- (26.5,39) -- (26,38.5) -- (26,35.5) -- (27.5,34) -- (31.5,34) -- (33,35.5) -- (33,36.5) -- (34.5,38) -- (37.5,38) -- (38,37.5) -- (38,36.5) -- (41,33.5) -- (41,31.5) -- (40,30.5) -- (40,27.5) -- (39,26.5) -- (39,25.5) -- (38,24.5) -- (38,23.5) -- (37,22.5) -- (37,19.5) -- (37.5,19) -- (38.5,20) -- (39.5,20) -- (41,21.5) -- (41,24.5) -- (42,25.5) -- (42,27.5) -- (43,28.5) -- (43,30.5) -- (44,31.5) -- (44,35.5) -- (45,36.5) -- (45,37.5) -- (45.5,38) -- (46.5,38) -- (47,38.5) -- (47,40.5) -- (48,41.5) -- (48,42.5) -- (49,43.5) -- (49,44.5) -- (50,45.5) -- (50,48.5) -- (53,51.5) -- (52,52.5) -- (52,55.5) -- (53.5,57) -- (56.5,57) -- (58.5,59) -- (59.5,59) -- (60.5,60) -- (61.5,60) -- (62,60.5) -- (62,61.5) -- (63,62.5) -- (63,64.5) -- (63.5,65) -- (64.5,65) -- (65.5,64) -- (68.5,64) -- (69.5,65) -- (73.5,65) -- (74,65.5) -- (74,66.5) -- (75,67.5) -- (75,68.5) -- (75.5,69) -- (76,68.5) -- (76,66.5) -- (76.5,66) -- (78.5,68) -- (80.5,68) -- (81.5,69) -- (84.5,69) -- (85,68.5) -- (84,67.5) -- (84,65.5) -- (83.5,65) -- (82.5,65) -- (82,64.5) -- (82.5,64) -- (84.5,64) -- (86,65.5) -- (86,66.5) -- (88.5,69) -- (89.5,69) -- (90,69.5) -- (90,70.5) -- (91,71.5) -- (91,72.5) -- (93,74.5) -- (91.5,76) -- (90.5,76) -- (88.5,74) -- (87.5,74) -- (86.5,73) -- (85.5,73) -- (83.5,71) -- (82.5,71) -- (81.5,70) -- (80.5,70) -- (79.5,69) -- (77.5,69) -- (74.5,72) -- (73.5,72) -- (73,72.5) -- (73,74.5) -- (73.5,75) -- (78.5,75) -- (79.5,74) -- (80.5,74) -- (81,74.5) -- (81,78.5) -- (80,79.5) -- (80,80.5) -- (78.5,82) -- (73.5,82) -- (73,81.5) -- (73,79.5) -- (72,78.5) -- (72,77.5) -- (71.5,77) -- (70.5,77) -- (69.5,76) -- (68.5,76) -- (68,76.5) -- (68,77.5) -- (67,78.5) -- (67,79.5) -- (66.5,80) -- (65.5,80) -- (65,80.5) -- (66,81.5) -- (64.5,83) -- (63,81.5) -- (63,79.5) -- (62,78.5) -- (62,77.5) -- (61,76.5) -- (61,74.5) -- (60,73.5) -- (60,72.5) -- (59,71.5) -- (59,68.5) -- (58,67.5) -- (58,64.5) -- (56,62.5) -- (56,61.5) -- (55.5,61) -- (54.5,62) -- (53.5,62) -- (52,63.5) -- (52,64.5) -- (51,65.5) -- (51,66.5) -- (50,67.5) -- (50,68.5) -- (49,69.5) -- (49,70.5) -- (48,71.5) -- (48,75.5) -- (46,77.5) -- (46,79.5) -- (44.5,81) -- (42.5,81) -- (42,80.5) -- (44,78.5) -- (44,77.5) -- (43.5,77) -- (41.5,77) -- (39.5,75) -- (38.5,76) -- (37.5,76) -- (37,76.5) -- (37,78.5) -- (36,79.5) -- (38,81.5) -- (38,82.5) -- (37.5,83) -- (34.5,83) -- cycle
    (65.5,22) -- (63.5,22) -- (63,21.5) -- (63.5,21) -- (65.5,21) -- (66,21.5) -- cycle
    (55.5,28) -- (53.5,28) -- (53,27.5) -- (53.5,27) -- (54.5,27) -- (57.5,24) -- (59.5,24) -- (60,24.5) -- (57.5,27) -- (56.5,27) -- cycle
    (64.5,26) -- (64,25.5) -- (64.5,25) -- (65,25.5) -- cycle
    (63.5,29) -- (63,28.5) -- (63,27.5) -- (63.5,27) -- (64.5,27) -- (65,27.5) -- cycle
    (58.5,55) -- (57,53.5) -- (57,52.5) -- (56,51.5) -- (56,49.5) -- (57,48.5) -- (57,46.5) -- (60,43.5) -- (60,42.5) -- (59.5,42) -- (58.5,42) -- (58,41.5) -- (59,40.5) -- (59,39.5) -- (61.5,37) -- (62.5,37) -- (63.5,36) -- (65.5,36) -- (66.5,37) -- (68.5,37) -- (70,35.5) -- (69,34.5) -- (69,32.5) -- (68,31.5) -- (68.5,31) -- (74.5,31) -- (75,31.5) -- (75,32.5) -- (76,33.5) -- (76,35.5) -- (75,36.5) -- (75,37.5) -- (72,40.5) -- (72,41.5) -- (71,42.5) -- (72,43.5) -- (72,44.5) -- (72.5,45) -- (73.5,45) -- (74.5,46) -- (75.5,46) -- (76.5,45) -- (78.5,45) -- (79.5,44) -- (80,44.5) -- (80,47.5) -- (79.5,48) -- (77.5,48) -- (76.5,49) -- (74.5,49) -- (73.5,50) -- (71.5,50) -- (70.5,51) -- (68.5,51) -- (67.5,52) -- (65.5,52) -- (64.5,53) -- (63.5,53) -- (62.5,54) -- (59.5,54) -- cycle
    (86.5,50) -- (85.5,50) -- (85,49.5) -- (86,48.5) -- (86,46.5) -- (85.5,46) -- (82.5,46) -- (81.5,47) -- (81,46.5) -- (82,45.5) -- (82,44.5) -- (82.5,44) -- (84.5,44) -- (86.5,42) -- (88.5,42) -- (89.5,41) -- (92.5,41) -- (93,41.5) -- (93,42.5) -- (94,43.5) -- (94,44.5) -- (93.5,45) -- (92.5,45) -- (88.5,49) -- (87.5,49) -- cycle
    (44,42.5) -- (43.5,42) -- (43,42.5) -- (43.5,43) -- cycle
    (45,43.5) -- (44.5,43) -- (44,43.5) -- (44.5,44) -- cycle
    (80.5,60) -- (79.5,60) -- (79,59.5) -- (79,56.5) -- (81,54.5) -- (81,53.5) -- (82.5,52) -- (83,52.5) -- (83,53.5) -- (82,54.5) -- (82,55.5) -- (81.5,56) -- (80.5,56) -- (80,56.5) -- (80,57.5) -- (81,58.5) -- (81,59.5) -- cycle
    (55.5,54) -- (55,53.5) -- (55.5,53) -- (56,53.5) -- cycle
    (43,65.5) -- (42.5,65) -- (42,65.5) -- (42.5,66) -- cycle
    (61,67.5) -- (60.5,67) -- (60,67.5) -- (60.5,68) -- cycle
    (47,69.5) -- (47,68.5) -- (46.5,68) -- (46,68.5) -- (46,69.5) -- (46.5,70) -- cycle
    (62,68.5) -- (61.5,68) -- (61,68.5) -- (61.5,69) -- cycle
    (67.5,97) -- (65.5,97) -- (64,95.5) -- (64,94.5) -- (61.5,92) -- (60.5,92) -- (54.5,86) -- (53.5,86) -- (49.5,90) -- (48.5,90) -- (47.5,91) -- (46.5,91) -- (45.5,92) -- (45,91.5) -- (45,90.5) -- (45.5,90) -- (46.5,90) -- (48.5,88) -- (49.5,88) -- (54.5,83) -- (55.5,83) -- (60.5,88) -- (61.5,88) -- (62.5,89) -- (63.5,89) -- (64,89.5) -- (64,90.5) -- (64.5,91) -- (66.5,91) -- (67,90.5) -- (66,89.5) -- (67,88.5) -- (65,86.5) -- (65,85.5) -- (65.5,85) -- (66.5,85) -- (67,85.5) -- (67,86.5) -- (68,87.5) -- (68,90.5) -- (69,91.5) -- (69,94.5) -- (70,95.5) -- (69.5,96) -- (68.5,96) -- cycle
    (40.5,96) -- (36.5,96) -- (36,95.5) -- (37,94.5) -- (37,91.5) -- (40.5,88) -- (41.5,88) -- (42,88.5) -- (41,89.5) -- (41,92.5) -- (42,93.5) -- (42,94.5) -- cycle;
  \end{tikzpicture}%
}
\providecommand{\noethpIIrosette}{\makebox[9.44444pt][c]{\raisebox{-3.85pt}[0pt][0pt]{\NoetherRosetteGlyph}}}
\providecommand{\NoetherRosetteInline}{\makebox[17.72671698pt][c]{\raisebox{-3.85pt}[10.87671698pt][3.85pt]{\NoetherRosetteGlyph}}}
\usepackage[unicode,hidelinks]{hyperref}
\renewcommand{\contentsname}{목차}
\hypersetup{pdftitle={에미 뇌터 — 수학 저작집 한국어판},pdfauthor={Emmy Noether},pdfsubject={논문·강의·보충문·참고문헌의 한국어 번역}}
\begin{document}
\begin{center}
{\LARGE\bfseries 에미 뇌터 — 수학 저작집}\par
\bigskip
{\Large 한국어판 — 교정 본문과 편집 주석}\par
\bigskip
논문 1–43, 1929/30년 겨울학기 강의,\par
Kapferer와 Noether의 논문 및 보충문, 참고문헌과 권말 목록\par
\end{center}
\vfill
\begingroup\small
\noindent\textbf{저자와 번역}\par
원저자는 Emmy Noether 및 각 저작에 명시된 공동 저자이다.
1929/30년 강의는 M. Deuring이 기록하였다.
한국어 본문은 AI가 번역하였다. 초기 한국어판의 제작 주체는 Manuscript Typesetting Project / OpenAI Codex이며, 당시의 정확한 모델은 현재 남아 있는 출처 자료로 확정되지 않는다.
이 판본의 AI 제작 수정과 조판: OpenAI Codex — GPT-6 Astra, Ultra 추론 설정; Claude Opus 5.5(Anthropic), 2026년 10월 26일.
발행자가 직접 번역하거나 독립적인 학술 심사를 수행했다는 뜻은 아니다.
\par\medskip
\noindent\textbf{출처}\par
한국어판 v1.0.0 (2026-08-13):
\url{https://doi.org/10.5281/zenodo.21921524}.
그 판본의 독일어 저본은 ED0011/v047 (2026-08-08)이다.
각 저작의 원래 출판 정보, 저자 표시 및 인용은 본문에 수록한다.
이 DOI는 이전 한국어판을 식별하며, 현재 GitHub판에 부여된 DOI가 아니다.
\par\medskip
\noindent\textbf{권리}\par
원저작과 보호되는 파생 기여의 권리는 각 권리자에게 있다.
이 판본은 권리자들이 반대하지 않을 것이라는 발행자의 추정에 따라 비상업적이고 순수한 교육 목적으로 무료 제공된다.
이는 추정이며 확인된 허락이 아니다. 무료 열람 가능성 자체가 파생판의 공개를 허용하지는 않는다.
권리자는 \url{https://github.com/KokunoYumeto/emmy-noether-ko/issues/new}에서 이의를 제기할 수 있으며, 발행자는 권리자의 요청에 따라 해당 자료를 내린다.
이 판본은 새로운 포괄적 이용허락이나 공공영역 귀속을 부여하지 않는다.
이전 판본의 실제 권리 표시는 그 판본과 실제로 허락된 권리의 범위에서 존속한다.
\endgroup
\clearpage
\pdfbookmark[0]{목차}{contents}
\tableofcontents
\clearpage
\setcounter{footnote}{0}
\phantomsection
\addcontentsline{toc}{section}{1. 삼원 4차 형식의 형식계 구성에 관하여}
%
%

%
%
\begin{center}
{\Large\bfseries 1. 삼원 4차 형식의 형식계 구성에 관하여.}\par
\vspace{1.2em}
에미 뇌터(Emmy Noether).\par
\vspace{0.8em}
(저자의 학위논문에서 발췌.)\par
\vspace{1.2em}
Sitz. Ber. d. Physikal.-mediz. Sozietät in Erlangen 39 (1907), 176--179쪽
\end{center}

삼원 4차 형식의 형식계는 고르단(Gordan), 마이사노(Maisano), 파스칼(Pascal)의 논문들\srcfn{1)}{\srcspaced{P. Gordan}, Über das volle Formensystem der ternären biquadratischen Form $f=x_1^3x_2+x_2^3x_3+x_3^3x_1$ (Math. Annalen Bd. XVII, S. 217 bis 233. 1880); \srcspaced{G. Maisano}, 1. Sistemi completi dei primi cinque gradi della forma ternaria biquadratica e degl' invarianti, covarianti e contravarianti di sesto grado (Giorn. di Battaglini XIX). 2. Sui covarianti indipendenti di $6^\circ$ grado nei coefficienti della forma biquadratica ternaria (Rend. Circ. Mat. di Palermo I. 1887); \srcspaced{E. Pascal}, Contributo alla teoria della forma ternaria biquadratica e delle sue varie decomposizioni in fattori (Memoria premiata dalla R. Accademia delle scienze fisiche e matematiche di Napoli. 1905).}에서 다루었다. 고르단은 특수한 자기동형 형식
$f=x_1^3x_2+x_2^3x_3+x_3^3x_1$의 54개 구성식으로 이루어진 완전한 형식계를 세웠는데, 그 바탕에는 그가 이원 형식론의 형식계에 사용한 것과 비슷한 원리들이 놓여 있다.

마이사노는 일반 4차 형식에 대하여 계수차수 5까지의 형식들\srcfn{2)}{여기서 ``계수차수(Ordnung)''는 계수에 관한 차수를, ``변수차수(Grad)''는 변수에 관한 차수를 뜻한다.}과 더 높은 계수차수의 몇몇 불변식, 공변식, 반변식을 제시했다. 그가 사용한 방법은 고르단이 \emph{Mathematische Annalen} 제1권에서 삼원 3차 형식에 적용한 방법이다. 파스칼은 마이사노의 결과를 이용하면서 주로 4차 형식이 인수들로 분해되는 문제를 다룬다\srcfn{1)}{두 번째 짧은 논문에서 마이사노는 계수차수 6, 변수차수 6인 세 공변식의 선형종속성을 증명하려 한다. 완전히 갖추어지지 않은 이 증명과 달리 파스칼은 계수차수 6인 세 공변식의 선형독립성과 계수차수 9인 세 불변식의 선형독립성을 증명했다고 본다. 그러나 실제로는 한편의 세 공변식 사이와 다른 한편의 세 불변식 사이에 각각 선형관계가 존재한다는 것이 나의 계산에서 드러났다. 파스칼의 표기법으로 쓴 명시적 공식은 다음과 같다.
\[
0=20\Omega_1+6\Omega_2-3\Omega_3+4f\cdot C_1-12f\cdot C_2-4A\cdot\Delta.
\]
\[
0=4A^3-15A\cdot B+30C-90D+9E.
\]}

\emph{내 연구의 목표는 일반 삼원 4차 형식의 형식계를 세우는 것이다. 여기서는 우선 그 주요 토대만을 제시하고, 이른바 ``상대적으로 완전한 계''를 세운다.}\srcfn{2)}{명칭에 관해서는 \srcspaced{Gordan-Kerschensteiner}, Vorlesungen über Invariantentheorie, Bd. II, S. 227을 보라.}

삼원 형식의 계를 구성하는 기본 발상은 이원 형식론에서와 같다. 첫 번째 상대적으로 완전한 계, 곧 이원 형식론에서 가져온 형식들의 계에서 출발하여 일정한 법칙에 따라 모듈이 갈수록 높은 계들로 나아간다. 이 과정은 어떤 모듈의 계가---그 모듈을 기본형식으로 삼았을 때---유한하고 알려진 것이 되거나, 또는 어떤 모듈이 불변식을 인수로 갖는 형식들로 환원될 수 있을 때까지 계속된다. 첫째 경우에는 상대적으로 완전한 계와 그 모듈의 계에 전이연산(Überschiebung, transvection)을 시행하여 절대적으로 완전한 계를 얻는다. 둘째 경우에는 상대적으로 완전한 계와 절대적으로 완전한 계가 일치한다. 힐베르트(Hilbert)가 형식계의 유한성을 일반적으로 증명했으므로, 이 절차는 반드시 끝나야 한다.

%
%
우리의 경우에는 첫 번째 상대적으로 완전한 계의 모듈 $(abc)$를 다음 모듈들로 환원할 수 있다.
\[
\Delta=(abc)^2a_x^2b_x^2c_x^2\quad\text{및}\quad \nu=(abu)^4.
\]
이로부터 mod $\nu$에 관한 상대적으로 완전한 계를 쉽게 얻는다. 이제 모듈열로 다음 형식들을 택한다.
\[
\nu;\qquad \nu^{(\nu)}=(\nu\nu_1x)^4=s_x^4,
\]
\[
\nu(s)=(ss'u)^4\ldots
\]
그리고 mod $s$에 관한 상대적으로 완전한 계를 구성할 때 나타나는 두 이차형식 $u_\varrho^2$와 $t_x^2$를 추가 모듈로 덧붙인다. 그러면 $s$ 다음에 오는 모듈 $(ss'u)^4$가 모듈 $(\varrho,t)$로 환원 가능함을 보일 수 있다. 그런데 두 이차형식의 동시 형식계는 유한하고 알려져 있으므로, 이로써 앞에서 말한 종결에 도달한다. mod $(\varrho,t)$에 관한 상대적으로 완전한 계로는 331개의 구성식이 얻어진다. ---

이제 사용한 방법에 관하여 몇 가지를 더 적어 둔다.

형식을 생성하는 기본 과정인 \emph{수축 과정(Faltungsprozeß)}은 삼원 형식론에서 다음과 같이 정의할 수 있다.

기호적 곱
\[
s_x^m t_x^n u_\sigma^\mu u_\tau^\nu
\]
에서 인수쌍들을
\[
\begin{array}{@{}c@{\qquad}c@{\qquad}c@{\qquad}c@{\qquad}c@{}}
 & s_xt_x & u_\sigma u_\tau & s_xu_\sigma\ \text{또는}\ s_xu_\tau & t_xu_\tau\ \text{또는}\ t_xu_\sigma\\[0.45em]
\text{각각 다음으로} & (stu) & (\sigma\tau x) & s_\sigma\ \text{또는}\ s_\tau & t_\tau\ \text{또는}\ t_\sigma\\[0.45em]
\text{수축} & \mathrm{I} & \mathrm{II} & \mathrm{III} & \mathrm{IV}
\end{array}
\]
바꾸면, 그 결과로 생기는 형식들은 원래 형식에서 수축을 통해 얻어진 것이다\srcfn{1)}{\srcspaced{Gordan}. Math. Annalen, Bd. XVII, S. 219.}.

이때 언제나 $s_\sigma=0;\ t_\tau=0$으로 둘 수도 있다. 각 수축 사이의 관계에 대해서는 다음 정리가 성립한다.

\emph{수축 I과 II는 기본 수축이다. 수축 III과 IV는 합성 순서와 무관하게 이 기본 수축들로 합성할 수 있다. 다시 말해 주어진 식에서 수축으로 생기는 모든 형식을 만들려면 수축 I과 II만 적용하면 된다.}\srcfn{2)}{$n$과 $\mu$, 또는 각각 $m$과 $\nu$가 동시에 0이 되는 경우에는 이 정리에 예외가 생긴다. 그때 ``형식열''은 대각 성분 하나로 줄어든다.}

%
%
``형식열(Formenreihe)''이란 하나의 시작형식과, 그 형식 자체에서 수축을 통해 생겨나는 모든 형식을 함께 이르는 것으로 정의한다. 위 정리에 따르면 형식열 $s_x^m t_x^n u_\sigma^\mu u_\tau^\nu$를 직사각형 도식으로 배열할 수 있다. 도식에서 나아갈 때에는 다음을 얻는다.
\begin{enumerate}
\item 오른쪽으로 한 열 이동하면, 바로 왼쪽의 형식들에서 수축 I로 생기는 모든 형식,
\item 아래쪽으로 한 행 이동하면, 바로 위의 형식들에서 수축 II로 생기는 모든 형식, 따라서 수축으로 생기는 모든 형식.
\end{enumerate}

이 전제 아래에서 환원자에 관한 정리들은 가장 일반적인 형태로 말할 수 있다. 먼저 다음과 같이 정의한다.
\begin{center}
``환원자(Reduzent)란 환원 가능한 형식열이다.''
\end{center}
그러면 다음이 성립한다.

\begin{center}
\emph{어떤 형식열의 시작형식이 환원 가능한 까닭이 그 성분 하나가 환원자와의 수축으로 생겨났기 때문이고, 그 형식열의 끝형식이 바로 그 성분에서 수축으로 생겨났다면, 전체 형식열은 환원 가능하다.}
\end{center}

그 밖의 환원 방법으로는 다음을 들 수 있다.
\begin{enumerate}
\item 이른바 ``이중 환원'', 곧 한 형식을 두 가지 방식으로 환원하여 더 높은 형식들 사이의 관계를 얻는 방법,
\item ``분해되는 형식과의 수축''. 이 방법은 일부는 환원자를 통한 환원의 성격을, 일부는 ``이중 환원''의 성격을 띤다.
\end{enumerate}

\clearpage
\setcounter{footnote}{0}
\phantomsection
\addcontentsline{toc}{section}{2. 삼원 4차 형식의 형식계 구성에 관하여}
%
%

%
%
\begin{center}
{\Large\bfseries 2. 삼원 4차 형식의 형식계 구성에 관하여\srcfn{*)}{서론과 제I장의 짧은 발췌문이 \glqq{}Sitzungsber. der phys.-med. Sozietät Erlangen 1907\grqq{}, 176--179쪽에 발표되었다.}.}\par
\vspace{0.8em}
에를랑겐의 에미 뇌터 양.\par
\vspace{1.2em}
Journal f. d. reine u. angew. Math. 134 (1908), 23--90쪽 및 표 2개
\end{center}

\vspace{3em}

%
%
\begin{center}
\textbf{서론.}
\end{center}

삼원 4차 형식의 형식계는 \emph{Gordan}, \emph{Maisano}, \emph{Pascal}의 연구에서 다루어졌다\srcfn{**)}{\emph{P. Gordan}, Über das volle Formensystem der ternären biquadratischen Form
\[
f=x_1^3x_2+x_2^3x_3+x_3^3x_1.
\]
(Math. Annalen, Bd. XVII (1880), S. 217--233.)
\emph{G. Maisano}, 1) \emph{Sistemi completi dei primi cinque gradi della forma ternaria biquadratica e degl' invarianti, covarianti e contravarianti di sesto grado}. (Giorn. di Battaglini XIX.)
2) \emph{Sui covarianti indipendenti di $6^\circ$ grado nei coefficienti della forma biquadratica ternaria}. (Rend. Circ. Mat. di Palermo I, 1887.)
\emph{E. Pascal}, \emph{Contributo alla teoria della forma ternaria biquadratica e delle sue varie decomposizioni in fattori}. (Memoria premiata dalla R. Accademia delle scienze fisiche e matematiche di Napoli. 1905.)}. \emph{Gordan}은 이항 영역의 형식계에 사용했던 것과 비슷한 원리를 바탕으로, 특수한 자기동형 형식 \(f\equiv x_1^3x_2+x_2^3x_3+x_3^3x_1\)의 54개 구성식으로 이루어진 완전한 형식계를 세웠다.

\emph{Maisano}는 일반 4차 형식에 대해 계수차수 5까지의 형식들\srcfn{***)}{여기서 \glqq{}Ordnung\grqq{}은 계수에 관한 차수를, \glqq{}Grad\grqq{}는 변수에 관한 차수를 뜻한다.}과 더 높은 계수차수의 몇몇 불변식, 공변식, 반공변식을 제시하였다. 그 방법은 \emph{Gordan}이 Math. Annalen 제1권에서 삼원 3차 형식에 적용한 것이다. \emph{Pascal}은 \emph{Maisano}의 결과를 이용하여 주로 4차 형식의 인수분해 문제를 다루었다\srcfn{*)}{두 번째 짧은 논문에서 \emph{Maisano}는 계수차수 6, 변수차수 6인 세 공변식의 선형 종속성을 보이려 하였다. 완전하지 못한 이 증명과 달리 \emph{Pascal}은 계수차수 6인 세 공변식의 선형 독립성과 계수차수 9인 세 불변식의 선형 독립성을 증명했다고 여겼다. 그러나 실제로 한편의 세 공변식 사이에, 다른 한편의 세 불변식 사이에 각각 선형 관계가 존재한다는 사실은 이 논문의 공식 (13), § 11과 공식 (22), § 17 주석에서 확인되며, 그곳에 이 관계들이 명시적으로 주어져 있다. \emph{Pascal}의 통지에 따르면 이 모순은 그의 논문 46쪽에 있는 반공변식 $p$(여기서는 $\varrho$라 부른다)의 식에 생긴 수치 오류로 설명된다.}.

%
%
\emph{이하 연구의 목적은 일반 삼원 4차 형식의 형식계를 세우는 것이다. 이 논문에서는 그 주된 기초만을 제시하고, 이른바 \glqq{}상대 완전계\grqq{}}\srcfn{**)}{명칭에 대해서는 § 3a를 참조하라.}\emph{를 구성한다.}

이 논문은 \emph{Gordan}의 연구와 밀접하게 이어져 있지만, 그곳에서 아주 일반적으로만 제시된 원리들을 먼저 세부적으로 전개해야 했다. 수축들 사이의 관계에 관한 한 정리를 이용하고 \glqq{}형식열\grqq{}(§ 1)을 도입함으로써 환원 정리를 엄밀하게 정식화하고 완성하며(§ 3), 특수한 급수 전개를 재귀적으로 구성하는 일(§ 2)은 환원을 실제로 수행하기 위한 계산 수단을 제공한다.

삼원 형식계를 구성하는 기본 발상은 이항 영역에서와 같다. 첫 번째 상대 완전계, 곧 이항 영역에서 넘겨받은 형식들의 계에서 출발하여 일정한 법칙에 따라 모듈이 점점 더 높은 계들로 나아간다. 이 과정은 어떤 모듈의 계가 --- 그 모듈을 기본형으로 삼았을 때 --- 마침내 유한하고 알려진 것이 될 때까지, 또는 한 모듈을 불변식을 인수로 갖는 형식들로 환원할 수 있을 때까지 계속된다. 첫째 경우에는 상대 완전계를 그 모듈의 계와 전이연산하여 절대 완전계를 얻고, 둘째 경우에는 상대 완전계와 절대 완전계가 일치한다. \emph{Hilbert}가 형식계의 유한성을 일반적으로 증명했으므로, 이 절차는 반드시 유한 단계에서 끝나야 한다.

%
%
우리의 경우 첫 번째 상대 완전계의 모듈 $(abc)$(§ 4)은 다음 모듈들로 환원된다.
\[
\Delta=(abc)^2a_x^2b_x^2c_x^2\quad\text{및}\quad \nu=(abu)^4
\]
(§ 5). 이어서 $\nu$를 법으로 하는 상대 완전계를 구한다(§ 6). 이 결과들은 일반 4차 형식에 대해 이미 \emph{Gordan}의 연구에 제시되어 있지만, 여기의 유도 방식은 더 높은 차수의 형식으로 더 쉽게 옮길 수 있다.

%
%
이제 모듈열로 다음 형식들을 택한다(§ 7).
\[
\nu,\qquad \nu(\nu)=(\nu\nu_1x)^4=s_x^4,\qquad \nu(s)=(ss'u)^4,\ldots .
\]

$s$를 법으로 하는 상대 완전계를 구성할 때, 그 계에 나타나는 두 이차형식 $u_\varrho^2$와 $t_x^2$를 모듈로 첨가한다(§ 9와 10). 따라서 $(s,\varrho,t)$를 법으로 하는 상대 완전계를 구성한다. 이어서 (§ 17) 모듈 $(ss'u)^4$가 모듈 $\varrho$와 $t$로 환원 가능함을 보인다. 이로써 그다음의 더 높은 계는 $(\varrho,t)$를 법으로 하는 상대 완전계로 넘어간다. 그런데 두 이차형식의 연립계는 유한하고 알려져 있으며, 일반적인 경우 20개의 구성식으로 이루어진다\srcfn{*)}{\emph{Clebsch-Lindemann}, Vorlesungen über Geometrie. (Bd. I, S. 288 ff.) 서로 동변인 두 형식의 계. \emph{Gordan}, Über Büschel von Kegelschnitten. (Math. Ann. Bd. XIX. S. 530). 서로 반변인 두 형식의 계.}. 그러므로 $(\varrho,t)$를 법으로 하는 상대 완전계를 세움으로써 위에서 말한 종결점에 이른다. 절대 완전계를 만들기 위하여 이 계를 $(\varrho,t)$의 계와 전이연산하는 일은 이후의 과제로 남겨 둔다.

$(\varrho,t)$를 법으로 하는 상대 완전계로 331개의 구성식을 얻으며, 첨부된 표에서는 이를 변수 $x$와 $u$에 관한 차수에 따라 배열하였다.

%
%
끝으로 도입한 기호들을 한눈에 볼 수 있도록 정리한다.
\begin{align*}
f&=a_x^4,\\
\theta&=\theta_x^4u_\vartheta^2=(abu)^2a_x^2b_x^2,\\
K&=K_x^6u_k^3=(a\theta u)u_\vartheta^2a_x^3\theta_x^3,\\
j&=u_j^6=(a\theta u)^4u_\vartheta^2,\\
\Delta&=\Delta_x^6=a_\vartheta^2a_x^2\theta_x^4=(abc)^2a_x^2b_x^2c_x^2,\\
N&=N_x^8u_n=(a\Delta u)a_x^3\Delta_x^5,\\
\nu&=u_\nu^4=(abu)^4,\\
H&=H_x^2u_\eta^4=(\nu\nu_1x)^2u_\nu^2u_{\nu_1}^2,\\
L&=L_x^3u_l^6=(\nu\eta x)H_x^2u_\nu^3u_\eta^3,\\
g&=g_x^6=(\nu\eta x)^4H_x^2,\\
\sigma&=u_\sigma^6=H_\nu^2u_\nu^2u_\eta^4,\\
s&=s_x^4=(\nu\nu_1x)^4,\\
Z&=Z_x^4u_\zeta^2=(ss'u)^2s_x^2s_x'^2,\\
\varrho&=u_\varrho^2=\theta_\nu^4u_\vartheta^2,\\
t&=t_x^2=a_\eta^4H_x^2,\\
i&=a_\nu^4,\\
J&=s_\nu^4.
\end{align*}

%
%
\begin{center}
{\Large\bfseries 제I장. \quad 삼원 형식에 관한 일반 정리.}
\end{center}

\begin{center}
\textbf{§ 1.}\quad \emph{수축 과정. 형식열.}
\end{center}

형식을 생성하는 기본 과정은 \emph{수축 과정}이며, 삼원 영역에서는 다음과 같이 정의할 수 있다. 기호곱
\[
s_x^m t_x^n u_\sigma^\mu u_\tau^\nu
\]
이 주어졌다고 하자. 인수쌍
\[
s_xt_x\qquad u_\sigma u_\tau\qquad s_xu_\sigma\ \text{또는}\ s_xu_\tau
\qquad t_xu_\tau\ \text{또는}\ t_xu_\sigma
\]
을 각각
\[
(stu)\qquad (\sigma\tau x)\qquad s_\sigma\ \text{또는}\ s_\tau
\qquad t_\tau\ \text{또는}\ t_\sigma
\]
로 바꾸는 것을 수축
\[
\text{I}\qquad\text{II}\qquad\text{III}\qquad\text{IV}
\]
라 하면, 이렇게 생기는 형식들은 원래 형식에서 수축으로 얻어진 것이다\srcfn{*)}{\emph{Gordan}, Math. Annalen Bd. XVII S. 219.}.

%
%
이때 식 $s_x^m t_x^n u_\sigma^\mu u_\tau^\nu$는 두 형식의 실제 곱 $S\cdot T$로 볼 수도 있고, 여러 기호열을 가진 하나의 형식으로 볼 수도 있다.
\[
s_x^m t_x^n u_\sigma^\mu u_\tau^\nu=A_x^{m+n}u_\lambda^{\mu+\nu}.
\]
첫째 경우에는 형식 $S$와 $T$의 수축이라 하고, 둘째 경우에는 \glqq{}형식 $A$의 자기 수축\grqq{}이라 한다. 간결함을 위해 이하에서는 언제나 다음을 가정한다. 이는 실제로 뒤에서 살필 모든 형식에 성립한다.
\srcnumdisplay{(a)}{\begin{gathered}
s_\sigma^\lambda=0,\\
t_\tau^\varkappa=0
\end{gathered}}
여기서 0이 아닌 임의의 $\lambda$와 $\varkappa$를 취하고, 이 식들을 임의의 다른 수축들과 결합한다. 이는 형식 $s_x^m t_y^n u_\sigma^\mu v_\tau^\nu$가 이른바 \glqq{}정규형\grqq{}임을 뜻한다. 즉 변수 $x,u$에 관해서든 변수 $y,v$에 관해서든 $\Omega$-과정을 적용하면 소멸한다.

따라서 조건 (a)는 아무런 제한도 주지 않으며 언제나 달성할 수 있다\srcfn{*)}{\emph{Gordan}, Math. Annalen Bd. V S. 104 참조.}.

개별 수축들 사이의 관계에 대해서는 이제 다음 정리가 성립한다.

\emph{정리 I. 수축 I과 II는 기본 수축이며, 합성 순서와 무관하게 이들로부터 수축 III과 IV를 합성할 수 있다. 다시 말해, 주어진 식에서 수축으로 생기는 모든 형식을 만들기 위해서는 수축 I과 II만 적용하면 된다.}

%
%
\textbf{증명:} 항등식 정리와 행렬의 곱정리를 각각 사용하고 (a)를 고려하면 다음이 따른다.
\[
(\widehat{st}\sigma x)=-t_\sigma,\qquad
(\widehat{\sigma\tau}su)=-s_\tau,
\]
\[
(\widehat{st}\tau x)=s_\tau,\qquad
(\widehat{\sigma\tau}tu)=t_\sigma,
\]
\[
(stu)(\sigma\tau x)=s_\tau+t_\sigma-u_x\cdot s_\tau t_\sigma .
\]
따라서 수축 I과 II를 합성하면 수축 III과 IV가 생긴다. 이는 수축 I, 곧 인수 $(stu)u_\sigma u_\tau$에 수축 II를 적용할 때에도, 수축 II, 곧 인수 $(\sigma\tau x)s_xt_x$에 수축 I을 적용할 때에도 마찬가지이다.

%
%
\emph{형식열}\srcfn{**)}{\emph{Clebsch}, Über eine Fundamentalaufgabe der Invariantentheorie (Abh. der Gött. Ges. d. Wiss. Bd. XVII), § 17 및 § 18 끝 참조. \glqq{}형식열\grqq{}은 더 높은 형식에 따른 배열 원리에 의해 그곳에서 도입한 \glqq{}본래적으로 환원된 동치계\grqq{}와 구별된다.}\setcounter{footnote}{9}을 하나의 시작형식과 이 시작형식의 자기 수축으로 생긴 모든 형식으로 이루어진 것으로 정의하고, 그 시작형식으로 형식열을 나타낸다. 여기서 자기 수축이 더 많은 형식을 \glqq{}더 높은 형식\grqq{}이라 하며, 서로 동등한 수축 $s_\tau$와 $t_\sigma$ 가운데 하나를 더 높은 것으로 정규화해야 한다. 형식열의 구성 법칙에서 다음이 따른다.

1) 형식열의 각 형식은 시작형식의 기호들에 관해 선형이다.

2) 서로 반변인 기호열을 각각 두 개씩 가진 시작형식에서는 형식열이 유일하게 정해진다. 더 많은 기호열을 가진 시작형식에서는 항등식 정리로 연결되는 수축들 가운데 특정 수축들을 우선시함으로써 형식열을 유일하게 정규화해야 한다.

%
%
정리 I에 따라 형식열 $s_x^m t_x^n u_\sigma^\mu u_\tau^\nu$ ($m>n;\ \mu>\nu$)는 다음 직사각형 도식으로 배열할 수 있다.
\begingroup\small
\[
\begin{array}{c|c|c|c|c}
s_x^m t_x^n u_\sigma^\mu u_\tau^\nu
& (stu)
& (stu)^2
& \cdots
& \noethpIIrosette\\[0.3em]
(\sigma\tau x)
& t_\sigma,\ s_\tau
& t_\sigma(stu),\ s_\tau(stu)
& \cdots
& \\[0.3em]
(\sigma\tau x)^2
& t_\sigma(\sigma\tau x),\ s_\tau(\sigma\tau x)
& t_\sigma^2,\ t_\sigma s_\tau,\ s_\tau^2
& \cdots
& \vdots\\[0.3em]
\vdots&\vdots&\vdots&\ddots&\vdots\\[0.3em]
(\sigma\tau x)^\nu
& t_\sigma(\sigma\tau x)^{\nu-1},\ s_\tau(\sigma\tau x)^{\nu-1}
& t_\sigma^2(\sigma\tau x)^{\nu-2},\ t_\sigma s_\tau(\sigma\tau x)^{\nu-2},\ s_\tau^2(\sigma\tau x)^{\nu-2}
& \cdots
& \vdots\\[0.3em]
\vdots
& t_\sigma(\sigma\tau x)^\nu
& t_\sigma^2(\sigma\tau x)^{\nu-1},\ t_\sigma s_\tau(\sigma\tau x)^{\nu-1}
& \cdots
& \vdots\\[0.3em]
\vdots
& \vdots
& t_\sigma^2(\sigma\tau x)^\nu
& \cdots
& \vdots
\end{array}
\]
\endgroup
\begingroup\small
\[
\begin{array}{r@{\quad}c}
\noethpIIrosette\srcfnmark{*)}
&
\begin{array}{c|c|c}
(stu)^n & \cdots & \cdots\\[0.3em]
t_\sigma(stu)^{n-1},\ s_\tau(stu)^{n-1}
& s_\tau(stu)^n & \cdots\\[0.3em]
t_\sigma^2(stu)^{n-2},\ t_\sigma s_\tau(stu)^{n-2},\ s_\tau^2(stu)^{n-2}
& t_\sigma s_\tau(stu)^{n-1},\ s_\tau^2(stu)^{n-1}
& s_\tau^2(stu)^n .
\end{array}
\end{array}
\]
\endgroup
\srcfntext{*)}{여기서 그리고 이하의 비슷한 곳에서, \NoetherRosetteInline{}가 붙은 아래쪽 부분은 도식의 오른쪽 위에 같은 기호로 표시한 자리에 이어 붙여야 한다.}

%
%
이 도식에서

1) 한 열씩 오른쪽으로 나아가면 바로 왼쪽의 형식들에 수축 I을 적용하여 생기는 모든 형식을 얻고,

2) 한 행씩 아래로 나아가면 바로 위의 형식들에 수축 II를 적용하여 생기는 모든 형식을 얻는다.

따라서 수축으로 생기는 모든 형식을 얻는다.

전체 형식열의 임의의 형식 $t_\sigma^\kappa s_\tau^\lambda(stu)^\mu$ 또는 $t_\sigma^\kappa s_\tau^\lambda(\sigma\tau x)^\nu$는 새로운 형식열의 시작을 이룬다. 이 새 형식열은 직사각형 도식에서 해당 형식 $t_\sigma^\kappa s_\tau^\lambda(stu)^\mu$ 또는 $t_\sigma^\kappa s_\tau^\lambda(\sigma\tau x)^\nu$를 왼쪽 위 꼭짓점으로 삼고, 인수 $t_\sigma^\kappa s_\tau^\lambda$를 갖는 모든 형식으로 이루어진다.

%
%
후기. 수 $n$과 $\mu$(각각 $m$과 $\nu$)가 0이 되는 경우에는 정리 I에 예외가 생긴다. 이때 수축 I, II 및 IV는 더 이상 존재하지 않지만 수축 III(각각 IV)는 여전히 존재한다. 이 경우에는 다음 도식의 대각 성분을 형식열이라 한다.
\[
\begin{array}{c@{\qquad}c@{\qquad}c}
\begin{array}{ccccc}
s_x^m u_\tau^\nu &&&&\\
&s_\tau&&&\\
&&s_\tau^2&&\\
&&&\ddots&\\
&&&&s_\tau^\nu
\end{array}
&\text{각각}&
\begin{array}{ccccc}
t_x^n u_\sigma^\mu &&&&\\
&t_\sigma&&&\\
&&t_\sigma^2&&\\
&&&\ddots&\\
&&&&t_\sigma^\mu
\end{array}
\end{array}
\]
수축 I과 II가 존재하지 않는 것과 일치하게, 이 경우 형식열의 편극에 따른 급수 전개는 언제나 시작 성분으로 환원된다. 그러나 환원자에 관한 정리들은 성립한다.

%
%
\begin{center}
\textbf{§ 2.}\quad \emph{형식열의 편극에 따른 급수 전개}\srcfn{*)}{Study, Methoden zur Theorie der ternären Formen (Teubner 1889) II. § 7 참조. 그곳의 유도와 달리 우리의 유도는 수치 계수 $C_{pr\varrho}$를 빠르게 계산하여 정하는 방법을 제공한다. 특수화 a), a')에서 Clebsch--Study 전개는 다음과 같을 것이다.
\[
(stu)^\mu u_\tau^{\nu-\varkappa}v_\tau^\varkappa s_x^m t_x^{\,n-\lambda}(tuv)^\lambda
=\sum_{p,r,\varrho} C'_{pr\varrho}
\bigl[s_\tau^\varrho(stu)^{\mu+r-\varrho}\bigr]_{\widehat{uv}^{\lambda-r+\varrho}v^{\varkappa-\varrho},\,w^p,\,(w=x\widehat{uv})}.
\]
따라서 인수 $u_x$를 가진 모든 항을 생략해야 한다는 사실을 알고 있을 때 우리의 전개에서는 계산 도중에 이미 나타나는 단순화를 그 전개에서는 할 수 없다.}.
\end{center}

$n$개 변수를 가진 형식에 대해서는 두 종류의 급수 전개가 명시적으로 주어져 있다\srcfn{**)}{Gordan, Über Kombinanten, § 2와 5 (Math. Ann. Bd. V) 참조.}.

1) 서로 반변인 변수열을 하나씩 가진 형식 $s_x^m u_\tau^\nu$에 대해서는 $u_x$의 거듭제곱에 따라 진행하는 급수이며, 그 계수들은 \glqq{}정규형\grqq{}이다.

2) 서로 동변인 변수열 두 개를 가진 형식 $s_x^m t_x^n$에 대해서는 기본 공변식의 편극(형식열 $s_x^m t_x^n$의 편극)에 따른 전개이며, 삼원 영역에서는 다음과 같다.
\[
\tag{I}
s_x^m t_x^{n-\lambda}t_y^\lambda
=\sum_{\varrho}
\frac{\binom{m}{\varrho}\binom{\lambda}{\varrho}}
 {\binom{m+n-\varrho+1}{\varrho}}\cdot
\bigl[(st\widehat{xy})^\varrho s_x^{m-\varrho}t_x^{n-\varrho}\bigr]_{y^{\lambda-\varrho}}.
\]

%
%
서로 반변인 변수열을 두 개씩 가진 형식
\[
s_x^m t_x^{n-\lambda}t_y^\lambda u_\sigma^\mu u_\tau^{\nu-\varkappa}v_\tau^\varkappa
\]
에 대해 두 전개를 결합한다. 즉 $u_x$의 거듭제곱에 따른 전개를 세우되, 그 계수는 변수 $x$와 $u$에 관해 취한 형식열 $s_x^m t_x^n u_\sigma^\mu u_\tau^\nu$의 합성 편극이 되게 한다.

서로 반변인 기호열(각각 변수열)을 두 개씩 가진 경우의 급수만 세우면 충분하다. 기호열(변수열)을 두 개씩 하나의 새로운 열로 묶으면, 더 많은 기호열(변수열)을 가진 형식들을 이 경우로 환원할 수 있기 때문이다.

형식열의 모든 형식이 기호열과 변수열을 각각 두 개만 포함하게 하기 위해 다음과 같이 특수화한다.
\[
\begin{array}{ll}
\text{a) } \sigma=\widehat{st}, & \text{a') } y=\widehat{uv},\\[0.3em]
\text{b) } s=\widehat{\sigma\tau}, & \text{b') } v=\widehat{xy}
\end{array}
\]
그리고 특수화 a), a')에 대해 전개를 수행한다.

%
%
전개 I을 직접 옮기면 형식 $(stu)^\mu s_x^m t_x^{n-\lambda}t_y^\lambda$에 대해서는 합성 편극을 얻는다. 그러나 형식 $(stu)^\mu u_\tau^{\nu-\varkappa}v_\tau^\varkappa s_x^m t_x^{n-\lambda}t_y^\lambda$에 대해서는 합성 편극 대신 그러한 편극들의 항들이 생긴다. 이 항들을 계산하려면 계속 새로운 보조변수를 도입하여 마지막에야 서로 같게 놓아야 하므로 계산이 불필요하게 복잡해진다.

따라서 우리는 간접적인 길을 택한다. 형식열의 모든 형식의 합성 편극, 곧 식
\[
\bigl[s_\tau^\varrho(stu)^{\mu-\varrho+r}\bigr]_{y^\lambda}^{v^\varkappa}
\]
을 이 편극들의 시작항들로 나타내고, 그 결과 생기는 재귀 방정식계를 역으로 풀어 원하는 편극에 따른 전개를 얻는다.

%
%
전이 원리에 따라 형식 $s_x^m t_x^n$의 $\lambda$번째 편극 $s_x^m t_x^n:[s_x^m t_x^n]_{y^\lambda}$ ($\lambda\le n$)에는 다음 전개가 성립한다\srcfn{*)}{Gordan, Die Resultante binärer Formen (Kap. I, § 5). Rend. Circ. Mat. di Palermo XXII. 1906.}. $\lambda>n$인 경우에는 $[s_y^m t_y^n]_{x^{m+n-\lambda}}$의 전개를 통하여 첫째 경우로 돌아간다.
\[
\bigl[s_x^m t_x^n\bigr]_{y^\lambda}
=\sum_r (-1)^r\,
\frac{\binom{m}{r}\binom{\lambda}{r}}{\binom{m+n}{r}}\cdot
(st\widehat{xy})^r s_x^{m-r}t_x^{n-\lambda}t_y^{\lambda-r}.
\]
이에 대응하여 $u_\sigma^\mu u_\tau^\nu$의 $\varkappa$번째 편극 $u_\sigma^\mu u_\tau^\nu:[u_\sigma^\mu u_\tau^\nu]^{v^\varkappa}$에는 다음이 성립한다.
\[
\bigl[u_\sigma^\mu u_\tau^\nu\bigr]^{v^\varkappa}
=\sum_\varrho (-1)^\varrho\,
\frac{\binom{\mu}{\varrho}\binom{\varkappa}{\varrho}}{\binom{\mu+\nu}{\varrho}}\,
(\sigma\tau\widehat{uv})^\varrho u_\sigma^{\mu-\varrho}u_\tau^{\nu-\varkappa}v_\tau^{\varkappa-\varrho}.
\]

%
%
두 식을 곱하고 특수화 a)와 a')를 도입하면 다음을 얻는다.
\[
\begin{aligned}
\bigl[s_x^m t_x^n(stu)^\mu u_\tau^\nu\bigr]_{\widehat{uv}^{\lambda}}^{v^\varkappa}
={}&
\sum_{r,\varrho}c_{r\varrho}\cdot
(st\widehat{x\widehat{uv}})^r(\widehat{st}\tau\widehat{uv})^\varrho
s_x^{m-r}t_x^{n-\lambda}(tuv)^{\lambda-r}(stu)^{\mu-\varrho}
u_\tau^{\nu-\varkappa}v_\tau^{\varkappa-\varrho},
\end{aligned}
\]
여기서 $u_x$의 거듭제곱에 따라 전개하면서 첫째 항을 따로 떼어 낸다. $\sum'_{r,\varrho}$는 항 $r=0$, $\varrho=0$을 생략한다는 뜻이다. 또 26쪽의 (a)를 고려하면 다음을 얻는다.
\[
\tag{II}
\begin{aligned}
&s_x^m t_x^{n-\lambda}(tuv)^\lambda(stu)^\mu u_\tau^{\nu-\varkappa}v_\tau^\varkappa\\
&\quad =
\bigl[s_x^m t_x^n(stu)^\mu u_\tau^\nu\bigr]_{\widehat{uv}^{\lambda}}^{v^\varkappa}\\
&\qquad
-\sum_{r,\varrho}' c_{r\varrho}\cdot s_\tau^\varrho(stu)^{\mu-\varrho+r}
u_\tau^{\nu-\varkappa}v_\tau^{\varkappa-\varrho}
s_x^{m-r}t_x^{n-\lambda}(tuv)^{\lambda-r+\varrho}\cdot v_x^r\\
&\qquad
+u_x\cdot\sum_{\varrho}\cdot\sum_{r=1}^{\lambda} c_{r\varrho}\cdot
s_\tau^\varrho(stu)^{\mu-\varrho+r-1}(stv)
u_\tau^{\nu-\varkappa}v_\tau^{\varkappa-\varrho}
s_x^{m-r}t_x^{n-\lambda}(tuv)^{\lambda-r+\varrho}\cdot v_x^{r-1}\\
&\qquad
\vdots\\
&\qquad
-(-1)^\lambda u_x^\lambda\cdot\sum_\varrho c_{\lambda\varrho}\cdot
s_\tau^\varrho(stu)^{\mu-\varrho}(stv)^\lambda
u_\tau^{\nu-\varkappa}v_\tau^{\varkappa-\varrho}
s_x^{m-\lambda}t_x^{n-\lambda}(tuv)^\varrho,
\end{aligned}
\]
여기서
\[
c_{r\varrho}=(-1)^{r+\varrho}\cdot
\frac{\binom{m}{r}\binom{\lambda}{r}}{\binom{m+n}{r}}\cdot
\frac{\binom{\mu}{\varrho}\binom{\varkappa}{\varrho}}{\binom{\mu+\nu}{\varrho}}
\qquad
\begin{array}{l}
\lambda\le n,\\
\varkappa\le\nu
\end{array}
\]
이고, 수 $\lambda,n;\varkappa,\nu$의 나머지 배열들에 대해서도 이에 대응하는 식을 얻는다.

%
%
따라서 식
\[
s_x^m t_x^{n-\lambda}(tuv)^\lambda(stu)^\mu u_\tau^{\nu-\varkappa}v_\tau^\varkappa
\]
은 하나의 합성 편극으로 환원된다. 또한 항등식
\[
(stv)u_\tau=(stu)v_\tau-s_\tau(tuv)
\]
을 적용하면, 형식열의 더 높은 형식들에 대응하는 같은 방식으로 구성된 식들의 합으로 환원된다.

이를 반복하면 다음 전개에 이른다.
\[
\tag{III}
s_x^m t_x^{n-\lambda}(tuv)^\lambda(stu)^\mu u_\tau^{\nu-\varkappa}v_\tau^\varkappa
=
\sum_{p,r,\varrho}u_x^p\cdot C_{pr\varrho}\cdot
\bigl[s_\tau^\varrho(stu)^{\mu-\varrho+r}\bigr]_{\widehat{uv}^{\lambda-r+\varrho}v^{\varkappa-\varrho+p}\cdot v_x^{r-p}}.
\]
수치 계수 $C_{pr\varrho}$는 (II.)를 반복하여 생기는 방정식계의 계수 $c_{r\varrho}$와 $c'_{r\varrho}$로부터 유일하고 재귀적으로 계산된다(42쪽의 예 참조). 나머지 특수화들에서도 이와 같은 전개를 얻는다.

%
%
\begin{center}
\textbf{§ 3.}\quad \emph{환원 정리.}
\end{center}

형식과 형식계의 환원에 관한 잘 알려진 정리들을 이제 § 1과 § 2에 연결하고자 한다. 이를 위해 이항 형식론에서 가져온 몇몇 정의가 필요하다.

a) 주어진 형식열을 법으로 하는 \emph{상대 완전계}란 다음 성질을 지닌 형식계를 말한다. 이 형식들을 임의로 곱한 것을 수축하여 생기는 모든 구성식은 그 계의 형식들로 된 다항식으로 나타낼 수 있으며, 다만 계의 형식들을 법의 계와 수축하여 생긴 식들로 이루어진 가법적 항은 제외한다\srcfn{*)}{Gordan-Kerschensteiner, Vorlesungen über Invariantentheorie. Bd. II. S. 227.}.

b) 어떤 형식이 불변식을 인수로 갖는 형식들 또는 \glqq{}더 높은 형식들\grqq{}로 나타낼 수 있을 때 그 형식을 \emph{환원 가능}하다고 한다. 다시 말해, 해당 형식이 속한 전체 형식열의 더 높은 형식들, 또는 법의 기호를 더 높은 차수로 포함하는 형식들로 나타낼 수 있을 때이다.

이제 환원자에 관한 정리들\srcfn{**)}{Gordan, Math. Annalen Bd. XVII. S. 222.}은 가장 일반적인 형태로 다음과 같이 진술할 수 있다.

\textbf{정의:} \emph{환원자란 환원 가능한 형식열이다.}

%
%
\textbf{정리 II.} \emph{형식열의 시작형식이 그 항들 가운데 하나가 환원자와 수축하여 생겼기 때문에(환원자를 인수로 갖기 때문에) 환원 가능하고, 그 형식열의 끝형식도 바로 이 항에서 수축하여 생겼다면, 전체 형식열은 환원 가능하다}\srcfn{*)}{정리 II가 가져오는 단순화는 다음 예에서 볼 수 있다. 특수 형식 $f=x_1^3x_2+x_2^3x_3+x_3^3x_1$에 대해서는 $a_\nu^2a_x^2u_\nu^2$가 환원자이다. 따라서 정리 II에 의해 형식열
\[
\theta_\nu(\vartheta\nu x)\theta_x^3u_\vartheta u_\nu^2
=
a_\nu^2(abu)-\underbrace{a_\nu b_\nu(abu)}_{0}
\]
의 환원이 따른다. 이는 $\theta_\nu^4=a_\nu^2b_\nu^2(abu)^2$가 항 $a_\nu^2(abu)$에서 수축하여 생겼기 때문이다. 반면 Gordan은 앞서 인용한 문헌 231쪽에서 형식들
\[
\theta_\nu(\vartheta\nu x),\quad \theta_\nu(\vartheta\nu x)^2,\quad
\theta_\nu^2,\quad \theta_\nu^2(\vartheta\nu x),\quad
\theta_\nu^2(\vartheta\nu x)^2
\]
을 하나씩 환원한다.}.

\textbf{증명:} 형식열은 시작항을 그 자체 안에서 수축하여 생긴다. 따라서 한편으로는 환원자와 더 높게 수축하여, 다른 한편으로는 환원자를 그 자체 안에서 수축하여 생긴다. 두 경우 모두 § 2에 따라 형식열의 모든 형식은 환원자의 형식열에 대한 편극(전이연산)으로 나타낼 수 있다.

%
%
다른 환원 방법들에서는 시작형식에서 전체 형식열로 나아가는 추론이 더 이상 성립하지 않는다. 그 방법들은 다음과 같다.

1) 이른바 \emph{\glqq{}이중 환원\grqq{}}.

이는 서로 독립적인 두 환원자를 인수로 갖는 식을 두 가지 방식으로 환원하는 것을 뜻하며, 이로써 환원의 도착점이 된 더 높은 형식들 사이의 관계가 생긴다.

한편 \glqq{}이중 환원\grqq{}을 체계적으로 적용하면 모든 관계를 얻어야 한다\srcfn{**)}{예를 들어 \glqq{}이중 환원\grqq{}으로 다음을 찾았다. 형식 $a_\eta^3$(§ 10)의 환원인데, 이 형식은 Maisano가 불필요하게 계에 포함시킨 유일한 형식이다. 이어서 서론의 주석에서 언급한 3개의 공변식과 3개의 불변식 사이의 관계도 찾았다.}. 다른 한편, 서로 다른 식들에 적용한 \glqq{}이중 환원\grqq{}이 새로운 관계를 내놓지 않을 때에는 더 높은 형식들의 독립성에 대한 판정 기준과 계산 검산을 동시에 얻게 된다.

%
%
2) \emph{분해 가능 형식과의 수축.}

여기서 \glqq{}분해 가능 형식\grqq{}이란 --- 통상적인 개념을 약간 일반화하여 --- 더 낮은 차수의 형식들의 곱과 \glqq{}더 높은\grqq{} 형식들로 나타낼 수 있는 형식을 뜻한다. 형식들의 곱은 한 번 수축하면 다시 곱으로 옮겨 간다. 마찬가지로 \emph{비선형} 형식들의 곱은 특수화 a')와 b')(§ 2)에서 곱 정리에 따라 두 번 수축하면 다시 곱으로 옮겨 간다.
\[
a_yb_y a_xb_x=\frac12a_y^2\cdot b_x^2+\frac12b_y^2\cdot a_x^2-\frac12(ab\widehat{yx})^2.
\]
따라서 분해 가능 형식의 1차 편극(전이연산)과 몇몇 2차 편극은 다시 분해 가능 형식이다. 그러므로 다음이 따른다.

분해 가능 형식에 대한 1차(2차) 전이연산의 한 \emph{항}이 분해 가능 형식과 한 번(두 번) 수축하여 생겼다면, 그 항 자체도 분해 가능 형식이 된다.

a) 분해 가능 형식의 형식열에서 바로 다음의 더 높은 형식들이 분해되거나 환원 가능해지는 경우, 또는 한 번 수축할 때 특수화 $y=\widehat{uv}$나 $v=\widehat{xy}$가 일어나는 경우,

b) 나머지 한 번의 수축들이 더 높은 형식으로 이어지는 경우(§ 12, B. $H_k$와 $H_k(k\eta x)$ 참조).

그러한 전이연산의 항은 다음 경우에도 분해될 수 있다.

c) 나머지 한 번의 수축들이 더 낮은 형식으로 이어지고, 이 낮은 형식들이 분해 가능 형식과의 수축과는 독립적으로 분해되는 경우이다. 즉 곱들과 환원하려는 형식으로 나타낼 수 있는 경우인데, 이때에는 해당 항의 수치 계수가 항등적으로 소멸하지 않는다는 사실을 매번 계산으로 먼저 확인해야 한다. 이는 낮은 형식의 \glqq{}이중 환원\grqq{}에 지나지 않는다(§ 23, C. $H_\vartheta(\theta H u)(\theta s u)$ 참조).

%
%
분해 가능 형식은 야코비 행렬식과의 모든 한 번의 수축으로 생기며\srcfn{*)}{Gordan, 앞서 인용한 문헌 § 3 참조.}, 그 밖에도 몇몇 더 높은 수축으로 생긴다. 다음이 성립한다.
\[
(st\widehat{yx})s_y=\frac12\,t\cdot s_y^2-\frac12\,s\cdot t_y^2+\frac12(st\widehat{yx})^2,
\]
\[
(st\widehat{yx})s_y^2=\frac13\,t\cdot s_y^3-\frac13\,s\cdot t_y^3+s_y(st\widehat{yx})^2-\frac13(st\widehat{yx})^3.
\]
유사한 공식이 쌍대 형식들
\[
(\sigma\tau\widehat{vu})v_\sigma
\quad\text{및}\quad
(\sigma\tau\widehat{vu})v_\sigma^2
\]
에 대해서도 성립한다.

%
%
이 절의 정리들에 따라 $S$와 $T$의 수축에서 생기는 모든 환원 불가능한 형식을 세우기 위한 규칙은 다음과 같다.

가장 낮은 수축부터 시작하여, 미리 임의로 정해 둔 경로(예를 들어 행별로 또는 대각 성분에 대칭이 되도록)를 따라 형식열 $ST$를 구성해 나가다가 환원자를 인수로 갖는 형식에 이를 때까지 계속한다. 이 형식이 정하는 형식열은 끝형식이 정리 II에서 세운 조건을 만족하는 즉시 제외한다. 환원 가능한 형식열을 모두 제거한 뒤에는 이중 환원을 적용하여 나머지 환원 가능한 형식들을 모두 제거하고, 분해 가능 형식들도 모두 제거한다. 서로 독립적으로 환원 가능한 형식열들이 겹쳐서 두 번 덮는 전체 형식열의 위치에서는 더 높은 형식들을 환원할 계기가 생긴다(§ 11, D 참조).

%
%
\begin{center}
\textbf{제II장. \quad 상대 완전 기본계.}

\textbf{§ 4.}\quad\emph{첫 번째 상대 완전계.}
\end{center}

이하에서는 첫 번째 정리를 제외하고 4차 형식에 한정한다. 그러나 모듈들을 적절히 선택하면 §§ 5와 6의 결과를 더 높은 차수의 형식으로 쉽게 옮길 수 있으며, 이 결과들은 3차 형식에도 마찬가지로 성립한다\srcfn{*)}{3차 형식의 경우 이와 유사하게 유도한 대응 공식은 다음과 같다:
\[
(abc)a_x^2b_y^2c_z^2=\frac13(sxy)(sxz)(syz)+\frac23\Delta_x\Delta_y\Delta_z\cdot(xyz),\qquad (\mathrm{IV}.)
\]
\[
\begin{gathered}
s=(abc)(abu)(acu)(bcu),\qquad \Delta=(abc)^2a_xb_xc_x,\\
t=\Delta_\vartheta(\Delta\theta u)^2u_\vartheta=\theta_s^2u_su_\vartheta^2,
\qquad T=a_t^3,
\end{gathered}
\]
\[
\left.
\begin{aligned}
(a\theta u)^2&=\frac13 u_xs,\\
a_\vartheta&=\frac13u_x\cdot\Delta,
& a_\vartheta(a\theta u)&=0,
& a_\vartheta(a\theta u)^2&=s,\\
a_\vartheta^2&=\Delta,
& a_\vartheta^2(a\theta u)&=0.
\end{aligned}
\right\}\tag{1.}
\]
}.

%
%
임의 차수의 형식에 대한 첫 번째 상대 완전계는 잘 알려진 다음 정리로 얻어진다.

\clearpage
\noindent\textbf{정리 III.}
\emph{이항 영역에서 넘겨받은 형식들, 즉 대응하는 이항 형식의 형식계에 속한 형식들로부터 Clebsch의 전이 원리에 따른 경계 첨가로 생기는 형식들은 $(abc)$를 법으로 하는 상대 완전계를 이룬다.}

\noindent\textbf{증명:}
형식 $Q'_1,\ldots,Q'_n$이 하나의 이항 기본형의 완전계를 이룬다는 다음 이항 관계
\[
Q'-F(Q')=0=\sum_\lambda\bigl\{(ab)c_x+(bc)a_x+(ca)b_x\bigr\}\varphi_\lambda\srcfnmark{*)}
\]
\srcfntext{*)}{Gordan--Kerschensteiner, 앞의 책 134쪽.}
에는 경계 첨가를 통하여 다음 삼원 관계가 대응한다.
\[
Q-F(Q_i)=\sum_\lambda u_x\cdot(abc)\varphi_\lambda .
\]
그러나 이것은 바로 $(abc)$를 법으로 하는 상대 완전계의 정의 방정식이다. 4차 형식에서는 넘겨받은 형식들의 계가 다음 형식들로 이루어진다.

%
%
\[
\begin{gathered}
f=a_x^4,\qquad
\theta=\theta_x^4u_\vartheta^2=(abu)^2a_x^2b_x^2,
\qquad
K=K_x^6u_k^3=(a\theta u)u_\vartheta^2a_x^3\theta_x^3,\\
j=u_j^6=(a\theta u)^4u_\vartheta^2,\qquad
\nu=u_\nu^4=(abu)^4,
\end{gathered}
\]
이들 가운데 처음 네 형식은 $((abc),\nu)$를 법으로 하는 상대 완전계를 이룬다.

%
%
\begin{center}
\textbf{§ 5.}\quad\emph{모듈 $(abc)$를 모듈 $\Delta$와 $\nu$로 환원.}
\end{center}

괄호 인자 하나로만 주어진 모듈 $(abc)$는 정의(§ 3, a)에 따라 계의 형식들로 환원해야 한다. 달리 말하면 형식열 $(abc)$를 유일하게 정규화해야 한다(§ 1 참조).

\[
a_s^2a_xu_s=\frac13a_s^3\cdot u_x=\frac13S\cdot u_x. \tag{2.}
\]
\[
\left.
\begin{aligned}
(a\Delta u)^2&=a_s-\frac{S}{6}u_x^2,\\
(a\Delta u)^3&=0.
\end{aligned}
\right\} \tag{3.}
\]
\[
\left.
\begin{aligned}
\theta(s)=(ss_1x)^2&=2a_t+\frac13S\cdot\theta-\frac13Tu_x^2,\\
\theta(t)=(tt_1x)^2&=-\frac13S\cdot a_t+\frac23T\cdot a_s+\frac1{18}S^2\cdot\theta-\frac1{18}u_x^2\cdot S\cdot T.
\end{aligned}
\right\} \tag{4.}
\]

%
%
모듈의 열: $(abc)$; $\Delta$; $s$; $t$ (모듈 $t$의 계는 유일한 형식 $t$로 이루어진다).

다음 형식들이
\[
\Delta=(abc)^2a_x^2b_x^2c_x^2,\qquad \nu=(abu)^4
\]
모듈로서 충분하다는 것, 즉 형식들
\[
\begin{gathered}
\Delta=(abc)^2a_x^2b_x^2c_x^2,
\qquad a_\nu=(abc)(bcu)^3a_x^3,\\
a_\nu^2=(abc)^2(bcu)^2a_x^2,
\qquad i=a_\nu^4=(abc)^4
\end{gathered}
\]
이 형식열 $(abc)$를 정한다는 것을 보일 것이다. 형식열 $a_\nu$에는 다음 항등식에 따라 형식 $a_\nu^3$이 빠진다.
\[
a_\nu^3=(abc)^3a_x(bcu)=\frac13u_x\cdot(abc)^4=\frac13i\cdot u_x.
\]

%
%
모듈 하나를 더 높은 모듈로 환원하려면 정의에 따라, 그 모듈과의 가장 일반적인 수축 또는 고려 대상인 모든 특수 수축을 더 높은 모듈과의 수축으로 환원해야 한다.

우리의 경우 가장 일반적인 수축은 다음 식들에서 생긴다.
\[
(abc)a_x^3b_y^3c_z^3,
\qquad
(abc)a_x^2b_y^2c_z^2a_xb_xc_x
\]
(3차 형식들에서 넘겨받음),
\[
(abc)a_xb_yc_za_x^2b_x^2c_x^2
\]
(2차 형식들에서 넘겨받음), 그리고 이 식들을 편극하여 생긴다.

이 식들을 $\Delta$와 형식열 $a_\nu$의 편극, 또는 각각 그 시작항들로 계산하기 위해 § 2에서와 같이 간접적인 길을 택한다. 편극항들
\[
(abc)^2a_x^2b_y^2c_z^2,
\qquad
 a_\nu(\nu xy)(\nu yz)(\nu zx)a_xa_ya_z=(abc)(bc\widehat{xy})(bc\widehat{yz})(bc\widehat{zx})a_xa_ya_z\quad\text{등}
\]
을 기호 인자 $(abc)$를 갖는 식들의 선형 결합으로 전개하고, 방정식계를 역으로 풀어 원하는 관계들과 서로 다른 변수 조합에 따라 취한 편극들 사이의 한 관계를 얻는다. 계산에는 항등식 정리와 행렬식의 곱 정리를 이용한다.

%
%
$(abc)a_x^3b_y^3c_z^3$에 대한 방정식계는 다음과 같다.
\begingroup\small
\[
\begin{aligned}
1)\quad
\sum_3a_\nu(\nu yz)^3a_x^3
&=6(abc)a_x^3b_y^3c_z^3-6\sum_3(abc)a_x^3b_y^2b_zc_z^2c_y\srcfnmark{*)},\\[0.4em]
2)\quad
3(abc)^2a_x^2b_y^2c_z^2\cdot(xyz)-\frac12\sum_3a_\nu^2(\nu yz)^2a_x^2\cdot(xyz)
&=2\sum_3(abc)a_x^3b_y^2b_zc_z^2c_y\\
&\quad-4\sum_3(abc)a_xa_ya_zb_y^2b_xc_z^2c_x,\\[0.4em]
3)\quad
a_\nu(\nu xy)(\nu yz)(\nu zx)a_xa_ya_z
&=2\sum_3(abc)a_xa_ya_zb_y^2b_xc_z^2c_x,\\[0.4em]
4)\quad
\sum_3a_\nu(\nu yz)^3a_x^3-\frac32\sum_3a_\nu^2(\nu yz)^2a_x^2\cdot(xyz)+\frac16i\cdot(xyz)^3
&=6\sum_3(abc)a_xa_ya_zb_y^2b_xc_z^2c_x.
\end{aligned}
\]
\endgroup
\srcfntext{*)}{$\sum_3$은 시작항에서 변수들을 순환 치환하여 생기는 세 항의 합을 뜻한다. 방정식들은 그 합의 각 항에 대해서도 개별적으로 성립한다.}

%
%
이로부터 $(abc)a_x^3b_y^3c_z^3$에 대해 다음을 얻고, $(abc)a_x^2b_y^2c_z^2a_xb_xc_x$와 $(abc)a_xb_yc_za_x^2b_x^2c_x^2$에 대해서도 유사한 계산으로 다음을 얻는다.
\[
\begin{aligned}
(abc)a_x^3b_y^3c_z^3
&=\frac32(abc)^2a_x^2b_y^2c_z^2\cdot(xyz)
+\frac32a_\nu(\nu xy)(\nu yz)(\nu zx)a_xa_ya_z
-\frac1{12}i\cdot(xyz)^3,\\[0.4em]
(\mathrm{IV}.)\quad
(abc)a_x^2b_y^2c_z^2a_xb_xc_x
&=\frac23(abc)^2a_x^2b_yb_xc_zc_x\cdot(xyz)
+\frac13a_\nu(\nu xy)(\nu yz)(\nu zx)a_x^3\\
&\quad-\frac16a_\nu^2(\nu xy)(\nu zx)a_x^2\cdot(xyz),\\[0.4em]
(abc)a_xb_yc_za_x^2b_x^2c_x^2
&=\frac16(abc)^2a_x^2b_x^2c_x^2\cdot(xyz).
\end{aligned}
\]

이로써 네 형식 $f,\theta,K,j$로 이루어진, $(\Delta,\nu)$를 법으로 하는 상대 완전계를 얻었다.

따라서 이항 영역에서 넘겨받지 않은 $f$와 $\theta$의 모든 전이연산, 그리고 이항 영역에서 $(ab)^4$로 환원 가능한 전이연산 $(a\theta u)^2$는 모두 기호 $\Delta$와 $\nu$로 나타낼 수 있다.

%
%
이하에서 기본이 되는 환원 공식들을 전개 IV의 특수화로, 또는 더 짧게는 직접 계산(기호의 교환)으로 얻는다. $a_\vartheta(a\theta u)^2$, $a_\vartheta^2(a\theta u)^2$, $(a\theta u)^2$에 대해서는 다음과 같이 둔다.
\[
\begin{aligned}
a_\vartheta(a\theta u)^2
&=(abc)(bcu)(abu)^2-\frac13(abc)(bcu)\{(bcu)a_x-u_x\cdot(abc)\}^2,\\
a_\vartheta^2(a\theta u)^2
&=(abc)^2(abu)^2-\frac13(abc)^2\{(bcu)a_x-u_x\cdot(abc)\}^2
\end{aligned}
\]
$(a\theta u)^2$에 대해서는\srcfn{*)}{Gordan--Kerschensteiner, 앞서 인용한 문헌 92쪽.} 다음을 쓴다.
\[
\begin{aligned}
(abu)a_y^3b_x^3&=\frac32u_\vartheta(\vartheta yx)\theta_y^2+\frac14(\nu yx)^3,\\
(abu)(acu)^3b_x^3&=\frac32u_\vartheta(\vartheta\widehat{au}x)(\theta au)^2+\frac14(\nu\widehat{au}x)^3:
\end{aligned}
\]

%
%
\begingroup\small
\[
\begin{array}{r@{\quad}|l@{\qquad}c|}
& \underline{(a\theta u)^2}=\frac16\nu\cdot f-\frac23u_x\cdot a_\nu+\frac23u_x^2\cdot a_\nu^2-\frac{i}{18}u_x^4
& \bigstar\\[0.55em]
\underline{a_\vartheta}=\frac13u_x\cdot\Delta
& \underline{a_\vartheta(a\theta u)}=0
& \\[0.55em]
& \underline{a_\vartheta^2}=\Delta
&
\end{array}
\]
\vspace{-0.5em}
\[
\begin{array}{r@{\quad}c@{\quad}|l@{\qquad}|l}
(1.) & \bigstar
& \underline{(a\theta u)^3}=0
& \underline{(a\theta u)^4}=j\\[0.55em]
&
& \underline{a_\vartheta(a\theta u)^2}=-\frac56a_\nu+\frac76u_x\cdot a_\nu^2-\frac{i}{9}u_x^3
& \underline{a_\vartheta(a\theta u)^3}=0\\[0.55em]
&
& \underline{a_\vartheta^2(a\theta u)}=0
& \underline{a_\vartheta^2(a\theta u)^2}=\frac23a_\nu^2-\frac{i}{9}u_x^2
\end{array}
\]
\endgroup

%
%
여기에 역시 모듈 $(abc)$의 환원에서 생긴 다음 공식도 취한다.
\[
\begin{array}{r@{\qquad}l}
(2.) & a_\nu^3a_xu_\nu=\frac13i\cdot u_x.
\end{array}
\]

\emph{공식 (1.)과 (2.), 그리고 기본형 $\nu$에 대해 이와 유사하게 만든 공식들로부터, 분해 가능 형식들에 관한 관계들을 제외한 이후의 모든 환원 공식이 편극을 통해 유도된다.} 공식 (1.)에서 다음을 알 수 있다.

형식열은 다음과 같이 환원된다.
\[
\begin{array}{rcl}
a_\vartheta(a\theta u)^2 &\text{는}& \text{기호 }\nu\text{만으로 환원되고},\\[0.2em]
a_\vartheta &\text{는}& \text{기호 }\Delta\text{와 }\nu\text{로 환원되고},\\[0.2em]
(a\theta u)^2 &\text{는}& \text{기호 }j,\Delta\text{ 및 }\nu\text{로 환원되고},\\[0.2em]
(a\theta u)^2\text{에 수축 I을 적용하면} &\text{}& \text{기호 }j\text{와 }\nu\text{로 환원되고},\\[0.2em]
(a\theta u)^2\text{에 수축 II를 적용하면} &\text{}& \text{기호 }\Delta\text{와 }\nu\text{로 환원된다}.
\end{array}
\]

%
%
따라서 다음과 같이 바꾸어 쓸 수 있다.
\begingroup
\renewcommand{\labelenumi}{\arabic{enumi})}
\begin{enumerate}
\item $(\Delta,\nu)$를 법으로 하여: $j$와의 수축을 $(a\theta u)^2$ 또는 $(a\theta u)^3$과의 대응하는 수축으로 바꾼다.
\item $(\nu)$를 법으로 하여: $\Delta$와의 수축을 $a_\vartheta$ 또는 $a_\vartheta(a\theta u)$와의 대응하는 수축으로 바꾸거나, $(a\theta u)^2$과의 특수한 수축으로도 바꾼다(후자는 한 번의 수축 I로만 이어진다).
\end{enumerate}
\endgroup
특히 다음이 성립한다.
\[
\begin{aligned}
(a\theta u)^2a_y^2\theta_x^2&=\frac13(jyx)^2\modu{\nu},&
(a\theta u)^2a_y\theta_y&=-\frac16(jyx)^2\modu{\nu},\\[0.4em]
(a\theta u)^2v_\vartheta^2&=\frac13(\Delta uv)^2\modu{\nu},&
(a\theta u)(a\theta v)u_\vartheta v_\vartheta&=-\frac16(\Delta uv)^2\modu{\nu},\\[0.4em]
(a\theta u)^2a_yv_\vartheta^2&=\frac13(\Delta uv)^2\Delta_y\modu{\nu}.
\end{aligned}
\]

%
%
\begin{center}
\textbf{§ 6.}\quad\emph{모듈 $(\Delta,\nu)$를 모듈 $(\nu)$로 환원.}
\end{center}

앞 절의 환원 공식들은 $\nu$를 법으로 하는 상대 완전계를 직접 세울 수단을 제공한다. 넘겨받은 형식들의 계에는 형식 $\Delta$와 $(a\Delta u)=N$만 덧붙이면 된다는 것을 보일 것이다(3차 형식의 경우와 유사하며, 사실 일반적으로 성립한다).

모듈 $\Delta$를 환원하기 위해 고려해야 할 모든 특수 수축, 즉 $(\Delta,\nu)$를 법으로 하는 상대 완전계와 $\Delta$의 계 사이의 수축을 살펴본다. 계수에 관한 차수가 가장 낮은 형식들, 곧 $f$와 $\Delta$의 수축에서 출발한다.

형식 $(a\Delta u)^2$, $(a\Delta u)^3$, $(a\Delta u)^4$를 환원하기 위해 § 5에 따라 기호 $\Delta$를 형식열의 더 낮은 형식으로 바꾼다. 즉 다음 식들을
\[
a_\vartheta b_\vartheta(b\theta u)^2,
\qquad a_\vartheta(a\theta u)b_\vartheta(b\theta u)^2,
\qquad a_\vartheta(a\theta u)b_\vartheta(b\theta u)^3
\]
1) 형식열 $a_\vartheta$의 편극(기호 $\Delta$와 $\nu$)에 따라,

2) 형식열 $b_\vartheta(b\theta u)^2$의 편극(기호 $\nu$)에 따라

전개하여 다음 환원 공식들을 얻는다(계산은 아래 참조).

%
%
\[
\begin{aligned}
\text{a)}\quad \underline{(a\Delta u)^2}
&=\frac12(\vartheta \nu x)^2-\frac25f\cdot a_\nu^2
-\frac45u_x\cdot\theta_\nu(\vartheta \nu x)^2
+u_x^2\left\{\frac16i f-\frac1{40}s\right\},\\
(3.)\quad\text{b)}\quad \underline{(a\Delta u)^3}
&=-\frac3{10}\theta_\nu(\vartheta \nu x),\\
\text{c)}\quad \underline{(a\Delta u)^4}
&=\frac{21}{10}\theta_\nu^2-\frac3{10}H
-\frac65u_x\cdot\theta_\nu^3+\frac1{10}u_x^2\cdot\varrho.
\end{aligned}
\]

($a_\vartheta^2$를 소거한 뒤 식 $a_\vartheta^2(b\theta u)^2$, $a_\vartheta^2(b\theta u)^3$, $a_\vartheta^2(a\theta u)^2(b\theta u)^2$, $a_\vartheta^2(a\theta u)(b\theta u)^3$을 이중 환원해도 같은 결과를 얻는다.)
공식 (3.)으로부터 $(a\Delta u)^2$가 모듈 $\nu$에 관한 환원자임이 따른다. 따라서 다음을 얻는다.

1) $\Delta$의 계의 환원: 형식열 $(\Delta\Delta' u)^2\equiv a_\vartheta^2(a\Delta u)^2$은 환원자 $(a\Delta u)^2$를 인수로 갖는다.

2) 모듈 $\Delta$와의 수축으로 생긴 계의 환원:

형식열 $(\theta\Delta u)^2$은 환원자 $(a\Delta u)^2$를 인수로 갖는다.

형식 $(\theta\Delta u)$와 $\Delta_\vartheta$에 대해서는 다음을 얻는다.
\[
(a\Delta u)^2(abu)=(\theta\Delta u)-u_x\cdot\Delta_\vartheta(\theta\Delta u),
\]
\[
(a\Delta u)(a\Delta b)=-\frac12\Delta_\vartheta+\frac12u_x\cdot\Delta_\vartheta^2.
\]

그러나 환원자 $(\theta\Delta u)$로부터 모듈 $\Delta$와의 수축으로 생기는 계의 환원이 따른다.

%
%
\bigskip

따라서 $\nu$를 법으로 하는 상대 완전계는 다음 여섯 형식으로 이루어진다.
\[
f,\ \theta,\ K,\ j,\ \A,\ N.
\]

§ 2의 급수 전개에 대한 예로 공식 (3.)a의 유도를 자세히 제시한다. 이후의 계산에서는 최종 공식 III을 직접 세울 것이다. (소멸하는 형식의 편극은 계산 도중에 이미 생략하였다. 첫째 줄은 전개 II에 따라 대입한 값을, 둘째 줄은 연립계의 마지막 방정식부터 시작하여 재귀적으로 구한 최종값을 나타낸다.)

%
%
전개 II에 따라 다음이 성립한다.
\[
\begin{aligned}
\text{1)}\quad&m=3.\quad n=4.\quad \lambda=2.\quad \mu=0.\quad \nu=\varkappa=1,\\[0.25em]
\underline{a_\vartheta b_\vartheta(b\theta u)^2}
&=[a_\vartheta]_{\widehat{ub^2}}b+\frac67\{a_\vartheta b_\vartheta(a\theta u)(b\theta u)-u_x\cdot a_\vartheta b_\vartheta(a\theta b)(b\theta u)\}\\
&\quad-\frac17\{a_\vartheta(a\theta u)^2b_\vartheta-2u_x\cdot a_\vartheta(a\theta u)(a\theta b)b_\vartheta+u_x^2\cdot a_\vartheta(a\theta b)^2b_\vartheta\}\\
&=\frac23(a\A u)^2+\frac15[a_\vartheta(a\theta u)^2]b+\frac1{30}f\cdot[a_\vartheta^2(a\theta u)^2]-\frac25u_x\cdot[a_\vartheta(a\theta u)^2]b^2\\
&\quad-\frac1{15}u_x\cdot[a_\vartheta^2(a\theta u)^2]b+\frac15u_x^2\cdot[a_\vartheta(a\theta u)^2]b^3+\frac1{30}u_x^2\cdot[a_\vartheta^2(a\theta u)^2]b^2;\\[0.7em]
\text{2)}\quad&m=2.\quad n=3.\quad \lambda=1.\quad \mu=1.\quad \nu=\varkappa=1,\\[0.25em]
\underline{a_\vartheta(a\theta u)b_\vartheta(b\theta u)}
&=\frac25\{a_\vartheta(a\theta u)^2b_\vartheta-u_x\cdot a_\vartheta(a\theta u)(a\theta b)b_\vartheta\}+\frac12a_\vartheta^2(b\theta u)^2\\
&\quad-\frac15\{a_\vartheta^2(b\theta u)(a\theta u)-u_x\cdot a_\vartheta^2(a\theta b)(b\theta u)\}\\
&=\frac12(a\A u)^2+\frac25[a_\vartheta(a\theta u)^2]b+\frac1{15}[a_\vartheta^2(a\theta u)^2]\cdot f\\
&\quad-\frac25u_x\cdot[a_\vartheta(a\theta u)^2]b^2-\frac1{10}u_x\cdot[a_\vartheta^2(a\theta u)^2]b+\frac1{30}u_x^2\cdot[a_\vartheta^2(a\theta u)^2]b^2;\\[0.7em]
\text{3)}\quad&m=2.\quad n=3.\quad \lambda=1.\quad \mu=1.\quad \nu=1.\quad \varkappa=2,
\end{aligned}
\]

%
%
\[
\begin{aligned}
\underline{a_\vartheta(a\theta b)b_\vartheta(b\theta u)}
&=\frac25\{a_\vartheta(a\theta b)(a\theta u)b_\vartheta-u_x\cdot a_\vartheta(a\theta b)^2b_\vartheta\}\\
&=\frac25[a_\vartheta(a\theta u)^2]b^2+\frac1{30}[a_\vartheta^2(a\theta u)^2]b-\frac25u_x\cdot[a_\vartheta(a\theta u)^2]b^3\\
&\quad-\frac1{30}u_x\cdot[a_\vartheta^2(a\theta u)^2]b^2;\\[0.7em]
\text{4)}\quad&m=1.\quad n=2.\quad \lambda=0.\quad \mu=2.\quad \nu=\varkappa=1,\\[0.25em]
\underline{a_\vartheta(a\theta u)^2b_\vartheta}
&=[a_\vartheta(a\theta u)^2]b+\frac23a_\vartheta^2(a\theta u)(b\theta u)\\
&=[a_\vartheta(a\theta u)^2]b+\frac16[a_\vartheta^2(a\theta u)^2]\cdot f-\frac16u_x\cdot[a_\vartheta^2(a\theta u)^2]b;\\[0.7em]
\text{5)}\quad&m=1.\quad n=2.\quad \lambda=0.\quad \mu=2.\quad \nu=1.\quad \varkappa=2.\quad(\varkappa>\nu),\\[0.25em]
\underline{a_\vartheta(a\theta u)(a\theta b)b_\vartheta}
&=[a_\vartheta(a\theta u)^2]b^2+\frac13a_\vartheta^2(a\theta b)(b\theta u)\\
&=[a_\vartheta(a\theta u)^2]b^2+\frac1{12}[a_\vartheta^2(a\theta u)^2]b-\frac1{12}u_x\cdot[a_\vartheta^2(a\theta u)^2]b^2;\\[0.7em]
\text{6)}\quad&m=1.\quad n=2.\quad \lambda=0.\quad \mu=2.\quad \nu=1.\quad \varkappa=3,\\[0.25em]
\underline{a_\vartheta(a\theta b)^2b_\vartheta}&=[a_\vartheta(a\theta u)^2]b^3;\\[0.7em]
\text{7)}\quad&m=2.\quad n=4.\quad \lambda=2.\quad \mu=\nu=\varkappa=0.\quad(\text{전개 I에 따라}),\\[0.25em]
\underline{a_\vartheta^2(b\theta u)^2}&=[a_\vartheta^2]_{\widehat{ub^2}}+\frac1{10}\{[a_\vartheta^2(a\theta u)^2]\cdot f-2u_x\cdot[a_\vartheta^2(a\theta u)^2]b+u_x^2\cdot[a_\vartheta^2(a\theta u)^2]b^2\};\\[0.7em]
\text{8)}\quad&m=1.\quad n=3.\quad \lambda=1.\quad \mu=1.\quad \nu=\varkappa=0.\quad(\text{전개 I}),\\[0.25em]
\underline{a_\vartheta^2(a\theta u)(b\theta u)}&=\frac14[a_\vartheta^2(a\theta u)^2]\cdot f-\frac14u_x\cdot[a_\vartheta^2(a\theta u)^2]b;\\[0.7em]
\text{9)}\quad&m=1.\quad n=3.\quad \lambda=1.\quad \mu=\varkappa=1.\quad \nu=0.\quad(\text{전개 I}),\\[0.25em]
\underline{a_\vartheta^2(a\theta b)(b\theta u)}&=\frac14[a_\vartheta^2(a\theta u)^2]b-\frac14u_x\cdot[a_\vartheta^2(a\theta u)^2]b^2.
\end{aligned}
\]

%
%
1)과 4)로부터 다음이 따른다.
\[
\begin{aligned}
(a\A u)^2
&=\frac65[a_\vartheta(a\theta u)^2]b+\frac15f\cdot[a_\vartheta^2(a\theta u)^2]+\frac35u_x\cdot[a_\vartheta(a\theta u)^2]b^2-\frac3{20}u_x\cdot[a_\vartheta^2(a\theta u)^2]b\\
&\quad-\frac3{10}u_x^2\cdot[a_\vartheta(a\theta u)^2]b^3-\frac1{20}u_x^2\cdot[a_\vartheta^2(a\theta u)^2]b^2,\\[0.4em]
(a\A u)^2&=-a_\nu b_\nu+\frac35f\cdot a_\nu^2+\frac45u_x\cdot a_\nu^2b_\nu-\frac3{20}u_x^2\cdot a_\nu^2b_\nu^2-\frac1{12}u_x^2\cdot i\cdot f.
\end{aligned}
\]

($\theta$의 기호로 환산하는 데 대해서는 §§ 8과 9를 참조하라.)
공식 (3.)a는 보조변수 $w=x\widehat{ub}$를 도입하여 전개 I을 직접 전이하면 더 짧게 계산할 수도 있다. 반면 공식 (3.)b와 (3.)c에서는 우리가 택한 방법이 이미 더 간단해진다. (3.)b에서는 § 9에 따라 인수 $u_x$를 가진 모든 항을 생략해야 한다는 것을 미리 알고 있다. (3.)b와 (3.)c의 계산에는 다음을 얻는다.

%
%
\[
\begin{aligned}
(3.)\mathrm{b)}\quad 1)&\quad a_\vartheta(a\theta u)b_\vartheta(b\theta u)^2=\frac12(a\A u)^3+\frac45[a_\vartheta(a\theta u)^2]_{\widehat{ub}}b+\frac{77}{360}[a_\vartheta^2(a\theta u)^2]_{\widehat{ub}},\\
2)&\quad a_\vartheta(a\theta u)b_\vartheta(b\theta u)^2=[a_\vartheta(a\theta u)^2]_{\widehat{ub}}b+\frac{13}{36}[a_\vartheta^2(a\theta u)^2]_{\widehat{ub}};\\[0.5em]
(3.)\mathrm{c)}\quad 1)&\quad a_\vartheta(a\theta u)b_\vartheta(b\theta u)^3=\frac12(a\A u)^4+\frac65[a_\vartheta(a\theta u)^2]_{\widehat{ub^2}}b+\frac{53}{120}[a_\vartheta^2(a\theta u)^2]_{\widehat{ub^2}}\\
&\qquad-\frac65u_x\cdot[a_\vartheta(a\theta u)^2]_{\widehat{ub^2}}b^2-\frac35u_x\cdot[a_\vartheta^2(a\theta u)^2]_{\widehat{ub^2}}b+\frac{19}{120}u_x^2\cdot[a_\vartheta^2(a\theta u)^2]_{\widehat{ub^3}}b^2,\\
2)&\quad a_\vartheta(a\theta u)b_\vartheta(b\theta u)^3=\frac34[a_\vartheta^2(a\theta u)^2]_{\widehat{ub^2}}.
\end{aligned}
\]

%
%
부기: § 5에 따라 환원 가능한 $f$와 $j$의 전이연산들도 같은 방식으로 계산된다. 다음을 얻는다.
\[
\tag{4.}
\begin{aligned}
a_j&=-\theta_\nu+u_x\cdot\left\{\frac92\theta_\nu^2-\frac12H\right\}-3u_x^2\cdot\theta_\nu^3+\frac12u_x^3\cdot\varrho,\\
a_j^2&=\frac{21}{10}\theta_\nu^2-\frac3{10}H-2u_x\cdot\theta_\nu^3+\frac12u_x^2\cdot\varrho,\\
a_j^3&=-\frac7{10}\theta_\nu^3+\frac12u_x\cdot\varrho,\\
a_j^4&=\frac35\varrho.
\end{aligned}
\]

(3.)과 (4.) 및 이항 영역으로부터의 전이에 의해 다음이 따른다.
\[
(\theta\theta'u)^2=\frac13 j\cdot f\modu{\nu}\srcfn{*)}{Gordan--Kerschensteiner, 앞의 책 181쪽.}.
\]
형식열 $(\theta\theta'u)^2$의 나머지 형식들과 형식 $\theta'_\vartheta$는 기호 $\nu$로 환원 가능하다. 따라서 다음과 같이 바꿀 수 있다.
\[
\modu{\nu}:\quad \text{$f\cdot j$와의 수축을 $(\theta\theta'u)^2$와의 대응 수축으로.}
\]
또한 다음을 얻는다.
\[
(\vartheta\vartheta'x)^2=\frac43 f\cdot\A,\qquad
\theta_{\vartheta'}(\vartheta\vartheta'x)=-N.
\]

%
%
\begin{center}
{\Large\bfseries 제III장.\quad $(s,\varrho,t)$를 법으로 하는 상대 완전계.}
\end{center}

\begin{center}
\textbf{\S\ 7.}\quad\emph{계 구성의 개요.}
\end{center}

$\nu$를 법으로 하는 상대 완전계에서 그다음으로 높은 모듈을 법으로 하는 상대 완전계로 넘어가려면, 정의 \S\ 3에 따라 이 더 높은 모듈에 관하여 $\nu$의 계를 구성하고 이를 $\nu$를 법으로 하는 상대 완전계의 형식들과 전이연산해야 한다. 그런데 $\nu=u_\nu^4$는 형식 $f$의 4차 반공변식으로서 $f$와 쌍대 관계에 있다. 따라서 $\nu$를 기본형으로 보면 $\nu(\nu)=(\nu\nu_1x)^4=s_x^4$를 법으로 하는 여섯 형식의 상대 완전계를 얻는다.

그러므로 $s$를 법으로 하는 상대 완전계(계 III)를 구성하려면 다음 두 계를 전이연산해야 한다.
\[
\begin{array}{ll}
\text{계 I:} & f,\ \theta,\ K,\ j,\ \A,\ N\quad \text{와}\vphantom{\dfrac11}\\[0.3em]
\text{계 II:} & \nu,\ H=(\nu\nu_1x)^2u_\nu^2u_{\nu_1}^2,\ L=(\nu\eta x)H_x^2u_\nu^3u_\eta^3,\ g=(\nu\eta x)^4H_x^2,\\
& \sigma=H_\nu^2u_\nu^2u_\eta^4,\ (\nu\sigma x).
\end{array}
\]
형식 $g$는 환원 가능하며, 따라서 계 II는 다섯 형식으로만 이루어진다는 것이 드러날 것이다(공식 (13.)).

%
%
계산 과정에서 이차형식들
\[
\varrho=u_\varrho^2=\theta_\nu^4u_\vartheta^2,\qquad t=t_x^2=a_\eta^4H_x^2
\]
을 모듈로 첨가하여 $(s,\varrho,t)$를 법으로 하는 상대 완전계를 얻는다. 이때 모듈 $(\varrho,t)$는 모듈 $s$보다 더 높은 모듈로 간주한다.

개별 형식들은 계수차수에 따라 배열한다.

수축
\[
\theta_\eta(K_\eta,\theta_\iota,\ldots)
\]
을 수축
\[
H_\vartheta(H_\kappa,L_\vartheta,\ldots).
\]
보다 더 높은 것으로 정규화한다. 이로써 \S\ 1과 \S\ 3b에 따라 같은 차수의 개별 형식들의 순서가 유일하게 정해진다.

%
%
같은 차수의 형식들은 표를 이용하여 구성한다.

I. 특정 계수차수를 얻기 위해 계 I과 계 II의 어떤 형식들을 서로 수축해야 하는지 제시한다.

II. 형식열의 도식에 따라 환원 불가능한 형식들을 제시한다.

III. 환원 가능하거나 분해 가능한 형식열 또는 형식을 환원한다. 즉 전체 형식열이 끊기는 위치를 밝히고, 이로써 제시한 것들이 모든 환원 불가능한 구성식을 소진한다는 것을 증명한다.

IV. 결론으로서 1) 새로 생긴 환원자들을 제시하고, 2) 같은 계의 형식들과 곱을 이루어 다른 계의 형식들과 수축하는 데 들어가지 않는 형식들을 제시한다.

%
%
다음 두 경우만 나타난다는 것이 드러날 것이다.

1) 계 I의 형식 하나와 계 II의 형식 하나를 각각 수축하는 경우.

2) $\theta$ 또는 $H$의 거듭제곱을 각각 수축하거나, 한 계의 두 형식을 다른 계의 형식 하나와 각각 수축하는 경우.

다음과 같은 수축 도식을 얻게 될 것이다.
\[
\begin{array}{r@{\ }c@{\ }l@{\qquad}c@{\qquad}r@{\ }c@{\ }l}
\multicolumn{3}{c}{\text{계 I}} & \text{와 수축} & \multicolumn{3}{c}{\text{계 II}}\\[0.25em]
1. & \text{차수} & f & & 2. & \text{차수} & \nu\\
2. & ,, & \theta & & 4. & ,, & H\\
3. & ,, & K,j,\A & & 6. & ,, & L,\sigma\\
4. & ,, & N,\theta^2,f\cdot j & & 8. & ,, & (\nu\sigma x),H^2\\
5. & ,, & \theta\cdot K & & 10. & ,, & H\cdot L\\
6. & ,, & \theta^3,\A\cdot j & & 12. & ,, & H^3.
\end{array}
\]

%
%
계산을 검산하기 위해 \glqq{}이중 환원\grqq{} 외에도 다음 두 특수 형식을 이용한다.
\[
\begin{array}{ll}
\text{I.}&\begin{array}{rlrlrl}
f&=\displaystyle\sum_3x_1^3x_2,& a_\nu^2&=\dfrac{i}{6}u_x^2,& s&=\dfrac{i}{3}f,\\
a_\eta^2&=-\dfrac19i\theta,& \A_\nu^2&=\dfrac{i}{6}\theta,& \sigma&=-\dfrac{i}{6}j,\\
g&=-\dfrac23i\A,& \varrho&=0,& t&=0,
\end{array}\\[1.2em]
\text{II.}&\begin{array}{rlrlrl}
f&=\displaystyle\sum_3x_1^4,& s&=\dfrac43if,& a_\eta^2&=\dfrac29i\theta,\\
\A_\nu^2&=\dfrac1{15}i\theta,& \sigma&=\dfrac43ij,& g&=\dfrac43i\A,\\
\varrho&=0,& t&=0.
\end{array}
\end{array}
\]

%
%
후기: 공식 (1.), (2.), (3.), (4.)에 대응하여 기본형 $\nu$에 대해서는 다음 공식들이 성립한다.
\[
\tag{5.}
\begin{array}{@{}c|c|c@{}}
\cdot & \underline{H_\nu}=\dfrac13u_x\cdot\sigma & \cdot\\[0.45em]
\underline{(\nu\eta x)^2}=A & \underline{H_\nu(\nu\eta x)}=0 & \underline{H_\nu^2}=\sigma\\[0.45em]
\underline{(\nu\eta x)^3}=0 & \underline{H_\nu(\nu\eta x)^2}=B & \underline{H_\nu^2(\nu\eta x)}=0\\[0.45em]
\underline{(\nu\eta x)^4}=g & \underline{H_\nu(\nu\eta x)^3}=0 & \underline{H_\nu^2(\nu\eta x)^2}=C,
\end{array}
\]
\[
\begin{aligned}
A&=\frac16s\cdot\nu-\frac23u_x\cdot s_\nu+\frac23u_x^2\cdot s_\nu^2-\frac{J}{18}u_x^4,\\
B&=-\frac56s_\nu+\frac76u_x\cdot s_\nu^2-\frac{J}{9}u_x^3,\\
C&=\frac23s_\nu^2-\frac{J}{9}u_x^2.
\end{aligned}
\]
\[
\tag{6.}
s_\nu^3u_\nu s_x=\frac13u_x\cdot s_\nu^4=\frac13J\cdot u_x.
\]

%
%
\[
\tag{7.}
\begin{aligned}
\text{a)}\quad (\nu\sigma x)^2
&=\frac12(Hsu)^2-\frac25\nu\cdot s_\nu^2-\frac45u_x\cdot s_\eta(Hsu)^2+u_x^2\left\{\frac16J\cdot\nu-\frac1{40}(ss'u)^4\right\},\\
\text{b)}\quad (\nu\sigma x)^3&=-\frac3{10}s_\eta(Hsu),\\
\text{c)}\quad (\nu\sigma x)^4&=\frac{21}{10}s_\eta^2-\frac3{10}Z-\frac65u_x\cdot s_\eta^3+\frac1{10}u_x^2\cdot s_\eta^4.
\end{aligned}
\]

%
%
\[
\tag{8.}
\begin{aligned}
\text{a)}\quad g_\nu&=-s_\eta+u_x\left\{\frac92s_\eta^2-\frac12Z\right\}-3u_x^2\cdot s_\eta^3+\frac12u_x^3\cdot s_\eta^4,\\
\text{b)}\quad g_\nu^2&=\frac{21}{10}s_\eta^2-\frac3{10}Z-2u_x\cdot s_\eta^3+\frac12u_x^2\cdot s_\eta^4,\\
\text{c)}\quad g_\nu^3&=-\frac7{10}s_\eta^3+\frac12u_x\cdot s_\eta^4,\\
\text{d)}\quad g_\nu^4&=\frac35s_\eta^4.
\end{aligned}
\]

%
%
\begin{center}
\textbf{\S\ 8.}\quad\emph{계수차수 3의 형식들(계 III).}
\end{center}

\noindent\textbf{$f$와 $\nu$의 수축.}

\noindent\emph{환원 불가능 형식들:}
\[
\begin{array}{cccc}
a_\nu&\cdot&\cdot&\cdot\\
\cdot&a_\nu^2&\cdot&\cdot\\
\cdot&\cdot&\cdot&\cdot\\
\cdot&\cdot&\cdot&a_\nu^4=i.
\end{array}
\]

\noindent\emph{환원:}
\[
a_\nu^3=\frac13iu_x\qquad\text{(공식 (2.)).}
\]
\noindent\emph{결론:} 1) $a_\nu^3$은 환원자이다. 2) $f$는 반변 형식 $(j)$과의 곱 안에서만 $\nu$와 수축된다. $f$에 대한 $\nu$의 경우에도 마찬가지이다.

%
%
\begin{center}
\textbf{\S\ 9.}\quad\emph{계수차수 4의 형식들(계 III).}
\end{center}

\noindent\textbf{$\theta$와 $\nu$의 수축.}

\noindent\emph{환원 불가능 형식들:}
\[
\begin{array}{cccc}
(\vartheta\nu x)&\theta_\nu&\cdot&\cdot\\
(\vartheta\nu x)^2&\theta_\nu(\vartheta\nu x)&\theta_\nu^2&\cdot\\
\cdot&\theta_\nu(\vartheta\nu x)^2&\cdot&\theta_\nu^3.
\end{array}
\]

\noindent\emph{환원:} 환원자 $a_\nu^3$으로부터 식 $a_\nu^3(abu)$, $a_\nu^3b_\nu$, $a_\nu^3b_\nu(abu)$을 다음과 같이 두고 이중 환원한다.

%
%
\[
\begin{aligned}
\theta_\nu^2(\vartheta\nu x)
&=a_\nu^2(\widehat{ab}\nu x)(abu)-\frac13(\nu\nu_1x)^3
=a_\nu b_\nu(\widehat{ab}\nu x)(abu)+\frac16(\nu\nu_1x)^3\\
&=a_\nu^3(abu)-a_\nu^2b_\nu(abu)=2a_\nu^2b_\nu(abu)=\frac23a_\nu^3(abu),\\[.4em]
\theta_\nu^2(\vartheta\nu x)^2
&=a_\nu^2(\widehat{ab}\nu x)^2-\frac13(\nu\nu_1x)^4
=a_\nu b_\nu(\widehat{ab}\nu x)^2+\frac16(\nu\nu_1x)^4\\
&=i\cdot f-2a_\nu^3b_\nu+a_\nu^2b_\nu^2-\frac13s
=2a_\nu^3b_\nu-2a_\nu^2b_\nu^2+\frac16s
=\frac23i\cdot f-\frac23a_\nu^3b_\nu-\frac16s,\\[.4em]
\theta_\nu^3(\vartheta\nu x)&=a_\nu^2b_\nu(\widehat{ab}\nu x)(abu)=a_\nu^3b_\nu(abu).
\end{aligned}
\]

다음 공식들이 생긴다.
\[
\tag{9}
\begin{array}{c|c|c}
\underline{\theta_\nu^2(\vartheta\nu x)}=0
&\underline{\theta_\nu^3(\vartheta\nu x)}=0
&\underline{\theta_\nu^4u_\vartheta^2}=u_\varrho^2=\varrho\quad(\text{첨가 모듈, \S\ 7}),\\[0.6em]
\underline{\theta_\nu^2(\vartheta\nu x)^2}=\dfrac49 i f-\dfrac16s
&&
\end{array}
\]

\noindent\emph{결론:} 1) $\theta_\nu^2(\vartheta\nu x)^2$은 환원자이다.
2) $\nu$는 반변 형식 $(g)$과의 곱 안에서만 $\theta$와 수축된다.

%
%
\begin{center}
\textbf{\S\ 10.}\quad\emph{계수차수 5의 형식들(계 III).}
\end{center}

\[
\text{다음의 수축: }\left\{\begin{array}{l}
K,j,\A \text{와 }\nu,\\
f\text{와 }H.
\end{array}\right.
\]

\noindent\emph{환원 불가능 형식들:}
\[
\begin{array}{c@{\qquad}c@{\qquad}c@{\qquad}c}
\text{A)} & \begin{array}{cccc}
\cdot&K_\nu&\cdot&\cdot\\
(k\nu x)^2&\cdot&K_\nu^2&\cdot\\
(k\nu x)^3&K_\nu(k\nu x)^2&\cdot&\cdot\\
\cdot&K_\nu(k\nu x)^3&\cdot&\cdot,
\end{array}
&
\text{B)}\ \begin{array}{c}(j\nu x)\\(j\nu x)^2\\(j\nu x)^3\\ \cdot\end{array}
&
\text{C)}\ \begin{array}{c}\A_\nu\\\cdot\\\cdot\\\cdot\end{array}
\end{array}
\]
\[
\text{D)}\qquad
\begin{array}{cccc}
(aHu)&(aHu)^2&\cdot&\\
a_\eta&a_\eta(aHu)&a_\eta(aHu)^2&\\
\cdot&a_\eta^2&\cdot&\cdot
\end{array}
\]

%
%
\noindent\emph{환원:}

\noindent A) 형식열 $K_\nu^3$: 환원자 $a_\nu^3$.

형식 $K_\nu^2(k\nu x)^2$: 환원자 $\theta_\nu^2(\vartheta\nu x)^2$.
\[
(k\nu x),\quad K_\nu(k\nu x),\quad K_\nu^2(k\nu x)
\]
은 \S\ 3에 따라 분해 가능 형식이다.

%
%
\noindent B) 형식
\[
\begin{aligned}
(j\nu x)^4&\equiv(\widehat{a\theta}\nu x)^4
=a_\nu^4\cdot\theta+\theta_\nu^4\cdot f-4a_\nu^3\theta_\nu
-4a_\nu\{a_\nu\cdot\theta_x-(\widehat{a\theta}\nu x)\}^3\\
&\quad+6a_\nu^2\{a_\nu\theta_x-(\widehat{a\theta}\nu x)\}^2
=\R\cdot f+\frac53 i\theta-6a_\nu^2(\widehat{a\theta}\nu x)^2
+4a_\nu(\widehat{a\theta}\nu x)^3,
\end{aligned}
\]
이고, 따라서
\[
(j\nu x)^4=\R\cdot f+\frac53 i\theta\cdot\modu{(\A,\nu)},\qquad (\text{\S\ 5 끝부분 참조}).
\]

%
%
\noindent C) 형식 $\A_\nu^4$: 환원자 $\theta_\nu^4$.

형식 $\A_\nu^2$와 $\A_\nu^3$: 형식 $\A$를 형식열의 더 낮은 형식들로 바꾸어 얻는 환원자 $a_\nu^3$ 또는 $\theta_\nu^4$(\S\ 5 끝부분), 즉 식들
\[
a_\vartheta a_\nu^3,\qquad
a_\vartheta a_\nu^3\theta_\nu,\qquad
a_\vartheta\theta_\nu^4.
\]
의 이중 환원으로 얻는다. $\A_\nu^3$을 두 가지로 계산하면 더 높은 형식 $a_\eta^3$을 환원할 계기가 생긴다.

%
%
\noindent\emph{환원 공식:}
\[
\tag{10}
\begin{aligned}
\text{a)}\quad \A_\nu^2
&=\frac35a_\eta^2-\frac3{10}(asu)^2+\frac13 i\theta
+\frac15u_x^2\cdot(2a_\R^2-t),\\
\text{b)}\quad \A_\nu^3
&=-\frac3{10}a_\R+\frac35u_x\left\{\frac{13}{10}a_\R^2-\frac25t\right\},\\
\text{c)}\quad \A_\nu^4&=a_\R^2-\frac25t.
\end{aligned}
\]

%
%
\noindent\emph{공식의 유도:}
\[
\begin{aligned}
\text{a)}\quad
 a_\vartheta a_\nu^3
&=\frac13 i\theta
=[a_\vartheta]_{\nu^3}+\frac67[a_\vartheta^2]_{\nu^2}
+\frac65[a_\vartheta(a\theta u)^2]_{\nu^2}\,\widehat{\nu x}^2
+\frac{11}{30}[a_\vartheta^2(a\theta u)^2]_{\nu}\,\widehat{\nu x}^2\\
&\quad-\frac67u_x\cdot[a_\vartheta^2]_{\nu^3}
-\frac16u_x\cdot[a_\vartheta^2(a\theta u)^2]_{\nu^2}\,\widehat{\nu x}^2,\\
\frac13 i\theta
&=\A_\nu^2-\frac23u_x\cdot\A_\nu^3-a_\nu a_{\nu_1}(\nu\nu_1x)^2
+\frac25a_\nu^2(\nu\nu_1x)^2
+\frac15u_x\cdot a_\nu^2a_{\nu_1}(\nu\nu_1x)^2;
\end{aligned}
\]

%
%
\[
\begin{aligned}
\text{b)}\quad
 a_\vartheta a_\nu^3\theta_\nu
&=0=[a_\vartheta]_{\nu^4}+\frac9{14}[a_\vartheta^2]_{\nu^3}
+\frac35[a_\vartheta(a\theta u)^2]_{\nu^2}\,\widehat{\nu x}^2
+\frac{17}{120}[a_\vartheta^2(a\theta u)^2]_{\nu}\,\widehat{\nu x}^2\\
&\quad-\frac9{14}u_x\cdot[a_\vartheta^2]_{\nu^4}
-\frac1{24}u_x\cdot[a_\vartheta^2(a\theta u)^2]_{\nu^2}\,\widehat{\nu x}^2,\\
0&=\frac56\A_\nu^3-\frac12u_x\cdot\A_\nu^4-\frac14a_\eta^3+\frac1{20}u_x\cdot a_\eta^4;
\end{aligned}
\]

%
%
\[
\begin{aligned}
\text{b')}\quad
 a_\vartheta\theta_\nu^4
&=a_\R=a_\vartheta\theta_\nu\{a_\nu-(a\theta\widehat{\nu x})\}^3\\
&=-\frac32[a_\vartheta^2]_{\nu^3}
+\frac35[a_\vartheta(a\theta u)^2]_{\nu^2}\,\widehat{\nu x}^2
-\frac{37}{120}[a_\vartheta^2(a\theta u)^2]_{\nu}\,\widehat{\nu x}^2\\
&\quad+\frac32u_x\cdot[a_\vartheta^2]_{\nu^4}
+\frac{49}{120}u_x\cdot[a_\vartheta^2(a\theta u)^2]_{\nu^2}\,\widehat{\nu x}^2,\\
a_\R&=-\frac32\A_\nu^3+\frac32u_x\cdot\A_\nu^4-\frac{11}{20}a_\eta^3+\frac7{20}u_x\cdot a_\eta^4;
\end{aligned}
\]

%
%
\[
\begin{aligned}
\text{c)}\quad
 a_\vartheta^2\theta_\nu^4
&=a_\R^2=[a_\vartheta^2]_{\nu^4}+\frac35[a_\vartheta^2(a\theta u)^2]_{\nu^2}\,\widehat{\nu x}^2,\\
a_\R^2&=\A_\nu^4+\frac25a_\eta^4.
\end{aligned}
\]

\noindent D) C)의 $\A_\nu^3$을 이중 환원하여 $a_\eta^3$을 환원한다.

나머지 환원들은 \S\ 9의 계산과 쌍대인 계산으로 얻는다.

%
%
\noindent\emph{환원 공식:}
\[
\tag{11}
\begin{array}{rl|rl}
\underline{a_\eta^2(aHu)}&=-\dfrac16(asu)^3,&
\underline{a_\eta^2(aHu)^2}&=\dfrac49 i \nu-\dfrac16(asu)^4,\\[0.7em]
\underline{a_\eta^3}&=-a_\R+\dfrac35u_x\cdot a_\R^2+\dfrac15u_x\cdot t,
&\underline{a_\eta^3(aHu)}&=0,\\[0.7em]
&&\underline{a_\eta^4}&=t\quad(\text{모듈로 첨가됨}).
\end{array}
\]

\noindent\emph{결론:}

1) $(j\nu x)^4$, $\A_\nu^2$, $a_\eta^2(aHu)$는 환원자이다.

2) $\nu$는 반변 형식들과의 곱 안에서만 $K,j,\A$와 수축된다. $f$가 $H$와 수축되는 경우와 $j$가 $\nu$와 수축되는 경우에도 마찬가지다. 왜냐하면 $\theta_\nu(\vartheta\cdot j,\nu,x)^3$이 \S\ 11 B에 따라 환원 가능해지기 때문이다.

%
%
\begin{center}
\textbf{\S\ 11.\quad 계수차수 6의 형식들(계 III).}
\end{center}
\[
\text{다음의 수축: }\left\{\begin{array}{l}
\theta^2,\ f\cdot j,\ N \text{과 }\nu,\\
\theta\text{와 }H.
\end{array}\right.
\]


\noindent\emph{환원 불가능 형식들:}
\[
\renewcommand{\arraystretch}{1.25}
\begin{array}{@{}r@{\quad}c@{\qquad\qquad}r@{\quad}c@{\qquad\qquad}r@{\quad}c@{}}
\text{A)}&(\vartheta^2\nu x)^3&\text{B)}&a_\nu(j\nu x)&\text{C)}&N_\nu,\\[0.45em]
&\cdot&&\cdot&&\\[-0.1em]
&\theta_\nu(\vartheta^2\nu x)^3,&&a_\nu^2(j\nu x),&&
\end{array}
\]
\[
\renewcommand{\arraystretch}{1.25}
\begin{array}{@{}r@{\quad}c|c|c|c@{}}
&\cdot&(\theta Hu)&(\theta Hu)^2&\cdot\\
\text{D)}&(\vartheta\eta x)&\theta_\eta&H_\vartheta(\theta Hu),\ \theta_\eta(\theta Hu)&\theta_\eta(\theta Hu)^2\\
&(\vartheta\eta x)^2&H_\vartheta(\vartheta\eta x),\ \theta_\eta(\vartheta\eta x)&\theta_\eta^2&\theta_\eta^2(\theta Hu)\\
&\cdot&\theta_\eta(\vartheta\eta x)^2&\cdot&\cdot
\end{array}
\]

%
%
\noindent\emph{환원:}

\noindent A) $(\vartheta^2\nu x)^4$는 공식 (3.)a에 따라 분해된다.
\[
(a\A\widehat{\vartheta}x)^2=\frac12(\vartheta^2\nu x)^4-\frac35 f\cdot a_\nu^2(\nu\vartheta x)^2
=\frac13\A^2+f\cdot\A_\vartheta^2.
\]

\noindent B) $a_\nu(j\nu x)^2$은 공식 (9.)a에 따라 분해된다. 이때 \S\ 6 끝부분에 따라 $f\cdot j$를 $(\theta\theta'u)^2$으로 바꾸고 $\theta_\vartheta'$의 환원 가능성을 이용한다.
\[
\frac13a_\nu(j\nu x)^2=(\widehat{\theta\theta'}\nu x)^2\theta_\nu
=\theta_\nu^3\cdot\theta-\theta_\nu^2\theta_\nu'
=\theta_\nu^3\cdot\theta+\theta_\nu^2(\widehat{\vartheta\nu}\theta'u)
=\theta_\nu^3\cdot\theta-\frac23\theta_\nu^3\cdot\theta.
\]
형식 $a_\nu(j\nu x)^3$과 $a_\nu^2(j\nu x)^2$: 환원자 $a_j$와 $(j\nu x)^4$.
\[
a_\nu(j\nu x)^3=a_j(j\nu x)^3-(j\nu x)^3(j\nu\widehat{au}),
\]
\[
a_\nu^2(j\nu x)^2=a_j^2(j\nu x)^2-2a_j(j\nu x)^2(j\nu\widehat{au})+(j\nu x)^2(j\nu\widehat{au})^2.
\]

\noindent C) 형식열 $N_\nu^2$: 환원자 $\A_\nu^2$.
$(n\nu x)$와 $N_\nu(n\nu x)$는 \S\ 3에 따라 야코비 행렬식에 대한 전이연산으로서 분해된다.

%
%
\noindent D) \emph{환원의 개요:}

a) 형식열 $\theta_\eta^2H_\vartheta$: 환원자 $a_\eta^2(aHu)$. (더 높은 기호를 갖는 형식들로 이어진다.)

b) 형식열 $\theta_\eta H_\vartheta$: 이중 환원을 통한 환원자 $a_\eta^2(aHu)$.
(더 높은 기호를 갖는 형식들과 전체 형식열의 더 높은 형식들, 즉 형식열 $\theta_\eta^2$로 이어진다.) $\theta_\eta H_\vartheta$ 대신 형식열
\[
\theta_\eta(\vartheta\eta x)(\theta Hu)=\theta_\eta H_\vartheta+\theta_\eta^2
\]
을 고찰한다. 여기서 기호 $\theta$를 더 낮은 형식 $(abu)$로 바꾸면 형식열
\[
a_\eta^2(aHu)(abu)=\theta_\eta H_\vartheta+\frac32\theta_\eta^2\modu{s}
\]
의 이중 환원을 거쳐 끝형식
\[
a_\eta^3b_\eta(bHu)(abu)=\theta_\eta^3H_\vartheta+\frac12\theta_\eta^4
\]
에 이른다.

%
%
c) 형식열 $\theta_\eta^3$: 형식열 $\theta_\eta H_\vartheta$와 $\theta_\eta^2H_\vartheta$에 동시에 속하는 형식들, 즉 형식열 $\theta_\eta^2H_\vartheta$의 형식들을 이중 환원하여 얻는다. 한편으로는 더 높은 기호를 갖는 형식들로 환원하고, 다른 한편으로는 더 높은 기호를 갖는 형식들과 더 높은 형식열 $\theta_\eta^3$으로 환원한다.

d) 형식 $\theta_\eta^2(\vartheta\eta x)$, $\theta_\eta^2(\vartheta\eta x)^2$, $\theta_\eta^3(\vartheta\eta x)$: \S\ 9의 계산과 마찬가지로 환원자 $a$를 이용한 이중 환원으로 얻는다.

e) 형식 $\theta_\eta^2(\theta Hu)^2$와 $\theta_\eta^2(\vartheta\eta x)^2$: 환원자 $\A_\nu^2$와 $(a\A u)^4$를 이용하여 식 $\A_\nu^2(a\A u)^4$와 $(a\A\widehat{\nu}x)^4$을 이중 환원함으로써 얻는다.

f) 형식 $H_\vartheta$와 $H_\vartheta^2$: 분해 가능 형식이다. $H_\vartheta^2$의 경우 이는 식 $\A_\nu^2(a\A u)^2$의 이중 환원으로, 또는 곱 $(a_\nu^2)^2$, $a_\nu^2\cdot b_{\nu_1}$을 계의 형식들로 직접 계산함으로써 얻는다. ($(a_\nu)^2$을 이와 같이 계산하면 분해 가능 형식들 사이의 관계가 생긴다.)

g) 형식열 $\theta\cdot H$의 형식들이 두 가지 환원 방법을 허용하기 때문에, 즉 두 환원 가능한 형식열에 속하기 때문에 다음과 같이 더 높은 형식들이 환원된다.
\[
\begin{array}{rl}
g&\text{는 }\theta_\eta^2(\vartheta\eta x)^2\text{의 이중 환원으로: 환원 d)와 e)가 가능하다;}\\[0.3em]
s_\vartheta^2(\theta s u)&\text{는 }\theta_\eta^3(\vartheta\eta x)\text{의 이중 환원으로: 환원 c)와 d)가 가능하다;}\\[0.3em]
s_\vartheta^2(\theta s u)^2&\text{은 }\theta_\eta^2H_\vartheta^2\text{의 이중 환원으로: 환원 a)가 가능하며}\\
&\text{환원자 }\theta_\nu^2(\vartheta\nu x)^2\text{을 인수로 갖는다;}\\[0.3em]
s_\nu^2&\text{은 }\theta_\eta^3H_\vartheta\text{의 이중 환원으로: 환원 a)와 c)가 가능하다.}
\end{array}
\]
나머지 형식들의 독립성은 형식열 $\A_\nu^2(a\A u)^2\equiv\theta_\eta^2$의 이중 환원으로 증명된다.

%
%
\noindent\emph{환원의 실행:}

제시한 구성 법칙에 따르면 이중 환원에서 생기는 관계들이 서로 독립이므로, 환원 a), b), c), d)의 존재는 선험적으로 명백하다. 환원 e), f), g)에서는 항등식이 생기지 않는다는 것을 계산으로 증명해야 한다.

형식열 $\theta\cdot H$의 대각 성분에 대칭을 이루도록 형식 $\theta_\eta$까지 나아가면서, 환원 a), b), c), d)로 얻은 공식들도 일부 계산의 개요와 함께 제시한다. 이 공식들은 일부는 환원 e), f), g)에, 일부는 뒤에서 사용되기 때문이다.

\noindent 1) f)에 따른 $H_\vartheta$와 $H_\vartheta^2$. 다음과 같이 둔다\srcfn{*)}{등호 대신 쓴 기호 $\sim$은 수축 III과 IV에서 생기는 형식들보다 더 높은 형식에 $u_x$를 곱한 곱들을 생략했다는 뜻이다.}.
\[
\begin{aligned}
a_\nu\cdot b_{\nu_1}^2
&=\nu\cdot a_\nu b_\nu^2-(aHu)(bHu)b_\eta
\sim \nu\cdot a_\nu b_\nu^2-f\cdot a_\eta(aHu)^2-(abu)(aHu)b_\eta+(abu)^2b_\eta,\\
a_\nu^2\cdot b_{\nu_1}^2
&=b_{\nu_1}^2\{a_{\nu_1}u_\nu-(a\widehat{\nu\nu_1}u)\}^2
\sim \nu\cdot a_\nu^2b_\nu^2-2a_\eta b_\eta(aHu)(bHu)-\frac16(asu)^2(bsu)^2\\
&\sim \nu\cdot a_\nu^2b_\nu^2-2a_\eta^2(aHu)(abu)-\frac16(asu)^2(bsu)^2+(\vartheta\eta x)^2(\theta Hu)^2
\end{aligned}
\]
$\theta_\eta H_\vartheta$의 값을 대입하면 다음 공식들이 생긴다.
\[
\tag{12}
\begin{aligned}
\text{a)}\quad H_\vartheta&\sim 2a_\nu\cdot a_\nu^2+2\nu\cdot\theta_\nu(\vartheta\nu x)^2+2f\cdot a_\eta(aHu)^2-2\theta_\eta,\\
\text{b)}\quad H_\vartheta^2&\sim (a_\nu^2)^2+2\theta_\eta^2+\frac23(\theta su)^2+\frac1{12}s\cdot\nu-\frac19 i f\cdot\nu-\frac16f\cdot(asu)^4.
\end{aligned}
\]

%
%
\noindent 2) b)에 따른 $\theta_\eta H_\vartheta$:
\[
\begin{aligned}
a_\eta^2(aHu)(abu)
&\sim [a_\eta^2(aHu)]_{\widehat{bu}}+\frac12 f\cdot[a_\eta^2(aHu)^2]
=a_\eta b_\eta(aHu)(abu)\\
&\quad+a_\eta(\widehat{ab}\eta x)(abu)(aHu)\sim\theta_\eta(\vartheta\eta x)(\theta Hu)+\frac12\theta_\eta^2+\frac14(\nu\eta x)^2.
\end{aligned}
\]
공식 (5.)와 (11.)에 따라 다음이 성립한다.
\[
\theta_\eta H_\vartheta\sim-\frac32\theta_\eta^2-\frac14(\theta su)^2-\frac1{12}s\cdot\nu+\frac29 i f\cdot\nu.
\]

\noindent 3) $\theta_\eta H_\vartheta(\theta Hu)$는 b)에 따라 $\theta_\eta H_\vartheta$와 같은 방식으로 계산된다.
\[
\theta_\eta H_\vartheta(\theta Hu)\sim-\theta_\eta^2(\theta Hu)-\frac16(\theta su)^3.
\]

\noindent 4) b)에 따른 $\theta_\eta H_\vartheta(\vartheta\eta x)$; d)에 따른 $\theta_\eta^2(\vartheta\eta x)$(\S\ 9 참조):
\[
\theta_\eta H_\vartheta(\vartheta\eta x)\sim-\theta_\eta^2(\vartheta\eta x)-\frac16s_\vartheta(\theta su)
\sim\frac13(\vartheta\R x)-\frac16s_\vartheta(\theta su),
\]
\[
\theta_\eta^2(\vartheta\eta x)\sim\frac23a_\eta^3(abu)\sim-\frac13(\vartheta\R x).
\]

%
%
\noindent 5)
\[
\begin{array}{ccc}
\theta_\eta H_\vartheta^2\text{은 b)에 따라},&
\theta_\eta^2H_\vartheta\text{는 a)에 따라},&
\theta_\eta^2H_\vartheta\text{는 b)에 따라, 이로써 }\theta_\eta^3\text{은 c)에 따라},\\
a_\eta^2(aHb)(bHu)\text{에서};&
a_\eta^2(Hab)(abu)\text{에서};&
a_\eta^3(bHu)(abu)\text{에서}:
\end{array}
\]
\[
\theta_\eta H_\vartheta^2\sim-\frac32\theta_\eta^2H_\vartheta-\frac14s_\vartheta(\theta su)^2+\frac16s_\nu-\frac49ia_\nu
\sim-\theta_\varrho-\frac16s_\nu-\frac19ia_\nu,
\]
\[
\theta_\eta^2H_\vartheta\sim\frac23\theta_\varrho-\frac16s_\vartheta(\theta su)^2+\frac29s_\nu-\frac29ia_\nu,
\]
\[
\theta_\eta^2H_\vartheta\sim\frac23\theta_\varrho-\frac23\theta_\eta^3+\frac5{36}s_\nu\quad\text{이고, 따라서}
\]
\[
\theta_\eta^3\sim\frac14s_\vartheta(\theta su)^2-\frac18s_\nu+\frac13i\cdot a_\nu.
\]

%
%
\noindent 6) e)에 따른 $\theta_\eta^2(\theta Hu)^2$:
\[
\A_\nu^2(a\A u)^4=[(a\A u)^4]_{\nu^2}=-\frac{12}{5}\theta_\eta^2(\theta Hu)^2-\frac3{10}\sigma-\frac1{20}(\theta su)^4+\nu\cdot\varrho
\quad(\text{공식 (3.) c에 따라}),
\]
\[
\A_\nu^2(a\A u)^4=[\A_\nu^2]^4_{\widehat{u a}}=\frac35\theta_\eta^2(\theta Hu)^2+\frac15\sigma-\frac3{10}(\theta su)^4+\frac13ij
\quad(\text{공식 (10.) a에 따라}),
\]
\[
\theta_\eta^2(\theta Hu)^2=-\frac16\sigma+\frac1{12}(\theta su)^4-\frac19ij+\frac13\nu\cdot\varrho.
\]

%
%
\noindent 7) d)와 e)에 따른 $\theta_\eta^2(\vartheta\eta x)^2$. 이로부터 더 높은 형식 $g$를 환원한다.

\S\ 9와 유사하게 계산하면 다음이 성립한다.
\[
\begin{aligned}
\theta_\eta^2(\vartheta\eta x)^2
&=\frac12t\cdot f-\frac12a_\eta^2b_\eta-\frac1{12}g=\frac23t\cdot f-\frac23a_\eta^3b_\eta-\frac16g\\
&=-\frac16g-\frac13(\vartheta\varrho x)^2+\frac4{15}f\cdot a_\varrho^2+\frac8{15}f\cdot t\quad(\text{공식 (11.)에 따라}).
\end{aligned}
\]

%
%
위와 같은 방식으로 대응하는 계산을 하면 다음이 성립한다.
\[
(a\A\widehat{\nu x})^4=[(a\A u)^4]\widehat{\nu x}^4=\frac{21}{10}\theta_\eta^2(\vartheta\eta x)^2+\frac7{10}s_\vartheta^2-\frac3{10}g
\quad(\text{공식 (3.) c에 따라}),
\]
\[
\begin{aligned}
(a\A\widehat{\nu x})^4&=-\frac13i\A+6[\A_\nu^2]a^2-4[\A_\nu^3]a+\A_\nu^4\cdot f\\
&=-\frac{36}{5}\theta_\eta^2(\vartheta\eta x)^2-\frac95s_\vartheta^2-\frac35g-\frac35(\vartheta\varrho x)^2\\
&\quad+\frac53i\A+\frac{37}{25}f\cdot a_\varrho^2+\frac{74}{25}f\cdot t,\quad(\text{공식 (10.)에 따라}).
\end{aligned}
\]
따라서
\[
\begin{aligned}
\frac{93}{10}\theta_\eta^2(\vartheta\eta x)^2
&=-\frac3{10}g-\frac52s_\vartheta^2-\frac35(\vartheta\varrho x)^2\\
&\quad+\frac53i\A+\frac{37}{25}f\cdot a_\varrho^2+\frac{74}{25}f\cdot t,
\end{aligned}
\]
이고, g)에 따라 $\theta_\eta^2(\vartheta\eta x)^2$의 두 값을 비교하면
\[
\tag{13.)\srcfnmark{*)}}
0=g-2s_\vartheta^2+2(\vartheta\varrho x)^2+\frac43i\A-\frac45f\cdot a_\varrho^2-\frac85f\cdot t.
\]
\srcfntext{*)}{이는 서론의 주석에서 언급한 계수차수 6, 변수차수 6인 세 공변식 사이의 관계이다.}

%
%
\noindent 8) a)와 b)에 따른 $\theta_\eta^2H_\vartheta(\theta Hu)$, 이로부터 c)에 따른 $\theta_\eta^3(\theta Hu)$(인수 $u_x$를 갖는 형식들은 소멸한다):
\[
\theta_\eta^2H_\vartheta(\theta Hu)=-\frac16s_\vartheta(\theta su)^3-\frac13(\nu\varrho x),
\]
\[
\theta_\eta^2H_\vartheta(\theta Hu)=-\frac13\theta_\eta^3(\theta Hu)-\frac13(\nu\varrho x),
\]
\[
\theta_\eta^3(\theta Hu)=\frac12s_\vartheta(\theta su)^3.
\]

%
%
\noindent 9) a)와 b)에 따른 $\theta_\eta^2H_\vartheta(\vartheta\eta x)$, 이로부터 c)에 따른 $\theta_\eta^3(\vartheta\eta x)$. d)에 따른 $\theta_\eta^3(\vartheta\eta x)$, 이로부터 g)에 따른 더 높은 형식 $s_\vartheta^2(\theta su)$의 환원(인수 $u_x$를 갖는 형식들은 소멸한다):
\[
\theta_\eta^2H_\vartheta(\vartheta\eta x)=\frac13(atu),
\]
\[
\theta_\eta^2H_\vartheta(\vartheta\eta x)=\frac13(atu)+\frac13\theta_\eta^3(\vartheta\eta x)-\frac16s_\vartheta^2(\theta su),
\]
\[
\theta_\eta^3(\vartheta\eta x)=\frac12s_\vartheta^2(\theta su),
\]
\[
\theta_\eta^3(\vartheta\eta x)=\frac25\theta_\varrho(\vartheta\varrho x)+\frac15(atu):
\]
\srcnumdisplay{(14.)}{%
s_\vartheta^2(\theta su)=\frac45\theta_\varrho(\vartheta\varrho x)+\frac25(atu).
}

%
%
\noindent 10) a)에 따른 $\theta_\eta^2H_\vartheta^2$ 및 환원자 $\theta_\nu^2(\vartheta\nu x)^2$을 통한 $\theta_\eta^2H_\vartheta^2$, a)와 c)에 따른 $\theta_\eta^3H_\vartheta$, c)에 따른 $\theta_\eta^4$. 이로부터 g)에 따라 더 높은 형식 $s_\vartheta^2(\theta su)^2$과 $s_\nu^2$을 환원한다.
\[
\begin{aligned}
\text{a)에 따라}\quad
\theta_\eta^2H_\vartheta^2&=[a_\eta^2(aHu)^2]b^2+\frac13[a_\eta^4]_{\widehat{bu^2}}-\frac13H_\nu^2(\nu\eta x)^2\\
&=\frac49i\cdot a_\nu^2-\frac16s_\vartheta^2(\theta su)^2-\frac5{18}s_\nu^2+\frac13(atu)^2+\frac1{54}J\cdot u_x^2.
\end{aligned}
\]
\[
\begin{aligned}
\text{환원자 }\theta_\nu^2(\vartheta\nu x)^2\text{을 통해}\quad
\theta_\eta^2H_\vartheta^2
&=[\theta_\nu^2(\vartheta\nu x)^2]_{\nu_1^2}+\frac13[\theta_\nu^4]\widehat{\nu_1}x^2-\frac13s_\vartheta^2(\theta su)^2\\
&=\frac49i\cdot a_\nu^2-\frac16s_\nu^2-\frac13s_\vartheta^2(\theta su)^2+\frac13(\nu\varrho x)^2.
\end{aligned}
\]

%
%
\[
\begin{aligned}
\text{a)에 따라}\quad
\theta_\eta^3H_\vartheta&=-[a_\eta^2(aHu)]_{\widehat{bu}}b^2-\frac12[a_\eta^2(aHu)^2]b^2-\frac13[a_\eta^3]_{\widehat{bu^2}}b+\frac1{18}[a_\eta^4]_{\widehat{bu^2}}\\
&\quad+\frac12u_x\cdot[a_\eta^2(aHu)^2]b^3\\
&=-\frac29ia_\nu^2+\frac1{12}s_\nu^2+\frac4{15}\theta_\varrho^2-\frac7{90}(\nu\varrho x)^2+\frac1{30}(atu)^2+\frac2{27}i^2\cdot u_x^2-\frac1{36}J\cdot u_x^2.
\end{aligned}
\]
\[
\begin{aligned}
\text{c)에 따라}\quad
\theta_\eta^3H_\vartheta:\quad a_\eta^3(Hab)(bHu)&=\frac12(atu)^2
=\theta_\eta^2H_\vartheta(\theta Hu)(\vartheta\eta x)+\frac14H_\nu^2(\nu\eta x)^2\\
&\quad+\frac12\theta_\eta^2H_\vartheta^2
=\frac32\theta_\eta^2H_\vartheta^2+\theta_\eta^3H_\vartheta-u_x\cdot[a_\eta^2(aHu)^2]b^3\\
&\quad+\frac14H_\nu^2(\nu\eta x)^2,
\end{aligned}
\]
따라서
\[
\begin{aligned}
\theta_\eta^3H_\vartheta&=-\frac23i\cdot a_\nu^2+\frac5{12}s_\nu^2+\frac12(\nu\varrho x)^2-\frac12(atu)^2\\
&\quad+\frac4{27}i^2\cdot u_x^2-\frac1{12}J\cdot u_x^2.
\end{aligned}
\]

%
%
c)에 따라
\[
\theta_\eta^4=\frac49ia_\nu^2-\frac16s_\nu^2+\frac{16}{15}\theta_\varrho^2-\frac{14}{45}(\nu\varrho x)^2+\frac2{15}(atu)^2.
\]

%
%
g)에 따라 $\theta_\eta^2H_\vartheta^2$의 두 값에서 다음을 얻는다.
\par\smallskip\noindent
\begin{minipage}[c]{3.6em}
\centering (15.)
\end{minipage}%
\begin{minipage}[c]{\dimexpr\linewidth-3.6em\relax}
\centering
\(\displaystyle
\begin{aligned}
s_\vartheta^2(\theta su)^2
&=\frac23s_\nu^2-2(atu)^2+2(\nu\varrho x)^2-\frac19J\cdot u_x^2\\
&=\frac89i\cdot a_\nu^2+\frac8{15}\theta_\varrho^2-\frac{14}{15}(atu)^2+\frac{38}{45}(\nu\varrho x)^2-\frac4{27}i^2\cdot u_x^2.
\end{aligned}
\)
\end{minipage}
\par\smallskip

%
%
g)에 따라 $\theta_\eta^3H_\vartheta$의 두 값에서 다음을 얻는다.
\par\smallskip\noindent
\begin{minipage}[c]{3.6em}
\centering (16.)
\end{minipage}%
\begin{minipage}[c]{\dimexpr\linewidth-3.6em\relax}
\centering
\(\displaystyle
s_\nu^2=\frac43i\cdot a_\nu^2+\frac45\theta_\varrho^2+\frac85(atu)^2-\frac{26}{15}(\nu\varrho x)^2+\frac16J\cdot u_x^2-\frac29i^2\cdot u_x^2.
\)
\end{minipage}
\par\smallskip

%
%
공식 (16.)은 모듈 $(ss'u)^4$의 환원을 위한 기초를, 공식 (13.)은 $s$의 계를 환원하기 위한 기초를 제공한다(\S\ 17.).

\noindent\emph{결론:}

1) $a_\nu(j\nu x)^3$, $\theta_\eta H_\vartheta$, $\theta_\eta^2(\vartheta\eta x)$, $\theta_\eta^2(\theta Hu)^2$은 환원자이며, $a_\nu(j\nu x)^2$, $H_\vartheta$, $H_\vartheta^2$은 분해 가능 형식이다.

2) 계 II의 형식 $j(\nu)=g=(\nu\eta x)^4H_x^2$은 환원 가능하다.

3) 유일한 반변 형식 $g$가 환원 가능해지므로, $\nu$는 같은 계의 형식들과의 거듭제곱이나 곱 안에서 계 I과 전혀 수축되지 않는다.

$\theta$는 반공변식으로 볼 때만 곱이나 거듭제곱 안에서 수축되고, 이에 대응하여 $H$는 공변식으로 볼 때만 그러하다(\S\ 13 B 참조).

%
%
\begin{center}
\textbf{\S\ 12.}\quad\emph{계수차수 7의 형식들.}
\end{center}
\[
\text{다음의 수축: }\left\{\begin{array}{l}
\theta\cdot K\text{와 }\nu,\\
K,j,\A\text{와 }H,\\
f\text{와 }L,\sigma.
\end{array}\right.
\]

\noindent\emph{환원 불가능 형식들:}
\[
\renewcommand{\arraystretch}{1.35}
\begin{array}{@{}r@{\quad}c|c|c|c@{}}
&\cdot&\cdot&(KHu)^2&\cdot\\
&\cdot&K_\eta&\cdot&K_\eta(KHu)^2\\
\text{B)}&(k\eta x)^2&\cdot&K_\eta^2&\cdot\\
&(k\eta x)^3&H_k(k\eta x)^2,\ K_\eta(k\eta x)^2&\cdot&\cdot\\
&\cdot&K_\eta(k\eta x)^3&\cdot&\cdot
\end{array}\, ,
\]
\[
\renewcommand{\arraystretch}{1.35}
\begin{array}{@{}r@{\quad}c@{\qquad}c@{\qquad}c@{\qquad\qquad}r@{\quad}c@{}}
\text{C)}&(j\eta x)&H_j&\cdot&\text{D)}&(\A Hu)\\
&\cdot&H_j(j\eta x)&H_j^2,&&\A_\eta,
\end{array}
\]
\[
\renewcommand{\arraystretch}{1.35}
\begin{array}{@{}r@{\quad}c@{\qquad}c@{\qquad}c@{\qquad}c@{}}
&\cdot&(aLu)^2&(aLu)^3&\cdot\\
\text{E)}&a_l&\cdot&a_l(aLu)^2&a_l(aLu)^3\\
&\cdot&a_l^2&\cdot&\cdot
\end{array}\, .
\]

%
%
\noindent\emph{환원:}

\noindent A) $(\vartheta\cdot k,\nu,x)^4$은 분해 가능 형식 $(\vartheta^2\nu x)^4$에 대한 한 번의 전이연산으로서 분해된다.

\noindent B) 형식열 $H_kK_\eta$: 환원자 $H_\vartheta\theta_\eta$.

형식열 $K_\eta^2(KHu)$: 환원자 $a_\eta^2(aHu)$.

형식열 $K_\eta^2(k\eta x)$: 환원자 $\theta_\eta^2(\vartheta\eta x)$.

여기서 생기는 형식열 $K_\eta^3$의 이중 환원은 더 높은 형식열 $(j\eta x)^3$을 환원할 계기를 준다.

$(KHu)$, $(k\eta x)$, $K_\eta(k\eta x)$는 \S\ 3에 따른 분해 가능 형식이다.

$H_k(KHu)$와 $K_\eta(KHu)$는 분해 가능 형식 $K_\nu(k\nu x)$ 및 야코비 행렬식과 한 번 수축하여 생기므로 분해된다.
\[
K_\nu^2\equiv\theta_\nu^2(a\theta u)\equiv a_\nu^2(a\theta u)\modu{\nu^2}.
\]
$K_\nu^2$의 이중 표현에는 분해 가능 형식들 사이의 한 관계가 대응한다. 다음을 얻는다.

%
%
\[
H_k(KHu)=2K_{\nu_1}(\nu\nu_1k)=2\theta_{\nu_1}(\nu\nu_1\widehat{a\theta})
=2\theta_\nu^2\cdot a_\nu-a_\nu^2\cdot\theta_\nu+f\cdot\theta_\eta(\theta Hu)^2-f\cdot \nu\cdot\theta_\nu^3\quad(\text{곱 정리})\modu{(\nu^3)},
\]
\[
\begin{aligned}
K_\eta(KHu)&=-K_\nu^2(\nu\nu_1x)=-a_\nu^2(a\theta u)(\nu\nu_1x)=-\theta_\nu^2(a\theta u)(\nu\nu_1x)\\
&=a_\eta(aHu)^2\cdot\theta+a_\nu^2\cdot\theta_\nu
=-\theta_\eta(\theta Hu)^2\cdot f-\theta_\nu^2\cdot a_\nu:
\end{aligned}
\]
\par\smallskip\noindent
\begin{minipage}[c]{4.8em}
\centering (17.)\srcfnmark{*)}
\end{minipage}%
\begin{minipage}[c]{\dimexpr\linewidth-4.8em\relax}
\centering
\(\displaystyle
0=a_\nu^2\cdot\theta_{\nu_1}+\theta_\nu^2\cdot a_{\nu_1}+f\cdot\theta_\eta(\theta Hu)^2+\theta\cdot a_\eta(aHu)^2\modu{(\nu^3)}.
\)
\end{minipage}
\par\smallskip
\srcfntext{*)}{이 경우의 $f\cdot\theta_\nu^3\cdot\nu$처럼 반공변식 하나를 포함하는 형식들의 곱은 분해 가능 형식들 사이의 모든 관계에서 생략한다. 계 $s_{IV}$에 공식을 적용할 때에는 이러한 곱들이 아무 의미도 갖지 않기 때문이다.}
형식 $H_k$, $H_k(k\eta x)$, $H_k^2$, $H_k^2(k\eta x)$는 분해 가능 형식 $H_\vartheta$와 $H_\vartheta^2$와 한 번 수축하여 생기므로 분해된다(\S\ 3. 2), a)와 b) 참조).

%
%
다음이 성립한다.
\[
H_k=H_\vartheta(a\theta u)\sim[H_\vartheta]_{\widehat{ua}}-f\cdot H_\vartheta(\theta Hu).\qquad(\text{2) a)에서 }y=\widehat{uv}\text{로 특수화}).
\]
\[
H_k(k\eta x)=H_\vartheta a_\eta-H_\vartheta\theta_\eta\cdot f,
\qquad
H_\vartheta a_\eta=\frac54[H_\vartheta]a-\frac14H_\vartheta a_\vartheta
\quad(\S\ 3. 2)b),
\]
\[
H_\vartheta^2(a\theta u)=[H_\vartheta^2]_{\widehat{ua}},
\qquad
H_\vartheta^2a_\eta=[H_\vartheta^2]a=H_k^2(k\eta x)+H_\vartheta^2\theta_\eta\cdot f.
\]

%
%
\noindent C) 형식열 $(j\eta x)^3$: B)에 따른 형식열 $K_\eta^3$의 이중 환원으로 환원 가능하다. 다음과 같이 둔다.
\[
\begin{aligned}
0&=a_\eta^3(a\theta u)=\theta_\eta^3(a\theta u)
=\{\theta_\eta+(\widehat{a\theta}\eta x)\}^3(a\theta u)-\theta_\eta^3(a\theta u)\\
&=\frac12(j\eta x)^3\modu{(\nu^3,s,\R,t,i,J)}.\qquad(\text{\S\ 5 끝부분 참조.})
\end{aligned}
\]
유사한 설정은 형식열의 최종형까지 성립한다.
\[
0=H_\vartheta^2(\widehat{a\theta}\eta x)a_\eta^3
=H_\vartheta^2(\widehat{a\theta}\eta x)\theta_\eta^3=\frac12H_j^2(j\eta x)^4\modu{(\nu^3,s,\R,t,i,J)}.
\]
형식 $(j\eta x)^2$: 1) 분해 가능 형식 $H_\vartheta^2$에 대한 관계 ((12.) b)를 두 번 편극하여 분해된다.

2) 곱 $a_\nu\cdot\theta_\nu^3$을 직접 계산하여 분해된다. 이로부터 더 높은 형식 $(\A Hu)^2$이 분해된다. 다음을 얻는다.

%
%
\par\smallskip\noindent
\begin{minipage}[c]{4.8em}
\centering (18.)
\end{minipage}%
\begin{minipage}[c]{\dimexpr\linewidth-4.8em\relax}
\centering
\(\displaystyle
\left.\begin{array}{ll}
\text{a)}&\dfrac13(\A Hu)^2+\dfrac13(j\eta x)^2=\theta_\nu^2\cdot a_{\nu_1}^2,\\[0.6em]
\text{b)}&(j\eta x)^2=\dfrac43\theta_\nu^3\cdot a_{\nu_1},
\end{array}\right\}\modu{(\nu^3,s,\R,t,i,J)}
\)
\end{minipage}
\par\smallskip
다음과 같이 두었다.
\[
\begin{aligned}
H_\vartheta^2(a\theta u)^2=\frac13(\A Hu)^2
&=2\theta_\eta^2(a\theta u)(aHu)+\theta_\nu^2\cdot a_{\nu_1}^2\quad(\text{곱 정리})\\
&=(a\theta u)(aHu)\{a_\eta-(\widehat{a\theta}\eta x)\}^2+\theta_\nu^2\cdot a_{\nu_1}^2,
\end{aligned}
\]
\[
\begin{aligned}
a_{\nu_1}\cdot\theta_\nu^3
&=-(a\widehat{\nu\nu_1}u)\theta_\nu^2(\theta\widehat{\nu\nu_1}u)-\frac12(a\widehat{\nu\nu_1}u)\theta_\nu\theta_{\nu_1}(\theta\widehat{\nu\nu_1}u)\\
&=-\frac32\theta_\eta^2(\theta Hu)(aHu)=-\frac32\theta_\eta^2(\theta Hu)(a\theta u)\\
&=-\frac32(a_\eta-(\widehat{a\theta}\eta x))^2\{(aHu)-(a\theta u)\}(a\theta u).
\end{aligned}
\]

%
%
\noindent D) 형식열 $\A_\eta(\A Hu)$: 환원자 $\A_\nu^2$.

$(\A Hu)^2$은 C)에 따른 분해 가능 형식이다.

\noindent E) 형식열 $a_l^2(aLu)$: 환원자 $a_\eta^2(aHu)$.
형식 $(aLu)$와 $a_l(aLu)$: \S\ 3에 따라 야코비 행렬식에 대한 전이연산으로서 분해된다.

\noindent F) 형식열 $a_\sigma^3$: 환원자 $a_\eta^3$.

형식 $a_\sigma^2$와 $a_\sigma^3$: 대응하는 식
\[
H_\nu a_\nu^3\quad\text{및}\quad H_\nu a_\eta^3,
\qquad
H_\nu a_\nu^3a_\eta\quad\text{및}\quad H_\nu a_\eta^4
\]
을 이중 환원하여 얻는 환원자 $a_\nu^3$과 $a_\eta^3$ (\S\ 10. C)의 $\A_\nu^2$와 $\A_\nu^3$의 환원 참조).

%
%
이로부터 (16.)을 고려하여 더 높은 형식 $a_\nu s_\nu(asu)^2$, $a_\nu^2s_\nu(asu)^2$ 및 기호 $\R,t(a_\nu,a_\varrho^2)$로 된 형식들에 대한 환원을 얻는다.
\par\smallskip\noindent
\begin{minipage}[c]{4.8em}
\centering (19.)
\end{minipage}%
\begin{minipage}[c]{\dimexpr\linewidth-4.8em\relax}
\centering
\(\displaystyle
\left.\begin{array}{rcl}
a_\nu s_\nu(asu)^2&=&\dfrac13 i\theta_\nu^2-\dfrac1{18}iH,\\[0.5em]
a_\nu^2s_\nu(asu)^2&=&0,
\end{array}\right\}\modu{(\R,t)}
\)
\end{minipage}
\par\smallskip
형식 $a_\sigma$: (12.)b를 편극하여 분해되며, 더 짧게는 다음과 같이 두고 곱 $a_\nu^2\cdot\theta_\nu^3$을 직접 전개하여 분해된다.

%
%
\[
a_\nu^2\cdot\theta_{\nu_1}^3=a_\nu^2\theta_{\nu_1}\{a_{\nu_1}-(\widehat{a\theta}\nu_1x)\}^2
=-2a_\eta(aHu)(a\theta H)(\widehat{a\theta}\eta x)+(\widehat{a\theta}\nu_1x)^2a_\nu^2\theta_{\nu_1}
\]
그리고 $H_j(j\eta x)^2$의 환원(C)을 고려하면 다음을 얻는다.
\par\smallskip\noindent
\begin{minipage}[c]{4.8em}
\centering (20.)
\end{minipage}%
\begin{minipage}[c]{\dimexpr\linewidth-4.8em\relax}
\centering
\(\displaystyle a_\sigma=2a_\nu^2\cdot\theta_{\nu_1}^3\modu{(s,\R,t,i)}.\)
\end{minipage}
\par\smallskip

%
%
\noindent\emph{결론:}

1) $(j\eta x)^3$, $\A_\eta(\A Hu)$, $a_\sigma^2$은 환원자이고, $(j\eta x)^2$, $(\A Hu)^2$, $a_\sigma$는 분해 가능 형식이다.

2) $f$는 반변 형식들과의 곱 안에서만 $L$과 수축되며, $f$와 따라서 $K$ 및 $N$은 $\sigma$와 전혀 수축되지 않고, 그러므로 $(\nu\sigma x)$와도 수축되지 않는다. $j$는 반변 형식들과의 곱 안에서만 $H$와 $L$에 수축되며, $\A$는 곱 안에서 $H$나 $L$과 전혀 수축되지 않는다($\A_\eta(\A Hu)$와 $(\A Hu)^2$으로부터 따른다). 마찬가지로 $\sigma$와 따라서 $(\nu\sigma x)$는 계 I의 어떤 형식과도 곱 안에서는 수축되지 않는다. 계 II의 곱들 가운데에는 \S\ 11의 결론들을 고려하면 $H$의 거듭제곱과 $H$와 $L$의 곱만 남는다.

%
%
\begin{center}
\textbf{\S\ 13.}\quad\emph{계수차수 8의 형식들(계 III).}
\end{center}
\[
\text{다음의 수축: }\left\{\begin{array}{l}
\A\cdot j\text{와 }\nu,\\
\theta^2,
 f\cdot j,
 N\text{과 }H,\\
\theta\text{와 }L,\sigma.
\end{array}\right.
\]


\noindent\emph{환원 불가능 형식들:}
\[
\renewcommand{\arraystretch}{1.28}
\begin{array}{@{}r@{\quad}c@{\qquad\qquad}r@{\quad}c@{\qquad}c@{}}
\text{A)}&\A_\nu(j\nu x),&\text{B)}&(\vartheta^2\eta x)^3&\cdot\\
&&&(\vartheta^2\eta x)^4&H_\vartheta(\vartheta^2\eta x)^3,\ \theta_\eta(\vartheta^2\eta x)^3,
\end{array}
\]
\[
\renewcommand{\arraystretch}{1.28}
\begin{array}{@{}r@{\quad}c@{\qquad}c@{\qquad\qquad}r@{\quad}c@{}}
\text{C)}&(aHu)H_j&\cdot&\text{D)}&H_\eta;\ N_\eta,\\
&\cdot&a_\eta(aHu)H_j,&&
\end{array}
\]
\[
\renewcommand{\arraystretch}{1.30}
\begin{array}{@{}r@{\quad}c@{\qquad}c@{\qquad}c@{\qquad}c@{\qquad}c@{}}
&\cdot&\cdot&(\theta Lu)^2&(\theta Lu)^3&\cdot\\
&\cdot&\theta_l&\cdot&L_\vartheta(\theta Lu)^2,\ \theta_l(\theta Lu)^2&\theta_l(\theta Lu)^3\\
\text{E)}&(\vartheta lx)^2&\cdot&\theta_l^2&\cdot&\cdot\\
&\cdot&\theta_l(\vartheta lx)^2&\cdot&\cdot&\cdot
\end{array}
\]

%
%
\noindent\emph{환원:}

\noindent A) $\A_\nu(j\nu x)^3$: 환원자 $a_\nu(j\nu x)^3$.

$\A_\nu(j\nu x)^2$은 환원자 $\theta_{\vartheta'}$를 고려할 때 \S\ 3. 2)b에 따라 분해 가능 형식 $a_\nu(j\nu x)^2$에 의해 환원 가능하다. 실제로 \S\ 11 B)에 따라 다음이 성립한다.
\[
\frac3{10}\A_\nu(j\nu x)^2+\frac35a_\nu(j\nu\vartheta)a_\vartheta(j\nu x)
+\frac1{10}a_\nu(j\nu\vartheta)^2
=\frac35\theta_\nu^3\cdot\theta_{\vartheta''}^{\prime2}+\frac25\theta_\nu^3\theta_{\vartheta''}\theta_{\vartheta''}'.
\]
(환원자 $a_j$와 $\theta_{\vartheta'}$).

%
%
\noindent B) 형식 $H_\vartheta(\vartheta'\eta x)^2$, $H_\vartheta^2(\vartheta'\eta x)$, $H_\vartheta^2(\vartheta'\eta x)^2$은 \S\ 3. 2)b)와 다음 관계에 따라 분해 가능 형식 $H_\vartheta$와 $H_\vartheta^2$, 각각 $H_\vartheta a_\eta$와 $H_\vartheta^2a_\eta$와의 연속적인 한 번의 수축으로 생기므로 분해된다.
\[
H_\vartheta(\vartheta'\eta x)^2=2H_\vartheta a_\eta^2\cdot f-2H_\vartheta a_\eta b_\eta.
\]

%
%
\noindent C) 형식열 $a_\eta(j\eta x)^2$: 환원자 $a_j$와 $(j\eta x)^3$.

형식 $a_\eta(j\eta x)$은 분해된다. 환원자 $a_j$와 \S\ 3. 2)a에 따른 분해 가능 형식 $(j\eta x)^2$과의 한 번의 수축을 이용한다(특수화 $y=uv$).

형식 $a_\eta H_j$는 다음 두 방식으로 분해된다. 1) 분해 가능 형식 $a_\nu(j\nu x)^2$과 한 번 수축한다.

2) 분해 가능 형식 $(j\eta x)^2$과 특수하게 두 번 수축한다. 이로부터 분해 가능 형식들 사이의 관계가 생긴다.

$a_\nu(j\nu x)^2$으로부터:
\[
0=\theta_\nu^3\cdot\theta_\nu+2a_\eta H_j\modu{(s,\R,t,i,J)}.
\]
$(j\eta x)^2$으로부터:
\[
0=\theta_\nu^3\cdot\theta_\nu-\frac32a_\eta H_j\modu{(s,\R,t,i,J)}
\]
(반공변식과의 곱은 생략한다). 다음과 같이 둔다.
\[
0=\frac12a_\nu(abu)^2\theta_\nu^3
+\frac12a_\nu(abu)\theta_{\nu_1}^3(\theta bu)
-\frac34(\widehat{j\eta}bu)^2+\frac34H_j(\widehat{j\eta}bu)
-\frac18f\cdot H_j^2
\]
그리고 $a_\nu(abu)(\theta bu)\theta_{\nu_1}^3$에서 기호들을 교환한다.

%
%
\noindent D) 형식열 $N_\eta(NHu)$: 환원자 $\A_\nu^2$.

$(NHu)$, $(n\eta x)$, $N_\eta(n\eta x)$, $H_n(NHu)$는 분해 가능 형식이며(\S\ 12 B 참조), $(NHu)^2$은 형식 $(\A Hu)^2$과 한 번 수축함으로써 분해된다.

%
%
\noindent E) 형식열 $\theta_lL_\vartheta$, $\theta_l^2(\theta Lu)^2$, $\theta_l^2(\vartheta lx)$:

환원자: $\theta_\eta H_\vartheta$, $\theta_\eta^2(\theta Hu)^2$, $\theta_\eta^2(\vartheta\eta x)$.

형식 $(\theta Lu)$, $(\vartheta lx)$, $L_\vartheta$, $L_\vartheta(\theta Lu)$, $\theta_l(\theta Lu)$, $L_\vartheta(\vartheta lx)$, $\theta_l(\vartheta lx)$, $L_\vartheta^2$, $L_\vartheta^2(\theta Lu)$, $\theta_l^2(\theta Lu)$는 분해된다. (\S\ 12 B의 분해 가능 형식들과 쌍대 관계에 있다.)

%
%
\noindent F) 형식열 $\theta_\sigma(\vartheta\sigma x)$: 환원자 $a_\sigma^2$.

형식 $(\vartheta\sigma x)$와 $(\vartheta\sigma x)^2$은 분해 가능 형식 $a_\sigma$와 한 번 수축하여 생기므로 분해된다. 두 번 수축하여 생긴 $\theta_\sigma$도 다음 형식열 때문에 분해된다.
\[
a_\nu^2(a\theta u)\theta_{\nu_1}^3\equiv H_j^2(j\eta x)\equiv0\;\text{법 }(s,\R,t,i,J):
\]
\[
\theta_\sigma=a_\sigma(abu)^2=a_\nu^2(abu)^2\cdot\theta_\nu^3.
\]

%
%
\noindent\emph{결론:}

$\theta^3$ 또는 $\theta^2K$는 $H$와 수축되지 않는다. $\theta$는 반공변식으로 볼 때만 거듭제곱이나 곱 안에서 $L$과 수축되며, $\sigma$와 $(\nu\sigma x)$과는 전혀 수축되지 않는다.

$N$은 곱 안에서 계 II의 어떤 형식과도 수축되지 않으며, 계 II 형식들의 거듭제곱이나 곱과도 마찬가지다. 앞 절들의 결론을 고려하면 계 I에 남는 곱은 $f\cdot j$, $\A\cdot j$, $\theta$의 거듭제곱, 그리고 $\theta\cdot K$의 곱뿐이다.

%
%
\begin{center}
\textbf{\S\ 14.}\quad\emph{계수차수 9의 형식들(계 III).}
\end{center}
\[
\text{다음의 수축: }\left\{\begin{array}{l}
\theta\cdot K\text{와 }H,\\
K,j,\A\text{와 }L,\\
j,\A\text{와 }\sigma,\\
f\text{와 }H^2.
\end{array}\right.
\]


\noindent\emph{환원 불가능 형식들:}
\[
\renewcommand{\arraystretch}{1.35}
\begin{array}{@{}r@{\quad}c@{\qquad}c@{\qquad\qquad}r@{\quad}c@{\qquad}c@{\qquad}c@{\qquad}c@{\qquad}c@{}}
&&&&\cdot&\cdot&\cdot&(KLu)^3&\cdot\\
\text{A)}&(\vartheta\cdot k,\eta,x)^4&\cdot&\text{B)}&\cdot&\cdot&\cdot&\cdot&K_l(KLu)^3\\
&\cdot&H_k(\vartheta\cdot k,\eta,x)^4,&&\cdot&\cdot&K_l^2&\cdot&\cdot\\
&&&& (klx)^3&\cdot&\cdot&\cdot&\cdot\\
&&&&\cdot&K_l(Klx)^3&\cdot&\cdot&\cdot,
\end{array}
\]
\[
\renewcommand{\arraystretch}{1.35}
\begin{array}{@{}r@{\quad}c@{\qquad}c@{\qquad}r@{\quad}c@{\qquad}r@{\quad}c@{\qquad}r@{\quad}c@{\qquad}c@{}}
\text{C)}&L_j&\cdot&\text{D)}&\A_l,&\text{E)}&(j\sigma x),&\text{G)}&(aH^2u)^3&(aH^2u)^4\\
&\cdot&L_j^2,&&&&&&\cdot&a_\eta(aH^2u)^3.
\end{array}
\]

%
%
\noindent\emph{환원:}

A) 형식 $H_\vartheta(k\eta x)^3$, $H_\vartheta^2(k\eta x)^2$, $H_\vartheta^2(k\eta x)^3$은 분해 가능 형식들과 연속적으로 한 번 수축하여 생기므로 분해된다.

B) 형식열 $K_lL_k$, $K_l^2(KLu)$, $K_l^2(klx)$:

환원자: $\theta_\eta H_\vartheta$, $a_\eta^2(aHu)$, $\theta_\eta^2(\vartheta\eta x)$.

형식 $K_l$, $K_l(KLu)$, $K_l(klx)$, $(KLu)^2$, $(klx)^2$은 야코비 행렬식에 대한 전이연산으로서 분해된다(\S\ 3).

형식 $L_k$, $L_k(KLu)$, $L_k(KLu)^2$, $L_k(klx)$, $L_k(klx)^2$, $L_k^2$, $L_k^2(KLu)$, $L_k^2(klx)$, $L_k^3$, $K_l(KLu)^2$, $K_l(klx)^2$은 다음 분해 가능 형식들과 한 번 수축하여 생기므로 분해된다.
\[
L_\vartheta,
\ H_k(KHu),
\ L_\vartheta(\vartheta lx),
\ H_k^2,
\ L_\vartheta^2,
\ L_\vartheta^2(\theta Lu),
\ K_\eta(KHu),
\ \theta_l(\vartheta lx)
\]
이는 \S\ 3. 2), a)와 b)에 따른다.

%
%
C) 형식열 $(jlx)^3$: 환원자 $(j\eta x)^3$.

형식 $(jlx)^2$과 $L_j(jlx)$은 분해 가능 형식 $(j\eta x)^2$과 한 번 수축하여 생긴다.

%
%
D) 형식열 $\A_l(\A Lu)$: 환원자 $\A_\nu^2$.

형식 $(\A Lu)^2$, $(\A Lu)^3$: 분해 가능 형식 $(\A Hu)^2$과 한 번 수축한다.

%
%
E) 형식열 $(j\sigma x)^2$: $j$를 형식열의 더 낮은 형식들로 바꾸어 얻는 환원자 $a_\sigma^2$(\S\ 5):
\[
a_\sigma^2(a\theta u)^2=\frac13(j\sigma x)^2\ldots,
\qquad
 a_\sigma^2(a\theta\sigma x)^2=\frac13(j\sigma x)^4.
\]

%
%
F) 형식열 $\A_\sigma^2$: 환원자 $a_\sigma^2$.

형식 $\A_\sigma$: $\A$를 형식열의 더 낮은 형식들로 바꾸어 얻는 환원자 $a_\sigma^2$:
\[
a_\vartheta a_\sigma^2=\frac23\A_\sigma\modu{\nu}.
\]

%
%
\noindent\emph{결론:} 형식 $j$만 $\sigma$와 $(\nu\sigma x)$과 수축된다.

$\A_l$에 대응하는 형식 $N_l$이 분해되므로 $N$은 $L$과 수축되지 않는다.

%
%
\begin{center}
\textbf{\S\ 15.}\quad\emph{계수차수 10의 형식들(계 III).}
\end{center}
\[
\text{다음의 수축: }\left\{\begin{array}{l}
\theta^2,
\ f\cdot j\text{와 }L,\\
\theta\text{와 }H^2.
\end{array}\right.
\]
\noindent\emph{환원 불가능 형식들:}
\[
\begin{array}{c@{\quad}l@{\qquad}c@{\quad}l}
\text{A)}&\begin{array}{cc}(\vartheta^2lx)^4&\cdot\\ \cdot&
\theta_l(\vartheta^2lx)^4,\end{array}&
\text{B)}&\begin{array}{cc}(aLu)^2L_j&\cdot\\ \cdot&
a_l(aLu)^2L_j,\end{array}\\[0.8em]
\text{C)}&\begin{array}{ccc}(\theta H^2u)^3
&(\theta H^2u)^4&\cdot\\ \cdot
&H_\vartheta(\theta H^2u)^3,
&\theta_\eta(\theta H^2u)^3.\end{array}
\end{array}
\]

\noindent\emph{환원:}

A)와 B): 나머지 형식들은 분해 가능 형식과 한 번 수축하여 분해된다.

C) 나머지 형식들은 \S\ 13 B)에 따라 쌍대적으로 반대되는 고찰을 통해 분해된다.

\noindent\emph{결론:}

$\theta^2\cdot K$는 $L$과 수축되지 않으며, 이에 대응하여 $H^2L$은 $K$와 수축되지 않는다. $H^3$과 $H^2L$은 $\theta$와 수축되지 않는다. 이로써 \S\ 7의 수축 도식이 증명된다.

%
%
\begin{center}
\textbf{\S\ 16.}\quad\emph{계수차수 11, 12, 13, 15의 형식들(계 III).}
\end{center}

\noindent\emph{계수차수 11.}
\[
\text{다음의 수축: }\left\{\begin{array}{l}
\theta\cdot K\text{와 }L,\\
K,j\text{와 }H^2,\\
j\text{와 }(\nu\sigma x),\\
f\text{와 }H\cdot L.
\end{array}\right.
\]
\noindent\emph{환원 불가능 형식들:}
\[
\begin{array}{c@{\quad}l@{\qquad}c@{\quad}l}
\text{A)}&\begin{array}{cc}(\vartheta\cdot k,l,x)^5&\cdot\\ \cdot&
\theta_l(\vartheta\cdot k,l,x)^5,\end{array}&
\text{B)}&\begin{array}{cc}(KH^2u)^4&\cdot\\ \cdot&
K_\eta(KH^2u)^4,\end{array}\\[0.8em]
\text{C)}&(\nu\sigma j),&
\text{D)}&(a,H\cdot L,u)^4.
\end{array}
\]

%
%
\noindent\emph{계수차수 12.}
\[
\text{다음의 수축: }\left\{\begin{array}{l}
\theta^3\text{와 }L,\\
\theta\text{와 }H\cdot L.
\end{array}\right.
\]
\noindent\emph{환원 불가능 형식들:}
\[
\text{E)}\ (\vartheta^3lx)^6,
\qquad
\text{F)}\ (\theta,H\cdot L,u)^4,
\qquad L_\vartheta(\theta,H\cdot L,u)^4.
\]

\noindent\emph{계수차수 13.}
\[
\text{다음의 수축: }K\text{와 }H\cdot L.
\]
\noindent\emph{환원 불가능 형식들:}
\[
\text{G)}\ (K,H\cdot L,u)^5,
\qquad K_\eta(K,H\cdot L,u)^5.
\]

\noindent\emph{계수차수 15.}
\[
\text{다음의 수축: }K\text{와 }H^3.
\]
\noindent\emph{환원 불가능 형식:}
\[
\text{H)}\ (KH^3u)^6.
\]

%
%
\noindent\emph{환원:}

형식들 $H_{j^3}^2$, $H_{j^4}^2$. 다음 관계식에서 얻는 환원자는 $L_j^3$이다:
\[
(\nu lx)L_x^3=-\frac12H^2\modu{(\sigma,s)}.
\]
나머지 형식들은 앞 절들의 고찰에 따라 분해되거나 환원 가능하다.

\noindent\emph{결론:}

\S\ 7의 수축 도식과 \S\S\ 8부터 15까지의 결론들에 따라, 이로써 계 III의 환원 불가능 형식들(117개 형식)이 모두 망라되었다.

\clearpage

%
%
\begin{center}
{\Large\bfseries 제 IV장.\quad $(\varrho,t)$를 법으로 하는 상대 완전계.}
\end{center}

\begin{center}
\textbf{\S\ 17.}\quad\emph{모듈 $(ss'u)^4$와 $s$의 계의 환원.}
\end{center}

%
%
바로 다음의 더 높은 상대 완전계를 구성하려면, \S\ 7과 유사하게 다음 모듈에 관하여 $s$의 계를 구성하고
\[
\nu(s)=(ss'u)^4
\]
$s$를 법으로 하는 상대 완전계의 형식들과 수축해야 한다. 더 높은 모듈로서 모듈 $(\varrho,t)$를 도입하면 다음이 성립함을 보이고자 한다.
\[
\begin{array}{ll}
\text{A)}&\text{모듈 }(ss'u)^4\text{가 환원 가능해지고,}\vphantom{\dfrac11}\\
\text{B)}&s\text{의 계는 유일한 형식 }s\text{로 이루어진다.}
\end{array}
\]

%
%
A) 모듈 $(ss'u)^4$의 환원은 \S\ 11의 공식 (16.)으로 구한 환원자 $s_\nu^2$에서 즉시 얻어진다. 다음 식에서
\[
 s_\nu^2=\frac43 i a_\nu^2+\frac45\theta_\varrho^2+\frac85(atu)^2-
 \frac{26}{15}(\nu\varrho x)^2+\frac16J\cdot u_x^2-
 \frac29i^2\cdot u_x^2
\]
$x$를 $\nu_1$에 따라 편극하면 다음을 얻는다.
\srcnumdisplay{(21.)\srcfnmark{*)}}{%
\begin{aligned}
(ss'u)^4&=-\frac23J\cdot \nu+6s_\nu^2s_{\nu_1}^2\\
&=\frac{24}{5}\theta_\nu^2\theta_\varrho^2+
\frac{48}{5}a_\nu^2(atu)^2-
\frac{52}{5}H_\varrho^2+
\frac43i\cdot(asu)^4+\frac13J\cdot \nu-
\frac49i^2\cdot \nu,
\end{aligned}
}
이로써 모듈 $(ss'u)^4$가 환원된다.

%
%
\srcfntext{*)}{공식 (21.)에서, 공식 (10.)c를 적용하면 서론의 주석에서 언급한 계수차수 9의 세 불변식 사이의 관계가 나온다:
\srcnumdisplay{(22.)}{%
\begin{aligned}
(ass')^4&=\frac{24}{5}\theta_\nu^2\theta_\varrho^2a_\nu^2a_\vartheta^2+
\frac{48}{5}a_\nu^2b_\nu^2(tab)^2-
\frac{52}{5}H_\varrho^2a_\eta^4+
\frac53J\cdot i-\frac49i^3,\\
(ass')^4&=\frac{24}{5}a_\varrho^2a_{\varrho'}^2-
\frac{12}{5}t_\varrho^2+\frac53J\cdot i-\frac49i^3.
\end{aligned}
}
}

%
%
B) 다음 형식의 환원과
\[
Z=(ss'u)^2=\theta(s)
\]
그에 따른 $s$의 계의 환원은, \S\ 7의 공식 (8.)b에 따라 형식 $Z$를 더 낮은 형식 $g_\nu^2$으로 바꾸고 다음 관계를 고려함으로써 얻어진다.
\[
 s_\eta^2=s_\nu^2(\nu\nu_1x)^2-\frac13Z.
\]
형식 $g_\nu^2$은 \S\ 11의 공식 (13.)에 따라 환원자 $s_\nu^2$를 인수로 갖는다. 공식 (8.)과 (16.)에 따라 다음이 성립한다.
\[
 g_\nu^2=-Z+\frac{14}{5}ia_\eta^2+\frac{14}{15}i(asu)^2
 \modu{(\varrho,t)},
\]
그리고 공식 (13.), (14.), (15.), (16.)에 따라 다음이 성립한다.
\[
\begin{aligned}
g_\nu^2&=2[s_\vartheta^2]_{\nu^2}-\frac43 i\A_\nu^2
      =2s_\vartheta^2s_\nu^2-\frac65[s_\vartheta^2(\theta su)^2]\widehat{\nu x}^{2}
        -\frac43i\A_\nu^2\\
&=\frac45i\cdot a_\eta^2-\frac25i\cdot(asu)^2+\frac13J\cdot\theta
\modu{(\varrho,t)}:
\end{aligned}
\]
따라서 다음을 얻는다.
\srcnumdisplay{(23.)}{%
Z=2i\cdot a_\eta^2+\frac43i\cdot(asu)^2-\frac13J\cdot\theta
 \modu{(\varrho,t)}.
}

%
%
따라서 $((ss'u)^4,\varrho,t)$를 법으로 하는 상대 완전계는 $(\varrho,t)$를 법으로 하는 상대 완전계로 넘어가며, 이 계를 알려진 두 이차형식의 계와 전이연산하면 절대 완전계가 생긴다. $(\varrho,t)$를 법으로 하는 상대 완전계를 구성하기 위해서는, 공식 (23.)에 따라 $s$의 계가 형식 $s$로 환원되므로, $(s,\varrho,t)$를 법으로 하는 상대 완전계를 $s$ 및 $s$의 거듭제곱들과 전이연산해야 한다. $s$의 거듭제곱들은 계 I과만 수축하게 됨을 보일 것이다.

%
%
같은 계수차수와 같은 차수를 지닌 형식들 가운데, $(s,\varrho,t)$를 법으로 하는 상대 완전계의 더 높은 형식들을 $s$와 수축하여 생긴 형식들을 \glqq{}더 높은 형식들\grqq{}로 정의한다. 이때 $s$와의 수축은 원래 형식들의 편극으로만 간주한다. 이로써 개별 형식들의 배열은 유일하게 정해진다(\S\ 1, 형식열 2 참조). 예를 들면 다음과 같다.
\[
(\theta_\eta(\theta Hu),s,u)^2\text{는 }s_\eta((\theta Hu)^2,s,u)\text{보다 더 높고}\srcfnmark{*)},
\]
\[
(\theta_\eta(\theta Hu)^2,s,u)\text{는 }(\theta_\eta(\theta Hu),s,u)^2\text{보다 더 높다.}
\]

%
%
생기는 계를 다음 네 부분계로 나눈다.
\[
\begin{array}{ll}
\text{계 }s_I:&s\text{ 및 }s\text{의 거듭제곱들을 계 I과 수축하여 생긴다,}\\
\text{계 }s_{II}:&s\text{를 계 II와 수축하여 생긴다,}\\
\text{계 }s_{III}:&s\text{를 계 III의 개별 형식들(곱은 제외)과 수축하거나,}\\
&\text{계 I과 II에서 각각 한 형식씩 취한 곱과 수축하여 생긴다,}\\
\text{계 }s_{IV}:&s\text{를 계 I과 III, II와 III, III과 III에서 각각 한 형식씩}\\
&\text{취한 곱과 수축하여 생긴다.}
\end{array}
\]
또한 이제부터 모든 공식은 $(\varrho,t)$를 법으로 이해하고, 공식 (27.)부터는 $(\varrho,t,i,J,\text{ 더 높은 형식들})$을 법으로 이해한다. 이를 매번 명시하지는 않는다.

%
%
환원을 위하여 앞서 얻은 공식들을 모아 적는다.
\srcnumdisplay{(24.)}{%
\begin{aligned}
s_\nu^2&=\frac43i\,a_\nu^2+\frac16J\cdot u_x^2-\frac29i^2\cdot u_x^2,\\
s_\nu^3&=\frac13J\cdot u_x,
& s_\nu^4&=J,\\
g&=2s_\vartheta^2-\frac43i\Delta,\\
s_\vartheta^2(\theta su)&=0,\\
s_\vartheta^2(\theta su)^2&=\frac89i\cdot a_\nu^2-\frac4{27}i^2\cdot u_x^2,\\
Z&=2i\cdot a_\eta^2+\frac43i(asu)^2-\frac13J\cdot\theta,\\
g_\nu&=-s_\eta+u_x\left\{2ia_\eta^2-\frac23i(asu)^2+\frac23J\cdot\theta\right\},\\
g_\nu^2&=\frac45ia_\eta^2-\frac25i(asu)^2+\frac13J\cdot\theta,
&g_\nu^3&=0,\quad g_\nu^4=0.
\end{aligned}
}

%
%
\srcfntext{*)}{간결하게 쓰기 위하여, 완전한 전이연산 $[(\theta Hu)^2]_{\widehat{su}}s$ 대신 전이연산의 항을 계의 형식으로 채택한다. 공식을 세울 때에는 모든 전이연산을 그 선두항으로 바꾼다. 예를 들어 $(\theta_\eta(\theta Hu),s,u)^2$을 $\theta_\eta(\theta Hu)(\theta su)^2$으로 바꾼다.}

%
%
\noindent\emph{결론:}
\[
 s_\nu^2,\quad s_\vartheta^2(\theta su),\quad s_\vartheta^2s_\nu,\quad s_\vartheta^2\theta_\nu
\]
은 환원자이다.

%
%
\begin{center}
\textbf{\S\ 18.}\quad\emph{계 $s_I$.}
\end{center}

\[
\text{다음의 수축: }\left\{\begin{array}{l}
s\text{와 } f;\ \theta;\ K,j,\Delta;\ f\cdot j,N;\ \Delta\cdot j,\\
s^2\text{와 }\theta,K.
\end{array}\right.
\]

%
%
\noindent\emph{환원 불가능 형식들.}

\noindent 계수차수 5:
\[
(asu);\qquad (asu)^2;\qquad (asu)^3;\qquad (asu)^4.
\]

%
%
\noindent 계수차수 6:
\[
\begin{array}{cccc}
(\theta su)&(\theta su)^2&(\theta su)^3&(\theta su)^4\\
s_\vartheta&s_\vartheta(\theta su)&s_\vartheta(\theta su)^2&s_\vartheta(\theta su)^3\\
\cdot&s_\vartheta^2&\cdot&\cdot
\end{array}
\]

%
%
\noindent 계수차수 7:
\[
\begin{array}{c@{\quad}c@{\quad}c@{\quad}c@{\qquad}l}
\cdot&(Ksu)^2&(Ksu)^3&(Ksu)^4&\\
\text{A)}\ s_k&s_k(Ksu)&s_k(Ksu)^2&s_k(Ksu)^3&
\text{B)}\ s_j,\\
\cdot&s_k^2&\cdot&\cdot,&
\text{C)}\ (\Delta su),\ (\Delta su)^2.
\end{array}
\]

%
%
\noindent 계수차수 8:
\[
\begin{array}{ll}
\text{A)}\ s_j(asu),\ s_j(asu)^2,\ s_j(asu)^3,
&
\text{B)}\ \begin{array}[t]{cc}
\cdot&(Nsu)^2\\
s_n&s_n(Nsu).
\end{array}
\end{array}
\]

%
%
\noindent 계수차수 10:
\[
\text{A)}\ s_j(\Delta su),\qquad \text{B)}\ s_\vartheta(\theta s'u)^4.
\]

%
%
\noindent 계수차수 11:
\[
\begin{array}{cc}
(Ks^2u)^6&\cdot\\
s_k(Ks^2u)^5&s_k(Ks^2u)^6.
\end{array}
\]

%
%
\noindent\emph{환원:}

다음 형식열들은
\[
 s_\vartheta^2(\theta su),\quad s_k^2(Ksu),\quad s_j^2,\quad (\Delta su)^3,\quad (Nsu)^3
\]
환원자 $s_\vartheta^2(\theta su)$를 인수로 가지며, $s_j^2$과 $(\Delta su)^3$은 $j$와 $\Delta$를 형식열의 더 낮은 형식들로 되돌려서 얻는다:
\[
\begin{aligned}
s_\vartheta^2(\theta su)(a\theta u)^3&=-\frac12s_j^2,\\
s_\vartheta^2(\theta su)(a\theta u)^2&=\frac13(\Delta su)^3,\\
s_\vartheta^2(\theta su)^2(a\theta u)^2&=\frac13(\Delta su)^4+\frac13s_j^2.
\end{aligned}
\]

%
%
형식 $s_\vartheta^2s_{\vartheta'}$와 $s_\vartheta^2s_{\vartheta'}^2$: 환원자 $s_\vartheta^2(\theta su)$와 $\theta_{\vartheta'}$:
\[
\begin{aligned}
s_\vartheta^2s_{\vartheta'}&=s_\vartheta^2\theta_{\vartheta'}-s_\vartheta^2(\theta s\widehat{\vartheta' x}),\\
s_\vartheta^2s_{\vartheta'}^2&=s_\vartheta^2s_{\vartheta'}\theta_{\vartheta'}-s_\vartheta^2s_{\vartheta'}(\theta s\widehat{\vartheta' x}).
\end{aligned}
\]
형식열 $s_j(\Delta su)^2$: 환원자 $\Delta_j$와 $(\Delta su)^3$.

%
%
\noindent\emph{결론:}

$s$는 거듭제곱으로는 $f,j,\Delta,N$과 수축되지 않는다. $f,j,\Delta,N$은 반변 형식들과의 곱 안에서만 수축되는데, $f$와의 곱에서는 그 반변 형식이 $x$에 관해 여전히 이차일 수 있고, $\Delta$와의 곱에서는 $x$에 관해 여전히 일차일 수 있다. 즉 다음 형식들의 곱을 고찰해야 한다:
\[
 f\cdot(0,\lambda),\quad f\cdot(1,\lambda),\quad f\cdot(2,\lambda),\quad
 j\cdot(\mu,0),\quad N,\quad
 \Delta\cdot(0,\lambda)\quad \Delta\cdot(1,\lambda)\quad(\lambda>0),
\]
여기서 $(\mu,\lambda)$는 $x$에 관해 $\mu$차이고 $u$에 관해 $\lambda$차인 형식을 뜻한다.

%
%
$\theta$ 또는 $K$는 임의의 형식과 이루는 거듭제곱이나 곱 안에서 $s$ 또는 $s^2$과 수축되지 않는다. 실제로 야코비 행렬식 $K$와의 곱 $K\cdot\varphi_x$ ($\varphi$는 $f$도 $\theta$도 아님)에 대해서는 분해 가능 형식들 사이의 다음 관계가 성립한다:
\[
K\cdot\varphi_x=(a\theta u)\varphi_x=f\cdot(\varphi\theta u)-\theta\cdot(\varphi a u).
\]
따라서 위에 제시한 계 $s_I$의 수축 도식이 증명되었다.

%
%
\begin{center}
\textbf{\S\ 19.}\quad\emph{계 $s_{II}$.}
\end{center}

$s$와 $\nu;\ H;\ L;\ H^2$의 수축.

\noindent\emph{환원 불가능 형식들:}
\[
\begin{array}{c@{\qquad}c@{\qquad}c@{\qquad}c}
\text{계수차수 6} & \text{계수차수 8} & \text{계수차수 10} & \text{계수차수 12}\\[0.3em]
s_\nu & (Hsu),\ (Hsu)^2 & (Lsu)^2,\ (Lsu)^3 & (H^2su)^3\\
\cdot & s_\eta & s_l & \cdot\\
\cdot & \cdot & \cdot & \cdot\\
s_\nu^4=J & \cdot & \cdot & \cdot
\end{array}
\]

%
%
\noindent\emph{환원:}

형식열 $s_\eta(Hsu)$와 $s_l(Lsu)$: 환원자 $s_\nu^2$.

형식열 $s_\sigma$: $H$에서 얻는 환원자 $s_\nu^2$,
$s_\nu^2=\dfrac23s_\sigma$.

$(H^2su)^4$은 공식 (7.)a에서 분해된다(\S\ 11. A 참조). 그로부터 $(H\cdot L,s,u)^4$의 분해가 따른다.

%
%
\noindent\emph{결론:}

$s^2$은 계 II와 수축되지 않는다.

$\nu$는 반변 형식과의 곱에서만 수축된다. $H$는 공변식으로 보았을 때 형식 $(1,\lambda)$, $(2,\lambda)$, $(3,\lambda)$ ($\lambda$는 임의)과의 곱에서 수축된다.

$\sigma$와 $(\nu\sigma x)$은 전혀 수축되지 않으며, 야코비 행렬식으로서의 $L$은 곱 안에서 수축되지 않는다(계 III의 공변식들에 대한 완전 전이연산은 환원 가능해진다).

계 I과 II의 곱으로 남는 것은 다음과 같다:
\[
f\cdot \nu,\qquad \Delta\cdot \nu.
\]

%
%
\begin{center}
\textbf{\S\ 20.}\quad\emph{계수차수 7의 형식들(계 $s_{III}$).}
\end{center}

$s$와 $f\cdot \nu$, $a_\nu$, $a_\nu^2$의 수축.

\noindent\emph{환원 불가능 형식들:}
\[
\begin{gathered}
s_\nu(asu),\quad a_\nu(asu),\quad s_\nu(asu)^2,\quad a_\nu(asu)^2,\quad
s_\nu(asu)^3,\quad a_\nu(asu)^3,\\
a_\nu s_\nu,\qquad a_\nu s_\nu(asu),\qquad a_\nu^2(asu),\qquad a_\nu^2(asu)^2.
\end{gathered}
\]

%
%
\noindent\emph{환원:}

환원자 $s_\nu^2$과 $s_\eta(Hsu)$와 직접 수축하여 생기는 자명한 환원들은 언급하지 않으며, 야코비 행렬식과의 전이연산이 분해되는 경우도 마찬가지로 언급하지 않는다.

형식 $a_\nu s_\nu(asu)^2$과 $a_\nu^2s_\nu(asu)^2$: 공식 (19.), \S\ 12 참조;
\[
a_\nu s_\nu(asu)^3=a_\nu(a\widehat{\nu_1\nu_2}u)^3(\nu_1\nu_2\nu)=0.
\]

%
%
형식열 $a_\nu^2s_\nu$: $a_\nu^2$을 더 낮은 형식들로 되돌림으로써 얻는 환원자 $s_\vartheta^2(\theta su)$, 즉 다음 표현들을 이중 환원함으로써 얻는다:
\[
s_\vartheta^2(a\theta s)(a\theta u),\qquad
s_\vartheta^2(a\theta s)(a\theta u)^2
\]
다음과 같이 둔다:
\[
\begin{aligned}
s_\vartheta^2(a\theta s)(a\theta u)
&=[(a\theta u)^2]s^3+\frac15[a_\vartheta(a\theta u)^2]s^2+\frac1{60}[a_\vartheta^2(a\theta u)^2]s\\
&=\frac12a_\nu^2s_\nu-\frac23a_\nu s_\nu^2,\\
s_\vartheta^2(a\theta s)(a\theta u)^2
&=\frac45[a_\vartheta(a\theta u)^2]_{\widehat{us}}s^2+\frac{23}{180}[a_\vartheta^2(a\theta u)^2]_{\widehat{us}}s\\
&=-\frac12a_\nu^2s_\nu(asu).
\end{aligned}
\]
이때 $(a_\nu^2b_\nu(abu)=0)$이다.

%
%
다음 공식들을 얻는다:
\srcnumdisplay{(25.)}{%
\begin{array}{rcl@{\qquad}rcl}
\underline{a_\nu^2s_\nu}&=&\dfrac43i\cdot a_\nu^2b_\nu,
& \underline{a_\nu s_\nu(asu)^2}&=&\dfrac13i\cdot\theta_\nu^2-\dfrac1{18}iH,\\[0.7em]
\underline{a_\nu s_\nu(asu)^3}&=&0,
& \underline{a_\nu^2s_\nu(asu)}&=&0,\\[0.7em]
\underline{a_\nu^2s_\nu(asu)^2}&=&0.
\end{array}
}

%
%
\noindent\emph{결론:}

1) $a_\nu s_\nu(asu)^2$, $a_\nu^2s_\nu$은 환원자이다.

2) \S\ 19에서 환원 가능하다고 판정된 형식들 $(\Delta su)^3$, $(\Delta su)^4$의 값들을 얻는다. 대응하는 형식열 $(agu)^2$에 대해서도 마찬가지이며, 이 형식열은 $\nu$와 수축할 때 환원자 $g_\nu$에 의해 이중 환원을 일으킨다.
\srcnumdisplay{(26.)}{%
\begin{aligned}
\text{a)}\quad (\Delta su)^3&=-\frac65a_\nu^2(asu)+3a_\nu s_\nu(asu)+2i\cdot\theta_\nu(\vartheta\nu x),\\
\text{b)}\quad (\Delta su)^4&=-\frac25a_\nu^2(asu)^2+\frac43i\theta_\nu^2+\frac19iH.
\end{aligned}
}

%
%
공식 (24.)와 (26.)에서, $(i,J)$를 법으로 하여(71쪽 참조):
\srcnumdisplay{(27.)}{%
\begin{aligned}
\text{a)}\quad (agu)^2&=\frac23(\Delta su)^2-\frac43a_\nu s_\nu+\frac35s\cdot a_\nu^2,\\
\text{b)}\quad (agu)^3&=2a_\nu s_\nu(asu)-a_\nu^2(asu)^2,\\
\text{c)}\quad (agu)^4&=-a_\nu^2(asu)^2.
\end{aligned}
}

%
%
\begin{center}
\textbf{\S\ 21.}\quad\emph{계수차수 8의 형식들(계 $s_{III}$).}
\end{center}

$s$와 계수차수 4의 형식들(계 III)의 수축. (\S\ 9.)

\noindent\emph{환원 불가능 형식들:}
\[
\begin{array}{c|c|c}
\cdot
%
& ((\vartheta \nu x),s,u)^2,\ (\theta_\nu su)
& \bigstar\\[0.45em]
(\vartheta \nu x)s_\nu,\ ((\vartheta \nu x)^2,s,u)
& ((\vartheta \nu x)^2,s,u)^2,\ \theta s_\nu
& \\[0.45em]
(\vartheta \nu x)^2s_\nu & ((\vartheta \nu x)^2,s,u)s_\nu,\ (\theta_\nu(\vartheta \nu x)^2,s,u)
&
\end{array}
\]

%
%
\[
\begin{array}{c|c|c|c}
\bigstar
& ((\vartheta \nu x),s,u)^3,\ (\theta_\nu su)^2
& ((\vartheta \nu x),s,u)^4,\ (\theta_\nu su)^3
& \cdot\\[0.45em]
& ((\vartheta \nu x)^2,s,u)^3,\ (\theta_\nu(\vartheta \nu x),s,u)^2,\ (\theta_\nu^2su)
& (\theta_\nu(\vartheta \nu x),s,u)^3,\ (\theta_\nu^2su)^2
& (\theta_\nu(\vartheta\nu x),s,u)^4.\\[0.45em]
& \cdot
& (\theta_\nu^3su)
& \cdot
\end{array}
\]

%
%
\noindent\emph{환원:}

$((\vartheta \nu x),s,u)s_\nu$은 \S\ 3에 따라 야코비 행렬식에 대한 전이연산으로 분해된다.

형식열 $((\vartheta \nu x),s,u)^2s_\nu$: 환원자 $s_\nu^2$과 $s_\vartheta^2\theta_\nu$:
\[
(\widehat{\vartheta \nu}su)s_\nu(\theta su)=s_\nu s_\vartheta(\theta su)=-\frac12(\widehat{\vartheta \nu}su)^2(\theta su),
\]
\[
\theta_\nu(\widehat{\vartheta \nu}su)s_\nu=\theta_\nu s_\nu s_\vartheta=-\frac12\theta_\nu(\widehat{\vartheta \nu}su)^2.
\]
형식열 $\theta_\nu s_\nu(\theta su)$: 다음 표현들을 이중 환원하여 얻는 환원자 $a_\nu s_\nu(asu)^2$:
\[
a_\nu s_\nu(asu)^2(abu)\quad\text{및}\quad a_\nu s_\nu(asu)^2(abu)(sbu);
\]
$\theta_\nu s_\nu(\theta su)^3$은 $(agu)^3(abu)(bgu)^3$을 이중 환원하여 얻는다.

%
%
형식열 $\theta_\nu(\vartheta \nu x)s_\nu$: 등식 $\theta_\nu(\vartheta \nu x)s_\nu=a_\nu^2(abu)s_\nu$에서 얻는 환원자 $a_\nu^2s_\nu$.

형식열 $(\theta_\nu(\vartheta \nu x)^2,s,u)$: 형식열 $\theta_\nu(\widehat{\vartheta \nu}su)s_\nu$의 형식들 가운데 동시에 형식열 $((\vartheta \nu x)su)^2s_\nu$과 $\theta_\nu(\vartheta \nu x)s_\nu$에 속하는 것들을 이중 환원한다.

형식 $(\theta_\nu(\vartheta \nu x),s,u)$은 야코비 행렬식 $\theta_\nu(\vartheta \nu x)=a_\nu^2(abu)$에 대한 한 번의 전이연산으로 분해된다.

형식 $((\vartheta \nu x)^2,s,u)^4$: 표현 $(a\A u)^2(\A su)^4$을 이중 환원한다. 또는 공식 (27.)에서 나온 $(\theta gu)^4$를 $(j\eta x)^4$의 환원(\S\ 12 C))과 유사하게 이중 환원한다.

%
%
이중 환원을 통하여 다음 공식들을 얻는다:
\srcnumdisplay{(28.)}{%
\begin{aligned}
0&=\theta_\nu s_\nu(\theta su)-\frac16(\widehat{\vartheta \nu}su)^2(\theta su)-\frac1{12}(Hsu),\\
0&=\theta_\nu s(\theta su)^2-\frac14(\widehat{\vartheta \nu}su)^2(\theta su)^2-\frac1{24}(Hsu)^2,\\
0&=(\widehat{\vartheta \nu}su)^2(\theta su)^2+2\theta_\nu(\widehat{\vartheta \nu}su)(\theta su)^2
  +3\theta_\nu^2(\theta su)^2-\frac13H(su)^2,\\
0&=\theta_\nu s_\nu(\theta su)^3+\frac12\theta_\nu(\widehat{\vartheta \nu}su)(\theta su)^3,\\
0&=\theta_\nu s_\nu s_\vartheta-\frac14s_\eta,
\quad 0=\theta_\nu s_\nu s_\vartheta(\theta su),
\quad 0=\theta_\nu s_\nu s_\vartheta(\theta su)^2.
\end{aligned}
}
마지막 것은 $(\theta_\nu(\vartheta \nu x)^2,s,u)^2\ldots$에 대응한다.

%
%
\noindent\emph{결론:}

1) $\theta(\vartheta \nu x)s_\nu$, $(\widehat{\vartheta \nu}su)(\theta su)s_\nu$, $\theta_\nu(\widehat{\vartheta \nu}su)^2$, $(\widehat{\vartheta \nu}su)^2(\theta su)^2$은 환원자이다.

2) 계수차수 4의 형식들 $(5,\lambda)$, $(6,\lambda)$ 등은 $s^2$과의 수축에 들어가지 않는다.

%
%
3) 형식열 $(\theta gu)^2$에 대해서는 공식 (27.)과 이 절의 환원들을 고려함으로써 다음 공식들을 얻는다:
\srcnumdisplay{(29.)}{%
\begin{aligned}
\text{a)}\quad (\theta gu)^2&=\frac43\theta_\nu s_\nu+\frac23s_\nu s_\vartheta-\frac13s\cdot\theta_\nu^2-\frac19s\cdot H-\frac13f\cdot a_\nu^2(asu)^2+\frac13a_\nu^2\cdot(asu)^2,\\
\text{b)}\quad (\theta gu)^3&=-\frac13(\widehat{\vartheta \nu}su)^2(\theta su)-\theta_\nu(\widehat{\vartheta \nu}su)(\theta su)-\theta_\nu^2(\theta su),\\
\text{c)}\quad (\theta gu)^4&=2\theta_\nu^2(\theta su)^2-\frac13(Hsu)^2,\\
g\vartheta(\theta gu)&=(\widehat{\vartheta \nu}su)(\vartheta \nu x)s_\nu-\theta_\nu(\vartheta \nu x)(\widehat{\vartheta \nu}su),\\
g\vartheta(\theta gu)^2&=\frac23s\cdot\theta_\nu^3,
\quad g\vartheta(\theta gu)^3=\frac13\theta_\nu^3(\theta su),
\quad g\vartheta(\theta gu)^4=0.
\end{aligned}
}

%
%
\begin{center}
\textbf{\S\ 22.}\quad\emph{계수차수 9의 형식들(계 $s_{III}$).}
\end{center}

$s$와 계수차수 5의 형식들(계 III) 및 곱 $\A\cdot \nu$의 수축(\S\ 10).

\noindent\emph{환원 불가능 형식들:}
\[
\begin{array}{c@{\quad}c|c|c|c}
\text{A)}& \cdot & \cdot & (K_\nu su)^2 & \bigstar\\[0.45em]
& \cdot & ((k\nu x)^2,s,u)^2,\ K_\nu s_\nu & ((k\nu x)^2,s,u)^3 & \\[0.45em]
& (k\nu x)^2s_\nu,\ ((k\nu x)^3,s,u) & ((k\nu x)^3,s,u)^2 & \cdot & \cdot\\[0.45em]
& (k\nu x)^3s_\nu & ((k\nu x)^3,s,u)s_\nu,\ (K_\nu(k\nu x)^3,s,u) & \cdot & \cdot\\[0.45em]
\bigstar & (K_\nu su)^3 & (K_\nu su)^4 & \cdot & \\[0.45em] & \cdot & \cdot & (K_\nu^2su)^4, &
\end{array}
\]
\[
\begin{array}{c@{\quad}c|c@{\qquad}c@{\quad}c}
\text{B)}& (j\nu x)s_\nu,\ ((j\nu x)^2,s,u) & ((j\nu x)^2,s,u)^2 && \\
& ((j\nu x)^3,s,u) & \cdot, && \text{C)}\quad s_\nu(\A su),\ (\A_\nu su),
\end{array}
\]
\[
\begin{array}{c@{\quad}c|c|c}
\text{D)}& \cdot & ((aHu),s,u)^2,\ ((aHu)^2,s,u) & \bigstar\\[0.45em]
& (aHu)s_\eta,\ (a_\eta su) & (aHu)^2s_\eta,\ (a_\eta su)^2 & \\[0.45em]
& \cdot & (a_\eta^2su) & \\
\bigstar& ((aHu),s,u)^3,\ ((aHu)^2,s,u)^2 & ((aHu),s,u)^4 & \\
& (a_\eta su)^3,\ (a_\eta(aHu),s,u)^2,\ (a_\eta(aHu)^2,s,u) & (a_\eta(aHu),s,u)^3. &
\end{array}
\]

%
%
\noindent\emph{환원:}

A) 환원들은 한 번 수축하여 \S\ 21의 환원들에서 직접 따른다. 형식열 $K_\nu s_\nu(Ksu)^2$은 환원자 $a_\nu s_\nu(asu)^2$과 $\theta_\nu s_\nu(\theta su)^2$을 인수로 갖는다. 이로부터 다음과 같이 두면 더 높은 형식 $K_\nu^2(Ksu)^2$, $K_\nu^2(Ksu)^3$이 분해된다:
\[
0=a_\nu s_\nu(asu)^2(a\theta u)=K_\nu s_\nu(Ksu)^2=\theta_\nu s_\nu(\theta su)^2(a\theta u).
\]

%
%
B) 형식열 $(j\nu x)^2s_\nu$: 다음 식에서 얻는 환원자 $a_\nu^2s_\nu$:
\[
(j\nu x)^2s_\nu=3a_\nu^2s_\nu(a\theta u)^2.
\]
$((j\nu x)^3,s,u)^2$: 다음 식에서 얻는 환원자 $a_\nu^2s_\nu$:
\[
((j\nu x)^3,s,u)^2=-6a_\nu^2s_\nu(a\theta u)^2(\theta su).
\]

%
%
C) 형식 $\A_\nu(\A su)^2$과 $s_\nu(\A su)^2$: 공식 (27.)a에서 얻는 환원자 $g_\nu$.

형식 $\A_\nu s_\nu$: 1) $a_\vartheta a_\nu^2s_\nu=\frac23\A_\nu s_\nu$에서 얻는 환원자 $a_\nu^2s_\nu$;

2) 공식 (27.)a에 따라 표현식 $(\A s\widehat{\nu x})^2$을 이중 환원함으로써 분해된다.

그로부터 더 높은 형식 $a_\eta s_\eta$의 분해가 따른다.

%
%
D) 형식열 $(aHu)^2s_\eta(asu)$: 형식열 $[a_\nu s_\nu(asu)^2]\widehat{\nu_1x}$을 이중 환원하여 얻는 환원자 $a_\nu s_\nu(asu)^2$; $a_\nu s_\nu(asu)^2$과 $a_\nu s_\nu(asu)^3$에서 $(aHu)^2s_\eta(asu)^2$의 환원을 얻는다. 그로부터 $(a_\eta,s,u)^4$의 환원이 따른다.

형식열 $a_\eta(aHu)s_\eta$: 표현식 $(\A s\widehat{\nu x})^3$을 이중 환원하여 얻는 환원자 $a_\nu^2s_\nu$, $a_\eta^2s_\eta$. 이는 $a_\eta^2s_\eta$가 형식 $a_\eta(aHu)s_\eta$의 환원 가능 항 $a_\nu^2s_\nu(\nu\nu_1x)$에서 수축으로 생기지 않기 때문이다.

형식열 $(a_\eta^2su)^2$: 다음 식에서 얻는 환원자 $a_\nu^2s_\nu$:
\[
a_\nu^2s_\nu s_{\nu_1}=-\frac12a_\eta^2(Hsu)^2.
\]
형식 $(a_\eta(aHu),s,u)$은 야코비 행렬식에 대한 전이연산으로서 분해된다.

%
%
다음 환원 공식들을 얻는다:
\srcnumdisplay{(30.)}{%
\begin{aligned}
\text{a)}\quad 0&=(aHu)^2(asu)s_\eta-a_\eta(asu)(Hsu)^2-a_\eta(aHu)(asu)(Hsu),\\
\text{b)}\quad 0&=(aHu)^2(asu)^2s_\eta-a_\eta(aHu)(asu)^2(Hsu),\\
\text{c)}\quad 0&=a_\eta(asu)^2(Hsu)^2,\\
0&=a_\eta s_\eta+\frac14a_\nu^2\cdot g-\frac14s\cdot a_\eta^2,\\
0&=a_\eta(aHu)s_\eta-\frac12a_\eta^2(Hsu),
\quad 0=a_\eta^2(Hsu)^2,\ldots,
\quad 0=a_\eta^2s_\eta.
\end{aligned}
}

%
%
\noindent\emph{결론:}

1) $(j\nu x)^2s_\nu$, $\A_\nu s_\nu$, $\A_\nu(\A su)^2$, 따라서 $N_\nu(Nsu)^2$, $(aHu)^2(asu)s_\eta$, $a_\eta(aHu)s_\eta$, $a_\eta^2(Hsu)^2$은 환원자이다. $a_\eta s_\eta$는 분해 가능 형식이며, 모든 한 번 및 두 번의 수축에서 분해 가능 형식으로 옮겨간다.

2) 계수차수 5의 형식들 $(5,\lambda)$, $(6,\lambda)$ 등은 $s^2$과 수축되지 않는다.

%
%
\begin{center}
\textbf{\S\ 23.}\quad\emph{계수차수 10의 형식들(계 $s_{III}$).}
\end{center}

$s$와 계수차수 6의 형식들(계 III)의 수축(\S\ 11).

\noindent\emph{환원 불가능 형식들:}
\[
\begin{array}{r@{\quad}c}
\text{A)}&
\begin{array}{c|c}
(\vartheta^2\nu x)^3s_\nu&\cdot\\[0.45em]
\cdot&(\vartheta^2\nu x)^3s_\nu s_\vartheta,\ \theta_\nu(\vartheta^2\nu x)^3s_\vartheta,
\end{array}
\end{array}
\]

%
%
\[
\begin{array}{r@{\quad}c}
\text{B)}&
\begin{array}{c|ccc}
\cdot&(a_\nu(j\nu x),s,u)^2&(a_\nu(j\nu x),s,u)^3&(a_\nu(j\nu x),s,u)^4\\[0.45em]
a_\nu(j\nu x)s_\nu&\cdot&(a_\nu^2(j\nu x),s,u)^2&(a_\nu^2(j\nu x),s,u)^3,
\end{array}
\end{array}
\]

%
%
\[
\begin{array}{r@{\quad}c}
\text{C)}&
\begin{array}{c|c|c|c}
\cdot&\cdot&((\theta Hu),s,u)^2,\ ((\theta Hu)^2,s,u)&\bigstar\\[0.35em]
\cdot&((\vartheta\eta x),s,u)^2,\ (\theta Hu)s_\eta,\ (\theta_\eta su)&((\vartheta\eta x),s,u)^3,\ (\theta_\eta su)^2&\\[0.35em]
((\vartheta\eta x)^2,s,u)&\cdot&(\theta_\eta^2su)&\\[0.35em]
\cdot&(\theta_\eta(\vartheta\eta x)^2,s,u)&\cdot&\cdot
\end{array}
\end{array}
\]
\[
\begin{array}{c@{\quad}|c|c}
\bigstar&((\theta Hu),s,u)^3,\ ((\theta Hu)^2,s,u)^2&((\theta Hu),s,u)^4\\[0.45em]
&(H_\vartheta(\theta Hu),s,u)^2,\ (\theta_\eta(\theta Hu),s,u)^2,\ (\theta_\eta(\theta Hu)^2,s,u)&(\theta_\eta su)^4,\ (\theta_\eta(\theta Hu),s,u)^3\\[0.45em]
&\cdot&(\theta_\eta^2(\theta Hu),s,u)^2.
\end{array}
\]

%
%
\noindent\emph{환원:}

A) 형식들 $((\vartheta^2\nu x)^3,s,u)^2$, $((\vartheta^2\nu x)^3,s,u)^3$은 분해된다:
\[
(\vartheta \nu x)(\widehat{\vartheta\nu}su)(\widehat{\vartheta'\nu}su)
=(\vartheta \nu x)s_\vartheta s_{\vartheta'}-(\vartheta\nu x)s_\nu s_{\vartheta'}
=-2\theta\cdot(\vartheta\nu x)s_\nu s_\vartheta+(\vartheta\nu x)\cdot s_\vartheta^2.
\]
나머지 형식들은 환원자를 인수로 갖거나 분해 가능 형식들과 한 번 수축하여 생긴다.

%
%
B) 환원자는 $a_\nu^2s_\nu$와 $(\widehat{j\nu}su)s_\nu$이다.

C) 1. 형식열 $(\theta Hu)^2s_\eta$: a) 환원자 $(aHu)^2(asu)s_\eta$; b) 형식열 $(\theta gu)^3g_\nu$와 $g_\nu(agu)^3(bgu)^2(abu)$의 이중 환원.

이로부터 더 높은 형식들 $(\theta_\eta su)^3$, $(H_\vartheta(\theta Hu),s,u)^3$, $(a_\nu\cdot b_{\nu_1}^2,s,u)^4$의 환원이 따른다[마지막 형식은 $(H_\vartheta su)^4$ 대신이다].

%
%
2. $((\vartheta\eta x),s,u)^4$: 환원자 $(a_\eta su)^4$.

3. $((\vartheta\eta x)^2,s,u)^3$: $(\widehat{\vartheta\eta}su)^2(Hsu)$에서 얻는 환원자 $s_\vartheta^2s_\nu$.
형식들 $(\vartheta\eta x)s_\eta$, $(\vartheta\eta x)^2s_\eta$, $((\vartheta\eta x)^2,s,u)^2$, $\theta_\eta s_\eta$는 분해 가능 형식 $a_\eta s_\eta$와 수축하여 분해된다.

%
%
4. 형식열 $H_\vartheta(\theta Hu)s_\eta$, $\theta_\eta(\theta Hu)s_\eta$, $H_\vartheta(\vartheta\eta x)s_\eta$, $\theta_\eta(\vartheta\eta x)s_\eta$:
환원자 $a_\eta(aHu)s_\eta$와 $a_\eta^2s_\eta$.

형식열 $(\theta_\eta(\vartheta\eta x),s,u)^2$: 환원자 $a_\eta^2(Hsu)^2$ [단, $(\theta_\eta^2(\theta Hu),s,u)^2$은 예외이며, 이는 항 $a_\eta^2(bsu)(Hsu)$에서 수축하여 생긴다].

%
%
5. $(H_\vartheta(\vartheta\eta x),s,u)$, $(H_\vartheta(\vartheta\eta x),s,u)^2$: $a_\eta(aHu)s_\eta(abu)$와 $g_\vartheta(\theta gu)g_\nu$의 이중 환원.

형식열 $(H_\vartheta(\vartheta\eta x),s,u)^3$: 환원자 $((\vartheta \nu x)^2,s,u)^2s_\nu$.

%
%
6. $(H_\vartheta(\theta Hu),s,u)$는 \S\ 3. 2), c)에 따라 분해 가능 형식 $(\theta_\nu(\vartheta \nu x),s,u)$과 한 번 수축하여 분해된다. 즉 더 낮은 분해 가능 형식 $(H_\vartheta su)^2$을 이중 환원한다:\srcfn{*)}{기호 $\equiv$는 더 낮은 차수의 형식들의 곱을 생략했음을 뜻한다.}
\[
\begin{aligned}
0&\equiv \theta_\nu(\vartheta\nu\nu_1)(\theta su)+\frac23\theta_\nu^2(\vartheta \nu_1 x)(\theta su)+\theta_\nu(\vartheta \nu x)(\theta su)s_{\nu_1}\\
&\equiv -\frac12H_\vartheta(\theta Hu)(\theta su)+\theta_\eta(\theta su)(Hsu)+\frac12H_\vartheta(\theta su)(Hsu)\\
&\equiv -\frac12H_\vartheta(\theta Hu)(\theta su)+\theta_\eta(\theta su)(Hsu)+\frac12[H_\vartheta]_{\widehat{su}^2}-\frac12\cdot\frac35H_\vartheta(\theta Hu)(\theta su)\\
&\equiv -\frac35H_\vartheta(\theta Hu)(\theta su).
\end{aligned}
\]

%
%
1과 5에 따라 이중 환원하면 다음 공식들이 생긴다:
\[
\begin{array}{@{}r@{\quad}l@{}}
&\displaystyle
\begin{aligned}
0={}&\theta_\eta(\theta su)(Hsu)^2+\frac14H_\vartheta(\theta Hu)(\theta su)^2-\frac34\theta_\eta(\theta Hu)(\theta su)(Hsu)\\
&-\frac54\theta_\eta(\theta Hu)^2(\theta su)-(asu)^3\cdot b_\eta(bHu)^2-a_\nu(asu)^3\cdot b_{\nu_1}^2,
\end{aligned}\\[0.55em]
(31.)&\displaystyle 0=H_\vartheta(\theta Hu)(\theta su)^3-7\theta_\eta(\theta Hu)(\theta su)^2(Hsu),\\[0.55em]
&\displaystyle 0=a_\nu(asu)^2b_{\nu_1}^{2}(bsu)^2-\theta_\eta(\theta su)^2(Hsu)^2+6\theta_\eta(\theta Hu)(\theta su)^2(Hsu),\\[0.55em]
&\displaystyle 0\equiv (H_\vartheta(\vartheta\eta x),s,u),\qquad
0\equiv H_\vartheta(\widehat{\vartheta\eta}su)(Hsu)-4\theta_\eta^2(Hsu).
\end{array}
\]

%
%
\noindent\emph{결론:}

1) $(\theta Hu)^2s_\eta$, $H_\vartheta(\theta Hu)s_\eta$, $\theta_\eta(\theta Hu)s_\eta$, $H_\vartheta(\vartheta\eta x)s_\eta$, $\theta_\eta(\vartheta\eta x)s_\eta$, $(H_\vartheta(\vartheta\eta x),s,u)$, $(\theta_\eta(\vartheta\eta x),s,u)^2$은 환원자이다.

2) 계수차수 6의 형식들 $(5,\lambda)$ 등은 $s^2$과의 수축에 들어가지 않는다.

%
%
\begin{center}\textbf{\S\ 24.}\quad\emph{계수차수 11의 형식들(계 $s_{III}$).}\end{center}
$s$와 계수차수 7의 형식들(계 III)의 수축(\S\ 12).

\noindent\emph{환원 불가능 형식들:}
\[
\text{A)}\quad
\begin{array}{c|c|c|c|c}
\cdot&\cdot&\cdot&((KHu)^2su)^2&\bigstar\\[0.35em]
\cdot&\cdot&(K_\eta su)^2&\cdot&\\[0.35em]
\cdot&\cdot&\cdot&\cdot&\cdot\\[0.35em]
((k\eta x)^3,s,u)&\cdot&\cdot&\cdot&\cdot\\[0.35em]
\cdot&(K_\eta(k\eta x)^3,s,u)&\cdot&\cdot&\cdot
\end{array}
\]
\[
\begin{array}{c@{\quad}|c|c}
\bigstar&((KHu)^2,s,u)^3&((KHu)^2,s,u)^4\\[0.45em]
&\cdot&(K_\eta(KHu)^2,s,u)^3,
\end{array}
\]
\[
\text{B)}\quad ((j\eta x),s,u)^2,\ (H_jsu)\ \big|\ ((j\eta x),s,u)^3,
\qquad \text{C)}\quad (\D_\eta su),
\]
\[
\text{D)}\quad
\begin{array}{c|c|c}
\cdot&((aLu)^2,s,u)^2,\ ((aLu)^3,s,u)&((aLu)^2,s,u)^3\\[0.5em]
(aLu)^2s_l,\ (a_lsu)^2&(a_lsu)^3&(a_l(aLu)^2,s,u)^2.
\end{array}
\]

%
%
\noindent\emph{환원:}

A) 기호 $\theta\cdot H$로 나타낸 대응 형식들과 한 번 수축하여 환원되거나 분해된다.

$(K_\eta(KHu)^2,s,u)$와 $(K_\eta(KHu)^2,s,u)^2$: $K_\nu^2(Ksu)$, $K_\nu^2(Ksu)^2$과 수축하여 \S\ 3. 2), a)에 따라 분해된다.

$(K_\eta(KHu),s,u)^4\equiv(a_\nu^2\cdot\theta_{\nu_1},s,u)^4$: 분해 가능 형식 $K_\nu^2(Ksu)^3$으로부터 \S\ 3. 2), b)에 따라 분해된다.

%
%
B) 형식열 $(j\eta x)s_\eta$: 환원자 $(j\nu x)^2s_\nu$.

형식 $H_j(\widehat{j\eta}su)(Hsu)$: 환원자 $(j\nu x)(\widehat{j\nu}su)^2$.

형식 $H_j(j\eta x)(Hsu)$: 환원자 $(j\nu x)^2s_\nu$를 사용하여 분해 가능 형식 $((j\eta x)^2,s,u)^2$을 환원함으로써 분해된다(공식 (18.)a):
\[
0=-2(j\nu x)^2s_\nu s_{\nu_1}=[(j\eta x)^2]_{\widehat{su}^2}+H_j(j\eta x)(Hsu).
\]

%
%
C) 환원자 $\D_\nu s_\nu$와 $\D_\nu(\D su)^2$.

D) 기호 $f\cdot H$로 나타낸 대응 형식들과 한 번 수축하여 환원되거나 분해된다.

%
%
E) (30.)c에서 분해 가능 형식열 $(a_\sigma su)^2$의 환원을 얻는다:
\[
\begin{aligned}
0&=a_\eta(Hsu)^2(as\widehat{\nu x})^2
=a_\eta(Hsu)^2(a\widehat{\nu\eta}u)^2+2a_\eta(a\widehat{\nu\eta}u)(\widehat{\nu\eta}su)(Hsu)^2
=\frac13a_\sigma(asu)^2,\\
0&=a_\eta a_\nu(Hsu)^2(as\widehat{\nu x})(asu)
=a_\eta(a\widehat{\nu\eta}u)^2(Hsu)^2(asu)+a_\eta(a\widehat{\nu\eta}u)(\widehat{\nu\eta}su)(Hsu)^2(asu)
=\frac13a_\sigma(asu)^3.
\end{aligned}
\]

%
%
\noindent\emph{결론:} $s^2$은 계수차수 7의 형식들과 수축되지 않는다.

%
%
\begin{center}\textbf{\S\ 25.}\quad\emph{계수차수 12, 13, 14, 15의 형식들(계 $s_{III}$).}\end{center}
$s$와 계수차수 8, 9, 10, 11의 형식들(계 III)의 수축(\S\S\ 13--16).

\noindent\emph{환원 불가능 형식들:}

%
%
\noindent 계수차수 12.
\[
\begin{adjustbox}{max width=0.98\linewidth}
$\displaystyle
\begin{array}{c@{\qquad}c@{\qquad}c}
\text{A)}&
\begin{array}{c|c}
((\vartheta^2\eta x)^4,s,u)&\cdot\\[0.35em]
\cdot&\theta_\eta(\vartheta^2\eta x)^3s_\vartheta,
\end{array}
&
\text{B)}\quad ((aHu)H_j,s,u)^2\;\big|\;((aHu)H_j,s,u)^3,
\\[1.0em]
\text{C)}&
\begin{array}{c|c|c}
\cdot&((\theta Lu)^2,s,u)^2,\ ((\theta Lu)^3,s,u)&((\theta Lu)^2,s,u)^3\\[0.35em]
(\theta_lsu)^2&\cdot&(\theta_l(\theta Lu)^2,s,u)^2.
\end{array}
&
\end{array}$
\end{adjustbox}
\]

%
%
\noindent\emph{환원:} 환원자 또는 분해 가능 형식과 직접 수축하여 생기는 환원은 앞으로 더 이상 언급하지 않는다.

B) $((aHu)H_j,s,u)$와 $(a_\eta(aHu)H_j,s,u)$는 야코비 행렬식에 대한 전이연산으로서 분해된다.

%
%
$(a_\eta(aHu)H_j,s,u)^2$의 분해: 분해 가능 형식 $[\theta_\nu\cdot\theta_{\nu_1}^{\prime 3}]_{\widehat{su}^3}$의 이중 환원(\S\ 13. C)
\begin{align*}
0&\equiv[\theta_{\nu_1}^{\prime 3}\cdot\theta_\nu]_{\widehat{su}^3}
\equiv-\frac34\theta_\nu(\theta su)^2(\theta\theta'u)\theta_{\nu_1}^{\prime 3}\\
&\equiv-\frac98(\widehat{\theta\theta'}\eta\widehat{su})(\theta\theta'u)(\theta Hu)(\theta'Hu)\theta_\eta'(\theta su)
\equiv-\frac38a_\eta(aHu)H_j(asu)^2.
\end{align*}

%
%
C) 공식 (31.)로부터 분해 가능 형식
\[
 [L_\vartheta(\theta Lu)]_{\widehat{su}^3}\quad\hbox{및}\quad [\theta_l(\theta Lu)]_{\widehat{su}^3}
\]
의 환원이 얻어진다. 즉 곱 $[\theta_\nu^2\cdot H]_{\widehat{su}^3}$와 $[a_\nu^2\cdot(bHu)^2]_{\widehat{su}^3}$에 대해 다음을 얻는다:
\[
\begin{array}{@{}r@{\quad}l@{}}
&\displaystyle 0\equiv\theta_\nu^2(\theta su)^2(Hsu),\qquad
0\equiv [L_\vartheta(\theta Lu)]_{\widehat{su}^4}-7[\theta_l(\theta Lu)]_{\widehat{su}^4},\\[0.55em]
&\displaystyle 0\equiv a_\nu^2(asu)^2(bHu)^2(bsu)+6\theta_l(\theta Lu)^2(\theta su)(Lsu),\\[0.55em]
(32.)&\displaystyle 0\equiv\theta_\nu^3(\theta su)(Hsu)^2
\quad\hbox{은 }0\equiv\theta_l^2(\theta Lu)(\theta su)(Lsu)^2\quad\hbox{에서 얻는다}\\[-0.05em]
&\displaystyle\hfill(\hbox{환원자 }a_\eta^2(Hsu)^2).
\end{array}
\]

%
%
\noindent\emph{환원 불가능 형식들.}

\noindent 계수차수 13.
\[
\begin{array}{rl}
\text{A)}&((KLu)^3,s,u)^2;\ ((KLu)^3,s,u)^3,\\[0.25em]
\text{B)}&(L_jsu)^2,\\[0.25em]
\text{C)}&((aH^2u)^3,s,u)^2.
\end{array}
\]

%
%
\noindent 계수차수 14.
\[
\begin{array}{rl}
\text{A)}&((aLu)^2L_j,s,u)^2,\\[0.25em]
\text{B)}&((\theta H^2u)^3,s,u)^2.
\end{array}
\]
\noindent 계수차수 15.
\[
 ((KH^2u)^4,s,u)^2.
\]

%
%
\noindent\emph{환원:} $(K_l(KLu)^3,s,u)^2$, $(H_\vartheta(\theta H^2u)^3,s,u)$, $((K,H\cdot L,u)^5,s,u)$는 \S\ 3. 2), c)에 따라 분해된다. 즉 분해 가능 형식 $(K_l(KLu)^2,s,u)^3$, $(H_\vartheta(\theta H'u)^2,s,u)^2$, $((K,H\cdot L,u)^4,s,u)^2$을 동시에 고찰한다.

\noindent\emph{결론:} $s^2$은 계 III의 개별 형식들과 수축되지 않는다.

%
%
\begin{center}\textbf{\S\ 26.}\quad\emph{계 $s_{IV}$.}\end{center}
\noindent 1) \emph{$s$와 계 I$\cdot$III의 수축.}

\noindent\emph{환원:} 앞 절들의 환원들 $(a_\nu^2s_\nu,\ (aHu)^2(asu)s_\eta$ 등)에 따라 형식 $(0,\lambda)$, $(1,\lambda)$, $(2,\lambda)$는 수축 I과 III을 동시에 허용하지 않는다. 따라서 \S\ 18에 따라 계 $s_{IV}$를 만들 때 $f,\D,N$은 곱 안에서 수축되지 않는다.

$j$는 수축되지 않는다. 같은 환원들에 따라 $j$에 반변인 다음 수축들
\[
(K_\nu(k\nu x)^3,s,u),\qquad ((\vartheta^2\eta x)^4,s,u)\quad\hbox{등}
\]
에 $j$와 동변인 다음 수축들이 대응하기 때문이다:
\[
K_\nu(k\nu x)^2s_k,
\qquad (\vartheta^2\eta x)^3s_\vartheta\quad\hbox{등}.
\]

%
%
\noindent 2) \emph{$s$와 계 II$\cdot$III의 수축},
즉 $s$와 $H\cdot(1,\lambda)$, $H\cdot(2,\lambda)$, $H\cdot(3,\lambda)$의 수축.

\noindent\emph{환원 불가능 형식들:}
\[
\begin{array}{rl}
\text{A)}&(a_\nu\cdot H,s,u)^4,
((aHu)^2\cdot H,s,u)^3,
((aHu)^2\cdot H,s,u)^4,\\[0.25em]
&((aLu)^2\cdot H,s,u)^4,
((aLu)^3\cdot H,s,u)^3,
((aH^2u)^3\cdot H,s,u)^4.\\[0.45em]
\text{B)}&(a_\nu^2\cdot H,s,u)^3,
(a_\eta(aHu)^2\cdot H,s,u)^3,
(a_l(aLu)^2\cdot H,s,u)^4.\\[0.45em]
\text{C)}&(\theta_\nu\cdot H,s,u)^4,
((\theta Hu)^2\cdot H,s,u)^3,
((\theta Hu)^2\cdot H,s,u)^4,\\[0.25em]
&((\theta Lu)^2\cdot H,s,u)^4,
((\theta Lu)^3\cdot H,s,u)^3,
((\theta H^2u)^3\cdot H,s,u)^4,\\[0.45em]
\text{D)}&(\theta_\eta(\theta Hu)^2\cdot H,s,u)^3,
(\theta_l(\theta Lu)^2\cdot H,s,u)^4,\\[0.45em]
\text{E)}&((KLu)^3\cdot H,s,u)^4,
((KH^2u)^4\cdot H,s,u)^4,\\[0.45em]
\text{F)}&((j\nu x)^2\cdot H,s,u)^3,
((j\nu x)^2\cdot H,s,u)^4,
((j\eta x)\cdot H,s,u)^4,\\[0.25em]
&(H_j\cdot H,s,u)^3,
(L_j\cdot H,s,u)^4.
\end{array}
\]

%
%
\noindent\emph{환원:} $\nu$는 $j$와 마찬가지로 수축되지 않는다.

B)와 D). $(a_\nu^2\cdot H,s,u)^4$과 $(\theta_\nu^2\cdot H,s,u)^4$은 $(gHu)^2=-\frac13s\cdot\sigma$를 고려하면 공식 (27.)c)와 (29.)c)에서 분해된다:
\[
(a_\nu^2\cdot H,s,u)^4\equiv\frac13(asu)^4\cdot\sigma,
\qquad
(\theta_\nu^2\cdot H,s,u)^4\equiv-\frac16(\theta su)^4\cdot\sigma;
\]
그로부터 $(a_\eta(aHu)\cdot H,s,u)^4$, $(\theta_\eta(\theta Hu)\cdot H,s,u)^4$의 분해가 따른다.

%
%
$(\theta_\nu^2\cdot H,s,u)^3$, $(\theta_\nu^3\cdot H,s,u)^3$, 그리고 그로부터 $(\theta_\eta^2(\theta Hu)\cdot H,s,u)^4$은 \S\ 25의 계수차수 12의 형식 C)에 따라 분해된다.

%
%
\noindent 3) \emph{$s$와 계 III$\cdot$III의 수축}, 즉 $s$와 계 III의 다음 형식들의 수축:
\[
(3,\lambda)\cdot(3,\lambda),\quad (3,\lambda)\cdot(2,\lambda),\quad (3,\lambda)\cdot(1,\lambda),\quad
(2,\lambda)\cdot(2,\lambda),\quad (2,\lambda)\cdot(1,\lambda),\quad
(1,\lambda)\cdot(1,\lambda).
\]
\noindent\emph{환원 불가능 형식들:}
\[
\begin{array}{rl}
\text{A)}&(a_\nu^2\cdot a_\nu^2,s,u)^3,
(a_\nu^2\cdot a_\nu^2,s,u)^4,
(a_\nu^2\cdot a_\eta(aHu)^2,s,u)^3,\\[0.45em]
\text{B)}&(a_\nu\cdot H_j,s,u)^4,
((aHu)^2\cdot H_j,s,u)^3,
((aLu)^2\cdot H_j,s,u)^4,\\[0.25em]
&((aLu)^3\cdot H_j,s,u)^2,
((aH^2u)^3\cdot H_j,s,u)^3.
\end{array}
\]

%
%
\noindent\emph{환원:}

기호 $\nu$와 $f$에서 계수차수가 같을 때, 곱 II$\cdot$III은 곱 III$\cdot$III보다 더 높은 것으로 간주한다.

환원들은 일부는 곱들 사이의 관계식(시지지)에서 생기고, 일부는 곱들을 대응하는 분해 가능 형식들로 바꾸고 이 형식들과 $s$의 수축을 환원함으로써 생긴다. (반공변식을 포함하는 곱들은 곱으로 넘어가므로 항상 생략한다.)

A) $f$에 관하여 이차이고, 기호 $\nu$는 임의의 계수차수이다. \S\ 11의 공식 (12.)와 유사한 계산에 따라 다음이 성립한다:
\begin{align*}
1)\quad 0&=a_\nu\cdot b_{\nu_1}+\frac12f\cdot(aHu)^2-\frac14\theta\cdot H+\frac12u_xH_\vartheta-\frac14u_x^2\cdot H_\vartheta^2,\\
2)\quad 0&=a_\nu\cdot b_{\nu_1}^2+f\cdot a_\eta(aHu)^2-\theta_\eta-\frac12H_\vartheta,\\
3)\quad 0&=a_\nu^2\cdot b_{\nu_1}^2-H_\vartheta^2+2\theta_\eta^2.
\end{align*}

%
%
그로부터 다음이 따른다:
\begin{align*}
4)\quad 0&=a_\nu^2\cdot b_{\nu_1}^2(j\nu_1x)\quad(3)\text{에서}),\\
5)\quad L_\vartheta(\theta Lu)&=-a_\nu\cdot b_\eta(bHu)^2+\frac34a_\nu^2\cdot(bHu)^2+\frac34\theta_\nu^2\cdot H\quad(2)\text{에서}),\\
6)\quad L_\vartheta(\theta Lu)&=-6a_\nu\cdot b_\eta(bHu)^2+2a_\nu^2\cdot(bHu)^2-\frac12\theta_\nu^2\cdot H\quad(1)\text{에서}),\\
7)\quad 0&=a_\nu\cdot(bHu)^2+f\cdot(aLu)^3+\theta_\nu\cdot H\quad(1)\text{에서}),
\end{align*}
\[
\left.
\begin{aligned}
8)\quad 0&=a_\nu\cdot(bLu)^3+\frac34(\theta Hu)^2\cdot H\\
9)\quad 0&=(aHu)^2\cdot(bHu)^2+2(\theta Hu)^2\cdot H
\end{aligned}
\right\}\quad(7)\text{에서}),
\]
\begin{align*}
10)\quad 0&=a_\eta(aHu)^2\cdot b_\eta(bHu)^2\quad(5)\text{와 }6)\text{에서}),\\
11)\quad 0&= [a_\nu^2\cdot b_\eta(bHu)]_{\widehat{su}^4}\quad(\text{공식 (30.)b에서}).
\end{align*}

%
%
다음 형식들의 환원
\[
(H_\vartheta su)^4,
\qquad (L_\vartheta(\theta Lu),s,u)^3,
\qquad (L_\vartheta(\theta Lu),s,u)^4
\]
(공식 (31.)과 (32.))은 공식 1)부터 11)까지와 함께, 기호 $f$에 관하여 이차이고 $\nu$에 관하여 임의인 곱들과 $s$의 모든 수축을 환원한다.

%
%
B) 기호 $f,\theta,j$ 중 각각 두 개의 곱이며, $\nu$는 임의의 계수차수이다. 공식 (17.), (18.), (20.)과 직접 계산에 따라 다음 관계식들이 성립한다:
\[
\left.
\begin{aligned}
1)\quad 0&=a_\nu\cdot\theta_{\nu_1}+\frac14\theta\cdot(aHu)^2+\frac14f\cdot(\theta Hu)^2,\\
2)\quad 0&=\theta_\nu\cdot\theta_{\nu_1}+\frac12\theta\cdot(\theta Hu)^2
\end{aligned}
\right\}\quad\text{(A 1)과 유사한 직접 계산).}
\]
그로부터 다음이 따른다:
\[
\left.
\begin{aligned}
3)\quad 0&=3a_\nu\cdot(\theta Hu)^2+\theta\cdot(aLu)^3+2f\cdot(\theta Lu)^3\\
4)\quad 0&=3\theta_\nu\cdot(aHu)^2+f\cdot(\theta Lu)^3+2\theta\cdot(aLu)^3
\end{aligned}
\right\}\quad\text{(1)과 A 6)에서),}
\]

%
%
\begin{align*}
5)\quad 0&=a_\nu\cdot(\theta Lu)^3,\quad 0=\theta_\nu\cdot(aLu)^3,\quad 0=(aHu)^2\cdot(\theta Hu)^2
\quad\text{(3), 4), A 8)에서),}\\
6)\quad 0&=a_\nu\cdot\theta_\nu^2+\theta_\nu\cdot a_\nu^2+f\cdot\theta_\eta(\theta Hu)^2+\theta\cdot a_\eta(aHu)^2\quad(\text{공식 (17.)}),\\
7)\quad 0&=\theta_\nu\cdot\theta_\nu^2+\theta\cdot\theta_\eta(\theta Hu)^2+\frac16f\cdot H_j\\
&\qquad [\text{을 }6)\text{에서 유도한 것과 유사한 유도}\theta_\nu^2(\theta\theta'\widehat{\nu}_1x)],\\
8)\quad 0&=a_\nu\cdot(j\nu x)^2+f\cdot H_j\quad(6)\text{에서}),\\
9)\quad 0&=(aHu)^2\cdot(j\nu x)^2-2a_\nu\cdot H_j\quad(8)\text{에서}),\\
10)\quad 0&=\theta_\nu\cdot(j\nu x)^2,\quad 0=\theta\cdot H_j\quad(7)\text{과 }8)\text{에서}),
\end{align*}
\[
\left.
\begin{aligned}
11)\quad 0&=a_\nu\cdot\theta_\nu^3-\frac34(j\eta x)^2,\\
12)\quad 0&=a_\nu^2\cdot\theta_\nu^2-\frac13(j\eta x)^2-\frac13(\D Hu)^2
\end{aligned}
\right\}\quad\text{(공식 (18.)),}
\]
\begin{align*}
13)\quad 0&=a_\nu^2\cdot\theta_\nu^3-\frac12a_\sigma\quad(\text{공식 (20.)}),\\
14)\quad 0&=\theta_\nu\cdot\theta_\nu^3,
&0&=a_\eta H_j\quad(\S\ 13.\ C)).
\end{align*}

%
%
공식들은 곱들이 편극 가능해질 때까지 전개하였다. 즉, 다음 이유로 수축을 통해 단 \emph{하나}의 새로운 곱만 생기는 때까지이다:
\begin{center}
\begin{tabular}{rl}
a)&곱 자체와의 수축이 환원 가능해지고,\\
b)&수축으로 생기는 나머지 곱들이 환원 가능해지거나
\end{tabular}
\end{center}
반공변식으로 이어진다(\S\ 3. 2), a)와 b) 참조).

공식 1)부터 5)까지는 $a_\nu\cdot\theta_\nu$에서 수축으로 생기는 곱들과 $s$의 모든 수축을 환원하고,

공식 6)부터 10)까지는 $a_\nu\cdot\theta_\nu^2$에서 수축으로 생기는 곱들을 환원한다.

\S\ 24 A)에 따르면, 공식 6)부터 10)까지를 고려할 때 $a_\nu^2\cdot\theta_\nu$에서 생기는 곱들과 $s$의 수축들은 분해된다.

%
%
다음을 얻는다:
\begin{align*}
0&\equiv a_\nu^2(asu)^2\theta_{\nu_1}(\theta su)^2,
&0&\equiv a_\nu^2(asu)\theta_{\nu_1}(\theta su)^3-K_\eta(KHu)^2(Ksu)^3,\\
0&\equiv a_\nu^2(asu)^2\theta_{\nu_1}(\theta su),
&0&\equiv a_\nu^2(asu)\theta_{\nu_1}(\theta su)^2,\\
0&\equiv a_\nu^2(asu)\theta_{\nu_1}(\theta su).
\end{align*}
환원자들
\[
((j\eta x)^2,s,u)^3
\ [\hbox{$(\vartheta\eta\widehat{su})^2(Hsu)$에서. \S\ 23.\ C\ 3)}],
\]
\[
(H_j(j\eta x),s,u)^2\ [\S\ 24,\ B],
\quad ((\D Hu)^2,s,u)^2\ [\S\ 24,\ C],
\]
\[
(a_\sigma su)^2\ [\S\ 24,\ E]
\]
은 공식 11)부터 14)까지와 함께 $a_\nu\cdot\theta_{\nu_1}^3$, $a_\nu^2\cdot\theta_{\nu_1}^2$, $a_\nu^2\cdot\theta_{\nu_1}^3$에서 생기는 모든 곱을 환원한다.

%
%
지금까지의 계산은 다음 \emph{결론}으로 요약된다:

상대 완전계 $\modu{(\R,t)}$는 다음 331개 형식과 부분계로 이루어지며, 이어서 이를 표로 배열하여 제시한다:
\[
\left.
\begin{array}{lr}
u_x&1\ \text{개 형식}\\
\text{계 I}&6\ \text{개 형식}\\
\text{계 II}&5\ \text{개 형식}\\
\text{계 III}&117\ \text{개 형식}
\end{array}
\right\}
\quad
\begin{gathered}
(s,\varrho,t)\text{를 법으로 하는 상대 완전계}\\
[\text{표 I 참조}].
\end{gathered}
\]

%
%
\[
\left.
\begin{array}{lr}
s&1\ \text{개 형식}\\
\text{계 }s_I&35\ \text{개 형식}\\
\text{계 }s_{II}&9\ \text{개 형식}\\
\text{계 }s_{III}&125\ \text{개 형식}\\
\text{계 }s_{IV}&32\ \text{개 형식}
\end{array}
\right\}
\quad [\text{표 II 참조}].
\]

%
%
\clearpage
\begin{landscape}
\thispagestyle{plain}
\begin{center}\bfseries 표 I (90쪽 참조).\end{center}
\vspace{-0.6em}
\begingroup
\setlength{\tabcolsep}{2pt}
\renewcommand{\arraystretch}{1.35}
\tiny
\begin{center}
\begin{adjustbox}{max width=0.98\linewidth,max totalheight=0.72\textheight,center}
\begin{tabular}{|c|c|c|c|c|c|c|c|c|c|c|c|c|c|c|c|c|c|c|c|c|c|c|}
\hline
{} & 0 & 1 & 2 & 3 & 4 & 5 & 6 & 7 & 8 & 9 & 10 & 11 & 12 & 13 & 14 & 15 & 16 & 17 & 18 & 19 & 20 & 21 \\ \hline
0 & \mcell{i} &  &  &  & \mcell{f} &  & \mcell{\A} &  & \mcell{K_{\nu}(k\nu x)^3} &  & \mcell{K_{\eta}(k\eta x)^3} &  &  &  & \mcell{(\vartheta^2\eta x)^4} & \mcell{H_k(k\cdot\vartheta,\eta x)^4} &  & \mcell{\theta_l(\vartheta\cdot k,lx)^5} &  &  &  & \mcell{(\vartheta^3lx)^6} \\ \hline
1 &  & \mcell{u_x} &  &  &  & \mcell{\theta_\nu(\vartheta\nu x)^2} &  & \mcell{\theta_\eta(\vartheta\eta x)^2} & \mcell{N} & \mcell{(k\nu x)^3} & \mcell{\theta_\nu(\vartheta^2\nu x)^3} & \mcell{(k\eta x)^3} & \mcell{H_{\vartheta}(\vartheta^2\eta x)^3\\ \theta_\eta(\vartheta^2\eta x)^3} &  & \mcell{\theta_l(\vartheta^2lx)^4} &  & \mcell{(\vartheta k,\eta,x)^4} &  & \mcell{(\vartheta k,l,x)^5} &  &  &  \\ \hline
2 &  &  & \mcell{a_\nu^2} &  & \mcell{\theta\\ a_\eta^2} &  & \mcell{(\vartheta\nu x)^2} & \mcell{K_\nu(k\nu x)^2} & \mcell{(\vartheta\eta x)^2} & \mcell{H_k(k\eta x)^2\\K_\eta(k\eta x)^2} &  & \mcell{(\vartheta^2\nu x)^3\\ K_l(klx)^3} &  & \mcell{(\vartheta^2\eta x)^3} &  & \mcell{(\vartheta^2lx)^4} &  &  &  &  &  &  \\ \hline
3 &  & \mcell{\theta_\nu^3} &  & \mcell{a_\nu} & \mcell{\theta_\nu(\vartheta\nu x)} & \mcell{\A_\nu\\a_\eta} & \mcell{K\\H_{\vartheta}(\vartheta\eta x)\\\theta_\eta(\vartheta\eta x)} & \mcell{\A_\eta} & \mcell{(k\nu x)^2\\\theta_l(\vartheta lx)^2} &  & \mcell{(k\eta x)^2} &  & \mcell{(klx)^3} &  &  &  &  &  &  &  &  &  \\ \hline
4 & \mcell{\nu} &  & \mcell{H\\\theta_\nu^2} & \mcell{(j\nu x)^3\\a_\eta(aHu)} & \mcell{\theta_\eta^2} & \mcell{(\vartheta\nu x)\\a_l^2} &  & \mcell{N_\nu\\(\vartheta\eta x)} &  & \mcell{H_n\\N_\eta\\(\vartheta lx)^2} &  &  &  &  &  &  &  &  &  &  &  &  \\ \hline
5 &  & \mcell{a_\eta(aHu)^2} & \mcell{\theta_\eta^2(\theta Hu)} & \mcell{\theta_\nu} & \mcell{K_\nu^2\\(aHu)} & \mcell{\theta_\eta} & \mcell{K_\eta^2\\(\A Hu)\\a_l} &  & \mcell{\A_l} &  &  &  &  &  &  &  &  &  &  &  &  &  \\ \hline
6 & \mcell{j\\\sigma} &  & \mcell{(j\nu x)^2\\(aHu)^2} & \mcell{L\\a_\nu^2(j\nu x)\\H_{\vartheta}(\theta Hu)\\\theta_\eta(\theta Hu)} &  & \mcell{K_\nu\\\theta_l^2} &  & \mcell{K_\eta} &  &  &  &  &  &  &  &  &  &  &  &  &  &  \\ \hline
7 &  & \mcell{\theta_\eta(\theta Hu)^2} & \mcell{H_j(j\eta x)\\a_l(aLu)^2} &  & \mcell{a_\nu(j\nu x)\\(\theta Hu)} &  & \mcell{\A_\nu(j\nu x)\\\theta_l} & \mcell{K_l^2} &  &  &  &  &  &  &  &  &  &  &  &  &  &  \\ \hline
8 & \mcell{H_j^2\\a_l(aLu)^3} & \mcell{(\nu\sigma x)\\(j\nu x)} & \mcell{(\theta Hu)^2} & \mcell{K_\eta(KHu)^2\\(j\eta x)\\(aLu)^2} &  &  &  &  &  &  &  &  &  &  &  &  &  &  &  &  &  &  \\ \hline
9 &  & \mcell{H_j\\(aLu)^3} & \mcell{a_\eta(aHu)H_j\\L_{\vartheta}(\theta Lu)^2\\\theta_l(\theta Lu)^2} &  & \mcell{(KHu)^2} &  &  &  &  &  &  &  &  &  &  &  &  &  &  &  &  &  \\ \hline
10 & \mcell{\theta_l(\theta Lu)^3} & \mcell{L_j^2\\(j\sigma x)\\a_\eta(aH^2u)^3} &  & \mcell{H_j(aHu)\\(\theta Lu)^2} &  &  &  &  &  &  &  &  &  &  &  &  &  &  &  &  &  &  \\ \hline
11 &  & \mcell{(\theta Lu)^3} & \mcell{K_l(KLu)^3\\L_j\\(aH^2u)^3} &  &  &  &  &  &  &  &  &  &  &  &  &  &  &  &  &  &  &  \\ \hline
12 & \mcell{(aH^2u)^4} & \mcell{a_l(aLu)^2L_j\\H_{\vartheta}(\theta H^2u)^3\\\theta_\eta(\theta H^2u)^3} &  & \mcell{(KLu)^3} &  &  &  &  &  &  &  &  &  &  &  &  &  &  &  &  &  &  \\ \hline
13 & \mcell{(\nu\sigma j)} &  & \mcell{(aLu)^2L_j\\(\theta H^2u)^3} &  &  &  &  &  &  &  &  &  &  &  &  &  &  &  &  &  &  &  \\ \hline
14 & \mcell{(\theta H^2u)^4} & \mcell{K_\eta(KH^2u)^4\\(a,H\!\cdot L,u)^4} &  &  &  &  &  &  &  &  &  &  &  &  &  &  &  &  &  &  &  &  \\ \hline
15 & \mcell{L_\vartheta(\theta,H\!\cdot L,u)^4} &  & \mcell{(KH^2u)^4} &  &  &  &  &  &  &  &  &  &  &  &  &  &  &  &  &  &  &  \\ \hline
16 &  & \mcell{(\theta,H\!\cdot L,u)^4} &  &  &  &  &  &  &  &  &  &  &  &  &  &  &  &  &  &  &  &  \\ \hline
17 & \mcell{K_\eta(K,H\!\cdot L,u)^5} &  &  &  &  &  &  &  &  &  &  &  &  &  &  &  &  &  &  &  &  &  \\ \hline
18 &  & \mcell{(K,H\!\cdot L,u)^5} &  &  &  &  &  &  &  &  &  &  &  &  &  &  &  &  &  &  &  &  \\ \hline
19 &  &  &  &  &  &  &  &  &  &  &  &  &  &  &  &  &  &  &  &  &  &  \\ \hline
20 &  &  &  &  &  &  &  &  &  &  &  &  &  &  &  &  &  &  &  &  &  &  \\ \hline
21 & \mcell{(KH^3u)^6} &  &  &  &  &  &  &  &  &  &  &  &  &  &  &  &  &  &  &  &  &  \\ \hline
\end{tabular}
\end{adjustbox}
\end{center}
\vspace{-0.2em}
\begin{center}\footnotesize (맨 위 가로행의 수들은 $x$에 관한 차수를, 첫 번째 세로행의 수들은 $u$에 관한 차수를 나타낸다.)\end{center}
\endgroup
\end{landscape}

%
%
\clearpage
\begin{landscape}
\thispagestyle{plain}
\begin{center}\bfseries 표 II (90쪽 참조).\end{center}
\vspace{-0.6em}
\begingroup
\setlength{\tabcolsep}{2pt}
\renewcommand{\arraystretch}{1.35}
\tiny
\begin{center}
\begin{adjustbox}{max width=0.98\linewidth,max totalheight=0.72\textheight,center}
\begin{tabular}{|c|c|c|c|c|c|c|c|c|c|c|c|c|c|c|c|c|c|}
\hline
{} & 0 & 1 & 2 & 3 & 4 & 5 & 6 & 7 & 8 & 9 & 10 & 11 & 12 & 13 & 14 & 15 & 16 \\ \hline
0 & \mcell{J} &  &  &  & \mcell{s} &  & \mcell{s_\vartheta^2} &  &  &  &  & \mcell{s_n} & \mcell{(k\nu x)^3s_\nu} & \mcell{(\vartheta^2\nu x)^3s_\nu s_\vartheta\\\theta_\nu(\vartheta^2\nu x)^3s_\vartheta} &  & \mcell{\theta_\eta(\vartheta^2\eta x)^3s_\vartheta} &  \\ \hline
1 &  &  &  &  &  &  & \mcell{(asu)} & \mcell{s_\vartheta} & \mcell{s_k^2\\(\Delta su)} & \mcell{s_n(Nsu)\\(\vartheta\nu x)^2s_\nu} & \mcell{((k\nu x)^3,s,u)s_\nu\\(K_\nu(k\nu x)^3,s,u)} &  & \mcell{(K_\eta(k\eta x)^3,s,u)} &  & \mcell{(\vartheta^2\nu x)^3s_\nu} &  & \mcell{((\vartheta^2\eta x)^4,s,u)} \\ \hline
2 &  &  &  &  & \mcell{(asu)^2} & \mcell{s_\vartheta(\theta su)} & \mcell{(\Delta su)^2\\a_\nu s_\nu} & \mcell{((\vartheta\nu x)^2,s,u)s_\nu\\(\theta_\nu(\vartheta\nu x)^2,s,u)} &  & \mcell{s_k\\(\theta_\eta(\vartheta\eta x)^2,s,u)} &  & \mcell{(k\nu x)^2s_\nu\\((k\nu x)^3,s,u)} &  & \mcell{((k\eta x)^3,s,u)} &  &  &  \\ \hline
3 &  &  & \mcell{(asu)^3} & \mcell{s_\vartheta(\theta su)^2\\s_\nu} & \mcell{a_\nu s_\nu(asu)\\a_\nu^2(asu)} & \mcell{s_\eta} & \mcell{(\theta su)\\(a_\eta^2su)} & \mcell{s_k(Ksu)} & \mcell{(Nsu)^2\\(\vartheta\nu x)s_\nu\\((\vartheta\nu x)^2,s,u)} & \mcell{((k\nu x)^3,s,u)^2} & \mcell{((\vartheta\eta x)^2,s,u)} &  &  &  &  &  &  \\ \hline
4 & \mcell{(asu)^4} & \mcell{s_\vartheta(\theta su)^3} & \mcell{a_\nu^2(asu)^2} & \mcell{(\theta_\nu^3su)} & \mcell{(\theta su)^2} & \mcell{s_k(Ksu)^2\\a_\nu(asu)\\s_\nu(asu)} & \mcell{((\vartheta\nu x)^2,s,u)^2\\\theta_\nu s_\nu} & \mcell{s_\nu(\Delta su)\\(\Delta_\nu su)\\(aHu)s_\eta\\(a_\eta su)} &  & \mcell{(\Delta_\eta su)} &  &  &  &  &  &  &  \\ \hline
5 &  &  & \mcell{(\theta su)^3} & \mcell{s_k(Ksu)^3,s_j,\\s_\vartheta(\theta s^\prime u)^4\\s_\nu(asu)^2\\a_\nu(asu)^2} & \mcell{(Hsu)\\((\vartheta\nu x)^2,s,u)^3\\(\theta_\nu(\vartheta\nu x),s,u)^2\\(\theta_\nu^2su)} & \mcell{((j\nu x)^3,s,u)\\(aHu)^2s_\eta\\(a_\eta su)^2} & \mcell{(Ksu)^2\\s_l\\(\theta_\eta^2su)} &  & \mcell{((k\nu x)^2,s,u)^2\\K_\nu s_\nu} &  &  &  &  &  &  &  &  \\ \hline
6 & \mcell{(\theta su)^4} & \mcell{s_\nu(asu)^3\\a_\nu(asu)^3} & \mcell{(Hsu)^2\\(\theta_\nu(\vartheta\nu x),s,u)^3\\(\theta_\nu^2su)^2} & \mcell{(a_\eta su)^3\\(a_\eta(aHu),s,u)^2\\(a_\eta(aHu)^2,s,u)} & \mcell{(Ksu)^3} & \mcell{s_j(asu)\\((\vartheta\nu x),s,u)^2\\(\theta_\nu su)} & \mcell{((k\nu x)^2,s,u)^3} & \mcell{s_j(\Delta su)\\a_\nu(j\nu x)s_\nu\\((\vartheta\eta x),s,u)^2\\(\theta Hu)s_\eta\\(\theta_\eta su)} &  &  &  &  &  &  &  &  &  \\ \hline
7 & \mcell{(\theta_\nu(\vartheta\nu x),s,u)^4} & \mcell{(a_\eta(aHu),s,u)^3} & \mcell{(Ksu)^4\\(\theta_\eta^2(\theta Hu),s,u)^2\\(a_\nu^2\cdot a_\nu^2,s,u)^3} & \mcell{s_j(asu)^2\\s_k(Ks^2u)^5\\((\vartheta\nu x),s,u)^3\\(\theta_\nu su)^2} & \mcell{(j\nu x)s_\nu\\((j\nu x)^2,s,u)\\((aHu),s,u)^2\\((aHu)^2,s,u)} & \mcell{((\vartheta\eta x),s,u)^3\\(\theta_\eta su)^2} & \mcell{(aLu)^2s_l\\(a_lsu)^2} &  &  &  &  &  &  &  &  &  &  \\ \hline

8 & \mcell{(a_\nu^2\cdot a_\nu^2,s,u)^4} & \mcell{s_j(asu)^3\\s_k(Ks^2u)^6\\((\vartheta\nu x),s,u)^4\\(\theta_\nu su)^3} & \mcell{((j\nu x)^2,s,u)^2\\((aHu),s,u)^3\\((aHu)^2,s,u)^2} & \mcell{(Lsu)^2\\(a_\nu^2(j\nu x),s,u)^2\\(H_\vartheta(\theta Hu),s,u)^2\\(\theta_\eta(\theta Hu),s,u)^2\\(\theta_\eta(\theta Hu)^2,s,u)} & \mcell{(a_lsu)^3} & \mcell{(K_\nu su)^2} &  & \mcell{(K_\eta su)^2} &  &  &  &  &  &  &  &  &  \\ \hline
9 & \mcell{(K_\nu^2su)^4\\((aHu),s,u)^4} & \mcell{(Lsu)^3\\(a_\nu^2(j\nu x),s,u)^3\\(\theta_\eta su)^4\\(\theta_\eta(\theta Hu),s,u)^3} & \mcell{(Ks^2u)^6\\(a_l(aLu)^2,s,u)^2\\(a_\nu^2\cdot H,s,u)^3} & \mcell{(K_\nu su)^3} & \mcell{(a_\nu(j\nu x),s,u)^2\\((\theta Hu),s,u)^2\\((\theta Hu)^2,s,u)} &  & \mcell{(\theta_lsu)^2} &  &  &  &  &  &  &  &  &  &  \\ \hline
10 &  & \mcell{(K_\nu su)^4\\(a_\nu^2\cdot a_\eta(aHu)^2,s,u)^3} & \mcell{(a_\nu(j\nu x),s,u)^3\\((\theta Hu),s,u)^3\\((\theta Hu)^2,s,u)^2} & \mcell{((j\eta x),s,u)^2\\(H_jsu)\\((aLu)^2,s,u)^2\\((aLu)^3,s,u)} &  &  &  &  &  &  &  &  &  &  &  &  &  \\ \hline
11 & \mcell{(a_\nu(j\nu x),s,u)^4\\((\theta Hu),s,u)^4} & \mcell{(K_\eta(KHu)^2,s,u)^3\\((j\eta x),s,u)^3\\((aLu)^2,s,u)^3\\(a_\nu\cdot H,s,u)^4} & \mcell{(H^2su)^3\\(\theta_l(\theta Lu)^2,s,u)^2} &  & \mcell{((KHu)^2,s,u)^2} &  &  &  &  &  &  &  &  &  &  &  &  \\ \hline
12 &  & \mcell{(a_\eta(aHu)^2\cdot H,s,u)^3} & \mcell{((KHu)^2,s,u)^3} & \mcell{((aHu)H_j,s,u)^2\\((\theta Lu)^2,s,u)^2\\((\theta Lu)^3,s,u)} &  &  &  &  &  &  &  &  &  &  &  &  &  \\ \hline
13 & \mcell{((KHu)^2,s,u)^4} & \mcell{((aHu)H_j,s,u)^3\\((\theta Lu)^2,s,u)^3\\(\theta_\nu\cdot H,s,u)^4} & \mcell{(L_jsu)^2\\((aH^2u)^3,s,u)^2\\((j\nu x)^2\cdot H,s,u)^3\\((aHu)^2\cdot H,s,u)^3} &  &  &  &  &  &  &  &  &  &  &  &  &  &  \\ \hline
14 & \mcell{((j\nu x)^2\cdot H,s,u)^4\\((aHu)^2\cdot H,s,u)^4} & \mcell{(\theta_\eta(\theta Hu)^2\cdot H,s,u)^3} &  & \mcell{((KLu)^3,s,u)^2} &  &  &  &  &  &  &  &  &  &  &  &  &  \\ \hline
15 & \mcell{(a_l(aLu)^2\cdot H,s,u)^4} & \mcell{((KLu)^3,s,u)^3} & \mcell{((aLu)^2L_j,s,u)^2\\((\theta H^2u)^3,s,u)^2\\((\theta Hu)^2\cdot H,s,u)^3} &  &  &  &  &  &  &  &  &  &  &  &  &  &  \\ \hline
16 & \mcell{((\theta Hu)^2\cdot H,s,u)^4\\(a_\nu\cdot H_j,s,u)^4} & \mcell{((j\eta x)\cdot H,s,u)^4\\(H_j\cdot H,s,u)^3\\((aLu)^2\cdot H,s,u)^4\\((aLu)^3\cdot H,s,u)^3} &  &  &  &  &  &  &  &  &  &  &  &  &  &  &  \\ \hline
17 & \mcell{(\theta_l(\theta Lu)^2\cdot H,s,u)^4} &  & \mcell{((KH^2u)^4,s,u)^2} &  &  &  &  &  &  &  &  &  &  &  &  &  &  \\ \hline
18 &  & \mcell{((\theta Lu)^2\cdot H,s,u)^4\\((\theta Lu)^3\cdot H,s,u)^3\\((aHu)^2\cdot H_j,s,u)^3} &  &  &  &  &  &  &  &  &  &  &  &  &  &  &  \\ \hline
19 & \mcell{((aH^2u)^3\cdot H,s,u)^4\\(L_j\cdot H,s,u)^4} &  &  &  &  &  &  &  &  &  &  &  &  &  &  &  &  \\ \hline
20 &  & \mcell{((KLu)^3\cdot H,s,u)^4} & \mcell{((aLu)^3\cdot H_j,s,u)^2} &  &  &  &  &  &  &  &  &  &  &  &  &  &  \\ \hline
21 & \mcell{((\theta H^2u)^3\cdot H,s,u)^4\\((aLu)^2\cdot H_j,s,u)^4} &  &  &  &  &  &  &  &  &  &  &  &  &  &  &  &  \\ \hline
22 &  &  &  &  &  &  &  &  &  &  &  &  &  &  &  &  &  \\ \hline
23 & \mcell{((KH^2u)^4\cdot H,s,u)^4} & \mcell{((aH^2u)^3\cdot H_j,s,u)^3} &  &  &  &  &  &  &  &  &  &  &  &  &  &  &  \\ \hline
\end{tabular}
\end{adjustbox}
\end{center}
\vspace{-0.2em}
\begin{center}\footnotesize (맨 위 가로행의 수들은 $x$에 관한 차수를, 첫 번째 세로행의 수들은 $u$에 관한 차수를 나타낸다.)\end{center}
\endgroup
\end{landscape}

\clearpage
\setcounter{footnote}{0}
\phantomsection
\addcontentsline{toc}{section}{3. n변수 형식의 불변식론에 관하여}
%
%

%
%
\begin{center}
{\Large\bfseries 3. $n$변수 형식의 불변식론에 관하여.\srcfn{1)}{1909년 잘츠부르크 자연과학자대회에서 행한 강연.}}\par
\vspace{1.2em}
에를랑겐의 \textsc{에미 뇌터}.\par
\medskip
Jahresbericht d. DMV 19 (1910), S. 101--104
\end{center}

\vspace{1em}
$n$변수 형식의 사영 불변식론에서는 주된 문제가 해결되었고, 형식계의 유한성이 증명되었다. 그러나 그 밖의 일련의 문제들, 곧 형식들의 연관관계를 연구하는 방법에 관한 문제들은 다루어지지 않았다. 그러한 문제들에 관해 내가 얻은 결과를 간략히 전하고자 한다. 다만 명료성을 위해 먼저 그 문제 설정을, 그것들이 잘 알려져 있는 삼원 영역에서 개략적으로 설명하겠다.

나는 먼저 기호적 방법의 두 기본정리를 생각한다. 첫째는 모든 불변 구성식을 기호적으로 표현할 수 있다는 정리이고, 둘째는 불변식들 사이에 성립하는 모든 관계를 소수의 기호적 항등식들을 차례로 적용하여 얻을 수 있다는 정리이다. 따라서 기호적 방법은 불변식의 영역을 벗어나지 않고도 모든 관계를 제공한다.\srcfn{1)}{Study: \emph{Methoden zur Theorie der ternären Formen}. II. \S\ 6. (Leipzig, Teubner 1889.)} 이를 토대로 형식을 생성하는 과정들, 곧 기호적 수축과정과 비기호적 미분과정, 급수전개 이론 등이 이어진다. 이와 관련하여 나의 학위논문\srcfn{2)}{\emph{Über die Bildung des Formensystems der ternären biquadratischen Form}. \S\ 1. (Crelle's Journal Bd. 134. 1908.)}에서 증명한 정리 하나를 더 언급하고자 한다. 그 정리는 가능한 두 ``동변'' 수축을 차례로 적용함으로써 모든 수축을 생성할 수 있다는 것이다. 여기서 내가 말하는 두 수축은 기호적 곱 $s_x t_x u_\sigma u_\tau$에 대한 $(stu)$와 $(\sigma\tau x)$이다. 이 정리를 토대로, 나의 학위논문에서 보인 것처럼 이원 영역과 유사하게 환원식들의 이론과 형식계의 환원 이론을 세울 수 있다.

%
%
이제 $n$원 영역으로 넘어가면, 기호적 방법의 첫 번째 정리조차도 제한된 의미에서만 성립한다. 잘 알려져 있듯이 4원 형식에서는 점좌표와 평면좌표 $x$와 $u$ 외에도 선좌표 $p_{ik}$가 나타나며, 이들은 각각 두 점좌표열로 만든 소행렬식이다. 이와 마찬가지로 $n$원 영역에서는 변수열 $p_\varrho$가 있는데, 그 각각의 원소는 $\varrho$개의 $x$열로 만든 소행렬식이며, 또 기호열 $S^\varrho$가 있는데, 그 원소들은 $u$와 같은 방식으로 변환되는 $\varrho$개의 $s$열로 만든 소행렬식처럼 거동한다. 잘 알려져 있듯이 행렬식에 의한 기호적 표현은 기호열 $S^\varrho$를 $s$열들로 풀어내고 이것들을 서로 다른 행렬식들에 분배할 때에만 성립한다. 그러면 $S^\varrho$에 의존하는 행렬식 결합식들만이 실질적인 의미를 갖는다. 그러나 어떤 행렬식 결합식들이 이를 충족하는지는 한눈에 알 수 없으며, 그 결과 3원 영역에서 언급한 나머지 정리들은 모두 사라진다.

이제 \emph{기호열 $S^\varrho$가 명시적으로 나타나는 기호적 표현}을 얻고 그와 함께 나머지 정리들도 다시 얻기 위한 간단한 원리를 제시하고자 한다. 이것은 대응 행렬에 관한 잘 알려진 정리의 응용이다: ``각 변수열 $p_\varrho$에는 다른 변수열 $q_{n-\varrho}$를 일대일로 대응시킬 수 있는데, 그 각각의 원소는 $n-\varrho$개의 $u$열로 만든 소행렬식이며, 이 $u$열들은 방정식 $(u\mid x)=0$에 의해 $\varrho$개의 $x$열과 연결된다.'' 이에 대응하여 각 기호열 $S^\varrho$에는 다른 기호열 $S_{n-\varrho}$가 대응되며, 이는 $x$와 같은 방식으로 변환되는 $n-\varrho$개의 열로 구성된 것처럼 거동한다.

이 정리는 응용에 적합한 또 다른 형태로도 진술할 수 있다: ``두 번째 행렬의 열들이 첫 번째 행렬의 열들과 반변 관계에 있는 모든 행렬곱 $(v^{(1)}\ldots v^{(\varrho)}\mid y^{(1)}\ldots y^{(\varrho)})$,\srcfn{1)}{즉, $\varrho$개의 $v$열로 이루어진 행렬과 $\varrho$개의 $y$열로 이루어진 행렬에서 얻은 서로 대응하는 행렬식들의 곱을 합한 것이다.} 에는 하나의 행렬식이 대응한다. 역으로 모든 행렬식은 미리 주어진 열의 개수 $\varrho$를 갖는 행렬곱으로 변환할 수 있다.'' 따라서 행렬식과 행렬곱은 완전히 동등한 것으로 나타나며, 기호적 방법의 첫 번째 정리는 다음과 같은 형태를 취한다: ``모든 불변 구성식은 행렬곱으로 기호적으로 나타낼 수 있다.'' 이제 두 기호열 $S^\sigma$, $T^\tau$와 변수열들을 이용하면 이 기호열들을 명시적으로 포함하는 행렬곱을 아주 쉽게 만들 수 있다. 따라서 이것은 $(S^\sigma\mid p_\sigma)(T^\tau\mid p_\tau)$ 형태의 불변 구성식을 나타내며, 다음과 같다:
\srcnumdisplay[3.2em]{(1)}{%
  \pair{q_{n-\sigma-\lambda}S^\sigma}{T_{n-\tau}p_{\tau-\lambda}},
  \qquad (\lambda=1,2\ldots\tau\ \text{또는}\ n-\sigma).
}

%
%
더 중요한 것은 그 역, 곧 \emph{모든 불변 구성식을 행렬곱으로 명시적으로 표현할 수 있다}는 사실이며, 따라서 위 과정은 이원 및 삼원 수축과정의 확장을 뜻한다. 이 역을 증명하려면 기호적 방법의 두 번째 정리를 \(n\)원 영역으로 옮겨야 한다. 즉 모든 열 \(S^\varrho\), \(p_\varrho\)를 비분해 대상으로 간주할 때의 서로 독립인 비분해 항등식 전체를 제시해야 한다. 이 항등식들은 \(x\)와 \(u\)의 열만을 포함하는 알려진 항등식들로부터\srcfn{2)}{E. Pascal in \emph{Memorie d. R. Acc. d. Lincei}. (4) 5 (1888), p. 375.} 얻어진다. 이때 행렬식의 곱정리를 행렬에 관한 곱정리 체계로 바꾸고, 바로 이 곱정리들을 이용하여 서로 다른 행렬곱들을 하나로 결합한다. 이렇게 하여 두 쌍대적 공식에 이르는데, 여기서는 그중 하나만을 제시한다.
\srcnumdisplay[3.2em]{(2)}{%
\begin{aligned}
 \pair{S^{\varrho_1}S^{\varrho_2}\ldots S^{\varrho_k}}{x p_{\varrho-1}}
 &=\sum_{i=1}^{k}\eps\,\bigl(S^{\varrho_1}\ldots S^{\varrho_{i-1}}S^{\varrho_{i+1}}\ldots S^{\varrho_k}q_{n-\varrho+1}\mid S_{n-\varrho_i}x\bigr);\\
 \varrho&=\sum_{i=1}^{k}\varrho_i\leqq n;\qquad \eps=\pm1.
\end{aligned}
}

이 공식들에 포함된 유일한 특수화, 곧 하나의 \(x\)열이 들어간다는 조건은 피할 수 없다. 완전히 일반적인 기호열과 변수열을 취하면 하나의 \(p_\varrho\)열을 개별 \(x\)열들로 분해하는 항등식인 ‘분해항등식’에 이르게 된다. Clebsch는 이미 자신의 ‘\emph{Fundamentalaufgabe der Invariantentheorie}’에서 급수전개의 기본 공변식을 유도하는 데 분해항등식의 가장 간단한 특수 경우를 사용하였다. 비명시적 행렬식 표현에서 행렬곱으로의 이행을 매개하는 이 항등식들을 이용하면, 행렬표현을 통해 원하는 명시적 기호표현이 실제로 얻어졌음을 보일 수 있다.

한편 수축과정에 대해서도 삼원 영역에서와 완전히 유사한 정리가 성립하며, 그 위에 환원정리들도 마찬가지 방식으로 세울 수 있다. (1)에서 수 \(\lambda\)를 수축의 결손이라 하고, \(\sigma=\tau\)인 수축을 동변이라 부르면, \emph{모든 수축은 결손 1인 동변 수축들을 차례로 적용하여 생성할 수 있다}. 이 정리의 증명은 가장 일반적인 형식들을 ‘정규형식’, 즉 모든 불변 미분연산을 적용하면 항등적으로 소멸하는 형식들로 바꿀 수 있다는 사실에 본질적으로 의존한다. 사원 영역에서는 Mertens가 이 마지막 사실을 증명하였다\srcfn{1)}{Über invariante Gebilde quaternärer Formen. Wiener Berichte. Bd. 98. 1889.}. 가장 일반적인 미분연산은 기호열을 미분기호로 바꾸어 수축과정으로부터 생긴다는 것이 알려져 있다. 이 미분연산들에 대한 지식을 이용하면, 분해항등식을 적용하여 Mertens의 증명을 거의 그대로 \(n\)원 영역으로 옮길 수 있다.

\clearpage
\setcounter{footnote}{0}
\phantomsection
\addcontentsline{toc}{section}{4. n개 변수의 형식에 관한 불변식론}
%
%

%
%
\begin{center}
{\Large\bfseries $n$개 변수의 형식에 관한 불변식론\footnote{잘츠부르크 자연과학자 대회에서 발표한 짧은 보고를 참조하라. \emph{Jahresbericht d. deutsch. Math.-Vereinigung}, Bd.\ 19, S.\ 101.}}\par
\vspace{0.5em}
에를랑겐의 \emph{Emmy Noether} 양.\par
\medskip
\rule{0.22\linewidth}{0.4pt}\par
\vspace{0.7em}
Journal f.\ d.\ reine u.\ angew.\ Math.\ 139 (1911), S.\ 118--154
\end{center}

\begin{center}
\textbf{서론.}
\end{center}
$n$개 변수의 형식에 관한 사영 불변식론에서는, 이항 및 삼항 영역을 넘어서면 주된 문제, 즉 형식계의 유한성 증명에 뒤따르는 문제들이 거의 다루어지지 않았다. 그것은 형식들 사이의 상호 의존성과 이 연관관계를 다루는 방법에 관한 문제들이다. 이 방법들은 본질적으로 --- 이항 및 삼항 영역에서는 이것이 완전히 증명되어 있듯이 --- Study가 \glqq{}기호적 방법의 기본정리\grqq{}라고 부른 두 정리, 곧 형식의 기호적 표현 가능성에 관한 정리와 기호 항등식에 관한 정리에 기초한다. 뒤의 정리는 정수 불변 구성식들 사이에 성립하는 모든 관계를 소수의 기호 항등식을 차례로 적용하여 얻는다는 것, 따라서 기호적 방법만이 불변식의 영역을 벗어나지 않고 형식들의 연관관계를 파악하게 한다는 것을 말한다.\footnote{E. Study, \emph{Methoden zur Theorie der ternären Formen} (Leipzig, Teubner, 1889). II. \S\ 6. --- 특수 변환군의 불변식론에서도 이 기본정리들이 나타나는 것에 관해서는 Study, \emph{Das Apollonische Problem} (Math. Ann. Bd. 49 [1897]) 및 \emph{Zur Differentialgeometrie der analytischen Kurven} (Transactions of the American Math. Society, Bd. 10 [1909])을 참조하라.}

$n$개 변수로 넘어갈 때 생기는 어려움의 근거는 Gordan이 이미 자신의 에를랑겐 강령에서 지적하였다.\footnote{P. Gordan, \emph{Über das Formensystem binärer Formen}. S. 46. (Leipzig, Teubner, 1875).} 그것은 기호적 방법의 첫째 정리가 그 일반적 형태로는 더 이상 성립하지 않는 데 있다. 잘 알려져 있듯이 $n$항 영역에서는 더 높은 단계의 변수열 $p_\rho$\footnote{표기는 \S\ 1을 참조하라.}가 나타나며, 이와 마찬가지로 더 높은 단계의 기호열 $S^\rho$도 나타난다. 실제로 처음 Clebsch가 했듯이\footnote{A. Clebsch, Über eine Fundamentalaufgabe der Invariantentheorie. (Abh. der Ges. d. Wiss. zu Göttingen, Bd. XVII, 1872).} 임의의 개수의 열 $p_\rho$를 지닌 형식에서 출발하든, F. Deruyts와 A. Capelli의 방식에 따라 서로 동변인 임의의 개수의 열 $x$만을 지닌 형식에서 출발하든, 이 열들에 도달한다.\footnote{A. Capelli, \emph{Lezioni sulla teoria delle forme algebriche} (Napoli 1902), \S\ 22--24 및 거기에 제시된 문헌을 참조하라.} 그런데 열 $p_\rho,S^\rho$를 명시적으로\footnote{명시적 표현의 더 정확한 정의는 \S\ 2에 있다.} 포함하는 기호적 표현은 존재하지 않는다. 따라서 $p_\rho,S^\rho$를 각각의 열 $x$와 이에 반변인 열 $s$로 분해하고, 이들을 여러 행렬식에 분배해야 한다. 물론 미분 조건(공식 (19.) 참조)을 이용하면 주어진 행렬식 결합식이 오직 열 $p_\rho,S^\rho$에만 의존한다는 성질을 검증할 수 있다. 그러나 이 방법으로는 불변 구성식 전체와 그것을 생성하는 과정 전체를 조망할 수 없다.\footnote{Study에서 유래한, 계산상 매우 유용한 기호법의 한 발전(전달: F. Engel, Ber. d. K. Sächs. Ges. d. Wiss., math.-phys. Kl., 1900, S. 126; 4항 영역에 대해서는 E. Study, \emph{Geometrie der Dynamen}, S. 135)도 여기서 말하는 의미의 명시적 표현은 아니다. 예컨대 그것으로부터 형식을 생성하는 비기호적 미분 과정을 얻기는 어려울 것이다.}

이제 우리는 \S\S\ 2와 6에서 \glqq{}대응 행렬에 관한 정리\grqq{}\footnote{1862년 Graßmann의 확장론 Nr. 112를 제외하면, A. Brill, Über zwei Berührungsprobleme, \S\ 2 (Math. Ann. Bd. 4. 1871)에 처음 나타난다.}를 적용하여, 행렬식에 의한 표현을 \glqq{}행렬곱\grqq{}에 의한 표현으로 바꿈으로써 첫째 기본정리를 일반적인 형태로 다시 얻는 방법을 보인다. \S\ 2에서는 열 $S^\rho,p_\rho$를 명시적으로 포함하는 모든 행렬곱이 기본형의 불변 구성식임을 보인다. \S\ 6에 주어진 역, 즉 모든 불변 구성식이 행렬곱으로 명시적으로 표현된다는 사실은 둘째 기본정리를 전이한 뒤에야(\S\S\ 3--5) 얻어진다.

이 정리는 변수열 $x$만을 포함하는 형식에 대해서는 알려져 있다.\footnote{E. Pascal, Giorn. di mat. 26 (1888), S. 33, 102; Rendic. d. Accad. d. Linc. (4) 4 (1888), S. 119; ib. Mem. (4) 5 (1888), S. 375.} 모든 열 $S^\rho,p_\rho$를 분해 불가능한 것으로 보았을 때의 분해 불가능한 항등식 전체를 세우는 일반 문제는 Study가 정식화하였다.\footnote{\glqq{}Methoden \dots\grqq{}, S. 204, Anmerkung 17).} 그는 동시에 나타나는 항등식의 개수를 $n$항 영역에서 방법의 어려움을 재는 척도로 본다. 이제 항등식 전체를 하나의 공식 (36.)에 모을 수 있으므로, 이 어려움은 적어도 최소한으로 줄어든다. 다만 실제 적용에서는 (36.)의 몇몇 두드러진 특수한 경우가 자주 나타난다. 예컨대 새 열을 대입하지 않고도 명시적 표현을 허용한다는 특징을 지닌 (20.), (21.), 또는 열 $p_\rho$가 겉보기에는 열 $y$들로 분해되어 나타나므로 \glqq{}분해 항등식\grqq{}이라 부르는 공식 (31.)이 그러하다.

이 두 기본정리 위에서 나머지 이론은 쉽게 세울 수 있다. 첫째 정리는 형식을 생성하는 과정, 즉 기호적 수축 과정과 비기호적 미분 과정을 직접 준다.\footnote{\glqq{}Invariante Gebilde von Nullsystemen\grqq{} (Sitzungsber. d. Ak. d. Wiss. Wien, Bd. 97, Abt. IIa, 1888)에서 Mertens는 다른 방식, 곧 직접 일반화할 수 없는 방식으로 4항 수축 과정에 도달하였다. 거기서 나타나는 여섯 열 $S^2$의 교대 불변식은 여기서 대입 (11.)을 한 번 적용하여 생기는 불변식과 짝수 불변식의 한 결합식만큼 다르다. 특수한 한 경우의 4항 수축 과정은 P. Gordan, \emph{Das simultane System von zwei quadratischen quaternären Formen} (Math. Ann. Bd. 56. 1903)에도 나타난다.} 한편 분해 항등식은 이 과정들 가운데 서로 동치인 것을 모두 제거한다(\S\ 6). 더 나아가 미분 과정을 알면 \glqq{}정규형\grqq{}이라는 개념을 $n$항 영역으로 옮길 수 있고, Mertens가 4항 영역에서 제시한 증명법\footnote{F. Mertens, Über invariante Gebilde quaternärer Formen. (Sitzber. d. Ak. d. Wiss. Wien. Bd. 98. Abt. IIa. 1889.)}을 써서 불변 구성식의 한 계를 그와 동치인 정규형의 계로 바꿀 수 있다는 정리를 얻는다(\S\ 7). 나는 이러한 마지막 가정 아래에서, 삼항 수축 과정을 두 기본 수축을 차례로 적용하여 생성할 수 있음을 학위논문에서 보였다.\footnote{Über die Bildung des Formensystems der ternären biquadratischen Form. (Dieses Journal, Bd. 134. 1908.)} 이 정리와 급수 전개 이론에 기초하여 \glqq{}형식열\grqq{} 개념을 도입하고 형식계의 주요 환원정리를 전개하였다. 완전히 마찬가지로 $n$항 영역에서는 수축 과정을 $n-1$개의 기본 수축으로부터 생성할 수 있고(\S\ 8), 환원정리를 세울 수 있다(\S\ 9).

\paragraph{후기.} 잘츠부르크 보고 때 빈의 E. Müller 교수는 이 논문의 바탕을 이루는 행렬곱 계산이 원리상 Graßmann의 확장론 계산법과 동일하다는 점을 내게 알려 주었다.\footnote{Hermann Graßmanns gesammelte mathematische und physikalische Werke. Ersten Bandes zweiter Teil: Die Ausdehnungslehre von 1862. In Gemeinschaft mit H. Graßmann d. Jüngeren herausgegeben v. Fr. Engel (Leipzig, Teubner 1896).} 실제로 Graßmann의 \glqq{}조합곱\grqq{} $[S^{\rho_1}S^{\rho_2}\cdots S^{\rho_i}]$은 이 열들로 주어진 행렬로 볼 수 있다(\S\ 2 참조). 이때 \glqq{}전진곱\grqq{}은 열 자체의 행렬에, \glqq{}후퇴곱\grqq{}은 대응하는 열 $S_{n-\rho_1},S_{n-\rho_2},\ldots S_{n-\rho_i}$의 행렬에 해당한다. \glqq{}혼합곱\grqq{}에 해당하는 행렬은 여러 열 $S$를 대입 (9.)로 새 열 $P,Q$에 모아 만들게 된다. 여기서 우리의 \glqq{}대응 열\grqq{}은 Graßmann의 \glqq{}보충\grqq{}에, 우리의 \glqq{}행렬곱\grqq{}은 \glqq{}단계수가 0인 혼합곱\grqq{}에 해당한다.

이 대응 아래에서 1862년 확장론 Nr. 113에 주어진 전개는 우리의 공식 (32.)와 쌍대인 공식과 동일해진다. 특수한 경우 $\rho=1$에서 (32.)는 (17.)로, 쌍대 공식은 (16.)으로 바뀐다. 최근 Müller는 Graßmann 전개의 이 특수한 경우들로부터\footnote{E. Müller, Beiträge zur Graßmannschen Ausdehnungslehre. 1. Mitteilung: Einige allgemeine Sätze. (Sitzungsber. d. Ak. d. Wiss. Wien; Bd. 118. Abt. IIa. Juli 1909), Satz 16 und 18.} 앞서 말한 항등식 (20.), (21.)을 유도하였다.

우리의 유도는 Graßmann--Müller의 유도와 달리 항등식의 완전성까지 동시에 증명한다. 여기서 다루는 문제에는 바로 이것이 가장 본질적이다. 또한 우리의 유도는 (20.), (21.)에서 출발하여 지금까지 주목되지 않은 듯한 가장 일반적인 분해 불가능 항등식 (36.)으로 나아간다.

%
%
\section*{\S\ 1. 표기와 정의. 알려진 결과의 요약.}
통상적인 방식대로 원소 $x_1,x_2,\ldots,x_n$의 열을 변수열 $x$라 부른다. 이와 마찬가지로 변수열 $p_\rho$ ($\rho=1,2,\ldots,n-1$)라 함은 $\binom n\rho$개의 원소 $p_{i_1i_2\ldots i_\rho}$로 이루어진 열을 뜻한다. 여기서 $p_{i_1i_2\ldots i_\rho}$는 $\rho$개의 열 $x^{(1)},x^{(2)},\ldots,x^{(\rho)}$의 행렬에서 얻는 다음 $\binom n\rho$개의 행렬식이다.
\[
\left|\begin{array}{cccc}
 x_{i_1}^{(1)}&x_{i_2}^{(1)}&\cdots&x_{i_\rho}^{(1)}\\
 x_{i_1}^{(2)}&x_{i_2}^{(2)}&\cdots&x_{i_\rho}^{(2)}\\
 \vdots&\vdots&&\vdots\\
 x_{i_1}^{(\rho)}&x_{i_2}^{(\rho)}&\cdots&x_{i_\rho}^{(\rho)}
\end{array}\right|,
\qquad (i_1<i_2<\cdots<i_\rho=1,2,\ldots,n).
\]
대응 행렬에 관한 정리\footnote{예를 들어 Gordan--Kerschensteiner, \emph{Vorlesungen über Invariantentheorie}, Bd. I, S. 99를 참조하라.}에 따르면, 각 열 $p_\rho$에는 그 열의 각 원소마다 성립하는 다음 관계를 통하여 둘째 열 $q_{n-\rho}$가 일대일로 대응한다.
\begin{equation}
 p_{i_1i_2\ldots i_\rho}=(-1)^{\frac{\rho(\rho+1)}{2}+i_1+i_2+\cdots+i_\rho}q_{j_1j_2\ldots j_{n-\rho}},
\tag{1}
\end{equation}
여기서 $i$들과 $j$들은 각각 오름차순으로 놓여 있고, 이들을 합하면 1부터 $n$까지의 수를 한 번씩 모두 갖는 순열이 된다. 열 $q_{n-\rho}$는 $(n-\rho)$차 행렬식 $q_{j_1j_2\ldots j_{n-\rho}}$ $\binom n\rho$개로 이루어진다. 이 행렬식들은 $x$에 반변인 $n-\rho$개의 열 $u^{(1)},u^{(2)},\ldots,u^{(n-\rho)}$의 행렬에서 만들어진다. 이 열들은 다음 $\rho\cdot(n-\rho)$개의 조건방정식을 만족한다.
\begin{equation}
 (u^{(k)}\mid x^{(l)})=\sum_{\alpha=1}^{n}u_\alpha^{(k)}x_\alpha^{(l)}=0;
 \qquad (k=1,2,\ldots,n-\rho;\ l=1,2,\ldots,\rho).
\tag{2}
\end{equation}
우리는 Clebsch를 따라\footnote{a. a. O. \S\ 5.} 이 열들을 정규화하여, (1.)에 원래 들어가는 비례인자를 1로 놓기로 한다.

%
%
Clebsch에 따르면\footnote{a. a. O. \S\ 12.}, 그리고 뒤에 나온 F. Deruyts와 A. Capelli의 연구에 따르면(서론 참조), 가장 일반적인 형식들은 $n-1$개의 각 열 $p_\rho$를 기껏해야 하나씩만 포함하는 형식들로 환원할 수 있다. 따라서 이들은 기호적으로 다음과 같이 표현된다.
\begin{equation}
 (S^1\mid p_1)^{m_1}(S^2\mid p_2)^{m_2}\cdots(S^{n-1}\mid p_{n-1})^{m_{n-1}},
\tag{3}
\end{equation}
여기서 (2.)와 유사하게 다음과 같이 놓았다.
\[
 (S^\rho\mid p_\rho)=\sum S^{i_1i_2\ldots i_\rho}p_{i_1i_2\ldots i_\rho}.
\]
한 열 $p_\rho$의 원소들 사이에 성립하는 비선형 항등관계 때문에, 기호열 $S^\rho$는 바로 그 항등관계들을 만족하도록 정규화할 수 있다. 즉 이 열은 $x$에 반변인 열 $s^{(1)},s^{(2)},\ldots s^{(\rho)}$로 이루어진 $\rho$행 행렬의 행렬식처럼 행동한다.\footnote{Clebsch, a. a. O. \S\ 4.} 따라서 각 열 $S^\rho$에는 다음 관계로 다른 한 열 $S_{n-\rho}$가 일대일로 대응한다.
\begin{equation}
 S_{i_1i_2\ldots i_{n-\rho}}=(-1)^{\frac{(n-\rho)(n-\rho+1)}2+i_1+i_2+\cdots+i_{n-\rho}}S^{j_1j_2\ldots j_\rho},
\tag{4}
\end{equation}
여기서 다시 $i$들과 $j$들은 각각 오름차순으로 놓여 있으며, 합쳐서 1부터 $n$까지의 수를 정확히 모두 소진한다. 열 $S_{n-\rho}$의 원소 $S_{i_1i_2\ldots i_{n-\rho}}$는 $x$에 동변인 $n-\rho$개의 열 $\sigma^{(1)},\sigma^{(2)},\ldots\sigma^{(n-\rho)}$로 만든 행렬식처럼 행동하며, 열 $s$와 $\sigma$는 다음 $\rho\cdot(n-\rho)$개의 조건으로 연결된다.
\begin{equation}
 (s^{(k)}\mid \sigma^{(l)})=0.\qquad(k=1,2,\ldots\rho;\ l=1,2,\ldots n-\rho).
\tag{5}
\end{equation}
관계 (1.)과 (4.)에 따라 (3.)의 각 인자는 다음 네 가지 서로 동치인 기호적 표현을 갖는다.
\begin{equation}
 (S^\rho\mid p_\rho)=(S^\rho q_{n-\rho})=(S_{n-\rho}p_\rho)=(-1)^{\rho(n-\rho)}(q_{n-\rho}\mid S_{n-\rho}\footnotemark),
\tag{6}
\end{equation}
\footnotetext{Clebsch, a. a. O. \S\ 5를 참조하라.}
이하에서 언제나 그렇듯이, $(S^\rho q_{n-\rho})$와 $(S_{n-\rho}p_\rho)$는 각각 행렬식
$(s^{(1)}\cdots s^{(\rho)}u^{(1)}\cdots u^{(n-\rho)})$와
$(\sigma^{(1)}\cdots\sigma^{(n-\rho)}x^{(1)}\cdots x^{(\rho)})$의 약기호이다. 라플라스 전개를 하면 이 행렬식들이 각각 열 $S^\rho,q_{n-\rho}$ 및 $S_{n-\rho},p_\rho$에만 의존한다는 것이 드러나며, 위 표기는 바로 이것을 나타내려는 것이다. 앞의 설명에 따라 $(S^\rho\mid p_\rho)$와 $(q_{n-\rho}\mid S_{n-\rho})$는 각각 열 $s$와 $x$의 행렬, 그리고 열 $u$와 $\sigma$의 행렬을 곱한 것을 뜻한다. 여기서 행렬곱이란 서로 대응하는 행렬식들의 곱을 모두 더한 것, 즉 곱행렬의 행렬식 값을 말한다.

%
%
기호열(또는 변수열)에 속하는 특성수 $\rho\cdot(n-\rho)$를 Clebsch를 따라 그 열의 무게라 부르고,\footnote{a. a. O. \S\ 11.} 열 $s$의 개수를 나타내는 수 $\rho$를 위 무게지수, 열 $\sigma$의 개수를 나타내는 수 $n-\rho$를 아래 무게지수라 부른다. 관계 (4.)에 따르면 하나의 기호열은 같은 관계 안에서 한 번은 위 무게지수를, 한 번은 아래 무게지수를 달고 나타날 수 있다. 서로 짝을 이루는 열 $p_\rho$와 $q_{n-\rho}$에도 같은 사실이 성립한다. 또한 일반적으로 다음과 같이 표기하기로 한다. 대응 행렬이 열 $x$들로 이루어진 더 높은 단계의 변수열은 $p$로, 그 행렬이 열 $u$들로 이루어진 것은 $q$로 나타낸다. $x$에 동변인 기호열은 작은 그리스 문자 $\sigma,\tau,\alpha,\beta$로, $u$에 동변인 기호열은 작은 라틴 문자 $s,t,a,b$로 나타낸다. 그 밖의 모든 기호열은 무게지수를 붙인 큰 라틴 문자 $S^\rho,T^\tau,A^\rho$ 또는 $S_{n-\sigma},T_{n-\tau},A_{n-\rho}$로 나타낸다. 위 무게지수 또는 아래 무게지수에 맞추어, (4.)에서처럼 한 열의 개별 원소도 위 첨자 또는 아래 첨자로 쓴다.

%
%
\section*{\S\ 2. 행렬곱에 의한 표현.}
이하에서 우리는 언제나 ``행렬곱''을, 같은 수의 열로 이루어진 두 행렬에서 서로 대응하는 행렬식들의 곱을 모두 합한 것으로 이해한다. 이때 둘째 행렬의 열들은 첫째 행렬의 대응 열들에 반변한다. 각 행렬의 열들은 1. 하나씩 주어질 수도 있고, 2. (6.)에서처럼 하나의 열 $S^\rho$로 묶일 수도 있으며, 3. 서로 다른 여러 열 $S^{\rho_1},S^{\rho_2}\ldots$로 합쳐질 수도 있다. 행렬곱은 (6.)에서 사용한 기호 $(\ \mid\ )$로 나타내기로 한다. 표현 (6.)은 대응 행렬에 관한 정리의 둘째 정식화를 특수하게 적용한 것이다. ``모든 $\rho$열 행렬곱 $(v^{(1)}v^{(2)}\cdots v^{(\rho)}\mid y^{(1)}\cdots y^{(\rho)})$에는 하나의 행렬식이 대응한다. 역으로 모든 행렬식에는 미리 정한 열의 수 $\rho$ ($\rho=1,2,\ldots,n-1$)를 가진 하나의 행렬곱을 대응시킬 수 있다.''\footnote{$\rho>n$이면 행렬곱은 잘 알려져 있듯이 소멸하고, $\rho=n$이면 두 행렬식의 곱이 된다.} 실제로 주어진 행렬곱에서 행렬식으로, 그리고 그 역으로 넘어가기 위해서는 열 $y$들을 하나의 열 $p_\rho$로 묶거나(각각 열 $v$들을 하나의 열 $q_\rho$로 묶고) (1.)에 주어진 대입을 실행하기만 하면 된다. 다만 이때 행렬식의 값은 그 개별 원소들에 의해 직접 주어지는 것이 아니라 라플라스 전개를 거쳐서야 얻어진다. 이와 같이 행렬식과 행렬곱이 동치이므로, 형식의 기호적 표현 가능성에 관한 정리(첫째 기본정리)는 다음과 같이도 표현된다.

\emph{정리 I:} ``모든 불변 구성식은 행렬곱으로 기호적으로 표현할 수 있으며, 역으로 모든 행렬곱은 불변 구성식이다.''\footnote{미리 정한 기본형들의 불변 구성식은 물론 기호열 $S^\rho\ldots$에 의존하는 행렬곱의 결합식뿐이다.}

행렬곱에 의한 표현을 쓰면, 행렬식 표현에 따르는 본질적 난점, 곧 표현이 명시적이지 않다는 난점을 극복할 수 있다.\footnote{서론을 참조하라.} 여기서 ``명시적 표현''이란 열 $S^\sigma,T^\tau,\ldots$ 자체나 이들로부터 유일하게 유도할 수 있는 열 $R^\rho$만이 들어가고, 개별 부분열 $s,t,\ldots$는 더 이상 나타나지 않는 표현을 뜻한다. 우리는 \S\ 6에서 기호열 $S^\sigma,T^\tau,\ldots$에만 의존하는 모든 불변 구성식에 대하여 이 기호열들만으로 된 명시적 표현을 얻고, 그와 함께 생성 과정인 수축 과정과 미분 과정을 $n$항 영역으로 옮길 것이다. 명시적 표현이 가능하다는 사실은 다음에서 분명해진다. 가장 단순한 행렬곱들만이 행렬식 표현에서 나타나는 행렬식들에 그대로 대응할 수 있으므로, 그 밖의 모든 행렬곱은 서로 다른 여러 행렬식을 하나의 식으로 묶은 것이다. 실제로 기호열과 변수열을 개별 열 $s,u,x$로 풀어 놓으면, 첫째 기본정리의 통상적인 형태에 따라 반드시 행렬식의 결합식이 생긴다. 역으로, 미리 정한 기본형들의 불변 구성식을 나타내는 모든 행렬식 결합식을 명시적 행렬곱으로 묶을 수 있다는 증명은 둘째 기본정리를 사용하여 \S\ 6에서 얻는다.

%
%
이제 두 기호열 $S^\sigma,T^\tau$와 변수열을 이용하여 기본형 $(S^\sigma\mid p_\sigma)(T^\tau\mid p_\tau)$에서 모든 열을 명시적으로 포함하는 불변 구성식을 얻으려면(수축 과정), 이 절의 첫머리에 제시한 정의에 따라 제3종 행렬곱을 하나 만들기만 하면 된다.
\begin{equation}
 \pair{q_{n-\sigma-\lambda}S^\sigma}{T_{n-\tau}p_{\tau-\lambda}},
 \qquad (\lambda=1,2,\ldots,\tau,n-\sigma\footnotemark).
\tag{7}
\end{equation}
\footnotetext{수 $\lambda$는 언제나 마지막에 적은 두 수 가운데 작은 것까지만 변한다.}
(7.)을 전개하면 다음을 얻는다.
\begingroup
\setlength{\arraycolsep}{2pt}
\begin{equation}
\begin{aligned}
&\pair{v^{(1)}\cdots v^{(n-\sigma-\lambda)}s^{(1)}\cdots s^{(\sigma)}}
       {\alpha^{(1)}\cdots\alpha^{(n-\tau)}y^{(1)}\cdots y^{(\tau-\lambda)}} \\
&=\sum
\left|\begin{array}{ccc}
 v_{i_1}^{(1)}&\cdots&v_{i_{n-\lambda}}^{(1)}\\
 \vdots&&\vdots\\
 v_{i_1}^{(n-\sigma-\lambda)}&\cdots&v_{i_{n-\lambda}}^{(n-\sigma-\lambda)}\\
 s_{i_1}^{(1)}&\cdots&s_{i_{n-\lambda}}^{(1)}\\
 \vdots&&\vdots\\
 s_{i_1}^{(\sigma)}&\cdots&s_{i_{n-\lambda}}^{(\sigma)}
\end{array}\right|
\left|\begin{array}{ccc}
 \alpha_{i_1}^{(1)}&\cdots&\alpha_{i_{n-\lambda}}^{(1)}\\
 \vdots&&\vdots\\
 \alpha_{i_1}^{(n-\tau)}&\cdots&\alpha_{i_{n-\lambda}}^{(n-\tau)}\\
 y_{i_1}^{(1)}&\cdots&y_{i_{n-\lambda}}^{(1)}\\
 \vdots&&\vdots\\
 y_{i_1}^{(\tau-\lambda)}&\cdots&y_{i_{n-\lambda}}^{(\tau-\lambda)}
\end{array}\right|\\
&=\sum_{i,k,l}\eps\,q_{k_1k_2\ldots k_{n-\sigma-\lambda}}
 S^{k_{n-\sigma-\lambda+1}\ldots k_{n-\lambda}}
 T_{l_1\ldots l_{n-\tau}}
 p_{l_{n-\tau+1}\ldots l_{n-\lambda}},
\end{aligned}
\tag{8}
\end{equation}
\endgroup
즉, 이는 열 $S,T,q,p$에만 의존하는 식이다. 여기서 $\eps=\pm1$이고\footnote{이 표기는 이하 줄곧 유지한다.}, 합은 다음과 같이 이해한다. 각 열 $S\ldots$의 지수들은 그 열 안에서 올바른 순서로 배열하고, 각 조합 $i_1i_2\ldots i_{n-\lambda}$에 대하여 $k$와 $l$은 각각 그때 가능한 수 $i$들의 모든 순열을 따로 돈다. $\sigma\geqq\tau$인 경우 (7.)에 들어가는 수 $\lambda$를 수축의 결손이라 부른다. 이 마지막 제한의 이유는 \S\ 6에서 드러난다.

%
%
(7.)에 들어가는 두 행렬의 행렬식들은 새로운 열 $Q^{n-\lambda},P_{n-\lambda}$의 원소로 볼 수 있으며, 이들에 대응하는 열을 $Q_\lambda,P^\lambda$라 하자. 이 열 $Q,P$는 (8.)에 따라 원래의 $S$와 $q$, 또는 각각 $T$와 $p$로 나타낼 수 있다. 이를 다음 등식으로 표시한다.
\begin{equation}
\begin{aligned}
 q_{n-\sigma-\lambda}S^\sigma&=Q^{n-\lambda}\sim Q_\lambda, &\quad\text{또는}\quad Q_\lambda&=[q_{n-\sigma-\lambda}S^\sigma]_\lambda,\\
 T_{n-\tau}p_{\tau-\lambda}&=P_{n-\lambda}\sim P^\lambda, &\quad\text{또는}\quad P^\lambda&=[T_{n-\tau}p_{\tau-\lambda}]^\lambda.
\end{aligned}
\tag{9}
\end{equation}
(4.)를 고려하면, (6.)과 마찬가지로 행렬곱 (7.)에는 다음과 같은 서로 동치인 네 가지 기호적 표현이 성립한다.
\begin{equation}
\pair{q_{n-\sigma-\lambda}S^\sigma}{T_{n-\tau}p_{\tau-\lambda}}
=(q_{n-\sigma-\lambda}S^\sigma P^\lambda)=(Q_\lambda T_{n-\tau}p_{\tau-\lambda})=(-1)^{\lambda(n-\lambda)}\pair{P^\lambda}{Q_\lambda}.
\tag{10}
\end{equation}
(7.)에는 다음과 같이 놓음으로써 새로운 열을 한 번 더 대입할 수 있다.
\begin{equation}
 \pair{q_{n-\sigma-\lambda}S^\sigma}{T_{n-\tau}p_{\tau-\lambda}}
 =\pair{R^{\sigma+\lambda}}{p_{\sigma+\lambda}}(R^{\tau-\lambda}p_{\tau-\lambda}).
\tag{11}
\end{equation}
여기서도 (8.)은 열 $R$의 원소들의 곱을 열 $S$와 $T$로 유일하게 나타낼 수 있음을 보여 준다. 대입 (9.)와 (11.)은 특히 (7.)에 다른 열들을 사용한 수축을 한 번 더 시행할 때 적용하게 된다.

%
%
행렬 표현을 이용하면 열 $p_\rho$의 원소들 사이의 이차 관계(잘 알려져 있듯이 나머지 모든 관계가 여기서 따른다), 따라서 \S\ 1에 따르면 기본형의 계수에 관해 선형인, 열 $S^\rho$가 만족하는 관계를 불변적으로 세울 수 있다. 그것은 다음 항등식들이다.
\begin{equation}
 \pair{q_{n-\rho-\lambda}S^\rho}{S_{n-\rho}p_{\rho-\lambda}}=0\footnotemark,
 \qquad (\lambda=1,2,\ldots,\rho,n-\rho),
\tag{12}
\end{equation}
\footnotetext{Graßmann의 저작 제1권 제2부, S. 510에 있는 Study의 후기를 참조하라($S^\rho$라는 양이 단순하기 위한 조건).}
여기서 $p$와 $q$는 임의의 양으로 보아야 한다. 결손수 $\lambda$가 홀수인 항등식들은 더 높은 결손의 항등식들로부터 따른다는 것을 \S\ 5에서 볼 것이다. 관계 (12.)는 관계 (5.)의 직접적인 귀결이다. 실제로
\[
 \pair{q_{n-\rho-\lambda}S^\rho}{S_{n-\rho}p_{\rho-\lambda}}
 =(v^{(1)}\cdots v^{(n-\rho-\lambda)}s^{(1)}\cdots s^{(\rho)}\mid \sigma^{(1)}\cdots\sigma^{(n-\rho)}y^{(1)}\cdots y^{(\rho-\lambda)}),
\]
이고, 곱정리를 적용하면 (5.) 때문에 소멸하는 $(n-\lambda)$열 행렬식을 얻는다. 항등식 (12.)이 말하는 것은, 모든 수축 유형을 찾기 위해서는 각 변수열에 관해 선형인 형식들의 곱을 수축하는 것으로 충분하다는 것이다. 또는 \S\ 6에 따르면, 정규화된 형식 $(S^\rho\mid p_\rho)^m$은 계수에 관해 선형인 불변 구성식을 갖지 않는다는 것이다.

%
%
조건 (12.)는 열 $p_\rho$를 특징짓는 데에도 충분하다. 실제로 $p_\rho$의 개별 원소가 어떤 행렬의 행렬식들이라는 성질은 불변적이므로, 그 성질도 불변적으로 표현할 수 있어야 한다. 그런데 정리 VI에 따르면 조건 (12.)는 $p_\rho$에 가능한 가장 큰 불변적 제한을 나타낸다. 즉, 열 $p_\rho$를 계수로 삼아 만든 선형형식은 어떠한 불변 구성식도 갖지 않는다는 제한이다.

%
%
\section*{\S\ 3. 기호 항등식\footnote{문제의 정식화와 문헌에 대해서는 서론을 참조하라.}.}
이제 기호적 방법의 둘째 기본정리를 옮겨야 한다. 즉, 모든 열 $S^\rho,p_\rho$를 분해 불가능한 것으로 보았을 때의 기약인 분해 불가능 항등식 전체를 세워야 한다. 이를 위하여 다음을 추가로 정한다. 기호열과 변수열의 의미에 맞추어, 변수열 $p_\rho$를 $p_{\rho-1}y$로 특수화한 뒤 생기는 식에 이 계의 항등식 가운데 하나를 적용하여 얻는 항등식은 합성된 것으로 본다. 반면 $k$개의 기호열을 포함하는 항등식은, 위 과정으로 더 적은 기호열을 가진 항등식에서 생기더라도 분해 불가능한 것으로 본다. 열 $S^\rho,p_\rho$가 항등식에 반드시 명시적으로 들어갈 필요는 없다. 나타나는 행렬곱들이 이 열들에만 의존한다는 것을 보이면, 이들을 분해 불가능한 양으로 이해하기에 충분하다. 그러나 우리는 이 경우에도, 일부에는 대입 (11.)을 적용하여, 언제나 명시적 표현을 얻을 수 있음을 보일 것이다. 일반 항등식은 Pascal이 제시한 특수 항등식들로부터 행렬 표현의 원리를 써서 얻는다. 이 특수 항등식들은 열 $x$와 $u$만을 포함하며, Study가 삼항 영역에서 제시한 항등식을 직접 옮긴 것이다.
\begin{align}
\sum_{i=1}^n(-1)^{i+1}\pair{a^{(i)}}{x}(a^{(1)}\cdots a^{(i-1)}a^{(i+1)}\cdots a^{(n)}u)&=(-1)^{n+1}\pair{u}{x}(a^{(1)}\cdots a^{(n)}),\tag{13}\\
\sum\pm\pair{a^{(1)}}{x^{(1)}}\pair{a^{(2)}}{x^{(2)}}\cdots\pair{a^{(n)}}{x^{(n)}}&=(a^{(1)}\cdots a^{(n)})(x^{(1)}\cdots x^{(n)}),\tag{14}\\
\sum_{i=1}^n(-1)^{i+1}\pair{u}{\alpha^{(i)}}(\alpha^{(1)}\cdots\alpha^{(i-1)}\alpha^{(i+1)}\cdots \alpha^{(n)}x)&=(-1)^{n+1}\pair{u}{x}(\alpha^{(1)}\cdots\alpha^{(n)}).\tag{15}
\end{align}
(13.)과 (15.)에 각각 $(a^{(i)}\mid x)=(a^{(i)}q_{n-1})$와 $(u\mid \alpha^{(i)})=(p_{n-1}\alpha^{(i)})$를 대입하면 두 항등식을 더 얻지만, 우리의 관점에서는 본질적으로 다른 것으로 보지 않는다. (14.)는 행렬식의 곱정리를 나타낸다. 이제 (13.)과 (14.)를(쌍대적으로는 (15.)와 (14.)를) 적절히 만든 행렬의 곱정리 계로 어떻게 바꿀 수 있는지 보이겠다. 곱정리를 한 번은 $\rho$열 행렬에, 한 번은 $(\rho-1)$열 행렬에 적용하면 다음을 얻는다.

%
%
\[
\pair{a^{(1)}a^{(2)}\cdots a^{(\rho)}}{xp_{\rho-1}}
=(a^{(1)}a^{(2)}\cdots a^{(\rho)}\mid xy^{(2)}\cdots y^{(\rho)})
=\sum\pm\pair{a^{(1)}}{x}\pair{a^{(2)}}{y^{(2)}}\cdots\pair{a^{(\rho)}}{y^{(\rho)}},
\]
\[
\sum\pm\pair{a^{(2)}}{y^{(2)}}\cdots\pair{a^{(\rho)}}{y^{(\rho)}}
=(a^{(2)}\cdots a^{(\rho)}\mid y^{(2)}\cdots y^{(\rho)})
=(a^{(2)}\cdots a^{(\rho)}\mid p_{\rho-1})=(a^{(2)}\cdots a^{(\rho)}q_{n-\rho+1});
\]
따라서 $\rho=2,3,\ldots n$일 때 다음과 같은 곱정리 계가 따른다.
\begin{equation}
\begin{aligned}
\pair{a^{(1)}a^{(2)}\cdots a^{(\rho)}}{xp_{\rho-1}}
 &=\sum_{i=1}^{\rho}(-1)^{i+1}\pair{a^{(i)}}{x}(a^{(1)}\cdots a^{(i-1)}a^{(i+1)}\cdots a^{(\rho)}\mid p_{\rho-1})\\
 &=\sum_{i=1}^{\rho}(-1)^{i+1}\pair{a^{(i)}}{x}(a^{(1)}\cdots a^{(i-1)}a^{(i+1)}\cdots a^{(\rho)}q_{n-\rho+1}).\footnotemark
\end{aligned}
\tag{16}
\end{equation}
\footnotetext{항등식 (16.)은 $\rho$개 변수의 영역에서 가져온 항등식 (13.)으로도 볼 수 있다.}
$\rho=n$이면 (16.)은 (13.)이 된다. 더 나아가 $p_{n-1}=y^{(2)}p_{n-2}$, $p_{n-2}=y^{(3)}p_{n-3}$ 등으로 특수화하고, 식 $(a^{(1)}\cdots a^{(i-1)}a^{(i+1)}\cdots a^{(n)}\mid y^{(2)}p_{n-2})$, $(a^{(1)}\cdots a^{(i-1)}a^{(i+1)}\cdots a^{(k-1)}a^{(k+1)}\cdots a^{(n)}\mid y^{(3)}p_{n-3})$ 등에 항등식 (16.)을 적용하면 (13.)에서 (14.)에 도달한다. 따라서 우리의 정의에 따르면 항등식 (14.)는 항등식 (16.)들로부터 합성된 것이다. 달리 말해, 계 (16.)은 항등식 (13.) 및 (14.)와 동치이다. 계 (16.)은 일반 항등식에 요구한 조건 가운데 하나를 만족한다. 즉, 열 $p_{\rho-1}$을 분해 불가능한 양으로 포함하는 동시에 명시적인 형태로 포함한다. 이것은 이 요구만을 만족하는 유일한 계이다. 실제로 열 $a,x,p$로부터 다른 행렬식이나 행렬곱을 만들 수 없으며, 이들 사이에는 (16.)과 독립인 관계가 있을 수 없다. 그런 관계가 있다면 $(a^{(1)}\cdots a^{(\rho)}\mid xp_{\rho-1})$를 소거하여 $\rho$개($\rho\leqq n$)의 식 $(a^{(i)}\mid x)(a^{(1)}\cdots a^{(i-1)}a^{(i+1)}\cdots a^{(\rho)}q_{n-\rho+1})$ 사이의 관계를 얻을 것이다. 그러나 (13.)에 따르면 적어도 $n+1$개의 그러한 식이 하나의 관계에 들어가야 한다. 이와 유사한 고찰로부터 (15.) 및 (14.)와 동치인 다음 계를 얻는다.

%
%
\begin{equation}
\begin{aligned}
\pair{u q_{\rho-1}}{\alpha^{(1)}\cdots \alpha^{(\rho)}}
 &=\sum_{i=1}^{\rho}(-1)^{i+1}\pair{u}{\alpha^{(i)}}(q_{\rho-1}\mid \alpha^{(1)}\cdots \alpha^{(i-1)}\alpha^{(i+1)}\cdots \alpha^{(\rho)})\\
 &=\sum_{i=1}^{\rho}(-1)^{i+1}\pair{u}{\alpha^{(i)}}(p_{n-\rho+1}\alpha^{(1)}\cdots \alpha^{(i-1)}\alpha^{(i+1)}\cdots \alpha^{(\rho)}).
\end{aligned}
\tag{17}
\end{equation}

%
%
(16.)에서 임의의 열 $A^{\rho_1},B^{\rho_2},\ldots,x,p$ 사이의 항등식으로 넘어가기 위하여 다음을 생각한다. 이 항등식들은 $A^{\rho_1}=a^{(1)}a^{(2)}\cdots a^{(\rho_1)}$, $B^{\rho_2}=b^{(1)}b^{(2)}\cdots b^{(\rho_2)}$ 등으로 풀어 놓는 순간 (16.)과 같아져야 한다. 또한 개별 식들이 (16.)에 들어가는 유형이라면, (16.)으로 정해지는 연산들만을 통해서도 종결식으로 묶일 수 있어야 한다. 이제
\begin{equation}
 \rho=\rho_1+\rho_2+\cdots+\rho_k\leqq n
\tag{18}
\end{equation}
로 놓으면 다음을 얻는다.
\begin{align*}
&\pair{a^{(1)}\cdots a^{(\rho_1)}b^{(1)}\cdots b^{(\rho_2)}\cdots k^{(1)}\cdots k^{(\rho_k)}}{xp_{\rho-1}}
  =\pair{A^{\rho_1}B^{\rho_2}\cdots K^{\rho_k}}{xp_{\rho-1}} \\
&\quad=\sum_{i=1}^{\rho_1}(-1)^{i+1}\pair{a^{(i)}}{x}(a^{(1)}\cdots a^{(i-1)}a^{(i+1)}\cdots a^{(\rho_1)}B^{\rho_2}\cdots K^{\rho_k}q_{n-\rho+1})\\
&\qquad+(-1)^{\rho_1\rho_2}\sum_{i=1}^{\rho_2}(-1)^{i+1}\pair{b^{(i)}}{x}(b^{(1)}\cdots b^{(i-1)}b^{(i+1)}\cdots b^{(\rho_2)}A^{\rho_1}\cdots K^{\rho_k}q_{n-\rho+1})\\
&\qquad\vdots\\
&\qquad+(-1)^{(\rho_1+\cdots+\rho_{k-1})\rho_k}\sum_{i=1}^{\rho_k}(-1)^{i+1}\pair{k^{(i)}}{x}(k^{(1)}\cdots k^{(i-1)}k^{(i+1)}\cdots k^{(\rho_k)}A^{\rho_1}\cdots (K-1)^{\rho_{k-1}}q_{n-\rho+1}).
\end{align*}
이제 각각의 합 전체는(그러나 그 합의 일부만으로는 안 된다) 실제로 열 $A^{\rho_1},B^{\rho_2}\ldots$를 분해 불가능한 양으로 포함한다. 그에 필요한 동시에 충분한 다음 조건들을 만족하기 때문이다.\footnote{Capelli, a. a. O. \S\ 23은 충분조건 $\bigl(a^{(k)}\mid \frac{\partial}{\partial a^{(i)}}\bigr)=0$, $(k>i,\ i=1,2,\ldots,\rho_1-1)$을 제시한다. 여기서 필요충분조건은 Wellstein, Von den Differentialgleichungen der projektiven Invarianten, Math. Ann. 67 (1909) \S\ 3에 따라 푸아송 괄호 과정을 적용하면 가장 간단히 얻어진다.}
\begin{equation}
 \left(a^{(i+1)}\middle|\frac{\partial}{\partial a^{(i)}}\right)=0;\ \ldots\ ;
 \left(k^{(j+1)}\middle|\frac{\partial}{\partial k^{(j)}}\right)=0.
 \quad (i=1,2,\ldots,\rho_1-1;\ j=1,2,\ldots,\rho_k-1)
\tag{19}
\end{equation}

%
%
이 합들이 모든 열을 명시적으로도 포함한다는 것을 보이자. $B^{\rho_2}\cdots K^{\rho_k}q_{n-\rho+1}=q_{n-\rho_1+1}$로 놓으면, (16.)과 (10.)에 따라 다음이 나온다.
\begin{align*}
\sum(-1)^{i+1}\pair{a^{(i)}}{x}(a^{(1)}\cdots a^{(i-1)}a^{(i+1)}\cdots a^{(\rho_1)}q_{n-\rho_1+1})
&=\pair{A^{\rho_1}}{xp_{\rho_1-1}}\\
&=(A_{n-\rho_1}xp_{\rho_1-1})=(-1)^{(\rho_1+1)(n+1)}\pair{q_{n-\rho_1+1}}{A_{n-\rho_1}x}\\
&=(-1)^{(\rho_1+1)(n+1)}\pair{B^{\rho_2}\cdots K^{\rho_k}q_{n-\rho+1}}{A_{n-\rho_1}x}.
\end{align*}
나머지 합들도 같은 방식으로 계산하고, 열들이 이미 무게지수로 서로 다르게 특징지어져 있으므로 나중에 $B^{\rho_2}=A^{\rho_2}$, $C^{\rho_3}=A^{\rho_3},\ldots$로 놓으면, 원하는 항등식을 얻는다.
\begin{equation}
\begin{aligned}
\pair{A^{\rho_1}A^{\rho_2}\cdots A^{\rho_k}}{xp_{\rho-1}}
&=\sum_{i=1}^{k}(-1)^{\delta_i}\pair{A^{\rho_1}\cdots A^{\rho_{i-1}}A^{\rho_{i+1}}\cdots A^{\rho_k}q_{n-\rho+1}}{A_{n-\rho_i}x},\\
\delta_i&=(\rho_1+\rho_2+\cdots+\rho_{i-1})\rho_i+(\rho_i+1)(n+1),
\end{aligned}
\tag{20}
\end{equation}
그리고 (17.)로부터 그 쌍대 항등식들을 얻는다.

%
%
\begin{equation}
\begin{aligned}
\pair{u q_{\sigma-1}}{A_{\sigma_1}A_{\sigma_2}\cdots A_{\sigma_k}}
&=\sum_{i=1}^{k}(-1)^{\eps_i}\pair{u A^{n-\sigma_i}}{p_{n-\sigma+1}A_{\sigma_1}\cdots A_{\sigma_{i-1}}A_{\sigma_{i+1}}\cdots A_{\sigma_k}},\\
\eps_i&=(\sigma_1+\sigma_2+\cdots+\sigma_{i-1})\sigma_i+(\sigma_i+1)(n+\sigma).
\end{aligned}
\tag{21}
\end{equation}
\emph{정리 II.} (20.)과 (21.)로 열 $A^{\rho_i},x,p$ 사이의, 또는 각각 $A_{\sigma_i},u,q$ 사이의 분해 불가능 항등식 전체가 소진된다. 동시에, 대입 (11.)을 하지 않고도 명시적 표현을 허용하는 분해 불가능 항등식 전체도 이들로 소진된다.

%
%
정리의 첫째 부분은 (20.)과 (21.)을 유도할 때의 고찰에서 나오며, 둘째 부분은 다음과 같은 두 기호열 사이의 항등식이 존재할 수 없다는 데서 나온다.
\[
\pair{q_\rho A^\sigma}{B_{\rho+\sigma-\tau}p_\tau}
=\eps_1\pair{q_\rho B^{n-\rho-\sigma+\tau}}{A_{n-\sigma}p_\tau}
+\eps_2\pair{q_\rho q_{n-\tau}}{B_{\rho+\sigma-\tau}A_{n-\sigma}},
\]
여기서 $\rho\leqq\tau\leqq\sigma$이고 어느 열도 무게 $1\cdot(n-1)$을 갖지 않는다. ($\rho,\tau,\sigma$에 관한 다른 가정은 열의 표기를 서로 바꾸면 이 경우로 환원된다.) 예를 들어 이 불가능성은 (8.)을 전개하되 $q_\rho=lM^{\rho-1}$, $q_{n-\tau}=lN^{n-\tau-1}$로 특수화하여 얻는다. 이때 $l_1=1$, $M^{2,3,\ldots\rho}=1$, $A^{\rho+1,\ldots\rho+\sigma}=1$로 놓고 열 $l,M,A$의 나머지 모든 원소를 0으로 놓는 한편, $N$과 $B$는 임의로 둔다. 이제 임의의 열 사이의 항등식은 하나의 열 $p_\tau$를 $y^{(1)}y^{(2)}\cdots y^{(\tau)}$로 풀어 놓음으로써 (20.)으로부터 합성되는 항등식으로 바뀐다. 그러므로 (20.)을 기호열 두 개만을 가진 항등식들로부터 생성할 수 있음을 보이기만 하면 정리의 둘째 부분이 증명된다. 실제로
\begin{equation}
\pair{A^{\rho_1}A^{\rho_2}}{xp_{\rho-1}}
=\eps\pair{A^{\rho_1}q_{n-\rho+1}}{A_{n-\rho_2}x}
+\pair{A^{\rho_1}}{xp_{\rho_1-1}},
\quad \text{여기서 }p_{\rho_1-1}\sim A^{\rho_2}q_{n-\rho+1},
\tag{20a}
\end{equation}
에서 $A^{\rho_1}=B^{\sigma_1}B^{\sigma_2}$로 특수화한다. 즉,
\[
\pair{B^{\sigma_1}B^{\sigma_2}}{xp_{\rho_1-1}}
=\eps\pair{B^{\sigma_1}q_{n-\rho_1+1}}{B_{n-\sigma_2}x}
+\pair{B^{\sigma_1}}{xp_{\sigma_1-1}},
\quad \text{여기서 } p_{\sigma_1-1}\sim B^{\sigma_2}q_{n-\rho_1+1},
\]
로 특수화하면 세 기호열에 관한 하나의 항등식 (20.)을 얻는다. 그리고 $B^{\sigma_1}=C^{\tau_1}C^{\tau_2}$ 등으로 계속 특수화하면 일반 항등식 (20.)을 얻는다. 따라서 두 기호열과 두 임의의 변수열 사이의 항등식에는 명시적 표현이 존재하지 않는다는 사실로부터, 변수열들 가운데 열 $x$나 $u$가 하나도 없다면 기호열이 임의로 많이 있는 경우에도 명시적 표현이 존재하지 않는다는 결론이 따른다.

%
%
\section*{\S\ 4. 분해 항등식.}
이제 임의의 기호열과 변수열 사이의 항등식 가운데 가장 일반적인 경우를 생각한다. 이 항등식의 종결식에는 대입 (11.)을 하지 않은 채로 하나의 열 $p_\tau$가 그 개별 열 $y^{(1)}\cdots y^{(\tau)}$로 분해되어 나타난다. 우리는 이러한 항등식을 ``분해 항등식''이라 부른다. 앞 절 끝의 주석에 이어서, 먼저 기호열이 두 개뿐인 경우에 이 항등식을 구한다. 즉, 식 $\pair{q_\rho A^\sigma}{B_{\rho+\sigma-\tau}p_\tau}$을 전개한다. 이때 다음과 같이 놓는다.
\begin{equation}
\begin{aligned}
q_{\rho-\alpha}^{(i_1\ldots i_\alpha)}&\sim p_{n-\rho}\,y^{(i_1)}\cdots y^{(i_\alpha)},\qquad
p_{\tau-\alpha}^{(i_1\ldots i_\alpha)}=y^{(i_{\alpha+1})}\cdots y^{(i_\tau)},\qquad
p_\tau=y^{(1)}\cdots y^{(\tau)}.
\end{aligned}
\tag{22}
\end{equation}
수 $\rho,\sigma,\tau$에 관해서는 네 경우를 구별해야 한다.
\[
\begin{array}{llll}
1)&\rho+\sigma\leqq n,& \tau\leqq\rho,& \tau\leqq\sigma;\\[2pt]
2)&\rho+\sigma\leqq n,& \tau\geqq\rho,& \tau\leqq\sigma;\\[2pt]
3)&\rho+\sigma\leqq n,& \tau\leqq\rho,& \tau\geqq\sigma;\\[2pt]
4)&\rho+\sigma\leqq n,& \tau\geqq\rho,& \tau\geqq\sigma,
\end{array}
\]
이 가운데 첫째 경우를 택하여 다룬 뒤 나머지 경우를 짧게 특징짓겠다. 열 $y^{(1)}y^{(2)}\cdots y^{(\tau)}B_{\rho+\sigma-\tau}$를 $y^{(1)}$과 $y^{(2)}\cdots y^{(\tau)}B_{\rho+\sigma-\tau}$로 분해하면, (20.), (10.), (9.)에 따라 다음이 나온다.

%
%
\begin{equation}
\begin{aligned}
\pair{q_\rho A^\sigma}{y^{(1)}y^{(2)}\cdots y^{(\tau)}B_{\rho+\sigma-\tau}}
&=\pair{q_{\rho-1}^{(1)}A^\sigma}{y^{(2)}\cdots y^{(\tau)}B_{\rho+\sigma-\tau}}\\
&\quad+(-1)^\rho\pair{q_\rho[A_{n-\sigma}y^{(1)}]^{\sigma-1}}{y^{(2)}\cdots y^{(\tau)}B_{\rho+\sigma-\tau}}.
\end{aligned}
\tag{23}
\end{equation}
(23.)의 마지막 항은 원래 식과 정확히 같은 형태를 가진다. 다만 $\sigma$가 $\sigma-1$로, $\tau$가 $\tau-1$로 바뀌었을 뿐이다. 따라서 이 항에도 다시 등식 (23.)을 적용할 수 있다. 그러므로 $\tau\leqq\sigma$이므로 연산 (23.)을 $\tau$번 적용하고 (10.)을 고려하면 다음 관계에 이른다.
\begin{equation}
\begin{aligned}
\pair{q_\rho A^\sigma}{p_\tau B_{\rho+\sigma-\tau}}
&=(-1)^{\rho\sigma+(\sigma+\tau)(n+1)}
  \pair{q_\rho B^{n-\rho-\sigma+\tau}}{A_{n-\sigma}p_\tau}\\
&\quad+\sum_{i_1=1}^{\tau}(-1)^{\rho(i_1-1)}
 \pair{q_{\rho-1}^{(i_1)}[A_{n-\sigma}y^{(1)}y^{(2)}\cdots y^{(i_1-1)}]^{\sigma-i_1+1}}
 {y^{(i_1+1)}\cdots y^{(\tau)}B_{\rho+\sigma-\tau}}.
\end{aligned}
\tag{24}
\end{equation}

%
%
합 기호 아래의 각 항은 다시 원래 식의 유형을 가진다. $\rho$는 $\rho-1$로, $\sigma$는 $\sigma-i_1+1$로, $\tau$는 $\tau-i_1$로 바뀌었으므로 조건 1)은 오히려 더 강하게 충족된다. 따라서 $\tau\leqq\rho$이므로 연산 (24.)을 $\tau$번 적용할 수 있고, 원하는 분해 항등식을 얻는다.
\begin{equation}
\begin{aligned}
\pair{q_\rho A^\sigma}{p_\tau B_{\rho+\sigma-\tau}}
&=\eps_{i_0}\pair{q_\rho B^{n-\rho-\sigma+\tau}}{A_{n-\sigma}p_\tau}\\
&\quad+\sum_{i_1=1}^{\tau}\eps_{i_1}
 \pair{q_{\rho-1}^{(i_1)}B^{n-\rho-\sigma+\tau}}{A_{n-\sigma}p_{\tau-1}^{(i_1)}}\\
&\quad+\sum_{i_2=i_1+1}^{\tau}\eps_{i_2}
 \pair{q_{\rho-2}^{(i_1i_2)}B^{n-\rho-\sigma+\tau}}{A_{n-\sigma}p_{\tau-2}^{(i_1i_2)}}\\ &\quad\vdots\\
&\quad+\eps_{i_\tau}\pair{q_{\rho-\tau}^{(i_1i_2\ldots i_\tau)}B^{n-\rho-\sigma+\tau}}{A_{n-\sigma}}\\
&=\sum_{\alpha=0}^{\tau}\sum_{i_\alpha=i_{\alpha-1}+1}^{\tau}\eps_{i_\alpha}
 \pair{q_{\rho-\alpha}^{(i_1\ldots i_\alpha)}B^{n-\rho-\sigma+\tau}}{A_{n-\sigma}p_{\tau-\alpha}^{(i_1\ldots i_\alpha)}}.
\end{aligned}
\tag{25}
\end{equation}
부호 $\eps$에 대해서는 (24.)로부터 다음을 얻는다.
\begin{equation}
\begin{aligned}
\eps_{i_\alpha}&=(-1)^{\delta_{i_\alpha}},\qquad
\delta_{i_\alpha}=\sum_{\beta=1}^{\alpha}(\rho-\beta+1)(i_{\beta-1}+i_\beta+1)\\
&\quad +(\rho-\alpha)(\sigma+i_\alpha+\alpha)+(\sigma+\tau+\alpha)(n+1),
\end{aligned}
\tag{26}
\end{equation}
여기서 $i_0=0$으로 놓아야 한다. (25.)의 합 기호는 다음과 같이 이해한다. $i_1<i_2\cdots<i_{\alpha-1}$인 각 조합 $i_1i_2\ldots i_{\alpha-1}$에 대하여 $i_{\alpha-1}<i_\alpha\leqq\tau$를 만족하는 $i_\alpha$의 모든 값을 덧붙인다. (25.)의 개별 합들은 (26.)을 고려하면 쉽게 확인할 수 있듯이 (19.)와 유사한 다음 미분방정식을 만족한다.
\begin{equation}
 \left(\frac{\partial}{\partial y^{(i)}}\middle|y^{(i+1)}\right)=0,
 \qquad (i=1,2,\ldots\tau-1),
\tag{27}
\end{equation}
따라서 이 합들은 열 $p_\tau$에만 의존한다. 그러므로 이 분해 항등식은 열 $A,B,q,p$ 사이의 분해 불가능 항등식이다. 명시적 표현은 다음 절에서 다시 다룬다.

%
%
경우 2)에서도 연산 (23.)은 $\tau$번 적용할 수 있지만, 연산 (24.)는 $\rho$번째 뒤에 중단된다. 따라서 (25.)의 종결식에서는 첫째 합 기호를 $\sum_{\alpha=0}^{\rho}$로 바꾸어야 한다. 경우 3)과 4)에서는 연산 (23.)을 $\sigma$번만 실행할 수 있다. 이때 (24.) 대신 다음 관계가 들어간다.
\begin{equation}
\begin{aligned}
\pair{q_\rho A^\sigma}{p_\tau B_{\rho+\sigma-\tau}}
&=(-1)^{\rho\sigma}
 (q_{\rho-(\tau-\sigma)}^{(\sigma+1\ldots\tau)}B^{n-\rho+\tau-\sigma})(A_{n-\sigma}p_{\tau-(\tau-\sigma)}^{(\sigma+1\ldots\tau)})\\
&\quad+\sum_{i_1=1}^{\sigma}(-1)^{\rho(i_1-1)}
 \pair{q_{\rho-1}^{(i_1)}[A_{n-\sigma}y^{(1)}\cdots y^{(i_1-1)}]^{\sigma-i_1+1}}
 {y^{(i_1+1)}\cdots y^{(\tau)}B_{\rho+\sigma-\tau}};
\end{aligned}
\tag{28}
\end{equation}
(28.)을 $(\tau-\sigma)$번 반복하면, 그동안 합 기호는 $\sum_{i_\beta=i_{\beta-1}+1}^{\sigma+\beta-1}$로 바뀌며, (25.)에서 지수 $\alpha=\tau-\sigma$에 해당하는 개별 합의 모든 항을 각각 알맞은 부호와 함께 얻는다. 이는 (28.)의 합 기호 아래 항들에 (23.)을 $(\sigma-i_1+1)$번 적용하면 알 수 있다. 그 뒤에는 경우 3)과 4)가 경우 1)과 2)로 넘어간다. 실제로 $\sigma,\tau$에 대응하는 수들은 $\sigma'=\sigma-i_{\tau-\sigma}+\tau-\sigma$, $\tau'=\tau-i_{\tau-\sigma}=\sigma'$가 되며, 절차를 계속하면 $\tau'<\sigma'$가 된다. 그러므로 (25.)의 종결식에서는 첫째 합 기호를 $\sum_{\alpha=\tau-\sigma}^{\tau,\rho}$로 바꾸어야 한다.

%
%
(21.)로부터 이와 유사하게 쌍대적인 분해 항등식을 얻는다. 여기서는 하나의 열 $q_\rho$가 개별 열 $v^{(1)},\ldots,v^{(\rho)}$로 분해된다. 다음과 같이 놓자.
\begin{equation}
 \dot p_{\tau-\beta}^{(i_1\ldots i_\beta)}\sim v^{(i_1)}\cdots v^{(i_\beta)}q_{n-\tau};\qquad
 \dot q_{\rho-\beta}^{(i_1\ldots i_\beta)}=v^{(i_{\beta+1})}\cdots v^{(i_\rho)};\qquad
 q_\rho=v^{(1)}\cdots v^{(\rho)},
\tag{29}
\end{equation}
그러면 다음이 나온다.
\begin{equation}
 \pair{q_\rho A^\sigma}{p_\tau B_{\rho+\sigma-\tau}}
 =\sum_{\substack{0,\tau-\sigma}}^{\substack{\tau,\rho}}
   \sum_{i_\beta=i_{\beta-1}+1}^{\rho}\eps
\pair{\dot q_{\rho-\beta}^{(i_1\ldots i_\beta)}B^{n-\rho-\sigma+\tau}}{A_{n-\sigma}\dot p_{\tau-\beta}^{(i_1\ldots i_\beta)}} ,
\tag{30}
\end{equation}
여기서 합 지수 $\beta$는 아래의 두 값 가운데 큰 값부터 위의 두 값 가운데 작은 값까지 변한다. (25.)와 (30.)의 한 개별 합에 속하는 모든 항은 하나의 형식을 다음 편극 과정들로부터 얻은 편극들이다.
\[
\left(\frac{\partial}{\partial p_{\tau-\alpha}}\middle|p_{\tau-\alpha}^{(i_1\ldots i_\alpha)}\right),\qquad
\left(q_{\rho-\alpha}^{(i_1\ldots i_\alpha)}\middle|\frac{\partial}{\partial q_{\rho-\alpha}}\right).
\]
이를 표현하기 위하여 상수와 곱한 이러한 편극 연산들의 결합식을 $\Pi$로 나타내기로 한다. 그러면 (25.)와 (30.) 대신 다음 관계를 얻는다.
\begin{equation}
 \pair{q_\rho A^\sigma}{p_\tau B_{\rho+\sigma-\tau}}
 =\sum_{\substack{0,\tau-\sigma}}^{\substack{\tau,\rho}}
 \Pi\pair{q_{\rho-\alpha}B^{n-\rho-\sigma+\tau}}{A_{n-\sigma}p_{\tau-\alpha}}.
\tag{31}
\end{equation}

%
%
\section*{\S\ 5. 명시적 형태의 분해 항등식.}
경우 4)의 (25.)에서 $\tau$를 $\sigma+\rho$로 특수화하면 다음이 나온다.
\begin{equation}
 \pair{q_\rho A^\sigma}{y^{(1)}\cdots y^{(\rho+\sigma)}}
 =\sum_{i_\rho=i_{\rho-1}+1}^{\tau}
 \eps_{i_\rho}\pair{q_\rho}{y^{(i_1)}\cdots y^{(i_\rho)}}
 \pair{A^\sigma}{y^{(i_{\rho+1})}\cdots y^{(i_{\rho+\sigma})}},
\tag{32}
\end{equation}
여기서 $\eps_{i_\rho}=(-1)^{\delta_{i_\rho}}$, $\delta_{i_\rho}=\rho(\rho+1)/2+i_1+\cdots+i_\rho$가 된다. 이것은 (16.)과 (17.)보다 더 일반적인 형태의 행렬 곱정리이며, 곱행렬의 행렬식
\[
 \sum\pm(v^{(1)}y^{(1)})\cdots(v^{(\rho)}y^{(\rho)})(a^{(1)}y^{(\rho+1)})\cdots(a^{(\sigma)}y^{(\rho+\sigma)})
\]
을 라플라스 전개하여 얻는다. 따라서 더 일반적인 곱정리는 항등식 (20.), 또는 각각 (21.)로부터 합성되며, 이로써 특수한 형태 (16.), (17.)을 따로 골라낸 이유가 설명된다.

%
%
이제 관계 (32.)는 분해 항등식을 명시적으로 표현하는 수단을 준다. 이를 위하여 (31.)에 나타나는 형식들을 기본형으로 본다. 즉, 대입 (11.)을 실행한다. (8.)에 따르면 새로 도입한 열 $R$은 열 $A,B$ 자체로 표현할 수 있고 그 부분열 $a^{(1)}a^{(2)}\ldots$, $b^{(1)}b^{(2)}\ldots$로 표현할 필요가 없으므로, 이 대입은 열 $A,B$에 관한 명시적 표현을 바꾸지 않는다. 따라서
\begin{equation}
 \pair{q_{\rho-\alpha}B^{n-\rho-\sigma+\tau}}{A_{n-\sigma}p_{\tau-\alpha}}
 =\bigl(R_{\rho-\alpha}p_{n-\rho+\alpha}\bigr)
  \pair{R^{\tau-\alpha}}{p_{\tau-\alpha}},
\tag{33}
\end{equation}
라 하면, (25.)에서 지수 $\alpha$에 해당하는 개별 합에 대하여 다음이 나온다.
\begin{equation}
\begin{aligned}
&\sum \eps_{i_\alpha}\bigl(R_{\rho-\alpha}p_{n-\rho}y^{(i_1)}\cdots y^{(i_\alpha)}\bigr)
 \pair{R^{\tau-\alpha}}{y^{(i_{\alpha+1})}\cdots y^{(i_\tau)}}\\
&\quad=\sum \eps_{i_\alpha}\bigl([R_{\rho-\alpha}p_{n-\rho}]^\alpha y^{(i_1)}\cdots y^{(i_\alpha)}\bigr)
 \pair{R^{\tau-\alpha}}{y^{(i_{\alpha+1})}\cdots y^{(i_\tau)}}\\
&\quad=\eps\pair{[R_{\rho-\alpha}p_{n-\rho}]^\alpha R^{\tau-\alpha}}{p_\tau}
 =\eps\pair{R^{\tau-\alpha}q_{n-\tau}}{R_{\rho-\alpha}p_{n-\rho}}.
\end{aligned}
\tag{34}
\end{equation}
이로써 실제로 명시적인 종결식이 주어진다. $q_\rho$와 $p_\tau$가 들어가는 대칭적인 형태는 (30.)도 같은 종결식에 이른다는 것을 보여 준다. 따라서 이 종결식은 처음의 분해 방식에 의존하지 않는다. (25.)와 (30.) 사이의 차이는 열 $y$와 $v$가 독립적인 의미를 가질 때에만 존재한다. 이때는 우리의 정의에 따른 분해 항등식이 하나의 분해 가능한 항등식으로 바뀐다. 더 많은 기호열을 가진 항등식으로 넘어가려면 \S\ 3 끝에서처럼 한 열 $q_\rho$를 $q_{\rho_1}C^{\sigma_1}$로, 또는 각각 $p_\tau$를 $p_{\tau_1}C_{\tau-\tau_1}$로 풀어 놓기만 하면 된다. 그러면 생기는 식
\begin{equation}
 \pair{q_{\rho_1}C^{\sigma_1}}{[R^{\tau-\alpha}q_{n-\tau}]_\alpha R_{\rho_1+\sigma_1-\alpha}},
 \qquad
 \pair{[R_{\rho-\alpha}p_{n-\rho}]^\alpha R^{\tau-\alpha}}{p_{\tau_1}C_{\tau-\tau_1}},
\tag{35}
\end{equation}
은 두 기호열의 경우에 나타나는 유형과 정확히 같다. 따라서 (33.)에 해당하는 대입을 하면 다시 명시적 표현을 허용하며, 이 과정을 반복하면 임의로 많은 열에 대해서도 명시적 표현을 얻는다. 또한 (25.)와 각각 (30.)의 첫째 합과 마지막 합은 (33.)을 적용하지 않고 공식 (32.)만으로 명시적 표현을 주거나, 아예 모든 열을 명시적으로 포함하는 단 하나의 항으로 줄어든다. (32.)에서 $q_\rho$를 $q_{\rho_1}C^{\sigma_1}$ 등으로 풀어 얻는 관계는 가장 일반적인 형태의 곱정리로 볼 수 있다.

%
%
요약하면 우리는 다음을 보일 것이다.

\emph{정리 III:} 항등식
\begin{equation}
 \pair{q_\rho A^\sigma}{p_\tau B_{\rho+\sigma-\tau}}
 =\sum_{\substack{0,\tau-\sigma}}^{\substack{\tau,\rho}}
 \eps\pair{R^{\tau-\alpha}q_{n-\tau}}{R_{\rho-\alpha}p_{n-\rho}},
\tag{36}
\end{equation}
에서 열 $R$은 (33.)으로 정해진다. 이 항등식들과, 여기서 $q_\rho,p_\tau$를 $q_{\rho_1}C^{\sigma_1}$ 등으로 풀어 얻는 항등식들로 분해 불가능 항등식 전체가 소진된다.

실제로 이렇게 생기는 항등식들은 그 종류 가운데 가장 일반적이다. (20.), (21.)에서는 아직 존재하던 무게와 개수에 관한 제한을 이제는 개별 열들이 더 이상 받지 않기 때문이다. 또한 이 열들 사이에는 다른 항등식도 존재하지 않는다. 한 열 $p_\tau$를 $y^{(1)}\cdots y^{(\tau)}$로 풀면 가장 일반적인 항등식은 (20.)으로부터, 따라서 \S\ 3 끝에 따르면 (20a.)로부터 합성된다. 그리고 (16.) 뒤에서 논한 합성은 정확히 \S\ 4에서 실행한 합성이다. 임의로 많은 열로 넘어가는 것은, (20a.)의 논의가 보여 주듯이, $q_\rho$를 $q_{\rho_1}C^{\sigma_1}$ 등으로 풀어 얻은 식들에 언제나 같은 연산을 적용함으로써 이루어진다. 이로부터 (20a.) 대신 (20.)에서 직접 넘어가더라도 새로운 것은 생기지 않는다는 결론도 따른다. 이는 한 번은 (25.)에서 열 $q_\rho$를 특수화하고, 다른 한 번은 \S\ 4의 방법으로 (20.)에서 넘어가서 계산하면 쉽게 확인할 수 있다. 두 경우 모두 같은 합들이 생기며, 일반 곱정리(위 참조)를 적용하면 (36.)에서 얻는 항등식으로 바뀐다. 항등식 (20.), (21.)은 특수한 경우 $\tau=1$, $\rho=1$에 해당한다. 이때는 각각 두 개의 개별 합만 나타나므로, 위 주석에 맞게 대입 (33.) 없이 명시적 표현을 얻으며, 이는 앞서 얻은 결과와 일치한다.

%
%
끝으로 분해 항등식의 두 특수한 경우를 짧게 언급하자. (31.)에서 $\tau=\sigma$, $\rho=n-\sigma$로 놓고, 열 $p_\tau$를 $p'_\sigma\sim q'_{n-\sigma}$로 나타내면 다음이 나온다.
\begin{equation}
\begin{aligned}
\pair{A^\sigma}{p_\sigma}\pair{B^\sigma}{p'_\sigma}
-\pair{A^\sigma}{p'_\sigma}\pair{B^\sigma}{p_\sigma}
=\sum_{\alpha=1}^{\sigma,n-\sigma}
\Pi\pair{q_{n-\sigma-\alpha}B^\sigma}{A_{n-\sigma}p_{\sigma-\alpha}}.
\end{aligned}
\tag{37}
\end{equation}
이것은 서로 동변인 두 열 $p_\sigma,p'_\sigma$를 가진 한 형식을 기본 공변식들로 전개하는 Clebsch 공식이다.\footnote{Clebsch, a. a. O., S. 25--32. Clebsch는 오른쪽의 개별 항들이 열 $A,B$에만 의존한다는 것을 보였다. 그의 식 $\Phi_h$에 (32.)를 적용하면 명시적 표현을 얻을 수 있을 것이다.} 형식 $\pair{A^\sigma}{p_\sigma}^m\pair{B^\sigma}{p'_\sigma}^n$에 속하는 기본 공변식들은 다음과 같다.
\begin{equation}
 \left\{\sum_{\alpha=1}^{\sigma,n-\sigma}
 \pair{q_{n-\sigma-\alpha}B^\sigma}{A_{n-\sigma}p_{\sigma-\alpha}}
 \right\}^{\rho}
 \pair{A^\sigma}{p_\sigma}^{m-\rho}\pair{B^\sigma}{p'_\sigma}^{n-\rho},
 \qquad (\rho=0,1,\ldots m,n).
\tag{38}
\end{equation}
이를 간단히 말하면, 기본형으로부터 ``결손 $\alpha$의 동변 수축''(\S\ 2 참조)을 하여 얻는다.

%
%
둘째로 $\tau=\sigma-\lambda$, $\rho=n-\sigma-\lambda$, $B^\sigma=A^\sigma$로 놓으면 다음이 나온다.
\begin{equation}
\begin{aligned}
\pair{q_{n-\sigma-\lambda}A^\sigma}{A_{n-\sigma}p_{\sigma-\lambda}}
&=(-1)^\lambda\pair{q_{n-\sigma-\lambda}A^\sigma}{A_{n-\sigma}p_{\sigma-\lambda}}\\
&\quad+(-1)^{(n-\sigma)(\sigma-\lambda)}
 \sum_{\alpha=1}^{\sigma-\lambda,n-\sigma-\lambda}
 \Pi\pair{q_{n-\sigma-\lambda-\alpha}A^\sigma}{A_{n-\sigma}p_{\sigma-\lambda-\alpha}}.
\end{aligned}
\tag{39}
\end{equation}
따라서 \S\ 2에서 언급했듯이, 기호열을 특징짓는 조건 (12.) 가운데 결손 $\lambda$가 홀수인 조건은 더 높은 결손의 조건들로부터 따른다. 선형형식 $\pair{A^\sigma}{p_\sigma}\equiv\pair{B^\sigma}{p_\sigma}$에 대해서는 (39.)로부터, 계수에 관해 이차인 그 기약 불변 구성식들이 짝수 결손의 구성식들로 환원된다는 사실을 얻는다. 이는 이항 및 삼항 영역에서 하나의 선형형식이 불변 구성식을 갖지 않는다는 알려진 사실과 일치한다.

일반 항등식들은 처음에 개별 열들이 조건 (12.)를 만족한다는 가정 아래 유도되었다는 점을 덧붙인다. 그러나 이 개별 열들은 항등식에 선형으로만 들어가므로, 이 항등식들은 조건 (12.)를 만족하지 않는 더 일반적인 열들에 대해서도 성립한다.

%
%
\section*{\S\ 6. 불변 구성식의 명시적 표현.\\수축 과정과 미분 과정.}
앞 절들의 결과를 이용하면, \S\ 2에서 진술한 기호적 표현 가능성에 관한 정리(첫째 기본정리)에 명시적 표현 가능성을 덧붙여 완결할 수 있다. 즉, 다음을 보인다.

\emph{정리 IV:} ``임의의 기본형들의 모든 불변 구성식은 행렬곱으로 명시적으로 표현할 수 있다. 다시 말해, 종결식에는 변수열 $p_\rho$ 이외에 원래 형식들의 열 $S^\sigma,\ldots T^\tau$나, 이들로부터 대입 (9.) 또는 (11.)로 유도한 열 $P,Q,R$만이 들어간다.''

증명을 위하여 한 불변 구성식이 변수열뿐 아니라 열 $S^\sigma,\ldots T^\tau$에도 의존한다고 하자. 이 열들은 행렬식에 의한 표현이 가능하도록 개별 열들로 충분히 세분되어 있다고 하자. 즉, $S^\sigma=S^{\sigma_1}\cdots S^{\sigma_i},\ldots T^\tau=T^{\tau_1}\cdots T^{\tau_k}$이다. 열들은 그 자체로는 임의이지만 확정해 둔 순서 $S^\sigma,\ldots T^\tau$로 배열되어 있다고 하자. 이제 첫째 전체 열 $S^\sigma=S^{\sigma_1}\cdots S^{\sigma_i}$에, 그 부분열 몇 개에만이 아니라 전체에 의존하는 행렬식 결합식 하나를 택한다. 그런데 이 의존성은 \S\ 3--\S\ 5에 따르면 일반 항등식들의 귀결이다. 이제 한 변수열 $p_\rho$를 열 $y,v$로 풀거나, 뒤에 오는 열 $T$ 가운데 하나를 개별 열 $t$, 또는 각각 $\tau$로 풀면, 가장 일반적인 항등식들은 (20.), (21.)로부터 합성된다. 그런데 이 항등식들은 언제나 열 $S^\sigma=S^{\sigma_1}\cdots S^{\sigma_i}$를 명시적으로 포함하는 식, 즉 (20.), (21.)의 왼쪽으로 이끈다. 실제로 (19.)를 보면 오른쪽의 완전한 합 전체만이 --- (20a.)의 한 항에 해당하는 그 일부가 아니라 --- $S^\sigma$에만 의존한다는 성질을 가짐을 알 수 있다. 절차를 계속하는 동안 한 번 명시적으로 들어간 모든 열은 계속 명시적으로 남는다. 항등식을 적용하면서 요구될 수 있는 것은 여러 열을 하나로 묶는 것, 즉 대입 (9.)뿐이다. 마지막으로 단 하나의 열 $T^\tau$만 명시적 형태로 바꾸면 되는 단계에서는 두 경우가 가능하다. 자유로운 변수열, 곧 대입 (9.)로 다른 열과 결합되지 않은 변수열이 아직 남아 있을 수 있다. 이를 부분열 $y$ 또는 $v$로 분해하면 절차를 끝낼 수 있고, (31.)에서처럼 모든 열을 명시적으로 포함하는 하나 이상의 형식의 편극들이 생긴다. 자유로운 변수열이 더 이상 없다면, 개별 열 $t$ 또는 $\tau$를 포함하면서 전체 열 $T^\tau$에 의존하는 식들의 총체에서, 그 생성 방식에 따라 분해 항등식의 항들을 알아볼 수 있다. 그런데 이 항들은 \S\ 5에 따라 대입 (11.)을 적용하면 언제나 모든 열에 관한 명시적 표현을 허용한다. 이로써 위 정리는 일반적으로 증명되었다. 기호열이 두 개만 나타날 때에는 대입 (11.)이 전혀 일어나지 않는다는 점도 덧붙인다. 실제로 $\sigma,n-\sigma,\tau,n-\tau<n$이므로, 변수열이 전혀 들어가지 않을 수 있는 것은 두 기호열이 하나의 행렬식에 결합되어 이미 명시적으로 나타나는 경우뿐이다. 반대로 변수열이 들어가면 (34.)와 (33.)에 따라 대입 (11.)은 불필요해진다.

%
%
방금 말한 사실에 따르면 모든 불변 구성식을 생성하기 위해서는 기호열들에 변수열을 덧붙여 가능한 모든 행렬곱으로 결합하는 것으로 충분하다. 가장 일반적인 행렬곱은 다음 형태이다.
\begin{equation}
 \pair{P^{\mu_1}\cdots P^{\mu_i}}{Q_{\nu_1}\cdots Q_{\nu_k}},\qquad \sum\mu=\sum\nu,
\tag{40}
\end{equation}
여기서 $P$와 $Q$는 원래 형식들의 열일 수도 있고, 원래 열들에 대입 (9.) 또는 (11.)을 차례로 적용하여 생긴 열일 수도 있으며, 마지막으로 변수열일 수도 있다. 그런데 식 (40.)은 변수열을 이용하면서 매번 두 기호열을 하나의 행렬곱으로 차례로 결합하여 생성할 수 있다. 실제로
\[
 \pair{q_{n-\mu-\lambda}P^\mu}{Q_{\nu_1}p_{n-\nu_1-\lambda}}
 =\bigl(R_\lambda Q_{\nu_1}p_{n-\nu_1-\lambda}\bigr)
\]
에 열 $Q_{\nu_2}$를 결합하면 식
\[
 \pair{[R_\lambda Q_{\nu_1}]^{n-\lambda-\nu_1}}{Q_{\nu_2}p_{n-\nu_1-\nu_2-\lambda}}
 =\pair{q_{n-\mu-\lambda}P^\mu}{Q_{\nu_1}Q_{\nu_2}p_{n-\nu_1-\nu_2-\lambda}}
\]
을 얻는다. 따라서 절차를 계속하면 식 (40.)을 얻는다. 이때 $P$와 $Q$가 원래 형식들의 열이 아니라면, (9.)와 (11.)을 방금 보인 사실과 함께 적용하여 이 열들도 매번 두 기호열을 결합하는 절차를 차례로 거쳐 생겼음을 알 수 있다. 따라서 다음이 성립한다.

\emph{정리 V:} 모든 불변 구성식을 생성하는 과정인 수축 과정은, 변수열을 이용하면서 매번 두 기호열을 하나의 행렬곱으로 차례로 결합하는 데 있다.

%
%
$\sigma\geqq\tau$인 두 기호열 $S^\sigma,T^\tau$로부터 다음 행렬곱들을 만들 수 있다.
\begin{align}
&\pair{q_{n-\sigma-\lambda}S^\sigma}{T_{n-\tau}p_{\tau-\lambda}}; &&(\lambda=1,2,\ldots\tau,n-\sigma)\tag{41}\\
&\pair{q_{n-\tau-\mu}T^\tau}{S_{n-\sigma}p_{\sigma-\mu}}; &&(\mu=\sigma-\tau+\lambda)\tag{42}\\
&\pair{q_{n-\tau-\nu}T^\tau}{S_{n-\sigma}p_{\sigma-\nu}}. &&(\nu=1,2,\ldots\sigma-\tau)
\tag{43}
\end{align}
그러나 (31.)에서 (42.)와 (43.)에 대하여 다음 값들을 얻는다.
\begin{align}
&\pair{q_{n-\sigma-\lambda}S^\sigma}{T_{n-\tau}p_{\tau-\lambda}}
 +\sum_{\alpha}\Pi\pair{q_{n-\sigma-\lambda-\alpha}S^\sigma}{T_{n-\tau}p_{\tau-\lambda-\alpha}};\tag{42a}\\
&\Pi(q_{n-\sigma}S^\sigma)(T_{n-\tau}p_\tau)
 +\sum_{\beta}\Pi\pair{q_{n-\sigma-\beta}S^\sigma}{T_{n-\tau}p_{\tau-\beta}}.
\tag{43a}
\end{align}
따라서 형식 (42.)는 (41.) 및 (41.)의 편극들의 선형식으로 나타낼 수 있고, 형식 (43.)는 기본형 및 형식 (41.)의 편극들로 나타낼 수 있다. 이와 마찬가지로 (31.)에서 형식 (41.)이 (42.)와 (42.)의 편극들에 선형으로 의존함을 얻는다. 그러므로 Clebsch의 정의에 따르면\footnote{Clebsch, a. a. O. \S\ 7.} 형식 (42.)는 형식 (41.)와 동치인 반면, (43.)는 기본형으로 되돌아간다. 따라서 생성 과정은 (41.) 또는 (42.)로 동등하게 주어진다. 우리는 수 $\lambda$에 관해 더 단순한 과정 (41.)을 수축 과정으로 정의하겠다. 수축의 결손인 수 $\lambda$(\S\ 2 참조)는 행렬식의 곱에 도달하기까지 부족한 열의 개수를 나타낸다. 정확히는 주어진 수축에서 최대 개수의 열을 포함하는 그 행렬곱에서 부족한 열의 수이다. $\sigma\geqq\tau$일 때 $\lambda$는 두 값 $\tau,n-\sigma$ 가운데 작은 값까지만 변하므로 다음이 성립한다.
\begin{equation}
 \lambda\leqq\frac n2,\ n\text{ 짝수};\qquad
 \lambda\leqq\frac{n-1}{2},\ n\text{ 홀수}.
\tag{44}
\end{equation}

%
%
추가 수축으로 열 $p,q$가 기호열로 바뀐 경우에도 (42.)와 (43.)을 (41.)로 환원하는 것은 여전히 성립한다. 실제로 (36.)에 따라
\[
 \pair{A^{n-\tau-\varkappa}T^\tau}{S_{n-\sigma}B_{\sigma-\varkappa}}
 =\sum_{\alpha=0,\varkappa-\sigma+\tau}^{\tau,n-\sigma}
 \eps\pair{R^{\tau-\alpha}B^{n-\sigma+\varkappa}}{A_{\tau+\varkappa}R_{n-\sigma-\alpha}},
\]
이고, 여기서 열 $R$은 (33.)에 따라 $S$와 $T$에서 생긴다. 따라서 기호열만을 포함하는 이 형식들은, $\sigma-\tau\gtreqless 2\varkappa$인지에 따라 각각 $\pair{q_\tau A^{n-\tau-\varkappa}}{B_{\sigma-\varkappa}p_{n-\sigma}}$ 또는 $\pair{q_{\tau-\alpha}B^{n-\sigma+\varkappa}}{A_{\tau+\varkappa}p_{n-\sigma-\alpha}}$를 기본형 및 형식 (41.)과 수축하여 얻는다.

%
%
기호적 수축 과정에서 불변 미분 과정을 얻는 방법은 잘 알려져 있다. 기호열 $S^\sigma,T_{n-\tau}$를 미분기호의 열 $\frac{\partial}{\partial p_\sigma}$, $\frac{\partial}{\partial q_{n-\tau}}$로 바꾸기만 하면 된다. 이때 각 곱 $\frac{\partial}{\partial p_{i_1\ldots i_\sigma}}\cdot\frac{\partial}{\partial q_{j_1\ldots j_{n-\tau}}}$에는 잘 알려져 있듯이 실제 의미 $\frac{\partial^2}{\partial p_{i_1\ldots i_\sigma}\partial q_{j_1\ldots j_{n-\tau}}}$가 대응한다. 열 $\frac{\partial}{\partial p_\sigma}$, $\frac{\partial}{\partial q_{n-\tau}}$는 열 $S^\sigma,T_{n-\tau}$에 동변하므로 이들과 같은 항등식들을 만족하며, 특히 (42.)와 (42a.), (43.)와 (43a.) 사이에 성립하는 항등식들도 만족한다. 요약하면 다음을 얻는다.

\emph{정리 VI:} 임의의 기본형들의 모든 불변 구성식은 이 기본형들에 기호적 수축 과정 (41.)을 반복 적용하거나, 다음 불변 미분 과정을 반복 적용하여 생긴다.
\begin{equation}
 \left(q_{n-\sigma-\lambda}\frac{\partial}{\partial p_\sigma}\middle|\frac{\partial}{\partial q_{n-\tau}}p_{\tau-\lambda}\right),
 \qquad(\sigma\geqq\tau,\ \lambda=1,2,\ldots,\tau,n-\sigma).
\tag{45}
\end{equation}
각 수축 (41.)에서 형식의 전체 무게, 즉 개별 변수열들의 무게의 합(\S\ 1 참조)은 양수 $2(\lambda^2+\lambda(\sigma-\tau))$만큼 감소한다는 점도 언급해야 한다. 따라서 일반적인 유한성 증명과 독립적으로 형식열, 즉 주어진 기본형의 계수에 관해 선형인 불변 구성식 전체의 유한성이 따른다.\footnote{Clebsch, a. a. O. \S\ 11 및 12를 참조하라. 기본 공변식 계의 유한성.}

%
%
더 나아가 정리 VI는 조건 (12.)를 만족하지 않는 형식들에 대해서도 성립한다. 실제로 (3.)의 각 인자 $\pair{S^\rho}{p_\rho}^{m_\rho}$를 $m_\rho$개의 선형 인자들의 곱 $\pair{A^\rho}{p_\rho}\pair{B^\rho}{p_\rho}\cdots\pair{M^\rho}{p_\rho}$로 바꿀 수 있다. 그러면 불변 구성식들은 선형형식 계의 불변 구성식들로 바뀌며, 이 선형형식들은 나중에 서로 같다고 놓기만 하면 된다. 그런데 각 개별 선형형식에 대해서는 조건 (12.)가 언제나 충족된다. 미분기호를 도입했을 때 이 조건들이 2계 미분계수만을 포함하기 때문이다.

%
%
\section*{\S\ 7. 일반 형식의 정규형에 의한 전개.}
이하(\S\ 8)에서 우리는 수축 과정 (41.)의 한 가지 주된 성질, 곧 그것이 기본 수축들로부터 합성될 수 있다는 사실을 유도할 것이다. 그러나 이를 위해서는 Mertens가\footnote{F. Mertens, \emph{\"Uber invariante Gebilde quatern\"arer Formen} (서론 참조), 이하 “Mertens”로 인용한다.} 4항 영역에서 제시한 중요한 정리를 $n$항 영역으로 옮길 필요가 있다. 이 정리에 따르면 임의의 유한한 형식계는 그와 동치인 정규형의 계로 바꿀 수 있다. 여기서 우리는 --- 이항 및 삼항 정의를 일반화하여\footnote{Study, a. a. O. S. 54.} --- “정규형”을 기본형 자체를 제외하고 계수에 관해 선형인 모든 불변 구성식이 항등적으로 소멸하는 형식으로 이해한다. 따라서 정리 VI(\S\ 6)에 의해 정규형 $\varphi$는 다음 미분방정식계를 만족한다.
\begin{equation}
 \left(q_{n-\sigma-\lambda}\frac{\partial}{\partial p_\sigma}\middle|\frac{\partial}{\partial q_{n-\tau}}p_{\tau-\lambda}\right)\varphi=0,
 \qquad(\sigma\geqq\tau;\ \lambda=1,2,\ldots\tau,n-\sigma),
\tag{46}
\end{equation}
여기서 열 $q_{n-\sigma-\lambda},p_{\tau-\lambda}$는 임의의 양으로 보아야 한다. 이 미분방정식들은 4항 영역에서 --- $\sigma=\tau=2$일 때 (39.)를 고려하면 --- Mertens가 임의의 열 $p,q$를 도입하지 않고 제시한 10개의 미분방정식(a. a. O., 공식 28))과 완전히 일치한다. 이 방정식들과 정리 VI를 알면 그곳의 증명법을 거의 그대로 $n$개 변수의 경우로 옮길 수 있다. 추가해야 할 유일한 것은 분해 항등식 (30.)으로 구성된 형식들이 정규형임을 보이는 증명이다. 4항 영역에서는 이 증명이 더 이상의 변형 없이도 얻어진다. 그러므로 증명을 간략히 서술하는 것으로도 충분하다. 또한 단순화를 위하여 이하에서 유일하게 고려해야 하는 하나의 기본형 $f$와 그 형식열의 경우(\S\ 6 끝 참조)만을 다룬다. 임의로 많은 기본형들의 계수에 관해 임의 차수인 불변 구성식에 대해서도 증명은 같기 때문이다.

증명의 기본 착상은 변환된 계수를 도입하여, $n-1$개의 열 $p_\rho$를 가진 한 형식을 $x$에 동변인 $n$개의 열 $\xi$, 곧 변환량의 열들을 가진 형식들로 바꾸는 것이다. 이 형식들로부터 불변 미분 과정에 의해 찾는 정규형들이 직접 유도된다. 일반 형식을 정규형으로 전개하는 것은 $x$에 동변인 $\rho$개의 열 $\xi$를 가진 형식들($\rho\leqq n$)에 대한 급수 전개를 통하여 이루어진다. 불변 미분 과정은 원래의 열 $p_\rho$를 가진 형식들로 되돌아가게 한다.

%
%
$\xi^{(1)},\ldots\xi^{(n)}$을 변환량의 열이라 하자. 그러면 다음이 성립한다.
\begin{equation}
\begin{aligned}
x_i&=\xi_i^{(1)}x'_1+\xi_i^{(2)}x'_2+\cdots+\xi_i^{(n)}x'_n,\\
p_{i_1\ldots i_\rho}&=\sum_j(\xi_{i_1}^{(j_1)}\cdots \xi_{i_\rho}^{(j_\rho)})p'_{j_1\ldots j_\rho}.
\end{aligned}
\tag{47}
\end{equation}
형식 $f,\varphi,\ldots$의 변환된 계수들을 $(f),(\varphi),\ldots$로 나타내고, 다음과 같이 놓자.
\[
 f\equiv\pair{S^1}{p_1}^{m_1}\pair{S^2}{p_2}^{m_2}\cdots\pair{S^{n-1}}{p_{n-1}}^{m_{n-1}}.
\]
그러면 다음을 얻는다.
\begin{equation}
\begin{aligned}
(f)&=\pair{S^1}{\xi^{(1)}}^{m_{11}}\cdots\pair{S^1}{\xi^{(n)}}^{m_{1n}}
\pair{S^2}{\xi^{(1)}\xi^{(2)}}^{m_{21}}\cdots
\pair{S^{n-1}}{\xi^{(2)}\cdots\xi^{(n)}}^{m_{n-1,n}}\\
&=\pair{s^{(1)}}{\xi^{(1)}}^{k_1}\pair{s^{(2)}}{\xi^{(2)}}^{k_2}\cdots\pair{s^{(n)}}{\xi^{(n)}}^{k_n},
\end{aligned}
\tag{48}
\end{equation}
여기서 $m_{11}+\cdots+m_{1n}=m_1$ 등이 성립한다. 다음 조건을 만족하는 각 계수 $(f)$로부터
\begin{equation}
 k_1\geqq k_2\geqq k_3\cdots\geqq k_n,
\tag{49}
\end{equation}
열 $\xi$를 정확히 모두 소진하는 다음 미분 과정을 거치면
\begin{equation}
\left(\frac{\partial}{\partial\xi^{(1)}}\middle|p_1\right)^{k_1-k_2}
\left(\frac{\partial}{\partial\xi^{(1)}}\frac{\partial}{\partial\xi^{(2)}}\middle|p_2\right)^{k_2-k_3}
\cdots
\left(\frac{\partial}{\partial\xi^{(1)}}\cdots\frac{\partial}{\partial\xi^{(n-1)}}\middle|p_{n-1}\right)^{k_{n-1}-k_n}
\left(\frac{\partial}{\partial\xi^{(1)}}\cdots\frac{\partial}{\partial\xi^{(n)}}\right)^{k_n}
\tag{50}
\end{equation}
다음 형식을 얻는다.
\begin{equation}
 \varphi=\bigl(s^{(1)}\mid p_1\bigr)^{k_1-k_2}
 \bigl(s^{(1)}s^{(2)}\mid p_2\bigr)^{k_2-k_3}\cdots
 \bigl(s^{(1)}\cdots s^{(n-1)}\mid p_{n-1}\bigr)^{k_{n-1}-k_n}
 \bigl(s^{(1)}s^{(2)}\cdots s^{(n)}\bigr)^{k_n}.
\tag{51}
\end{equation}
형식 $\varphi$들이 바로 찾던 정규형들이다. 실제로 이들은 정규형을 정의하는 다음 성질들을 가진다.

%
%
1. 이들은 기본형 $f$의 계수에 관해 선형인 불변 구성식이다. 이는 이들이 $f$에 불변 미분 과정 $\left(\frac{\partial}{\partial p_\rho}\middle|\xi^{(i_1)}\xi^{(i_2)}\cdots\xi^{(i_\rho)}\right)$ 및 (50.)을 적용하여 생기기 때문이다. 따라서 이들은(\S\ 6 끝에 따라) $f$의 유한 형식열의 편극들의 선형식으로 나타낼 수 있다. 이러한 편극들은 $\varphi$에 들어가는 변수 $p_\rho$를, 변수열 $\mathfrak{p}_\sigma$를 가진 다음 형식, 곧 선형형식들의 곱으로 분해되는 형식의 계수로 간주하여 얻는다.
\begin{equation}
 H=(q_{n-1}\mid \mathfrak{p}_{n-1})^{k_1-k_2}(q_{n-2}\mid \mathfrak{p}_{n-2})^{k_2-k_3}\cdots(q_1\mid \mathfrak{p}_1)^{k_{n-1}-k_n}.
\tag{52}
\end{equation}
형식열 $f$와 $H$의 연립불변식들은 고려해야 하는 편극들을 모두 준다. 이들 편극은 각 개별 열 $p_\rho$에 관하여 $\varphi$ 자체와 같은 차수를 가져야 하기 때문이다.

2. 이들은 미분방정식 (46.)을 만족한다. 실제로 $\varphi$에 미분 과정 (45.)을 적용하면 다음이 나온다.
\[
 \varphi'=\bigl(q_{n-\sigma-\lambda}s^{(1)}\cdots s^{(\sigma)}\mid [s^{(1)}\cdots s^{(\tau)}]_{n-\tau}p_{\tau-\lambda}\bigr).\qquad (\sigma\geqq\tau).
\]
분해 항등식 (30.)에서 열 $s$로 만들어진 열 $q_\sigma=s^{(1)}\cdots s^{(\sigma)}$, $p_{n-\tau}=[s^{(1)}\cdots s^{(\tau)}]_{n-\tau}$를 그곳의 $q_\rho,p_\tau$와 동일시하면 다음을 얻는다.
\[
 \varphi'=\sum_{\beta=\sigma-\tau+\lambda}^{n-\tau,\sigma}\sum_{i_\beta}\eps\,
 \pair{\dot{q}_{\sigma-\beta}^{(i_1\ldots i_\beta)}q_{n-\tau+\lambda}}{p_{\sigma+\lambda}\dot{p}_{n-\tau-\beta}^{(i_1\ldots i_\beta)}}.
\]
이때 다음이 성립한다.
\[
 \dot{p}_{n-\tau-\beta}^{(i_1\ldots i_\beta)}\sim s^{(1)}\cdots s^{(\tau)}s^{(i_1)}\cdots s^{(i_\beta)},
\]
여기서 $i_1,\ldots i_\beta$는 1부터 $\sigma$까지의 수 가운데서 고른, 서로 다른 임의의 $\beta$개 수를 뜻한다. 이제 $\beta>\sigma-\tau$이므로 이 수들 $i_1,\ldots i_\beta$ 중에는 반드시 1부터 $\tau$까지의 수가 들어 있다. 그러면 $\dot{p}_{n-\tau-\beta}^{(i_1\ldots i_\beta)}$는 항등적으로 소멸하고, 따라서 각 개별 행렬곱 및 $\varphi'$도 소멸한다.

%
%
정규형의 계가 형식열 $f$와 동치임을 보이기 위해 이제 남은 것은, 형식열의 모든 형식을 1.에서 정한 의미에서 정규형들의 편극들에 따라 전개할 수 있음을 보이는 것뿐이다. $\Pi$를 다음 편극 연산들의 결합식이라 하자.
\srcnumdisplay[4.2em]{(53.)}{%
 \left(\frac{\partial}{\partial\xi^{(i)}}\middle|\xi^{(k)}\right),
 \qquad (i\ne k;\ i=1,2,\ldots,\rho;\ k=1,2,\ldots,\rho)}
또한 $\Pi^\sigma$를 $i=\sigma$이고 $k<\sigma$인 특수 연산 (53.)들의 결합식이라 하자. 열 $\xi^{(\rho)}$에 관해 차수가 $k_\rho$인, $\rho$개의 열 $\xi^{(1)},\ldots,\xi^{(\rho)}$의 형식 $F$에 대해서는 다음 전개가 성립한다\srcfn{*)}{F. Mertens, \emph{\"Uber eine Formel der Determinantentheorie}. Sitzungsber. d. Ak. d. Wiss., Wien, math.-naturw. Kl., Bd. 91. 2. Abt. (1885). Formel 20).}:
\srcnumdisplay[4.2em]{(54.)}{%
 F=\sum_{\alpha=0}^{k_\rho}\Pi F_\alpha,
 \qquad(\rho\leqq n)}
여기서 $F_\alpha$는 마지막 열 $\xi^{(\rho)}$에 관해 $\alpha$차이고, $F$에 다음 불변 미분 연산
\srcnumdisplay[4.2em]{(55.)}{%
 \left(\frac{\partial}{\partial\xi^{(1)}}\cdots\frac{\partial}{\partial\xi^{(\rho)}}\middle|\xi^{(1)}\cdots\xi^{(\rho)}\right)^\alpha}
및 $\Pi^\rho$를 적용하여 생긴다. $F_\alpha$를 생성할 때 과정 (55.)를
\srcnumdisplay[4.2em]{(56.)}{%
 \left(\frac{\partial}{\partial\xi^{(1)}}\cdots\frac{\partial}{\partial\xi^{(\rho)}}\middle|p_\rho\right)^\alpha,}
로 바꾸면, $\rho-1$개의 열 $\xi^{(1)},\ldots,\xi^{(\rho-1)}$만을 가진 형식 $F_\alpha(p_\rho)$가 생기며, 여기에 다시 (54.)를 적용할 수 있다. 이 절차를 계속하면 형식 $G(p_\rho,p_{\rho-1},\ldots,p_1)$에 이른다. 이 형식들로부터 (54.)에 따른 $F$의 전개를 얻으려면, 각 열 $p_\sigma$를 $\xi^{(1)},\ldots,\xi^{(\sigma)}$로 바꾸고 생기는 형식들에 편극 과정 $\Pi$ (53.)를 적용해야 한다. 그러면 (48.)과 비교할 때 변환된 계수 $(G)$들의 합을 얻는다. $\rho=n$이면 첫 단계에서 과정 (55.)를 그대로 둔다. $c$가 수치 상수들을 뜻한다고 하고, 더 나아가 다음과 같이 줄여 쓰자.
\srcnumdisplay[4.2em]{(57.)}{%
 (\xi^{(1)}\xi^{(2)}\cdots\xi^{(n)})=\Delta;
 \qquad
 \left(\frac{\partial}{\partial\xi^{(1)}}\frac{\partial}{\partial\xi^{(2)}}\cdots\frac{\partial}{\partial\xi^{(n)}}\right)=\nabla,}
그러면 $\rho=n$인 경우에는
\srcnumdisplay[4.2em]{(58.)}{%
 F=\sum c\Delta^\alpha(G),}
을 얻는 반면, $\rho<n$이면 오른쪽에는 $\alpha=0$인 항, 곧 $\sum c(G)$만 나타난다.

%
%
(58.)을 변환된 계수 $(f)$ (48.)에 적용하자. 편극 과정 $\Pi^\sigma$는 $\rho\geqq\sigma$인 모든 연산 (56.)과 교환할 수 있고, 변환된 계수의 편극은 다시 변환된 계수가 되므로 다음을 얻는다.
\[
\begin{aligned}
 G={}&\left(\frac{\partial}{\partial\xi^{(1)}}\middle|p_1\right)^{\beta_1}
 \left(\frac{\partial}{\partial\xi^{(1)}}\frac{\partial}{\partial\xi^{(2)}}\middle|p_2\right)^{\beta_2}
 \cdots
 \left(\frac{\partial}{\partial\xi^{(1)}}\cdots\frac{\partial}{\partial\xi^{(n-1)}}\middle|p_{n-1}\right)^{\beta_{n-1}}
 \nabla^{\beta_n}\Pi(f)\\
 ={}&\sum c\varphi
\end{aligned}
\]
따라서 (58.)로부터 다음이 따른다.
\srcnumdisplay[4.2em]{(59.)}{%
 (f)=\sum c\Delta^\alpha(\varphi)\qquad\text{(Mertens, Formel 23)).}}
(59.)에서 $f$ 자체의 전개로 넘어가기 위해, 다시 변수 $p_\rho$를 $f$에 대응하는 차수 $m_1,\ldots m_{n-1}$을 가진 형식 $H$ (52.)의 계수로 간주하고, $m=m_1+\cdots+m_{n-1}$이라 하자. 그러면 불변식의 정의에 따라 다음이 성립한다.
\[
 f\cdot\Delta^m=\sum(f)(H)=\sum c\Delta^\alpha(\varphi)(H),
\]
열 $\xi$를 정확히 모두 소진하는 과정 $\nabla^m$을 적용하면 다음을 얻는다.
\srcnumdisplay[4.2em]{(60.)}{%
 f=\sum c\Phi,}
여기서 형식 $\Phi$는 $\varphi$와 $H$로부터 변수들에 관한 불변 미분 과정을 통해 유도되므로, 형식열 $\varphi$와 $H$의 연립불변식이다\srcfn{*)}{직접 이 결론을 내리는 대신 Mertens, \S\ 7에서는 본질적으로 형식열 $H$를 만드는 데 목적을 둔 일련의 추가 미분 과정들을 도입하는데, 이는 우리의 정리 VI가 수행하는 일이다.}. 그런데 $\varphi$가 정규형이므로 형식열 $\varphi$는 형식 $\varphi$ 자체로 제한되고, 형식열 $H$는 가능한 모든 불변 변수 결합을 포함한다. 따라서 $\varphi$와 $H$의 연립불변식들은 1.에서 말한 의미의 $\varphi$의 편극들이다. 즉, (60.)은 $f$를 정규형들의 편극들에 따라 전개한 것이다. 이제 형식열 $f$의 임의의 형식 $\theta$에 대해서도 같은 전개를 증명해야 한다. $\Gamma$를 $(\theta)$로부터 (50.)에 의해 유도된 정규형들이라 하자. 그러면 다음이 성립한다.

%
%
\srcnumdisplay[4.2em]{(61.)}{%
 (\theta)=\sum c\Delta^\beta(\Gamma).}
형식 $\theta$는 각 변환된 계수에 대해 성립하는 다음 관계로 특징지어진다.
\[
 (\theta)\Delta^\sigma=\sum c(f)\qquad\text{(Mertens, Formel 9)).}
\]
그런데 $\Delta^\sigma$를 인자로 포함하는 $n$개의 열 $\xi$의 임의의 형식 $F$는 둘째 것인 $\nabla^\sigma\Pi F$도 포함하므로(Mertens, Formel 20)), 다음이 따른다.
\[
 (\theta)=\sum c\nabla^\sigma(f),\qquad\text{따라서}\qquad \Gamma=\sum c\varphi,
\]
그리고 (61.)로부터 (59.)에 대응하는 다음 관계를 얻는다.
\[
 (\theta)=\sum c\Delta^\beta(\varphi),
\]
이 관계는 $\theta$에 대해 정규형 $\varphi$의 편극들에 따른 전개 (60.)을 이끌어 낸다. 요약하면 다음을 얻는다.

\emph{정리 VII. 모든 형식열 $f$는 그와 동치인 정규형의 계로 바꿀 수 있다.}

%
%
\begin{center}
\emph{\S\ 8. 수축 과정의 기본 수축으로의 환원.}
\end{center}
정리 VII에 이어 이제 \S\ 7의 첫머리에서 언급한 수축 과정의 기본 수축으로의 환원을 수행한다. 이때(\S\ 5 참조) 서로 동변인 두 열 $S^\rho,T^\rho$의 모든 수축을 “\emph{동변 수축}”이라 부르고, 특히 두 열 $S^\rho,T^\rho$의 결손 1인 동변 수축을 “\emph{유형 $\rho$의 기본 수축}”이라 부른다. 그러면 정리 VI를 보완하여 다음을 증명한다.

\emph{정리 VIII: 일반 수축 과정 (41.)은 $\rho=1,2,\ldots n-1$인 유형 $\rho$의 $(n-1)$개 기본 수축으로 합성할 수 있다. 다시 말해, 모든 불변 구성식을 생성하려면 $n-1$개의 기본 수축을 차례로 적용하는 것으로 충분하다.}

이항 영역에서는 수축 과정 (41.)이 가능한 유일한 기본 수축과 곧바로 일치한다. 삼항 영역에서는 내가 학위논문의 \S\ 1에서(서론 참조) 정리 VIII을 증명하였다. 그곳에서는 (44.) 때문에 일반적인 동변 수축들이 여전히 기본 수축들과 동일하다. 주목할 점은 정리 VIII이 $n$개 변수에 대하여 기본 수축들로의 환원을 이루어 내며, 훨씬 약한 동변 수축들로의 환원에 그치지 않는다는 것이다. 증명은 원리상 삼항 경우를 본뜬 것으로, 기호적 항등식들을 적용하여 가장 일반적인 수축들을 기본 수축들로부터 구성한다. 우리는 다음을 생성할 것이다.
\begin{enumerate}
\renewcommand{\labelenumi}{\arabic{enumi})}
\item 결손 1의 임의의 수축을 결손 1의 동변 수축, 즉 기본 수축으로부터 생성한다(삼항 경우의 대응물),
\item 결손 $\lambda$의 임의의 수축을 결손 $\lambda$의 동변 수축과 결손 1의 임의의 수축으로부터 생성한다,
\item 결손 $\lambda$의 동변 수축을 결손 $\lambda-1$의 동변 수축과 결손 1의 임의의 수축으로부터 생성한다.
\end{enumerate}
이때 본질적으로 중요한 것은 정리 VII에 따라 수축하려는 두 형식 $S$와 $T$를 정규형으로 가정한다는 점이다. 따라서 (46.)에 따라 다음 관계들이 성립한다.
\srcnumdisplay[4.2em]{\raisebox{-1.7ex}{(62.)}}{%
\begin{aligned}
 \pair{q_{n-\sigma_1-\lambda}S^{\sigma_1}}{S_{n-\sigma_2}p_{\sigma_2-\lambda}}&=0,
 && (\sigma_1\geqq\sigma_2,\ \lambda=1,2,\ldots,\sigma_2,n-\sigma_1),\\
 \pair{q_{n-\tau_1-\lambda}T^{\tau_1}}{T_{n-\tau_2}p_{\tau_2-\lambda}}&=0,
 && (\tau_1\geqq\tau_2,\ \lambda=1,2,\ldots,\tau_2,n-\tau_1).
\end{aligned}}
더 나아가 \S\ 2의 끝부분에 따라 형식 $S$와 $T$가 모든 변수열에 관해 선형이라고 가정하면 충분하다. 즉, 구성되는 형식들은 다음 형식열에 속한다.
\[
 f=\pair{S^1}{p_1}\pair{S^2}{p_2}\cdots\pair{S^{n-1}}{p_{n-1}}
   \pair{T^1}{p_1}\pair{T^2}{p_2}\cdots\pair{T^{n-1}}{p_{n-1}}.
\]

%
%
\emph{1) $\sigma>\tau$일 때 $\pair{q_{n-\sigma-1}S^\sigma}{T_{n-\tau}p_{\tau-1}}$를 유형 $\tau+\alpha$의 기본 수축들로부터 생성하기; $\alpha=0,1,\ldots\sigma-\tau$.} 유형 $\tau$의 기본 수축
\[
 \pair{q_{n-\tau-1}S^\tau}{T_{n-\tau}p_{\tau-1}}
\]
에서 대입
\[
 w\sim T_{n-\tau}p_{\tau-1}
\]
을 통해 위 무게지수가 $\tau+1$인 세 열을 가진 형식열 $f$의 한 형식을 얻는다(\S\ 1 참조).
\srcnumdisplay[4.2em]{(63.)}{%
 \pair{wS^\tau}{p_{\tau+1}}\pair{S^{\tau+1}}{p_{\tau+1}}\pair{T^{\tau+1}}{p_{\tau+1}}.}
(63.)에 유형 $\tau+1$의 기본 수축을 적용하면 다음을 얻는다.
\[
 \pair{q_{n-\tau-2}wS^\tau}{S_{n-\tau-1}p_\tau};
\]
여기에서 (31.)과 대입 $q_{n-\tau-2}w=q'_{n-\tau-1}$을 통해 다음 값을 얻는다.
\[
 \eps\bigl(q_{n-\tau-2}wS^{\tau+1}\bigr)\pair{S^\tau}{p_\tau}
 +\sum_{\alpha=1}^{\tau,n-\tau-1}\Pi\pair{q'_{n-\tau-1-\alpha}S^{\tau+1}}{S_{n-\tau}p_{\tau-\alpha}}.
\]
합 기호 아래의 식은 (62.)에 따라 항등적으로 소멸한다. 따라서 다음 형식에 이른다.
\[
 \pair{wS^{\tau+1}}{p_{\tau+2}}\pair{S^{\tau+2}}{p_{\tau+2}}\pair{T^{\tau+2}}{p_{\tau+2}},
\]
이는 (63.)에서 $\tau$를 $\tau+1$로 바꾸어 생기는 형식이다. 그러므로 이 과정을 $(\sigma-\tau)$번 반복하면 찾던 수축에 이른다.
\[
 \pair{wS^\sigma}{p_{\sigma+1}}=\eps\pair{q_{n-\sigma-1}S^\sigma}{T_{n-\tau}p_{\tau-1}}.
\]
방금 수행한 과정을 다음 공식으로 간단히 나타낸다.
\srcnumdisplay[4.2em]{(64.)}{%
 T^\tau S^\tau,\quad S^{\tau+1},\quad S^{\tau+2},\ldots S^\sigma;}
여기서는 $S$의 정규형 성질만 사용되었고 $T$의 정규형 성질은 사용되지 않았음을 알 수 있다. (64.)와 쌍대인 과정은 $T$의 정규형 성질만 사용하여 역순으로 수행할 수 있다. 곧 다음 과정이다.
\srcnumdisplay[4.2em]{(65.)}{%
 S^\sigma T^\sigma,\quad T^{\sigma-1},\quad T^{\sigma-2},\ldots T^\tau}
이는 다음 관계들에 따른다.
\begin{align*}
\pair{q_{n-\sigma-1}S^\sigma}{T_{n-\sigma}p_{\sigma-1}}&=\eps\bigl(T_{n-\sigma}y p_{\sigma-1}\bigr),\quad \text{여기서 }y\sim q_{n-\sigma-1}S^\sigma,\\
\pair{q_{n-\sigma}T^{\sigma-1}}{T_{n-\sigma}yp_{\sigma-2}}&=\eps\bigl(T_{n-\sigma+1}y p_{\sigma-2}\bigr),\\
&\vdots\\
\pair{q_{n-\tau-1}T^\tau}{T_{n-\tau-1}yp_{\tau-1}}&=\eps\pair{q_{n-\sigma-1}S^\sigma}{T_{n-\tau}p_{\tau-1}}.
\end{align*}
끝으로 \S\ 6의 끝부분에서 제시한 전체 무게의 감소와의 관련성을 지적해 두자. 그 값은 결손 1의 수축에서는 $2(1+(\sigma-\tau))$이고, 기본 수축에서는 $2\cdot1$이다. 따라서 기본 수축들로의 합성이 가능하기만 하다면, 위에서 수행한 것처럼 반드시 $\sigma-\tau+1$개의 기본 수축이 필요하다.

%
%
\emph{2) 일반 수축을 동변 수축과 기본 수축들로부터 생성하기.} 일반 수축 (41.)의 무게 감소는 $2(\lambda^2+\lambda(\sigma-\tau))=2[\lambda^2+(\sigma-\tau)(1+(\lambda-1))]$이다. 따라서 이를 결손 $\lambda$의 동변 수축(무게 감소 $2\lambda^2$)과, 결손이 1이고 무게지수의 차가 $\lambda-1$인 $\sigma-\tau$개의 수축(무게 감소 $2\{1+(\lambda-1)\}$)으로 분해할 가능성이 생긴다. $\lambda=1$인 경우 이 분해는 1)에서 제시한 분해와 일치한다. 이제 이 분해를 실제로 실현할 수 있음을 보이겠다.

결손 $\lambda$의 동변 수축
\[
 \pair{q_{n-\tau-\lambda}S^\tau}{T_{n-\tau}p_{\tau-\lambda}}
\]
에서 대입 $R^\lambda\sim T_{n-\tau}p_{\tau-\lambda}$을 통해 위 무게지수가 각각 $\tau+\lambda$와 $\tau+1$인 두 열을 가진 형식을 얻는다.
\srcnumdisplay[4.2em]{(66.)}{%
 \pair{R^\lambda S^\tau}{p_{\tau+\lambda}}\pair{S^{\tau+1}}{p_{\tau+1}}.}
더 낮은 무게지수 $\tau+1$을 가진 열은 정규형에 속한다. 따라서 (66.)에는 (65.)의 순서에 따라 $\lambda$개의 기본 수축으로 합성된 결손 1의 수축을 적용할 수 있다.
\[
 (R^\lambda S^\tau)S^{\tau+\lambda},\quad S^{\tau+\lambda-1},\quad S^{\tau+\lambda-2},\ldots S^{\tau+1}.
\]
(31.)과 (62.)를 고려하면 이는 다음을 준다.
\[
 \pair{q_{n-\tau-\lambda-1}R^\lambda S^\tau}{S_{n-\tau-1}p_\tau}=\eps\pair{R^\lambda S^{\tau+1}}{p_{\tau+\lambda+1}},
\]
즉, 1)에서 (63.)에 대해 그러했던 것과 마찬가지로 (66.)에서 $\tau$를 $\tau+1$로 바꾸어 생기는 형식을 준다. 이 과정을 $(\sigma-\tau)$번 반복하면 찾던 수축에 이른다.
\[
 \pair{R^\lambda S^\sigma}{p_{\sigma+\lambda}}=\eps\pair{q_{n-\sigma-\lambda}S^\sigma}{T_{n-\tau}p_{\tau-\lambda}}.
\]
이 전체 과정에서는 $S$의 정규형 성질만 사용되었다. $T$를 정규형으로 보면 순서가 완전히 뒤집힌 쌍대 과정에 이르며, 그 공식들은 다음과 같다.
\begin{align*}
\pair{q_{n-\sigma-\lambda}S^\sigma}{T_{n-\sigma}p_{\sigma-\lambda}}&=\eps\bigl(T_{n-\sigma}R_\lambda p_{\sigma-\lambda}\bigr),\quad \text{여기서 }R_\lambda\sim q_{n-\sigma-\lambda}S^\sigma,\\
\pair{q_{n-\sigma}T^{\sigma-1}}{T_{n-\sigma}R_\lambda p_{\sigma-\lambda-1}}&=\eps\bigl(T_{n-\sigma+1}R_\lambda p_{\sigma-\lambda-1}\bigr),\\
&\vdots\\
\pair{q_{n-\tau-1}T^\tau}{T_{n-\tau-1}R_\lambda p_{\tau-\lambda}}&=\eps\pair{q_{n-\sigma-\lambda}S^\sigma}{T_{n-\tau}p_{\tau-\lambda}}.
\end{align*}

%
%
\emph{3) 결손 $\lambda-1$의 동변 수축과 기본 수축들로부터 결손 $\lambda$의 동변 수축을 생성하기.} 결손 $\lambda$의 동변 수축에서 무게 감소는 $2\lambda^2=2((\lambda-1)^2+2\lambda-1)$이므로, 이를 결손 $\lambda-1$의 동변 수축과 $2\lambda-1$개의 기본 수축으로 분해할 가능성이 생긴다. 먼저 특별히 $n=2\lambda$로 놓으면 이 분해를 가장 쉽게 찾을 수 있다. 유일하게 가능한 결손 $\lambda$의 동변 수축은 $\pair{S^{n/2}}{T_{n/2}}=\pair{S^\lambda}{T_\lambda}$가 된다. 이 수축은 같은 열들의 결손 $\lambda-1$인 수축 $\pair{uS^\lambda}{T_\lambda x}$보다 열 $u$와 $x$를 하나씩 덜 포함한다. 역으로 결손 $\lambda$의 수축을, 열 $u$와 $x$를 소멸시키고(무게 감소 $2(n-1)=2(2\lambda-1)$) 동시에 새로운 기호열 $R^\lambda$를 생성하는 기본 수축들과 결손 $\lambda-1$의 동변 수축으로 합성한다. 이를 수행하는 가장 단순한 수축은 다음 결손 1의 수축이다.
\[
 \pair{sT^\lambda}{\sigma p_\lambda}=\eps\pair{S^{2\lambda-1}}{R_{\lambda-1}p_\lambda}=\eps\pair{R^\lambda}{p_\lambda},
\]
여기서 다음과 같이 놓았다.
\[
 R^{\lambda+1}=sT^\lambda,
 \quad s=S^1,
 \quad \sigma=S_1,
 \quad R^\lambda\sim\sigma R_{\lambda-1}.
\]
이 수축은 (65.)와 (64.)에 따라 다음 $2\lambda-1$개의 기본 수축으로 합성된다.
\srcnumdisplay[4.2em]{(67.)}{%
 T^\lambda S^\lambda,S^{\lambda-1},\ldots S^1;
 \qquad R^{\lambda+1}S^{\lambda+1},S^{\lambda+2},\ldots S^{2\lambda-1}.}
결손 $\lambda-1$의 수축을 적용하면 (21.)과 (62.)를 고려하여 다음을 얻는다.
\[
 \pair{uS^\lambda}{\sigma R_{\lambda-1}x}=\eps\pair{R^{\lambda+1}}{S_\lambda x}
 =\eps\pair{sT^\lambda}{S_\lambda x}=\eps\pair{S^\lambda}{T_\lambda}.
\]

%
%
이제 두 열 $S^\rho,T^\rho$와 임의의 $n$에 관한 일반적인 경우는 특수한 경우로부터 직접 옮길 수 있다. (67.) 대신 다음 $2\lambda-1$개의 기본 수축이 들어간다.
\srcnumdisplay[4.2em]{(68.)}{%
 T^\rho S^\rho,S^{\rho-1},\ldots S^{\rho-\lambda+1};
 \qquad R^{\rho+1}S^{\rho+1},S^{\rho+2},\ldots S^{\rho+\lambda-1}.}
수축들의 전반부는 다음 형식을 준다.
\[
 \pair{q_{n-\rho-1}T^\rho}{S_{n-\rho+\lambda-1}p_{\rho-\lambda}},
\]
여기에 대입
\[
 S_{n-\rho+\lambda-1}p_{\rho-\lambda}\sim w,
 \qquad T^\rho w=R^{\rho+1}
\]
을 하고 수축들의 후반부를 적용하면 다음 형식이 생긴다.
\[
 \pair{q_{n-\rho-\lambda}S^{\rho+\lambda-1}}{R_{n-\rho-1}p_\rho}.
\]
다음과 같이 대입하면
\[
 q_{n-\rho-\lambda}S^{\rho+\lambda-1}\sim y,
\]
결손 $\lambda-1$의 동변 수축에 의해 다음을 얻는다.
\srcnumdisplay[4.2em]{(69.)}{%
 \pair{q_{n-\rho-\lambda+1}S^\rho}{yR_{n-\rho-1}p_{\rho-\lambda+1}}.}
이것이 찾던 결손 $\lambda$의 수축임을 보이자. 다음과 같이 대입하고
\[
 R_{n-\rho-1}p_{\rho-\lambda+1}=R_{n-\lambda}
\]
$y$를 풀어 놓으면 다음을 얻는다는 점을 고려하면
\[
 \pair{q_{n-\rho-\lambda+1}R^\lambda}{S_{n-\rho}y}
 =\eps\pair{q_{n-\rho-\lambda}S^{\rho+\lambda-1}}{S_{n-\rho}p'_{\rho-1}}=0
\]
(20.)에 따라 (69.)의 값은 다음과 같다.
\[
 \eps\pair{S^\rho R^\lambda}{p_{\rho+\lambda-1}y}
 =\eps\pair{q_{n-\rho-\lambda}S^{\rho+\lambda-1}}{[S^\rho R^\lambda]_{n-\rho-\lambda}p_{\rho+\lambda-1}}.
\]
그런데 이 수축은 유형 (43.)이므로 다음과 동치이다.
\[
 \pair{q_{n-\rho-\lambda}S^\rho}{R_{n-\rho-1}p_{\rho-\lambda+1}}
 =\eps\pair{wT^\rho}{R_\lambda p_{\rho-\lambda+1}},
 \qquad \text{여기서 }R_\lambda\sim q_{n-\rho-\lambda}S^\rho.
\]
여기서 $R_\lambda$는 이미 원하는 종결형을 갖고 있으며, 이 식은 특수한 경우의 형식 $\pair{sT^\lambda}{S_\lambda x}$에 대응한다. 이제 $w$와 $R_\lambda$를 풀어 놓으면 다음을 얻는다는 점을 고려하자.
\[
 \pair{wR^{n-\lambda}}{T_{n-\rho}p_{\rho-\lambda+1}}
 =\eps\pair{q'_{\lambda-1}R^{n-\lambda}}{S_{n-\rho+\lambda-1}p_{\rho-\lambda}}
 =\eps\pair{q''_{n-\rho-1}S^\rho}{S_{n-\rho+\lambda-1}p_{\rho-\lambda}}=0,
\]
그러면 (21.)에 따라 다음 값을 얻는다.
\[
 \pair{wq_{n-\rho+\lambda-1}}{T_{n-\rho}R_\lambda}
 =\eps\pair{q_{n-\rho+\lambda-1}[T_{n-\rho}R_\lambda]^{\rho-\lambda}}{S_{n-\rho+\lambda-1}p_{\rho-\lambda}}
 \quad\text{동치}\quad
 \pair{q_{n-\rho-\lambda}S^\rho}{T_{n-\rho}p_{\rho-\lambda}}.
\]
이 전체 과정에서는 $S$의 정규형 성질만 사용되었고 $T$의 정규형 성질은 사용되지 않았다. 따라서 (69.)를 생성하기 위한 결손 $\lambda-1$의 동변 수축은 $2\lambda-3$개의 기본 수축과 결손 $\lambda-2$의 동변 수축으로 바꿀 수 있고, 이 후자의 수축도 다시 기본 수축들로 분해할 수 있으며, 이와 같이 계속할 수 있다. 마찬가지로 $T$의 정규형 성질을 사용하고 쌍대적 과정, 즉 다음 기본 수축들을 적용하면 생성이 이루어진다.
\[
 S^\rho T^\rho,\quad T^{\rho+1},\ldots T^{\rho+\lambda-1};
 \qquad R^{\rho-1}T^{\rho-1},\quad T^{\rho-2},\ldots T^{\rho-\lambda+1},
\]
그리고 (69.)와 쌍대적인 수축을 적용한다. 그러면 앞의 수축과 동치인 다음 수축이 생긴다.
\[
 \pair{q_{n-\rho-\lambda}T^\rho}{S_{n-\rho}p_{\rho-\lambda}}.
\]

%
%
2)에서 얻은 분해를 덧붙이면 실제로 \emph{가장 일반적인 수축들을 기본 수축들로부터 생성하는 방법}을 얻는다. 이제 이 생성 방법이 순서의 교환을 제외하면 \emph{유일함}을 증명해야 한다. 즉, 각 수축에는 $\rho=1,2,\ldots n-1$인 유형 $\rho$의 기본 수축 수 $\alpha_\rho$가 오직 하나만 대응함을 보여야 한다.

$m_\rho$를 출발 형식 $f$의 변수열 $p_\rho$에 관한 차수, $l_\rho$를 수축으로 생긴 형식의 차수라 하자. 유형 $\rho$의 각 기본 수축은 차수 $m_{\rho-1},m_\rho,m_{\rho+1}$을 각각 $m_{\rho-1}+1,m_\rho-2,m_{\rho+1}+1$로 바꾼다. 따라서 $\alpha$들은 다음 방정식계를 만족한다.
\srcnumdisplay[4.2em]{(70.)}{%
 m_\rho-l_\rho=-\alpha_{\rho-1}+2\alpha_\rho-\alpha_{\rho+1},
 \qquad (\rho=1,2,\ldots n-1)}
여기서 $\alpha_0$와 $\alpha_n=0$은 0으로 놓아야 하며, 이 방정식계의 행렬식은 값 $n$을 갖는다\srcfn{*)}{$\Delta_n$이 $n$행 행렬식을 나타내면 다음 점화식이 성립한다: $\Delta_n=2\Delta_{n-1}-\Delta_{n-2}$; 즉 $\Delta_n-\Delta_{n-1}=\Delta_{n-1}-\Delta_{n-2}=\cdots=\Delta_2-\Delta_1=1$이다. 왜냐하면 $\Delta_2=3$, $\Delta_1=2$이기 때문이다.}. 따라서 $\alpha$들은 유일하게 결정되며, 정리 VIII에 따라 $0$ 이상의 정수이다. 이로부터 차 $m_\rho-l_\rho$에 대한 조건들이 생긴다.

이 절의 결과들은 \emph{관계 (62.)에 의거하여} 성립한다. 그러나 관계 (62.)가 성립하지 않는 경우에는 종결식들이 시작식들과 동치인 것이 아니다. 실제로 \S\ 6에서 제시한 동치의 정의는 다음과 같이도 표현할 수 있다\srcfn{**)}{이에 대해서는 다음 절 끝의 고찰을 참조하라.}: “두 형식이 더 높은 수축의 형식들, 즉 더 큰 무게 감소를 보이는 형식들만큼 서로 다를 때, 그리고 오직 그때에만 두 형식은 동치이다.” 이 더 큰 무게 감소는 수축에 들어가는 두 열 $q_{n-\sigma-\lambda},p_{\tau-\lambda}$가 실제로 변수열일 때에는 언제나 일어난다. 이는 (41.), (42.), 이 절에서 동치인 것으로 나타나는 수축들, 그리고 정리 VII의 동치인 형식들의 경우에도 그러했다. 그러나 이 열들이 기호열들로부터 대입을 통해 유도되었을 때에는 그러한 무게 감소가 일어날 필요가 없다. 관계 (62.)에 의거하여 서로 표현할 수 있었던 이 절의 형식들이 바로 그러한 경우였다. 소멸하는 형식들 가운데 일부는 오히려 더 큰 전체 무게를 가짐을 쉽게 알 수 있다.

%
%
\begin{center}
\emph{\S\ 9. 형식열. 환원 정리}\srcfn{*)}{삼항 영역에서의 대응하는 고찰을 참조하라. 학위논문 \S\S\ 1--3.}.
\end{center}
\S\ 6(끝)에서 도입한 \emph{형식열}을 이제 조금 더 엄밀하게 다음과 같이 정의한다. 즉 “주어진 출발형식의 계수에 관해 선형인 불변 구성식 전체, 다시 말해 출발형식을 자기 자신에 대해 수축하여 생긴 형식 전체를 \emph{더 높은 형식에 따른 순서로 배열한 것}”이다.

%
%
더 높은 형식을 설명하기 위해, 정리 VII에 따라 출발형식이 두 정규형의 곱 $S\cdot T$로 주어졌다고 하자. 그러면 자기 자신에 대한 더 높은 차수의 수축으로 생긴 형식을 \emph{더 높은 형식}이라 부르고, 자기수축의 차수가 같은 형식들 가운데서는 특별히 정규화한 형식을 그렇게 부른다. 더 정확히 말하면, 형식 $A$가 유형 $\rho$의 기본 수축 $\alpha_\rho$개로, 형식 $B$가 그와 같은 기본 수축 $\beta_\rho$개로 생겼다고 하자. 다음 경우에 $B$는 $A$보다 더 높다.
\begin{enumerate}
\renewcommand{\labelenumi}{\arabic{enumi})}
\item 부등식계 $\beta_\rho\geqq\alpha_\rho$, $\rho=1,2,\ldots n-1$에서 적어도 하나의 $\rho$에 대해 엄격한 부등호가 성립하는 경우\srcfn{**)}{이 부등식계에서 일부에는 부호 $>$가, 일부에는 부호 $<$가 성립하면 두 형식은 서로 비교할 수 없다. 한 형식이 다른 형식에서 수축으로 생긴 경우에는 언제나 양의 기본 수축 값계가 대응하므로, 이 경우에는 $\beta_\rho-\alpha_\rho$가 양수이거나 0이어야 하기 때문이다.};
\item 모든 $\rho$에 대해 등호가 성립하지만 하나의 행렬곱으로 묶인 기호열의 수가 더 적은 경우. 이러한 정규화는 \S\ 8의 환원 때문에 필요함을 알 수 있다. 예를 들어 이에 따라 $\pair{S^{n-1}}{T_{n-1}}$은 $\pair{S^{n-1}}{[S^1T^\rho]_{n-\rho-1}p_\rho}$보다 더 높아진다.
\item 모든 $\rho$에 대해 등호가 성립하고 하나의 행렬곱으로 묶인 기호열의 수마저 같을 때에는, 그 자체로는 임의적이지만 한 번 정한 뒤에는 계속 유지하는 정규화에 의해 $B$를 $A$보다 더 높게 정한다. 하나의 값계 $\alpha$에 속하는 형식의 수는 유한하므로 이러한 정규화는 언제나 가능하다. 예를 들어 한 형식 $S$를 우선시하는 방식(기호열 $S$의 수가 더 많음, 열 $S$의 위쪽 무게지수 합이 더 큼 등)으로 이를 이룰 수 있다.
\end{enumerate}
형식열은 격자점들로 이루어진 $(n-1)$차원 다양체로 배열되어 있다고 생각하는 것이 편리하다. 이때 $n-1$개의 방향은 $n-1$개의 기본 수축에 대응한다. 그러면 각각의 형식에는 하나의 값계 $\alpha$, 곧 하나의 $\alpha_\rho\ge0$인 정수 격자점이 유일하게 대응하고, 하나의 값계 $\alpha$에 속하는 여러 형식의 순서는 위의 정규화로 유일하게 정해진다. 형식열의 임의의 한 형식은 새로운 형식열의 출발형식이 되며, 이 새 형식열은 전체 형식열 안에 한 부분으로 포함된다. 더 정확히는 새 출발형식에서 시작하여 $n-1$개의 양의 방향으로 나아가 얻는 형식들의 부분으로 포함된다.

%
%
정리 VIII과 급수 전개의 일반 이론에 따라 이제 형식열에는 이항 및 삼항 영역의 환원 정리들이 그대로 성립한다(학위논문 \S\ 3 참조). 그 가운데 주정리를 간단히 유도하기에 앞서 다음 정의를 둔다.

“주어진 형식계에서 어떤 유한 개의 형식들이 하나의 모듈로 묶여 있다고 하자. 불변식을 인수로 갖는 형식들을 통해, 또는 ‘더 높은 형식들’을 통해 표현할 수 있는 형식을 \emph{환원 가능}이라 부른다. 여기서 더 높은 형식들이란 그 형식이 속하는 전체 형식열의 더 높은 형식들, 또는 그 모듈의 기호들을 더 높은 차수로 포함하는 형식들을 뜻한다. 환원 가능한 형식열을 \emph{환원자}라 부른다.” 그러면 다음이 성립한다.

\emph{정리 IX: 한 형식열의 출발형식이 어떤 환원자와의 수축으로 생겼기 때문에 환원 가능하다면, 이 출발형식에서 생기는 형식열 전체가 환원 가능하다.}

%
%
증명을 위해, 형식 $f$와 또 다른 형식 $g$를 수축하여 생긴 형식 $h$는 (45.)에 의해 서로 동변인 여러 변수열 $p_\rho,p'_\rho$를 포함하는 형식 $f$로 볼 수 있음을 생각하자. 급수 전개의 일반 이론에 따르면 이러한 형식들은, 서로 동변인 변수열 $p_\rho,p'_\rho$ 두 개씩에 급수 전개를 차례로 적용함으로써 $f$의 기본 공변식들의 편극으로 나타낼 수 있다. 이때 (38.)에 따라 나타나는 기본 공변식들은 동변 수축으로 생성된 형식들이다. 그러나 급수 전개를 적용하는 순서는 여전히 임의적이므로, 특히 기본 수축들로부터 일반 수축을 생성하는 순서에 대응하는 순서를 택할 수 있다. 그러면 생기는 기본 공변식계는 형식열 $f$와 동일해지며, 결손이 더 높은 동변 수축으로 생긴 형식들은 기본 수축으로 생긴 형식들 속에 배열된다. 한편 급수 전개를 임의의 순서로 적용하여도 형식열 $f$의 편극들에 도달하며, 그 순서에는 두 번째 기본 공변식계가 대응한다고 하자. 실제로 형식열은 불변 구성식 전체를 소진하므로, 두 번째 계의 각 형식은 형식열의 편극들, 그중에서도 무게 감소가 같거나 더 큰 형식들의 편극들을 통해 선형으로 표현될 수 있다. 마지막 진술은 다음 사실에서 따른다. 곧 편극들은(\S\ 7 참조) 형식 $h$에 대응하는 차수들을 가진 형식열 $H$ (52.)와 형식열 $f$의 연립불변식으로 볼 수 있다. $H$를 자기 자신에 대해 수축할 때마다 하나의 무게 감소가 대응하며, 이 감소는 치환 (11.)을 시행한 뒤 기호열로, 따라서 $f$의 형식들로 옮겨간다\srcfn{*)}{이 고찰은 동치 개념의 두 번째 표현(\S\ 8 끝)의 바탕이기도 하다. 서로 다른 기본 공변식계에 관한 더 자세한 논의는 삼항 영역에 관한 Study, a. a. O., II. \S\ 7, 정리 7)과 8)에 있다. 우리의 정리 VII에 따라 같은 논의가 $n$개 변수에도 성립한다.}.

%
%
형식열 $h$의 임의의 형식은 $h$를 자기 자신에 대해 수축하여 생긴다. 다시 말해 한편으로는 $g$와 $f$의 더 높은 수축을 통해, 다른 한편으로는 $f$ 또는 $g$를 자기 자신에 대해 수축하여 생긴다. 두 경우 모두 급수 전개는 출발형식 $h$에 관해 위에서 보인 것과 마찬가지로 형식열 $f$의 편극들로 이끈다. 특히 $f$가 환원자라면 환원자의 편극들에 따른 전개가 얻어지므로, 형식열 $h$는 환원 가능하다.

%
%
완전 형식계의 환원에 이 정리를 적용하는 방법은 이항 및 삼항 영역에서와 마찬가지로 얻어진다.

\clearpage
\section*{1911년 논문에 대한 교정 주석 (원본 118--154쪽)}
\addcontentsline{toc}{section}{1911년 논문에 대한 교정 주석 (원본 118--154쪽)}
\begin{enumerate}
\item 원본 118쪽, Study 문헌. 원문은 제1권으로 적고 있다. Eduard Study의 \emph{Zur Differentialgeometrie der analytischen Curven}은 \emph{Transactions of the American Mathematical Society} 제10권 제1호(1909), 1--49쪽에 실렸다. DOI: \url{https://doi.org/10.1090/S0002-9947-1909-1500824-2}.
\item 원본 129쪽, (16) 뒤의 논증. 원문에서는 $a^{(i+1)}$와 $a^{(\rho)}$ 사이의 이어짐 표시가 빠져 있다. $a^{(i+1)}\cdots a^{(\rho)}$라는 표기는 같은 쪽의 항등식 (16)에 근거하며 그 연속된 인수열을 나타낸다.
\item 원본 151쪽, (70) 뒤. $\alpha_\rho$는 각 유형의 기본 수축 횟수를 세므로 영일 수 있다. 따라서 조건은 $\alpha_\rho\in\mathbb Z_{\ge0}$이다. 원본 152쪽의 각주도 해당 횟수에 영을 명시적으로 포함한다. 원문의 양의 정수라는 표현은 여기서 영을 포함하는 조건으로 명확히 한다. 원본 153쪽도 대응하는 격자점을 양의 정수 격자점이라고 부른다. 같은 횟수를 나타내므로 여기서도 $\alpha_\rho\ge0$이며, 수축하지 않은 출발형식에서는 $\alpha=0$이다. 이 명시는 횟수에 관한 것으로, 문단 뒤에서 말하는 양의 방향에 관한 것이 아니다.
\end{enumerate}
\smallskip
\noindent{\small AI가 제작한 교정 주석과 정정: OpenAI Codex, GPT-6 Astra, Ultra 추론 강도.}

\clearpage
\setcounter{footnote}{0}
\phantomsection
\addcontentsline{toc}{section}{5. 유리 함수체}
%
%

%
%
\begin{center}
{\Large\bfseries 5. 유리 함수체.\srcfn{1)}{1913년 빈 자연과학자대회에서 행한 강연.}}\par
\vspace{1.2em}
에를랑겐의 \textsc{에미 뇌터}.\par
\medskip
J. Ber. d. DMV 22 (1913), S. 316--319
\end{center}

\vspace{2em}

다음 문제들은 원래 E. Fischer와 나눈 대화에서 비롯되었다. 다만 이 문제들 가운데 몇 가지는 부정원이 하나인 특수한 경우에 대하여, 그것도 계수영역에는 더 일반적인 가정을 두고, 이미 E. Steinitz가 제기하고 해결하였다.\srcfn{2)}{E. Steinitz: Algebraische Theorie der Körper (특히 \S\ 24). Crelles Journal 137. 1910.}

%
%
\glqq{}유리 함수체\grqq{}란, 주어진 수체의 계수를 갖는 $n$개 부정원의 유리함수들을 원소로 하는 체를 뜻한다. 여기서 그 수체는 특히 모든 복소수를 포함하도록 택할 수도 있다. 이러한 체의 예로는 $n$개 변수의 대칭함수체가 있다. 더 일반적으로는 어떤 군의 치환들을 허용하는 $n$개 부정원의 유리함수 전체, 곧 라그랑주 종영역(Gattungsbereich)이 있으며, 그 밖에 불변식체가 있다.

앞에서 처음 말한 두 체와 마지막 체 사이에는 특징적인 차이가 있다. 앞의 체들은 \emph{$n$개의 대수적으로 독립인 함수}, 곧 기본대칭함수들을 포함하는 반면, 대수적으로 독립인 불변식의 수는 언제나 부정원의 수보다 작다. 따라서 부정원 몇 개를 수로 치환함으로써 \emph{일반적으로 이러한 둘째 종류의 체를 첫째 종류의 체와 일대일로 대응시킬 수 있다}는 사실이 중요하다. 그러므로 이하에서는 \emph{첫째 종류의 체, 즉 $n$개의 대수적으로 독립인 함수를 갖는 체}로 논의를 한정하겠다. 그러면 그 대응에 의해 결과들은 일반적으로 성립한다.

문제들은 유리기저, 최소기저, 정수성 기저라는 세 가지 기본개념을 중심으로 묶인다.

1. \glqq{}유리기저\grqq{}란 체의 유한개의 함수로 이루어진 집합으로서, 체의 모든 함수를 이 유한개의 함수에 대한 유리결합으로 나타낼 수 있고 그 계수는 주어진 계수영역에 속하는 것을 뜻한다. 간단한 고찰을 통해 모든 유리 함수체에 유리기저가 존재함을 보일 수 있다. 부정원이 하나인 경우에는 이러한 고찰이 E. Steinitz의 앞에서 인용한 문헌에도 나온다. 예를 들어 라그랑주의 정리에 따르면 기본대칭함수들과 그 군에 속하는 하나의 함수가 앞에서 언급한 라그랑주 종영역의 유리기저를 이룬다.

2. 첫째 종류의 체에서 모든 유리기저는 적어도 $n$개의 함수를 포함해야 한다. 정확히 $n$개의 함수를 포함하면 그 함수들은 대수적으로 독립이어야 하며, 나는 그러한 기저를 \glqq{}최소기저\grqq{}라 부른다. 최소기저의 존재 문제는 Lüroth, Castelnuovo 및 Enriques의 정리들을 통해 부분적으로 답할 수 있다.\srcfn{1)}{J. Lüroth: Beweis eines Satzes über rationale Kurven. Math. Ann. 9. 1875. --- G. Castelnuovo: Sulla razionalità delle involuzioni piane. Math. Ann. 44. 1893. --- F. Enriques: Sopra una involuzione non razionale dello spazio. Rend. Acc. Linc. vol. 21. 21. Jan. 1912.} 이에 따르면 부정원이 하나인 체에는 언제나 최소기저가 존재한다. 그것은 일차분수변환을 제외하면 유일하게 정해지는 그 체의 Lüroth 함수이다(Steinitz \S\ 24 참조). 부정원이 두 개인 체에서는 계수영역의 대수적 확대를 허용할 때 언제나, 그리고 일반적으로는 그때에만 최소기저가 존재한다. 반면 부정원이 세 개 이상인 경우에는 일반적으로 최소기저가 전혀 존재하지 않는다. 최소기저를 갖는 특수한 체들을 특징짓는 일은 아직 시도되지 않았다.

나는 이제 라그랑주 종영역의 최소기저 문제를 더 추적하였다. 여기에서 이 문제는 다음과 같은 의미를 얻는다. \glqq{}군 $G$에 속하는 라그랑주 종영역이 최소기저를 가지면, 지정된 군 $G$를 갖는 방정식 전체를 유리적 매개변수표현으로 구성할 수 있다. 이는 최소기저의 계수영역을 포함하는 모든 수체에 대하여 성립한다. 따라서 그 계수영역이 유리수체이면 특히 임의의 모든 수체에 대하여 성립한다.\grqq{}

%
%
가장 간단한 예는 대칭군이다. 여기에서는 기본대칭함수들이 유리수 계수를 갖는 최소기저를 이룬다. 이로부터 매개변수표현으로서 부정계수를 갖는 방정식을 곧바로 얻는다. 이 방정식에서 출발하여 각 수체에 대하여 임의로 많은 아펙트 없는(affektlos) 방정식으로 어떻게 넘어가는지는 힐베르트의 불가약성 정리가 보여 준다.

이제 Lüroth와 Castelnuovo의 정리들로부터 3차 및 4차 방정식에 나타나는 모든 군에 대하여 최소기저가 존재한다는 사실이 직접 따라온다. 더 나아가 이 최소기저는 유리수 계수를 갖는다는 것도 드러난다. 따라서 임의의 수체마다 지정된 군을 갖는 3차 및 4차 방정식 전체를 유리적으로 구성할 수 있다. 순환군과 이면체군에 대해서는 $n$차 단위근의 체를 계수영역으로 갖는 최소기저가 존재하며, 몇몇 다른 메타순환군에 대해서도 마찬가지이다. 최소기저가 어떤 군들을 특징짓는 성질인지 여부는 미결로 남겨 두어야 한다.

3. 마지막 기저개념에 이르기 위하여, 체에 들어 있는 모든 다항식의 집합을 생각한다. \glqq{}정수성 기저(Integritätsbasis)\grqq{}란 이 다항식들 가운데 유한개로 이루어진 집합으로서, 체의 모든 다항식을 이 유한개에 대한 다항식적 결합으로 나타낼 수 있고 그 계수는 주어진 계수영역에 속하는 것을 뜻한다. 정수성 기저의 존재 문제는, 예를 들어 불변식계의 유한성을 특수한 경우로 포함하며, 힐베르트가 자신의 \glqq{}수학 문제들\grqq{}\srcfn{1)}{D. Hilbert: Mathematische Probleme. Vortrag Paris 1900. Problem 14. Göttinger Nachr. 1900.}에서 상대적으로 정수적인 함수의 문제라는 조금 다른 형태로 제기하였다.

%
%
현재로서는 정수성 기저가 존재하는 체의 한 부류만을 제시할 수 있다. 즉 그 체가 $n$개의 다항식으로 이루어진 한 계를 포함하고, 이 다항식들의 최고차항들의 동차 종결식이 0이 아니면 정수성 기저가 존재한다. 이 기저는 다음 선형형식에 대하여 그 체에서 불가약인 방정식의 모든 $u$에 관한 계수들로 주어진다:
\[
  u_0=u_1x_1+u_2x_2+\cdots+u_nx_n.\srcfnmark{2)}
\]
\srcfntext{2)}{부정원이 하나인 경우, 이 불가약 방정식은 Lüroth 정리를 증명하기 위한 E. Steinitz의 논의에 나온다.}
특히 종결식의 값이 계수영역의 단위원과 같으면 표현의 정수성도 보장된다. 이 체의 부류에 속하는 한 예는 라그랑주 종영역이다. 기본대칭함수들의 종결식 값이 $\pm 1$이기 때문이다. 또 다른 예는 정수성은 없지만, 대수적으로 독립인 다항식을 단 하나만 포함하는 모든 체이다. 이때 정수성 기저는 모든 다항식을 포함하는 가장 작은 부분체의 Lüroth 함수와 동일하다. 이와 같이 $n$개 변수의 이차형식에 대한 모든 다항 사영불변식은 판별식의 다항함수라는 것이 알려져 있다.

위의 조건은 정수성 기저가 존재하기 위한 충분조건일 뿐, 결코 필요조건은 아니다. 임의의 $n$개 다항식의 종결식이 0이면서도 저 불가약 방정식의 계수들이 주는 정수성 기저를 갖는 체를 쉽게 구성할 수 있다. 따라서 가장 일반적인 경우에도 혹시 저 방정식의 계수들이 정수성 기저를 주는가라는 문제가 생긴다.

\clearpage
\setcounter{footnote}{0}
\phantomsection
\addcontentsline{toc}{section}{6. 유리 함수의 체와 계}
%
%

%
%
\begin{center}
{\Large\bfseries 6. 유리 함수의 체와 계}\par
\vspace{0.6em}
Emmy Noether (Erlangen).\par
\vspace{1.2em}
Math. Ann. 76 (1915), S. 161--196
\end{center}

\vspace{2em}

이 논문은 임의의 유리함수계와 다항함수계에서의 기저 문제를 다룬다. 그런데 여기에 적용되는 체론의 방법에 따르면, 유리 함수들로 이루어진 체 --- 유리 함수체\srcfn{*)}{빈 자연과학자대회 때 발표한 예비보고를 참조하라: Rationale Funktionenkörper --- Jahresber. d. D. Math.-Ver. 22 (1913) ---; 여기에는 문제 설정과 함수체에 관한 결과의 개요가 제시되어 있다.} --- 에서 이 문제들을 해결하는 일이 본질을 이루며, 그 결과를 임의의 계로 일반화하는 일은 그 따름으로 얻어진다.

%
%
일반적인 계의 기저 문제 가운데 지금까지 알려진 것은 힐베르트 정리(Math. Ann. 36)가 보장하는 모듈기저의 존재뿐이다. 이하에서는 유리적 표현가능성의 문제가 임의의 모든 계에 대한 유리기저의 존재와 함께 완전히 해결되며(\S\ 7), 체의 유리기저는 이미 \S\ 4에서 얻어진다. 이 유리기저의 존재 덕분에 처음부터 끝까지 추상적으로 정의된 체나 계에서 출발할 수 있고, 그리하여 계 자체가 아니라 오직 유리기저의 특수한 선택 때문에 생기는 어려움, 예컨대 특수한 분모나 기저함수들의 기본점이 나타나는 일을 피할 수 있다.

체에 대해서는 더 나아가 최소기저, 즉 대수적으로 독립인 함수들로 이루어진 유리기저의 문제도 다룬다(\S\ 6). 체론의 방법은 또한 특별히 구별되는 유리기저인 인볼루션기저(Involutionsbasis)\srcfn{**)}{이 명칭은 유리 함수체가 기하학적 해석에서 선형공간의 가장 일반적인 인볼루션을 나타내기 때문에 택하였다.}(\S\ 5)로 이끈다. 이 기저는 특히 다항식들로 이루어진 정역에서 본질적인 역할을 한다(\S\ 8). 여기서 이 기저는 고정된 분모를 갖는 표현을 주는데, 그 분모는 정역이 정하는 한 함수의 거듭제곱이다. 이러한 표현은 가령 불변식의 전형표현이라는 특수한 경우에 알려져 있다.

%
%
이제 유리적 표현 가능성으로부터 다항식적 표현 가능성, 즉 좁은 의미의 유한성 문제에 대하여 가능한 한 폭넓은 결론을 이끌어 낸다.

지금까지 알려진 유한성 정리들은 모두 다음 사실에 근거한다. 곧 모듈기저에 관한 힐베르트 정리는
\[
  F=A_1f_1+A_2f_2+\cdots+A_kf_k
\]
라는 표현이
\[
  F=B_1f_1+B_2f_2+\cdots+B_kf_k
\]
라는 또 하나의 표현을 수반하고, 여기서 $B_i$가 그 계에 속하며 $F$보다 낮은 차수를 갖는 모든 계에 대하여 유한성을 보장한다.

이에 반해 유리기저 정리와 체론의 방법을 이용하면 전혀 다른 종류의 가정 아래에서도 유한성을 결론지을 수 있다.

%
%
그리하여 힐베르트의 한 문제(Math. Probleme 14)에 대한 부분적인 답으로서, 유한 정역이며 앞서 말한 인볼루션기저가 그 정기저를 이루는 상대적으로 정수적인 함수들의 한 부류(제1종 상대정역)를 얻게 된다(\S\ 10). 대수적 정수인 양들의 기본계에 관한 크로네커 정리에 대한 힐베르트의 증명(Volle Invariantensysteme, \S\ 2; Math. Ann. 42)과 유사하게, 더 나아가 다항식으로 이루어진 모든 정칙계 --- 즉 기본점이 없는 계로 사상할 수 있는 모든 계 --- 가 정기저를 가짐을 보일 수 있다(\S\ 12). 한편 정기저가 없는 비정칙계도 쉽게 제시할 수 있다. 이 사실의 간단한 예로 다음 정리를 들 수 있다. \glqq{}한 부정원에 관한 무한히 많은 다항식으로 이루어진 임의의 계에는 유한개의 다항식이 존재하여, 그 계의 모든 다항식을 이 유한개에 대한 다항식 결합으로 나타낼 수 있다.\grqq{} 더 많은 예는 \S\ 13에 있다.

끝으로 마지막 절들(\S\S\ 14 및 15)은 정수계수성에 관하여 기저정리들을 강화한다.

이 연구의 본질적인 보조수단은 \S\ 3의 전이 원리이다. 이에 따르면 여기에서 제기한 모든 기저 문제는 부정원의 수와 그 계의 대수적으로 독립인 함수의 수가 일치하는 계들에 대해서만 고찰하면 충분하다.

%
%
본 연구의 계기는 E. Fischer와 나눈 대화, 특히 그가 라그랑주 종영역(Gattungsbereich)의 최소기저에 관하여 제기한 질문이었다. 그와 관련된 연구는 지정된 군을 갖는 방정식의 구성으로 이어졌으며(앞에서 인용한 \glqq{}유리 함수체\grqq{}에 관한 보고의 제2부를 참조하라), 추후 발표로 남겨 둔다.

본 논문이 최초의 보고에 비해 갖는 본질적인 확장은, 그곳에서는 체나 특수한 정역에 대해서만 진술되었던 기저정리들을 임의의 계로 옮겼다는 데 있다. 이러한 이행은 임의의 계를 포함하는 최소체라는 개념, 곧 정역의 분수체를 일반화한 개념에 힘입어 간단하게 이루어진다(\S\ 7). 체의 유리기저가 알려진 뒤 K. Hentzelt가 힐베르트의 모듈기저를 거치는 우회로를 통해 임의의 계에도 유리기저가 존재함을 증명할 수 있었던 사실이 이 개념을 세우도록 나를 이끌었다. 또한 내가 정역에 적용한 추론들이 동일한 전제를 만족하는 임의의 계에도 유효하다는 K. Hentzelt의 지적은 정칙계의 정기저에 관한 정리로도 나를 이끌었다. 그 밖의 몇 가지 개별적인 작은 지적 역시 Hentzelt에게 빚지고 있다.

%
%
끝으로, 부정원이 하나인 경우 체의 유리기저와 최소기저가 E. Steinitz의 \glqq{}체의 대수적 이론\grqq{}, \S\ 24 (Crelles Journ. 137)에 나와 있음을 언급해 둔다. 다만 그곳에서는 계수영역을 가장 일반적인 의미의 체로 가정한다. 여기에서 제시한 정리들이 이러한 더 일반적인 전제 아래 어디까지 유지되는지는 이하에서 다루지 않는다. 어쨌든 증명에는 수정이 필요하다. 예컨대 표수 $p$인 체에서는 야코비 행렬식 이론에서 가져온 보조수단부터 이미 쓸 수 없기 때문이다.

%
%
\section*{\S\ 1. 체 $\Kfield_{n\rho}$와 계 $\Ssys_{n\rho}$.}

이하에서는 $n$개 부정원의 유리함수계, 특히 유리함수체를 고찰한다.

%
%
$f(x_1\cdots x_n)$, $g(x_1\cdots x_n)$, \ldots\ 또는 줄여서 $f(x)$, $g(x)$, \ldots\ 로 적은 것은 언제나 $x_1\cdots x_n$의 유리함수를 뜻한다. 이 함수들은 기약표현 --- 즉 분자와 분모가 서로소인 표현 --- 으로 주어진다고 가정한다. 다항함수(다항식)는 대문자 $F(x)$, $G(x)$, \ldots\ 로 표시한다. 계수영역 $\Omega$는 어떤 수체로 가정하며, 따라서 특히 모든 복소수를 포함할 수도 있다. $\Omega$에는 또한 유한개의 매개변수가 들어 있을 수도 있다.

$\Omega$의 원소들, 즉 모든 0차 함수는 언제나 그 계에 속하는 것으로 간주한다.

이러한 약정 아래 유리함수들로 이루어진 체 --- 유리함수체 --- 를 추상적으로 정의할 수 있다.

%
%
\emph{정의 I: 유리함수계가 다음 조건들을 만족할 때 이를 체라 한다:}
\begin{enumerate}
\item \emph{그 계가 $f(x)$를 포함하면, $\Omega$의 모든 원소 $c$에 대하여 $c\cdot f(x)$도 포함한다.}
\item \emph{그 계가 $f(x)$와 $g(x)$를 포함하면 언제나 $f(x)+g(x)$, $f(x)\cdot g(x)$도 포함하며, $g(x)\ne0$일 때에는 몫 $f(x):g(x)$도 포함한다.}\srcfn{*)}{$f(x)$와 $g(x)$는 0차일 수도, 즉 $\Omega$의 원소일 수도 있다.}
\end{enumerate}
특수한 체란 $\Omega$에 유리함수 유한개
\[
  f_1(x),\ f_2(x),\ \cdots,\ f_k(x)
\]
를 첨가하여 생기는 체를 말하며, 이를
\[
  \Omega\bigl(f_1(x)\cdots f_k(x)\bigr)
  \quad\text{또는}\quad
  \Omega(f_1\cdots f_k)
\]
로 표시한다.\srcfn{**)}{이 \glqq{}특수한 체들\grqq{}로 이미 모든 체가 망라된다는 사실은 \S\ 4에서 도출된다.} 이 특수한 체들 가운데에는 $\Omega$에 $x_1\cdots x_n$을 첨가하여 생기는 체
\[
  \Omega(x_1\cdots x_n)
\]
도 언급해 둔다. 이 체는 $\Omega$의 계수를 갖는 $x_1\cdots x_n$의 모든 유리함수 전체와 동일하다.

%
%
이제 (E. Steinitz를 따라) 중간체의 개념도 도입한다:
\begin{quote}
\glqq{}체 $\Mfield$가 $\mathfrak{L}_1$과 $\mathfrak{L}_2$ 사이에 놓인다는 것은, 그 체가 $\mathfrak{L}_1$을 포함하고 $\mathfrak{L}_2$에 포함된다는 뜻이다\grqq{};
\end{quote}
그러면 정의 I은 다음과 같이도 표현할 수 있다:

$n$개 부정원의 유리함수체는 $\Omega$와 $\Omega(x_1\cdots x_n)$의 중간체로서, 그 체의 적어도 하나의 함수가 이 $n$개의 부정원을 모두 실질적으로 포함한다.

%
%
부정원의 수 외에 유리함수계에는 또 하나의 특성수, 곧 대수적으로 독립인 함수의 수 또는 대수적 랭크가 있으며, 이는 다음 정의로 주어진다.

\emph{정의 II: 유리함수계에서 대수적으로 독립인 $\rho$개의 함수를 제시할 수 있지만, 그 계에서 택한 임의의 $(\rho+1)$개 함수는 대수적으로 종속일 때, 그 계는 대수적 랭크 $\rho$를 갖는다.}\srcfn{***)}{$\rho$개의 함수 $f_1(x)\cdots f_\rho(x)$는 모든 관계
\[
F\bigl(f_1(x)\cdots f_\rho(x)\bigr)=0\quad\text{에 관한 항등식}x_1\cdots x_n
\]
로부터 다음이 따라 나올 때 대수적으로 독립이라 한다.
\[
F(\lambda_1\cdots \lambda_\rho)=0\quad\text{에 관한 항등식}\lambda_1\cdots \lambda_\rho.
\]
반대의 경우에는 이 함수들을 대수적으로 종속이라 한다.}

%
%
$n$개의 부정원에 관한 (유한개 또는 무한개의) 유리함수로 이루어지고 대수적 랭크가 $\rho$인 계는 이중 첨자를 사용하여
\[
  \Ssys_{n\rho}
\]
로 나타내기로 한다. 특히
\[
  \Kfield_{n\rho}
\]
는 $n$개의 부정원에 관한 대수적 랭크 $\rho$의 체를 뜻한다. 다음 부등식이 성립한다.
\[
  1\leq \rho\leq n.
\]

%
%
체 $\Kfield_{nn}$과 $\Kfield_{n\rho}$의 몇 가지 예를 더 제시한다.

1. 체 $\Kfield_{nn}$은 라그랑주 종영역이며, 이는 다음과 같이 정의할 수 있다.
\begin{quote}
\glqq{}$x_1\cdots x_n$의 유리함수(및 다항함수) 가운데 $x_1\cdots x_n$의 한 순열군 $G$의 순열들을 허용하는 함수 전체.
\grqq{}
\end{quote}
이 체들은 항상 대칭함수체를 포함하며, 따라서 대수적으로 독립인 $n$개의 기본대칭함수도 포함한다.

%
%
2. 체 $\Kfield_{n\rho}$의 예로는 하나 또는 여러 기본형식의 유리불변식(및 다항불변식) 전체로 이루어진 사영불변식체가 있다.\srcfn{*)}{행렬식이 $1$인 변환만을 고려하는 편이 적절하다. 그러면 임의로 많은 불변식의 모든 유리결합은 다시 그 변환을 허용하며, 실제로 체를 얻게 된다. 동차인 등중(等重) 불변식만 따로 취하면 체를 이루지 않지만, 계 $\Ssys_{n\rho}$는 이룬다.} 대수적으로 독립인 불변식의 수는 언제나 계수의 수보다 작으므로 $\rho<n$이다.

사영군의 부분군에 대한 불변식체에도 같은 사실이 성립한다.

%
%
\section*{\S\ 2. 대수적으로 종속인 함수들에 관한 보조정리.}

이 절의 보조정리에 의해 \S\ 3에서 한 계의 대수적 랭크가 진정한 특성수임이 드러날 것이다. 즉 이 보조정리는, $\rho$개의 대수적으로 독립인 함수가 대수적으로 독립인 채로 남도록 몇몇 부정원에 잇따라 수치를 특수화할 때, 이 $\rho$개 함수에 대수적으로 종속인 임의의 함수가 항등적으로 영이 되거나 부정형이 되는 일이 없음을 보여 준다.

%
%
그러므로 계
\[
\begin{array}{@{}r@{\quad}l@{}}
(1)& g_1(x),\ g_2(x),\ldots,g_\rho(x)
\end{array}
\]
가 대수적 랭크 $\rho$ --- 여기서 $\rho<n$ --- 를 갖고, $h(x)$가 함수들 (1)에 대수적으로 종속인 유리함수라고 하자. 알려진 정리들에 따르면 (1)의 야코비 행렬의 랭크는 $\rho$이다. 따라서 부정원들은 다음이 성립하도록 번호가 매겨져 있다고 하자:
\[
\begin{array}{@{}r@{\quad}l@{}}
(2)& \displaystyle
\frac{\partial(g_1\ g_2\cdots g_\rho)}{\partial(x_1\ x_2\cdots x_\rho)}
=\frac{\Gamma(x)}{G(x)^{\rho+1}}\ne0,
\end{array}
\]
여기서 잘 알려진 바와 같이 $G(x)$는 함수들 (1)의 공통분모를 뜻한다.

%
%
\emph{이제 수 $a$를 곱 $G(x)\cdot\Gamma(x)$가 $(x_i-a)$로 나누어떨어지지 않도록 택하자}. 여기서 $i$는 $\rho+1,\rho+2,\ldots,n$ 가운데 하나를 뜻한다. 그러면 $x_i=a$일 때 함수들 (1)의 공통분모는 0이 될 수 없고, $\rho$개의 함수
\[
\begin{array}{@{}r@{\quad}l@{}}
(3)& [g_1(x)]_{x_i=a},\ldots,[g_\rho(x)]_{x_i=a}
\end{array}
\]
도 대수적으로 독립이다. 앞으로 계 (1)의 대수적 랭크는 대입 $x_i=a$에 의해 보존된다고 말하기로 한다.

%
%
함수들 (1)에 관하여 $h(x)$가 만족하는 기약방정식이 다음과 같이 주어졌다고 하자:
\[
\begin{array}{@{}r@{\quad}l@{}}
(4)& A_0(g_1\cdots g_\rho)\cdot h(x)^\tau
 +A_1(g_1\cdots g_\rho)\cdot h(x)^{\tau-1}
 +\cdots+A_\tau(g_1\cdots g_\rho)=0,
\end{array}
\]
여기서 $A_0(g_1\cdots g_\rho)$와 $A_\tau(g_1\cdots g_\rho)$는 항등적으로 영이 되지 않는다. 이 식으로부터 다항식들 사이의 항등식을 얻기 위해 다음과 같이 둔다:
\[
\begin{array}{@{}r@{\quad}l@{}}
(5)& g_i(x)=G_i(x):G(x);\qquad h(x)=H_1(x):H_2(x),
\end{array}
\]
그리고 $h(x)$가 기약표현으로 주어졌다는 가정에 의해 $H_1(x)$과 $H_2(x)$가 서로소임을 주목한다. (5)에 의해 (4)는 다음과 같이 바뀐다:
\[
\begin{array}{@{}r@{\quad}l@{}}
(6)& \begin{aligned}[t]
&B_0(G\,G_1\cdots G_\rho)\cdot G(x)^{\sigma_0}\cdot H_1(x)^\tau
+B_1(G\,G_1\cdots G_\rho)\cdot G(x)^{\sigma_1}\cdot H_1(x)^{\tau-1}H_2(x)\\
&\quad+\cdots+B_\tau(G\,G_1\cdots G_\rho)\cdot G(x)^{\sigma_\tau}\cdot H_2(x)^\tau=0,
\end{aligned}
\end{array}
\]
여기서 $B_k$는 $G,G_i$의 동차형식이고, 음이 아닌 지수 $\sigma_k$ 가운데 적어도 하나는 0이어야 한다.

%
%
이제 $H_1(x)$가 $(x_i-a)$로 나누어떨어진다고 가정하면, (6)에 의해
\[
  B_\tau(G\,G_1\cdots G_\rho)\cdot G(x)^{\sigma_\tau}\cdot H_2(x)^\tau
\]
도 $(x_i-a)$로 나누어떨어질 것이다. 가정에 따라 $G(x)$에 대해서는 이것이 배제되어 있다. 따라서 $B_\tau$가 나누어떨어진다면
\[
  A_\tau=B_\tau:G^\lambda
\]
도 나누어떨어지게 되고, 함수들 (3) 사이에 관계 $A_\tau=0$이 성립하게 되는데, 이는 가정한 대수적 독립성에 어긋난다.
끝으로 $H_1(x)$과 서로소인 $H_2(x)$도 $(x_i-a)$로 나누어떨어질 수 없다\srcfn{*)}{기약표현에 근거한 이 마지막 결론은, 나머지 가정들을 유지하면서 $x_i=a$, $x_k=b$를 대입하는 경우에는 더 이상 내릴 수 없을 것이다. 그 경우 $H_1(x)$과 $H_2(x)$가 동시에 영이 될 수 있으므로 몫은 대입에 의해 부정형 $=\frac00$이 될 것이다. 예컨대
\[
 h(x)=\frac{x_3(\alpha x_1+\beta x_2)+x_4(\gamma x_1+\delta x_2)}
 {x_3(\alpha' x_1+\beta' x_2)+x_4(\gamma' x_1+\delta' x_2)}
\]
에 $x_3=0$, $x_4=0$을 대입하는 경우가 그러하다.}; 따라서 $H_1(x)$가 $(x_i-a)$로 나누어떨어진다는 가정은 모순에 이른다. 같은 방식으로 $H_2(x)$도 $(x_i-a)$로 나누어떨어질 수 없음을 보인다. 그러므로 함수 $[h(x)]_{x_i=a}=k(x)$는 항등적으로 영이 될 수도, 부정형이 될 수도 없다.

%
%
계 (3)과 다시 기약표현으로 주어졌다고 가정하는 함수 $k(x)$에는 위의 논법을 다시 적용할 수 있으며, 이 절차를 계속하면 마침내 다음 보조정리에 이른다.

\emph{보조정리: 두 계}
\[
\begin{array}{c}
g_1(x),\ g_2(x),\ldots,\ g_\rho(x),\\
g_1(x),\ g_2(x),\ldots,\ g_\rho(x),\ h(x)
\end{array}
\]
\emph{는 각각 대수적 랭크 $\rho$ --- 여기서 $\rho<n$ --- 를 갖는다고 하고, 예컨대 다음이 성립한다고 하자.}
\[
 \frac{\partial(g_1\cdots g_\rho)}
 {\partial(x_1\cdots x_\rho)}
 =\frac{\Gamma(x)}{G(x)^{\rho+1}}\ne0.
\]

%
%
\emph{수들}
\[
 a_n,\ a_{n-1},\ldots,\ a_{\rho+1}
\]
\emph{은 다음 대입에 의해}
\[
 x_n=a_n,\quad x_{n-1}=a_{n-1},\ldots,\quad x_{\rho+1}=a_{\rho+1}
\]
\emph{곱 $G(x)\cdot\Gamma(x)$가 영이 되지 않도록 --- 즉 계 (1)의 대수적 랭크가 보존되도록 --- 택한다고 하자. 잇따른 대입들}
\[
\begin{array}{c}
[h(x)]_{x_n=a_n}=h^{(n-1)}(x);\quad
[h^{(n-1)}(x)]_{x_{n-1}=a_{n-1}}=h^{(n-2)}(x);\quad\cdots;\\[0.25em]
[h^{(\rho+1)}(x)]_{x_{\rho+1}=a_{\rho+1}}=h^*(x_1,x_2,\ldots,x_\rho)
\end{array}
\]
\emph{에서는 함수들 $h(x),h^{(n-1)}(x),\ldots,h^{(\rho+1)}(x)$을 각각 기약표현으로 만든다고 하자. 이 가정들 아래에서는 함수 $h^*(x_1\cdots x_\rho)$의 분자도 분모도 항등적으로 영이 될 수 없으며, 이 함수는 $\frac00$이 될 수도 없다.}\srcfn{**)}{잇따른 대입은 예컨대 앞 주석의 함수 $h(x)$에서 다음을 준다:

%
%
\[
 [h(x)]_{x_4=0}=h^{(3)}(x)=\frac{(\alpha x_1+\beta x_2)}{(\alpha' x_1+\beta' x_2)}.\qquad
 h^*(x_1 x_2)=\frac{\alpha x_1+\beta x_2}{\alpha' x_1+\beta' x_2}.
\]}

%
%
\section*{\S\ 3. 계 $\Ssys_{n\rho}$에서 계 $\Ssys^*_{\rho\rho}$로의 일대일 대응.}

%
%
보조정리로부터 몇몇 부정원을 수로 바꿈으로써 계 $\Ssys_{n\rho}$를 계 $\Ssys^*_{\rho\rho}$로 보내는 사상을 곧바로 얻게 된다. 이에 따라 대수적 랭크 $\rho$는 진정한 특성수로 나타난다.

%
%
$\Ssys_{n\rho}$의 대수적 랭크가 $\rho$이므로, $\Ssys_{n\rho}$에는 대수적으로 독립인 $\rho$개의 함수
\[
\begin{array}{@{}l@{\quad}c@{}}
(1)& g_1(x)\cdots g_\rho(x),
\end{array}
\]
가 존재한다. 또한 $f(x)$가 $\Ssys_{n\rho}$의 임의의 함수를 나타내면, 계
\[
\begin{array}{@{}l@{\quad}c@{}}
(2)& g_1(x)\cdots g_\rho(x),\ f(x)
\end{array}
\]
도 대수적 랭크 $\rho$를 가지며, 계 (1)과 (2)는 보조정리의 조건을 충족한다.\srcfn{*)}{필요하다면 부정원들의 번호를 다시 매긴다.}

%
%
그러므로 보조정리에 따라 수
\[
 a_n,\ a_{n-1},\ldots,\ a_{\rho+1}
\]
를, 계 (1)의 대수적 랭크 $\rho$가 대입에 의해 보존되도록 정하면 --- 이는 임의로 많은 방식으로 가능하다 --- 다음과 같이 정의되는 함수에서
\[
\begin{array}{c}
[f(x)]_{(x_n=a_n)}=f^{(n-1)}(x);\quad
[f^{(n-1)}(x)]_{(x_{n-1}=a_{n-1})}=f^{(n-2)}(x);\\[0.25em]
[f^{(\rho+1)}(x)]_{(x_{\rho+1}=a_{\rho+1})}=f^*(x_1\cdots x_\rho)
\end{array}
\]
곧
\[
\begin{array}{@{}l@{\quad}c@{}}
(3)& f^*(x_1\cdots x_\rho)
\end{array}
\]
에서는 분자도 분모도 항등적으로 사라질 수 없다.

%
%
이제 $f(x)$가 $\Ssys_{n\rho}$에 속하는 모든 함수(유한개일 수도 무한개일 수도 있다)를 두루 나타내게 하자. 모든 함수 (3)의 전체로 이루어진 계를 생각하자. 이 계는 다음 성질을 갖는다:

%
%
1. 이 계는 대수적 랭크 $\rho$를 갖는다. 실제로 이 계는 (1)에서 생기는 함수들
\[
g_1^*(x_1\cdots x_\rho)\cdots g_\rho^*(x_1\cdots x_\rho),
\]
을 포함하며, 이 함수들은 $a_n,a_{n-1},\ldots,a_{\rho+1}$의 선택으로 인해 대수적으로 독립이다. 그러므로 이 계를
\[
\Ssys^*_{\rho\rho}
\]
로 나타내기로 한다.\srcfn{**)}{여기서 $f^*(x)$, $g^*(x)$, \ldots 는 이하 각각 모두 사상 함수, 즉 $x_1\cdots x_\rho$의 함수로 이해한다.}

%
%
2. 함수 $f(x)$와 $f^*(x)$ 사이의 대응은 예외 없이 일대일이다.

실제로 $f_1(x)$와 $f_2(x)$가 $\Ssys_{n\rho}$에 속하면, 그 차
\[
d(x)=f_1(x)-f_2(x)
\]
도 계 (1)에 대수적으로 종속이며, 또한 다음이 성립한다:
\[
d^*(x)=f_1^*(x)-f_2^*(x).
\]
$d(x)$는 계 (1)과 함께 보조정리의 조건을 충족하므로,
\[
d(x)\ne0
\]
이면 반드시
\[
d^*(x)\ne0
\]
이다. 즉 $\Ssys_{n\rho}$의 서로 다른 함수들에는 $\Ssys^*_{\rho\rho}$의 서로 다른 함수들이 대응한다.

%
%
3. $\Ssys^*_{\rho\rho}$에 속하는 함수들 사이의 임의의 관계식
\[
F\bigl(f_1^*(x),f_2^*(x),\ldots,f_k^*(x)\bigr)=0
\quad\text{$x_1\cdots x_\rho$에 관하여 항등적으로}
\]
으로부터 $\Ssys_{n\rho}$에 속하는, 각각 유일하게 대응하는 함수들에 대해 다음이 따라 나온다:
\[
F\bigl(f_1(x),f_2(x),\ldots,f_k(x)\bigr)=0
\quad\text{$x_1\cdots x_n$에 관하여 항등적으로}.
\]
실제로 $f_1(x),f_2(x),\ldots,f_k(x)$와 마찬가지로 이 함수들의 모든 유리 결합도 계 (1)에 대수적으로 종속이다.

%
%
또한 (4)의
\[
\begin{array}{@{}l@{\quad}l@{}}
(4)& h(x)=F\bigl(f_1(x)\cdots f_k(x)\bigr):\\
& h^*(x)=F\bigl(f_1^*(x)\cdots f_k^*(x)\bigr).
\end{array}
\]
로부터, 보조정리에 따라
\[
h(x)\ne0\quad\text{이면 또한}\quad h^*(x)\ne0;
\]
이 따른다.
q. e. d.

%
%
요약하면 다음 정리를 얻는다.

\emph{정리 I: (유한 또는 무한히 많은) 유리함수들로 이루어진 각 계 $\Ssys_{n\rho}$에는, $\Ssys^*_{\rho\rho}$의 함수들 사이에 성립하는 모든 관계가 $\Ssys_{n\rho}$에서도 보존되고 그 역도 성립하도록, 계 $\Ssys^*_{\rho\rho}$를 임의로 많은 방식으로 일대일 대응시킬 수 있다. 이때 체(體) $\Kfield_{n\rho}$에는 (4)에 따라 특히 다시 체(體) $\Kfield^*_{\rho\rho}$가 대응한다.}

%
%
정리 I로부터 앞으로의 모든 증명을 단순하게 하는 다음 따름결과도 얻는다. \emph{계 $\Ssys_{n\rho}$의 함수들 사이의 유한 또는 무한히 많은 관계로 특징지어지는 계 $\Ssys_{n\rho}$의 성질은, 그것이 하나의 사상계 $\Ssys^*_{\rho\rho}$에 대하여 성립하기만 하면 계 $\Ssys_{n\rho}$에 대해서도 성립한다.}

이 따름결과를 줄여서 \glqq{}전이 원리\grqq{}라고 부른다.

%
%
이 사상의 가장 간단한 예는 대수적 불변식 이론이다. 실제로 선형변환을 통해 기본형식의 몇몇 계수를 언제나 수치적으로 정하여, 특수형식에 대해 성립하는 모든 관계가 일반적으로도 보존되도록 할 수 있다.

%
%
\section*{\S\ 4. 체(體) $\Kfield_{n\rho}$의 유리기저.}

%
%
이하에서 기저 개념을 중심으로 전개되는 연구에서는 시종일관 \S\ 3의 \glqq{}전이 원리\grqq{}를 이용한다. 우선 체(體)에 대하여 기저의 존재를 증명하지만, 정의들은 일반적으로 제시한다:

%
%
\emph{정의 III: $\Ssys_{n\rho}$에 속하는 유한개의 함수는, $\Ssys_{n\rho}$의 모든 함수를 $\Omega$의 계수를 갖는 이 유한개 함수의 유리결합으로 나타낼 수 있을 때 유리기저라 한다.}

%
%
유리기저로부터 유도된 계의 대수적 랭크는 그 유리기저의 대수적 랭크와 같으므로, 유리기저는 대수적으로 독립인 함수 $\rho$개를 포함해야 한다.

%
%
먼저 유리기저의 존재 증명에 전이 원리를 적용할 수 있음을 보인다. 실제로 정의 III는 다음과 같이 말할 수 있다:

$\Ssys_{n\rho}$에 속하는 함수들
\[
g_1(x),\ g_2(x),\ldots,g_k(x)
\]
는 다음 무한히 많은 관계가 성립할 때 그리고 그때에만 유리기저를 이룬다:
\[
\begin{array}{@{}l@{\quad}c@{}}
(1)& f(x)=\gamma\bigl(g_1(x)\cdots g_k(x)\bigr),
\end{array}
\]
여기서 $f(x)$는 $\Ssys_{n\rho}$의 모든 함수를 두루 나타내고, $\gamma$는 $f$에 따라 달라지는 $k$개 인수의 유리함수를 뜻한다.

%
%
그러므로 사상계 $\Ssys^*_{\rho\rho}$에
\[
g_1^*(x),\ g_2^*(x),\ldots,g_k^*(x)
\]
로 주어진 유리기저가 있고, 이것이 관계
\[
f^*(x)=\gamma\bigl(g_1^*(x)\cdots g_k^*(x)\bigr)
\]
를 귀결한다면, 정리 I에 따라 $\Ssys_{n\rho}$에서 관계 (1)이 성립한다. 즉 \emph{유리기저에 대한 전이 원리}는 다음과 같다:

\emph{사상계 $\Ssys^*_{\rho\rho}$에 유리기저가 존재하면 원래 계 $\Ssys_{n\rho}$에도 유리기저가 존재한다.}\srcfn{*)}{여기서 다루는 각 기저 문제에 대해서도 전이 원리는 이에 상응하게 진술된다.}

%
%
모든 체(體) $\Kfield_{n\rho}$에 대한 유리기저의 존재를 증명하기 위해, 따라서 우리는 체(體) $\Kfield^*_{\rho\rho}$만을 다루면 충분하다. 여기서는 E. Steinitz가 사용한 논법을 직접 일반화하면 목적에 이른다.\srcfn{*)}{E. Steinitz, Algebraische Theorie der Körper, \S\ 24, Crelles Journ. 137 (1909).}

%
%
$\Kfield^*_{\rho\rho}$에서 택한 $\rho$개의 대수적으로 독립인 함수로 이루어진 다음 계를 생각하자.
\[
\begin{array}{@{}l@{\quad}c@{}}
(2)& g_1^*(x_1\cdots x_\rho),\ldots,g_\rho^*(x_1\cdots x_\rho)
\end{array}
\]
$x_1\cdots x_\rho$를 소거하면 $\rho$개의 관계식이 얻어진다.
\[
\begin{array}{@{}l@{\quad}c@{}}
(3)& F_1\bigl(g_1^*(x)\cdots g_\rho^*(x),x_1\bigr)=0,\ldots,
F_\rho\bigl(g_1^*(x)\cdots g_\rho^*(x),x_\rho\bigr)=0,
\end{array}
\]
각 관계식 $F_i=0$에는 부정원 $x_i$가 실제로 나타난다. 그렇지 않으면 가정과 달리 계 (2)의 함수들 사이에 대수적 종속관계가 생길 것이기 때문이다. 관계식 (3)은 다음을 뜻한다.

%
%
\[
x_1,x_2,\ldots,x_\rho
\]
가
\[
\Omega\bigl(g_1^*(x)\cdots g_\rho^*(x)\bigr)
\]
에 관하여 대수적 원소라는 것이다. 따라서 $x_1\cdots x_\rho$를 첨가하여 생기는 체(體)
\[
\Omega\bigl(g_1^*\cdots g_\rho^*;\ x_1\cdots x_\rho\bigr)=\Omega(x_1\cdots x_\rho)
\]
는 $\Omega(g_1^*\cdots g_\rho^*)$ 위의 (유한) 대수적 확대체이다. 알려진 정리\srcfn{**)}{예컨대 Weber, Lehrbuch d. Algebra 1, \S\ 144 (제1판) 또는 \S\ 55 (축약판) 참조.}에 따라 $\Kfield^*_{\rho\rho}$는 $\Omega(g_1^*\cdots g_\rho^*)$와 $\Omega(x_1\cdots x_\rho)$의 중간체로서 그 자체가 $\Omega(g_1^*\cdots g_\rho^*)$ 위의 대수적 확대체이며, 한편 $\Omega(x_1\cdots x_\rho)$는 $\Kfield^*_{\rho\rho}$ 위의 대수적 확대체이다. 그러므로 $\Kfield^*_{\rho\rho}$는 $\Omega(g_1^*\cdots g_\rho^*)$에 $\Kfield^*_{\rho\rho}$의 유한 개 원소를 첨가하여 얻어지며, 또는 적절히 택한 함수 \emph{하나}를 첨가하여도 얻어진다. 즉,
\[
\Kfield^*_{\rho\rho}
=\Omega(g_1^*\cdots g_\rho^*g_{\rho+1}^*\cdots g_k^*)
=\Omega(g_1^*\cdots g_\rho^*g_0^*).
\]

%
%
전이 원리에 따라 다음을 얻는다.

\emph{정리 II: 모든 체(體) $\Kfield_{n\rho}$는 대수적 랭크 $\rho$의 유리기저를 갖는다. 특히 이 유리기저는 많아야 $\rho+1$개의 함수를 포함하도록 택할 수 있다.}

%
%
$\Kfield_{n\rho}$에서 택한 $\rho$개의 대수적으로 독립인 함수로 이루어진 임의의 계가 $h_1(x)\cdots h_\rho(x)$로 주어졌다고 하자. 이 계는 위의 방법으로 다음 유리기저까지 확장할 수 있다.
\[
h_1(x)\cdots h_\rho(x);\ h_{\rho+1}(x)\cdots h_\sigma(x)
\]
정의 성질에 따라 다음 관계식들이 존재한다.
\[
\begin{array}{ll}
h_i(x)=\chi_i\bigl(g_1(x)\cdots g_k(x)\bigr)&(i=1,2,\ldots,\sigma),\\
g_j(x)=\psi_j\bigl(h_1(x)\cdots h_\sigma(x)\bigr)&(j=1,2,\ldots,k),
\end{array}
\]
여기서 $\chi_i$, $\psi_j$는 각각 그 인수들의 유리함수를 뜻한다. 즉,

임의의 한 유리기저로부터 다른 모든 유리기저가 쌍유리 변환으로 얻어진다. 또한,

%
%
만일
\[
h_1(x)\cdots h_\rho(x)
\]
가 $\Kfield_{n\rho}$에서 택한 $\rho$개의 대수적으로 독립인 함수로 이루어진 임의의 계라면, $\Kfield_{n\rho}$는
\[
\Omega\bigl(h_1(x)\cdots h_\rho(x)\bigr)
\]
위의 (유한) 대수적 확대체이다.

%
%
\section*{\S\ 5. 체(體) $\Kfield_{n\rho}$의 인볼루션형식과 인볼루션기저.}

%
%
\S\ 4에서 구성된 모든 유리기저는 임의로 택한 출발함수들에 의존한다는 단점을 지녔으나, 이제 우리는 \emph{그 체(體)에 의해 유일하게 결정되는 유리기저}가 존재함을 보인다.

%
%
\S\ 4에 따르면 $\Omega(x_1\cdots x_\rho)$는 $\Kfield^*_{\rho\rho}$ 위의 --- 가령 차수가 $i$인 --- 대수적 확대체이며, 다음 등식 때문에
\[
\Kfield^*_{\rho\rho}(x_1\cdots x_\rho)=\Omega(x_1\cdots x_\rho)
\]
$\Kfield^*_{\rho\rho}$에 관하여 대수적인 원소 $x_1\cdots x_\rho$를 첨가함으로써 얻어진다. 따라서 우리는 $\Kfield^*_{\rho\rho}$에 관하여 기약다항식 $\Phi^*(z,u)$를 생각하며, 이 다항식은 다음을 근으로 갖는다.
\[
z=u_1x_1+u_2x_2+\cdots+u_\rho x_\rho=u(x),
\]
그리고 이 다항식은 $z$에 관하여 $i$차이다. 더 나아가 $\Phi^*(z,u)$는 $[z-u(x)]$와 켤레인 양들의 곱이므로 $z,u_1\cdots u_\rho$에 관한 $i$차 동차형식이다; 따라서 $z^i$의 계수로서 $u$에 무관한 $\Kfield^*_{\rho\rho}$의 함수를 포함한다. 이 최고차계수를 나누어 $1$로 만들면 기약다항식 $\Phi^*(z,u)$는 유일하게 결정된다. 이와 같이 정의한 동차형식

%
%
\[
\begin{array}{@{}l@{\quad}l@{}}
(1)& \displaystyle \Phi^*(z,u)=z^i+\sum{}'g^*_{\alpha\alpha_1\ldots\alpha_\rho}(x_1\cdots x_\rho)\cdot
z^\alpha u_1^{\alpha_1}\cdots u_\rho^{\alpha_\rho}\\
& \displaystyle\qquad(\alpha+\alpha_1+\cdots+\alpha_\rho=i)
\end{array}
\]
이 형식을 $\Kfield^*_{\rho\rho}$의 \emph{인볼루션형식}\srcfn{*)}{이 명칭의 이유는 이 절의 끝에서 드러난다; $\sum'$은 $z^i$를 제외한 모든 항에 걸쳐 합을 취한다는 뜻이다.}이라 부르기로 한다; (1)로부터 다음 형식을 얻는다.

%
%
\[
\begin{array}{@{}l@{\quad}l@{}}
(2)& \displaystyle \Phi(z,u)=z^i+\sum{}'g_{\alpha\alpha_1\ldots\alpha_\rho}(x)\cdot
z^\alpha u_1^{\alpha_1}\cdots u_\rho^{\alpha_\rho}\\
& \displaystyle\qquad(\alpha+\alpha_1+\cdots+\alpha_\rho=i)
\end{array}
\]
이는 $\Kfield_{n\rho}$의 한 \emph{인볼루션형식}이다.

%
%
인볼루션형식 $\Phi^*(z,u)$의 계수(係數)들은 $\Kfield^*_{\rho\rho}$의 유리기저를 이루게 된다. 즉 다음 관계가 성립한다.
\[
\begin{array}{@{}l@{\quad}c@{}}
(3)& \Omega\bigl(g^*_{\alpha\alpha_1\ldots\alpha_\rho}(x_1\cdots x_\rho)\bigr)=\Kfield^*_{\rho\rho}\qquad(\alpha+\alpha_1+\cdots+\alpha_\rho=i).
\end{array}
\]
이를 증명하기 위해, $\Phi^*(z,u)$가 $\Omega(g^*_{\alpha\alpha_1\ldots\alpha_\rho}(x))$에 속하는 계수(係數)들을 가지며, $\Kfield^*_{\rho\rho}$에 관하여 기약이므로 이 부분체에 관해서는 더욱 기약임을 주목한다. 따라서 $\Phi^*(z,u)$는 $\Omega(g^*_{\alpha\alpha_1\ldots\alpha_\rho}(x))$에 관한 $u(x)$의 기약방정식을 나타낸다. 여기서 $\Omega(x_1\cdots x_\rho)$는 $\Omega(g^*_{\alpha\alpha_1\ldots\alpha_\rho}(x))$ 위의 차수 $i$인 대수적 확대체임이 따른다. $\Kfield^*_{\rho\rho}$가 $\Omega(g^*_{\alpha\alpha_1\ldots\alpha_\rho}(x))$와 $\Omega(x_1\cdots x_\rho)$ 사이에 놓이므로, 알려진 바와 같이 차수의 동일성에서 체들의 동일성, 곧 관계 (3)이 따른다. 전이 원리에 따라 $\Phi(z,u)$의 계수(係數)들도 이에 대응하여 $\Kfield_{n\rho}$의 유리기저를 이룬다. 따라서 다음 정리를 얻는다.

%
%
\emph{정리 III: 인볼루션형식의 계수(係數)들은 체(體)에 의해 정해지는 유리기저, 곧 인볼루션기저를 이룬다. 체 $\Kfield_{\rho\rho}$에 대해서는 이 기저가 유일하게 정해진다.}\srcfn{*)}{체 $\Kfield_{n\rho}$의 경우 인볼루션형식은 사상에 의존할 수 있다. 미정 계수(係數)를 가진 사상을 사용하면 유일성을 얻을 수 있다.}

%
%
이제 이 사실과 인볼루션의 기하학적 이론 사이의 연관을 드러내는 무리적 해석을 제시한다.

%
%
$\Phi^*(z,u)$를 일차인수로 분해한 것이 어떤 확대체에서 다음과 같이 주어졌다고 하자.
\[
\begin{array}{@{}l@{\quad}l@{}}
(4)& \Phi^*(z,u)=(z-u_1x_1-\cdots-u_\rho x_\rho)(z-u_1x_1^{(1)}-\cdots-u_\rho x_\rho^{(1)})\cdots\\[0.25em]
& \qquad\qquad\quad (z-u_1x_1^{(i-1)}-\cdots-u_\rho x_\rho^{(i-1)}).
\end{array}
\]
(1)과 비교하면 함수들
\[
g^*_{\alpha\alpha_1\ldots\alpha_\rho}(x_1\cdots x_\rho)\qquad(\alpha+\alpha_1+\cdots+\alpha_\rho=i)
\]
은 다음 양들의 열들에 대한 기본대칭함수와 항등적으로 같음을 알 수 있다.
\[
\begin{array}{@{}l@{\quad}cccc@{}}
(5)&x_1&x_2&\cdots&x_\rho\\
&x_1^{(1)}&x_2^{(1)}&\cdots&x_\rho^{(1)}\\
&\vdots&&&\\
&x_1^{(i-1)}&x_2^{(i-1)}&\cdots&x_\rho^{(i-1)}.
\end{array}
\]
그러므로 $\Kfield^*_{\rho\rho}$의 모든 함수는 이 기본대칭함수들의 유리 결합으로서 양들의 열 (5)에 대하여 대칭이며, 거꾸로 양들의 열 (5)의 모든 대칭함수는 $\Kfield^*_{\rho\rho}$에 속한다.

%
%
따라서
\[
f^*(x)=\frac{F^*(x)}{H^*(x)}
\]
가 $\Kfield^*_{\rho\rho}$의 모든 함수를 두루 취하면, 다음과 같은 무한히 많은 관계로 이루어진 계가 성립한다.
\[
\begin{array}{@{}l@{\quad}c@{}}
(6)&\displaystyle \frac{F^*(x)}{H^*(x)}=\frac{F^*(x^{(1)})}{H^*(x^{(1)})}=\cdots=\frac{F^*(x^{(i-1)})}{H^*(x^{(i-1)})};
\end{array}
\]
또는 이를 한데 쓰면 다음과 같다.
\[
F^*(x^{(k)})H^*(x^{(l)})-F^*(x^{(l)})H^*(x^{(k)})=0\qquad(k,l=0,1,\cdots,(i-1));
\]
여기서 $H^*(x^{(k)})$는 항등적으로 영이 아니며, $f^*\ne0$이면 $F^*(x^{(k)})$도 항등적으로 영이 아니다. 이 식들은 각각 $H^*(x)$와 $F^*(x)$에 켤레이며 형식적으로 대칭이다.

%
%
거꾸로, 무한히 많은 관계
\[
\begin{array}{@{}l@{\quad}c@{}}
(7)&F^*(x)H^*(\xi)-H^*(x)F^*(\xi)=0,\qquad H^*(\xi)\ne0
\end{array}
\]
를 만족하는 양들의 열 $\xi$ 각각은 양들의 열 (5)에 속한다.

%
%
실제로 $G^*(x)$가 $\Phi^*(z,u)$의 계수들의 공통분모를 뜻한다고 하자. 그러면 가정 (7)에 따라 $\xi$에 관해서도 다항식인
\[
G^*(\xi)\cdot\Phi^*(z,u;\xi)
\]
는 항등적으로 소멸하지 않는다. 즉 $\Phi^*(z,u;\xi)$의 계수들은 부정형이 되지 않으며, 이 함수는 다음 근을 갖는다.
\[
z=u_1\xi_1+u_2\xi_2+\cdots+u_\rho\xi_\rho.
\]
따라서 (7)에 의해 $\Phi^*(z,u;x)$ 또한 $z=u(\xi)$를 근으로 갖는다. 즉 $\xi$는 양들의 열 (5)에 속한다. (6)에 따라, $\xi$가 켤레인 양들의 열에 관한 방정식계를 만족할 때에도 이에 상응하는 사실이 성립한다:

%
%
\[
F^*(x^{(k)})H^*(\xi)-H^*(x^{(k)})F^*(\xi)=0;\qquad H^*(\xi)\ne0,
\]
즉 양들의 열 (5)는 다음 방정식계의 해 $\xi$들로 정의된다.
\[
F^*(x)H^*(\xi)-H^*(x)F^*(\xi)=0;\qquad H^*(\xi)\ne0,
\]
여기서 $\frac{F^*(x)}{H^*(x)}$는 $\Kfield^*_{\rho\rho}$의 모든 함수를 두루 취한다. 이때 양들의 열 $x$는 어떤 켤레인 열로도 바꿀 수 있다.

%
%
기하학적으로 각 열 $x$를 하나의 점으로 해석하면 이는 다음을 뜻한다. (7)의 해들은 $\rho$차원 선형공간 $\infty^\rho$ 안에서 각각 $i$개의 점으로 이루어진 군들을 나타내며, 그 군의 어느 한 점으로부터도 개별 군이 유일하게 정해진다. 이러한 점들의 군으로 이루어진 계를 \glqq{}\emph{선형공간의 인볼루션}\grqq{}이라 한다.\srcfn{*)}{(7)의 해들 가운데 $F^*(\xi)=0$, $H^*(\xi)=0$이 되는 까닭에 제외된 것들 --- 그리고 인볼루션기저에 대응하는 방정식계에서 제외된 해들도 마찬가지로 --- 을 인볼루션의 기본점이라 한다. 이때 동차화로 생기는 \glqq{}무한원에 놓인\grqq{} 점들도 함께 세어야 한다.}

%
%
이에 대응하여 체 $\Kfield_{n\rho}$에 의해, 선형공간 안에서 곡선, 곡면 등으로 이루어진 하나의 인볼루션이 특징지어진다.

%
%
위에서 본 바에 따라, $\Kfield^*_{\rho\rho}$에 관한 $\Omega(x_1\cdots x_\rho)$의 차수 $i$는 부가조건 $H^*(\xi)\ne0$을 갖는 (7)의 해계의 수 --- 이를 $\Kfield^*_{\rho\rho}$의 해계라 부를 수 있다 --- 와 같다. 그러므로 인볼루션기저의 존재 증명에 쓰인 논법은 다음과 같은 무리적 표현도 허용한다.

\glqq{}$\Lsys^*$가 $\Kfield^*_{\rho\rho}$의 한 부분체이고, 체 $\Lsys^*$의 모든 해계가 $\Kfield^*_{\rho\rho}$의 해계이기도 하면, $\Lsys^*$와 $\Kfield^*_{\rho\rho}$는 동일하다.\grqq{}

%
%
전이 원리에 따라 $\Kfield_{n\rho}$의 부분체 $\Lsys$에도 이에 대응하는 명제가 성립한다.

%
%
\section*{\S\ 6. 체(體) $\Kfield_{n\rho}$의 최소기저.}

%
%
이 절에서는 이미 알려진 정리들을 개관하되, 전이 원리와 유리기저의 존재로 가능해진 확장된 형태로 제시한다.

%
%
\emph{정의 IV: 체(體) $\Kfield_{n\rho}$의 유리기저가 함수의 최소 개수인 $\rho$개만을 포함하고, 따라서 이 함수들이 대수적으로 독립이어야 할 때 이를 최소기저라 한다.}

%
%
뤼로트, 카스텔누오보, 엔리케스의 정리들\srcfn{*)}{J. Lüroth: Beweis eines Satzes über rationale Kurven, Math. Ann. 9 (1875). --- G. Castelnuovo: Sulla razionalità delle involuzioni piane, Math. Ann. 44 (1893). --- F. Enriques: Sopra una involuzione non razionale dello spazio, Rend. Acc. Linc, Vol. 21, 21. Jan. 1912.}로부터 다음 사실들이 따른다:

%
%
I. 모든 체(體) $\Kfield_{n1}$에는 최소기저가 존재하며, 그것은 일차분수변환까지 유일하게 정해지는 그 체의 뤼로트 함수이다.

%
%
II. 체(體) $\Kfield_{n2}$에서는 계수영역 $\Omega$의 대수적 확대를 허용할 때 언제나 최소기저가 존재하고, 일반적으로는 그러한 경우에만 존재한다.

%
%
III. 체(體) $\Kfield_{n\rho}$ ($\rho\ge3$)에는 일반적으로 최소기저가 존재하지 않는다. 최소기저를 갖는 특수한 체 $\Kfield_{n\rho}$의 특성화는 아직 시도되지 않았다.

%
%
E. 슈타이니츠가 체(體) $\Kfield^*_{11}$에 대한 뤼로트 정리에 관하여 제시한 증명\srcfn{**)}{앞의 문헌, \S\ 24.}을 간단히 개략한다:

$\Phi^*(z,u)$를 $\Kfield^*_{11}$의 인볼루션형식이라 하고, $\psi^*(x)$를 $\Phi^*(z,u)$의 계수들 가운데 부정원 $x$를 실제로 포함하는 임의의 하나라 하자. $\Omega(\psi^*(x))$에 대한 $\Omega(x)$의 차수는 $\Kfield^*_{11}$에 대한 $\Omega(x)$의 차수와 같음이 드러나며, 이로부터 두 체가 동일함이 따른다. 또한 $\Omega(\chi^*(x))$가 $\Omega(\psi^*(x))$와 같다고 하자. 그러면 $\chi^*(x)$와 $\psi^*(x)$가 서로 유리적으로 종속되므로, 두 함수가 일차분수적으로 종속됨이 따른다.

%
%
그러므로 전이 원리에 따라 뤼로트 정리는 다음과 같이도 진술할 수 있다:

\emph{체(體) $\Kfield_{n1}$의 최소기저는 인볼루션기저의 임의의 한 함수로 이루어진다. 또한 $\Kfield_{n1}$의 인볼루션기저에 속하는 임의의 두 함수는 서로 일차분수적으로 종속된다.}

%
%
G. 카스텔누오보는 세 함수로 이루어진 유리기저에서 유도되는 체(體) $\Kfield^*_{22}$에 관하여, 인볼루션의 기하학적 이론을 사용하여 사실 II를 증명한다. 계수영역 $\Omega$의 유한 대수적 확대에서도 전이 원리가 유지되므로, 유리기저의 존재에 의해 이 증명은 모든 체 $\Kfield_{n2}$에 대하여 성립한다.

%
%
III에 관해서는 엔리케스가 3차원 공간에서 비유리 인볼루션을 구성했다는 점, 즉 최소기저를 갖지 않는 체(體) $\Kfield_{33}$을 구성했다는 점에 주목해야 한다.

%
%
\section*{\S\ 7. 유리함수의 임의의 계 $\Ssys_{n\rho}$, 선형족 $\Lsys_{n\rho}$ 및 정역(整域) $\Jdom_{n\rho}$. 유리기저의 존재.}

%
%
체(體) $\Kfield_{n\rho}$의 유리기저가 존재한다는 사실로부터 이제 유리함수와 다항함수의 임의의 계에도 유리기저가 존재한다는 결과가 나올 것이다. 우리는 이 계들을 덧셈, 곱셈, 나눗셈이라는 기본 연산을 거쳐 차례로 서로로부터 생겨나는 순서대로 배열한다.

%
%
I. 언제나와 같이 $\Ssys_{n\rho}$는 $n$개의 부정원에 관한, 대수적 랭크가 $\rho$인 유리함수(또는 다항함수)의 임의의 계를 나타낸다. $\Omega$의 원소들도 이 계에 속하는 것으로 센다.

%
%
II. 계 $\Lsys_{n\rho}$가 다음 조건들을 만족할 때 이를 \emph{선형족}이라 한다:

%
%
1. $\Lsys_{n\rho}$는 $f(x)$와 함께 언제나 $c\cdot f(x)$도 포함한다. 여기서 $c$는 $\Omega$의 임의의 원소이다.

%
%
2. $\Lsys_{n\rho}$는 $f(x)$와 $g(x)$와 함께 언제나 $f(x)+g(x)$도 포함한다.

%
%
III. 선형족 $\Jdom_{n\rho}$가 다음의 추가 조건을 만족할 때 이를 \emph{정역(整域)}이라 한다:

%
%
3. $\Jdom_{n\rho}$는 $f(x)$와 $g(x)$와 함께 언제나 $f(x)\cdot g(x)$도 포함한다.

%
%
IV. 정역(整域) $\Kfield_{n\rho}$가 다음의 추가 조건을 만족할 때 이를 \emph{체(體)}라 한다:

%
%
4. $\Kfield_{n\rho}$는 $f(x)$와 $g(x)$ --- 여기서 $g(x)\ne0$ --- 와 함께 언제나 몫 $f(x):g(x)$도 포함한다.

%
%
조건 1부터 4까지는 \S\ 1에서 체에 대하여 제시한 조건들이다.\srcfn{*)}{$f(x)$와 $g(x)$는 0차일 수도 있다. 즉 $\Omega$의 원소일 수도 있다.}

%
%
이제 이 영역들을 차례로 서로로부터 유도한다.

%
%
a) 임의의 계 $\Ssys_{n\rho}$에 다음 조건을 부과하여 이를 선형족 $\Lsys_{n\rho}$로 확대한다:

%
%
$\Lsys_{n\rho}$는 $\Ssys_{n\rho}$와 함께, $\Ssys_{n\rho}$의 유한 개 함수로부터 $\Omega$의 계수를 써서 선형적으로 나타낼 수 있는 모든 함수 $h(x)$를, 그리고 오직 그러한 함수들만을 포함한다:
\[
h(x)=c_1f_1(x)+c_2f_2(x)+\cdots+c_kf_k(x),
\]
여기서 $f_1(x)\cdots f_k(x)$는 $\Ssys_{n\rho}$의 모든 함수를 두루 나타내고, $c_1\cdots c_k$는 $\Omega$의 모든 원소를 두루 나타낸다.

%
%
이는 계의 함수들에 대한 유리 결합을 첨가해도 계의 대수적 랭크가 유지되기 때문이다. $\Lsys_{n\rho}$는 조건 1과 2를 만족하므로 선형족이다. 더 나아가 $\Ssys_{n\rho}$를 포함하는 모든 선형족은 조건 1과 2에 따라 $\Lsys_{n\rho}$도 포함하므로, $\Lsys_{n\rho}$는 $\Ssys_{n\rho}$를 포함하는 최소 선형족이다.

%
%
b) 이에 대응하여 선형족 $\Lsys_{n\rho}$로부터 이를 포함하는 최소 정역(整域) $\Jdom_{n\rho}$을 다음 조건으로 얻는다:

%
%
$\Jdom_{n\rho}$는 $\Lsys_{n\rho}$와 함께, $\Lsys_{n\rho}$의 유한 개 함수 $f_1(x)\cdots f_k(x)$의 $\Omega$-계수 \emph{다항식 결합}으로 나타낼 수 있는 모든 함수 $h(x)$를, 그리고 오직 그러한 함수만을 포함한다. 여기서 $f_1(x)\cdots f_k(x)$는 $\Lsys_{n\rho}$의 모든 함수를 두루 나타낸다.

%
%
c) 끝으로 정역(整域) $\Jdom_{n\rho}$로부터 이를 포함하는 최소 체(體) $\Kfield_{n\rho}$를 다음 조건으로 얻는다. $\Kfield_{n\rho}$는 $\Jdom_{n\rho}$와 함께, $\Jdom_{n\rho}$의 두 함수의 몫으로 나타낼 수 있는 모든 함수 $h(x)$를, 그리고 오직 그러한 함수만을 포함한다.

%
%
이 세 확대를 한 단계로 합치면 다음을 얻는다:

%
%
\emph{모든 계 $\Ssys_{n\rho}$는 다음 조건에 따라 이를 포함하는 최소 체 $\Kfield_{n\rho}$로 확대할 수 있다:}

\emph{$\Kfield_{n\rho}$는 $\Ssys_{n\rho}$와 함께, $\Ssys_{n\rho}$의 유한 개 함수 $f_1(x)\cdots f_k(x)$의 $\Omega$-계수 유리 결합으로 나타낼 수 있는 모든 함수 $h(x)$를, 그리고 오직 그러한 함수만을 포함한다. 여기서 $f_1(x)\cdots f_k(x)$는 $\Ssys_{n\rho}$의 모든 함수를 두루 나타낸다.}

%
%
마지막 사실로부터 임의의 영역에 대한 유리기저의 존재가 곧바로 따른다:

%
%
$\Kfield_{n\rho}$를 $\Ssys_{n\rho}$를 포함하는 최소 체라 하고, $\Kfield_{n\rho}$의 한 유리기저가 다음으로 주어진다고 하자:
\[
\begin{array}{@{}l@{\quad}c@{}}
(1)&h_1(x),\ h_2(x),\ldots,h_\sigma(x).
\end{array}
\]

%
%
$\Kfield_{n\rho}$의 정의에 따라 다음 관계들이 존재한다:
\[
\begin{array}{@{}l@{\quad}c@{}}
(2)&h_1(x)=\gamma_1(g_1(x)\cdots g_\lambda(x));\quad\ldots;\quad h_\sigma(x)=\gamma_\sigma(g_\mu(x)\cdots g_\nu(x)),
\end{array}
\]

%
%
여기서 $\gamma_1\cdots\gamma_\sigma$는 각자의 인수에 관한 유리함수이고,
\[
\begin{array}{@{}l@{\quad}c@{}}
(3)&g_1(x)\cdots g_\lambda(x),\ldots,g_\mu(x)\cdots g_\nu(x)
\end{array}
\]
는 모두 $\Ssys_{n\rho}$에 속한다. 관계 (2)에 따라 (3)도 (1)과 더불어 $\Kfield_{n\rho}$의 유리기저를 이룬다. 또한 (3)의 함수들이 $\Ssys_{n\rho}$에 속하므로, 이는 동시에 $\Ssys_{n\rho}$의 유리기저이다. 따라서 다음을 얻는다.

%
%
\emph{정리 IV: 유리함수(또는 다항함수)의 임의의 계 $\Ssys_{n\rho}$에는 그 계의 유한 개 함수로 이루어지고 대수적 랭크가 $\rho$인 유리기저가 존재한다.}

%
%
이제 특별히 $\Lsys_{n\rho}$가 선형족이고,
\[
\begin{array}{@{}l@{\quad}c@{}}
(4)&g_1(x)\cdots g_\rho(x),\ g_{\rho+1}(x)\cdots g_\nu(x)
\end{array}
\]
가 $\Lsys_{n\rho}$의 유리기저라고 하자. 이를테면

%
%
\[
g_1(x)\cdots g_\rho(x)
\]
가 대수적으로 독립이라고 하자. 이제 다음과 같이 두면

%
%
\[
g(x)=c_1g_{\rho+1}(x)+c_2g_{\rho+2}(x)+\cdots+c_{\nu-\rho}g_\nu(x),
\]
$c_i$들을 $\Omega$의 수들로 다음이 성립하도록 선택할 수 있다:

%
%
\[
\begin{array}{@{}l@{\quad}c@{}}
(5)&g_1(x)\cdots g_\rho(x),\ g(x)
\end{array}
\]
가 이를 포함하는 최소 체 $\Kfield_{n\rho}$의 유리기저가 된다. 그런데 $\Lsys_{n\rho}$는 선형족이므로 $g(x)$는 $\Lsys_{n\rho}$에 속하고, (5)는 $\Lsys_{n\rho}$의 유리기저를 이룬다. 즉,

%
%
\emph{정리 V: 유리함수(또는 다항함수)의 모든 선형족 $\Lsys_{n\rho}$의 유리기저, 따라서 특히 모든 정역 $\Jdom_{n\rho}$과 모든 체 $\Kfield_{n\rho}$의 유리기저는 많아야 $\rho+1$개의 함수(그리고 적어도 $\rho$개)를 포함하도록 특별히 선택할 수 있다.}

%
%
따라서 \S\ 4에서 체에 대하여 얻은 유리기저 함수 수의 제한은 그 체가 선형족이라는 성질에 근거한다. 반면 최소기저가 존재하는 경우 $\rho$개의 함수로 제한되는 것은 체에 관한 모든 가정을 사용한다.

%
%
끝으로 \S\ 3 (4)에 따라 계와 그 계를 포함하는 최소 계의 모든 특징적 성질은 사상에서도 보존된다는 점을 언급한다.

%
%
\section*{\S\ 8. 다항식으로 이루어진 정역(整域) $\Jdom_{n\rho}$의 인볼루션기저.}

%
%
이하의 절들에서는 원소가 다항함수(다항식)인 계 $\Ssys_{n\rho}$를 다룬다. 우선 다항식으로 이루어진 정역(整域) $\Jdom_{n\rho}$를 다룬다. 이 정역 $\Jdom_{n\rho}$에 대해서는 먼저 나눗셈 관계의 개념을 정한다.

%
%
$F(x),G(x)$를 $\Jdom_{n\rho}$의 두 다항식이라 하고, $L(x)$를 이 다항식들의 공약수라 하여 다음이 성립한다고 하자:
\[
F(x)=L(x)\cdot F_1(x);\qquad G(x)=L(x)\cdot G_1(x).
\]
세 다항식 $L(x),F_1(x),G_1(x)$가 $\Jdom_{n\rho}$에 속할 때, 그리고 그때에만 $L(x)$를 \emph{$\Jdom_{n\rho}$에서의 공약수}라 한다.

%
%
$\Jdom_{n\rho}$의 두 다항식이 $\Jdom_{n\rho}$에서 공약수를 갖지 않을 때 이들을 \emph{$\Jdom_{n\rho}$에서 서로소}라 한다.

%
%
$\Kfield_{n\rho}$를 $\Jdom_{n\rho}$를 포함하는 최소 체(體)라 하자. 그러면 정의에 따라 $\Kfield_{n\rho}$의 모든 함수는 다음과 같이 나타낼 수 있다:
\[
f(x)=F_1(x):F_2(x),
\]
여기서 $F_1(x),F_2(x)$는 $\Jdom_{n\rho}$에 속하며 $\Jdom_{n\rho}$에서 서로소이다. 이에 따라 $\Kfield_{n\rho}$의 여러 함수도 $\Jdom_{n\rho}$에서 공통분모로 통분할 수 있다:

%
%
\[
f_1(x)=\frac{F_1(x)}{F(x)}\cdots f_k(x)=\frac{F_k(x)}{F(x)}.
\]
$\Jdom_{n\rho}$의 다항식 $F(x),F_1(x),\ldots,F_k(x)$가 $\Jdom_{n\rho}$에서 서로소일 때, 그리고 그때에만 $F(x)$를 \emph{$\Jdom_{n\rho}$에서의 최소공통분모}라 한다.

%
%
$\Jdom_{n\rho}$에서는 일반적으로 기약다항식으로의 유일분해 법칙이 성립하지 않으므로, 이러한 표현은 여러 방식으로 가능할 수 있다.

%
%
이제 인볼루션형식과 인볼루션기저의 개념을 정역에 옮김으로써, $\Kfield_{n\rho}$의 인볼루션기저가 $\Jdom_{n\rho}$에 대하여 갖는 의미를 조사한다. 다음을 두자:
\[
\begin{gathered}
\Phi(z,u)=z^i+\sum{}' g_{\alpha\alpha_1\ldots\alpha_\rho}(x)z^\alpha u_1^{\alpha_1}\cdots u_\rho^{\alpha_\rho}\\
(\alpha+\alpha_1+\cdots+\alpha_\rho=i)
\end{gathered}
\]
이를 $\Kfield_{n\rho}$의 인볼루션형식이라 하고, $G(x)$를 $\Phi(z,u)$의 계수들에 대한 $\Jdom_{n\rho}$에서의 최소공통분모라 하자. $G(x)$를 곱하면 $\Phi(z,u)$는 형식 $\Psi(z,u)$로 바뀌며, 그 계수들은 $\Jdom_{n\rho}$의 다항식이고 $\Jdom_{n\rho}$에서 서로소이다:

%
%
\[
\begin{gathered}
G(x)\cdot\Phi(z,u)=\Psi(z,u)=G(x)z^i+\sum{}'G_{\alpha\alpha_1\ldots\alpha_\rho}(x)z^\alpha u_1^{\alpha_1}\cdots u_\rho^{\alpha_\rho}\\
(\alpha+\alpha_1+\cdots+\alpha_\rho=i).
\end{gathered}
\]

%
%
\emph{$\Kfield_{n\rho}$의 인볼루션형식 $\Phi(z,u)$에 $\Jdom_{n\rho}$의 한 다항식을 곱하여 얻어지는 모든 형식 $\Psi(z,u)$ 가운데, $\Psi(z,u)$의 계수들이 $\Jdom_{n\rho}$의 다항식이고 $\Jdom_{n\rho}$에서 서로소가 되도록 얻어진 형식을 $\Jdom_{n\rho}$의 인볼루션형식이라 하고, $\Psi(z,u)$의 계수들을 그에 대응하는 인볼루션기저라 한다.}

%
%
우선 $\Kfield_{n\rho}$의 인볼루션기저가 갖는 의미로부터, $\Jdom_{n\rho}$의 모든 인볼루션기저가 동시에 유리기저임은 분명하다. 더 조사하기 위하여 사상영역 $\Jdom^*_{\rho\rho}$를 생각하자. 이를 포함하는 최소 체 $\Kfield^*_{\rho\rho}$에 대하여, 근
\[
z=u_1x_1+\cdots+u_\rho x_\rho
\]
을 갖는 기약방정식은 $\Phi^*(z,u)=0$으로 주어진다.

%
%
\S\ 5에 따르면
\[
\begin{gathered}
\Phi^*(z,u)=z^i+\sum{}'g^*_{\alpha\alpha_1\ldots\alpha_\rho}(x)z^\alpha u_1^{\alpha_1}\cdots u_\rho^{\alpha_\rho}\\
(\alpha+\alpha_1+\cdots+\alpha_\rho=i)
\end{gathered}
\]
의 계수들은 $\Kfield^*_{\rho\rho}$에 관하여 서로 켤레인 양들의 열
\[
x,\ x^{(1)}\cdots x^{(i-1)}
\]
의 기본대칭함수들을 나타낸다.

%
%
$F^*(x)$를 $\Kfield^*_{\rho\rho}$의 임의의 다항식이라 하면,
\[
F^*(x)=\frac1i\left(F^*(x)+F^*(x^{(1)})+\cdots+F^*(x^{(i-1)})\right),
\]
이므로 $F^*(x)$는 이 양들의 열에 대한 대칭다항식이다. 따라서 이는 기본대칭함수들, 즉 함수
\[
g^*_{\alpha\alpha_1\ldots\alpha_\rho}(x)
\]
의 다항식 결합으로 나타낼 수 있다.

%
%
이제
\[
\begin{gathered}
G^*(x)\cdot\Phi^*(z,u)=\Psi^*(z,u)=G^*(x)z^i+
\sum{}'G^*_{\alpha\alpha_1\ldots\alpha_\rho}(x)z^\alpha u_1^{\alpha_1}\cdots u_\rho^{\alpha_\rho}\\
(\alpha+\alpha_1+\cdots+\alpha_\rho=i)
\end{gathered}
\]
가 $\Jdom^*_{\rho\rho}$의 인볼루션형식이면, 이에 따라 $\Kfield^*_{\rho\rho}$의 모든 다항식, 특히 $\Jdom^*_{\rho\rho}$의 모든 다항식은 몫
\[
G^*_{\alpha\alpha_1\ldots\alpha_\rho}(x):G^*(x)
\]
의 다항식 결합이 된다.

%
%
전이 원리를 고려하면 이로부터 다음이 얻어진다.

%
%
\emph{정리 VI: 다항식으로 이루어진 모든 정역(整域) $\Jdom_{n\rho}$에는 적어도 하나의 특별한 유리기저, 곧 인볼루션기저가 존재한다. 이 기저는 $\Jdom_{n\rho}$의 모든 다항식을 인볼루션기저의 함수들로 유리적으로 나타낼 때 분모에는 하나의 고정된 함수(그에 대응하는 인볼루션형식의 최고차항의 계수)의 거듭제곱만이 나타나도록 하는 성질을 갖는다.}

%
%
\section*{\S\ 9. 다항식으로 이루어진 상대정역(相對整域) $\Gdom_{n\rho}$.}

%
%
이하에서는 정역과 체 사이의 중간 단계를 이루는, 다항식으로 된 특수한 정역들을 고찰한다.

%
%
다항식으로 이루어진 정역 $\Gdom_{n\rho}$가 다음 조건을 만족할 때 이를 
\glqq{}상대정역\grqq{}이라 한다:

%
%
\emph{4a) $\Gdom_{n\rho}$는 $F(x)$와 $G(x)$ --- 여기서 $G(x)\ne0$ --- 와 함께, 몫 $F(x):G(x)$가 부정원들에 관한 정식(整式), 즉 다항식일 때, 그리고 그때에만 언제나 이 몫도 포함한다.}

%
%
우선 상대정역이 유리함수체 $\Kfield_{n\sigma}$ ($\sigma\ge \rho$)의 다항식 전체와 동일함을 증명해야 한다.

%
%
$\Kfield_{n\rho}$를 $\Gdom_{n\rho}$를 포함하는 최소 체라 하자. $\Kfield_{n\rho}$의 모든 함수 $f(x)$는 $\Gdom_{n\rho}$의 두 다항식의 몫으로 나타낼 수 있다:
\[
f(x)=F_1(x):F_2(x),
\]
여기서는 0차 다항식도 허용한다. $f(x)$가 다항식이면 조건 4a)에 따라 $\Gdom_{n\rho}$에 속한다. 또한 $\Gdom_{n\rho}$는 그 밖의 다항식을 포함할 수 없으므로, \emph{모든 상대정역 $\Gdom_{n\rho}$는 이를 포함하는 최소 체의 다항식 전체와 동일하다.}

%
%
거꾸로 $\Kfield_{n\sigma}$를 임의의 유리함수체(즉 주어진 어떤 계를 포함하는 최소 체일 필요는 없는 체)라 하고, $\mathfrak J$를 $\Kfield_{n\sigma}$에 포함되는 다항식 전체라 하자. $\mathfrak J$는 정역의 조건 1, 2, 3과 조건 4a)를 만족하므로 상대정역이다. 즉 \emph{체 $\Kfield_{n\sigma}$에 포함되는 다항식 전체는 상대정역 $\Gdom_{n\rho}$ ($\rho\le\sigma$)를 이룬다.}\srcfn{*)}{$\rho=0$일 수도 있다. 즉 $\Kfield_{n\sigma}$의 다항식 전체가 계수영역 $\Omega$의 원소들로 이루어질 수도 있다.}

%
%
더 나아가 상대정역이 힐베르트가 도입한 
\glqq{}상대적으로 정수적인 함수들의 정역\grqq{}과 동일함을 보인다.\srcfn{**)}{D. Hilbert: Mathematische Probleme. Vortrag Paris 1900. Problem 14 (Göttinger Nachrichten 1900).}

%
%
다음을
\[
\begin{array}{@{}l@{\quad}c@{}}
(1)& G_1(x)\cdots G_k(x)
\end{array}
\]
$\Gdom_{n\rho}$의 유리기저라 하자. 조건 1, 2, 3, 4a)에 따라 다항식 (1)의 유리 결합 가운데 부정원들에 관한 정식, 즉 다항식이 되는 모든 것은 $\Gdom_{n\rho}$에 속한다. 거꾸로 $\Gdom_{n\rho}$의 모든 다항식은 유리기저의 개념에 따라 다항식 (1)로 유리적으로 나타낼 수 있다. 따라서 $\Gdom_{n\rho}$는 \emph{다항식 (1)의 유리 결합 가운데 부정원들에 관한 정식, 즉 다항식인 것들의 전체와 동일하다.} 이로써 특수한 유리기저에서 출발하는 힐베르트의 정의를 얻는다. 반면 우리의 정의는 이 영역의 추상적 성질들만을 사용한다.

%
%
상대정역의 경우에는 \S\ 8에서 일반 정역에 대하여 제시한 공약수의 정의가 간단해진다:

%
%
$F(x),G(x)$를 $\Gdom_{n\rho}$의 두 다항식이라 하고, $L(x)$를 이 다항식들의 공약수라 하자. $L(x)$가 $\Gdom_{n\rho}$에 속할 때, 그리고 그때에만 $L(x)$를 \emph{$\Gdom_{n\rho}$에서의 공약수}라 한다.

%
%
몫 $F(x):L(x)$와 $G(x):L(x)$가 속한다는 것은 이때 조건 4a)에서 따르기 때문이다.

%
%
상대정역의 또 하나의 본질적 성질은 
\glqq{}$\Gdom_{n\rho}$에 관하여 대수적으로 정수적인 것\grqq{}의 개념을, 대수적 수체나 체 $\Omega(x_1\cdots x_n)$에 관하여 알려진 것과 똑같이 임의의 방정식을 통해 정의할 수 있다는 점이다.

%
%
\emph{정의 V: 한 양 $\xi$가 최고차항의 계수는 1이고 나머지 계수들은 $\Gdom_{n\rho}$의 다항식인 어떤 방정식을 만족할 때, $\xi$를 $\Gdom_{n\rho}$에 관하여 대수적으로 정수적이라고, 또는 상대대수적으로 정수적이라고도 한다.}

%
%
상대대수적으로 정수적인 양 $\xi$가 이를테면 방정식
\[
L(z)=z^\lambda+G_1(x)z^{\lambda-1}+\cdots+G_\lambda(x)=0
\]
을 만족한다고 하자. $L(z)$는 이를 포함하는 최소 체 $\Kfield_{n\rho}$에 관하여 기약다항식

%
%
\[
M(z)=z^\mu+H_1(x)z^{\mu-1}+\cdots+H_\mu(x)
\]
로 나누어진다. 그런데 이 나눗셈 관계로부터 잘 알려진 가우스 정리의 확장\srcfn{*)}{Vgl. etwa J. König: Einleitung in die allgemeine Theorie der algebraischen Größen (Leipzig, B. G. Teubner, 1903), Kap. IX, \S\ 2 oder Weber (kleine Ausgabe), \S\ 20.}에 따라 유리함수 $H_1(x)\cdots H_\mu(x)$가 $\Gdom_{n\rho}$의 다항식이며, 따라서 $\Gdom_{n\rho}$에 속함이 따른다. 이로써 상대대수적으로 정수적인 양 $\xi$의 정의를 기약방정식으로도 얻었다. 특히 $\Gdom_{n\rho}$의 모든 다항식 $F(x)$는 방정식 $z-F(x)=0$을 만족하므로 $\Gdom_{n\rho}$에 관하여 대수적으로 정수적이다.

%
%
또한 다항식으로 이루어진 정역 $\Jdom_{n\rho}$에 대해서는, 이를 포함하는 최소 체와 마찬가지로 이를 포함하는 최소 상대정역 $\Gdom_{n\rho}$가 다음 조건으로 정의된다는 점을 언급한다:

%
%
$\Gdom_{n\rho}$는 $\Jdom_{n\rho}$와 함께, $\Jdom_{n\rho}$의 두 다항식의 몫으로 나타낼 수 있는 모든 다항식 $H(x)$를, 그리고 오직 그러한 다항식만을 포함한다.

%
%
이 정의로부터 더 나아가 다음이 따른다. $\Jdom_{n\rho}$의 모든 인볼루션형식과 인볼루션기저로부터, 기저의 모든 다항식이 $\Gdom_{n\rho}$에서 갖는 공약수를 하나까지 분리해 냄으로써 $\Gdom_{n\rho}$의 대응하는 인볼루션형식과 인볼루션기저가 생긴다.

%
%
$\Jdom_{n\rho}$와 $\Gdom_{n\rho}$는 이를 포함하는 동일한 최소 체 $\Kfield_{n\rho}$를 갖기 때문이다. 또한 $\Jdom_{n\rho}$의 모든 인볼루션형식은 $\Kfield_{n\rho}$의 한 인볼루션형식에 $\Jdom_{n\rho}$의 다항식을 곱하여 생긴다. 따라서 그 계수들은 $\Gdom_{n\rho}$에 속하지만, $\Gdom_{n\rho}$에서 공약수를 가질 수 있다.\srcfn{*)}{예를 들어 $x_1^2$과 $x_1x_2$의 모든 다항식 결합으로 이루어진 정역 $\Jdom_{22}$에 대해서는 다음 인볼루션형식이 얻어진다:
\[
\Psi(z,u)=x_1^2\cdot z^2-x_1^4u_1^2-2x_1^3x_2u_1u_2-x_1^2x_2^2u_2^2.
\]
$x_1^2$은 $\Jdom_{22}$에, 따라서 더더욱 이를 포함하는 최소 상대정역 $\Gdom_{22}$에 속하고, 그러므로 $\Gdom_{22}$에서 공약수이다. 이에 따라 $\Gdom_{22}$의 인볼루션형식으로 다음이 얻어진다:
\[
\Psi(z,u)=z^2-x_1^2u_1^2-2x_1x_2u_1u_2-x_2^2u_2^2.
\]}

%
%
끝으로 상대정역 $\Gdom_{n\rho}$의 \emph{사상영역}은 \S\ 7에 따라 다시 다항식으로 이루어진 정역 $\Jdom^*_{\rho\rho}$이지만, 반드시 상대정역일 필요는 없다는 점을 언급한다. 실제로 $F(x)$와 $G(x)$가 $\Gdom_{n\rho}$의 두 다항식이고 그 몫이 다항식이 아니어서 $\Gdom_{n\rho}$에 속하지 않으면, 몫 $F^*(x):G^*(x)$도 $\Jdom^*_{\rho\rho}$에 속하지 않는다. 그러나 몇몇 부정원의 특수화로 생기는 이 몫은 얼마든지 다항식일 수 있으며, 그러면 $\Jdom^*_{\rho\rho}$에서는 조건 4a)가 성립하지 않는다.\srcfn{**)}{예를 들어 $\Gdom_{32}$가 $F=x_1^2x_2$, $G=x_1^2(1+x_3)$의 유리 결합인 모든 다항식으로 이루어져 있다고 하자. 허용되는 특수화 $x_3=0$은 $F^*:G^*=x_2$를 주지만, $F:G$는 다항식이 아니다. 따라서 $\Jdom^*_{22}$는 상대정역이 아니다.}

%
%
\section*{\S\ 10. 제1종 상대정역과 그 정기저(整基底).}

%
%
\S\ 9 (정의 V)에서 도입한 
\glqq{}상대대수적으로 정수적인 것\grqq{}의 개념을 이용하면 유한 정역인 상대정역의 한 부류를 제시할 수 있고, 따라서 D. 힐베르트가 제기한 문제(Mathematische Probleme, Probl. 14)에 적어도 이 부류에 대해서는 긍정적으로 답할 수 있다. 다음 정의를 둔다.

%
%
\emph{정의 VI: $\Gdom_{n\rho}$의 유한 개 다항식으로 된 집합을 사용하여 $\Gdom_{n\rho}$의 모든 다항식을 이 유한 개 다항식의 $\Omega$-계수 다항식 결합으로 나타낼 수 있을 때, 이 유한 개 다항식을 정기저라 한다. 다항식으로 이루어진 정역이 정기저를 가지면 이를 유한 정역이라 한다.}

%
%
인볼루션기저가 정기저임을 보이게 될 상대정역의 부류를 다음 정의에 따라 제1종 상대정역이라 부른다.

%
%
\emph{정의 VII: 상대정역 $\Gdom_{n\rho}$가 사상영역으로 상대정역 $\Gdom^*_{\rho\rho}$를 가지며, $x_1\cdots x_\rho$의 임의의 모든 다항식($\Omega$의 계수를 갖는)이 $\Gdom^*_{\rho\rho}$에 대하여 대수적으로 정수적으로 종속될 때, $\Gdom_{n\rho}$를 제1종이라 한다.}

%
%
실제로 이 정의에 따르면 부정원 $x_1\cdots x_\rho$와 따라서 선형형식 $u(x)=u_1x_1+\cdots+u_\rho x_\rho$도 $\Gdom^*_{\rho\rho}$에 대하여 대수적으로 정수적으로 종속된다. 그러므로 근 $z=u(x)$를 갖는 기약다항식 --- 이를 포함하는 최소 체 $\Kfield^*_{\rho\rho}$의 인볼루션형식 --- 은 최고차항의 계수가 1이고 나머지 계수들은 $\Gdom^*_{\rho\rho}$의 다항식이다. 따라서 이는 동시에 $\Gdom^*_{\rho\rho}$의 인볼루션형식이며, $\Gdom^*_{\rho\rho}$에는 다른 인볼루션형식이 존재할 수 없다. 전이 원리에 따라 $\Gdom_{n\rho}$도 최고차항의 계수가 1인 인볼루션형식을 갖는다. 정리 VI에 따라 $\Gdom_{n\rho}$의 모든 다항식을 그에 대응하는 인볼루션기저로 유리적으로 나타낼 때 분모에는 이 최고차항의 계수만 들어가므로, 인볼루션기저는 정기저가 된다. 즉,

%
%
\emph{정리 VII. 모든 제1종 상대정역은 유한 정역이다. 정기저는 그 상대정역을 특징짓는 사상영역에서 얻은 인볼루션기저로 주어진다. 특히 제1종 상대정역 $\Gdom_{\rho\rho}$의 정기저는 $\Gdom_{\rho\rho}$의 유일하게 정의되는 인볼루션기저로 주어진다.}

%
%
제1종 상대정역에 대한 한 판정법을 더 제시한다. $\Gdom_{n\rho}$가 제1종이고 그 정기저가 다음으로 주어진다고 하자:
\[
\begin{array}{@{}l@{\quad}c@{}}
(1)& F_1(x),F_2(x),\ldots,F_k(x).
\end{array}
\]

%
%
그러면 힐베르트의 보조정리\srcfn{*)}{D. Hilbert: Über die vollen Invariantensysteme, Math. Ann. 42 (1893), \S\ 1. (거기서는 동차다항식에 대해서만 진술된 보조정리이지만 비동차다항식에도 똑같이 성립한다.)}에 따라 $\Gdom_{n\rho}$에는 대수적으로 독립인 $\rho$개의 다항식
\[
\begin{array}{@{}l@{\quad}c@{}}
(2)& G_1(x),G_2(x),\ldots,G_\rho(x)
\end{array}
\]
가 존재한다.

%
%
다항식 (1), 따라서 $\Gdom_{n\rho}$의 모든 다항식은 다항식 (2)에 대하여 대수적으로 정수적으로 종속된다. 사상영역 $\Gdom^*_{\rho\rho}$의 모든 다항식에 대해서도, 그에 대응하는 다항식
\[
\begin{array}{@{}l@{\quad}c@{}}
(3)& G_1^*(x),G_2^*(x),\ldots,G_\rho^*(x).
\end{array}
\]
에 관하여 같은 사실이 성립한다.

%
%
따라서 정의 VII에 의해 $x_1\rcdots x_\rho$의 임의의 모든 다항식도 다항식 (3)에 대하여 대수적으로 정수적으로 종속된다.

%
%
이제 거꾸로 임의의 상대정역 $\Gdom_{n\rho}$의 한 사상영역 $\Jdom^*_{\rho\rho}$가 다항식 (3)을 $\rho$개 포함하고, $x_1\rcdots x_\rho$의 모든 다항식이 이들에 대하여 대수적으로 정수적으로 종속된다고 하자. 그러면 이로부터 $\Jdom^*_{\rho\rho}$가 상대정역 $\Gdom^*_{\rho\rho}$임이 나오고, 더 나아가 $\Gdom^*_{\rho\rho}$와 따라서 $\Gdom_{n\rho}$가 제1종임을 보이겠다.

%
%
실제로 $F^*(x)$를 $\Jdom^*_{\rho\rho}$를 포함하는 최소 체 $\Kfield^*_{\rho\rho}$의 한 다항식이라 하자. 가정에 따라 이는 다항식 (3)에 대하여 대수적으로 정수적으로 종속된다. 그러면 $\Kfield_{n\rho}$의 대응하는 함수는 $\Gdom_{n\rho}$의 $\rho$개 다항식에 대하여 대수적으로 정수적으로 종속되므로 다항식 $F(x)$이고, 따라서 $\Gdom_{n\rho}$에 속한다. 그러므로 $F^*(x)$는 $\Jdom^*_{\rho\rho}$에 속한다. 따라서 사상영역 $\Jdom^*_{\rho\rho}$는 $\Kfield^*_{\rho\rho}$의 모든 다항식을 포함하여 상대정역 $\Gdom^*_{\rho\rho}$이다. 그런데 $\Gdom^*_{\rho\rho}$와 따라서 $\Gdom_{n\rho}$에 대해서는 가정과 정의 V 및 VII로부터 이들이 제1종임이 곧바로 따른다. 즉,

%
%
\emph{제1종 상대정역 $\Gdom_{n\rho}$는 한 사상영역 $\Jdom^*_{\rho\rho}$에 $\rho$개의 다항식이 존재하여 $x_1\rcdots x_\rho$의 임의의 모든 다항식이 이들에 대하여 대수적으로 정수적으로 종속된다는 성질로 특징지어진다. 이 가정으로부터 사상영역은 상대정역임이 나온다.}\srcfn{*)}{예를 들어 \S\ 9의 끝 주석에서 언급한 영역 $\Gdom_{32}$에서는 허용되는 특수화 $x_1=1$이 두 다항식 $F^*=x_2$, $G^*=1+x_3$을 주며, $x_2,x_3$의 임의의 모든 다항식은 이들에 대하여 대수적으로 정수적으로 종속된다. 따라서 $\Gdom_{32}$는 제1종이고, 이때 $F(x)$와 $G(x)$가 정기저이다.}

%
%
\section*{\S\ 11. 대수적 정수적 종속에 관한 보조정리.}

%
%
이제 $x_1\rcdots x_\rho$의 임의의 모든 다항식이 $x_1\rcdots x_\rho$의 $\rho$개 다항식
\[
\begin{array}{@{}l@{\quad}c@{}}
(1)& G_1^*(x)\rcdots G_\rho^*(x)
\end{array}
\]
에 대하여 대수적으로 정수적으로 종속되기 위한 판정법을 제시한다.

%
%
다항식 (1)이 특별히 동차형식들이라면, (1)을 0으로 만드는 모든 특수한 값의 계에 대하여 대수적으로 정수적으로 종속되는 모든 형식도 동시에 0이 된다. 특히 $x_1\rcdots x_\rho$도 모두 0이 된다.

%
%
따라서 형식 (1)은
\[
x_1=0,\ldots,x_\rho=0
\]
이외의 어떤 값의 계에서도 동시에 0이 될 수 없다. 그러므로 그 종결식은 0이 아니어야 한다.

%
%
거꾸로 이 조건은 충분하기도 하며, 다음 보조정리에 따라 비동차다항식으로 확장할 수 있다.

%
%
\emph{보조정리: $x_1\rcdots x_\rho$의 임의의 모든 다항식이 $x_1\rcdots x_\rho$의 $\rho$개 다항식
\[
\tag{1} G_1^*(x)\rcdots G_\rho^*(x)
\]
에 대하여 대수적으로 정수적으로 종속되기 위한 충분조건은, 다항식 (1)의 최고차항들}

%
%
\[
H_1^*(x)\rcdots H_\rho^*(x)
\]
\emph{의 $x_1\rcdots x_\rho$에 관한 종결식이 0이 아닌 것이다:}
\[
\begin{array}{@{}l@{\quad}c@{}}
(2)& \displaystyle R_0=\left[\begin{array}{c} H_1^*\rcdots H_\rho^*\\ x_1\rcdots x_\rho\end{array}\right]\ne0.
\end{array}
\]
\emph{동차형식 (1)에 대해서는 조건 (2)가 필요조건이다.}\srcfn{*)}{종결식 이론에 관해서는 F. Mertens, Zur Theorie der Elimination (I. u. II. Teil), Sitzungsber. d. k. Ak. d. Wiss. Wien, Math.-naturw. Kl. Abt. IIa, Bd. 108 (1899) 참조. (J. König, Theorie d. algebr. Größen, Kap. VI에 일부 재수록.)}

%
%
증명을 위하여 다음 다항식들을 만든다:
\[
\begin{array}{@{}l@{\quad}c@{}}
(3)& \begin{gathered}
\bar G_1^*(x)=G_1^*(x)-\lambda_1,\ldots,\bar G_\rho^*(x)=G_\rho^*(x)-\lambda_\rho,\\
\bar\Phi^*(x)=\Phi^*(x)-\mu,
\end{gathered}
\end{array}
\]

%
%
여기서 $\lambda_1\rcdots\lambda_\rho,\mu$는 부정원이고, $\Phi^*(x)$는 $x_1\rcdots x_\rho$의 임의의 다항식이다. 다항식 (3)의 $x_1\rcdots x_\rho,1$에 관한 종결식\srcfn{**)}{즉 다른 부정원 $x_0$을 사용하여 다항식 (3)을 $x$에 관한 차수가 유지되도록 동차화했을 때 $x_1\rcdots x_\rho x_0$에 관한 동차 종결식이다.}이 다음으로 주어진다고 하자:
\[
\begin{array}{@{}l@{\quad}c@{}}
(4)& \displaystyle R(\lambda,\mu)=\left[\begin{array}{cc} \bar G_1^*\rcdots \bar G_\rho^*\bar\Phi^*\\ x_1\rcdots x_\rho\quad 1\end{array}\right];
\end{array}
\]

%
%
이는 $\lambda_1\rcdots\lambda_\rho,\mu$의 다항함수이다. 이제 F. 메르텐스에 따르면 (4)를 $\mu$의 거듭제곱으로 전개했을 때 최고차항의 계수를 명시적으로 제시할 수 있다.\srcfn{***)}{A. a. O. \S\ 3 (Beweis von Satz I) 및 명시적으로 \S\ 6 (Bestimmung des Inbegriffs aller Glieder der Resultante $R$, welche ein gegebenes Potenzprodukt $a_{\lambda\lambda}^{p_\lambda}a_{\mu\mu}^{p_\mu}\cdots a_{\sigma\sigma}^{p_\sigma}$ als Faktor enthalten). 여기서는 $p_\lambda=0$, $p_\mu=0,\ldots,p_\sigma=\tau$, $a_{\sigma\sigma}=(-\mu)$로 둔다. König, Kap. VI, \S\ 9에 재수록.} 다음을 얻는다:
\[
\begin{array}{@{}l@{\quad}c@{}}
(5)& (-1)^\tau\cdot R(\lambda,\mu)=R_0^\sigma\cdot\mu^\tau+A_1(\lambda)\mu^{\tau-1}+\cdots+A_\tau(\lambda),
\end{array}
\]

%
%
여기서 $A_1(\lambda)\rcdots A_\tau(\lambda)$는 $\lambda_1\rcdots\lambda_\rho$와 다항식 (3)의 계수들에 관한 정수계수 다항함수이다. 이 전개는 우선 가정 (2) $R_0\ne0$ 아래에서 $R(\lambda,\mu)$도 항등적으로 0일 수 없음을 보여 준다. 또한 $x_1\rcdots x_\rho,\lambda_1\rcdots\lambda_\rho\mu$에 관한 항등식
\[
R(\lambda,\mu)\equiv0\pmod{\bar G_1^*(x),\ldots,\bar G_\rho^*(x),\bar\Phi^*(x)}
\]

%
%
은 치환
\[
\lambda_i=G_i^*(x),\qquad \mu=\Phi^*(x)
\]
에 따라 다음으로 바뀐다:

%
%
\[
\begin{array}{@{}l@{\quad}l@{}}
(6)& \displaystyle R(G_i^*(x),\Phi^*(x))=R_0^\sigma\cdot\Phi^{*\tau}(x)+A_1(G_i^*(x))\Phi^{*\tau-1}(x)+\cdots\\
& \displaystyle\qquad +A_\tau(G_i(x))=0,
\end{array}
\]
이로써 보조정리가 증명되었다. 실제로 가정에 따라 $R_0$은 $\Omega$의 한 수이다.

%
%
\section*{\S\ 12. 다항식으로 이루어진 정칙계 $\Ssys_{n\rho}$의 정기저.}

%
%
\S\ 11의 보조정리는 \S\ 10의 결론과 함께 다음 사실을 곧바로 준다.

%
%
\emph{$\Gdom_{n\rho}$의 한 사상영역 $\Jdom^*_{\rho\rho}$ 안에 최고차항들의 동차 종결식이 항등적으로 0이 아닌 --- $R_0\ne0$ --- $\rho$개의 다항식이 존재한다면, $\Gdom_{n\rho}$는 제1종이고 따라서 유한 정역이며, 인볼루션기저가 그 정기저이다.}

%
%
이제 조건 $R_0\ne0$은 보조정리에 힘입어 정기저의 존재에 대한 두 번째 증명을 가능하게 한다. 이 증명은 본질적으로 D. 힐베르트가 제시한 것으로서\srcfn{*)}{Über die vollen Invariantensysteme, \S\ 2. 힐베르트의 경우에는 불변식체의 정기저가 다른 방법으로 이미 존재한다고 가정하고, 이 기저를 한 체의 대수적 정수인 양들에 대한 크로네커 기본계로 해석하는 일만을 다룬다.}, 인볼루션기저 자체를 식별해 주지는 않지만 그 대신 위 성질을 갖는 모든 계 $\Ssys_{n\rho}$에 한결같이 적용된다.

%
%
실제로 $R_0\ne0$이고 따라서 대수적으로 독립인 $\Jdom^*_{\rho\rho}$의 $\rho$개 다항식이 다음과 같이 주어졌다고 하자.
\[
\begin{array}{@{}l@{\quad}c@{}}
(1)& G_1^*(x)\rcdots G_\rho^*(x).
\end{array}
\]

%
%
보조정리에 따르면 $x_1\rcdots x_\rho$의 임의의 다항식, 특히 $\Ssys^*_{\rho\rho}$를 포함하는 최소 체 $\Kfield^*_{\rho\rho}$의 모든 다항식은 다항식 (1)에 대하여 대수적으로 정수적으로 종속된다. 따라서 이에 대응하는 $\Ssys_{n\rho}$의 다항식이
\[
\begin{array}{@{}l@{\quad}c@{}}
(2)& G_1(x)\rcdots G_\rho(x)
\end{array}
\]
라면, 전이 원리에 따라 $\Ssys_{n\rho}$를 포함하는 최소 체 $\Kfield_{n\rho}$의 모든 다항식, 따라서 특히 $\Ssys_{n\rho}$의 모든 다항식은 다항식 (2)에 대하여 대수적으로 정수적으로 종속된다. 거꾸로 $\Kfield_{n\rho}$의 함수 가운데 다항식 (2)에 대하여 대수적으로 정수적으로 종속되는 것은 모두 다항식이다. 그러므로 (\S\ 4에 따라) $\Kfield_{n\rho}$를 $\Omega(G_1\rcdots G_\rho)$ 위의 대수적 체로 보면, $\Kfield_{n\rho}$의 다항식 전체 --- 상대정역 $\Gdom_{n\rho}$ --- 는 이 체의 대수적 정수인 양들 전체와 일치한다. $G(x)$를 $\Kfield_{n\rho}$의 다항식으로서 (2)를 유리기저로 보충하는 것이라 하고, $D(x)$를 $G(x)$가 $\Omega(G_1\rcdots G_\rho)$에 관하여 만족하는 $k$차 기약방정식의 판별식이라 하자. 그러면 $\Kfield_{n\rho}$의 모든 다항식, 특히 $\Ssys_{n\rho}$의 모든 다항식 $F(x)$은 대수적 정수인 양으로서 다음 표현을 갖는다.
\[
\begin{array}{@{}l@{\quad}c@{}}
(3)& \displaystyle F(x)=\frac{\Gamma_1(G_i)\cdot G(x)^{k-1}+\Gamma_2(G_i)\cdot G(x)^{k-2}+\cdots+\Gamma_k(G_i)}{D(x)}.
\end{array}
\]

%
%
이제 $F(x)$가 $\Ssys_{n\rho}$의 모든 다항식을 두루 달리게 하자. 이때 해당 함수 $v_1\Gamma_1(G_i)+v_2\Gamma_2(G_i)+\cdots+v_k\Gamma_k(G_i)$로 이루어진 계에 힐베르트의 모듈기저 정리 I (Math. Ann. 36)를 적용하면, $\Ssys_{n\rho}$의 유한개의 다항식
\[
\begin{array}{@{}l@{\quad}c@{}}
(4)& F_1(x),F_2(x),\ldots,F_\lambda(x)
\end{array}
\]
이 존재함을 알 수 있다.

%
%
그리하여 $\Ssys_{n\rho}$의 모든 다항식은 다음 표현을 갖는다.
\[
\begin{array}{@{}l@{\quad}c@{}}
(5)& F(x)=A_1(G_i)F_1(x)+A_2(G_i)\cdot F_2(x)+\cdots+A_\lambda(G_i)F_\lambda(x),
\end{array}
\]
여기서 $A_1\rcdots A_\lambda$는 $G_i(x)$의 다항식이다. 따라서 다항식 (2)와 (4)는 $\Ssys_{n\rho}$의 한 정기저를 이룬다. 즉,

%
%
\emph{정리 VIII: 다항식으로 이루어진 계 $\Ssys_{n\rho}$의 한 사상계 $\Ssys^*_{\rho\rho}$ 안에 최고차항들의 종결식이 항등적으로 0이 아닌 --- $R_0\ne0$ --- $\rho$개의 다항식이 존재한다면, $\Ssys_{n\rho}$는 한 정기저를 갖는다.}

%
%
이 정리는 쉽게 일반화할 수 있다. $R_0\ne0$인 $\rho$개의 다항식 (1)이 $\Ssys^*_{\rho\rho}$를 포함하는 최소 정역 $\Jdom^*_{\rho\rho}$에 속하기만 하면 충분하다. 실제로 정리 VIII에 이르는 추론은 이 경우에도 표현 (5)가 성립함을 보여 준다. 그런데 정의에 따라 최소 포함 정역 $\Jdom_{n\rho}$의 모든 다항식 $L(x)$은 $\Ssys_{n\rho}$의 유한개의 다항식 --- 이 다항식들은 $L(x)$에 따라 달라진다 --- 으로 다항식적으로 나타낼 수 있다. 따라서 $\Ssys_{n\rho}$에는 $\sigma$개의 다항식
\[
\begin{array}{@{}l@{\quad}c@{}}
(6)& K_1(x),K_2(x),\ldots,K_\sigma(x)
\end{array}
\]
이 존재하여, 다항식 $G_i(x)$는 (6)의 다항식 결합이 된다. 그러므로 (6)과 (4)가 한 정기저를 이룬다.

%
%
\emph{정의 VIII: 다항식으로 이루어진 계 $\Ssys_{n\rho}$를 포함하는 최소 정역의 한 사상영역에 최고차항들의 종결식이 항등적으로 0이 아닌 --- $R_0\ne0$ --- $\rho$개의 다항식이 존재할 때, $\Ssys_{n\rho}$를 정칙이라고 한다.}

%
%
따라서 다음이 성립한다.

\emph{정리 IX: 다항식으로 이루어진 모든 정칙계 $\Ssys_{n\rho}$는 한 정기저를 갖는다.}\srcfn{*)}{정기저가 없는 비정칙계가 존재한다는 것을 힐베르트는 한 예로 보였다(Gött. Nachr. 1891, S. 233). 다음 예는 조금 더 간단하다.
\[
\Ssys_{22}=x_1,
x_1x_2,
x_1x_2^2,
\ldots,
x_1x_2^\nu,
\ldots\quad\text{무한히 계속}.
\]
여기서 정기저의 존재는 다음 항등식들이 성립하도록 하는 양의 정수 $\nu$가 존재한다는 것과 동치이다.
\[
x_1x_2^{\nu+\sigma}=\sum c\cdot x_1^{i_0}(x_1x_2)^{i_1}(x_1x_2^2)^{i_2}\cdots(x_1x_2^\nu)^{i_\nu},
\qquad(\sigma=1,2,\ldots).
\]
$x_1,x_2$에 관한 차수를 비교하면 이미 $\sigma=1$에서 지수들에 대한 두 조건방정식
\[
1=i_0+i_1+i_2+\cdots+i_\nu,
\qquad
\nu+1=i_1+2i_2+\cdots+\nu\cdot i_\nu
\]
을 얻는데, 이는 명백히 음이 아닌 정수해를 갖지 않는다. 따라서 $\Ssys_{22}$는 정기저를 갖지 않는다.}

%
%
이제 인볼루션에 대해 \S\ 5에서 한 것과 유사하게 정칙계를 기하학적으로 해석한다. 이를 위해 그곳의 (7)에서와 같이 방정식계
\[
\begin{array}{@{}l@{\quad}c@{}}
(7)& L^*(x)-L^*(\xi)=0
\end{array}
\]
을 고찰하되, $L^*(x)$가 $\Jdom^*_{\rho\rho}$의 모든 다항식을 두루 달리게 한다.

%
%
$x_0,\xi_0$을 사용하여 동차화하면, (7)이 동차형식들 사이의 다음 방정식계로 옮겨간다고 하자.
\[
\begin{array}{@{}l@{\quad}c@{}}
(8)& \xi_0^m\cdot M^*(x)-x_0^m\cdot M^*(\xi)=0.
\end{array}
\]

%
%
여기서 $H^*(x)$가 $L^*(x)$의 최고차항들을 나타낼 때 다음 관계가 성립한다.
\[
\begin{array}{@{}l@{\quad}c@{}}
(9)& H^*(x)\equiv M^*(x)\pmod{x_0}.
\end{array}
\]

%
%
$\Jdom^*_{\rho\rho}$의 \emph{기본점}이라 함은(\S\ 5의 주석 참조) (8)의 해 가운데 $x$와 독립인 것, 곧 $x$에 켤레인 열들과 다른 것을 말한다. 따라서 기본점은 $\xi_1=0,\ldots,\xi_\rho=0$과 다른 값의 계로서 다음을 동시에 만족하는 것이다.
\[
\xi_0^m=0,
\qquad M^*(\xi)=0
\]
또는 (9)에 따라
\[
\xi_0=0,
\qquad H^*(\xi)=0.
\]

%
%
이제 $R_0\ne0$인 $\rho$개의 형식 $H^*(\xi)$, 곧 공통해를 갖지 않는 형식들이 존재한다면, $\Jdom^*_{\rho\rho}$는 기본점을 가질 수 없다. 특히 $\Ssys^*_{\rho\rho}$가 동차형식들로 이루어져 있다면, $R_0\ne0$일 때 $\Ssys^*_{\rho\rho}$ 역시 기본점을 가질 수 없다. 그러한 기본점은 동시에 $\Jdom^*_{\rho\rho}$의 기본점일 것이기 때문이다. 그러나 그 역도 성립한다. 이를 위해 모든 형식 $H^*(x)$로 이루어진 계 $\Jdom(H^*)$를 고찰하면, 이 역시 한 정역을 이룬다.

%
%
그러므로 힐베르트의 모듈기저 정리 I와 힐베르트 보조정리(\S\ 10의 인용 참조)로부터, 그 소멸이 $\Jdom(H^*)$의 모든 형식의 소멸을 수반하는 $\Jdom(H^*)$의 $\rho$개 형식이 존재함을 얻는다. 이제 $\Jdom^*_{\rho\rho}$가 기본점을 갖지 않는다면 이 $\rho$개 형식은 동시에 소멸할 수 없으므로, 그 종결식은 $R_0\ne0$이다.

%
%
따라서 다음이 성립한다.

\emph{다항식으로 이루어진 정칙계 $\Ssys_{n\rho}$는, 이를 포함하는 최소 정역의 한 사상영역 $\Jdom^*_{\rho\rho}$가 기본점을 갖지 않는다는 성질로 특징지어진다. 동차형식으로 이루어진 정칙계 $\Ssys^*_{\rho\rho}$의 경우에는 계 자체가 이미 기본점을 갖지 않는다.}\srcfn{*)}{앞 주석에서 제시한 비정칙계 $\Ssys_{22}$는 다음 기본점을 갖는다.
\[
\xi_1:\xi_2:\xi_0=0:1:0.
\]}

%
%
\section*{\S\ 13. 제1종 상대정역과 정칙계의 예.}

%
%
1. 제1종 상대정역 $\Gdom_{nn}$의 첫 번째 예로, \emph{주어진 순열군 $G$를 허용하는 $x_1\rcdots x_n$의 다항식 전체}, 곧 한 라그랑주 종영역의 다항식 전체를 고찰한다. $x_1\rcdots x_n$은 항등식
\[
\begin{array}{@{}l@{\quad}c@{}}
(1)& x_k^n-\sigma_1(x)x_k^{n-1}+\sigma_2(x)x_k^{n-2}+\cdots\pm\sigma_n(x)=0\quad(k=1,2,\ldots,n)
\end{array}
\]
에 따라 $\Gdom_{nn}$에 속하는 기본대칭함수들에 대하여 대수적으로 정수적이므로, 정의 V에 따라 $\Gdom_{nn}$에 대하여 대수적으로 정수적으로 종속된다. 따라서 $\Gdom_{nn}$은 제1종이다. 더 나아가 $\Gdom_{nn}$의 원소들이 동차형식이고 $n=\rho$이므로, $\Gdom_{nn}$은 제1종 영역으로서(\S\ 10의 결론에 따라) 한 정칙계를 이룬다. 실제로 (1)에 의해 기본대칭함수들의 종결식 $R_0$은 항등적으로 소멸할 수 없으며, 종결식 이론의 계산 규칙에 따르면 $R_0=\pm1$이다. $\Gdom_{nn}$의 \emph{인볼루션형식}은 --- $\Gdom_{nn}$이 제1종이고 $n=\rho$이므로 --- 대응하는 라그랑주 종영역의 인볼루션형식으로서 유일하게 정해진다. 이 영역에서 $x_1\rcdots x_n$에 켤레인 양들은 군 $G$의 순열들로 생기므로, 이 인볼루션형식은 다음과 같다.
\[
\begin{array}{@{}l@{\quad}c@{}}
(2)&
\begin{aligned}
  \Psi(z,u)&=\prod_G
  (z-u_1x_{i_1}-u_2x_{i_2}-\cdots-u_nx_{i_n})\\
  &=z^i+\sum{}'G_{\alpha\alpha_1\ldots\alpha_n}(x)
  z^\alpha u_1^{\alpha_1}\cdots u_n^{\alpha_n},
\end{aligned}
\end{array}
\]
여기서 곱은 $G$의 모든 순열에 걸쳐 취한다. (2)는 \glqq{}군 $G$의 갈루아형식\grqq{}을 나타낸다. 따라서 다음 사실이 성립한다.

\emph{순열군 $G$를 허용하는 $x_1\rcdots x_n$의 모든 다항식은 그 군의 갈루아형식의 계수 $G_{\alpha\alpha_1\ldots\alpha_n}(x)$들의 다항식 결합이다.}

%
%
2. \emph{정칙계}의 예로 다항식으로 이루어진 계 $\Ssys_{n1}$, 곧 대수적으로 독립인 다항식을 하나만 포함하는 계를 고찰한다.

%
%
실제로 $\Ssys^*_{11}$을 $\Ssys_{n1}$의 한 사상계라 하고
\[
  F^*(x)=a_0x^k+a_1x^{k-1}+\cdots+a_k\qquad(a_0\ne0)
\]
를 $\Ssys^*_{11}$의 한 다항식이라 하면,
\[
\begin{array}{@{}l@{\quad}c@{}}
(3)& R_0=a_0\ne0\text{\srcfnmark{*)}}
\end{array}
\]
\srcfntext{*)}{Mertens, a. a. O. \S\ 5.}
가 된다. 그러므로 $\Ssys_{n1}$은 정의 VIII에 따라 정칙이며, 따라서 다음이 성립한다.

\emph{다항식으로 이루어진 모든 계 $\Ssys_{n1}$, 특히 한 부정원의 무한히 많은 다항식으로 이루어진 모든 계는 한 정기저를 갖는다.}\srcfn{**)}{따라서 정기저가 없는 가장 단순한 비정칙계는 계 $\Ssys_{22}$이다. 그러한 계가 실제로 존재한다는 것은 힐베르트의 예와 \S\ 12의 주석에 든 예가 보여 준다.}

%
%
3. 이제 계 $\Ssys_{n1}$을 상대정역 $\Gdom_{n1}$로 특수화한다. 이들은 정칙계이므로 \S\ 12에 따라 제1종 상대정역이다. 따라서 그 정기저는 인볼루션기저로 주어지며, 이 기저는 동시에 이를 포함하는 최소 체 $\Kfield_{n1}$의 인볼루션기저이다. 그런데 \S\ 6에 따르면 $\Kfield_{n1}$의 인볼루션기저의 임의의 함수는 동시에 최소기저이고 --- 우리는 이를 그 체의 뤼로트함수라 불렀다 --- 인볼루션기저의 모든 함수는 이 뤼로트함수에 일차분수식으로 종속된다. 우리의 경우 뤼로트함수는 $\Gdom_{n1}$에 속하므로 다항식이다. 따라서 인볼루션기저의 다른 모든 다항식은 이 뤼로트함수 $L(x)$에 일차분수식으로 종속되고, (3)에 따라 대수적으로 정수적이므로 정수적이며 선형적으로 종속된다. 그러므로 $L(x)$가 정기저가 된다. 여기에 $(u_1=1)$을 놓으면 $\Gdom^*_{11}$의 인볼루션형식은
\[
  \Psi^*(z)=L^*(z)-L^*(x)
\]
가 되며, 여기서도 $L(x)$가 정기저임이 드러난다. 따라서 다음이 성립한다.

\emph{상대정역 $\Gdom_{n1}$은 제1종이다. 그 정기저는 하나의 다항식, 곧 이를 포함하는 최소 체의 뤼로트함수로 이루어진다.}

%
%
4. 상대정역 $\Gdom_{n1}$의 특수한 예로, $s$개 변수의 이차형식에 대한 다항 사영불변식들을 들 수 있다. 앞의 결과와 일치하게, 이들은 잘 알려진 대로 그 형식의 행렬식 $|a_{ik}|$의 다항식 결합으로 나타낼 수 있다. 그리고 이 행렬식은 동시에 대응하는 불변식체의 뤼로트함수이다.

%
%
\section*{\S\ 14. 정수계수 다항식으로 이루어진 계.}

%
%
이제 지금까지 얻은 기저정리들을 정수성에 관하여 강화하고자 한다.

%
%
따라서 이하에서는 계수영역 $\Omega$를 (유한) 대수적 수체라 가정한다. 또한 $[\Omega]$는 $\Omega$ 안의 대수적 정수 전체를 나타내며, 이를 줄여서 정수라 부르기로 한다. 정수계수 다항식이란 $[\Omega]$의 계수를 갖는 다항식을 뜻한다. 이제부터 $F(x),G(x),\ldots$는 오직 정수계수 다항식만을 나타낸다.

%
%
\S\ 7과 유사하게 정수계수 다항식으로 이루어진 여러 계를 고찰한다.

%
%
정수계수 다항식으로 이루어진 임의의 계를 정수계수계 $\Ssys_{n\rho}$라 부른다.

%
%
정수계수 선형족 $\Lsys_{n\rho}$은 $F_1(x)$과 $F_2(x)$를 포함할 때마다 $c_1,c_2$가 $[\Omega]$의 임의의 정수인 $c_1F_1(x)+c_2F_2(x)$도 항상 포함하는 정수계수계이다.

%
%
정수계수 정역 $\Jdom_{n\rho}$은 $F(x)$과 $G(x)$를 포함할 때마다 그 곱 $F(x)\cdot G(x)$도 항상 포함하는 정수계수 선형족이다.

%
%
정수계수 상대정역 $\Gdom_{n\rho}$은 $F(x)$과 $G(x)$를 포함하고 $G(x)\ne0$일 때, 그 몫 $F(x):G(x)$이 정수계수 다항식인 경우에 그리고 그 경우에만 이 몫도 포함하는 정역이다.

%
%
\S\ 7에서와 마찬가지로, 정수계수계 $\Ssys_{n\rho}$를 포함하는 최소 정수계수 정역 $\Jdom_{n\rho}$은 $\Ssys_{n\rho}$의 유한개의 다항식에 대한 정수계수 다항식 결합으로 나타나는 모든 다항식 $H(x)$로 이루어진다.

%
%
이를 포함하는 최소 정수계수 상대정역 $\Gdom_{n\rho}$은 $\Ssys_{n\rho}$의 유한개의 다항식에 대한 유리결합으로 나타나는 모든 정수계수 다항식 $H(x)$로 이루어진다.

%
%
이제 기저정리들을 옮기자. 유리기저에 관해서는 오직 \glqq{}유리적인 것\grqq{}에만 기초하는 체, 최소 포함 체, 유리기저의 개념은 아무런 변화도 겪지 않는다는 점에 유의해야 한다. 더 나아가 정수계수계 $\Ssys_{n\rho}$의 사상계 $\Ssys^*_{\rho\rho}$를 다시 정수계수계가 되도록 임의로 많은 방식으로 선택할 수 있다. 정수계수 선형족에서도 개수의 제한을 이끌어 내는 추론이 그대로 유지되므로 다음이 성립한다.

%
%
\emph{정리 V$'$: 모든 정수계수계 $\Ssys_{n\rho}$는 한 유리기저를 갖는다. 특히 정수계수 선형족 $\Lsys_{n\rho}$의 유리기저는 많아야 $\rho+1$개의 다항식을 포함하도록 선택할 수 있다.}

%
%
이제 다항식으로 이루어진 정역의 인볼루션기저(\S\ 8)에 대응하는 것을 고찰한다. 이를 위해 \emph{양들의 열에 대한 대칭함수}의 정리를 \emph{정수성}에 관하여 다음 보조정리로 확장해야 한다.

%
%
\emph{보조정리. $n$개의 양들의 열 $(yz\cdots t)$}
\[
\begin{array}{@{}l@{\quad}c@{}}
(1)&
\begin{array}{cccc}
  y_1&z_1&\cdots&t_1,\\
  y_2&z_2&\cdots&t_2,\\
  \vdots&\vdots&&\vdots\\
  y_n&z_n&\cdots&t_n
\end{array}
\end{array}
\]
\emph{에 대한 정수계수 대칭다항식들은 그 가운데 유한개의 함수로 된 정수계수 다항함수이다. 구체적으로 그 유한개의 함수는 기본대칭함수}
\[
\begin{array}{@{}l@{\quad}c@{}}
(2)&
\begin{gathered}
  \sigma_1(y),\sigma_2(y),\ldots,\sigma_n(y);\quad
  \sigma_1(z),\sigma_2(z),\ldots,\sigma_n(z);\quad\ldots;\\
  \sigma_1(t),\ldots,\sigma_n(t);
\end{gathered}
\end{array}
\]
\emph{와 함수}
\[
\begin{array}{@{}l@{\quad}c@{}}
(3)&
  \chi_1(y,z,\ldots,t),\ldots,\chi_k(y,z,\ldots,t),
\end{array}
\]
\emph{이며, 함수 (3)은 함수 (2)에 대하여 대수적으로 정수적이고 정수계수적으로 종속된다.}

%
%
실제로 $n$개의 양 $y_1,\ldots,y_n$은 $\sigma_1(y),\ldots,\sigma_n(y)$에 대하여 대수적으로 정수적이고 정수계수적으로 종속된다. $z_1,\ldots,z_n$과 $\sigma_1(z),\ldots,\sigma_n(z)$ 사이에도 같은 관계가 성립하며, 그 밖의 열들에도 마찬가지이다. 따라서 양들의 열 (1)에 대한 대칭함수체는 함수 (2)가 정하는 체 위의 대수적 체를 이루고, 이 체의 대수적으로 정수적이며 정수계수적인 양들은 양들의 열 (1)에 대한 정수계수 대칭다항식들과 일치한다. 그러므로 (3)을 크로네커 기본계로 정하면 보조정리가 증명된다.

%
%
이제 $\Jdom_{n\rho}$을 정수계수 정역, $\Jdom^*_{\rho\rho}$을 정수계수 사상영역이라 하고
\[
  \Psi^*(z,u)=G^*(x)\cdot z^i+
  \sum{}'G^*_{\alpha\alpha_1\ldots\alpha_\rho}(x)\cdot
  z^\alpha u_1^{\alpha_1}\cdots u_\rho^{\alpha_\rho}
\]
를 $\Jdom^*_{\rho\rho}$의 인볼루션형식이라 하자. 여기서 $G^*(x)$는 0차, 곧 $[\Omega]$의 한 수일 수도 있다. 그러면 함수 (2)에 대응하는 기본대칭함수들은 다음 몫들 가운데 몇몇으로 주어진다.
\[
\begin{array}{@{}l@{\quad}c@{}}
(4)&
  G^*_{\alpha\alpha_1\ldots\alpha_\rho}(x):G^*(x).
\end{array}
\]

%
%
보조정리에 따라 함수 (3)은 (2)에 대하여 대수적으로 정수적이고 정수계수적으로 종속된다. 그러므로 대수적 정수인 계수를 갖는 다항식에 대한 가우스 정리의 확장에 의해, (3)에 대응하는 $\Kfield^*_{\rho\rho}$의 함수들은
\[
\begin{array}{@{}l@{\quad}c@{}}
(5)&
  K_\lambda^*(x):(G^*(x))^\sigma
\end{array}
\]
로 주어진다. 여기서 $K_\lambda^*(x)$는 $\Kfield^*_{\rho\rho}$의 정수계수 다항식이다. 따라서 상대정역을 다루는 경우 이 함수들은 그 영역에 속하고, 일반적으로는 $\Jdom^*_{\rho\rho}$의 다항식들의 몫이다.

%
%
이제 $F^*(x)$가 $\Jdom^*_{\rho\rho}$의 임의의 정수계수 다항식이면
\[
  F^*(x)=\frac1i\Bigl(F^*(x)+F^*(x^{(1)})+\cdots+F^*(x^{(i-1)})\Bigr)
\]
이고, 따라서 $i\cdot F^*(x)$는 양들의 열 $x,x^{(1)},\ldots,x^{(i-1)}$에 대한 정수계수 대칭다항식이 된다.

%
%
보조정리와 전이 원리에 따라 다음이 성립한다.

\emph{정리 VI$'$: 정수계수 정역 $\Jdom_{n\rho}$에서는 각 인볼루션형식이 하나의 특별한 유리기저를 정하며, $\Jdom_{n\rho}$의 모든 다항식 $F(x)$은 다음 표현을 갖는다.}
\[
  F(x)=\frac{\Gamma(H(x),H_1(x),\ldots,H_\sigma(x))}{i\cdot H(x)^k},
\]
\emph{여기서 $\Gamma$는 정수계수 다항함수이다. 정수계수 상대정역 $\Gdom_{n\rho}$의 사상영역도 정수계수 상대정역 $\Gdom^*_{\rho\rho}$인 경우, 분모 $H(x)$는 대응하는 인볼루션형식의 최고차항의 계수로 주어진다.}

%
%
\section*{\S\ 15. 제1종 정수계수 상대정역과 정수계수 정칙계.}

%
%
정수계수계에서는 정기저 대신 정수계수 정기저를 사용한다. 이는 그 계의 유한개의 정수계수 다항식으로서, 계의 모든 정수계수 다항식을 이 유한개에 대한 정수계수 다항식 결합으로 나타낼 수 있게 하는 것을 말한다.

%
%
정수계수 정기저에 관한 정리들은 본질적으로 앞의 정리들과 나란히 진행되며, 증명만 약간 수정하면 된다.

%
%
먼저 정의 V(\S\ 9)를 그대로 옮길 수 있다.

\emph{정의 V$'$: 한 양 $\xi$가 최고차항의 계수가 $[\Omega]$의 한 단원이고 나머지 계수들이 $\Gdom_{n\rho}$의 다항식인 어떤 방정식을 만족할 때, $\xi$를 정수계수 상대정역 $\Gdom_{n\rho}$에 관하여 대수적으로 정수적이고 정수계수적이라고 한다.}

%
%
실제로 대수적 정수인 계수를 갖는 다항식에 대한 가우스 정리의 확장은 \S\ 9에서와 마찬가지로 이로부터 기약방정식에 의한 정의를 얻을 수 있음을 보여 준다.

%
%
정의 V$'$로부터 정의 VII도 옮길 수 있다.

\emph{정의 VII$'$: 정수계수 상대정역 $\Gdom_{n\rho}$이 어떤 상대정역 $\Gdom^*_{\rho\rho}$을 사상영역으로 가지며, $x_1\rcdots x_\rho$의 임의의 정수계수 다항식이 모두 $\Gdom^*_{\rho\rho}$에 대하여 대수적으로 정수적이고 정수계수적으로 종속될 때, $\Gdom_{n\rho}$을 제1종이라 한다.}

%
%
이 제1종 영역들의 유한성을 보이기 위해, 먼저 \S\ 10에서와 마찬가지로 정의로부터 인볼루션형식의 최고차항의 계수가 $[\Omega]$의 한 단원임이 따른다는 점에 유의한다.

%
%
따라서 정리 VI$'$ (\S\ 14)에 의해 $\Gdom_{n\rho}$의 모든 다항식 $F(x)$은 다음 표현을 갖는다.
\[
  F(x)=\frac{\Gamma(H_1(x)\rcdots H_\sigma(x))}{i}.
\]
이제 정수계수 다항식 $\Gamma(H_1\rcdots H_\sigma)$들로 이루어진 계에, 정수 계수에 대하여 진술된 힐베르트의 정수계수 모듈기저 정리 II (Math. Ann. 36)를 대수적 정수계수 다항식으로 확장한 형태로 적용한다.\srcfn{*)}{$w_1\rcdots w_k$가 $\Omega$의 대수적 정수들에 대한 한 기본계이고
\[
  \Gamma(H_1\rcdots H_\sigma)=\Gamma_1(H)w_1+\cdots+\Gamma_k(H)w_k
\]
로 놓으면, $\Gamma_1(H)\rcdots \Gamma_k(H)$는 보통 정수 계수를 갖는다. 형식 $v_1\Gamma_1(H)+\cdots+v_k\Gamma_k(H)$로 이루어진 계에 정리 II를 적용하면 이 확장의 타당성이 곧바로 따른다.} 그러면 정수계수 정기저의 존재가 드러난다. 즉,

\emph{정리 VII$'$: 모든 제1종 정수계수 상대정역은 한 정수계수 정기저를 갖는다.}

%
%
이제 정칙계의 정의를 옮긴다.

\emph{정의 VIII$'$: 정수계수계 $\Ssys_{n\rho}$을 포함하는 최소 정수계수 정역의 한 사상영역 $\Jdom^*_{\rho\rho}$ 안에 최고차항들의 종결식이 $[\Omega]$의 한 단원인 --- $R_0=\eps$ --- $\rho$개의 다항식이 존재할 때, $\Ssys_{n\rho}$을 정수계수 정칙이라 한다.}

%
%
정수계수 정기저의 존재증명을 옮기기 위해, \S\ 11의 보조정리가 다음과 같은 더 강한 형태를 갖는다는 점에 유의한다. 이는 \S\ 11의 방정식 (5), (6)과 $A_1(\lambda)\rcdots A_\tau(\lambda)$가 $\lambda_1\rcdots\lambda_\rho$의 정수계수 다항함수라는 사실에서 따른다.

\emph{$x_1\rcdots x_\rho$의 임의의 모든 정수계수 다항식이 $\rho$개의 정수계수 다항식에 대하여 대수적으로 정수적이고 정수계수적으로 종속되기 위한 충분조건은, 이 다항식들의 최고차항들의 종결식이 $[\Omega]$의 한 단원인 것이다 --- $R_0=\eps$.}

%
%
이 보조정리로부터 \S\ 12에서와 똑같이 정수계수 정칙계의 모든 다항식 $F(x)$에 대하여 \S\ 12의 표현 (3)을 얻는다. 이제 $\Gamma_1(G_i)\rcdots\Gamma_k(G_i)$는 $G_i$의 정수계수 다항식이다. 더 나아가 \S\ 12에서와 같이, 대수적 정수계수 다항식으로 확장한 힐베르트 정리 II (Math. Ann. 36)를 적용하고 이를 최소 포함 정수계수 정역의 정의와 결합하면 정수계수 정기저의 존재를 얻는다. 즉,

\emph{정리 IX$'$: 모든 정수계수 정칙계는 한 정수계수 정기저를 갖는다.}

%
%
제1종 정수계수 상대정역의 예로, \S\ 13의 예 1과 유사하게 한 라그랑주 종영역의 정수계수 다항식 전체 $\Gdom_{nn}$을 들 수 있다. $\Gdom_{nn}$이 정수계수 제1종이라는 사실은 항등식
\[
  x_k^n-\sigma_1(x)x_k^{n-1}+\cdots\pm\sigma_n(x)=0
  \qquad(k=1,2,
\cdots,n)
\]
에서 따른다. 기본대칭함수들의 종결식 값이 $\pm1$이므로, $\Gdom_{nn}$은 정수계수 정칙계이기도 하다.

%
%
또 다른 예는 \S\ 14의 보조정리에 따라 $n$개의 양들의 열에 대한 정수계수 대칭다항식 전체로 주어진다.

%
%
에를랑겐, 1914년 5월.

\clearpage
\section*{P06 교정 주석}
\addcontentsline{toc}{section}{P06 교정 주석}
원문 인쇄면 174쪽, 식(6) 뒤. 분자가 항등적으로 영이 아니라는 주장에는 $f^*\ne0$이라는 조건이 필요하다. 허용되는 영함수에 대해 $F^*=0$, $H^*=1$로 놓으면 모든 켤레 분자는 영이고 분모는 여전히 영이 아니다. 모든 함수에 대한 양화와 식(7)의 조건 $H^*(\xi)\ne0$은 유지된다. 영함수는 그 식에 $0=0$만을 준다.
\smallskip
\noindent{\small AI가 제작한 교정 주석과 정정: OpenAI Codex, GPT-6 Astra, Ultra 추론 강도.}

\clearpage
\setcounter{footnote}{0}
\phantomsection
\addcontentsline{toc}{section}{7. 유한군 불변식의 유한성 정리}
%
%

%
%
\section*{7. 유한군 불변식의 유한성 정리}
\begin{center}
Math. Ann. 77 (1916), S. 89--92
\end{center}

\begin{center}
{\Large 유한군 불변식의 유한성 정리.}\par
\vspace{1.2em}
지음\par
\vspace{0.8em}
에를랑겐의 {\scshape Emmy Noether}.
\end{center}

이하에서는 대칭함수 이론에만 기초한 완전히 초등적인 유한군 불변식의 유한성 증명을 제시한다. 이 증명은 동시에 완전계를 실제로 제시해 준다. 이에 반해 힐베르트의 가군 기저 정리(Ann. 36)에 의존하는 통상적인 증명은 존재 증명에 그친다.\NoetherSrcNote{*)}{예를 들어 Weber, \emph{Lehrbuch der Algebra} (2. Aufl.) 2. Band, \S\ 57을 보라.}

%
%
유한군 $\Hgrp$가 행렬식이 0이 아닌 $h$개의 선형변환 $A_1\cdots A_h$로 이루어진다고 하자. 여기서 $A_k$는 선형변환
\[
  x_1^{(k)}=\sum_{\nu=1}^{n}a_{1\nu}^{(k)}x_\nu,\cdots,
  x_n^{(k)}=\sum_{\nu=1}^{n}a_{n\nu}^{(k)}x_\nu
\]
을 나타내며, 이를 줄여서 $(x^{(k)})=A_k(x)$로 쓴다. 따라서 군 $\Hgrp$는 원소가 $x_1\cdots x_n$인 열 $(x)$를 원소가 $x_1^{(k)}\cdots x_n^{(k)}$인 열들 $(x^{(k)})$로 옮긴다. $A_1\cdots A_h$ 가운데 항등변환이 반드시 포함되므로, 열들 $(x^{(k)})$ 가운데에도 열 $(x)$가 포함된다. 군의 다항 (절대)불변식이란 $x_1\cdots x_n$의 다항함수 가운데 $A_1\cdots A_h$를 적용해도 항등적으로 변하지 않는 것을 뜻한다. 그러므로 그러한 불변식 $f(x)$에 대해서는 다음이 성립한다:
\begin{equation}
  f(x)=f(x^{(1)})=\cdots=f(x^{(k)})=
  \frac1h\cdot\sum_{k=1}^{h}f(x^{(k)}).
\tag{1}
\end{equation}

%
%
\subsection*{1.}
공식 (1)은 $f(x)$가 변수열 $(x^{(1)})\cdots(x^{(h)})$의 다항 대칭함수임을 나타낸다. 더구나 각 항 $f(x^{(k)})$에는 오직 하나의 열 $(x^{(k)})$만 들어 있으므로, 여기서는 보통 \emph{단일형}이라고 부르는 가장 간단한 경우를 다룬다. 변수열의 대칭함수에 관한 잘 알려진 정리\NoetherSrcNote{*)}{2절의 주를 보라.}에 따라 $f(x)$는 이 열들의 기본대칭함수들로 다항적으로 표현된다. 다시 말해, 다음 ‘갈루아 레졸벤트’의 계수 $G_{\alpha\alpha_1\cdots\alpha_n}(x)$로 표현된다:
\[
\begin{aligned}
  \Phi(z,u)
  &=\prod_{k=1}^{h}\bigl(z+u_1x_1^{(k)}+\cdots+u_nx_n^{(k)}\bigr) \\
  &=z^h+
    \sum G_{\alpha\alpha_1\cdots\alpha_n}(x)
    z^{\alpha}u_1^{\alpha_1}\cdots u_n^{\alpha_n}
\end{aligned}
\qquad
\left(\begin{array}{c}
(\alpha+\alpha_1+\cdots+\alpha_n=h)\\
\alpha\ne h
\end{array}\right)
\]
여기서 $G_{\alpha\alpha_1\cdots\alpha_n}(x)$는 $x$들에 관하여 차수가 $\alpha_1+\cdots+\alpha_n$인 불변식이다. 따라서 다음이 증명되었다:

\emph{갈루아 레졸벤트의 계수 $G_{\alpha\alpha_1\cdots\alpha_n}(x)$들은 군의 완전 불변식계를 이룬다. 즉 모든 불변식은 이 유한 개의 불변식들로 다항적으로 표현된다.}

%
%
\subsection*{2.}
(1)에서 출발하여 다음과 같은 한층 더 초등적인 고찰을 할 수도 있다.\NoetherSrcNote{**)}{이는 앞에서 언급한 단일형의 경우에 변수열의 대칭함수에 관한 정리를 증명하는 것과 같다.} 이 고찰은 동시에 두 번째 완전계를 준다. 다음과 같이 두자:
\[
  f(x)=a+b x_1^{\mu_1}\cdots x_n^{\mu_n}
       +\cdots+c x_1^{\nu_1}\cdots x_n^{\nu_n},
\]
여기서 $a,b,\cdots,c$는 상수이다. 그러면 (1)에 따라 다음을 얻는다:
\[
\begin{aligned}
 h\cdot f(x)
 &=h\cdot a+b\cdot\sum_{k=1}^{h}x_1^{(k)^{\mu_1}}\cdots x_n^{(k)^{\mu_n}}
     +\cdots
     +c\cdot\sum_{k=1}^{h}x_1^{(k)^{\nu_1}}\cdots x_n^{(k)^{\nu_n}} \\
 &=h\cdot a+b\cdot J_{\mu_1\cdots\mu_n}+\cdots+c\cdot J_{\nu_1\cdots\nu_n}.
\end{aligned}
\]

%
%
따라서 모든 불변식은 다음의 특수한 불변식들을 선형결합하여 얻을 수 있으므로, 이 특수한 불변식들에 대해서만 보이면 충분하다:
\[
  J_{\mu_1\cdots\mu_n}
  =\sum_{k=1}^{h}x_1^{(k)^{\mu_1}}\cdots x_n^{(k)^{\mu_n}}.
\]
그런데 수치 인자를 제외하면 $J_{\mu_1\cdots\mu_n}$는 식
\[
  S_\mu=\sum_{k=1}^{h}
  \bigl(u_1x_1^{(k)}+\cdots+u_nx_n^{(k)}\bigr)^{\mu}
\]
에서 $u_1^{\mu_1}\cdots u_n^{\mu_n}$의 계수이다. 여기서 $\mu=\mu_1+\cdots+\mu_n$으로 두었으며, 이 식은 $h$개의 선형형식
\[
  \xi_1=u_1x_1^{(1)}+\cdots+u_nx_n^{(1)},\cdots,
  \xi_h=u_1x_1^{(h)}+\cdots+u_nx_n^{(h)}
\]
의 $\mu$차 거듭제곱합을 나타낸다. 잘 알려져 있듯이 무한히 많은 거듭제곱합 $S_\mu$는
\[
  S_1,S_2,\cdots,S_h
\]
의 다항식으로 표현된다. 이들의 계수는 불변식
\[
  J_{\mu_1\cdots\mu_n}\qquad(\mu_1+\cdots+\mu_n\leqq h)
\]
으로 주어진다. 따라서 두 번째 완전계가 얻어진다:

\emph{군의 완전 불변식계는 $\mu_1+\cdots+\mu_n\leqq h$인 모든 불변식 $J_{\mu_1\cdots\mu_n}$으로 주어진다. 여기서 $h$는 군의 위수이다.}

%
%
1절과의 관계는 거듭제곱합 $S_1\cdots S_h$를 $\xi_1\cdots\xi_h$의 기본대칭함수들로 바꿀 수 있다는 관찰로 얻어진다. 이렇게 하면 그곳에서 제시한 $G_{\alpha\alpha_1\cdots\alpha_n}(x)$의 완전계에 이른다. 두 결과는 모두 모든 불변식을 $x$들에 관한 차수가 군의 위수 $h$를 넘지 않는 불변식들로 다항적으로 표현할 수 있음을 보여 준다.

%
%
\subsection*{3.}
방금 증명한 사실로부터 유리 표현에 관한 결론을 이끌어 낼 수 있다. 모든 유리 절대불변식은, 분자와 분모에 군에 관한 분모의 켤레들을 곱해 보면 곧 알 수 있듯이, 서로소일 필요가 없는 두 다항 절대불변식의 몫으로 나타낼 수 있다. 따라서 모든 유리 절대불변식은 갈루아 레졸벤트 $\Phi(z,u)$의 계수 $G_{\alpha\alpha_1\cdots\alpha_n}(x)$로 유리적으로 표현된다. 이 정리는 유한성 정리를 거치지 않고도 여러 방식으로 쉽게 증명할 수 있으며, 이미 Weber II, \S\ 58에 진술되어 있다.\NoetherSrcNote{*)}{그러나 그곳의 증명은 타당하지 않다. 공식 (7)의 함수 $\Psi(t)$가 불변식들을 계수로 갖는다는 것만 보였을 뿐, 그 불변식들이 $\Phi(t)$의 계수들로 유리적으로 표현된다는 것은 보이지 않았기 때문이다. $x_i^{(k)}$를 $\xi_k$로 표현할 때 라그랑주의 보간공식 대신 잘 알려진 미분법을 사용하면 이 빈틈을 피할 수 있다. 이는 항등식 $\Phi(-\xi_k,u)=0$을 모든 $u$에 관하여 미분해서 얻는 관계식
\[
 \left[\frac{\partial\Phi}{\partial u_i}
      -x_i^{(k)}\cdot\frac{\partial\Phi}{\partial z}\right]_{z=-\xi_k}=0
\]
이다. 이에 따라 Weber의 공식 (7), (8) 대신 임의의 유리함수 $\omega(x)$에 대해 다음 공식을 써야 한다:
\[
 \omega(x_1^{(k)}\cdots x_n^{(k)})=
 \left.
 \omega\!\left(
   \frac{\frac{\partial\Phi}{\partial u_1}}{\frac{\partial\Phi}{\partial z}}\cdots,
   \frac{\frac{\partial\Phi}{\partial u_n}}{\frac{\partial\Phi}{\partial z}}
 \right)\right|_{z=-\xi_k}.
\]
이 식에는 이제 $\Phi(z,u)$의 계수들만 들어 있으며, $\omega(x)$가 불변식일 때 모든 $k$에 대해 합하면 원하는 표현을 얻는다.}

%
%
\subsection*{4.}
끝으로 여기서 도출한 결과들이 나의 논문 \emph{Körper und Systeme rationaler Funktionen}에 암묵적으로 포함되어 있음을 언급해 둔다.\NoetherSrcNote{*)}{Math. Ann. 76, S. 161 (1915).} 구체적으로 1절에서 제시한 완전계는 정리 VI와 VII에, 이 유한성 정리와 무관한 유리 표현의 증명은 정리 III에 들어 있다.\NoetherSrcNote{**)}{상대불변식에 대해서는 정리 VIII과 IX에 유한성 증명이 들어 있다. 그것 역시 통상적인 증명과는 다른 기초에 서 있지만 존재 증명에 그친다. 이는 상대불변식들이 체를 이루지 않기 때문이다.} 그러나 여기의 특수한 경우에는 증명의 진행과 결과가 일반 이론에 비해 현저히 단순해진다. 일반적인 고찰을 유한군의 불변식에 적용한다는 생각은 E. Fischer 선생과의 대화에서 얻었다. 2절에서 제시한 증명의 내용---일반적인 경우에는 거기에 나타나는 완전계를 유리적으로 정의할 수 없다---과 3절의 주도 그에게서 비롯되었다.

에를랑겐, 1915년 5월.

\clearpage
\setcounter{footnote}{0}
\phantomsection
\addcontentsline{toc}{section}{8. 임의로 많은 기본형식들의 계의 불변식에 대한 다항식적 표현}
%
%

%
%
\section*{8. 임의로 많은 기본형식들의 계의 불변식에 대한 다항식적 표현}
\begin{center}
Math. Ann. 77 (1916), 93--102쪽
\end{center}

\begin{center}
{\Large 임의로 많은 기본형식들의 계의 불변식에 대한\\
다항식적 표현에 관하여.}\par
\vspace{1.2em}
저자\par
\vspace{0.8em}
에를랑겐의 {\scshape Emmy Noether}.
\end{center}

%
%
힐베르트는 슈바르츠 기념논문집에 기고한 「임의로 많은 기본형식들의 계의 불변식에 관하여」의 끝에서, \emph{이 불변식들을 계 $(J,PJ)$의 유한개의 불변식으로 다항식적으로 나타낼 수 있다}고 추측한다. 여기서 $J$는 기본형식들 가운데 선형독립인 것들의 완전 불변식계를 뜻하고, 불변식 $PJ$는 $J$에서 편극연산을 통해 얻어진다.

아래 I에서 나는 이 추측이 실제로 성립하며, 그것이 다음 환원정리의 간단한 귀결임을 보인다. \emph{각각 $n$개의 변수를 갖는 임의로 많은 변수열의 모든 형식은, 고정된 $n$개의 변수열만을 포함하고 원래 형식에서 편극연산으로 유도된 형식들의 편극형들의 합으로 나타낼 수 있다.} 이 환원정리는 F. Mertens가 처음 정식화한 특수한 경우로서, A. Capelli, F. Mertens, J. Deruyts가 클렙슈--고르단 급수전개를 각각 $n$개의 변수를 갖는 임의로 많은 변수열의 형식으로 일반화한 결과에 속한다. 그 일반화는 미분연산을 형식적으로 변형하여 증명되었다.\NoetherSrcNote{*)}{F. Mertens: 「Über eine Formel der Determinantentheorie」, 공식 5, Sitzb. d. Ak. d. Wiss. Wien, Bd. 91, II. Abt. (1885); 더 상세한 설명과 응용은 「Über ganze Funktionen von $m$ Systemen von je $n$ Unbestimmten」, Monatsh. f. Math. u. Phys. 4. Jahrg. (1893) 참조. Capelli의 몇몇 증명(1882년 이후)에 대해서는 예컨대 Math. Ann. 37 또는 등사판 「Lezioni sulla teoria delle forme algebriche」(1902) 참조. J. Deruyts: \emph{Essai d'une theorie générale des formes algébriques}, Mém. d. l. Soc. Roy. des Sciences de Liège (2) 17 (1892).}

%
%
II에서는 이 문제를 선형 형식족의 동치 문제로 파악하여 환원정리의 좀 더 개념적인 증명도 제시한다. 이러한 관점과 사용되는 논의의 중요한 부분은 E. Fischer의 논문 「대수적 가군계에 관하여」\NoetherSrcNote{*)}{E. Fischer: Über algebraische Modulsysteme und lineare homogene partielle Differentialgleichungen mit konstanten Koeffizienten, \S\ 1--4, J. f. M. 140 (1911).} 및 그 논문에 이어지는 그의 미발표 정리들에서 가져온 것이다. Fischer는 이 정리들을 사용할 수 있도록 매우 친절하게 허락해 주었다.

마지막으로 III에서는 이에 대응하는 논의가 어떻게 가장 일반적인 급수전개도 이끌어 내는지 보인다.

%
%
\subsection*{I.}
$\theta$를 $N>n$개의 각 변수열
\[
  A_1,A_2,\ldots,A_N
\]
각각에 대해 동차인 형식이라 하자. 일반적으로 변수열 $A_k$는 $n$개의 원소
\[
  A_k^{(1)},A_k^{(2)},\ldots,A_k^{(n)}
\]
로 이루어진다고 하자. $\theta$의 계수들은 부정원일 수도 있고, 주어진 유리성 영역의 원소일 수도 있다. 또한 $P$는 편극연산자들
\[
  P_{hk}=\left(A_h\frac{\partial}{\partial A_k}\right)
  =A_h^{(1)}\frac{\partial}{\partial A_k^{(1)}}+
   A_h^{(2)}\frac{\partial}{\partial A_k^{(2)}}+
   \cdots+
   A_h^{(n)}\frac{\partial}{\partial A_k^{(n)}}
\]
로 이루어진, \emph{유리수 계수를 갖는} 다항식이라 하자. 여기서 곱 $P_{hk}P_{ij}\theta$는 $P_{ij}\theta$에 $P_{hk}$를 적용한 것을 뜻한다.

%
%
그러면 앞서 말한 \emph{환원정리}는 항등식
\begin{equation}
  \theta=\sum PZ
\tag{1}
\end{equation}
으로 주어진다. 여기서 형식 $Z$는 변수열
\[
  A_1,A_2,\ldots,A_n
\]
만을 포함하며, $\theta$에서 편극연산 $P$를 통해 유도된 형식이다.

%
%
이제 같은 차수와 같은 변수 개수를 갖는 $N$개의 형식
\[
  F_1,F_2,\ldots,F_N
\]
으로 이루어진 기본형식계 $(F)$를 생각하자. 이 형식들의 계수는 각각 $N$개의 변수열
\[
  A_1,A_2,\ldots,A_N
\]
으로 취한다. $(J)$는 유한개의 불변식으로 이루어진 $F_1,\ldots,F_n$의 완전계라 하자. 즉 $F_1,\ldots,F_n$의 모든 불변식을 $(J)$의 불변식들로 된 다항식으로 나타낼 수 있다고 하자. \emph{그러면 증명해야 할 힐베르트의 추측은 다음과 같다. $F_1,\ldots,F_N$의 동시불변식 $S$에 대한 완전계는 유한개의 불변식 $(J,PJ)$로 주어진다. 여기서 $PJ$는 $J$에서 편극연산 $P$를 통해 유도할 수 있는 모든 불변식을 뜻한다.}

%
%
실제로 $(F)$의 유리수 계수 동시불변식 $S$는 $N$개의 변수열 $A_1,\ldots,A_N$ 각각에 대해 동차인, 유리수 계수의 형식이 된다. 따라서 항등식 (1)을 $S$에 적용할 수 있다. 편극연산 $P$를 $S$에 적용하면 다시 불변식이 생긴다는 것은 잘 알려져 있으므로, 형식 $Z$도 다시 불변식이 된다. 더구나 $Z$는 변수열 $A_1,
\ldots,A_n$만을 포함하므로 $F_1,
\ldots,F_n$의 동시불변식이다. 따라서 가정에 의해 $Z$는 불변식 $(J)$로 된 다항식이다. 그러므로 항등식 (1)은
\[
  S=\sum PY
\]
로 바뀐다. 여기서 $Y$는 불변식 $(J)$의 거듭제곱 단항곱을 뜻한다. 그러나 $Y$에 연산 $P$를 실제로 수행하는 것은, $P$의 정의에 따라 단순 편극연산 $P_{hk}$를 유한번 연속해서 수행하는 것이다. 곱의 미분법칙에 의해 이 연산은 불변식 $(J)$와 $(PJ)$의 곱들의 합을 낳는다. $(J)$와 $(PJ)$의 곱에 $P_{hk}$를 적용하는 경우에도 마찬가지이므로, $PY$ 역시 $(J,PJ)$의 다항식적 결합만을 낳는다. \emph{따라서 모든 불변식 $S$를 $(J,PJ)$로 다항식적으로 나타낼 수 있음이 증명되었다.}

%
%
\subsection*{II.}
환원정리의 증명에 앞서 선형 형식족에 관한 몇 가지 사실을 먼저 살펴본다. 여기서 \emph{$K$ 위의 선형 형식족}이란 같은 차수의 형식들로 이루어진 계로서, $\theta_1$과 $\theta_2$가 계에 속하면 $c_1\theta_1+c_2\theta_2$도 언제나 계에 속한다는 닫힘성을 갖는 것을 뜻한다. 여기서 $c_1,c_2$는 주어진 유리성 영역 $K$의 원소이고, $K$는 $\theta$들의 계수영역을 포함한다. 이 형식들 가운데에는 $K$ 위에서 선형독립인 형식이 언제나 유한개, 이를테면 $\rho$개 있고, 임의의 $\rho+1$개 형식은 $K$의 계수를 갖는 하나의 선형관계를 만족한다. 이때 그 형식족은 $K$에 관하여 \emph{랭크 $\rho$}를 갖는다고 한다.\NoetherSrcNote{*)}{같은 차수라는 제한을 버리거나, $K$가 $\theta$들의 계수영역을 포함한다는 제한을 버리면, 랭크가 유한하지 않을 수도 있는 더 일반적인 선형 형식족을 얻는다.} 다음 \emph{사실}은 쉽게 알 수 있다. $\mathfrak L$과 $\mathfrak T$가 $K$ 위에서 같은 랭크 $\rho$를 갖는 두 선형 형식족이고 $\mathfrak T$가 $\mathfrak L$의 부분족이면, $\mathfrak L$과 $\mathfrak T$는 동일하다(동치이다).

%
%
이를 전제로 환원정리로 넘어가자. 이제 I의 항등식
\[
  \theta=\sum PZ
\]
의 양변에 같은 편극연산 $P$를 적용하면 $P\theta$에 대한 이와 유사한 항등식을 얻는다. 따라서 환원정리는 동시에 $\theta$의 \emph{모든} 편극형을 특수한 편극형 $PZ$들의 합으로 나타내는 표현을 준다. 한편 이 특수한 편극형들은 분명 전체 편극형들 가운데 포함되므로, \emph{환원정리는 이 두 선형 형식족이 동치임을 말한다}. 특히 각 \emph{개별} 변수열에서 $\theta$와 같은 차수를 갖는 형식들로 이루어진 부분족들도 동치이다. 이 동치를 증명하기 위하여 형식 $Z$들로 정해지는 하나의 중간 형식족을 더 도입한다. 간단히 하기 위해 먼저 \emph{세 개의 이진 변수열인 경우 $(N=3,n=2)$}에 대한 증명을 제시한 뒤, 이와 나란히 진행되는 일반적인 증명을 짧게 설명한다.

%
%
그러므로 $\theta(x,y,z)$가 변수열
\[
  x_1x_2;\quad y_1y_2;\quad z_1z_2
\]
에서 각각 차수 $\alpha,\beta,p$를 갖는다고 하자. 우선 $\theta$의 계수는 부정원 $a$라고 하며, 유리수라고 해도 좋다. 이제 다음 세 선형 형식족을 생각한다.

%
%
\noindent
1. 형식족 $\mathfrak L$. 이는 $\theta$에서 편극연산 $P$를 통해 얻어지고 각 변수열 $x,y,z$에서 $\theta$와 같은 차수를 갖는 모든 형식 $f(x,y,z)$로 이루어진다. 특히 $\mathfrak L$은 $\theta$도 포함한다. $f(x,y,z)$를 $x,y,z$와 부정원 $a$에 관한 형식으로 보면, 그 계수영역은 유리수체 $R$이다. $K$ 역시 $R$로 잡아야 한다. $c_1,c_2$가 유리수일 때에만 $c_1P_1+c_2P_2$도 다시 하나의 편극연산 $P$가 되기 때문이다. $\mathfrak L$은 $R$에 관하여 랭크 $\rho$를 갖는다고 하자.

%
%
\noindent
2. 형식족 $\mathfrak S$. 이는 다음 식으로 정의되는 모든 형식 $g(\lambda;x,y)$로 이루어진다.
\[
  g(\lambda;x,y)=f(x,y,\lambda_1x+\lambda_2y),
\]
이를 편극연산으로 나타내면 다음과 같다.
\[
\begin{aligned}
  g(\lambda;xy)
  &=\frac1{p!}\left[(\lambda_1x+\lambda_2y)\frac{\partial}{\partial z}\right]^p f(x,y,z)\srcfnmark{*)} \\
  &=g_0(xy)\lambda_1^p+g_1(x,y)\lambda_1^{p-1}\lambda_2+\cdots+g_p(xy)\lambda_2^p
\end{aligned}
\]
이때 다음 식이 성립한다.
\[
  i!(p-i)!\,g_i(xy)=
  \left(y\frac{\partial}{\partial z}\right)^i
  \left(x\frac{\partial}{\partial z}\right)^{p-i}f(xyz).
\]
\srcfntext{*)}{I에서와 마찬가지로
\[
 \left(t\frac{\partial}{\partial z}\right)=t_1\frac{\partial}{\partial z_1}+t_2\frac{\partial}{\partial z_2}
\]
를 뜻한다. 따라서 특히
\[
 \left[(\lambda_1x+\lambda_2y)\frac{\partial}{\partial z}\right]
 = (\lambda_1x_1+\lambda_2y_1)\frac{\partial}{\partial z_1}
 + (\lambda_1x_2+\lambda_2y_2)\frac{\partial}{\partial z_2},
\]
\[
 \left[z\left(\frac{\partial}{\partial\lambda_1}\frac{\partial}{\partial x}
 +\frac{\partial}{\partial\lambda_2}\frac{\partial}{\partial y}\right)\right]
 =z_1\left(\frac{\partial}{\partial\lambda_1}\frac{\partial}{\partial x_1}
 +\frac{\partial}{\partial\lambda_2}\frac{\partial}{\partial y_1}\right)
 +z_2\left(\frac{\partial}{\partial\lambda_1}\frac{\partial}{\partial x_2}
 +\frac{\partial}{\partial\lambda_2}\frac{\partial}{\partial y_2}\right).
\]}
따라서 $g_i(xy)$들은 항등식 (1)의 형식 $Z$들을 준다. $g$들은 $x,y,\lambda$와 부정원 $a$에 관한 유리수 계수 형식이다. $\mathfrak S$는 $R$에 관하여 랭크 $\sigma$를 갖는다고 하자.

%
%
\noindent
3. 형식족 $\mathfrak T$. 이는 $\mathfrak S$가 $\mathfrak L$에서 생겨나는 것과 같은 방식으로 $\mathfrak S$에서 역편극연산을 통해 생겨난다. $\mathfrak T$의 형식 $h$는 다음 식으로 정의된다.
\[
\begin{aligned}
 h(x,y,z)
 &=\frac1{p!}\left[z\left(\frac{\partial}{\partial\lambda_1}\frac{\partial}{\partial x}
 +\frac{\partial}{\partial\lambda_2}\frac{\partial}{\partial y}\right)\right]^p g(\lambda;xy)\\
 &=h_0(xyz)+h_1(xyz)+\cdots+h_p(xyz).
\end{aligned}
\]
2에 따라
\[
  h_i(xyz)=
  \left(z\frac{\partial}{\partial x}\right)^{p-i}
  \left(z\frac{\partial}{\partial y}\right)^i g_i(xy)
\]
이다. 그러므로 각 $h_i$와 이에 따라 $h$도 형식 $PZ$ 및 그 합이다. $\mathfrak T$는 $R$에 관하여 랭크 $\tau$를 갖는다고 하자.

%
%
$\mathfrak T$의 형식 $h(x,y,z)$를 구성한 방식에서 알 수 있듯이, 이 형식들은 $\theta$에서 편극연산 $P$를 통해 얻어졌으며 각 변수열 $x,y,z$에서 $\theta$와 같은 차수를 갖는다. 따라서 \emph{형식족 $\mathfrak T$는 $\mathfrak L$의 부분족}을 이룬다. 그러므로 위에서 말한 \emph{사실}에 따라 두 형식족의 동치를 증명하려면 이제 랭크가 같다는 것만 보이면 된다. 이 랭크를 비교할 때에는 $f(xyz)$의 특수한 구조, 곧 모두 같은 형식 $\theta$의 편극형이라는 사실을 더 이상 사용하지 않는다.

%
%
\emph{a) $\mathfrak L$과 $\mathfrak S$의 랭크를 비교하는 데에는 다음 보조정리 a)가 쓰인다. $f(x,y,\lambda_1x+\lambda_2y)=0$이면 반드시 $f(x,y,z)=0$이다.}

이는 세 개의 이진 변수열 사이에 성립하는 항등관계
\[
  (xy)z+(yz)x+(zx)y=0,
  \qquad \text{여기서 }(xy)=(x_1y_2-x_2y_1)\cdots,
\]
에서 곧바로 따른다. 이 관계로부터
\[
  f[x,y,(zy)x+(xz)y]=(xy)^p f(x,y,z)
\]
를 얻는다. 따라서 $f(x,y,\lambda_1x+\lambda_2y)=0$이 $\lambda$에 관한 항등식으로 성립하면, $\lambda_1=(zy)$, $\lambda_2=(xz)$를 대입하고 $(xy)\ne0$임을 이용하여 반드시 $f(x,y,z)=0$임을 얻는다.

%
%
이 보조정리에 따라 \emph{$\mathfrak S$의 형식들과 $\mathfrak L$의 형식들 사이에는 일대일 대응이 존재한다}. 실제로
\[
  f_1(x,y,\lambda_1x+\lambda_2y)=g(\lambda;xy),\qquad
  f_2(x,y,\lambda_1x+\lambda_2y)=g(\lambda;xy)
\]
이면 반드시
\[
  f_1(xyz)=f_2(xyz)
\]
이다.

%
%
더 나아가 \emph{$\mathfrak S$에서 성립하는 모든 관계는 그것과 유일하게 대응하는 $\mathfrak L$의 형식들 사이에서도 보존되며, 그 역도 성립한다}. 왜냐하면
\[
  c_1g^{(1)}(\lambda;xy)+c_2g^{(2)}(\lambda;xy)+\cdots+c_\varkappa g^{(\varkappa)}(\lambda;xy)=0
\]
이면 보조정리에 따라 반드시
\[
  c_1f^{(1)}(x,y,z)+c_2f^{(2)}(x,y,z)+\cdots+c_\varkappa f^{(\varkappa)}(x,y,z)=0
\]
이기 때문이다. \emph{따라서 $\mathfrak L$과 $\mathfrak S$의 랭크가 같음이 증명되었다. 즉 $\rho=\sigma$이다.}

%
%
\emph{b) $\mathfrak S$와 $\mathfrak T$의 랭크를 비교하는 데에는 이에 대응하는 보조정리 b)가 쓰인다. $h(x,y,z)=0$이면 반드시 $g(\lambda;xy)=0$이다.}

이를 증명하기 위해, 3에 따라 $h(x,y,z)=0$이라는 것은 다음과 같은 각 미분연산자 거듭제곱의 곱이 0이 되는 것과 동치임을 생각하자.
\[
 \left(\frac{\partial}{\partial\lambda_1}\frac{\partial}{\partial x_1}
 +\frac{\partial}{\partial\lambda_2}\frac{\partial}{\partial y_1}\right)^{p_1}
 \left(\frac{\partial}{\partial\lambda_1}\frac{\partial}{\partial x_2}
 +\frac{\partial}{\partial\lambda_2}\frac{\partial}{\partial y_2}\right)^{p_2}g(\lambda;x,y),
 \qquad (p_1+p_2=p).
\]

%
%
여기에 다시 미분을 적용하면 다음 연산자곱들도 0이 된다.
\[
\begin{aligned}
&\left(\frac{\partial}{\partial x_1}\right)^{\alpha_1}
 \left(\frac{\partial}{\partial x_2}\right)^{\alpha_2}
 \left(\frac{\partial}{\partial y_1}\right)^{\beta_1}
 \left(\frac{\partial}{\partial y_2}\right)^{\beta_2} \\
&\quad\cdot
 \left(\frac{\partial}{\partial\lambda_1}\frac{\partial}{\partial x_1}
 +\cdots\right)^{p_1}
 \left(\frac{\partial}{\partial\lambda_1}\frac{\partial}{\partial x_2}
 +\cdots\right)^{p_2}
 g(\lambda,xy)
\end{aligned}
\]
따라서 이 연산자곱들의 모든 선형결합도 0이 되며, 그 결과 다음 꼴의 모든 식도 0이 된다.
\[
  f\left(\frac{\partial}{\partial x},\frac{\partial}{\partial y},
  \frac{\partial}{\partial\lambda_1}\frac{\partial}{\partial x}
  +\frac{\partial}{\partial\lambda_2}\frac{\partial}{\partial y}\right)g(\lambda,xy).
\]

%
%
$f$를 특별히
\[
  f(x,y,\lambda_1x+\lambda_2y)=g(\lambda;xy)
\]
가 되도록 택하면
\[
  g\left(\frac{\partial}{\partial\lambda};
    \frac{\partial}{\partial x},
    \frac{\partial}{\partial y}\right)g(\lambda,x,y)=0
\]
도 얻는다. 또는
\[
  g(\lambda,xy)=\sum G_i\lambda_1^{\varkappa_1}\lambda_2^{\varkappa_2}
  x_1^{\varkappa_3}x_2^{\varkappa_4}y_1^{\varkappa_5}y_2^{\varkappa_6}
\]
로 놓으면, 여기서 $G_i$들은 부정원 $a$의 유리수 계수 선형형식이므로,
\[
  \sum \varkappa_1!\varkappa_2!\cdots\varkappa_6!\,G_i^2=0
\]
이고 따라서 $g(\lambda;xy)=0$이다.

%
%
보조정리 b)에서 a)와 똑같은 방식으로 \emph{$\mathfrak S$와 $\mathfrak T$의 랭크가 같음, 즉 $\sigma=\tau$임}이 따른다.

%
%
그러므로 a)와 b)에 의해 $\rho=\tau$이다. 또한 $\mathfrak T$는 $\mathfrak L$의 부분족이므로, 이로써 두 선형 형식족의 동치가 증명되었다. 따라서 $\mathfrak L$의 모든 형식, 특히 $\theta$도 선형독립인 $\rho$개의 형식 $h(x,y,z)$의 유리수 계수 선형결합으로 나타낼 수 있다.
\[
  \theta=\sum c_i h^{(i)}(x,y,z)=\sum PZ.
\]
$\theta$의 부정계수 $a$에 대해 유도한 이 항등식은 부정원 $a$를 임의의 유리성 영역의 원소들로 바꾸어도 그대로 성립한다. 따라서 \emph{세 개의 이진 변수열인 경우에 환원정리가 일반적으로 증명되었다.}

%
%
각각 $n$개의 변수를 갖는 $N$개의 변수열에 대한 일반적인 증명도 완전히 나란히 진행된다. $\theta(A)$가 변수열 $A_i$에서 각각 차수 $\alpha_i$를 갖는다고 하자. 다시 다음 세 선형 형식족을 생각한다.

\noindent
1. 형식족 $\mathfrak L$은 $\theta(A)$에서 편극연산 $P$를 통해 얻어지고 각 개별 변수열 $A_i$에서 $\theta(A)$와 같은 차수를 갖는 모든 형식 $f(A)$로 이루어진다.

%
%
\noindent
2. 형식족 $\mathfrak S$는 다음 대입을 통해 $f(A)$에서 얻어지는 모든 형식 $g(\lambda;A_1\cdots A_n)$으로 이루어진다.
\[
\begin{aligned}
 A_{n+1}&=\lambda_{n+1}^{(1)}A_1+\lambda_{n+1}^{(2)}A_2+
        \cdots+\lambda_{n+1}^{(n)}A_n\\
 &\vdotswithin{=}\\
 A_N&=\lambda_N^{(1)}A_1+\lambda_N^{(2)}A_2+
        \cdots+\lambda_N^{(n)}A_n
\end{aligned}
\]
이때 $g(\lambda;A_1\cdots A_n)$에서 $\lambda$의 각 거듭제곱 단항곱의 계수들은 항등식 (1)의 형식 $Z$들을 준다.

%
%
\noindent
3. 형식족 $\mathfrak T$는 다음 식으로 정의되는 모든 형식 $h(A)$로 이루어진다.
\[
\begin{aligned}
 h(A)=&\left[A_{n+1}\left(
 \frac{\partial}{\partial\lambda_{n+1}^{(1)}}\frac{\partial}{\partial A_1}
 +\cdots+
 \frac{\partial}{\partial\lambda_{n+1}^{(n)}}\frac{\partial}{\partial A_n}\right)\right]^{\alpha_{n+1}} \\
&\cdots
\left[A_N\left(
 \frac{\partial}{\partial\lambda_N^{(1)}}\frac{\partial}{\partial A_1}
 +\cdots+
 \frac{\partial}{\partial\lambda_N^{(n)}}\frac{\partial}{\partial A_n}\right)\right]^{\alpha_N}
 g(\lambda;A_1\cdots A_n).
\end{aligned}
\]
$h(A)$는 형식 $PZ$들의 합이다.

%
%
이제 임의의 $n+1$개 변수열 사이에 성립하는 항등관계 때문에 보조정리 a)는 그대로 성립한다. 보조정리 b)도 마찬가지로 성립하는데, 그 증명에서는 변수의 개수가 전혀 고려되지 않았기 때문이다. 따라서 $\mathfrak L$과 $\mathfrak T$는 같은 랭크를 가지며, $\mathfrak T$가 $\mathfrak L$의 부분족이므로 두 형식족은 동일하다. 그러므로 항등식
\[
  \theta=\sum PZ
\]
은 부정계수에 대해 성립하고, 따라서 임의의 유리성 영역의 원소들을 계수로 하는 경우에도 성립한다. \emph{이로써 환원정리가 일반적으로 증명되었다.}

%
%
\subsection*{III.}
이와 비슷한 고찰을 통해 \emph{일반 급수전개}도 얻을 수 있음을 간단히 보이겠다. 여기서는 $N\le n$이다. $Z$를 각각 $n$개의 원소를 갖는 $n$개 변수열 $A_1,\ldots,A_n$에 대해 개별적으로 동차인 형식이라 하고, $\Delta$를 이 변수열들의 행렬식이라 하자. 먼저 (1)에 대응하는 항등식
\begin{equation}
  Z\equiv\sum PH \mod{\Delta}
\tag{2}
\end{equation}
을 증명한다. 여기서 형식 $H$는 변수열 $A_1\cdots A_{n-1}$만을 포함하고, $Z$에서 연산 $P$를 통해 유도된다.

%
%
증명은 II에서 유도한 관계들을 $\Delta$를 법으로 하는 합동식으로 바꾸는 데에 기초한다. 간단히 하기 위해 \emph{세 개의 삼항 변수열} $x,y,z$를 전제로 하고, 그 계수들은 부정원이라고 하자.

세 선형 형식족 $\mathfrak L,\mathfrak S,\mathfrak T$를 II의 이진 변수열의 경우와 같이 정의하고 그 랭크를 각각 $\rho,\sigma,\tau$라 하자. 한편 $\rho^*$와 $\tau^*$는 $\Delta$를 법으로 한 $\mathfrak L$과 $\mathfrak T$의 랭크를 뜻한다. 즉 $\mathfrak L$ 안에는 $\Delta$를 법으로 선형독립인 형식이 $\rho^*$개 있지만 $\rho^*+1$개는 존재하지 않는다는 뜻이다. 그러면 II에서와 같이 다음 보조정리가 성립한다.

%
%
\emph{보조정리 a): $f(x,y,\lambda_1x+\lambda_2y)=0$이면 반드시 $f(xyz)\equiv0\pmod{\Delta}$이며, 그 역도 성립한다.}

실제로 네 개의 삼항 변수열 사이의 관계
\[
  \Delta_1x-\Delta_2y+\Delta_3z=\Delta t,
  \qquad \text{여기서 }\Delta_1=(yzt)\cdots,
\]
에 따라 세 개의 삼항 변수열 사이에는 언제나 $\Delta$를 법으로 한 선형관계
\[
  \Delta_1x-\Delta_2y+\Delta_3z\equiv0\mod{\Delta}
\]
가 존재한다. 따라서 II에서와 같이
\[
  f(x,y,-\Delta_1x+\Delta_2y)\equiv\Delta_3^p f(x,y,z)\pmod{\Delta}
\]
를 얻는다. $f(x,y,\lambda_1x+\lambda_2y)=0$이 $\lambda$에 관한 항등식으로 성립하면, $\Delta$가 기약이고 $\Delta_3$가 $\Delta$로 나누어지지 않으므로 반드시 $f(x,y,z)\equiv0\pmod{\Delta}$이다. 그 역은 $\Delta(x,y,\lambda_1x+\lambda_2y)=0$이므로 분명하다. 여기서 II와 같이 $\rho^*=\sigma$를 얻는다.

%
%
보조정리 b)는 그 증명과 함께 그대로 성립하므로 $\sigma=\tau$이다.

%
%
$\tau=\tau^*$를 보이기 위해 다음 보조정리를 사용한다.

\emph{보조정리 c): $h(x,y,z)\equiv0\pmod{\Delta}$이면 반드시 $h(x,y,z)=0$이다.}

실제로 $h(x,y,z)$는 편극형이므로 잘 알려진 $\Omega$-연산을 적용하면 소멸한다. 이는 다음 미분방정식으로 나타낼 수 있다.
\[
  \Delta\!\left(\frac{\partial}{\partial x},\frac{\partial}{\partial y},\frac{\partial}{\partial z}\right)h(xyz)
  =\left(\frac{\partial}{\partial x}\frac{\partial}{\partial y}\frac{\partial}{\partial z}\right)h(x,y,z)=0.\srcfnmark{*)}
\]
\srcfntext{*)}{3에 따라
\[
 c\cdot\Omega(h)=\left(
\frac{\partial}{\partial \xi}\frac{\partial}{\partial \eta}
\left(\frac{\partial}{\partial\lambda_1}\frac{\partial}{\partial \xi}
+\frac{\partial}{\partial\lambda_2}\frac{\partial}{\partial \eta}\right)
\right)
\left(x\frac{\partial}{\partial \xi}\right)^{\alpha-1}
\left(y\frac{\partial}{\partial \eta}\right)^{\beta-1}
\left[
z\left(\frac{\partial}{\partial\lambda_1}\frac{\partial}{\partial \xi}
+\frac{\partial}{\partial\lambda_2}\frac{\partial}{\partial \eta}\right)
\right]^{p-1}\cdot g(\lambda;\xi,\eta)
\]
이다. 따라서 행렬식 인자가 소멸하므로 $\Omega(h)$도 소멸한다.}

이제 $h(x,y,z)=\Delta(x,y,z)\cdot k(x,y,z)$라 하면, 다시 미분함으로써 다음 식도 소멸한다.
\[
  k\!\left(\frac{\partial}{\partial x},\frac{\partial}{\partial y},\frac{\partial}{\partial z}\right)
  \Delta\!\left(\frac{\partial}{\partial x},\frac{\partial}{\partial y},\frac{\partial}{\partial z}\right)
  =
  h\!\left(\frac{\partial}{\partial x},\frac{\partial}{\partial y},\frac{\partial}{\partial z}\right)h(x,y,z).
\]
따라서 $h(x,y,z)$ 자체도 소멸한다. 그러므로 $\rho^*=\sigma=\tau=\tau^*$이고, $\mathfrak T$가 $\mathfrak L$의 부분족이므로 \emph{항등식 (2)가 증명되었다}. 이와 같은 고찰로 변수가 $n$개인 경우도 증명된다.

%
%
항등식 (2)를 $\Delta$의 승수에 적용하면 다음 전개를 얻는다.
\[
  Z=\varphi_0+\Delta\varphi_1+\Delta^2\varphi_2+\cdots+\Delta^\mu\varphi_\mu.
\]
각 $\varphi_i$는 $\Omega$-연산을 적용하면 개별적으로 소멸한다. 일반적으로 $\Omega(\Delta\cdot\psi)=c_1\psi+c_2\Delta\cdot\Omega(\psi)$이므로, 이 연산을 $i$번 반복하면
\[
  c\cdot \Omega^i(Z)\equiv\varphi_i\pmod{\Delta}
\]
이다. 한편 (2)에 의해
\[
  c\cdot\Omega^i(Z)\equiv\sum PH_i\pmod{\Delta}
\]
이다. 여기서 $H_i$는 $H$가 $Z$에서 유도되는 것과 같은 방식으로 형식 $\Omega^i(Z)$에서 유도된다. 따라서 변수 $n$개로 확장한 보조정리 c)에 의해 $\varphi_i=\sum PH_i$이고, 그러므로
\begin{equation}
  Z=\sum PH+\Delta\cdot\sum PH_1+\Delta^2\cdot\sum PH_2+
  \cdots\Delta^\mu\cdot\sum PH_\mu
\tag{3}
\end{equation}
이다. 이는 \emph{각각 $n$개의 변수를 갖는 $n$개 변수열의 형식을 $n-1$개 변수열만을 갖는 형식들의 편극형으로 나타내는 가장 일반적인 급수전개}이다.

%
%
$n$개의 변수를 갖는 $N$개 변수열의 형식에 대한 급수전개 $(N<n)$는 잘 알려진 대로 간단한 대입을 통해 여기서 얻는다. $N=3$, $n>3$인 경우
\[
  u=\xi_1x+\xi_2y+\xi_3z,
  \quad v=\eta_1x+\eta_2y+\eta_3z,
  \quad w=\zeta_1x+\zeta_2y+\zeta_3z
\]
로 놓고 $H(u,v,w)=Z(\xi,\eta,\zeta)$를 (3)에 따라 전개한다. $\xi,\eta,\zeta$는 결합 $u,v,w$ 안에서만 나타나므로
\[
  \left(\eta\frac{\partial}{\partial \xi}\right)=
  \left(v\frac{\partial}{\partial u}\right)\cdots,
  \qquad
  \Omega=\sum\left(\frac{\partial}{\partial u}
  \frac{\partial}{\partial v}
  \frac{\partial}{\partial w}\right)_{ikl}(xyz)_{ikl}
  =\nabla_{uvw}(xyz)
\]
이다. $\xi_1=\eta_2=\zeta_3=1$이고 $\xi_2=\xi_3=\cdots=\zeta_2=0$인 특수화에 의해 $H(u,v,w)$는 $H(x,y,z)$로 바뀐다. 또한 $(y\partial/\partial x)$ 등을 행렬식 $(xyz)_{ikl}$에 적용하면 이 행렬식들이 소멸하므로, $\nabla_{uvw}(xyz)$는 수치인자를 제외하면 $\nabla_{xyz}(xyz)=\nabla_{xyz}$로 바뀐다. $N$개의 변수열에서도 같은 사실이 성립한다.

%
%
따라서 (3)으로부터 다음 전개를 얻는다.
\begin{equation}
  H=\sum P\Phi+\sum P\Phi_1+\cdots+\sum P\Phi_\mu
\tag{4}
\end{equation}
여기서 $\Phi_i$는 $\nabla^i_{A_1\ldots A_N}H$에서 연산 $P$를 통해 유도되며 변수열 $A_1,\ldots,A_{N-1}$만을 포함한다. 다만 $\nabla^i$ 안에 나타나는 행렬식 결합들은 예외인데, 앞에서 말했듯이 이들은 편극연산에서 고려되지 않는다.

%
%
따라서 $\nabla^i$ 안의 이 행렬식들을 부정원 $p$로 바꾸면, 그렇게 생기는 형식들을 다시 (4)에 따라 전개할 수 있다. 이로써 마침내 다음 전개를 얻는다.
\begin{equation}
  H=\left[\sum P\psi(p)\right]_{p=\text{$A_1\cdots A_N$의 행렬식들}A_1\cdots A_N}
\tag{5}
\end{equation}
여기서 $\psi$들은 원래 형식에서 연산 $P$와 $\nabla_{A_1\ldots A_N}(p)$, $\nabla_{A_1\ldots A_{N-1}}(p)$ 등을 통해 생겨난다. 이 전개들과 그 결합들, 예를 들어 (1)에 나타나는 형식 $Z$에 (3)을 적용한 다음 (4)와 (5)를 적용하여 얻는 전개들을 사용하면, \emph{각각 $n$개의 동변 변수를 갖는 임의로 많은 변수열의 형식에 대해 알려진 모든 전개를 망라한다}.

%
%
에를랑겐, 1915년 1월 5일.



\clearpage
\setcounter{footnote}{0}

\clearpage
\setcounter{footnote}{0}
\phantomsection
\addcontentsline{toc}{section}{9. 정수적 초월수로 이루어진 가장 일반적인 영역}
%
%

%
%
\section*{9. 정수적 초월수로 이루어진 가장 일반적인 영역}
\begin{center}
{\Large\bfseries 9. 정수적 초월수로 이루어진 가장 일반적인 영역}\par
\vspace{0.75em}
에를랑겐의 에미 뇌터.\par
\vspace{1.2em}
Math. Ann. 77 (1916), 103--128쪽
\end{center}

%
%
E. 체르멜로는 정렬정리를 이용하여 ``정수적 초월수''의 한 영역, 곧 다음과 같은 \emph{추상적} 성질을 지닌 수의 영역 $\mathfrak G$를 구성하였다.
\begin{enumerate}
\item[I.] $\mathfrak G$의 두 수의 합, 차, 곱은 다시 $\mathfrak G$의 수이다.
\item[II.] 모든 실수 또는 복소수는 $\mathfrak G$의 두 수의 몫이다.
\item[III.] 모든 유리정수(또는 대수적 정수)는 $\mathfrak G$의 원소이다.
\item[IV.] 정수가 아닌 유리수(또는 대수적 정수가 아닌 대수적 수)는 어느 것도 $\mathfrak G$에 속하지 않는다.\srcfn{*)}{E. Zermelo, 「Über ganze transzendente Zahlen」, Math. Ann. 75, 434쪽 (1914).}
\end{enumerate}

%
%
이 구성은 정렬정리로 모든 수의 \emph{초월기저}가 존재함을 알 수 있다는 사실에 바탕을 둔다. 다시 말해, 이는 수 $\eta$들의 \emph{계} $H$로서, \emph{그 수들 사이에는 아무런 대수적 관계도 없지만 그 수들을 써서 다른 모든 수를 대수적으로 나타낼 수 있는} 것이다. 이 기저원소 $\eta$들은 대수적으로 독립이므로 \emph{부정원}으로 볼 수 있고, 따라서 찾는 영역은 계수가 모든 대수적 수의 체 $K$에 속하는, 부정원 $\eta$의 \emph{유리함수와 대수함수로 이루어진 정역}이 된다.\srcfn{**)}{이 때문에 본 논문의 고찰은 저자의 앞선 논문 「Körper und Systeme rationaler Funktionen」, Math. Ann. 76, 161쪽 (1915)의 고찰과 부분적으로 나란히 진행된다.} 그러면 체르멜로가 구성한 특수한 영역 $\mathfrak G_\eta$는 다음 \emph{기저 성질}들로 특징지어진다.

%
%
$\mathfrak G_\eta$는 다음을 포함한다.
\begin{enumerate}
\item[1)] 모든 기저원소 $\eta$,
\item[2)] 정수적인\srcfn{***)}{이 글에서 ``정수적''이라는 말은 언제나 계수가 모든 대수적 정수의 정역 $[\mathsf{K}]$에 속한다는 뜻으로 이해한다. $\mathfrak G_\eta$를 정의할 때에는 1)부터 5)까지에서 유리정수 계수만을 고려해도 충분하다. 그러나 더 일반적인 영역에서는 그렇게 하면 기저 성질들이 덜 단순해진다.} 계수를 갖는 $\eta$의 모든 다항식,
\item[3)] 정수적인 계수를 갖는 $\eta$의 모든 정수적 대수함수.
\end{enumerate}
$\mathfrak G_\eta$는 다음을 배제한다.
\begin{enumerate}
\item[4)] 정수적인 계수를 갖는 $\eta$의 모든 분수형 유리함수,
\item[5)] 정수적인 계수를 갖는 $\eta$의 모든 분수형 대수함수.\srcfn{*)}{Zermelo, 앞의 글, § 3.}
\end{enumerate}

%
%
그런데 체르멜로의 영역 $\mathfrak G_\eta$와 \emph{본질적으로 다른} 영역 $\mathfrak G$가 존재한다는 사실은 쉽게 드러난다(§ 2와 § 3의 예 참조). 즉, \emph{어떠한 정렬에서도} 기저 성질 1)부터 5)까지를 모두 갖지 않아, 따라서 어떠한 정렬에서도 체르멜로의 구성으로 만들어 낼 수 없는 영역 $\mathfrak G$가 존재한다. 달리 말하면, $\mathfrak G_\eta$ 위로 \emph{일대일이며 동형인} 사상을 갖지 않는 영역 $\mathfrak G$가 존재한다. 이로써 다음에서 완전히 답할 \emph{정수적 초월수로 이루어진 가장 일반적인 영역}에 관한 문제가 생긴다.

%
%
이를 위해서는 먼저 어느 기저 성질이 추상적인 정의 성질의 귀결이고, 어느 것이 특수한 구성에서 비롯되는지를 확정해야 한다. 모든 주어진 영역 $\mathfrak G$에 대하여 기저 성질 1)과 2)가 성립하도록 하는 정렬이 존재한다는 것이 § 1에서 드러난다. 따라서 \emph{모든 영역 $\mathfrak G$는 모든 수의 초월기저를 이용하여 구성할 수 있다}. 나머지 기저 성질들은 어느 것도 반드시 성립할 필요가 없다(§ 2와 § 3). 특히 체르멜로 영역의 또 하나의 추상적 성질인 \emph{대수적 정수성에 대한 폐쇄성}은 모든 영역 $\mathfrak G$에 성립하지 않는다. 이 성질은 다음과 같이 표현할 수 있다.
\begin{enumerate}
\item[V.] $\mathfrak G$의 유한개의 수에 대하여 대수적으로 정수이고, 그 종속방정식의 계수가 정수적인 모든 수는 $\mathfrak G$에 속한다.
\end{enumerate}

%
%
I--IV와 더불어 V도 만족하는 특수한 영역을 \emph{대수적으로 정수인 초월수}의 영역 $\mathfrak H$라 부르자. 가장 일반적인 영역 $\mathfrak H$에 대해서는 다음의 단순한 결과가 얻어진다(§ 4).
\begin{enumerate}
\item[a)] \emph{같은 기저 $\mathsf{H}$를 이용하여 구성할 수 있는 모든 영역 $\mathfrak H$의 교집합은 그 기저의 모든 정수적 대수함수로 주어진다. 따라서 그 교집합은 체르멜로 영역 $\mathfrak G_\eta$이며, 이 영역 자체도 영역 $\mathfrak H$이다.} (이로써 동시에 체르멜로 영역의 추상적 특징짓기도 얻어진다.)
\item[b)] \emph{모든 영역 $\mathfrak H$는 임의의 계 $\mathfrak S(\eta)$를 $\mathfrak G_\eta$에 대수적으로 정수인 방식으로 첨가하여 얻어진다.\srcfn{**)}{이 표현의 설명은 § 4에 있다.} 여기서 $\mathfrak S(\eta)$가 만족해야 하는 유일한 조건은 그 첨가로 어떠한 분수형 대수수도 그 영역 안으로 들어오지 않는다는 것이다.}
\end{enumerate}

%
%
초월기저를 바탕에 둘 때 가장 일반적인 영역 $\mathfrak G$에 관한 결과는 그보다 덜 단순하다. 이 경우 모든 영역 $\mathfrak G$의 교집합은 그 자체가 영역 $\mathfrak G$가 아니며, 그 결과 가장 일반적인 영역의 구성도 한눈에 파악하기가 더 어렵다(§ 5).

%
%
영역 $\mathfrak H$에 관한 것과 비슷한 결과를 얻기 위해 초월기저를 \emph{유리기저} $\Theta$로 바꾼다. 즉, 이는 모든 수를 유리적으로 나타낼 수 있게 하는 수 $\vartheta$들의 \emph{계} $\Theta$이면서, 적어도 하나의 정렬순서에서는 어떠한 기저원소도 유한개의 앞선 기저원소를 써서 유리적으로 나타낼 수 없는 것이다. 이 유리기저 $\Theta$에는 $\vartheta$의 모든 정수 계수 다항식으로 이루어진 정역 $[\Theta]$가 어떠한 분수형 대수수도 포함하지 않는다는 부가조건도 부과해야 한다. (이 부가조건을 만족하는 기저의 구성은 § 7에서 보인다.) 그러면 \emph{유리적으로 정수인 초월수}의 영역이라고도 부를 수 있는 가장 일반적인 영역 $\mathfrak G$에 대해서도 다시 다음의 단순한 결과가 얻어진다(§ 6).
\begin{enumerate}
\item[a)] \emph{같은 유리기저 $\Theta$를 이용하여 구성할 수 있는 모든 영역 $\mathfrak G$의 교집합은 그 기저의 모든 다항함수로 주어진다. 따라서 그 교집합은 영역 $[\Theta]$이며, 이 영역 자체도 영역 $\mathfrak G$이다.}
\item[b)] \emph{모든 영역 $\mathfrak G$는 임의의 계 $\mathfrak S(\vartheta)$를 $[\Theta]$에 유리적으로 정수인 방식으로 첨가하여 얻어진다. 여기서 $\mathfrak S(\vartheta)$가 만족해야 하는 유일한 조건은 그 첨가로 어떠한 분수형 대수수도 그 영역 안으로 들어오지 않는다는 것이다.}
\end{enumerate}

%
%
영역 $\mathfrak H$와 완전한 유사성을 얻으려면 영역 $\mathfrak G$의 정의 성질 III과 IV에 대수적 수를 포함하지 않아야 한다. III과 IV를 이렇게 바꾸면 유리기저를 이용한 정수적 원소의 구성을 \emph{임의의} 체로 옮길 수 있지만(§ 10), 체르멜로의 구성은 그 체가 대수적으로 닫혀 있음을 전제로 한다.

%
%
\emph{$\mathsf{H}$가 모든 초월기저를, 그리고 $\Theta$가 부가조건을 만족하는 모든 유리기저를 각각 두루 취하면, 결과 a)와 b)로 모든 영역 $\mathfrak H$와 $\mathfrak G$의 전체가 얻어진다.} 서로 동형인 모든 영역을 하나의 동형류로 묶으면, 이 전체에서 부분집합 하나를 골라 § 9에 따라 적어도 가산무한개인 모든 동형류를 대표하게 할 수 있다(§ 8).

%
%
\subsection*{§ 1. 모든 영역 $\mathfrak G$에 대한 기저 성질 1)과 2)의 증명}

%
%
$\mathfrak G$를 추상적 성질 I--IV를 만족하는, 정수적 초월수로 이루어진 임의의 주어진 영역이라 하자. E. 체르멜로가 모든 복소수의 체에 적용한 절차에 따라 $\mathfrak G$에서 \emph{$\mathfrak G$의 모든 수의 초월기저} $\mathsf{H}$를 골라낸다. II에 따르면 \emph{모든} 실수 또는 복소수를 $\mathfrak G$의 두 수의 몫으로 나타낼 수 있으므로, $\mathsf{H}$는 동시에 \emph{모든} 수의 초월기저이다. 따라서 주어진 모든 영역 $\mathfrak G$에 대하여 모든 수의 초월기저를 $\mathfrak G$에 속하도록 선택할 수 있고, 그 결과 \emph{기저 성질 1)이 성립한다.} 또한 III에 따라 $\mathfrak G$는 모든 대수적 정수의 정역 $[\mathsf{K}]$를 포함하며, 따라서 I에 따라 기저원소 $\eta$에 관한, 계수가 $[\mathsf{K}]$에 속하는 모든 다항식, 즉 $\eta$의 모든 정수 계수 다항식도 포함한다. 이 다항식들은 정역 $[\mathsf{H}]$를 이룬다. 그러므로 \emph{기저 성질 2)도 언제나 성립한다}. 다시 말해, \emph{모든 영역 $\mathfrak G$는 모든 수의 초월기저 $\mathsf{H}$를 이용하여, $\mathfrak G$가 정역 $[\mathsf{H}]$를 부분정역으로 포함하도록 구성할 수 있다.}

%
%
\textbf{주:} 그렇더라도 기저 $\mathsf{H}$에서 출발하여, 바로 $H$에 대해서가 아니라 다른 어떤 기저 $\mathsf{H}'$에 대해서 기저 성질 1)과 2)를 갖는 영역 $\mathfrak G$를 구성할 수 있음은 물론이다. 예를 들어 체르멜로 영역 $\mathfrak G_\eta$의 전체 구성에서 어떤 기저원소 $\eta$를 다항식 $f(\eta)$로 바꾸어 영역 $\mathfrak G'$를 얻는다고 하자. $\mathfrak G'$는 $\eta$도 $[\mathsf{H}]$도 포함하지 않는다. 그러나 기저 $\mathsf{H}$에서 기저원소 $\eta$를 $\zeta=f(\eta)$로 바꾸면, $\mathsf{H}$는 다시 초월기저 $\mathsf{H}'$가 되며, $\mathfrak G'$에 대해서는 $\mathsf{H}'$에 관한 기저 성질 1)과 2)가 성립한다.

%
%
\subsection*{§ 2. 기저 성질 4)와 5)의 배제}

%
%
이제 영역 $\mathfrak G$가 나머지 기저 성질 가운데 어느 것도 반드시 만족할 필요가 없음을 보이겠다. 먼저 체르멜로의 구성에서 정역 $[\mathsf{H}]$를 정수적 유리 ``함수식''으로 이루어진 영역으로 바꿈으로써 4)와 5)가 성립하지 않게 한다.

%
%
베버에 따르면,\srcfn{*)}{H. Weber, \emph{Lehrbuch der Algebra}, Kleine Ausgabe, § 96 및 97.} \emph{유리 함수식}이란 유리수 계수를 갖는 유리함수를 다음과 같이 유일하게 정해지는 꼴로 나타낸 것을 말한다.
\[
 A(\eta_1\cdots \eta_t)
   = a\cdot\frac{E_1(\eta_1\cdots \eta_t)}{E_2(\eta_1\cdots \eta_t)}.
\]
여기서 $E_1$, $E_2$는 서로소인 원시다항식으로서 유리정수 계수를 가지며, $a$는 양의 유리수, 곧 $A$의 절댓값을 뜻한다. 특히 $a$가 \emph{유리정수}이면 $A$를 \emph{정수적 유리 함수식}이라 한다. 정수적 유리 함수식의 특수한 경우에는 유리정수와 유리정수 계수의 다항식이 포함된다. 반면 정수가 아닌 유리수와 정수가 아닌 유리수 계수의 다항식은 분수형 유리 함수식으로 분류해야 한다.

%
%
\emph{이제 영역 $\mathfrak G$가 각각 유한개의 기저원소 $\eta$에 관한 모든 정수적 유리 함수식의 전체 $\mathfrak J$와, $\mathfrak J$에 대하여 대수적으로 정수인 모든 함수로 이루어진다고 하자.} 여기서 후자의 함수란 최고차항의 계수가 $1$이고 나머지 계수들이 정수적 유리 함수식인 어떤 방정식을 만족하는 함수이다.

%
%
$\mathfrak G$가 성질 I--IV를 만족한다는 증명은 체르멜로의 해당 증명과 정확히 나란히 진행되므로 간단히만 적겠다. 우선 $\mathfrak G$가 체르멜로 영역 $\mathfrak G_\eta$를 부분정역으로 포함하므로 II와 III은 명백하다. 또한 유리정수 계수 다항식에 관한 가우스 정리에 따르면 정수적 유리 함수식들의 합, 차, 곱도 다시 \emph{정수적} 유리 함수식이다. 따라서 $\mathfrak J$는 정역을 이룬다. 알려진 논법에 따라\srcfn{*)}{예컨대 Weber, § 97, 8 참조.} $\mathfrak G$도 정역이 되므로 I이 성립한다. 더 나아가 확장된 가우스 정리로부터\srcfn{**)}{Weber, § 20, 5.} 다음 두 보조정리가 따른다. ``$z$가 \emph{어떤} 방정식을 통해 $\mathfrak J$에 대하여 대수적으로 정수이면, 모든 유리 함수식의 체에 대하여 $z$가 만족하는 \emph{기약}방정식을 통해서도 그러하다.''\srcfn{***)}{Weber, § 97, 4 참조.} 그리고 ``어떤 대수적 수가 만족하는, 모든 유리수의 체에 대하여 기약인 방정식은 모든 유리 함수식의 체에 대해서도 기약이다.'' 그러므로 $\mathfrak G$는 분수형 대수수를 포함할 수 없다. 이로써 IV와 따라서 추상적 성질 I--IV가 모두 성립함이 증명되었다.

%
%
\emph{그러나 어떠한 정렬에서도 $\mathfrak G$로부터, $\mathfrak G$가 기저원소들의 모든 분수형 유리함수와 분수형 대수함수를 배제하도록 기저를 골라낼 수는 없다.} 실제로 영역의 어떤 함수 $z(\eta)$를 기저원소로 고른다고 하자. 이 함수는 다음 방정식으로 정의된다.
\[
 z^k+A_1(\eta)z^{k-1}+\cdots+A_k(\eta)=0,
\]
여기서
\[
 A_i(\eta)=a_i\cdot\frac{E_1(\eta)}{E_2(\eta)}
\]
는 정수적 유리 함수식이다. 그러면 함수 $\frac{a_k}{z}$는 $\mathfrak G$에 속하며, $\frac{a_k}{\sqrt z}$, $\frac{a_k}{\sqrt[3]z}$ 등도 마찬가지이다.\srcfn{*)}{완전히 같은 방식으로, $\eta$의 모든 대수함수에 어떤 유리정수를 곱하여 함수식 영역 $\mathfrak G$의 원소로 만들 수 있다. 따라서 이 영역은 \emph{모든 복소수를 $\mathfrak G$의 한 원소와 한 유리정수의 몫으로 나타낼 수 있다}는 성질을 갖는다. 이는 대수적 수의 경우와 유사하다.} 따라서 \emph{영역 $\mathfrak G$ 전체에 대하여 기저 성질 4)와 5)가 배제된다.}

%
%
\textbf{주:} § 3의 한 예는 $\mathfrak G$가 분수형 유리수나 분수형 대수수를 포함하지 않으면서도, 모든 정렬에서 분수형 유리 함수식을 포함할 수 있음을 보일 것이다.

%
%
\subsection*{§ 3. 기저 성질 3)의 배제}

%
%
기저 성질 3)을 배제하기 위해 일반화된 \emph{모듈 영역}을 사용한다. 이 영역에서는 영역의 다항식에 들어가는 인수가 더 이상 부정원 자체가 아니라 부정원들의 모든 정수적 대수함수이며,\srcfn{**)}{여기서 ``정수적 대수함수''는 언제나 \emph{정수적인} 계수를 갖는 것을 뜻한다. 즉, 최고차항의 계수가 $1$이고 나머지 계수들이 대수적 정수 계수를 갖는 $\eta$의 다항식, 곧 $[\mathsf{H}]$의 다항식인 어떤 방정식을 만족하는 함수이다. 특히 $[\mathsf{H}]$의 모든 다항식은 정수적 대수함수에 속한다.} 부정원 자체는 모듈 기저의 역할을 맡는다.

%
%
$\mathfrak G$는 다음과 같은 \emph{모든 양}으로 이루어진다고 하자.
\[
\begin{array}{@{}l@{\quad}c@{}}
(1)& z=\alpha+g_1(\eta)\cdot\eta_1+g_2(\eta)\cdot\eta_2+\cdots+g_t(\eta)\cdot\eta_t.
\end{array}
\]
\emph{여기서 $\alpha$는 모든 대수적 정수를, $\eta_1\cdots\eta_t$는 모든 기저원소를 두루 취하며, $g_1(\eta)\cdots g_t(\eta)$는 $\eta_1\cdots\eta_t$의 각 고정된 조합마다 서로 독립적으로 $\eta$의 모든 정수적 대수함수를 두루 취한다.}\srcfn{***)}{모듈 성질은 양 $z-\alpha$에 성립한다.}

%
%
이 모듈 영역이 성질 I--IV를 만족함은 쉽게 증명할 수 있다. 우선 꼴 (1)의 두 양을 더하거나 빼거나 곱해도 다시 꼴 (1)의 양이 되므로 I이 성립한다. 또한 $\mathfrak G$는 모든 기저원소 $\eta$를 포함하고, $g(\eta)$가 $\eta$의 모든 정수적 대수함수를 두루 취할 때 모든 양 $\eta\cdot g(\eta)$도 포함한다. 따라서 $\eta$의 모든 \emph{정수적} 대수함수, 나아가 $\eta$의 모든 대수함수를 $\mathfrak G$의 두 양의 몫으로 나타낼 수 있으므로 II가 증명된다. 더 나아가 표현 (1)은 특히 $g_1=0\cdots g_t=0$으로 놓을 때 모든 대수적 정수를 포함한다. 그러나 $z$는 어떤 분수형 대수수와도 같아질 수 없다. 실제로 $z$가 대수적 수 $\beta$를 나타낸다면
\[
 \beta=\alpha+g_1(\eta)\eta_1+\cdots+g_t(\eta)\eta_t
\]
가 $\eta$에 관하여 항등적으로 성립하므로 반드시 $\beta=\alpha$이다. 이는 다음 특수화에서 알 수 있다.
\[
 \eta_1=0\cdots \eta_t=0.
\]
이때 승수 $g_1(\eta)\cdots g_t(\eta)$는 정수적 대수함수이므로 유한한 값을 유지한다. 곧 정수적 대수함수나 수로 바뀐다.\srcfn{*)}{IV에 대한 이 더 자세한 증명은 이 절 끝의 확장된 예 때문에 제시하였다. 여기서는 $z$가 구성상 정수적 대수함수라는 지적만으로도 충분하다.} 따라서 $\mathfrak G$는 조건 I--IV를 만족한다.

%
%
\emph{반대로, 어떠한 정렬에서도 기저원소들의 모든 정수적 대수함수가 $\mathfrak G$에 속하도록 $\mathfrak G$에서 기저를 골라낼 수는 없다.} 실제로 영역에서 기저원소로 고를 임의의 함수 $z(\eta)$를 생각하자. 이 함수는 부정원 $\eta$를 실제로 포함해야 한다. 그러면
\[
 \xi=\sqrt[\nu]{z-\alpha}
\]
로 주어지는 함수들은 $\nu$가 $z$에 따라 정해지는 어떤 유한한 경계를 넘으면 더 이상 $\mathfrak G$에 속할 수 없음을 보이겠다.

%
%
가령 $\xi$가 $\mathfrak G$에 속한다고 하자. 그러면 (1)에 따라 다음 꼴이 된다.
\[
 \zeta=\gamma+h_1(\eta)\eta_{i_1}+h_2(\eta)\eta_{i_2}
      +\cdots+h_\tau(\eta)\eta_{i_\tau}.
\]
여기서 $h_1(\eta)\cdots h_\tau(\eta)$는 다시 $\eta$의 정수적 대수함수이고, $\eta_{i_1}\cdots\eta_{i_\tau}$는 임의의 기저원소로서 특히 $\eta_1\cdots\eta_t$와 일치할 수도 있다. 먼저 대수적 정수 $\gamma$가 $0$임을 보인다. $h_1(\eta)\cdots h_\tau(\eta)$가 정수적 대수함수이므로, 나머지 $\eta$의 값이 어떠하든 $\eta_{i_1}=0\cdots\eta_{i_\tau}=0$이면 $\xi=\gamma$이다. 특히 다음과 같이 놓아도 $\xi=\gamma$이다.
\[
 \eta_{i_1}=0\cdots \eta_{i_\tau}=0; \qquad
 \eta_1=0\cdots \eta_t=0.
\]
그러나 이 특수화에서는 (1)에 따라 함수 $(z-\alpha)$가 사라지고, 따라서 $\xi$도 사라진다. 그러므로 $\gamma=0$이며
\[
 \zeta=h_1(\eta)\eta_{i_1}+h_2(\eta)\eta_{i_2}
      +\cdots+h_\tau(\eta)\eta_{i_\tau}
\]
를 얻는다.

%
%
이를 $\nu$제곱하면
\[
 \zeta^\nu=z-\alpha=F_\nu(\eta)
\]
를 얻는다. 여기서 $F_\nu(\eta)$는 $\eta_{i_1}\cdots\eta_{i_\tau}$에 관한 $\nu$차 동차형식으로서 항등적으로 사라지지 않으며, 그 계수들은 $\eta$의 정수적 대수함수이다. 한편 (1)에 따라 $\eta$의 \emph{정수적} 대수함수인 $z-\alpha$는 $[\mathsf{H}]$에 대하여 기약인 다음 방정식을 만족한다.
\[
 (z-\alpha)^\varkappa+B(\eta)(z-\alpha)^{\varkappa-1}+\cdots+C(\eta)=0.
\]
여기서 마지막 계수 $C(\eta)$는 $(z-\alpha)$와 그 켤레들의 곱인 다항식이며, 그 \emph{차수}를 $\lambda$라 하자. 다른 한편 $z-\alpha=F_\nu(\eta)$에서 $F_\nu(\eta)$와 그에 대응하는 켤레들을 곱하면
\[
 C(\eta)=G_{\nu\varkappa}(\eta)
\]
를 얻는다. 여기서 $G_{\nu\varkappa}(\eta)$는 $\eta_{i_1}\cdots\eta_{i_\tau}$에 관한 $\nu\varkappa$차 동차형식이고, 그 계수들은 $\eta$의 다항식이다. 따라서 $G_{\nu\varkappa}$의 $\eta$에 관한 차수는 적어도 $\nu\varkappa$이므로 $\lambda\geqq\nu\varkappa$이다. 다시 말해, \emph{$\nu>\frac{\lambda}{\varkappa}$인 함수 $\sqrt[\nu]{z-\alpha}$는 $\mathfrak G$에 속하지 않는다. 따라서 영역 $\mathfrak G$들의 전체에 대하여 기저 성질 3)이 배제된다.}

%
%
\textbf{주:} 방금 구성한 영역을 조금 일반화하면 어떠한 정렬에서도 \emph{분수형 함수식}을 포함하는 영역 $\mathfrak G$를 얻는다. $\mathfrak G$가 다음과 같은 모든 양으로 이루어진다고 하자.
\[
 z=\alpha+\gamma_1(\eta)\eta_1+\gamma_2(\eta)\eta_2
      +\cdots+\gamma_t(\eta)\eta_t.
\]
다만 이제 $\gamma(\eta)$는 $\eta$의 대수적으로 정수이되 반드시 \emph{정수적인 계수}를 갖지는 않는 모든 함수를 두루 취한다. 성질 I--IV의 증명은 위와 같다. 그런데 $\mathfrak G$는 각 함수 $z$와 더불어 $a\cdot(z-\alpha)$도 포함하며, 여기서 $a$는 임의의 유리수 또는 분수형 대수수이다. 따라서 어떠한 정렬에서도 분수형 함수식이 나타난다.

%
%
\subsection*{§ 4. 대수적으로 정수인 초월수로 이루어진 가장 일반적인 영역}

%
%
다음을 더 편리하게 서술하기 위해 다음 표현을 도입하자. ``\emph{양 $z$는 $\mathfrak T$에 대하여 대수적으로 정수이다}.'' 이는 $\mathfrak T$가 임의로 주어진 영역일 때, $z$가 최고차항의 계수는 $1$이고 나머지 계수들은 $\mathfrak T$의 양들로 된 정수 계수 다항식인 \emph{어떤} 방정식을 만족한다는 뜻이다.\srcfn{*)}{$\mathfrak T$가 정역이면 최고차항의 계수는 $1$이고 나머지 계수들은 $\mathfrak T$의 양이 된다.}

%
%
영역 $\mathfrak T$가 다음 조건을 만족하면 이를 \emph{대수적 정수성에 관하여 닫혀 있다}고 한다(서론 참조).
\begin{enumerate}
\item[V.] $\mathfrak T$에 대하여 대수적으로 정수인 모든 양은 $\mathfrak T$에 속한다.
\end{enumerate}

%
%
먼저 체르멜로 영역 $\mathfrak G_\eta$가 \emph{추상적 성질} V를 가짐을 보이겠다. 어떤 양 $z$가 방정식 $f(z)=0$을 통하여 $\mathfrak G_\eta$에 대하여 대수적으로 정수인 것으로 하자. $f(z)$에 $[\mathsf{H}]$에 관한 켤레들을 곱하면 $z$가 $[\mathsf{H}]$에 대해서도 대수적으로 정수임을 알 수 있으므로, $z$는 $\mathfrak G_\eta$에 속한다. § 2에서 살펴본 함수식 영역도 마찬가지로 대수적 정수성에 관하여 닫혀 있다.

%
%
반면 § 3에서 살펴본 모듈 영역은 기저 성질 3)을 갖지 않으므로 성질 V가 모든 영역 $\mathfrak G$에 성립하지는 않음을 보여 준다. 더 정확히 말하면, § 3에서는 그 모듈 영역이 영역의 어떠한 \emph{초월 원소}에 관해서도 대수적 정수성에 관하여 닫혀 있지 않음이 드러났다. 물론 III과 IV에 따라 영역의 \emph{각 대수적 원소}에 관해서는 닫혀 있다.

%
%
따라서 먼저 영역 $\mathfrak G$들의 전체 가운데 추상적 성질 V까지 갖는 것들을 골라내고, 이를 ``\emph{대수적으로 정수인 초월수로 이루어진 영역 $\mathfrak H$}''라 부르기로 한다.

%
%
$\mathfrak H$를 그러한 영역이라 하고, $\mathsf{H}$를 $\mathfrak H$의 모든 수의 초월기저라 하자. II에 따라 $\mathsf{H}$는 동시에 \emph{모든} 수의 초월기저이다. III, I, V에 따라 $\mathfrak H$는 $\mathsf{H}$와 더불어 $\eta$의 모든 정수적 대수함수, 곧 체르멜로 영역 $\mathfrak G_\eta$도 포함한다. 따라서 § 1과 마찬가지로 다음이 성립한다. \emph{대수적으로 정수인 초월수로 이루어진 모든 영역 $\mathfrak H$는 모든 수의 초월기저를 이용하여, 체르멜로 영역 $\mathfrak G_\eta$를 부분정역으로 포함하도록 구성할 수 있다.}

%
%
이제 $\mathsf{H}$를 모든 수의 임의의 초월기저라 하고, $\mathsf{H}$를 포함하는 모든 영역 $\mathfrak H$의 전체를 $\mathfrak M(\mathfrak H)$로 나타내자. 방금 증명한 바에 따라 $\mathfrak M(\mathfrak H)$의 각 영역 $\mathfrak H$는 체르멜로 영역 $\mathfrak G_\eta$를 부분정역으로 포함한다. 또한 $\mathfrak G_\eta$ 자체도 영역 $\mathfrak H$이므로, $\mathfrak G_\eta$는 $\mathfrak M(\mathfrak H)$의 모든 영역 $\mathfrak H$의 \emph{최대 공통 부분정역, 곧 교집합}이다. 이로써 $\mathfrak G_\eta$가 추상적으로 정의된다. 동시에 $\mathfrak M(\mathfrak H)$가 교집합을 취하는 연산에 관하여 닫혀 있음도 곧바로 따른다. 즉, $\mathfrak M(\mathfrak H)$의 어떠한 부분계에 속하는 모든 영역 $\mathfrak H$의 교집합도 $\mathfrak M(\mathfrak H)$에 속한다. 실제로 교집합에는 I과 V가 성립하고, $\mathfrak G_\eta$가 부분정역이므로 II와 III도 성립하며, 그 교집합이 $\mathfrak H$의 부분정역이므로 IV도 성립한다.

%
%
§ 2의 함수식 영역 예가 보여 주듯이, 영역 $\mathfrak H$는 일반적으로 $\mathfrak G_\eta$ 외에도 다른 함수들을 포함한다. $\psi(\eta)$를 그러한 $\eta$의 유리함수 또는 대수함수라 하자. 그러면 I에 따라 $\mathfrak H$는 $\psi$와 $\mathfrak G_\eta$의 함수들 가운데 각각 유한개를 택해 만든 모든 다항식 결합도 포함한다. 이렇게 생기는 정역, 곧 계수가 $\mathfrak G_\eta$에 속하는 $\psi$의 모든 다항식으로 이루어진 정역을 $[\mathfrak G_\eta,\psi]$로 나타내고, 이 과정을 ``\emph{$\mathfrak G_\eta$에 $\psi$를 다항식적으로 첨가하는 것}''이라 하자. 더 나아가 V에 따라 $[\mathfrak G_\eta,\psi]$에 대하여 대수적으로 정수인 모든 함수도 $\mathfrak H$에 속한다. 이렇게 생기는 영역을 $\{\mathfrak G_\eta,\psi\}$로 나타내고, 이 과정을 ``\emph{$\mathfrak G_\eta$에 $\psi$를 대수적 정수로 첨가하는 것}''이라 하자. $\{\mathfrak G_\eta,\psi\}$는 계수가 $\mathfrak G_\eta$에 속하는 $\psi$의 모든 정수적 대수함수로 이루어지므로, 알려진 논법에 따라\srcfn{*)}{예컨대 Weber, § 97, 6과 7 참조.} 정역이다.

%
%
$\{\mathfrak G_\eta,\psi\}$도 대수적 정수성에 관하여 닫혀 있어 조건 V를 만족함을 보이겠다. 실제로 $z$가 다음 방정식을 통하여 $\{\mathfrak G_\eta,\psi\}$에 대하여 대수적으로 정수인 것으로 하자.
\[
 f(z;a\cdots c)=z^\varkappa+a z^{\varkappa-1}+\cdots+c=0.
\]
여기서 $a\cdots c$는 $\{\mathfrak G_\eta,\psi\}$의 양으로서 계수가 $[\mathfrak G_\eta,\psi]$에 속하는 다음 방정식들을 만족한다.
\[
\begin{gathered}
 a^\lambda+A_1a^{\lambda-1}+\cdots+A_\lambda=0,\\
 \vdots\\
 c^\mu+C_1c^{\mu-1}+\cdots+C_\mu=0.
\end{gathered}
\]
그 근들을 $a,a'\cdots a^{(\lambda-1)}\cdots c,c'\cdots c^{(\mu-1)}$라 하자. 이제 근들의 모든 조합에 걸쳐 곱을 만들면
\[
 g(z)=f(z;a\cdots c)f(z;a'\cdots c')\cdots
      f(z;a^{(\lambda-1)}\cdots c^{(\mu-1)})
\]
를 얻는다. $g(z)$의 최고차항 계수는 $1$이고 나머지 계수들은 $[\mathfrak G_\eta,\psi]$의 양이므로, $z$는 $\{\mathfrak G_\eta,\psi\}$에 속한다.\srcfn{*)}{따라서 \emph{어떤} 방정식에 의해 대수적으로 정수인 것을 정의하면, I과 V를 증명할 때 기초체에 관한 기약성 개념이 필요하지 않다. § 2도 참조하라. 거기서는 IV를 증명하기 위해서만, 임의의 방정식에 의한 대수적 정수성의 정의에서 기약방정식으로 넘어가는 보조정리가 필요하다. 이 보조정리는 현재의 경우 일반적으로 더 이상 성립하지 않으므로, IV를 함수 $\psi$에 대한 별도의 조건으로 부과해야 한다.}

%
%
따라서 영역 $\{\mathfrak G_\eta,\psi\}$는 성질 I과 V를 가지며, $\mathfrak G_\eta$가 그 부분정역이므로 II와 III도 가진다. IV의 성립 여부는 $\psi$의 선택에 달려 있다. 예를 들어 $\psi=\frac{1}{n\cdot\eta}$이면 $\frac1n$이 영역에 속하므로 IV가 성립하지 않는다. 반면 § 2에 따라 $\psi=\frac1\eta$일 때나 $\psi=\frac{\eta}{n}$일 때에는 IV가 성립한다. 실제로 후자의 경우 $\{\mathfrak G_\eta,\psi\}$의 어떤 함수가 대수적 수를 나타내면 $\eta=0$으로 놓아도 그 값이 유지된다. 그러면 그 함수는 $\mathfrak G_\eta$의 함수로 바뀌고, 따라서 대수적 정수로 바뀐다.

%
%
\emph{그러므로 대수적 정수 첨가를 통해 어떠한 분수형 대수수도 영역에 들어오지 않는다는 조건을 $\psi$에 부과하면, $\{\mathfrak G_\eta,\psi\}$는 영역 $\mathfrak H$가 된다.}

%
%
$\eta$의 유리함수 또는 대수함수로 이루어진 임의의 계 $\mathfrak S$를 $\mathfrak G_\eta$에 다항식적으로 또는 대수적 정수로 첨가하는 것도 완전히 같은 방식으로 정의한다. 다항식적 첨가는 정역 $[\mathfrak G_\eta,\mathfrak S]$를 낳는다. 이는 $\mathfrak G_\eta$와 $\mathfrak S$에서 각각 유한개씩 택한 양으로 만든 모든 정수 계수 다항식으로 이루어진다. 대수적 정수 첨가는 영역 $\{\mathfrak G_\eta,\mathfrak S\}$를 낳는다. 이는 $[\mathfrak G_\eta,\mathfrak S]$에 대하여 대수적으로 정수인 모든 양으로 이루어진다. 위와 글자 그대로 같은 방식으로 정역 $\{\mathfrak G_\eta,\mathfrak S\}$가 대수적 정수성에 관하여 닫혀 있고 II와 III도 만족함을 알 수 있다. 따라서 위와 마찬가지로 다음이 성립한다.

\emph{대수적 정수 첨가를 통해 어떠한 분수형 대수수도 영역에 들어오지 않는다는 조건을 $\mathfrak S$에 부과하면, $\{\mathfrak G_\eta,\mathfrak S\}$는 영역 $\mathfrak H$가 된다.} 다시 말해 $\mathfrak S$는 허용 가능한 계가 된다.\srcfn{*)}{더 일반적으로도 같은 방식으로 다음을 알 수 있다. $\mathfrak T$를 임의의 영역이라 하고 $\{\mathfrak T\}$가 $\mathfrak T$에 대하여 대수적으로 정수인 모든 양의 전체를 나타낸다면, $\{\mathfrak T\}$는 대수적 정수성에 관하여 닫혀 있다.}

%
%
이제 그 역도 성립한다는 사실이 중요하다. 따라서 $\mathfrak S$가 모든 허용 가능한 계를 두루 취하면 영역 $\{\mathfrak G_\eta,\mathfrak S\}$가 실제로 $\mathfrak M(\mathfrak H)$의 모든 영역 $\mathfrak H$를 빠짐없이 준다. 실제로 $\mathfrak H$를 $\mathfrak M(\mathfrak H)$의 임의의 영역이라 하고, $\mathfrak L=\mathfrak H-\mathfrak G_\eta$가 $\mathfrak H$에 속하지만 $\mathfrak G_\eta$에는 속하지 않는 모든 함수로 이루어진 계를 나타낸다고 하자. V에 따라 $\mathfrak H$는 영역 $\{\mathfrak G_\eta,\mathfrak L\}$을 부분정역으로 포함한다. 반대로 $\{\mathfrak G_\eta,\mathfrak L\}$도 $\mathfrak H$를 부분정역으로 포함한다. 실제로 이는 $\mathfrak G_\eta$와 $\mathfrak L$을 모두 포함하며, 두 계가 공통 원소를 갖지 않으므로 그 합집합인 $\mathfrak H$도 포함한다. 따라서
$\mathfrak H=\{\mathfrak G_\eta,\mathfrak L\}$이다. $\mathfrak H$가 IV를 만족하므로 $\mathfrak L$은 물론 허용 가능한 계이다.\srcfn{**)}{$\mathfrak L$의 한 부분계를 첨가하는 것만으로도 충분할 수 있다. 예를 들어 § 2의 함수식 영역은 모든 정수적 유리 함수식 가운데 다항식을 제외한 것들을 $\mathfrak G_\eta$에 대수적 정수로 첨가하여 얻어진다.}

%
%
$\mathsf{H}$가 가능한 모든 초월기저를 두루 취하면, 처음에 증명한 대로 $\mathfrak M(\mathfrak H)$는 모든 영역 $\mathfrak H$의 전체를 두루 취한다. 따라서 대수적으로 정수인 초월수로 이루어진 가장 일반적인 영역 $\mathfrak H$에 관한 물음은 이 절의 다음 결과들로 완전히 답해진다.

%
%
\begin{enumerate}
\item[a)] \emph{$\mathfrak M(\mathfrak H)$의 모든 영역 $\mathfrak H$의 교집합은 체르멜로 영역 $\mathfrak G_\eta$이며, $\mathfrak G_\eta$ 자체도 영역 $\mathfrak H$이다. $\mathfrak M(\mathfrak H)$는 교집합을 취하는 연산에 관하여 닫혀 있다.}
\item[b)] \emph{$\mathfrak M(\mathfrak H)$의 각 영역 $\mathfrak H$는 허용 가능한 계 $\mathfrak S$를 $\mathfrak G_\eta$에 대수적 정수로 첨가하여 생긴다. $\mathsf{H}$가 가능한 모든 초월기저를 두루 취하면 $\mathfrak M(\mathfrak H)$는 모든 영역 $\mathfrak H$의 전체를 두루 취한다.}\srcfn{***)}{허용 가능한 계 $\mathfrak S$는 다음과 같이 찾을 수 있다. $\eta$의 모든 분수형 대수함수에 임의의 정렬순서 $\Omega$를 부여하고, 그 함수를 $\mathfrak G_\eta$와 앞서 받아들인 함수들에 대수적 정수로 첨가했을 때 분수형 대수수를 만들어 내는 경우를 제외하고 나머지 모든 함수를 $\mathfrak S$에 넣는다. 이 절차는 임의의 위치에서 멈출 수 있다. $\Omega$가 모든 정렬순서를 두루 취하고, 각 고정된 $\Omega$에서 임의의 위치마다 절차를 멈추게 하면 모든 허용 가능한 계 $\mathfrak S$가 빠짐없이 얻어진다.}
\end{enumerate}

%
%
\section*{§ 5. 초월기저를 바탕으로 한 정수적 초월수로 이루어진 가장 일반적인 영역}

%
%
$\mathsf{H}$를 모든 수의 임의의 초월기저라 하고, $H$와 따라서 $[\mathsf{H}]$를 포함하는 모든 영역 $\mathfrak G$의 전체를 $\mathfrak M(\mathfrak G)$로 나타내자. $H$가 임의의 모든 기저를 두루 취하면 $\mathfrak M(\mathfrak G)$는 모든 영역 $\mathfrak G$의 전체를 두루 취한다. 그러므로 § 4에서와 마찬가지로 $\mathfrak M(\mathfrak G)$에 속하는 영역 $\mathfrak G$만 살펴보면 된다.

%
%
$\mathfrak M(\mathfrak G)$의 영역 $\mathfrak G$들은 모두 $[\mathsf{H}]$를 부분정역으로 포함한다. 그러나 $[\mathsf{H}]$ 자체는 영역 $\mathfrak G$가 아니므로, § 4에서처럼 $[\mathsf{H}]$가 최대 공통 부분정역이라고 결론지을 수는 없다. 그럼에도 실제로 그러하다는 사실은 § 3의 모듈 영역을 조금 일반화하여 얻는 $\mathfrak M(\mathfrak G)$의 가산개 영역 $\mathfrak G$가 보여 준다. 이 영역들의 공통 원소는 실제로 $[\mathsf{H}]$의 원소뿐이다.

%
%
\emph{각 고정된 $\nu$에 대하여 영역 $\mathfrak G_\nu$ $(\nu=1,2,\ldots)$가 다음과 같은 모든 양으로 이루어진다고 하자.}
\[
\begin{array}{@{}l@{\quad}c@{}}
(1)& z=f(\eta)+g_1(\eta)\eta_1^\nu+g_2(\eta)\eta_2^\nu+\cdots+g_t(\eta)\eta_t^\nu.
\end{array}
\]
\emph{여기서 $f(\eta)$는 $\eta$의 모든 정수 계수 다항식을, $\eta_1\cdots\eta_t$는 모든 기저원소를 두루 취한다. 또한 $g_1(\eta)\cdots g_t(\eta)$는 $\eta_1^\nu\cdots\eta_t^\nu$의 각 고정된 조합마다 서로 독립적으로 $\eta$의 모든 정수적 대수함수를 두루 취한다.}

%
%
§ 3에서와 같은 방식으로 영역 $\mathfrak G_\nu$가 추상적 성질 I--IV를 가짐을 알 수 있다. $\mathfrak G_\nu$는 정수적 대수함수만을 포함하며, $\nu=1$일 때에는 § 3의 영역과 일치한다.

%
%
이제 영역 $\mathfrak G_\nu$들이 $[\mathsf{H}]$의 다항식 $f(\eta)$ 이외에는 어떠한 공통 원소도 갖지 않음을 보이겠다. $z$를 $\eta$의 임의의 정수적 대수함수이되 유리함수는 아닌 것으로 하고, $[\mathsf{H}]$에 대하여 기약인 다음 방정식을 만족한다고 하자.
\[
\begin{array}{@{}l@{\quad}c@{}}
(2)& z^\varkappa+a_1(\eta)z^{\varkappa-1}+a_2(\eta)z^{\varkappa-2}+\cdots+a_\varkappa(\eta)=0.
\end{array}
\]
따라서 $\varkappa>1$이다. 다항식 $a_1(\eta)\cdots a_\varkappa(\eta)$의 차수 가운데 가장 큰 것을 $\lambda$라 하자. $z$가 $\mathfrak G_\nu$에 속하면 양 $\zeta=z-f(\eta)$도 $[\mathsf{H}]$에 대하여 기약인 다음 방정식을 만족한다.
\[
\begin{array}{@{}l@{\quad}c@{}}
(3)& \zeta^\varkappa+b_1(\eta)\zeta^{\varkappa-1}+b_2(\eta)\zeta^{\varkappa-2}+\cdots+b_\varkappa(\eta)=0.
\end{array}
\]
$b_1(\eta),b_2(\eta)\cdots b_\varkappa(\eta)$는 $\zeta$의 대칭함수이므로, (1)에 따라 그 최저차항의 차수는 각각 적어도 $\nu,2\nu,\cdots,\varkappa\nu$이다. (2)와 (3)을 비교하면
\[
 b_1(\eta)=a_1(\eta)+\varkappa\cdot f(\eta)
\]
를 얻는다.

%
%
그러므로 양 $\xi=z+\frac{a_1(\eta)}{\varkappa}$가 방정식
\srcnumdisplay[3.2em]{(4)}{%
\xi^\varkappa+c_2(\eta)\xi^{\varkappa-2}+\cdots+c_\varkappa(\eta)=0
}
을 만족하면
\srcnumdisplay[3.2em]{(5)}{%
\xi-\zeta=\dfrac{b_1(\eta)}{\varkappa};\qquad
c_i(\eta)\equiv b_i(\eta)\;\mathrm{mod.}\;\dfrac{b_1(\eta)}{\varkappa}
}
가 된다. 여기서 $b_1(\eta)$은 차수가 적어도 $\nu$인 항들로 시작한다. 한편 (2)와 (4)를 비교하면 $c_i(\eta)$는 (2)의 계수 $a_i(\eta)$에 관하여 차수가 $\chi$이므로, $\eta$에 관해서는 차수가 기껏해야 $\varkappa\cdot\lambda$이다. 반면 모든 $b_i(\eta)$는 차수가 적어도 $\nu$인 항들로 시작한다. 따라서 $\nu>\varkappa\cdot\lambda$로 택하면 (5)에서
\[
 c_i(\eta)=0;\qquad
 \xi^\varkappa=0\quad\text{또는}\quad
 \left(z+\frac{a_1(\eta)}{\varkappa}\right)^\varkappa=0
\]
을 얻는다. 이는 가정에 반하여 $z$가 $\eta$의 유리함수가 된다는 뜻이다. 따라서 다음이 성립한다.

%
%
\emph{$z$가 $\eta$의 정수적 대수함수이되 유리함수가 아니라면, $\nu>\varkappa\lambda$인 어떠한 영역 $\mathfrak G_\nu$에도 $z$가 속할 수 없다.} 다시 말해 다음과 같다.

%
%
\emph{$\mathfrak M(\mathfrak G)$의 모든 영역 $\mathfrak G$의 교집합은 정역 $[\mathsf{H}]$이다. $[\mathsf{H}]$ 자체는 영역 $\mathfrak G$가 아니므로, $\mathfrak M(\mathfrak G)$는 교집합을 취하는 연산에 관하여 닫혀 있지 않다.}

%
%
따라서 초월기저를 이용해 가장 일반적인 영역 $\mathfrak G$를 구성하는 일은 영역 $\mathfrak H$의 경우보다 파악하기 어렵다. 영역 $[\mathsf{H}]$는 성질 I, III, IV를 가진다. 그러므로 임의의 계 $\mathfrak S$를 다항식적으로 첨가하면 생기는 정역 $[\mathsf{H},\mathfrak S]$에는 성질 I과 III이 그대로 남지만, IV는 $\mathfrak S$에 부과해야 할 별도의 조건이 된다. 영역 $\mathfrak G$가 생기려면 II도 $\mathfrak S$에 대한 조건으로 부과해야 한다. 달리 말하면, $(\mathsf{H})$에 대한\srcfn{*)}{$(\mathsf{H})$는 계수가 대수적 수인 $\eta$의 모든 유리함수로 이루어진 체를 뜻한다.} 모든 유한확대체 $\mathfrak K$마다 $\mathfrak K$의 유한확대체 $\mathfrak L$이 존재하여, $[\mathsf{H},\mathfrak S]$가 $\mathfrak L$의 한 원시원소를 포함해야 한다.\srcfn{**)}{이에 관해서는 § 7 참조.} 역으로 $\mathfrak M(\mathfrak G)$의 각 영역 $\mathfrak G$도 $[\mathsf{H},\mathfrak S]$ 꼴로 나타낼 수 있다. 따라서 $\mathfrak S$가 이와 같이 제한된 모든 계를 두루 취하게 하면 $\mathfrak M(\mathfrak G)$의 모든 영역 $\mathfrak G$를 얻는다. 이러한 계 $\mathfrak S$의 구조는 다음 절에서 ``유리기저''를 이용하면 더 간단하게 드러나며, 동시에 § 4의 정리들과 완전한 유사성이 얻어진다.

%
%
\section*{\centering § 6.\\[0.35em]
유리기저를 바탕으로 한 정수적 초월수로 이루어진 가장 일반적인 영역.}

%
%
초월기저와 유사하게 유리기저를 다음과 같이 정의한다. \glqq{}수 $\vartheta$들의 계 $\Theta$가 모든 수를 유리적으로 --- $\mathsf{K}$의 계수를 써서\srcfn{*)}{유리함수는 언제나 대수적 수를 계수로 갖는 것을 뜻한다. 이 대수적 수 계수들은 III와 IV 때문에 정의에 들어 있다. III와 IV에서도, 그리고 유리기저의 정의에서도 유리수 계수를 요구하는 편이 더 상응할 것이다. 이 경우 모든 대수적 수의 체 $\mathsf{K}$를 모든 유리수의 체로 바꾸기만 하면 § 6과 § 7의 결론은 글자 그대로 유지된다.} --- 나타낼 수 있게 하면서, 적어도 하나의 정렬순서에서는 어떠한 기저원소도 유한개의 앞선 기저원소를 써서 --- $\mathsf{K}$의 계수로 --- 유리적으로 나타낼 수 없을 때, $\Theta$를 모든 수의 유리기저라 한다.\grqq{}\srcfn{**)}{선형기저와 초월기저의 경우와 달리 여기서는 원소의 순서가 구실을 한다. 예를 들어 $\eta,\sqrt{\eta}$의 순서에서는 두 원소를 모두 넣어야 하지만, 역순에서는 $\sqrt{\eta}$만 넣으면 된다.}

%
%
정수적 초월수의 영역 $\mathfrak G$에 정렬순서 $\Omega$가 주어졌다고 하자. 그러면 $\mathfrak G$의 어떤 양 $z$가 $\mathfrak G$의 유한개의 앞선 양을 써서 유리적으로 나타내어질 수 있다. 다시 말해, $z=\varphi(\zeta)$라는 방정식을 만족할 수 있다. 여기서 $\varphi(\zeta)$는 $\xi_1,\ldots,\xi_t$의, $K$에 속하는 계수를 가진 유리함수이다. 특히 이 영역의 모든 대수적 수 $\alpha$는 방정식 $z=\alpha$를 만족한다.

%
%
$\mathfrak G$의 나머지 양 $\vartheta$들은 계 $\Theta$를 이루며, 이 계가 $\mathfrak G$의 모든 양의 유리기저임을 보이겠다. 실제로 가정에 따라 각 양 $\vartheta$는 $\Theta$의 앞선 양들에 유리적으로 종속되지 않는다. 또한 귀납 논법\srcfn{***)}{Zermelo, 앞의 논문 § 1 참조.}은 모든 관계 $z=\varphi(\zeta)$가 관계 $z=\psi(\vartheta)$를 수반함을 보여 준다. 그렇지 않다면 $\vartheta$들로 유리적으로 나타낼 수 없는 첫 원소 $z_0=\varphi_0(\xi_1,\ldots,\xi_t)$가 존재해야 하지만, 그보다 앞선 $\xi_1,\ldots,\xi_t$는 $\vartheta$들의 유리함수이므로 모순이 생긴다. $\Theta$는 II에 따라 동시에 모든 수의 유리기저이다. III와 I에 따라 $\mathfrak G$는 $\Theta$와 더불어 $\vartheta$의 모든 정수 계수 다항식으로 이루어진 정역 $[\Theta]$를 부분정역으로 포함하며, IV에 따라 $[\Theta]$는 어떠한 분수형 대수수도 포함하지 않는다. 따라서 모든 수의 유리기저 $\Theta$는, 그로부터 생기는 정역 $[\Theta]$가 어떠한 분수형 대수수도 포함하지 않는다는 부가조건을 만족한다.\srcfn{*)}{이것이 실제 조건이라는 사실은, 예를 들어 두 양 $\eta,\frac1{2\sqrt{\eta}}$는 기저 원소로 허용되지 않지만 $\eta,\frac1{\sqrt{\eta}}$는 허용된다는 데서 드러난다.} 즉 $\Theta$는 모든 수에 대한 \emph{부가조건을 만족하는 유리기저}이다. 그러므로 § 2와 유사하게 다음이 성립한다. \emph{각 영역 $\mathfrak G$는 부가조건을 만족하는 유리기저를 이용하여 구성할 수 있으며, 이때 $\mathfrak G$는 정역 $[\Theta]$를 부분정역으로 포함한다.}

%
%
$\Theta$를 모든 수에 대한 부가조건을 만족하는 임의의 유리기저라고 하자. 이러한 기저의 존재는 영역 $\mathfrak G$의 존재로 보장된다. $\Theta$와 따라서 $[\Theta]$를 포함하는 모든 영역 $\mathfrak G$의 전체를 $\mathfrak N(\mathfrak G)$로 나타내자. $\Theta$가 부가조건을 만족하는 임의의 모든 기저를 두루 취하면, 위의 결과에 따라 $\mathfrak N(\mathfrak G)$는 모든 영역 $\mathfrak G$의 전체를 두루 취한다. 따라서 다시 $\mathfrak N(\mathfrak G)$에 속하는 영역 $\mathfrak G$만 살펴보면 된다.

%
%
이제 $[\Theta]$ 자체가 영역 $\mathfrak G$이다. 실제로 모든 수는 $\Theta$로 유리적으로, 곧 $[\Theta]$의 두 다항식의 몫으로 나타낼 수 있으며, $[\Theta]$는 그 구성에 따라 나머지 조건 I, III, IV도 만족한다. \emph{따라서 $[\Theta]$는 $\mathfrak N(\mathfrak G)$의 모든 영역 $\mathfrak G$의 최대 공통 부분정역, 곧 교집합이다. 또한 $\mathfrak N(\mathfrak G)$는 교집합을 취하는 연산에 관하여 닫혀 있다. 즉 $\mathfrak N(\mathfrak G)$의 임의의 부분계에 속하는 모든 영역 $\mathfrak G$의 교집합도 $\mathfrak N(\mathfrak G)$에 속한다.} 실제로 그 교집합에는 I이 성립하고, II와 III은 $[\Theta]$가 부분정역이라는 데서 따르며, IV는 그 교집합 영역이 어떤 영역 $\mathfrak G$에 부분정역으로 포함된다는 데서 따른다.

%
%
$\vartheta$의 유리함수로 이루어진 임의의 계 $\mathfrak S$를 $[\Theta]$에 다항식적으로 첨가하자. $\vartheta$의 모든 대수함수는 유리기저의 성질에 따라 어떤 다른 $\vartheta$들을 써서 유리적으로 나타낼 수 있다. 그러면 정역 $[\Theta,\mathfrak S]$가 생기며, 이 정역에는 성질 I, II, III이 성립한다. 따라서 $\mathfrak S$가 \glqq{}허용 가능한 계\grqq{}, 곧 다항식적 첨가로 어떠한 분수형 대수수도 영역 안에 들어오지 않는다는 조건을 만족하면 $[\Theta,\mathfrak S]$는 영역 $\mathfrak G$가 된다. 역으로 $\mathfrak N(\mathfrak G)$의 각 영역 $\mathfrak G$도 $\mathfrak S$를 나머지 계 $\mathfrak G-[\Theta]$로 잡으면 $[\Theta,\mathfrak S]$ 꼴로 나타낼 수 있다. 따라서 가장 일반적인 영역 $\mathfrak G$에 대해서도 § 4와 완전히 유사한 다음 결과들이 성립한다.

%
%
\begin{enumerate}
\item[a)] \emph{$\mathfrak N(\mathfrak G)$의 모든 영역 $\mathfrak G$의 교집합은 정역 $[\Theta]$이며, 이 정역 자체도 영역 $\mathfrak G$이다. $\mathfrak N(\mathfrak G)$는 교집합을 취하는 연산에 관하여 닫혀 있다.}
\item[b)] \emph{$\mathfrak N(\mathfrak G)$의 각 영역 $\mathfrak G$는 허용 가능한 계 $\mathfrak S$를 $[\Theta]$에 다항식적으로 첨가하여 생긴다. $\Theta$가 부가조건을 만족하는 가능한 모든 유리기저를 두루 취하면, $\mathfrak N(\mathfrak G)$는 모든 영역 $\mathfrak G$의 전체를 두루 취한다.}\srcfn{*)}{가장 일반적인 허용 가능한 계 $\mathfrak S$에 관해서는 § 4에서 말한 내용에서 \glqq{}대수적\grqq{}을 \glqq{}유리적\grqq{}으로 바꾸기만 하면 된다.}
\end{enumerate}

%
%
\textbf{주:} $\Theta$에서 $\Theta$의 모든 양에 대한 초월기저 $\mathsf{H}$를 골라내면, $\mathsf{H}$는 동시에 모든 수의 초월기저가 된다. 역으로 모든 초월기저는, 정렬순서에서 $\mathsf{H}$를 맨 앞에 놓으면 유리기저로 확장할 수 있다. 이에 따라 § 5에서 초월기저를 이용하여 영역 $\mathfrak G$를 구성한 방식을 다음과 같이 해석할 수 있다. $[\mathsf{H}]$에 첨가할 계 $\mathfrak S$를 두 부분계 $\mathfrak S_1$과 $\mathfrak S_2$로 나눈다. $\mathfrak S_1$을 첨가하는 것은 $\mathsf{H}$를 부가조건을 만족하는 유리기저로 확장하는 것과 같고, 그 뒤 허용 가능한 계 $\mathfrak S_2$를 다시 다항식적으로 첨가한다.

%
%
\section*{\centering § 7.\\[0.35em]
부가조건을 만족하는 가장 일반적인 유리기저의\\
구성.}

%
%
앞 절의 결과는 모든 수에 대하여 부가조건을 만족하는 유리기저의 \emph{존재}는 증명했지만, 주어진 정렬순서에 따른 그 구성법은 제시하지 않았다는 점에서 보충이 필요하다. 이 구성은 영역 $\mathfrak G$ 대신 모든 복소수의 체를 다룰 때 부가조건 때문에 더 복잡해진다.\srcfn{**)}{§ 6 b)에서는 $\Theta$가 초월기저와 기저원소들의 대수적으로 정수인 함수들로 이루어진 그러한 특수한 유리기저들만 두루 취하게 해도 충분하다. 그러면 부가조건은 저절로 만족된다. 이를 보려면 § 6에 따라 $\mathfrak G$에서 고른 초월기저를 유리기저로 확장하고, 그렇게 얻은 기저에서 $\eta$의 각 분수형 대수함수에 적절히 고른 $\eta$의 다항식을 곱하기만 하면 된다. 그러면 기저 성질은 보존된다. 이 구성은 모든 복소수의 체에도 변함없이 옮길 수 있다. 그러나 § 6에서 생기는, 부가조건을 만족하는 가장 일반적인 유리기저의 문제는 그 자체로도 독립적인 관심을 지닌다.}

%
%
모든 복소수의 체에 정렬순서 $\Omega$를 부여하자. 그러면 다음 세 가지 유형의 양이 나타날 수 있다.
\begin{enumerate}
\item[1.] 앞선 --- 3)에서 재귀적으로 정의되는 --- 기저원소들\srcfn{***)}{귀납 논법이 보여 주듯이 그러한 재귀적 정의는 가능하다.}에 $y$를 다항식적으로 첨가하면, 생성되는 정역 안에 분수형 대수수가 들어오게 되는 양 $y$; 앞선 기저원소가 하나도 없으면 이 정역은 $y$의 모든 정수 계수 다항식으로 이루어진다. 특히 모든 분수형 대수수는 유형 1의 양 $y$이다. 생성되는 정역이 양 $y=\beta$를 포함하기 때문이다.
\item[2.] 성질 1)을 갖지는 않지만, 앞선 수들 가운데 $y$가 아닌 유한개의 수를 써서 유리적으로 나타낼 수 있는 양 $z$. 즉 이들은
$z=\varphi(\zeta_1\cdots\zeta_t)$라는 관계를 만족한다. 여기서 $\varphi$는 계수가 $\mathsf{K}$에 속하는 유리함수이고, $\xi$들 가운데에는 어떠한 $y$도 없다. 특히 모든 대수적 정수 $\alpha$가 여기에 속한다. 이들은 성질 1)을 갖지 않으며 관계 $z=\alpha$를 만족하기 때문이다.
\item[3.] 1)도 2)도 성립하지 않으며, 앞으로 기저원소라 부를 양 $\vartheta$. 특히 정렬순서에서 첫 번째 초월수 $\vartheta_0$는 그러한 기저원소이다. $\vartheta_0$의 정수 계수 다항식으로는 어떠한 분수형 대수수도 나타낼 수 없고, $\vartheta_0$는 앞선 대수적 수들에 유리적으로 종속될 수도 없기 때문이다.
\end{enumerate}

%
%
\emph{이제 모든 기저원소 $\vartheta$로 이루어진 계 $\Theta$가 모든 수에 대하여 부가조건을 만족하는 유리기저임을 보이겠다.}

%
%
어떠한 $\vartheta$도 양 $y$가 아니므로 정역 $[\Theta]$는 분수형 대수수를 포함할 수 없다. 따라서 부가조건이 만족된다. 또한 어떠한 $\vartheta$도 양 $z$가 아니므로 각 $\vartheta$는 앞선 기저원소들에 유리적으로 종속되지 않는다. 더 나아가 § 6에서와 똑같은 귀납 논법으로, $\xi$들 가운데 양 $y$가 하나도 없는 모든 관계 $z=\varphi(\zeta)$는 관계 $z=\psi(\vartheta)$를 수반함을 알 수 있다. 따라서 모든 양 $z$는 $\vartheta$들을 써서 유리적으로 나타낼 수 있다.

%
%
이제 더 많은 수고를 요하는 것은 양 $y$도 $\vartheta$들을 써서 유리적으로 나타낼 수 있음을 증명하는 일뿐이다. 먼저 $y$가 $\vartheta$들에 대수적으로 종속됨을 보이겠다. 가정에 따라 각 $y$마다 다음 표현을 허용하는 분수형 대수수 $\beta$가 적어도 하나 존재한다.
\[
\begin{array}{@{}l@{\quad}c@{}}
(1)& \beta=f_0(\vartheta)+f_1(\vartheta)y+\cdots+f_\varkappa(\vartheta)y^\varkappa,
\end{array}
\]
여기서 $f_0(\vartheta)\cdots f_\varkappa(\vartheta)$는 $[\Theta]$의 다항식이므로 어떠한 분수형 대수수도 나타낼 수 없다. 만약 $y$가 $\vartheta$들에 대수적으로 독립이라면 (1)은 $y$에 관한 항등식이어야 하며, 따라서 $\beta=f_0(\vartheta)$가 된다. 이는 $[\Theta]$의 성질과 모순된다.

%
%
이 대수적 종속으로부터 유리적 종속을 이끌어 내기 위하여 $\Theta$에서 $\Theta$의 모든 양에 대한 초월기저 $\mathsf{H}$를 골라낸다. 그러면 각 $y$는 $\vartheta$들의 대수함수이므로 $\eta$들에도 대수적으로 종속된다. 따라서 $y=y(\eta)$는 $(\mathsf{H})$ 위의 대수적 확대체 $(\mathsf{H},y(\eta))$를 정한다. 모든 $\eta$가 동시에 양 $\vartheta$이므로, 유리적 종속의 증명은 그러한 각 체에 대하여 $\vartheta$들의 유리함수인 원시원소가 존재함을 보이는 일로 귀착된다. 더 정확히는 $y(\eta)$에 $\eta$의 어떤 다항식을 곱하면 성질 1)을 잃고, 따라서 그러한 원시원소로 바뀜을 보이겠다.

%
%
알려져 있듯이 $(\mathsf{H})$와 $(\mathsf{H},y)$ 사이에는 유한개의 중간체만 존재한다. 그 가운데 $\vartheta$들로 유리적으로 나타낼 수 있는 원시원소를 가진 것들을 $\Omega_0=(H),\Omega_1,\ldots,\Omega_\sigma$라 하자. 그러면 그 체의 각 원소는 $\vartheta$들의 유리함수가 된다. 이 조건은 적어도 $\mathfrak L_0=(\mathsf{H})$에 대해서는 언제나 만족된다. $\mathfrak L_i$에 속하는 원시함수가
$\frac{g_i(\vartheta)}{h_i(\vartheta)}$로 주어졌다고 하자. 여기서 적절한 정수 $\alpha_i$를 택하면 $[\Theta]$의 다항식
\[
 \zeta_i=g_i(\vartheta)+\alpha_i h_i(\vartheta)
\]
는 $\mathfrak L_i$ 또는 $\mathfrak L_i$ 위의 어떤 대수적 확대체의 원시함수가 된다.\srcfn{*)}{이에 관해서는 § 5의 끝을 참조하라.} 따라서 $\mathfrak L_i$에 속하는 각 함수 $\psi(\eta)$는 다음과 같이 나타낼 수 있다.
\[
\begin{array}{@{}l@{\quad}l@{}}
(2)& \displaystyle \psi(\eta)=
\frac{a_1(\eta)\zeta_i^{\varkappa-1}+a_2(\eta)\zeta_i^{\varkappa-2}+\cdots+a_\varkappa(\eta)}
     {a(\eta)}
=\frac{A(\eta,\zeta_i)}{a(\eta)}.
\end{array}
\]
여기서 $a(\eta),a_1(\eta),\ldots,a_\varkappa(\eta)$는 $\eta$의 정수 계수 다항식이며, 따라서 $A(\eta,\zeta_i)$와 $a(\eta)$는 $[\Theta]$의 양이다.

%
%
$y$가 $\mathfrak L_i$에 관하여 만족하는 기약방정식이 (2)에 따라 다음 꼴이라고 하자.
\[
 f^{(i)}(\eta)y^{\lambda_i}+f_1^{(i)}(\eta,\zeta_i)y^{\lambda_i-1}
+\cdots+f_{\lambda_i}^{(i)}(\eta,\zeta_i)=0,
\]
이 방정식의 모든 계수는 $[\Theta]$에 속한다. 그러면 함수 $y_i=f^{(i)}(\eta)\cdot y$는 $\mathfrak L_i$에 관하여 마찬가지로 기약인 다음 방정식을 만족한다.
\[
 y_i^{\lambda_i}+f_1^{(i)}(\eta,\zeta_i)y_i^{\lambda_i-1}
+\cdots+\bigl(f^{(i)}(\eta)\bigr)^{\lambda_i-1}
 f_{\lambda_i}(\eta,\zeta_i)=0.
\]
따라서 $y_i$는 동시에 $[\mathsf{H},\zeta_i]$에 대하여 대수적으로 정수이며, $y_i$에 $\eta$의 임의의 다항식을 곱한 것에도 같은 사실이 성립한다. 특히 함수
\[
\begin{array}{@{}l@{\quad}l@{}}
(3)& \displaystyle Y=f^{(0)}(\eta)f^{(1)}(\eta)\cdots f^{(\sigma)}(\eta)\cdot y
\end{array}
\]
는 각각 $\mathfrak L_i$에 관한 다음 기약방정식에 의하여 모든 정역 $[\mathsf{H},\zeta_i]$에 대하여 대수적으로 정수이다.
\[
\begin{array}{@{}l@{\quad}l@{}}
(4)& \displaystyle Y^{\lambda_i}+F_1^{(i)}(\eta,\zeta_i)Y^{\lambda_i-1}
+\cdots+F_{\lambda_i}^{(i)}(\eta,\zeta_i)=0
\qquad (i=0,1,\ldots,\sigma).
\end{array}
\]

%
%
이제 $Y$가 우리가 찾는, $\vartheta$들로 유리적으로 나타낼 수 있는 원시원소임을 보이겠다. $Y$와 $[\Theta]$를 써서 어떤 대수적 수 $\gamma$가 다음과 같이 나타난다고 하자.
\srcnumdisplay[3.2em]{(5)}{%
\gamma=k_0(\vartheta)+k_1(\vartheta)Y+\cdots+k_\nu(\vartheta)Y^\nu
}
여기서 $k_0(\vartheta)\cdots k_\nu(\vartheta)$는 $[\Theta]$의 다항식이다. 이제 (5)의 계수 영역에서 생기는 체
$\mathfrak K=(\mathsf{H};k_0(\vartheta)\cdots k_\nu(\vartheta))$에 관하여 $Y$가 만족하는 기약방정식을 정하자. 알려져 있듯이 이 방정식은 $(\mathsf{H})$와 $(\mathsf{H},Y)$ 사이에 놓인 $\mathfrak K$와 $(\mathsf{H},Y)$의 교집합체에 관하여 $Y$가 만족하는 기약방정식으로 주어진다. 한편 (3)에 따라 $(\mathsf{H},Y)$는 $(\mathsf{H},y)$와 같다. 더 나아가 $\mathfrak K$의 모든 원소, 따라서 교집합체의 모든 원소는 유한개의 $\vartheta$를 써서 유리적으로 나타낼 수 있으므로, 이 교집합체는 반드시 $\mathfrak L_0,\mathfrak L_1,\cdots,\mathfrak L_\sigma$ 가운데 하나, 이를테면 $\mathfrak L_\tau$이다. 따라서 찾는 기약방정식의 차수는 (4)에 따라 $\lambda_\tau$이며, 다음과 같다.
\srcnumdisplay[3.2em]{(6)}{%
Y^{\lambda_\tau}+F_1^{(\tau)}(\eta,\zeta_\tau)Y^{\lambda_\tau-1}
+\cdots+F_{\lambda_\tau}^{(\tau)}(\eta,\zeta_\tau)=0
}
그 계수들은 $[\Theta]$의 다항식이다. (6)에 따라 (5)를 다음과 같이 바꿀 수 있다.
\srcnumdisplay[3.2em]{(7)}{%
\gamma=K_0(\vartheta)+K_1(\vartheta)Y+\cdots+
K_{\lambda_\tau-1}(\vartheta)Y^{\lambda_\tau-1},
}
여기서 $K_0(\vartheta),K_1(\vartheta),\cdots$는 $\mathfrak K$에 속하면서 동시에 $[\Theta]$의 다항식이다. 그런데 $\gamma$는 대수적 수이므로 관계 (7)의 $Y$에 관한 차수는 $\lambda_\tau$에 이르지 못하며, 따라서 (7)은 $Y$에 관한 항등식으로 성립한다. 그러므로
\[
 \gamma=K_0(\vartheta),
\]
즉 $\gamma$는 $[\Theta]$에 속하고, 따라서 대수적 정수이다.

%
%
이제 $\gamma$가 $Y$와 $[\Theta]$를 써서 나타낼 수 있는 모든 대수적 수를 두루 취하게 하자. 그러면 $(\mathsf{H},Y)$와 그에 딸린 체 $\mathfrak K$의 교집합체는 언제나 $\mathfrak L_0,\mathfrak L_1,\cdots,\mathfrak L_\sigma$ 가운데 하나일 뿐이며, 위의 모든 결론은 그대로 유지된다. 다시 말해 \emph{$Y$를 $[\Theta]$에 다항식적으로 첨가해도 어떠한 분수형 대수수도 나타낼 수 없다}. 따라서 $Y$는 더 이상 성질 1)을 갖지 않는다. 그러므로 2) 또는 3)에 따라 양 $z$나 $\vartheta$이며, 따라서 유한개의 $\vartheta$를 써서 유리적으로 나타낼 수 있다. (3)에 따라 $y$에도 같은 사실이 성립한다. 이로써 \emph{$\Theta$가 모든 수의 유리기저임이 증명되었다.} 역으로 정렬순서에서 $\Theta$를 맨 앞에 두면, 부가조건을 만족하는 모든 유리기저를 1), 2), 3)에 따라 구성할 수 있다. 따라서 다음이 증명되었다. \emph{1), 2), 3)에 주어진 구성은 모든 수에 대하여 부가조건을 만족하는 가장 일반적인 유리기저를 산출한다.}

%
%
\section*{§ 8. 영역 $\mathfrak G$와 $\mathfrak H$의 동형류 분류.\srcfn{*)}{E. Steinitz, \emph{Algebraische Theorie der Körper}, Crelles Journal 137 (1910), 특히 § 23과 § 24 참조. (§ 23에서는 이미 초월기저의 존재가 증명되어 있다.)}}

가장 일반적인 영역 $\mathfrak G$와 $\mathfrak H$에 관한 물음과 더불어, 아래에서 더 정확히 규정할 의미에서 같은 구성 원리로 만들어 낼 수 없는 \emph{본질적으로 서로 다른} 영역들에 관한 물음이 제기된다. 이는 이 영역들을 동형류로 분류하는 데 이른다.

\emph{두 영역의 원소들을 일대일로 대응시켜, 한 영역의 임의의 두 원소의 합과 곱에 언제나 대응하는 원소들의 합과 곱이 대응하게 할 수 있을 때, 두 영역은 같은 동형류에 속한다거나 서로 동형이라고 하기로 한다.} 이 정의에 따르면 모든 동형 대응에서 영과 단위원은 각각 자기 자신에 대응하고, 유한개의 원소 사이에 성립하는 모든 대수적 관계는 대응하는 원소들에 대해서도 보존되며 그 역도 성립한다. 따라서 단위원뿐 아니라 모든 유리수도 자기 자신에 대응한다. 반면 대수적 수 전체와 마찬가지로 대수적 정수 전체는 각각 자기 자신 안으로 옮겨지지만, 개별 대수적 수에는 그 켤레가 대응할 수도 있다.

먼저, 고정된 초월기저 $H$를 포함하는 모든 영역 $\mathfrak G$의 전체 $\mathfrak M(\mathfrak G)$에 속하는 각 영역은 다른 초월기저 $\mathsf{H}^*$에 딸린 전체 $\mathfrak M^*(\mathfrak G)$의 한 영역과 동형임을 지적해 두자. 실제로 모든 복소수의 체는 어느 초월기저에서든 대수적 확대를 거쳐 생기므로 그 기저와 같은 기수를 가진다.\srcfn{**)}{Steinitz, § 23과 § 24.} 따라서 $\mathsf{H}$와 $\mathsf{H}^*$는 같은 기수를 가지며, 기저원소 $\eta_i$와 $\eta_i^*$를 서로 대응시키면 원하는 동형사상을 얻는다.\srcfn{***)}{물론 이 논법은 두 유리기저 $\Theta$와 $\Theta^*$에 딸린 집합 $\mathfrak N(\mathfrak G)$와 $\mathfrak N^*(\mathfrak G)$에는 적용되지 않는다. $\Theta$와 $\Theta^*$를 서로 유리적으로 나타낼 수 있다는 사실만으로는 $\mathfrak N(\mathfrak G)$와 $\mathfrak N^*(\mathfrak G)$가 동형이라고 결론지을 수 없다. 다만 이 사실은 체 $(\Theta)$와 $(\Theta^*)$ 사이의 동형 대응을 준다.} 그러나 § 2, § 3, § 5의 예에서 알 수 있듯이 $\mathfrak M(\mathfrak G)$는 서로 동형이 아닌, 곧 본질적으로 서로 다른 영역 $\mathfrak G$들을 포함한다. 동형사상은 모든 대수적 관계를 보존하므로 초월기저에는 다시 초월기저가 대응해야 한다. 그런데 § 2와 § 3에서는 함수식 영역과 모듈 영역에 체르멜로 영역의 모든 기저 성질을 만족하는 초월기저가 존재하지 않음을 보였다. § 2의 함수식 영역은 대수적 정수성에 관하여 닫혀 있으므로, 부분집합 $\mathfrak M(\mathfrak H)$도 본질적으로 서로 다른 영역 $\mathfrak H$들을 포함한다.

이제 기저를 구성할 때와 마찬가지로 $\mathfrak M(\mathfrak G)$에서 다음 성질을 가진 부분집합 $\mathfrak L(\mathfrak G)$를 골라낼 수 있음을 보이겠다. $\mathfrak L(\mathfrak G)$에 속하는 영역 $\mathfrak G$들은 모두 본질적으로 서로 다르지만, $\mathfrak M(\mathfrak G)$의 각 영역, 따라서 모든 영역 $\mathfrak G$는 $\mathfrak L(\mathfrak G)$의 한 영역과 동형이다. $\mathfrak M(\mathfrak G)$에 정렬순서 $\Omega$를 부여하자. 그러면 $\mathfrak M(\mathfrak G)$의 영역 $\mathfrak G$는 앞선 어떤 영역과 동형일 수도 있고 그렇지 않을 수도 있다. 앞선 어느 영역과도 동형이 아닌 영역 $\mathfrak G$들은 집합 $\mathfrak L(\mathfrak G)$를 이룬다. 구성상 이 집합에는 본질적으로 서로 다른 영역 $\mathfrak G$만 들어 있다. 또한 귀납 논법에 따르면 $\mathfrak M(\mathfrak G)$의 각 영역은 $\mathfrak L(\mathfrak G)$의 한 영역과 동형이다. 실제로 이 성질을 만족하지 않는 첫 영역이 $\mathfrak G_1$이라고 하자. $\mathfrak G_1$은 $\mathfrak L(\mathfrak G)$에 속하거나, 앞선 영역 $\mathfrak G_0$와 동형이다. 후자의 경우 가정에 따라 $\mathfrak G_0$가 $\mathfrak L(\mathfrak G)$의 한 영역과 동형이므로 $\mathfrak G_1$도 그러하다. 모든 영역 $\mathfrak G$는 어떤 집합 $\mathfrak M^*(\mathfrak G)$에 속하고, 따라서 $\mathfrak M(\mathfrak G)$의 한 영역과 동형이다. \emph{그러므로 $\mathfrak L(\mathfrak G)$는 영역 $\mathfrak G$의 모든 동형류를 대표한다.}

$\mathfrak L(\mathfrak G)$의 영역 $\mathfrak G$에서 기저원소 $\eta$를 다른 임의의 기저의 기저원소 $\eta^*$로 바꾸면, 위에서 보았듯이 $\mathfrak G$와 동형인 영역이 생긴다. 역으로 $\mathfrak G^*$를 임의의 영역, 따라서 $\mathfrak L(\mathfrak G)$의 어떤 영역 $\mathfrak G$와 동형인 영역이라 하자. 이때 기저원소 $\eta$에 $\mathfrak G^*$의 원소 $\eta^*$가 대응한다면, 이 원소들은 다시 초월기저를 이룬다. 또한 $\mathfrak G^*$는 $\eta^*$의 대수함수 가운데 $\mathfrak G$가 $\eta$에 관하여 포함하는 것과 같은 함수들, 또는 그 켤레들을 포함한다. \emph{따라서 $\mathfrak G$와 동형인 모든 영역은 $\mathfrak G$와 $(\mathsf{H})$에 관한 그 켤레 영역들에서}\srcfn{*)}{$\mathfrak G$의 켤레 영역에 속하는 원소들은 II에 따라 분명 $\mathfrak G$의 원소들을 써서 유리적으로 나타낼 수 있다. 그러나 반드시 다항식적으로 나타낼 수 있는 것은 아니므로, 켤레 영역은 $\mathfrak G$와 다를 수 있다.} \emph{--- 켤레 영역이 $\mathfrak G$와 다를 경우에는 이들도 포함하여 --- 초월기저 $\mathsf{H}$가 가능한 모든 초월기저를 두루 취하게 함으로써 얻어진다.} $\mathfrak L(\mathfrak G)$의 고정된 각 영역 $\mathfrak G_i$에서 이 방법으로 동형류 $\mathfrak K_i$가 생긴다. 이 동형류는 $\mathfrak G_i$와 동형인 모든 영역 $\mathfrak G$만을 빠짐없이 포함하지만, 같은 영역을 여러 번, 심지어 무한히 여러 번 포함할 수도 있다. 위의 방법으로 $\mathfrak K_i$에서 $\mathfrak G_i$와 동형인 각 영역을 정확히 한 번씩만 포함하는 부분집합 $\mathfrak K_i^*$를 골라낼 수 있다. $\mathfrak K_i^*$가 모든 동형류를 두루 취하게 하면, 모든 영역 $\mathfrak G$를 각각 정확히 한 번씩 얻는다. 물론 동형류 $\mathfrak K_i$를 특징짓는 데에는 $\mathfrak G_i$ 대신 그 동형류의 다른 어느 영역을 써도 된다. 그 영역에서도 위와 같이 $\mathfrak K_i$의 모든 영역을 이끌어 낼 수 있기 때문이다. 요약하면 다음과 같다.

\emph{$\mathfrak M(\mathfrak G)$에서 부분집합 $\mathfrak L(\mathfrak G)$를 골라내어, 그 영역 $\mathfrak G$들을 모든 동형류의 전체와 중복 없이 일대일로 대응시킬 수 있다. 이때 $\mathfrak G_i$에 딸린 동형류 $\mathfrak K_i$는 $\mathfrak G_i$와 $(\mathsf{H})$에 관한 그 켤레 영역들에서 $\mathsf{H}$가 가능한 모든 초월기저를 두루 취하게 함으로써 생긴다. $\mathfrak K_i^*$가 $\mathfrak K_i$의 서로 다른 모든 영역을 뜻하고, $\mathfrak K_i^*$가 모든 동형류를 두루 취하게 하면, 모든 영역 $\mathfrak G$를 각각 정확히 한 번씩 얻는다.}\srcfn{*)}{$\mathfrak M(\mathfrak G)$는 § 6과 § 7에 따라 다음과 같이 구성한다. § 7에 따라 초월기저 $\mathsf{H}$를 그것에서 생기는, 부가조건을 만족하는 각 유리기저 $\Theta$로 확장하고, 각 $\Theta$에 대하여 § 6 b)에 따라 $\Theta$를 포함하는 모든 영역 $\mathfrak G$의 전체 $\mathfrak N(\mathfrak G)$를 만든다.}

대수적으로 정수인 초월수로 이루어진 영역 $\mathfrak H$에 대해서도 같은 사실이 성립한다.

\section*{§ 9. 영역 $\mathfrak G$의 가산무한히 많은 동형류의 한 예.}

§ 4 b)와 § 6 b)에 제시한 구성들의 기수는 그곳의 주석에서 알 수 있듯이 모든 정렬순서의 기수와 같고, 따라서 연속체의 기수를 $c$라 할 때 $2^c$와 같다. 그러나 여기서 각 영역을 정확히 한 번씩 센 모든 영역 $\mathfrak G$의 전체가 갖는 기수를 역으로 알아낼 수는 없으며, 모든 동형류의 기수는 더더욱 알아낼 수 없다. 동형류의 수가 유한하지 않다는 사실을 § 5의 영역들과 비슷한 가산무한히 많은 영역 $\mathfrak G$의 예로 보이겠다. 이 영역들은 모두 본질적으로 서로 다름이 드러날 것이므로 가산무한히 많은 동형류를 준다.\srcfn{**)}{§ 6의 영역 $\mathfrak G$들도 마찬가지로 가산무한히 많은 동형류를 주지만, 그 증명은 조금 더 복잡하다.}

§ 5의 (1)과 마찬가지로 영역 $\mathfrak G_\sigma$가 다음 모든 양으로 이루어진다고 하자.
\[
\begin{array}{@{}l@{\quad}l@{}}
(1)& \displaystyle z\equiv a_0+a_1(\eta)+\cdots+a_{\sigma-1}(\eta)
 \;\mathrm{mod.}\;\mathfrak M_\sigma(\eta),
\end{array}
\]
\emph{여기서 $a_i(\eta)$는 $\eta$의 $i$차 동차형식이고, 모듈 $\mathfrak M_\sigma$는 $\eta$의 $\sigma$차 멱단항식 전부를 기저로, $\eta$의 모든 정수적 대수함수를 승수로 가진다.}\srcfn{***)}{물론 개별 $z$마다 모듈에 나타나는 $\eta$의 멱단항식은 유한개뿐이다.} 두 영역 $\mathfrak G_\sigma$와 $\mathfrak G_\tau$ $(\sigma\ne\tau)$ 사이에 동형 대응이 있으면 모듈 $\mathfrak M_\sigma$와 $\mathfrak M_\tau$도 반드시 서로 동형으로 대응함을 보이겠다. 그러면 그러한 동형 대응이 불가능하다는 결론이 쉽게 따른다.

$\mathfrak G_\tau$의 부정원들을 $\xi$라 하고, 동형 대응에 따라 이 $\xi$들에 $\mathfrak G_\sigma$의 양 $f(\eta)$들이 대응한다고 하자. 동형 대응은 몫을 취해도 보존되므로, $\xi$의 모든 정수적 대수함수로 이루어진 영역 $\mathfrak G_\zeta$에는 대수적 정수성에 관하여 닫힌 영역 $\mathfrak T$가 대응한다. $\mathfrak T$는 $f(\eta)$에 대하여 대수적으로 정수인 모든 함수로 이루어지며, 따라서 $\eta$에 관해서도 대수적으로 정수이다.

(1)의 한 양에 $\mathfrak M_\tau$의 한 양이 대응한다고 하자. $\mathfrak M_\tau$의 모듈 성질에 따라 (1)에 $\mathfrak T$의 임의의 양 $\psi(\eta)$를 곱한 것도 $\mathfrak G_\sigma$에 속해야 한다. 식으로 쓰면
\[
\begin{array}{@{}l@{\quad}l@{}}
(2)& \displaystyle (a_0+a_1(\eta)+\cdots+a_{\sigma-1}(\eta))\cdot\psi(\eta)
\equiv b_0+b_1(\eta)+\cdots+b_{\sigma-1}(\eta)
\;\mathrm{mod.}\;\mathfrak M_\sigma.
\end{array}
\]
모든 $\eta$를 영으로 놓으면 $a_0\cdot\psi(0)=b_0$를 얻으므로, (2)는 다음과 같이 바뀐다.
\[
\begin{array}{@{}l@{\quad}l@{}}
(3)& \displaystyle (a_0+a_1(\eta)+\cdots+a_{\sigma-1}(\eta))(\psi(\eta)-\psi(0))
\equiv c_1(\eta)+\cdots+c_{\sigma-1}(\eta)
\;\mathrm{mod.}\;\mathfrak M_\sigma.
\end{array}
\]

$\psi(\eta)$와 더불어 각 $\nu$에 대하여 함수
\[
 \chi_\nu(\eta)=\sqrt[\nu]{\psi(\eta)-\psi(0)}
\]
도 $\mathfrak T$에 속하고 $\chi_\nu(0)=0$이므로, (3)과 마찬가지로 다음 식들이 성립한다.
\[
\begin{array}{@{}l@{\quad}l@{}}
(4)& \displaystyle (a_0+a_1(\eta)+\cdots+a_{\sigma-1}(\eta))\chi_\nu(\eta)
\equiv c_1^{(\nu)}(\eta)+\cdots+c_{\sigma-1}^{(\nu)}(\eta)
\;\mathrm{mod.}\;\mathfrak M_\sigma
\quad(\nu=1,2,\cdots).
\end{array}
\]
$\chi_1=\psi(\eta)-\psi(0)$와 그 $\varkappa-1$개 켤레의 곱을 차수 $\lambda$인 $A(\eta)$라 하자. $\chi_1(0)=0$이므로 $A(\eta)$는 차수가 적어도 1인 항들로 시작해야 한다. (4)에서
\[
\begin{array}{@{}l@{\quad}l@{}}
(5)& \displaystyle a_0\cdot\chi_\nu(\eta)\equiv0\;\mathrm{mod.}\;\mathfrak M_1;\qquad
a_0^\nu\cdot\chi_1(\eta)\equiv0\;\mathrm{mod.}\;\mathfrak M_\nu;\qquad
a_0^{\nu\varkappa}\cdot A(\eta)\equiv0\;\mathrm{mod.}\;\mathfrak M_{\nu\varkappa},
\end{array}
\]
를 얻는다. 따라서 $a_0\ne0$이면 § 3에서와 마찬가지로 $\lambda>\nu\varkappa$이다. 그러나 $\nu<\frac{\lambda}{\varkappa}$라는 제한은 $\mathfrak T$가 대수적 정수성에 관하여 닫혀 있다는 사실과 모순되므로, 반드시 $a_0=0$이다. 이제 (4)에서
\[
\begin{array}{@{}l@{\quad}l@{}}
(6)& \displaystyle a_1(\eta)\cdot\chi_\nu(\eta)\equiv c_1^{(\nu)}(\eta)\;\mathrm{mod.}\;\mathfrak M_2;\qquad
[a_1(\eta)]^{\nu\varkappa}\cdot A(\eta)
\equiv[c_1^{(\nu)}(\eta)]^{\nu\varkappa}\;\mathrm{mod.}\;\mathfrak M_{\nu\varkappa+1};
\end{array}
\]
을 얻는다. $A(\eta)$가 차수 적어도 1인 항들로 시작하므로 여기서 $c_1^{(\nu)}(\eta)=0$이 따른다. 그러므로 (5)와 완전히 마찬가지로 각 $\nu$에 대하여
\[
 a_1(\eta)\cdot\chi_\nu(\eta)\equiv0\;\mathrm{mod.}\;\mathfrak M_2;\qquad
 [a_1(\eta)]^{\nu\varkappa}A(\eta)\equiv0\;\mathrm{mod.}\;\mathfrak M_{2\nu\varkappa}
\]
이고, 여기서 $a_1(\eta)=0$을 얻는다. 이 과정을 계속하여 (5)와 (6)을 번갈아 적용하면
\[
 a_0=0,\qquad a_1(\eta)=0,\quad \ldots,\quad a_{\sigma-1}(\eta)=0
\]
을 얻는다. $\mathfrak G_\tau$에 관해서도 같은 논법이 성립하므로, \emph{$\mathfrak G_\sigma$와 $\mathfrak G_\tau$ $(\sigma\ne\tau)$ 사이의 모든 동형 대응에서 모듈 $\mathfrak M_\sigma$와 $\mathfrak M_\tau$는 서로 동형으로 대응한다.}

이를테면 $\sigma>\tau$라고 하자. 부정원 $\xi$에는 $\mathfrak G_\sigma$의 양 $f(\eta)$가 대응했다. $\mathfrak G_\tau$의 모든 양은 $\mathfrak M_\tau$를 법으로 하여 $\xi$의 다항식으로 나타낼 수 있고, $\mathfrak M_\tau$와 $\mathfrak M_\sigma$가 서로 대응하므로 $\mathfrak G_\sigma$의 모든 양은 $\mathfrak M_\sigma$를 법으로 하여 $f(\eta)$의 다항식으로 나타난다. $\sigma>\tau\geqq1$이므로 부정원 $\eta$는 분명 $\mathfrak M_\sigma$에 속하지 않는다. 따라서 $\eta$를 $\mathfrak M_\sigma$를 법으로 하여 $f(\eta)$의 다항식으로 나타낼 수 있으려면, 선형항이 사라지지 않는 $f(\eta)$가 있어야 한다. 그러한 대응을
\[
 \zeta:f(\eta)\equiv a_0+a_1(\eta)\;\mathrm{mod}\;\mathfrak M_2\quad[a_1(\eta)\ne0],
\]
\[
 \xi^\tau:[f(\eta)]^\tau\equiv a_0^\tau\pmod{\mathfrak M_1}\quad\text{$a_0\ne0$일 때},
\]
\[
 \hphantom{\zeta^\tau:[f(\eta)]^\tau}\equiv [a_1(\eta)]^\tau
 \pmod{\mathfrak M_{\tau+1}}\quad\text{$a_0=0$일 때}
\]
라고 쓰자. $\zeta^\tau$는 $\mathfrak M_\tau$에 속한다. 그러나 $[f(\eta)]^\tau$는 사라지지 않는 0차항 또는 $\tau$차항으로 시작하고 $\tau<\sigma$이므로 $\mathfrak M_\sigma$에 속할 수 없다. 이는 앞서 얻은 $\mathfrak M_\sigma$와 $\mathfrak M_\tau$의 대응과 모순이다. $\tau>\sigma$인 경우는 $\eta$와 $\xi$를 맞바꾸면 함께 처리되므로, 영역 $\mathfrak G_\sigma$와 $\mathfrak G_\tau$가 \emph{서로 동형이 아니며, 곧 본질적으로 서로 다름}이 증명되었다. \emph{따라서 영역 $\mathfrak G_i$ $(i=1,2,\cdots)$는 영역 $\mathfrak G$의 가산무한히 많은 동형류의 한 예를 준다.}

\emph{주.} 이처럼 서로 다른 두 영역 $\mathfrak G_i$는 동형이 아니지만, 두 영역 각각이 다른 영역의 한 부분정역과 동형일 수는 있다.\srcfn{*)}{이러한 현상은 체에서도 나타날 수 있다. Steinitz, 앞의 논문 끝부분 참조.} 이를테면 $\tau=1$이고 $\sigma$는 임의라고 하자.
\[
\mathfrak G_1:a_0+\mathfrak M_1(\zeta);\qquad
\mathfrak G_\sigma:a_0+a_1(\eta)+\cdots+a_{\sigma-1}(\eta)+\mathfrak M_\sigma(\eta).
\]
대응 $\zeta=\eta$에 따라 $\mathfrak G_\sigma$는 $\mathfrak G_1$의 진부분정역이 된다. 역으로 $\mathfrak G_1$과 $\mathfrak G_\sigma$에서 일대일인 대응 $\zeta=\eta^\sigma$, $\eta=\sqrt[\sigma]{\zeta}$에 따라 $\mathfrak G_1$은 $\mathfrak G_\sigma$의 진부분정역이 된다. 따라서 각 $\mathfrak G_\sigma$는 자기 자신의 진부분정역과도 동형이다. 그러므로 교집합이라는 개념과 달리 교집합 동형류, 곧 그 동형류의 각 영역이 다른 모든 동형류의 각 영역의 진부분정역과 동형인 그러한 동형류는 존재하지 않는다. 적어도 유일하게 정해지는 것으로는 존재하지 않는다.

\section*{§ 10. 임의의 체의 정수적 양.}

앞 절들에서 전개한 내용은, 영역 $\mathfrak G$의 정의 성질 III과 IV에 대수적 수를 포함시키지 않는 한, $\mathfrak G$에 관해서는 모든 복소수의 전체가 체를 이룬다는 사실에만 의존했다.\srcfn{**)}{§ 7에서만 유리수의 ``소체''가 무한히 많은 원소를 갖는다는 추가 가정을 사용했다. 원소가 유한개뿐인 소체의 경우에는, 뒤에서 보듯이 이 절이 단순히 빠진다.} 영역 $\mathfrak H$에 관해서는 이 체가 대수적으로 닫혀 있다는 사실을 추가로 사용했다. 따라서 얻은 결과들은 Weber와 E. Steinitz가\srcfn{*)}{앞의 논문 서론.} ``결합법칙과 교환법칙을 따르고 분배법칙으로 연결되며, 제한 없이 유일하게 역연산을 할 수 있는 두 연산, 곧 덧셈과 곱셈을 갖춘 원소들의 계''라고 추상적으로 정의한 \emph{가장 일반적인 체}로 곧바로 일반화된다.\srcfn{**)}{영으로 나누는 것만은 제외해야 한다.}

그러한 모든 체는 단위원을 포함하고, 따라서 단위원에서 덧셈과 곱셈 및 그 역연산으로 이끌어 낸 ``소체''를 포함한다. 이 소체는 유리수형(표수 0)이거나 소수 $p$에 대한 잉여류계형(표수 $p$)이다.\srcfn{***)}{Steinitz, § 4.} 이때 \emph{소체의 정수적 양}은 단위원에서 덧셈과 뺄셈으로 생기는 양으로 잘 정의된다. 표수 $p$의 소체에서는 정수적 양이 이미 소체 전체를 소진하므로 소체의 분수형 양은 존재하지 않는다. 모든 대수적으로 닫힌 체는 소체에 관하여 대수적인 모든 원소를 소체에 첨가하여 얻는 ``절대대수체''를 포함한다. 따라서 표수 0인 체에서는 이 체가 모든 대수적 수의 체와 같은 유형이다. \emph{절대대수체의 대수적 정수인 양}은 소체의 정수적 양에 대하여, 또는 같은 말로 단위원에 대하여 대수적으로 정수인 모든 양으로 잘 정의된다. 그러므로 표수 $p$의 체에서는 절대대수체의 분수형 대수적 양도 존재하지 않는다. 이러한 예비 고찰에 따라 ``정수적'' 양과 ``대수적으로 정수인'' 양의 정의를 각각 임의의 체와 임의의 대수적으로 닫힌 체로 옮길 수 있다. \emph{체 $\Omega$의 정수적 양은 $\Omega$의 양들로 이루어진 정역 $\mathfrak G$로 정의한다. 이 정역의 분수체}\srcfn{$\dagger$)}{즉 $\mathfrak G$의 두 양의 모든 몫으로 이루어진 체.} \emph{는 $\mathfrak K$ 전체와 같고, 이 정역은 단위원을 포함하지만 소체의 분수형 양은 포함하지 않는다.}

\emph{대수적으로 닫힌 체 $\mathfrak L$의 대수적 정수인 양은 $\Omega$의 양들로 이루어지고 대수적 정수성에 관하여 닫힌 정역 $\mathfrak H$로 정의한다. 이 정역의 분수체는 $\mathfrak L$ 전체와 같고, 이 정역은 단위원을 포함하지만 절대대수체의 분수형 대수적 양은 포함하지 않는다.}

표수 0인 체에 대해서는 앞 절들에서 얻은 결과들이 글자 그대로 성립한다. 다만 ``대수적 수''를 각각 소체 또는 절대대수체의 양으로 바꾸면 된다. 표수 $p$인 체에서는 소체에 분수형 양이 없기 때문에 결과들이 훨씬 더 단순해진다. 유리기저와 첨가할 계에 관한 부가조건이 완전히 사라지는 것이다.\srcfn{*)}{이 절의 첫 주석 참조.} 따라서 결과는 다음과 같다. \emph{추상적 정의에 따라 표수 $p$인 체의 정수적 양으로 허용되는 것은 임의의 유리기저에서 생기는 각 정역, 체 자체, 그리고 이 정역들과 체 사이에 놓인 모든 정역이다. 이에 대응하여 소수 표수의 대수적으로 닫힌 체에서 대수적 정수인 양으로 간주되는 것은, 초월기저에서 생긴 정역과 체 자체 사이에 놓인 대수적 정수성에 관하여 닫힌 각 정역이며, 두 경계 영역도 포함된다.} 따라서 소수 표수의 가장 일반적인 체에서는 그 소체에서와 마찬가지로 정수적인 것과 분수형인 것의 차이가 희미해진다.

\medskip
\noindent 에를랑겐, 1915년 3월 30일.

\clearpage
\setcounter{footnote}{0}
\phantomsection
\addcontentsline{toc}{section}{10. 동형사상의 함수방정식}
%
%
%
%

\section*{10. 동형사상의 함수방정식}
\begin{center}
{\Large\bfseries 10. 동형사상의 함수방정식.}\par
\vspace{0.75em}
저자\par
\vspace{0.75em}
괴팅겐의 Emmy Noether.\par
\vspace{1em}
Math. Ann. 77 (1916), 536--545쪽
\end{center}

\vspace{3em}

데데킨트에 따르면\srcfn{*)}{예를 들어 Dirichlet--Dedekind, \emph{Vorlesungen über Zahlentheorie}, 제4판, 부록 XI, § 161 참조. 군론을 본뜬 ``동형사상'' 또는 ``동형''이라는 명칭은 후대 문헌에서야 나타난다. 데데킨트는 ``체의 치환''이라고 말한다.}, 체 \(\Afield\)를 한 계 \(\Bfield\) 위로 보내는 \emph{동형사상}이란 다음과 같은 대응을 뜻한다. 뒤에 가서 \(\Bfield\) 역시 체임이 드러나며, \emph{\(\Afield\)의 각 원소에는 \(\Bfield\)의 원소가 하나씩 유일하게 대응하고, \(\Afield\)의 임의의 두 원소의 합·차·곱·몫에는 언제나 \(\Bfield\)에서 유일하게 대응하는 원소들의 합·차·곱·몫이 대응한다.}

그러한 사상이 위의 조건에 따라 단일값인 함수 \(f(z)\)로 주어진다고 하자. \(z\)가 체 \(\Afield\)의 모든 원소 \(x,y,\cdots\)를 달릴 때, \(f(z)\)는 다음 함수방정식으로 특징지어진다.
\[
\begin{array}{ll}
(1)\quad f(x+y)=f(x)+f(y); & (2)\quad f(x-y)=f(x)-f(y);\\[0.45em]
(3)\quad f(x\cdot y)=f(x)\cdot f(y); & (4)\quad f\!\left(\dfrac{x}{y}\right)=\dfrac{f(x)}{f(y)}.
\end{array}
\]
\emph{이 글에서는 \(x,y,\cdots\)가 모든 실수와 복소수를 달릴 때 이 함수방정식들의 가장 일반적인 단일값 해 \(f(z)\), 즉 모든 복소수의 체를 동형적으로 사상하는 가장 일반적인 함수를 제시한다.}\srcfn{**)}{이 문제는 란다우 선생이 언젠가 내게 제기한 것이다. 나중에 알게 되었지만, 해의 일부는 이미 H. Lebesgue, ``Sur les transformations ponctuelles, transformant les plans en plans'', Atti Acad. Torino, 1906/07에 있으며, A. Ostrowski, ``Über einige Fragen der allgemeinen Körpertheorie'', Crelles Journal 143 (1913), § 2에도 암묵적으로 들어 있다.} 임의의 체에 대한 해도 완전히 같은 방식으로 얻어진다.

G. Hamel은\srcfn{*)}{G. Hamel, ``Eine Basis aller Zahlen und die unstetigen Lösungen der Funktionalgleichung: \(f(x+y)=f(x)+f(y)\)'', Math. Ann. 60 (1905), 459쪽.} 모든 수의 \emph{선형}기저를 이용하여 함수방정식 (1)의 가장 일반적인 해를 구성했다. 여기서 제시하는 (1)--(4)의 해, 곧 유리연산과 대수연산을 보존하는 사상함수의 구성은 모든 수의 \emph{유리기저}와 \emph{초월기저}를 이용한다. 1에서 동형사상을 좀 더 정확히 논한 뒤, 2에서는 \(f(z)\)에 필요한 조건을 제시하고, 3에서는 이 조건들이 충분함을 밝혀 실제 구성에 이른다.\srcfn{**)}{내 논문 ``Die allgemeinsten Bereiche aus ganzen transzendenten Zahlen'', Math. Ann. 77 (1916), 103쪽, 특히 § 8과 § 10의 병행하는 논의를 참조하라.} 잘 알려져 있듯이 (1)--(4)의 해는 \(f(z)=z\)와 \(\overline z\)를 제외하면 불연속이다. 4에서는 더 나아가 그 해들이 ``극도로 불연속''임을 보인다. 이는 G. Hamel이 (1)의 실수값 해에 대하여 증명한 사실이지만, 복소수의 경우에는 일반적인 덧셈방정식의 해가 반드시 이 성질을 갖는 것은 아니다.\srcfn{***)}{Hamel이 쓴 ``완전 불연속''이라는 명칭은 다른 문헌에서는 다른 뜻으로 쓰인다.}

\medskip
\noindent\textbf{1.}\quad 먼저 함수방정식 (1)--(4)에서 나오는 몇 가지 귀결을 일반적인 체 \(\Afield\)에 대하여 직접 살펴보자. (1)과 단일값성에서 \(f(0)=0\)이 따른다. 따라서 \(f(y)\)이면 원래의 \(y\)도 0일 수 없고, 함수방정식 (1)--(4)에 의해 \emph{모든 값 \(f(z)\)로 이루어진 상집합 \(\Bfield\)는 다시 체를 이룬다.} 거꾸로 초기조건 \(f(0)=0\)을 부과하면 (2)에서 함수 \(f(z)\)의 단일값성이 따른다. 그러므로 \emph{단일값성의 요구와 초기조건 \(f(0)=0\)은 동치이다.} (4)에서 \(f(z)\)가 항등적으로 0일 수 없음이 따르고, 다시 (3)에서 \(f(z)\)는 \(z=0\)일 때에만 0임이 따른다. 실제로 \(f(y)=0\), \(y\ne0\)이라면 \(\frac{z}{y}\) 역시 \(\Afield\)에 속하므로
\[
f(z)=f\!\left(\frac{z}{y}\cdot y\right)=f\!\left(\frac{z}{y}\right)\cdot f(y)=0
\]
이 되어, (4)에 모순되게 \(f(z)\)가 항등적으로 0이 된다. 따라서 \(f(x)=f(y)\)이면 \(x=y\)이다. 즉 \(f(z)\)의 각 값에도 \(z\)의 값이 정확히 하나씩 대응한다. \emph{따라서 \(f(z)\)는 \(z\)의 일대일 함수이다.} 함수방정식에서 알 수 있듯이 역함수가 주는 사상도 다시 동형사상이다. 모든 유리연산은 유한번의 덧셈과 곱셈 및 그 역연산으로 합성되므로, 다음과 같이 말할 수도 있다. \emph{동형사상은 \(\Afield\)의 원소와 \(\Bfield\)의 원소 사이에 일대일 대응을 세우며, \(\Afield\)의 유한개의 원소 사이에 성립하는 모든 유리관계}
\[
F(x,y,\cdots,t)=0
\]
\emph{는 \(\Bfield\)에서도 보존되고 그 역도 성립한다.}

체에 대해서는 함수방정식 (1)과 (3), \(f(z)\)가 항등적으로 0이 아니라는 조건, 초기조건 \(f(0)=0\)만으로도 \(f(z)\)가 이미 특징지어진다는 점도 덧붙여 둔다. 실제로 (1)에서 \(x\)를 \(\Afield\)에 속하는 \((x-y)\)로 바꾸면 (2)가 나오고, (2)와 \(f(0)=0\)에서 \(f(z)\)의 단일값성이 나온다. 위에서 보았듯이 (3)에서 \(x\)를 \(\Afield\)에 속하는 \(\frac{x}{y}\)로 바꾸면 \(y\ne0\)일 때 \(f(y)\ne0\)임이 나오므로 (4)와 일대일성이 함께 성립한다. 체 대신 정역 \(\Jdom\)을 바탕으로 삼으면 \(\frac{x}{y}\)가 \(\Jdom\)에 속하지 않을 수 있으므로 (3)에서 일대일성을 결론 내릴 수 없다. 이 경우에는 다음과 같은 가장 널리 알려진 서술이면 충분하다. 동형사상이란 \(\Jdom\)과 \(\Jdom'\)의 원소 사이의 \emph{일대일} 대응으로서, 합과 곱에 언제나 합과 곱이 대응하는 것이다.\srcfn{*)}{1에 대해서는 Dedekind, 앞의 책 § 161 참조.}

\medskip
\noindent\textbf{2.}\quad 이제부터는 간단히 하기 위하여 \(\Afield\) 대신 모든 실수와 복소수의 체 \(\Cfield\)를 바탕으로 삼는다. 우선 \(f(z)\)가 취하는 값들의 집합을 논하자. (3) 또는 (4)에서 \(f(1)=1\)이 나오고, 함수방정식들에 의해 모든 \emph{유리수} \(\alpha\)에 대하여 \(f(\alpha)=\alpha\)가 나온다. 따라서 \emph{모든 유리수는 사상에서 자기 자신에 대응한다.} \emph{대수적 수} \(\beta\)는 유리수 계수의 기약방정식 \(F(\beta,\alpha)=0\)을 만족하므로, 1과 방금 한 관찰에 따라 \(F(f(\beta),\alpha)=0\)이다. 따라서 \emph{각 대수적 수는 자기 자신 또는 자기의 한 켤레수로 옮겨가며}, 일대일성에 의해 대수적 수 전체에는 다시 모든 대수적 수의 체가 대응한다. \emph{초월수}의 거동을 알아보고 동시에 켤레값들을 구별하기 위하여 모든 수의 \emph{유리기저}\srcfn{**)}{위에서 인용한 ``ganze transzendente Zahlen''에 관한 내 논문의 § 6 참조.}를 택한다. 이는 모든 수를 유리수 계수의 유리식으로 표현하게 하는 수 \(\vartheta\)들의 계 \(\Theta\)로서, 적어도 하나의 정렬순서에서는 어떤 기저수도 그보다 앞선 유한개의 수로 유리적으로 표현되지 않는 것이다.\srcfn{***)}{여기서 ``유리함수''는 그 계수도 유리수인 것만을 뜻한다.} 이 기저의 존재는 선형기저와 초월기저의 경우와 꼭 마찬가지로 정렬정리에서 따른다.

\(\Cfield\)에 정렬순서 \(\Omega\)를 주자. 그러면 각 수는 앞선 유한개의 수로 유리적으로 표현되거나, 그렇지 않으면 앞선 수들에 대하여 유리적으로 독립이다. 후자의 수들이 기저 \(\Theta\)를 이루며, 초한귀납법으로 모든 수가 실제로 유한개의 \(\vartheta\)로 유리적으로 표현됨을 알 수 있다.

각 기저수 \(\vartheta\)에 대하여 ``절단체 \(\Kfield(\vartheta)\)''를 \(\vartheta\)보다 정렬순서에서 앞선 기저수들의 모든 유리적 결합으로 정의한다. 첫 번째 기저수의 절단체는 모든 유리수의 체 \(\Rfield\)로 본다. \(\vartheta\)는 \(\Kfield(\vartheta)\)에 대하여 두 가지 거동을 보일 수 있다. 곧 \(\Kfield(\vartheta)\) 위에서 대수적으로 종속이거나 대수적으로 독립이다. 1에 따라 \(\Cfield\)에서 성립하는 모든 관계가 상의 체에서도 성립하고 그 역도 성립하므로, \(\vartheta\)들에 대응하는 값 \(f(\vartheta)\)의 전체는 상의 영역의 기저를 이루며 \(\Theta\)의 정렬순서에 의해 그 자체도 정렬된다. 특히 절단체 \(\Kfield(\vartheta)\)에는 절단체 \(\Kfield(f(\vartheta))\)가 대응한다. \(\vartheta\)가 \(\Kfield(\vartheta)\) 위에서 대수적으로 종속인지 독립인지에 따라 \(f(\vartheta)\)도 \(\Kfield(f(\vartheta))\) 위에서 같은 성질을 가지며, 종속인 경우에는 \(\vartheta\)와 \(f(\vartheta)\)의 기약방정식들도 서로 대응한다. 각각 \(\Kfield(\vartheta)\) 위에서 대수적으로 독립인 값 \(\vartheta\)들의 전체는, 다시 초한귀납법으로 알 수 있듯이, 모든 수의 \emph{초월기저}\srcfn{*)}{E. Zermelo, ``Über ganze transzendente Zahlen'', § 1, Math. Ann. 75 (1914) 참조.}를 이룬다. 즉 이는 모든 수를 대수적으로 표현할 수 있게 하면서 그 원소들 사이에는 아무 대수관계도 없는 수 \(\eta\)들의 계 \(H\)이다. 이 초월기저 \(H\)에는 상의 영역에서 다시 초월기저 \(\mathsf{Z}\)가 대응하며, 일대일성 때문에 \(\mathsf{Z}\)의 기수는 \(\mathsf{H}\)와 같다. 이 기수는 연속체의 기수인데, 대수적 확대는 기수를 증가시키지 않는다는 것이 잘 알려져 있기 때문이다.\srcfn{**)}{E. Steinitz, ``Algebraische Theorie der Körper'', Crelles Journal 137 (1910), § 24 참조.} 값 \(f(\vartheta)\)를 알면 모든 값의 집합에서 \(f(z)\)가 곧바로 정해진다. 실제로 \(z=\psi(\vartheta)\)이면 언제나 \(f(z)=\psi(f(\vartheta))\)이다.

\medskip
\noindent\textbf{3.}\quad 이제 앞에서 찾은 필요조건들이 실제로 \(f(z)\)를 구성하기에도 충분함을 보이자. 초월기저 \(\mathsf{H}\)에 일대일로 대응시킬 기저 \(\mathsf{Z}\)를 다음과 같이 구성하여 \(\mathsf{H}\)와 짝짓는다. 두 번째 정렬순서 \(\Omega'\)를 택하여 모든 수의 초월기저 \(\mathsf{Z}'\)를 얻는다. 원소를 \(\xi\)라 쓰는 \(Z\)를 \(\mathsf{Z}'\)와 같은 기수, 따라서 \(H\)와 같은 기수를 갖는 \(\mathsf{Z}'\)의 부분집합으로 잡는다. 이제 \(\mathsf{H}\)의 각 원소 \(\eta_i\)에 같은 첨자의 \(\xi_i\)를 다음 규칙으로 일대일 대응시킨다. \(\xi_i\)는 \(\Omega'\)에서, \(H\)에서 \(\eta_i\)보다 앞선 어떤 원소에도 아직 대응하지 않은 \(Z\)의 원소들 가운데 \emph{첫 번째} 원소이다.

이제 \(f(z)\)를 다음과 같이 \emph{정의한다.}
\begin{enumerate}
\item[(a)] \emph{모든 유리수} \(\alpha\)\emph{에 대하여} \(f(0)=0\), \(f(\alpha)=\alpha\)로 둔다.
\item[(b)] \emph{초월기저} \(\mathsf{H}\)\emph{의 모든 원소에 대하여} \(f(\eta_i)=\xi_i\)로 둔다.
\item[(c)] \emph{유리기저} \(\Theta\)\emph{의 나머지 모든 원소} \(\vartheta\)\emph{에 대하여 다음과 같이 둔다. 이 원소들은 절단체} \(\Kfield(\vartheta)\) \emph{위에서 대수적으로 종속이므로} \(\Kfield(\vartheta)\)\emph{에서 기약인 방정식} \(F(\vartheta;\vartheta_{i_1}\cdots\vartheta_{i_\tau})=0\)\emph{을 만족한다. 이때} \(f(\vartheta)\)\emph{를 방정식} \(F(f(\vartheta);f(\vartheta_{i_1})\cdots f(\vartheta_{i_\tau}))=0\)\emph{의 한 근으로 정한다.}
\item[(d)] \emph{나머지 모든 값} \(z\)\emph{는} \(\Theta\)\emph{의 정의에 따라} \(z=\psi(\vartheta_{k_1}\cdots\vartheta_{k_\sigma})\)\emph{로 나타낼 수 있으므로,} \(f(z)=\psi(f(\vartheta_{k_1})\cdots f(\vartheta_{k_\sigma}))\)\emph{로 정한다.}\srcfn{*)}{(d)는 (a)를 특수한 경우로 포함한다.}
\end{enumerate}
이렇게 정의한 \(f(z)\)가 실제로 함수방정식 (1)--(4)를 만족함을 보이려면, 1에 따라 (1)과 (3)만 확인하면 충분하다. \(\Cfield\)는 체이고, \(f(z)\)는 (a)와 (b)에 의해 항등적으로 0이 아니며 초기조건 \(f(0)=0\)도 만족하기 때문이다. 유리기저 \(\Theta\)를 사용하여 \(x,y,x+y\)를 다음과 같이 나타내자.
\[
x=\psi_1(\vartheta_1\cdots\vartheta_t),\quad
y=\psi_2(\vartheta_1\cdots\vartheta_t),\quad
(x+y)=\psi_3(\vartheta_1\cdots\vartheta_t).
\]
\(f(z)\)가 함수방정식 (1), 곧
\[
f(x)+f(y)-f(x+y)=0
\]
을 만족한다는 것은 (d)에 따라 다음 함의와 같다.
\[
\psi_1(\vartheta)+\psi_2(\vartheta)-\psi_3(\vartheta)
=\frac{H_1(\vartheta)}{H_2(\vartheta)}=0
\]
이면 언제나
\[
\psi_1(f(\vartheta))+\psi_2(f(\vartheta))-\psi_3(f(\vartheta))
=\frac{H_1(f(\vartheta))}{H_2(f(\vartheta))}=0
\]
이다. 여기서 \(H_1,H_2\)는 그 인수들의 다항식이다. 함수방정식 (3)도 정확히 같은 꼴의 조건으로 귀착된다. 따라서 모든 함수방정식이 성립한다는 것은 다음 조건, 곧 1과 2에 따라 필요하기도 한 조건과 같다. \emph{모든 다항식 \(H(\vartheta)\)에 대하여 \(H(\vartheta)=0\)이면 언제나 \(H(f(\vartheta))=0\)이고, \(H(\vartheta)\ne0\)이면 언제나 \(H(f(\vartheta))\ne0\)이다.} 뒤의 조건은 분모가 나타나기 때문에 필요하다.

초한귀납법으로 이 조건이 실제로 성립함을 보일 수 있다. \(H(\vartheta_1\cdots\vartheta_t)\)에서 \(\vartheta_t\)가 \(\vartheta_1,\cdots,\vartheta_t\) 가운데 정렬순서상 마지막인 기저원소라고 하자.\srcfn{**)}{\(\vartheta\)들에 관하여 0차인 식 \(H(\vartheta)\)는 사상에서 자기 자신으로 옮겨가므로 더 살필 필요가 없다.} 그러면 \(H(\vartheta_1\cdots\vartheta_t)=0\)은 \(\vartheta_t\)가 절단체 \(\Kfield(\vartheta_t)\) 위에서 만족하는 관계로 볼 수 있고, \(H(f(\vartheta_1)\cdots f(\vartheta_t))=0\)도 \(f(\vartheta_t)\)가 \(\Kfield(f(\vartheta_t))\) 위에서 만족하는 관계로 볼 수 있다. 모든 관계 \(H(\vartheta)=0\)이 \(f(\vartheta)\)에 대해서도 성립하고 그 역도 성립하는 것이 아니라면, 적어도 한 관계가 상의 체에서 보존되지 않거나 그 역이 보존되지 않는 최초의 기저원소 \(\vartheta_0\)가 있어야 한다. 그보다 앞선 절단체 \(\Kfield(\vartheta_0)\)와 \(\Kfield(f(\vartheta_0))\)는 서로 동형이다. 특히 (a)에 따라 \(\Theta\)의 첫 기저원소의 절단체, 즉 모든 유리수의 체 \(\Rfield\)는 그 상의 체와 동형이다. 이제 \(H(\vartheta_0)\)를
\[
H(\vartheta_0)=a_0(\vartheta)\cdot\vartheta_0^{\varkappa}
+a_1(\vartheta)\cdot\vartheta_0^{\varkappa-1}+\cdots+a_{\varkappa}(\vartheta)
\]
의 꼴로 쓰자. 여기서 \(a_i(\vartheta)\)는 \(\Kfield(\vartheta_0)\)의 원소이다. 두 경우가 있다.

1) \(\vartheta_0\)가 \(\Kfield(\vartheta_0)\) 위에서 대수적으로 독립인 경우. 그러면 \(H(\vartheta_0)=0\)은 \(\vartheta_0\)에 관한 항등식으로 성립하고,\srcfn{*)}{\(x,y,\cdots\)를 \(\vartheta\)들로 표현할 수 있기 때문에 이런 꼴의 관계도 충분히 나타날 수 있다.} 모든 계수에 대하여 \(a_i(\vartheta)=0\)이다. 따라서 \(\Kfield(\vartheta_0)\)와 \(\Kfield(f(\vartheta_0))\)의 동형성에 의해 \(a_i(f(\vartheta))=0\)이다. 그러므로 \(H(\vartheta_0)=0\)이면 언제나 \(H(f(\vartheta_0))=0\)이다. 거꾸로 \(H(f(\vartheta_0))=0\)이라고 하자. \(\vartheta_0\)는 \(\Kfield(\vartheta_0)\) 위에서 대수적으로 독립이므로 초월기저 \(\mathsf{H}\)에 속한다. 따라서 (b)에 의해 \(f(\vartheta_0)\)는 초월기저 \(\mathsf{Z}\)에 속하고 \(\Kfield(f(\vartheta_0))\) 위에서 대수적으로 독립이다. 그러므로 모든 \(a_i(f(\vartheta))\)가 0이고, 동형성에 의해 모든 \(a_i(\vartheta)\)도 0이다. 따라서 \(H(f(\vartheta_0))=0\)에서 \(H(\vartheta_0)=0\)이 따르며, 동치로 \(H(\vartheta_0)\ne0\)이면 \(H(f(\vartheta_0))\ne0\)이다.

2) \(\vartheta_0\)가 \(\Kfield(\vartheta_0)\) 위에서 대수적으로 종속인 경우. 그러면 \(H(\vartheta_0)\)는 \(\Kfield(\vartheta_0)\) 위에서 기약인 방정식 \(F(\vartheta_0)\)으로 나누어지며, 다시 말해 \(t\)에 관하여 항등적으로
\[
H(t;a_i(\vartheta))=F(t,\vartheta_i)\cdot G(t,\vartheta_i)
\]
가 성립한다. 나타나는 계수들이 모두 \(\Kfield(\vartheta_0)\)에 속하므로, 그와 동형인 체 \(\Kfield(f(\vartheta_0))\)에서는 대응하는 관계
\[
H(t;a_i(f(\vartheta)))=F(t;f(\vartheta))\cdot G(t;f(\vartheta))
\]
가 성립한다. 그런데 (c)에서 \(f(\vartheta_0)\)를 \(F(t;f(\vartheta))\)의 근으로 택했으므로, \(H(\vartheta_0)=0\)에서 \(H(f(\vartheta_0))=0\)이 따른다.

거꾸로 \(H(f(\vartheta_0))=0\)이라고 하자. 그러면 \(H(t;a_i(f(\vartheta)))\)가 \(\Kfield(f(\vartheta_0))\) 위에서 기약인 다항식 \(F(t;f(\vartheta))\)으로 나누어진다는 데서, 동형인 체 \(\Kfield(\vartheta_0)\)에 대하여 \(H(t;a_i(\vartheta))\)가 \(F(t;\vartheta)\)로 나누어짐이 똑같이 따른다. \(F(t;f(\vartheta))\)는 \(\vartheta_0\)의 기약방정식에서 동형관계로 생겨났고 이 관계는 일대일이므로, \(F(t;\vartheta)\)는 반드시 \(\vartheta_0\)를 근으로 갖는 기약다항식이다. 따라서 \(H(f(\vartheta_0))=0\)에서 \(H(\vartheta_0)=0\)이 따르며, 동치로 \(H(\vartheta_0)\ne0\)이면 \(H(f(\vartheta_0))\ne0\)이다.

따라서 서로 동형인 절단체 \(\Kfield(\vartheta_0)\)와 \(\Kfield(f(\vartheta_0))\)에 각각 \(\vartheta_0\)와 \(f(\vartheta_0)\)를 첨가하면 다시 서로 동형인 체가 생긴다. 어떤 \(\vartheta_0\)에서 이것이 실패한다는 가정은 모순에 이른다. 이로써 \emph{(a)--(d)로 정의한 함수 \(f(z)\)가 실제로 함수방정식 (1)--(4)의 해이며, 2에 따라 가장 일반적인 해임}이 증명되었다.

\(f(z)\)를 구성하는 데 쓴 모든 논의는 모든 수의 집합 \(\Cfield\)가 체라는 사실만을 사용했다. 따라서 그 논의는 추상적으로 정의된 임의의 체 \(\Afield\)에도 적용된다. 이때 유리수체 \(\Rfield\)에는 \(\Afield\)의 ``소체'', 즉 \(\Afield\)의 단위원에서 덧셈과 곱셈 및 그 역연산으로 생성되는 체가 대응한다. 그러므로 (a)--(d)에서 유리수를 소체의 원소로 바꾸면, \emph{그렇게 얻은 함수 \(f(z)\)는 추상적으로 정의된 임의의 체의 가장 일반적인 동형사상을 준다.}

\medskip
\noindent\textbf{4.}\quad 이제 모든 실수와 복소수의 체 \(\Cfield\)에 대하여 정의한 함수 \(f(z)\)는, 잘 알려진 연속함수 \(f(z)=z\)와 \(\overline z\)를 제외하면, \emph{극도로 불연속}임을 보이자. 이는 모든 실수 또는 복소수 \(z_0\)의 근방에, \(f(z)\)가 임의로 미리 지정한 값 \(Z_0\)에 임의로 가까워지는 값 \(z\)가 존재한다는 뜻이다. G. Hamel이 앞의 논문에서 보였듯이, 이런 극도의 불연속성은 모든 \emph{실수}에 정의된 함수방정식 (1)의 \emph{실수값} 해 \(\varphi(x)\)가 갖는다. 그러나 아래의 예들이 보이듯이 복소수에 정의된 해는 반드시 그렇지 않다. Hamel의 논법은 다음과 같다. \(\varphi(x)\)가 \(C\cdot x\)와 다르면, 선형기저의 적어도 두 값\srcfn{*)}{모든 실수의 선형기저란 모든 실수를 유리수 계수의 선형결합으로 표현할 수 있게 하면서, 임의의 유한개의 기저수 사이에는 유리수 계수의 선형관계가 없는 실수들의 계이다.}, 이를테면 \(x_1,x_2\)가 있어서 \(\varphi(x_1):x_1\ne\varphi(x_2):x_2\)이다. 그러면 두 방정식
\[
\alpha x_1+\beta x_2=a;\qquad \alpha\varphi(x_1)+\beta\varphi(x_2)=b
\]
은 실수 \(\alpha,\beta\)에 관한 해를 갖는다. 따라서 유리수 \(\alpha,\beta\)를 사용하여 \(a,b\)를 임의로 가깝게 근사할 수 있다. 그런데 \(\alpha,\beta\)가 유리수이면 \(\alpha\varphi(x_1)+\beta\varphi(x_2)\)는 실제로 \(\varphi(\alpha x_1+\beta x_2)\)이다. 복소수값의 집합에서는 이 논법이 실패한다. 실제로 복소수에 정의된 (1)의 해 가운데 불연속이지만 극도로 불연속은 아닌 것을 줄 수 있다. 예를 들면
\[
\varphi(z)=\varphi(x+iy)=\varphi(x)+i(y+\varphi(x))
\]
이며, 여기서 \(\varphi(x)\)는 실수에서 극도로 불연속인 실수값 해이다. 더 간단한 예는 \(\varphi(z)=\varphi(x)\)이다. 두 경우 모두 \(\varphi(z)\)의 실수부는 극도로 불연속이지만 허수부는 실수부에 의해 함께 결정된다.

랭크의 개념을 도입하면 이 거동과 함수 \(f(z)\)의 거동이 모두 분명해진다. \(z=x+iy\), \(\varphi(z)=X+iY\)로 쓰고, 실수 \(c_1,c_2,c_3,c_4\)에 관한 선형독립인 관계
\[
c_1x+c_2y+c_3X+c_4Y=0
\]
가 각각 0개, 1개, 2개 또는 3개 존재함에 따라 \(\varphi(z)\)의 랭크를 각각 4, 3, 2 또는 1이라고 부르자.

1) \(\varphi(z)\)의 \emph{랭크가 4}인 경우. 선형기저의 적어도 네 값 \(z_1,z_2,z_3,z_4\)를 택하여 행렬식
\[
\begin{vmatrix}
x_1&x_2&x_3&x_4\\
y_1&y_2&y_3&y_4\\
X_1&X_2&X_3&X_4\\
Y_1&Y_2&Y_3&Y_4
\end{vmatrix}
\]
이 0이 아니게 할 수 있다. 따라서 방정식
\[
\begin{aligned}
\alpha x_1+\beta x_2+\gamma x_3+\delta x_4&=a,\\
\alpha y_1+{}\cdot{}\ \cdot{}\ \cdot{}+\delta y_4&=b,\\
\alpha X_1+{}\cdot{}\ \cdot{}\ \cdot{}+\delta X_4&=A,\\
\alpha Y_1+{}\cdot{}\ \cdot{}\ \cdot{}+\delta Y_4&=B
\end{aligned}
\]
은 실수 \(\alpha,\cdots,\delta\)에 관한 해를 가지며, 유리수 \(\alpha,\beta,\gamma,\delta\)를 사용하여 \(a,b,A,B\)를 임의로 가깝게 근사할 수 있다. 또한
\[
\alpha(X_1+iY_1)+\beta(X_2+iY_2)+\gamma(X_3+iY_3)+\delta(X_4+iY_4)
=\varphi(\alpha z_1+\beta z_2+\gamma z_3+\delta z_4)
\]
이다. 따라서 ``일반적으로'' 해 \(\varphi(z)\)는 복소수에서도 극도로 불연속이며, \(z_0=a+ib\)의 근방에서 \(\varphi(z)\)의 \emph{불연속값들은 이차원 선형다양체를 이룬다.}

2) \(\varphi(z)\)의 \emph{랭크가 3}인 경우. 선형기저에서 적어도 세 값 \(z_1,z_2,z_3\)를 택하여 위 행렬식의 랭크가 3이 되게 할 수 있다. 관계
\[
c_1a+c_2b+c_3A+c_4B=0
\]
을 만족하는 값 \(a,b,A,B\), 그리고 그런 값들만 임의로 가깝게 근사할 수 있다. \emph{불연속값들은 \(z_0=a+ib\)가 정하는 일차원 선형다양체를 이룬다.} 위의 예들이 보이듯이 이 경우는 실제로 나타날 수 있다.

3) \(\varphi(z)\)의 \emph{랭크가 2}인 경우. 언제나 두 복소수 \(z_1,z_2\)를 택하여
\[
\begin{vmatrix}x_1&x_2\\ y_1&y_2\end{vmatrix}\ne0
\]
이 되게 할 수 있으므로, 관계들은
\[
X=\lambda_1x+\lambda_2y;\qquad Y=\mu_1x+\mu_2y
\]
의 꼴이다. 이로부터 얻는
\[
\varphi(z)=(\lambda_1x+\lambda_2y)+i\cdot(\mu_1x+\mu_2y)
\]
는 복소수 영역에서 함수방정식 (1)의 가장 일반적인 \emph{연속}해이다.

4) \(\varphi(z)\)의 \emph{랭크가 1}인 경우. 이 경우는 복소수 영역에서는 나타날 수 없고, \emph{실수} 영역에서는 연속해 \(\varphi(x)=C\cdot x\)를 뜻한다.

이제 모든 실수와 복소수에 정의된 함수방정식 (1)--(4)의 \emph{해 \(f(z)\)는 랭크 2 또는 4만을 가질 수 있음}을 보이자. 랭크 2인 경우에는 3)의 결과에 \(f(1)=1\), \(f(i)=\pm i\)를 대입하면 두 \emph{연속}함수 \(f(z)=z\)와 \(\overline z\)만 나온다. 랭크 1은 이미 위에서 배제했으므로 랭크 3을 배제하면 된다. 따라서 적어도 하나의 관계
\[
c_1x+c_2y+c_3X+c_4Y=0
\]
가 존재하여 \(f(z)\)의 랭크가 3 이하라고 가정하고, 이 가정 아래에서는 언제나 실제 랭크가 더 낮음을 보이자. 다시 \(f(1)=1\), \(f(i)=\pm i\)를 대입하면 이 관계는
\[
c_3(X-x)+c_4(Y\mp y)=0
\]
이 되며, 여기서 \(c_3,c_4\)는 동시에 0일 수 없다. 먼저 \(c_4=0\)이라고 하자. 관계 \((X-x)=0\)은 순허수 \(z\)에는 순허수 \(f(z)\)가 대응함을 뜻한다. 따라서 모든 실수 \(\mathfrak{y}\)에 대하여 \(f(i\mathfrak{y})=i\cdot\Yreal\)이며, 여기서 \(\Yreal\)도 실수이다. 그런데
\[
f(i\mathfrak{y})=\pm i f(\mathfrak{y})\quad\text{이므로}\quad f(\mathfrak{y})=\pm\Yreal
\]
이다. 따라서 실수 \(z\)에는 실수 \(f(z)\)가 대응하고, \(X-x=0\) 때문에 모든 실수는 자기 자신으로 옮겨간다. 그러므로 \(f(z)=z\) 또는 \(\overline z\)뿐이며, 실제 랭크는 2이다. \(c_3=0\), \(c_4\ne0\)인 경우도 완전히 같은 방식으로 결론을 얻는다.

이제 \(c_3,c_4\)가 모두 0이 아니라고 하자. 그러면 관계를
\[
(Y\mp y)=c\cdot(X-x)
\]
로 쓸 수 있으며, \(c\)는 0이 아닌 유한한 실수이다. 이 관계는 \(y=\pm cx\)이면 \(Y=c\cdot X\)가 됨을 뜻한다. 따라서 함수방정식들을 함께 사용하면 모든 실수 \(\mathfrak{x}\)에 대하여
\[
f(\mathfrak{x}\cdot(1\pm ic))=\Xreal(1+ic)=f(\mathfrak{x})\cdot(1+i f(c))
\]
이다. 여기서 \(\Xreal\)도 실수이다. 1에 따라 \(f(\mathfrak{x})\)에 곱해진 인자는 항등적으로 0일 수 없으므로 나누어
\[
f(\mathfrak{x})=\Xreal(\sigma+i\tau)
\]
를 얻는다. \(\sigma,\tau\)는 실수이며 \(x\)와 \(\Xreal\)에 무관하다. 특히 \(\mathfrak{x}=1\)에도 실수 \(\mathfrak{X}_0\)가 대응하고, 조건 \(1=\mathfrak{X}_0(\sigma+i\tau)\)는 반드시 \(\tau=0\)임을 보인다. 따라서 \(f(\mathfrak{x})=\sigma\cdot\Xreal\)이다. 즉 실수 \(z\)에는 실수 \(f(z)\)가 대응하며, \(y=0\)이면 반드시 \(Y=0\)이다. 그러면 실수 \(z\)에 대한 선형관계는 \(X-x=0\)이 된다.\srcfn{*)}{여기서는 \(c\ne0\), 즉 \(c_3,c_4\)가 모두 0이 아니라는 사실을 쓴다. 앞에서는 \(c_4\ne0\)만을 썼으므로 둘 중 하나가 0인 경우를 먼저 따로 다루어야 했다.} 모든 실수는 자기 자신에 대응하므로 \(f(z)=z\) 또는 \(\overline z\)뿐이고, 실제 랭크는 2이다. 이로써 모든 경우에 랭크 3이 배제되었다. 그런데 앞 절들의 논의가 보이듯이 불연속해는 실제로 존재하고, 따라서 그 랭크는 반드시 4이다. 그러므로 1)에 의해 다음 정리가 성립한다. \emph{``함수방정식 (1)--(4)의 모든 해 \(f(z)\)는, 연속해 \(f(z)=z\)와 \(\overline z\)를 제외하면, 실수와 복소수에서 극도로 불연속이다.''}\srcfn{**)}{Lebesgue는 앞의 논문에서 \(f(z)\)의 실수부와 허수부가 모두 \(x,y\)의 ``비가측'' 함수임을 보인다. 그러나 4의 처음에 든 예 \(\varphi(z)=\varphi(x)+i(y+\varphi(x))\)가 보이듯이, 이 사실만으로는 복소수에서의 극도의 불연속성이 따라오지 않는다.}

\medskip
\noindent 에를랑겐, 1915년 10월 30일.

\clearpage
\setcounter{footnote}{0}

\clearpage
\setcounter{footnote}{0}
\phantomsection
\addcontentsline{toc}{section}{11. 지정된 군을 갖는 방정식}
%
%
%
%

\section*{11. 지정된 군을 갖는 방정식}
\begin{center}
{\Large\bfseries 11. 지정된 군을 갖는 방정식.}\par
\vspace{0.75em}
저자\par
\vspace{0.75em}
괴팅겐의 Emmy Noether,\par
\vspace{1.5em}
Math. Ann. 78 (1918), 221--229쪽
\end{center}

\vspace{4em}

지정된 군을 갖는 방정식을 구성하는 문제에는 두 방향에서 접근할 수 있다. 간단히 말하면, 근을 특징짓는 ``무리적'' 방향과 계수를 특징짓는 ``유리적'' 방향이다.

함수론적·산술적으로 연구하는 ``무리적'' 방향에는, 유리수 영역의 모든 아벨체가 원분체라는 크로네커의 정리와 이차 수체에 대한 상대 아벨체의 대응하는 정리들이 있다. 이 정리들은 한편으로 이러한 특수한 유리성 영역에서 모든 아벨군을 실현할 수 있음을 보이고, 다른 한편으로 방정식 전체를 주면서 그 방정식들이 정의하는 수체의 구조를 더 깊이 드러낸다. 그러나 각각의 정리는 하나의 특수한 유리성 영역에 한정되며, 새로운 유리성 영역마다 완전히 새로운 처리가 필요하다.

대수적으로 연구하는 ``유리적'' 방향은 힐베르트의 불가약성 정리에 의존한다. 이 정리에 따르면 계수의 표현에서는 매개변수표현만을 생각해도 된다. 예를 들어 임의의 유리성 영역에 대하여 교대군을 갖는 방정식이 임의로 많이 존재한다는 사실이 이 방향에 속한다. 여기서는 방정식이 정의하는 체의 구조에 대한 위와 같은 통찰을 얻을 수 없는 것이 당연하다. 그러나 이 ``유리적'' 방향의 장점은 군의 실현가능성이 \emph{모든 유리성 영역에 대하여 동시에} 증명된다는 데 있다. 다만 지금까지 알려진 매개변수표현은 일반적으로 방정식 전체를 주지 않는다.\srcfn{*)}{가해방정식에서는 계수의 매개변수표현을 근의 매개변수표현으로 바꿀 수 있다. 최근 F. Mertens는 몇몇 8차 군에 대하여 가장 일반적인 근의 표현을 제시했다. ``Gleichungen 8ten Grades mit Quaternionengruppen'', Sitzb. d. Ak. d. Wiss., Wien, Abt. IIa, Bd. 125, 735쪽. 이 논문에 따르면 G. Bucht는 그보다 앞서 언제나 만들 수 있는 3차·4차 방정식의 정규형과 앞서 말한 군을 갖는 8차 방정식의 정규형에 대하여 가장 일반적인 근의 식을 제시했다. ``Über einige algebraische Körper achten Grades'', Arkiv for Math., Astron. och Physik, Bd. 6, Nr. 30. [교정 때의 추가.]}

이 글은 ``유리적'' 방향에 대한 한 기여이다. 곧 \emph{임의의 유리성 영역에서 지정된 군을 갖는 방정식 전체를 주는 매개변수표현이 존재하기 위한 충분조건}을 제시한다. 이 조건은 군에 속하는 불변체, 즉 라그랑주 종영역을 그 영역의 \emph{독립인} 함수들로 유리적으로 나타낼 수 있어서 ``최소기저''를 갖는다는 것이다. 달리 말하면, 그 불변체가 \(n\)개 부정원의 모든 유리함수체와 동형이라는 것이다. § 2에서 보겠지만, 예를 들어 3차와 4차 방정식에서 나타나는 모든 군이 이 조건을 만족한다. 이 3차·4차 방정식의 실제 작성과 더 정확한 논의는 F. Seidelmann의 학위논문\srcfn{*)}{``Die Gesamtheit der kubischen und biquadratischen Gleichungen mit Affekt bei beliebigem Rationalitätsbereich'', Erlangen, 1916.}에 있으며, 그 요약이 이 글 바로 뒤에 실려 있다.

라그랑주 종영역의 최소기저 문제는, 앞에서 말한 연관관계와는 독립적으로, E. Fischer가 내게 제기한 것이다. ``Körper und Systeme rationaler Funktionen''(Math. Ann. 76)에서도 이미 언급했듯이 이 질문이 이 연구의 계기가 되었다.\srcfn{**)}{``Rationale Funktionenkörper''에 관한 예비보고의 제2부, Jahresb. d. d. Math. Vereinig., Bd. 22 (1913) 참조.}

\begin{center}
§ 1.\\[0.75em]
\textbf{매개변수표현을 위한 충분조건.}
\end{center}

1. 지정된 군을 갖는 방정식의 구성을 매개변수표현으로 환원할 수 있음은 힐베르트의 불가약성 정리에서 따른다. 이 정리는 다음을 말한다. 방정식
\[
\begin{array}{@{}l@{\quad}l@{}}
(1)& \displaystyle f(x)=x^n+F_1(\lambda_1\cdots\lambda_r)x^{n-1}+\cdots+F_n(\lambda_1\cdots\lambda_r)=0
\end{array}
\]
을 생각하자. 여기서 \(F_i(\lambda_1\cdots\lambda_r)\)는 독립인 매개변수 \(\lambda_1\cdots\lambda_r\)의 유리함수이고, 그 계수는 수체\srcfn{***)}{수체 대신 일반적인 유리성 영역을 놓아도 된다.} \(\Omega\)에 속한다. 이 방정식이 \(\Omega\)의 계수를 갖는 \(\lambda_1,\ldots,\lambda_r\)의 모든 유리함수체 \(\Omega(\lambda_1\cdots\lambda_r)\)에 대하여 군 \(\Gamma\)를 갖는다면, 매개변수 \(\lambda\)를 임의로 많은 방식으로 유리수로 바꾸어 그 결과 생기는 수치방정식
\[
\begin{array}{@{}l@{\quad}l@{}}
(2)& \displaystyle g(x)=x^n+a_1x^{n-1}+\cdots+a_n=0
\end{array}
\]
이 \(\Omega\)에 대하여 바로 군 \(\Gamma\)를 갖게 할 수 있다. 그런 \(g(x)\)를 더 정확히 \(g_\Gamma\)라고 쓰자.

이제 다음과 같은 \emph{매개변수표현} (1)을 줄 수 있는지가 문제이다. \(\lambda_1,\ldots,\lambda_r\)가 \(\Omega\)의 모든 원소를 달릴 때, 그 결과 생기는 수치방정식 (2)가 모든 \(g_\Gamma\)를 달리게 하되, 군을 축소시키는 \(\Omega\)의 값들은 제외하고자 한다.\srcfn{*)}{\(g(x)=0\)에 중근이 나타나는 경우도 여기에 포함된다.} 달리 말하면 비선형 방정식계
\[
\begin{array}{@{}l@{\quad}l@{}}
(3)& \displaystyle F_1(\lambda_1\cdots\lambda_r)=a_1;\ \cdots\ F_n(\lambda_1\cdots\lambda_r)=a_n
\end{array}
\]
이 \(g_\Gamma\)에 대응하는 모든 값의 계 \(a_1,\ldots,a_n\)에 대하여 \(\Omega\)의 원소 \(\lambda_1,\ldots,\lambda_r\)로 된 해를 가져야 한다. (3)은 비선형이므로 모든 \(g_\Gamma\)를 하나의 매개변수표현으로 얻으려는 요구에는 처음부터 제한이 필요하다. 이러한 비선형 방정식계에는 대수관계 \(H(a_1\cdots a_n)=0\)을 만족하면서도 \emph{해가 전혀 없는} 특이한 값의 계 \(a_1,\ldots,a_n\)이 언제나 나타날 수 있기 때문이다.\srcfn{**)}{이는 ``비선형 방정식계의 양자택일''에 관한 논문에서 곧 일반적으로 전개할 예정이다.} 그때는 매개변수표현 (1)이 실패하므로 보충표현이 필요하다. 그러나 여기서는 뒤에서 보듯이 이 특이한 값들을 간단히 특징지을 수 있다.

2. 그런 매개변수표현을 얻기 위하여 우리는 더 강하게, \(\lambda_i\)들이 \emph{부정원으로 간주한 근들의 항등적인 함수이며 군에 속하는 자연스러운 무리량}일 것을 요구한다. 그 뜻은 다음과 같다. (1)의 매개변수 \(\lambda_1,\ldots,\lambda_r\)를 부정원 \(x_1,\ldots,x_n\)의 적절한 \emph{유리}함수 \(\varphi_1(x),\ldots,\varphi_r(x)\), 곧 \(\Omega\)의 계수를 갖는 함수로 바꾸었을 때, (1)이
\[
\begin{array}{@{}l@{\quad}l@{}}
(4)& \displaystyle h(x)=x^n-\sigma_1(x)x^{n-1}+\sigma_2(x)x^{n-2}\cdots\pm\sigma_n(x)
\end{array}
\]
으로 옮겨가야 한다. 여기서 \(\sigma_1(x)\ldots\sigma_n(x)\)는 \(x_1\ldots x_n\)의 기본대칭함수이다. 또 \(h(x)\)는 그렇게 \(\Omega(\lambda_1\ldots\lambda_r)\)에서 생기는 유리성 영역 \(\Omega(\varphi_1(x),\ldots,\varphi_r(x))\)에 대하여 군 \(\Gamma\)를 가져야 한다. 매개변수들이 독립이므로 \(\varphi_1(x)\ldots\varphi_r(x)\)도 대수적으로 독립인 함수여야 한다.

이 요구를 정확히 하려면 \(\Omega(\varphi_1(x)\ldots\varphi_r(x))\)와 군 \(\Gamma\)의 관계를 밝혀야 한다. 먼저 \(h(x)\)가 \(h_\Gamma\)가 되는 가장 일반적인 유리성 영역 \(K\)를 정하자. \(\mathfrak L_\Gamma\)를 \(x_1\ldots x_n\)의 유리수 계수 유리함수 가운데 \(\Gamma\)를 허용하는 모든 함수의 체라고 하자. 이는 \(\Gamma\)에 속하는 ``라그랑주 종영역'' 또는 \(\Gamma\)의 불변체라고도 부른다. 또 \(\Omega(\mathfrak L_\Gamma)\)를 \(\Omega\)에 \(\mathfrak L_\Gamma\)를 첨가하여 생기는 대응하는 체라고 하자. 따라서 \(\mathfrak L_\Gamma\)는 \(x_1,\ldots,x_n\)의 모든 유리수 계수 유리함수체 \(\mathfrak R(x_1\ldots x_n)\)의 부분체이다. 군의 정의에 따르면 가장 일반적인 영역 \(K\)는 \emph{\(\mathfrak L_\Gamma\)를 부분체로 포함하며 \(\mathfrak R(x_1\ldots x_n)\)과의 교집합이 정확히 \(\mathfrak L_\Gamma\)인 임의의 체}이다. 이에 따라 \(\Omega\)를 포함하는 모든 영역 \(K\)에 대해서는 \(\Omega(x_1\ldots x_n)\)과의 교집합이 \(\Omega(\mathfrak L_\Gamma)\)이다. 특히 우리의 요구에 따라 \(\Omega(\varphi_1(x)\ldots\varphi_r(x))\)는 그러한 영역 \(K\)이므로, 이것과 \(\Omega(x_1\ldots x_n)\)의 교집합은 \(\Omega(\mathfrak L_\Gamma)\)이다. 다른 한편 \(\Omega(\varphi_1(x)\ldots\varphi_r(x))\) 자체가 \(\Omega(x_1\ldots x_n)\)의 부분체이므로, \(\Omega(\varphi_1(x)\ldots\varphi_r(x))\)는 정확히 \(\Omega(\mathfrak L_\Gamma)\)와 같다. 따라서 우리의 요구는 \emph{라그랑주 종영역 \(\Omega(\mathfrak L_\Gamma)\)가 대수적으로 독립인 함수 \(\varphi_1(x)\ldots\varphi_r(x)\)를 \(\Omega\)에 첨가하여 생겨야 한다}는 뜻이다. 이때 \(\mathfrak L_\Gamma\)는 대수적으로 독립인 \(n\)개의 기본대칭함수를 포함하고, 다른 한편 대수적으로 독립인 함수는 \(n\)개를 넘을 수 없으므로 \(r=n\)이어야 한다. 이제 \(\varphi_1(x),\ldots,\varphi_n(x)\)가 \(\Omega(\mathfrak L_\Gamma)\)의 \emph{최소기저}를 이룬다고 말하자. 같은 뜻으로 \(\mathfrak L_\Gamma\)가 \(\Omega\)의 계수를 갖는 최소기저를 가진다고도 말한다.

3. 이제 중요한 점은 이 논의 전체를 거꾸로 돌릴 수 있으며, 그 결과 실제로 원하는 매개변수표현의 충분조건을 얻는다는 것이다.\srcfn{*)}{2의 요구는 1의 요구보다 강해졌으므로 이 조건들이 필요조건일 필요는 없다.} 따라서 \(\varphi_1(x)\ldots\varphi_n(x)\)가 \(\Omega(\mathfrak L_\Gamma)\)의 최소기저라고 가정하자. \(\Omega\)는 \(\varphi_1(x)\ldots\varphi_n(x)\)의 계수를 포함하는 가장 작은 체라고 하자. 그러면 \(\Omega\)는 반드시 대수적 수체이고, 특히 모든 유리수의 체 \(\mathfrak R\)일 수 있다. \(\sigma_1(x),\ldots,\sigma_n(x)\)가 \(\Omega(\mathfrak L_\Gamma)\)에 속하므로, \(F_i\)를 그 인수들의 유리함수라고 할 때 \(x_1,\ldots,x_n\)에 관한 항등식
\[
\begin{array}{@{}l@{\quad}l@{}}
(5)& \displaystyle \sigma_1(x)=F_1(\varphi_1(x)\ldots\varphi_n(x));\quad \ldots;\quad \sigma_n(x)=F_n(\varphi_1(x)\ldots\varphi_n(x));
\end{array}
\]
이 성립한다. 거꾸로 \(\varphi_1(x)\ldots\varphi_n(x)\), 따라서 \(\varphi(x)=\mu_1\varphi_1(x)+\cdots+\mu_n\varphi_n(x)\)는 \(\sigma_1(x)\ldots\sigma_n(x)\)에 대수적으로 종속이며, \(\Omega(\sigma_1(x)\ldots\sigma_n(x))\) 위에서 기약인 방정식
\[
\begin{array}{@{}l@{\quad}l@{}}
(6)& \displaystyle G(\varphi(x);\sigma_1(x)\ldots\sigma_n(x))=0,
\end{array}
\]
을 만족한다. 이 역시 \(x\)들에 관한 항등식이다.

(5)에 의해 항등식 (6)은
\[
\begin{array}{@{}l@{\quad}l@{}}
(7)& \displaystyle G(\varphi(x);F_1(\varphi_1\ldots\varphi_n);\ldots F_n(\varphi_1\ldots\varphi_n)) =H(\varphi_1(x)\ldots\varphi_n(x))=0;
\end{array}
\]
으로 바뀐다. \(\varphi_1(x)\ldots\varphi_n(x)\)가 대수적으로 독립이므로 (7)은 \(x\)에 관해서뿐 아니라 \(\varphi_i\)에 관해서도 항등적으로 성립한다. 즉 독립인 매개변수 \(\lambda\)에 관하여
\[
\begin{array}{@{}l@{\quad}l@{}}
(8)& \displaystyle H(\lambda_1\ldots\lambda_n)=G(\mu_1\lambda_1+\cdots+\mu_n\lambda_n;F_1(\lambda_1\ldots\lambda_n);\ldots F_n(\lambda_1\ldots\lambda_n))=0.
\end{array}
\]
이 항등적으로 성립한다.

항등식 (8)에서 방정식
\[
\begin{array}{@{}l@{\quad}l@{}}
(9)& \displaystyle f(x)=x^n-F_1(\lambda_1\ldots\lambda_n)x^{n-1}+F_2(\lambda_1\ldots\lambda_n)x^{n-2}\ldots\pm F_n(\lambda_1\ldots\lambda_n)=0
\end{array}
\]
이 \emph{\(\Omega(\lambda_1\ldots\lambda_n)\)에 대하여 군 \(\Gamma\)를 가짐}을 추론할 수 있다. 실제로 이 항등식은 유리적으로 알려진 \(\lambda=\mu_1\lambda_1+\cdots+\mu_n\lambda_n\)가 \(\Gamma\)에 속하는 근의 함수임을 보인다. 한편 \(\lambda_i=\varphi_i(x)\)를 대입해 보면 \(\Gamma\)의 진부분군에 속하는 함수가 추가로 유리적으로 알려질 수 없음을 알 수 있고, 같은 대입에서 \(\lambda\)가 자기의 모든 켤레값과 다름도 알 수 있다. 더 나아가 \(\Omega^*\)가 \(\Omega\)를 포함하는 임의의 수체이면 \(f(x)\)는 \emph{\(\Omega^*(\lambda_1\ldots\lambda_n)\)에 대하여 군 \(\Gamma\)를 갖는다.} 2에 따라 \(\lambda_i=\varphi_i(x)\)를 대입했을 때 \(\Omega^*\)를 첨가해도 군이 축소되지 않기 때문이다.

이제 1에서 수정한 의미로 \(f(x)\)가 실제로 모든 \(g_\Gamma\)를 달림을 보이면 된다. (5)에 의해 \(F_i(\lambda)\)는 그 인수 \(\lambda\)의 대수적으로 독립인 함수들이다. 따라서 방정식계
\[
\begin{array}{@{}l@{\quad}l@{}}
(10)& \displaystyle F_1(\lambda_1\ldots\lambda_n)=-a_1;\quad F_2(\lambda_1\ldots\lambda_n)=a_2;\quad \ldots;\quad F_n(\lambda_1\ldots\lambda_n)=\pm a_n
\end{array}
\]
은 아직 명시해야 할 특이한 경우를 제외하면 임의의 값의 계 \(a_1,\ldots,a_n\)에 대하여 해를 갖는다. \(\Omega^*\)에 대한 어떤 \(g_\Gamma\)에 대응하는 값의 계 \(a_1,\ldots,a_n\)에서는, 그때 생기는 \(\lambda_i\)가 군에 속하는 자연스러운 무리량이므로\srcfn{*)}{아래도 참조하라.} \(\Omega^*\)의 원소가 된다. 따라서 \(\lambda_1\ldots\lambda_n\)가 모든 값의 계를 달리면 \(g(x)\)는 수정된 의미에서 모든 방정식을 달리고, 그 가운데에는 \(\Omega^*\)의 값에 대응하는 모든 \(g_\Gamma\)가 포함된다. 거꾸로 힐베르트의 불가약성 정리에 따라 \(\Omega^*\)의 \(\lambda\) 값에는 ``일반적으로'' 역시 방정식 \(g_\Gamma\)가 대응한다. 그러므로 1에서 수정한 의미로 매개변수표현 (9)는 \(\Omega^*\)에 대한 모든 \(g_\Gamma\)를 실제로 준다.

이제 매개변수표현이 실패하는 특이한 값의 계를 정하면 된다. 최소기저의 함수들을 분자와 분모로 나누어
\[
\varphi_i(x)=\frac{\psi_i(x)}{\chi_i(x)}
\]
라고 하자. 이것을 항등식 (5)에 대입하고 약분하지 않은 채
\[
\begin{array}{@{}l@{\quad}l@{}}
(11)& \displaystyle \sigma_k(x)=\frac{G_k(x)}{H_k(x)}\quad\text{또는}\quad H_k(x)\cdot\sigma_k(x)=G_k(x)
\end{array}
\]
을 얻는다고 하자. 그러면 모든 \(H_k(x)\)의 곱은 모든 분모 \(\chi_i(x)\)의 곱으로 나누어진다. 근의 계 \(\alpha_1\ldots\alpha_n\)가
\[
H_1(\alpha)\cdot H_2(\alpha)\cdots H_n(\alpha)\neq0
\]
을 만족하면 어느 분모 \(\chi_i(\alpha)\)도 0이 아니다. 따라서 (11)을 \(H_k(\alpha)\)로 나누어
\[
\begin{array}{@{}l@{\quad}l@{}}
(12)& \displaystyle \sigma_k(\alpha)=\frac{G_k(\alpha)}{H_k(\alpha)}=F_k(\varphi_1(\alpha);\cdots\varphi_n(\alpha)),
\end{array}
\]
을 얻는다. 분모가 0이 아니므로 \(\varphi_i(\alpha)\)는 유한한 값을 취하며, \(g_\Gamma\)의 경우에는 \(\Omega^*\)의 원소가 된다. \(a_k=(-1)^k\sigma_k(\alpha)\)가 되도록 \(\alpha_1,\ldots,\alpha_n\)를 정하면 (12)는 (10)의 해를 준다.\srcfn{*)}{이로써 방정식을 푸는 문제는 독립인 함수들로 이루어진 방정식계
\[
\varphi_1(\alpha_1\cdots\alpha_n)=\lambda_1;\quad \cdots;\quad \varphi_n(\alpha_1\cdots\alpha_n)=\lambda_n
\]
을 푸는 문제로 환원된다.}

따라서 \emph{특이할 수 있는} 계수의 계 \(a_k=(-1)^k\sigma_k(\alpha)\)는
\[
H_1(\alpha)\cdot H_2(\alpha)\cdots H_n(\alpha)=0
\]
을 만족하는 것뿐이다. 그런 경우 (11)은 표현식이 아니라 \(0=0\)으로 바뀐다. 이렇게 정의된 특이한 \(a_1\ldots a_n\) 가운데 어떤 수체에 대하여 \(g_\Gamma\)에 대응하는 것이 있는지는 따로 조사해야 한다.\srcfn{**)}{예를 들어 F. Seidelmann이 앞의 논문에서 보였듯이, 3차·4차 방정식의 교대군에서만 그러한 특이한 값의 계에 대응하는 방정식이 나타나며, 그것도 3차 단위근의 체를 포함하는 수체 \(\Omega^*\)에서만 나타난다. 이때는 매번 간단한 보충표현을 줄 수 있다.} 이를 요약하면 다음 정리를 얻는다.

\emph{정리. 군 \(\Gamma\)에 속하는 라그랑주 종영역이 \(\Omega\)의 계수를 갖는 최소기저를 가지면}\srcfn{***)}{하나의 최소기저가 존재하면, 원래의 기저에서 가역 유리변환으로 유도되는 최소기저가 임의로 많이 존재함이 곧바로 따른다.} \emph{\(\Omega\)를 포함하는 모든 수체 \(\Omega^*\)에 대하여, \(\Omega^*\)에 대한 지정된 군 \(\Gamma\)를 갖는 방정식 전체를 매개변수표현 (2)로 유리적으로 구성할 수 있다. 다만 매개변수표현에 관하여 특이한 방정식들은 예외일 수 있으며, 이들은 따로 제시할 수 있다. 매개변수는 군을 축소시키지 않는 \(\Omega^*\)의 모든 값을 달려야 한다. 최소기저의 계수영역 \(\Omega\)가 모든 유리수의 체 \(\mathfrak R\)이면 이 구성은 임의의 모든 수체에 대하여 가능하다.}

힐베르트의 불가약성 정리에 따르면, 군을 축소시키는 \(\Omega^*\)의 값들을 제외해도 매개변수다양체의 차원은 낮아지지 않는다. 따라서 다음도 성립한다. \emph{최소기저가 존재하면 \(\Omega^*\)에 대한 군 \(\Gamma\)의 방정식들의 다양체는 아펙트가 없는 방정식들의 다양체와 같다. 아펙트 때문에 다양체의 차원이 낮아지지 않는다.}

\begin{center}
§ 2.\\[0.75em]
\textbf{특수한 방정식에 대한 응용.}
\end{center}

이제 3차와 4차 방정식에서 나타나는 모든 군에 최소기저가 존재함을 보이자. 이를 위하여 먼저 일반적으로 \(n\)차 방정식의 문제를 \(n-2\)개 부정원의 문제로 환원할 수 있음을 보인다. 첫 번째 환원은 두 번째로 높은 항을 없애는 선형 치른하우스 변환에 해당한다. 첫 기저함수로 \(\varphi_1(x)=\sigma_1(x)\)를 택한다. 또 \(\mathfrak L_\Gamma\)는 \(x_1,\ldots,x_n\)의 \emph{동차}형식 가운데 \(\Gamma\)를 허용하는 모든 형식으로부터 유리적으로 생성됨을 유의하자. \(p(x_1\ldots x_n)\)가 그러한 형식이라고 하자. 그러면
\[
q(x)=p\left(x_1-\frac{\sigma_1(x)}{n};\cdots;x_n-\frac{\sigma_1(x)}{n}\right)=p\left(x_i-\frac{\sigma_1(x)}{n}\right)
\]
도 \(\Gamma\)를 허용하므로 \(\mathfrak L_\Gamma\)에 속한다. 더 나아가
\[
p(x)\equiv p\left(x_i-\frac{\sigma_1(x)}{n}\right)\operatorname{mod}\sigma_1(x)
\]
또는
\[
p(x)=q(x)+\sigma_1(x)\cdot p_1(x)
\]
이며, 여기서 \(p_1(x)\)는 \(\mathfrak L_\Gamma\)에 속하고 그 차수는 \(p(x)\)보다 낮다. 이 과정을 계속하면 \(\mathfrak L_\Gamma\)의 모든 동차형식, 따라서 \(\mathfrak L_\Gamma\)의 모든 유리함수를 \(\varphi_1(x)=\sigma_1(x)\)와 동차형식
\[
q(x)=p\left(x_i-\frac{\sigma_1(x)}{n}\right)=p(y_i)
\]
으로 유리적으로 나타낼 수 있다. 그런데 \(y_i\)들 사이에는 선형관계 \(\sum y_i=0\)이 있으므로, \(q(x)\)들은 \emph{\((n-1)\)개의 인수에만 의존한다.}

두 번째 환원은 \(q(x)\)들이 그 인수들의 동차형식이라는 사실에 바탕을 둔다. 두 번째 기저함수로
\[
\varphi_2(x)=\frac{\tau_3(x)}{\tau_2(x)},\quad \text{여기서}\quad
\tau_3(x)=\sigma_3\left(x_i-\frac{\sigma_1(x)}{n}\right)=\sigma_3(y_i),\quad
\tau_2(x)=\sigma_2\left(x_i-\frac{\sigma_1(x)}{n}\right)=\sigma_2(y_i)
\]
를 택한다. 따라서 \(\varphi_2(x)\)는 \(\sum y_i=0\)을 만족하는 \(y_i\), 곧 \(n-1\)개 인수의 1차 동차분수함수이다. \(\varphi_2(x)\)를 사용하면 모든 \(q(x)\)를 역시 \(\mathfrak L_\Gamma\)에 속하는
\[
r(x)=\frac{q(x)}{\varphi_2(x)^\lambda}=\frac{q(x)\cdot\tau_2(x)^\lambda}{\tau_3(x)^\lambda}=\frac{p(y_i)\cdot\sigma_2(y_i)^\lambda}{\sigma_3(y_i)^\lambda}
\]
로 유리적으로 나타낼 수 있다. 여기서 \(\lambda\)는 \(q(x)\)의 차수이다. 그러므로 \(r(x)\)는 0차 동차함수이고, \(\sum y_i=0\)을 만족하는 비 \(y_1:y_2:\cdots:y_n\), 곧 \(n-2\)개의 인수에만 의존한다. 따라서 \(\mathfrak L_\Gamma\)는
\[
\varphi_1(x)=\sigma_1(x);\quad \varphi_2(x)=\frac{\tau_3(x)}{\tau_2(x)}
\]
와 \emph{이 \(n-2\)개 인수에 의존하는 \(\mathfrak L_\Gamma\)의 모든 함수 \(r(x)\)의 체}\srcfn{*)}{두 번째 항이 없는 \(n\)차 방정식
\[
f(x)=x^n+a_2x^{n-2}+a_3x^{n-3}+\cdots+a_n=0
\]
에서는 이 두 번째 환원에, \(a_2\neq0\), \(a_3\neq0\)일 때 언제나 가능한 치환 \(y=\frac{x\cdot a_2}{a_3}\)가 대응한다. 이 치환으로 \(f(x)\)는 한 인자를 제외하면
\[
g(y)=y^n+\frac{a_2^3}{a_3^2}y^{n-2}+\frac{a_2^3}{a_3^2}y^{n-3}+\frac{a_4\cdot a_2^4}{a_3^4}y^{n-4}+\cdots+a_n\cdot\frac{a_2^n}{a_3^n}=0
\]
으로, 따라서 \(n-2\)개의 매개변수만을 포함하는 방정식
\[
g(y)=y^n+c\cdot y^{n-2}+c\cdot y^{n-3}+c_4y^{n-4}+\cdots+c_n=0
\]
으로 옮겨간다.}으로 유리적으로 표현된다. 이로써 원하는 환원이 이루어졌다.

특히 3차와 4차 방정식에서는 각각 하나와 두 개의 인수를 갖는 함수체로 환원된다. 뤼로트와 카스텔누오보의 정리에 따라 이러한 함수체에는 언제나 최소기저가 존재한다.\srcfn{**)}{J. Lüroth, ``Beweis eines Satzes über rationale Kurven'', Math. Ann. 9 (1875); G. Castelnuovo, ``Sulla razionalità delle involuzioni piane'', Math. Ann. 44 (1893). 세 개보다 많은 부정원의 체에는 최소기저가 존재하지 않을 수 있다는 사실은 F. Enriques가 반례로 보였다. Rend. Acc. Linc., Vol. 21, 1912년 1월 21일.} 부정원이 하나인 경우 최소기저는 유리연산으로 유도할 수 있으므로, 특히 \(\mathfrak L_\Gamma\)에 대해서는 유리수 계수만을 갖는다. 부정원이 둘이면 카스텔누오보의 정리에 따라 대수적 수체 \(\Omega\)의 계수를 갖는 기저가 필요할 수도 있다. 그러나 실제 조사(F. Seidelmann, 앞의 논문 참조)에 따르면 이 경우에도 각 \(\mathfrak L_\Gamma\)에 유리수 계수의 기저가 존재한다. 클라인 4원군에 대해서는 그러한 유리수 계수 기저
\[
\varphi_1(x)=\sigma_1(x);\quad \varphi_2(x)=\frac{v w}{u};\quad \varphi_3(x)=\frac{w\cdot u}{v};\quad \varphi_4(x)=\frac{u\cdot v}{w}
\]
를 택할 수 있다는 사실에서 거의 계산 없이 이를 알 수 있다. 여기서
\[
u=x_1+x_2-x_3-x_4;\quad v=x_1-x_2+x_3-x_4;\quad w=x_1-x_2-x_3+x_4
\]
이다. 교대군은 세 기저함수 \(\varphi_2,\varphi_3,\varphi_4\)의 순환군으로, 각 8원군은 이 함수들 가운데 두 개씩의 대칭군으로 옮겨가므로, 문제가 3차와 2차 방정식의 군으로 환원된다. 한편 순환군에 대해서는 아벨이 이미 유리수 계수의 기저를 알고 있었지만 그러한 말로 정식화하지는 않았다.\srcfn{*)}{Weber, \emph{Algebra}, 소판, § 93, 공식 (12) 참조.} 따라서 다음이 증명되었다. \emph{지정된 군을 갖는 3차와 4차 방정식 전체는, 매개변수표현에 관하여 특이한 방정식들은 예외일 수 있지만, 임의의 모든 유리성 영역에서 매개변수표현으로 유리적으로 구성할 수 있다. 그 다양체는 아펙트가 없는 방정식들의 다양체와 같다.}

실제 조사(F. Seidelmann, 앞의 논문 참조)에 따르면, 3차 단위근의 체를 포함하는 유리성 영역에 대한 교대군의 경우에만 보충표현이 필요하다. 그 밖의 모든 경우에는 매개변수표현 하나가 예외 없이 방정식 전체를 준다.

끝으로 모든 \emph{아벨군}에 대해서도 최소기저가 알려져 있음을 덧붙인다.\srcfn{**)}{E. Fischer, ``Die Isomorphie der Invariantenkörper der endlichen Abelschen Gruppen linearer Transformationen'', Gött. Nachr. 1915; ``Zur Theorie der endlichen Abelschen Gruppen'', Math. Ann. 77 (1915).} 그러나 이때의 기저는 군의 지표들에서 유도되는 원분체의 계수를 갖는다. 따라서 지정된 군을 갖는 방정식의 구성에 관하여 이 경우에는 새로운 결과를 얻을 수 없다.

괴팅겐, 1916년 7월.

\clearpage
\setcounter{footnote}{0}

\clearpage
\setcounter{footnote}{0}
\phantomsection
\addcontentsline{toc}{section}{12. 임의의 미분식의 불변식}
%
%
%
%
%
%

\section*{12. 임의의 미분식의 불변식}
\begin{center}
{\Large\bfseries 12. 임의의 미분식의 불변식.}\par
\vspace{0.8em}
괴팅겐의 에미 뇌터.\par
\vspace{1.5em}
Nachr. v. d. Ges. d. Wiss. zu Göttingen 1918, 37--44쪽
\end{center}

\vspace{3em}
\begin{center}
1918년 1월 25일 회의에서 F. Klein이 제출.
\end{center}

미분식이란 \(n\)개의 변수 \(x_1\ldots x_n\)과 그 미분
\(dx_1\ldots dx_n\)의 함수
\[
f(x;dx)=f(x_1,\ldots,x_n;\,dx_1,\ldots,dx_n)
\]
를 뜻한다. 이 함수는 두 인수계 각각에 관하여 해석적이라고 가정하지만 반드시 실함수일 필요는 없다. 이제 변수의 모든 해석적 변환으로 이루어진 군을 바탕에 두면, 이에 미분의 모든 선형변환으로 이루어진 군이 대응한다.
\srcnumdisplay{(1)}{%
x_i=x_i(y_1\ldots y_n);\qquad
dx_i=\sum\frac{\partial x_i}{\partial y_k}\,dy_k;\qquad
\delta x_i=\sum\frac{\partial x_i}{\partial y_k}\,\delta y_k.
}
이 변환으로 \(f(x,dx)\)가 \(g(y,dy)\)로 옮겨 간다고 하자. \(f(x,dx)\)의 불변식이란 이 군과, 이 군이 고차 미분 \(d^2x,d\delta x\ldots\) 및 \(f(x,dx)\)의 도함수들에 유도하는 군에 대한 불변식을 뜻한다. 다시 말해 임의의 \(f\)에 대하여 (1)에 따라 변수들과 나타나는 모든 미분에 관하여 항등적으로
\[
\begin{aligned}
&J\!\left(f,\frac{\partial f}{\partial dx}\cdots
\frac{\partial^{\varrho+\sigma}f}{\partial x^\varrho\partial dx^\sigma}
\cdots dx,\delta x,d^2x,\ldots\right)\\
&\qquad =
J\!\left(g,\frac{\partial g}{\partial dy}\cdots
\frac{\partial^{\varrho+\sigma}g}{\partial y^\varrho\partial dy^\sigma}
\cdots dy,\delta y,d^2y\ldots\right)\srcfnmark{1)}.
\end{aligned}
\]
을 만족하는 해석함수 \(J\)이다.
\srcfntext{1)}{\(\displaystyle
\frac{\partial^{\rho+\sigma}f}{\partial x^\rho\partial dx^\sigma}\)
는 이하에서 줄여 쓴 표기로서
\[
\frac{\partial^{\varrho+\sigma}f}
{\partial x_{i_1}\ldots\partial x_{i_\varrho}
 \partial dx_{k_1}\ldots\partial dx_{k_\sigma}}
\]
를 뜻하며, 지표들은 \(1\)부터 \(n\)까지의 임의의 수이다.}
여러 미분식으로 이루어진 연립계의 불변식도 완전히 같은 방식으로 정의한다. 그러한 불변식이 특히 일차 미분 \(dx,\delta x\)와 이 미분들에 관한 도함수만을 포함하고 변수에 관한 도함수는 포함하지 않는다면, 함수에 변수가 나타난다는 사실과 미분의 선형변환은 고려할 필요가 없다. 이 특수한 불변식들은 부정계수를 가진 \(dx,\delta x\)의 모든 선형변환군에 대한 불변식이 되며, 이하에서는 연립계의 사영불변식이라 부르겠다.

Christoffel(Crelle 70)과 Ricci(Math. Annalen 54)는 이차 동차 미분형식에 대하여, 모든 불변식이 그 연립계의 사영불변식이 되도록 하는 불변 구성들의 한 계를 세웠다. 이로써 불변식 전체와 동치성의 문제는 선형불변식론의 문제로 환원된다. 나는 이를 \glqq{}환원정리\grqq{}라 부르고자 한다. 완전계는 가산무한개의 형식으로 이루어지지만, 각 유한 차수 \(\varrho\)의 불변식, 곧 변수에 관한 \(\varrho\)차보다 높은 도함수를 포함하지 않는 불변식에는 유한개의 형식만 필요하다.

이하에서는 임의의 미분식에 대하여 \glqq{}환원정리\grqq{}가 성립함을 보인다. 여기서도 완전계는 가산무한개의 불변 구성으로 이루어지되, 각 유한 차수에는 유한개만 필요하다. Christoffel과 Ricci는 불변식을 만들고 환원정리를 증명할 때 조망하기 어려운 소거에 의존한다. 이에 반해 나는 불변적인 미분과정과 변분과정으로 불변식을 직접 만든다. 변분과정도 다른 매개변수에 관한 미분과정으로 간단히 볼 수 있다. Lagrange(\emph{Mécanique analytique}, 제2부 제IV절)에 이어 Riemann(Werke XXII)과 Lipschitz(Crelle 70, 72)는 이 불변 미분과정의 알고리즘을 이차 미분형식에 대하여 전개하였으나 환원정리에까지 이르지는 않았다. 환원정리의 증명에 쓰이는 도구는 Riemann이 이차형식에 도입한 \glqq{}정규좌표\grqq{}이다(Werke XIII). 이 좌표는 대응하는 변분문제에서 한 점으로부터 나가는 극값곡선들을 직선으로 바꾸지만, 동차 미분식
\[
f(x,\varkappa\cdot dx)=\varPhi(\varkappa)\cdot f(x,dx)
\]
에 대해서만 존재한다.\srcfn{1)}{\(\varPhi(\varkappa)=\varkappa^c\)가 됨을 쉽게 알 수 있으며, 여기서 \(c\)는 임의의 실수 또는 복소수이다.} 그러나 비동차 경우는 이 경우로 환원할 수 있다.

\begin{center}
\textbf{I. 불변식의 생성}
\end{center}

일반적으로 끝나지 않는 \(f(x,dx)\)의 사영불변식 열은 극화로 얻어진다. 미분의 변환이 선형이므로 \(f(x,dx)=g(y,dy)\)에서
\[
f(x,dx+\lambda\delta x)=g(y,dy+\lambda\delta y)
\]
도 따라온다. 이를 \(\lambda\)에 관하여 전개하면 불변성을 특징짓는 다음 관계들이 생긴다.
\srcnumdisplay{(2)}{%
\begin{gathered}
f(dx)=g(dy),\\
f_\delta=\sum\frac{\partial f(dx)}{\partial dx}\,\delta x
       =\sum\frac{\partial g(dy)}{\partial dy}\,\delta y,\\
\vdots\\
f_{\delta^\sigma}=
\sum\frac{\partial^\sigma f(dx)}{\partial dx^\sigma}\,\delta x^\sigma
=
\sum\frac{\partial^\sigma g(dy)}{\partial dy^\sigma}\,\delta y^\sigma,\\
\vdots
\end{gathered}
}

각 극형에 변분과정 \(\bdelta\)를 적용하면 끝나지 않는 일반 불변식 열을 얻는다.
\srcnumdisplay{(3)}{%
\begin{gathered}
\bdelta^\varrho f(x,dx)=\bdelta^\varrho g(y,dy),\\
\vdots\\
\bdelta^\varrho f_{\delta^\sigma}
=\bdelta^\varrho\sum\frac{\partial^\sigma f(dx)}{\partial dx^\sigma}\,
\delta x^\sigma
=\bdelta^\varrho\sum\frac{\partial^\sigma g(dy)}{\partial dy^\sigma}\,
\delta y^\sigma,\\
\vdots\\[0.2em]
(\varrho=0,1,2\ldots\ldots)
\end{gathered}
}
관계 (2)와 (3)에서는 (1)을 이용하여 \(f(x,dx)\)의 도함수들에 대한 변환법칙을 읽어 낼 수 있지만, 이하에서는 이를 쓰지 않는다. 이제 불변식 (3)을 선형결합하면 \glqq{}\(\varrho\)차 변분의 정규형\grqq{}을 얻을 수 있다. \(\varrho=1\)인 경우는 Lagrange에게, \(\varrho=2\)인 경우는 Riemann에게서 찾을 수 있다. 이 정규형은 \((\varrho+1)\)차 혼합미분
\[
d^\varrho\delta,\ d^{\varrho-1}\delta^2,\ldots
\]
이 소거된다는 성질로 특징지어진다. 구체적으로 다음을 얻는다.
\srcnumdisplay{(4)}{%
\begin{gathered}
\varOmega_1=\delta f(dx)-d f_\delta(dx),\\
\varOmega_2=\delta^2 f-\delta d f_\delta+\frac12 d^2 f_{\delta^2},\\
\vdots\\
\varOmega_\varrho=\delta^\varrho f-\delta^{\varrho-1}d f_\delta
+\frac12\delta^{\varrho-2}d^2 f_{\delta^2}
+\cdots+\frac{(-1)^\varrho}{\varrho!}\,d^\varrho f_{\delta^\varrho},\\
\vdots
\end{gathered}
}

이제 Riemann의 한 착상을 일반화하면 불변식 (4)에서 일차 미분만을 포함하는 불변식으로 나아갈 수 있다. 이러한 불변식을 \glqq{}기본함수\grqq{}라 부르겠다. 실제로 \(\varOmega_1\)에서 미분들을 미분몫으로 바꾸면, \(\varOmega_1=0\)은 \(\bdelta x\)에 관하여 항등적으로
\[
f\!\left(\frac{dx}{dt}\right)
\]
에 딸린 변분문제의 라그랑주 방정식과 정확히 같아진다. 이제 방향 \((d+\lambda\delta)\)가 이 라그랑주 방정식들을 만족한다는 불변조건을 세운다.
\srcnumdisplay{(5)}{%
\begin{gathered}
\varOmega_1(d+\lambda\delta)=
\bdelta f(dx+\lambda\delta x)
-(d+\lambda\delta)f_{\bdelta}(dx+\lambda\delta x)=0,\\
(\bdelta x\text{에 관하여 항등적으로}).
\end{gathered}
}
동차 \(f(dx)\)에서는 \(\bdelta=d+\lambda\delta\)를 대입하여 다음도 따라온다.
\srcnumdisplay{(6)}{%
(d+\lambda\delta)f(dx+\lambda\delta x)=0.
}
다만 일차 동차인 경우는 예외이다. 이 경우 방정식계 (5)는 독립인 조건을 \((n-1)\)개만 주므로, (6)을 새 조건으로 덧붙인다. (5), 또는 (5)와 (6)에서 \(\lambda\)의 0차, 1차, 2차 계수를 각각 0으로 두면 세 불변 방정식계
\[
\varphi=0,\qquad\psi=0,\qquad\chi=0
\]
을 얻는다. 이들은 \(\bdelta x\)에 관하여 항등적이며, 이를 이용하면 \(d^2x,d\delta x,\delta^2x\)를 일차 미분으로 나타낼 수 있다. 더 나아가 다음 불변 방정식계는 모든 고차 미분을 일차 미분으로 나타내는 데 정확히 충분하다.
\srcnumdisplay{(7)}{%
\begin{gathered}
d^\sigma\varphi=0,\qquad d^\sigma\psi=0,\qquad
d^\sigma\chi=0,\qquad d^{\sigma-1}\delta\chi=0,\\
d^{\sigma-2}\delta^2\chi=0\ldots\delta^\sigma\chi=0
\qquad(\sigma=0,1,2\ldots).
\end{gathered}
}
\srcfn{1)}{선형형식 \(\sum A_i\,dx_i\)만은 따로 다루어야 한다. 이 경우 방정식계 (5)에 이차 미분이 들어 있지 않기 때문이다.}
따라서 함수 (4)를 불변 방정식계 (7)과 결합하면 일차 미분만을 포함하는 기본함수 \([\varOmega_\varrho]\)를 얻으며, 이때 \([\varOmega_1]\)은 항등적으로 사라진다.

마찬가지로 기본함수 \(h\)의 미분 \(\delta h(dx,\delta x)\)에서도 (7)을 이용하여 다시 하나의 기본함수를 얻을 수 있다. 이를 \(h\)의 \glqq{}공변미분\grqq{} \([h^{(1)}]\)이라 부르겠다. 일반적으로 \([h^{(\sigma)}]\)는 \([h^{(\sigma-1)}]\)의 공변미분을 뜻한다. 이 두 과정은 다음의 이중 무한 기본함수 열을 준다.
\srcnumdisplay{(8)}{%
\begin{array}{cccccc}
[\varOmega_2], & [\varOmega_2^{(1)}], & [\varOmega_2^{(2)}] & \ldots
& [\varOmega_2^{(\nu)}] & \ldots\\
\vdots &&&&&\\
{}[\varOmega_\varrho], & [\varOmega_\varrho^{(1)}], & \ldots & \ldots
& [\varOmega_\varrho^{(\nu)}] & \ldots\\
\vdots & \vdots &&& \vdots&
\end{array}
}

고차 미분을 일차 미분으로 나타내는 방정식들의 좌변은 일차 미분과 동변임도 알 수 있다. 즉 같은 선형변환을 따른다. 따라서 고차 미분 대신 이 좌변들 \(p,q,r\ldots\)을 불변식의 인수로 도입하면, 그 불변식은 도함수들 이외에는 서로 동변인 양들에만 의존한다. 위에서 제시한 함수 (8)의 구성은 바로 이 새 양들을 도입하는 것과 같다. 각 불변식은 이로써 불변식들의 합으로 바뀌며, 그 가운데 \(dx,\delta x\)만을 포함하고 \(p,q,r\ldots\)은 포함하지 않는 항을 골라낸다.

이하에서는 동차성 가정
\[
f(\varkappa\cdot dx)=\varPhi(\varkappa)\cdot f(dx)
\]
아래에서 기본함수 (8)이 찾고 있는 완전계 전체를 이룸을 보인다.

\begin{center}
\textbf{II. 정규좌표, 동치성과 환원정리}
\end{center}

\(f\!\left(\frac{dx}{dt}\right)\)에 딸린 변분문제를 적분하면 정규좌표에 이른다.
\srcnumdisplay{(9)}{%
\begin{gathered}
\varOmega_1=\delta f\!\left(\frac{dx}{dt}\right)
-\frac{d}{dt}f_\delta\!\left(\frac{dx}{dt}\right)=0
\quad(\delta x\text{에 관하여 항등적으로}),\\
\frac{d}{dt}f\!\left(\frac{dx}{dt}\right)=0
\quad(\text{귀결 또는 추가조건}).
\end{gathered}
}
\(f(dx)\)가 동차라는 가정 때문에 이 방정식계는 매개변수 치환 \(t=\varkappa\cdot\tau\) 아래에서 자기 자신으로 옮겨진다. 따라서 (9)를 이용하여
\(\frac{d^\varrho x_i}{dt^\varrho}\)를 일차 도함수의 함수
\(\varphi_\varrho^{(i)}\!\left(\frac{dx}{dt}\right)\)로 나타내면,
\(\varphi_\varrho^{(i)}\)는 \(\varrho\)차 동차가 된다.
\srcnumdisplay{(10)}{%
\varphi_\varrho^{(i)}\!\left(\varkappa\cdot\frac{dx}{dt}\right)
=\varkappa^\varrho\varphi_\varrho^{(i)}\!\left(\frac{dx}{dt}\right).
}
이제 (9)를 적분할 때 \(t=t_0\)에서 초깃값
\[
x=(x)_0,\qquad \frac{dx}{dt}=\left(\frac{dx}{dt}\right)_0
\]
을 주고 다음과 같이 둔다.
\srcnumdisplay{(11)}{%
u_i=(t-t_0)\left(\frac{dx_i}{dt}\right)_0.
}
여기서 \(f\!\left(\frac{dx}{dt}\right)=\mathrm{const.}\)이므로 매개변수 \((t-t_0)\)는 \glqq{}극값곡선의 길이\grqq{}에 비례한다. 그러면 \((t-t_0)\)의 거듭제곱으로 전개하여 (10)에 따라 다음을 얻는다.
\srcnumdisplay{(12)}{%
x_i-(x_i)_0=u_i+\frac12\varphi_2^{(i)}(u)+\cdots+
\frac{1}{\varrho!}\varphi_\varrho^{(i)}(u)+\cdots.
}
이 전개는 변수변환 \(x_i=x_i(u)\)로 볼 수 있으며, 이때 \(u\)가 바로 리만 정규좌표이다. 따라서 이 좌표는 (11)과 (12)에 따라 \((x_0)\)를 지나는 극값곡선들을 직선으로 바꾼다. 또 임의의 변수변환 \(x=x(y)\)에는 (11)에 따라 대응하는 정규좌표들의 선형변환 \(u=u(v)\)가 대응한다. 그 계수들은 임의로 택한 값의 계 \((x)_0\)에만 의존하므로, 미분 \(du,\delta u\)도 \(u\) 자체와 같은 변환을 따른다.

(12)에 따라 \(f(x,dx)\)가 \(F(u,du)\)로 옮겨 간다고 하자. 이를 \(u\)의 거듭제곱으로 전개하면
\srcnumdisplay{(13)}{%
F(u,du)=F_0(u,du)+F_1(u,du)+\cdots F_\varrho(u,du)+\cdots,
}
여기서 \(F_\varrho(u,du)\)는 \(u\)에 관한 \(\varrho\)차 동차형식이다. 변환 \(u=u(v)\)가 선형이므로 \(F(u,du)=G(v,dv)\)이면 각 동차 성분도 서로 대응한다.
\srcnumdisplay{(14)}{%
F_\varrho(u,du)=G_\varrho(v,dv).
}
따라서 \((x)_0\)를 임의의 고정된 상수 값의 계로 택하면, (13)과 (14)는 \(f(dx)\)와 \(g(dy)\)의 동치성이 대응하는 함수계 \(F_\varrho(u,du)\)들의 선형변환에 관한 동치성과 같은 뜻임을 보여 준다. 더 나아가 \(F_\varrho(u,du)\), 또는 대응하는 \(F_\varrho(\delta u,du)\)는 점 \(x=(x)_0\)에서 취한 \(f(dx)\)의 불변식이며, 이 점에 대한 기본함수의 완전계를 이룬다. 실제로
\[
F_\varrho(\delta u,du)=\frac1{\varrho!}\sum
\frac{\partial F_\varrho}{\partial \delta u^\varrho}\,\delta u^\varrho;
\qquad
\frac{\partial F_\varrho}{\partial \delta u^\varrho}
=\left(\frac{\partial F(du)}{\partial u^\varrho}\right)_0
\]
이고, 이 마지막 도함수는 (12)를 이용하여 점 \(x=(x)_0\)에서 취한 \(f(dx)\)의 도함수들로 나타낼 수 있다. 대응하게 만든 \(G(dv)\)의 도함수도 \(g(dy)\)의 도함수들로 나타낼 수 있다. 또한
\[
\left(\frac{\partial^{\varrho+\sigma}F(du)}
{\partial u^\varrho\partial du^\sigma}\right)_0
=
\frac{\partial^{\varrho+\sigma}F_\varrho(\delta u,du)}
{\partial\delta u^\varrho\partial du^\sigma}
\]
이므로, 점 \(x=(x)_0\)에서 취한 \(f(dx)=F(du)\)의 모든 불변식은 고차 미분들을 \(dx,\delta x\)와 동변인 \(p,q,r\ldots\)로 바꾸기만 하면 \(F_\varrho(\delta u,du)\)들의 사영불변식이 된다.

이제 \((x)_0\)를 매개변수로 볼 수 있다. 점 \(x=(x)_0\)에서의 각 불변식에는 \((x)_0\)를 다시 \(x\)로 바꾸어 얻는 \(f(dx)\)의 불변식이 대응한다. 따라서 \(F_\varrho(\delta u,du)\)들이 점 \(x=(x)_0\)에서 취한 불변식 \(\varPsi_\varrho(\delta x,dx)\)가 된다면 --- 이 점에서 \(du,\delta u\)는 \(dx,\delta x\)로 옮겨 간다 --- \(\varPsi_\varrho\) 자체가 한 완전계의 기본함수들을 이룬다. 이제 이 기본함수계를 \(f(dx)\)의 명시적인 불변식들, 곧 함수 (8)로 표현해야 한다. 미분과정에 의한 생성법은 이것이 가능함을 보여 준다. 최고차항에 대해서는 이 과정들이 단순한 극화과정이 되며, Clebsch--Gordan의 급수전개와 비슷한 변형에 생략한 항들에 관한 귀납을 결합하면 명제가 증명된다. 연립계의 불변식을 다루는 경우에는, \(f(dx)\)에 딸린 정규좌표를 다른 미분식들에 도입해 보면 알 수 있듯이 기본함수 (8)에 그 다른 미분식들의 공변미분도 덧붙여야 한다.

\(f(dx)\)와 \(g(dy)\)의 동치성은 (14)의 동치성으로 환원되었으므로, 이제 \(\varPsi_\varrho\)들의 동치성, 또는 함수 (8)의 미분에 대한 선형변환 아래의 동치성과도 같은 뜻이 된다. (14)로 환원할 수 있다는 사실에서 알 수 있듯이 이때 별도의 적분가능조건은 필요하지 않다. 이로써 환원정리의 모든 부분이 증명되었다. 특히 함수열
\[
[\varOmega_2],[\varOmega_3],\ldots,[\varOmega_\varrho],\ldots
\]
이 항등적으로 사라지는 것은 \(f(dx)\)를 상수계수를 가진 식으로 변환할 수 있기 위한 필요충분조건이다.

마지막으로 모든 해석적 변환의 군 대신 한 부분군을 바탕에 두면, 대응하는 \(u\)의 선형변환군도 사영군의 한 부분군으로 바뀐다. \(f(dx)\)의 불변식들은 다시 함수 (8)의 선형변환에 대한 불변식이 되지만, 이제는 이 부분군에 대한 불변식이다. 따라서 동차화하면 비동차 함수 \(f(dx)\)의 경우를 아핀군으로 환원할 수 있고, 그 완전계는 변수를 하나 더 두고 만든 함수 (8)에서 이끌어 낼 수 있다.

이차형식, 더 일반적으로 \(p\)차 형식의 특수 경우는 다음과 같이 환원정리에서 얻어진다. 이 경우 함수계는 \([\varOmega_2]\), 더 일반적으로는 \([\varOmega_p]\)와 그 공변미분들에서 끝나며, 그 뒤의 모든 기본함수는 이들의 사영불변식이 된다. 이로써 이차형식에서 리만 곡률형식 \([\varOmega_2]\)가 갖는 특별한 위치가 설명된다. 이차형식, 또는 \(p\)차 형식의 동치성 문제는 \([\varOmega_2]\), 또는 \([\varOmega_2]\ldots[\varOmega_p]\)와 그 공변미분들의 선형변환 아래의 동치성으로 환원된다. 또 \([\varOmega_2]\), 또는 \([\varOmega_2]\ldots[\varOmega_p]\)가 항등적으로 사라지는 것은 상수계수를 가진 형식으로 변환하기 위한 필요충분조건이다. 이차형식에 대해서는 이 명제가 가장 잘 알려진 결과들 가운데 하나이다.

더 자세한 서술은 곧 Math. Annalen에 출판할 예정이다.

\clearpage
\setcounter{footnote}{0}
\phantomsection
\addcontentsline{toc}{section}{13. 불변 변분문제}
%
%
%
%
%
%

\section*{13. 불변 변분문제}

\begin{center}
(F. 클라인의 박사학위 50주년에)\\[0.6em]
에미 뇌터 (괴팅겐)\\[1em]
\emph{Nachr. v. d. Ges. d. Wiss. zu Göttingen 1918, 235--257쪽}\\[2em]
F. 클라인이 1918년 7월 26일 회의에 제출함\srcfn{1)}{원고의 최종본은 9월 말에야 제출되었다.}.
\end{center}

여기서는 리의 의미에서 연속군을 허용하는 변분문제를 다룬다. 이에 대응하는 미분방정식들에 대한 귀결은 § 1에서 정식화하고 뒤의 절들에서 증명하는 정리들에 가장 일반적으로 표현되어 있다. 변분문제에서 생기는 이 미분방정식들에 대해서는 군을 허용하는 임의의 미분방정식, 곧 리가 연구한 방정식들보다 훨씬 더 정밀한 명제를 세울 수 있다. 따라서 이하의 논의는 형식적 변분법의 방법과 리 군론의 방법을 결합한다. 특수한 군과 변분문제에 대해서 이 방법들의 결합은 새롭지 않다. 특수한 유한매개변수군에 대해서는 하멜과 헤르글로츠를, 특수한 무한군에 대해서는 로런츠와 그의 제자들(예컨대 포커), 바일과 클라인을 들 수 있다\srcfn{2)}{Hamel, \emph{Math. Ann.} 59 및 \emph{Zeitschrift f. Math. u. Phys.} 50. Herglotz, \emph{Ann. d. Phys.} (4) 36, 특히 § 9, 511쪽. Fokker, Verslag d. Amsterdamer Akad., 1917년 1월 27일. 다른 문헌은 클라인의 둘째 글, Göttinger Nachrichten, 1918년 7월 19일을 보라. 막 출판된 Kneser의 논문 \emph{Math. Zeitschrift} 2에서는 비슷한 방법으로 불변식을 세운다.}. 특히 클라인의 둘째 글과 이 논의는 서로 영향을 주고받았으며, 이에 관해서는 클라인 논문의 맺음말을 참조할 수 있다.

\section*{§ 1. 예비사항과 정리의 정식화}

이하에 나타나는 모든 함수는 해석적이거나 적어도 연속이고 필요한 유한차수만큼 연속미분가능하며, 생각하는 영역에서 단일값이라고 가정한다.

잘 알려져 있듯이 \glqq{}변환군\grqq{}이란 각 변환의 역변환도 그 계에 들어 있고, 그 계의 임의의 두 변환을 합성한 것도 다시 그 계에 들어 있는 변환계를 뜻한다. 변환들이 본질적인 매개변수 \(\varepsilon\) \(\varrho\)개에 해석적으로 의존하는 가장 일반적인 변환 안에 포함될 때, 다시 말해 이 \(\varrho\)개 매개변수를 더 적은 매개변수들의 \(\varrho\)개 함수로 나타낼 수 없을 때 그 군을 유한매개변수 연속군 \(\mathfrak G_\varrho\)라 한다. 이에 대응하여 무한 연속군 \(\mathfrak G_{\infty\varrho}\)란 가장 일반적인 변환이 본질적인 임의함수 \(p(x)\) \(\varrho\)개와 그 도함수들에 해석적으로, 또는 적어도 연속이며 필요한 유한차수만큼 연속미분가능하게 의존하는 군이다. 그 중간에는 임의함수에는 의존하지 않지만 무한히 많은 매개변수에 의존하는 군이 있다. 끝으로 임의함수와 매개변수 모두에 의존하는 군을 혼합군이라 한다\srcfn{1)}{리는 \emph{Grundlagen für die Theorie der unendlichen kontinuierlichen Transformationsgruppen}, Ber. d. K. Sächs. Ges. der Wissensch. 1891 [이하 \glqq{}Grundlagen\grqq{}]에서 무한 연속군을, 변환들이 편미분방정식계의 가장 일반적인 해로 주어지며 그 해들이 유한개의 매개변수에만 의존하지 않는 변환군으로 정의한다. 이 정의로 위에서 말한 유한군과 다른 유형들 가운데 하나를 얻는다. 반대로 무한히 많은 매개변수라는 극한적인 경우가 반드시 미분방정식계를 만족해야 하는 것은 아니다.}.

\(x_1,\ldots,x_n\)을 독립변수라 하고 \(u_1(x),\ldots,u_\mu(x)\)를 그에 종속하는 함수들이라 하자. \(x,u\)에 한 군의 변환들을 적용하면 변환의 가역성 때문에 변환된 양들 가운데 다시 정확히 \(n\)개의 독립변수 \(y_1,\ldots,y_n\)이 들어 있어야 한다. 그에 종속하는 나머지 함수들을 \(v_1(y),\ldots,v_\mu(y)\)라 하자. 변환에는 \(u\)의 \(x\)에 대한 도함수
\[
\frac{\partial u}{\partial x},\quad
\frac{\partial^2u}{\partial x^2},\ldots
\]
도 나타날 수 있다\srcfn{1)}{가능한 곳에서는 합을 쓸 때도 지표를 생략한다. 예를 들어 \(\partial^2u/\partial x^2\)는 \(\partial^2u_\alpha/\partial x_\beta\partial x_\gamma\) 등을 뜻한다.}. 관계
\[
P\!\left(x,u,\frac{\partial u}{\partial x},
\frac{\partial^2u}{\partial x^2},\ldots\right)
=
P\!\left(y,v,\frac{\partial v}{\partial y},
\frac{\partial^2v}{\partial y^2},\ldots\right)
\]
가 성립하면 그 함수를 군의 불변식이라 한다.

특히 적분 \(I\)에 대하여
\begin{equation}
\tag{1}
\begin{aligned}
I&=\int\!\cdots\!\int
f\!\left(x,u,\frac{\partial u}{\partial x},
\frac{\partial^2u}{\partial x^2},\ldots\right)\dd x\\
&=\int\!\cdots\!\int
f\!\left(y,v,\frac{\partial v}{\partial y},
\frac{\partial^2v}{\partial y^2},\ldots\right)\dd y
\end{aligned}
\end{equation}
가 성립하면 \(I\)는 군의 불변식이다\srcfn{2)}{\(d x,d y\)는 각각 \(d x_1\ldots d x_n,d y_1\ldots d y_n\)의 줄임표기이다.}. 적분은 임의의 실수 \(x\)-영역과 그에 대응하는 \(y\)-영역에서 한다\srcfn{3)}{변환에 나타나는 모든 인수 \(x,u,\varepsilon,p(x)\)는 실수라고 가정하되 계수는 복소수여도 된다. 그러나 최종결과들은 \(x,u\), 매개변수, 임의함수에 관한 항등식이므로 모든 함수가 해석적이면 복소수 값에 대해서도 성립한다. 결과의 큰 부분은 적분 없이도 근거 지을 수 있으므로 그 부분의 증명에는 실수로의 제한도 필요 없다. 반면 § 2 끝과 § 5 시작의 고찰은 적분 없이 수행하기 어려워 보인다.}.

한편 불변일 필요가 없는 임의의 적분 \(I\)에 대하여 제1변분 \(\delta I\)를 만들고, 변분법의 규칙에 따라 부분적분으로 바꾼다. \(\delta u\)와 나타나는 모든 도함수가 경계에서 사라지고 그 밖에서는 임의라고 가정하면 잘 알려진 바와 같이
\begin{equation}
\tag{2}
\delta I=\int\!\cdots\!\int \delta f\,d x
=\int\!\cdots\!\int
\left(\sum_i\psi_i\!\left(x,u,\frac{\partial u}{\partial x},\cdots\right)
\delta u_i\right)d x
\end{equation}
를 얻는다. 여기서 \(\psi_i\)들은 라그랑주 식, 곧 대응 변분문제 \(\delta I=0\)의 라그랑주 방정식의 좌변이다. 경계항들도 함께 쓰면 이 적분관계에는 \(\delta u\)와 그 도함수에 관한 적분 없는 항등식이 대응한다. 부분적분이 보여 주듯 이 경계항들은 발산식의 적분, 곧
\[
\operatorname{Div}A=
\frac{\partial A_1}{\partial x_1}+\cdots+
\frac{\partial A_n}{\partial x_n}
\]
꼴의 식의 적분이다. 여기서 \(A\)는 \(\delta u\)와 그 도함수에 관하여 선형이다. 따라서
\begin{equation}
\tag{3}
\sum_i\psi_i\,\delta u_i=\delta f+\operatorname{Div}A
\end{equation}
를 얻는다.

특히 \(f\)가 \(u\)의 일차도함수만 포함하면, 단일적분의 경우 항등식 (3)은 호인이 \glqq{}라그랑주 중심방정식\grqq{}이라 부른 식과 같다.
\begin{equation}
\tag{4}
\sum_i\psi_i\,\delta u_i
=\delta f-\frac{\dd}{\dd x}
\left(\sum_i\frac{\partial f}{\partial u_i'}\delta u_i\right),
\qquad u_i'=\frac{\dd u_i}{\dd x}.
\end{equation}
\(n\)중적분의 경우 (3)은
\begin{equation}
\tag{5}
\begin{aligned}
\sum_i\psi_i\,\delta u_i
={}&\delta f
-\frac{\partial}{\partial x_1}
\left(\sum_i\frac{\partial f}
{\partial\left(\frac{\partial u_i}{\partial x_1}\right)}
\delta u_i\right)\\
&-\cdots-
\frac{\partial}{\partial x_n}
\left(\sum_i\frac{\partial f}
{\partial\left(\frac{\partial u_i}{\partial x_n}\right)}
\delta u_i\right)
\end{aligned}
\end{equation}
이 된다.

단일적분이고 \(u\)의 도함수가 \(\varkappa\)차까지 나타나는 경우 (3)은
\begin{equation}
\tag{6}
\begin{aligned}
\sum_i\psi_i\,\delta u_i
={}&\delta f\\
&-\frac{\dd}{\dd x}\left\{\sum_i\left[
\binom11\frac{\partial f}{\partial u_i^{(1)}}\delta u_i+
\binom21\frac{\partial f}{\partial u_i^{(2)}}\delta u_i^{(1)}
+\cdots+
\binom{\varkappa}{1}\frac{\partial f}{\partial u_i^{(\varkappa)}}
\delta u_i^{(\varkappa-1)}\right]\right\}\\
&+\frac{d^2}{d x^2}\left\{\sum_i\left[
\binom22\frac{\partial f}{\partial u_i^{(2)}}\delta u_i+
\binom32\frac{\partial f}{\partial u_i^{(3)}}\delta u_i^{(1)}
+\cdots+
\binom{\varkappa}{2}\frac{\partial f}{\partial u_i^{(\varkappa)}}
\delta u_i^{(\varkappa-2)}\right]\right\}\\
&+\cdots+
(-1)^\varkappa\frac{d^\varkappa}{d x^\varkappa}
\left\{\sum_i\binom{\varkappa}{\varkappa}
\frac{\partial f}{\partial u_i^{(\varkappa)}}\delta u_i\right\}.
\end{aligned}
\end{equation}
\(n\)중적분에도 이에 대응하는 항등식이 성립한다. 특히 \(A\)에는 \(\delta u\)의 \((\varkappa-1)\)차 도함수까지 들어 있다. (4), (5), (6)이 실제로 라그랑주 식 \(\psi_i\)들을 정의한다는 사실은, 우변의 결합에서 \(\delta u\)의 모든 고차도함수가 소거되는 반면 부분적분이 유일하게 이끄는 관계 (2)는 충족된다는 데서 따라온다.

이하에서 다룰 두 정리는 다음과 같다.

\textbf{I.} 적분 \(I\)가 \(\mathfrak G_\varrho\)에 관하여 불변이면 라그랑주 식들의 선형독립인 결합 \(\varrho\)개가 발산식이 된다. 역으로 그러한 관계들에서 \(I\)의 \(\mathfrak G_\varrho\)-불변성이 따라온다. 이 정리는 매개변수가 무한히 많은 극한적인 경우에도 성립한다.

\textbf{II.} 적분 \(I\)가 임의함수들이 \(\sigma\)차 도함수까지 나타나는 \(\mathfrak G_{\infty\varrho}\)에 관하여 불변이면, 라그랑주 식들과 그 \(\sigma\)차까지의 도함수들 사이에 \(\varrho\)개의 항등관계가 성립한다. 여기서도 역이 성립한다\srcfn{1)}{몇 가지 자명한 예외는 § 2, 2의 주를 보라.}.

혼합군에 대해서는 두 정리의 명제가 모두 성립한다. 따라서 종속관계와, 그와 독립인 발산관계가 함께 나타난다.

이 항등식들에서 대응 변분문제로 넘어가 \(\psi=0\)으로 두면\srcfn{2)}{조금 더 일반적으로는 \(\psi_i=T_i\)로 둘 수도 있다. § 3의 첫째 주를 보라.}, 정리 I은 일차원의 경우 발산식이 전미분이 되므로 \(\varrho\)개의 제1적분이 존재함을 말한다. 다만 이 제1적분들 사이에는 비선형 종속관계가 있을 수 있다\srcfn{3)}{§ 3의 끝을 보라.}. 다차원의 경우에는 최근 흔히 \glqq{}보존법칙\grqq{}이라 부르는 발산방정식을 얻는다. 정리 II는 라그랑주 방정식 가운데 \(\varrho\)개가 나머지 방정식들의 귀결임을 말한다.

역을 제외한 정리 II의 가장 간단한 예는 바이어슈트라스의 매개변수표시이다. 잘 알려져 있듯 일차동차성 아래에서는 독립변수 \(x\)를 \(u\)는 바꾸지 않는 \(x\)의 임의함수로 바꾸어도 적분이 불변이다.
\[
y=p(x),\qquad v_i(y)=u_i(x).
\]
따라서 도함수 없이 임의함수 하나가 나타나며, 이에 대응하여 라그랑주 식들 자체 사이에 잘 알려진 선형관계
\[
\sum_i\psi_i\,\frac{\dd u_i}{\dd x}=0
\]
가 성립한다.

또 다른 예는 물리학자들의 \glqq{}일반상대성이론\grqq{}이다. 여기서는 모든 \(x\)-변환의 군
\[
y_i=p_i(x)
\]
을 다루고, \(u\), 곧 \(g_{\mu\nu}\)와 \(q\)라 표시되는 양들은 이 변환이 이차 및 일차 미분형식의 계수에 유도하는 변환을 받는다. 이 변환들에는 임의함수 \(p(x)\)의 일차도함수가 들어 있다. 이에 대응하여 라그랑주 식들과 그 일차도함수들 사이에 잘 알려진 \(n\)개의 종속관계가 성립한다\srcfn{4)}{예컨대 클라인의 서술을 보라.}.

특히 변환에서 \(u(x)\)의 도함수를 허용하지 않고, 변환된 독립변수들은 \(u\)가 아니라 \(x\)에만 의존하게 하여 군을 특수화하면 다음이 성립한다. § 5에서 보이듯 \(I\)의 불변성에서 \(\sum_i\psi_i\delta u_i\)의 상대불변성이 따라오고\srcfn{1)}{즉 \(\sum_i\psi_i\delta u_i\)는 변환할 때 한 인자를 얻는다.}, 매개변수에 적절한 변환을 가하면 정리 I에 나타나는 발산식들도 상대불변이다. 여기서 앞에서 말한 제1적분들도 그 군을 허용함이 따라온다. 정리 II에서도 임의함수들로 묶은 종속관계의 좌변들이 상대불변이다. 또 여기서 발산이 항등적으로 사라지면서 군을 허용하는 함수가 하나 생기는데, 물리학자들의 상대성이론에서는 이것이 종속관계와 에너지법칙의 연관을 매개한다\srcfn{2)}{클라인의 둘째 글을 보라.}. 끝으로 정리 II는 \glqq{}일반상대성\grqq{}에서 고유한 에너지법칙이 실패한다는 힐베르트의 관련 주장을 군론적으로 증명한다. 이러한 부가 설명까지 포함하면 정리 I은 역학 등에서 알려진 제1적분에 관한 모든 정리를 포함하고, 정리 II는 \glqq{}일반상대성이론\grqq{}의 가능한 가장 넓은 군론적 일반화라고 할 수 있다.

\bigskip
\section*{§ 2. 발산관계와 종속관계}

\(\mathfrak G\)를 유한 또는 무한 연속군이라 하자. 그러면 항등변환에 매개변수 \(\varepsilon\), 또는 임의함수 \(p(x)\)의 값 0이 대응하게 언제나 만들 수 있다\srcfn{3)}{예컨대 Lie, \emph{Grundlagen}, 331쪽을 보라. 임의함수를 다루는 경우에는 매개변수의 특수값 \(a^\sigma\)를 고정함수 \(p^\sigma,\partial p^\sigma/\partial x,\ldots\)로 바꾸고, 이에 따라 \(a^\sigma+\varepsilon\)을 \(p^\sigma+p(x),\partial p^\sigma/\partial x+\partial p/\partial x,\ldots\)로 바꾸어야 한다.}. 따라서 가장 일반적인 변환은
\[
\begin{aligned}
y_i&=A_i\!\left(x,u,\frac{\partial u}{\partial x},\ldots\right)
=x_i+\varDelta x_i+\cdots,\\
v_i(y)&=B_i\!\left(x,u,\frac{\partial u}{\partial x},\ldots\right)
=u_i+\varDelta u_i+\cdots
\end{aligned}
\]
꼴이다. 여기서 \(\varDelta x_i,\varDelta u_i\)는 \(\varepsilon\), 또는 \(p(x)\)와 그 도함수들에 관한 최저차항이다. 이 항들은 그 양들에 관하여 선형이라고 가정한다. 뒤에서 보이듯 이는 일반성을 제한하지 않는다.

이제 적분 \(I\)가 \(\mathfrak G\)에 관하여 불변이어서 관계 (1)을 만족한다고 하자. 그러면 특히 \(I\)는 \(\mathfrak G\)에 들어 있는 무한소변환
\[
y_i=x_i+\varDelta x_i,\qquad v_i(y)=u_i+\varDelta u_i
\]
에 관하여도 불변이다. 이 변환에 대하여 (1)은
\begin{equation}
\tag{7}
\begin{aligned}
0=\varDelta I
&=\int\!\cdots\!\int
f\!\left(y,v(y),\frac{\partial v}{\partial y},\cdots\right)d y\\
&\quad-\int\!\cdots\!\int
f\!\left(x,u(x),\frac{\partial u}{\partial x},\cdots\right)d x
\end{aligned}
\end{equation}
가 된다. 첫 적분은 \(x\)-영역에 대응하는 \(x+\varDelta x\)-영역에서 해야 한다. 그러나 무한소 \(\varDelta x\)에 성립하는 변형
\begin{equation}
\tag{8}
\begin{aligned}
\int\!\cdots\!\int
f\!\left(y,v(y),\frac{\partial v}{\partial y},\ldots\right)\dd y
&=\int\!\cdots\!\int
f\!\left(x,v(x),\frac{\partial v}{\partial x},\ldots\right)\dd x\\
&\quad+\int\!\cdots\!\int
\operatorname{Div}(f\cdot\varDelta x)\,\dd x
\end{aligned}
\end{equation}
을 사용하면 이를 \(x\)-영역에서의 적분으로 바꿀 수 있다.

따라서 무한소변환 \(\varDelta u\) 대신 변분
\begin{equation}
\tag{9}
\bar\delta u_i=v_i(x)-u_i(x)
=\varDelta u_i-\sum_\lambda
\frac{\partial u_i}{\partial x_\lambda}\varDelta x_\lambda
\end{equation}
를 도입하면 (7), (8)은
\begin{equation}
\tag{10}
0=\int\!\cdots\!\int
\left\{\bar\delta f+\operatorname{Div}(f\cdot\varDelta x)\right\}\dd x
\end{equation}
가 된다. 관계 (10)은 모든 임의의 영역에서 적분할 때 성립하므로 피적분함수는 항등적으로 사라져야 한다. 따라서 \(I\)의 불변성을 나타내는 리의 미분방정식은 관계
\begin{equation}
\tag{11}
\bar\delta f+\operatorname{Div}(f\cdot\varDelta x)=0
\end{equation}
이 된다. (3)에 따라 \(\bar\delta f\)를 라그랑주 식으로 나타내면
\begin{equation}
\tag{12}
\sum_i\psi_i\bar\delta u_i=\operatorname{Div}B,
\qquad B=A-f\cdot\varDelta x
\end{equation}
를 얻는다. 이 관계는 각 불변적분 \(I\)에 대하여 나타나는 모든 인수에 관한 항등식이고, 바로 \(I\)에 대한 리의 미분방정식의 구하는 형태이다\footnote{자명한 경우에는 (12)가 \(0=0\)이 된다. 이는 \(\varDelta x,\varDelta u\)가 \(u\)의 도함수에도 의존할 때에만 일어날 수 있으며, \(\operatorname{Div}(f\cdot\varDelta x)=0,\ \bar\delta u=0\)인 경우이다. 이러한 무한소변환들은 언제나 군에서 분리하고, 정리를 정식화할 때에는 나머지 매개변수 또는 임의함수만 세어야 한다. 나머지 무한소변환들이 언제나 다시 군을 이루는지는 여기서 결정하지 않는다.}.

먼저 \(\mathfrak G\)가 유한매개변수 연속군 \(\mathfrak G_\rho\)라고 하자. 가정에 따라 \(\Delta u,\Delta x\)가 매개변수 \(\varepsilon_1,\ldots,\varepsilon_\rho\)에 관하여 선형이므로, (9)에 따라 \(\bar\delta u\)와 그 도함수들도 그러하다. 따라서 \(A,B\)도 \(\varepsilon\)에 관하여 선형이다. 그러므로
\[
B=B^{(1)}\varepsilon_1+\cdots+B^{(\rho)}\varepsilon_\rho,\qquad
\bar\delta u=\bar\delta u^{(1)}\varepsilon_1+\cdots+
\bar\delta u^{(\varrho)}\varepsilon_\varrho
\]
로 두자. 여기서 \(\bar\delta u^{(1)},\ldots\)는 \(x,u,\partial u/\partial x,\ldots\)의 함수들이다. 그러면 (12)에서 구하는 발산관계
\begin{equation}
\tag{13}
\sum\psi_i\bar\delta u_i^{(1)}=\operatorname{Div}B^{(1)}
\quad\ldots,\quad
\sum_i\psi_i\bar\delta u_i^{(\varrho)}=\operatorname{Div}B^{(\varrho)}
\end{equation}
가 따라온다.

따라서 라그랑주 식들의 선형독립인 결합 \(\varrho\)개가 발산식이 된다. 선형독립성은 다음에서 따라온다. (9)에 따라 \(\bar\delta u=0,\ \varDelta x=0\)이면 \(\varDelta u=0,\ \varDelta x=0\)이 되어 무한소변환들 사이에 종속관계가 생긴다. 그러나 어떠한 매개변수값에도 그러한 종속관계가 성립할 수 없다. 그렇지 않으면 이 무한소변환들을 적분하여 다시 얻는 \(\mathfrak G_\varrho\)가 \(\varrho\)보다 적은 본질매개변수에 의존할 것이기 때문이다. 다른 가능성인 \(\bar\delta u=0,\ \operatorname{Div}(f\varDelta x)=0\)은 이미 제외하였다. 이 추론은 매개변수가 무한히 많은 극한적인 경우에도 성립한다.

이제 \(\mathfrak G\)가 무한 연속군 \(\mathfrak G_{\infty\rho}\)라고 하자. 그러면 \(\bar\delta u\)와 그 도함수들, 따라서 \(B\)도 임의함수 \(p(x)\)와 그 도함수들에 관하여 다시 선형이 된다\footnote{\(p\)가 \(u,\partial u/\partial x,\ldots\)와 무관하다고 가정해도 일반성을 잃지 않는다는 사실은 역정리가 보여 준다.}. (12)와는 아직 무관하게 \(\bar\delta u\)의 값을 대입하여
\[
\sum\psi_i\bar\delta u_i
=\sum_{\lambda,i}\psi_i\left\{
a_i^{(\lambda)}(x,u,\ldots)p^{(\lambda)}(x)
+b^{(\lambda)}(x,u,\ldots)
\frac{\partial p^{(\lambda)}}{\partial x}
+\cdots+c_i^{(\lambda)}(x,u,\ldots)
\frac{\partial^\sigma p^{(\lambda)}}{\partial x^\sigma}
\right\}
\]
라고 하자. 이제 부분적분 공식과 유사하게 항등식
\[
\varphi(x,u,\ldots)\frac{\partial^\tau p(x)}{\partial x^\tau}
=(-1)^\tau\frac{\partial^\tau\varphi}{\partial x^\tau}\,p(x)
\quad\bmod\text{ 발산식}
\]
에 따라 \(p\)의 도함수들을 \(p\) 자체와, \(p\) 및 그 도함수들에 관하여 선형인 발산식들로 바꿀 수 있다. 따라서
\begin{equation}
\tag{14}
\begin{aligned}
\sum_i\psi_i\bar\delta u_i
={}&\sum_\lambda\left\{
a_i^{(\lambda)}\psi_i
-\frac{\partial}{\partial x}\bigl(b_i^{(\lambda)}\psi_i\bigr)
+\cdots+(-1)^\sigma
\frac{\partial^\sigma}{\partial x^\sigma}
\bigl(c_i^{(\lambda)}\psi_i\bigr)
\right\}p^{(\lambda)}+\operatorname{Div}\varGamma
\end{aligned}
\end{equation}
를 얻고, 이를 (12)와 결합하면
\begin{equation}
\tag{15}
\sum\left\{
a_i^{(\lambda)}\psi_i
-\frac{\partial}{\partial x}\bigl(b_i^{(\lambda)}\psi_i\bigr)
+\cdots+(-1)^\sigma
\frac{\partial^\sigma}{\partial x^\sigma}
\bigl(c_i^{(\lambda)}\psi_i\bigr)
\right\}p^{(\lambda)}
=\operatorname{Div}(B-\varGamma)
\end{equation}
을 얻는다. 이제 임의의 영역에 걸쳐 (15)의 \(n\)중적분을 취하고, \(B-\varGamma\)에 나타나는 모든 도함수와 함께 \(p(x)\)가 경계에서 사라지도록 \(p(x)\)를 택한다. 발산식의 적분은 경계적분으로 환원되므로, (15)의 왼쪽 적분도 경계에서 충분히 많은 도함수와 함께 사라지는 임의의 \(p(x)\)에 대하여 영이다. 잘 알려진 추론에 따라 그 피적분함수는 모든 \(p(x)\)에 대하여 사라져야 하므로 다음 \(\varrho\)개의 관계를 얻는다.
\begin{equation}
\tag{16}
\sum\left\{
a_i^{(\lambda)}\psi_i
-\frac{\partial}{\partial x}\bigl(b_i^{(\lambda)}\psi_i\bigr)
+\cdots+(-1)^\sigma
\frac{\partial^\sigma}{\partial x^\sigma}
\bigl(c_i^{(\lambda)}\psi_i\bigr)
\right\}=0,
\qquad(\lambda=1,2,\ldots,\varrho).
\end{equation}
이것들이 \(I\)가 \(\mathfrak G_{\infty\varrho}\)에 대하여 불변일 때 라그랑주 식들과 그 도함수들 사이에 성립하는 구하는 종속관계들이다. 선형독립성은 위와 마찬가지로 드러난다. 역방향의 논증이 (12)로 되돌아가기 때문이며, 또한 제4절에서 더 자세히 설명하듯 무한소변환들로부터 다시 유한변환들을 추론할 수 있기 때문이다. 따라서 \(\mathfrak G_{\infty\varrho}\)에서는 이미 무한소변환들 속에 언제나 \(\varrho\)개의 임의함수가 나타난다. (15)와 (16)에서는 또한 \(\operatorname{Div}(B-\varGamma)=0\)이 따라온다.

``혼합군''에 맞추어 \(\varDelta x\)와 \(\varDelta u\)를 \(\varepsilon\)들과 \(p(x)\)들에 관하여 선형으로 놓으면, 한 번은 \(p(x)\)들을 영으로, 다른 한 번은 \(\varepsilon\)들을 영으로 둠으로써 발산관계 (13)과 종속관계 (16)이 모두 성립함을 알 수 있다.

\bigskip
\section*{\S{}3. 유한군의 경우 역정리}

역을 보이려면 우선 앞의 고찰을 본질적으로 역순으로 밟아야 한다. (13)이 성립하면 이를 \(\varepsilon\)들로 곱하고 더하여 (12)를 얻고, 항등식 (3)에 따라 여기서
\[
\bar\delta f+\operatorname{Div}(A-B)=0
\]
이라는 관계를 얻는다. 따라서
\[
\varDelta x=\frac1f\,(A-B)
\]
로 두면 (11)에 도달한다. 이를 적분하면 마침내 (7), 즉 \(\varDelta I=0\)을 얻으므로, \(I\)는 \(\varDelta x,\varDelta u\)로 정해지는 무한소변환에 대하여 불변이다. 여기서 \(\varDelta u\)는 (9)에 따라 \(\varDelta x\)와 \(\bar\delta u\)로부터 정해지고, \(\varDelta x\)와 \(\varDelta u\)는 매개변수들에 관하여 선형이 된다. 잘 알려진 바와 같이 \(\varDelta I=0\)이면 다음 연립계를 적분하여 생기는 유한변환들에 대해서도 \(I\)가 불변이다.
\begin{equation}
\tag{17}
\frac{\dd x_i}{\dd t}=\varDelta x_i,\qquad
\frac{\dd u_i}{\dd t}=\varDelta u_i,\qquad
\left(\begin{array}{l}
x_i=y_i,\\
u_i=v_i
\end{array}\right\}
\quad\text{단, }t=0.
\end{equation}
이 유한변환들은 \(\rho\)개의 매개변수 \(a_1,\ldots,a_\rho\), 즉 \(t\varepsilon_1,\ldots,t\varepsilon_\rho\)라는 결합들을 포함한다. 선형독립인 발산관계 (13)이 정확히 \(\rho\)개뿐이라는 가정으로부터, 유한변환들이 도함수 \(\partial u/\partial x\)를 포함하지 않는 한 언제나 군을 이룬다는 것도 따라온다. 그렇지 않다면 리 괄호과정으로 생기는 무한소변환들 가운데 적어도 하나는 나머지 \(\rho\)개의 선형결합이 아닐 것이다. \(I\)는 이 변환도 허용하므로 선형독립인 발산관계가 \(\rho\)개보다 많아질 것이다. 아니면 이 무한소변환은 \(\bar\delta u=0,\ \operatorname{Div}(f\Delta x)=0\)인 특수한 형태일 것이다. 그러나 그러면 가정에 어긋나게 \(\Delta x\) 또는 \(\Delta u\)가 도함수에 의존할 것이다. \(\Delta x\)나 \(\Delta u\)에 도함수가 나타날 때 이 경우가 실제로 생길 수 있는지는 여기서 미결로 남겨야 한다. 이때 군의 성질을 다시 얻으려면 위에서 정한 \(\Delta x\)에 \(\operatorname{Div}(f\Delta x)=0\)을 만족하는 모든 함수 \(\Delta x\)를 더해야 한다. 그러나 약속에 따라 이로써 추가되는 매개변수들은 세지 않는다. 이로써 역정리가 증명되었다.

이 역정리로부터 실제로 \(\varDelta x\)와 \(\varDelta u\)를 매개변수들에 관하여 선형이라고 가정할 수 있다는 것도 따라온다. 만일 \(\varDelta u\)와 \(\varDelta x\)가 \(\varepsilon\)에 관하여 더 높은 차수의 형식들이라면, \(\varepsilon\)의 거듭제곱곱들이 선형독립이므로 전적으로 같은 종류의 관계 (13)이 단지 더 많이 생긴다. 역정리에 따라 이 관계들로부터, 무한소변환들이 매개변수를 선형으로 포함하는 한 군에 대한 \(I\)의 불변성이 따라온다. 이 군이 정확히 \(\varrho\)개의 매개변수를 포함해야 한다면, 원래 \(\varepsilon\)의 고차항들로부터 얻은 발산관계들 사이에 선형 종속관계가 있어야 한다.

\(\varDelta x\)와 \(\varDelta u\)가 \(u\)의 도함수들도 포함하는 경우, 유한변환들은 \(u\)의 무한히 많은 도함수에 의존할 수 있다는 점도 덧붙여 두자. 실제로 이 경우 (17)을 적분하면서
\[
\frac{\dd^2x_i}{\dd t^2},\qquad
\frac{\dd^2u_i}{\dd t^2}
\]
를 정하려 하면
\[
\varDelta\!\left(\frac{\partial u}{\partial x_\varkappa}\right)
=\frac{\partial\,\varDelta u}{\partial x_\varkappa}
-\sum_\lambda\frac{\partial u}{\partial x_\lambda}
\frac{\partial\,\varDelta x_\lambda}{\partial x_\varkappa}
\]
가 나타나므로, 일반적으로 각 단계에서 \(u\)의 도함수 차수가 증가한다. 한 예는 다음과 같다.
\[
f=\frac12u'^2,\qquad
\psi=-u'',\qquad
\psi\cdot x=\frac{\dd}{\dd x}\{u-u'x\},\qquad
\bar\delta u=x\cdot\varepsilon;
\]
\[
\Delta x=-\frac{2u}{u'^2}\,\varepsilon,\qquad
\Delta u=\left(x-\frac{2u}{u'}\right)\varepsilon.
\]
발산식의 라그랑주 식들은 항등적으로 사라지므로, 역정리는 끝으로 다음 사실도 보여 준다. \(I\)가 한 \(\mathfrak G_\varrho\)를 허용한다면, \(I\)와 경계적분, 즉 발산식의 적분만큼만 다른 모든 적분도 같은 \(\bar\delta u\)를 갖는 한 \(\mathfrak G_\varrho\)를 허용한다. 그 무한소변환들은 일반적으로 \(u\)의 도함수들을 포함한다. 이를테면 위의 예에 따라
\[
f^*=\frac12\left\{u'^2-\frac{\dd}{\dd x}\left(\frac{u^2}{x}\right)\right\}
\]
는 무한소변환 \(\Delta u=x\varepsilon,\ \Delta x=0\)을 허용하지만, \(f\)에 대응하는 무한소변환들에는 \(u\)의 도함수들이 나타난다.

변분문제로 넘어가, 즉 \(\psi_i=0\)으로 두면\footnote{물리학에서는 \(\psi_i=0\), 또는 좀 더 일반적으로 새로 추가되는 함수들을 \(T_i\)라 할 때 \(\psi_i=T_i\)를 ``장방정식''이라 부른다. \(\psi_i=T_i\)인 경우 항등식 (13)은 방정식 \(\operatorname{Div}B^{(\lambda)}=\sum_iT_i\bar\delta u_i^{(\lambda)}\)로 바뀌며, 물리학에서는 이것들도 보존법칙이라 부른다.}, (13)은 방정식
\[
\operatorname{Div}B^{(1)}=0,\ldots,\qquad
\operatorname{Div}B^{(\varrho)}=0
\]
으로 바뀌는데, 이들은 흔히 ``보존법칙''이라 불린다. 일차원인 경우에는
\[
B^{(1)}=\mathrm{const.},\ldots,\qquad
B^{(\varrho)}=\mathrm{const.}
\]
가 따라온다. \(\Delta u\)와 \(\Delta x\)가 \(f\)에 나타나는 \(\varkappa\)차 도함수보다 더 높은 도함수를 포함하지 않는 한, (6)에 따라 여기서 \(B\)들은 기껏해야 \(u\)의 \((2\varkappa-1)\)차 도함수를 포함한다. 일반적으로 \(\psi\)에는 \(2\varkappa\)차 도함수가 나타나므로\footnote{\(f\)가 \(\varkappa\)차 도함수들에 관하여 비선형인 한 그러하다.}, 이로써 \(\rho\)개의 제1적분이 존재한다. 이들 사이에 비선형 종속관계가 있을 수 있다는 사실은 위의 \(f\)가 다시 보여 준다. 선형독립인 \(\Delta u=\varepsilon_1,\ \Delta x=\varepsilon_2\)에는 선형독립인 관계
\[
u''=\frac{\dd}{\dd x}u',\qquad
u''u'=\frac12\frac{\dd}{\dd x}(u')^2
\]
가 대응하지만, 제1적분 \(u'=\mathrm{const.},\ u'^2=\mathrm{const.}\) 사이에는 비선형 종속관계가 있다. 여기서는 \(\varDelta u,\varDelta x\)가 \(u\)의 도함수를 포함하지 않는 초등적인 경우를 다룬다\footnote{그렇지 않으면 모든 \(\lambda\)에 대하여 \(u'^\lambda=\mathrm{const.}\)도 얻는다. 이는
\[
u''(u')^{\lambda-1}=\frac1\lambda\frac{\dd}{\dd x}(u')^\lambda
\]
에 대응한다.}.

\bigskip
\section*{\S{}4. 무한군의 경우 역정리}

먼저 \(\varDelta x\)와 \(\varDelta u\)의 선형성을 가정해도 제한이 없음을 보이자. 여기서는 역정리를 쓰지 않고도, \(\mathfrak G_{\infty\varrho}\)가 형식적으로 정확히 \(\varrho\)개의 임의함수에 의존한다는 사실에서 이를 얻는다. 실제로 비선형인 경우에는 변환들을 합성할 때 최저차항들이 서로 더해지므로 임의함수의 수가 늘어남을 알 수 있다. 이를테면
\[
\begin{aligned}
y&=A\!\left(x,u,\frac{\partial u}{\partial x},\cdots;p\right)\\
&=x+\sum a(x,u,\ldots)p^\nu
+b(x,u,\ldots)p^{\nu-1}\frac{\partial p}{\partial x}
+c\,p^{\nu-2}\left(\frac{\partial p}{\partial x}\right)^2+\cdots
+d\left(\frac{\partial p}{\partial x}\right)^\nu+\cdots,
\end{aligned}
\]
\[
p^\nu=(p^{(1)})^{\nu_1}\cdots(p^{(\varrho)})^{\nu_\varrho},
\]
이고 이에 대응하여
\[
v=B\!\left(x,u,\frac{\partial u}{\partial x},\ldots;p\right)
\]
라고 하자. 이를
\[
z=A\!\left(y,v,\frac{\partial v}{\partial y},\ldots;q\right)
\]
와 합성하면 최저차항에 대하여
\[
\begin{aligned}
z={}&x+\sum a(p^\nu+q^\nu)
+b\left\{p^{\nu-1}\frac{\partial p}{\partial x}
+q^{\nu-1}\frac{\partial q}{\partial x}\right\}\\
&+c\left\{p^{\nu-2}\left(\frac{\partial p}{\partial x}\right)^2
+q^{\nu-2}\left(\frac{\partial q}{\partial x}\right)^2\right\}+\cdots
\end{aligned}
\]
를 얻는다. 여기서 \(a,b\)와 다른 어떤 계수가 영이 아니어서 \(\sigma>1\)인 항
\[
p^{\nu-\sigma}\left(\frac{\partial p}{\partial x}\right)^\sigma
+q^{\nu-\sigma}\left(\frac{\partial q}{\partial x}\right)^\sigma
\]
이 실제로 나타난다면, 이를 한 함수의 미분계수나 그러한 미분계수의 거듭제곱곱으로 쓸 수 없다. 따라서 가정에 어긋나게 임의함수의 수가 늘어난다. \(a,b\)와 다른 모든 계수가 영이라면, 지수 \(\nu_1\ldots\nu_\varrho\)의 값에 따라 둘째 항이 첫째 항의 미분계수가 되어 실제로 선형성이 생기거나(예를 들어 \(\mathfrak G_{\infty1}\)에서는 언제나 그러하다), 이 경우에도 임의함수의 수가 늘어난다. 그러므로 무한소변환들은 \(p(x)\)의 선형성 때문에 선형 편미분방정식계를 만족한다. 또한 군의 성질이 성립하므로, 이들은 리의 정의에 따른 ``무한 무한소변환군''을 이룬다(\emph{Grundlagen}, \S{}10).

이제 역정리는 유한군의 경우와 비슷하게 얻어진다. 종속관계 (16)이 성립하면 이를 \(p^{(\lambda)}(x)\)로 곱하고 더한 뒤 항등변형 (14)를 적용하여
\[
\sum_i\psi_i\bar\delta u_i=\operatorname{Div}\varGamma
\]
를 얻는다. 여기서 제3절과 마찬가지로 \(\Delta x,\Delta u\)가 정해지고, 이 무한소변환들에 대한 \(I\)의 불변성이 따라온다. 이 변환들은 실제로 \(\rho\)개의 임의함수와 그 \(\sigma\)차까지의 도함수들에 선형으로 의존한다. 이 무한소변환들이 \(\partial u/\partial x,\ldots\)를 포함하지 않을 때 반드시 군을 이룬다는 사실도 제3절에서와 같이 따라온다. 그렇지 않다면 합성에 의해 더 많은 임의함수가 나타날 것이지만, 가정에 따르면 종속관계 (16)은 \(\rho\)개뿐이기 때문이다. 따라서 이들은 ``무한 무한소변환군''을 이룬다. 그러나 그러한 군은(\emph{Grundlagen}, 정리 VII, 391쪽) 이로써 정의되는, 리의 의미에서의 어떤 ``유한변환들의 무한군 \(\mathfrak G\)''의 가장 일반적인 무한소변환들로 이루어진다. 각 유한변환은 무한소변환들로부터 생성되며(\emph{Grundlagen}, \S{}7)\footnote{특히 이로부터 \(\mathfrak G_{\infty\rho}\)의 무한소변환 \(\Delta x,\Delta u\)로부터 생성된 군 \(\mathfrak G\)가 다시 \(\mathfrak G_{\infty\rho}\)로 돌아온다는 것이 따라온다. 실제로 \(\mathfrak G_{\infty\rho}\)는 임의함수에 의존하면서 \(\Delta x,\Delta u\)와 다른 무한소변환들을 포함하지 않는다. 또한 매개변수에 의존하면서 이들과 독립인 변환들도 포함할 수 없다. 그렇다면 혼합군이 될 것이기 때문이다. 한편 위에서 본 바와 같이 무한소변환들이 유한변환들을 정한다.}, 따라서 다음 연립계를 적분하여 생긴다.
\[
\frac{\dd x_i}{\dd t}=\varDelta x_i,\qquad
\frac{\dd u_i}{\dd t}=\varDelta u_i,\qquad
\left(\begin{array}{l}
x_i=y_i,\\
u_i=v_i
\end{array}\right\}
\quad\text{단, }t=0.
\]
다만 여기서는 임의의 \(p(x)\)가 \(t\)에도 의존한다고 가정해야 할 수 있다. 따라서 \(\mathfrak G\)는 실제로 \(\rho\)개의 임의함수에 의존한다. 특히 \(p(x)\)를 \(t\)와 무관하다고 해도 충분하다면, 이 의존성은 임의함수 \(q(x)=t\,p(x)\)에 관하여 해석적이다\footnote{언제나 이 후자의 경우가 성립하는가라는 문제는 리가 다른 형식으로 제기하였다(\emph{Grundlagen}, \S{}7 및 \S{}13 끝).}. \(\partial u/\partial x,\ldots\)가 나타나는 경우에는 같은 결론을 내리기 전에 \(\bar\delta u=0,\ \operatorname{Div}(f\Delta x)=0\)인 무한소변환들을 더해야 할 수도 있다.

리의 한 예(\emph{Grundlagen}, \S{}7)에 이어, 명시적 공식들까지 나아갈 수 있는 꽤 일반적인 경우를 하나 더 들자. 이 공식들은 동시에 임의함수들의 도함수가 정확히 \(\sigma\)차까지 나타난다는 것을 보여 주므로 이 경우 역정리는 완전하다. 이는 모든 \(x\)의 변환과, 그로써 ``유도된'' \(u\)의 변환으로 이루어진 군에 대응하는 무한소변환군이다. 다시 말해 \(\Delta u\), 따라서 \(u\)가 \(\Delta x\)에 나타나는 임의함수들에만 의존하는 \(u\)의 변환이다. 또한 \(\Delta u\)에는 도함수 \(\partial u/\partial x,\ldots\)가 나타나지 않는다고 가정한다. 그러면
\[
\varDelta x_i=p^{(i)}(x);\qquad
\varDelta u_i=
\sum_{\lambda=1}^{n}\left\{
a^{(\lambda)}(x,u)p^{(\lambda)}
+b^{(\lambda)}\frac{\partial p^{(\lambda)}}{\partial x}
+\cdots+c^{(\lambda)}
\frac{\partial^\sigma p^{(\lambda)}}{\partial x^\sigma}
\right\}
\]
이다. 무한소변환 \(\varDelta x=p(x)\)는 임의의 \(g(y)\)를 갖는 모든 변환 \(x=y+g(y)\)를 생성한다. 특히 \(p(x)\)의 \(t\) 의존성을 정하여 일항군
\begin{equation}
\tag{18}
x_i=y_i+t\cdot g_i(y)
\end{equation}
을 생성할 수 있다. 이 군은 \(t=0\)일 때 항등변환이고 \(t=1\)일 때 구하는 \(x=y+g(y)\)가 된다. 실제로 (18)을 미분하면
\begin{equation}
\tag{19}
\frac{d x_i}{d t}=g_i(y)=p^{(i)}(x,t)
\end{equation}
를 얻는다. 여기서 \(p(x,t)\)는 (18)을 역으로 풀어 \(g(y)\)로부터 정해진다. 거꾸로 (19)에 보조조건 \(t=0\)일 때 \(x_i=y_i\)를 붙여 적분을 유일하게 정하면 (18)을 얻는다. (18)에 따라 \(\varDelta u\) 안의 \(x\)들을 적분상수 \(y\)와 \(t\)로 바꿀 수 있다. 이때 \(g(y)\)는 정확히 \(\sigma\)차 도함수까지 나타난다. 즉
\[
\frac{\partial p}{\partial x}
=\sum\frac{\partial g}{\partial y_\varkappa}
\frac{\partial y_\varkappa}{\partial x}
\]
에서 \(\partial y/\partial x\)를 \(\partial x/\partial y\)로 나타내고, 일반적으로 \(\partial^\sigma p/\partial x^\sigma\)를
\[
\frac{\partial g}{\partial y},\ldots,
\frac{\partial x}{\partial y},\ldots,
\frac{\partial^\sigma x}{\partial y^\sigma}
\]
로 나타낸 값으로 바꾼다. 따라서 \(u\)를 정하기 위한 방정식계는
\[
\frac{d u_i}{d t}
=F_i\!\left(g(y),\frac{\partial g}{\partial y},\ldots,
\frac{\partial^\sigma g}{\partial y^\sigma},u,t\right),
\qquad(u_i=v_i\ \text{단, }t=0)
\]
가 된다. 여기서는 \(t,u\)만 변수이고 \(g(y),\ldots\)는 계수체에 속한다. 따라서 적분하면
\[
u_i=v_i+B_i\!\left(v,g(y),\frac{\partial g}{\partial y},\ldots,
\frac{\partial^\sigma g}{\partial y^\sigma},t\right)_{t=1}
\]
을 얻는다. 즉 임의함수들의 도함수에 정확히 \(\sigma\)차까지 의존하는 변환들이다. (18)에 따라 \(g(y)=0\)이면 그 안에 항등변환이 포함된다. 또한 주어진 절차는 모든 변환 \(x=y+g(y)\)를 낳고, 이로부터 유도된 \(u\)의 변환이 유일하게 정해지므로, 이 절차가 군 \(\mathfrak G\)를 모두 소진한다. 따라서 군의 성질도 따라온다.

역정리로부터, 임의함수들이 \(u,\partial u/\partial x,\ldots\)가 아니라 \(x\)에만 의존한다고 가정해도 제한이 없다는 사실도 부수적으로 따라온다. 후자의 경우라면 항등변형 (14), 따라서 (15)에는 \(p^{(\lambda)}\)들 외에도 \(\partial p^{(\lambda)}/\partial u,\ \partial p^{(\lambda)}/\partial(\partial u/\partial x),\ldots\)가 나타날 것이다. 이제 \(p^{(\lambda)}\)들을 차례로 \(u,\partial u/\partial x,\ldots\)의 0차, 1차, \(\ldots\) 형식으로 잡고 그 계수들을 \(x\)의 임의함수로 두면, 다시 종속관계 (16)을 단지 더 많이 얻는다. 그러나 위의 역정리에 따라 이들을 \(x\)에만 의존하는 임의함수들과 함께 묶으면 앞의 경우로 돌아간다. 마찬가지로 종속관계와, 이들과 독립인 발산관계가 동시에 나타나면 혼합군이 대응함을 보일 수 있다\footnote{제3절에서와 마찬가지로, 역정리로부터 \(I\)뿐 아니라 발산식의 적분만큼 다른 모든 적분 \(I^*\)도 같은 \(\bar\delta u\)를 갖는 무한군을 허용한다는 것이 여기서도 따라온다. 다만 \(\Delta x\)와 \(\Delta u\)에는 일반적으로 \(u\)의 도함수들이 나타난다. 아인슈타인은 일반상대성이론에서 에너지 정리들을 더 간단히 쓰기 위해 그러한 적분 \(I^*\)를 도입하였다. 둘째 클라인 논문의 표기를 그대로 따라 이 \(I^*\)가 허용하는 무한소변환들을 제시하겠다. 적분 \(I=\int\!\cdots\!\int K\,\dd\omega=\int\!\cdots\!\int\mathfrak R\,\dd S\)는 모든 \(w\)의 변환과, 그로써 유도되는 \(g_{\mu\nu}\)의 변환으로 이루어진 군을 허용한다. 이에 종속관계(클라인의 (30))
\[
\sum\mathfrak R_{\mu\nu}g_\tau^{\mu\nu}
+2\sum\frac{\partial(g^{\mu\nu}\mathfrak R_{\mu\tau})}{\partial w^\sigma}=0
\]
가 대응한다. 이제 \(I^*=\int\!\ldots\!\int \mathfrak{R}^*\,d S\)이며, 여기서 \(\mathfrak R^*=\mathfrak R+\operatorname{Div}\)이다. 따라서 \(\mathfrak R_{\mu\nu}^*=\mathfrak R_{\mu\nu}\)이다. 여기서 \(\mathfrak R_{\mu\nu}^*,\mathfrak R_{\mu\nu}\)는 각각 라그랑주 식을 뜻한다. 그러므로 위의 종속관계들은 \(\mathfrak R_{\mu\nu}^*\)에 대해서도 성립한다. 이를 \(p^\tau\)로 곱하고 더한 뒤 곱의 미분법칙을 역으로 적용하면
\[
\sum\mathfrak R_{\mu\nu}p^{\mu\nu}
+2\operatorname{Div}\!\left(\sum g^{\mu\sigma}\mathfrak R_{\mu\tau}p^\tau\right)=0,
\]
\[
\delta\mathfrak R^*+\operatorname{Div}\!\left(
\sum2g^{\mu\sigma}\mathfrak R_{\mu\tau}p^\tau
-\frac{\partial\mathfrak R^*}{\partial g^{\mu\nu}_\sigma}p^{\mu\nu}
\right)=0
\]
을 얻는다. 이를 리의 미분방정식 \(\delta\mathfrak R^*+\operatorname{Div}(\mathfrak R^*\varDelta w)=0\)과 비교하면
\[
\varDelta w^\sigma=\frac1{\mathfrak R^*\cdot}\left(
\sum2g^{\mu\sigma}\mathfrak R_{\mu\tau}p^\tau
-\frac{\partial\mathfrak R^*}{\partial g^{\mu\nu}_\sigma}p^{\mu\nu}
\right),
\qquad
\varDelta g^{\mu\nu}=p^{\mu\nu}+\sum g^{\mu\nu}_\sigma\varDelta w^\sigma
\]
를 \(I^*\)가 허용하는 무한소변환들로 얻는다. 따라서 이 무한소변환들은 \(g^{\mu\nu}\)의 1차와 2차 도함수에 의존하며, 임의의 \(p\)를 1차 도함수까지 포함한다.}.

\bigskip
\section*{\S{}5. 관계를 이루는 각 성분의 불변성}

변환들에 \(u\)의 도함수가 나타나지 않고 변환된 독립변수들이 \(u\)가 아니라 \(x\)에만 의존한다고 함으로써 군 \(\mathfrak G\)를 통상 고찰하는 가장 단순한 경우로 특수화하면, 공식들을 이루는 각 성분의 불변성을 추론할 수 있다. 먼저 잘 알려진 추론에 따라
\[
\int\!\cdots\!\int\left(\sum_i\psi_i\delta u_i\right)\dd x
\]
가 불변이고, 따라서 \(\delta\)가 임의의 변분을 뜻할 때 \(\sum_i\psi_i\delta u_i\)가 상대불변임을 얻는다\footnote{즉 \(\sum_i\psi_i\delta u_i\)는 변환할 때 한 인자를 얻는데, 대수적 불변식론에서는 이것을 언제나 상대불변성이라 불러 왔다.}. 실제로 한편으로는
\[
\begin{aligned}
\delta I
&=\int\!\cdots\!\int
\delta f\!\left(x,u,\frac{\partial u}{\partial x},\ldots\right)\dd x\\
&=\int\!\cdots\!\int
\delta f\!\left(y,v,\frac{\partial v}{\partial y},\ldots\right)\dd y
\end{aligned}
\]
이고, 다른 한편으로 경계에서 사라지는 \(\delta u,\delta(\partial u/\partial x),\ldots\)에 대하여 다음이 성립한다. \(\delta u,\delta(\partial u/\partial x),\ldots\)의 변환이 선형이고 동차이므로 이에 대응하는 \(\delta v,\delta(\partial v/\partial y),\ldots\)도 경계에서 사라진다.
\[
\int\!\cdots\!\int
\delta f\!\left(x,u,\frac{\partial u}{\partial x},\ldots\right)\dd x
=\int\!\cdots\!\int
\left(\sum_i\psi_i(u,\ldots)\delta u_i\right)\dd x,
\]
\[
\int\!\cdots\!\int
\delta f\!\left(y,v,\frac{\partial v}{\partial y},\ldots\right)\dd y
=\int\!\cdots\!\int
\left(\sum_i\psi_i(v,\ldots)\delta v_i\right)\dd y.
\]
따라서 경계에서 사라지는 \(\delta u,\delta(\partial u/\partial x),\ldots\)에 대하여
\[
\int\!\cdots\!\int
\left(\sum_i\psi_i(u,\ldots)\delta u_i\right)\dd x
=\int\!\cdots\!\int
\left(\sum_i\psi_i(v,\ldots)\delta v_i\right)\dd y
\]
\[
=\int\!\cdots\!\int
\left(\sum_i\psi_i(v,\ldots)\delta v_i\right)
\left|\frac{\partial y_i}{\partial x_\varkappa}\right|\dd x
\]
이다. 셋째 적분에서 \(y,v,\delta v\)를 \(x,u,\delta u\)로 나타내고 이를 첫째 적분과 같게 두면, 경계에서는 사라지지만 그 밖에서는 임의인 \(\delta u\)에 대하여
\[
\int\!\cdots\!\int
\left(\sum_i\chi_i(u,\ldots)\delta u_i\right)d x=0
\]
이라는 관계를 얻는다. 잘 알려진 바와 같이 여기서 임의의 \(\delta u\)에 대하여 피적분함수가 사라짐이 따라온다. 따라서 \(\delta u\)에 관한 다음 항등관계가 성립한다.
\[
\sum_i\psi_i(u,\ldots)\delta u_i
=\left|\frac{\partial y_i}{\partial x_\varkappa}\right|
\left(\sum_i\psi_i(v,\ldots)\delta v_i\right).
\]
이 관계는 \(\sum_i\psi_i\delta u_i\)의 상대불변성, 따라서
\[
\int\!\cdots\!\int
\left(\sum_i\psi_i\delta u_i\right)\dd x
\]
의 불변성을 말한다\footnote{\(y\)도 \(u\)에 의존하면 이 추론은 성립하지 않는다. 그때에는 \(\delta f(y,v,\partial v/\partial y,\ldots)\)가 \(\sum(\partial f/\partial y)\delta y\)라는 항들도 포함하여, 발산변형이 라그랑주 식들로 이어지지 않기 때문이다. \(u\)의 도함수들을 허용할 때도 마찬가지이다. 그러면 \(\delta v\)는 \(\delta u,\delta(\partial u/\partial x),\ldots\)의 선형결합이 되므로, 한 번 더 발산변형을 거친 뒤에야 항등식 \(\int\!\cdots\!\int(\sum\chi_i(u,\ldots)\delta u_i)\dd x=0\)을 얻고, 따라서 오른쪽에는 다시 라그랑주 식들이 나타나지 않는다. \(\int\!\cdots\!\int(\sum\psi_i\delta u_i)\dd x\)의 불변성만으로 발산관계의 존재를 추론할 수 있는가라는 문제는, 역정리에 따르면 동일한 \(\Delta u,\Delta x\)를 줄 필요는 없지만 동일한 \(\bar\delta u\)를 주는 어떤 군에 대하여 \(I\)가 불변임을 그로부터 추론할 수 있는가라는 문제와 동치이다. 단일적분이고 \(f\)에 1차 도함수만 나타나는 특수한 경우에는 유한군에 관하여 라그랑주 식들의 불변성에서 제1적분의 존재를 추론할 수 있다(예: Engel, \emph{Gött. Nachr.} 1916, 270쪽 참조).}.

이것을 유도된 발산관계와 종속관계에 적용하려면, 먼저 \(\varDelta u,\varDelta x\)로부터 유도된 \(\bar\delta u\)가 실제로 변분 \(\delta u\)의 변환법칙을 만족함을 보여야 한다. 다만 \(\bar\delta v\)에 나타나는 매개변수 또는 임의함수는 \(y,v\)에서의 닮은 무한소변환군에 맞게 정해야 한다. \(\mathfrak T_q\)가 \(x,u\)를 \(y,v\)로 옮기는 변환이고 \(\mathfrak T_p\)가 \(x,u\)에서의 무한소변환이라 하자. 그러면 \(y,v\)에서 이에 닮은 변환은 \(\mathfrak T=\mathfrak T_q\mathfrak T_p\mathfrak T_q^{-1}\)로 주어지며, 그 매개변수 또는 임의함수 \(r\)는 \(p,q\)로부터 정해진다. 공식으로 쓰면
\[
\mathfrak T_p:\quad
\xi=x+\varDelta x(x,p),\qquad
u^*=u+\varDelta u(x,u,p);
\]
\[
\mathfrak T_q:\quad
y=A(x,q);\qquad v=B(x,u,q);
\]
\[
\mathfrak T_q\mathfrak T_p:\quad
\eta=A(x+\varDelta x(x,p),q),\qquad
v^*=B(x+\varDelta x(p),u+\varDelta u(p),q).
\]
여기서 \(\mathfrak T_r=\mathfrak T_q\mathfrak T_p\mathfrak T_q^{-1}\), 즉
\[
\eta=y+\varDelta y(r),\qquad v^*=v+\varDelta v(r)
\]
를 얻는다. 이때 \(\mathfrak T_q\)의 역변환에 따라 \(x\)를 \(y\)의 함수로 보고 무한소항들만 고려한다. 따라서 다음 항등식이 성립한다.
\begin{equation}
\tag{20}
\begin{gathered}
\eta=y+\varDelta y(r)
=y+\sum\frac{\partial A(x,q)}{\partial x}\varDelta x(p),\\[0.55em]
v^*=v+\varDelta v(r)
=v+\sum\frac{\partial B(x,u,q)}{\partial x}\varDelta x(p)
+\sum\frac{\partial B(x,u,q)}{\partial u}\varDelta u(p).
\end{gathered}
\end{equation}
여기서 \(\xi=x+\varDelta x\)를 \(\xi-\varDelta \xi\)로 바꾸면 \(\xi\)는 다시 \(x\)가 되고 \(\varDelta x\)는 사라진다. 따라서 (20)의 첫째 공식에 따라 \(\eta\)도 다시 \(y=\eta-\varDelta\eta\)가 된다. 이 치환으로 \(\varDelta u(p)\)가 \(\bar\delta u(p)\)로 바뀌므로 \(\varDelta v(r)\)도 \(\bar\delta v(r)\)로 바뀌고, (20)의 둘째 공식은
\[
v+\bar\delta v(y,v,\ldots,r)
=v+\sum\frac{\partial B(x,u,q)}{\partial u}\,\bar\delta u(p),
\]
\[
\bar\delta v(y,v,\ldots,r)
=\sum\frac{\partial B}{\partial u_\varkappa}
\,\bar\delta u_\varkappa(x,u,p)
\]
를 준다. 따라서 \(\bar\delta v\)가 매개변수 또는 임의함수 \(r\)에 의존한다고 가정하기만 하면 실제로 변분의 변환공식들이 성립한다\footnote{이 추론이 성립하려면 \(y\)가 \(u\)와 무관해야 한다는 것 등이 다시 드러난다. 한 예로 클라인이 제시한 \(\delta g^{\mu\nu}\)와 \(\delta q_\varrho\)를 들 수 있다. \(p\)에 벡터변환을 가하면 이들은 변분의 변환법칙을 만족한다.}.

특히 \(\sum_i\psi_i\bar\delta u_i\)가 상대불변임이 따라온다. 따라서 발산관계들이 \(y,v\)에서도 성립하므로 (12)에 따라 \(\operatorname{Div} B\)도 상대불변이다. 나아가 (14)와 (13)에 따라 \(\operatorname{Div}\varGamma\)와, 종속관계들의 왼쪽을 \(p^{(\lambda)}\)들과 묶은 식도 상대불변이다. 변환된 공식들에서는 언제나 임의의 \(p(x)\), 또는 매개변수들을 \(r\)로 바꾸어야 한다. 여기서 \(\operatorname{Div}(B-\varGamma)\)의 상대불변성도 따라온다. 이는 항등적으로 사라지지 않는 함수계 \(B-\varGamma\)의 발산식이지만, 그 발산 자체는 항등적으로 사라진다.

일차원이고 유한군인 경우에는 \(\operatorname{Div}B\)의 상대불변성으로부터 제1적분의 불변성도 추론할 수 있다. (20)에 따라 무한소변환에 대응하는 매개변수변환은 선형이고 동차이다. 또한 모든 변환이 가역이므로 \(\varepsilon\)들도 변환된 매개변수 \(\varepsilon^*\)의 선형 동차함수이다. (20)에는 \(u\)의 도함수가 나타나지 않으므로 \(\psi=0\)으로 두어도 이 가역성은 확실히 유지된다. 식
\[
\operatorname{Div}B(x,u,\ldots,\varepsilon)
=\frac{\dd y}{\dd x}\,
\operatorname{Div}B(y,v,\ldots,\varepsilon^*)
\]
에서 \(\varepsilon^*\)의 계수들을 비교하면, \(\frac{d}{d y}B^{(\lambda)}(y,v,\ldots)\)들도 \(\frac{d}{d x}B^{(\lambda)}(x,u,\ldots)\)들의 선형 동차함수가 된다. 따라서
\[
\frac{\dd}{\dd x}B^{(\lambda)}(x,u,\ldots)=0
\quad\text{또는}\quad
B^{(\lambda)}(x,u)=\mathrm{const.}
\]
이면
\[
\frac{\dd}{\dd y}B^{(\lambda)}(y,v,\ldots)=0
\quad\text{또는}\quad
B^{(\lambda)}(y,v)=\mathrm{const.}
\]
도 따라온다. 따라서 \(\mathfrak G_\rho\)에 대응하는 \(\rho\)개의 제1적분은 각각 이 군을 허용하므로, 이어지는 적분도 단순해진다. 이에 대한 가장 단순한 예는 \(f\)가 \(x\), 또는 한 \(u\)와 무관한 경우이며, 이는 각각 무한소변환 \(\Delta x=\varepsilon,\Delta u=0\)과 \(\Delta x=0,\Delta u=\varepsilon\)에 대응한다. 그러면 \(\bar\delta u\)는 각각 \(-\varepsilon\,\dd u/\dd x\)와 \(\varepsilon\)가 된다. \(B\)는 \(f\)와 \(\bar\delta u\)로부터 미분과 유리결합을 거쳐 유도되므로, \(B\)도 각각 \(x\)와 \(u\)에 무관하고 따라서 대응하는 군들을 허용한다\footnote{\(\int(\sum\psi_i\delta u_i)\dd x\)의 불변성만으로 이미 제1적분들의 존재가 따라오는 경우, 이 제1적분들은 완전한 군 \(\mathfrak G_\rho\)를 허용하지 않는다. 예를 들어 \(\int(u''\delta u)\dd x\)는 무한소변환 \(\Delta x=\varepsilon_2,\Delta u=\varepsilon_1+x\varepsilon_3\)을 허용한다. 그러나 \(\Delta x=0,\Delta u=x\varepsilon_3\)에 대응하는 제1적분 \(u-u'x=\mathrm{const.}\)는 \(u\)와 \(x\)를 모두 명시적으로 포함하므로 다른 두 무한소변환을 허용하지 않는다. 이 제1적분에 대응하는 \(f\)의 무한소변환들이 도함수들을 포함하는 것이다. 따라서 \(\int\!\cdots\!\int(\sum\psi_i\delta u_i)\dd x\)의 불변성이 \(I\)의 불변성보다 아무튼 더 적은 것을 준다는 사실을 알 수 있다. 이는 앞의 주석에서 제기한 문제와 관련하여 덧붙일 점이다.}.

\bigskip
\section*{\S{}6. 힐베르트의 한 주장}

앞의 결과들로부터 끝으로 ``일반상대성''과 고유한 에너지법칙의 실패 사이의 연관에 관한 힐베르트의 주장(첫째 클라인 논문, \emph{Göttinger Nachr.} 1917, 답변 제1문단)을 일반화된 군론적 형태로 증명할 수 있다.

적분 \(I\)가 \(\mathfrak G_{\infty\varrho}\)를 허용하고, \(\mathfrak G_\sigma\)는 임의함수들을 특수화하여 생기는 임의의 유한군, 즉 \(\mathfrak G_{\infty\varrho}\)의 부분군이라고 하자. 그러면 무한군 \(\mathfrak G_{\infty\varrho}\)에는 종속관계 (16)이, 유한군 \(\mathfrak G_\sigma\)에는 발산관계 (13)이 대응한다. 거꾸로 어떤 발산관계들이 성립하면 한 유한군에 대한 \(I\)의 불변성이 따라오며, 그 군은 \(\bar\delta u\)들이 \(\mathfrak G_\sigma\)로부터 생기는 것들의 선형결합일 때 그리고 그때에만 \(\mathfrak G_\sigma\)와 같다. 따라서 \(\mathfrak G_\sigma\)에 대한 불변성은 (13)과 다른 발산관계들을 낳을 수 없다. 한편 (16)이 성립하면 임의의 \(p(x)\)에 대하여 \(\mathfrak G_{\infty\varrho}\)의 무한소변환 \(\varDelta u,\varDelta x\)에 대한 \(I\)의 불변성이 따라온다. 특히 그로부터 특수화하여 얻는 \(\mathfrak G_\sigma\)의 무한소변환들, 따라서 \(\mathfrak G_\sigma\)에 대한 불변성도 이미 따라온다. 그러므로 발산관계
\[
\sum_i\psi_i\bar\delta u_i^{(\lambda)}
=\operatorname{Div}B^{(\lambda)}
\]
는 종속관계 (16)의 귀결이어야 한다. 후자는 다음과 같이 쓸 수도 있다.
\[
\sum_i\psi_i a_i^{(\lambda)}
=\operatorname{Div}\chi^{(\lambda)},
\]
여기서 \(\chi^{(\lambda)}\)는 라그랑주 식들과 그 도함수들의 선형결합이다. \(\psi\)는 (13)과 (16) 양쪽에서 선형으로 들어가므로, 발산관계들은 특히 종속관계 (16)의 선형결합이어야 한다. 따라서
\[
\operatorname{Div}B^{(\lambda)}
=\operatorname{Div}\left(\sum\alpha\,\chi^{(\varkappa)}\right)
\]
가 된다. 그러므로 \(B^{(\lambda)}\) 자체는 \(\chi\), 즉 라그랑주 식들과 그 도함수들의 선형결합, 그리고 발산이 항등적으로 사라지는 함수들로 이루어진다. 후자의 예로 제2절 끝의 \(B-\varGamma\)가 있는데, 여기서는 \(\operatorname{Div}(B-\varGamma)=0\)이고 그 발산은 동시에 불변식의 성질을 갖는다. \(B^{(\lambda)}\)들을 말한 방식으로 라그랑주 식들과 그 도함수들로 구성할 수 있는 발산관계를 ``비고유'', 나머지를 모두 ``고유''라고 부르겠다.

거꾸로 발산관계들이 종속관계 (16)의 선형결합, 즉 ``비고유''라면 \(\mathfrak G_\sigma\)에 대한 불변성은 \(\mathfrak G_{\infty\varrho}\)에 대한 불변성에서 따라온다. 따라서 \(\mathfrak G_\sigma\)는 \(\mathfrak G_{\infty\varrho}\)의 부분군이 된다. 그러므로 유한군 \(\mathfrak G_\sigma\)에 대응하는 발산관계들은 \(\mathfrak G_\sigma\)가 \(I\)를 불변으로 두는 어떤 무한군의 부분군일 때 그리고 그때에만 비고유이다.

군들을 특수화하면 여기서 힐베르트의 원래 주장을 얻는다. ``평행이동군''은 다음 유한군을 뜻한다고 하자.
\[
y_i=x_i+\varepsilon_i,\qquad v_i(y)=u_i(x),
\]
따라서
\[
\varDelta x_i=\varepsilon_i,\qquad
\varDelta u_i=0,\qquad
\bar\delta u_i=-\sum_\lambda
\frac{\partial u_i}{\partial x_\lambda}\varepsilon_\lambda.
\]
잘 알려진 바와 같이 평행이동군에 대한 불변성은
\[
I=\int\!\cdots\!\int
f\!\left(x,u,\frac{\partial u}{\partial x},\ldots\right)\dd x
\]
에서 \(f\)에 \(x\)가 명시적으로 나타나지 않음을 뜻한다. 이에 대응하는 \(n\)개의 발산관계
\[
\sum_i\psi_i\frac{\partial u_i}{\partial x_\lambda}
=\operatorname{Div}B^{(\lambda)},
\qquad(\lambda=1,2,\ldots n)
\]
를 ``에너지 관계''라고 부르자. 변분문제에 대응하는 ``보존법칙'' \(\operatorname{Div}B^{(\lambda)}=0\)은 ``에너지법칙''에, \(B^{(\lambda)}\)는 ``에너지 성분''에 대응하기 때문이다. 그러면 다음이 성립한다. \(I\)가 평행이동군을 허용할 때, 에너지 관계들은 \(I\)가 평행이동군을 부분군으로 포함하는 한 무한군에 대하여 불변일 때 그리고 그때에만 비고유이다\footnote{고전역학의 에너지법칙과 옛 ``상대성이론''의 에너지법칙(여기서는 \(\sum d x^2\)가 자기 자신으로 옮겨진다)은 무한군이 나타나지 않으므로 ``고유''이다.}.

그러한 무한군의 한 예는 모든 \(x\)의 변환과 그로써 유도되는 \(u(x)\)의 변환으로 이루어진 군이다. 이때 유도된 변환에는 임의함수 \(p(x)\)의 도함수들만 나타나며, \(p^{(i)}(x)=\varepsilon_i\)로 특수화하면 평행이동군을 얻는다. 그러나 이것들과, \(I\)를 경계적분만큼 바꾸어 생기는 군들로 이미 가장 일반적인 모든 군을 얻는지는 미결로 남겨야 한다. 말한 종류의 유도된 변환들은, 예컨대 \(u\)에 ``전미분형식''의 계수변환을 가하면 생긴다. 여기서 전미분형식이란
\[
\sum a\,\dd^\lambda x_i
+\sum b\,\dd^{\lambda-1}x_i\,\dd x_\varkappa+\cdots
\]
처럼 \(d x\) 외에 고차 미분들도 포함하는 형식이다. \(p(x)\)가 1차 도함수로만 나타나는 더 특수한 유도변환들은 보통의 미분형식 \(\sum c\,d x_{i_1}\ldots d x_{i_\lambda}\)의 계수변환들로 주어지며, 보통은 이것들만 고찰해 왔다.

말한 종류의 또 다른 군으로 다음을 들 수 있다. 로그항이 나타나므로 이것은 계수변환일 수 없다.
\[
y=x+p(x),\qquad
v_i=u_i+\lg(1+p'(x))
=u_i+\lg\frac{d y}{d x},
\]
\[
\varDelta x=p(x),\qquad
\varDelta u_i=p'(x)\footnote{이 무한소변환들로부터 유한변환들을 역으로 계산하는 데에는 제4절 끝에서 제시한 방법을 쓴다.};\qquad
\bar\delta u_i=p'(x)-u_i'p(x).
\]
여기서 종속관계 (16)은
\[
\sum_i\left(\psi_i u_i'+\frac{d\psi_i}{d x}\right)=0
\]
이 되고, 비고유 에너지 관계는
\[
\sum\left(\psi_i u_i'+
\frac{d(\psi_i+\mathrm{const.})}{d x}\right)=0
\]
이 된다. 이 군의 가장 단순한 불변적분 하나는
\[
I=\int \frac{e^{-2u_1}}{u_1'-u_2'}\,d x
\]
이다. 가장 일반적인 \(I\)는 리의 미분방정식 (11)
\[
\bar\delta f+\frac{\dd}{\dd x}(f\varDelta x)=0
\]
을 적분하여 정해진다. \(f\)가 \(u\)의 1차 도함수들에만 의존한다고 가정하고 \(\varDelta x\)와 \(\bar\delta u\)의 값을 대입하면 이 방정식은
\[
\frac{\partial f}{\partial x}p(x)
+\left\{\sum_i\frac{\partial f}{\partial u_i}
-\frac{\partial f}{\partial u_i'}u_i'+f\right\}p'(x)
+\left\{\sum_i\frac{\partial f}{\partial u_i'}\right\}p''(x)=0
\]
으로 바뀐다. 이는 \(p(x),p'(x),p''(x)\)에 관한 항등식이다. 이 방정식계는 두 함수 \(u(x)\)만으로도 도함수를 실제로 포함하는 해
\[
f=(u_1'-u_2')\,
\varPhi\!\left(u_1-u_2,\frac{e^{-u_1}}{u_1'-u_2'}\right)
\]
를 갖는다. 여기서 \(\varPhi\)는 표시된 변수들의 임의함수이다.

힐베르트는 고유한 에너지법칙의 실패가 ``일반상대성이론''의 특징이라는 형태로 자신의 주장을 진술한다. 이 주장이 문자 그대로 성립하려면 ``일반상대성''이라는 명칭을 통상적인 것보다 넓게 이해하여, 앞에서 다룬 \(n\)개의 임의함수에 의존하는 군들까지 포함해야 한다\footnote{이로써 물리학에서 통용되는 ``상대성''이라는 명칭을 ``한 군에 관한 불변성''으로 바꾸어야 한다는 클라인의 논평이 옳다는 사실도 다시 확인된다. F. Klein, ``Über die geometrischen Grundlagen der Lorentzgruppe'', \emph{Jhrber. d. d. Math. Vereinig.} 19 (1910), 287쪽; \emph{Physikalische Zeitschrift}에도 재수록.}.

\clearpage
\setcounter{footnote}{0}
\phantomsection
\addcontentsline{toc}{section}{14. 한 변수 대수함수의 산술이론}
%
%
%
%
%
%

\section*{14. 한 변수 대수함수의 산술이론과 다른 이론들 및 수체론과의 관계}

\begin{center}
괴팅겐의 에미 뇌터 보고\\[0.6em]
\emph{Jahresber. d. D. Math.-Verein. 28 (1919), 182--203쪽}
\end{center}

이 보고서는 대수함수에 관한 개별 이론들을 서로 이어 주는 고리를 마련하려는 뜻에서 생겼다. 동시에 산술이론에 대한 논의가 본질적으로 빠져 있는 브릴--뇌터의 보고\srcfn{1)}{\emph{Die Entwicklung der Theorie der algebraischen Funktionen in älterer und neuerer Zeit}, \emph{Jahresb. d. D. Math. Verein.} 3 (1894) [이하 브릴--뇌터].}를 보완하는 글로 볼 수 있다\srcfn{2)}{크로네커의 판별식 연구에 대한 논의는 예외이다(\emph{J. f. M.} 91).}.

브릴--뇌터 외에도 각 이론에 대해서는 다음 백과사전 논문들과 거기에 실린 문헌을 참고할 수 있다.

리만 이론: Wirtinger (II B 2), 1--19번;

바이어슈트라스 이론: Wirtinger (II B 2), 23, 32, 33, 34번;

브릴--뇌터의 대수기하학적 이론: Wirtinger (II B 2), 21--31번; Berzolari (III C 4), 23--38번;

데데킨트--베버 및 헨젤--란츠베르크의 산술이론: Hensel (II C 5) [이하 헨젤]\srcfn{3)}{데데킨트--베버에 관한 부분은 A. Ostrowski가 썼다.};

더 일반적으로 대수적 원소의 산술이론: Landsberg (I C 5), 여러 변수 대수함수의 산술이론: Jung (II C 6).

수체론에 대해서는 힐베르트의 \glqq{}수론 보고\grqq{}\srcfn{4)}{\emph{Die Theorie der algebraischen Zahlkörper}, \emph{Jahresb. d. D. Math. Verein.} 4 (1894/95).}와 디리클레--데데킨트의 수론 강의 [이하 데데킨트--수론]\srcfn{5)}{Braunschweig, Vieweg, 제4판, 1894.}를 참조하라.

\begin{center}
\textbf{목차}
\end{center}
\begin{enumerate}
\item 여러 이론의 기초.
\item 대수적 수와 대수함수의 평행성(켤레원소, 기저정리, 모듈 개념).
\item 대수적 수와 대수함수의 아이디얼론에서 나타나는 평행성과 차이.
\item 급수전개의 산술적 기초와 아이디얼론의 관계.
\item 대수함수의 산술이론과 기하학적 이론의 비교. 서로 다른 불변성 개념.
\item 비교의 계속. 특이점.
\item 잔여정리와 아이디얼론 정리들의 비교.
\item 대수함수와 대수적 수에서의 초월적 문제.
\end{enumerate}

\begin{center}
\textbf{1. 여러 이론의 기초}
\end{center}

리만의 이론은 \emph{순전히 초월적}이며 \emph{존재정리}(디리클레 원리, 슈바르츠--노이만 방법)에 바탕을 둔 유일한 이론이다. 여기서는 대수함수의 \emph{적분}이 일차적이고 대수함수 자체가 이차적이다. 리만곡면 위의 함수는 영점과 극점으로 특징지어지는데, 이는 산술이론의 폴리곤(또는 인자)에 대응한다. 한편 바이어슈트라스도 함수 자체를 그 영점 및 극점과 함께 다룬다. 존재정리 외의 주된 문제는 주어진 극점들을 갖는 모든 함수를 세우는 것이며, 이는 제1종 및 제2종 적분을 통해 해결된다. 더 나아가 극점들을 지정한 뒤에도 들어가는 임의상수의 수를 묻는다. 이 마지막 문제에는 리만--로흐 정리가 답한다. 산술이론에서는 폴리곤류(또는 인자류)의 차원을 묻는 문제가, 대수기하학적 이론에서는 완전 선형계의 차원을 묻는 문제가 이에 대응하며, 여기서도 리만--로흐 정리가 답을 준다. 함수의 실제 구성은 산술이론에서는 기저정리로, 대수기하학적 이론에서는 잔여정리로 이루어진다. 특히 제1종 적분의 피적분함수의 영점들에는 산술이론의 미분류가 대응하고, 대수기하학적 이론의 특수점군 이론, 즉 \((n-3)\)차 수반곡선인 \(\varphi\)-곡선이 잘라 내는 점군들의 선형계가 대응한다.

리만 이론은 역사적으로 다른 이론들보다 앞섰지만 존재정리들의 엄밀한 증명은 뒤에 이루어졌다. 마찬가지로 훗날에도 그 초월적 도구로, 순수 대수적 또는 산술적 방법으로는 오늘날까지 접근할 수 없는 대수적 문제들을 해결하였다. 여기서는 무엇보다 특이대응에 관한 후르비츠의 이론을 들어야 한다\srcfn{1)}{A. Hurwitz, \emph{Über algebraische Korrespondenzen und das erweiterte Korrespondenzprinzip}, \emph{Ann.} 28. 브릴--뇌터 제10절 참조.}.

대수적 수체론의 초월적 문제들과의 유추는 마지막 절에서 논의한다.

나머지 이론들에서는 \emph{대수함수}가 일차적이고 적분이 이차적이다. \emph{바이어슈트라스 이론}은 거듭제곱급수 전개와 유수정리라는 초월적 기초 위에 서지만, 그 뒤에는 고려되는 모든 함수를 명시적으로 제시하면서 대수적으로 전개된다.

헨젤--란츠베르크의 형태를 취한 \emph{산술이론}도 \emph{급수전개라는 초월적 기초}를 쓴다. 그러나 이 기초 위에서 각 점에 인자를 대응시킨 뒤에는 순수 산술적으로 진행하므로, 바이어슈트라스와 데데킨트--베버의 융합으로 볼 수 있다. 후자의 이론과 달리, 고려되는 모든 기저함수를 실제로 구성하는 방법도 함께 준다. 세부적으로는 특히 \glqq{}분수\grqq{} 인자를 도입함으로써 데데킨트--베버보다 여러 점에서 단순해진다. 최근 헨젤\srcfn{1)}{\emph{Mathematische Zeitschrift} 4; 그의 보고도 이 기초에 바탕을 둔다.}은 수렴증명을 위한 우세급수법만 제외하고 급수전개 이론에 산술적 기초를 주었으며, 이로써 이 이론은 거의 순수 산술적이 되었다.

브릴--뇌터의 \emph{대수기하학적 이론}은 한 가지의 \glqq{}연속하는 점들\grqq{}이라는 개념만을 전제로 한다. 이는 \emph{보통점}에서의 급수전개로 환원되고\srcfn{2)}{연속하는 \glqq{}특이\grqq{}점에 대해서는 브릴--뇌터 제6절, 특히 10, 11번을 참조하라. 보통의 연속점은 이를테면 접선과 곡선의 교점 또는 고차 접촉점(분기점)이다.} 바이어슈트라스의 원소 개념에 대응한다. 그 뒤에는 대수적 유리방법, 본질적으로 삼원 소거론의 방법을 써서 진행하며, 이 이론도 명시적 표현에까지 이른다.

데데킨트--베버의 원래 \emph{순수 산술이론}\srcfn{3)}{\emph{J. f. M.} 92 [이하 D.--W.].}은 그 정리들이 다시 존재증명의 성격을 띤다는 점에서 리만 이론과 가깝고, 방법의 완전한 순수성을 지킨다\srcfn{4)}{그 때문에 이후의 비교에서는 가끔만 나타나는 헨젤--란츠베르크 이론보다 다른 이론들과 비교하기에 더 알맞다. 다만 아벨 적분과 그 역에 관한 이론은 연속성을 전제하는 한 D.--W.에서도, 대수이론의 후속 발전(M. Noether, \emph{Ann.} 37)에서도 다루어지지 않았다.}. 변수들이 단지 부정원의 역할만 하는 이 산술적 기초는 제2절에서 더 자세히 밝히듯 \emph{대수적 수}의 이론과 평행하게 진행된다\srcfn{1)}{헨젤에서는 헨젤--란츠베르크 이론과 수론, 특히 \(p\)-진수론의 비교를 볼 수 있다.}. 이 기초는 점 개념을 산술적으로 도입할 수 있게 하며(헨젤 10a 참조), 그 결과 연속성에 대한 어떤 고찰도 없이 절대 리만곡면의 개념을 세울 수 있게 한다. 이로써 이론은 형식적인 성격을 얻는다. 예를 들어 미분의 개념도 순전히 형식적으로 도입된다. 점 개념으로 얻은 기초 위에서는, 원소 개념을 제외하면 역시 초월적 도구를 쓰지 않는 대수기하학적 이론의 방법 및 개념과 가까운 친족관계를 밝힐 수 있다. 이를 뒤의 절들에서 수행하겠다.

\begin{center}
\textbf{2. 대수적 수와 대수함수의 평행성\newline
(켤레원소, 기저정리, 모듈 개념)}
\end{center}

대수적 수체는 유리수체 \(R\) 위의 \(n\)차 체로 정의되고, 한 변수 대수함수체는 임의의 복소계수를 갖는 한 부정원 \(z\)의 모든 유리함수체 \(K(z)\) 위의 \(n\)차 체로 정의된다. 따라서 기초체의 특수한 성질을 완전히 도외시하는 모든 정리, 즉 기저, 노름, 대각합, 판별식에 관한 정리들은 두 체에 공통으로 성립한다\srcfn{2)}{D.--W., \S{}1, 2; 헨젤 4번(주 5 포함).}. 데데킨트--베버에서는 이 개념들을 모두 켤레원소 없이 도입한다. 그러나 슈타이니츠가 크로네커를 이어 보였듯, \emph{켤레체의 개념은 완전히 형식적으로 정의할 수 있다}.\srcfn{3)}{Steinitz, \emph{Algebraische Theorie der Körper}. (\emph{J. f. M.} 137, \S{}6, 8.)} 실제로 \(K\)가 임의의 체이고 \(\varphi(x)\)가 \(K\)에서 기약인 \(x\)의 다항식이면, \(\varphi(x)\)를 법으로 한 잉여류계는 체를 이룬다. \(x\)의 잉여류에 원소 \(j\)를 대응시키면, 동형대응 아래 다항식 \(f(x)\)가 대표하는 잉여류에는 원소 \(f(j)\)가 대응한다. \(\varphi(x)\)로 나누어지는 모든 다항식은 영류를 대표하므로 \(\varphi(j)\)는 영류에 대응한다. 따라서 \(j\)를 기호적으로 \(\varphi(x)=0\)의 근으로 볼 수 있고, 이것이 대수체 \(K(j)\)를 생성한다. 이 절차를 반복하면 \(n\)개의 켤레체 \(K(j),K(j'),\ldots,K(j^{(n-1)})\)를 얻는다. 그러므로 특히 대수함수에서도 \emph{순수 산술이론을 벗어나지 않고} 켤레원소를 말할 수 있다.

두 기초체 \(R\)과 \(K(z)\)에는 모든 \emph{정수적} 원소의 계가 정의되어 있다는 공통점이 있다. 각각 유리정수와 한 부정원의 다항함수이다. 이 정수적 원소들은 임의의 두 원소 \(a,b\)가 \emph{최대공약원} \(d\)를 가지며
\[
d=ma+nb
\]
로 쓸 수 있다는 특징을 갖는다. 여기서 \(d,m,n\)도 각각 \(R\) 또는 \(K(z)\)의 정수적 원소이고, \(d\)는 이 형태로 나타낼 수 있는 절댓값이 가장 작은 수 또는 차수가 가장 낮은 다항식이다. 대수체 또는 상체의 정수적 원소에 관한 정리들은 오직 이 성질에만 바탕을 두며, 이 성질은 유리정수 또는 한 부정원의 다항식에 대한 \emph{소인수분해의 유일성}을 낳는다. 최대공약원의 존재를 보이는 논증은 대수체의 정수적 원소들에 대한 \emph{기본계}, 즉 \emph{기저}의 존재도 준다. 실제로 \emph{정수적} 원소들로 이루어진 체의 기저들만 생각하면 그 판별식은 유리정수 또는 한 부정원의 다항식이 된다. 절댓값이 가장 작은 수 또는 차수가 가장 낮은 함수를 주는 모든 기저가 기본계를 이룬다.

더 일반적인 체까지 포함하여 기저 문제를 더 정확히 이해하려면 \emph{모듈 개념}을 도입한다. 문헌에서 원소의 성질에 따라 서로 다르게 나타나는 정의들을 포괄하도록 다음처럼 정하겠다. \emph{기초정역 \(\mathfrak J\)에 관한 모듈}이란, \emph{두 원소의 차가 다시 그 계에 속하고 한 원소와 \(\mathfrak J\)의 임의의 원소의 곱도 그 계에 속하는 원소계}를 말한다. \(\mathfrak J\)가 유리정수 전체이면 둘째 조건은 첫째 조건에서 따라온다\srcfn{1)}{모듈 개념은 데데킨트에게서 비롯되었다. 그는 『수론』 제2판에서 수모듈(\(\mathfrak J=\) 유리정수)을 처음 도입했으며, 정수계수 선형형식의 모듈도 거기에 암묵적으로 포함되어 있다. 데데킨트--베버에서는 원소가 대수함수이고 \(\mathfrak J\)가 한 부정원의 모든 다항식으로 이루어진 \glqq{}함수모듈\grqq{}이 나타난다. \(n\)개 부정원의 동차형식, 더 일반적으로 다항식의 모듈은 힐베르트(\emph{Annalen} 36)에게서 처음 나타난다. 여기서는 \(\mathfrak J\)가 모든 다항식으로 이루어지고 원소들 역시 다항식이므로 이 모듈은 우리의 아이디얼 정의에 들어간다. 크로네커의 다항식 \glqq{}모듈계\grqq{}는 추상적 정의가 아니라 기저에서 출발한다.}.

유한개의 원소로 이루어진 기저를 가져서 모든 원소를 \(\mathfrak J\)의 계수로 선형표현할 수 있는 모듈을 \emph{유한}(유한생성) 모듈이라 한다. 특히 일생성 모듈은 \(c\cdot\alpha\) 꼴이며, 여기서 \(\alpha\)는 기저원소이고 \(c\)는 \(\mathfrak J\)의 모든 원소를 돈다. \emph{원소 자체가 \(\mathfrak J\)에 속하는 모듈을 \(\mathfrak J\)의 아이디얼이라 한다.} 여기서 유한 아이디얼, 특히 일생성 아이디얼인 주아이디얼의 개념이 따라온다.

정수적 원소와 아이디얼에 관한 모든 기저정리는 다음 정리에 바탕을 둔다.

\emph{\(M\)이 기초정역 \(\mathfrak J\)에 관한 모듈이고, 그 원소들이 \(\mathfrak J\)의 계수를 갖는 \(n\)개 부정원의 선형형식들이라고 하자. \(\mathfrak J\)의 모든 아이디얼이 유한이면 \(M\)도 유한모듈이다. 특히 \(\mathfrak J\)의 모든 아이디얼이 주아이디얼이면 \(M\)은 선형독립인 형식들로 이루어진 기저를 갖는다.}\srcfn{1)}{이 정리는 여러 계수정역 \(\mathfrak J\)에 대하여 문헌에 따로 나타난다. 유리정수에 대해서는 힐베르트, \emph{Zahlbericht} \S{}3, 4를, 여러 부정원의 다항식에 대해서는 힐베르트, \emph{Über die vollen Invariantensysteme}, \S{}2 끝(\emph{Math. Ann.} 42)을 참조하라.}

실제로 \(M\)의 원소들이
\[
l(x)=a_1x_1+\cdots+a_nx_n
\]
으로 주어지면, 계수 \(a_n\) 전체는 \(\mathfrak J\)의 한 아이디얼을 돈다. 가정에 따라 이는 \(b_1a_n^{(1)}+\cdots+b_\varrho a_n^{(\varrho)}\) 꼴이다. 따라서 \(M\)에 속하는 선형형식
\[
m(x)=l(x)-b_1l^{(1)}(x)-\cdots-b_\varrho l^{(\varrho)}(x)
\]
은 \((n-1)\)개의 부정원에만 의존한다. 이를 유한번 반복하면 명제가 따라온다. 매 단계에서 \(\varrho=1\)이면 기저는 각각 \(n,\ n-1\ldots1\)개의 부정원에 의존하는 \(n\)개의 형식으로 이루어지고, 따라서 영이 아닌 형식들은 선형독립이다.

이로부터 여러 경우의 기본계 정리를 다음과 같이 얻는다.

대수체가 정수적 대수원소 \(\delta\)를 기초체에 첨가하여 생기고, \(\mathfrak J\)가 기초체의 모든 정수적 원소로 이루어진다고 하자. 그러면 상체의 모든 정수적 원소는 다음 표현을 갖는다(예: \emph{Zahlbericht} \S{}3).
\[
\omega=\frac{a_1\delta^{\,n-1}+a_2\delta^{\,n-2}\cdots+a_n}{\Delta(\delta)},
\]
여기서 \(a_i\)들은 \(\mathfrak J\)의 원소이고 \(\Delta(\delta)\)는 \(\delta\)의 판별식이다. 치환
\[
x_i=\frac{\delta^{n-i}}{\Delta(\delta)}
\]
에 의해 모든 정수적 원소의 모듈은 선형형식의 모듈로 바뀐다. 따라서 \(\mathfrak J\)의 모든 아이디얼이 유한인 모든 경우에 기본계가 존재한다. 같은 논증은 더 일반적으로 상체의 정수적 원소 정역에 있는 모든 아이디얼이 유한임을 보여 준다. 특히 유리정수 또는 한 부정원의 다항식에서는 모든 아이디얼이 주아이디얼이다. 그러므로 대수적 수체 또는 한 변수 대수함수체의 차수가 \(n\)이면 그 기본계는 \emph{정확히 \(n\)개의 원소}로 이루어진다. 또한 앞에서 본 바와 같이 이 체들의 모든 아이디얼이 유한이므로 상대체에도 기본계가 존재하며, 일반적으로 그 원소 수는 상대차수보다 많다\srcfn{1)}{\(\mathfrak J\)가 한 유한 대수적 수체의 정수환일 때 선형형식의 모듈에 관한 슈타이니츠의 연구에 따르면, 상대차수가 \(\nu\)일 때 \(\nu+1\)개의 기저원소면 충분하다(\emph{Math. Ann.} 71, \S{}4).}. 더 나아가 \(\mathfrak J\)가 여러 부정원의 모든 다항식으로 이루어진 정역일 때도 모든 아이디얼은 유한이므로\srcfn{2)}{힐베르트의 모듈기저정리(\emph{Math. Ann.} 36)에 따른다. 여기서 일관성을 위해 아이디얼이라 부른 다항식 모듈은 보통 \glqq{}형식모듈\grqq{} 또는 단순히 \glqq{}모듈\grqq{}이라 한다. 그러한 모듈 \(N\)이 유한이라는 힐베르트의 정리도 선형형식 정리로 환원할 수 있다. 실제로 \(f\)가 \(N\)에 속하는 \(\varrho+1\)개 부정원 \(\xi_1\ldots\xi_\varrho,\eta\)의 \(n\)차 다항식이고 항 \(\eta^n\)을 포함한다고 하자. 이는 선형변환으로 언제나 달성할 수 있다. 그러면 \(N\)의 모든 다항식은 \(f\)를 법으로 하여 \(a_1(\xi)\eta^{n-1}+\ldots a_n(\xi)\) 꼴인 것과 합동이다. 따라서 \(\eta^{n-i}=x_i\)로 두면 선형형식의 모듈을 이룬다. 한 부정원에 대해서는 명제가 성립하므로 귀납적으로 \(\varrho\)개 부정원의 정역에서 모든 아이디얼을 유한이라고 가정할 수 있다.}, 여러 변수 대수함수체에도 기본계가 존재한다. 이 경우에도 그 원소 수는 차수보다 많다.

\begin{center}
\textbf{3. 대수적 수와 대수함수의 아이디얼론에서 나타나는 평행성과 차이}
\end{center}

모든 아이디얼이 소아이디얼들의 거듭제곱곱으로 유일하게 표현된다는 아이디얼론의 기본정리는 대수적 수와 대수함수에 대하여 함께 증명할 수 있다. 유리정수 또는 한 부정원의 다항식에서는 모든 아이디얼이 주아이디얼이라는 성질만을 쓰면 된다\srcfn{3)}{Hurwitz, \emph{Über die Theorie der Ideale}, \emph{Gött. Nachr.} 1894의 서론에 있는 논평을 참조하라. 베버의 대수학 교과서에서는 크로네커를 따라 함수적을 도입하여 두 체에 평행한 증명을 준다. 크로네커 이론과 아이디얼론 사이의 연결은 \glqq{}확장된 가우스 정리\grqq{}가 담당한다(Hurwitz, 앞의 글).}. 증명의 핵심은 아이디얼 \(\mathfrak c\)가 아이디얼 \(\mathfrak a\)로 나누어진다는 것, 즉 \(\mathfrak c\)의 모든 원소가 \(\mathfrak a\)에 들어 있다는 것으로부터 곱표현 \(\mathfrak c=\mathfrak a\cdot\mathfrak b\)가 따라옴을 보이는 데 있다\srcfn{4)}{데데킨트의 『수론』, 보유 XI, 정리 VII, \S{}177에서 특히 강조되어 있다. 554쪽의 주도 참조하라. 이 정리는 대수체를 다룬다는 사실을 쓰며, 다항식의 아이디얼(모듈)에는 더 이상 성립하지 않는다. 그 경우 곱 대신 최소공배수가 나타난다.}. 후르비츠의 증명은 \glqq{}확장된 가우스 정리\grqq{}에 바탕을 둔다. 이 정리에 따르면 두 다항식 \(\varphi,\psi\)의 곱의 모든 계수를 나누는 정수적 원소 \(\omega\)는 \(\varphi\)와 \(\psi\)의 모든 계수의 곱도 나눈다. 데데킨트의 증명은 이 정리와 동치인 모듈정리에 바탕을 둔다\srcfn{1)}{『수론』, 보유 XI, 정리 VI, \S{}173. 두 정리의 동치성은 데데킨트가 지적하였다(\emph{Über die Begründung der Idealtheorie}, \emph{Göttinger Nachr.} 1895).}.

더 나아가 여모듈의 개념(헨젤 12번)과 그로부터 분기아이디얼(기초아이디얼, 디퍼런트) \(\mathfrak z\)을 함께 정의할 수 있다. 그 노름은 체판별식 \(D\)가 된다. 또한 다음 기본정리가 성립한다. \emph{분기아이디얼 \(\mathfrak z\)는 모든 주아이디얼 \(F'(\delta)\)의 최대공약아이디얼이다.} 여기서 \(\delta\)는 체의 모든 정수적 원소를 돌고, \(F(\delta)=0\)은 그때마다 대응하는 기약방정식이다. 여기서 \emph{체판별식은 소아이디얼의 제곱으로 나누어지는 유리소수 또는 \(z\)의 일차인자만을, 그리고 그 모두를 포함한다}는 것이 따라온다. \(F'(\delta)\) 자체에 대해서는
\[
F'(\delta)=\mathfrak f\mathfrak D,\qquad
\Delta(\delta)=R^2D
\]
라는 표현을 얻는다. 여기서 \(\mathfrak f\)는 \(\delta\)로부터 생성되는 환(오더, 정역)의 도체이다. 즉 \(\mathfrak f\)는 임의의 정수적 원소와 곱했을 때 모듈 \([1,\delta,\ldots\delta^{n-1}]\)로 나누어지는 모든 정수적 원소로 이루어진다. \(R^2=\operatorname{Norm}\mathfrak f\)는 기저 \(1,\delta\ldots\delta^{n-1}\)을 정수적 원소들의 기저로 옮기는 치환행렬식의 제곱이다\srcfn{2)}{Dedekind, \emph{Über die Diskriminante endlicher Körper}, \emph{Abhandl. der Gött. Ges. d. Wissensch.} 29 (1882). 거기서 대수적 수에 대하여 준 정의들은 대수함수로 곧바로 옮길 수 있고, 증명들도 거의 그대로 옮길 수 있다. D.--W.의 증명 배열은 이와 다르다. 아래를 보라. 인용한 정리들에 대해서는 \emph{Zahlbericht}, \S{}31, 32도 참조하라. 거기서는 \(F'(\delta)\)를 정수 \(\delta\)의 \glqq{}디퍼런트\grqq{}라 부른다.}.

뒤의 전개에서 생기는 \emph{차이}는 기초체의 특수한 성질들 때문이다. 대수적 수의 경우에는 \emph{기초체의 원소들이 동시에 지수}라는 사실 때문에, \(p\)에서 \(e\)제곱으로 나타나는 소아이디얼 \(\mathfrak p\)의 \((e-1)\)제곱보다 높은 거듭제곱이 분기아이디얼을 나눌 수 있다. 이는 \(e\)가 \(p\)로 나누어질 때 일어난다. 그 결과 관성체와 분기체라는 더 나아간 개념들\srcfn{3)}{\emph{Zahlbericht}, \S{}39--45 참조.}이 생기지만, 대수함수 이론에는 그 유사물이 없다. 무엇보다 대수함수에서는 기초체에 \emph{상수가 무한히 많이 존재}하므로 아이디얼류의 수가 유한하다는 성질이 사라지고, 이에 따라 유수의 결정과 관련된 모든 문제도 사라진다. 다만 어떤 의미에서는 모든 초월적 소형식의 정역\srcfn{1)}{Klein, \emph{Zur Theorie der Abelschen Funktionen}, \emph{Ann.} 36 (1890).}을 유체의 초월적 유사물로 볼 수 있다. 여기서는 각 점이 하나의 형식, 즉 한 주아이디얼로 생성되기 때문이다.

반대로 기초체 \(K(z)\)에 \emph{상수가 무한히 많이 존재}한다는 사실은 대수함수에서 대수적 수에는 유사물이 없는 정리와 문제를 낳는다. 먼저 증명의 배열이 어느 정도 단순해진다. 예를 들어 D.--W.가 아이디얼의 소아이디얼 분해 유일성을 증명할 때, 각 일차형식 \((z-c)\)가 모든 상수라는 무한히 많은 잉여류를 갖는다는 사실을 이용하여 \glqq{}확장된 가우스 정리\grqq{}나 그와 동치인 정리를 피할 수 있다. 반면 소수 \(p\)를 법으로 한 잉여류의 수는 유한하다\srcfn{2)}{D.--W. \S{}9. 212, 213쪽의 주를 참조하라.}.

실제로 새로운 결과를 낳는 것은 다음 사실이다. 계수가 \(z\)의 다항식인 모든 다항식은 \((z-c)\)를 법으로 일차인자들로 분해된다. 정수계수 다항식을 \(p\)를 법으로 볼 때에는 성립하지 않는 성질이다. 이 사실은 모든 상수 대신 모든 대수적 수만을 취해도 여전히 성립한다\srcfn{3)}{이 경우에도 이론 전체가 유지된다는 점을 D.--W.가 서론에서 지적한다.}. 여기서 \emph{모든 소아이디얼은 일차 아이디얼}이라는 것이 따라오고\srcfn{4)}{D.--W., \S{}9, 7번.}, 더 나아가 \emph{체판별식 \(D\)는 모든 판별식 \(\Delta(\delta)\)의 최대공약원}이 된다. 여기서 \(\delta\)는 모든 정수적 원소를 돈다\srcfn{5)}{D.--W., 224쪽.}. 두 사실 모두 대수적 수의 경우와 대조된다. 그 뒤에는 다시 분기아이디얼이 모든 주아이디얼 \(F'(\delta)\)의 최대공약아이디얼임이 따라온다.

또 다른 차이는 대수적 수에서는 유리수체 \(R\)가 \emph{절대적으로 특별한} 기초체, 즉 그 체에 포함된 \emph{소체}(단위원으로 생성되는 체)인 반면, \emph{기초체 \(K(z)\)는 상대적}이라는 데서 생긴다. 상수가 아닌 임의의 함수 \(\xi\)를 독립\emph{변수}로 택할 수 있으며, 그러면 원래 체는 \(K(\xi)\) 위의 대수체가 된다. 모든 소아이디얼 \(\mathfrak p\)가 일차 아이디얼이라는 정리, 즉 모든 함수가 \(\mathfrak p\)를 법으로 한 상수와 합동이라는 정리를 임의의 함수를 독립변수로 택할 수 있다는 사실과 결합하면, 대수적 수의 경우를 넘어서는 가장 중요한 개념인 \emph{점 개념}을 순수 산술적 형태로 얻는다(헨젤 10a 참조). \emph{각 점은 함수체와 동형인 한 상수체로 표현된다}\srcfn{1)}{잘 알려져 있듯 두 체가 \glqq{}동형\grqq{}이라는 것은 합, 차, 곱, 몫도 대응하도록 일대일로 서로 대응시킬 수 있다는 뜻이다.}. 이렇게 정의한 모든 점이 \glqq{}절대 리만곡면\grqq{}을 이룬다. 이 점 정의는 체에만 의존하고 특정 변수에는 의존하지 않는다. 다만 존재증명에서는 어떤 함수 하나를 독립변수 \(z\)로 특별히 택한다. 각 함수에 \(\mathfrak p\)를 법으로 한 상수 잉여를 대응시키면 소아이디얼 \(\mathfrak p\)가 점 \(P\)를 생성한다. 특히 \(\mathfrak p\)로 나누어지는 함수들은 \(P\)에서 사라지고, 그 역도 성립한다. \(\mathfrak p\)가 \(z\)에 관한 모든 소아이디얼을 돌면 \(z\)가 유한인 모든 점을 얻는다. 나머지 점들은 독립변수 \(\zeta=\frac{1}{z}\)를 도입하여 얻으며, 주아이디얼 \((\xi)\)를 나누는 소아이디얼들에 대응한다. 당연히 점 개념에 이어지는 폴리곤, 폴리곤류, 폴리곤류의 차원 등의 개념도 대수적 수에는 유사물이 없다.

독립변수를 자유롭게 택할 수 있다는 사실은 분해 \(F'(\vartheta)=\mathfrak f\mathfrak z\)에 또 다른 해석도 준다. \(\vartheta\)를 독립변수로 보면, \(\vartheta\)로 생성되는 환은 \(z,\vartheta\)에 대하여 대칭적으로 정의되었으므로 이에 대응하여
\[
\left(\frac{\partial F}{\partial z}\right)=\mathfrak f\mathfrak a
\]
를 얻는다. 그리고 \(\mathfrak f\)는 두 주아이디얼 \(\partial F/\partial z\)와 \(\partial F/\partial\vartheta\)의 최대공약아이디얼이 된다. 즉 \emph{도체 \(\mathfrak f\)는 동시에 \(\vartheta,z\)에 관한 \glqq{}이중점 아이디얼\grqq{}이다}\srcfn{2)}{D.--W., 263쪽.}.

\begin{center}
\textbf{4. 급수전개의 산술적 기초와 아이디얼론의 관계}
\end{center}

헨젤--란츠베르크에서는 급수전개라는 초월적 도구와 함수론에서 가져온 점 개념이 아이디얼론을 대신한다. 최근 헨젤이 대수적 수와 대수함수의 급수전개에 부여한 산술적 기초\srcfn{3)}{\emph{Eine neue Theorie der algebraischen Zahlen}, \emph{Math. Zeitschr.} 2 (1918), 그리고 \emph{Neue Begründung der arithmetischen Theorie der algebraischen Funktionen einer Variablen}, 같은 곳 4 (1919). 더 자세한 서술은 헨젤 44번 이하를 참조하라.}는 따라서 아이디얼론과 동등한 것으로 볼 수 있다. 실제로 이 산술적 기초는 각 소아이디얼마다 그 소아이디얼이 주아이디얼이 되는 한 초월적 확대체로 넘어가서 통상적인 분해정리들이 성립하게 하는 데 귀착된다. 이 고찰은 한 소수 \(p\), 또는 일차함수 \(z-a\)를 나누는 모든 소아이디얼에 대하여 수행하면 충분하다.

먼저 \(z-a\), 또는 \(p\)의 정수 거듭제곱에 따른 모든 급수전개로 이루어지고 음의 거듭제곱은 매번 유한개만 허용하는 초월적 확대체를 도입한다. 후자는 \(p\)-진수체이다. 이 체에서 대수적 수체 또는 함수체를 정의하는 기약방정식은 \(\varrho\)개의 기약인자의 곱으로 분해되며, 이는 \(p\) 또는 \(z-a\)가 \(\varrho\)개의 \glqq{}으뜸\grqq{} 아이디얼 \(\mathfrak q_1,\ldots\mathfrak q_\varrho\)의 곱으로 분해되는 데 대응한다. 각 으뜸아이디얼을 다시 소아이디얼의 거듭제곱
\[
\mathfrak q_i=\mathfrak p_i^{\,l_i}
\]
으로 나타내는 데에는, 헨젤의 방법에서 초월적 확대체 위의 또 다른 대수체를 도입하는 것이 대응한다. 대수함수의 경우 이 체는 순수방정식으로 정의할 수 있고, 대수적 수의 경우에는 더 일반적으로 대수적으로 풀 수 있는 방정식으로 정의할 수 있다. 대수함수에서의 단순화는 이 경우 관성체와 분기체가 나타나지 않는다는 사실에 기인한다(제3절 참조).

이 연구들에서 소원소 \(\pi\)에 관하여 나타나는 \glqq{}종결급수\grqq{}\srcfn{1)}{\emph{Neue Begründung \(\ldots\)}, 식 (11).}
\[
v=c_0+c_1\pi+\ldots+c_{\sigma-1}\pi^{\sigma-1}-v_\sigma\pi^\sigma,
\]
여기서 \(c_0\ldots c_{\sigma-1}\)는 상수이고 \(v_\sigma\)는 체의 정수적 원소인데, 이와 정확히 대응하는 식을 D.--W.에서도 볼 수 있다\srcfn{2)}{D.--W., 231쪽 및 \S{}15.}. 거기서는 아이디얼론으로부터 이를 유도하며, 이 식이 한 점에서 함수가 영이 되는 차수 또는 무한대가 되는 차수를 정의하는 기초를 준다.

\begin{center}
\textbf{5. 대수함수의 산술이론과 기하학적 이론의 비교\newline
서로 다른 불변성 개념}
\end{center}

\glqq{}불변\grqq{}이라는 것은 \emph{오직 체 자체로 특징지어지는} 모든 개념을 말한다. 산술이론에서는 일반적으로 정의부터 체의 \emph{모든} 원소에 관하여 진술되고 어떤 기저나 특별히 선택한 변수에 관하여 진술되지 않으므로, 이 불변성이 이미 정의에 들어 있다. 앞에서 제시한 점 개념이 그 예이다. 반면 다른 이론들은 특정한 변수나 기저에서 출발하여 그 뒤에 그 개념이 이 선택과 무관함을 보인다. 여기서 불변성은 \emph{쌍유리변환에 대한 불변성}이 된다.

실제로 한편으로 체가 \(K(z)\)에 대수원소 \(s\)를 첨가하여 생기고, 다른 한편으로 \(K(\zeta)\)에 \(\sigma\)를 첨가하여 생긴다고 하자. \(f(z,s)=0\)과 \(F(\zeta,\sigma)=0\)은 각각 \(s,\sigma\)의 기약방정식이다. 더 일반적으로는 \(K(\eta)\)에 유한개의 원소를 첨가하고 대응하는 기약방정식들을
\[
g_1(\eta,\tau_1)=0,\qquad
g_2(\eta,\tau_1,\tau_2)=0,\ldots
\]
로 둘 수도 있다. 정의에 따라 모든 원소는 \(z,s\)로도, \(\xi,\sigma\)로도, \(\eta,\tau_1,\tau_2\ldots\)로도 유리적으로 표현된다. 특히 \(z,s\)는 \(\xi,\sigma\)로 유리적으로 표현되고 그 역도 성립한다. 마찬가지로 \(z,s\)와 \(\eta,\tau_1,\tau_2\ldots\) 사이에도 양방향의 유리표현이 있다. 따라서 평면곡선 \(f(z,s)=0\)은 \glqq{}쌍유리변환\grqq{}으로 평면곡선 \(F(\zeta,\sigma)=0\)으로 옮겨지며, 또는 변수 \(\eta,\tau_1,\tau_2\ldots\)의 공간에서 위 방정식들이 정의하는 곡선으로 옮겨진다. 불변성은 특정 곡선에 대하여 진술한 정의가 그 곡선에서 쌍유리변환으로 생기는 모든 곡선, 즉 임의 차원의 공간에 있는 모든 곡선에 대하여 유지된다는 방식으로 나타난다. 이들은 리만의 의미에서 한 \glqq{}류\grqq{}를 이룬다.

여기서 \glqq{}점\grqq{}은 \(f(a,b)=0\)을 만족하는 서로 결부된 값의 계 \(z=a,s=b\)로 정의된다. 이 점은 쌍유리변환으로 일반적으로 다시 \(F(\zeta,\sigma)=0\) 위의 한 결부된 값의 계 \(\alpha,\beta\)로 옮겨진다. 그러나 \(a,b\)가 \glqq{}특이\grqq{}점인 특별한 경우에는 \(F(\zeta,\sigma)=0\) 위의 여러 \glqq{}점\grqq{}에 대응할 수 있다. 또한 \(f(z,s)=0\)의 한 \glqq{}특이\grqq{}점이 유한개의 단순점, 즉 서로 떨어져 있거나 이웃한 점들에 대응하는 곡선을 언제나 줄 수 있다. 이것이 특이점 해소이다\srcfn{1)}{예를 들어 브릴--뇌터 제6절을 참조하라. \glqq{}특이\grqq{}점에 관한 자세한 내용은 다음 절에 있다.}. 마찬가지로 점군은 다시 점군으로 옮겨진다. 적절한 변환에서 특이점에 대응하는 단순점들의 수만큼 그 특이점을 세면 점의 수는 불변이다. 모든 곡선은 \glqq{}보통 다중\grqq{}점들만 갖는 곡선으로 쌍유리변환할 수 있으므로\srcfnmark{1)}, 기하학적 이론은 그러한 특별한 곡선들에 대해서만 개념들을 전개한 뒤, 가장 일반적인 특이점들을 낳는 변환까지 포함한 모든 쌍유리변환에 대하여 불변임을 보인다.

특히 초타원인 경우를 제외하면 이 \glqq{}류\grqq{}에서 특별한 곡선으로 \((p-1)\)차원 공간의 이른바 \glqq{}\(\varphi\)들의 정규곡선\grqq{}을 도입할 수 있다. 이 곡선은 특이점을 갖지 않는다. 여기서 \(\varphi\)들은 동차화한 \(f(z,s)=0\)에 대한 \(p\)개의 \glqq{}수반 \((n-3)\)차 곡선\grqq{}이다. \(f(z,s)=0\)을 \(F(\zeta,\sigma)=0\)으로 옮기는 쌍유리변환에는 \(\varphi\)들의 선형변환이 대응하므로, 여기서는 사영적 의미의 불변성을 다루게 된다. 체의 모든 함수를 \(\varphi\)들의 유리함수로 표현하는 것을 기하학적 이론에서는 \glqq{}불변표현\grqq{}이라 한다\srcfn{1)}{예를 들어 브릴--뇌터 제8절을 참조하라. 여러 변수 대수함수에서는 불변성을 선형변환에 대한 불변성으로 환원할 수 있는가라는 문제가 다루어지지 않았다.}.

따라서 \(\varphi\)들의 \emph{선형계}, 또는 이들이 잘라 내는 \glqq{}특수점군\grqq{}의 선형계는 모든 쌍유리변환 아래 자기 자신으로 옮겨지며, 오직 체로만 특징지어진다. 이는 산술이론의 미분류에 대응한다. 산술이론에서는 다시 체의 \emph{모든} 함수를 독립변수로 보았을 때의 성질들로 정의하기 때문에 불변성이 정의에 직접 나타난다.

\begin{center}
\textbf{6. 비교의 계속. 특이점}
\end{center}

점 개념을 서로 다르게 이해한다는 사실과 관련하여, 대수기하학적 이론에서 특별한 어려움을 일으키는 곡선의 \glqq{}특이점\grqq{}은 산술이론의 기초에는 전혀 나타나지 않는다. 따라서 기하학적 이론에서 이론을 세우는 데 필요한 \glqq{}수반함수\grqq{}도 산술이론의 구성에는 필요하지 않고, 더 높은 특수 부분에만 속한다. 반대로 기하학적 이론의 기초에는 분기점이 나타나지 않는다. 실제로 아론홀트 정규형으로 적분을 표현할 때에도 분기점은 사라진다(헨젤 27번, 식 80).

제5절에서 준 특이점의 정의는 다음과 같이 쓸 수 있다. \(f(z,s)=0\), 또는 더 높은 차원 공간의 한 곡선이 \(z=a,s=b\), 또는 \(z=a; s_i=b_i\)에서 \emph{특이점을 갖는다는 것은, 이 값의 계를 절대 리만곡면의 둘 이상의 점}(산술적 의미의 점)\emph{이 취하거나, 하나 이상의 점이 더 높은 차수로 취한다는 뜻이다.} 후자의 경우에는 \(z\)의 분기폴리곤에도 \(s\)의 분기폴리곤에도 속하는 한 점을 다룬다. 이는 기하학적 이론의 \glqq{}연속하는 점들\grqq{}에 대응하며, 그 점은 \(f(z,s)=0\), 또는 더 높은 차원 공간의 곡선의 고차 첨점이 된다.

가정에 따라 체의 임의의 함수 \(\eta\)는 \(z,s\)로 유리적으로 표현된다. 그러므로 동형성 때문에 \(z=a,s=b\)가 되는 절대 리만곡면의 각 점은, \(z,s\)로 나타낸 각 함수 \(\eta\)에 \(z=a,s=b\)에서의 값을 대응시킴으로써 생긴다. 그 점이 \emph{특이}하려면 \(z=a,s=b\)에서 여러 값을 취하거나 어떤 값의 계를 여러 번 취하는 함수 \(\eta\)가 적어도 하나 있어야 한다. \(z,s\)에 의한 유리표현 때문에 이는 그 표현의 \emph{분자와 분모가 \(z=a,s=b\)에서 동시에 사라질 때에만} 가능하다. 여기서 \emph{위 정의는 \(z=a,s=b\)에서 \(f,\partial f/\partial z,\partial f/\partial s\)가 모두 사라진다는 통상적인 정의와 동일하다}는 것이 따라온다. 실제로 이 경우
\[
\eta=\frac{\partial f/\partial z}{\partial f/\partial s}
\]
는 말한 종류의 함수로서 \(z=a,s=b\)에서 여러 값을 갖는다. 반대의 경우에는 어떤 함수가 \(0/0\) 꼴이 되더라도 오직 한 값만을 취한다.

산술적 점 개념에서는 어떤 한 함수에 서로 다른 값의 계가 대응하기만 해도 두 점을 서로 다르다고 보기 때문에 특이점이 나타날 수 없다. 산술적 점 개념의 기초가 되는 아이디얼론에서도, 체의 정수적 원소 \emph{전체}를 언제나 동시에 고려하기 때문에 특이점 개념이 나타나지 않는다. 실제로 \(\delta\)가 \(z\)에 관하여 정수적인 \(F(z,\delta)=0\)의 유한영역 특이점들은 모두, 둘째 정의와 제3절 끝에 따라 \glqq{}이중점 아이디얼\grqq{} \(\mathfrak f\)로 생성된다. 한편 기본정리에 따라 분기아이디얼 \(\mathfrak z\)는 모든 주아이디얼 \(F'(\delta)\)의 최대공약아이디얼이다. 따라서 각 소아이디얼 \(\mathfrak p\)에 대하여 대응하는 \(\mathfrak f\)가 \(\mathfrak p\)로 나누어지지 않도록 \(\delta\)를 택할 수 있다. 또 \(\mathfrak f\)는 \(F'(\delta)\)와 \(F'(z)\)의 최대공약아이디얼이므로, 두 주아이디얼 \(F'(\delta),F'(z)\) 가운데 적어도 하나는 \(\mathfrak p\)로 나누어지지 않는다. 그러므로 \emph{모든 \(F'(\delta)\)와 \(F'(z)\)가 동시에 사라지는 값의 계는 없다}. 따라서 값의 계 \(z=a,\delta_i=b_i\)는 절대 리만곡면의 \emph{한} 점에서만 취해진다. 즉 무한차원 공간의 곡선으로 해석할 수 있는 \emph{체의 모든 정수적 대수함수 전체는 유한영역에서 특이점을 갖지 않는다}. 아이디얼론에서는 유한값만을 다룬다.

모든 정수적 원소의 기저함수 \(\omega_1\ldots\omega_n\) 역시 유한영역에서 특이점을 갖지 않는다. 앞에서 본 바와 같이 모든 \emph{정수적} 함수와 동형으로 대응하는 상수계만으로도 산술적 의미의 각 점은 이미 \emph{유일하게} 정해진다. 따라서 임의의 \emph{정수적} 함수 \(s_i\)들로 정의되는 \glqq{}공간곡선\grqq{}이 \(s_i=\alpha_i\)에서 유한영역의 특이점을 갖는다면, \(s_i=\alpha_i\)에서 여러 값의 계를 취하는 \emph{정수적} 함수 \(\eta\)가 적어도 하나 있어야 한다. 이제 \(\omega_1\ldots\omega_n\)이 모든 정수적 원소의 기저이면
\[
\delta=a_1\omega_1+\cdots+a_n\omega_n
\]
이고 \(a_i\)들이 정수적 유리함수이므로, \(\omega\)의 각 값의 계에는 \(\delta\)의 값의 계 하나만이 대응한다. 따라서 \(\omega_1\ldots\omega_n\)이 정의하는 공간곡선은 유한영역에 특이점을 가질 수 없다. 또한 \emph{\(z\)가 유한인 절대 리만곡면의 각 점은 \(\omega\)들과 동형으로 대응하는 상수계 하나만으로 이미 유일하게 정의된다. 기저함수들은 바로 \(s_i=\alpha_i\)에서 여러 값을 갖게 되는 함수 \(\eta\)들이다.} 그러므로 예컨대 기저 \(1,\delta\ldots\delta^{n-1}\) 대신 정수적 원소들의 기저를 도입하면 특이점의 어려움을 피할 수 있다.

기하학적 이론에서는 절대 리만곡면의 모든 점을 유일하게 생성하는 문제, 더 일반적으로 주어진 점들에서 극을 갖는 모든 함수를 구성하는 문제를 수반형식과 그 몫을 통해 해결한다. 형식 \(\psi\)가 수반이라는 것은 기초곡선 \(f(z,s)=0\)의 각 \(i\)중점에서 \(\psi=0\)이 적어도 \((i-1)\)중점을 갖는다는 뜻이다. 고차 특이점은 먼저 보통 다중점들로 해소해야 한다. \emph{\(f(z,s)=0\)의 한 특이점에서 가장 일반적인 함수 \(\eta=\psi/\chi\)의 분자와 분모가 모두 사라져야 한다}는 사실은 위에서 보았다. 그것들이 수반이어야 함은 전역적으로 유한인 미분들을 고찰하면 드러난다. 거꾸로 수반조건들이 가장 일반적인 함수를 구성하는 데 충분하다는 것은 뒤에서 더 자세히 다룰 \glqq{}잔여정리\grqq{}가 보여 준다. 따라서 수반함수들은 산술이론의 기저함수들을 대신한다.

모든 특이점이 유한영역에 놓인다고 가정하자. 이는 변환으로 언제나 달성할 수 있다. 그러면 산술적으로 한 \emph{수반형식}에는 체에서 \emph{\glqq{}이중점 아이디얼\grqq{} \(\mathfrak f\)로 나누어지는 함수의 분자}가 대응한다. 이로부터 형식적으로 \emph{기저원소 \(\omega_1\ldots\omega_n\)을 수반형식의 몫으로 나타내는 표현}과, 그 결과 서로 다른 표현방식들의 관계를 다음처럼 얻는다.

\(R(z)\)가 기저 \(1,\delta\ldots\delta^{n-1}\)을 기저 \(\omega_1\ldots\omega_n\)으로 옮기는 치환행렬식이라 하자. 거꾸로 각 \(\omega\)는 \(z,\delta\)의 다항함수를 \(R(z)\)로 나눈 몫으로 표현된다.
\[
\omega_i=\frac{G_i(z,\delta)}{R(z)}.
\]
따라서 \(R(z)\omega_i\)는 \(\delta\)로 생성되는 환에 속하므로, 도체의 정의에 따라 \(R(z)\)는 \(\mathfrak f\)로 나누어져야 한다. \((R(z))^2=\operatorname{Norm}\mathfrak f\)에서는 \((R(z))^2\)에 대해서만 이 나눗셈이 따라온다. 한편 \(\omega_i\)는 정수적 함수이므로 유한한 모든 \(z\)에서 유한하게 남는다. 따라서 분자함수 \(G_i(z,\delta)\)도 \(\mathfrak f\)로 나누어져야 한다. 즉 \emph{모든 \(\omega_i\)의 분자와 분모는 수반이다.} \(\omega\)들의 이 표현으로부터 체의 모든 함수를 수반형식의 몫으로 나타내는 표현이 따라온다.

앞의 결과에 따르면 기하학적 이론은 \emph{한 원소 \(\delta\)로 생성되는 환}(오더)의 이론으로 특징지을 수 있다. 이는 \emph{모든 정수적 원소}를 동시에 고찰하는 산술이론과 대조된다. 따라서 환의 도체, 즉 특이점들이 결정적인 역할을 한다는 것이 분명하다. 환(오더)의 산술이론과 기하학적 이론 사이에는 더 많은 유추가 있다. 예를 들어 기하학적 이론은 수반곡선과의 교차에서 특이점으로 들어가는 교점들을 어디에서도 세지 않고, 그와 다른 점군들만을 고찰한다. 이에 대응하여 데데킨트\srcfn{1)}{\emph{Über die Anzahl der Idealklassen in den verschiedenen Ordnungen eines endlichen Körpers}, Braunschweig 기념논문 (1877), \S{}3--5. 이 절들에서 대수적 수에 관하여 진술한 정리들은 한 변수 대수함수에도 문자 그대로 성립한다.}는 도체아이디얼 \(\mathfrak f\)와 서로소인 환의 아이디얼들만을 고찰하고, 이 범위의 아이디얼에 대하여 소아이디얼 분해의 유일성을 포함한 모든 나눗셈법칙을 세운다.

기하학적 이론의 기초가 되는 \glqq{}대수함수의 기본정리\grqq{}\srcfn{2)}{M. Noether, \emph{Über einen Satz aus der Theorie der algebraischen Funktionen}, \emph{Math. Ann.} 6 (1873).}, 또는 적어도 그 직접적인 한 귀결도 어떤 의미에서는 환의 산술이론에서 유사물을 갖는다. 이 정리는
\[
\psi(z,s)=A(z,s)f(z,s)+B(z,s)\varphi(z,s)
\]
로 표현할 수 있기 위한 조건들을 제시한다. 여기서 나타나는 모든 양은 \(s,z\)의 다항식이며, 더 일반적으로 \(x_1x_2x_3\)의 정수적 동차형식이다. 이 조건들에서\srcfn{3)}{브릴--뇌터, 55번.}, \(f(z,s)=0\)을 이용할 때 몫 \(\psi/\varphi\)가 \(f\)에 수반이면 언제나 \(s,z\)의 다항함수 \(B(z,s)\)와 같아진다는 것이 따라온다. 이에 대응하는 것은 이 모든 비교의 기초가 되는 제3절의 정리, 즉 \(\delta\)로 생성되는 환의 도체 \(\mathfrak f\)가 \(f(z,\delta)=0\)의 이중점 아이디얼과 같다는 정리이다. 실제로 도체의 정의에 따라 \(\mathfrak f\)로 나누어지는 모든 정수적 대수원소는 환에 속하므로 \(z,\delta\)의 다항식으로 표현된다. 반면 방금 말한 정리는 그러한 원소가 수반함수라고 말한다\srcfn{4)}{여기서는 데데킨트를 따라 \(\delta\)가 정수적 대수원소라는 가정을 계속 두었다. 이는 변환으로 언제나 달성할 수 있다. 기하학에 더 정확히 대응하여 \(s\)가 \(z\)에 대하여 대수적 분수함수일 수도 있는 \(f(z,s)=0\)의 경우는 헨젤--란츠베르크가 다룬다. 헨젤 26, 29번 참조.}.

\begin{center}
\textbf{7. 잔여정리와 아이디얼론 정리들의 비교}
\end{center}

방금 말한 \glqq{}대수함수의 기본정리\grqq{}의 귀결은 기하학적 이론에서 곧바로 \glqq{}잔여정리\grqq{}로 이어지며, 그 결과 기하학적 이론 전체가 세워지는 기초로 이어진다. 그러나 어떤 면에서 잔여정리는 다른 기초들보다 더 초등적이다. 산술적으로는 소아이디얼 분해의 유일성, 더 정확히는 그 분해정리의 바탕에 있는 사실들의 직접적인 귀결로 표현할 수 있으므로 더 높은 환론에 의존하지 않기 때문이다. 다음과 같이 하면 이 산술적 표현에 도달한다.

변수들을 분수선형변환하여 기초곡선 \(f(z,\delta)=0\)의 특이점들과 고려하는 유한개의 점군이 모두 유한영역에 놓이게 하자. 또한 그 뒤에 변수들의 정수계수 선형변환을 가하여 \(\delta\)가 \(z\)에 관하여 정수적이게 하자. 그러면 절대 리만곡면에서 \(z=a,\delta=b\)가 되는 점 \(P\)를 생성하는 소아이디얼 \(\mathfrak p\)에는 \(f(z,\delta)=0\) 위의 점 \(z=a,\delta=b\)가 대응한다. 더 일반적으로 아이디얼
\[
\mathfrak a=\mathfrak p_1^{\varrho_1}\cdots\mathfrak p_\sigma^{\varrho_\sigma}
\]
은 각 점 \(z=a_i,\delta=b_i\)를 대응하는 중복도만큼 포함하는 점군을 생성한다. \(f(z,\delta)=0\) 위의 유한영역에 있는 모든 점이 이렇게 생성된다. 한 점 \(P\)에서는 \(\mathfrak p\)로 나누어지는 모든 함수 \(g_1,g_2\ldots\)가 사라지고, 거꾸로 이 함수들 전체가 아이디얼 \(\mathfrak p\)를 생성한다. 따라서 \(f\) 위의 대응점은 \(g_1(z,\delta)=0,\ g_2(z,\delta)=0\ldots\)과 \(f=0\)의 교점이다. 아이디얼 \(\mathfrak a\)에 대응하는 점군도 마찬가지이다. 모든 아이디얼이 두 주아이디얼의 최대공약아이디얼이라는 정리에 따라, 한 점군을 잘라 내는 데에는 언제나 그러한 곡선 두 개면 충분하다. 특히 한 주아이디얼에 대응하는 점군은 \emph{하나의} 곡선 \(g(z,\delta)=0\)이 잘라 내며 그 역도 성립한다. 주아이디얼은 완전교차에 대응한다.

기하학적 이론에서는 두 점군 \(A,B\)가 같은 점군 \(R\)과 합쳐져 완전교차를 이룰 때 이들을 \emph{공잔여}라고 부른다. 이 점군들을 각각 아이디얼 \(\mathfrak a,\mathfrak b,\mathfrak r\)가 생성한다면
\[
\mathfrak a\mathfrak r=(\alpha),\qquad
\mathfrak b\mathfrak r=(\beta)
\]
라는 관계가 성립하며, 이는 아이디얼 \(\mathfrak a,\mathfrak b\)가 동치임을 보여 준다. 거꾸로 두 아이디얼의 동치 \(\mathfrak a=\eta\mathfrak b\)로부터도, 모든 아이디얼은 다른 한 아이디얼을 곱하면 주아이디얼이 된다는 정리에 따라 언제나 이 관계가 따라온다\srcfn{1)}{데데킨트, 『수론』 \S{}181 참조. 실제로 \(\mathfrak a\mathfrak r=(\alpha)\)가 되도록 \(\mathfrak r\)을 택하면 \(\mathfrak b\mathfrak r=\eta^{-1}(\alpha)\)이다. \(\mathfrak b\mathfrak r\)은 정수적 원소들만을 포함하므로 \(\eta^{-1}(\alpha)\)도 그러하며, 따라서 이것도 주아이디얼 \((\beta)\)가 된다.}. 즉 \emph{동치인 아이디얼과 공잔여 점군은 동일한 개념이다.}

이제 잔여정리는 다음과 같이 말한다. \emph{\(A,B\)가 공잔여이면, 같은 잔여를 갖는 수반곡선들로 두 점군을 잘라 낼 수 있다. 단 특이점들은 제외한다.} 한 함수의 영점과 극점은 언제나 공잔여이므로, 이 정리는 동시에 체의 가장 일반적인 함수를 수반형식의 몫으로 표현할 수 있음을 증명한다. 아이디얼론의 언어로 쓰면, 잔여정리는 \(\mathfrak a,\mathfrak b\)뿐 아니라 \(\mathfrak a\mathfrak f,\mathfrak b\mathfrak f\)도 동치라는 사실에 불과하다. 여기서 \(\mathfrak f\)는 이중점 아이디얼이다. 실제로 위의 결과에 따라
\[
\mathfrak a\mathfrak f\mathfrak m=(\alpha),\qquad
\mathfrak b\mathfrak f\mathfrak m=(\beta)
\]
가 주아이디얼이 되게 하는 아이디얼 \(\mathfrak m\)이 언제나 존재한다. \(\alpha=0,\beta=0\)은 점군 \(A,B\), 더 정확히는 그 가운데 특이점과 다른 부분군 \(A',B'\)을 잘라 내는 수반곡선들이다\srcfn{2)}{\(\mathfrak a=\mathfrak f'\mathfrak a_1,\ \mathfrak b=\mathfrak f'\mathfrak b_1,\ \mathfrak f=\mathfrak f'\mathfrak f_1\)이고 \(\mathfrak a_1,\mathfrak b_1,\mathfrak f_1\)가 서로소라면, 이미 \(\mathfrak a_1\mathfrak f\mathfrak n,\mathfrak b_1\mathfrak f\mathfrak n\)이 주아이디얼이 된다. 여기서 \(\mathfrak n\)은 \(\mathfrak f\)와 서로소이다. 실제로 모든 아이디얼은 주어진 한 아이디얼과 서로소인 아이디얼을 곱하여 주아이디얼로 만들 수 있다(D.--W. 214쪽 또는 『수론』 \S{}178, 정리 XI).}. 따라서 잔여정리의 산술적 증명은 소아이디얼들에 점들을 대응시킬 수 있다는 사실과, 동치인 아이디얼들을 같은 하나의 아이디얼로 곱하여 주아이디얼들로 만들 수 있다는 사실에 있다.

공잔여 점군들은 같은 수의 점으로 이루어질 필요가 없다. 마찬가지로 동치인 아이디얼들도 같은 차수일 필요가 없다. 예를 들어 모든 주아이디얼은 서로 동치이고, 완전교차가 생성하는 모든 점군은 서로 공잔여이다. 동치인 폴리곤(인자)과의 이러한 차이\srcfn{3)}{폴리곤 \(A\)는 절대 리만곡면의 유한개의 점으로 이루어진다. 두 폴리곤 \(A,B\)가 한 함수 \(\eta\)의 영점과 극점이면 이들을 동치라고 한다. 그러면 \(\eta\)는 폴리곤의 몫 \(\eta=\frac{A}{B}\)으로 표현된다(D.--W. \S{}17, 18).}는 \(z\)가 무한대가 되는 점들이 \(z\)에 관한 소아이디얼로 생성되지 않는다는 사실에서 설명된다. 아이디얼의 몫
\[
\eta=\frac{\mathfrak a}{\mathfrak b}
\]
으로 함수를 표현하는 것은 \(\mathfrak a,\mathfrak b\)의 동치를 말하지만, 여기에는 \(z\)가 무한대가 되는 곳에 있는 \(\eta\)의 영점이나 극점이 나타나지 않는다. 예를 들어 주아이디얼 \((\alpha)\)는 \emph{정수적} 함수 \(\alpha\)의 \emph{영점}에 대응하는 소아이디얼들의 곱이다. \(\alpha\)의 극점들은 나타나지 않는다. \(\alpha\)는 \(z\)가 무한대가 되는 점에서만 무한대가 되기 때문이다. 이런 의미에서 아이디얼론은 기하학적 이론의 \emph{형식론}과 정확히 대응한다. 형식론에도 영점만 있고 극점은 없다.

기하학에서는 같은 차수의 형식들로 몫을 만들어 형식에서 함수로 넘어간다. 산술적으로 이에 대응하는 것은 아이디얼류에서 불변적으로 정의된 폴리곤류(인자류)로 넘어가는 일이다. 이제 더 이상 어떤 변수도 특별히 선택되지 않으므로, 함수를 폴리곤의 몫으로 나타내면 모든 영점과 극점이 드러난다. 아이디얼류, 즉 공잔여 점군의 류와 폴리곤류의 관계는 다음과 같다. 아이디얼 \(\mathfrak a,\mathfrak b\)가 생성하는 폴리곤을 \(\mathfrak A,\mathfrak B\)라 하고, \(z\)가 무한대가 되는 점들의 폴리곤, 즉 \(z\)와 \(z\)의 정수적 함수들의 분모폴리곤을 \(\mathfrak N\)이라 하자. \(\mathfrak a,\mathfrak b\)가 동치일 때 \(\mathfrak A,\mathfrak B\)가 동치인 것은 두 아이디얼의 차수가 같을 때 그리고 그때에만 성립한다. 일반적으로는 \(\mathfrak A\mathfrak N^\rho\)와 \(\mathfrak B\mathfrak N^\sigma\), \(\rho\ne\sigma\), 의 동치만이 따라온다. 따라서 아이디얼들, 또는 공잔여 점군들을 류로 나누는 것은 무한히 많은 폴리곤류를 한 류로 합치는 것이다. 이것은 폴리곤들을 \(\mathfrak N\)의 거듭제곱을 법으로 한 류들로 나누는 것으로도 볼 수 있다\srcfn{1)}{헨젤--란츠베르크, 제25강, 438쪽. 두 변수 \(z,s\)를 폴리곤의 몫으로 표현하면 각 변수를 이원 동차화하는 것으로 볼 수 있다. 특히 \(z,s\)가 같은 분모폴리곤을 갖도록 택하면 기하학에서 통상 쓰는 삼원 동차화가 생긴다.}.

서로 다른 수의 점으로 이루어진 공잔여 점군에 대한 잔여정리는 순수 기하학적 문제에서만 나타난다. 대수함수에 적용할 때에는 실제로 같은 수의 공잔여 점들, 즉 폴리곤류에 대응하는 경우만 나타난다. 한 점군이 생성하는 \glqq{}완전 선형계\grqq{}, 곧 같은 수의 점으로 이루어지고 그 점군과 공잔여인 모든 점군은 대응하는 폴리곤이 생성하는 폴리곤류와 정확히 일치한다. \glqq{}두 완전계의 합\grqq{}과 \glqq{}두 폴리곤류의 곱\grqq{}이라는 개념들도 일치한다. 그러한 완전계가 유한개의 선형독립인 점군만을 포함한다는 사실은 잔여정리의 직접적인 귀결이다. 잔여정리는 그 선형계의 차원, 즉 선형독립인 점군의 수를, 잔여군을 지나는 임의의 고정된 차수의 모든 수반곡선의 차원으로 환원한다. 이 차원을 실제로 정하는 리만--로흐 정리에 도달하기 위해 기하학적 이론은 잔여정리로부터 먼저 \glqq{}환원정리\grqq{}, 즉 고정점에 관한 정리\srcfn{1)}{브릴--뇌터, 58번.}를 이끌어 내고, 다시 거기서 리만--로흐 정리를 얻는다.

완전히 이에 대응하여 산술이론은 각 폴리곤류의 차원이 유한함을 보이고, 정수적 원소들의 특수한 기저인 정규기저를 통해 그 차원을 정한다. 이 기저는 여모듈을 이용하여 정의된다. 이로써 여모듈과 분기폴리곤의 관계에서 리만--로흐 정리를 얻는다. D.--W.의 서술에서는 환원정리도 나타나지만\srcfn{2)}{D.--W., \S{}30, 2번.}, 기하학적 이론에서처럼 기초로 쓰이지는 않는다. 따라서 제6절에서 언급한 기저와 수반함수의 대응 때문에, 여기에서도 두 이론의 보조수단들 사이에 어느 정도의 대응이 있음을 확인할 수 있다. 요약하면 두 이론의 개념 구성은 표현과 배열이 다르기는 해도 본질적으로 일치한다. 두 이론은 서로 독립적으로 생겼으며, 기하학적 이론이 약 10년 앞선다. 끝으로 리만--로흐 정리는 원래 리만과 로흐에게서, 또 바이어슈트라스에게서 계수정리로 나타났음을 덧붙인다. 여기서 다루는 것은 행렬
\[
(\varphi_i(x_j))
\]
의 계수이다. \(\varphi\)들은 \(p\)개의 수반 \(n-3\)차 곡선을 돌고, \(x\)들은 잔여군의 점들을 돈다.

\begin{center}
\textbf{8. 대수함수와 대수적 수에서의 초월적 문제}
\end{center}

대수함수와 대수적 수의 초월적 문제들 사이의 유추는 특히 힐베르트가 지적하였다\srcfn{3)}{\emph{Mathematische Probleme}, 문제 12, \emph{Göttinger Nachrichten} 1900.}. 본질적으로 존재문제를 다룬다. 주어진 리만곡면, 따라서 주어진 분기점을 갖는 대수함수가 존재하는가라는 리만의 문제에는, 주어진 군과 판별식을 갖는 대수적 수체가 존재하는가라는 문제가 대응한다\srcfn{4)}{그로부터 자연스럽게 생기는, 군에서 출발하여 대수함수를 다루는 방법은 R. König이 더 일반적인 문제들의 특수한 경우로 수행하였다. 특히 \emph{Riemannsche Funktionen- und Differentialsysteme in der Ebene}, \emph{J. f. M.} 148, 그리고 \emph{Math. Ann.} 78, 79를 참조하라. König에서는 대수체 대신 \(K(z)\)에 관한 특수한 유한모듈이 나타나며, \glqq{}류의 성질\grqq{}만이 성립한다.}. 대수함수의 경우에는 존재증명을 일반적으로 수행할 수 있지만, 대수적 수체에서는 특수한 기초체, 즉 유리수체와 이차수체에 대해서, 그리고 아벨군에 대해서만 성공하였다\srcfn{1)}{삼차와 사차 방정식에서 나타나는 군들의 단순한 경우와 팔차 방정식에서의 몇몇 매우 특수한 군들만은 임의의 기초체에 대하여 다룰 수 있다. G. Bucht, \emph{Über einige algebraische Körper 8. Grades}, \emph{Arkiv f. Math., Astron. och Physik} 6, 30번. E. Noether, \emph{Gleichungen mit vorgeschriebener Gruppe}와 거기서 인용한 Seidelmann의 학위논문, \emph{Math. Ann.} 78도 참조하라.}. 여기서 문제되는 것은 유리수체 위의 모든 아벨수체가 단위근체들로부터 합성된다는 크로네커의 정리와, 상대 아벨수체에 대한 퓌터\srcfn{2)}{\emph{Math. Ann.} 75.} 및 헤케\srcfn{3)}{\emph{Math. Ann.} 76.}의 대응 정리들이다. 따라서 구하는 수체들은 지수함수와 모듈러함수라는 초월적 수단을 통해 실제로 주어진다. 이는 대수함수에서의 순수 존재증명과 대조된다.

대수함수의 존재증명에는 전역적으로 유한인 적분들의 존재가 선행한다. 이에 대한 유사물은 한 대수적 수체에서 소수 \(l\)에 관하여 정의되는 \glqq{}특이 일차수\grqq{}의 존재이다. 그 \(l\)제곱근들은 상대적으로 비분기이고 유체를 구성하는 최초의 벽돌이 된다\srcfn{4)}{Furtwängler, \emph{Reziprozitätsgesetze für Potenzreste mit Primzahlexponenten in algebraischen Zahlkörpern}, \emph{Math. Ann.} 67, 72.}. 이 존재증명에는 수체 안에 주어진 잉여지표를 갖는 소아이디얼이 언제나 존재한다는 정리가 본질적으로 쓰인다. 따라서 이 정리는 전역적으로 유한인 적분의 존재가 바탕을 두는 경곗값문제의 유사물로 볼 수 있다.

전역적으로 유한인 적분들은 두 폴리곤이 같은 류에 속하는지, 즉 같은 수의 점으로 이루어진 두 점군이 공잔여인지를 판정하는 초월적 기준을 준다. 이것은 \glqq{}공잔여 경로\grqq{}에 걸쳐 취한 제1종 적분들이 사라진다고 말하는 아벨 정리, 또는 그 역을 포함한 대응 미분정리이다\srcfn{5)}{예를 들어 M. Noether, \emph{Ann.} 37의 표현을 참조하라. 역사적으로 잔여정리가 아벨 정리에서 생겨났기 때문에 기하학적 이론은 적분의 합을 본떠 점군과 선형계의 \glqq{}합\grqq{}을 말한다. 이는 산술이론에서 류의 \glqq{}곱\grqq{}을 말하는 것과 대조된다.}.

대수적 수체에서 아벨 정리의 유사물은 \glqq{}\(l\)제곱잉여의 상호법칙\grqq{}이다. 이는 수체 자체에서, 또는 적어도 적절한 상체에서 두 아이디얼이 같은 류에 속한다는 판정기준을 준다. 그 내용은 이 상체에서 특이 일차수들이 같은 류의 모든 아이디얼에 관하여 같은 잉여지표를 갖는다는 식으로도 쓸 수 있다. 이 마지막 사실은 원래 체에서도 성립하지만, 그곳에서는 상호법칙과 일치하지 않는다\srcfn{1)}{Furtwängler, 앞의 글, \emph{Annalen} 67, 특히 서론 참조.}.

끝으로 헨젤의 \(p\)-진수론은 대수함수의 거듭제곱급수 전개에 대한 유사물을 준다. 이에 대해서는 헨젤의 보고를 참조하라.

\clearpage
\setcounter{footnote}{0}
\phantomsection
\addcontentsline{toc}{section}{15. 이원형식의 정수 계수 불변식계의 유한성}
%
%
%
%
%
%

\section*{15. 이원형식의 정수 계수 불변식계의 유한성}

\begin{center}
\emph{Nachr. v. d. Ges. d. Wiss. zu Göttingen 1919, 138--156쪽}\\[5em]
에미 뇌터 (괴팅겐)\\[0.8em]
F. 클라인이 1919년 3월 27일 회의에 제출함.
\end{center}

이 글에서는 잘 알려진 유한성 정리를 확장하여 다음 정리를 증명한다.

\emph{한 이원 기본형식계의 모든 정수 계수 다항불변식은 그 가운데 유한개에 관한 정수 계수 다항함수이다.}

통상적인 용법과 같이 여기서 \glqq{}불변식\grqq{}은 언제나 한 기본형식 또는 기본형식계의 다항 유리불변식, 더 정확히는 모든 선형변환의 군이 유도하는 계수변환에 관한 불변식을 뜻한다.

증명은 메르텐스--힐베르트의 유한성 증명을 정수성에 관하여 강화하는 데에 바탕을 둔다\srcfn{1)}{F. Mertens, Wiener Sitzungsberichte, 1889년 1월; D. Hilbert, \emph{Math. Ann.} 33 (1889) [1888년 3월 날짜]. 메르텐스의 Crelle 100권 논문에 이어 서로 독립적으로 나온 두 증명은 내용상 일치한다.}.

이 유한성 증명은 다음과 같이 특징지을 수 있다. 기본형식들을 일차인수들로 분해한다. 더 정확히 말하면 각 기본형식에 일차형식들의 곱을 대응시킨다. 그러면 각 불변식 \(I\)는 이 일차형식들의 불변식, 곧 동차화한 근들의 함수로 넘어간다. 이 함수는 세 조건을 만족해야 한다. 1) 불변성 조건, 2) 가중치 조건, 3) 대칭성 조건이다. 조건 2), 3)은 \(I\)가 기본형식의 계수들에 관한 다항 유리함수가 되기 위하여 필요하다. 조건 1)은 \(I\)를 일차형식들의 행렬식에 관한 다항함수로 나타내어 만족시킨다. 조건 2)는 개별 행렬식 대신 그 행렬식들의 어떤 거듭제곱곱들을 새로운 인수로 삼아 만족시킨다. 이제 조건 3)은 \(I\)가 변수열들의 다항 대칭함수가 된다는 것뿐이며, 역으로 변수열들의 그러한 각 대칭함수는 기본형식들의 불변식이 된다. \emph{따라서 이 변수열들의 기본대칭함수들이 완전계를 이룬다.}

이제 정수 계수 불변식만을 생각하자. 일차형식의 불변식에 관하여 \emph{이원 일차형식의 모든 정수 계수 불변식은 이 일차형식들의 행렬식에 관한 정수 계수 다항함수이다}라는 더 강한 정리를 증명하면 조건 1)은 위와 똑같이 충족된다. 이것이 실제로 옳다는 사실을 § 3에서 보인다. 그곳에서는 이 행렬식들 사이에 어떤 소수를 법으로 성립하는 모든 관계의 전체가 항등관계의 전체와 일치한다는 형태로 증명한다. 즉 이 관계들에 대응하는 모듈은 각 소수를 법으로 해도 소모듈로 남는다. 증명은 행렬식들 사이의 이차관계에 대응하는 이차형식들의 모듈에 관하여 잉여류들을 정규화하여 이루어진다. 이 이차형식들은 각 소수를 법으로 해도 모듈의 정수 계수 기저를 이룬다. 이전에는 정수성을 무시할 때에만 모듈 기저를 이룬다는 사실이 알려져 있었다.

조건 2)를 충족하는 데에는 아무 수정도 필요 없다. 거듭제곱곱을 도입해도 수계수는 영향을 받지 않기 때문이다. 조건 3)은 이제 \(n_1!\ldots n_\mu! I\), 여기서 \(n_1,\ldots n_\mu\)는 개별 기본형식들의 차수, 가 변수열들의 정수 계수 대칭함수가 된다는 뜻이고, 역으로 이 변수열들의 각 정수 계수 대칭함수는 정수 계수 불변식이 된다는 뜻이다. 그런데 쉽게 보일 수 있듯이, 변수열들의 이 정수 계수 대칭함수들은 정수 계수 기저를 가지며 여기에는 기본함수들만으로는 충분하지 않다. 힐베르트의 정수 계수 모듈 기저 정리를 적용하면 여기서 \(I\)들 자체에 대해서도 정수 계수 기저를 얻는다.

§ 1에서는 조건 1), 2), 3)에 각각 대응하는 특수 유한성 정리들을 모은다. § 2에서는 이 정리들로 유한성 증명을 이끌어 낸다. 조건 3)에 대응하는 특수 유한성 정리들은 동시에 \emph{정수 계수 또는 대수적 정수 계수 치환들로 이루어진 유한군의 정수 계수 불변식의 유한성}을 주며, § 4에서 이를 간단히 설명한다.

\begin{center}
\textbf{§ 1. 특수 유한성 정리}
\end{center}

서론의 불변성, 가중치, 대칭성 조건을 만족시키는 데 사용할 다음 특수 유한성 정리들을 먼저 둔다.

\textbf{I.} \emph{이원 일차형식계의 각 정수 계수 불변식은 이 일차형식들의 행렬식에 관한 정수 계수 다항함수이다.} 증명은 § 3에서 한다.

\textbf{II.} \emph{음이 아닌 지수를 갖는 \(n\)개 변수의 거듭제곱곱들로 이루어진 계 \(\mathfrak S\)가 다음 성질을 가진다고 하자. \(f,g\in\mathfrak S\)이고 그 몫이 변수들에 관한 다항식이면 그 몫도 \(\mathfrak S\)에 속한다. 그러면 \(\mathfrak S\)에는 유한 곱셈기저 \(f_1,\ldots,f_k\)가 존재하여, 각 \(f\in\mathfrak S\)를 음이 아닌 지수 \(\lambda_1,\ldots,\lambda_k\)로}
\[
f=f_1^{\lambda_1}\cdots f_k^{\lambda_k}
\]
\emph{와 같이 나타낼 수 있다.}

실제로 힐베르트의 모듈 기저 정리에 따라 \(\mathfrak S\)의 함수 가운데 \(x\)가 실제로 나타나는 유한개의 함수 \(f_1,\ldots,f_k\)가 존재하여, \(\mathfrak S\)의 상수가 아닌 각 \(f\)는 이들이 생성하는 모듈에 속한다.
\[
f\equiv0\pmod{(f_1,\ldots,f_k)}.
\]
따라서 각 \(f=x_1^{i_1}\cdots x_n^{i_n}\)에 대하여 \(f_\nu=x_1^{\nu_1}\cdots x_n^{\nu_n}\)인 적어도 하나의 \(f_\nu\)가 존재하여
\[
\nu_1\le i_1,\ldots,\nu_n\le i_n
\]
이다. 표현
\[
f=x_1^{i_1-\nu_1}\cdots x_n^{i_n-\nu_n}f_\nu=Af_\nu\srcfnmark{1)}
\]
\srcfntext{1)}{이 표현은 직접 증명할 수도 있고, 그러면 역으로 힐베르트의 정리도 다시 따라온다. P. Gordan, \emph{Neuer Beweis des Hilbertschen Satzes über homogene Funktionen}, Gött. Nachr. 1899를 보라. 이 표현을 직접 증명하고 여기서 정리 II도 이끌어 내는 또 다른 증명은 A. Ostrowski, \emph{Über die Existenz einer endlichen Basis bei gewissen Funktionensystemen}, Math. Ann. 78 (1916), § 2, 1 및 § 3, 2에 있다. § 2에서 쓰는 정리 II의 특수한 경우, 곧 지수들이 어떤 디오판토스 방정식계의 음이 아닌 모든 해를 두루 취하는 경우는 Gordan--Kerschensteiner, \emph{Vorlesungen über Invariantentheorie}, 1권, 199쪽에 이미 나온다.}
에서 가정에 따라 \(A\in\mathfrak S\)이고, \(A\)의 변수 차수는 \(f\)보다 낮다. 따라서 이 논법을 유한번 반복하면 주장이 따라온다. \(f=1\)에도 표현이 성립하게 하려면 모든 지수 \(\lambda\)를 0으로 두면 된다.

\textbf{III.} \emph{\(n\)개 변수열의 모든 정수 계수 다항 대칭함수는 그 가운데 유한개에 관한 정수 계수 다항함수이다.}

\(\tau_1(x),\ldots,\tau_n(x)\)를 부정원 \(x_1,\ldots,x_n\)의 기본대칭함수들이라 하자. 그러면 각 지수 \(k\)에 대하여 항등식
\[
x_i^k=G_0^{(k)}(\tau)+x_iG_1^{(k)}(\tau)+\cdots+x_i^{n-1}G_{n-1}^{(k)}(\tau),
\qquad (i=1,2,\ldots,n)
\]
이 성립한다. 여기서 \(G_0^{(k)},\ldots,G_{n-1}^{(k)}\)는 \(\tau(x)\)의 정수 계수 다항함수이다. 이제 \(x_i,y_i,\ldots,z_i\)를 \(i\)번째 열의 원소들이라 하자. 그러면 이 항등식과 \(y,\ldots,z\)에 대한 대응 항등식들에 따라 특수한 단일형 대칭함수
\[
\sum_{i=1}^n x_i^\alpha y_i^\beta\cdots z_i^\gamma
\]
는 유한개의 함수
\[
\sum_{i=1}^n x_i^a y_i^b\cdots z_i^c,
\qquad
(0\le a<n;\;0\le b<n;\;\ldots\;0\le c<n)
\]
와 기본대칭함수 \(\tau(x),\tau(y),\ldots,\tau(z)\)의 정수 계수 다항함수가 된다\srcfn{1)}{힐베르트는 앞의 논문에서 \(\tau(x),\ldots,\tau(z)\) 대신 거듭제곱합들을 기저에 넣는다. 그러면 기저도 단일형 함수들만으로 이루어지지만 표현의 정수성이 사라진다.}.

이하에서는 나타나지 않는 가장 일반적인 \(n\)개 열의 대칭함수의 경우에도 같은 방법으로
\[
\begin{gathered}
\sum x_{i_1}^{a_1}y_{i_1}^{b_1}\cdots z_{i_1}^{c_1}
      x_{i_2}^{a_2}y_{i_2}^{b_2}\cdots z_{i_2}^{c_2}
      \cdots
      x_{i_n}^{a_n}y_{i_n}^{b_n}\cdots z_{i_n}^{c_n},\\
(0\le a_\lambda<n;\;0\le b_\lambda<n;\;\ldots\;0\le c_\lambda<n)
\end{gathered}
\]
여기서 \(i_1,i_2,\ldots,i_n\)은 모든 순열을 두루 취한다, 와 \(\tau(x),\tau(y),\ldots,\tau(z)\)가 함께 정수 계수 기저를 이룸을 알 수 있다.

\textbf{IV.} \emph{\(F_1(x),\ldots,F_k(x)\)를 \(x_1,\ldots,x_n\)의 정수 계수 다항함수들이라 하고 \(a\)를 고정된 정수라 하자. \(G\)가 정수 계수 다항함수이고 \(\frac{G(F)}a\)가 \(x\)에 관한 정수 계수 다항함수가 되는 모든 함수 \(\frac{G(F)}a\)는 그 가운데 유한개에 관한 정수 계수 다항함수이다.}

실제로 힐베르트의 정수 계수 모듈 기저 정리에 따라 각 \(\frac{G(F)}a\)는
\[
\frac{G(F)}a=
A_1(F_1,\ldots,F_k)\frac{G_1(F)}a+\cdots+
A_\varrho(F_1,\ldots,F_k)\frac{G_\varrho(F)}a
\]
로 나타난다. 여기서 \(A_1,\ldots,A_\varrho\)는 \(F\)의 정수 계수 다항함수이고,
\[
\frac{G_1(F)}a,\ldots,\frac{G_\varrho(F)}a
\]
도 해당 계에 속한다. 따라서 함수들
\[
F_1=\frac{aF_1}{a},\ldots,F_k=\frac{aF_k}{a},
\qquad
\frac{G_1(F)}a,\ldots,\frac{G_\varrho(F)}a
\]
가 구하는 기저를 이룬다.

\begin{center}
\textbf{§ 2. 정수 계수 이원불변식의 유한성}
\end{center}

\(x,y\)를 부정계수 변환
\begin{equation}
\tag{1}
x=s_{11}x'+s_{12}y';\qquad y=s_{21}x'+s_{22}y'
\end{equation}
을 받는 이원변수라 하고, \(f\)를 계수가 부정원인 이원 기본형식이라 하자. 그러면 (1)에 따라
\begin{equation}
\tag{2}
f=a_0x^n+a_1x^{n-1}y+\cdots+a_ny^n
=a'_0x'^n+a'_1x'^{\,n-1}y'+\cdots+a'_ny'^n
\end{equation}
이다. 잘 알려져 있듯이 \(f\)의 정수 계수 불변식이란
\begin{equation}
\tag{3}
I(a')=|s_{ik}|^\varrho I(a)
\qquad\text{\(a,s_{ik}\)에 관한 항등식}
\end{equation}
을 만족하는 정수 계수 다항함수 \(I(a)\)를 뜻한다\srcfn{1)}{불변식론에서 흔히 하듯 이항계수를 써서 \(f=b_0x^n+\binom n1b_1x^{n-1}y+\cdots+b_ny^n\)로 쓰면, 모든 정수 계수 불변식 \(I(b)\)의 영역은 \(I(a)\)의 영역보다 넓어지고 여기서 쓰는 유한성 논법은 실패한다. \(I(b)\)가 더는 근들의 정수 계수 다항함수가 되지 않기 때문이다.}. 기본형식계의 정수 계수 불변식도 이에 맞추어 정의한다. 각 기본형식의 계수들에 관한 동차성을 정의에 함께 넣어도 된다. 각 형식에 관하여 따로 동차인 불변식들의 정수 계수 선형결합으로 나머지 불변식들을 만들 수 있기 때문이다.

이제 (2)와 함께 하나의 \emph{특수한} 이원형식 \(g\), 곧 계수들이 다시 \emph{부정원}인 \(n\)개 일차형식의 곱을 생각하자.
\begin{equation}
\tag{4}
\begin{aligned}
g&=(\alpha_1x+\beta_1y)(\alpha_2x+\beta_2y)\cdots(\alpha_nx+\beta_ny)\\
 &=\sigma_0(\alpha,\beta)x^n+\sigma_1(\alpha,\beta)x^{n-1}y+\cdots+\sigma_n(\alpha,\beta)y^n\\
 &=(\alpha'_1x'+\beta'_1y')(\alpha'_2x'+\beta'_2y')\cdots(\alpha'_nx'+\beta'_ny')\\
 &=\sigma'_0x'^n+\sigma'_1x'^{\,n-1}y'+\cdots+\sigma'_ny'^n.
\end{aligned}
\end{equation}
여기서 \(\sigma_0(\alpha,\beta),\sigma_1(\alpha,\beta),\ldots,\sigma_n(\alpha,\beta)\)는 동차화한 기본대칭함수들이고, 변환된 계수 \(\sigma'_i\)는 \(\sigma_i(\alpha',\beta')\)와 같다.

이 기본대칭함수들이 대수적으로 독립이므로 각 관계
\begin{equation}
\tag{5}
G(\sigma_0,\sigma_1,\ldots,\sigma_n)=0
\qquad\text{\(\alpha,\beta\)에 관한 항등식}
\end{equation}
에는 관계
\begin{equation}
\tag{6}
G(\sigma_0,\sigma_1,\ldots,\sigma_n)=0
\qquad\text{\(\sigma_0,\ldots,\sigma_n\)에 관한 항등식}
\end{equation}
이 대응하고, 따라서 부정원 \(a_0,\ldots,a_n\)에 대해서도
\begin{equation}
\tag{7}
G(a_0,a_1,\ldots,a_n)=0
\qquad\text{\(a_0,\ldots,a_n\)에 관한 항등식}
\end{equation}
이 대응한다.

\(f\)의 각 정수 계수 불변식 \(I(a)\)에는 치환 \(a_i=\sigma_i(\alpha,\beta)\)을 통하여 \(g\)의 정수 계수 불변식
\[
I(\sigma)=H(\alpha,\beta)
\]
가 대응한다. 여기서 \(H\)는 \(\alpha,\beta\)의 정수 계수 다항함수이다. 이 불변식에 대하여 관계 (3)은
\begin{equation}
\tag{8}
I(\sigma')=|s_{ik}|^\varrho I(\sigma)
\qquad\text{\(\sigma,s_{ik}\)에 관한 항등식}
\end{equation}
으로 넘어간다. \(\sigma'=\sigma(\alpha',\beta')\)이므로 (8)에서
\begin{equation}
\tag{9}
H(\alpha',\beta')=|s_{ik}|^\varrho H(\alpha,\beta)
\qquad\text{\(\alpha,\beta;s_{ik}\)에 관한 항등식}
\end{equation}
이 따라온다. 따라서 \(I(\sigma)=H(\alpha,\beta)\)는 \(n\)개 일차형식 \((\alpha_i x+\beta_i y)\)의 정수 계수 불변식이다.

역으로 \(H(\alpha,\beta)\)가 이 \(n\)개 일차형식의 정수 계수 불변식이어서 (9)를 만족하고, 동시에 \(\sigma\)의 정수 계수 다항함수 \(I(\sigma)\)와 같다고 하자. 그러면
\[
H(\alpha',\beta')=I(\sigma(\alpha',\beta'))=I(\sigma')
\]
이므로 (9)는
\begin{equation}
\tag{9a}
I(\sigma')=|s_{ik}|^\varrho I(\sigma)
\qquad\text{\(\alpha,\beta;s_{ik}\)에 관한 항등식}
\end{equation}
으로도 쓸 수 있다. (5), (6)에 따라 여기서 (8)이, 다시 (7)에 따라 (3)이 따라온다. \emph{따라서 각 정수 계수 불변식 \(I(a)\)에는 \(n\)개 일차형식의 정수 계수 불변식 \(H(\alpha,\beta)=I(\sigma)\)가 대응하고, 그 역도 성립한다.}

더 나아가 \(I_1(a),\ldots,I_k(a)\)가 모든 \(I(\alpha)\)의 정수 계수 기저이면 \(H_1(\alpha,\beta),\ldots,H_k(\alpha,\beta)\)는 동시에 \(\sigma\)의 정수 계수 다항함수 \(I(\sigma)\)인 모든 \(H(\alpha,\beta)\)의 정수 계수 기저가 된다. (5), (6), (7)에 따라 여기서도 역이 성립한다. \emph{그러므로 모든 \(I(a)\)의 정수 계수 기저가 존재함을 보이려면, 동시에 정수 계수 다항함수 \(I(\sigma)\)인 모든 정수 계수 불변식 \(H(\alpha,\beta)\)에 대하여 이를 보이면 충분하다. 기저}
\[
H_1(\alpha,\beta)=I_1(\sigma),\ldots,H_k(\alpha,\beta)=I_k(\sigma)
\]
\emph{는 \(I(a)\)들에 기저 \(I_1(a),\ldots I_k(a)\)를 준다.}

이제 \(H(\alpha,\beta)=I(\sigma)\)가 (4)의 모든 정수 계수 불변식을 두루 취하게 하자. (3)에 따라 각 불변식은 \(\sigma\)들에 관하여 동차이고, (4)에 따라 각 \(\sigma\)는 각 일차형식의 계수에 관하여 동차이며 일차이다. 따라서 \(H(\alpha,\beta)\)는 각 열 \(\alpha_i\beta_i\)에 관하여 동차이고 모든 열에서 차수가 같으며, 물론 \(n\)개 열에 관한 대칭함수이다. 역으로 \(n\)개 열의 정수 계수 대칭함수 \(H(\alpha,\beta)\)가 각 열에 관하여 따로 동차이고 모든 열에서 차수가 같으면, 대칭함수의 기본정리에 따라 \(\sigma\)들의 정수 계수 다항함수이다. 따라서 \(H(\alpha,\beta)\)가 불변이면 위의 결과에 따라 불변식 \(I(\sigma)\)가 된다. 그러므로 (4)의 정수 계수 불변식 전체는 다음 성질들을 가진 정수 계수 다항함수 \(H(\alpha,\beta)\) 전체로 특징지어진다.
\[
\begin{array}{ll}
1)& n\text{개 일차형식 }\alpha_i x+\beta_i y\text{의 정수 계수 불변식이다};\\
2)& n\text{개 열 }\alpha_i,\beta_i\text{ 각각에 관하여 동차이고 모든 열에서 차수가 같다};\\
3)& n\text{개 열 }\alpha_i,\beta_i\text{에 관하여 대칭이다}.
\end{array}
\]

특수 유한성 정리 I에 따라 조건 1) 때문에 \(H(\alpha,\beta)\)는 행렬식
\[
(ik)=\alpha_i\beta_k-\beta_i\alpha_k
\]
의 정수 계수 다항함수이다. 즉
\[
H(\alpha,\beta)=G((ik)).
\]
이 표현에 \(n\)개 열의 모든 순열을 적용하면 대칭성을 형식적으로도 나타낼 수 있다.
\begin{equation}
\tag{10}
n!\,H(\alpha,\beta)
G^{(1)}((ik))+G^{(2)}((ik))+\cdots+G^{(n!)}((ik))
=G^*((ik)).
\end{equation}
\(G^*\)는 \((ik)\)들의 각 거듭제곱곱 \(P\)와 함께 \(P\)의 모든 순열에 대한 합
\[
Q=P^{(1)}+\cdots+P^{(n!)}
\]
를 포함하며, 유한개의 \(Q\)에 관한 정수 계수 선형결합 \(G^*=\sum cQ\)가 된다. \(Q\)들은 조건 1), 2), 3)을 만족하여 그 자체가 불변식이므로 \(G^*((ik))\) 대신 \(Q\)들을 생각하면 충분하다.

이제 \(P\)가 \(Q\)들에 나타나는 모든 거듭제곱곱, 곧 조건 2)를 만족하는 \((ik)\)의 모든 거듭제곱곱으로 이루어진 계 \(\mathfrak S\)를 두루 취하게 하자. 두 \(P\)의 몫이 \((ik)\)들의 다항식이면 그 몫도 조건 2)를 만족하므로 \(\mathfrak S\)에 속한다. 정리 II에 따라 그러한 \(P\) 가운데 유한개 \(P_1,\ldots,P_\varrho\)가 존재하여 각 \(P\)를 음이 아닌 지수로
\[
P=P_1^{\lambda_1}\cdots P_\varrho^{\lambda_\varrho}
\]
와 같이 나타낼 수 있다. 모든 순열을 적용하면
\[
Q=
P_1^{(1)\lambda_1}\cdots P_\varrho^{(1)\lambda_\varrho}
+\cdots+
P_1^{(n!)\lambda_1}\cdots P_\varrho^{(n!)\lambda_\varrho}.
\]
따라서 \(Q\)는 각각 \(\varrho\)개 원소로 이루어진 다음 \(n!\)개 변수열의 정수 계수 대칭함수이다.
\[
P_1^{(1)},\ldots,P_\varrho^{(1)};\quad
P_1^{(2)},\ldots,P_\varrho^{(2)};\quad\ldots;\quad
P_1^{(n!)},\ldots,P_\varrho^{(n!)}.
\]
그러므로 정리 III에 따라 \(Q\)는 이 변수열들의 유한개의 대칭함수에 관한 정수 계수 다항함수이다. 이 기저함수들은 그 자체가 \(Q\)이거나 \(n!\)개 원소 \(P_i^{(1)},P_i^{(2)},\ldots,P_i^{(n!)}\)의 기본대칭함수이므로 조건 1), 2), 3)을 만족하며 다음 정수 계수 불변식들이 된다.
\begin{equation}
\tag{11}
H_1(\alpha,\beta)=I_1(\sigma);\ldots;\quad
H_l(\alpha,\beta)=I_l(\sigma).
\end{equation}

(10)에서 \(G^*=\sum cQ\)이므로 각 \(H(\alpha,\beta)=I(\sigma)\)는
\[
H(\alpha,\beta)=I(\sigma)=\frac1{n!}\Phi(I_1,\ldots,I_l)
\]
꼴이다. 여기서 \(\Phi\)는 정수 계수 다항함수이다. 역으로 \(\frac1{n!}\Phi\)가 \(\alpha,\beta\)의 정수 계수 다항함수가 되면 정수 계수 불변식이다. 따라서 정리 IV에 따라 불변식 (11)에 유한개의 \(I_{l+1}(\sigma),\ldots,I_k(\sigma)\)를 더하여
\[
\begin{gathered}
H_1(\alpha,\beta)=I_1(\sigma),\ldots,H_l(\alpha,\beta)=I_l(\sigma),\\
H_{l+1}(\alpha,\beta)=I_{l+1}(\sigma),\ldots,H_k(\alpha,\beta)=I_k(\sigma)
\end{gathered}
\]
\emph{가 모든 불변식 \(H(\alpha,\beta)=I(\sigma)\)의 정수 계수 기저를 이루게 할 수 있다. 이로써 하나의 기본형식인 경우의 정리가 증명되었다.}

기본형식계를 다룰 때에는 위에서 말했듯 각 기본형식의 계수에 관한 불변식의 동차성을 가정할 수 있다. 다시 일반 기본형식들을 각각 일차형식들의 곱인 형식으로 바꾼다. 그러면 그 불변식들은 1) 이 일차형식들의 불변식이고, 2) 일차형식의 계수들로 이루어진 각 열에 관하여 동차이며 한 기본형식에 속한 열들에서는 차수가 같고, 3) 각 기본형식에 속한 열들에 관하여 대칭이다. 역으로 이 조건들을 만족하는 모든 일차형식의 정수 계수 불변식에는 다시 기본형식계의 정수 계수 불변식이 대응한다. 증명은 위의 증명과 완전히 나란히 진행된다. 개별 행렬식들의 가중치 조건을 만족하는 각 거듭제곱곱
\[
P=P_1^{\lambda_1}\cdots P_\varrho^{\lambda_\varrho}
\]
(정리 I, II)에 모든
\[
N=n_1!n_2!\cdots n_\mu!
\]
개의 순열을 적용한다. 여기서 \(n_1,\ldots n_\mu\)는 기본형식들의 차수이다. 그러면 \(N\cdot I\)는 각각 \(\varrho\)개 원소로 이루어진 \(N\)개 열의 정수 계수 대칭함수로 넘어간다. 여기서 모든 \(N\cdot I\)의 정수 계수 기저를 얻고, 이어서 정리 IV에 따라 모든 \(I\)의 정수 계수 기저를 얻는다.

\bigskip
\begin{center}
\textbf{§ 3. 이원 일차형식의 정수 계수 불변식의 기저}
\end{center}

이제 특수 유한성 정리 I의 증명을 보충해야 한다. 클렙슈가 처음 증명한 잘 알려진 정리에 따르면, \(n\)개 일차형식 \(\alpha_i x+\beta_i y\)의 모든 다항 유리불변식은 행렬식 \((ik)\)들의 다항함수이다. 따라서 이 일차형식들의 모든 \emph{정수 계수} 불변식도 이 행렬식들의 다항함수이지만 반드시 정수 계수 다항함수인 것은 아니다.

\emph{이하에서는 \((ik)\)들 사이에 성립하는 항등관계들을 사용하면 이 표현을 언제나 정수 계수로도 만들 수 있음을 보인다.}

이를 공식과 모듈론의 언어로 쓰면 다음과 같다. \(i<k\)인 \(\xi_{ik}\)들을 부정원이라 하자. \(F((ik))\)가 \(\alpha,\beta\)에 관하여 항등적으로 사라지는 모든 정수 계수 다항식 \(F(\xi_{ik})\)의 전체는 모듈 \(\mathfrak M\)을 이룬다. 두 \(F\)의 합과 한 \(F\)에 \(\xi_{ik}\)들의 임의의 정수 계수 다항식을 곱한 것도 이 전체에 속하기 때문이다. 공식으로 쓰면
\[
F((ik))=0\qquad\text{\(\alpha,\beta\)에 관한 항등식}
\]
에서
\[
F(\xi_{ik})\equiv0\pmod{\mathfrak M}
\qquad\text{\(\xi_{ik}\)에 관한 항등식}
\]
이 따라오고, 그 역도 성립한다.

이제 완전히 대응하여 다음이 성립함을 보이겠다.
\[
F((ik))\equiv0\pmod p
\qquad\text{\(\alpha,\beta\)에 관한 항등식}
\]
에서
\begin{equation}
\tag{1}
F(\xi_{ik})\equiv0\pmod{(\mathfrak M,p)}
\qquad\text{\(\xi_{ik}\)에 관한 항등식}
\end{equation}
이 따라오고 그 역도 성립하며, 이는 \emph{모든} 소수 \(p\)에 대하여 옳다. \emph{따라서 \(\mathfrak M\)은 \((ik)\)들의 모든 \(p\)를 법으로 하는 관계에 대응하는 모듈 \(\mathfrak M_p\)와도 같다.} 여기서 \(\mathfrak M_p\)는 \(F((ik))\)가 \(p\)를 법으로 \(\alpha,\beta\)에 관한 항등식으로 사라지는 모든 정수 계수 다항식 \(F(\xi_{ik})\)로 이루어진다.

(1)이 \emph{정수 계수 표현 가능성}을 뜻한다는 사실은 다음과 같이 알 수 있다. (1)은 이렇게도 쓸 수 있다. \(F((ik))=p\cdot H(\alpha,\beta)\)가 \(\alpha,\beta\)에 관한 항등식이면
\[
F(\xi_{ik})-p\cdot G(\xi_{ik})\equiv0\pmod{\mathfrak M}
\qquad\text{\(\xi_{ik}\)에 관한 항등식}
\]
이고, 따라서 \(\mathfrak M\)의 정의에 따라
\[
F((ik))=p\cdot G((ik))
\qquad\text{\(\alpha,\beta\)에 관한 항등식}\srcfn{1)}{이 마지막 추론은 동시에 (1)의 역을 증명한다.}
\]
이다. 그러므로 \(p\cdot H(\alpha,\beta)=p\cdot G((ik))\), 즉
\[
H(\alpha,\beta)=G((ik))
\]
이다. 따라서
\[
H(\alpha,\beta)=
\frac{1}{p_1^{\lambda_1}\cdots p_\varrho^{\lambda_\varrho}}F((ik))
\]
와 같은 표현이 있고 \(p_1,\ldots p_\varrho\)가 소수이면
\[
H(\alpha,\beta)=
\frac{1}{p_1^{\lambda_1-1}\cdots p_\varrho^{\lambda_\varrho}}G((ik))
\]
와 같은 표현도 있다. 이를 유한번 반복하면 표현의 정수성이 따라온다.

(1)을 증명하기 위하여
\[
f_{ik,rs}=\xi_{ik}\xi_{rs}-\xi_{ir}\xi_{ks}+\xi_{is}\xi_{kr},
\qquad i<k<r<s
\]
를 \((ik)\)들 사이의 이차관계
\[
(ik)(rs)-(ir)(ks)+(is)(kr)=0
\]
에 대응하는 부정원 \(\xi_{ik}\)들의 형식이라 하고, \(\mathfrak N=(f_{ik,rs})\)를 이 형식들이 생성하는 정수 계수 다항식 모듈이라 하자. 그러면
\[
F(\xi_{ik})\equiv0\pmod{\mathfrak N}
\quad\Longrightarrow\quad
F((ik))=0
\qquad\text{\(\alpha,\beta\)에 관한 항등식}
\]
이고,
\begin{equation}
\tag{2}
\begin{array}{c}
F(\xi_{ik})\equiv0\pmod{(\mathfrak N,p)}\\
\Longrightarrow\quad
F((ik))\equiv0\pmod p\quad\text{\(\alpha,\beta\)에 관한 항등식}
\end{array}
\end{equation}
이다.

\(\mathfrak N\)에 관한 잉여류의 정규화를 이용하여 (2)의 역도 성립함을 보이겠다. 그러면 \(\mathfrak N=\mathfrak M=\mathfrak M_p\)이고 (1)이 증명된다. 여기서 \(\mathfrak N=\mathfrak M\)은 \(p=0\)으로 둔 \(\mathfrak N=\mathfrak M_p\)의 특수한 경우로 얻어진다. \emph{이하의 고찰에는 \(p=0\)인 경우도 언제나 포함된다.} 달리 말하면, \emph{\(f_{ik,rs}\)들은 \(\mathfrak M_p\), 특히 \(\mathfrak M_0=\mathfrak M\)의 정수 계수 모듈 기저를 이룬다.}

\textbf{a) 두 열 또는 세 열 \(\alpha_i\beta_i\)인 경우.}
서로 다른 네 지표가 아직 나타나지 않으므로 \(\mathfrak N\)은 영모듈이다. 따라서
\begin{equation}
\tag{3}
\begin{array}{c}
F((ik))\equiv0\pmod p\quad\text{\(\alpha,\beta\)에 관한 항등식}\\[2pt]
\Longrightarrow\quad
F(\xi_{ik})\equiv0\pmod p\quad\text{\(\xi_{ik}\)에 관한 항등식}
\end{array}
\end{equation}
을 보이면 된다. 그러면 \(\mathfrak M=\mathfrak M_p=0\)인 (1)이 증명된다.

\(F(ik)=p\cdot H(\alpha,\beta)\)라는 조건은 각 열에 관하여 \emph{동차인} 성분마다 따로 충족되어야 한다. 따라서 언제나 \(F((ik))\)가 각 열에 관하여 동차이고, 곧 \(F(\xi_{ik})\)가 각 지표에 관하여 동차라고 가정하면 충분하다. 두 열인 경우에는 행렬식 \((12)\) 하나와 이에 대응하는 부정원 \(\xi_{12}\) 하나만 있다. 동차성에 따라 \(F=c\cdot\xi_{12}^{\varkappa}\)이다. 그런데
\[
c\cdot(12)^\varkappa\equiv0\pmod p
\qquad\text{\(\alpha,\beta\)에 관한 항등식}
\]
이고 \((12)\not\equiv0\pmod p\)이므로 언제나 \(c\equiv0\pmod p\)이다. 따라서
\[
F(\xi_{ik})\equiv0\pmod p
\qquad\text{\(\xi_{ik}\)에 관한 항등식}
\]
이고, 이로써 (3), 따라서 (1)이 증명된다.

세 열인 경우에는 행렬식 \((12),(13),(23)\)과 그에 대응하는 부정원 \(\xi_{12},\xi_{13},\xi_{23}\)이 있다. 따라서
\[
F=\sum c\,\xi_{12}^{\tau_{12}}\xi_{13}^{\tau_{13}}\xi_{23}^{\tau_{23}}.
\]
\(F\)가 지표 \(1,2,3\)에 관하여 각각 차수 \(\sigma_1,\sigma_2,\sigma_3\)라고 하자. 그러면
\[
\sigma_1=\tau_{12}+\tau_{13},\qquad
\sigma_2=\tau_{12}+\tau_{23},\qquad
\sigma_3=\tau_{13}+\tau_{23}
\]
이고, 역은 유일하게
\[
\tau_{12}=\frac12(\sigma_1+\sigma_2-\sigma_3),\quad
\tau_{13}=\frac12(\sigma_1-\sigma_2+\sigma_3),\quad
\tau_{23}=\frac12(-\sigma_1+\sigma_2+\sigma_3)
\]
로 주어진다. 따라서 각 지표에 관하여 동차인 \(F\)는 많아야 항 하나만 포함한다.
\[
F=c\,\xi_{12}^{\tau_{12}}\xi_{13}^{\tau_{13}}\xi_{23}^{\tau_{23}}.
\]
여기서 위와 마찬가지로 (3), 따라서 (1)이 성립한다.

\textbf{b) 네 열 \(\alpha_i,\beta_i\)인 경우.}
여섯 부정원 \(\xi_{12},\xi_{13},\xi_{14},\xi_{23},\xi_{24},\xi_{34}\)와 대응 행렬식들이 나타나고, \(\mathfrak N=(f_{12,34})\)이다. 관계
\begin{equation}
\tag{4}
\xi_{14}\xi_{23}\equiv\xi_{13}\xi_{24}-\xi_{12}\xi_{34}
\pmod{\mathfrak N}
\end{equation}
에 따라 \(\xi_{ik}\)들의 각 거듭제곱곱은 \(\mathfrak N\)을 법으로 할 때 곱 \(\xi_{14}\xi_{23}\)을 포함하지 않는 거듭제곱곱들의 정수 계수 선형결합과 합동이다. 따라서 각 다항식 \(F(\xi_{ik})\)에 대하여
\begin{equation}
\tag{5}
F(\xi_{ik})\equiv F_0(\xi_{ik})\pmod{\mathfrak N}
\end{equation}
이고, 여기서 \(F_0\)는 곱 \(\xi_{14}\xi_{23}\)을 포함하지 않는다. \emph{\(F_0\)를 \(F\)의 잉여류의 정규화 다항식이라 부르겠다.} (4)가 각 지표에 관하여 동차이므로, 각 지표에 관하여 동차인 \(F\)에는 같은 성질의 \(F_0\)가 대응한다.

이제
\begin{equation}
\tag{6}
F_0(\xi_{ik})=G(\xi_{ik})+\xi_{34}\cdot F_1(\xi_{ik})
\end{equation}
로 두자. 여기서 \(G\)는 \(\xi_{34}\)로 나누어지는 거듭제곱곱을 하나도 포함하지 않는다. 또
\begin{equation}
\tag{7}
\left[G\right]_{\xi_{14}=\xi_{13},\,\xi_{24}=\xi_{23}}
=K(\xi_{12},\xi_{13},\xi_{23})
\end{equation}
로 둔다.

\[
F((ik))\equiv0\pmod p
\qquad\text{\(\alpha,\beta\)에 관한 항등식}
\]
에서 (5)와 (2)에 따라
\begin{equation}
\tag{8}
F_0((ik))\equiv0\pmod p
\qquad\text{\(\alpha,\beta\)에 관한 항등식}
\end{equation}
이 따라온다. 여기에 \(\alpha_4=\alpha_3,\ \beta_4=\beta_3\)을 특수화하고 (6), (7)을 사용하면
\[
K((ik))\equiv0\pmod p
\qquad\text{\(\alpha,\beta\)에 관한 항등식}
\]
이다. 따라서 a)에서 증명한 결과에 따라
\[
K(\xi_{ik})\equiv0\pmod p
\qquad\text{\(\xi_{ik}\)에 관한 항등식}
\]
이고, 여기서 아래에서 보충하여 보일 바와 같이
\[
G(\xi_{ik})\equiv0\pmod p
\qquad\text{\(\xi_{ik}\)에 관한 항등식}
\]
이 따라온다. 그러므로 (6)에 따라
\[
F_0(\xi_{ik})\equiv\xi_{34\cdot}F_1(\xi_{ik})\pmod p
\qquad\text{\(\xi_{ik}\)에 관한 항등식}
\]
이다. \((34)\not\equiv0\pmod p\)이므로 (8)에서
\[
F_1((ik))\equiv0\pmod p
\qquad\text{\(\alpha,\beta\)에 관한 항등식}
\]
도 따라온다. 여기서 \(F_1(\xi_{ik})\)는 다시 정규화되어 있고 \(\xi_{34}\)에 관한 차수가 더 낮다. 이 절차를 유한번 반복하면
\[
F_0(\xi_{ik})\equiv\xi_{34}^\lambda F_\lambda(\xi_{ik})\pmod p
\]
를 얻는다. \(F_\lambda\)의 \(\xi_{34}\)에 관한 차수는 0이므로 방금 증명한 바에 따라 \(p\)를 법으로 사라진다. 따라서
\begin{equation}
\tag{9}
\begin{array}{rcl}
F_0(\xi_{ik})&\equiv&0\pmod p
\quad\text{\(\xi_{ik}\)에 관한 항등식},\quad\text{또는}\\
F(\xi_{ik})&\equiv&0\pmod{(\mathfrak N,p)}
\quad\text{\(\xi_{ik}\)에 관한 항등식}
\end{array}
\end{equation}
이다. \emph{이로써 (2)의 역, 따라서 (1)이 증명되었다.} 동시에 \(F_0(\xi_{ik})\)에 관한 더 강한 명제도 얻었다\srcfn{1)}{이 더 강한 명제는 각 잉여류가 정규화 다항식을 정확히 하나만 포함한다는 뜻이다.}.

\(K(\xi_{ik})\equiv0\pmod p\)에서 \(G(\xi_{ik})\equiv0\pmod p\)이 따라옴, 곧 \(G(\xi_{ik})\not\equiv0\pmod p\)에서 \(K(\xi_{ik})\not\equiv0\pmod p\)이 따라옴을 보충해야 한다. \emph{귀납추론의 이 부분에서만 정규화가 필요하다.} 치환 \(\xi_{14}=\xi_{13},\,\xi_{24}=\xi_{23}\)으로 \(G\)의 개별 거듭제곱곱은 사라지지 않는다. 따라서 \(K\)가 사라지려면 \(G\)의 서로 다른 거듭제곱곱들이 \(K\)의 같은 거듭제곱곱에 대응해야 한다. 이는 그 거듭제곱곱들이 지표 3과 4의 맞바꿈만으로 서로 달라질 때에만 가능하다. 예컨대 가정에 따라 정규화된 한 거듭제곱곱이
\[
\xi_{12}^{\tau_{12}}\xi_{13}^{\tau_{13}}
\xi_{14}^{\tau_{14}}\xi_{24}^{\tau_{24}}
\]
로 주어졌다고 하자. 각 지표에 관한 동차성 때문에 한 \(\xi_{ik}\)에서 3을 4로 바꾸면 다른 \(\xi_{ik}\)에서는 4를 3으로 바꾸어야 한다. 따라서 소거될 수 있는 항은 \(\xi_{13}\xi_{24}\) 대신 곱 \(\xi_{14}\xi_{23}\)을 포함해야 하지만, 정규화에 따라 이는 배제된다. 완전히 같은 방식으로
\[
\xi_{12}^{\tau_{12}}\xi_{13}^{\tau_{13}}
\xi_{23}^{\tau_{23}}\xi_{24}^{\tau_{24}}
\]
와 소거될 수 있는 항도 인수 \(\xi_{14}\xi_{23}\)을 포함해야 하므로 정규화에 따라 배제된다. \emph{따라서 \(K(\xi_{ik})\equiv0\pmod p\)에서 \(G(\xi_{ik})\equiv0\pmod p\)이 따라온다.}

\textbf{c) \(n\)개 열 \(\alpha_i,\beta_i\)인 경우.}
이제 네 열의 경우를 그대로 일반화하여 완전귀납법으로 처리할 수 있다\srcfn{2)}{이 증명은 \(n=3\)인 경우에도 성립한다.}.

먼저 각 다항식 \(F\)에 대하여
\begin{equation}
\tag{10}
F(\xi_{ik})\equiv F_0(\xi_{ik})\pmod{\mathfrak N}
\qquad\text{\(\xi_{ik}\)에 관한 항등식}
\end{equation}
인 합동식이 성립하고 \(F_0(\xi_{ik})\)가 정규화되어 있음을 다시 보이겠다. 여기서 \(\mathfrak N=(f_{ik,rs})\)이고, 정규화는 다음을 뜻한다. \emph{\(i,k,r,s\)가 서로 다른 네 지표이고 \(i<k<r<s\)이면 \(F_0\)는 곱 \(\xi_{is}\xi_{kr}\)을 포함하지 않는다. 즉 한 \(\xi\)의 두 지표가 다른 \(\xi\)의 두 지표 사이에 함께 놓이지 않는다.} 두 열이나 세 열인 경우에는 이 조건이 아무 제한도 주지 않으므로 모든 다항식이 정규화되어 있다. 네 열인 경우에도 이미 보았듯 각 잉여류에는 정규화 다항식이 존재한다. 따라서 \emph{\((n-1)\)개 열의 각 잉여류에 정규화 다항식이 존재함}을 귀납가정으로 둘 수 있다.

\(F\)의 임의의 거듭제곱곱을
\[
\xi_{i_1n}\cdots\xi_{i_\sigma n}\,P
\]
꼴로 쓰자. 여기서 \(P\)에는 지표 \(n\)이 더는 나타나지 않는다. 귀납가정에 따라 \(P\)는 \(\mathfrak N\)을 법으로 한 정수 계수 정규화 다항식 \(P_0\)와 합동이다.
\[
\xi_{i_1n}\cdots\xi_{i_\sigma n}P
\equiv
\xi_{i_1n}\cdots\xi_{i_\sigma n}P_0
\pmod{\mathfrak N}.
\]
\(P_0\)가 다음 성질의 \(\xi_{rs}\)를 하나도 포함하지 않는다고 하자. 적어도 하나의 \(i\)에 대하여 \(r,s\)가 \(i_\lambda\)과 \(n\) 사이에 놓이는 것, 곧 \(i_\lambda<r<s<n\)인 것이다. 그러면 \(P_0\) 안의 임의의 두 \(\xi\)의 곱은 이미 정규화 조건을 만족하므로 정규화가 끝난다.

그렇지 않으면 \(i_\lambda<r<s<n\)인 \(i_\lambda\)를 택하고, \(P_0^*\)를 \(P_0\)에서 \(\xi_{rs}\)를 포함하는 거듭제곱곱이라 하자. 관계
\[
\xi_{i_\lambda n}\xi_{rs}\equiv
\xi_{i_\lambda s}\xi_{rn}-\xi_{i_\lambda r}\xi_{sn}
\pmod{\mathfrak N}
\]
에 따라
\begin{align*}
&\xi_{i_1n}\cdots\xi_{i_\lambda n}\cdots\xi_{i_\sigma n}P_0^*\\
&\equiv
\xi_{i_1n}\cdots\xi_{rn}\cdots\xi_{i_\sigma n}Q
+
\xi_{i_1n}\cdots\xi_{sn}\cdots\xi_{i_\sigma n}R
\pmod{\mathfrak N}\\
&\equiv
\xi_{i_1n}\cdots\xi_{rn}\cdots\xi_{i_\sigma n}Q_0
+
\xi_{i_1n}\cdots\xi_{sn}\cdots\xi_{i_\sigma n}R_0
\pmod{\mathfrak N}.
\end{align*}
여기서 \(Q_0,R_0\)는 다시 \(Q,R\)의 잉여류의 정규화 다항식들이다. \(Q,R\)에는 지표 \(n\)이 없으므로 귀납가정에 따라 \(Q_0,R_0\)가 존재한다. 또한 \(Q,R\)가 단순한 거듭제곱곱이므로 이들은 확실히 정수 계수 다항식이다. \(i_\lambda<r<s<n\)이므로
\[
((n-i_1)+\cdots(n-r)+\cdots(n-i_\sigma))
<
((n-i_1)+\cdots(n-i_\lambda)+\cdots(n-i_\sigma)),
\]
\[
((n-i_1)+\cdots(n-s)+\cdots(n-i_\sigma))
<
((n-i_1)+\cdots(n-i_\lambda)+\cdots(n-i_\sigma)).
\]
따라서 \(P_0^*\)에서 \(Q_0,R_0\)로 넘어가면 이 차이합이 작아진다. 이 합은 반드시 \(\ge \sigma\) 이상이므로 각 거듭제곱곱 \(P_0^*\)에 대한 절차는 유한 단계 뒤에 끝난다. 그러므로
\[
\begin{aligned}
F(\xi_{ik})
&=\sum \xi_{i_1n}\cdots\xi_{i_\sigma n\cdot}P^{(\nu)}\\
&\equiv
\sum \xi_{j_1n}\cdots\xi_{j_\sigma n}\cdot S_0^{(\nu)}
=F_0(\xi_{ik})\pmod{\mathfrak N}
\end{aligned}
\]
을 얻는다. 정수 계수 다항식 \(F_0(\xi_{ik})\)는 반드시 정규화되어 있다. 그렇지 않으면 차이합을 더 낮출 수 있기 때문이다. 이로써 (10)이 증명되었다. 더 나아가 \((n-1)\)개 지표에서 각 지표에 관하여 동차인 \(F\)에 같은 성질의 정규화 다항식 \(F_0\)가 대응한다는 귀납가정 아래, 위 절차는 \(n\)개 지표에서도 그러함을 보인다. 따라서 이하에서는 각 지표에 관하여 동차인 다항식만 다룬다.

b)에서와 비슷하게 다시
\begin{equation}
\tag{11}
F_0(\xi_{ik})=G(\xi_{ik})+\xi_{n-1,n}F_1(\xi_{ik})\srcfnmark{1)}
\end{equation}
\srcfntext{1)}{(10), (11)에 대응하는 항등식
\[
F((ik))=F_0((ik))=G((ik))+(n-1,n)\cdot F_1((ik))
\]
은 두 열 \((n-1)\), \((n)\)을 변수열 \(x,y\), 따라서 \((n-1,n)\)을 \((xy)\)로 바꾸면 클렙슈--고르단 급수전개에 대응한다. 클렙슈--고르단에서는 극형에 따라 전개하지만, \(F_0\)의 정규화는 극형의 \glqq{}선두항\grqq{}에 대응한다. 이 때문에 클렙슈--고르단에서는 성립하지 않는 정수성을 강제할 수 있다. 정수성을 무시하면 \(f_{ik,rs}\)들이 \(\mathfrak M\)의 기저를 이룬다는 잘 알려진 증명은 바로 클렙슈--고르단 급수전개를 사용한다. Gordan--Kerschensteiner, \emph{Vorlesungen über Invariantentheorie}, 132쪽.}
로 둔다. 여기서 \(G\)는 \(\xi_{n-1,n}\)으로 나누어지는 거듭제곱곱을 하나도 포함하지 않는다. 또
\begin{equation}
\tag{12}
\left[G\right]_{\xi_{in}=\xi_{i,n-1}}=K(\xi_{ik})
\end{equation}
로 둔다. 정규화의 정의에서 알 수 있듯 \(G\)뿐 아니라 \(K\)도 정규화되어 있다.

그러면 \(G\)와 \(K\) 사이에는 다시 \emph{일대일 대응}이 성립한다. 실제로 \(G\)의 서로 다른 거듭제곱곱들이 \(K\)에서 같은 거듭제곱곱에 대응할 수 있는 것은 오직 지표 \(n\)과 \((n-1)\)의 맞바꿈으로만 서로 다를 때이다. 예컨대 \(G\)의 지표 \(n-1\)에 관한 차수가 \(\varrho\), 지표 \(n\)에 관한 차수가 \(\sigma\)라 하고,
\[
u_1\le u_2\le\cdots\le u_\varrho
\le u_{\varrho+1}\le\cdots\le u_{\varrho+\sigma}
\]
를 \(G\)의 한 거듭제곱곱에서 \(n-1\) 및 \(n\)과 연결된 지표들이라 하자. (11)에 따라 부정원 \(\xi_{n-1,n}\)이 \(G\)에 나타나지 않으므로 그 개수는 \(\varrho+\sigma\)이다. 정규화에 따라 그 거듭제곱곱은
\[
\xi_{\varkappa_1,n-1}\cdots\xi_{\varkappa_\varrho,n-1}
\xi_{\varkappa_{\varrho+1},n}\cdots\xi_{\varkappa_{\varrho+\sigma},n}
P(\xi_{ik})
\]
꼴이고, \(P(\xi_{ik})\)에는 지표 \(n,n-1\)이 더는 나타나지 않는다. \(i\ne j\)일 때 \(n\)과 \(n-1\)을 맞바꾸면 인수 \(\xi_{in}\xi_{j,n-1}\), \(i<j\)가 생긴다. 그러면 \(j,n-1\)이 \(i,n\) 사이에 놓여 그 항은 정규화되어 있지 않다. 따라서 고정된 \(u_1,u_2,\ldots,u_{\varrho+\sigma}\)와 고정된 \(P\)에 대하여 \(G\)는 거듭제곱곱을 \emph{하나만} 포함한다. 이로써 \(G\)의 서로 다른 거듭제곱곱들은 \(K\)의 서로 다른 거듭제곱곱에 대응함이 증명되었다.
\begin{equation}
\tag{13}
K(\xi_{ik})\equiv0\pmod p
\quad\Longrightarrow\quad
G(\xi_{ik})\equiv0\pmod p.
\end{equation}

\[
F((ik))\equiv0\pmod p
\qquad\text{\(\alpha,\beta\)에 관한 항등식}
\]
에서 (10)과 (2)에 따라
\begin{equation}
\tag{14}
F_0((ik))\equiv0\pmod p
\qquad\text{\(\alpha,\beta\)에 관한 항등식}
\end{equation}
이 따라온다. 여기에 \(\alpha_n=\alpha_{n-1},\ \beta_n=\beta_{n-1}\)을 특수화하고 (11), (12)를 쓰면
\begin{equation}
\tag{15}
K((ik))\equiv0\pmod p
\qquad\text{\(\alpha,\beta\)에 관한 항등식}
\end{equation}
이다. \(K(\xi_{ik})\)는 정규화되어 있고 지표가 \(n-1\)개뿐이다. 따라서 a)와 b)의 (9)에서 증명한 결과를 (15)에 적용하면
\[
K(\xi_{ik})\equiv0\pmod p
\qquad\text{\(\xi_{ik}\)에 관한 항등식}
\]
이라고 둘 수 있고, (13)에 따라
\[
G(\xi_{ik})\equiv0\pmod p
\qquad\text{\(\xi_{ik}\)에 관한 항등식}
\]
이 따라온다. 그러므로 (11)에 따라
\[
F_0(\xi_{ik})\equiv\xi_{n-1,n}F_1(\xi_{ik})\pmod p
\qquad\text{\(\xi_{ik}\)에 관한 항등식}
\]
이다. \((n,n-1)\not\equiv0\pmod p\)이므로 (14)에서
\[
F_1((ik))\equiv0\pmod p
\qquad\text{\(\alpha,\beta\)에 관한 항등식}
\]
도 따라온다. 여기서 \(F_1(\xi_{ik})\)는 다시 정규화되어 있고, \(\xi_{n-1,n}\)에 관한 차수가 \(F_0\)보다 낮다. 이 절차를 유한번 반복하면
\[
F_0(\xi_{ik})\equiv\xi_{n-1,n}^{\lambda}F_\lambda(\xi_{ik})\pmod p
\]
이고, \(F_\lambda\)의 \(\xi_{n-1,n}\)에 관한 차수는 0이므로 방금 증명한 결과에 따라 \(p\)를 법으로 사라진다. 따라서
\begin{equation}
\tag{16}
\begin{array}{rcl}
F_0(\xi_{ik})&\equiv&0\pmod p
\quad\text{\(\xi_{ik}\)에 관한 항등식},\quad\text{또는}\\
F(\xi_{ik})&\equiv&0\pmod{(\mathfrak N,p)}
\quad\text{\(\xi_{ik}\)에 관한 항등식}
\end{array}
\end{equation}
이다. 이로써 (2)의 역, 따라서 (1)이 증명되었다. 동시에 귀납추론에서 사용한 \(F_0(\xi_{ik})\)에 관한 더 강한 명제도 \(n\)개 지표의 경우에 증명되었으므로 일반적으로 성립한다.

\bigskip
\begin{center}
\textbf{§ 4. 유한군의 정수 계수 불변식의 유한성}
\end{center}

나의 글 \glqq{}유한군 불변식의 유한성 정리\grqq{}\srcfn{1)}{\emph{Math. Ann.} 77, 89쪽 (1915).}에서 제시한 초등적인 유한성 증명은 § 1의 특수 유한성 정리 III, IV를 사용하여 정수성까지 강화할 수 있다. 반면 힐베르트의 모듈 기저 정리에 근거한 통상적 증명은 이 정수성을 주지 않는다. 정수 계수 모듈 기저를 사용해도 군의 위수인 정수 \(h\)의 임의로 높은 거듭제곱이 분모에 나타나는 표현만 얻기 때문이다.

이하에서는 위 글의 기호를 그대로 따른다. 유한군 \(\mathfrak H\)가 행렬식이 0이 아닌 \(h\)개의 선형변환 \(A_1,\ldots,A_h\)로 이루어진다고 하자\srcfn{2)}{\(\mathfrak H\)가 한 추상 유한군과 한 단계 동형이라고 가정하는 것으로도 충분하다. 즉 \(A\) 가운데 어떤 변환행렬식이 사라지면 \(\mathfrak H\)가 \(\begin{pmatrix}\mathfrak S&0\\0&0\end{pmatrix}\) 꼴의 군과 상사이고, \(\mathfrak S\)의 치환행렬식들은 사라지지 않는다고 하면 된다. 이 경우 \(\mathfrak H\)의 불변식은 \(\mathfrak S\)의 불변식과 일치하므로 행렬식의 비소멸 조건은 실제로 일반성을 제한하지 않는다.}. \(A_k\)는 선형변환
\[
x_1^{(k)}=\sum_{\nu=1}^{n}a_{1\nu}^{(k)}x_\nu,
\quad\ldots,\quad
x_n^{(k)}=\sum_{\nu=1}^{n}a_{n\nu}^{(k)}x_\nu
\]
또는 줄여서 \((x^{(k)})=A_k(x)\)를 나타낸다. 따라서 군 \(\mathfrak H\)는 원소 \(x_1,\ldots,x_n\)으로 이루어진 열 \((x)\)를 원소 \(x_1^{(k)},\ldots,x_n^{(k)}\)으로 이루어진 열 \((x^{(k)})\)들로 보낸다. 계수 \(a_{i\nu}^{(k)}\)는 \emph{대수적 정수}라고 가정한다. J. 슈어의 정리\srcfn{3)}{\emph{Über Gruppen linearer Substitutionen mit Koeffizienten aus einem algebraischen Zahlkörper}, 정리 VII, Math. Ann. 71 (1912), 355쪽.}에 따라 각 유한 선형변환군은 \(a_{i\nu}^{(k)}\)들이 대수적 정수가 되는 군과 상사이므로, 이는 일반성을 제한하지 않는다.

군의 \emph{정수 계수 (절대)불변식}이란 \(x_1,\ldots,x_n\)의 다항함수로서 \(a_{i\nu}^{(k)}\)들이 생성하는 대수적 수체 \(\mathfrak K\)의 \emph{대수적 정수}들을 계수로 가지며 \(A_1,\ldots,A_h\)를 적용해도 항등적으로 바뀌지 않는 것을 뜻한다고 하자. 그러한 불변식 \(f(x)\)에 대해서는
\begin{equation}
\tag{1}
f(x)=f(x^{(1)})=\cdots=f(x^{(h)})
=\frac1h\cdot\sum_{k=1}^{h}f(x^{(k)})
\end{equation}
이다. 역으로 \(h\)개 변수열 \((x^{(k)})\)의 각 정수 계수 대칭함수도 군의 정수 계수 불변식이 된다.

\(\omega_1,\omega_2,\ldots,\omega_\varrho\)를 \(\mathfrak K\)의 모든 대수적 정수의 기저라 하자. 그러면 각 불변식은
\[
f(x)=\omega_1f_1(x)+\cdots+\omega_\varrho f_\varrho(x)
\]
로 나타낼 수 있고, \(f_1(x),\ldots,f_\varrho(x)\)는 유리정수 계수를 갖는다. (1)에 따라
\[
h.\,f(x)=
\omega_1\sum_{k=1}^{h}f_1(x^{(k)})+\cdots+
\omega_\varrho\sum_{k=1}^{h}f_\varrho(x^{(k)})
\]
이다. 여기서 각 합은 유리정수 계수를 갖는 \(h\)개 변수열 \((x^{(k)})\)의 대칭함수이다. 따라서 정리 III에 따라, 이 합들은 변수열들의 유한개의 유리정수 계수 대칭함수에 관한 유리정수 계수 다항함수로 나타난다. 이 기저함수 \(I_1(x),\ldots,I_\sigma(x)\)는 위의 결과에 따라 군의 정수 계수 불변식이 된다. 다만 \(\mathfrak H\)가 일반적으로 \(A\)와 함께 대수적으로 켤레인 모든 변환을 포함하지는 않으므로 그 계수는 일반적으로 \emph{대수적 정수}이다.

따라서
\begin{equation}
\tag{2}
h\cdot f(x)=\omega_1G_1(I)+\cdots+\omega_\varrho G_\varrho(I)
\end{equation}
인 표현이 성립하고, \(G_1,\ldots,G_\varrho\)는 유리정수 계수를 갖는다.

이제 \(w_1,\ldots,w_\varrho\)를 부정원이라 하고,
\[
L=w_1G_1(I)+\cdots+w_\varrho G_\varrho(I)
\]
가 \(w_i=\omega_i/h\)라는 치환으로 표현 (2)의 불변식 \(f(x)\)에 대응하는 모든 \(w\)의 선형형식을 두루 취하게 하자. 힐베르트의 정수 계수 모듈 기저 정리에 따라
\begin{equation}
\tag{3}
L=A_1(I)L_1+\cdots+A_\tau(I)L_\tau
\end{equation}
이고, \(A_i\)들은 유리정수 계수를 갖는다. 모든 \(L\)이 \(w\)에 관하여 선형이므로 \(A_i\)들에는 \(w\)가 더는 나타나지 않는다. 치환 \(\omega_i=\frac{w_i}{h}\)에 따라 \(L_1,\ldots,L_\tau\)에 각각 불변식 \(\mathfrak I_1(x),\ldots,\mathfrak I_\tau(x)\)가 대응한다고 하자. 이 치환으로 (3)은
\begin{equation}
\tag{4}
f(x)=A_1(I)\cdot\mathfrak I_1(x)+\cdots+
A_\tau(I)\mathfrak I_\tau(x)
\end{equation}
으로 넘어간다. 표현 (4)는
\[
I_1(x),\ldots,I_\sigma(x);
\quad\mathfrak I_1(x),\ldots,\mathfrak I_\tau(x)
\]
\emph{가 군 \(\mathfrak H\)의 모든 정수 계수 불변식의 기저를 이루며, 각 불변식은 이 유한개의 불변식에 관한 유리정수 계수 다항함수로 나타남}을 보여 준다.

특히 군 \(\mathfrak H\)가 각 변환 \(A\)와 함께 대수적으로 켤레인 계수를 갖는 모든 변환 \(A',A'',\ldots\)도 포함한다고 하자. 그러면 \(h\)개 열 \((x^{(k)})\)의 유리정수 계수 대칭함수들은 \(x\)의 유리정수 계수 다항함수가 되고, 대칭함수의 기저함수들도 그러하다. 따라서 유리정수 계수를 갖는 불변식 \(f(x)\)들로 제한하면 기저함수들
\[
I_1(x),\ldots,I_\sigma(x);
\quad\mathfrak I_1(x),\ldots,\mathfrak I_\tau(x)
\]
도 유리정수 계수를 갖는다. 그러므로 일반적으로 군 \(\mathfrak H\)가 만족하지는 않는 이 특수한 가정 아래에서는 유리정수 계수 불변식들의 부분영역에 대해서도 정수 계수 유한성이 성립한다.

\clearpage
\section*{논문 15에 대한 편집 주석}
\addcontentsline{toc}{section}{논문 15에 대한 편집 주석}
\begin{enumerate}
\item 원문 149쪽의 (7)과 150쪽의 이어지는 증명. \(\alpha_4=\alpha_3,\ \beta_4=\beta_3\)를 대입하면 \(\xi_{i4}=\xi_{i3}\)과 \(\xi_{24}=\xi_{23}\)을 모두 얻는다. \(K\)의 정의에는 두 조건이 필요하다. 원문은 두 곳 모두 첫 번째 조건만 제시한다. \(G=\xi_{24}\)로 놓으면 첫 번째 조건만으로는 충분하지 않다.
\item 원문 150쪽의 두 번째 별도 표시 단항식. 원문은 첫 번째 단항식을 반복한다. 두 번째에서는 \(\xi_{14}^{\tau_{14}}\) 대신 \(\xi_{23}^{\tau_{23}}\)을 쓴다. 정규화는 인수 \(\xi_{14}\xi_{23}\)을 배제한다. 두 단항식은 각각 \(\xi_{23}\)이 없는 경우와 \(\xi_{14}\)가 없는 경우에 해당한다.
\item 원문 156쪽, (3)과 (4) 사이의 치환. 155쪽과 마찬가지로 관계는 \(w_i=\omega_i/h\)이며, 원문은 여기서 이를 거꾸로 적었다. (2)의 \(hf=\sum_i\omega_iG_i\)에서 \(L=\sum_iw_iG_i=f\)를 얻는다. 원문에 적힌 역방향 관계로는 \(L=h^2f\)가 된다.
\end{enumerate}
\smallskip
\noindent{\small AI가 제작한 교정 주석과 정정: OpenAI Codex, GPT-6 Astra, Ultra 추론 강도.}

\clearpage
\setcounter{footnote}{0}
\phantomsection
\addcontentsline{toc}{section}{16. 형식론의 급수전개에 관하여}
%
%
%
%

\section*{16. 형식론의 급수전개에 관하여}
\begin{center}
{\large 형식론의 급수전개에 관하여.}\\[0.5em]
\textbf{(논문 ``임의로 많은 기본형식들의 계의 불변식에 대한\\
다항식적 표현에 관하여'', Math. Ann. 77, 93쪽의 III에 대한\\
추가와 정정.)}\\[0.65em]
저자\\[0.25em]
괴팅겐의 Emmy Noether.\\[0.65em]
Math. Ann. 81 (1920), 25--30쪽
\end{center}

잘 알려져 있듯이 ``형식론의 급수전개''에는 두 가지가 있다. 첫째는 클렙슈--고르단 급수전개를 각각 \(n\)개의 변수를 갖는 임의로 많은 변수열의 형식으로 일반화하는 것이다. 클렙슈--고르단 급수전개에서는 두 이진 변수열을 포함하는 형식을 두 변수열의 행렬식의 거듭제곱들에 따라 전개하며, 그 계수들은 하나의 변수열만을 갖는 형식들의 편극형이다. 같은 뜻으로, 그 계수들은 \(\Omega\)-연산을 적용하면 소멸한다. 둘째는 ``콤플렉스 좌표'' \(t_{ik}\)에서 나타나는 정규형을 일반화하여, 서로 다른 단계의 양 \(t_i, t_{ik}, t_{ikl}\ldots\)를 동시에 포함하는 형식들을 정규형들에 따라 전개하는 것이다. 그러나 이 두 번째 전개는 불변 미분연산으로 첫 번째 전개에서 얻을 수 있다. 특히 Mertens와 Deruyts가 이를 전개하였다.\footnote{4원 경우에는 F. Mertens, ``Über invariante Gebilde quaternärer Formen'', Sitzber. d. Ak. d. Wiss. Wien 98, Abt. IIa (1889). 일반 정규형은 J. Deruyts, ``Sur la réduction des fonctions invariantes dans le système des variables géométriques'', Bulletins de l'Ac. R. de Belgique 24 (1892)에 있다. 나는 이 논문을 알지 못한 채 Mertens를 정확히 따라 ``Zur Invariantentheorie der Formen von \(n\) Variabeln'', J. f. M. 139 (1910), § 7에서 정규형에 따른 일반 전개를 제시했다.}

첫째 종류의 여러 전개는 모두, 이를테면 Math. Ann. 77의 앞서 인용한 논문 III에서 간단히 개괄했듯이, 각각 \(n\)개의 변수를 갖는 \(n\)개 변수열의 형식을 그 변수열들의 행렬식의 거듭제곱에 따라 전개하는 데서 나온다. 이때 계수들은 \emph{\((n-1)\)개 변수열만을 갖는 형식들의 편극형이며, 그 형식들 자체도 원래 형식에서 편극연산과 \(\Omega\)-연산으로 유도된다}.\footnote{앞의 논문에 든 참고문헌을 보라. 그 뒤에 출판되었고 역시 형식적 방법에 바탕을 둔 논문으로 Schouten, ``Over reeksontwikkelingen van algebraische vormen met verschillende rijen van variabelen van verschillenden graad'', Verslag Ak. Amsterdam 27 (1919), 1481쪽이 있다.} 원래의 클렙슈--고르단 전개, 곧 \((n=2)\)의 경우는 한 걸음 더 나아가 나타나는 특수한 편극연산들을 수치계수까지 포함하여 \emph{명시적으로} 제시하였다.

이제 나는 위의 논문, 곧 급수전개를 선형 형식족의 문제로 좀 더 개념적으로 파악한 논문에, 수치계수를 제외하고 \emph{명시적인 표현}을 덧붙이고자 한다. 이 표현은 일반 모듈계에 관한 E. Fischer의 연구\footnote{E. Fischer, ``Über die Differentiationsprozesse der Algebra'', J. f. M. 148 (1917).}를 특수화하여 얻어진다.

\(t_{ik}\)의 정규형을 해석하면 이 포함관계에 이른다. ``콤플렉스 좌표'' \(t_{ik}\)를 ``선좌표''
\[
p_{ik}=(x_i y_k-x_k y_i)
\]
로 바꾸면, \(p_{ik}\)들 사이의 항등관계 때문에 형식 \(F(p_{ik})\)의 표현은 더 이상 유일하지 않다. \(F\)의 계수들 사이에 \(F\)에 나타나는 \(p_{ik}\)의 거듭제곱 단항곱들 사이와 같은 선형관계가 성립하도록 요구하면 유일성을 얻을 수 있다.\footnote{Clebsch, ``Über eine Fundamentalaufgabe der Invariantentheorie'', § 4, Abh. d. K. G. d. Wissensch. zu Göttingen 17 (1872) 참조.} 그러므로 이 정규형 \(\Phi\)에 대하여
\srcnumdisplay[4.0em]{[1]}{%
\begin{aligned}
F(p_{ik})&=\Phi(p_{ik}) &&[\text{\(x,y\)에 관한 항등식}x,y],\quad\text{또는}\notag\\
F(t_{ik})&\equiv\Phi(t_{ik}) &&\modu{M}
\end{aligned}}
가 성립한다. 여기서 \(M\)은 \(p_{ik}\)들 사이의 모든 항등관계의 모듈이다. 즉 \(G(p_{ik})\)가 \(x,y\)에 관하여 항등적으로 소멸하는 모든 형식 \(G(t_{ik})\)의 전체이다. \(\Phi(t_{ik})\)는 콤플렉스 좌표의 ``정규형''이고, \(M\)의 기저함수들은 \(t_{ik}\)의 이차형식이다. 이 기저함수들의 인자들에 대하여 절차를 반복하면 완전한 급수전개가 생긴다. \(t_i,t_{ik},t_{ikl},\ldots\)가 동시에 나타나는 경우에도 같은 논의가 성립한다.

편극형에 따른 전개도 비슷한 논의로 환원할 수 있다. 예를 들어 형식 \(F(x,y,z)\)의 세 삼항 변수열 \(x=(x_1,x_2,x_3),y,z\)를 두 변수열의 선형결합으로 나타낼 수 있는 변수열 \(\xi,\eta,\zeta\), 이를테면 \(\xi=x,\eta=y;\zeta=\lambda_1x+\lambda_2y\)로 바꾸면, \(F(\xi,\eta,\zeta)\)의 표현은 다시 유일하지 않다. 여기서도 계수들 사이에 \(F\)에 나타나는 \(\xi,\eta,\zeta\)의 거듭제곱 단항곱들 사이와 같은 선형관계를 부과할 수 있다. 이제 모든 항등관계의 모듈 \(M\)은 세 변수열의 행렬식 \(\Delta=(xyz)\)를 기저함수로 가지므로(앞의 논문 III, 보조정리 a)), [1] 대신
\srcnumdisplay[4.0em]{[2]}{%
\begin{aligned}
F(\xi,\eta,\zeta)&=\Phi(\xi,\eta,\zeta)&&[\text{\(x,y\)에 관한 항등식}x,y],\quad\text{또는}\notag\\
F(x,y,z)&\equiv\Phi(x,y,z)&&\modu{\Delta}
\end{aligned}}
를 얻는다. \(\Phi\)가 편극연산으로 생겨남을 보일 수 있고(앞의 논문 공식 (2) 참조), 다시 반복을 통해 완전한 급수전개를 얻는다.

따라서 두 경우 모두 정규형 \(\Phi\)의 구성은 다음 문제에 포함된다. \(F(z)\)의 부정원 \(z_i\)를 매개변수 \(\xi\)의 다항식 \(f_i(\xi)\)로 바꾸어 \(F(f_i(\xi))\)의 표현이 더 이상 유일하지 않게 되었다고 하자. 그러면 \(F\)의 계수들 사이에 \(F\)에 나타나는 \(f_i(\xi)\)의 거듭제곱 단항곱들 사이와 같은 선형관계가 성립하도록 하여 유일성을 정하고자 한다. 따라서
\srcnumdisplay[4.0em]{[3]}{%
\begin{aligned}
F(f_i(\xi))&=\Phi(f_i(\xi))&&[\text{\(\xi\)에 관한 항등식}\xi],\quad\text{또는}\notag\\
F(z)&\equiv\Phi(z)&&\modu{M}
\end{aligned}}
이며, 여기서 \(M\)은 \(f_i(\xi)\)들 사이의 모든 항등관계의 모듈이다.

E. Fischer는 앞에서 인용한 논문의 § 11, I에서 이 정규형 \(\Phi(z)\)를 수치계수를 제외하고 선형연산으로 명시적으로 전개하였다. 동시에 서론 III과 § 6 II에서는 \(M\)에 속하는 차 \(F-\Phi\)도, 모듈 기저를 안다는 가정 아래, 같은 뜻에서 명시적인 식으로 주었다. § 11 II에서 이 전개들을 결합하여 공식 [3]을 명시적인 표현으로 바꾸었다. 아직 정해지지 않은 수치계수를 구하는 문제는 같은 논문의 § 12에 따라 선형연산들의 고윳값과 고유함수를 구하는 문제(공식 82)와 같다.

1에서는 Fischer의 공식들을 편극형에 따른 급수전개의 경우로 특수화하고, 반복을 통하여 급수전개를 얻는다. 이때 \(\Phi\)를 정하는 연산은 \emph{특정한 편극연산}으로 바뀐다. 이는 위에서 말한 종래의 서술보다 훨씬 정확하다. \(F-\Phi\)를 정하는 연산에는 행렬식 \(\Delta\)와의 곱셈과 \(\Omega\)-연산이 섞여 나타난다. 통상적인 비명시적 꼴을 얻으려면 \(\Omega\Delta\)를 편극연산들로 나타낼 수 있다는 점도 사용해야 한다.

여기에 정정 하나를 덧붙인다. 나는 앞의 논문 101쪽 위에서 여섯째 줄에서, 이진 영역에서만 성립하는 항등변형
\[
\Omega(\Delta\psi)=c_1\psi+c_2\Delta\cdot\Omega(\psi)
\]
을 실수로 일반적으로 성립한다고 가정하였다. 그 때문에 [2]의 정규형 \(\Phi\)를 주는 기본공식 (2)에서 급수전개 (3)으로 넘어갔다. 따라서 그 논문의 논의로부터는 공식 (2)와 그 특수한 경우인 환원정리 (1)만 얻어진다. 차 \(F-\Phi\)의 명시적인 식이 없으므로 그 논의만으로는 급수전개를 정당화하기에 충분하지 않다.

2에서도 마찬가지로 특수화를 통하여 [1]의 정규형 \(\Phi(t_{ik})\), 더 일반적으로 정규형 \(\Phi(t_i,t_{ik},t_{ikl},\ldots)\)를 얻으며, 이는 알려진 명시적 공식들로 이어진다. 그러나 완전한 급수전개를 작성하려면 \(M\)의 모듈 기저를 알아야 하므로, 여기서는 위에서 든 문헌의 불변식론적 방법이 낫다. 이 방법은 그 지식을 요구하지 않고, 오히려 급수전개를 이용하여 모듈 기저를 주기 때문이다.

1. Fischer의 § 11, 공식 (19)와 그 뒤의 줄들에 따르면 [3]의 정규형 \(\Phi(z)\), 그 논문에서는 \(N(\zeta)\), 는
\srcnumdisplay[4.0em]{[4]}{%
\Phi(z)=(a_1S+a_2S^2+\cdots+a_rS^r)F(z)}
로 주어진다. 상수 \(a_1\ldots a_r\)는 \(f(\xi)\)의 계수영역에 속하며 \(F(z)\)의 차수에만 의존한다. 연산자 \(S\)는
\srcnumdisplay[4.0em]{[5]}{%
\begin{gathered}
S=AB;\qquad BF(z)=F(f(\xi));\\
AF(f(\xi))=\left\{\sum z_i f_i\!\left(\frac{\p}{\p\xi}\right)\right\}^{p}F(f(\xi))\footnotemark
\end{gathered}}
로 정의된다.\footnotetext{공식 (75) 참조. [3]의 경우 합 \(\alpha\)는 \(p\)차인 모든 거듭제곱 단항곱을 달리므로 (75)는 [5]로 특수화된다.}
\([2]\)에서 \(F\)를 \(z\)에 대해서만 \(p\)차인 형식으로 보면 \(f_i(\xi)\)는 \(\lambda_1x_i+\lambda_2y_i\)로 특수화된다. 따라서
\[
\begin{gathered}
BF=F(x,y,\lambda_1x+\lambda_2y),\\
SF=\left\{\sum z_i\left(\frac{\p}{\p\lambda_1}\frac{\p}{\p x_i}+
\frac{\p}{\p\lambda_2}\frac{\p}{\p y_i}\right)\right\}^{p}
F(x,y,\lambda_1x+\lambda_2y)
\end{gathered}
\]
이다. 이를 편극연산으로 쓰면(앞의 논문 II, 2와 3 및 그곳의 약기법 참조)
\srcnumdisplay[4.0em]{[6]}{%
SF(x,y,z)=\frac1{p!}\left[z\left(\frac{\p}{\p\lambda_1}\frac{\p}{\p x}+
\frac{\p}{\p\lambda_2}\frac{\p}{\p y}\right)\right]^p
\left[(\lambda_1x+\lambda_2y)\frac{\p}{\p z}\right]^p F(x,y,z)}
이다. \emph{[4]와 [6]은 원하는 \(\Phi(x,y,z)\)의 명시적인 식을 준다.} 동시에 \(\Phi(x,y,z)\)는 처음에 말한 일반형, 곧 앞의 논문 공식 (2)의 \(\sum PH\)에 포함된다. 다시 말해 두 변수열에만 의존하고 \(F\)에서 편극연산으로 얻는 식들의 편극형이다. 아직 정해지지 않은 계수 \(a_i\)는 유리수이다. 여기서는 \(f_i(\xi)\)가 정수 계수 다항식이기 때문이다. [2]를
\srcnumdisplay[4.0em]{[2a]}{%
F=\Phi+\Delta F_1}
로 쓰면, Fischer의 서론 III 및 § 6 II, 41쪽의 공식을 특수화하여
\srcnumdisplay[4.0em]{[7]}{%
F_1=(\alpha_1\Omega+\alpha_2\Omega\Delta\Omega+\cdots+\alpha_\lambda\Omega\Delta\Omega\cdots\Delta\Omega)F\footnotemark,
\qquad \left[\Omega=\left(\frac{\p}{\p x}\frac{\p}{\p y}\frac{\p}{\p z}\right)\right]}
를 얻는다.\footnotetext{Kerim은 곧 출판될 에를랑겐 학위논문에서 이 특수화된 공식을 ``관성형식''에 적용한다.} 여기서 \(\alpha_i\)들도 유리수이며 \(F\)의 차수에만 의존한다.

반복하면
\[
F_1=\Phi_1+\Delta F_2;\quad F_2=\Phi_2+\Delta F_3;\quad\ldots;\quad
F_i=\Phi_i+\Delta F_{i+1};\quad\ldots
\]
를 얻으며, 각 \(\Phi_i\)는 \(F_i\)에 대응하는 정규형이다. \(F_i\)는 \(z\)에 관하여 \((p-i)\)차이므로 [4]와 [6]에 의해
\srcnumdisplay[4.0em]{\vspace*{1.45em}[8]}{%
\begin{gathered}
\Phi_i=(a_{i1}S_i+a_{i2}S_i^2+\cdots)F_i,\\
S_iF_i=\frac1{(p-i)!}
\left[z\left(\frac{\p}{\p\lambda_1}\frac{\p}{\p x}+
\frac{\p}{\p\lambda_2}\frac{\p}{\p y}\right)\right]^{p-i}
\left[(\lambda_1x+\lambda_2y)\frac{\p}{\p z}\right]^{p-i}F_i
\end{gathered}}
이다. 공식 [7]에는
\[
F_i=(\alpha_{i1}\Omega+\alpha_{i2}\Omega\Delta\Omega+\cdots)F_{i-1},
\quad\text{그리고 반복하면}
\]
\srcnumdisplay[4.0em]{[9]}{%
\begin{gathered}
{}=(\beta_1\Omega^i+\beta_2\Omega^i\Delta\Omega+\cdots
+\beta_\rho\Omega^{k_1}\Delta\Omega^{k_2}\cdots\Delta\Omega^{k_\lambda}+\cdots)F,\\
(k_1+k_2+\cdots+k_\lambda-\lambda+1=i)
\end{gathered}}
가 대응한다. \emph{[8]과 [9]에 의해 급수전개}
\[
F=\Phi+\Delta\Phi_1+\Delta^2\Phi_2+\cdots+\Delta^i\Phi_i+\cdots+\Delta^p\Phi_p
\]
\emph{에서 위의 의미로 명시적인 표현을 얻었다.}

[7]에 항등변형 \(\Omega\Delta F=PF\)를 적용하자. 여기서 \(P\)는 편극연산을 뜻한다.\footnote{예를 들어 F. Mertens, ``Über ganze Funktionen von \(m\) Systemen von je \(n\) Unbestimmten'', Monatsh. f. Math. u. Phys. 4, 공식 (5) 참조.} 이 연산들이 \(\Omega\)-연산과 교환하므로 알려진 공식
\[
F_1=P\Omega F;\quad\ldots;\quad F_i=P\Omega^iF
\]
를 얻는다. 이것과 [8] 또는 앞의 논문의 논의를 결합하면, 그 논문 (3)과 같은 알려진 급수전개의 서술이 나온다. 각각 \(n\)개의 변수를 갖는 \(n\)개 변수열 \(x^{(1)},x^{(2)},\ldots,x^{(n-1)},z\)을 다룰 때에는 이에 대응하여
\[
f_i(\xi)=\lambda_1x_i^{(1)}+\lambda_2x_i^{(2)}+\cdots+\lambda_{n-1}x_i^{(n-1)}
\]
이고, \(\Delta\)와 \(\Omega\)도 \(n\)개 변수열에 맞추어 정의한다.

2. \(F(t_i,t_{ik},t_{ikl},\ldots)\)의 정규형에서는 \(f_i(\xi),f_{ik}(\xi),\ldots\)가 1열, 2열, \((n-1)\)열 행렬식
\[
\xi_i^{(1)},\quad (\xi^{(1)}\xi^{(2)})_{ik},\quad\ldots,\quad
(\xi^{(1)}\xi^{(2)}\cdots\xi^{(n-1)})_{i_1i_2\cdots i_{n-1}}
\]
으로 주어진다. 따라서 [5]를 특수화하면
\[
\begin{gathered}
S=AB;\qquad
BF=F\bigl(\xi_i^{(1)},(\xi^{(1)}\xi^{(2)})_{ik},\ldots\bigr),\\
ABF=\left(\sum t_i\frac{\p}{\p\xi_i^{(1)}}\right)^{\lambda_1}
\left(\sum t_{ik}\left(\frac{\p}{\p\xi^{(1)}}\frac{\p}{\p\xi^{(2)}}\right)_{ik}\right)^{\lambda_2}
\cdots F\bigl(\xi_i^{(1)},(\xi^{(1)}\xi^{(2)})_{ik},\ldots\bigr)
\end{gathered}
\]
를 얻는다. 여기서 \(\lambda_1,\lambda_2\ldots\)는 \(t_i,t_{ik},\ldots\)에 관한 차수이다.

불변식론적 논의에서 다음 합동식이 나온다.
\[
\Phi\equiv a_1SF\modu{M}.
\]
실제로 \(F\)의 계수에 선형인 모든 불변식은 \(F\)와 \(M\)의 형식들의 선형결합으로 나타낼 수 있으며, 특히 \(M\)에 속하지 않는 \(SF\)도 그렇다.\footnote{예를 들어 내가 인용한 논문 ``Zur Invariantentheorie der Formen von \(n\) Variabeln''의 § 6 및 § 7, 2에서 따른다.} 정규형의 성질을 함께 쓰면 [4] 대신 간단한 관계
\srcnumdisplay[4.0em]{[10]}{\Phi=a_1SF=a_1S\Phi}
를 얻는다. 두 번째 등식은 \(S\)를 적용하면 \(M\)의 모든 형식이 소멸하기 때문에 성립한다. 공식 [10]은 정규형 \(\Phi\) 자체가 동시에 고유함수임을 보인다. 위의 불변식론적 논의에 따라, 동시에 \(F\)의 불변식인 다른 고유함수는 없다([10]에 대해서는 Deruyts의 공식 (10), (10')도 참조).

\begin{center}
\textbf{정정.}
\end{center}

1. 논문 ``Die allgemeinsten Bereiche aus ganzen transzendenten Zahlen''(Math. Ann. 77)의 121쪽 위에서 여덟째 줄, 위에서 아홉째 줄, 아래에서 열한째 줄에서 체 \((H,Y)\)를 ``체 \((H,Y,Y',\ldots,Y^{(\lambda_0-1)})\), 여기서 \(Y',\ldots,Y^{(\lambda_0-1)}\)는 \(H\)에 대한 \(Y\)의 켤레들이다''로 바꾸어야 한다. 이에 따라 120쪽 위에서 다섯째 줄과 121쪽 위에서 아홉째 줄에서도 \((H,y)\)를 \((H,y,y',\ldots,y^{(\lambda_0-1)})\)로 바꾸어야 한다.

2. 논문 ``Gleichungen mit vorgeschriebener Gruppe''(Math. Ann. 78)에서 힐베르트의 불가약성 정리를 서술한 222쪽 아래에서 셋째 줄의 ``수체 \(\Omega\)''는 더 정확히 ``유한 대수적 수체 \(\Omega\)''라고 해야 한다. Ostrowski 선생이 이 점을 알려 주었다. 실제로 모든 대수적 수의 체에 대해서는 대수적 수 계수의 모든 방정식의 군이 항등군이다. 이에 따라 서론의 ``임의의 유리성 영역''도 유한 대수적 수체에 유한개의 매개변수를 첨가하여 얻는 체 또는 그렇게 생긴 체의 한 중간체를 뜻한다.

\begin{center}
(1919년 9월 채택.)
\end{center}

\clearpage
\setcounter{footnote}{0}

\clearpage
\setcounter{footnote}{0}
\phantomsection
\addcontentsline{toc}{section}{17. 비가환 영역의 모듈}
%
%
%
%
%
%

\providecommand{\mM}{\mathfrak M}
\providecommand{\mN}{\mathfrak N}
\providecommand{\mL}{\mathfrak L}
\providecommand{\mG}{\mathfrak G}
\providecommand{\mA}{\mathfrak A}
\providecommand{\mB}{\mathfrak B}
\providecommand{\mU}{\mathfrak U}
\providecommand{\mP}{\mathfrak P}
\providecommand{\mD}{\mathfrak D}
\providecommand{\mS}{\mathfrak S}
\providecommand{\mx}{\mathfrak x}
\providecommand{\my}{\mathfrak y}
\providecommand{\mz}{\mathfrak z}
\providecommand{\me}{\mathfrak e}
\providecommand{\mbarA}{\overline{\mathfrak A}}
\providecommand{\mbarB}{\overline{\mathfrak B}}
\providecommand{\bara}{\overline{\mathfrak a}}

\begingroup
\makeatletter
\renewcommand{\thefootnote}{\arabic{footnote}}
\renewcommand{\@makefnmark}{\hbox{\textsuperscript{\normalfont\thefootnote}}}
\long\def\@makefntext#1{\parindent 1em\noindent
  \hb@xt@2.4em{\normalfont\thefootnote)\hfil}#1}
\makeatother
\setcounter{footnote}{0}
\newcommand{\pseventeenkoheading}[2]{%
  \par\medskip
  \begin{center}
    \textnormal{\S~#1.}\\[.3ex]
    \textbf{#2}
  \end{center}
  \nopagebreak
}

\begin{center}
{\Large\bfseries 비가환 영역의 모듈, 특히\\
미분식과 차분식으로 이루어진 모듈}\par
\medskip
에미 뇌터와 베르너 슈마이들러, 괴팅겐\par
\medskip
\rule{2.2em}{0.25pt}
\end{center}

\begin{center}
\textbf{목차}
\end{center}
\begin{center}
\begin{tabular}{@{}r@{\hspace{0.7em}}p{0.77\linewidth}@{}}
& 서론.\\
\S~1. & 미분식과 차분식의 형식적 성질.\\
\S~2. & 일반 다항식영역.\\
\S~3. & 한쪽 모듈과 그 잉여군.\\
\S~4. & 잉여군의 두 부분군으로의 가약성과 모듈과의 관계.\\
\S~5. & 단위원을 이용한 계산과 \(k\)개 부분군으로의 가약성.\\
\S~6. & 유한성정리.\\
\S~7. & 유일분해가 성립하는 한 경우.\\
\S~8. & 완전가약군의 소군들로의 가법적 분해. 동형정리.\\
\S~9. & ‘동형’과 ‘같은 종류’라는 개념의 관계.\\
\S~10. & 한 군에 무한히 많은 분해가 존재함.\\
\S~11. & 미분방정식계와 그 적분해들과의 관계.\\
\S~12. & 예.
\end{tabular}
\end{center}

\begin{center}\textbf{서론}\end{center}

한 변수 미분식의 형식이론은 어떤 의미에서 \emph{유일한} 기약인자분해까지 이른다.\footnote{E. Landau, \emph{Ein Satz über die Zerlegung homogener linearer Differentialausdrücke in irreduzible Faktoren}, Journal für die reine und angewandte Mathematik 124 (1902), 115--120쪽과, 이에 이어지는 A. Loewy, \emph{Über reduzible lineare homogene Differentialgleichungen}, Mathematische Annalen 56 (1903), 549--584쪽을 보라.} 그러나 대수학의 대응 정리는 여러 변수의 다항식으로 옮길 수 있는 반면, 편미분식에는 더 이상 옮길 수 없다.\footnote{H. Blumberg, \emph{Über algebraische Eigenschaften von linearen homogenen Differentialausdrücken}, 괴팅겐 박사학위논문, 1912, 52쪽에 실린 Landau의 반례(51--52쪽)를 보라.} 그런데 한 변수의 미분식에는 곱표현과 나란히 최소공배수로서의 표현도 있으며, 여기에도 비슷한 법칙이 성립한다. 곧 \emph{서로 다른 두 분해의 개별 성분들을 두 개씩 서로 대응시킬 수 있고}, 그렇게 대응하는 성분들은 ‘같은 종류’이다.\footnote{앞서 든 논문들과 함께 A. Loewy, \emph{Über vollständig reduzible lineare homogene Differentialgleichungen}, Mathematische Annalen 62 (1906), 89--117쪽도 보라. 이하 Loewy로 인용한다.}

이 최소공배수 개념은 모듈개념으로 이해할 수 있고, 그렇게 하면 \(n\)개의 변수로 옮겨진다. 이 논문은 이 방식으로 위 정리들의 \emph{자연스러운} 일반화를 얻는다. 더 나아가 미분식에서 실제로 쓰는 것은 다음 두 성질뿐임도 주목한다. 미분식들을 비가환 곱셈을 갖는 다항식으로 볼 수 있고, 곱의 차수가 인자들 차수의 합과 같다는 것이다. 예컨대 차분식도 이 조건들을 만족하므로, 지금까지 한 변수에 대해서만 알려져 있던 정리들은 차분식에도 성립한다.\footnote{G. Wallenberg--A. Guldberg, \emph{Theorie der linearen Differenzengleichungen}, 라이프치히·베를린, 1911을 보라.}

따라서 우리는 일반적으로 원소가 비가환 곱셈을 갖는 다항식인 모듈들의 이론을 전개하며, \emph{모듈의 고찰을 잉여류의 고찰로 환원한다}.\footnote{가환영역의 경우에는 W. Schmeidler, \emph{Über Moduln und Gruppen hyperkomplexer Größen}, Mathematische Zeitschrift 3 (1919), 29--42쪽을 보라.} 가환인 특수한 경우와 달리 여기서는 잉여류들의 합은 말할 수 있지만 곱은 말할 수 없고, 잉여류와 다항식의 곱만 말할 수 있다. 따라서 ‘잉여군’의 개념을 가환인 경우와 다르게 고쳐야 한다. 한 변수 다항식의 특수한 경우에는 잉여군이 ‘유한’, 곧 선형독립인 잉여류를 유한개만 갖는다. 일반적으로는 그러한 유한한 ‘차수’가 존재하지 않아 군이 ‘무한’이다. 그러나 우리의 방법에서는 이 차이가 문제가 되지 않는다.

가환인 경우와 마찬가지로, 한 모듈을 서로 소인 모듈들의 최소공배수로 나타내는 데에는 잉여군을 가법적으로 분해하는 더 간단한 과정이 대응한다. 여기서 핵심이 되는 것은 잉여군들의 \emph{동형}이다. 이를 한 변수 미분식으로 특수화하면 대응 모듈들의 같은 종류라는 개념이 되고, 가환모듈로 특수화하면 동일성이 된다. 같은 종류라는 개념 자체도 \(n\)개의 변수로 옮길 수 있고 동형개념과 일치함을 보일 수 있다. 잉여군의 가법적 분해에는 언제나 \emph{유한개}의 합항만 나타난다는 사실은 고르단의 증명법에 근거한 힐베르트의 모듈기저정리를 옮겨서 보인다. 힐베르트의 원래 증명법을 쓰면 곱셈법칙에 제한이 필요할 것이다. 공식 (6)의 오른쪽에 \(a_i\) 대신 \(a\)가 와야 하는데, 이는 미분식에서는 성립하지만 차분식에서는 성립하지 않는다.

이제 \emph{유일성정리}에 이르기 위해서는 Loewy와 마찬가지로 모듈을 ‘완전가약’, 곧 소모듈들의 최소공배수로 나타낼 수 있는 모듈로 특수화해야 한다. \emph{서로 다른 두 분해가 있으면 그 성분들은 두 개씩 서로 동일하거나, 적어도 한 분해의 각 성분에 다른 분해의 동형인 성분이 대응하고 그 역도 성립한다. 이 둘째 경우에는 각 분해가 동형인 성분을 적어도 두 개씩 포함한다.} 더 나아가 같은 분해 안에 동형인 두 성분이 나타나면 무한히 많은 분해가 존재한다. 이는 소모듈로 제한하지 않아도 성립하며, 지금까지 한 변수 미분식의 특수한 경우에도 알려지지 않았던 사실이다.

끝으로 미분식으로 이루어진 모듈에서는 적분해와의 관계로 돌아간다. 유한 잉여군의 차수는 선형독립인 적분해들의 최대수와 같다. 따라서 잉여군이 유한일 필요충분조건은 일반적분에 임의함수가 나타나지 않는 것이다. 두 모듈이 같은 종류라는 것은 한 변수의 경우와 마찬가지로 그 적분해들이 ‘쌍유리적’ 관계에 있다는 뜻이다. 마지막에는 몇 가지 예를 든다. ---

이 논문을 쓰게 된 최초의 계기는 몇 해 전 E. Landau가 E. Noether에게 곱정리의 일반화를 물은 것이었다. 당시 E. Noether는 이것이 비가환 다항식영역의 모듈이론 문제라는 점만 알아보았지만, 문제에 손댈 수는 없었다. 그때 유일하게 알려져 있던 Lasker의 모듈정리들은 이 영역으로 직접 옮길 수 없었다. Lasker의 증명법은 한 다항식이 기약인자들로 유일하게 분해된다는 정리에 근거하기 때문이다.\footnote{최근 E. Noether는 여기에 쓰인 방법과 가까운 방법으로 이 정리를 쓰지 않고도 Lasker의 정리들을 정초할 수 있음을 알아보았다.} 비로소 Schmeidler의 방법들이 전이 가능성을 열어 주었고, 우리는 이를 함께 수행하였다. 세부적으로 말하면 잉여군의 일반화, 힐베르트 유한성정리의 전이와 응용, 무한히 많은 분해의 존재는 E. Noether에게서 나왔다. Loewy의 완전가약 개념의 일반화, 동형개념과 이것이 같은 종류라는 개념과 맺는 관계는 W. Schmeidler에게서 나왔다.

\pseventeenkoheading{1}{미분식과 차분식의 형식적 성질}

\textbf{1.} 한 변수 미분식에는 기호적인 곱셈을 도입해 왔다. 다음과 같이 두자.
\[
 a_0(x)\frac{\dd^n y}{\dd x^n}+a_1(x)\frac{\dd^{n-1}y}{\dd x^{n-1}}+\cdots+a_n(x)y=A(y),
\]
\[
 b_0(x)\frac{\dd^m y}{\dd x^m}+b_1(x)\frac{\dd^{m-1}y}{\dd x^{m-1}}+\cdots+b_m(x)y=B(y).
\]
그러면 기호적 곱 \(AB(y)\)를 다음 미분식으로 정의한다.
\[
 a_0(x)\frac{\dd^n B(y)}{\dd x^n}+a_1(x)\frac{\dd^{n-1}B(y)}{\dd x^{n-1}}+\cdots+a_n(x)B(y).
\]
이 곱셈은 실제로 두 연산자의 곱셈이다. 통상적으로 \(\frac{\dd^n}{\dd x^n}\)을 \(\left(\frac{\dd}{\dd x}\right)^n\)으로 바꾸면 이를 가장 간단히 쓸 수 있다. 그러면 \(AB\)는
\[
 \left(a_0\left(\frac{\dd}{\dd x}\right)^n+\cdots+a_n\right)
 \left(b_0\left(\frac{\dd}{\dd x}\right)^m+\cdots+b_m\right)y
\]
가 되며, 곱을 실행할 때에는 합과 곱의 통상적인 미분법칙을 쓴다. 연산자가 작용하는 함수를 명시적으로 쓰지 않고 변수 \(\frac{\dd}{\dd x}=\xi\)에 곱의 미분법칙에 대응하는 곱셈법칙을 부여하면, 각 연산자를 이 한 변수의 다항식으로 볼 수 있다.
\[
 \frac{\dd}{\dd x}\bigl(a(x)\cdot y\bigr)=a(x)\frac{\dd}{\dd x}(y)+a'(x)y.
\]
우리 표기에서는 (\((n=1, m=0)\))
\begin{equation}
 \xi\cdot a=a\cdot\xi+a'.
\tag{1}
\end{equation}

한 함수 \(y\)의 선형 편미분식
\[
 \left(\sum a_{\alpha_1\ldots\alpha_n}(x_1,\ldots,x_n)
 \frac{\p^{\alpha_1+\cdots+\alpha_n}}
 {\p x_1^{\alpha_1}\cdots\p x_n^{\alpha_n}}\right)y
\]
도 완전히 같은 방식으로 다룬다. \(n\)개 연산자 \(\frac{\p}{\p x_1},\ldots,\frac{\p}{\p x_n}\)에 변수 \(\xi_1,\ldots,\xi_n\)을 도입하면 다항식
\[
 \sum a_{\alpha_1\ldots\alpha_n}(x_1,\ldots,x_n)
 \xi_1^{\alpha_1}\cdots\xi_n^{\alpha_n}
\]
을 얻는다. 두 다항식의 곱은 다음 법칙으로 계산한다.
\begin{equation}
 \xi_i a=a\xi_i+a_i \qquad (i=1,\ldots,n).
\tag{2}
\end{equation}
여기서 \(a_i\)는 \(x_i\)에 관한 \(a\)의 편도함수이고, \(\xi\)들끼리의 곱은 가환적이다. 이렇게 정의한 다항식에는 곱셈의 교환법칙을 제외한 통상적인 대수학의 모든 덧셈법칙과 곱셈법칙이 성립한다. 특히 두 그러한 다항식의 곱은 전체차수와 각 개별 변수에 관한 차수가 모두 두 인자의 대응 차수의 합과 같다.

\textbf{2.} 다음 꼴의 차분식(앞서 든 Wallenberg의 책을 보라)
\[
 a_x^{(0)}y_{x+n}+a_x^{(1)}y_{x+n-1}+\cdots+a_x^{(n)}y_x=A_x,
\]
\[
 b_x^{(0)}y_{x+m}+b_x^{(1)}y_{x+m-1}+\cdots+b_x^{(m)}y_x=B_x
\]
에는 이에 대응하여 다음 곱을 정의한다.
\[
 (AB)_x=a_x^{(0)}B_{x+n}+a_x^{(1)}B_{x+n-1}+\cdots+a_x^{(n)}B_x.
\]
이 곱은 곱셈법칙
\[
 (ay)_{x+1}=a_{x+1}y_{x+1}
\]
을 이용하여 실행할 수 있다. \(x\)에서 \(x+1\)로 가는 연산자를 \(\xi\)라 하면 \(\xi\{ay\}=a^*\xi\{y\}\)이고, 여기서 \(a_x^*=a_{x+1}\)이다. 다시 연산자가 작용하는 함수 \(y\)를 생략하면 다음 법칙을 얻는다.
\begin{equation}
 \xi a=a^*\xi.
\tag{3}
\end{equation}
여기서 \(a\)와 함께 \(a^*\)도 항등적으로 사라지지 않는다. \(n\)개의 변수에서도 완전히 같은 방식으로
\begin{equation}
 \xi_i a=a_i\xi_i\qquad (i=1,\ldots,n)
\tag{4}
\end{equation}
이다. 여기서 \(\xi_i\)는 \(x_i\)에서 \(x_i+1\)로 가는 연산자이고, \(a_i\)는 \(a\)에서 \(x_i\)를 \(x_i+1\)로 바꾸어 얻은 함수이다. 따라서 \(a\)가 항등적으로 사라지지 않으면 \(a_i\)도 그러하다. 이 규정 아래에서도 교환법칙을 제외한 덧셈과 곱셈의 모든 법칙, 그리고 곱의 차수에 관한 법칙이 성립한다.

다음 절에서는 이 두 경우를 임의의 계수를 갖고 아주 일반적으로 정의한 하나의 다항식영역 아래에 포함시킨다.

\pseventeenkoheading{2}{일반 다항식영역}

임의의 추상적으로 정의된 체 \(\mathsf{P}\)에서 시작하자. 그 원소들은 작은 라틴문자로 나타내며, 이 체에는 모든 대수법칙, 특히 곱셈의 교환법칙과 유일한 역연산이 성립한다. 이 체에 \(n\)개의 부정원 \(\xi_1,\ldots,\xi_n\)을 첨가한다. 곧 곱 \(a\xi_{\nu_1}b\xi_{\nu_2}\ldots c\xi_{\nu_\rho}\)과 그 유한합들로 이루어진 정역 \(J\)를 고찰하되, 곱셈의 교환법칙을 제외한 덧셈과 곱셈의 모든 법칙이 성립한다고 한다. 특히 다음 꼴의 유한합을 ‘다항식’이라고 부르고, 언제나 큰 라틴문자로 나타낸다.
\[
 F(\xi_1,\ldots,\xi_n)=\sum a_{\nu_1\ldots\nu_\rho}\xi_{\nu_1}\cdots\xi_{\nu_\rho}.
\]
이제 교환법칙 대신 \(J\)의 원소들에 \emph{두 다항식의 곱이 다시 다항식이 되는} 곱셈법칙을 부여한다. 곧 \(J\)의 모든 다항식들의 영역이 이미 영역 \(J\) 전체가 되게 한다. 이를 위한 필요충분조건은 \(\xi_1,\ldots,\xi_n\)과 \(\mathsf{P}\)의 임의의 원소 \(a\)에
\begin{equation}
 \xi_i a=F_i(\xi_1,\ldots,\xi_n)\qquad (i=1,\ldots,n)
\tag{5}
\end{equation}
꼴의 곱셈법칙이 성립하는 것이다. 여기서 \(F_i(\xi_1,\ldots,\xi_n)\)는 \(\xi_i\)와 \(a\)에 자연스럽게 의존하는 어떤 다항식이다. 특히 미분식과 차분식에 관한 법칙 (2), (4)가 이 꼴이다. 실제로 \(\xi_i a\) 자체는 다항식이 아니지만 가정상 영역에 속하므로 이 등식이 필요하다. 역으로 공식 (5)를 유한번 적용하면 영역 \(J\)의 각 원소를 다항식으로 바꿀 수 있다.

1로 나타내는 체의 단위원에는 \(\xi_i\cdot1=1\cdot\xi_i=\xi_i\)가 성립한다고 하자. 또 \(\xi_1,\ldots,\xi_n\)끼리의 곱셈은 가환적이라고 규정한다.\footnote{이 가정은 § 6부터 유한성정리와 그 귀결들을 위해서만 쓰인다.} 그러면 \(J\)의 각 다항식은 \(\mathsf{P}\)의 계수를 갖는 \(\xi_1,\ldots,\xi_n\)의 통상적인 다항식처럼 쓰이며, 두 다항식은 모든 계수가 일치할 때에만 일치한다.

이제 이하에서는 더 특별히,\footnote{이 가정도 § 6부터만 쓰인다.} (5) 대신
\begin{equation}
 \xi_i a=a_i\xi_i+b_i\qquad (a_i\ne0)\quad(i=1,\ldots,n)
\tag{6}
\end{equation}
이 성립한다고 요구한다. 여기에도 특히 법칙 (2), (4)가 포함된다. 그러면 곱의 차수는 인자들 차수의 합과 같으며, 이는 전체차수와 각 개별 변수에 관한 차수 모두에 성립한다. 따라서 항등적으로 사라지지 않는 두 다항식의 곱도 영과 다르다.

\pseventeenkoheading{3}{한쪽 모듈과 그 잉여군}

우리 영역 \(J\)에서는 곱셈이 비가환이므로 다음 종류의 모듈들을 구별할 수 있다.\footnote{모듈은 큰 독일문자로 나타낸다.}\footnote{A. Hurwitz도 이미 한쪽 아이디얼들을 고찰하였다. \emph{Vorlesungen über die Zahlentheorie der Quaternionen}, 베를린, Springer, 1919, 제6강을 보라. 그러나 거기서는 본질적으로 주아이디얼을 다루며, 그곳의 주 1에 인용된 Du Pasquier의 논문들도 마찬가지이다.}

\textbf{1. 우측모듈}은 \(M\)과 함께 임의의 다항식 \(F\)에 대한 \(FM\)을 포함하고, \(M_1,M_2\)와 함께 \(M_1+M_2\)도 포함한다.

\textbf{2. 좌측모듈}은 \(M\)과 함께 \(MF\)를 포함하고, \(M_1,M_2\)와 함께 \(M_1+M_2\)도 포함한다.

이 두 종류를 한데 \emph{한쪽 모듈}이라고 부르며, 우측이면서 좌측인 모듈을 \emph{양쪽 모듈}이라 한다. 이하에서는 우측모듈로 제한한다. 좌측모듈의 이론도 완전히 같은 방식으로 세울 수 있다.

두 다항식 \(F,G\)의 차가 \(\mM\)으로 나누어지는, 곧 \(\mM\)에 속하는 경우 이 두 다항식을 법 \(\mM\)에 관하여 합동이라고 하고 \(F\equiv G(\mM)\)라고 쓴다. 한 다항식과 합동인 모든 다항식의 집합은 법 \(\mM\)의 한 \emph{잉여류}를 이룬다. 잉여류는 작은 독일문자로 나타낸다. 합동식
\[
 F_1\equiv F_2(\mM)\qquad\text{및}\qquad G_1\equiv G_2(\mM)
\]
에서 \(F_1+G_1\equiv F_2+G_2(\mM)\)가 따르므로, 두 잉여류 \(\mx,\my\)의 합은 정의되며 다시 잉여류 \(\mx+\my\)이다. 반면 모듈이 한쪽 모듈이므로 두 잉여류의 곱은 말할 수 없다. 실제로
\[
 F_1=F_2+M_1,\qquad G_1=G_2+M_2
\]
에서
\[
 F_1G_1=F_2G_2+F_2M_2+M_1G_2+M_1M_2
\]
가 따르기는 한다. 그러나 \(M_1\)이 \(\mM\)에 속하더라도 \(M_1G_2\)는 \(\mM\)에 속하지 않을 수 있으므로, 일반적으로 \(F_1G_1\)과 \(F_2G_2\)는 합동일 필요가 없다. 이에 비해 \(M_1=0\), 곧 \(F_1=F_2=F\)이면 \(G_1\equiv G_2\)에서 \(FG_1\equiv FG_2\)가 따른다. 따라서 임의의 잉여류 \(\my\)의 앞에 다항식 \(F\)를 곱할 수 있고, 그렇게 해서 잘 정의된 잉여류 \(F\my=\mz\)를 얻는다.

잉여류의 계산에서는 덧셈의 교환법칙과 결합법칙, 덧셈의 유일한 역연산법칙이 성립한다. 또 잉여류와 다항식의 곱과 덧셈에는 분배법칙 \(F(\my_1+\my_2)=F\my_1+F\my_2\)이 성립한다. 이는 법 \(\mM\)의 다항식 계산에서 곧바로 따른다.

단위원이 속하는 잉여류를 단위원류, 또는 단위원 \(\me\)라고 하자. 그러면 \(F\me=\mx\)는 다항식 \(F\)가 속하는 잉여류이다. 일반적인 잉여류와 달리 여기서는 곱 \(\mx\me\)도 잘 정의되며 \(\mx\)와 같다. 실제로
\[
 F_1=F_2+M_1,\qquad G_1=1+M_2
\]
이면 \(M_1\cdot1=M_1\)이므로
\[
 F_1G_1=F_2+M
\]
도 성립한다. 법 \(\mM\)의 모든 잉여류의 집합을 \emph{\(\mM\)의 잉여군}이라 한다. 따라서 한 잉여군은 다음 성질을 갖는다.

\textbf{1.} 각 잉여류 \(\mx\)와 임의의 다항식 \(F\)에 대하여 잉여류 \(F\mx\)를 포함하고, 잉여류 \(\mx_1,\mx_2\)와 함께 \(\mx_1+\mx_2\)도 포함한다.\footnote{이 성질을 더 일반적인 의미에서 잉여군의 ‘모듈 성질’이라고 부를 수도 있다.}

\textbf{2.} 잉여군에는 단위원 \(\me\)가 있어서, 각 다항식 \(F\)에 대하여 \(F\me\)가 \(F\)의 잉여류 \(\mx\)를 나타낸다. \emph{따라서 \(F\)가 모든 다항식을 돌면 \(F\me\)는 잉여군 전체를 돈다.}

\pseventeenkoheading{4}{잉여군의 두 부분군으로의 가약성과 모듈과의 관계}

다음 정의들이 보여 주듯, 한쪽 모듈에서도 \emph{나눗셈}, \emph{최대공약인자}, \emph{최소공배수}라는 개념은 양쪽 모듈의 경우와 똑같이 유지된다.

한 모듈 \(\mM\)의 \(M\equiv0(\mM)\)에서 언제나 \(M\equiv0(\mN)\)이 따르면 \(\mM\)은 다른 모듈 \(\mN\)으로 나누어진다고 한다.

두 모듈 \(\mM,\mN\)의 \emph{최대공약인자} \((\mM,\mN)\)는 \(M\equiv0(\mM)\), \(N\equiv0(\mN)\)인 \(M,N\)에 대하여 \(M+N\) 꼴로 나타나는 모든 다항식의 집합으로 정의되며, 다시 모듈이다. 여러 모듈에 대해서도 마찬가지이다.

\(\mM\)과 \(\mN\)으로 모두 나누어지는 모든 다항식의 집합은 모듈 \([\mM,\mN]\), 곧 \(\mM,\mN\)의 \emph{최소공배수}를 이룬다. 이 정의도 여러 모듈에 대하여, 나아가 임의의 기수를 갖는 모듈들의 집합에 대하여 그대로 성립한다.

두 모듈의 최대공약인자가 모든 다항식으로 이루어진 단위원모듈이면 두 모듈이 서로 소라고 한다.

이제 한 모듈 \(\mM\)이 서로 소인 두 진인자 \(\mN,\mL\)의 최소공배수로 나타난다고 하자.
\begin{equation}
 \mM=[\mN,\mL].
\tag{7}
\end{equation}
\(\mM\)의 잉여군 \(\mG\)를 조사하자. 가정에 따라, \(N\equiv0(\mN)\), \(L\equiv0(\mL)\)인 두 다항식 \(N,L\)에 대하여
\begin{equation}
 1\equiv N+L(\mM)
\tag{8}
\end{equation}
이다. 여기서 다음을 보이겠다.
\begin{equation}
 \mL=(\mM,L),\qquad \mN=(\mM,N).
\tag{9}
\end{equation}
예컨대 \((\mM,L)\)는 \(M+FL\) 꼴의 모든 다항식으로 이루어진 모듈이다. 우선 \((\mM,L)\equiv0(\mL)\)이다. 한편 (8)에서 임의의 다항식 \(F\)에 대하여
\begin{equation}
 F\equiv FN(\mL)
\tag{10}
\end{equation}
을 얻는다. 이제 \(F\equiv0(\mL)\)이면 \(FN\equiv0(\mL)\), 따라서 \(\equiv0(\mM)\)이고, 결국 \(F\equiv FL(\mM)\)이다. 같은 방식으로 \(\mN=(\mM,N)\)을 증명한다. 또 \(\mN,\mL\)이 \(\mM\)의 진인자이므로 \(L\equiv0\)도 \(N\equiv0(\mM)\)도 성립할 수 없다는 것이 여기서 따라온다.

\(\ma\)를 법 \(\mM\)의 \(N\)의 잉여류, \(\mb\)를 \(L\)의 잉여류라 하자. 그러면 (8)에 따라
\[
 \me=\ma+\mb\qquad(\ma\ne0,\ \mb\ne0).
\]
각 잉여류는 \(F\me\) 꼴이므로, 분배법칙에 따라 모든 잉여류에 대하여
\begin{equation}
 F\me=F\ma+F\mb\qquad(\ma\ne0,\ \mb\ne0)
\tag{11}
\end{equation}
이다. 임의의 잉여류를 \(G\ma\) 꼴의 잉여류와 \(H\mb\) 꼴의 잉여류의 합으로 나타내는 이 표현은 \emph{유일하다}. 실제로 관계 \(0=G\ma+H\mb\)가 있고 \(K\)가 잉여류 \(G\ma=-H\mb\)의 한 다항식이면, 분명 \(K\equiv GN(\mM)\), \(K\equiv-HL(\mM)\)이다. 따라서 \(K\equiv0(\mN)\), \(K\equiv0(\mL)\)이고, 결국 \(K\equiv0(\mM)\), 곧 \(G\ma=0\), \(H\mb=0\)이다.

역으로 이제 관계 (11)이 각 다항식 \(F\)에 대하여 위 의미에서 \emph{유일하게} 성립한다고 하자. \(N\)을 \(\ma\)의 한 다항식, \(L\)을 \(\mb\)의 한 다항식이라 하고, (9)의 두 모듈을 만들자. 그러면 (7), (8)이 성립한다. (8)은 \(F=1\)인 (11)에서 따른다. 또 \(\mN,\mL\)의 정의에 따라 \(\mM\)은 이 두 모듈의 공배수이고, 더구나 최소공배수이다. 실제로 \(F\equiv0(\mL)\), \(F\equiv0(\mN)\)이면 \(F\equiv HL\equiv GN(\mM)\)이다. 따라서 \(G\ma=H\mb\), \(G\ma-H\mb=0\)이고, 표현의 유일성 때문에 \(G\ma=0\), \(H\mb=0\), \(F\equiv0(\mM)\)이다. 이로써 \emph{한 모듈 \(\mM\)을 서로 소인 두 모듈의 최소공배수로 나타내는 것과, 공식 (11)이 주는 \(\mM\)의 잉여군 \(\mG\)의 유일한 가법적 분해가 같은 뜻}임을 보였다.

이제 \(F\ma\) 꼴 잉여류들의 계 \(\mA\)를 조사하고, 이를 \(\mL\)의 잉여군 \(\mbarA\)와 관계시키자. \(\bara\)를 법 \(\mL\)의 \(N\)의 잉여류, 따라서 \(\mbarA\)의 단위원류라고 하면, (11)에 따라 \(\mbarA\)의 각 다항식 \(F\)는 잉여류 \(F\ma\)에 속한다. \(\mA\)의 잉여류 \(F\ma\)와 \(\mbarA\)의 잉여류 \(F\bara\)를 서로 대응시키자. 그러면 \(\mA\)의 각 잉여류에 \(\mbarA\)의 잉여류가 대응하고, \(\mA\)의 두 잉여류의 합에는 대응하는 \(\mbarA\)의 잉여류들의 합이, \(\mA\)의 한 잉여류와 한 다항식의 곱에는 대응하는 \(\mbarA\)의 잉여류와 같은 다항식의 곱이 대응한다. 이 대응은 일대일이다. \(F\ma=0\)이면 (11)에 따라 \(F\equiv0(\mL)\), 따라서 \(F\bara=0\)이기 때문이다. 역으로 \(F\bara=0\), 곧 \(F\equiv0(\mathfrak L)\)이면 (10)에 따라 \(FN\equiv0(\mL)\), 따라서 \(FN\equiv0(\mM)\), 곧 \(F\ma=0\)이다.

이로써 다음과 같이 정의할 근본개념을 얻었다.

\emph{두 잉여류계 \(\mA,\mbarA\) 사이의 일대일 대응에서 \(\mA\)의 두 잉여류의 합에 \(\mbarA\)의 대응 잉여류들의 합이 대응하고, \(\mA\)의 임의의 잉여류와 한 다항식의 곱에 \(\mbarA\)의 대응 잉여류와 같은 다항식의 곱이 대응하면, 이 대응을 동형이라고 한다.}

\(\mG\) 안의 잉여류 부분계 \(\mA\)가 한 모듈의 잉여군과 동형이면 \(\mA\)를 \(\mG\)의 \emph{부분군}이라 한다. 위 증명이 보여 주듯, 여기서는 특별히 \(\bar{\mathfrak A}\)의 각 잉여류가 \(\mbarA\)의 대응 잉여류에 어떤 다항식들, 곧 \(\mL\)에는 속하지만 \(\mM\)에는 속하지 않는 모든 다항식을 덧붙여 생긴다.

마찬가지로 \(F\mb\) 꼴의 모든 잉여류가 \(\mG\)의 한 부분군 \(\mB\)를 이루고, 이는 \(\mN\)의 잉여군 \(\mbarB\)와 동형임을 보일 수 있다.

잉여류들이 유일한 가법적 분해 (11)을 갖는 잉여군 \(\mG\)를 두 부분군 \(\mA,\mB\)로 \emph{가약}이라고 하고 \(\mG=\mA+\mB\)라고 쓴다. 가약이 아닌 군을 \emph{기약}이라고 한다. 대응 모듈들도 때때로 각각 가약, 기약이라고 부른다. 따라서 다음 정리가 성립한다.

\textbf{정리 I}\footnote{앞서 든 Schmeidler의 논문, 39쪽을 보라.}. \emph{한 모듈 \(\mM\)을 서로 소인 두 진인자 \(\mL,\mN\)의 최소공배수로 나타낼 수 있을 필요충분조건은 그 잉여군 \(\mG\)가 두 부분군 \(\mA,\mB\)로 가약인 것이다. 이때 \(\mA,\mB\)는 각각 \(\mL,\mN\)의 잉여군과 동형이다.}

\pseventeenkoheading{5}{단위원을 이용한 계산과 \(k\)개 부분군으로의 가약성}

\(\mG\)의 잉여류 \(\ma,\mb\)는 단위원 \(\me\)와 마찬가지로 임의의 잉여류 \(\mx\)와의 곱이 정의된다. 따라서 \(\ma,\mb\)도 \(\mG\)의 단위원이라고 부른다. 실제로 (11)에 \(F=M\)을 넣으면
\[
 0=M\me=M\ma+M\mb
\]
이고, 유일성 때문에 \(M\ma=0\), \(M\mb=0\)이다. 따라서
\[
 (F+M)\ma=F\ma=\mx\ma
\]
이다. 여기서 \(\mx\)는 \(F\)의 잉여류이다. 그러므로 관계 (11) 대신 그와 동치인 잉여류 관계
\[
 \mx\me=\mx\ma+\mx\mb
\]
를 쓸 수 있다. 특히 \(\mx=\ma\)를 넣으면
\[
 \ma=\ma^2+\ma\mb
\]
이고, 유일성 때문에 \(\ma^2=\ma\), \(\ma\mb=0\)이다. 마찬가지로 \(\mb^2=\mb\), \(\mb\ma=0\)이다.

잉여군을 두 성분의 합으로 나타낸 것에 대응하여 \(k\)개 성분의 합도 정의할 수 있다. 곧 유일한 표현
\[
 F\me=F\ma_1+F\ma_2+\cdots+F\ma_k
 \qquad(\ma_1\ne0,\ldots,\ma_k\ne0)
\]
가 주어졌다고 하자. 여기서 \(0=G_1\ma_1+\cdots+G_k\ma_k\)이면 각 \(G_i\ma_i=0\)이다. 특히 \(F=M\)으로 두면 \(M\ma_i=0\)이므로 이 표현을 다음 꼴로도 쓸 수 있다.\footnote{잉여군의 가법적 분해에는 교환법칙과 결합법칙이 성립한다. 그러나 가환인 경우와 달리 \(\mathfrak A_1+\mathfrak A_2=\mathfrak A_1+\mathfrak A_3\)에서 \(\mathfrak A_2=\mathfrak A_3\)가 따르지는 않는다. § 12의 예 1을 보라.}
\begin{equation}
 \mx\me=\mx\ma_1+\cdots+\mx\ma_k
 \qquad(\ma_1\ne0,\ldots,\ma_k\ne0).
\tag{12}
\end{equation}
이 표현은 두 성분으로의 분해를 차례로 반복하여 생긴 것으로도 볼 수 있다. 모든 \(\ma_i\ne0\)인 채로 이를 제한 없이 계속할 수 있으면 무한히 많은 성분의 합이라고 한다.

(12)에 \(\mx=\ma_i\)를 넣으면
\begin{equation}
 \ma_i^2=\ma_i,\qquad
 \ma_i\ma_j=0\quad(i\ne j),\qquad
 (i,j=1,\ldots,k),\qquad
 (\ma_i\ne0)
\tag{13}
\end{equation}
을 얻는다. 역으로 \(\mG\) 안에 직교조건 (13)을 만족하고
\begin{equation}
 \me=\ma_1+\cdots+\ma_k
\tag{14}
\end{equation}
인 잉여류계 \(\ma_1\ldots\ma_k\)가 있다고 하자. 그러면 군을 \(k\)개 성분의 합으로 나타내는 표현이 생김을 보이겠다. 먼저 \(\ma_i\)의 임의의 다항식 \(A_i\)에 대하여 \(A_i\ma_i=(A_i+M)\ma_i\)이므로 \(M\ma_i=0\)이다. 따라서 \(\ma_i\)는 각 잉여류와 곱할 수 있다. 이제 (14)에서 표현 (12)가 따르며, 이 표현은 유일하다. 실제로 \(0=\mx_1\ma_1+\cdots+\mx_k\ma_k\)이면 오른쪽에 \(\ma_i\)를 곱하여 (13)에 따라 \(0=\mx_i\ma_i\)를 얻는다. 이로써 유일성과 주장을 증명하였다. 단위원류의 분해에 나타나는 잉여류 \(\ma_1\ldots\ma_k\)를 이 분해에 속하는 단위원들이라고 한다. 따라서 이 단위원들은 직교관계 (13)과 조건 (14)로도 정의된다. 이후의 전개는 본질적으로 단위원들을 이용한 계산으로 귀착될 것이다.

이제 군을 \(k\)개 성분으로 가법적으로 분해하는 것과 대응 모듈분해 사이의 관계를 밝혀야 한다. \(k=2\)의 경우처럼 모듈을 서로 완전히 소인 \(k\)개 모듈의 최소공배수로 나타내는 표현이 생기며, 귀납법을 쓴다.

분해 (12)가 주어졌다고 하고, 이를
\begin{align}
 \mx\me&=\mx\mb+\mx\ma_k,\tag{15}\\
 \mx\mb&=\mx\ma_1+\cdots+\mx\ma_{k-1}\tag{16}
\end{align}
로 쓰자. (15)와 정리 I에 따라 잉여류 \(\mx\mb\)는 모듈
\begin{equation}
 \mN=(\mM,A_k)
\tag{17}
\end{equation}
의 잉여군과 동형이다. 따라서 그 잉여류 \(\mx\bar{\mb}\)에는 (16)에 대응하는 표현
\begin{equation}
 \mx\bar{\mb}=\mx\bar{\ma}_1+\cdots+\mx\bar{\ma}_{k-1}
\tag{18}
\end{equation}
도 성립한다. 이제 이미 다음이 증명되었다고 하자.
\begin{equation}
 \mN=[\mM_1,\ldots,\mM_{k-1}],
\tag{19}
\end{equation}
여기서
\begin{equation}
\left\{
\begin{aligned}
 \mM_1&=(\mN,A_2+\cdots+A_{k-1}),\\
 \mM_2&=(\mN,A_1+A_3+\cdots+A_{k-1}),\\[-0.2ex]
 &\hspace{0.25cm}\cdot\quad\cdot\quad\cdot\quad\cdot\quad\cdot\quad\cdot\quad\cdot\quad\cdot\\[-0.2ex]
 \mM_{k-1}&=(\mN,A_1+\cdots+A_{k-2})
\end{aligned}\right.
\tag{20}
\end{equation}
합항이 두 개일 때 이 명제는 앞의 결과와 일치한다. 또 § 4에 따라 뒤 잉여류의 다항식들이 앞 잉여류에 포함되므로, \(A_i\)들은 \(\bar{\ma}_i\)에 속하면서 동시에 \(\ma_i\)에도 속하는 다항식들로 특별히 택할 수 있다. 그러면 (15)에 따라
\begin{equation}
 \mM=[\mN,\mM_k]
\tag{21}
\end{equation}
이고,
\[
 \mM_k=(\mM,A_1+\cdots+A_{k-1})
\]
이다. (19), (21)에서 최소공배수의 정의에 따라
\[
 \mM=[[\mM_1,\ldots,\mM_{k-1}],\mM_k]
     =[\mM_1,\ldots,\mM_{k-1},\mM_k]
\]
가 따른다. (17), (20)에 따라
\[
 \mM_1=(\mM,A_k,A_2+\cdots+A_{k-1})
\]
이고, 따라서
\[
 \mM_1=(\mM,A_2+\cdots+A_k).
\]
실제로 뒤 모듈은 (13)에 따라 다항식
\[
 A_k\equiv A_k(A_2+\cdots+A_k)(\mM)
\]
도 포함한다. \(\mM_2,\ldots,\mM_{k-1}\)에 대해서도 대응하는 표현들을 얻는다. 또 (14)에 따라
\[
 \frac{1}{k-1}\Bigl((A_2+\cdots+A_k)
 +(A_1+A_3+\cdots+A_k)+\cdots+(A_1+\cdots+A_{k-1})\Bigr)
 \equiv1(\mM)
\]
이므로 \(\mM_1,\ldots,\mM_k\)는 서로 소이다. 더 나아가 이 모듈들은 \emph{서로 완전히 소인 계}를 이룬다. 곧 그 가운데 임의로 고른 모듈들의 최소공배수는 나머지 모듈들의 최소공배수와 서로 소이다.\footnote{이는 Loewy가 도입한 개념의 일반화이다.} 공식 (12)의 합항들을 알맞은 두 무리로 묶으면 이 사실이 곧바로 따른다.

이로써 잉여군을 \(k\)개 합항으로 분해하는 데에는 모듈을 서로 완전히 소인 계를 이루는 \(k\)개 모듈의 최소공배수로 나타내는 표현이 대응함을 보였다. 이제 이 정리의 역도 귀납법으로 보이겠다.

다음과 같이 두자.
\[
 \mM=[\mM_1,\ldots,\mM_k]
     =[[\mM_1,\ldots,\mM_{k-1}],\mM_k]
     =[\mN,\mM_k],
\]
여기서 \(\mM_1,\ldots,\mM_k\)는 서로 완전히 소인 계를 이룬다. 정리 I을 함께 쓰면 잉여군에 대한 표현 (15)가 따른다. \(\mM_1,\ldots,\mM_{k-1}\)도 서로 완전히 소인 계를 이룬다. 실제로 이 계의 모듈들을 어떤 식으로든 두 무리로 나누었을 때 한 무리의 최소공배수와 다른 무리의 최소공배수가 공통인자를 가진다면, 한 무리에 \(\mM_k\)를 덧붙여 최소공배수를 만든 뒤에도 이 인자가 남을 것이다.\footnote{이 논의를 계속하면 \(\mM_1,\ldots,\mM_k\) 가운데 임의의 두 모듈이 서로 소임을 알 수 있다. 그 역, 곧 이 사실에서 \(\mM_1,\ldots,\mM_k\)가 서로 완전히 소인 계를 이룬다는 결론은 가환인 경우에만 일반적으로 성립한다. § 12의 반례 1을 보라.} \(k=2\)에서는 ‘서로 완전히 소’와 ‘서로 소’가 같은 뜻이므로, 주장은 \(k-1\)까지 증명되었다고 가정할 수 있다. 따라서 잉여류 \(\mx\bar{\mb}\)에는 유일한 표현 (18)이 성립하고, 동형에 따라 (16), 나아가 주장이 따른다.

\textbf{정리 II.} \emph{한 모듈 \(\mM\)을 서로 완전히 소인 계를 이루는 \(k\)개 모듈 \(\mM_1,\ldots,\mM_k\)의 최소공배수로 나타낼 수 있을 필요충분조건은 그 잉여군 \(\mG\)가 \(k\)개 부분군 \(\mU_1,\ldots,\mU_k\)로 가약인 것이다. 적절히 번호를 붙이면 \(\mU_i\)는 \(\mM_i\)의 잉여군과 동형이다.}

\pseventeenkoheading{6}{유한성정리}

이 절에서는 모든 잉여군을 \emph{유한개}의 기약성분들의 합으로 나타낼 수 있음을 보인다. 이를 위해 먼저 힐베르트의 모듈기저정리를 옮긴다.

이제부터는 § 2에서 세운 더 특별한 가정들을 쓴다. 곧 \(\xi_1,\ldots,\xi_n\)끼리의 곱셈은 가환적이고, 다항식의 곱셈법칙은
\begin{equation}
 \xi_i a=a_i\xi_i+b_i,\qquad
 (a_i\ne0)\quad(i=1,\ldots,n)
\tag{6}
\end{equation}
꼴이다.

모듈에 대해서만 증명하면 충분하다. 실제로 \(\mS\)가 임의의 다항식계이면, \(\mS\)의 유한개의 다항식 \(G_1,\ldots,G_k\)에서 \(A_1G_1+\cdots+A_kG_k\) 꼴로 합성할 수 있는 모든 다항식을 \(\mS\)에 덧붙여 여기서 유도된 모듈 \(\mM\)을 만든다. \(F_1,\ldots,F_\varrho\)이 \(\mM\)의 한 모듈기저이고
\[
 F_i=\sum_j A_{ij}G_j\qquad(i=1,\ldots,l)
\]
이면, 이 표현들에 나타나는 유한개의 다항식 \(G_j\)가 \(\mS\)의 한 기저를 이룬다.

모듈기저의 증명은 Gordan의 증명(Göttinger Nachrichten, 1899)을 따른다. 증명은 두 부분으로 갈린다. 첫 부분은 무한히 많은 거듭제곱곱으로 이루어진 한 계 \(\mP\)에 관한 것이며, \(\xi_1,\ldots,\xi_n\)의 가환성 때문에 본질적으로 그대로 유지된다. 각 거듭제곱곱은 계에 한 번만 넣고,
\[
 P=\xi_1^{a_1}\cdots\xi_n^{a_n}
\]
들을 차수가 커지는 순서로, 같은 차수 안에서는 임의의 고정된 순서로 배열하자. 이제 앞에 오는 어떤 \(P\)로도 나누어지지 않는 모든 거듭제곱곱 \(P=Q\)로 이루어진 \(\mP\)의 부분계 \(\mathfrak{Q}\)를 고찰한다. 어떤 \(Q\)도 뒤에 오는 \(Q\)로 나누어지지 않는다. 뒤의 것들은 차수가 더 높거나, 같은 차수이면 앞의 것과 다르기 때문이다. 반면 \(\mP\)의 각 \(P\)는 적어도 하나의 \(Q\)로 나누어진다. 그렇지 않다고 하고, 어떤 \(Q\)로도 나누어지지 않는 첫 \(P\)를 \(P_m\)이라 하자. \(P_m\) 자체는 \(Q\)일 수 없으므로 적어도 하나의 앞선 \(P\)로 나누어진다. 그런데 이 앞선 \(P\)는 가정에 따라 어떤 \(Q\)로 나누어지므로 \(P_m\)도 \(Q\)로 나누어져 모순이다.

이제 \emph{\(Q\)의 수가 유한함}을 주장한다. 첫 \(Q\)를
\[
 Q_0=\xi_1^{h_1}\cdots\xi_n^{h_n}
\]
이라 하자. 임의의 \(Q_\rho=\xi_1^{h_{\rho1}}\cdots\xi_n^{h_{\rho n}}\)는 \(Q_0\)로 나누어지지 않으므로, 각 \(Q\)마다 적어도 하나의 \(h_{\rho\lambda}<h_\lambda\)가 있다. \(\lambda\)는 매번 이 부등식이 처음 성립하는 첨자라고 하자. 이제 \(h_{\rho\lambda}\)가 같은 고정된 값 \(c_\lambda\)인 모든 \(Q\)를 한 류 \(K(c_\lambda)\)로 묶는다. \(c_\lambda<h_\lambda\)이므로 이 류들의 수는 유한하다. 각 류의 원소 수도 유한하다. 실제로 \(K(c_\lambda)\)의 한 \(Q\)를 \(\xi_\lambda^{c_\lambda}Q'\)로 쓰면 \(Q'\)들도 어느 하나가 다른 하나로 나누어지지 않는 거듭제곱곱들의 계를 이루고 변수는 \(n-1\)개뿐이다. \(n=2\)이면 \(Q'\)들은 변수 하나만 포함하므로 그 수가 \(=1\), 따라서 유한하다. 그러므로 변수 \(n-1\)개에 대해서도 그 수가 유한하다고 가정할 수 있고, 주장이 증명된다.

둘째로 이제 \(\mM\)이 다항식들의 한 모듈이고 그 항들을 사전식으로 배열했다고 하자. \(\mP\)는 각 최고항에 나타나는 거듭제곱곱들의 계이며, 각각을 한 번만 쓴다. 그러면 계 \(\mathfrak{Q}=(Q_1,\ldots,Q_\varrho)\)에 유한개의 다항식 \(F_1,\ldots,F_\varrho\)이 대응한다. 각 \(Q_i\)에 최고항이 \(b_iQ_i\)인 \(\mM\)의 어떤 다항식을 대응시켜 \(F_i=b_iQ_i+\)로 두며, 생략한 항들은 더 낮은 항들이다.

\(F=bP+\cdots\)가 \(\mM\)의 임의의 다항식이고 \(P=\xi_1^{\beta_1}\cdots\xi_n^{\beta_n}Q_i\)라 하자. 곱
\[
 \xi_1^{\beta_1}\cdots\xi_n^{\beta_n}F_i
 =\xi_1^{\beta_1}\cdots\xi_n^{\beta_n}(b_iQ_i+\cdots)
 =c\,\xi_1^{\beta_1}\cdots\xi_n^{\beta_n}Q_i+\cdots
 =cP+\cdots
\]
을 만들면 (6)에 따라 \(c\ne0\)이다. 또 곱셈은 지수에 관해서 더 낮은 항들을 제외하면 가환적이므로 이어지는 항들은 모두 사전식 순서에서 더 낮은 항들이다.\footnote{여기서 법칙 (6) 대신 조금 덜 제한적인
\[
 \xi_i a=a_i\xi_i+a_{i,i+1}\xi_{i+1}+\cdots+a_{in}\xi_n+b_i
 \qquad(a_i\ne0)
\]
를 써도 됨을 알 수 있다. 여기서는 \(\xi\)들에 어떤 고정된 번호를 붙인다.} 따라서 차
\[
 F-\frac bc\,\xi_1^{\beta_1}\cdots\xi_n^{\beta_n}F_i
\]
는 더 낮은 항들만 포함하고 다시 \(\mM\)에 속한다. 이 절차는 유한번 뒤에 끝나므로 \(F_1,\ldots,F_\varrho\)이 모듈기저임이 드러난다.

\textbf{정리 III.} \emph{무한히 많은 다항식으로 이루어진 모든 계 \(\mS\)는 유한 모듈기저 \(G_1,\ldots,G_k\)를 갖는다. 곧 \(\mS\)의 각 다항식 \(G\)는 \(G=A_1G_1+\cdots+A_kG_k\)로 나타낼 수 있다.}

여기서 다음이 곧바로 따른다.

\textbf{따름명제.} \emph{법 \(\mM\)의 무한히 많은 잉여류로 이루어진 모든 계 \(\mS\)는 기저 \(\mx_1,\ldots,\mx_k\)를 갖는다. 곧 \(\mS\)의 각 \(\mx\)는 \(\mx=A_1\mx_1+\cdots+A_k\mx_k\)로 나타낼 수 있다.}

실제로 각 잉여류에 임의의 대표원을 대응시키고, 이 대표원들과 \(\mM\)의 모든 다항식으로 이루어진 계를 고찰하면 된다.

이제 다음을 증명한다.

\textbf{정리 IV.} \emph{모든 잉여군은 유한개의 기약성분들의 합으로 나타낼 수 있다. 따라서 정리 II에 따라 모든 모듈은 서로 완전히 소인 계를 이루는 유한개의 기약모듈들의 최소공배수로 나타낼 수 있다.}

반대로 \(\mG\)가 무한히 많은 성분 \(\mU_1,\mU_2,\ldots\)의 합이라고 하자(§ 5를 보라). 그에 속하는 단위원들을 \(\ma_1,\ma_2,\ldots\)라 하면 정의상 모두 영과 다르다. 이 단위원들 전체의 계 \(\mS\)를 만들자. 정리 III의 따름명제에 따라 이 계에는 \(i_1<\cdots<i_\nu\)인 유한 기저 \(\ma_{i_1},\ldots,\ma_{i_\nu}\)가 있다. \(\mu>i_\nu\)라 하면
\[
 \ma_\mu=\mx_1\ma_{i_1}+\cdots+\mx_\nu\ma_{i_\nu}
\]
이다. 이는 단위원들 사이의 선형관계인데, 정의상 그러한 관계는 있을 수 없다. 따라서 한 잉여군의 각 가법적 분해에는 유한번 뒤에 기약성분이 하나 나타나며, 분해가 가환적이므로 이를 맨 앞으로 옮길 수 있다. 일반성을 잃지 않고 \(\mU_i\)들 자체를 기약이라고 할 수 있다. 위 논의에 따라 성분 \(\mU_i\)는 유한개만 나타날 수 있으므로 주장이 증명된다.

\pseventeenkoheading{7}{유일분해가 성립하는 한 경우}

\[
 \mG=\mU_1+\cdots+\mU_k=\mB_1+\cdots+\mB_l
\]
가 한 잉여군이고 그 성분 \(\mU_1,\ldots,\mU_k\), \(\mB_1,\ldots,\mB_l\)이 기약이라고 하자. 그러면
\begin{equation}
 \mx\me=\mx\ma_1+\cdots+\mx\ma_k
       =\mx\mb_1+\cdots+\mx\mb_l
\tag{22}
\end{equation}
이다. 여기서 \(\ma_\varkappa\)는 \(\mU_\varkappa\)의 단위원류에 동형으로 대응하는 \(\mG\)의 잉여류이고, \(\mb_\lambda\)도 \(\mB_\lambda\)에 대하여 같은 뜻이다. 이제 \(\mx\)를 \(\mx\mb_\lambda\)로 바꾸자. 여기서 \(\lambda\)는 \(1,\ldots,l\) 가운데 고정된 첨자이다. 그러면
\begin{equation}
 \mx\mb_\lambda
 =\mx\mb_\lambda\ma_1+\cdots+\mx\mb_\lambda\ma_k
\tag{23}
\end{equation}
이다. 군 \(\mB_\lambda\)의 이 표현은 유일하지만, 오른쪽의 개별 잉여류들이 \(\mB_\lambda\)에 속할 필요가 없으므로 \(\mB_\lambda\)의 가법적 분해는 아니다. \(\mb_\lambda\ne0\)이므로 모든 곱 \(\mb_\lambda\ma_\varkappa\)가 영일 수는 없다. 일반성을 잃지 않고
\[
 \mb_\lambda\ma_1\ne0,\ldots,\mb_\lambda\ma_\rho\ne0,\qquad
 \mb_\lambda\ma_{\rho+1}=0,\ldots,\mb_\lambda\ma_k=0
\]
이라 하자. 그러면
\[
 \mx\mb_\lambda
 =\mx\mb_\lambda\ma_1+\cdots+\mx\mb_\lambda\ma_\rho
\]
이다.

이제 \(1,\ldots,\rho\) 가운데 고정된 \(\varkappa\)에 대하여, 영과 다르고 \(\mx\mb_\lambda\ma_\varkappa\ne0\)도 만족하는 잉여류 \(\mx\mb_\lambda\)들의 전체를 고찰한다. \emph{두 계 \(\mx\mb_\lambda\)와 \(\mx\mb_\lambda\ma_\varkappa\)는 동형이다.} 우선 대응은 일대일이다. 가정에 따라 \(\mx\mb_\lambda\ma_\varkappa=0\)에서 \(\mx\mb_\lambda=0\)이 따르고, 역으로 뒤 등식에서 표현 (23)의 유일성에 따라 \(\mx\mb_\lambda\ma_\varkappa=0\)이 따르기 때문이다. 또 합에는 합이, 임의의 다항식을 곱한 것에는 대응하는 곱이 대응한다. 이제 \(\varkappa\)를 고정하고 \(\mb_\lambda\ma_\varkappa\ne0\)인 각 \(\lambda\)에 대하여 같은 논의를 할 수 있다. 각 \(\ma_\varkappa\)에는 그러한 \(\mb_\lambda\)가 적어도 하나 있어야 한다. 실제로 \(\ma_\varkappa=(\mb_1+\cdots+\mb_l)\ma_\varkappa\ne0\)이다.

\emph{먼저 각 \(\varkappa\)에 대하여 \(\mb_\lambda\ma_\varkappa\ne0\)인 \(\lambda\)가 하나뿐이고, 각 \(\lambda\)에 대하여 그러한 \(\varkappa\)도 하나뿐이라고 가정하자.}

이 대응에서 \(k=l\)이 따른다. 번호를 알맞게 붙여 \(\mb_\varkappa\ma_\varkappa\ne0\)이고 \(\lambda\ne\varkappa\)이면 \(\mb_\lambda\ma_\varkappa=0\)이 되게 할 수 있다. 이때 곱 \(\mb_\lambda\ma_\varkappa\)들이 대각도식을 이룬다고 한다. 이 경우 (23)에 따라 단위원 \(\mb_\varkappa=\mb_\varkappa\ma_\varkappa\)이고,
\begin{equation}
 \ma_\varkappa
 =\mb_1\ma_\varkappa+\cdots+\mb_l\ma_\varkappa
 =\mb_\varkappa\ma_\varkappa
 =\mb_\varkappa
\tag{24}
\end{equation}
이므로 일반적으로 \(\mU_\varkappa=\mB_\varkappa\)이다. \emph{따라서 분해는 유일하고}, 곱 \(\ma_\varkappa\mb_\lambda\)들도 대각도식을 이룬다.

\emph{특별히 잉여류 \(\ma_\varkappa,\mb_\lambda\)의 곱에 교환법칙이 성립하면 위 가정은 언제나 충족된다.} 실제로 이때 (23)의 각 성분 \(\mx\mb_\lambda\ma_\varkappa=\mx\ma_\varkappa\mb_\lambda\)는 \(\mB_\lambda\)에 속하므로 (23)은 군 \(\mB_\lambda\)의 가법적 분해이다. \(\mB_\lambda\)가 기약이므로 하나를 제외한 모든 \(\mb_\lambda\ma_\varkappa\)가 영이다. 또 (24)에 따라 고정된 \(\varkappa\)에 대하여도 영과 다른 \(\mb_\lambda\ma_\varkappa\)는 하나뿐이다. 그렇지 않으면 \(\ma_\varkappa\)와 \(\mb_\lambda\)의 교환성 때문에 (24)가 기약이어야 할 \(\mU_\varkappa\)의 가법적 분해를 줄 것이기 때문이다.\footnote{앞서 든 Schmeidler의 논문에서 다룬 비가환 양쪽 모듈의 경우에도 같은 사실이 성립한다. 그 경우에는 한 부분군의 각 잉여류와 다른 부분군의 각 잉여류의 곱이 사라진다. 실제로 \(\mb_\lambda\ma_\varkappa=\mb_\lambda\ma_\varkappa\mb_1+\cdots+\mb_\lambda\ma_\varkappa\mb_l\)이고, \(\lambda\ne\mu\)이면 \(\mb_\lambda(\ma_\varkappa\mb_\mu)=0\)이므로 \(\mb_\lambda\ma_\varkappa=\mb_\lambda\ma_\varkappa\mb_\lambda\), 따라서 이 곱은 \(\mB_\lambda\)에 속한다. \(\mB_\lambda\)의 기약성에서 \(\mb_\lambda\ma_\varkappa\ne0\)인 \(\varkappa\)가 하나뿐, 따라서 정확히 하나임이 따른다. 같은 방식으로 각 \(\ma_\varkappa\)에 대하여 \(\ma_\varkappa\mb_\lambda\ne0\)인 \(\mb_\lambda\)가 정확히 하나 있음을 보이고, 여기서 \(k=l\)이 따른다. 따라서 \(\mb_\lambda\ma_\varkappa\)들은 대각도식을 이루며, 이로써 그곳에서 증명한 유일성이 나온다.}

모든 \(\mb_\lambda\)가 아니라 그 일부 \((\lambda=1,\ldots,\sigma)\)에 대해서만 \(\mb_\lambda\ma_\lambda\ne0\), \(\lambda\ne\varkappa\), \(\varkappa=1,\ldots,k\)이면 \(\mb_\lambda\ma_\varkappa=0\)이고, 마찬가지로 각 \(\mu=1,\ldots,l\)에 대하여 \(\varkappa=1,\ldots,\sigma\)이면 \(\mb_\mu\ma_\varkappa=0\), \(\mb_\varkappa\ma_\varkappa\ne0\)이라는 조건들이 성립하는 경우에도 같은 결론이 적용된다. 이때 \(\mU_\varkappa=\mB_\varkappa\) (\(\varkappa=1,\ldots,\sigma\))를 얻는다. 또
\[
 \mU_{\sigma+1}+\cdots+\mU_k=\mU,\qquad
 \mB_{\sigma+1}+\cdots+\mB_l=\mB
\]
로 두면, \(\mU,\mB\)의 단위원 \(\ma,\mb\)에 대해서도 \(\mb\ma_\varkappa=0\), \(\mb_\varkappa\ma=0\) (\(\varkappa=1,\ldots,\sigma\)), \(\mb\ma\ne0\)이고, 여기서 \(\mU=\mB\)가 따른다.

\pseventeenkoheading{8}{완전가약군의 소군들로의 가법적 분해. 동형정리}

이제 유일성은 성립하지 않지만 서로 다른 분해들의 동형이 성립하는 다른 경우를 고찰한다. 이는 Loewy가 다룬 경우의 일반화이다.

\textbf{정의.} \emph{자기 자신과 모든 다항식을 포함하는 단위원모듈 외에는 인자를 갖지 않는 모듈을 ‘소모듈’이라고 한다. 한 소모듈의 잉여군을 ‘소군’이라고 한다.}

모든 소군은 당연히 기약이다. 무한 소군도 존재한다(§ 11의 정의와 § 12의 예 3를 보라).

주어진 잉여군 \(\mG\)를 두 가지 방식으로 각각 유한개의 소군 \(\mU_1,\ldots,\mU_k\), 또는 \(\mB_1,\ldots,\mB_l\)의 합으로 나타낼 수 있다고 하자. \emph{소군들의 합으로 나타내는 그러한 표현을 적어도 하나 갖는 잉여군과 그 대응 모듈 \(\mM\)을 Loewy의 용어에 따라 완전가약이라고 한다.}

이제 곱 \(\ma_\varkappa\mb_\lambda\) (\(\varkappa=1,\ldots,k\), \(\lambda=1,\ldots,l\))를 고찰한다. \(\ma_\varkappa\mb_\lambda=0\)이거나, 아니면 \(\mx\)가 \(\mG\)의 모든 잉여류를 돌 때 \(\mx\ma_\varkappa\mb_\lambda\)가 군 \(\mB_\lambda\) 전체를 돈다는 것을 보이겠다. 실제로 \(\mB_\lambda\)에 속하는 모듈
\[
 (\mM,\mb_1+\cdots+\mb_{\lambda-1}
       +\mb_{\lambda+1}+\cdots+\mb_l)=\mM_{\mb_\lambda},
\]
곧 \(\mb_1+\cdots+\mb_{\lambda-1}+\mb_{\lambda+1}+\cdots+\mb_l\)이라는 잉여류의 모든 다항식을 \(\mM\)에 덧붙여 얻는 모듈에는 인자
\[
 (\mM,\mb_1+\cdots+\mb_{\lambda-1}
       +\mb_{\lambda+1}+\cdots+\mb_l,
       \ma_\varkappa\mb_\lambda)
\]
가 있다. 그런데 가정상 \(\mM_{\mb_\lambda}\)는 소모듈이므로 이 인자는 \(\mM_{\mb_\lambda}\) 자신이거나 단위원모듈이다. 첫 경우에는 \(\ma_\varkappa\mb_\lambda\)가 \(\mM_{\mb_\lambda}\)에 속하므로 관계
\[
 \ma_\varkappa\mb_\lambda
 =\mx(\mb_1+\cdots+\mb_{\lambda-1}
       +\mb_{\lambda+1}+\cdots+\mb_l)
\]
가 있고, 여기서 \(\ma_\varkappa\mb_\lambda=0\)이 따른다. 다른 경우에는 두 잉여류 \(\mx_0,\mx_1\)가 있어서
\[
 \mx_0\ma_\varkappa\mb_\lambda
 +\mx_1(\mb_1+\cdots+\mb_{\lambda-1}
        +\mb_{\lambda+1}+\cdots+\mb_l)
 =\me
 =\mb_1+\cdots+\mb_{\lambda-1}
  +\mb_{\lambda+1}+\cdots+\mb_l+\mb_\lambda
\]
이다. 유일성 때문에 \(\mx_0\ma_\varkappa\mb_\lambda=\mb_\lambda\)이다. 따라서 \(F\), 또는 \(\mx\)가 각각 모든 다항식, 또는 \(\mG\)의 모든 잉여류를 돌면 \(F\mx_0\ma_\varkappa\mb_\lambda\), 따라서 \(\mx\ma_\varkappa\mb_\lambda\)가 군 \(\mB_\lambda\) 전체를 돈다. 이로써 주장을 증명하였다.

\(\ma_\varkappa\mb_\lambda\ne0\)이면 \(\mU_\varkappa\)는 \(\mB_\lambda\)와도 동형이다. 실제로 \(\mx\ma_\varkappa\ne0\)이고 \(\mx\ma_\varkappa\mb_\lambda\ne0\)인 모든 류의 집합은 § 7에 따라 \(\mx\ma_\varkappa\mb_\lambda\) 꼴 류들의 계와 동형이고, 뒤의 류들은 부분군 \(\mB_\lambda\)를 모두 소진한다. 따라서 앞서 말한 류 \(\mx\ma_\varkappa\)들이 군 \(\mU_\varkappa\)를 모두 소진함만 보이면 된다. 이를 위해 \(\mx\ma_\varkappa\mb_\lambda=0\)인 모든 잉여류 \(\mx\ma_\varkappa\)를 고찰하자. 이 계의 두 류의 합과 한 류에 임의의 다항식을 곱한 것은 다시 이 계에 속한다. 따라서 이 모든 잉여류를 모듈 \(\mM_{\ma_\varkappa}\)에 덧붙이면 다시 \(\mM_{\ma_\varkappa}\)의 인자인 모듈이 생기며, 이는 \(\mM_{\ma_\varkappa}\) 자신이거나 단위원모듈이다. 뒤의 경우라면 \(\ma_\varkappa\) 자체도 이 모듈에 속하여 가정과 달리 \(\ma_\varkappa\mb_\lambda=0\)일 것이다. 그러므로 \(\mx\ma_\varkappa\ne0\)이면서 \(\mx\ma_\varkappa\mb_\lambda=0\)인 류는 없고, \(\mU_\varkappa\)와 \(\mB_\lambda\)의 동형이 증명된다.

어떤 고정된 \(\varkappa\)에 대하여 \(\ma_\varkappa\mb_{\lambda_1}\ne0,\ldots,\ma_\varkappa\mb_{\lambda_i}\ne0\) (\(i>1\))이라고 하자. 그러면 \(\mU_\varkappa\)는 군 \(\mB_{\lambda_1},\ldots,\mB_{\lambda_i}\)와 동형이므로, 이 군들은 서로 동형이다.

이때 \(\mU\)들 가운데도 적어도 두 개가 동형임을 보이겠다. 곱 \(\mb_\lambda\ma_\varkappa\)에서 각 \(\mb_\lambda\)에 \(\mb_\lambda\ma_\varkappa\ne0\)인 \(\ma_\varkappa\)가 하나씩만 대응한다면 \(\mb_\lambda=\mb_\lambda\ma_\varkappa\)일 것이다. 그러면 \(\mU_\varkappa\)는 소군이어서 진부분군을 포함할 수 없으므로 \(\mB_\lambda=\mU_\varkappa\)이다. 따라서 \(k=l\)이고, 알맞게 번호를 붙이면 \(\ma_\varkappa=\mb_\varkappa\)이다. 그러면 도식 \(\ma_\varkappa\mb_\lambda\)도 대각도식이어야 하지만, 지금은 그렇지 않다. 이로써 다음 정리를 증명하였다.

\textbf{정리 V.} \emph{완전가약군 \(\mG\)를 두 가지 방식으로 소군들의 합
\[
 \mG=\mU_1+\cdots+\mU_k=\mB_1+\cdots+\mB_l
\]
로 나타냈다고 하자. 그러면 각 군 \(\mU\)는 적어도 하나의 군 \(\mB\)와 동형이고, 그 역도 성립한다. 각 \(\mU\)가 정확히 하나의 \(\mB\)와 동형이면 두 분해는 동일하다. 그렇지 않으면 \(\mU\)들 가운데 적어도 두 개가 서로 동형이고, \(\mB\)들 가운데도 적어도 두 개가 서로 동형이다.}

\pseventeenkoheading{9}{‘동형’과 ‘같은 종류’라는 개념의 관계}

이 절에서는 한 변수 미분식에 대하여 알려진 ‘같은 종류’라는 개념을 뜻에 맞게 일반화하면 대응 잉여군들의 동형에 이르게 됨을 보인다.

\textbf{정의}\footnote{한 변수 미분식에 대해서는 앞서 든 Blumberg의 논문, 25쪽의 형식적 정의를 보라.}. \emph{두 모듈 \(\mM,\mN\)이 같은 종류라는 것은, \(\mM\)과 서로 소인 다항식 \(P\), \(\mN\)과 서로 소인 다항식 \(Q\)가 있어서 모든 \(M\equiv0(\mM)\), \(N\equiv0(\mN)\)에 대하여}
\begin{align}
 MQ&\equiv0(\mN),\tag{25}\\
 NP&\equiv0(\mM)\tag{26}
\end{align}
\emph{이 성립한다는 뜻이다. 여기서 서로 소라는 것은 두 다항식 \(X,Y\)가 있어서}
\begin{align}
 YQ&\equiv1(\mN),\tag{27}\\
 XP&\equiv1(\mM)\tag{28}
\end{align}
\emph{이 성립한다는 뜻이다.}

이제 다음 정리가 성립한다.

\textbf{정리 VI.} \emph{두 모듈이 같은 종류일 필요충분조건은 그 잉여군들이 동형인 것이다.}

\textbf{1.} \(\mA\)가 \(\mM\)의 잉여군, \(\mB\)가 \(\mN\)의 잉여군이고, \(\mA,\mB\)가 동형이라고 하자. \(\ma,\mb\)를 각각 \(\mA,\mB\)의 단위원류라 하자. 또 \(Q\mb\)는 \(\mA\)의 단위원 \(\ma\)에 대응하는 \(\mB\)의 류이고, \(P\ma\)는 \(\mB\)의 단위원 \(\mb\)에 대응하는 \(\mA\)의 류라 하자. 이 대응을
\begin{align}
 \ma&\sim Q\mb,\tag{29}\\
 P\ma&\sim\mb\tag{30}
\end{align}
로 쓴다. 여기서
\[
 0=M\ma\sim MQ\mb,\qquad\text{곧}\qquad MQ\equiv0(\mN)
\]
가 따르므로 (25)가 증명된다. 마찬가지로
\[
 NP\ma\sim N\mb=0,\qquad\text{곧}\qquad NP\equiv0(\mM)
\]
이어서 관계 (26)이 성립한다. 또 (29)에 따라 \(P\ma\sim PQ\mb\)이고, (30)에서는 \(P\ma\sim\mb\)라고 가정했으므로 대응의 유일성에 따라 \(PQ\mb=\mb\), 곧 \(PQ\equiv1(\mN)\)이다. 마찬가지로 \(QP\equiv1(\mM)\)이다. 따라서 \(X=Q,Y=P\)인 (27), (28)도 증명되었다.

\textbf{2.} \(\mM,\mN\)이 같은 종류여서 관계 (25)--(28)이 성립한다고 하자. \(\mA\)의 임의의 류 \(F\ma\)에 \(\mB\)의 류 \(FQ\mb\)를 대응시킨다. 이때 \(\mA\)의 각 류에는 \(\mB\)의 류 하나만 대응한다. 실제로 \(F\ma=0\), 곧 \(F=M\)이면 (25)에 따라 \(FQ\mb=0\)이다. 또 이 대응은 군 \(\mB\) 전체를 소진한다. 특별한 류 \(Y\ma\)에는 (27)에 따라 \(\mb\)와 같은 류 \(YQ\mb\)가 대응하기 때문이다. 따라서 \(\mA\)에서 \(\mB\) 위로 가는 \emph{일의적} 사상 \(\Phi\)를 얻었지만, 이것이 일대일임은 아직 보이지 않았다. 똑같이 (26), (28)에서 \(\mB\)에서 \(\mA\) 위로 가는 일의적 사상 \(\Psi\)를 얻지만, 역시 일대일임은 아직 보이지 않았다.

두 사상 가운데 하나가 일대일이면 합과 곱에 관한 요구가 충족되므로 동형이 증명된다. 그러나 \(\Phi,\Psi\)가 모두 일대일이 아니라고 하자. 그러면 이를테면 \(\mA\)에는 \(\mB\)에서 영류 \(GQ\mb=0\)에 대응하는 류 \(G\ma\)가 있다. 나머지 잉여류 \(F\ma\)에 대응하는 류 \(FQ\mb\)들만으로 이미 군 \(\mB\)를 소진하므로, 사상 \(\Psi\)를 적용하면 류 \(FQP\ma\)들로 군 \(\mA\) 전체를 다시 얻는다. 또 \(\mB\)에는 \(\Psi\)에 따라 \(\mA\)의 영류에 대응하는 영과 다른 류들이 있다. \(G_1\ma\ne0\)이지만 \(G_1QP\ma=0\)인 모든 류를 \(G_1\ma\)라고 나타내자. 다른 모든 다항식 \(F\)는 \(FQP\ma\)가 군 \(\mA\)를 소진한다는 성질을 갖는다. \(FQP\ma=G_1\ma\)가 되는 다항식 \(F\)들을 \(G_2\)라고 하자. \(G_1\)이 존재하면 그러한 다항식은 언제나 존재한다. 각 \(G_2\)에 대하여 \(G_2QP\ma\ne0\)이지만 \(G_2(QP)^2\ma=0\)이다. 마찬가지로 각 양의 정수 \(\nu\)에 대하여 \(G_\nu(QP)^{\nu-1}\ma\ne0\), \(G_\nu(QP)^\nu\ma=0\)인 류 \(G_\nu\ma\)가 있다.

모든 다항식 \(G_1,G_2,\ldots\)로 이루어진 계 \(\mS\)를 고찰하자. 이는 유한 모듈기저 \(G_{\lambda_1},\ldots,G_{\lambda_\varrho}\)를 갖는다. \(\lambda\)를 첨자 \(\lambda_1,\ldots,\lambda_\varrho\) 가운데 가장 큰 것으로 두고 \(\nu>\lambda\)를 택하면
\[
 G_\nu=A_1G_{\lambda_1}+\cdots+A_\varrho G_{\lambda_\varrho}.
\]
여기에 \((QP)^\lambda\)를 곱하면
\[
 G_\nu(QP)^\lambda\equiv0(\mM)
\]
을 얻어 가정과 모순이다. 따라서 동형이 성립하고, 1에 따라 \(X=Q,Y=P\)인 더 강한 관계 (27), (28)도 성립한다.

정리 VI를 이용하면 정리 V를 다음과 같이도 쓸 수 있다.

\textbf{정리 VII}\footnote{Loewy, 100--101쪽을 보라.}. \emph{완전가약모듈 \(\mM\)을 두 가지 방식으로, 각각 서로 완전히 소인 계를 이루는 소모듈 \(\mP_1,\ldots,\mP_k\), 또는 \(\mathfrak{Q}_1,\ldots,\mathfrak{Q}_l\)의 최소공배수로 나타냈다고 하자. 그러면 각 \(\mP\)는 적어도 하나의 \(\mathfrak{Q}\)와 같은 종류이고, 그 역도 성립한다. 각 \(\mP\)가 정확히 하나의 \(\mathfrak{Q}\)와 같은 종류이면 두 분해는 동일하다. 그렇지 않으면 \(\mP\)들 가운데 적어도 두 개가 서로 같은 종류이고, \(\mathfrak{Q}\)들 가운데도 적어도 두 개가 서로 같은 종류이다.}

\pseventeenkoheading{10}{한 군에 무한히 많은 분해가 존재함}

완전가약군의 합표현이 유일하지 않으면, 정리 V에 따라 각 표현에는 서로 동형인 부분군이 적어도 두 개씩 있다. 그러한 동형 부분군들의 한 계에서 이루어진 부분군
\[
 \mathfrak H=\mathfrak A_1+\cdots+\mathfrak A_\alpha
\]
를 고찰하자. 여기서 \(\mathfrak A_1,\mathfrak A_2,\ldots,\mathfrak A_\alpha\)는 모두 서로 동형이다. 유리성영역 \(\mathsf{P}\)가 ‘상수’, 곧 \(\xi_i a=a\xi_i\)인 원소 \(a\)를 무한히 많이 포함하기만 하면 \(\mathfrak H\)가 그러한 서로 다른 분해를 무한히 많이 가짐을 주장한다. 이 정리는 군 \(\mathfrak A_i\)들에 더 아무것도, 심지어 기약성조차 가정하지 않아도 성립한다. 따라서 다음이 성립한다.

\textbf{정리 VIII.} \emph{\(\mathfrak G\)를 유한개의 서로 동형인 부분군들의 합으로 나타낼 수 있고, 유리성영역 \(\mathsf{P}\)가 상수를 무한히 많이 포함하면, 이 표현은 서로 다른 방식으로 무한히 많이 가능하다. 여기서 상수란 \(\xi_i a=a\xi_i\)인 영역의 원소 \(a\)를 뜻한다.}

\textbf{증명.} 다음과 같이 두자.
\[
 \mathfrak G=\mathfrak A_1+\cdots+\mathfrak A_\alpha,
\]
그리고 \(i,\varkappa=1,\ldots,\alpha\), \(\alpha\geq2\)에 대하여 \(\mathfrak A_i\)가 \(\mathfrak A_\varkappa\)와 동형이라고 하자. 하나의 특정한 동형대응을 주기 위해 먼저 한 부분군, 이를테면 \(\mathfrak A_1\)을 특별히 택한다. \(\mathfrak A_1\)과 \(\mathfrak A_i\) 사이의 동형을
\begin{equation}
 a_1\sim t_{1\varkappa}a_\varkappa;\qquad
 a_i\sim t_{i1}a_1;\qquad
 t_{11}a_1=a_1
 \quad(i,\varkappa=1,\ldots,\alpha)
\tag{31}
\end{equation}
로 고정하자. 군 \(\mathfrak A_1\) 안에서는 각 잉여류가 자기 자신에 대응한다. \(M\mathfrak{a}_1=0\), \(M\mathfrak{a}_i=0\)에서 \(M\mathfrak{t}_{1\varkappa}\mathfrak{a}_\varkappa=0\), \(M\mathfrak{t}_{i1}\mathfrak{a}_1=0\)이 따른다. 따라서 류 \(t_{1\varkappa}a_\varkappa,t_{i1}a_1\)은 단위원과 마찬가지로 다항식뿐 아니라 잉여류와도 곱할 수 있다. 동형의 정의에 따라 서로 대응하고 이러한 잉여류 곱셈을 허용하는 두 류 \(\mathfrak y,\bar{\mathfrak y}\)에 대해서는 \(\mathfrak x\mathfrak y\)와 \(\mathfrak x\bar{\mathfrak y}\)도 서로 대응한다. 그러므로 관계 \(\mathfrak{a}_r\mathfrak{a}_s=0\) (\(r\ne s\), \(r,s=1,\ldots,\alpha\))에서 \(s=1\)로 두면
\begin{equation}
\left\{
\begin{aligned}
 a_r t_{1\varkappa}a_\varkappa&=0
 &&(\varkappa=1,\ldots,\alpha,\ r\ne1)\\
 \text{그리고 \(s>1\)이면}\qquad
 \mathfrak{a}_r \mathfrak{t}_{s1}\mathfrak{a}_1&=0 &&(r\ne s),\\
 \text{따라서}\qquad
 \mathfrak{t}_{1\varkappa}\mathfrak{a}_\varkappa&=\mathfrak{a}_1\mathfrak{t}_{1\varkappa}\mathfrak{a}_\varkappa;
 & \mathfrak{t}_{i1}\mathfrak{a}_1&=\mathfrak{a}_i\mathfrak{t}_{i1}\mathfrak{a}_1
\end{aligned}\right.
\tag{32}
\end{equation}
을 얻는다.

또 (31)과 대응의 유일성에서 (27), (28)의 일반화인
\begin{equation}
 \mathfrak{t}_{1i}\mathfrak{t}_{i1}\mathfrak{a}_1=\mathfrak{a}_1
 \mathfrak{t}_{\varkappa1}\mathfrak{t}_{1\varkappa}\mathfrak{a}_\varkappa=\mathfrak{a}_\varkappa
\tag{33}
\end{equation}
가 따른다. 이제 \(\mathfrak A_i,\mathfrak A_\varkappa\) 사이의 동형은 (31)의 합성으로 생기는 관계로 고정하자. 개별 군은 자기동형을 가질 수 있으므로 이렇게 고정할 필요가 있다. (31)--(33)에서
\begin{equation}
 a_i\sim t_{i1}t_{1\varkappa}a_\varkappa
      =t_{i\varkappa}a_\varkappa;\qquad
 \mathfrak{t}_{ii}\mathfrak{a}_i=\mathfrak{a}_i
 \mathfrak{a}_r\mathfrak{t}_{s\varkappa}\mathfrak{a}_\varkappa=0\quad (r\ne s)
\tag{34}
\end{equation}
를 얻는다. 여기서 \(\mathfrak{t}_{i1}\mathfrak{t}_{1\varkappa}=\mathfrak{t}_{i\varkappa}\)로 두었다. (33), (34)에서 더 나아가
\[
 t_{i\varkappa}t_{\varkappa l}a_l
 =t_{i1}t_{1\varkappa}t_{\varkappa1}t_{1l}a_l
 =\mathfrak{t}_{i1}\mathfrak{t}_{1l}\mathfrak{a}_l=\mathfrak{t}_{il}\mathfrak{a}_l
\]
이다. 이로써 \(\mathfrak A_1\)을 특별히 택한 것은 다시 사라진다. 정의에 \(\mathfrak A_1\) 대신 어느 군 \(\mathfrak A_l\)를 써도 \(\mathfrak A_i,\mathfrak A_\varkappa\) 사이의 동형관계는 같다. 따라서 \(\mathfrak{t}_{i\varkappa}\)들은 한데 다음 관계들을 만족하는 류들이다.
\begin{equation}
 M\mathfrak{t}_{i\varkappa}\mathfrak{a}_\varkappa=0
 \quad \mathfrak{a}_r\mathfrak{t}_{i\varkappa}\mathfrak{a}_\varkappa=0\ (r\ne i)
 \quad \mathfrak{t}_{i\varkappa}\mathfrak{t}_{\varkappa l}\mathfrak{a}_l=\mathfrak{t}_{il}\mathfrak{a}_l
 t_{\varkappa i}a_i
 =t_{\varkappa i}t_{i\varkappa}a_\varkappa
 =a_\varkappa
\tag{35}
\end{equation}
이다. 이 관계들은 \(\mathfrak A_i,\mathfrak A_\varkappa\)에 속하는 모듈 \(\mathfrak M_i,\mathfrak M_\varkappa\)가 같은 종류라는 것을 정확히 말한다(§ 9의 정의를 보라). 그곳의 \(P,Q\) 대신 여기서는 각각 류 \(t_{i\varkappa},t_{\varkappa i}\)가 오며, (35)의 첫 두 관계를 합친 것은 각각 공식 (25), (26)에 대응한다. 역으로 이 관계들에서는 정리 VI에 따라 동형이 따른다.\footnote{여기서는 유한성정리를 쓰지 않아도 된다. 공식들이 이미 사상 \(\varphi_{i\varkappa}\)와 그 역 \(\varphi_{\varkappa i}\)을 각각 나타내기 때문이다. 실제로 \(\mathfrak{x}\mathfrak{a}_i=\mathfrak{x}\mathfrak{a}_i \mathfrak{t}_{i\varkappa}\mathfrak{a}_\varkappa \mathfrak{t}_{\varkappa i}\mathfrak{a}_i\)이므로 \(\mathfrak{x}\mathfrak{a}_i\ne0\)에서 \(\mathfrak{x}\mathfrak{a}_i \mathfrak{t}_{i\varkappa}\mathfrak{a}_\varkappa\ne0\)도 따른다.}
이제 류 \(\mathfrak{t}_{i\varkappa}\mathfrak{a}_\varkappa\)들에서 다시 군 \(\mathfrak B_\rho\)의 단위원이 되는 류 \(\mathfrak{b}_\rho\)들을 합성할 수 있다. 그러면
\[
 \mathfrak G=\mathfrak B_1+\cdots+\mathfrak B_\alpha
\]
라는 표현이 따라온다. 실제로 \(\lambda_{i\varkappa}\)들을 행렬식 \(|\lambda_{i\varkappa}|\ne0\)인 상수들이라고 하고, \(\lambda^{i\varkappa}\)를 \(|\lambda_{i\varkappa}|\)로 나눈 대응 여인자라고 하자. 그러면
\begin{equation}
\left\{
\begin{aligned}
 \sum_{r=1}^{\alpha}\lambda_{ri}\lambda^{r\varkappa}
 &=\delta_{i\varkappa}
 =\begin{cases}0,&i\ne\varkappa,\\1,&i=\varkappa,\end{cases}\\
 \sum_{r=1}^{\alpha}\lambda_{ir}\lambda^{\varkappa r}
 &=\delta_{i\varkappa}
\end{aligned}\right.
\tag{36}
\end{equation}
이다. 이제
\begin{equation}
 b_\rho=\sum_{i,\varkappa}
 \lambda_{\rho i}\lambda^{\rho\varkappa}
 t_{i\varkappa}a_\varkappa
 \qquad(\rho=1,\ldots,\alpha)
\tag{37}
\end{equation}
로 정의한다. (35)--(37)에서 \(M\mathfrak{b}_\rho=0\)이고
\begin{equation}
 \sum_{\rho=1}^{\alpha}\mathfrak{b}_\rho
 =\sum_{i,\varkappa,\rho}
  \lambda_{\rho i}\lambda^{\rho\varkappa}
  t_{i\varkappa}a_\varkappa
 =\sum_{i,\varkappa}\delta_{i\varkappa}
  t_{i\varkappa}a_\varkappa
 =\sum_i \mathfrak{t}_{ii}\mathfrak{a}_i=\sum_i \mathfrak{a}_i=\mathfrak{e}
\tag{38}
\end{equation}
이다. 더 나아가
\[
\begin{aligned}
 \mathfrak{b}_\rho \mathfrak{b}_\sigma
 &=\sum_{\substack{i,\varkappa\\ \mu,\nu}}
  \lambda_{\rho i}\lambda^{\rho\varkappa}
  \lambda_{\sigma\mu}\lambda^{\sigma\nu}
  \mathfrak{t}_{i\varkappa}\mathfrak{a}_\varkappa \mathfrak{t}_{\mu\nu}\mathfrak{a}_\nu \\
 &=\sum_{i,\mu,\nu}
  \lambda_{\rho i}\lambda^{\rho\mu}
  \lambda_{\sigma\mu}\lambda^{\sigma\nu}
  t_{i\nu}a_\nu
 =\delta_{\rho\sigma}\sum_{i,\nu}
  \lambda_{\rho i}\lambda^{\sigma\nu}t_{i\nu}a_\nu\\
 &=\begin{cases}0&(\rho\ne\sigma),\\ b_\rho&(\rho=\sigma)\end{cases}
\end{aligned}
\]
이다 (\(\rho,\sigma=1,\ldots,\alpha\)). 따라서 표현 \(\mathfrak r e=\mathfrak r b_1+\cdots+\mathfrak r b_\alpha\)는 군 \(\mathfrak G\)의 한 가법적 분해이다. \(\lambda_{i\varkappa}\)가 대각행렬만을 이루지 않으면 군 \(\mathfrak B_\rho\)들은 군 \(\mathfrak A_\sigma\)들과 다르다. 실제로 \(a_\sigma b_\rho\)를 만들면 (35)에 따라
\[
 a_\sigma b_\rho
 =\sum_{i,\varkappa}\lambda_{\rho i}\lambda^{\rho\varkappa}
  a_\sigma t_{i\varkappa}a_\varkappa
 =\lambda_{\rho\sigma}\sum_\varkappa
  \lambda^{\rho\varkappa}\mathfrak{t}_{\sigma\varkappa}\mathfrak{a}_\varkappa
 \ne0
 \quad\text{\(\lambda_{\rho\sigma}\ne0\)이면}
\]
이다. 마지막 합이 영이라면 각 합항이 따로 사라져야 하지만, 그렇지 않기 때문이다. 따라서 고정된 \(\sigma\)에 대하여 \(\lambda_{\rho\sigma}\ne0\)인 \(\rho\)가 적어도 두 개이면 \(\mathfrak A_\sigma\)는 어떤 \(\mathfrak B_\rho\)와도 동일할 수 없다.

그런데 군 \(\mathfrak B_\rho\)들은 서로 동형이고 군 \(\mathfrak A_\sigma\)들과도 동형이다. 이를 위해 먼저 \(\mathfrak A_\sigma,\mathfrak B_\rho\)의 동형을 매개하는 류 \(\mathfrak{r}_{\rho\sigma}\mathfrak{a}_\sigma\), \(\mathfrak{r}^{\sigma\rho}\mathfrak{b}_\rho\)의 존재를 보이자 (\(\rho,\sigma=1,\ldots,\alpha\)). 곧 류
\[
 r_{\rho\sigma}a_\sigma
 =\sum_i\lambda_{\rho i}t_{i\sigma}a_\sigma,\qquad
 r^{\sigma\rho}
 =\sum_\varkappa\lambda^{\rho\varkappa}
  t_{\sigma\varkappa}a_\varkappa
 =\mathfrak{r}^{\sigma\rho}\mathfrak{b}_\rho\text{\footnotemark[21]}
\]
\footnotetext[21]{(37)에 따라 \(\mathfrak{r}^{\sigma\rho}\mathfrak{b}_\rho=\sum_{\varkappa,i,\mu}\lambda^{\rho\varkappa}\mathfrak{t}_{\sigma\varkappa}\mathfrak{a}_\varkappa\lambda_{\rho i}\lambda^{\rho\mu}\mathfrak{t}_{i\mu}\mathfrak{a}_\mu=\sum_{\varkappa,\mu}\lambda^{\rho\varkappa}\lambda_{\rho\varkappa}\lambda^{\rho\mu}\mathfrak{t}_{\sigma\mu}\mathfrak{a}_\mu=\sum_\mu\lambda^{\rho\mu}\mathfrak{t}_{\sigma\mu}\mathfrak{a}_\mu=\mathfrak{r}^{\sigma\rho}\)이다.}
가 모듈
\[
 (\mathfrak M,b_1+\cdots+b_{\rho-1}
              +b_{\rho+1}+\cdots+b_\alpha)
\]
과
\[
 (\mathfrak M,a_1+\cdots+a_{\sigma-1}
              +a_{\sigma+1}+\cdots+a_\alpha)
\]
사이의 ‘같은 종류’ 관계
\begin{equation}
\left\{
\begin{aligned}
 b_\varkappa r_{\rho\sigma}a_\sigma&=0
 &&(\varkappa\ne\rho);&
 \mathfrak{r}^{\sigma\rho}\mathfrak{r}_{\rho\sigma}\mathfrak{a}_\sigma
 &=a_\sigma\text{\footnotemark[22]},\\
 a_\varkappa r^{\sigma\rho}b_\rho&=0
 &&(\varkappa\ne\sigma);&
 \mathfrak{r}_{\rho\sigma}\mathfrak{r}^{\sigma\rho}\mathfrak{b}_\rho
 &=b_\rho
\end{aligned}\right.
\tag{39}
\end{equation}
\footnotetext[22]{역관계는 여기서처럼 첨자들을 서로 바꾸어 얻어지지 않으므로, 공식의 수가 (35)에 비해 두 배이다.}
\setcounter{footnote}{22}
를 만족함을 보인다. 그러면 (35)의 경우와 마찬가지로 대응
\begin{equation}
 \mathfrak{b}_\rho\sim \mathfrak{r}_{\rho\sigma}\mathfrak{a}_\sigma;\qquad
 \mathfrak{a}_\sigma\sim \mathfrak{r}^{\sigma\rho}\mathfrak{b}_\rho
\tag{40}
\end{equation}
에 따라 동형이 따른다.

실제로
\[
\begin{aligned}
 \mathfrak{b}_\varkappa \mathfrak{r}_{\rho\sigma}\mathfrak{a}_\sigma
 &=\sum_{\mu,\nu,i}
  \lambda_{\varkappa\mu}\lambda^{\varkappa\nu}
  t_{\mu\nu}a_\nu\cdot\lambda_{\rho i}t_{i\sigma}a_\sigma\\
 &=\sum_{\mu,i}
  \lambda_{\varkappa\mu}\lambda^{\varkappa i}
  \lambda_{\rho i}t_{\mu\sigma}a_\sigma
 =\delta_{\varkappa\rho}\sum_\mu
  \lambda_{\varkappa\mu}t_{\mu\sigma}a_\sigma\\
 &=\delta_{\varkappa\rho}r_{\varkappa\sigma}a_\sigma
 =0\quad(\varkappa\ne\rho)
\end{aligned}
\]
이다. 또
\[
 \mathfrak{r}^{\sigma\rho}\mathfrak{r}_{\rho\sigma}\mathfrak{a}_\sigma
 =\sum_{\varkappa,i}\lambda^{\rho\nu}t_{\sigma\nu}a_\nu
  \lambda_{\rho i}t_{i\sigma}a_\sigma
 =\sum_i\lambda^{\rho i}\lambda_{\rho i}a_\sigma
 =a_\sigma.
\]
역으로 마찬가지로
\[
 \mathfrak{a}_\varkappa \mathfrak{r}^{\sigma\rho}\mathfrak{b}_\rho
 =a_\varkappa\sum_\mu
  \lambda^{\rho\mu}t_{\sigma\mu}a_\mu
 =0\quad(\varkappa\ne\sigma)
\]
이고
\[
 \mathfrak{r}_{\rho\sigma}\mathfrak{r}^{\sigma\rho}\mathfrak{b}_\rho
 =\sum_{i,\mu}\lambda_{\rho i}t_{i\sigma}a_\sigma
  \lambda^{\rho\mu}t_{\sigma\mu}a_\mu
 =\sum_{i,\mu}\lambda_{\rho i}\lambda^{\rho\mu}
  t_{i\mu}a_\mu=b_\rho
\]
이다. 이제 증명한 같은 종류 관계에서 (다시 유한성정리를 쓰지 않고) \(\mathfrak A_\sigma,\mathfrak B_\rho\)의 동형이 따른다.

\(\mathfrak B_\rho,\mathfrak B_\tau\) 사이의 동형은 이를 합성하여 얻으며, 그 공식을 적어 두자. 다음과 같이 둔다.
\[
 \ell_{\rho\tau}b_\tau
 =r_{\rho\varkappa}r^{\varkappa\tau}b_\tau
 =\sum_{i,j}\lambda_{\rho i}t_{i\varkappa}a_\varkappa
  \lambda^{\tau j}t_{\varkappa j}a_j
 =\sum_{i,j}\lambda_{\rho i}\lambda^{\tau j}\mathfrak{t}_{ij}\mathfrak{a}_j.
\]
그러면 대응은 \(\mathfrak{b}_\rho\sim\mathfrak{s}_{\rho\tau}\mathfrak{b}_\tau\)이다. 여기에서도 같은 종류 관계를 형식적으로 다시 확인할 수 있다. \(\mathfrak{s}_{\rho\tau}\)들을 \(\mathfrak{t}_{i\varkappa}\) 대신 넣으면 관계 (35)를 만족한다. 실제로
\[
 b_\sigma\ell_{\rho\tau}b_\tau
 =b_\sigma r_{\rho\varkappa}r^{\varkappa\tau}b_\tau
 =0\quad(\sigma\ne\rho)
\]
이고
\[
 \ell_{\rho\sigma}\ell_{\sigma\tau}b_\tau
 =r_{\rho\varkappa}r^{\varkappa\sigma}
  r_{\sigma\lambda}r^{\lambda\tau}b_\tau
 =r_{\rho\varkappa}a_\varkappa r^{\varkappa\tau}b_\tau
 =\ell_{\rho\tau}b_\tau
\]
이다.

\(\lambda_{i\varkappa}\)들을 대각행렬이 아니고, 대각행렬을 곱하는 변환으로 서로 옮길 수도 없게 무한히 많은 방식으로 택하기만 하면 되므로 정리 VIII이 증명되었다.

정리 VIII을 정리 V, 또는 VII과 함께 쓰면 다음이 따른다.

\textbf{정리 IX}\footnote{Loewy, 106--107쪽을 보라.}. \emph{모듈 \(\mathfrak M\)이 완전가약이고, 서로 완전히 소인 계를 이루는 소모듈 \(\mathfrak P_1,\ldots,\mathfrak P_k\)의 최소공배수이며, 유리성영역 \(P\)가 상수를 무한히 많이 포함한다고 하자. \(\mathfrak M\)을 소모듈들의 최소공배수로 나타내는 서로 다른 표현이 무한히 많이 존재할 필요충분조건은 \(\mathfrak P_1,\ldots,\mathfrak P_k\) 가운데 적어도 두 모듈이 같은 종류인 것이다.}

이제 한 모듈 \(\mathfrak M\)이 무한히 많은 소모듈로 나누어지는 일반적인 판정기준을 얻기 위해 먼저 다음을 증명한다.

\textbf{보조정리}\footnote{Loewy, 96쪽을 보라.}. \emph{완전가약모듈 \(\mathfrak M=[\mathfrak P_1,\ldots,\mathfrak P_k]\)이 모든 \(\mathfrak P_i\)와 다른 소모듈 \(\mathfrak P'\)으로 나누어지면, \(\mathfrak P_i\)들 가운데 적어도 두 개는 같은 종류이다.}

먼저 \(\mathfrak P_1,\ldots,\mathfrak P_k\) 가운데 아직 지우지 않은 모듈들의 최소공배수에 포함되는 소모듈을 차례로 지워 서로 완전히 소인 계를 골라낸다. 따라서 이제 \(\mathfrak P_1,\ldots,\mathfrak P_k\) 자체가 서로 완전히 소인 계를 이룬다고 가정하자. \(a'\)를 법 \(\mathfrak P'\)의 단위원류에 대응하는 \(\mathfrak M\)의 한 류라 하자. 그러면 각 다항식 \(F\)에 대하여
\[
 F\mathfrak{a}'=F\mathfrak{a}'\mathfrak{a}_1+\cdots+F\mathfrak{a}'\mathfrak{a}_k
\]
이다. 곱 \(\mathfrak{a}'\mathfrak{a}_i\) 가운데 적어도 두 개, 이를테면 \(a'a_1,a'a_2\)는 영과 다르다. 이를테면 \(\mathfrak{a}'\mathfrak{a}_1\) 하나만 영과 다르면 \(F\mathfrak{a}'=F\mathfrak{a}'\mathfrak{a}_1\)이고, 잉여류 \(F\mathfrak{a}'\)들은 \(\mathfrak A_1\)의 부분군을 이룬다. \(\mathfrak A_1\)은 소군이므로 이 부분군은 \(\mathfrak A_1\)과 같고, 따라서 모듈 \(\mathfrak P',\mathfrak P_1\)도 일치할 것이다. 그러므로 § 8의 논법에 따라 \(\mathfrak A'\)은 \(\mathfrak A_1,\mathfrak A_2\)와 동형이다.

이제 \(\mathfrak M\)을 임의의 모듈이라 하자. \(\mathfrak M\)을 나누는 모든 소모듈의 최소공배수 \(\mathfrak N\)을 고찰한다.\footnote{그러한 소모듈이 무한히 많으면 정리 IV에 따라 이 표현에는 잉여군의 가법적 분해가 대응할 수 없다. § 12의 예 4도 보라.} 두 경우가 가능하다. 첫째, \(\mathfrak N\)을 유한개의 소모듈의 최소공배수로 나타낼 수 없다. 이 경우에는 \(\mathfrak N\), 따라서 \(\mathfrak M\)이 무한히 많은 소모듈을 인자로 갖는다. 둘째, \(\mathfrak N\)을 유한개의 소모듈의 최소공배수로 나타낼 수 있다. 이때 \(\mathfrak N\)을 Loewy의 대응 구성에 따라 \(\mathfrak M\)의 \emph{최대 완전가약모듈}이라고 한다. 이 유한개의 소모듈 가운데 같은 종류인 두 개가 있으면 정리 IX에 따라 \(\mathfrak N\), 따라서 \(\mathfrak M\)은 무한히 많은 소모듈로 나누어진다. 역으로 뒤의 사실이 성립하면, \(\mathfrak N\)을 최소공배수로 나타내는 그 유한개의 소모듈 모두와 다른 \(\mathfrak N\)의 소인자가 적어도 하나 있다. 보조정리에 따라 \(\mathfrak N\), 따라서 \(\mathfrak M\)은 같은 종류의 소인자 두 개를 갖는다. 이로써 다음 정리를 증명하였다.

\textbf{정리 X.} \emph{유리성영역 \(\mathsf{P}\)가 상수를 무한히 많이 포함하면, 한 모듈이 무한히 많은 서로 다른 소모듈로 나누어질 필요충분조건은 최대 완전가약모듈이 존재하지 않거나, 그러한 모듈이 존재할 경우 같은 종류의 서로 다른 소모듈 적어도 두 개로 나누어지는 것이다.}

\pseventeenkoheading{11}{미분방정식계와 그 적분해들과의 관계}

유한성정리에 따라 미분식들의 각 모듈에는 유한개의 미분식으로 이루어진 한 기저를 대응시킬 수 있다. 모듈의 한 미분식을 영으로 놓아 생기는 모든 미분방정식의 공통 적분해 전체는, 기저미분식들을 영으로 놓아 얻는 유한 미분방정식계의 적분해 전체와 일치한다. 이로써 선형 편미분방정식계의 이론과 연결된다. 모듈의 잉여군이 ‘유한’인 경우에 이 관계를 조금 더 추적하자.

\textbf{정의.} \emph{한 잉여군에 수 \(\mu\)가 있어서, 임의의 \(\mu+1\)개 잉여류 사이에는 \(\mathsf{P}\)의 계수를 갖는 선형관계가 있지만 아무 관계도 없는 \(\mu\)개 잉여류의 계가 적어도 하나 있으면 이 잉여군을 유한이라고 한다. \(\mu\)를 군의 차수라고 하고, 뒤의 계를 잉여군의 선형기저라고 한다. 그렇지 않으면 잉여군을 무한이라고 한다.}

두 유한 잉여군의 합은 유한하고, 그 차수는 두 차수의 합이다. 두 무한군의 합이나 한 유한군과 한 무한군의 합은 무한군이다. 한 변수 모듈의 모든 잉여군이 유한군의 예이다. 여러 변수에서도 차수는 유한할 수 있다. 예컨대 기저가 \(\xi_1^2,\xi_2^2\)인 모듈에서는 모든 잉여류가 \(1,\xi_1,\xi_2,\xi_1\xi_2\)의 잉여류들의 선형결합으로 만들어진다. 무한군의 예는 § 12를 보라.

이제 다음 정리가 성립한다.

\textbf{정리 XI.} \emph{한 모듈의 잉여군이 유한이면, 그 모듈의 한 기저를 영으로 놓아 생기는 모든 미분방정식계는 잉여군의 차수와 정확히 같은 수의 선형독립인 적분해를 갖는다. 역으로 유한개의 선형 편미분방정식으로 이루어진 주어진 계가 선형독립인 적분해를 유한개만 가지면, 그 수는 미분식들이 생성하는 모듈의 잉여군 차수와 같다.}

여기서 유리성영역 \(\mathsf{P}\)는 \(n\)개 변수 \(x_1,\ldots,x_n\)의 해석함수들로 이루어진다고 하자. 나타나는 특이점들은 어떤 \(n\)차원 영역 안에서 기껏해야 더 낮은 차원의 다양체들을 이루므로, 정칙인 \(n\)차원 근방을 갖는 점들이 얼마든지 존재한다고 한다.

\(\mathfrak M\)을 유한 잉여군을 갖는 모듈이라 하자. 그러면 다음 성질을 갖는 거듭제곱곱들의 한 선형기저가 존재함을 보이겠다. \emph{임의의 잉여류를 이 기저로 나타낼 때 분모에는 \(\mathsf{P}\)의 유한개의 서로 다른 원소들의 거듭제곱만 나타날 수 있다.}

이를 위해 모든 거듭제곱곱 \(\xi_1^{a_1}\cdots\xi_n^{a_n}\)을 사전식으로 배열하고, 앞선 것들과 법 \(\mathfrak M\)에서 선형독립인 것들을 모두 잉여류의 대표원 \(Q_1,\ldots,Q_m\)으로 택한다. 가정에 따라 이는 유한개뿐이다. 나머지는 모두 이 대표원들의 선형결합으로 나타낼 수 있다.

\(\mathfrak S\)를 선형기저에 포함되지 않은 모든 거듭제곱곱의 계라고 하자. 정리 III의 증명 첫 부분에 따라 \(\mathfrak S\)에는 유한 부분계 \(\mathfrak T\)가 있어서, \(\mathfrak S\)의 각 거듭제곱곱은 \(\mathfrak T\)의 적어도 하나의 거듭제곱곱으로 나누어진다. \(P_1,\ldots,P_\rho\)를 \(\mathfrak T\)의 거듭제곱곱들이라 하고, \(M_1\equiv0,\ldots,M_\rho\equiv0\;(\mathfrak M)\)을 각각 \(P_1,\ldots,P_\rho\)를 법 \(\mathfrak M\)에서 더 낮은 거듭제곱곱들의 선형결합으로 나타내는 관계들이라 하자. 예컨대
\[
 M_i=b_iP_i+\text{낮은 항들}
\]
이다. \(P\)가 \(\mathfrak S\)의 임의의 거듭제곱곱이면
\[
 \xi_1^{h_1}\cdots\xi_n^{h_n}P_i
\]
꼴이고,
\[
\begin{aligned}
 \xi_1^{h_1}\cdots\xi_n^{h_n}M_i
 &=b_i\xi_1^{h_1}\cdots\xi_n^{h_n}P_i
   +\text{낮은 항들}\\
 &=b_iP+\text{낮은 항들}
\end{aligned}
\]
이 따른다. \(P\)보다 낮은 모든 거듭제곱곱에서는 분모에 \(b_1,\ldots,b_\rho\)의 거듭제곱만 나타나고, 분자는 \(M_i\)의 계수들에서 덧셈·곱셈·미분으로 만들어진다고 이미 증명되었다고 하자. 그러면 \(P\) 자체에도 이 사실이 성립한다. \(\mathfrak S,\mathfrak T\)의 첫 원소 \(P_1\)에 대해서는 이 가정이 충족되므로 이러한 귀납가정이 허용된다.

또 정리 III에 따라 \(M_1,\ldots,M_\rho\)는 \(\mathfrak M\)의 모듈기저이다. 이제 \(b_1\ne0,\ldots,b_\rho\ne0\)이고 \(M_i\)의 나머지 계수들이 정칙인 \(x_1,\ldots,x_n\)의 한 값계 \(\beta_1,\ldots,\beta_n\), 곧 미분방정식계 \(M_1=0,\ldots,M_\rho=0\)의 한 정칙점을 고찰하자. 거듭제곱곱 \(Q_i=\xi_1^{h_{i1}}\cdots\xi_n^{h_{in}}\) (\(i=1,\ldots,m\))에는 편도함수
\[
 \frac{\partial^{h_{i1}+\cdots+h_{in}}}
 {\partial x_1^{h_{i1}}\cdots\partial x_n^{h_{in}}}y
\]
가 대응한다. 점 \((\beta_1,\ldots,\beta_n)\)에서 이 편도함수들에 임의의 상수값 \((\gamma_1,\ldots,\gamma_m)\)을 준다. 그러면 미분방정식 \(M_i=0\)에 따라 이 점에서 \(y\)의 모든 고계도함수 값이 유한한 값으로 유일하게 정해진다. 분모에는 \(b_i\)들만 나타나기 때문이다. 여기서 선형독립인 해석적 적분해 \(m\)개가 존재함을 보이는 것은 잘 알려져 있다.

한편 우리의 가정 아래에서는 선형독립인 적분해가 \(m\)개보다 많을 수 없다. 그렇다면 한 정칙점에서 적절히 고른 \(m+1\)개 도함수의 값을 임의로 줄 수 있다는 알려진 결론이 나오는데, 이는 임의의 \((m+1)\)개 거듭제곱곱 사이에 법 \(\mathfrak M\)의 선형종속관계가 있어야 한다는 사실과 모순이다.

역으로 한 모듈의 잉여군이 무한이면 같은 방식으로 선형독립인 적분해가 무한히 많이 존재함을 보인다. 따라서 선형독립인 적분해가 \(m\)개뿐이면 잉여군은 유한이고, 위에서 본 바에 따라 그 차수는 \(m\)이다.

이 절의 끝으로, 미분식들로 이루어진 두 모듈의 동형개념, 또는 같은 뜻인 같은 종류라는 개념이 대응 미분방정식계의 적분해들에 대하여 Poincaré가 한 변수에서 쓴 종류개념과 같은 뜻을 가짐을 보인다(앞서 든 Blumberg의 논문, 부록 I을 보라).

\textbf{정리 XII.} \emph{미분식들로 이루어진 두 모듈 \(\mathfrak M,\mathfrak N\)이 같은 종류이면, 두 미분식 \(P,Q\)가 있어서 \(y\)가 \(\mathfrak M\)의 모든 적분해를 돌 때 \(P(y)\)가 \(\mathfrak N\)의 모든 적분해를 돌고, \(z\)가 \(\mathfrak N\)의 모든 적분해를 돌 때 \(Q(z)\)가 \(\mathfrak M\)의 모든 적분해를 돈다.}

실제로 가정과 § 9에 따라 알맞게 고른 두 다항식 \(P,Q\)에 대하여
\[
 MQ\equiv0(\mathfrak N)\qquad PQ\equiv1(\mathfrak N)
\]
이고
\[
 NP\equiv0(\mathfrak M)\qquad QP\equiv1(\mathfrak M)
\]
이다. 그러므로 \(NP(y)=0\)에 따라 함수 \(P(y)\)는 \(\mathfrak N\)의 적분해이다. 마찬가지로 \(Q(z)\)는 \(\mathfrak M\)의 적분해이다. \(PQ(z)=z\), \(QP(y)=y\)이므로 \(P(y)\)는 \(\mathfrak N\)의 모든 적분해를 돌고, \(Q(z)\)는 \(\mathfrak M\)의 모든 적분해를 돈다. 이때 영과 다른 각 \(y\)에는 영과 다른 \(z\)가 대응하고, 그 역도 성립한다.

\pseventeenkoheading{12}{예}

\textbf{1.} 첫 예로 계수들이 적분해들의 대칭함수로 나타나는 통상적인 2계 선형미분방정식을 고찰한다. 여기서 유리성영역 \(\mathsf{P}\)는 적분해들과 그 도함수들의 모든 유리함수로 이루어지고, 미분방정식의 한 정칙점의 충분히 작은 근방으로 제한한다. \(\mathfrak M\)을 \(y_1,y_2\)를 적분해로 갖는 모든 미분식의 집합으로 정의하고, 이 식들을 다시 한 부정원 \(\xi\)의 다항식으로 쓴다. 모듈기저는 한 원소로 이루어지고, 다항식
\[
 M=
 \begin{vmatrix}y_1&y_2\\y_1'&y_2'\end{vmatrix}\xi^2
 -\begin{vmatrix}y_1&y_2\\y_1''&y_2''\end{vmatrix}\xi
 +\begin{vmatrix}y_1'&y_2'\\y_1''&y_2''\end{vmatrix}
\]
으로 주어진다. 이는 \(\mathsf{P}\)의 한 원소를 인자로 곱하는 차이까지만 유일하게 정해진다. 이제
\begin{equation}
 \mathfrak M=[\mathfrak M_1,\mathfrak M_2]
\tag{41}
\end{equation}
이다. 여기서 \(\mathfrak M_i\)는 \(y_i\)를 적분해로 갖는 모든 미분식으로 이루어지고, 그 기저는 \(\mathsf{P}\)의 한 인자를 곱하는 차이까지만
\[
 M_i=y_i\xi-y_i'\qquad(i=1,2)
\]
로 정해진다. 실제로 \(\mathfrak M\)은 적분해 \(y_1,y_2\)에서 사라지므로 \(\mathfrak M_1,\mathfrak M_2\)로, 따라서 \([\mathfrak M_1,\mathfrak M_2]\)로 나누어진다. \([\mathfrak M_1,\mathfrak M_2]\)의 기저다항식의 차수는 적어도 2여야 하므로 \(\mathfrak M=[\mathfrak M_1,\mathfrak M_2]\)이다. 또
\begin{equation}
 \frac{y_2}{D}(y_1\xi-y_1')
 -\frac{y_1}{D}(y_2\xi-y_2')=1,\qquad
 D=D(y)=\begin{vmatrix}y_1&y_2\\y_1'&y_2'\end{vmatrix}
\tag{42}
\end{equation}
이므로 \(\mathfrak M_1,\mathfrak M_2\)는 서로 소이다. 그 잉여군의 차수가 1이므로 이들은 소모듈이기도 하다.

그런데 \(\mathfrak M\)은 \(y_1,y_2\)뿐 아니라
\[
 z_1=\alpha_{11}y_1+\alpha_{12}y_2,\qquad
 z_2=\alpha_{21}y_1+\alpha_{22}y_2,\qquad
 A=|\alpha_{ik}|\ne0
\]
도 모두 적분해로 갖는다. 따라서 각각 \(z_1,z_2\)를 적분해로 갖는 서로 소인 두 모듈 \(\mathfrak N_1,\mathfrak N_2\)의 최소공배수와 동일하다. 그 기저다항식은
\[
 N_i=z_i\xi-z_i'
 =(\alpha_{i1}y_1+\alpha_{i2}y_2)\xi
  -(\alpha_{i1}y_1'+\alpha_{i2}y_2')
\]
이다. 그러므로 \(\mathfrak M\)은 서로 소인 소모듈들의 최소공배수로 나타내는 서로 다른 표현을 무한히 많이 갖는다. 따라서 정리 VII에 따라 \(\mathfrak N_1,\mathfrak N_2\)는 서로 같은 종류이고 \(\mathfrak M_1,\mathfrak M_2\) 각각과도 같은 종류이다. 실제로 이 모든 모듈의 잉여군은 차수가 1이므로 동형이다.

이제 표현 (41)에 대응하는 잉여군의 가법적 분해 \(\mathfrak G=\mathfrak A_1+\mathfrak A_2\)를 세우자. § 4에 따라 \(\mathfrak M_1,\mathfrak M_2\)가 서로 소임을 말하는 관계 (42)에서 시작한다. 단위원 \(a_1,a_2\)로는 각각 다항식
\[
 -\frac{y_1}{D}(y_2\xi-y_2')
 \quad\text{및}\quad
 \frac{y_2}{D}(y_1\xi-y_1')
\]
이 생성하는 잉여류들을 얻는다. 이로써
\[
 \mathfrak M_1
 =\left(\mathfrak M,\frac{y_2}{D}(y_1\xi-y_1')\right),
 \qquad
 \mathfrak M_2
 =\left(\mathfrak M,-\frac{y_1}{D}(y_2\xi-y_2')\right).
\]
마찬가지로 \(y_i\)를 \(z_i\)로 바꾸어 분해 \(\mathfrak M=[\mathfrak N_1,\mathfrak N_2]\)에 대응하는 단위원 \(b_1,b_2\)를 만들면
\[
 -\frac{z_1}{D(z)}(z_2\xi-z_2')
 =-\frac{\alpha_{11}y_1+\alpha_{12}y_2}{AD(y)}
 \bigl((\alpha_{21}y_1+\alpha_{22}y_2)\xi
      -(\alpha_{21}y_1'+\alpha_{22}y_2')\bigr),
\]
\[
 \frac{z_2}{D(z)}(z_1\xi-z_1')
 =\frac{\alpha_{21}y_1+\alpha_{22}y_2}{AD(y)}
 \bigl((\alpha_{11}y_1+\alpha_{12}y_2)\xi
      -(\alpha_{11}y_1'+\alpha_{12}y_2')\bigr)
\]
이고, 따라서
\begin{equation}
\left\{
\begin{aligned}
 \mathfrak{b}_1&=\frac{\alpha_{11}y_1+\alpha_{12}y_2}{y_1}
 \left(\frac{\alpha_{22}}{A}\right)\mathfrak{a}_1
 +\frac{\alpha_{11}y_1+\alpha_{12}y_2}{y_2}
 \left(-\frac{\alpha_{21}}{A}\right)\mathfrak{a}_2\\
 &=\frac{z_1}{y_1}\alpha^{11}a_1
  +\frac{z_1}{y_2}\alpha^{12}a_2,\\
 \mathfrak{b}_2&=\frac{\alpha_{21}y_1+\alpha_{22}y_2}{y_1}
 \left(-\frac{\alpha_{12}}{A}\right)\mathfrak{a}_1
 +\frac{\alpha_{21}y_1+\alpha_{22}y_2}{y_2}
 \left(\frac{\alpha_{11}}{A}\right)\mathfrak{a}_2\\
 &=\frac{z_2}{y_1}\alpha^{21}a_1
  +\frac{z_2}{y_2}\alpha^{22}a_2
\end{aligned}\right.
\tag{43}
\end{equation}
이다. 여기서 \((\alpha^{ik})\)는 행렬 \((\alpha_{ik})\)의 역행렬이다.

역으로 \(z_i\)와 \(y_i\)를 서로 바꾸면
\begin{equation}
\left\{
\begin{aligned}
 a_1&=\frac{y_1}{z_1}\alpha_{11}b_1
     +\frac{y_1}{z_2}\alpha_{12}b_2,\\
 a_2&=\frac{y_2}{z_1}\alpha_{21}b_1
     +\frac{y_2}{z_2}\alpha_{22}b_2
\end{aligned}\right.
\tag{44}
\end{equation}
를 얻는다. 이 공식들에서 군 \(\mathfrak A,\mathfrak B\)의 동형을 알아볼 수 있다. 실제로 우리의 일반적인 논의에 따라 \(\mathfrak{a}_1\mathfrak{b}_\sigma\ne0\)이면 대응은 \(\mathfrak{a}_1\sim \mathfrak{a}_1\mathfrak{b}_\sigma\)이고, 따라서 (44)에 따라
\[
 \mathfrak{a}_1\sim\frac{y_1}{z_\sigma}\alpha_{1\sigma}\mathfrak{b}_\sigma,
 \qquad\text{\(\alpha_{1\sigma}\ne0\)이면}
\]
이다. (43)에 따라 마찬가지로
\[
 \mathfrak{b}_\sigma\sim\frac{z_\sigma}{y_2}\alpha^{\sigma2}\mathfrak{a}_2,
 \qquad\text{\(\alpha^{\sigma2}\ne0\)이면}
\]
이고, 따라서
\[
 \frac{y_1}{z_\sigma}\alpha_{1\sigma}\mathfrak{b}_\sigma
 \sim
 \frac{y_1}{z_\sigma}\alpha_{1\sigma}
 \frac{z_\sigma}{y_2}\alpha^{\sigma2}a_2
 =\frac{y_1}{y_2}\alpha_{1\sigma}\alpha^{\sigma2}\mathfrak{a}_2.
\]
그러므로
\[
 \mathfrak{a}_1\sim\alpha_{1\sigma}\alpha^{\sigma2}\frac{y_1}{y_2}\mathfrak{a}_2
 \quad\text{\(\alpha_{1\sigma}\alpha^{\sigma2}\ne0\)이면},
 \qquad
 \frac{1}{\alpha_{1\sigma}\alpha^{\sigma2}}
 \frac{y_2}{y_1}a_1\sim a_2,
\]
이고, 뒤 관계는 역대응이다. \((\alpha_{i\varkappa})\)가 대각행렬이 아니면 어떤 \(\sigma\)에 대해서는 조건 \(\alpha_{1\sigma}\alpha^{\sigma2}\ne0\)이 언제나 성립한다. 대각행렬인 경우에는 모듈들이 동일하므로 제외해야 한다. 또 예컨대 \(\alpha_{12}=0\)만 성립하면 \(\mathfrak A_1=\mathfrak B_1\)이지만 \(\mathfrak A_2=\mathfrak B_2\)는 아니다. 따라서 \(\mathfrak A_1+\mathfrak A_2=\mathfrak B_1+\mathfrak B_2\), \(\mathfrak A_1=\mathfrak B_1\)에서 \(\mathfrak A_2=\mathfrak B_2\)가 따르지 않는다(주 12를 보라). 이때 세 모듈 \(\mathfrak M_1,\mathfrak M_2,\mathfrak N_2\)는 두 개씩 서로 소이지만 서로 완전히 소인 계를 이루지는 않는다. \([\mathfrak M_1,\mathfrak M_2]\)가 \(\mathfrak N_2\)로 나누어지기 때문이다(주 15를 보라).

또
\[
 \mathfrak{t}_{12}=\alpha_{1\sigma}\alpha^{\sigma2}\frac{y_1}{y_2}\mathfrak{a}_2,
 \qquad
 t_{21}=\frac{1}{\alpha_{1\sigma}\alpha^{\sigma2}}
 \frac{y_2}{y_1}a_1
\]
로 두면, 직접 계산하여 확인할 수 있듯 관계 (35)가 충족된다.

\textbf{2.} 둘 이상의 변수를 가지면서 기저가 다항식 하나뿐인 모든 모듈은 무한군의 예이다. 다음 예는 조금 더 복잡하며, 다항식을 둘 이상 포함하는 모듈기저를 갖는 여러 원소 모듈에서도 무한군이 나타날 수 있음을 보여 준다.\footnote{앞서 든 Blumberg의 논문에 실린 Landau의 예, 51--52쪽을 보라. 여기와 이하에서는 독립변수를 \(x,y\)로 나타낸다.}
\[
 \mathfrak M=[(\xi+x\eta),(\xi+1)]=[(P),(Q)]
 \qquad\left(\xi=\frac{\partial}{\partial x},
 \ \eta=\frac{\partial}{\partial y}\right).
\]
곱셈법칙은 미분식의 법칙으로 한다. \((\xi+x\eta)\)와 \((\xi+1)\)의 잉여군들이 모두 무한이고, 모듈 \((\xi+x\eta)\), \((\xi+1)\)은
\[
 1=(-x\xi-x^2\eta+1)Q+(x\xi+x-1)P
\]
에 따라 서로 소이므로 여기서 잉여군은 무한이다.

다음과 같이 놓으면 \(\mathfrak M\)이 2차 다항식을 갖지 않음을 아주 쉽게 보일 수 있다.
\[
 (u_1\xi+u_2\eta+u_3)(\xi+x\eta)
 =(v_1\xi+v_2\eta+v_3)(\xi+1).
\]
여기서 \(u_1=u_2=u_3=v_1=v_2=v_3=0\)을 얻는다. 또 \(\mathfrak M\)의 3차 다항식을 찾으면 선형독립인 것 두 개, 예컨대
\[
 (\xi^2+x\xi\eta+\xi+(2+x)\eta)Q
 =(\xi^2+2\xi+1)P
\]
와
\[
 (\xi^2+2x\xi\eta+x^2\eta^2+\eta)Q
 =(\xi^2+x\xi\eta+\xi+(x-1)\eta)P
\]
를 얻는다. \(\mathfrak M\)이 한 원소 모듈이면 두 다항식은 모두 기저다항식으로 나누어져야 한다. 그런데 \(\mathfrak M\)에는 2차 다항식이 없으므로 3차 다항식으로 나누어져야 하고, 따라서 서로 비례해야 하지만 그렇지 않다. \(\mathfrak M\)이 여러 원소 모듈이라는 사실이 한 변수의 경우에 성립하는 Landau의 곱정리들이 실패하는 내적인 이유이다.

\textbf{3.} 다음 예는 소모듈에서도 무한군이 나타날 수 있음을 보여 준다.\footnote{우리의 ‘소모듈’ 정의는 가환인 경우 소모듈이라고 부르는 것과 다르다. 그 경우에는 다양체가 0차원, 곧 잉여군이 유한인 경우를 제외하면 언제나 더 낮은 차원의 인자가 존재하기 때문이다.} \(P\)를 두 변수 \(x,y\)의 모든 유리함수의 영역이라고 하자. 기저
\[
 \xi-\frac{y}{x}=\xi-a
\]
를 갖는 모듈 \(\mathfrak M\)의 잉여군은 \(\eta\)의 각 거듭제곱의 잉여류를 포함하므로 무한이다. 한편 \(\mathfrak M\)은 소모듈이다. 실제로 각 인자
\[
 \mathfrak M_1=(\xi-a,F(\xi,\eta))
\]
가 단위원모듈임을 보이겠다. 모든 다항식 \(F(\xi,\eta)\)은 법 \((\xi-a)\)에서 한 다항식 \(G(\eta)\)으로 나타낼 수 있다.
\[
 F(\xi,\eta)\equiv
 c_0\eta^\nu+c_1\eta^{\nu-1}+\cdots+c_\nu
 \quad(\mathfrak M).
\]
일반성을 잃지 않고 \(c_0=1\)이라 하자. 곱셈법칙에 따라 \(\xi\eta^\nu-\eta^\nu\xi=0\)이다. 따라서
\[
 \xi\equiv a(\mathfrak M_1),\qquad
 \eta^\nu\equiv F_1(\mathfrak M_1)
\]
이고 \(F_1=-(c_1\eta^{\nu-1}+\cdots+c_\nu)\)로 두면
\[
 \xi F_1-\eta^\nu a\equiv0(\mathfrak M_1)
\]
이다.

\[
 \eta^\nu a
 =a\eta^\nu+\nu\frac{\partial a}{\partial y}\eta^{\nu-1}
 +\cdots
\]
로 전개하고 다시 \(\eta^\nu\)를 \(F_1\)으로 바꾸면
\[
 \xi(c_1\eta^{\nu-1}+\cdots+c_\nu)
 -a(c_1\eta^{\nu-1}+\cdots+c_\nu)
 +\nu\frac{\partial a}{\partial y}\eta^{\nu-1}
 +\cdots\equiv0(\mathfrak M_1),
\]
또는
\[
 c_1\eta^{\nu-1}a
 +\frac{\partial c_1}{\partial x}\eta^{\nu-1}+\cdots
 -ac_1\eta^{\nu-1}+\cdots
 +\nu\frac{\partial a}{\partial y}\eta^{\nu-1}+\cdots
 \equiv0(\mathfrak M_1),
\]
마지막으로
\[
 \left(\frac{\partial c_1}{\partial x}
 +\nu\frac{\partial a}{\partial y}\right)\eta^{\nu-1}
 +\cdots\equiv0(\mathfrak M_1)
\]
을 얻는다. 왼쪽은 정확히 \(\nu-1\)차이고 항등적으로 사라지지 않는 \(\eta\)의 다항식이다. 실제로 최고계수
\[
 \frac{\partial c_1}{\partial x}
 +\nu\frac{\partial a}{\partial y}
 =\frac{\partial c_1}{\partial x}+\nu\frac1x
 \ne0
\]
이다. 그렇지 않으면 \(c_1=-\nu\log x+\varphi(y)\)가 되어 유리함수가 아닐 것이기 때문이다. 따라서 이 절차를 계속할 수 있고, 마침내 항등적으로 사라지지 않는 0차 다항식에 이른다. 이로써 단위원도 \(\mathfrak M_1\)에 속함을 보였다.\footnote{이 절차는 한 미분방정식계에 필요한 적분가능조건들이 충족되지 않음을 보이는 것으로 귀착된다.}

\textbf{4.} 반면 유리성영역에 함수 \(\log x\)를 더 첨가하면 \(\mathfrak M\)은 인자
\[
 \left(\xi-\frac{y}{x},\ \eta-\log x+\varphi(y)\right)
\]
들을 갖는다. 여기서 \(\varphi(y)\)는 \(y\)의 모든 유리함수를 돈다. 이 인자들은 모두 서로 다르고, 각각 적분가능조건을 만족하며, 대응 적분해는
\[
 \exp\!\left(y\log x-\int\varphi(y)\,dy+\alpha\right)
\]
이다. 그러므로 단위원은 모듈에 속할 수 없다. 그렇다면 영함수가 유일하게 가능한 적분해일 것이기 때문이다. 한편 \(\xi,\eta\)는 이 모듈에 관하여 영역의 어떤 원소들과 합동이다. 따라서 잉여군은 유한이고 차수 1이다. 또 무한히 많은 인자 가운데 임의의 두 개는 서로 소이며, 주어진 모듈 \(\mathfrak M\)은 이 모든 인자의 최소공배수 \(\mathfrak N\)과 같다. 실제로 \(\mathfrak M\)은 모든 인자로, 따라서 \(\mathfrak N\)으로 나누어진다. \(\mathfrak N\)이 \(\mathfrak M\)의 진인자이면 \(\mathfrak N\)은 \(\eta\)에만 의존하는 다항식 \(F\), 이를테면 차수 \(\rho\)인 것을 적어도 하나 더 포함해야 한다. 이제 \(\varphi(y)\)가 함수 \(y,\ldots,y^{\rho+1}\)을 돌게 하면 대응 적분해
\[
 \exp\!\left(y\log x-\frac{y^{i+1}}{i+1}+\alpha_i\right)
 \qquad(i=1,\ldots,\rho+1)
\]
사이에는 \(y\)와 무관한 계수를 갖는 선형관계가 없다. 그러나 \(\rho\)계 미분방정식 \(F=0\)은 선형독립인 적분해 \(\rho+1\)개를 가질 수 없다.

그런데 모듈 \(\mathfrak M=\mathfrak N\)은 완전가약이 아니다.\footnote{이로써 정리 X의 첫 경우에 해당하는 예를 얻었다.} 그렇다면 유한개의 이 소모듈의 최소공배수일 것이고 그 잉여군도 유한차수일 것이다. 그러나 \(\xi-\frac yx\)의 잉여군은 \(\eta\)의 모든 거듭제곱의 잉여류를 포함하므로 무한이다.

괴팅겐, 1919년 8월 1일.

\begin{center}
(1919년 8월 4일 접수.)
\end{center}

\endgroup

\clearpage

\clearpage
\section*{논문 17에 대한 편집 주석}
\addcontentsline{toc}{section}{논문 17에 대한 편집 주석}
\begin{enumerate}
\item 원문 19쪽, 제8절: 제12절 예 2에 대한 참조를 예 3에 대한 참조로 읽는다. 여기에서 필요한 소모듈의 무한군은 예 3에서 다룬다. 예 2는 무한군을 다루지만 이 추가 주장을 포함하지 않는다.
\item 원문 34쪽, 예 3: 인쇄본의 $y$의 거듭제곱을 이 교정판에서는 $\eta$의 거듭제곱으로 읽는다. 제11절의 유한성은 $P$ 위에서의 선형종속에 관한 것이다. $y\in P$이므로 $y$의 거듭제곱은 길이가 임의로 긴 $P$ 위의 선형독립인 족을 줄 수 없다. $\eta$라는 독법은 $G(\eta)$로의 환원 및 35쪽 마지막 논증의 $\eta$의 거듭제곱과 일치한다.
\item 원문 35쪽, 예 4: 제시된 두 퍼텐셜을 각각 지수함수의 지수로 읽는다. 국소적으로 로그의 가지와 원시함수를 택하고 다음과 같이 둔다.\[
 F=y\log x-\Phi(y)+c,\quad \Phi'(y)=\varphi(y),\quad u=e^F,
\]
\[
 A=\partial_x-y/x,\qquad B=\partial_y-\log x+\varphi(y),
 \qquad Au=Bu=0.
\]
\[
 AF=(y/x)(1-F),\qquad BF=(\log x-\varphi(y))(1-F).
\]
따라서 $u$는 두 방정식을 만족하지만 $F$는 일반적으로 만족하지 않는다. $\varphi(y)=y^i$이면 $\Phi(y)=y^{i+1}/(i+1)$로 둔다. 예를 들어 $x>0$이고 극점을 포함하지 않는 실수 직사각형 영역에서 이 계산이 성립한다. 지수함수는 영이 되지 않으며, $u=0$도 해이다. 더 일반적으로 임의의 상수 $C$에 대하여 $C\exp(y\log x-\Phi(y))$는 해이다. 이 항등식들은 두 지수함수 독법의 근거이며, 뒤따르는 선형독립성 논증을 증명하지는 않는다.
\end{enumerate}
\smallskip
\noindent{\small AI가 제작한 교정 주석과 정정: OpenAI Codex, GPT-6 Astra, Ultra 추론 강도.}

\clearpage
\setcounter{footnote}{0}
\phantomsection
\addcontentsline{toc}{section}{18. K. 헨첼트의 소거론 논문에 관하여}
%
%
%

\begin{center}
\textbf{제11차 회의: 1921년 9월 23일 금요일 오전 11시.}\\
(의장: 뢰비(Loewy).)
\end{center}

\begin{center}
\textbf{1. E. 뇌터(Noether): 전쟁 중 전사한 K. 헨첼트(Hentzelt)의 소거론 논문에 관하여.}
\end{center}

쿠르트 헨첼트는 1914년 10월 이래 딕스뮈데 전선에서 실종되었으며, 전사한 것으로 보아야 한다. 그는 에를랑겐 학위논문의 핵심 부분을 빈틈없이 구성해 놓았지만 거의 이해할 수 없게 서술한 채 남겼다. 나는 이를 자유롭게 다시 정리하여 조만간 \emph{Mathematische Annalen}에 발표하고자 한다.

문제는 소거론의 \emph{일반 문제}, 곧 아이디얼 $\mathfrak M$에 속하는 모든 다항식의 공통 영점을 결정하는 것이다. 먼저 변수들에 부정계수를 갖는 선형변환을 가하면, 주된 결과는 다음과 같이 말할 수 있다.

아이디얼 $\mathfrak M$에는 다음과 같은 종결식형이 유일하게 대응한다.
\par\smallskip
\noindent\makebox[\linewidth][c]{\(\displaystyle
R_{\mathfrak M}=R^{(1)}(x_1\ldots x_n)\cdots R^{(i)}(x_i\ldots x_n)\cdots R^{(n)}(x_n)\equiv 0\pmod{\mathfrak M}.
\)}
\par\smallskip\noindent
이 종결식형 $R_{\mathfrak M}$은 $\mathfrak M$의 모든 영점에서, 그리고 오직 그 영점들에서만 사라진다. $\mathfrak N$이 $\mathfrak M$의 약수이고 $R_{\mathfrak N}$과 $R_{\mathfrak M}$이 일치하면, $\mathfrak N$과 $\mathfrak M$ 자체도 일치한다.

정리의 둘째 부분은 종결식형이 각 영점을 그 고유한 중복도와 함께 준다는 것을 보여 준다. 그 증명은 $\mathfrak M$을 $x_1,\ldots x_{i-1}$의 단항식들에 관한 선형형식 모듈로 볼 수 있다는 사실에 기초한다. 이때 $x_{i+1},\ldots x_n$을 계수영역에 첨가하면, $R_{\mathfrak M}$의 각 인자 $R^{(i)}(x_i\ldots x_n)$는 이 모듈에 대한 해당 ``기초모듈''의 노름이 된다. 추상 아이디얼론을 이용하면 이를 다음과 같이 해석할 수 있다. $\mathfrak M=[\mathfrak Q,\mathfrak Q_1,\ldots\mathfrak Q_r]=[\mathfrak Q,\mathfrak L]$이 $\mathfrak M$의 으뜸아이디얼 분해이고 $\mathfrak L$이 $\mathfrak Q$의 보충부분이라 하자. 그러면 아이디얼 $\mathfrak Q$에는 $R^{(i)}(x_i\ldots x_n)$의 으뜸인자 $Q^{(i)}(x_i\ldots x_n)$가 대응하며, 이 인자는 위의 의미에서 $\mathfrak Q$에 대한 $\mathfrak L$의 노름을 나타낸다.

\clearpage
\setcounter{footnote}{0}
\phantomsection
\addcontentsline{toc}{section}{19. 환영역의 아이디얼론}
%
%
%
%

\begin{center}
{\Large\bfseries 환영역의 아이디얼론.}\par
\vspace{0.7em}
지은이\par
\vspace{0.25em}
괴팅겐의 에미 뇌터.\par
\vspace{0.9em}
\rule{4em}{0.4pt}\par
\vspace{0.9em}
{\large\bfseries 목차.}
\end{center}
\noindent 서론.\par
\noindent\hangindent=3.0em\hangafter=1 \S~1. \textbf{환영역, 아이디얼, 유한조건.}\par
\noindent\hangindent=3.0em\hangafter=1 \S~2. \textbf{한 아이디얼을 유한개의 기약아이디얼의 최소공배수로 나타내기.}\par
\noindent\hangindent=3.0em\hangafter=1 \S~3. \textbf{기약아이디얼로의 서로 다른 두 분해에서 성분수의 일치.}\par
\noindent\hangindent=3.0em\hangafter=1 \S~4. \textbf{으뜸아이디얼. 기약아이디얼로의 서로 다른 두 분해에서 대응 소아이디얼의 유일성.}\par
\noindent\hangindent=3.0em\hangafter=1 \S~5. \textbf{한 아이디얼을 극대 으뜸아이디얼들의 최소공배수로 나타내기. 대응 소아이디얼의 유일성.}\par
\noindent\hangindent=3.0em\hangafter=1 \S~6. \textbf{한 아이디얼을 상대소-기약 아이디얼들의 최소공배수로 유일하게 나타내기.}\par
\noindent\hangindent=3.0em\hangafter=1 \S~7. \textbf{고립아이디얼의 유일성.}\par
\noindent\hangindent=3.0em\hangafter=1 \S~8. \textbf{한 아이디얼을 서로소-기약 아이디얼들의 곱으로 유일하게 나타내기.}\par
\noindent\hangindent=3.0em\hangafter=1 \S~9. \textbf{모듈로의 고찰 확장. 기약모듈로의 분해에서 성분수의 일치.}\par
\noindent\hangindent=3.0em\hangafter=1 \S~10. \textbf{다항식영역의 특수한 경우.}\par
\noindent\hangindent=3.0em\hangafter=1 \S~11. \textbf{수론과 미분식 이론의 예.}\par
\noindent\hangindent=3.0em\hangafter=1 \S~12. \textbf{초등인자론의 예.}\par

\begin{center}\textbf{서론.}\end{center}
이 논문의 내용은 \emph{유리정수의 분해정리, 또는 대수적 수체에서 아이디얼의 분해정리를 임의의 정역, 더 일반적으로는 환영역의 아이디얼로 옮기는 것}이다. 이 옮김을 이해하기 위해 먼저 유리정수의 분해정리를 통상적인 정식화와 조금 다른 꼴로 적어 두자.

식
\[
 a=p_1^{\varrho_1}p_2^{\varrho_2}\cdots p_\sigma^{\varrho_\sigma}=q_1q_2\cdots q_\sigma
\]
에서 소수의 거듭제곱 $q_i$들을 분해의 성분으로 보면, 이 성분들은 다음과 같은 특징을 갖는다.

1. 이들은 \emph{쌍마다 서로소}이다. 그러나 어느 $q$도 쌍마다 서로소인 수들의 곱으로 나타낼 수 없으므로, 이 뜻에서 기약성도 성립한다. 쌍마다 서로소라는 사실에서 곱 $q_1\cdots q_\sigma$가 최소공배수 $[q_1\cdots q_\sigma]$와 같다는 것도 따른다.

2. 임의의 두 성분 $q_i,q_k$는 \emph{상대소}이다. 곧 $b\cdot q_i$가 $q_k$로 나누어지면 $b$가 $q_k$로 나누어진다. 이 뜻에서도 기약성이 성립한다.

3. 각 $q$는 \emph{으뜸}이다. 곧 곱 $b\cdot c$가 $q$로 나누어지지만 $b$는 나누어지지 않으면, $c$의 어떤 거듭제곱\footnote{이 거듭제곱을 언제나 첫째 거듭제곱으로 잡을 수 있으면, 잘 알려진 대로 소수를 얻는다.}이 $q$로 나누어진다. 더 나아가 서로 다른 두 $q$의 곱은 더 이상 으뜸이 아니므로, 이 표현은 \emph{극대 으뜸성분들}에 의한 표현이다. 극대 으뜸성분들로의 분해에 관해서도 $q$들은 기약이다.

4. 각 $q$는 두 진약수의 최소공배수로 나타낼 수 없다는 뜻에서 \emph{기약}이다.

이 으뜸수 $q$와 소수 $p$의 관계는 다음과 같다. 각 $q$마다 $q$의 약수이면서 어떤 거듭제곱이 $q$로 나누어지는 $p$가 하나, 그리고 부호를 제외하면 오직 하나 존재한다. 이것이 대응 소수이다. $p^\varrho$가 그런 거듭제곱 가운데 가장 낮은 것, 곧 $\varrho$가 $q$의 지수이면, 여기서는 특히 $p^\varrho=q$이다. 이제 \emph{유일성 정리}는 다음과 같이 말할 수 있다.

\emph{한 유리정수를 기약인 극대 으뜸성분 $q$들로 나타내는 서로 다른 두 분해에서는 성분수, 대응 소수들(부호를 제외하고), 그리고 지수들이 일치한다. $p^\varrho=q$이므로 $q$들 자체도 부호를 제외하고 일치한다.}

수 대신 그 수에서 유도되는 아이디얼, 곧 $a$로 나누어지는 모든 수를 고찰하면 부호에서 생기는 부정성이 사라진다는 것은 잘 알려져 있다. 그러면 같은 정식화가 (유한) 대수적 수체의 아이디얼을 소아이디얼의 거듭제곱으로 유일하게 분해하는 정리에 그대로 적용된다.

이하에서는 (\S~1) 각 아이디얼이 유한 아이디얼기저를 갖는다는 \emph{유한조건}만을 만족하는 일반 환영역을 바탕으로 삼는다. 이러한 유한조건이 없으면 기약아이디얼과 소아이디얼이 존재하지 않을 수도 있다. 실제로 \emph{모든 대수적 정수}의 영역에는 \emph{소아이디얼로의 분해가 없다}.

성분 $q$의 네 가지 특징에 대응하여 일반적으로는 \emph{서로 구별되는 네 분해가 존재하며, 각 분해는 앞의 분해를 더 세분하여 얻어짐}이 드러난다. 서로소-기약 아이디얼로의 분해는 곱표현이고, 나머지 세 분해는 최소공배수에 의한 축약표현(\S~2)이다. 으뜸아이디얼과 대응 소아이디얼 사이의 \emph{관계}도 유지된다. 기약아이디얼도 으뜸아이디얼이다. 각 으뜸아이디얼 $\ideal Q$에는 $\ideal Q$의 약수이고 그 어떤 거듭제곱이 $\ideal Q$로 나누어지는 대응 소아이디얼 $\ideal P$가 유일하게 정해진다. $\ideal P^\varrho$가 그러한 거듭제곱 가운데 가장 낮은 것, 곧 $\varrho$가 $\ideal Q$의 지수라고 해도 여기서는 $\ideal P^\varrho$와 $\ideal Q$가 일치할 필요가 없다. 이제 \emph{유일성 정리}는 다음과 같이 말한다.

\emph{분해 1과 2는 유일하다. 분해 3 또는 4의 서로 다른 두 표현에서는 성분수와 대응 소아이디얼들이 일치한다}\footnote{아마 지수들도 일치하며, 더 일반적으로 대응 성분들이 동형이라는 사실도 성립할 것이다.}\emph{. 성분들 가운데 나타나는 고립아이디얼(\S~7)은 유일하게 정해진다.}

분해정리를 증명하기 위해 먼저 유한조건에서 데데킨트가 유한 수모듈에 대하여 처음 말한 ``유한사슬 정리''를 이끌어 내고, 그로부터 각 아이디얼을 유한개의 기약아이디얼의 최소공배수로 나타내는 분해 4를 얻는다. 한 성분의 가약성 개념을 바꾸어 쓰면 \emph{기약아이디얼로의 분해 4에 관한 기본 유일성 정리}가 나온다. 이어서 매번 유한개의 성분을 묶어 나머지 분해들로 올라가며, 그 유일성 정리들은 분해 4의 유일성 정리에서 따라온다.

끝으로 (\S~9) 유한개의 기약성분에 의한 표현은 더 약한 가정 아래에서도 성립함을 보인다. 환영역의 가환성은 필요하지 않고, 아이디얼 대신 그 영역에 관한 모듈을 고찰해도 충분하다. 이 더 일반적인 경우에도 서로 다른 두 분해에서 성분수는 일치한다. 반면 소와 으뜸이라는 개념은 가환성과 아이디얼 개념에 매여 있다. 그러나 비가환 영역의 아이디얼에서는 서로소라는 개념이 그대로 남는다.

실제로 서로 구별되는 네 분해가 나타나는 가장 간단한 환영역은 임의의 복소수 계수를 갖는 $n$변수 다항식 전체의 영역이다. 여기서는 각 분해를 대수적 형상의 거동으로 무리적으로 해석할 수 있고, 대응 소아이디얼에 관한 유일성 정리는 대수적 형상을 기약성분들로 유일하게 분해할 수 있다는 소거이론의 기본정리에 대응한다. 더 일반적인 예는 다항식으로 이루어진 모든 유한 정역이다(\S~10). 그러나 모든 짝수의 단순한 영역, 더 일반적으로 한 고정된 수로 나누어지는 모든 수의 영역에서도 이미 일부 분리된 분해의 예를 볼 수 있다(\S~11). 비가환 영역의 아이디얼론의 한 예는 초등인자론(\S~12)이 준다. 거기서는 아이디얼 또는 클래스가 기약아이디얼 또는 기약클래스로 유일하게 분해된다. 이 기약클래스들은 초등인자의 기약성분을 완전히 특징짓고, 보통의 초등인자론이 성립하지 않는 영역에서는 그 등가물로 볼 수도 있을 것이다.

기존 문헌에 관해서는 다음을 말해 둘 필요가 있다. 임의의 복소수 계수 또는 정수 계수를 갖는 다항식영역에서 극대 으뜸아이디얼로의 분해는 라스커가 제시했고, 매콜리가 몇몇 점에서 더 발전시켰다.\footnote{E. Lasker, \emph{Zur Theorie der Moduln und Ideale}. Math. Ann. 60 (1905), S. 20, Satz VII und XIII. -- F. S. Macaulay, \emph{On the Resolution of a given Modular System into Primary Systems including some Properties of Hilbert Numbers}. Math. Ann. 74 (1913), S. 66.} 두 사람은 모두 소거이론, 곧 다항식이 기약다항식들의 곱으로 유일하게 표현된다는 사실을 이용한다. 그러나 대수적 수체의 아이디얼론에서 짐작할 수 있고 이 논문이 보여 주듯, 아이디얼의 분해정리들은 실제로 이 가정과 무관하다. 라스커와 매콜리는 으뜸아이디얼도 소거이론의 개념들을 바탕으로 정의한다.

기약아이디얼로의 분해와 상대소-기약 아이디얼로의 분해는 다항식영역에 대해서도 문헌에서 주목되지 않은 듯하다. 매콜리에게서 고립된 으뜸아이디얼의 유일성에 관한 한 논평만을 찾을 수 있다.

다항식영역에서 서로소-기약 아이디얼로의 분해는 슈마이들러\footnote{W. Schmeidler, \emph{Über Moduln und Gruppen hyperkomplexer Größen}. Math. Zeitschr. 3 (1919), S. 29.}가 유한성 증명에 소거이론을 사용하여 제시했다. 그러나 거기서는 유일성 정리가 아이디얼 자체가 아니라 아이디얼의 클래스에 대해서만 진술된다. 아이디얼 자체에 관한 이 유일성 정리는 비가환 다항식영역의 아이디얼을 다룬 공동논문\footnote{E. Noether -- W. Schmeidler, \emph{Moduln in nichtkommutativen Bereichen, insbesondere aus Differential- und Differenzenausdrücken}. Math. Zeitschr. 8 (1920), S. 1.}에 있다. 그곳에서는 유한 아이디얼기저만을 사용하므로 정리와 방법은 일반 환영역에서도 유지되며, 이 논문에 의해 성분수의 일치라는 점에서 더 날카롭게 만들 수 있다(\S~11). 이 논문의 고찰은 두 앞선 논문의 개념구성을 크게 일반화하고 발전시킨 것이다. 두 논문의 핵심은 최소공배수에 의한 표현에서 잉여류계의 가법적 분해로 넘어가는 데 있다. 여기서는 서술을 간단히 하려고 다시 최소공배수를 사용한다. 그러면 가법적 분해에는 가약성 개념을 보충부분의 성질로 바꾸어 쓰는 것이 대응한다(\S~3). 그러나 본질적으로 그 공동논문과 같은 고찰을 거치면, 여기에 제시한 모든 정리를 잉여류계와 어떤 부분계의 가법적 분해정리로 볼 수도 있다. 이 잉여류계는 처음에 바탕으로 삼은 환과 똑같이 일반적인 환을 이룬다. 실제로 임의의 환은 환원소들 사이의 모든 항등관계 전체에 대응하는 아이디얼의 잉여류계로 볼 수 있다. 또는 나머지 관계들을 영역 안에서도 성립한다고 가정하면, 이 관계들의 한 부분계에 대응하는 아이디얼의 잉여류계로도 볼 수 있다.

이 논평은 프렝켈의 연구들이 어디에 놓이는지도 보여 준다.\footnote{A. Fraenkel, \emph{Über die Teiler der Null und die Zerlegung von Ringen}. J. f. M. 145 (1914), S. 139. \emph{Über gewisse Teilbereiche und Erweiterungen von Ringen}. Habilitationsschrift, Leipzig, Teubner, 1916. \emph{Über einfache Erweiterungen zerlegbarer Ringe}. J. f. M. 151 (1920), S. 121.} 프렝켈은 네 분해가 일치하도록 제한조건(정칙원소의 존재, 그 원소들에 의한 나눗셈, 분해가능성 조건)을 부과한 환의 가법적 분해를 고찰한다. 네 분해가 일치하기 때문에, 아이디얼이 유한개의 진약수만을 가져야 한다는 그의 유한조건도 --- 그는 일부 경우 이 조건을 쓰지 않기도 한다 --- 우리의 조건보다 더 강한 제한이 아니다. 프렝켈의 출발점은 본질적으로 대수적인 다른 목표들 때문에 정해진 것이며, 그는 대수적 확대를 통해 제한조건이 더 적은 일반적인 환들에 이른다.

\begin{center}
\S~1.\\[0.25em]
\textbf{환영역, 아이디얼, 유한조건.}
\end{center}

\textbf{1.} 바탕 영역 $\Sigma$는 추상적으로 정의된 (가환) \emph{환}이라 하자.\footnote{이 정의는 프렝켈의 교수자격논문에서 제한조건 6, I, II를 빼고 가져온 것이다. 그 대신 덧셈의 교환법칙을 포함해야 했다. 따라서 이는 곱셈의 가역성을 제외한 체의 정의법칙들이다.} 곧 $\Sigma$는 원소 $a,b,c,\ldots,f,g,h,\ldots$의 계이며, 그 안에는 통상적인 조건들을 만족하는 한 관계가 \emph{등식}으로 정의되어 있다. 또 \emph{덧셈}과 \emph{곱셈}이라는 \emph{두} 연산(결합방식)에 의해 임의의 두 환원소 $a,b$에서 각각 합 $a+b$와 곱 $a\cdot b$라는 제3의 원소가 언제나 유일하게 얻어진다. \emph{환}과, 그 밖에는 \emph{완전히} 임의적인 이 \emph{연산들}은 다음 \emph{법칙들}을 만족해야 한다.
\begin{enumerate}
\item \emph{덧셈의 결합법칙: $(a+b)+c=a+(b+c)$.}
\item \emph{덧셈의 교환법칙: $a+b=b+a$.}
\item \emph{곱셈의 결합법칙: $(a\cdot b)\cdot c=a\cdot(b\cdot c)$.}
\item \emph{곱셈의 교환법칙: $a\cdot b=b\cdot a$.}
\item \emph{분배법칙: $a\cdot(b+c)=a\cdot b+a\cdot c$.}
\item \emph{제한 없이 유일한 뺄셈의 법칙.} 방정식 $a+x=b$를 만족하는 원소 $x$가 $\Sigma$에 오직 하나 존재한다. (이를 $x=b-a$로 나타낸다.)
\end{enumerate}
이 성질들에서 영원의 존재가 따른다. 그러나 환에 단위원이 있을 필요는 없으며, 어느 인자도 영이 아닌데 두 원소의 곱이 영이 될 수도 있다. 곱이 영이면 언제나 한 인자가 영이고, 그뿐 아니라 단위원도 갖는 환을 \emph{진정한 정역}이라 한다. 유한합 $a+a\ldots+a$에는 통상적인 약기호 $na$를 도입한다. 이때 \emph{정수 $n$은 환원소가 아니라 단지 약기호로 보아야 하며}, $a=1\cdot a$, $na+a=(n+1)a$로 재귀적으로 정의된다.

\textbf{2.} $\Sigma$의 한 \emph{아이디얼} $\ideal M$\footnote{아이디얼은 대문자 독일체로 나타낸다. $\ideal M$은 보통 ``모듈'' 또는 형식모듈이라 불리는 다항식 아이디얼의 예를 떠올리게 하려고 택한 기호이다. 덧붙여 \S\S~1--3에서는 아이디얼 성질이 아니라 모듈 성질만을 사용한다. 이에 관해서는 \S~9를 보라.}이란 다음 두 조건을 만족하는 $\Sigma$ 원소들의 계를 뜻한다.
\begin{enumerate}
\item \emph{$\ideal M$이 $f$를 포함하고 $a$가 $\Sigma$의 임의의 원소이면, $\ideal M$은 $a\cdot f$도 포함한다.}
\item \emph{$\ideal M$이 $f,g$를 포함하면 차 $f-g$도 포함한다.} 따라서 $f$와 함께 모든 정수 $n$에 대한 $nf$도 포함한다.
\end{enumerate}
$f$가 $\ideal M$의 원소이면 통상처럼 $f\equiv0(\ideal M)$으로 나타내고, \emph{$f$가 $\ideal M$로 나누어진다}고 말한다. $\ideal N$의 모든 원소가 동시에 $\ideal M$의 원소, 곧 $\ideal M$로 나누어질 때 \emph{$\ideal N$이 $\ideal M$로 나누어진다}고 말하며 $\ideal N\equiv0(\ideal M)$으로 쓴다. $\ideal M$이 $\ideal N$에 없는 원소를 포함하여 역으로 $\ideal N$으로 나누어지지 않으면, $\ideal M$을 $\ideal N$의 \emph{진약수}라 한다. $\ideal N\equiv0(\ideal M)$과 $\ideal M\equiv0(\ideal N)$에서 $\ideal N=\ideal M$이 따른다.

그 밖의 익숙한 개념들도 문자 그대로 유지된다. 두 아이디얼 $\ideal A,\ideal B$의 \emph{최대공약아이디얼} $\ideal D=(\ideal A,\ideal B)$는 $a$가 $\ideal A$의 모든 원소를, $b$가 $\ideal B$의 모든 원소를 돌 때 $a+b$ 꼴로 나타나는 원소 전체를 뜻한다. $\ideal D$는 다시 아이디얼이다. 무한히 많은 아이디얼의 최대공약아이디얼 $\ideal D=(\ideal A_1,\ideal A_2,\ldots,\ideal A_\nu,\ldots)$도 매번 유한개의 아이디얼에서 택한 원소들의 합 $d=a_{i_1}+a_{i_2}+\cdots+a_{i_n}$으로 나타나는 모든 원소 $d$의 전체로 정의된다. 여기서도 $\ideal D$는 다시 아이디얼이다.

특히 아이디얼 $\ideal M$이 유한개의 원소 $f_1,f_2,\ldots,f_\varrho$를 포함하고
\[
 \ideal M=(f_1\ldots f_\varrho);\qquad \text{곧}\qquad
 f=a_1f_1+\cdots+a_\varrho f_\varrho+n_1f_1+\cdots+n_\varrho f_\varrho
\]
가 모든 $f\equiv0(\ideal M)$에 대하여 성립한다고 하자. 여기서 $a_i$들은 \emph{환영역}의 원소이고 $n_i$들은 정수이다. 그러면 $\ideal M$을 \emph{유한아이디얼}, $f_1,\ldots,f_\varrho$를 한 \emph{아이디얼기저}라 한다.

\emph{이하에서는 다음 유한조건을 만족하는 환 $\Sigma$만을 바탕으로 삼는다. $\Sigma$의 각 아이디얼은 유한아이디얼, 곧 아이디얼기저를 갖는 아이디얼이다.}

\textbf{3.} 유한조건에서 이후의 모든 고찰의 바탕이 되는 다음 정리가 바로 따른다.

\textbf{정리 I (유한사슬 정리).}\footnote{데데킨트가 수모듈에 대하여 처음 진술하였다. \emph{Zahlentheorie}, Suppl. XI, \S~172, Satz VIII (4판). 우리의 증명과 ``사슬''이라는 명칭은 그곳에서 가져왔다. 다항식 아이디얼에 대해서는 라스커, 앞의 글 56쪽(보조정리)에 있다. 그러나 두 경우 모두 이 정리는 산발적으로만 쓰였다. 우리의 적용은 모두 선택공리에 바탕을 둔다.} \emph{$\ideal M,\ideal M_1,\ideal M_2,\ldots,\ideal M_\nu,\ldots$가 $\Sigma$의 가산무한 아이디얼계이고 각 아이디얼이 바로 다음 아이디얼로 나누어진다면, 어떤 유한한 지수 $n$부터 모든 아이디얼은 동일하다: $\ideal M_n=\ideal M_{n+1}=\cdots$. 달리 말하면 $\ideal M,\ideal M_1,\ideal M_2,\ldots,\ideal M_\nu,\ldots$가 각 아이디얼이 바로 앞 아이디얼의 진약수인 단순순서 아이디얼사슬을 이루면, 이 사슬은 유한 단계에서 끝난다.}

실제로 $\ideal D=(\ideal M_1,\ideal M_2,\ldots,\ideal M_\nu,\ldots)$를 이 계의 최대공약아이디얼이라 하고, 유한조건에 의해 언제나 존재하는 $\ideal D$의 기저를 $f_1\ldots f_k$라 하자. 나눗셈 가정에 따라 $\ideal D$의 각 원소는 동시에 사슬의 한 아이디얼에 속한다. 실제로
\[
 f=g+h,\qquad g\equiv0(\ideal M_r),\qquad h\equiv0(\ideal M_s),\qquad (r\le s)
\]
이면 $g\equiv0(\ideal M_s)$이고, 따라서 $f\equiv0(\ideal M_s)$이다. $f$가 더 많은 성분의 합일 때도 마찬가지이다. 그러므로 어떤 유한한 지수 $n$이 존재하여
\[
 f_1\equiv0(\ideal M_n);\quad\ldots;\quad f_k\equiv0(\ideal M_n);\quad
 \ideal D=(f_1\ldots f_k)\equiv0(\ideal M_n)
\]
이 된다. 역으로 $\ideal M_n\equiv0(\ideal D)$이므로 $\ideal M_n=\ideal D$이다. 더 나아가
\[
 \ideal M_{n+i}\equiv0(\ideal D);\qquad
 \ideal D=\ideal M_n\equiv0(\ideal M_{n+i})
\]
이므로 모든 $i$에 대하여 $\ideal M_{n+i}=\ideal D=\ideal M_n$이다. 이로써 정리가 증명되었다.

역으로 이 정리에서 다시 아이디얼기저의 존재가 따라온다는 점도 주목하자. 따라서 유한조건을 이렇게 기저를 쓰지 않는 꼴로 진술할 수도 있었다.

\begin{center}
\S~2.\\[0.25em]
\textbf{한 아이디얼을 유한개의 기약아이디얼의 최소공배수로 나타내기.}
\end{center}

아이디얼 $\ideal B_1,\ideal B_2,\ldots,\ideal B_k$의 최소공배수 $[\ideal B_1,\ideal B_2,\ldots,\ideal B_k]$를 통상처럼 $\ideal B_1$로도, $\ideal B_2$로도, \ldots, $\ideal B_k$로도 나누어지는 모든 원소의 전체로 정의하자. 기호로 쓰면
\[
 \text{$f\equiv0(\ideal B_i)$,}\quad (i=1,2,\ldots,k)
 \quad\text{이면}\quad
 f\equiv0([\ideal B_1,\ideal B_2,\ldots,\ideal B_k])
\]
이고, 그 역도 성립한다. 최소공배수는 다시 아이디얼이다. 아이디얼 $\ideal B_i$들을 이 분해의 성분이라고도 한다.

\textbf{정의 I.} \emph{표현 $\ideal M=[\ideal B_1\cdots\ideal B_k]$에서 어느 $\ideal B_i$도 나머지 아이디얼들의 최소공배수 $\ideal A_i$의 약수가 아니고, 어느 $\ideal B_i$도 진약수로 바꿀 수 없으면 이 표현을 축약표현이라 한다.\footnote{축약되지 않은 표현의 한 예는 모든 지수 $\lambda\ge2$에 대한 $(x^2,xy)=[(x),(x^2,xy,y^\lambda)]$이다. $\lambda=1$에 대응하는 표현 $[(x),(x^2,y)]$은 그에 대응하는 축약표현이다. (전쟁에서 전사한 K. 헨첼트가 으뜸아이디얼로의 분해가 유일하지 않음을 보여 주는 가장 간단한 예로 이 표현을 알려 주었다.)} 두 조건이 아이디얼 $\ideal B_i$에 대해서만 만족되면 이 표현을 $\ideal B_i$에 관하여 축약되었다고 한다. 최소공배수 $\ideal A_i=[\ideal B_1,\ldots,\ideal B_{i-1},\ideal B_{i+1},\ldots,\ideal B_k]$를 $\ideal B_i$의 보충부분이라 한다. 첫째 조건만을 만족하는 표현을 최단표현이라 한다.}

이제 한 아이디얼을 최소공배수로 나타낼 때에는 다음 정리에 따라 축약표현만 고찰하면 충분하다.

\textbf{보조정리 I.} \emph{한 아이디얼을 유한개의 아이디얼의 최소공배수로 나타낸 모든 표현은 적어도 한 가지 방식으로 축약표현으로 바꿀 수 있다. 특히 연속적인 분해로 그러한 표현을 얻을 수 있다.}

실제로 $\ideal M=[\ideal B_1^*\cdots\ideal B_l^*]$가 $\ideal M$의 임의의 표현이라고 하자. 남겨 둔 아이디얼들의 최소공배수로 나누어지는 $\ideal B_i^*$들을 차례로 버린다. 나머지 아이디얼들도 여전히 $\ideal M$을 주므로, 이렇게 생긴 표현
\[
 \ideal M=[\ideal B_1\cdots\ideal B_k]=[\ideal A_i,\ideal B_i]
\]
에서는 첫째 조건이 만족되고, 따라서 이 표현은 최단표현이다. 어느 $\ideal B_i$를 진약수로 바꾸어도 이 조건은 유지된다. 둘째 조건은 유한사슬 정리(정리 I)에 따라 언제나 만족시킬 수 있다. 실제로
\[
 \ideal M=[\ideal A_i,\ideal B_i]=[\ideal A_i,\ideal B_i^{(1)}]
 =\cdots=[\ideal A_i,\ideal B_i^{(\nu)}],\ldots
\]
라고 하자. 여기서 각 $\ideal B_i^{(\nu)}$는 바로 앞 아이디얼의 진약수이다. 그러면 이 정리에 따라 사슬 $\ideal B_i,\ideal B_i^{(1)},\ldots,\ideal B_i^{(\nu)},\ldots$는 유한 단계에서 끝나야 한다. 따라서 표현 $\ideal M=[\ideal A_i,\ideal B_i^{(n)}]$에서는 $\ideal B_i^{(n)}$을 그 진약수로 바꿀 수 없다. $\ideal A_i$를 진약수로 바꾸면 이 사실은 더욱 분명히 유지된다. 그러므로 이미 축약된 $\ideal B$들을 포함하여 매번 보충부분을 만들면서 이 절차를 각 $\ideal B_i$에 차례로 적용하면 축약표현을 얻는다.\footnote{앞의 예는 주어진 표현이 그러한 축약표현을 유일하게 정하지 않음을 보여 준다. $\lambda\ge2$일 때 $(x^2,xy)=[(x),(x^2,xy,y^\lambda)]$에는 $[(x),(x^2,y)]$뿐 아니라 임의의 $\mu$에 대한 $[(x),(x^2,\mu x+y)]$도 대응하는 축약표현이다.}

그러한 표현을 연속적으로 얻으려면, 각각 축약된 표현들
\[
 \ideal M=[\ideal B_1,\ideal C_1];\quad
 \ideal C_1=[\ideal B_2,\ideal C_2],\quad\ldots,\quad
 \ideal C_{l-1}=[\ideal B_l,\ideal C_l]
\]
에서 생기는 표현 $\ideal M=[\ideal B_1\cdots\ideal B_l\ideal C_l]$도 축약되었음을 보이면 된다. 이를 위해서는 축약표현
\[
 \ideal M=[\ideal B_1\cdots\ideal B_\varrho\ideal C];\qquad
 \ideal C=[\ideal C_1,\ideal C_2]
\]
에서 $\ideal M=[\ideal B_1\cdots\ideal B_\varrho\ideal C_1\ideal C_2]$도 축약되었음을 보이면 충분하다. 실제로 가정에 따라 어느 $\ideal B$도 그 보충부분의 약수가 아니다. 어느 $\ideal C_i$가 그 보충부분의 약수라면, 둘째 표현이 축약되어 $\ideal C_1,\ideal C_2$가 $\ideal C$의 진약수라는 사실 때문에 첫째 표현에서 $\ideal C$를 진약수로 바꿀 수 있게 되어 가정에 어긋난다. 따라서 이 표현은 최단이다. 더 나아가 가정에 따라 어느 $\ideal B$도 진약수로 바꿀 수 없다. 어느 $\ideal C_i$를 진약수로 바꿀 수 있다면, $\ideal C$의 표현이 축약되어 있으므로 이는 $\ideal C$를 진약수로 바꾸는 것에 대응하여 다시 가정에 어긋난다. 이로써 보조정리가 증명되었다.

\textbf{정의 II.} \emph{한 아이디얼 $\ideal M$을 두 진약수의 최소공배수로 나타낼 수 있으면 $\ideal M$을 가약이라 하고, 그렇지 않으면 $\ideal M$을 기약이라 한다.}

이제 축약표현을 이용하여 유한사슬 정리인 정리 I로 다음을 증명한다.

\textbf{정리 II.} \emph{각 아이디얼은 유한개의 기약아이디얼의 최소공배수로 나타낼 수 있다.}\footnote{그러한 표현이 일반적으로 유일하지 않다는 것은 임의의 $\mu$에 대한 앞의 예 $(x^2,xy)=[(x),(x^2,\mu x+y)]$가 보여 준다. 두 성분은 임의의 $\mu$에 대하여 기약이다. 실제로 $(x)$의 모든 약수는 $(x,g(y))$ 꼴이며, 여기서 $g(y)$는 $y$의 다항식이다. 따라서 그러한 두 약수의 최소공배수도 이 꼴이고 $(x)$의 진약수가 된다. $(x^2,\mu x+y)$는 유일한 약수 $(x,y)$만을 가지므로 반드시 역시 기약이다.}

실제로 임의의 아이디얼 $\ideal M$은 기약이거나 가약이다. 기약이면 $\ideal M=[\ideal M]$가 정리 II가 요구하는 표현이다. 가약이면 $\ideal M=[\ideal B_1,\ideal C_1]$이고, 여기서 $\ideal B_1,\ideal C_1$은 $\ideal M$의 진약수이며 보조정리 I에 따라 이 표현을 축약되었다고 가정할 수 있다. $\ideal C_1$에도 같은 양자택일이 성립한다. 곧 $\ideal C_1$이 기약이거나, 축약표현 $\ideal C_1=[\ideal B_2,\ideal C_2]$가 존재한다.

계속하면 축약표현들의 열
\srcnumdisplay[3.1em]{(1)}{\begin{aligned}
 \ideal M&=[\ideal B_1,\ideal C_1];\qquad
 \ideal C_1=[\ideal B_2,\ideal C_2];\quad\ldots;\\
 \ideal C_{\nu-1}&=[\ideal B_\nu,\ideal C_\nu];\quad\ldots
\end{aligned}}
을 얻는다. 사슬 $\ideal C_1,\ideal C_2,\ldots,\ideal C_\nu,\ldots$에서 각 $\ideal C_i$은 바로 앞 아이디얼의 진약수이므로 사슬은 유한 단계에서 끝난다. 따라서 $\ideal C_n$이 기약인 지수 $n$이 존재한다. 보조정리 I에 따라 표현 $\ideal M=[\ideal B_1\cdots\ideal B_n\ideal C_n]$도 축약되어 있다. 따라서 $\ideal C_n$은 보충부분 $\ideal A_n$의 약수가 될 수 없고, 표현 $\ideal M=[\ideal A_n,\ideal C_n]$에서 $\ideal C_n$을 진약수로 바꿀 수 없다. 필요하면 $\ideal A_n$을 진약수로 바꾸어\footnote{실제로 이 표현은 $\ideal A_n$에 관해서도 축약되어 있다. 이는 \S~3의 보조정리 IV에서 보조정리 I의 역으로 증명한다.} 전체 표현을 축약하면, 모든 가약아이디얼이 한 기약아이디얼과 그 보충아이디얼의 최소공배수로 된 축약표현을 가짐을 알 수 있다. 그러므로 (1)에서는 일반성을 잃지 않고 모든 $\ideal B_i$를 기약이라고 가정할 수 있다. 위의 논증을 되풀이하면 기약인 $\ideal C_n$의 존재를 얻고, 이로써 정리 II가 증명된다.

\begin{center}
\S~3.\\[0.25em]
\textbf{서로 다른 두 기약아이디얼 분해에서 성분 수의 동일성.}
\end{center}

성분 수의 동일성을 증명하려면 먼저 한 아이디얼의 가약성 또는 기약성을 그 보충부분의 성질로 나타내야 한다. 이는 다음 정리가 해 준다.

\textbf{정리 III.}\footnote{정리 III은 슈마이들러와 뇌터--슈마이들러의 논문들에서 모듈로부터 잉여군으로 넘어가는 과정에 대응한다(서론 참조). $\ideal A$는 잉여군에, $\ideal N_1$과 $\ideal N_2$는 그 잉여군이 분해되는 부분군들에 대응한다.} \emph{최단표현 $\ideal M=[\ideal A,\ideal C]$가 $\ideal C$에 관하여 축약되어 있다고 하자. 이때 $\ideal C$가 가약이기 위한 필요충분조건은 $\ideal M$의 진약수인 두 아이디얼 $\ideal N_1,\ideal N_2$가 존재하여 다음을 만족하는 것이다.}
\srcnumdisplay[3.1em]{(2)}{\ideal N_1\equiv0(\ideal A);\qquad
 \ideal N_2\equiv0(\ideal A);\qquad
 [\ideal N_1,\ideal N_2]=\ideal M.}
\emph{또한 이로부터 다음이 따른다. 조건 (2)가 성립하고 $\ideal C$가 기약이면, $\ideal N_i$ 중 적어도 하나는 $\ideal M$의 진약수가 아니며 $\ideal N_i=\ideal M$이다.}

$\ideal C=[\ideal C_1,\ideal C_2]$이고 $\ideal C_1,\ideal C_2$가 $\ideal C$의 진약수라고 하자. 그러면
\[
 \ideal M=[\ideal A,\ideal C]=[\ideal A,\ideal C_1,\ideal C_2]
       =[[\ideal A,\ideal C_1],[\ideal A,\ideal C_2]].
\]
여기서 아이디얼 $[\ideal A,\ideal C_i]$는 $\ideal M$의 진약수이다. 그렇지 않으면 $[\ideal A,\ideal C]$가 $\ideal C$에 관하여 축약되어 있지 않기 때문이다. 또한 이 아이디얼들은 $\ideal A$로 나누어지므로 조건 (2)의 필요성이 증명되었다. (표현 (2)는 축약되어 있지 않다. $[\ideal A,\ideal C_i]$를 $\ideal C_i$로 바꿀 수 있기 때문이다.)

이제 거꾸로 (2)가 성립한다고 하자. 아이디얼들을
\[
 \ideal C_1=(\ideal C,\ideal N_1),\qquad
 \ideal C_2=(\ideal C,\ideal N_2),\qquad
 \ideal C^*=[\ideal C_1,\ideal C_2]
\]
로 만든다. 그러면 $\ideal C$는 $\ideal C_1$과 $\ideal C_2$ 각각으로 나누어지고, 따라서 그 최소공배수 $\ideal C^*$로도 나누어진다. $\ideal C^*$가 $\ideal C$로 나누어짐을 보이기 위하여
\[
 f\equiv0(\ideal C^*);\quad f\equiv0(\ideal C_1);\quad f\equiv0(\ideal C_2),
 \quad\text{또는}\quad f=c+n_1;\quad f=\bar c+n_2
\]
라고 하자. 여기서 $c,\bar c,n_1,n_2$는 각각 $\ideal C,\ideal N_1,\ideal N_2$에 속하는 원소들이며, 특히 $n_i$는 $\ideal N_i$로 나누어진다. 따라서 차
\[
 g=c-\bar c=n_2-n_1
\]
는 $\ideal C$와 $\ideal A$ 양쪽으로 나누어지고, 따라서 $\ideal M$으로 나누어진다. 더 나아가 $n_1=n_2+m$이므로 $n_1$은(마찬가지로 $n_2$도) $\ideal N_1$과 $\ideal N_2$ 양쪽으로 나누어지고, 따라서 $\ideal M$으로 나누어진다. 그러므로
\[
 f=c+m;\qquad f\equiv0(\ideal C);\qquad \ideal C^*=\ideal C
\]
를 얻는다. 이때 $\ideal C_1$과 $\ideal C_2$는 $\ideal C$의 진약수이다. 실제로 $\ideal C_i=(\ideal C,\ideal N_i)=\ideal C$라면 $\ideal N_i$이 $\ideal C$로 나누어질 것이고, $\ideal A$로도 나누어지므로 가정에 어긋나게 $\ideal N_1=\ideal M$이 될 것이다. 따라서 $\ideal C=[\ideal C_1,\ideal C_2]$로부터 $\ideal C$가 가약임을 알 수 있으며, 정리 III이 증명되었다.

거의 같은 논증이 다음 결과도 준다는 점을 덧붙여 둔다.

\textbf{보조정리 II.} \emph{최단표현 $\ideal M=[\ideal A,\ideal C]$에서 아이디얼 $\ideal C$를 진약수로 바꿀 수 있으면 $\ideal C$는 가약이다.}

실제로
\[
 \ideal M=[\ideal A,\ideal C]=[\ideal A,\ideal C_1]
\]
이라고 하고
\[
 \ideal C^*=[\ideal C_1,(\ideal A,\ideal C)]
\]
로 놓자. 그러면 다시 $\ideal C$가 $\ideal C^*$로 나누어진다. 또한 $f\equiv0(\ideal C^*)$로부터
\[
 f=c_1=a+c
\]
를 얻는다. 따라서 차 $a=c_1-c$는 $\ideal A$와 $\ideal C_1$ 양쪽으로 나누어지고, 따라서 $\ideal M$으로 나누어진다. 그러므로 $c_1=c+m$, $f\equiv0(\ideal C)$, $\ideal C^*=\ideal C$이다. 가정에 따라 $\ideal C_1$과 $(\ideal A,\ideal C)$가 모두 $\ideal C$의 진약수이므로, $\ideal C=\ideal C^*$로부터 $\ideal C$가 가약임을 알 수 있다.

따라서 기약아이디얼 $\ideal C$는 진약수로 바꿀 수 없다.

이제 $\ideal M$을 유한개의 기약아이디얼의 최소공배수로 나타내는 서로 다른 두 최단표현
\[
 \ideal M=[\ideal B_1\cdots\ideal B_k]=[\ideal D_1\cdots\ideal D_l]
\]
이 주어졌다고 하자. 보조정리 II 뒤의 주석에 따라 이 표현들은 동시에 축약되어 있다. 먼저 다음을 증명한다.

\textbf{보조정리 III.} \emph{각 보충부분 $\ideal A_i=[\ideal B_1\cdots\ideal B_{i-1}\ideal B_{i+1}\cdots\ideal B_k]$에 대하여 $\ideal M=[\ideal A_i,\ideal D_j]$가 되게 하는 아이디얼 $\ideal D_j$가 존재한다.}

실제로 $\ideal M=[\ideal D_1,\ideal C_1]$, $\ideal C_1=[\ideal D_2,\ideal C_{12}]$ 등으로 놓으면
\[
 \ideal M=[\ideal A_i,\ideal M]=[\ideal A_i,\ideal D_1,\ideal C_1]
 =[[\ideal A_i\ideal D_1],[\ideal A_i\ideal C_1]]
\]
을 얻는다. 여기서 $\ideal N_1=[\ideal A_i,\ideal D_1]$, $\ideal N_2=[\ideal A_i,\ideal C_1]$에 대하여 정리 III의 조건 (2)가 성립한다. 실제로 $\ideal M=[\ideal A_i,\ideal B_i]$는 $\ideal B_i$에 관하여 축약되어 있고 최단표현이다. 그런데 $\ideal B_i$가 기약이라고 가정했으므로 반드시 $\ideal N_i$ 중 하나가 $\ideal M$과 같아야 한다.

$\ideal N_1=\ideal M$이면 보조정리가 증명된다. $\ideal N_2=\ideal M$이면 그에 따라 $\ideal M=[[\ideal A_i\ideal D_2],[\ideal A_i\ideal C_{12}]]$가 되고, 같은 논증에 의하여 여기서도 한 성분이 $\ideal M$과 같아야 한다. 이 과정을 계속하면 $j<l$인 어떤 $j$에 대하여 $\ideal M=[\ideal A_i,\ideal D_j]$가 되거나, 아니면 $\ideal M=[\ideal A_i,\ideal C_{1\ldots l-1}]$가 된다. $\ideal C_{1\ldots l-1}=\ideal D_l$이므로 후자의 경우에도 보조정리가 증명된다.

이제 이로부터 다음을 얻는다.

\textbf{정리 IV.} \emph{한 아이디얼을 기약아이디얼들의 최소공배수로 나타내는 서로 다른 두 최단표현에서 성분의 수는 같다.}

실제로 $i=1$에 대하여 보조정리로부터
\[
 \ideal M=[\ideal A_1,\ideal B_1]=[\ideal A_1,\ideal D_{j_1}]
 = [\ideal D_{j_1},\ideal B_2,\ldots,\ideal B_k]
\]
를 얻는다. 이제 두 분해
\[
 \ideal M=[\ideal D_{j_1},\ideal B_2,\ldots,\ideal B_k]
          =[\ideal D_1,\ideal D_2,\ldots,\ideal D_l]
\]
를 생각하고, $\ideal B_2$의 보충부분 $\bar{\ideal A}_2=[\ideal D_{j_1},\ideal B_3,\ldots,\ideal B_k]$에 관하여 위의 논증을 반복하면
\[
 \ideal M=[\bar{\ideal A}_2,\ideal B_2]=[\bar{\ideal A}_2,\ideal D_{j_2}]
           =[\ideal D_{j_1},\ideal D_{j_2}\ideal B_3,\ldots,\ideal B_k]
\]
를 얻으며, 이 절차를 계속하면
\[
 \ideal M=[\ideal D_{j_1},\ideal D_{j_2},\ldots,\ideal D_{j_k}]
\]
가 된다. 가정에 따라 표현 $\ideal M=[\ideal D_1\cdots\ideal D_l]$은 최단이므로 어느 $\ideal D$도 생략할 수 없다. 따라서 $\ideal D_{j_i}$들 가운데 서로 다른 것들이 모든 $\ideal D$를 빠짐없이 포함해야 하며, 이에 따라 $k\ge l$이다. 보조정리와 이어지는 모든 논증에서 $\ideal B$와 $\ideal D$를 서로 바꾸면 마찬가지로 $l\ge k$를 얻는다. 그러므로 $k=l$이고, 이로써 성분 수의 동일성이 증명되었다. 또한 이로부터 아이디얼 $\ideal D_{j_i}$들이 모두 서로 다름을 알 수 있다. 그렇지 않으면 $\ideal D_{j_i}$들로 된 최단표현에 $k$개보다 적은 성분이 나타날 것이기 때문이다. 따라서 지표를 $\ideal D_{j_i}=\ideal D_i$가 되도록 정할 수 있다. 같은 이유로 모든 중간표현 $\ideal M=[\ideal D_{j_1}\cdots\ideal D_{j_i},\ideal B_{i+1}\cdots\ideal B_k]$도 최단이고, 따라서 보조정리 II 뒤의 주석에 따라 축약되어 있다.

성분 수의 동일성은 다음 보조정리를 통하여 보조정리 I의 역을 준다.

\textbf{보조정리 IV.} \emph{축약표현의 성분들을 여러 묶음으로 모으고 각 묶음의 최소공배수를 취하면, 그 결과로 생기는 표현도 축약되어 있다. 다시 말해 축약표현}
\[
 \ideal M=[\ideal C_{11}\cdots\ideal C_{1\mu_1};\;\ldots;\;
 \ideal C_{\sigma1}\cdots\ideal C_{\sigma\mu_\sigma}]
\]
\emph{에서 $\ideal N_i=[\ideal C_{i1}\cdots\ideal C_{i\mu_i}]$로 놓으면 $\ideal M=[\ideal N_1\cdots\ideal N_\sigma]=[\ideal N_i,\ideal L_i]$도 축약되어 있다.}

먼저 $\ideal N_i$는 그 보충부분 $\ideal L_i$의 약수가 될 수 없다. 실제로 그 약수 $\ideal C_{ij}$ 가운데 어느 것도 $\ideal L_i$의 약수가 아니기 때문이다. 따라서 이 표현은 최단이다. $\ideal N_i$를 진약수로 바꿀 수 없음을 보이기 위하여 각 $\ideal C$를 그 기약아이디얼 $\ideal B$들로 분해한다.\footnote{여기서는 언제나 $\ideal B$들로 된 최단표현, 따라서 축약표현을 뜻한다.} 그러면 보조정리 I에 따라 축약된 표현들
\[
 \ideal M=[\ideal B_{11}\cdots\ideal B_{1\lambda_1};\;\ldots;\;
 \ideal B_{\sigma1}\cdots\ideal B_{\sigma\lambda_\sigma}];
 \qquad
 \ideal N_i=[\ideal B_{i1}\cdots\ideal B_{i\lambda_i}]
\]
을 얻는다. 이제 $\ideal M=[\ideal N_i^*,\ideal L_i]$가 $\ideal N_i^*$에 관하여 축약되어 있고, $\ideal N_i^*$가 $\ideal N_i$의 진약수라고 하자. 보조정리 II에 따라
\[
 \ideal N_i=[\ideal N_i^*,(\ideal N_i,\ideal L_i)]
\]
이다. 이 표현은 반드시 $\ideal N_i^*$에 관하여 축약되어 있다. 그렇지 않으면 $\ideal M$에서도 $\ideal N_i^*$를 진약수로 바꿀 수 있기 때문이다. 이제 필요하면 $(\ideal N_i,\ideal L_i)$도 진약수로 바꾸어 $\ideal N_i$의 축약표현을 얻는다. $\ideal N_i$의 이 두 성분을 기약아이디얼들로 분해하면, $\ideal N_i$의 기약아이디얼 수 $\lambda_i$는 두 성분의 기약아이디얼 수의 합이 된다. 따라서 진약수 $\ideal N_i^*$에 대응하는 기약아이디얼의 수는 반드시 $\lambda_i$보다 작다. 그러면 $\ideal M=[\ideal N_i^*,\ideal L_i]$를 기약아이디얼들로 분해할 때에도 아이디얼의 수가 $\sum_i\lambda_i$보다 작아지므로, 성분 수의 동일성에 모순된다. 특수한 경우 $\sigma=2$로부터 표현 $\ideal M=[\ideal N_i,\ideal L_i]$가 보충부분 $\ideal L_i$에 관해서도 축약되어 있음이 또한 따른다.

\begin{center}
\S~4.\\[0.25em]
\textbf{으뜸아이디얼. 서로 다른 두 기약아이디얼 분해에서 대응 소아이디얼의 유일성.}
\end{center}

이하에서는 으뜸아이디얼과 기약아이디얼 사이의 관계를 다룬다.

\textbf{정의 III.} \emph{$a\cdot b\equiv0(\ideal Q)$, $a\not\equiv0(\ideal Q)$로부터 반드시 $b^\varkappa\equiv0(\ideal Q)$가 따르고 지수 $\varkappa$가 유한한 수이면, 아이디얼 $\ideal Q$를 으뜸아이디얼이라 한다.}

이 정의는 다음과 같이 말할 수도 있다. 곱 $a\cdot b$가 $\ideal Q$로 나누어지면 한 인자가 $\ideal Q$로 나누어지거나, 각 인자의 어떤 거듭제곱이 $\ideal Q$로 나누어진다. \emph{특히 $\varkappa$를 언제나 1로 취할 수 있으면 그 아이디얼을 소아이디얼이라 한다.}

으뜸아이디얼(또는 소아이디얼)의 정의와 기저의 존재로부터, 아이디얼의 곱만을 사용하는 다음 정의를 얻는다.\footnote{두 아이디얼의 곱 $\ideal A\cdot\ideal B$는 통상적으로 원소 $a\cdot b$ 전체와 그 유한합들로 이루어진 아이디얼을 뜻한다.}

\textbf{정의 IIIa.} \emph{$\ideal A\cdot\ideal B\equiv0(\ideal Q)$, $\ideal A\not\equiv0(\ideal Q)$로부터 반드시 $\ideal B^\lambda\equiv0(\ideal Q)$가 따르면 아이디얼 $\ideal Q$를 으뜸아이디얼이라 한다. $\lambda$를 언제나 1로 취할 수 있으면 그 아이디얼을 소아이디얼이라 한다. 따라서 소아이디얼 $\ideal P$에 대해서는 $\ideal A\cdot\ideal B\equiv0(\ideal P)$, $\ideal A\not\equiv0(\ideal P)$로부터 언제나 $\ideal B\equiv0(\ideal P)$가 따른다.}

IIIa에서 $\ideal A=(a)$, $\ideal B=(b)$로 놓으면 정의 III이 특수한 경우로 포함되므로, IIIa에 따라 으뜸인 모든 아이디얼은 III에 따라서도 으뜸이다. 거꾸로 $\ideal Q$가 III에 따라 으뜸이고 IIIa의 가정 $\ideal A\cdot\ideal B\equiv0(\ideal Q)$가 성립한다고 하자. 그러면 $\ideal A\equiv0(\ideal Q)$이거나, 아니면 $a\cdot\ideal B\equiv0(\ideal Q)$이면서 $a\not\equiv0(\ideal Q)$인 원소 $a\equiv0(\ideal A)$가 적어도 하나 존재한다. 이제 $b_1\ldots b_r$가 $\ideal B$의 아이디얼기저이면 $a\cdot b_i\equiv0(\ideal Q)$이므로 정의 III에 따라
\[
 b_1^{\varkappa_1}\equiv0(\ideal Q);\quad\ldots;\quad
 b_r^{\varkappa_r}\equiv0(\ideal Q)
\]
를 얻는다. 또한
\[
 b=f_1b_1+\cdots+f_r b_r+n_1b_1+\cdots+n_r b_r
\]
이므로 $\lambda=\varkappa_1+\cdots+\varkappa_r$로 놓으면 $\ideal B$의 원소 $\lambda$개로 이루어진 임의의 곱이 $\ideal Q$로 나누어진다. 이로써 III에 따라 으뜸인 아이디얼이 정의 IIIa도 만족함이 증명되었다. 특히 소아이디얼 $\ideal P$의 경우에는 $a\cdot\ideal B\equiv0(\ideal P)$, 따라서 모든 $b\equiv0(\ideal B)$에 대하여 $a\cdot b\equiv0(\ideal P)$이고 $a\not\equiv0(\ideal P)$이면 $b\equiv0(\ideal P)$, 따라서 $\ideal B\equiv0(\ideal P)$가 따른다. 이로써 두 정의의 동치가 증명되었다.

으뜸아이디얼과 소아이디얼의 관계는 다음 관찰로 얻어진다. 어떤 거듭제곱이 $\ideal Q$로 나누어지는 성질을 가진 모든 원소 $p$의 전체 $\ideal P$는 소아이디얼을 이룬다. 먼저 $p_1,p_2$와 함께 $a p_1$과 $p_1-p_2$도 이 성질을 가지므로 $\ideal P$가 아이디얼임은 분명하다. 또한 정의 IIIa에서 사용한 기저 논증에 따라 어떤 수 $\lambda$가 존재하여 $\ideal P^\lambda\equiv0(\ideal Q)$가 된다.

이제
\[
 a\cdot b\equiv0(\ideal P);\qquad a\not\equiv0(\ideal P)
\]
라고 하자. $\ideal P$의 정의에 따라
\[
 a^\lambda\cdot b^\lambda\equiv0(\ideal Q);\qquad a^\lambda\not\equiv0(\ideal Q)
\]
이고, 따라서 $\ideal Q$의 정의에 따라
\[
 b^{\lambda\varkappa}\equiv0(\ideal Q)\qquad\text{따라서}\qquad b\equiv0(\ideal P)
\]
이다. 이로써 $\ideal P$가 소아이디얼임이 증명되었다. 또한 $\ideal P$는 어떤 거듭제곱이 $\ideal Q$로 나누어지는 성질을 가진 모든 아이디얼 $\ideal B$의 최대공약아이디얼로 정의할 수 있다. 실제로 그러한 각 $\ideal B$는 $\ideal P$로 나누어지므로, 이 $\ideal B$들의 최대공약아이디얼 $\ideal D$도 $\ideal P$로 나누어진다. 거꾸로 $\ideal P$ 자체가 그러한 아이디얼 $\ideal B$이므로 $\ideal D$로 나누어지고, 따라서 $\ideal P=\ideal D$이다. 그러므로 $\ideal P$는 $\ideal Q$의 약수이고 그 어떤 거듭제곱이 $\ideal Q$로 나누어지는 소아이디얼이며, 이 성질로 유일하게 정해진다. 실제로
\[
 \ideal Q\equiv0(\ideal P);\qquad \ideal P^\lambda\equiv0(\ideal Q),
 \qquad
 \ideal Q\equiv0(\bar{\ideal P});\qquad \bar{\ideal P}^{\mu}\equiv0(\ideal Q)
\]
이면
\[
 \ideal P^\lambda\equiv0(\bar{\ideal P}),\qquad
 \bar{\ideal P}^{\mu}\equiv0(\ideal P)
\]
이고, 소아이디얼의 성질에 따라
\[
 \ideal P\equiv0(\bar{\ideal P}),\qquad
 \bar{\ideal P}\equiv0(\ideal P),\qquad
 \ideal P=\bar{\ideal P}
\]
를 얻는다. 정리하면 다음과 같다.

\textbf{정리 V.} \emph{각 으뜸아이디얼 $\ideal Q$에 대하여, $\ideal Q$의 약수이면서 그 어떤 거듭제곱이 $\ideal Q$로 나누어지는 소아이디얼 $\ideal P$가 유일하게 존재한다. $\ideal P$를 $\ideal Q$에 ``대응 소아이디얼''이라 한다.}\footnote{역이 성립하지 않음은 예 $\ideal M=(x^2,xy)$가 보여 준다. 소아이디얼 $(x)$는 모든 조건을 만족하지만 $\ideal M$은 으뜸아이디얼이 아니다.} \emph{$\ideal P$는 어떤 거듭제곱이 $\ideal Q$로 나누어지는 성질을 가진 모든 아이디얼 $\ideal B$의 최대공약아이디얼로 정의된다. $\ideal P^\varrho\equiv0(\ideal Q)$가 되게 하는 가장 작은 수를 $\varrho$라 하면, $\varrho$를 $\ideal Q$의 지수라 한다.}\footnote{유리정수 또는 대수적 정수의 환영역에서와 달리 일반적으로 $\ideal P^\varrho=\ideal Q$가 되지 않음은 다음 예가 보여 준다. $\ideal Q=(x^2,y)$, $\ideal P=(x,y)$, $\ideal P^2=(x^2,xy,y^2)\equiv0(\ideal Q)$이지만 $\ideal Q\not\equiv0(\ideal P^2)$이다. 따라서 $\ideal Q$는 $\ideal P^2$와 다르다.}

이제 으뜸성과 기약성 사이의 관계로서 다음을 증명한다.

\textbf{정리 VI.} \emph{으뜸아이디얼이 아닌 모든 아이디얼은 가약이다. 다시 말해 모든 기약아이디얼은 으뜸아이디얼이다.}\footnote{역이 성립하지 않음은 예를 들어 $\lambda\ge2$일 때 $\ideal Q=(x^2,xy,y^\lambda)=[(x^2,y),(x,y^\lambda)]$가 보여 준다. 여기서 $\ideal Q$는 으뜸이지만 가약이다. $\ideal Q$가 으뜸임은 이 아이디얼이 $x,y$의 모든 $\lambda$차 멱곱을 포함한다는 데서 따른다. 따라서 상수항이 없는 모든 다항식의 어떤 거듭제곱은 $\ideal Q$로 나누어진다. 그런데 $a\cdot b\equiv0(\ideal Q)$에서 다항식 $b$가 상수항을 포함한다면, 즉 모든 $\varkappa$에 대하여 $b^\varkappa\not\equiv0(\ideal Q)$이면, $\ideal Q$의 기저다항식들이 동차이므로 $a\cdot b$의 각 동차성분이 $\ideal Q$로 나누어진다는 사실에서 $a$가 $\ideal Q$로 나누어져야 한다.}

$\ideal K$가 으뜸이 아닌 아이디얼이라고 하자. 그러면 정의 III에 따라 적어도 한 원소쌍 $a,b$가 존재하여
\srcnumdisplay[3.1em]{(3)}{a\cdot b\equiv0(\ideal K);\qquad
 a\not\equiv0(\ideal K);\qquad
 b^\varkappa\not\equiv0(\ideal K)\quad\text{모든 }\varkappa\text{에 대하여}}
를 만족한다. 이제 두 아이디얼
\[
 \ideal L_0=(\ideal K,a);\qquad \ideal N_0=(\ideal K,b)
\]
를 만든다. 이들은 (3)에 따라 $\ideal K$의 진약수이며, 또한 (3)에 따라
\srcnumdisplay[3.1em]{(4)}{\ideal L_0\cdot\ideal N_0\equiv0(\ideal K)}
이다. 최소공배수 $\ideal K_0=[\ideal L_0,\ideal N_0]$의 원소 $f$에 대해서는 다음 둘 중 하나가 성립한다.

\emph{첫째,}
\[
 f\equiv0(\ideal L_0);\qquad f\equiv0(\ideal N_0),
 \qquad\text{즉}\qquad f\equiv a_1\cdot b\;(\ideal K)
\]
로부터 언제나
\[
 f\equiv l_0\cdot b\;(\ideal K);\qquad l_0\equiv0(\ideal L_0)
\]
인 표현이 따른다. 그러면 (4)에 따라 $f\equiv0(\ideal K)$이고, 따라서 $\ideal K_0\equiv0(\ideal K)$이다. 또한 $\ideal K\equiv0(\ideal K_0)$이므로 $\ideal K=\ideal K_0$이다. 이로써 $\ideal K$가 가약임이 증명된다.

\emph{둘째,} 그러한 $l_0$가 존재하지 않는 $f\equiv0(\ideal K_0)$가 적어도 하나 있다. 이때 이 $f$에 대응하는 $a_1$을 이용하여
\[
 \ideal L_1=(\ideal L_0,a_1)=(\ideal K,a,a_1);\qquad
 \ideal N_1=(\ideal K,b^2)
\]
를 만든다. $a_1\cdot b\equiv0(\ideal L_0)$이므로 (4)에 따라
\srcnumdisplay[3.1em]{(4')}{\ideal L_1\cdot\ideal N_1\equiv0(\ideal K)}
도 성립하고, $\ideal L_1$은 $\ideal L_0$의 진약수가 된다.

$\ideal K_1=[\ideal L_1,\ideal N_1]$의 원소 $f$에 대해서도 같은 양자택일이 성립한다.

\emph{첫째,}
\[
 f\equiv0(\ideal L_1);\qquad f\equiv0(\ideal N_1),
\]
즉
\[
 f\equiv a_2\cdot b^2\;(\ideal K)
 \quad\text{이면 언제나}\quad
 f\equiv l_1\cdot b^2;\qquad l_1\equiv0(\ideal L_1)
\]
가 따른다. 그러면 (4')에 따라 $\ideal K=\ideal K_1$이다.

\emph{둘째,} 적어도 한 $f$에 대하여 그러한 $l_1$이 존재하지 않는다. 이 경우 $\ideal L_2=(\ideal L_1,a_2)$, $\ideal N_2=(\ideal K,b^{2^2})$를 만들게 되며, $\ideal L_2\cdot\ideal N_2\equiv0(\ideal K)$이고 $\ideal L_2$는 $\ideal L_1$의 진약수이다. 이 과정을 계속하여 일반적으로
\[
 \ideal L_0=(\ideal K,a);\quad
 \ideal L_1=(\ideal L_0,a_1);\quad\ldots;\quad
 \ideal L_\nu=(\ideal L_{\nu-1},a_\nu);\quad\ldots,
\]
\[
 \ideal N_0=(\ideal K,b);\quad
 \ideal N_1=(\ideal K,b^2);\quad\ldots;\quad
 \ideal N_\nu=(\ideal K,b^{2^\nu});\quad\ldots
\]
로 정의한다. 여기서 $a_i$들은 어떤 $f$가 존재하여
\[
 f\equiv0(\ideal L_{i-1});\qquad f\equiv0(\ideal N_{i-1}),
\]
즉
\[
 f\equiv a_i\cdot b^{2^{i-1}}\;(\ideal K);
 \qquad\text{그러나}\qquad a_i\not\equiv0(\ideal L_{i-1})
\]
가 되도록 정해진다. 그러면 일반적으로 $\ideal L_i\cdot\ideal N_i\equiv0(\ideal K)$이고, (3)에 따라 $\ideal N_i$는 $\ideal K$의 진약수이며 $\ideal L_i$는 $\ideal L_{i-1}$의 진약수이다. 따라서 유한사슬 정리인 정리 I에 따라 $\ideal L$들의 사슬은 이를테면 $\ideal L_n$에서 유한하게 끝나야 한다. 그러면 모든 $f\equiv0(\ideal L_n)$, $f\equiv0(\ideal N_n)$에 대하여 $l_n\equiv0(\ideal L_n)$인 $l_n$이 존재하여 $f\equiv l_n\cdot b^{2^n}\;(\ideal K)$가 된다. 따라서 위의 논증에 따라 $\ideal K=[\ideal L_n,\ideal N_n]$이며, 이로써 $\ideal K$가 가약임이 증명되었다.

방금 증명한 결과로부터 \emph{대응 소아이디얼의 유일성}을 다음과 같이 얻는다.

\[
 \ideal M=[\ideal B_1\cdots\ideal B_k]=[\ideal D_1\cdots\ideal D_k]
\]
가 $\ideal M$을 기약아이디얼들의 최소공배수로 나타내는 두 최단표현이라고 하자. 이들은 축약되어 있으며, 정리 IV에 따라 성분의 수도 같다. 또한 정리 IV에 따라 거기서 나타난 중간표현들은(그곳에서 언급했듯 지표를 $j_i=i$로 놓을 수 있다)
\[
 \ideal M=[\ideal D_1\cdots\ideal D_{i-1}\ideal B_i\ideal B_{i+1}\cdots\ideal B_k]
       =[\ideal D_1\cdots\ideal D_{i-1}\ideal D_i\ideal B_{i+1}\cdots\ideal B_k]
       =[\bar{\ideal A}_i,\ideal B_i]=[\bar{\ideal A}_i,\ideal D_i]
\]
처럼 최단표현이다. 따라서
\[
 \bar{\ideal A}_i\cdot\ideal B_i\equiv0(\ideal D_i),\qquad
 \bar{\ideal A}_i\not\equiv0(\ideal D_i),
 \qquad
 \bar{\ideal A}_i\cdot\ideal D_i\equiv0(\ideal B_i),\qquad
 \bar{\ideal A}_i\not\equiv0(\ideal B_i)
\]
이다. 이제 정리 VI에 따라 기약아이디얼 $\ideal B_i$와 $\ideal D_i$는 으뜸이므로, 두 수 $\lambda_i,\mu_i$가 존재하여
\srcnumdisplay[3.1em]{(5)}{\ideal B_i^{\lambda_i}\equiv0(\ideal D_i);\qquad
 \ideal D_i^{\mu_i}\equiv0(\ideal B_i)}
가 된다. $\ideal B_i$와 $\ideal D_i$에 대응 소아이디얼을 각각 $\ideal P_i$, $\bar{\ideal P}_i$라 하자. 곧 $\ideal P_i^{\varrho_i}\equiv0(\ideal B_i)$, $\bar{\ideal P}_i^{\sigma_i}\equiv0(\ideal D_i)$라 하자. 그러면 (5)에 따라
\[
 \ideal P_i^{\lambda_i\varrho_i}\equiv0(\bar{\ideal P}_i),
 \qquad
 \bar{\ideal P}_i^{\mu_i\sigma_i}\equiv0(\ideal P_i)
\]
이고, 소아이디얼의 성질에 따라
\[
 \ideal P_i\equiv0(\bar{\ideal P}_i),
 \qquad
 \bar{\ideal P}_i\equiv0(\ideal P_i),
 \qquad
 \ideal P_i=\bar{\ideal P}_i
\]
를 얻는다. 이로써 다음이 증명되었다.

\textbf{정리 VII.} \emph{한 아이디얼을 기약아이디얼들의 최소공배수로 나타내는 서로 다른 두 최단표현에서 대응 소아이디얼들은 서로 일치한다. 그 가운데 같은 소아이디얼이 나타날 수도 있으며, 각 분해에서 같은 횟수로 나타난다.\footnote{예를 들어 각주 19의 예 $\lambda\ge2$일 때 $(x^2,xy,y^\lambda)=[(x^2,y),(x,y^\lambda)]$가 이를 보여 준다. 두 성분에 대응 소아이디얼은 모두 $(x,y)$이다.} 따라서 두 표현의 아이디얼들은 적어도 한 가지 방식으로 쌍을 지어, 각 아이디얼 $\ideal B_i$의 어떤 거듭제곱이 짝지어진 $\ideal D_i$로 나누어지고 그 역도 성립하게 할 수 있다. 그 수는 정리 IV에 따라 같다.}\footnote{기약아이디얼들 가운데 들어 있는 ``고립'' 아이디얼의 유일성에 대해서는 \S~7을 참조하라.}

\begin{center}
\S~5.\\[0.25em]
\textbf{한 아이디얼을 극대 으뜸아이디얼들의 최소공배수로 나타내기. 대응 소아이디얼의 유일성.}
\end{center}

\textbf{정의 IV.} \emph{최단표현 $\ideal M=[\ideal Q_1\cdots\ideal Q_\alpha]$에서 모든 $\ideal Q$가 으뜸이지만 어느 두 $\ideal Q$의 최소공배수도 더 이상 으뜸이 아니면, 이 표현을 극대 으뜸아이디얼들의 최소공배수라고 한다.}

그러한 표현이 적어도 하나 언제나 존재한다는 사실은 $\ideal M$을 기약아이디얼들의 최소공배수로 나타내는 표현에서 따른다. 실제로 기약아이디얼들은 으뜸이다. 이미 극대 으뜸아이디얼들로 된 표현이거나, 아니면 어떤 두 아이디얼의 최소공배수가 다시 으뜸이다. 후자의 경우 아이디얼의 수가 하나 줄어들므로, 이 절차를 반복하면 유한 단계 뒤에 원하는 표현을 얻는다.

이 표현은 보조정리 IV에 따라 축약되어 있다. 거꾸로 극대 으뜸아이디얼들로 된 모든 축약표현은 $\ideal Q$들을 기약아이디얼들로 분해해 보면 이러한 방식으로 생긴다.

여기서 정리 VII로부터 그에 대응하는 유일성 정리를 얻으려면 대응 소아이디얼과의 관계를 다음과 같이 조사해야 한다.

\textbf{정리 VIII.} \emph{으뜸아이디얼 $\ideal N_1,\ideal N_2,\ldots,\ideal N_\lambda$가 모두 같은 대응 소아이디얼 $\ideal P$를 가지면, 그 최소공배수 $\ideal Q=[\ideal N_1,\ideal N_2,\ldots,\ideal N_\lambda]$도 으뜸이고 $\ideal P$를 대응 소아이디얼로 가진다. 거꾸로 $\ideal Q=[\ideal N_1\cdots\ideal N_\lambda]$가 으뜸아이디얼 $\ideal Q$의 축약표현이면, 모든 $\ideal N_i$는 으뜸이고 $\ideal Q$에 대응 소아이디얼 $\ideal P$를 자신의 대응 소아이디얼로 가진다.}

주장의 첫 부분을 증명하기 위해 먼저 각 $i$에 대하여 $\ideal P^{\varrho_i}\equiv0(\ideal N_i)$이면 $\ideal P^\tau\equiv0(\ideal Q)$도 성립함을 관찰하자. 여기서 $\tau$는 지수 $\varrho_i$들 가운데 가장 큰 것이다. 또한 $\ideal P$는 $\ideal Q$의 약수이므로, $\ideal Q$가 으뜸이기만 하면 $\ideal P$는 반드시 그 대응 소아이디얼이다. 이제
\[
 \ideal A\cdot\ideal B\equiv0(\ideal Q),\qquad
 \ideal B^k\not\equiv0(\ideal Q)\quad\text{(모든 }k\text{에 대하여)}
\]
이면
\[
 \ideal B\not\equiv0(\ideal P);
 \quad\text{따라서}\quad
 \ideal B^k\not\equiv0(\ideal N_i)\quad\text{(모든 }k\text{에 대하여)}
\]
이다. 그러므로 $\ideal A\equiv0(\ideal N_i)$이고, 따라서 $\ideal A\equiv0(\ideal Q)$이다. 이로써 $\ideal Q$가 으뜸이고 $\ideal P$가 그 대응 소아이디얼임이 증명되었다.

이제 거꾸로 먼저
\[
 \ideal Q=[\ideal N_1\cdots\ideal N_\lambda]=[\ideal N_i,\ideal L_i]
\]
가 으뜸아이디얼 $\ideal N_i$들로 된 $\ideal Q$의 최단표현이고, 각 $\ideal P_i$가 $\ideal N_i$에 대응 소아이디얼이라고 하자. 그러면
\[
 \ideal L_i\cdot\ideal N_i\equiv0(\ideal Q),\qquad
 \ideal L_i\not\equiv0(\ideal Q)
\]
(최단표현이므로)에서
\[
 \ideal N_i^{\sigma_i}\equiv0(\ideal Q),\qquad
 \text{따라서}\qquad
 \ideal P_i^{\varrho_i\sigma_i}\equiv0(\ideal Q)
\]
를 얻는다. 동시에 $\ideal P_i$는 $\ideal Q$의 약수이므로, 각 $i$에 대하여 $\ideal P_i$는 $\ideal Q$에 대응 소아이디얼 $\ideal P$와 일치한다.

이제 모든 축약표현 $\ideal Q=[\ideal N_1\cdots\ideal N_\lambda]$에서 $\ideal N_i$들이 으뜸임을 보이는 일만 남았다.\footnote{여기서 축약표현이라는 가정이 본질적임은 축약되지 않은 다음 표현이 보여 준다. $\lambda\ge2$일 때 $\ideal Q=[x^2,xy,y^\lambda]=[(x^2,xy,y^2,yz),(x,y^\lambda)]$이다. 여기서 $(x^2,xy,y^2,yz)=[(x^2,y),(x,y^2,z)]$는 위에서 증명한 바에 따라 으뜸이 아니다. 실제로 뒤의 표현은 으뜸아이디얼들로 된 최단표현이지만 대응 소아이디얼 $(x,y)$와 $(x,y,z)$가 서로 다르다. ($\ideal Q$가 으뜸임은 각주 19에 따른다.)} 이를 위해 $\ideal N_i$들을 그 기약아이디얼 $\ideal B$들로 분해하자. 이렇게 생긴 최단표현 $\ideal Q=[\ideal B_1\cdots\ideal B_\mu]$에서는 각 아이디얼이 으뜸이고, 방금 증명한 결과에 따라 $\ideal P$를 대응 소아이디얼로 가진다. 그러면 위에서 증명한 정리의 첫 부분에 따라 각 $\ideal N_i$에도 같은 사실이 성립한다. 이로써 정리 VIII이 완전히 증명되었다.

\textbf{따름정리.} 이로부터 소아이디얼은 반드시 기약임이 또한 따른다. 실제로 축약표현 $\ideal P=[\ideal N,\ideal L]$에서 정리 VIII에 따라 $\ideal N\equiv0(\ideal P)$이고, $\ideal P\equiv0(\ideal N)$이기도 하므로 $\ideal P=\ideal N$이며 마찬가지로 $\ideal P=\ideal L$이다. $\ideal P$의 기약성은 직접 얻을 수도 있다. 실제로 $\ideal P=[\ideal N,\ideal L]$이면 $\ideal N\cdot\ideal L\equiv0(\ideal P)$이지만 $\ideal N\not\equiv0(\ideal P)$, $\ideal L\not\equiv0(\ideal P)$여서 소아이디얼의 정의 성질에 모순된다.

이제
\[
 \ideal M=[\ideal Q_1\cdots\ideal Q_\alpha]
          =[\bar{\ideal Q}_1\cdots\bar{\ideal Q}_\beta]
\]
가 $\ideal M$을 극대 으뜸아이디얼들의 최소공배수로 나타내는 두 축약표현이라고 하자. $\ideal Q$들을 기약아이디얼 $\ideal B$들로, $\bar{\ideal Q}$들을 그에 대응하여 기약아이디얼 $\ideal D$들로 분해하면, $\ideal M$을 기약아이디얼들의 최소공배수로 나타내는 두 축약표현을 얻는다. 정리 VII에 따라 두 표현에서는 성분의 수와 대응 소아이디얼들이 모두 일치한다. 정리 VIII에 따라 고정된 $\ideal Q_i$에서 나타나는 모든 기약아이디얼 $\ideal B$는 같은 대응 소아이디얼 $\ideal P_i$를 가진다. 한편 $\ideal Q_k$에 속하는 $\ideal P_k$는 반드시 $\ideal P_i$와 다르다. 그렇지 않으면 정리 VIII에 따라 극대 으뜸아이디얼들로 된 표현이 아니기 때문이다. 따라서 $\ideal Q$들의 수 $\alpha$는 $\ideal B$들에 결부된 서로 다른 소아이디얼 $\ideal P$들의 수와 같고, 이 서로 다른 $\ideal P$들이 $\ideal Q$들에 대응 소아이디얼들이다. $\bar{\ideal Q}$들을 $\ideal D$들로 분해한 것에 대해서도 같은 사실이 성립한다. 그러므로 정리 VII로부터 $\ideal Q$와 $\bar{\ideal Q}$의 수가 같고 그 대응 소아이디얼들이 일치함을 얻는다. 동시에 이 절의 첫머리에서 설명한, 기약아이디얼들을 극대 으뜸아이디얼들로 모으는 과정은 같은 대응 소아이디얼을 가진 것들을 모두, 그리고 오직 그것들만 한데 모으는 과정임을 알 수 있다. 더 나아가 정리 VIII은 극대 으뜸아이디얼의 기약성 성질을 보여 준다. 즉 극대 으뜸아이디얼은 극대 으뜸아이디얼들의 최소공배수로 된 축약표현을 허용하지 않는다.

정리하면 다음이 증명되었다.

\textbf{정리 IX.} \emph{한 아이디얼을 극대 으뜸아이디얼들의 최소공배수로 나타내는 두 축약표현에서 성분의 수와 대응 소아이디얼들이 일치하며, 이 소아이디얼들은 모두 서로 다르다. 다시 말해 각 $\ideal Q$에 $\bar{\ideal Q}$ 하나를 일대일로 대응시켜, $\ideal Q$의 어떤 거듭제곱이 $\bar{\ideal Q}$로 나누어지고 그 역도 성립하게 할 수 있다.}\footnote{서로 다른 표현의 예는 정리 II의 각주 12에 나온다. 임의의 $\mu$에 대하여 $(x^2,xy)=[(x),(x^2,\mu x+y)]$이다. 대응 소아이디얼 $\ideal P_1=(x)$, $\ideal P_2=(x,y)$가 서로 다르므로 이들은 극대 으뜸아이디얼이다. ``고립'' 극대 으뜸아이디얼의 유일성에 대해서는 \S~7을 참조하라.} \emph{또한 $\ideal Q$들과 $\bar{\ideal Q}$들은 극대 으뜸아이디얼로의 분해에 관하여 기약성 성질을 가진다.}

\textbf{따름정리.} 정리 IX는 축약표현 대신 최단표현만을 가정해도 본질적으로 그대로 성립한다는 점을 덧붙여 둔다. 예를 들어
\[
 \ideal M=[\ideal Q_1\cdots\ideal Q_i^*\cdots\ideal Q_\alpha]
\]
가 $\ideal Q_i^*$에 관하여 축약되어 있고, $\ideal Q_i^*$가 $\ideal Q_i$의 진약수이며, $\ideal L_i$가 $\ideal Q_i$의 보충부분이라고 하자. 그러면 보조정리 IV에 따라
\[
 \ideal Q_i=[\ideal Q_i^*;(\ideal L_i,\ideal Q_i)]
\]
이고, 이 표현은 $\ideal Q_i^*$에 관하여 축약되어 있다. 정리 VIII을 적용할 때 필요하면 $(\ideal L_i,\ideal Q_i)$를 진약수로 바꾸어도 되므로, $\ideal Q_i^*$는 으뜸이고 $\ideal Q_i$와 같은 대응 소아이디얼 $\ideal P_i$를 가진다. 이 절차를 계속하면 그러한 각 표현에 극대 으뜸아이디얼들로 된 축약표현을 대응시킬 수 있으며, 성분의 수와 대응 소아이디얼들은 일치한다. 따라서 최단표현의 경우에도 서로 다른 두 표현에서 성분의 수와 대응 소아이디얼들이 일치한다.

이렇게 유일하게 정의되는 서로 다른 대응 소아이디얼들을 줄여서 ``$\ideal M$에 대응 소아이디얼들''이라 부르기로 한다.

\begin{center}
\S~6.\\[0.25em]
\textbf{한 아이디얼을 상대소-기약아이디얼들의 최소공배수로 유일하게 나타내기.}
\end{center}

\textbf{정의 V.} \emph{$\ideal T\cdot\ideal R\equiv0(\ideal S)$로부터 반드시 $\ideal T\equiv0(\ideal S)$가 따르면 아이디얼 $\ideal R$을 $\ideal S$에 대하여 상대소라고 한다. $\ideal R$이 $\ideal S$에 대하여 상대소이고 $\ideal S$도 $\ideal R$에 대하여 상대소이면, $\ideal R$과 $\ideal S$를 서로 상대소라고 한다.}\footnote{상대소 조건은 대칭적이지 않다. 예를 들어 $\ideal R=(x^2,y)$는 $\ideal S=(x)$에 대하여 상대소이지만, $\ideal S$는 $\ideal R$에 대하여 상대소가 아니다. 실제로 $\ideal S^2\equiv0(\ideal R)$이지만 $\ideal T=\ideal S\not\equiv0(\ideal R)$이다.}

\emph{한 아이디얼을 서로 상대소인 두 진약수의 최소공배수로 나타낼 수 없으면 그 아이디얼을 상대소-기약이라 한다.}

특히 $\ideal T\cdot\ideal R\equiv0(\ideal S)$인 모든 $\ideal T$의 최대공약아이디얼 $\ideal T_0$를 취하면, $\ideal T_0\cdot\ideal R\equiv0(\ideal S)$이고 $\ideal S\equiv0(\ideal T_0)$이다. 따라서 $\ideal R$이 $\ideal S$에 대하여 상대소이면 $\ideal T_0=\ideal S$이고, 상대소가 아니면 $\ideal T_0$는 $\ideal S$의 진약수이다.\footnote{이렇게 정의한 $\ideal T_0$는 라스커의 ``잔여모듈''과 일치하며, 아이디얼 대신 수모듈로 확장하면 데데킨트의 두 모듈의 ``몫''과 일치한다. Lasker, 앞의 논문, 49쪽; Dedekind, \emph{Zahlentheorie}, 504쪽.}

유일성 증명의 바탕은 다음 정리이다.

\textbf{정리 X.} \emph{1. $\ideal R$이 아이디얼 $\ideal S_1\ldots\ideal S_\lambda$ 각각에 대하여 상대소이면, $\ideal R$은 그 최소공배수 $\ideal S$에 대해서도 상대소이다.}

\emph{2. 아이디얼 $\ideal S_1,\ldots,\ideal S_\lambda$가 각각 $\ideal R$에 대하여 상대소이면, 그 최소공배수 $\ideal S$도 $\ideal R$에 대하여 상대소이다.}

\emph{3. $\ideal R$이 $\ideal S$에 대하여 상대소이고 $\ideal S=[\ideal S_1\cdots\ideal S_\lambda]$가 $\ideal S$의 축약표현이면, $\ideal R$은 각 $\ideal S_i$에 대해서도 상대소이다.}

\emph{4. $\ideal S$가 $\ideal R$에 대하여 상대소이면, $\ideal S$의 모든 약수 $\ideal S_i$도 $\ideal R$에 대하여 상대소이다.}

1. $\ideal T\cdot\ideal R\equiv0(\ideal S)$로부터 반드시 $\ideal T\cdot\ideal R\equiv0(\ideal S_i)$가 따르므로, 가정에 따라 $\ideal T\equiv0(\ideal S_i)$이고 따라서 $\ideal T\equiv0(\ideal S)$이다.

2. 다음과 같이 놓자.
\[
 \ideal C_1=[\ideal S_2\cdots\ideal S_\lambda],\quad
 \ideal C_{12}=[\ideal S_3\cdots\ideal S_\lambda];\ \ldots
 \ideal C_{12\ldots\lambda-1}=\ideal S_\lambda.
\]
$\ideal T\cdot\ideal S\equiv0(\ideal R)$로부터 $\ideal T\cdot\ideal C_1\cdot\ideal S_1\equiv0(\ideal R)$가 따르고, 따라서 가정에 따라 $\ideal T\cdot\ideal C_1\equiv0(\ideal R)$이다. 여기서 다시 $\ideal T\cdot\ideal C_{12}\cdot\ideal S_2\equiv0(\ideal R)$, 따라서 $\ideal T\cdot\ideal C_{12}\equiv0(\ideal R)$를 얻는다. 마침내 $\ideal T\cdot\ideal C_{12\ldots\lambda-1}=\ideal T\cdot\ideal S_\lambda\equiv0(\ideal R)$이고, 따라서 $\ideal T\equiv0(\ideal R)$이다.

3. $\ideal C_i$를 $\ideal S_i$의 보충부분이라 하자. $\ideal T_0\cdot\ideal R\equiv0(\ideal S_i)$이면 $\ideal C_i\cdot\ideal R\equiv0(\ideal C_i)$이기도 하므로 $[\ideal T_0,\ideal C_i]\cdot\ideal R\equiv0(\ideal S)$이다. 만일 $\ideal T_0$가 $\ideal S_i$의 진약수라면, 축약표현 $\ideal S=[\ideal S_i,\ideal C_i]$ 때문에 $[\ideal T_0,\ideal C_i]$도 $\ideal S$의 진약수가 되어 가정에 어긋난다.

4. $\ideal T_0$가 $\ideal R$의 진약수이고 $\ideal T_0\cdot\ideal S_i\equiv0(\ideal R)$이면 $\ideal T_0\cdot\ideal S\equiv0(\ideal R)$도 성립한다. 그런데 $\ideal S$가 $\ideal R$에 대하여 상대소이므로 $\ideal T\equiv0(\ideal R)$이어야 하고, 이는 $\ideal T_0$가 $\ideal R$의 진약수라는 가정과 모순된다.

정의 V에 따르면 특히 모든 아이디얼은 $\Sigma$의 모든 원소로 이루어진 단위아이디얼 $\ideal O$에 대하여 상대소이다.\footnote{$\ideal O$는 나눗셈과 최소공배수에 대해서만 단위의 구실을 하고, 곱셈에 대해서는 그렇지 않다. 예를 들어 상수항이 없는 $x$의 정수계수 다항식 전체의 환영역에서는 $\ideal O=(x)$이고, 모든 짝수의 환영역에서는 $\ideal O=(2)$이다.} 그러나 다음 정리들은 $\ideal O$와 다른 아이디얼들에 대해서만 성립한다.

방금 증명한 정리 X로부터 먼저 다음을 얻는다.

\textbf{보조정리 V.} \emph{한 아이디얼을 $\ideal O$와 다르며 서로 상대소인 아이디얼들의 최소공배수로 나타내는 모든 표현은 축약되어 있다.}

\[
 \ideal M=[\ideal R_1\ideal R_2\cdots\ideal R_\sigma]=[\ideal R_i,\ideal L_i]
\]
가 그러한 표현이라고 하자. 정리 X의 2에 따라 $\ideal L_i$도 $\ideal R_i$에 대하여 상대소이다. 따라서 $\ideal R_i$는 $\ideal L_i$의 약수가 될 수 없고, 이 표현은 최단이다. 실제로 $\ideal L_i\equiv0(\ideal R_i)$이면서 $\ideal L_i$가 $\ideal R_i$에 대하여 상대소라면 $\ideal O\cdot\ideal L_i\equiv0(\ideal R_i)$로부터 $\ideal O\equiv0(\ideal R_i)$, 따라서 $\ideal R_i=\ideal O$가 따르지만 이는 가정에서 배제되었다.

이제 $\ideal R_i$를 진약수 $\ideal R_i^*$로 바꿀 수 있다고 하자. 그러면 보조정리 II에 따라 $\ideal R_i$는 가약이고
\[
 \ideal R_i=[\ideal R_i^*,(\ideal R_i,\ideal L_i)]
\]
이다. 여기서 필요하면 $(\ideal R_i,\ideal L_i)$를 진약수 $(\ideal R_i,\ideal L_i)^*$로 바꾸고 $\ideal R_i^*$도 진약수로 바꾸어 $\ideal R_i$의 축약표현을 얻는다. 그러면 정리 X의 3에 따라 $\ideal L_i$는 $(\ideal R_i,\ideal L_i)^*$에 대해서도 상대소이고, 이 아이디얼은 최단표현 때문에 $\ideal O$와 다르다. 그런데 $(\ideal R_i,\ideal L_i)^*$는 $(\ideal R_i,\ideal L_i)$의 약수이고 $(\ideal R_i,\ideal L_i)$는 $\ideal L_i$의 약수이므로 모순이다.

정리 X로부터 대응 소아이디얼과의 관계를 주는 다음 정리도 얻는다.

\textbf{정리 XI.} \emph{$\ideal R$이 $\ideal S$에 대하여 상대소이고 $\ideal S$가 $\ideal O$와 다르면, $\ideal R$에 결부된 어느 소아이디얼도\footnote{$\ideal R$에 대응 소아이디얼의 정의는 \S~5의 끝을 보라.} $\ideal S$에 대응 소아이디얼로 나누어지지 않는다. 거꾸로 그러한 나눗셈 관계가 하나도 없으면 $\ideal R$은 $\ideal S$에 대하여 상대소이고, 물론 $\ideal S$는 $\ideal O$와 다르다.}

\[
 \ideal R=[\ideal Q_1\cdots\ideal Q_\alpha],\qquad
 \ideal S=[\ideal Q_1^*\cdots\ideal Q_\beta^*]
\]
가 $\ideal R$과 $\ideal S$를 극대 으뜸아이디얼들로 나타내는 축약표현이고, $\ideal P_1,\ldots,\ideal P_\alpha$와 $\ideal P_1^*,\ldots,\ideal P_\beta^*$가 각각 대응 소아이디얼들이라고 하자. 우리는 다음과 같은 형태로 주장을 보인다. 어떤 $\ideal P$가 어떤 $\ideal P^*$로 나누어지면 $\ideal R$은 $\ideal S$에 대하여 상대소일 수 없으며, 그 역도 성립한다.

실제로 $\ideal P_\mu\equiv0(\ideal P_\nu^*)$라고 하자. 그러면 $\ideal Q_\mu\equiv0(\ideal P_\nu^*)$이기도 하고, $\ideal P_\nu^*$의 정의에 따라 $\ideal Q_\mu^{\sigma_\nu}\equiv0(\ideal Q_\nu^*)$, 따라서 $\ideal R^{\sigma_\nu}\equiv0(\ideal Q_\nu^*)$이다. $\ideal R$의 거듭제곱 가운데 $\ideal Q_\nu^*$로 나누어지는 가장 낮은 것을 $\ideal R^\tau$라고 하자. $\tau=1$이면 $\ideal R$이 $\ideal Q_\nu^*$로 나누어진다. 그런데 가정에 따라 $\ideal S$가 $\ideal O$와 다르므로 $\ideal R$은 $\ideal Q_\nu^*$에 대하여 상대소가 아니다. $\tau\ge2$이면
\[
 \ideal R^{\tau-1}\cdot\ideal R\equiv0(\ideal Q_\nu^*),\qquad
 \ideal R^{\tau-1}\not\equiv0(\ideal Q_\nu^*)
\]
이므로 역시 $\ideal R$은 $\ideal Q_\nu^*$에 대하여 상대소가 아니다.\footnote{$\Sigma$가 단위원을 포함할 필요가 없으므로 $\ideal R^0$은 정의되지 않는다. 따라서 $\tau=1$인 경우를 먼저 따로 다루어야 했다. $\Sigma$가 단위원을 포함하는 경우에도 $\tau=0$은 $\ideal S$에 관한 가정 때문에 배제된다.} 그러므로 두 경우 모두 정리 X의 3에 따라 $\ideal R$은 $\ideal S$에 대해서도 상대소가 아니다.

거꾸로 $\ideal R$이 $\ideal S$에 대하여 상대소가 아니면, 정리 X의 1에 따라 $\ideal R$은 적어도 하나의 $\ideal Q_\nu^*$에 대하여 상대소가 아니다. 따라서
\[
 \ideal T_0\cdot\ideal R\equiv0(\ideal Q_\nu^*),\qquad
 \ideal T_0\not\equiv0(\ideal Q_\nu^*)
\]
이고, $\ideal Q_\nu^*$가 으뜸이므로
\[
 \ideal R^\tau\equiv0(\ideal Q_\nu^*),\qquad
 \text{따라서}\qquad
 \ideal Q_1^\tau\cdots\ideal Q_\alpha^\tau\equiv0(\ideal Q_\nu^*)
\]
이다. 그러면 대응 소아이디얼들에 대하여
\[
 \ideal P_1^{\tau\varrho_1}\cdots\ideal P_\alpha^{\tau\varrho_\alpha}
 \equiv0(\ideal P_\nu^*)
\]
를 얻는다. 소아이디얼의 성질에 따라 적어도 하나의 $\ideal P$가 $\ideal P_\nu^*$로 나누어진다. 이로써 정리 XI가 증명되었다.

이제 정리 X와 XI로부터 상대소-기약아이디얼로의 분해가 존재하고\footnote{분해의 존재는 \S~2에서 유한개의 기약아이디얼로의 분해가 존재함을 증명한 것과 정확히 같은 방식으로 직접 증명할 수도 있다.} 유일함을 다음과 같이 얻는다.

$\ideal M=[\ideal Q_1\cdots\ideal Q_\alpha]$가 $\ideal M$을 극대 으뜸아이디얼들로 나타내는 축약표현(또는 적어도 최단표현)이고, $\ideal P_1,\ldots,\ideal P_\alpha$가 대응 소아이디얼들이라고 하자. 어느 한 군의 아이디얼도 다른 군의 아이디얼로 나누어지지 않도록 $\ideal P$들을 군들로 묶되, 각 군 자체는 이 성질을 갖는 두 부분군으로 더 나눌 수 없게 한다.

그러한 군 분할을 구성하려면, 정의상 각 군 $G$가 한 아이디얼 $\ideal P$와 함께 $\ideal P_1,\ldots,\ideal P_\alpha$ 가운데 나타나는 그 모든 약수와 배수(즉 $\ideal P$로 나누어지는 아이디얼)를 포함해야 함을 관찰하면 된다. 이를테면 $\ideal P^{(i_1)}$들을 $\ideal P$의 모든 배수, $\ideal P_{j_1}^{(i_1)}$들을 $\ideal P^{(i_1)}$의 모든 약수, $\ideal P_{j_1}^{(i_1i_2)}$들을 $\ideal P_{j_1}^{(i_1)}$의 모든 배수 등으로 정한다. 일반적으로 $\ideal P_{j_1\ldots j_{\lambda-1}}^{(i_1\ldots i_\lambda)}$들을 $\ideal P_{j_1\ldots j_{\lambda-1}}^{(i_1\ldots i_{\lambda-1})}$의 모든 배수로, $\ideal P_{j_1\ldots j_\lambda}^{(i_1\ldots i_\lambda)}$들을 $\ideal P_{j_1\ldots j_{\lambda-1}}^{(i_1\ldots i_\lambda)}$의 모든 약수로 정한다. 전체적으로 유한개의 아이디얼 $\ideal P$만을 다루므로 이 절차는 유한 단계 뒤에 끝나야 한다. 곧 앞서 나온 모든 아이디얼과 다른 새 아이디얼을 더 이상 내놓지 않는다. 이렇게 얻은 아이디얼 $\ideal P$들의 전체는 실제로 원하는 성질을 가진 한 군 $G$를 이룬다.

실제로 정의에 따라 $G$는 각 $\ideal P_{j_1\ldots j_{\lambda-1}}^{(i_1\ldots i_{\lambda-1})}$와 함께 그 모든 배수 $\ideal P_{j_1\ldots j_{\lambda-1}}^{(i_1\ldots i_\lambda)}$를 포함하고, 각 $\ideal P_{j_1\ldots j_{\lambda-1}}^{(i_1\ldots i_\lambda)}$와 함께 그 모든 약수 $\ideal P_{j_1\ldots j_\lambda}^{(i_1\ldots i_\lambda)}$를 포함한다. 여기에는 $\ideal P_{j_1\ldots j_{\lambda-1}}^{(i_1\ldots i_{\lambda-1})}$의 모든 약수도 들어 있고, $\ideal P_{j_1\ldots j_{\lambda-1}}^{(i_1\ldots i_\lambda)}$의 배수들도 다시 같은 방식으로 들어 있다. 따라서 $G$에 들어 있지 않은 아이디얼은 $G$에 들어 있는 아이디얼의 약수도 배수도 될 수 없다.

$G$는 기약성 조건도 만족한다. 실제로 $G$가 두 부분군 $G^{(1)}$, $G^{(2)}$로 나뉘고, 이를테면 $G^{(1)}$이 $\ideal P_{j_1\ldots j_\lambda}^{(i_1\ldots i_\lambda)}$(또는 $\ideal P_{j_1\ldots j_{\lambda-1}}^{(i_1\ldots i_\lambda)}$)를 포함한다고 하자. 그러면 앞선 아이디얼들은 번갈아 그 배수와 약수(또는 약수와 배수)이므로 $G^{(1)}$은 그 모두를 포함하고, 따라서 $\ideal P$와 군 $G$ 전체를 포함한다. $G$에 들어 있지 않은 아이디얼들에 같은 절차를 적용하면 모든 $\ideal P$의 군 분할 $G_1,\ldots,G_\sigma$를 얻으며, 이 분할은 원하는 성질들을 가진다. 그러한 군 분할은 유일하다. 실제로 $G'_1,\ldots,G'_\tau$가 두 번째 분할이고 $\ideal P_{j_1\ldots j_\lambda}^{(i_1\ldots i_\lambda)}$(또는 $\ideal P_{j_1\ldots j_{\lambda-1}}^{(i_1\ldots i_\lambda)}$)가 $G'_i$의 원소이면, 위의 논증에 따라 $G'_i$는 군 $G$ 전체를 포함하고, 기약성 성질에 따라 $G$와 일치한다.

$\ideal P_{i\mu}$($\mu=1,\ldots,\lambda_i$)가 군 $G_i$에 함께 묶인 아이디얼 $\ideal P$들을 나타낸다고 하자. $\ideal P$들이 모두 서로 다르므로 최단표현에서 이들에 대응하는 으뜸아이디얼 $\ideal Q_{i\mu}$들은 유일하게 정해진다. 이제
\[
 \ideal R_i=[\ideal Q_{i1}\cdots\ideal Q_{i\lambda_i}],\qquad
 \text{그러면}\qquad
 \ideal M=[\ideal R_1\cdots\ideal R_\sigma]
\]
로 놓자. 이것이 $\ideal M$의 상대소-기약아이디얼 분해임을 보이겠다. 먼저 $\ideal R_i=[\ideal Q_{i_1}\cdots\ideal Q_{i\lambda_i}]$가 최단표현일 뿐이어도 정리 XI를 적용할 수 있음을 관찰하자. 정리 IX의 따름정리에 따라 대응 소아이디얼들이 이미 유일하게 정의되기 때문이다. 이제 $\ideal R_i$에 결부된 어느 소아이디얼 $\ideal P_{i\mu}$도 $\ideal R_j$에 대응 소아이디얼 $\ideal P_{j\nu}$로 나누어지지 않고, 그 역도 성립하므로 정리 XI에 따라 $\ideal R_i$와 $\ideal R_j$는 서로 상대소이다. 또한 각 $\ideal R_i$는 군 $G_i$의 기약성 성질 때문에 정리 XI에 따라 상대소-기약이다. 가약인 경우에는 언제나 $\ideal O$와 다른 아이디얼들만을 다루기 때문이다. 더 나아가 처음에 $\ideal Q$들의 최단표현만에서 출발했더라도 보조정리 V에 따라 이 표현은 축약되어 있다. 거꾸로 정리 XI에 따라 상대소-기약아이디얼로의 모든 분해는 위에서 설명한 $\ideal P_{i\mu}$들의 군 분할을 낳는다.

이제
\[
 \ideal M=[\bar{\ideal R}_1\cdots\bar{\ideal R}_\tau]
\]
가 $\ideal M$을 상대소-기약아이디얼들로 나타내는 두 번째 표현이라고 하자. 이 표현은 보조정리 V에 따라 축약되어 있다. $\bar{\ideal R}$들을 극대 으뜸아이디얼들로 분해해 보면 대응 소아이디얼들이 일치하므로, 위에서 유일성을 증명한 이 소아이디얼들의 군 분할도 일치한다. 따라서 $\tau=\sigma$이고, $\ideal R_i$와 $\bar{\ideal R}_i$가 같은 군에 속하도록 지표를 정할 수 있다. 이제 아이디얼과 보충부분으로 된 표현을
\[
 \ideal M=[\ideal R_i,\ideal L_i]
          =[\bar{\ideal R}_i,\bar{\ideal L}_i]
\]
라 하자. $\bar{\ideal R}_i$가 $\ideal R_i$와 같은 군에 대응하므로 정리 XI에 따라 $\ideal L_i$는 $\bar{\ideal R}_i$에 대하여 상대소이고, $\bar{\ideal L}_i$는 $\ideal R_i$에 대하여 상대소이다. 또한
\[
 \ideal R_i\cdot\ideal L_i\equiv0(\bar{\ideal R}_i),
 \qquad
 \bar{\ideal R}_i\cdot\bar{\ideal L}_i\equiv0(\ideal R_i)
\]
이므로
\[
 \ideal R_i\equiv0(\bar{\ideal R}_i),
 \qquad
 \bar{\ideal R}_i\equiv0(\ideal R_i),
 \qquad
 \ideal R_i=\bar{\ideal R}_i
\]
를 얻는다. 이로써 다음이 증명되었다.

\textbf{정리 XII.} \emph{각 아이디얼은 유한개의 서로 상대소인 상대소-기약아이디얼들의 최소공배수로 유일하게 나타낼 수 있다.}

\begin{center}
\S~7.\\[0.25em]
\textbf{고립아이디얼의 유일성.}
\end{center}

\textbf{정의 VI.} \emph{최단표현 $\ideal M=[\ideal R,\ideal L]$가 $\ideal L$에 관하여 축약되어 있다고 하자. $\ideal R$에 결부된 어느 소아이디얼도 $\ideal L$에 대응 소아이디얼의 약수가 아니면, 다시 말해 $\ideal L$이 $\ideal R$에 대하여 상대소이면 $\ideal R$을 고립아이디얼이라 한다.}

그러면 표현 $\ideal M=[\ideal R,\ideal L]$은 보조정리 V의 표현 $\ideal M=[\ideal R_i,\ideal L_i]$에 대한 조건을 만족하므로, 보조정리 V와 함께 다음이 증명된다.

\textbf{보조정리 VI.} \emph{$\ideal R$이 $\ideal L$에 관하여 축약된 최단표현 $\ideal M=[\ideal R,\ideal L]$의 고립아이디얼이면, 이 표현은 $\ideal R$에 관해서도 축약되어 있다.}

따라서 고립아이디얼의 경우에는 축약표현을 다루므로, $\ideal R$과 $\ideal L$을 기약아이디얼들로 분해할 때 나타나는 대응 소아이디얼들을 합치면 그에 대응하는 $\ideal M$의 분해에서 나타나는 유일하게 정해진 대응 소아이디얼들이 된다. 그러므로 $\ideal R$의 기약아이디얼 분해에 속하는 어느 소아이디얼도, 그에 대응하는 $\ideal M$의 분해에 속하는 나머지 소아이디얼의 약수가 아니다. 거꾸로 이 조건이 성립하고 $\ideal R$이 $\ideal R$에 관하여 축약된 표현 $\ideal M=[\ideal R,\ideal L]$에 적어도 한 번 나타나면, 따라서 어떤 축약표현 $\ideal M=[\ideal R,\ideal L^*]$에도 나타나면, 정의 VI에 따라 $\ideal R$은 고립되어 있다. 이로부터 특정 보충부분 $\ideal L$에 의존하지 않는 다음 정의를 얻는다.

\textbf{정의 VIa.} \emph{$\ideal R$의 기약아이디얼 분해에 속하는 소아이디얼들이 그에 대응하는 $\ideal M$의 분해에 속하는 나머지 소아이디얼들의 약수가 아니고, $\ideal R$이 $\ideal R$에 관하여 축약된 표현 $\ideal M=[\ideal R,\ideal L]$에 적어도 한 번 나타나면 $\ideal R$을 고립아이디얼이라 한다.}\footnote{통상적인 대응 소아이디얼들(\S~5 끝)을 사용한다면, $\ideal L$에 대응 소아이디얼들이 모두 $\ideal R$에 대응 소아이디얼들과 다르다는 조건을 따로 덧붙여야 한다. 따라서 정의 VIa에서는 표현이 보충부분에 관하여 축약되어 있다고 더 이상 가정할 필요가 없다.}

이제
\[
 \ideal M=[\ideal R,\ideal L]=[\bar{\ideal R},\bar{\ideal L}]
\]
가 고립아이디얼 $\ideal R,\bar{\ideal R}$와 보충부분 $\ideal L,\bar{\ideal L}$로 된 $\ideal M$의 두 표현이고, $\ideal R$과 $\bar{\ideal R}$에 대응 소아이디얼들이 일치한다고 하자. 표현들이 축약되도록 $\ideal L$과 $\bar{\ideal L}$을 각각 약수 $\ideal L^*$와 $\bar{\ideal L}^{*}$로 바꾸면, $\ideal L^*$와 $\bar{\ideal L}^{*}$에 대응 소아이디얼들도 일치한다. 따라서 정리 XI에 따라 $\ideal L^*$는 $\bar{\ideal R}$에 대하여 상대소이고 $\bar{\ideal L}^{*}$는 $\ideal R$에 대하여 상대소이다. 또한
\[
 \ideal R\cdot\ideal L^*\equiv0(\bar{\ideal R}),
 \qquad
 \bar{\ideal R}\bar{\ideal L}^{*}\equiv0(\ideal R)
\]
이므로
\[
 \ideal R\equiv0(\bar{\ideal R}),
 \qquad
 \bar{\ideal R}\equiv0(\ideal R),
 \qquad
 \ideal R=\bar{\ideal R}
\]
를 얻는다. 따라서 고립아이디얼들은 대응 소아이디얼들에 의해 유일하게 정해진다. 특히 이것은 기약아이디얼 또는 극대 으뜸아이디얼로의 분해에 관한 정리 VII과 IX를 강화한다. 정리 IX의 따름정리에 따라 여기서는 최단표현만을 가정해도 된다. 정리하면 다음과 같다.

\textbf{정리 XIII.} \emph{한 아이디얼을 기약아이디얼들 또는 극대 으뜸아이디얼들의 최소공배수로 나타내는 모든 최단표현에서, 고립된 기약아이디얼들 또는 고립된 극대 으뜸아이디얼들은 유일하게 정해진다. 다의성은 고립되지 않은 기약아이디얼들 또는 극대 으뜸아이디얼들에만 관계한다.}\footnote{다항식의 아이디얼을 극대 으뜸아이디얼들로 분해하는 경우에 이 정리는 이미 매콜리가 증명 없이 제시하였다. 고립 및 비고립(매입된) 으뜸아이디얼에 대한 그의 정의는 아래에서 제시하는 것의 비유리적 정식화로 볼 수 있다.} \emph{일반적으로 고립아이디얼들은 대응 소아이디얼들에 의해 유일하게 정해진다.}

그러한 기약아이디얼 또는 극대 으뜸아이디얼의 최단표현에서 아이디얼 $\ideal B_i$ 또는 $\ideal Q_j$가 고립되어 있지 않다고 하자. 그러면 정의에 따라 보충부분 $\ideal A_i$ 또는 $\ideal L_j$는 $\ideal P_i$ 또는 $\ideal P_j$로 나누어지고, 따라서
\[
 \ideal A_i^{\varrho_i}\equiv0(\ideal B_i)
 \qquad\text{또는}\qquad
 \ideal L_j^{\sigma_j}\equiv0(\ideal Q_j)
\]
이다. 거꾸로 이 관계들이 성립하면 $\ideal P_i$ 또는 $\ideal P_j$는 보충부분에 결부된 적어도 하나의 소아이디얼의 약수이다. 따라서 정리 IX의 따름정리에 따라, $\ideal L_j$의 약수로서 $\ideal L_j^*$에 관하여 축약된 표현을 주는 $\ideal L_j^*$에 대응 소아이디얼 하나의 약수이기도 하다. 그러므로 $\ideal B_i$ 또는 $\ideal Q_j$는 고립되어 있지 않다. 특히 $\ideal M$의 분해에서 대응 소아이디얼이 한 번보다 많이 나타나는 기약아이디얼 $\ideal B_i$는 언제나 고립되어 있지 않다. \emph{따라서 고립되지 않은 으뜸아이디얼은 모든 보충부분의 어떤 거듭제곱이 그 아이디얼로 나누어진다는 성질로도 특징지어지고, 고립된 으뜸아이디얼은 이 성질이 성립할 수 없다는 것으로 특징지어진다.}

\begin{center}
\S~8.\\[0.25em]
\textbf{한 아이디얼을 서로소-기약아이디얼들의 곱으로 유일하게 나타내기.}
\end{center}

바탕 환영역 $\Sigma$가 단위원, 즉 $\Sigma$의 모든 원소 $a$에 대하여 $\varepsilon\cdot a=a$가 되는 원소 $\varepsilon$을 가진다고 하자.\footnote{곱셈이 가환이므로 $\Sigma$는 단위원을 둘 이상 가질 수 없다. 실제로 두 단위원 $\varepsilon_1,\varepsilon_2$에 대하여 $\varepsilon_1\varepsilon_2=\varepsilon_2=\varepsilon_1$이다.} 그러면 서로소 아이디얼들을 다음과 같이 정의할 수 있다.

\textbf{정의 VIII.} \emph{두 아이디얼 $\ideal R,\ideal S$의 최대공약아이디얼이 $\Sigma$의 모든 원소로 이루어진 단위아이디얼 $\ideal O=(\varepsilon)$와 같으면 $\ideal R$과 $\ideal S$를 서로소라고 한다. 한 아이디얼을 쌍마다 서로소인 아이디얼들의 최소공배수로 나타낼 수 없으면 그 아이디얼을 서로소-기약이라 한다.}

서로소인 두 아이디얼은 언제나 서로 상대소임을 관찰하자. 실제로 정의에 따라 $r\equiv0(\ideal R)$, $s\equiv0(\ideal S)$이고 $\varepsilon=r+s$인 두 원소 $r,s$가 존재한다. $\ideal T\cdot\ideal R\equiv0(\ideal S)$이면 $\ideal T\cdot r\equiv0(\ideal S)$이고, 따라서 $\ideal T\cdot\varepsilon=\ideal T\equiv0(\ideal S)$이다. 마찬가지로 $\bar{\ideal T}\cdot\ideal S\equiv0(\ideal R)$이면 $\bar{\ideal T}\equiv0(\ideal R)$이다.\footnote{그러나 역은 성립하지 않는다. 예를 들어 아이디얼 $\ideal R=(x)$와 $\ideal S=(y)$는 서로 상대소이지만 서로소는 아니다.} 따라서 보조정리 V에 따라 쌍마다 서로소인 아이디얼들로 된 모든 표현은 동시에 축약되어 있다.

유일성 증명의 바탕은 정리 X에 대응하는 다음 정리이다.

\textbf{정리 XIV.} \emph{$\ideal R$이 아이디얼 $\ideal S_1,\ldots,\ideal S_\lambda$ 각각과 서로소이면, $\ideal R$은 $\ideal S=[\ideal S_1\cdots\ideal S_\lambda]$와도 서로소이다. 거꾸로 $\ideal R$과 $\ideal S$가 서로소이면 $\ideal R$은 각 $\ideal S_j$와도 서로소이다. $\ideal R=[\ideal R_1\cdots\ideal R_\mu]$이고 각 $\ideal R_i$가 각 $\ideal S_j$와 서로소이면 $\ideal R$과 $\ideal S$는 서로소이며, 여기서도 역이 성립한다.}

실제로 $\ideal R$이 각 $\ideal S_j$와 서로소이면
\[
 s_j\equiv0(\ideal S_j),
 \qquad
 s_j\equiv\varepsilon(\ideal R)
\]
인 원소 $s_j$들이 존재한다. 따라서
\[
 s_1\cdot s_2\cdots s_\lambda\equiv0(\ideal S),
 \qquad
 s_1\cdot s_2\cdots s_\lambda\equiv\varepsilon(\ideal R),
 \qquad
 (\ideal R,\ideal S)=(\varepsilon)
\]
이다. 한편 $(\ideal R,\ideal S)$는 각 $(\ideal R,\ideal S_j)$로 나누어지므로 역도 성립한다. 이 논증을 여러 번 적용하면 주장의 둘째 부분을 얻는다. 실제로 고정된 $i$에 대하여 $\ideal R_i$가 $\ideal S_1,\ldots,\ideal S_\lambda$ 각각과 서로소이면 $\ideal R_i$는 $\ideal S$와 서로소이다. 이 사실이 모든 $i$에 대하여 성립하면 서로소 개념이 상호적이므로 $\ideal S$는 $\ideal R$과 서로소이다. 거꾸로 $\ideal R$과 $\ideal S$가 서로소이면 $\ideal S$는 각 $\ideal R_i$와 서로소이고, 따라서 $\ideal R_i$는 $\ideal S$와 서로소이다.

서로소-기약아이디얼로의 분해가 존재하고 유일함을 보이는 증명은\footnote{분해의 존재는 다시 \S~2와 비슷하게 직접 증명할 수 있고, 유일성도 직접 증명할 수 있다(서론에서 슈마이들러와 뇌터--슈마이들러에 관하여 말한 내용을 참조하라). 여기서 제시하는 증명은 서로소-기약아이디얼의 구조도 함께 밝혀 준다.} 상대소아이디얼의 경우에 대응하는 증명처럼 유일한 군 분할로 귀착된다. 그러나 $\ideal M$의 상대소-기약아이디얼 $\ideal R_1,\ldots,\ideal R_\sigma$는 정리 XII에 따라 유일하게 정해져 있으므로, 여기서는 대응 소아이디얼들로 거슬러 올라갈 필요가 없다.

$\ideal M$의 유일하게 정해진 상대소-기약아이디얼 $\ideal R_1,\ldots,\ideal R_\sigma$를 다음과 같이 군들로 묶는다. \emph{한 군의 각 아이디얼이 다른 군의 각 아이디얼과 서로소가 되게 하되, 각 군 자체는 한 부분군의 각 아이디얼이 다른 부분군의 각 아이디얼과 서로소가 되도록 두 부분군으로 나눌 수 없게 한다.} 그러한 군 분할은 다음과 같이 얻는다. 정의상 각 군 $G$는 한 아이디얼 $\ideal R$과 함께 $\ideal R$과 서로소가 아닌 모든 아이디얼도 포함해야 한다. 이들을 $\ideal R_{i_1}$이라 하자. 다시 이들과 서로소가 아닌 아이디얼들을 $\ideal R_{i_1i_2}$라 하고, 일반적으로 $\ideal R_{i_1\ldots i_\lambda}$가 $\ideal R_{i_1\ldots i_{\lambda-1}}$과 서로소가 아니라고 하자. 전체적으로 유한개의 아이디얼만을 다루므로 이 절차는 유한 단계 뒤에 끝나야 한다. 곧 앞서 나온 모든 아이디얼과 다른 새 아이디얼을 더 이상 내놓지 않는다. 이렇게 얻은 아이디얼들의 전체는 원하는 성질을 가진 한 군 $G$를 이룬다. 실제로 $G$에 들어 있지 않은 모든 아이디얼은 구성에 따라 $G$에 들어 있는 모든 아이디얼과 서로소이다. 더 나아가 $G$가 두 부분군 $G^{(1)}$, $G^{(2)}$로 나뉘고 이를테면 $\ideal R_{i_1\ldots i_\lambda}$가 $G^{(1)}$의 원소라고 하자. 서로소성은 상호적이므로 $G^{(1)}$은 $\ideal R_{i_1\ldots i_{\lambda-1}},\ldots,\ideal R_{i_1}$과 $\ideal R$, 따라서 군 $G$ 전체도 포함해야 한다. 이로써 기약성이 증명되었다. $G$에 들어 있지 않은 아이디얼들에 같은 절차를 적용하면 모든 $\ideal R$의 군 분할 $G_1,\ldots,G_\tau$를 얻는다.

이 군 분할은 유일하다. 실제로 $G'_1,\ldots,G'_\tau$가 두 번째 분할이고 $\ideal R_{i_1\ldots i_\lambda}$가 $G'_i$의 원소라고 하자. 그러면 $G'_i$는 $\ideal R$과 따라서 $G_1$도 포함하며, 기약성 조건 때문에 $G_1$과 다른 원소를 포함할 수 없다. 그러므로 $G'_i=G_1$이다.

이제 $\ideal T_i$를 군 $G_i$에 모인 아이디얼 $\ideal R$들의 최소공배수라 하자. 그러면
\[
 \ideal M=[\ideal T_1\cdots\ideal T_\tau]
\]
가 $\ideal M$을 서로소-기약아이디얼들로 나타내는 표현임을 보이겠다. 먼저 $\ideal T$들을 $\ideal R$들로 분해해 보면 $[\ideal T_1\cdots\ideal T_\tau]$가 실제로 $\ideal M$을 나타냄을 알 수 있다. 또한 정리 XIV에 따라 $\ideal T$들은 쌍마다 서로소이고, 각 $\ideal T$는 서로소-기약이다. 이로써 그러한 표현의 존재가 증명되었다.

유일성을 증명하기 위하여 $\ideal M=[\bar{\ideal T}_1\cdots\bar{\ideal T}_{\bar\tau}]$가 그러한 두 번째 표현이라고 하자. $\bar{\ideal T}$들을 상대소-기약아이디얼 $\ideal R$들로 분해하면 서로 다른 $\bar{\ideal T}$에서 나타나는 $\ideal R$들은 정리 XIV에 따라 서로소이고, 따라서 서로 상대소이기도 하다. 그러므로 이들은 $\ideal M$의 유일하게 정해진 상대소-기약아이디얼 $\ideal R$들과 일치한다. 또한 정리 XIV에 따라 $\bar{\ideal T}_i$들은 앞서 말한 성질을 가진 $\ideal R$들의 군 분할 $G'_i$를 만든다. 그런데 이 군 분할은 유일하고 각 $\bar{\ideal T}_i$는 군 $G'_i=G_i$에 의해 유일하게 정해지므로 $\bar{\ideal T}_i=\ideal T_i$이다. 이로써 유일성이 증명되었다.

쌍마다 서로소인 아이디얼들에서는 최소공배수가 곱과 같다. 실제로 정리 XIV에 따라 $\ideal M=[\ideal T_1\cdots\ideal T_\tau]$에서 보충부분 $\ideal L_i$도 $\ideal T_i$와 서로소이다. 따라서
\[
 t_i\equiv0(\ideal T_i),
 \qquad
 l_i\equiv0(\ideal L_i),
 \qquad
 \varepsilon=t_i+l_i
\]
인 원소들이 존재한다. $f\equiv0(\ideal T_i)$, $f\equiv0(\ideal L_i)$이면
\[
 f=f\varepsilon=ft_i+fl_i
 \quad\text{이므로}\quad
 f\equiv0(\ideal L_i\cdot\ideal T_i)
\]
이다. 거꾸로 $\ideal L_i\cdot\ideal T_i$는 $[\ideal L_i,\ideal T_i]$로 나누어지므로 $[\ideal L_i,\ideal T_i]=\ideal L_i\cdot\ideal T_i$이다. 이 절차를 $\ideal L_i$에 계속 적용하면 마침내
\[
 \ideal M=[\ideal T_1\cdots\ideal T_\tau]
       =\ideal T_1\cdot\ideal T_2\cdots\ideal T_\tau
\]
를 얻는다. 이로써 다음이 증명되었다.

\textbf{정리 XV.} \emph{각 아이디얼은 유한개의 쌍마다 서로소인 서로소-기약아이디얼들의 곱으로 유일하게 나타낼 수 있다.}

\begin{center}
\S~9.\\[0.25em]
\textbf{모듈로의 논의 확장. 기약모듈 분해에서 성분 수의 동일성.}
\end{center}

이제 으뜸아이디얼이나 소아이디얼이 아니라 \emph{기약아이디얼}에 관한 앞의 세 절의 내용이 더 약한 가정 아래에서도 성립함을 보이겠다. 실제로 그 절들은 곱셈의 가환법칙을 사용하지 않고 아이디얼이 모듈이라는 성질만을 사용한다. 따라서 이제 정의할 비가환 환영역에 관한 모듈에도 그대로 성립한다. 모듈을 정의하려면 다음 성질들을 가진 \emph{이중환영역 $(\Sigma,T)$}를 바탕으로 삼아야 한다.

\emph{$\Sigma$는 추상적으로 정의된 비가환환이다.} 즉 $\Sigma$는 원소 $a,b,c,\ldots$의 체계이고, 이 원소들에는 환의 덧셈 $(\#)$과 환의 곱셈 $(\times)$이라는 두 연산이 정의되어 있다. 이 연산들은 \S~1에서 세운 법칙들 가운데 환의 곱셈에 관한 가환법칙 4를 제외한 모든 법칙을 만족한다.

$T$는 원소 $\alpha,\beta,\gamma,\ldots$의 체계이며, $\Sigma$와 결합하여 여기에도 두 연산이 정의되어 있다. \emph{덧셈}은 두 원소 $\alpha,\beta$에서 세 번째 원소 $\alpha+\beta$를 유일하게 만들고, \emph{$T$의 원소 $\alpha$와 $\Sigma$의 원소 $c$의 곱셈}은 $T$의 원소 $c\cdot\alpha$를 유일하게 만든다.\footnote{따라서 여기서는 ``우측'' 곱셈, ``우측'' 환영역 $T$, 그리고 그에 따라 ``우측'' 모듈과 아이디얼을 다룬다. $T$에 좌측 곱셈 $\alpha\cdot c$를 바탕으로 삼으면 좌측모듈과 좌측아이디얼의 대응하는 이론을 얻는다. 즉 $M$은 $\alpha$와 함께 $\alpha\cdot c$도 포함한다. 이때 결합법칙은 $(\gamma\cdot b)\cdot a=\gamma\cdot(b\times a)$ 꼴이다.}

이 연산들에는 다음 법칙들이 성립한다.
\begin{enumerate}
\item 덧셈의 결합법칙: $(\alpha+\beta)+\gamma=\alpha+(\beta+\gamma)$.
\item 덧셈의 가환법칙: $\alpha+\beta=\beta+\alpha$.
\item 제한 없고 유일한 뺄셈의 법칙: 방정식 $\alpha+\xi=\beta$를 만족하는 원소 $\xi$가 $T$에 유일하게 존재한다($\xi=\beta-\alpha$로 쓴다).
\item 곱셈의 결합법칙: $a\cdot(b\cdot\gamma)=(a\times b)\cdot\gamma$.
\item 분배법칙:
\[
 (a\# b)\cdot\gamma=a\cdot\gamma+b\cdot\gamma,
 \qquad c\cdot(\alpha+\beta)=c\cdot\alpha+c\cdot\beta.
\]
\end{enumerate}
이 조건들로부터 영원소의 존재와 뺄셈과 곱셈의 결합에 대해서도 분배법칙
\[
 (a-b)\cdot\gamma=a\cdot\gamma-b\cdot\gamma,
 \qquad c\cdot(\alpha-\beta)=c\cdot\alpha-c\cdot\beta
\]
이 성립함은 잘 알려져 있다. 여기서 $(\dot{-})$는 $\Sigma$에서의 뺄셈이다. \emph{$\Sigma$가 단위원 $\varepsilon$을 포함하면 $T$의 모든 원소 $\alpha$에 대하여 $\varepsilon\cdot\alpha=\alpha$라고 한다.}

$(\Sigma,T)$의 모듈 $M$은 다음 두 조건을 만족하는 $T$의 원소들의 체계를 뜻한다.
\begin{enumerate}
\item \emph{$M$은 $\alpha$와 함께 $c\cdot\alpha$도 포함한다. 여기서 $c$는 $\Sigma$의 임의의 원소이다.}
\item \emph{$M$은 $\alpha,\beta$와 함께 차 $\alpha-\beta$도 포함한다.} 따라서 $\alpha$와 함께 모든 정수 $n$에 대한 $n\alpha$도 포함한다.\footnote{정수는 다시 환의 원소가 아니라 줄여 쓰는 기호로 이해해야 한다.}
\end{enumerate}
이 정의에 따라 \emph{$T$ 자체도 $(\Sigma,T)$의 모듈을 이룬다.} 특히 환영역 $T$와 거기서 정한 연산들이 환영역 $\Sigma$ 및 그 연산들과 일치하면 모듈 $M$은 $\Sigma$의 우측아이디얼 $M$이 된다. 더 나아가 $\Sigma$가 가환이라고 가정하면 통상적인 아이디얼 개념을 얻는다. 따라서 아이디얼은 모듈 개념의 특수한 경우이다.\footnote{가장 간단한 모듈의 예는 정수계수 선형형식들의 모듈이다. 여기서 $\Sigma$는 모든 유리정수로, $T$는 모든 정수계수 선형형식으로 이루어진다. $\Sigma$와 $T$에서 유리정수 대신 대수적 정수를 쓰거나, 이를테면 모든 짝수를 쓰면 조금 더 일반적인 모듈을 얻는다. 선형형식 대신 매번 계수 전체를 하나의 원소로 보면 $\Sigma$와 $T$의 연산들은 실제로 서로 다르다. 비가환 다항식 환영역의 아이디얼은 뇌터--슈마이들러의 공동논문에서 다룬다. 다른 특수한 비가환 환영역의 아이디얼은 후르비츠의 \emph{사원수의 수론 강의}(베를린, 슈프링어, 1919)와 거기서 인용한 뒤 파스키에의 논문들에서 다룬다.}

모듈에 대해서도 \S~1의 모든 정의가 유지된다. 따라서 $\alpha\eqzero{M}$ 또는 $N\eqzero{M}$은 $\alpha$ 또는 $N$의 모든 원소가 $M$의 원소임을 뜻한다. 달리 말해 $\alpha$ 또는 $N$은 $M$으로 나누어진다. $M$이 $N$과 다른 원소들을 포함하면 $M$은 $N$의 진약수가 된다. $N\eqzero{M}$, $M\eqzero{N}$이면 $M=N$이다. 최대공약과 최소공배수의 정의도 글자 그대로 유지된다. 특히 모듈 $M$이 유한개의 원소 $\alpha_1\ldots\alpha_\varrho$를 포함하고 $M=(\alpha_1\ldots\alpha_\varrho)$, 즉 모든 $\alpha\eqzero{M}$에 대하여
\[
 \alpha=c_1\alpha_1+\cdots+c_\varrho\alpha_\varrho+n_1\alpha_1+\cdots+n_\varrho\alpha_\varrho
\]
가 되며, 여기서 $c_i$가 $\Sigma$의 원소이고 $n_i$가 정수이면, $M$을 \emph{유한모듈}, $\alpha_1\ldots\alpha_\varrho$를 \emph{모듈기저}라고 한다.

\emph{이하에서는 \S~1과 마찬가지로 다음 유한성 조건을 만족하는 환영역 $(\Sigma,T)$만을 바탕으로 삼는다. $(\Sigma,T)$의 모든 모듈은 유한모듈, 따라서 모듈기저를 가진다.}

그러면 정리 I의 유한사슬 정리는 그곳의 증명이 보여 주듯 이 환영역 $(\Sigma,T)$의 모듈에도 성립하고, 따라서 \S\S~2, 3의 모든 전제가 충족된다. 최단표현과 축약표현에 관한 정의 I과 보조정리 I을 직접 옮길 수 있다. 각 모듈을 유한개의 기약모듈의 최소공배수로 나타낼 수 있다는 정리 II도 마찬가지이며, 보조정리 II는 그러한 최단표현이 모두 동시에 축약되어 있음을 보여 준다. 모듈의 가약성을 그 보충부분의 성질로 나타내는 정리 III도 유지되고, 여기서 보조정리 III과 마침내 \emph{한 모듈을 기약모듈들의 최소공배수로 나타내는 서로 다른 두 최단표현에서 성분 수가 같다}는 정리 IV가 따른다. 또한 정리 IV는 축약표현에 관한 보조정리 I의 역인 보조정리 IV를 준다.

\emph{따름정리.} 같은 논증은 양쪽 환영역 $T$, 즉 우측이면서 좌측인 환영역에 대하여 모듈을 일관되게 \emph{양쪽모듈}, 곧 $\alpha$와 함께 $c\cdot\alpha$와 $\alpha\cdot c$를 모두 포함하고 $\alpha,\beta$와 함께 $\alpha-\beta$를 포함하는 모듈로 이해해도 이 모든 정리와 정의가 유지됨을 보여 준다.

이 모든 정리는 \emph{나눗셈}과 \emph{최소공배수}라는 개념에만 의존하지만, 뒤의 유일성 정리들은 본질적으로 \emph{곱}의 개념에 의존하므로 직접 옮길 수 없다. 예를 들어 $T$의 두 원소의 곱이 정의되어 있지 않으므로 으뜸아이디얼과 소아이디얼의 정의를 모듈로 옮길 수 없다. 형식적으로 가능한 비가환환으로의 이식도 의미를 잃는다. 여기서는 대응 소아이디얼의 존재를 증명할 수 없고\footnote{실제로 여기서는 $p_1^{\lambda_1}\eqzero{\ideal Q}$, $p_2^{\lambda_2}\eqzero{\ideal Q}$로부터 $(p_1-p_2)^{\lambda_1+\lambda_2}\eqzero{\ideal Q}$가 따르지 않으며, $(a\cdot b)^\lambda\eqzero{\ideal Q}$로부터 $a^\lambda\cdot b^\lambda\eqzero{\ideal Q}$도 따르지 않는다. 따라서 $\ideal P$가 아이디얼임도 증명할 수 없고, 소아이디얼의 성질을 가짐도 증명할 수 없다. 양쪽아이디얼의 경우에는 $\ideal P$가 아이디얼임은 증명할 수 있지만 소아이디얼임은 증명할 수 없다.} 기약아이디얼이 으뜸임을 보이는 증명도 실패하기 때문이다. 그런데 이 두 사실이 뒤의 유일성 정리들의 토대이다. 반면 비가환환이 단위원을 가지면 서로소 및 서로소-기약아이디얼을 정의할 수 있고, 정리 II의 증명에 사용한 논증은 각 아이디얼을 \emph{유한개의} 쌍마다 서로소인 서로소-기약아이디얼들의 최소공배수로 나타낼 수 있음을 보여 준다.\footnote{특수한 ``완전가약'' 아이디얼에 대해서는 여기서도 유일성 정리를 세울 수 있다. 뇌터--슈마이들러의 공동논문을 참조하라.}

끝으로 $(\Sigma,T)$가 유한성 조건을 만족하기 위한 충분조건을 하나 덧붙인다. \emph{$\Sigma$가 단위원을 가지고 유한성 조건을 만족하며, $T$ 자체가 $(\Sigma,T)$의 유한모듈이면 $(\Sigma,T)$의 모든 모듈은 유한모듈이다.}

실제로 $\Sigma$에 단위원이 존재하므로
\[
 \ideal M=(f_1,\ldots,f_e),\qquad
 f=\bar b_1f_1+\cdots+\bar b_\varrho f_\varrho+n_1f_1+\cdots+n_\varrho f_\varrho,
\]
여기서 $\bar b_i$는 $\Sigma$의 원소이고 $n_i$는 정수인 표현은
\[
 f=b_1f_1+\cdots+b_\varrho f_\varrho
\]
라는 표현도 준다. 여기서 $b_i$는 $\Sigma$의 원소이다. 실제로 $f_i=\varepsilon f_i$이므로 $n_if_i=n_i\varepsilon f_i$이고, $n_i\varepsilon=(\varepsilon+\cdots+\varepsilon)$는 $\Sigma$에 속하므로 $b_i=\bar b_i+n_i\varepsilon$도 $\Sigma$의 원소이다. 이제 $T$에 관한 가정은 $T$의 모든 원소 $\alpha$가
\[
 \alpha=\bar a_1\tau_1+\cdots+\bar a_k\tau_k
       +n_1\tau_1+\cdots+n_k\tau_k
\]
꼴의 표현, 따라서
\[
 \alpha=a_1\tau_1+\cdots+a_k\tau_k
\]
꼴의 표현을 가진다는 뜻이다. 둘째 표현은 $\varepsilon\tau_i=\tau_i$이므로 앞에서와 같이 첫째 표현에서 얻어진다.

이제 $\alpha$가 $(\Sigma,T)$의 한 모듈 $M$의 원소들을 두루 달리면, $\tau_k$의 계수 $a_k$들은 $\Sigma$의 한 아이디얼 $\ideal M_k$를 두루 달린다. 따라서 위의 결과에 따라 모든 $a_k\eqzero{\ideal M_k}$는
\[
 a_k=b_1a_k^{(1)}+\cdots+b_\varrho a_k^{(\varrho)}
\]
로 나타난다. $\alpha^{(i)}$가 $M$의 원소이고 그 $\tau_k$의 계수가 $a_k^{(i)}$라고 하자. 그러면 $\alpha-b_1\alpha^{(1)}\ldots-b_\varrho\alpha^{(\varrho)}$는 $M$에 속하면서 $\tau_1,\ldots,\tau_{k-1}$에만 의존하는 원소이다. $\tau_k$를 더 이상 포함하지 않는 이 원소들 전체는 모듈 $M'$을 이루며, 여기에 같은 절차를 반복할 수 있다. 유한번 반복하면 주장이 증명된다.

\begin{center}
\S~10.\\[0.25em]
\textbf{다항식 환영역이라는 특수한 경우.}
\end{center}

1. 바탕 환영역 $\Sigma$가 임의의 복소수 계수를 갖는 $x_1,\ldots,x_n$의 모든 다항식으로 이루어진다고 하자. 이 환영역은 힐베르트의 모듈기저 정리(\emph{Math. Ann.} 36)에 따라 유한성 조건을 만족한다. 여기서는 \emph{우리 정리들과 소거이론 및 모듈론의 알려진 정리들 사이의 관계}를 다룬다.

이 관계는 잘 알려진 힐베르트 정리의 다음 특수한 경우로 얻어진다.\footnote{``완전 불변식계에 관하여'', \emph{Math. Ann.} 42 (1893), \S~3, 313쪽.}

$f$가 소아이디얼 $\ideal P$의 모든 다항식의 영점인 $x_1,\ldots,x_n$의 모든 유한 값계($\ideal P$의 영점)에서 0이 되면 $f$는 $\ideal P$로 나누어진다. 다시 말해 소아이디얼 $\ideal P$는 그 영점들에서 0이 되는 다항식 \emph{전체}로 이루어진다.\footnote{라스커[\emph{Math. Ann.} 60 (1905), 607쪽]가 보였듯 이 특수한 경우는 동차형식의 경우 직접 증명할 수도 있다. 거꾸로 여기서 동차 및 비동차 경우의 힐베르트 정리가 다시 따르며, 이를 다음과 같이 말할 수 있다. 아이디얼 $\ideal R$이 $\ideal M$의 모든 영점에서 0이 되면 $\ideal R$의 어떤 거듭제곱은 $\ideal M$으로 나누어진다. 그러나 이 정리와 위의 특수한 경우는 $x$의 값영역이 \emph{대수적으로 닫혀 있을 때}에만 성립한다. 따라서 우리 정리들만으로는 따를 수 없고 근의 존재를 이용해야 한다. 이를테면 영점이 $x_1=0,\ldots,x_n=0$뿐인 아이디얼은 어떤 차수 이상인 $x$들의 모든 멱곱을 포함한다는 대목에서 그러하다. 나머지 증명은 우리 정리들을 이용하면 라스커의 증명보다 조금 단순해진다. 라스커는 한 아이디얼을 극대 으뜸아이디얼들로 분해하는 일에도 힐베르트 정리를 사용해야 했기 때문이다.}

따라서 곱 $f\cdot g$가 $\ideal P$의 모든 영점에서 0이면 적어도 한 인자가 그 모든 영점에서 0이 된다. 그 영점들은 \emph{기약 대수도형}을 이룬다. 거꾸로 이것을 기약도형의 정의로 삼으면, 한 기약도형에서 0이 되는 모든 다항식은 소아이디얼을 이룬다. 따라서 소아이디얼과 기약도형은 일대일로 대응한다. 또한 $\ideal Q\eqzero{\ideal P}$, $\ideal P^e\eqzero{\ideal Q}$이므로 으뜸아이디얼의 영점들은 그에 대응 소아이디얼의 영점들과 일치한다.\footnote{으뜸아이디얼이 기약도형 하나를 가진다는 이 성질을 매콜리는 정의로 삼는다(서론 참조). 반면 라스커는 도형의 ``다양도''라는 개념만 정의에 넣고, 나머지는 \emph{추상적으로} 정의한다. $x_1=0,\ldots,x_n=0$에서만 0이 되는 으뜸아이디얼은 라스커에게 특수한 지위를 가진다.} 그러므로 한 아이디얼을 극대 으뜸아이디얼들의 최소공배수로 나타내면 그 아이디얼의 모든 영점을 기약도형들로 분해하게 되며, 라스커가 보였듯 그 역도 성립한다. \emph{따라서 대응 소아이디얼의 유일성을 증명하는 것은 여기서 대수도형을 기약도형들로 유일하게 분해할 수 있다는 소거이론의 기본정리에 대응한다.} 환영역 안의 다항식들을 기약다항식들의 곱으로 유일하게 나타낼 수 없고 따라서 소거이론도 존재하지 않는 특수한 다항식 환영역에서는, 우리의 정리를 그 소거이론 정리의 동치물로 볼 수 있다.

고립된 으뜸아이디얼들에 대응하는 기약도형들은 정확히 ``최소 레졸벤트''\footnote{예를 들어 J. König, \emph{Einleitung in die allgemeine Theorie der algebraischen Größen}(라이프치히, 토이프너, 1903), 235쪽을 참조하라.}에 나타나는 것들이다. 실제로 고립되지 않은 각 극대 으뜸아이디얼의 영점들은 적어도 하나의 고립아이디얼의 영점이기도 하다. 곧 그 고립아이디얼에 대응 소아이디얼은 고립되지 않은 아이디얼에 대응 소아이디얼로 나누어진다. 따라서 고립된 으뜸아이디얼의 유일성은 그 지수들에서 새로운 불변 중복도수를 준다. 으뜸아이디얼들을 기약아이디얼들로 분해하는 데 관한 유일성 정리들도 중복도의 관점에서 소거이론을 보완하는 것으로 볼 수 있다.

이 관찰들에 따라 여러 분해가 대수도형에 대하여 갖는 의미를 해석할 수 있다. 쌍마다 서로소인 아이디얼들에는 공통영점을 갖지 않는 도형들이 대응한다. 서로 상대소인 아이디얼들의 경우 한 아이디얼의 어느 기약도형도 다른 아이디얼의 공통영점도형이 아니다. 극대 으뜸아이디얼들은 모두 서로 다른 기약도형에서만 0이 된다. 기약아이디얼로 분해할 때에는 같은 기약도형이 여러 번 나타날 수도 있다.

일반 다항식 환영역 대신 모든 동차형식의 환영역을 바탕으로 삼아도 된다는 점을 덧붙인다. 거기서는 같은 차수의 형식들에 대해서만 덧셈이 정의되지만, 이 연산들 아래에서도 일반 정리들이 유지됨을 쉽게 확인할 수 있다.\footnote{그러나 다의성이 있는 경우 동차분해뿐 아니라 비동차분해도 존재할 수 있음은 다음 예가 보여 준다. $ (x^3,xy,y^3)=[(x^3,y);(y^3,x)]=[(xy,x^3,y^3,x+y^2);(xy,x^3,y^3,y+x^2)]$.}

서로 다른 네 분해의 가장 간단한 예 하나는, 아래에 공식을 적을 것인데, 위의 논의에 따라 한 직선과 그 직선과 각각 꼬인 위치에 있으면서 서로 교차하는 두 직선으로 주어진다. 이 두 직선 가운데 하나는 두 직선의 교점과 다른 한 점을 더 높은 중복도로 포함한다. 서로소-기약아이디얼 분해에는 첫 직선과 그에 꼬인 위치에 있는 도형으로의 분해가 대응한다. 이 도형은 상대소-기약아이디얼 분해에서 두 직선으로 나뉜다. 극대 으뜸아이디얼 분해에는 더 높은 중복도의 점을 분리해 내는 일이 대응하고, 기약아이디얼 분해에는 이 점의 분해가 수반된다.

이 점을 원점으로, 그 점을 지나는 직선을 $y$축으로, 이 직선과 교차하는 직선을 $x$축에 평행하게, 그에 꼬인 위치의 직선을 $z$축에 평행하게 잡으면, 그러한 배치는 다음 \emph{기약아이디얼들}로 나타낼 수 있다.\footnote{이 가운데 처음 세 아이디얼은 소아이디얼이므로 기약이다. $\ideal B_4$는 약수 $(x^2,y,z)$와 $(x,y,z)$만 가지므로 기약이고, $\ideal B_5$는 그 모든 약수가 다항식 $xy$를 포함하므로 기약이다.}
\par\medskip
\noindent\makebox[\linewidth][c]{\(\displaystyle
 \ideal B_1=(x-1,y);\qquad
 \ideal B_2=(y-1,z);\qquad
 \ideal B_3=(x,z);\qquad
 \ideal B_4=(x^3,y,z);
\)}\par
\noindent\makebox[\linewidth][c]{\(\displaystyle
 \ideal B_5=(x^2,y^2,z).
\)}\par\medskip
\emph{대응 소아이디얼들}은 다음과 같다.
\par\smallskip
\noindent\makebox[\linewidth][c]{\(\displaystyle
 \ideal P_1=\ideal B_1;\qquad
 \ideal P_2=\ideal B_2;\qquad
 \ideal P_3=\ideal B_3;\qquad
 \ideal P_4=\ideal P_5=(x,y,z).
\)}\par\medskip
\emph{극대 으뜸아이디얼들}은 다음과 같다.
\par\smallskip
\noindent\makebox[\linewidth][c]{\(\displaystyle
 \ideal Q_1=\ideal B_1;\qquad
 \ideal Q_2=\ideal B_2;\qquad
 \ideal Q_3=\ideal B_3;\qquad
 \ideal Q_4=[\ideal B_4,\ideal B_5]=(x^3,y^2,x^2y,z).
\)}\par\medskip
\emph{상대소-기약아이디얼들}은 다음과 같다.
\par\smallskip
\noindent\makebox[\linewidth][c]{\(\displaystyle
 \ideal R_1=\ideal Q_1;\qquad
 \ideal R_2=\ideal Q_2;\qquad
 \ideal R_3=[\ideal Q_3,\ideal Q_4]=(x^3,x^2y,xy^2,z).
\)}\par\medskip
\emph{서로소-기약아이디얼들}은 다음과 같다.
\par\smallskip
\noindent\makebox[\linewidth][c]{\(\displaystyle
 \ideal S_1=\ideal R_1;\qquad
 \ideal S_2=[\ideal R_2,\ideal R_3]
 =((y-1)x^3,(y-1)x^2y,(y-1)xy^2,z).
\)}\par\medskip
이로부터 \emph{전체아이디얼}은
\par\smallskip
\noindent\(\displaystyle \ideal M=[\ideal S_1,\ideal S_2]=\ideal S_1\cdot\ideal S_2\)\par
\noindent\(\displaystyle
 =((x-1)(y-1)x^3,(y-1)x^2y,(y-1)xy^2,(x-1)z,y(y-1)x^3,yz),
\)\par\smallskip
\noindent\begin{tabular*}{\linewidth}{@{}l@{\extracolsep{\fill}}c@{}}
여기서 & \(\displaystyle 1=-(y-1)x^3+(y-1)(x^3-1)+y,\)
\end{tabular*}\par
\noindent\makebox[\linewidth][c]{\(\displaystyle
 (y-1)(x^3-1)+y\eqzero{\ideal S_1};\qquad
 -(y-1)x^3\eqzero{\ideal S_2}
\)}\par
이다. 이때 아이디얼 $\ideal B_1,\ideal B_2,\ideal B_3$은 고립되어 있으므로 기약아이디얼 및 극대 으뜸아이디얼 분해에서도 유일하게 정해진다. 아이디얼 $\ideal B_4,\ideal B_5$와 $\ideal Q_4$는 고립되어 있지 않다. 이들은 유일하게 정해지지 않으며, 예를 들어
\par\smallskip
\noindent\makebox[\linewidth][c]{\(\displaystyle
 \ideal D_4=(x^3,y+\lambda x^2,z);\qquad
 \ideal D_5=(x^2+\mu xy,y^2,z)
\)}\par
로 바꿀 수 있다. 마찬가지로 $\ideal Q_4$는
\par\smallskip
\noindent\makebox[\linewidth][c]{\(\displaystyle
 \overline{\ideal Q}_4=(x^3,x^2y,y^2+\lambda xy,z)
\)}\par
로 바꿀 수 있다.

2. 일반 다항식 환영역과 정수계수 다항식 환영역뿐 아니라, 다항식으로 이루어진 모든 \emph{유한 정역}도 힐베르트의 모듈기저 정리가 보여 주듯 유한성 조건을 만족한다.\footnote{거꾸로 한 다항식 환영역이 유한성 조건을 만족하고, 각 다항식이 $x$에 관하여 더 낮은 차수의 승수를 사용하는 표현을 적어도 하나 허용하면 그 환영역은 유한 정역이다.} 이때 계수들은 임의의 체에 속하게 할 수 있다. 여기서는 \emph{모든 짝다항식의 환영역}에 대한 예를 하나 더 제시한다. 이는 $x^2\cdot y^2=(xy)^2$ 때문에 환영역의 다항식들을 환영역의 기약다항식들의 곱으로 유일하게 나타낼 수 없는 가장 간단한 환영역이다. 앞의 예와 같은 배치를 이제 다음 \emph{기약아이디얼들}로 나타낸다.
\par\medskip
\noindent\makebox[\linewidth][c]{\(\displaystyle
 \ideal B_1=(x^2-1,xy,y^2,yz);\qquad
 \ideal B_2=(y^2-1,xz,yz,z^2);
\)}\par
\noindent\makebox[\linewidth][c]{\(\displaystyle
 \ideal B_3=(x^2,xy,xz,yz,z^2);\qquad
 \ideal B_4=(x^4,xy,y^2,xz,yz,z^2);
\)}\par
\noindent\makebox[\linewidth][c]{\(\displaystyle
 \ideal B_5=(x^2,y^2,xz,yz,z^2).
\)}\par\medskip
\emph{대응 소아이디얼들}은 다음과 같다.
\par\smallskip
\noindent\makebox[\linewidth][c]{\(\displaystyle
 \ideal P_1=\ideal B_1;\qquad
 \ideal P_2=\ideal B_2;\qquad
 \ideal P_3=\ideal B_3;\qquad
 \ideal P_4=\ideal P_5=(x^2,xy,y^2,xz,yz,z^2).
\)}\par\medskip
$\ideal P_1$이 소아이디얼이라는 사실은 환영역의 각 다항식이
\par\smallskip
\noindent\makebox[\linewidth][c]{\(\displaystyle
 f\equiv\varphi(z^2)+xz\,\psi(z^2)\;(\ideal P_1)
\)}\par\smallskip
꼴이 된다는 데서 따른다. 실제로
\par\smallskip
\noindent\makebox[\linewidth][c]{\(\displaystyle
 f_1\equiv\varphi_1(z^2)+xz\,\psi_1(z^2);\qquad
 f_2\equiv\varphi_2(z^2)+xz\,\psi_2(z^2)
\)}\par\smallskip
라고 하면 $f_1\cdot f_2\eqzero{\ideal P_1}$로부터 방정식들
\par\smallskip
\noindent\makebox[\linewidth][c]{\(\displaystyle
 \varphi_1\varphi_2+z^2\psi_1\psi_2=0;\qquad
 \varphi_2\psi_1+\varphi_1\psi_2=0
\)}\par
이 따르고, 여기서 $f_1\eqzero{\ideal P_1}$ 또는 $f_2\eqzero{\ideal P_1}$를 얻는다. 꼭 같은 방식으로 $\ideal P_2$가 소아이디얼임을 알 수 있다. 환영역의 각 다항식은 $\ideal P_3$을 법으로 $y^2$의 다항식과 합동이므로 $\ideal P_3$도 소아이디얼이다. $\ideal P_4$는 환영역의 모든 다항식으로 이루어진다. 따라서 $\ideal B_1,\ideal B_2,\ideal B_3$은 소아이디얼이므로 기약이다. 한편 $\ideal B_4$와 $\ideal B_5$는 각각 진약수 $\ideal P_4$ 하나만을 가지므로 반드시 기약이다.

기약아이디얼들로부터 다음 \emph{극대 으뜸아이디얼들}을 얻는다.
\par\medskip
\noindent\makebox[\linewidth][c]{\(\displaystyle
 \ideal Q_1=\ideal B_1;\qquad
 \ideal Q_2=\ideal B_2;\qquad
 \ideal Q_3=\ideal B_3;\qquad
 \ideal Q_4=[\ideal B_4,\ideal B_5]=(x^4,x^3y,y^2,xz,yz,z^2);
\)}\par\medskip
다음은 \emph{상대소-기약아이디얼들}이다.
\par\smallskip
\noindent\makebox[\linewidth][c]{\(\displaystyle
 \ideal R_1=\ideal Q_1;\qquad
 \ideal R_2=\ideal Q_2;\qquad
 \ideal R_3=[\ideal Q_3,\ideal Q_4]
 =(x^4,x^3y,x^2y^2,xy^3,xz,yz,z^2),
\)}\par\medskip
다음은 \emph{서로소-기약아이디얼들}이다.
\par\smallskip
\noindent\makebox[\linewidth][c]{\(\displaystyle
 \ideal S_1=\ideal R_1;\qquad
 \ideal S_2=[\ideal R_2,\ideal R_3]
\)}\par
\noindent\makebox[\linewidth][c]{\(\displaystyle
 =((y^2-1)x^4,(y^2-1)x^3y,(y^2-1)x^2y^2,(y^2-1)xy^3,xz,yz,z^2).
\)}\par\medskip
첫째 예에서와 마찬가지로
\par\smallskip
\noindent\makebox[\linewidth][c]{\(\displaystyle
 [\ideal B_4,\ideal B_5]=[\ideal D_4,\ideal D_5]
\)}\par
이고, 여기서
\noindent\(\displaystyle
 \ideal D_4=(x^4,xy+\lambda x^2,\ldots);\qquad
 \ideal D_5=(x^2+\mu xy,\ldots),\qquad \lambda\cdot\mu\ne1.
\)\par
또한
\noindent\(\displaystyle [\ideal Q_3,\ideal Q_4]=[\ideal Q_3,\overline{\ideal Q}_4],\)\quad
여기서 \(\displaystyle \overline{\ideal Q}_4=(x^4,x^3y,y^2+\lambda xy,\ldots).\)\par
나머지 기약아이디얼 또는 극대 으뜸아이디얼 $\ideal B_1,\ideal B_2,\ideal B_3$은 고립아이디얼이므로 유일하게 정해진다.

\begin{center}
\S~11.\\[0.25em]
\textbf{수론과 미분식 이론의 예.}
\end{center}

1. 바탕 환영역 $\Sigma$가 모든 \emph{짝수 유리정수}로 이루어진다고 하자. 그러면 $\Sigma$의 각 수 $2a$에 수 $a$를 대응시킴으로써 $\Sigma$를 모든 유리정수의 환영역과 일대일로 대응시킬 수 있다. 여기서 곧바로 $\Sigma$의 모든 아이디얼이 주아이디얼 $(2a)$임이 따른다. 아이디얼의 각 원소 $2c$에 대한 기저표현 $2c=n\cdot2a$에서 홀수 $n$은 유한합을 줄여 쓰는 기호일 뿐이다.

이 환영역의 \emph{소아이디얼들}은 $\ideal P_0=\ideal O=(2)$와 $\ideal P=(2p)$로 주어지며, 여기서 $p$는 홀수 소수이다. 따라서 각 소아이디얼은 $\ideal P_0$로 나누어지지만 다른 어느 소아이디얼로도 나누어지지 않는다. \emph{으뜸아이디얼들}은
\[
 \ideal Q_{\varrho_0}=(2\cdot2^{\varrho_0})
 \qquad\text{및}\qquad
 \ideal Q_\varrho=(2p^\varrho)
\]
로 주어진다. 이들은 동시에 \emph{기약아이디얼들}이다. 소아이디얼들에 관하여 방금 말한 사실에 따라 서로 다른 홀수 소수에 속하는 어느 두 아이디얼도 서로 상대소이지만, 어느 $\ideal Q$도 어떤 $\ideal Q_{\varrho_0}$에 대해서는 상대소가 아니다.\footnote{실제로 홀수 $p_1\ne p_2$에 대하여 $2b\cdot2p_1^{\varrho_1}\eqzero{2p_2^{\varrho_2}}$이면 언제나 $2b\eqzero{2p_2^{\varrho_2}}$이다. 그러나 $2b\cdot2p^\varrho\eqzero{2\cdot2^{\varrho_0}}$에서는 $2b\equiv2\cdot2^{\varrho_0-1}\,(2\cdot2^{\varrho_0})$만이 따른다.} 따라서 $\ideal Q_{\varrho_0}$들이 유일한 비고립 으뜸아이디얼들이다. $a$를 소수 거듭제곱들로 유일하게 분해하는 것에는 아이디얼 $(2a)$을 극대 으뜸아이디얼이면서 동시에 기약아이디얼인 것들로 나타내는 \emph{유일한} 표현
\[
 (2a)=[(2\cdot2^{\varrho_0}),(2p_1)^{\varrho_1},\ldots,(2p_\alpha)^{\varrho_\alpha}]
\]
이 대응한다. 따라서 다항식 환영역의 예들과 달리 여기서는 \emph{비고립} 극대 으뜸아이디얼들도 \emph{유일하게} 정해진다. $\varrho_0=0$이면 이 표현은 동시에 서로 상대소인 아이디얼들로 된 표현이고, $\varrho_0>0$이면 그 아이디얼은 상대소-기약이다.

그러므로 \emph{모든 정수}의 환영역에서는 서로 다른 네 분해가 일치하지만, 여기서는 한편으로 극대 으뜸아이디얼 분해와 기약아이디얼 분해만이 일치한다. 다른 한편 $\varrho_0>0$일 때에는 서로소-기약 분해(환영역이 단위원을 갖지 않으므로 모든 아이디얼이 서로소-기약이다)와 상대소-기약 분해가 일치하지만, $\varrho_0=0$일 때에는 서로소-기약 분해와 상대소-기약 분해가 서로 달라진다. 동시에 여기서는 한 소아이디얼이 다른 소아이디얼로 나누어지면서도 그와 같지 않을 수 있다는 예, 더 일반적으로 \emph{나눗셈으로부터 곱표현이 따르지 않는다는} 예를 이미 얻는다. 후자의 사실은 환영역에 단위원이 없다는 데서 생기며, 모든 아이디얼이 주아이디얼임에도 환영역의 수들을 환영역의 기약수들의 곱으로 유일하게 나타낼 수 없는 이유이기도 하다. 따라서 여기서는 최소공배수를 도입하는 일이 필요하다. 모든 짝수 대신 \emph{고정된 소수 또는 소수 거듭제곱으로 나누어지는 모든 수}를 바탕으로 삼아도 사정은 정확히 \emph{똑같다}는 점을 덧붙인다.

반면 $\Sigma$가 \emph{합성수} $g=\mathfrak p_1^{\sigma_1}\ldots\mathfrak p_\nu^{\sigma_\nu}$로 나누어지는 모든 수로 이루어지면 \emph{기약아이디얼과 으뜸아이디얼도 서로 달라진다.} 다시 모든 아이디얼은 주아이디얼 $(g\cdot a)$이다. \emph{소아이디얼들}은 다시 $\ideal P_0=\ideal O=(g)$와 $\ideal P=(g\cdot p)$로 주어지는데, 여기서 $p$는 $g$를 나누는 소수들과 다른 소수이다. 반면 \emph{기약아이디얼들}은 $\ideal B_{\lambda_i}=(g\cdot p_i^{\lambda_i})$와 $\ideal B_e=(g\cdot p^e)$로 주어진다. \emph{으뜸아이디얼들}은 $\ideal Q_e=\ideal B_e$와, 기약아이디얼들과는 다른
\[
 \ideal Q_{\lambda_1\ldots\lambda_\nu}
 =(g\cdot p_1^{\lambda_1}\cdots p_\nu^{\lambda_\nu})
\]
로 주어진다. 여기서 $\ideal B_{\lambda_i}$들과 $\ideal Q_{\lambda_1\ldots\lambda_\nu}$는 모두 같은 대응 소아이디얼 $\ideal P_0=(g)$를 가진다. 여기서도 기약아이디얼로의 분해는 유일하고, 따라서 극대 으뜸아이디얼로의 분해도 유일하다. 그러므로 여기서도 비고립아이디얼들까지 유일하게 정해진다.

2. \emph{비가환 환영역}의 한 예는 뇌터--슈마이들러의 논문에서 다룬 비가환 다항식 환영역의 아이디얼론이다. 특히 그 논문은 ``완전가약'' 아이디얼, 즉 분해의 성분들이 쌍마다 서로소이고 진약수를 갖지 않는 아이디얼을 다룬다. 따라서 성분들은 더더욱 기약이다. 그러므로 거기서 증명한 동형성에 더하여 \S~9로부터 서로 다른 두 분해에서 \emph{성분 수가 같다}는 결과를 얻는다. 이로써 그 논문의 특수한 경우로 나오는 편미분 또는 상미분 선형미분식들의 연립계 분해에 관한 결과를 얻는데, 이 결과는 상미분 선형미분식 하나라는 잘 알려진 경우에도 지적된 적이 없는 듯하다.

동시에 고정된 아이디얼 $\ideal M$의 모든 잉여류로 이루어진 체계 $T$는 비가환 다항식 환영역 $\Sigma$와 함께 이중환영역 $(\Sigma,T)$를 이루며, $T$는 $\Sigma$에 관한 모듈 성질을 가진다. 실제로 두 잉여류의 차는 다시 잉여류이고, 한 잉여류와 임의의 다항식의 곱도 잉여류이지만 두 잉여류의 곱은 존재하지 않는다(앞의 논문 \S~3). 그곳에서 ``부분군''이라고 부른 잉여류의 체계들은 따라서 바탕 환영역 $\Sigma$가 비가환인 이중환영역 $(\Sigma,T)$의 모듈 예를 이룬다.

\begin{center}
\S~12.\\[0.25em]
\textbf{초등인자론의 예.}
\end{center}

여기서는 일반 전개에서 도출되는 초등인자론의 관점을 제시하지만, 초등인자론 자체는 \emph{알려져 있다}고 가정한다.

$\Sigma$를 $n^2$개의 원소를 가진 모든 정수행렬의 환영역이라 하고, 덧셈과 곱셈은 행렬에 통상적인 뜻으로 정의하자. 그러면 $\Sigma$는 \emph{비가환 환영역}이다. 따라서 아이디얼들은 일반적으로 한쪽이고, 양쪽아이디얼들은 그 특수한 경우로만 포함된다.\footnote{이 환영역의 아이디얼론은 뒤 파스키에의 다음 논문들에서 다룬다. ``테타리온의 수론'', 취리히 학위논문, \emph{Vierteljahrsschr. d. Naturf. Ges. Zürich} 51 (1906); ``테타리온아이디얼의 이론에 관하여'', 같은 곳 52 (1907). 둘째 논문의 내용은 모든 아이디얼이 주아이디얼임을 증명하는 것이다.}

먼저 \emph{모든 아이디얼이 주아이디얼임}을 보이겠다. 이를 위하여 우측아이디얼의 경우 각 행렬 $A=(a_{ik})$에 정수계수 선형형식들의 모듈
\[
 \mathsf{A}=(a_{11}\xi_1+\cdots+a_{1n}\xi_n,\ldots,
 a_{n1}\xi_1+\cdots+a_{nn}\xi_n)
\]
을 대응시킨다. 거꾸로 이 모듈에는 $\mathsf{A}$의 기저를 주는 모든 행렬, 곧 $A$뿐 아니라 $U$가 유니모듈러인 $UA$도 대응한다. 더 일반적으로 곱 $PA$에는 $A$의 배수인 모듈 $\mathsf{B}$가 대응한다. $\mathsf{A}$의 한 선형형식은 0이 아닌 행을 하나만 가진 $P$들로 주어진다.

이제 $A_1,A_2,\ldots,A_\nu,\ldots$가 아이디얼 $\ideal M$의 모든 원소이고, 같은 기호 $\mathsf{A}_1,\mathsf{A}_2,\ldots,\mathsf{A}_\nu,\ldots$가 그에 대응하는 모듈들이라고 하자. 이 모듈들의 최대공약을 $\mathsf{A}$라 하고, 이 모듈 $A$에 대응하는 가장 일반적인 행렬을 $UA$라 하자. 최대공약의 정의에 따라 $\mathsf{A}$의 각 선형형식에는 행렬 $P_1A_{i_1}+\cdots+P_\sigma A_{i_\sigma}$가 대응하며, 위에서와 같이 각 $P$는 0이 아닌 행을 하나만 가진다. 따라서 $\mathsf{A}$의 한 기저에 대응하는 행렬 $A$도 이제 일반적인 $P$를 사용한 그러한 표현을 가지며, 그 결과 $\ideal M$의 원소가 된다. 더 나아가 각 모듈 $\mathsf{A}_i$는 $\mathsf{A}$로 나누어지므로 각 행렬 $A_i$도 $A$로 나누어진다. \emph{따라서 $A$, 더 일반적으로 $UA$가 $\ideal M$의 기저를 이룬다.} 좌측아이디얼의 경우에는 그에 맞추어 각 행렬의 열들을 한 모듈의 기저로 보아야 한다. 모든 아이디얼은 주아이디얼이고, $A$와 함께 $AV$도 그 기저가 되며, 여기서 $V$는 임의의 유니모듈러 행렬이다.

이제 초등인자론과 연결하기 위하여 이하에서는 \emph{양쪽아이디얼}을 바탕으로 삼는다. 따라서 특히 $A$와 함께 $PAQ$도 아이디얼에 속한다. 그러면 위의 결과에 따라 그러한 아이디얼의 가장 일반적인 기저는 $UAV$로 주어지며, 여기서 $U,V$는 유니모듈러이다. 따라서 기저원소들은 \emph{서로 동치인 행렬들의 한 류}를 모두 이룬다.\footnote{한쪽아이디얼의 기저원소들에는 각각 우측류 또는 좌측류가 대응한다.} 아이디얼과 류 사이에는 일대일 관계가 있고, 따라서 아이디얼과 그 류의 초등인자계 $(a_1\mid a_2\mid\cdots\mid a_n)$ 사이에도 일대일 관계가 있다. 여기서 $a_i$들은 잘 알려진 대로 음이 아닌 정수이며 각 수는 그 다음 수를 나눈다. 초등인자들이 정하는 표준형에 나타나는 그 류의 행렬을 따라서 아이디얼의 특별한 기저로 볼 수 있다. 아이디얼 또는 류의 나눗셈 관계에서 초등인자들의 나눗셈 관계가 따르고, 그 역도 성립한다.

그런데 뒤 파스키에는 앞의 논문 \S~11에서 모든 양쪽아이디얼의 랭크가 $n$이고 모든 초등인자가 일치함을 보였다. 따라서 일반적인 초등인자를 다루려면 아이디얼이 아니라 양쪽류 자체에서 직접 출발해야 한다(이 류들도 큰 독일문자로 나타낸다). 이때 한 류 $\mathfrak A=UAV$는 $A=PBQ$가 될 때 다른 류 $\mathfrak B$로 나누어진다.

일반적으로 다음이 성립한다. \emph{두 류의 최소공배수(최대공약)는 대응하는 초등인자계들의 최소공배수(최대공약)를 취하여 얻는다.}\footnote{E. 슈타이니츠는 ``모듈론에 관하여'', \emph{Math. Ann.} 52 (1899), 1쪽에서 초등인자계들의 최소공배수와 최대공약으로 류의 최소공배수와 최대공약을 정의한다. 이와 독립적으로 류의 최소공배수는 H. 브란트, ``한 기본형식에 관한 이차 이항형식의 합성'', \emph{J. f. M.} 150 (1919), 1쪽에 ``합동합성''이라는 이름으로 나타난다.} 실제로 $(a_1\mid a_2\mid\cdots\mid a_n)$, $(b_1\mid b_2\mid\cdots\mid b_n)$, $(c_1\mid c_2\mid\cdots\mid c_n)$가 각각 $\mathfrak A,\mathfrak B,\mathfrak C^*$의 초등인자계이고 $c_i=[a_i,b_i]$라고 하자. 또한 $\mathfrak C=[\mathfrak A,\mathfrak B]$라 하자. 그러면 $\mathfrak C^*$는 $\mathfrak A$와 $\mathfrak B$로, 따라서 $\mathfrak C$로 나누어진다. 거꾸로 $\mathfrak C$의 초등인자들은 $\mathfrak C^*$의 초등인자들로 나누어지므로 $\mathfrak C$는 $\mathfrak C^*$로 나누어진다. 따라서 $\mathfrak C=\mathfrak C^*$이다. 최대공약의 증명도 그에 맞추어 얻어진다.

초등인자 $a_i$를 소수 거듭제곱들의 최소공배수로 유일하게 나타내는 것에는 $\mathfrak A$를 류 $\mathfrak Q$들의 최소공배수로 나타내는 표현이 대응한다. 이 $\mathfrak Q$들의 초등인자는 한 소수의 거듭제곱들로 주어지며, 기호로는
\[
 \mathfrak Q\sim(p^{r_1}\mid p^{r_2}\mid\cdots\mid p^{r_\varrho}\mid0\mid\cdots\mid0);
 \qquad r_1\leqq r_2\leqq\cdots\leqq r_\varrho
\]
이다. 특히 랭크 $\varrho$가 $n$과 같으면 이는 서로소인 서로소-기약류들로의 분해이며, 비가환 환영역임에도 이 분해는 유일하다.

류 $\mathfrak Q$는 다음 초등인자계들에 대응하는 류들로 더 분해할 수 있다.
\[
\begin{gathered}
 \mathfrak B_1\sim(p^{r_1}\mid\cdots\mid p^{r_1});\quad
 \mathfrak B_2\sim(1\mid p^{r_2}\mid\cdots\mid p^{r_2});\quad \ldots;\\
 \mathfrak B_\varrho\sim(1\mid\cdots\mid1\mid p^{r_\varrho}\mid\cdots\mid p^{r_\varrho});\quad
 \mathfrak B_{\varrho+1}\sim(1\mid\cdots\mid1\mid0\mid\cdots\mid0).
\end{gathered}
\]
여기서 $\mathfrak B_\nu$에는 매번 수 $1$이 $(\nu-1)$번 나타난다. 이를테면 $r_1=r_2=\cdots=r_\mu=0$이면 $\mathfrak B_1,\ldots,\mathfrak B_\mu$는 단위류와 같아지므로 분해에서 생략해야 한다. $\varrho=n$이면 $\mathfrak B_{\varrho+1}$에도 같은 말이 성립한다. 더 나아가 이를테면 $r_\nu=r_{\nu+1}=\cdots=r_{\nu+\lambda}$이면 $\mathfrak B_{\nu+1},\ldots \mathfrak B_{\nu+\lambda}$는 $\mathfrak B_\nu$의 진약수가 되므로 역시 제외해야 한다. 남은 류들을 $\mathfrak B_{i_1},\ldots,\mathfrak B_{i_k}$라 하면 이들은 이제 최단표현을 주고,
\[
 \mathfrak Q=[\mathfrak B_{i_1}\ldots\mathfrak B_{i_k}]
\]
는 \emph{$\mathfrak Q$를 기약류들로 나타내는 유일한 분해}이다.

실제로
\[
 \mathfrak B_\nu\sim(1\mid\cdots\mid1\mid p^{r_\nu}\mid\cdots\mid p^{r_\nu})
\]
를
\[
 \mathfrak C\sim(1\mid\cdots\mid1\mid p^{s_1}\mid\cdots\mid p^{s_\lambda})
 \quad\text{및}\quad
 \mathfrak D\sim(1\mid\cdots\mid1\mid p^{t_1}\mid\cdots\mid p^{t_\mu})
\]
의 최소공배수로 나타낼 수 있다고 하자. 그러면 $\mathfrak C$와 $\mathfrak D$의 초등인자계에서 처음 $(\nu-1)$자리에는 수 $1$이 놓여야 한다. $\nu$번째 자리에서는 지수 $s_1$ 또는 $t_1$, 이를테면 $s_1$이 $r_\nu$와 같아야 한다. 그런데
\[
 r_\nu=s_1\leqq s_2\leqq\cdots\leqq s_\lambda\leqq r_\nu
\]
이므로 $\mathfrak C=\mathfrak B_\nu$이다. 따라서 $\mathfrak B_\nu$는 \emph{기약}이다.\footnote{반면 $\mathfrak B$들은 한쪽류들로는 더 분해될 수 있다. 여기서는 초등인자와 류 사이의 관계가 더 이상 일대일이 아니고, 따라서 기약 한쪽류들로의 분해도 유일하지 않다. H. 브란트가 알려 준 다음 우측류 분해의 예가 이를 보여 준다. 여기서 류는 한 기저행렬 또는 그에 대응하는 모듈로 나타낸다.
\[
\begin{gathered}
 \mathfrak B=[\mathfrak C_1,\mathfrak C_2]=[\mathfrak D_1,\mathfrak D_2];\quad
 \mathfrak B\sim\begin{pmatrix}p&0\\0&p\end{pmatrix};\quad
 \mathfrak C_1\sim\begin{pmatrix}1&0\\0&p\end{pmatrix};\quad
 \mathfrak C_2\sim\begin{pmatrix}p&0\\0&1\end{pmatrix};\\
 \mathfrak D_1\sim\begin{pmatrix}1&0\\0&p\end{pmatrix};
 \quad
 \mathfrak D_2\sim\begin{pmatrix}p&0\\(p-1)&1\end{pmatrix}.
\end{gathered}
\]
실제로 모듈 $(\xi,p\eta)$, $(p\xi,\eta)$, $(p\xi,(p-1)\xi+\eta)$는 각각 진약수로 모듈 $(\xi,\eta)$만을 가지므로 기약이며, 더 나아가 서로 다르다.} $p$를 $0$으로 바꾸면 $\mathfrak B_{\varrho+1}$에도 같은 사실이 성립한다. 그런데 이 기약류들은 각각 최소공배수를 만들 때 하나의 특정한 지수와 그 지수가 $\mathfrak Q$ 또는 $\mathfrak A$의 초등인자계에서 처음 나타나는 자리를 주며, $\mathfrak B_{\varrho+1}$은 랭크를 준다. 이 수들은 $\mathfrak Q$ 또는 $\mathfrak A$의 초등인자계에 의해 유일하게 정해지고 초등인자와 류 사이의 관계는 일대일이므로, $\mathfrak Q$의 기약류 분해, 마찬가지로 임의의 류의 기약류 분해는 \emph{유일하다}.

정리하면 다음과 같다. \emph{원소 수가 제한된 정수행렬들의 모든 양쪽류 $\mathfrak A$는 유한개의 기약 양쪽류들의 최소공배수로 유일하게 나타낼 수 있다. 각 기약류는 $\mathfrak A$의 초등인자계의 고정된 소인수 하나, 그에 딸린 지수 하나, 그리고 그 지수가 처음 나타나는 자리를 나타낸다. 인수 $0$에 대응하는 기약류는 $\mathfrak A$의 랭크를 나타낸다.}

\begin{flushleft}
에를랑겐, 1920년 10월.
\end{flushleft}
\begin{center}
(1920년 10월 16일 접수.)
\end{center}

\clearpage
\setcounter{footnote}{0}

\clearpage
\setcounter{footnote}{0}
\phantomsection
\addcontentsline{toc}{section}{20. 절대기약성의 한 대수적 판정기준}
%
%
%
%
%
%

\section*{20. 절대기약성의 한 대수적 판정기준}
\begin{center}
에미 뇌터, 괴팅겐.\\[0.6em]
Math. Ann. 85 (1922), 26--33쪽
\end{center}

임의의 추상적으로 정의된 체에 계수를 둔 다항식은, 그 계수영역을 확대하여 얻을 수 있는 대수적으로 닫힌 체에서도 기약으로 남을 때 \emph{절대기약}이라고 한다\srcfnmark{1)}.
\srcfntext{1)}{모든 체는 본질적으로 유일한 방식으로 대수적으로 닫힌 체, 곧 한 변수의 모든 다항식이 일차인자들로 분해되는 체로 확대될 수 있음을 E. Steinitz가 \emph{Algebraische Theorie der Körper}(J. f. M. 137 (1910), 167쪽)에서 보였다. 이로써 그는 대수학의 기본정리에 대한 유리적 등가물을 제시하였다.}
\addtocounter{footnote}{1}
이하에서는 주어진 체에 대한 기약성과 달리 절대기약성은 \emph{대수적으로} 포착할 수 있음을 보이겠다. 더 정확히는 다음 정리가 성립한다.

\emph{부정계수를 가진 \(n\geqq2\)개 변수의 각 다항식에는 그 계수들과 추가 부정원들의 정수계수 다항함수 하나, 곧 가약성형식을 대응시킬 수 있다. 같은 차수의 각 특수 다항식이 절대기약이기 위한 필요충분조건은 그 가약성형식이 사라지지 않는 것이다.} 다시 말해, \emph{가약성형식이 부정원들에 관하여 항등적으로 사라지고 차수저하를 일으키지 않는 계수의 각 특수 값의 계에 대하여 특수 다항식은 가약이며, 역도 성립한다.}

다항식 대신 변수를 하나 더 가진 동차형식을 생각하면, 차수조건은 계수계가 항등적으로 사라지지 않는다는 더 간단한 조건으로 바뀐다.

이 정리에서는 A. Ostrowski가 처음 세운 다음 명제도 직접 따라 나온다.\footnote{\emph{Zur arithmetischen Theorie der algebraischen Größen}, Gött. Nachr. 1919, 279쪽(296쪽의 보조정리). 통상적인 수로 이루어진 체의 경우에는 Ostrowski도 위의 주정리에 도달한 것으로 알고 있으며, 그 연구를 곧 출판할 예정이다. 나의 증명 배열에서 주정리가 임의의 추상적으로 정의된 체에 대하여 성립하는 것은 부정원들에 대하여 만든 가약성형식이 \emph{특별히 바탕에 둔 체와 무관한 수계수}를 갖기 때문이다.}

\emph{대수적 수들을 계수로 가진 모든 절대기약 다항식은, 계수영역을 포함하는 유한 대수적 수체 \(\Omega\)의 소아이디얼들 가운데 많아야 유한개를 제외하면 각 소아이디얼을 법으로 하여 기약으로 남는다. 특히 \(\Omega\)를 유리수체로 택할 수도 있다.}

상대적으로 정수적인 함수의 이론에 대한 더 많은 응용\footnote{D. Hilbert, \emph{Mathematische Probleme}, Gött. Nachr. 1900, 253쪽, 문제 14.}은 다른 곳에서 다룰 예정이다.

\textbf{1.} 우선 \emph{부정계수를 가진 \(n\geqq2\)개 변수의 모든 다항식은 절대기약}임을 보이자. 다음과 같이 두자.
\srcnumdisplay{(1)}{%
\begin{aligned}
F(x)&=\sum A_{i_1\ldots i_n}x_1^{i_1}\ldots x_n^{i_n}
&&(i_1+\ldots+i_n\leqq l),\\
G(x)&=\sum B_{j_1\ldots j_n}x_1^{j_1}\ldots x_n^{j_n}
&&(j_1+\ldots+j_n\leqq m),\\
H(x)&=F(x)G(x)=\sum C_{k_1\ldots k_n}(A,B)x_1^{k_1}\ldots x_n^{k_n}
&&(k_1+\ldots+k_n\leqq l+m).
\end{aligned}
}
여기서 \(A\)와 \(B\)는 부정원이고 \(C(A,B)\)는 \(A,B\)의 정수계수 쌍선형 결합이다. 따라서 몫
\[
D=C(A,B)/C_{0\ldots0}(A,B)
\]
들은 몫 \(A/A_{0\ldots0}\)와 \(B/B_{0\ldots0}\)의 유리함수이다. 그런데 이 함수들의 개수는 인수들의 개수보다 크다. 각각
\[
\binom{l+m+n}{n}-1,\qquad
\binom{l+n}{n}-1,\qquad
\binom{m+n}{n}-1
\]
개이며
\[
\binom{l+m+n}{n}+1>
\binom{l+n}{n}+\binom{m+n}{n}
\quad(n\geqq2,\ l>0,\ m>0)
\]
이기 때문이다. 그러므로 \(D(A,B)\)들 사이에는 적어도 \emph{하나의} 대수적 관계가 있고, 따라서 \(C(A,B)\)들 사이에도 그러한 관계가 있다.
\srcnumdisplay{(2)}{%
\varPhi(Z)\ne0\quad[\text{에 관하여 항등적으로}Z_{k_1\ldots k_n}],\
\varPhi(C(A,B))=0\quad[\text{에 관하여 항등적으로}A,B],
}
여기서 \(\varPhi\)는 정수계수 다항식이다.

따라서 \emph{부정원 \(Z\)}를 계수로 가진 다항식
\srcnumdisplay{(3)}{%
E(x)=\sum Z_{k_1\ldots k_n}x_1^{k_1}\ldots x_n^{k_n}
\qquad(k_1+\ldots+k_n\leqq l+m)
}
은 \(n\geqq2\)일 때 반드시 \emph{절대기약}이다\srcfnmark{4)}.
\srcfntext{4)}{2절의 표현을 쓰면, 부정원 \(Z\)가 \(\mathfrak P\)의 영점이 아니기 때문이다.}
\addtocounter{footnote}{1}

\textbf{2.} \(Z=C(A,B)\)를 대입했을 때 \(A,B\)에 관하여 항등적으로 사라지는 이 모든 다항식 \(\varPhi(Z)\)의 전체는 \emph{다항식 아이디얼}\srcfnmark{5)}을 이루며, 더 정확히는 \emph{소아이디얼 \(\mathfrak P\)}을 이룬다.
\srcfntext{5)}{이러한 전체는 보통 가군이라 부른다. 그러나 이를 특징짓는 성질은 \(\varPhi_1,\varPhi_2\)와 함께 \(\varPhi_1\mp\varPhi_2\)도 포함하고, \(\varPhi\)와 함께 \(\Psi\varPhi\)도 포함한다는 것이다. 여기서 \(\Psi\)는 \(Z\)의 임의의 정수계수 다항식이다. 이는 바로 아이디얼의 성질이다.}
\addtocounter{footnote}{1}
실제로 대입에 따라 곱 \(\varPhi_1\varPhi_2\)가 사라지면 적어도 한 인자가 사라진다. (2)는 \(\mathsf{A},\mathsf{B}\)에 관하여 항등적으로 성립하므로, \(\mathsf{A},\mathsf{B}\), 따라서 \(\mathsf{\Gamma}=C(\mathsf{A},\mathsf{B})\)가 임의의 체의 원소가 되는 각 특수 값의 계에서도 그대로 성립한다. \emph{그러므로 어떤 확대체에서 가약인 모든 다항식의 계수 \(\mathsf{\Gamma}\)는 조건 \(\varPhi(\mathsf{\Gamma})=0\)을 만족하여 \(\mathfrak P\)의 영점이 되며, 따라서 \(\mathfrak P\)의 유한개 기저다항식의 영점이 된다.} 이제 \emph{역을 증명}해야 한다.

\(y\)라는 새 변수를 도입하여 (1)의 다항식들을 차수 \(l,m,l+m\)의 동차형식 \(F(x,y),G(x,y),H(x,y)\)로 바꾸어도 \(A,B,C\) 사이의 모든 관계는 보존된다. 따라서 예를 들어 \(\bar F(x,y)\)\footnote{\(F,G,\ldots,f,g,\ldots\)의 부정원들을 특수한 값의 계로 바꾸어 얻는 다항식과 형식은 이하에서 언제나 \(\bar F,\bar G,\ldots,\bar f,\bar g,\ldots\)처럼 윗줄을 그어 표시한다.}가 \(y\)의 한 거듭제곱이 되는 계수 \(\Gamma\)도 \(\mathfrak P\)의 영점이다. 이때 비동차 다항식으로 넘어가면 \(\bar H(x)\)의 차수저하만 생기고 인수분해는 생기지 않는다. 이 차수저하를 제외하면 역은 언제나 성립함을 보이겠다. 특히 \emph{동차형식에 대해서는 모든 \(\mathsf{\Gamma}\)가 사라지는 경우만 제외하면 역이 예외 없이 성립한다. 다시 말해 \(\mathfrak P\)의 각 영점은 어떤 확대체의 원소 \(A,B\)를 써서 매개표현 \(\mathsf{\Gamma}=C(\mathsf{A},\mathsf{B})\)으로 주어진다}\srcfnmark{7)}.
\srcfntext{7)}{\(\mathfrak P\)가 기약 대수다양체를 정의하는 소아이디얼이라는 성질에서, \(\mathfrak P\)의 영점들이 \glqq{}일반적으로\grqq{} 매개표현으로 주어진다는 사실은 따라온다. 그러나 이로부터 매개표현이 어느 특수 값의 계에서도 실패하지 않는다는 사실은 따라오지 않는다. 나는 처음에 이 점을 놓쳤고 A. Ostrowski가 오래전에 이를 지적해 주었다. 여기서 제시하는 증명은 더 초등적인 개념들에 바탕을 둔다.}
\addtocounter{footnote}{1}

이를 증명하기 위하여 유한개의 함수 \(\varPhi(Z)\)를 직접 구성한다. 즉 Kronecker 치환
\[
x_\lambda=\xi^{d^{\,n-\lambda}}\srcfnmark{8)}
\]
을 이용하여 (1)의 다항식들을 \emph{한} 변수의 다항식들에 대응시킨다. 그 지수들에 결락이 생김을 보이고, 대응하는 계수의 노름을 만들어 함수 \(\varPhi(Z)\)와 찾는 판정기준을 얻는다.
\srcfntext{8)}{Festschrift, 11쪽.}
\addtocounter{footnote}{1}

\textbf{3.} 다음과 같이 두자.
\srcnumdisplay{(4)}{%
\begin{aligned}
d&=l+m+1; &\qquad i&=i_1d^{n-1}+i_2d^{n-2}+\ldots+i_n;\\
j&=j_1d^{n-1}+\ldots+j_n; &\qquad k&=k_1d^{n-1}+\ldots+k_n.
\end{aligned}
}
Kronecker 치환에 따라 (1)의 \(F,G,H\)에는 각각 다음 다항식들이 대응한다.
\srcnumdisplay{(5)}{%
\begin{aligned}
f(\xi)&=\sum A_{i_1\ldots i_n}\xi^i
&&(i_1+\ldots+i_n\leqq l),\\
g(\xi)&=\sum B_{j_1\ldots j_n}\xi^j
&&(j_1+\ldots+j_n\leqq m),\\
h(\xi)&=\sum C_{k_1\ldots k_n}(A,B)\xi^k
&&(k_1+\ldots+k_n\leqq l+m).
\end{aligned}
}
이 대응은 잘 알려져 있듯이 \emph{일대일}이다. (5)에 나타나는 모든 \emph{유도지수}에 대하여 \(k=0\)이면 \(k_1=0,\ldots,k_n=0\)이기 때문이다. 따라서 \(i+j=i'+j'\)이면
\[
i_1+j_1=i'_1+j'_1;\ldots;i_n+j_n=i'_n+j'_n
\]
이다. 그러므로 서로 다른 곱 \(\xi^i\xi^j\)가 같은 지수 \(\xi^k\)를 주는 것은, 대응하는
\[
x_1^{i_1}\ldots x_n^{i_n}x_1^{j_1}\ldots x_n^{j_n}
\]
도 같은 지수를 줄 때뿐이다. 따라서
\srcnumdisplay{(6)}{%
h(\xi)=f(\xi)g(\xi)
}
를 얻으며, \(f(\xi)\)와 \(g(\xi)\)에 유도지수 이외의 지수가 나타나지 않는 관계 (6)가 성립하면 거꾸로
\srcnumdisplay{(7)}{%
H(x)=F(x)G(x)
}
를 얻는다. 이때 \emph{\(f(\xi)\)와 \(g(\xi)\)의 유도지수에는 실제로 결락이 생긴다.} \(f(\xi)\)의 최고지수는 \(l\cdot d^{n-1}\), 그다음 지수는 \((l-1)d^{n-1}+d^{n-2}\)이고, 그 차는
\[
d^{n-2\cdot}(d-1)>1
\]
이다. 실제로 \(d=l+m+1\geqq3\), \(n\geqq2\)이다. \(g(\xi)\)에 대해서도 같다.

부정원 \(\mathsf{A},\mathsf{B}\)에 대하여 성립하는 관계 (6)과 (7)은 각 특수 값의 계에서도 보존된다. 이때 차수가 보존되도록 변수 \(\eta\)와 \(y\)를 써서 (6)과 (7)을 동차화한다. \emph{따라서 동차형식 \(\bar{H}(x,y)\)가 한 대수적 확대체에서 인수분해되는 것은, 그 체에서 대응하는 이원형식 \(\bar{h}(\xi,\eta)\)가 두 인자로 분해되고 그 인자들에 유도지수 이외의 지수가 나타나지 않을 때 그리고 그때뿐이다.}

\textbf{4.} 2절 끝에서 말한 노름의 구성을 수행하기 위하여
\[
a_1,b_1;\ldots;a_p,b_p;\quad
a_{p+1},b_{p+1};\ldots;a_{p+q},b_{p+q}
\]
를 새 부정원들이라 하자. 다음과 같이 둔다.
\srcnumdisplay{(8)}{%
\begin{aligned}
\varphi(\xi,\eta)&=(a_1\xi+b_1\eta)\ldots(a_p\xi+b_p\eta)
 =r_p\xi^p+r_{p-1}\xi^{p-1}\eta+\ldots+r_0\eta^p,\\
\psi(\xi,\eta)&=(a_{p+1}\xi+b_{p+1}\eta)\ldots(a_{p+q}\xi+b_{p+q}\eta)
 =s_q\xi^q+s_{q-1}\xi^{q-1}\eta+\ldots+s_0\eta^q,\\
\chi(\xi,\eta)&=\varphi(\xi,\eta)\psi(\xi,\eta)
 =t_{p+q}\xi^{p+q}+t_{p+q-1}\xi^{p+q-1}\eta+\ldots+t_0\eta^{p+q}.
\end{aligned}
}
여기서 \(r,s,t\)는 각각 동차화한 기본대칭함수이다. 즉 각 열에 관하여 선형인 \(a,b\) 열들의 대칭함수로서, 각 열에 관하여 따로 같은 차수로 동차인 모든 정수계수 다항 대칭함수는 이들로 정수계수 다항식으로, 그리고 오직 한 가지 방식으로 나타난다.

이제 추가 부정원 \(u_{\mu\nu}\)를 써서 선형형식
\srcnumdisplay{(9)}{%
\zeta(u,a,b)=\sum u_{\mu\nu}r_\mu s_\nu
}
를 만들되, 합은 미리 정한 지표 \(\mu,\nu\)에 걸쳐 취한다. \(\zeta(u)\)와 \(a,b\) 열의 치환으로 생기는, \(\zeta(u)\) 및 서로와 다른 모든 켤레들의 곱을 생각하자. 이는 \(\zeta(u)\)를 \(r_\mu(a,b),s_\nu(a,b)\)의 체의 함수로 볼 때 \(t(a,b)\)의 체에 대한 \(\zeta(u)\)의 노름 \(N(\zeta(u))\)이다. 이 곱은 모든 \(a,b\) 열의 정수계수 다항 대칭함수이며 각 열에 관하여 같은 차수로 동차이다. 그러므로 위에서 본 바에 따라 \(t(a,b)\)와 \(u_{\mu\nu}\)의 유일하게 정해지는 정수계수 다항함수가 된다.
\srcnumdisplay{(10)}{%
N(\zeta(u,a,b))=T(t(a,b),u)
\qquad[a,b,u\text{에 관하여 항등적으로}].
}
같은 추론으로
\[
N(z-\zeta(u))=
z^\delta+T_{\delta-1}(t,u)z^{\delta-1}+\ldots+T(t,u)
\qquad[a,b,u,z\text{에 관하여 항등적으로}]
\]
도 성립하며, 여기서 모든 \(T\)는 \(t,u\)의 정수계수 다항함수이다. 이 항등식의 양변은 \(z=\zeta(u)\)일 때 사라진다. 기본대칭함수 \(r(a,b),s(a,b)\)가 대수적으로 독립이므로, \(t(a,b)\)를 (8)에 따라 \(r,s\)의 정수계수 쌍선형 결합 \(w(r,s)\)로 두고 \(\zeta(u)\)를 \(r,s\)의 함수로 보면 이 소거는 \(r,s\)에 관하여도 항등적이다. 따라서
\srcnumdisplay{(11)}{%
\zeta(u)^\delta+
T_{\delta-1}(w(r,s),u)\zeta(u)^{\delta-1}
+\ldots+T(w(r,s),u)=0
\quad[r,s,u\text{에 관하여 항등적으로}].\,\text{\textsuperscript{9)}}
}
\srcfntext{9)}{(9)의 합을 모든 지표에 걸쳐 취하면 (11)은 가우스 정리에 대한 Kronecker의 확장이 되며, 증명의 배열도 본질적으로 Kronecker의 것이다. \emph{Zur Theorie der Formen höherer Stufen}, Berliner Ber. 1883, Werke 2, 417쪽.}
\addtocounter{footnote}{1}

항등식 (10)과 (11)은 \(a,b\)의 모든 특수 값의 계 \(\alpha,\beta\), 또는 \(r,s\)의 특수 값의 계 \(\varrho,\sigma\)에서도 유지된다. 특히 (9)에 나타나는 지표 \(\mu,\nu\)에 대하여 모든 곱 \(\varrho_\mu\sigma_\nu\)가 사라지도록 \(\varrho,\sigma\)를 택하면, 대입에 따라 \(\zeta(u)\)가 사라진다. \(w(\varrho,\sigma)=\tau\)로 두면 (11)에서
\srcnumdisplay{(12)}{%
T(\tau,u)=0\qquad[u\text{에 관하여 항등적으로}]
}
를 얻는다.

역으로 \(\tau_0,\ldots,\tau_{p+q}\)가 모두 사라지지는 않으면서 (12)를 만족하는 \emph{값의 계}라고 하자. 한 대수적 확대체에서는
\srcnumdisplay{(12)}{%
\bar\chi(\xi,\eta)=
\tau_{p+q}\xi^{p+q}+\ldots+\tau_0\eta^{p+q}
=(\alpha_1\xi+\beta_1\eta)\ldots
(\alpha_{p+q}\xi+\beta_{p+q}\eta)
}
라는 분해가 존재한다. 각 인자에서 \(\alpha\)와 \(\beta\)가 동시에 사라지지는 않지만\srcfnmark{10)}, 몇몇 \(\alpha\) 또는 \(\beta\)는 사라질 수 있다.
\srcfntext{10)}{이 유리적 \glqq{}대수학의 기본정리\grqq{}가 2절 끝에서 말한 역이 예외 없이 성립한다는 사실의 바탕이 되는 존재정리이다.}
\addtocounter{footnote}{1}
따라서 \(\tau=t(\alpha,\beta)\)이다. 이 값들을 (10)에 넣으면 (12)에 따라 \(\zeta(u,\alpha,\beta)\) 또는 그 켤레인자 하나가 \(u\)에 관하여 항등적으로 사라진다. \(\alpha,\beta\)의 번호를 알맞게 매긴 뒤 \(r(\alpha,\beta)=\varrho\), \(s(\alpha,\beta)=\sigma\)로 두면, 이는 \(\bar\chi(\xi,\eta)\)가 \emph{적어도 하나의 다음 분해를 허용한다}는 뜻이다.
\srcnumdisplay{(13)}{%
\bar\chi(\xi,\eta)=\bar\varphi(\xi,\eta)\bar\psi(\xi,\eta)
=\sum \varrho_\varkappa\xi^\varkappa\eta^{p-\varkappa}
\cdot\sum \sigma_\lambda\xi^\lambda\eta^{q-\lambda},
}
\emph{여기서 \(\zeta(u)\)에 나타나는 모든 곱 \(\varrho_\mu\sigma_\nu\)는 사라지지만, \(\varrho\) 전체나 \(\sigma\) 전체가 사라지지는 않는다.}

\textbf{5.} 이제 (8)의 차수 \(p,q\)를 각각 \(ld^{n-1}\), \(md^{n-1}\), 곧 (5)의 \(f(\xi),g(\xi)\)의 차수와 같게 둔다. 지표 \(\mu,\nu\)를 \(\mu^{(1)},\nu^{(1)},\mu^{(2)},\nu^{(2)}\)로 나눈다. 이때 \(\mu^{(1)}\)은 \(f(\xi)\)에서 빠진 모든 지수를, \(\nu^{(1)}\)은 \(\psi(\xi,\eta)\)에서 \(\xi\)의 모든 지수를 돈다. 대응하게 \(\nu^{(2)}\)는 \(g(\xi)\)에서 빠진 모든 지수를, \(\mu^{(2)}\)는 \(\varphi(\xi,\eta)\)에서 \(\xi\)의 모든 지수를 돈다. 또
\[
e(\xi)=\sum Z_{k_1\ldots k_n}\xi^k
\qquad(k_1+\ldots+k_n\leqq l+m)
\]
를 다항식 (3)에 대응하는 한 변수 다항식이라 하자. \(S(Z,u)\)는 \(T(t,u)\)에서 다음과 같이 얻는 \(Z,u\)의 정수계수 다항식이라 하자. 여기서 \(T(t,u)\)는 (10)의 유일하게 정해지는 우변이고, (9)의 합은 지정한 지표 \(\mu^{(1)},\nu^{(1)},\mu^{(2)},\nu^{(2)}\)에 걸쳐 취한다. 이제 \(t\)들을 \(e(\xi)\)의 계수 \(Z\)와 대응하는 0들로 바꾼다.

\(r,s\)를 \(f(\xi),g(\xi)\)의 계수 \(A,B\)와 대응하는 0들로 바꾸면, 위에서 지정한 합에서는 \(\zeta(u)\)가 \emph{사라진다}. 또 \(t=w(r,s)\)는 \(h(\xi)\)의 계수 \(C(A,B)\)와 대응하는 0들로 바뀌므로, \(T(t,u)\)는 \(S(C(A,B),u)\)로 바뀐다. 따라서 항등식 (11)에서
\srcnumdisplay{(14)}{%
S(C(A,B),u)=0
\qquad[A,B,u\text{에 관하여 항등적으로}]
}
을 얻는다. (14)는 각 특수 값의 계에서도 유지된다. 그러므로 한 \emph{확대체에서 두 인자로 분해되는} 특수 \(\bar H(x)\), 또는 일반적인 \(\bar H(x,y)\)의 계수 \(\mathsf{\Gamma}=C(\mathsf{A},\mathsf{B})\)는 조건
\srcnumdisplay{(15)}{%
S(\mathsf{\Gamma},u)=0\qquad[\text{에 관하여 항등적으로}u]
}
을 만족한다.

역으로 모두 사라지지는 않는 다항식 \(\bar H(x)\), 또는 형식 \(\bar H(x,y)\)의 계수 \(\Gamma\)가 조건 (15)를 만족한다고 하자. 위에서 본 바에 따라 이는 \(T(\tau,u)=0\)과 동치이며, 여기서 \(\tau\)는 \(\bar h(\xi,\eta)\)의 모든 계수를 뜻한다. (13)에 따라 \(\mathsf{\Gamma}\)의 한 확대체에는
\[
\bar h(\xi,\eta)=\bar f(\xi,\eta)\cdot\bar g(\xi,\eta)
\]
라는 분해가 존재한다. 지정한 합 때문에 \(\bar f,\bar g\)에는 유도지수 이외의 지수가 나타나지 않는다. 따라서
\[
\bar H(x,y)=\bar F(x,y)\cdot\bar G(x,y)
\]
라는 관계가 성립한다. \(y=1\)일 때 어느 인자도 0차가 되지 않으면, 특히 \(\mathsf{\Gamma}\)가 \(\bar H(x)\)의 차수저하를 일으키지 않으면
\[
\bar H(x)=\bar F(x)\cdot\bar G(x)
\]
도 성립한다.

\emph{그러므로 조건 (15)는 \(\bar H(x,y)\)가 한 확대체에서 두 인자로 분해되기 위한 필요충분조건이며, 차수저하를 제외하면 \(\bar H(x)\)에 대해서도 필요충분조건이다.}

1절에서 \(n\geqq2\)일 때 부정원들을 계수로 가진 다항식, 따라서 변수 하나를 더 둔 대응하는 동차형식이 절대기약임을 보였다. 그러므로 \(S(Z,u)\)는 \(Z,u\)에 관하여 항등적으로 사라질 수 없다. \(U\)가 \(u\)들의 거듭제곱곱들을 뜻하면
\[
S(Z,u)=\sum \varPhi_\lambda(Z)U_\lambda
\]
이고, (14)에 따라 \(\varPhi_\lambda(Z)\)들은 \(\mathfrak P\)의 다항식이 된다. \emph{이로써 \(\mathfrak P\)에는 \(A,B\)가 \(\mathsf{\Gamma}\)의 한 대수적 확대체에 속하는 매개표현 \(\mathsf{\Gamma}=C(\mathsf{A},\mathsf{B})\)으로 주어진 것 이외의 영점이 없음을 보였다.} 따라서 2절 끝에서 말한 역이 증명되었다.

\textbf{6.} 이제 \emph{절대기약성의 조건을 직접 정식화할 수 있다.} (3)의 \(E(x)\)의 차수를 가능한 모든 방식으로 두 양의 정수의 합 \(l+m\)으로 나누고, 각 분해에 대하여 형식 \(S(Z,u)\)를 만들면 된다. 이 형식들을 모두 곱한
\srcnumdisplay{(16)}{%
R(Z,u)=\prod S(Z,u)
}
이 \(E(x)\)의 \emph{가약성형식}이다. 실제로 위에서 본 바에 따라 \(\bar E(x)\), 또는 \(\bar E(x,y)\)가 \emph{분해}될 때마다 (16)의 한 인자가 사라지고, 역도 성립한다. 이로써 \emph{서론에서 세운 주정리가 증명되었고}, 동시에 가약성형식을 유한 단계로 실제 구성하는 방법도 보였다.

\textbf{7.} 가약성형식 (16)의 \emph{존재}에서 서론에서 말한 \emph{Ostrowski 정리}가 곧바로 따라 나온다. 대수적 정수 \(\Gamma\)를 계수로 가진 절대기약 다항식을
\[
\bar E(x)=\sum \mathsf{\Gamma}_{k_1\ldots k_n}x_1^{k_1}\ldots x_n^{k_n}
\qquad(k_1+\ldots+k_n\leqq v)
\]
라 하자. \(\mathfrak K\)는 모든 \(\mathsf{\Gamma}\)를 포함하는 유한 대수적 수체이고, \(\mathfrak p\)는 \(\mathfrak K\)의 소아이디얼이라 하자.

\(\bar E(x)\)의 정확한 차수에 맞추어 만든 가약성형식 \(R(Z,u)\)는 \(Z=\mathsf{\Gamma}\)를 대입해도 \emph{0이 아니며},
\[
R(\Gamma,u)=\sum P_\lambda U_\lambda\ne0
\qquad[u\text{에 관하여 항등적으로}]
\]
가 된다. 따라서 아이디얼 \(\mathfrak r=(\mathsf{P}_1,\mathsf{P}_2,\ldots)\)은 \(\mathfrak K\)의 소아이디얼들 가운데 \emph{유한개}로만 나누어진다. 그러므로 \(\mathfrak p\)가 이 유한개와 다르면 \(R(\mathsf{\Gamma},u)\)는 \(\mathfrak p\)를 법으로 하여 항등적으로 사라지지 않는다. \(\mathfrak p\)를 법으로 한 잉여류들이 다시 체를 이루므로, \(\bar E(x)\)는 \emph{\(\mathfrak p\)를 법으로 하여 기약}, 더 정확히는 절대기약이다. \(\bar E(x,y)\)도 마찬가지이므로 \(\mathfrak p\)를 법으로 한 차수저하도 일어나지 않는다.

\(\bar E(x)\)의 계수들이 대수적 분수라면, 공통분모 \(\varDelta\)를 나누는 유한개의 소아이디얼도 \(\mathfrak p\)에서 제외해야 한다. 나머지 모든 소아이디얼을 법으로 할 때에는 \(\bar E(x)\)를 \(\varDelta\cdot\bar E(x)\)로 바꿀 수 있으므로 위 가정들이 성립한다. 이로써 정리가 증명되었다.

\begin{flushleft}
에를랑겐, 1921년 8월.
\end{flushleft}
\begin{center}
(1921년 8월 28일 접수.)
\end{center}

\clearpage
\setcounter{footnote}{0}
\phantomsection
\addcontentsline{toc}{section}{21. 형식적 변분법과 미분불변식}
%
%
%
%
%
%

\section*{21. 형식적 변분법과 미분불변식}
\begin{center}
Encyklopädie d. math. Wiss. III, 3 (1922), 68--71쪽
(R. Weitzenböck, \emph{Differentialinvarianten}에 수록)
\end{center}

\begingroup
\renewcommand{\thefootnote}{\arabic{footnote})}
\setcounter{footnote}{148}
\subsection*{28. 형식적 변분법과 미분불변식.\footnote{이 항목은 E. 뇌터가 집필하였다.}}

모든 미분불변식을 생성하고 주요 정리들을 유도하는 일은 변분법의 형식적 방법을 사용하면 가장 통일적으로 수행할 수 있다. 이 원리는 리만\textsuperscript{78)}에게서 비롯되었고, 형식론의 관점에서는 립시츠\textsuperscript{73)}가 이를 가장 본격적으로 전개하였다.

동차형식 $f(dx)$ --- 여기서
\[
\left|\frac{\partial^2 f}{\partial dx_i\,\partial dx_k}\right|\ne0
\]
이라고 하자 --- 에 딸린 변분문제
\srcnumdisplay{(140)}{\delta\int_{t_0}^{t_1} f(x')\,dt=0,
\qquad \left(x_i'=\frac{dx_i}{dt}\right)}
에서 출발하면, 부분적분을 거쳐 다음 불변 방정식계에 이른다.
\[
 2\psi_i=\frac{d}{dt}\frac{\partial f}{\partial x_i'}-\frac{\partial f}{\partial x_i}=0.
\]
라그랑주 식 $\psi_i$들은 $dx$에 반변이다. 이를 형식적으로 가장 간단히 확인하려면, 부분적분에 대응하면서 적분을 포함하지 않는 다음 형식적 항등식으로 $\psi_i$를 정의하면 된다(라그랑주의 중심방정식\footnote{특수한 $f$에 관한 이 명칭은 Heun에게서 비롯되었다. 그의 논문 IV, 1 II, 447쪽 참조.}).
\[
 \delta f-df_\delta=-2\sum\psi_i(d,d)\,\delta x_i,
 \qquad
 \left(f_\delta=\sum\frac{\partial f}{\partial dx_i}\,\delta x_i\right).
\]
$f$와 $df_\delta$가 모두 불변식이므로 우변도 불변식이다. 특히 이차형식의 경우에는
\[
 \delta\sum g_{ik}dx_i dx_k-2d\sum g_{ik}dx_i\delta x_k
 =-2\sum\psi_\mu(d,d)\,\delta x_\mu
\]
를 얻는다. 여기서 $d$를 그것과 동변인 $(d+\lambda\delta)$로 바꾸고 $\lambda$의 거듭제곱별로 계수를 비교하면, 이에 대응하는 다음 항등식계를 얻는다.
\srcnumdisplay{(141)}{\left\{
\begin{aligned}
D\sum g_{ik}dx_i dx_k-2d\sum g_{ik}dx_iDx_k
   &=-2\sum\psi_\mu(d,d)\,Dx_\mu,\\
D\sum g_{ik}dx_i\delta x_k-\delta\sum g_{ik}dx_iDx_k
   &\phantom{={}}{}-d\sum g_{ik}\delta x_iDx_k\\
   &=-2\sum\psi_\mu(d,\delta)\,Dx_\mu,\\
D\sum g_{ik}\delta x_i\delta x_k-2\delta\sum g_{ik}\delta x_iDx_k
   &=-2\sum\psi_\mu(\delta,\delta)\,Dx_\mu.
\end{aligned}
\right.}
여기서 $\psi_\mu(d,\delta)$가 반변벡터임이, 따라서 특히 $\psi_\mu(d,d)$와 $\psi_\mu(\delta,\delta)$도 반변벡터임이 따른다. (141)에서 $\psi_\mu(d,\delta)$를 계산하면
\[
 \psi_\mu(d,\delta)=\sum g_{i\mu}\,d\delta x_i+
 \sum\left[\begin{smallmatrix}ik\\ \mu\end{smallmatrix}\right]dx_i\delta x_k
\]
이고, 여기서 다음 동변벡터를 얻는다.
\srcnumdisplay{(142)}{p^\sigma(d,\delta)=\sum g^{\sigma\mu}\psi_\mu
=d\delta x_\sigma+
\sum\left\{\begin{smallmatrix}ik\\ \sigma\end{smallmatrix}\right\}dx_i\delta x_k.}
$p^\sigma(d,d)=0$으로 놓으면 측지선의 방정식이 됨은 분명하다.

동변벡터 $p^\sigma(d,\delta)$를 알면 공변미분을 곧바로 얻을 수 있다(19번 참조). 이를테면 $h(d,\delta)$가 두 열 $dx$와 $\delta x$의 형식이라면,
\srcnumdisplay{(143)}{h^{(1)}(dx,\delta x)=\delta h-
\sum\frac{\partial h}{\partial dx_\sigma}p^\sigma(d,\delta)-
\sum\frac{\partial h}{\partial \delta x_\sigma}p^\sigma(\delta,\delta)}
은 불변식들의 차이이므로 그 자체도 불변식이며, 이차 미분들이 소거되도록 만들어져 있다. 따라서 $h^{(1)}(dx,\delta x)$는 일차 미분만을 포함하는 형식, 곧 $h$의 공변미분이다. $h$가 더 많은 열 $d^{(1)}x,\ldots,d^{(\lambda)}x$를 포함할 때에도 같은 사실이 성립한다.\footnote{Lipschitz, J. f. Math. 72 (1870), 17쪽.}

리만과 립시츠의 곡률형식도 같은 원리에 따라 구성된다. $\delta^2f$에서 출발하여 ``제2변분의 정규형''
\srcnumdisplay{(144)}{\Omega=\delta^2\sum g_{ik}dx_i dx_k
-2d\delta\sum g_{ik}dx_i\delta x_k
+d^2\sum g_{ik}\delta x_i\delta x_k}
으로 넘어간다. 이는 불변식들의 결합이므로 다시 불변식이며, 라그랑주의 중심방정식을 일반화하여 삼차 미분들이 이미 소거되어 있다. 이차 미분은 공변미분의 경우와 마찬가지로 다음 차이를 만들어 소거한다.
\srcnumdisplay{(145)}{K=\Omega-2\sum\{p^\sigma(d,\delta)\psi_\sigma(d,\delta)
-p^\sigma(\delta,\delta)\psi_\sigma(d,d)\}.}
이는 다음과 같이 쓸 수도 있다.\footnote{Lipschitz, J. f. Math. 82 (1876), § 1.}
\srcnumdisplay{(146)}{\begin{aligned}
\frac12 K={}&d\sum\psi_\sigma(d,\delta)\delta x_\sigma
-\delta\sum\psi_\sigma(d,d)\delta x_\sigma\\
&-\sum\{p^\sigma(d,\delta)\psi_\sigma(d,\delta)
-p^\sigma(\delta,\delta)\psi_\sigma(d,d)\}.
\end{aligned}}
$K$의 구성법은 $\Omega$에서 $p(d,\delta)$와 $p(\delta,\delta)$, 또는 (141)의 우변들을 영으로 놓아 이차 미분을 소거하는 것이라고도 표현할 수 있다. 이것이 리만의 곡률형식 정의이다(21번 참조). 다만 (146)이 특히 분명히 보여 주듯이, 미분을 수행한 \emph{뒤에야} 영으로 놓아야 한다는 점을 명시적으로 유의해야 한다. 립시츠의 형식적 구성법은 이 어려움을 피한다. 마찬가지로 공변미분 $h^{(1)}$ (143)도 $\delta h$에서 $p(d,\delta)\ldots$를 영으로 놓아 이차 미분을 소거한 것으로 정의할 수 있다.

곡률형식에 공변미분을 차례로 취하면 형식열 $[\Phi_4,\Phi_5,\ldots\text{ (Christoffel)}]$을 얻는다. 크리스토펠과 리치에 따르면 여기에 $f$를 덧붙였을 때 이 형식열의 사영불변식들은 이차형식의 모든 미분불변식을 빠짐없이 주며, 별도의 적분가능성 조건 없이 선형변환에 관한 이 형식열의 동치는 두 이차 미분형식의 동치를 뜻한다(22번 참조). 이 환원정리의 내적인 근거는 변분문제 (140)에서 생기는 리만 정규좌표의 존재에 있다(20번 참조). 이 좌표는 한 점을 지나는 측지선들을 직선으로 바꾸므로, $x$를 임의로 변환할 때 선형적으로 변환된다. 따라서 모든 불변식을 구성하는 문제와 동치 문제는 $f$를 정규좌표로 전개했을 때의 각 동차 성분에 관한 대응 문제로 환원된다. 그리고 이 성분들은 $f,\Phi_4,\Phi_5,\ldots$의 선형결합으로 이루어짐을 보일 수 있다. 이차형식의 경우에는 E. 뇌터의 주\textsuperscript{93)}에 의거하여 Vermeil\textsuperscript{92)}이 이를 상세히 수행하였다. 그 주에 따르면 임의의 미분식을 바탕으로 삼아도 정리와 방법은 그대로 유지된다(22번과 23번 참조). 이로써 순수 소거론의 방법에 견주어 형식적 변분법의 방법이 갖는 효력이 드러난다.\footnote{이 설명과 더 많은 기하학적 해석에 관해서는 ``문헌'' 아래에 적힌 F. Klein의 세미나 강연들을 참조하라.}

Levi-Civita\textsuperscript{134)}는 $p(d,\delta)$를 영으로 놓는 것을 ``벡터의 평행이동''이라고 기하학적으로 해석하였다(Hessenberg\textsuperscript{133)}도 참조). Weyl이 공리적으로 정식화한 이 개념은 본질적으로 이차형식에 한정되지만, 다른 방향의 일반화를 허용한다. 이차형식뿐 아니라 선형형식 하나를 함께 바탕으로 삼으면, 군을 확장하여 $g_{ik}$에 임의의 함수를 곱하더라도 이 개념은 보존된다. 이때 $p(d,\delta)$는 유일하게 정해지고, 위와 같은 방식으로 불변식을 구성하고 측지좌표를 정의할 수 있다. 그러나 이제 $p(d,d)=0$은 더 이상 변분문제에서 생기지 않는다!\footnote{H. Weyl, Reine Infinitesimalgeometrie, Math. Ztschr. 2 (1918), 384--411쪽 및 Raum, Zeit, Materie (제3판 및 제4판).} 여기서 환원정리와 동치에 관한 문제는 아직 일반적으로 다루어지지 않았다. 가장 낮은 차수의 불변식에 대해서만 R. Weitzenböck이 환원정리를 제시하였다.\footnote{Wien. Ber. 129 (1920), 683쪽 및 697쪽.}

마지막으로 일반적인 ``불변 변분문제'', 곧 (단일 또는 다중) 적분 $J$가 리 군의 의미에서 어떤 군에 대하여 불변인 문제를 지적해 두자. 물리학적으로 흥미로운 특수한 경우는 Hamel\footnote{Math. Ann. 59 (1904), 416--434쪽; Ztschr. Math. Phys. 50 (1904), 1--54쪽.}, Herglotz\footnote{Ann. d. Phys. (4) 36 (1911), 511쪽, 특히 § 9.}, Lorentz\footnote{Verslag der Amsterd. Akad., 1916년 4월ㆍ9월ㆍ10월 및 1917년 5월.}, Fokker\footnote{같은 곳, 1917년 1월.}, Weyl\footnote{Ann. d. Phys. (4) 54 (1917), 117쪽, § 2.}, Klein\footnote{Gött. Nachr. 1917, 469--482쪽; 같은 곳 1918, 171--189쪽.}의 논문들에서 다루어진다. E. 뇌터의 원리적 정식화\footnote{Gött. Nachr. 1918, 235--257쪽.}에 따르면, $J$가 $\varrho$개의 본질적 매개변수를 갖는 유한군 $G_\varrho$에 대하여 불변이면 이에 $\varrho$개의 선형독립인 발산식이 대응한다. 또 $J$가 $\sigma$차 도함수까지 포함하는 $\varrho$개의 임의함수에 의존하는 무한군에 대하여 불변이면, 라그랑주 식들과 그 $\sigma$차까지의 도함수 사이에 $\varrho$개의 종속관계가 대응한다. 두 경우 모두 역도 성립한다. 라그랑주 식들이 그 군의 (상대) 불변식이 되므로, 동시에 하나의 불변식 생성과정을 얻는다.

\begin{center}
(1921년 3월 탈고.)
\end{center}
\endgroup

\clearpage
\setcounter{footnote}{0}
\phantomsection
\addcontentsline{toc}{section}{22. 다항식 아이디얼과 종결식 이론}
%
%
%
%
%
%

\providecommand{\Amod}{\mathfrak A}
\providecommand{\Bmod}{\mathfrak B}
\providecommand{\Gmod}{\mathfrak G}
\providecommand{\Mmod}{\mathfrak M}
\providecommand{\Nmod}{\mathfrak N}
\providecommand{\Rmod}{\mathfrak R}
\providecommand{\Smod}{\mathfrak S}
\providecommand{\Tmod}{\mathfrak T}
\providecommand{\mideal}{\mathfrak m}
\providecommand{\nideal}{\mathfrak n}
\providecommand{\gideal}{\mathfrak g}
\providecommand{\hideal}{\mathfrak h}
\providecommand{\jideal}{\mathfrak j}
\providecommand{\tideal}{\mathfrak t}
\providecommand{\aideal}{\mathfrak a}
\providecommand{\bideal}{\mathfrak b}
\providecommand{\bb}{\mathfrak b}
\providecommand{\cc}{\mathfrak c}
\providecommand{\zero}[1]{\equiv0\pmod{#1}}
\providecommand{\Norm}{\operatorname N}

\section*{22. 다항식 아이디얼과 종결식 이론}

\begin{center}
\emph{Math. Ann. 88 (1923), 53--79쪽}\\[0.9em]
쿠르트 헨첼(\(\dagger\)) 지음\\[0.6em]
괴팅겐의 에미 뇌터 편집\srcfn{1)}{쿠르트 헨첼은 1914년 10월 이래 딕스뮈드 전선에서 실종되었으며 사망자로 보아야 한다. 이 글은 그가 1914년 여름 에를랑겐에서 E. Fischer의 지도 아래 제출한 학위논문 \glqq{}\emph{Zur Theorie der Formenmoduln und Resultanten}\grqq{}의 가장 본질적인 부분을 완전히 자유롭게 다시 쓴 것이다. 전적으로 독창적인 착상으로 쓴 그 학위논문은 빈틈없이 구성되었지만, 보조정리가 보조정리를 잇고 모든 개념이 네댓 개 지표를 가진 공식으로만 서술되며 본문이 거의 전혀 없다. 그래서 이해하기가 극도로 어렵다. 본인은 계획한 개작을 완수하지 못했다. 나는 이 연구를 순전히 개념적인 형태로 다시 제시한다. 이로써 기본착상이 한결같이 헨첼에게서 비롯된 증명들이 크게 단순해지고, 바라건대 이 연구의 아름다움도 드러날 것이다. 주어진 기저에서 유한개의 단계로 나타나는 함수들을 실제로 구성하는 문제를 다룬 학위논문의 부분은 별도 출판을 위해 남겨 둔다. (E. N.)}
\end{center}

이하에서는 \emph{다항식 아이디얼에 속하는 모든 다항식의 공통영점을 정한다는 소거론의 일반문제}를 다룬다. 이는 그 아이디얼의 기저를 이루는 임의의 유한개 다항식의 영점을 정하는 것과 같다. 동시에 나타나는 중복도에 대한 개념적 해석도 얻을 것이다. 다항식의 계수는 임의의 체에 속해도 되며, 영점은 그 체에서 유도되는 대수적으로 닫힌 체에 속한다. 또한 일반성을 잃지 않고 아이디얼을 변환된 아이디얼(제3절)이라고 가정할 수 있다. 그러면 본질적인 결과는 다음과 같다.

\emph{각 다항식 아이디얼 \(\mideal\)에 종결식형
\[
R_{\mideal}
=R^{(1)}(x_1,\ldots,x_n)\cdots
R^{(i)}(x_i,\ldots,x_n)\cdots R^{(n)}(x_n)
\zero{\mideal}
\]
을 대응시킬 수 있으며, \(R_{\mideal}\)은 \(\mideal\)의 모든 영점에서, 그리고 오직 그 영점들에서만 사라진다. \(\nideal\)이 \(\mideal\)의 인자이고 종결식형 \(R_{\nideal},R_{\mideal}\)이 같으면 아이디얼 \(\nideal,\mideal\)도 같다.}

주정리의 둘째 부분은 종결식형이 각 영점을 특징적인 중복도와 함께 준다는 것을 보여 준다. 이는 한 아이디얼이 다른 아이디얼의 인자인 두 대수적 수체 아이디얼이 그 \emph{노름}이 같을 때 그리고 그때에만 같다는 사실의 유사물이며, 두 정리의 내적 근거도 같다.

실제로 각 다항식을 \(x_1,\ldots,x_{i-1}\)의 거듭제곱곱들에 관한 선형형식으로 보고 \(x_{i+1},\ldots,x_n\)을 계수영역에 첨가하면, \(\mideal\)을 선형형식 모듈 \(\mathfrak M_{i-1}\)로 볼 수 있다. 종결식인자 \(R^{(i)}\)는 그때 대응하는 기초모듈(제1절) \(\Gmod_{i-1}\)의 \(\mathfrak M_{i-1}\)에 대한 \emph{노름}과 같아진다. 따라서 \emph{\(R^{(i)}\)의 한 인자의 차수와 지수의 곱으로 주어지는 중복도를 선형독립인 잉여류의 수로 해석}할 수 있다\srcfn{2)}{이 마지막 결과들의 진정한 의미는 아이디얼의 으뜸아이디얼 분해를 함께 고려할 때 드러난다. 한 으뜸아이디얼에 대응하는 중복도는 그 아이디얼에 대한 보충부분의 선형독립인 잉여류 수로 정의된다. 이는 다른 곳에서 다루겠다. (E. N.)}. 지금까지 이 사실은 부정계수에 대하여 종결식을 정의할 수 있는 특수한 경우에만 알려져 있었다\srcfn{3)}{Macaulay, \emph{The algebraic theory of modular systems}, Cambridge Tracts 19 (1916), 67번; Lasker, \emph{Zur Theorie der Moduln und Ideale}, \emph{Math. Ann.} 60 (1905), 20--116쪽, 특히 98쪽. 이 경우 종결식과 종결식형이 일치한다는 것은 주 1에서 말한 아직 출판되지 않은 헨첼의 정리들에 들어 있다. (E. N.)}.

이 결과들을 유도하기 위해 제1절과 제2절에서는 계수가 유리정수 또는 한 부정원의 다항식인 선형형식 모듈에 관한 정리들을 제시한다. 제3절에서 도입하는 다항식 아이디얼과의 연결은 제4절의 동형정리가 맡는다. 제5절은 위에서 말한 주정리의 둘째 부분을, 제7절은 영점에 관한 정리를 준다. 그에 앞선 제6절에서는 일반 소거론을 전개한다. 여기서는 새 부정원들을 도입하여 레졸벤트가 이 부정원들의 일차형식들로 분해됨도 증명한다. 이로써 여기서 제시하는 소거법 안에서 크로네커의 한 주장에 엄밀한 증명을 준다\srcfn{4)}{König, \emph{Algebraische Größen} (Leipzig: Teubner, 1903), V, \S{}4가 크로네커 소거론에 대하여 시도한 증명은 성공하지 못한 것으로 알려져 있다. König이 크로네커 소거론에 도입한 중복도들에 대해서도 헨첼 주정리의 둘째 부분이 성립하지 않는다는 사실은 Macaulay, 앞의 책 23쪽에 나와 있다. 거기서는 크로네커 소거론이 으뜸아이디얼 분해와 대응하지 않는다는 것도 보인다. (E. N.)}.

\section*{\S{}1. 선형형식의 모듈}

앞의 두 절에서 정수적 원소 \(a,b,c,\ldots\)란 유리정수 또는 추상적으로 정의된 임의의 체 \(\mathsf{P}\)를 계수체로 갖는 한 부정원의 다항식을 뜻한다. 선형형식 \(a(\xi),b(\xi),\ldots\)란 정수적 원소를 계수로 갖는 유한개의 부정원 \(\xi_1\ldots\xi_k\)의 선형형식을 뜻한다.

\emph{선형형식의 모듈 \(\Amod\)}은 다음 조건으로 정의한다. \(\Amod\)가 \(a(\xi),b(\xi)\)를 포함하면 차 \(a(\xi)-b(\xi)\)도 포함한다. \(a(\xi)\)를 포함하면 임의의 정수적 원소 \(c\)에 대하여 \(c\,a(\xi)\)도 포함한다.

\(\Amod\)의 \emph{랭크}는 통상처럼 \(\Amod\)에 속하는 선형독립인 선형형식의 최대수로 정의한다. \(\Amod\)의 \(\lambda\)번째 \emph{행렬식인자} \(d_\lambda\)는 \(\Amod\)에서 얻는 모든 \(\lambda\)차 행렬식의 최대공약원이다. 즉 \(\Amod\)의 선형형식 \(\lambda\)개씩의 계수행렬에서 만들 수 있는 모든 \(\lambda\)차 행렬식을 취한다. \(\varrho\)가 \(\Amod\)의 랭크여서 \(d_1,\ldots,d_\varrho\)가 영이 아닌 행렬식인자들이라면, \(\Amod\)의 \emph{초등인자}는
\[
e_1=d_1,\quad e_2=\frac{d_2}{d_1},\ldots,\quad
e_\varrho=\frac{d_\varrho}{d_{\varrho-1}},\qquad
e_{\varrho+i}=0
\]
으로 정의한다. 잘 알려져 있듯 \(\Amod\)은 정확히 \(\varrho\)개의 선형형식 \(a_1(\xi),\ldots,a_\varrho(\xi)\)로 이루어진 모듈기저를 갖는 유한모듈이다. \(\Amod\)의 모든 선형형식은 이 기저를 정수적 원소의 계수로 선형표현할 수 있다.
\[
\Amod=(a_1(\xi),\ldots,a_\varrho(\xi)).
\]
\(A\)가 이 선형형식들의 계수행렬이면 \(\Amod\)의 행렬식인자와 초등인자는 \(A\)의 것들과 같다. 먼저 여기서 \emph{각 초등인자 \(e_\lambda\)가 다음 초등인자 \(e_{\lambda+1}\)을 나눈다}는 것이 따라온다. 초등인자론의 행렬방정식
\[
PAQ=
\begin{pmatrix}
e_1&&&0\\
&e_2&&\\
&&\ddots&\\
0&&&e_\varrho
\end{pmatrix}
\]
에서 단모듈러변환 \(\xi=Q(\eta)\)로 새 부정원 \(\eta_1\ldots\eta_k\)를 도입하면 모듈등식
\begin{equation}
\Amod=(e_1\eta_1,e_2\eta_2,\ldots,e_\varrho\eta_\varrho)
\tag{1}
\end{equation}
을 얻는다.

이하에서 가장 중요한 개념은 기초모듈의 개념이다\srcfn{5)}{Steinitz, \emph{Rechteckige Systeme und Moduln in algebraischen Zahlkörpern II}, \emph{Math. Ann.} 72 (1912), 297--345쪽, 36번 참조. 다른 접점들도 있지만, 적어도 헨첼은 더 단순한 자신의 경우에서 슈타이니츠와 독립적으로 식 (4)의 형태로 이 개념에 도달한 듯하다. (E. N.)}.

\medskip
\noindent\textbf{정의 I.} \emph{같은 랭크의 진인자를 갖지 않는 모듈 \(\Gmod\)을 기초모듈이라 한다}\srcfn{6)}{여기서 인자는 모듈의 의미이다. \(\Amod\)의 모든 원소가 \(\Bmod\)에 들어갈 때 \(\Amod\)은 \(\Bmod\)로 나누어진다고 쓰며, \(\Amod\zero{\Bmod}\)이다. \(\Bmod\)가 \(\Amod\)에 없는 원소를 포함하면 진인자라고 한다.}.

따라서 \(\Gmod\)이 기초모듈인 것은
\begin{equation}
c\cdot\,g(\xi)\zero{\Gmod},\qquad c\ne0
\quad\Longrightarrow\quad g(\xi)\zero{\Gmod}
\tag{2}
\end{equation}
가 언제나 성립할 때 그리고 그때뿐이다.

임의의 모듈과 기초모듈의 관계는 다음 정리가 준다.

\medskip
\noindent\textbf{정리 I.} \emph{모든 모듈 \(\Amod\)은 같은 랭크의 유일한 기초모듈 \(\Gmod\)로 나누어진다. \(\Gmod\)은 \(\Amod\)의 같은 랭크를 가진 최소인자로 정의되며
\begin{equation}
\Gmod=(\eta_1,\ldots,\eta_\varrho)
\tag{3}
\end{equation}
로 표현된다. \(\Gmod\)을 \(\Amod\)의 기초모듈이라 부른다.}

\(\Gmod\)이 같은 랭크의 최소인자라는 조건은 다음처럼 표현된다. \(\Gmod\)은 관계
\begin{equation}
b\,g(\xi)\zero{\Amod},\qquad b\ne0
\tag{4}
\end{equation}
를 만족하는 선형형식 \(g(\xi)\)만을, 그리고 그 모두를 포함한다. 여기서 먼저 \(\Gmod\)이 모듈임을 알 수 있다. 실제로
\[
b_1g_1(\xi)\zero{\Amod},\quad b_1\ne0,\qquad
b_2g_2(\xi)\zero{\Amod},\quad b_2\ne0
\]
이면
\[
b_1b_2\bigl(g_1(\xi)-g_2(\xi)\bigr)\zero{\Amod},
\qquad b_1b_2\ne0
\]
이고, 물론 \(b_1\cdot c\,g_1(\xi)\zero{\Amod}\)도 성립한다. \(\Gmod\)은 또한 기초모듈이다. 실제로
\[
c\cdot g(\xi)\zero{\Gmod};\quad c\ne0
\]
이면 (4)에 따라
\[
b c g(\xi)\zero{\Amod};\quad bc\ne0
\]
이고, 다시 (4)에 따라 \(g(\xi)\zero{\Gmod}\)이다.

조건들을 만족하면서 같은 랭크의 최소인자 \(\Gmod\)과 다른 기초모듈 \(\Gmod_1\)은 존재하지 않는다. \(\Gmod_1\)은 \(\Gmod\)로 나누어져야 하지만 기초모듈로서 같은 랭크의 진인자를 갖지 않으므로 \(\Gmod\)과 같다. 끝으로 (3)의 오른쪽이 나타내는 모듈은 기초모듈이다. 그 진인자는 또 다른 부정원 \(\eta\)를 포함해야 하므로 더 높은 랭크를 갖기 때문이다. 따라서 (3)은 \(\Amod\)의 유일하게 정의된 기초모듈을 주며, 이로써 \emph{정리 I의 모든 부분이 증명되었다}.

이하에서 또 하나 중요한 개념은 두 모듈의 몫이라는 데데킨트의 개념이다\srcfn{7)}{데데킨트는 곱셈도 정의되는 수모듈을 다루므로 그의 몫은 여기서 정의하는 몫과 일치하지 않는다. (E. N.)}.

\medskip
\noindent\textbf{정의 II.} \emph{두 선형형식 모듈의 몫 \(\cc=\Amod/\Bmod\)은 조건
\[
c\,\Bmod\zero{\Amod}
\]
을 만족하는 정수적 원소 \(c\) 전체로 정의한다. \(\cc\)는 정수적 원소들의 아이디얼, 따라서 주아이디얼이다. 이에 대응하여 선형형식 모듈 \(\Amod\)을 정수적 원소의 아이디얼 \(\bb\)로 나눈 몫 \(\mathfrak C=\Amod/\bb\)는 조건
\[
c(\xi)\,\bb\zero{\Amod}
\]
을 만족하는 선형형식 \(c(\xi)\) 전체로 정의한다. \(\mathfrak C\)는 다시 선형형식 모듈이며, \(\bb\)가 영아이디얼이 아니면 \(\Amod\)과 같은 랭크를 갖는다.}

정의에 따라 각 몫이 모듈의 성질을 갖는다는 것은 분명하다. 한편 정수적 원소들로 이루어지고 모듈의 성질을 갖는 계는 아이디얼이다.

몫 개념은 기초모듈과 가장 높은 초등인자 사이에 다음 관계를 준다.

\medskip
\noindent\textbf{정리 II.} \emph{\(\Amod\)의 기초모듈과 가장 높은 초등인자 사이에는 역관계
\begin{equation}
(e_\varrho)=\Amod/\Gmod,\qquad
\Gmod=\Amod/(e_\varrho)
\tag{5}
\end{equation}
가 성립한다. 여기서 \((e_\varrho)\)는 \(e_\varrho\)가 생성하는 주아이디얼이다.}

모든 초등인자가 \(e_\varrho\)를 나누므로 (1), (3)에 따라
\begin{equation}
e_\varrho\Gmod\zero{\Amod},
\qquad\text{따라서}\qquad
(e_\varrho)\cdot\Gmod\zero{\Amod}
\tag{6}
\end{equation}
이다. 그러므로 \((e_\varrho)\)는 몫 \(\Amod/\Gmod\)으로 나누어진다. 역도 성립한다. 실제로
\[
b\Gmod\zero{\Amod}
\quad\Longrightarrow\quad
b\eta_\varrho\zero{\Amod}
\quad\Longrightarrow\quad
b\equiv 0(e_\varrho).
\]
이로써 (5)의 첫째 공식이 증명된다\srcfn{8)}{\(\Amod_0=\Amod,\ldots,\Amod_\lambda=(\Amod,\eta_\varrho,\ldots,\eta_{\varrho-\lambda+1})\)로 두자. 여기서 \(\Amod_\lambda\)의 가장 높은 초등인자는 \(e_{\varrho-\lambda}\)이다. 같은 추론으로
\[
(e_\varrho)=\Amod_0/\Amod_1,\ldots,\quad
(e_{\varrho-\lambda})=\Amod_\lambda/\Amod_{\lambda+1},\ldots,\quad
(e_1)=\Amod_{\varrho-1}/\Amod_\varrho
\]
를 얻는다. 더 일반적으로 \(\zeta_\lambda=b_{\lambda1}\eta_1+\cdots+b_{\lambda\lambda}\eta_\lambda\)로 두고 \((b_{\lambda\lambda},e_\lambda)=1\)이라 하며
\[
\Bmod_0=\Amod,\ldots,\qquad
\Bmod_\lambda=(\Amod,\zeta_\varrho,\ldots,\zeta_{\varrho-\lambda+1})
\]
로 두자. 그러면 \(\Bmod_\lambda\)의 가장 높은 초등인자도 \(e_{\varrho-\lambda}\)이므로
\[
(e_{\varrho-\lambda})=\Bmod_\lambda/\Bmod_{\lambda+1}
\]
이다.}. 또한 (6)은 \(\Gmod\)이 몫 \(\Amod/(e_\varrho)\)으로 나누어짐을 보인다. 이 몫은 \(\Amod\)의 같은 랭크를 가진 인자이므로, \(\Amod\)의 기초모듈 정의에 따라 거꾸로 \(\Gmod\)로도 나누어진다. 이로써 (5)의 둘째 공식도 증명되었다.

\(d\)가 \(\Amod\)에서 얻는 항등적으로 영이 아닌 임의의 \(\varrho\)차 행렬식이라 하자. 그러면 \(d\)는 \(\varrho\)번째 행렬식인자 \(d_\varrho\), 따라서 \(e_\varrho\)로 나누어진다. 그러므로 \(d\Gmod\zero{\Amod}\)이고 앞의 추론에서 \(\Gmod=\Amod/(d)\)가 따라온다. 뒤에서 중요한 점은 \(\Gmod\)의 이 몫표현을 초등인자로 돌아가지 않고도 증명할 수 있어서 더 일반적으로 성립한다는 사실이다.

\medskip
\noindent\textbf{정리 III.} \emph{정수적 원소를 여러 부정원의 다항식으로 이해하더라도}\srcfn{9)}{실제로 쓰는 것은 정수적 원소들이 통상적인 계산법칙 아래 덧셈, 뺄셈, 곱셈으로 닫혀서 환을 이룬다는 것, 그리고 이 환에서 곱이 사라지면 한 인자가 사라져서 원소쌍을 첨가하는 몫구성으로 체까지 확대할 수 있다는 것뿐이다. 반면 모듈의 기저표현은 쓰지 않는다. 따라서 예를 들어 계수들이 모든 정수적 대수적 수의 정역에 속할 때에도 (7)이 성립한다. (E. N.)}, \emph{몫표현
\begin{equation}
\Gmod=\Amod/(d)
\tag{7}
\end{equation}
이 여전히 성립한다. 여기서 \(d\)는 \(\Amod\)에서 얻는 항등적으로 영이 아닌 임의의 \(\varrho\)차 행렬식이다.}

이 확대 아래에서도 랭크, 기초모듈, 모듈몫의 개념은 유지된다. 다만 이제 모듈몫은 주아이디얼일 필요가 없다. 또한 기저표현을 제외하면 정리 I도 유지된다.

이제
\[
\begin{aligned}
a_1(\xi)&=a_{11}\xi_1+\cdots+a_{1k}\xi_k,\ldots,\\
a_\varrho(\xi)&=a_{\varrho1}\xi_1+\cdots+a_{\varrho k}\xi_k
\end{aligned}
\]
가 \(\Amod\)에 속하는 \(\varrho\)개의 선형독립인 선형형식이라고 하자. 이들은 일반적으로 모듈기저를 이루지 않는다. 부정원 \(\xi\)들의 번호를
\[
d=|a_{ij}|\ne0,\qquad i,j=1,2,\ldots,\varrho
\]
가 되도록 정하자. 그러면 다음 특별한 모양의 \(\varrho\)개 선형형식이 \(\Amod\)에 속한다.
\begin{equation}
d\xi_\lambda+b_{\lambda,\varrho+1}\xi_{\varrho+1}
+\cdots+b_{\lambda,k}\xi_k\zero{\Amod},
\qquad(\lambda=1,2,\ldots,\varrho).
\tag{8}
\end{equation}
그러나 \(\Amod\)은 \(\xi_{\varrho+1}\ldots\xi_k\)에만 의존하는 선형형식
\[
c_{\varrho+1}\xi_{\varrho+1}+\cdots+c_k\xi_k
\]
을 포함할 수 없다. 그렇다면 적어도 하나의 \((\varrho+1)\)차 행렬식 \(d\,c_\sigma\)가 영이 아닐 것이기 때문이다.

이제 몫 \(\Amod/(d)\)는 \(\Amod\)의 같은 랭크를 가진 인자이므로 \(\Gmod\)로 나누어진다. 그러나 역도 성립한다. 각 \(g(\xi)\zero{\Gmod}\)에 대하여 정의상
\[
c\cdot\,g(\xi)\zero{\Amod},\qquad c\ne0,
\]
따라서
\[
d\cdot c\cdot\,g(\xi)\zero{\Amod}
\]
이다. 한편 (8)에 따라
\[
d\cdot\,g(\xi)\equiv g_{\varrho+1}\xi_{\varrho+1}+\cdots+g_k\xi_k
\zero{\Amod}.
\]
그러므로
\[
cg_{\varrho+1}\xi_{\varrho+1}+\cdots+cg_k\xi_k\zero{\Amod}.
\]
방금 지적한 사실에 따라 \(c\cdot g_{\varrho+1}=0,\ldots,c\cdot g_k=0\)이고, \(c\ne0\)이므로 \(g_{\varrho+1}=0,\ldots,g_k=0\)이다. 따라서 \(d\cdot\,g(\xi)\zero{\Amod}\)이며 (7)이 증명되었다.

\(d\cdot\,g(\xi)\zero{\Amod}\)은 두 모듈
\[
(a_1(\xi),\ldots,a_\varrho(\xi)),\qquad
(g(\xi),a_1(\xi),\ldots,a_\varrho(\xi))
\]
이 같은 랭크 \(\varrho\)를 갖는다는 사실에서도 직접 따라온다. 실제로 \(g(\xi)=c_1\xi_1+\cdots+c_k\xi_k\)로 두면 \((\varrho+1)\)차 행렬식
\[
\begin{vmatrix}
a_1(\xi)&a_{11}&\cdots&a_{1\varrho}\\
\vdots&\vdots&&\vdots\\
a_\varrho(\xi)&a_{\varrho1}&\cdots&a_{\varrho\varrho}\\
g(\xi)&c_1&\cdots&c_\varrho
\end{vmatrix}
\]
이 사라진다. 그러나 위에서 준 증명은 제4절에서 식 (30)에 이어 다시 나타난다.

\section*{\S{}2. 노름과 모듈의 분해정리}

이제 다시 유리정수 또는 \(\mathsf{P}\)를 계수체로 갖는 한 부정원의 다항식에 관한 선형형식 모듈을 다룬다.

\medskip
\noindent\textbf{정의 III.} \emph{통상처럼 \(\xi\)의 선형형식 가운데 \(\Amod\)을 법으로 서로 합동인 것들을 한 잉여류로 묶자. 기호 \(\Gmod\mid\Amod\)은 \(\Amod\)에 관한 \(\Gmod\)의 잉여류계, 즉 \(\Gmod\)의 원소들로 이루어진 \(\Amod\)의 모든 잉여류의 계를 뜻한다. \(\Amod\)에 대한 \(\Gmod\)의 노름, 기호로 \(\Norm(\Gmod\mid\Amod)\), 은 \(\Gmod\)에서 \(\Amod\)으로 넘어가는 치환의 행렬식으로 정의한다. 즉 \(\Amod\)의 임의의 선형독립인 모듈기저를 \(\Amod\)의 기초모듈 \(\Gmod\)의 그러한 기저로 나타내는 치환의 행렬식이다.}

(1), (3), 또는 정리 II 뒤의 주에서 다음을 얻는다.
\begin{equation}
\Norm(\Gmod\mid\Amod)=e_1e_2\cdots e_\varrho=d_\varrho,
\qquad
\Gmod=\Amod/\bigl(\Norm(\Gmod\mid\Amod)\bigr).
\tag{9}
\end{equation}

(1), (3), (9)에서 잘 알려진 방식으로 다음이 따라온다. 유리정수의 경우 \(\Norm(\Gmod\mid\Amod)\)은 \(\Amod\)에 관한 \(\Gmod\)의 잉여류 수를 나타낸다. 한 부정원의 다항식인 경우에는 잉여류의 수가 일반적으로 무한이지만, \(\Norm(\Gmod\mid\Amod)\)의 차수는 \emph{\(\mathsf{P}\)에 관하여 선형독립인 잉여류의 수}와 같다. \(\Bmod\)이 \(\Amod\)의 같은 랭크를 가진 인자이면, \(\Bmod\)은 그러한 한 잉여류의 대표원과 함께 그 잉여류의 모든 원소를 포함한다. 따라서 \(\Bmod\)은 \(\Amod\)에 유한개의 잉여류 또는 그 선형결합들을 덧붙여 얻는다. 여기서 곧바로 다음이 따라온다.

\medskip
\noindent\textbf{정리 IV.} \emph{\(\Bmod\)이 \(\Amod\)의 같은 랭크를 가진 인자라서 두 기초모듈이 같고, 또한
\[
\Norm(\Gmod\mid\Bmod)=\Norm(\Gmod\mid\Amod)
\]
이면 \(\Bmod=\Amod\)이다.}

노름의 분해와 모듈의 분해 사이의 관계에 대해서는 다음이 성립한다.

\medskip
\noindent\textbf{정리 V.} \emph{
\begin{equation}
\Norm(\Gmod\mid\Amod)=rs,\qquad(r,s)=1
\tag{10}
\end{equation}
이라고 하자. 그러면 \(\Amod\)의 같은 랭크를 가진 인자 \(\Rmod\)과 \(\Smod\)이 각각 유일하게 존재하여
\[
r=\Norm(\Gmod\mid\Rmod),\qquad
s=\Norm(\Gmod\mid\Smod)
\]
가 된다. \(\Rmod,\Smod\)은
\[
\Rmod=\Amod/(s),\qquad
\Smod=\Amod/(r)
\]
로 정의되며, \(\Amod\)은 \(\Rmod,\Smod\)의 최소공배모듈과 같다.}\srcfn{10)}{최소공배모듈 \([\Rmod,\Smod]\)은 \(\Rmod\)로도 \(\Smod\)으로도 나누어지는 모든 선형형식의 전체이다. 이에 대응하여 최대공약모듈 \((\Rmod,\Smod)\)은 \(\Rmod\)의 한 원소와 \(\Smod\)의 한 원소의 합으로 표현되는 모든 선형형식의 전체이다.}
\emph{\(r_\varrho=(r,e_\varrho)\), \(s_\varrho=(s,e_\varrho)\)로 두면 다음 역관계가 성립한다.
\begin{equation}
(s_\varrho)=\Amod/\Rmod,\quad
\Rmod=\Amod/(s_\varrho);\qquad
(r_\varrho)=\Amod/\Smod,\quad
\Smod=\Amod/(r_\varrho).
\tag{11}
\end{equation}}

실제로 일반적으로 \(r_i=(r,e_i)\), \(s_i=(s,e_i)\)로 두면 (9), (10)에 따라
\[
(r_i,s_i)=1,\qquad e_i=r_is_i,\qquad
r=r_1\cdots r_\varrho,\qquad s=s_1\cdots s_\varrho,
\]
이고 각 \(r_i\), 또는 \(s_i\)는 다음 \(r_{i+1}\), 또는 \(s_{i+1}\)을 나눈다. 따라서
\[
\Rmod=(r_1\eta_1,\ldots,r_\varrho\eta_\varrho),\qquad
\Smod=(s_1\eta_1,\ldots,s_\varrho\eta_\varrho)
\]
로 두면 \(\Rmod,\Smod\)은 \(\Amod\)의 같은 랭크를 가진 인자가 된다. \(r_i,s_i\)가 서로소이므로 \(\Amod\)은 \(\Rmod,\Smod\)의 최소공배모듈이고, 또한
\[
\Norm(\Gmod\mid\Rmod)=r_1\cdots r_\varrho=r,\qquad
\Norm(\Gmod\mid\Smod)=s_1\cdots s_\varrho=s
\]
이다.

더 나아가
\[
s_\varrho\Rmod\zero{\Amod},\qquad
r_\varrho\Smod\zero{\Amod}.
\]
\(c\Rmod\zero{\Amod}\)이면 \(cr_\varrho\eta_\varrho\zero{\Amod}\), 따라서 \(c\zero{(s_\varrho)}\)이다. 이로써 \((s_\varrho)=\Amod/\Rmod\)이 증명되고, 마찬가지로 \((r_\varrho)=\Amod/\Smod\)이 증명된다. 또
\[
b(\eta)=b_1\eta_1+\cdots+b_\varrho\eta_\varrho
\zero{\Amod/(s_\varrho)}
\]
이면 모든 \(r_i\)가 \(s_\varrho\)와 서로소이므로 \(b_i\zero{(r_i)}\)이다. 따라서 \(\Rmod=\Amod/(s_\varrho)\)이고, 마찬가지로 \(\Smod=\Amod/(r_\varrho)\)이다. 이로써 역관계 (11)이 증명되었다. 특수한 경우 \(r=1\)이면 (11)은 (5)가 된다.

이제 \(\Rmod,\Smod\)이 \emph{유일하게 정해짐}을 보여야 한다. \(\overline{\Rmod},\overline{\Smod}\)이 정리 V의 조건을 만족하는 모듈들이고, \(\bar r_i,\bar s_i\)가 각각의 초등인자라고 하자. \(\bar r_i,\bar s_i\)는 \(e_i\)의 인자이고
\[
r=\bar r_1\cdots\bar r_\varrho,\qquad
s=\bar s_1\cdots\bar s_\varrho,\qquad
(\bar r_i,\bar s_i)=1
\]
이므로
\[
e_i=\bar r_i\bar s_i,\qquad
\bar r_i=(r,e_i)=r_i,\qquad
\bar s_i=(s,e_i)=s_i
\]
이다. 따라서 \(\overline{\Rmod},\overline{\Smod}\)은 \(\Rmod,\Smod\)과 초등인자와 기초모듈이 각각 같다. 이제 단모듈러치환 \(\eta=Q(\zeta)\)로 새 부정원 \(\xi\)들을 도입하면
\[
\overline{\Rmod}=(r_1\zeta_1,\ldots,r_\varrho\zeta_\varrho)
\]
이고
\[
\Amod=\bigl(
r_1s_1(q_{11}\zeta_1+\cdots+q_{1\varrho}\zeta_\varrho),\ldots,
r_\varrho s_\varrho(q_{\varrho1}\zeta_1+\cdots+q_{\varrho\varrho}\zeta_\varrho)
\bigr).
\]
\(\Amod\zero{\overline{\Rmod}}\)에서
\[
r_i s_i(q_{i1}\zeta_1+\cdots+q_{i\varrho}\zeta_\varrho)\zero{\overline{\Rmod}}
\]
이고, 따라서 \(r_is_iq_{ij}\zero{(r_j)}\)이다. \((r,s)=1\)이므로 \(r_iq_{ij}\zero{(r_j)}\), 즉 \(\Rmod\zero{\overline{\Rmod}}\)이다. 그런데
\[
\Norm(\Gmod\mid\Rmod)=\Norm(\Gmod\mid\overline{\Rmod})
\]
이므로 정리 IV에 따라 \(\Rmod=\overline{\Rmod}\)이다. 마찬가지로 \(\Smod=\overline{\Smod}\)이다. 이로써 정리 V가 모두 증명되었다.

\section*{\S{}3. 다항식의 아이디얼}

\(\overline{\mideal}\)이 추상적으로 정의된 임의의 체 \(\mathsf{P}\)를 계수체로 갖는 \(n\)개 부정원 \(y_1\ldots y_n\)의 \emph{다항식 아이디얼}이라고 하자. 즉 \(\overline{\mideal}\)은 \(\Phi_1(y),\Phi_2(y)\)와 함께 차 \(\Phi_1(y)-\Phi_2(y)\)를 포함하고, \(\Phi(y)\)와 함께 \(\mathsf{P}\)의 계수를 갖는 임의의 다항식 \(A(y)\)에 대하여 \(A(y)\Phi(y)\)도 포함한다\srcfn{11)}{개념상 자연스러운 \glqq{}아이디얼\grqq{}이라는 명칭은 예전에 널리 쓰던 \glqq{}모듈\grqq{} 또는 \glqq{}형식모듈\grqq{} 대신 쓰인다. 여기서는 이미 선형형식 모듈과 구별하기 위해서도 이 명칭이 필요하다. (E. N.)}.

부정원들을 계수 자체가 부정원인 변환
\begin{equation}
y=U(x),\qquad\text{예를 들어}\qquad
\begin{cases}
y_1=x_1,\\
y_2=u_{21}x_1+x_2,\\
\vdots\\
y_n=u_{n1}x_1+\cdots+u_{n,n-1}x_{n-1}+x_n
\end{cases}
\tag{12}
\end{equation}
에 놓고, 부정원 \(u_{\mu\nu}\)들을 다항식의 계수영역에 첨가하자. 그러면 계수영역은 체 \(\mathsf{P}(u)\)로 바뀐다. 이 첨가로 \(\overline{\mideal}\)이 \(\overline{\mideal}_{(u)}\)로 바뀐다고 하자. 그러면 \(\overline{\mideal}_{(u)}\)은 \(\overline{\mideal}\)의 다항식 \(\Phi_\lambda(y)\)들과 함께, 매번 유한개의 \(\Phi_\lambda\)를 취한 모든 선형결합
\[
\sum_\lambda\alpha_\lambda(u)\Phi_\lambda(y)
\]
도 포함한다. 여기서 \(\alpha_\lambda(u)\)는 \(P\)의 계수를 갖는 \(u_{\mu\nu}\)들의 유리함수이다. 알맞은 \(u\)의 다항식을 곱하면 일반성을 잃지 않고 \(\alpha_\lambda(u)\)들을 \(u\)의 거듭제곱곱이라고 할 수 있다. 따라서
\begin{equation}
\sum_\lambda\alpha_\lambda(u)\Phi_\lambda(y)
\zero{\overline{\mideal}_{(u)}},
\qquad
\Phi_\lambda(y)\zero{\overline{\mideal}},
\quad\text{그리고 그 역도 성립한다.}
\tag{13}
\end{equation}
(12)에 따라 \(\overline{\mideal}_{(u)}\)은 \(\mathsf{P}(u)\)를 계수체로 갖는 \(x_1\ldots x_n\)의 다항식 아이디얼 \(\mideal\)로 옮겨진다.
\begin{equation}
\overline{\mideal}_{(u)}=\mideal
\qquad\text{단, }y=U(x).
\tag{14}
\end{equation}
관계 (14), (13)을 결합하면 다음을 얻는다.
\begin{equation}
\begin{gathered}
F(x)\zero{\mideal},\qquad
F(x)=\Phi(y)=\sum_i\alpha_i(u)\Phi_i(y)
=\sum_i\alpha_i(u)F_i(x)\\
\Longrightarrow\qquad F_i(x)\zero{\overline{\mideal}}.
\end{gathered}
\tag{15}
\end{equation}
반대로 (14), (15)에서 다시 (13)을 얻는다.

\medskip
\noindent\textbf{정의 IV.} \emph{\(\mathsf{P}(u)\)를 계수체로 갖는 \(x_1\ldots x_n\)의 다항식 아이디얼 \(\mideal\) 가운데 관계 (15)를 만족하는 것, 즉 \(y=U(x)\)에 따라 \(\overline{\mideal}_{(u)}\)에서 생기는 것을 변환된 아이디얼이라 한다.}

변환된 아이디얼의 특징은 정칙 다항식의 존재이다. \emph{\(r\)차 다항식이 영이 아닌 계수의 항 \(x_i^r\)을 포함하면 \(x_i\)에 관하여 정칙이라고 한다.} 다항식 \(F_\lambda(x)\)들이 \(x_1\)에 관하여 정칙임을 알 수 있다. 이하에서는 주로 이 변환된 아이디얼들을 \(x\)의 특정한 거듭제곱곱들에 관한 선형형식 모듈로 볼 것이다. 이를 위해 나중에 기초모듈로 옮겨갈 기초아이디얼의 개념을 정해야 한다. 먼저 이제부터 계속 쓸 표기를 정한다.

\medskip
\noindent\textbf{표기.} \emph{\(F^{(i)},f^{(i)},a^{(i)},\ldots\)는 언제나 \(x_1,\ldots,x_{i-1}\)과 무관한 \(x\)의 다항식을 뜻한다.}

\medskip
\noindent\textbf{정의 V.} \emph{각 아이디얼 \(\mideal\)에는 \(n\)개의 기초아이디얼 \(\gideal_0,\gideal_1,\ldots,\gideal_{n-1}\)이 다음처럼 정의된다. \((i-1)\)단계 기초아이디얼 \(\gideal_{i-1}\)은, 일반적으로 \(G(x)\)에 따라 달라지는 다항식 \(b^{(i)}\)가 존재하여
\begin{equation}
b^{(i)}G(x)\zero{\mideal},\qquad b^{(i)}\ne0
\tag{16}
\end{equation}
이 되는 다항식 \(G(x)\)만을, 그리고 그 모두를 포함한다.}

\(\gideal\)들이 실제로 아이디얼이라는 증명은 식 (4) 뒤에서 \(\Gmod\)의 모듈 성질을 증명한 것과 정확히 대응하며, (16)은 (4)의 유사물이다. 각 다항식 \(b^{(i+1)}\)은 동시에 \(b^{(i)}\)이므로 \(\gideal_i\)는 \(\gideal_{i-1}\)로 나누어진다.

기초아이디얼에 대해서는 다음이 성립한다.

\medskip
\noindent\textbf{정리 VI.} \emph{변환된 아이디얼의 기초아이디얼들도 변환된 아이디얼이다}\srcfn{12)}{정리 VI은 제7절에서 처음 쓴다.}.

증명은 잘 알려진 데데킨트--메르텐스 정리\srcfn{13)}{Dedekind, \emph{Über einen arithmetischen Satz von Gauß}, \emph{Mitt. d. deutsch. math. Ges. zu Prag} (1892); F. Mertens, \emph{Über einen algebraischen Satz}, \emph{Ber. d. Ak. d. Wissensch. Wien} 101 (1892), 1560--1566쪽. 크로네커가 가우스 정리를 부정계수 다항식으로 확장한 것은 이 정리의 직접적인 귀결이다.}에 바탕을 둔다. \(a,b\)가 부정원이고
\[
\begin{aligned}
&\sum a_{i_1\ldots i_s}z_1^{i_1}\cdots z_s^{i_s}\,
\sum b_{j_1\ldots j_s}z_1^{j_1}\cdots z_s^{j_s}\\
&\hspace{8em}=
\sum c(a,b)_{k_1\ldots k_s}z_1^{k_1}\ldots z_s^{k_s},
\end{aligned}
\]
\[
\sum i\leqq \nu
\sum j\leqq \mu
\sum k\leqq \nu+\mu
\]
라고 하자. 또 \(\Amod,\Bmod,\mathfrak C\)가 각각 \(a,b,c\)들로부터 유리정수에 의해 생성되는 모듈이라고 하자. 그러면 어떤 지수 \(q\)가 존재하여 항등식
\begin{equation}
\Amod^q\Bmod=\Amod^{q-1}\mathfrak C
\tag{17}
\end{equation}
이 성립한다. 여기서 \(\Amod^q\)는 \(a\)들의 \(q\)차 거듭제곱곱들과 그 정수계수 선형결합들로 이루어진다.

정리 VI을 증명하기 위해
\begin{equation}
\begin{cases}
G(x)=\Gamma(y)=\sum_\lambda\alpha_\lambda(u)\Gamma_\lambda(y)
=\sum_\lambda\alpha_\lambda(u)G_\lambda(x),\\
b^{(i)}(x)=\varphi(y)=\sum_\mu\beta_\mu(u)\varphi_\mu(y)
\end{cases}
\tag{18}
\end{equation}
로 두자. 여기서 \(\alpha_\lambda(u),\beta_\mu(u)\)는 \(u\)의 거듭제곱곱들이다. 그러면 (16), (14), (18)에 따라
\begin{equation}
\sum_\mu\beta_\mu(u)\varphi_\mu(y)\,
\sum_\lambda\alpha_\lambda(u)\Gamma_\lambda(y)
\zero{\overline{\mideal}_{(u)}},
\qquad
\sum_\mu\beta_\mu(u)\varphi_\mu(y)\ne0.
\tag{19}
\end{equation}
왼쪽은 \(u\)에 관한 두 다항식의 곱이고, 그 계수는 다항식 \(\varphi_\mu(y),\Gamma_\lambda(y)\)이다. (19)의 오른쪽에 따라 곱다항식의 \(u\)-계수들은 \(\overline{\mideal}\)에 속하는 다항식이다. 따라서 (17)에서 \(a,b,c\)를 이 \(y\)의 다항식들로 바꾸면
\[
\left\{\sum_\mu\beta_\mu(u)\varphi_\mu(y)\right\}^{q}
\Gamma_\lambda(y)\zero{\overline{\mideal}_{(u)}
}
\]
를 얻는다. (18)에 따라 이는
\begin{equation}
b^{(i)}(x)^q\cdot G_\lambda(x)\zero{\mideal},
\qquad
G_\lambda(x)\zero{\gideal_{i-1}}
\tag{20}
\end{equation}
과 같다. 관계 (16), (18), (20)은 \emph{기초아이디얼 \(\gideal_{i-1}\)들이 실제로 변환된 아이디얼}임을 보여 준다.

\section*{\S{}4. 다항식 아이디얼과 선형형식 모듈의 관계}

언제나 변환된 것이라고 가정하는 다항식 아이디얼 \(\mideal\)을 선형형식 모듈 \(\Mmod_{i-1}\), \((i=1\ldots n)\), 로 보기 위해 각 다항식 \(F(x)\)를 \(x_1,\ldots,x_{i-1}\)의 거듭제곱곱 \(\xi_\lambda\)들에 관한 선형형식으로 본다.
\[
F(x)=\sum_\lambda a_\lambda^{(i)}\xi_\lambda.
\]
제3절의 표기에 따라 \(a_\lambda^{(i)}\)는 \(x_i,\ldots,x_n\)의 다항식이다. 아이디얼의 정의에서 \(\Mmod_{i-1}\)이 이 정수적 원소 \(a^{(i)},b^{(i)},\ldots\)에 관하여 모듈 성질을 가짐이 곧바로 따라온다. 또한 \(x_1,\ldots,x_{i-1}\)의 무한히 많은 거듭제곱곱 \(\xi\)는 오직 \emph{선형}으로 나타나므로, 서로 선형독립이라는 사실 때문에 \(\Mmod_{i-1}\)에 대하여 부정원의 역할을 한다. 각 모듈 \(\Mmod_{i-1}\)은 무한히 많은 부정원 \(\xi\)를 포함하지만, 한 선형형식에는 그 가운데 유한개만 나타난다. 마찬가지로 각 기초아이디얼 \(\gideal_{i-1}\)을 선형형식 모듈 \(\Gmod_{i-1}\)로 보며, 부정원은 다시 \(x_1\ldots x_{i-1}\)의 거듭제곱곱이다. \(i=1\)일 때에는 단위원에 대응하는 부정원 \(\xi_0\) 하나만 있다.

이하의 모든 내용은 다음 정리에 따라, 이처럼 부정원 \(\xi\)가 무한히 많아도 유한개의 부정원에 관한 선형형식 모듈로 제한할 수 있다는 사실에 바탕을 둔다.

\medskip
\noindent\textbf{정리 VII.} \emph{잉여류계 \(\Gmod_{i-1}\mid\Mmod_{i-1}\)은 잉여류계 \(\Gmod^*_{i-1}\mid\Mmod^*_{i-1}\)과 동형이다. 여기서 \(\Mmod^*_{i-1}\)은 유한개의 부정원에 관한 모듈이고 \(\Gmod^*_{i-1}\)은 그 기초모듈이다. \(x_{i+1},\ldots,x_n\)을 \(\mathsf{P}(u)\)에 첨가하면 \(\Mmod_{i-1}\)은 \(1\)과 다른 초등인자를 유한개만 가지며, 이들은 \(\Mmod^*_{i-1}\)의 \(1\)과 다른 초등인자들이다.}

여기서 \(\Mmod_{i-1}\)의 초등인자는 주 8과 같이 정의한다. 정리 VII의 동형은 다음 정의에 따른다.

\medskip
\noindent\textbf{정의 V.} \emph{두 잉여류계를 일대일로 대응시킬 수 있고, 두 류의 차에는 대응하는 두 류의 차가 대응하며, 한 류와 한 정수적 원소의 곱에는 대응하는 류와 같은 정수적 원소의 곱이 대응하면 두 잉여류계를 동형이라고 한다.}

정리 VII은 완전귀납법으로 증명한다. \(i=1\)일 때에는 분명하다. \(\Mmod_0\)은 \(x\)의 0차 거듭제곱곱, 즉 단위원에 대응하는 부정원 \(\xi_0\) 하나에 관한 선형형식 모듈이고, 정수적 원소들은 \(x_1\ldots x_n\)의 모든 다항식이다. 기초아이디얼 \(\gideal_0\)은 \(x_1\ldots x_n\)의 모든 다항식으로 이루어진 단위원아이디얼이고, \(\Gmod_0\)은 \(\Mmod_0\)의 기초모듈인 모듈 \((\xi_0)\)과 같다. 따라서 \(\Gmod^*_0\mid\Mmod^*_0\)은 \(\Gmod_0\mid\Mmod_0\)과 같다. 끝으로 랭크 \(1\)인 \(\Mmod_0\)은 \(1\)과 다른 초등인자를 기껏해야 하나만 갖는다.

먼저 \(i=1\)에서 \(i=2\)로 넘어간다. 이때 \(\Gmod_1\mid\Mmod_1\)의 구조를 파악한 뒤 일반적인 귀납단계를 수행할 것이다. 우선 \(\Gmod_0\mid\Mmod_0\)에 대하여 \(x_1\)의 차수가 제한된 대표원계를 택할 수 있음을 보인다. 그러면 \(\gideal_1\)이 \(\gideal_0\)으로 나누어진다는 사실에 따라 \(\Gmod^*_1\)의 부정원은 유한개가 된다.

제3절에 따라 변환된 아이디얼 \(\mideal\)은 \(x_1\)에 관하여 정칙인 다항식 \(C^{(1)}(x)\)을 적어도 하나 포함한다. 그 차수를 \(k\)라 하자. 따라서 \(\Mmod_0\)은 선형형식 \(C^{(1)}\xi_0\)을 포함한다. \(C^{(1)}(x)\)의 정칙성 때문에 모든 다항식 \(H(x)\)은
\begin{equation}
H(x)\equiv b_0^{(2)}+b_1^{(2)}x_1+\cdots+
b_{k-1}^{(2)}x_1^{k-1}
\pmod{C^{(1)}(x)}
\tag{21}
\end{equation}
로 표현된다. 그러므로 \(C^{(1)}\xi_0\)이 \(\Mmod_0\)에 속한다는 사실에 따라, \(\Gmod_0\mid\Mmod_0\)의 대표원계를 \(x_1\)에 관한 차수가 \(k\)에 미치지 않는 것들로 실제로 택할 수 있다.

특히 \(\gideal_1\)의 다항식 \(G(x)\)와 \(\mideal\)의 다항식 \(F(x)\)에 대하여
\begin{equation}
\left\{
\begin{aligned}
G(x)&\equiv g_0^{(2)}+g_1^{(2)}x_1+\cdots+g_{k-1}^{(2)}x_1^{k-1}
&&\pmod{C^{(1)}(x)},\\
F(x)&\equiv a_0^{(2)}+a_1^{(2)}x_1+\cdots+a_{k-1}^{(2)}x_1^{k-1}
&&\pmod{C^{(1)}(x)}
\end{aligned}\right.
\tag{22}
\end{equation}
로 두자. \(\Gmod^*_1\)은 모든 선형형식
\[
g(\xi)=g_0^{(2)}\xi_0+\cdots+g_{k-1}^{(2)}\xi_{k-1}
\]
의 계, \(\Mmod^*_1\)은 이에 대응하는 모든
\[
a(\xi)=a_0^{(2)}\xi_0+\cdots+a_{k-1}^{(2)}\xi_{k-1}
\]
의 계라고 하자. \(\gideal_1,\mideal\)의 아이디얼 성질로부터 \(\Gmod^*_1,\Mmod^*_1\)은 정수적 원소 \(b^{(2)}, c^{(2)}\ldots\)에 관한, \(\xi_0,\ldots,\xi_{k-1}\)의 선형형식 모듈임이 따라온다. \(C^{(1)}(x)\)가 생성하는 주아이디얼도 이 정수적 원소들에 관한 모듈
\[
\mathfrak C_1=(\zeta_0,\zeta_1,\ldots,\zeta_\nu,\ldots)
\]
을 주며, 여기서 \(\zeta_\nu=x_1^\nu C^{(1)}(x)\)이다. \(\zeta_\nu\)들은 정수적 원소 \(b^{(2)},c^{(2)},\ldots\)에 관하여 서로 선형독립이고, 유한개의 \(\xi\)들과도 선형독립이다. \(\mathfrak C_1\)이 \(\Mmod_1\), 따라서 \(\Gmod_1\)로 나누어진다는 사실에서 \(\Gmod^*_1\)도 \(\Gmod_1\)로, \(\Mmod^*_1\)도 \(\Mmod_1\)로 각각 나누어짐이 따라온다. (22)를 함께 고려하면
\begin{equation}
\Gmod_1=(\Gmod^*_1,\mathfrak C_1),\qquad
\Mmod_1=(\Mmod^*_1,\mathfrak C_1)
\tag{23}
\end{equation}
을 얻는다.

(23)과 \(\xi,\zeta\)의 선형독립성에서 \(i=2\)에 대한 정리 VII이 다음처럼 따라온다. 먼저 \(\mathfrak C_1\)이 \(\Mmod_1\)로 나누어지므로
\[
\Gmod_1\mid\Mmod_1=\Gmod^*_1\mid\Mmod_1.
\]
\(\Gmod^*_1\mid\Mmod^*_1\)과의 동형을 보이기 위해, \(\Gmod^*_1\)의 어떤 선형형식 \(g^*(\xi)\)들이 \(\Gmod^*_1\mid\Mmod^*_1\)의 대표원계를 이룬다고 하자. \(\xi,\zeta\)의 선형독립성 때문에 \(g^*(\xi)\equiv0\pmod{\Mmod_1}\)이면 \(g^*(\xi)\equiv0\pmod{\Mmod^*_1}\)이다. 따라서 \(g^*(\xi)\)들은 \(\Mmod_1\)을 법으로도 서로 합동이 아니고, \(\Gmod^*_1\mid\Mmod_1\)의 대표원계를 이룬다. 이 대표원계를 통해 \(\Gmod^*_1\mid\Mmod_1\)과 \(\Gmod^*_1\mid\Mmod^*_1\) 사이에 생기는 대응은 정의 V의 의미에서 동형이다.

똑같이 \(b^{(2)}g(\xi)\equiv0(\Mmod_1)\), \(b^{(2)}\ne0\)이면 \(b^{(2)}g(\xi)\equiv0(\Mmod^*_1)\), \(b^{(2)}\ne0\)이다. 이 조건은 \(\Gmod^*_1\)에 속하는 모든 \(g(\xi)\)에 대하여, 그리고 오직 그 형식들에 대해서만 성립한다. (23)에 따라 \(\Gmod^*_1\)은 기초아이디얼 \(\gideal_1\)의 다항식 가운데 동시에 \(\xi\)의 선형형식인 것들 전체이기 때문이다. 따라서 \(\Gmod^*_1\)은 \(\Mmod^*_1\)의 기초모듈로 특징지어진다.

끝으로 \(x_3,\ldots,x_n\)을 \(\mathsf{P}(u)\)에 첨가하면
\begin{equation}
\left\{
\begin{aligned}
\Gmod^*_1&=(\eta_1\ldots\eta_\rho);&
\Mmod^*_1&=(e_1\eta_1,\ldots,e_\rho\eta_\rho),\\
\Gmod_1&=(\eta_\rho\ldots\eta_1,\zeta_0\ldots\zeta_\nu\ldots);&
\Mmod_1&=(e_\rho\eta_\rho,\ldots,e_1\eta_1,\zeta_0\ldots\zeta_\nu\ldots)
\end{aligned}\right.
\tag{24}
\end{equation}
이고, 따라서
\[
(e_\rho)=\Mmod_1/\Gmod_1
=\Mmod_1/(\Mmod_1,\eta_\rho),\qquad
(e_{\rho-1})=(\Mmod_1,\eta_\rho)/\Gmod_1,\quad\text{등등}
\]
이다. 주 8까지 고려하면 \(\Mmod_1\)의 초등인자는
\[
e_\rho,e_{\rho-1},\ldots,e_1,1,\ldots,1,\ldots
\]
임을 알 수 있다. 이로써 \(i=2\)에 대한 정리 VII의 모든 부분이 증명되었다.

(22), (23)은 \(C^{(1)}(x)\)가 \(x_1\)에 관하여 정칙이라는 사실만을 썼다는 점을 주목하자. 따라서 (12) 대신 특수한 변환
\begin{equation}
y_1=x_1,\qquad
y_2=u_{21}x_1+z_2,\quad\ldots,\quad
y_n=u_{n1}x_1+z_n
\tag{25}
\end{equation}
을 기초로 해도 이 식들과 정리 VII로 이어지는 논증은 그대로 성립한다. 이를
\begin{equation}
\left\{
\begin{array}{l}
z=U_1(x),\quad\text{또는}\\[1mm]
z_2=x_2,\quad z_3=u_{32}x_2+x_3,\quad\ldots,\quad
z_n=u_{n2}x_2+\cdots+u_{n,n-1}x_{n-1}+x_n
\end{array}\right.
\tag{26}
\end{equation}
과 합성하면 다시 (12)를 얻는다. \(\overline{\mideal}_{(u)}\)가 (25)에 따라 \(\widetilde{\mideal}\)로 옮겨지고, \(\widetilde{\gideal}_1,\widetilde{\Gmod}_1\ldots\)가 \(\widetilde{\mideal}\)에 대하여 \(\gideal_1,\Gmod_1\ldots\)과 같은 뜻을 갖는다면
\begin{equation}
\widetilde{\Gmod}_1=(\widetilde{\Gmod}^{*}_1,\widetilde{\mathfrak C}_1),
\qquad
\widetilde{\Mmod}_1=(\widetilde{\Mmod}^{*}_1,\widetilde{\mathfrak C}_1)
\tag{27}
\end{equation}
이다.

(27)은 (26)의 역변환에 따라 (23)에서 단순히 생긴다. 실제로 이는 \(\mideal,C^{(1)}(x)\), 따라서 \(\mathfrak C_1,\Mmod_1\)에 대해서 분명히 성립한다. \(\gideal_1\)에 대해서도 성립한다. 관계
\[
b^{(2)}(x)G(x)\equiv0\pmod{\mideal},\quad b^{(2)}\ne0,
\]
\[
\widetilde b^{(2)}(z)\widetilde G(x,z)
\equiv0\pmod{\widetilde\mideal},\quad
\widetilde b^{(2)}\ne0
\]
은 \(z=U_1(x)\)에 따라 서로를 함의하기 때문이다. 그러므로 (26)에 따라 \(\widetilde\gideal_1=\gideal_1\), 따라서 \(\widetilde\Gmod_1=\Gmod_1\)이다. (22)가 주는 \(\Gmod_1^*,\Mmod_1^*\)의 정의에 따라 \(\widetilde\Gmod_1^*,\widetilde\Mmod_1^*\), 따라서 표현 (27) 전체에도 같은 사실이 성립한다.

일반적인 귀납단계를 위해 다음을 가정하자.
\begin{equation}
\left\{
\begin{aligned}
\Gmod_{i-1}&=(\Gmod_{i-1}^*,\mathfrak C_{i-1}),&
\Mmod_{i-1}&=(\Mmod_{i-1}^*,\mathfrak C_{i-1}),\\
\mathfrak C_{i-1}&=(\zeta_0,\zeta_1,\ldots,\zeta_\nu,\ldots).
\end{aligned}\right.
\tag{28}
\end{equation}
여기서 \(\Gmod^*_{i-1},\Mmod^*_{i-1}\)은 정수적 원소 \(b^{(i)},c^{(i)},\ldots\)에 관한, 유한개의 부정원 \(\xi\)의 선형형식 모듈이다. 이 \(\xi\)들은 \(x_1,\ldots,x_{i-1}\)의 특정한 거듭제곱곱들이다. 또한 \(\zeta\)들은 이 정수적 원소들에 관하여 서로 선형독립이고 \(\xi\)들과도 선형독립이라고 하자. (12) 대신 특수한 변환
\begin{equation}
\begin{aligned}
y_1&=x_1,\quad\ldots,\quad
y_{i-1}=u_{i-1,1}x_1+\cdots+x_{i-1},\\
y_i&=u_{i,1}x_1+\cdots+u_{i,i-1}x_{i-1}+z_i,\quad\ldots,\quad
y_n=u_{n,1}x_1+\cdots+u_{n,i-1}x_{i-1}+z_n
\end{aligned}
\tag{29}
\end{equation}
을 기초로 할 때에도 (28)에 대응하는 표현과 \(\xi,\zeta\)의 선형독립성이 성립한다고 하자. 이 변환은
\[
z=U_{i-1}(x),
\]
즉
\[
z_i=x_i;\quad z_{i+1}=u_{i+1,i}x_i+x_{i+1};\quad\ldots;\quad
z_n=u_{n,i}x_i+\cdots+u_{n,n-1}x_{n-1}+x_n
\]
을 통해 (12)와 합성할 수 있다. 그리고 이 특수한 표현은 \(z=U_{i-1}(x)\)의 역변환에 따라 (28)에서 단순히 생긴다고 하자.

가정 (28)과 \(\xi,\zeta\)의 선형독립성에서, 위에서 (23)에 이어 쓴 것과 똑같은 논증으로 같은 대표원계를 대응시켜 \(\Gmod_{i-1}\mid\Mmod_{i-1}\)과 \(\Gmod^*_{i-1}\mid\Mmod^*_{i-1}\)의 동형을 얻는다. 마찬가지로 \(\Gmod^*_{i-1}\)이 \(\Mmod^*_{i-1}\)의 기초모듈이고, \(\Mmod_{i-1}\)이 \(1\)과 다른 초등인자를 유한개만 가지며 그것들이 \(\Mmod^*_{i-1}\)의 \(1\)과 다른 초등인자들임을 얻는다.

귀납단계를 수행하려면 먼저 \(\Gmod^*_{i-1}\mid\Mmod^*_{i-1}\), 따라서 \(\Gmod_{i-1}\mid\Mmod_{i-1}\)에 대하여 \(x_i\)의 차수가 제한된 대표원계를 택할 수 있음을 다시 보여야 한다. 이는 정리 III의 (7)에서 증명한 몫표현
\begin{equation}
\Gmod^*_{i-1}=\Mmod^*_{i-1}/(C^{(i)})
\tag{30}
\end{equation}
에서 따라온다. 여기서 \(C^{(i)}\)는 \(\Mmod^*_{i-1}\)에서 얻는 항등적으로 영이 아닌 임의의 \(\rho\)차 행렬식이고, \(\rho\)는 \(\Mmod^*_{i-1}\)의 랭크이다. 특수한 변환 (29)에 관한 가정에서 \(C^{(i)}\)를 언제나 \(x_i\)에 관하여 정칙으로 택할 수 있다는 것이 따라온다. 실제로 \(\Mmod^*_{i-1}\)에 대응하는 모듈 \(\widetilde\Mmod^*_{i-1}\)은 같은 랭크를 갖는다. \(\widetilde C^{(i)}\)가 \(\widetilde\Mmod^*_{i-1}\)에서 얻는 항등적으로 영이 아닌 임의의 \(\rho\)차 행렬식이고 \(z=U_{i-1}(x)\)에 따라 정칙인 \(C^{(i)}\)로 옮겨진다면, 가정에 따라 \(C^{(i)}\)는 \(\Mmod^*_{i-1}\)에서 얻는 \(\rho\)차 행렬식이다.

\(C^{(i)}\)가 존재하면 (8)에 대응하여 \(\xi\)들의 번호를 알맞게 정했을 때 특별한 모양의 \(\rho\)개 선형형식
\[
L_\lambda(\xi)
=C^{(i)}\xi_\lambda+c_{\lambda,1}^{(i)}\xi_{\rho+1}
+\cdots+c_{\lambda,s-\rho}^{(i)}\xi_s,
\qquad(\lambda=1,\ldots,\rho)
\]
이 \(\Mmod^*_{i-1}\)에 속한다. 이제 \(\xi\)의 모든 선형형식은
\begin{equation}
\begin{aligned}
H(\xi)={}&a_1^{(i)}\xi_1+\cdots+a_\rho^{(i)}\xi_\rho\\
&+b_1^{(i)}\xi_{\rho+1}+\cdots+b_{s-\rho}^{(i)}\xi_s
\pmod{L_1(\xi),\ldots,L_\rho(\xi)}
\end{aligned}
\tag{31}
\end{equation}
로 만들 수 있다. 여기서 \(a_1^{(i)},\ldots,a_\rho^{(i)}\)는 \(x_i\)의 차수가 제한되어 \(C^{(i)}\)의 차수 \(k\)에 미치지 않는다. 이 표현은 유일하다. 실제로 왼쪽이 \(L_1(\xi),\ldots,L_\rho(\xi)\)를 법으로 영이면, \(x_i\)의 차수를 고려하면서 \(\xi_1,\ldots,\xi_\rho\)의 계수를 비교하여 모든 \(a^{(i)},b^{(i)}\)가 항등적으로 사라짐을 얻는다.

특히
\[
H(\xi)\equiv0\pmod{\Gmod^*_{i-1}},
\qquad\text{따라서 (30)에 따라}\qquad
C^{(i)}H(\xi)\equiv0\pmod{\Mmod^*_{i-1}}
\]
라고 하자. 정리 III의 증명에서처럼
\[
C^{(i)}H(\xi)
-\sum_{\tau=1}^{s-\rho}
\left\{C^{(i)}b_\tau^{(i)}
-\sum_\lambda a_\lambda^{(i)}c_{\lambda,\tau}^{(i)}\right\}
\xi_{\rho+\tau}
\equiv0\pmod{\Mmod^*_{i-1}}
\]
을 얻는다. 정리 III에서와 같이 \(\xi_{\rho+1},\ldots,\xi_s\)의 계수들이 사라져야 하므로
\[
C^{(i)}b_\tau^{(i)}
=\sum_\lambda a_\lambda^{(i)}c_{\lambda,\tau}^{(i)},
\qquad(\tau=1,\ldots,s-\rho).
\]
따라서 \(b_\tau^{(i)}\)들의 차수도 제한되며 \(c_{\lambda,\tau}^{(i)}\)들의 최대차수에 미치지 않는다. 이로써 \(\Gmod^*_{i-1}\mid\Mmod^*_{i-1}\), 즉 \(\Gmod_{i-1}\mid\Mmod_{i-1}\)에 대하여 \(x_i\)의 차수가 어떤 고정된 차수 \(r\)에 미치지 않는 대표원계가 존재함이 증명되었다.

이제 유한개의 거듭제곱곱
\[
x_i^\alpha\xi_\lambda
\quad(\alpha=0,1,\ldots,k-1;\ \lambda=1,2,\ldots,\rho),
\]
\[
x_i^\beta\xi_{\rho+\tau}
\quad(\beta=0,1,\ldots,r-1;\ \tau=1,2,\ldots,s-\rho)
\]
를 \(\vartheta_1,\ldots,\vartheta_t\)라 하자. 또한 가산무한개의 식
\[
x_i^\gamma L_\lambda
\quad(\gamma=0,1\ldots\ldots;\ \lambda=1,2,\ldots,\rho),
\qquad
x_i^\gamma\zeta_\nu
\quad(\gamma,\nu=0,1\ldots\ldots)
\]
을 단순무한열 \(\omega_0,\omega_1,\ldots,\omega_\mu\ldots\)로 배열하자. 그러면 \(\vartheta\)들과 \(\omega\)들은 정수적 원소 \(b^{(i+1)},c^{(i+1)},\ldots\)에 관하여 각 계 안에서도, 두 계 사이에서도 선형독립이다. 실제로 각 합이 유한개의 항에만 걸치는 관계
\[
\begin{aligned}
&\sum c_{\alpha\lambda}^{(i+1)}x_i^\alpha\xi_\lambda
+\sum c_{\beta\tau}^{(i+1)}x_i^\beta\xi_{\rho+\tau}\\
&\quad+\sum d_{\gamma\lambda}^{(i+1)}x_i^\gamma L_\lambda(\xi)
+\sum e_{\gamma\nu}^{(i+1)}x_i^\gamma\zeta_\nu=0
\end{aligned}
\]
에서, 가정한 \(\xi,\zeta\)의 선형독립성과 정수적 원소 \(b^{(i)},\ldots\)에 관한 \(\xi\)들의 상호 선형독립성 때문에 모든 \(e_{\gamma\nu}^{(i+1)}\)가 사라진다. 표현 (31)의 유일성에서 모든 \(c_{\alpha\lambda}^{(i+1)},c_{\beta\tau}^{(i+1)}\)가 사라지고, 끝으로 \(L_\lambda(\xi)\)들의 선형독립성에서 모든 \(d_{\gamma\lambda}^{(i+1)}\)가 사라진다.

이제 모듈 \((\mathfrak C_{i-1},L_1(\xi),\ldots,L_\rho(\xi))\)을 정수적 원소 \(b^{(i+1)}\)에 관한 모듈 \(\mathfrak C_i\)로 보자. \(\mathfrak C_{i-1},L_1(\xi),\ldots,L_\rho(\xi)\)가 \(\Mmod_{i-1}\)로 나누어지므로 \(\mathfrak C_i\)는 \(\Mmod_i\)로 나누어지고,
\[
\mathfrak C_i=(\omega_0,\omega_1,\ldots,\omega_\mu\ldots)
\]
이다. 각 \(H(x)\equiv0\pmod{\gideal_{i-1}}\)에 대하여 (28), (31)과 방금 증명한 결과에서
\[
H(x)\equiv
b_1^{(i+1)}\vartheta_1+\cdots+b_t^{(i+1)}\vartheta_t
\pmod{\mathfrak C_i}
\]
라는, 정수적 원소 \(b^{(i+1)},c^{(i+1)}\)에 관한 모듈등식을 얻는다. 한편 \(\gideal_i\)는 \(\gideal_{i-1}\)로 나누어지고 \(\mideal\)은 \(\gideal_i\)로 나누어진다. 특히 \(\mathfrak C_i\)의 각 \(G(x)\)와 \(\Mmod_i\)의 각 \(F(x)\)에 대하여
\[
G(x)\equiv g_1^{(i+1)}\vartheta_1+\cdots+
g_t^{(i+1)}\vartheta_t\pmod{\mathfrak C_i},
\]
\[
F(x)\equiv a_1^{(i+1)}\vartheta_1+\cdots+
a_t^{(i+1)}\vartheta_t\pmod{\mathfrak C_i}
\]
로 두자. 모든 \(g(\vartheta)\)의 모듈을 \(\Gmod_i^*\), 모든 \(a(\vartheta)\)의 모듈을 \(\Mmod_i^*\)라 하면, (23)으로 이끈 것과 정확히 같은 논증으로
\begin{equation}
\Gmod_i=(\Gmod_i^*,\mathfrak C_i),\qquad
\Mmod_i=(\Mmod_i^*,\mathfrak C_i),\qquad
\mathfrak C_i=(\omega_0,\omega_1,\ldots,\omega_\mu\ldots)
\tag{32}
\end{equation}
을 얻는다.

(32)를 유도할 때에도 \(C^{(i)}\)가 \(x_i\)에 관하여 정칙이라는 사실만을 썼다. 따라서 지표를 하나 올려 특수한 변환 (29)을 기초로 해도 논증 전체가 유지된다. (27)에 이어진 것과 정확히 같은 논증은 이 특수한 표현이 \(z=U_i(x)\)의 역변환에 따라 (32)에서 단순히 생김을 보여 준다.

이로써 일반적인 귀납단계의 가정 (28) 등이 한 단계 높은 지표에 대하여 모두 증명되었다. 앞에서 보았듯 이 가정들에서 정리 VII이 직접 따라오므로, 정리 VII은 일반적으로 증명되었다. 동시에 \(C^{(i)}\)들의 선택에 따른 \(\Gmod^*_{i-1},\Mmod^*_{i-1}\)의 임의성은 잉여류계 \(\Gmod^*_{i-1}\mid\Mmod^*_{i-1}\)의 동형을 통해 다시 사라진다는 것도 보였다.

특히 제3절에서 계수체로 삼은 \(\mathsf{P}\)의 표수가 영, 즉 \(\mathsf{P}\)에 포함되고 단위원으로 생성되는 소체가 유리수체의 유형이면, 각 \(i=1,\ldots,n\)에 대하여 \(C^{(i)}\)가 \(x_i\)에 관하여 정칙으로 남도록 부정원 \(u_{\mu\nu}\)들을 \(P\)의 원소 \(\bar u_{\mu\nu}\)로 특수화할 수 있다. 위의 논증은 이 특수화 아래에서도 정리 VII이 유지됨을 보여 준다\srcfn{14)}{헨첼은 임의의 \(\mathsf{P}\) 대신 복소수체만을 택하고 주로 이 마지막 경우를 고찰했으며, 본래의 소거론에서만 부정원들을 기초로 삼았다. 그는 이어서, 본문에서 제시한 추론들을 본질적으로 사용하여 종결식형을 구성할 때 각 \(\Mmod_{i-1}\)에서 \(x\)에 관하여 어떤 고정된 경계를 넘는 유한한 차수까지만 진행하면 충분함을 보였다. 그러나 정리 VII이 주는 \(\Gmod_{i-1}\mid\Mmod_{i-1}\)과 \(\Gmod^*_{i-1}\mid\Mmod^*_{i-1}\)의 동형이 빠져 있어서, 차수들에 대한 이 독립성의 근거는 드러나지 않는다. (E. N.)}. 다만 기초모듈과 초등인자들은 일반적인 경우를 \(u_{\mu\nu}=\bar u_{\mu\nu}\)로 특수화하여 반드시 생길 필요는 없다\srcfn{15)}{\(\mideal=(x^2,ux+y)\)이면 \(\gideal_0=(1)\), \(\gideal_1=(1)\)이다. \(C^{(1)}=x^2\)을 택하면 \(\Gmod_1^*=(1,x)\), \(\Mmod_1^*=(ux+y,xy)\)이고, \(\Mmod_1^*\)의 초등인자는 \(y^2,1\)이며 \(C^{(2)}=y^2\)으로 택할 수 있다. \(u=0\)일 때에도 \(C^{(1)},C^{(2)}\)은 각각 \(x,y\)에 관하여 정칙으로 남지만, \(\widetilde{\Mmod}_1^*=(y,xy)\)의 초등인자는 \(y,y\)이다. 따라서 이 경우에도 \(\widetilde{\Gmod}_1\mid\widetilde{\Mmod}_1\)은 유한모듈의 잉여류계와 동형이지만 일반적인 경우의 잉여류계와 동형은 아니다. 반대로 \(C^{(1)}=ux+y\)를 택하고 \(C^{(1)}\)이 정칙으로 남도록 특수화했다면 동형도 유지되었을 것이다. (E. N.)}.

\section*{\S{}5. 다항식 아이디얼의 종결식형과 초등인자형}

\(x_{i+1},\ldots,x_n\)을 \(\mathsf{P}(u)\)에 첨가하자. 그러면 \(\Gmod_{i-1},\Mmod_{i-1}\), 따라서 \(\Gmod^*_{i-1},\Mmod^*_{i-1}\)도 정수적 원소들을 한 부정원 \(x_i\)의 다항식으로 볼 수 있는 선형형식 모듈이 된다. 정리 VII에 따라 이 경우 \(\Mmod_{i-1}\)의 \(1\)과 다른 초등인자들과 기껏해야 유한개의 초등인자 \(1\)은 \(\Mmod^*_{i-1}\)의 초등인자들과 일치한다. 이 초등인자들의 곱, 즉 \(\Mmod^*_{i-1}\)의 \(\rho\)번째 행렬식인자는 \(\Gmod^*_{i-1}\)에서 \(\Mmod^*_{i-1}\)로 넘어가는 치환의 행렬식과 같으며, 정의 III에 따라 \(\Norm(\Gmod^*_{i-1}\mid\Mmod^*_{i-1})\)이다. (24), 또는 일반적인 \(i\)에 대하여 (28)에서 생기는 유사한 공식은 이 초등인자들의 곱이 \(\Gmod_{i-1}\)에서 \(\Mmod_{i-1}\)로 넘어가는 치환의 행렬식과도 같음을 보여 준다. 따라서 이를 \(N(\Gmod_{i-1}\mid\Mmod_{i-1})\)이라 써야 한다. 즉
\[
N(\Gmod_{i-1}\mid\Mmod_{i-1})
=N(\Gmod^*_{i-1}\mid\Mmod^*_{i-1})
=R^{(i)}(x).
\]
다항식 \(R^{(i)}(x)\)은 \(\Mmod^*_{i-1}\)에서 얻는 모든 \(\rho\)차 행렬식의 다항식 의미의 최대공약원이다. 이를 \(x_{i+1},\ldots,x_n\)에 관하여 정수적이고 원시적인 다항식이라고 가정할 수 있으므로, \(C^{(i)}\)의 인자로서 \(x_i\)에 관하여 정칙이다. 마찬가지로 \(\Mmod_{i-1}\), 또는 \(\Mmod^*_{i-1}\)의 가장 높은 초등인자 \(E^{(i)}(x)\)도 \(x_{i+1},\ldots,x_n\)에 관하여 정수적이고 원시적이며, 따라서 \(x_i\)에 관하여 정칙이라고 가정할 수 있다. 이로써 다음 정의에 이른다.

\medskip
\noindent\textbf{정의 VI.} \emph{다항식 \(R^{(i)}(x)\)을 \(\mideal\)의 \(i\)단계 종결식이라 하고 \(E^{(i)}(x)\)을 \(i\)단계 초등인자라 한다. 따라서 \(i\)단계 종결식은 기초아이디얼 \(\gideal_{i-1}\)의 \(\mideal\)에 대한 노름, 즉 \(x_{i+1},\ldots,x_n\)을 \(\mathsf{P}(u)\)에 첨가했을 때 모듈 \(\Gmod_{i-1}\)의 \(\Mmod_{i-1}\)에 대한 노름이다. 같은 첨가 아래 \((E^{(i)}(x))\)은 몫 \(\Mmod_{i-1}/\Gmod_{i-1}\)과 같다. 곱
\[
R_\mideal=R^{(1)}\cdot R^{(2)}\cdot\ldots\cdot R^{(n)};\qquad
E_\mideal=E^{(1)}\cdot E^{(2)}\cdot\ldots\cdot E^{(n)}
\]
을 각각 \(\mideal\)의 종결식형과 초등인자형이라 한다.}

종결식형과 초등인자형에 대해서는 다음이 성립한다.

\medskip
\noindent\textbf{정리 VIII.} \emph{종결식형은 초등인자형으로 나누어지고, 거꾸로 \(E_\mideal\)의 어떤 거듭제곱은 \(R_\mideal\)로 나누어진다. \(E_\mideal,R_\mideal\)은 \(\mideal\)로 나누어진다. 더 일반적으로
\[
E^{(i)}\cdots E^{(n)}\gideal_{i-1}\zero{\mideal},
\]
따라서
\[
R^{(i)}\cdots R^{(n)}\gideal_{i-1}\zero{\mideal}.
\]}

노름은 가장 높은 초등인자로 나누어지고, 거꾸로 이 초등인자의 어떤 거듭제곱은 노름으로 나누어진다. 따라서 \(x_{i+1},\ldots,x_n\)을 첨가했을 때 \(R^{(i)}(x),E^{(i)}(x)\)에 대하여 이 나눗셈이 성립한다. 그런데 \(R^{(i)}(x),E^{(i)}(x)\)을 \(x_{i+1},\ldots,x_n\)에 관한 원시다항식으로 가정했으므로, 모든 부정원 \(x_i,x_{i+1},\ldots,x_n\)에 관하여도 이 나눗셈이 성립한다. 따라서 이 표현들의 정의에 따라 모든 부정원 \(x\)에 관하여 \(R_\mideal\)은 \(E_\mideal\)로 나누어지고 \(E_\mideal\)의 어떤 거듭제곱은 \(R_\mideal\)로 나누어진다.

정리 VII과 정리 II에 따라 \(x_{i+1},\ldots,x_n\)을 첨가했을 때 가장 높은 초등인자 \(E^{(i)}(x)\)은 몫 \(\Mmod_{i-1}/\Gmod_{i-1}\)으로 정의된다. 따라서 다시 모든 부정원 \(x\)의 다항식으로 돌아오면 각
\begin{equation}
G(x)\zero{\gideal_{i-1}}
\quad\text{에 대하여 어떤}\quad b^{(i+1)}\ne0
\quad\text{가 존재하여}\quad
b^{(i+1)}E^{(i)}(x)G(x)\zero{\mideal}
\tag{33}
\end{equation}
이 된다. 특히 \(\gideal_0\)은 단위원아이디얼과 같으므로
\[
b^{(2)}E^{(1)}(x)\zero{\mideal};\quad b^{(2)}\ne0
\qquad\text{따라서}\qquad
E^{(1)}(x)\zero{\gideal_1}.
\]
일반적으로 (33)에 따라
\[
E^{(1)}(x)\cdots E^{(i-1)}(x)\zero{\gideal_{i-1}}
\]
이면 어떤 \(b^{(i+1)}\ne0\)가 존재하여
\[
b^{(i+1)}E^{(1)}\cdots E^{(i)}\zero{\mideal}
\]
이고, 따라서
\[
E^{(1)}\cdots E^{(i)}\zero{\gideal_i}.
\]
\(\gideal_n=\mideal\)이므로
\begin{equation}
E_\mideal=E^{(1)}\cdots E^{(n)}\equiv0(\mideal),
\qquad\text{따라서}\qquad
R_\mideal=R^{(1)}\cdots R^{(n)}\equiv0(\mideal).
\tag{34}
\end{equation}

(33)의 \(G(x)\)가 \(\gideal_{i-1}\)의 한 아이디얼기저에 속하는 유한개의 다항식을 돌게 하고, \(c^{(i+1)}(x)\)를 이에 대응하는 \(b^{(i+1)}(x)\)들의 곱이라고 하자. 그러면
\[
c^{(i+1)}E^{(i)}(x)\gideal_{i-1}\equiv0(\mideal),
\qquad c^{(i+1)}\ne0,
\]
따라서
\[
E^{(i)}(x)\gideal_{i-1}\zero{\gideal_i}.
\]
이를 위와 같이 유한번 반복하면
\[
E^{(n)}(x)\cdots E^{(i)}(x)\gideal_{i-1}\zero{\mideal}
\]
이고, 따라서
\[
R^{(n)}(x)\cdots R^{(i)}(x)\gideal_{i-1}\zero{\mideal}.
\]
이로써 정리 VIII의 모든 부분이 증명되었다.

정리 VIII은 본질적으로 초등인자형의 성질에 바탕을 둔다. 그러나 종결식형이 \(\mideal\)에 대하여 특징적인 의미를 갖는다는 것은 다음 정리가 보여 준다.

\medskip
\noindent\textbf{정리 IX.} \emph{\(\nideal\)이 아이디얼의 의미에서 \(\mideal\)의 인자이고 종결식형 \(R_\nideal,R_\mideal\)이 같으면 아이디얼 \(\nideal,\mideal\)도 같다.}

\(\gideal_0,\ldots,\gideal_{n-1},\gideal_n=\mideal\)과 \(\hideal_0,\ldots,\hideal_{n-1},\hideal_n=\nideal\)이 각각 \(\mideal,\nideal\)의 기초아이디얼들이라고 하자. 먼저 \(\gideal_0,\hideal_0\)은 둘 다 단위원아이디얼이므로 \(\gideal_0=\hideal_0\)이다. 이제 \(\gideal_{i-1}=\hideal_{i-1}\)이라고 가정하자. 그러면 \(\Mmod_{i-1},\Nmod_{i-1}\)은 같은 기초모듈 \(\Gmod_{i-1}\)을 갖는다. 가정 \(R_\nideal=R_\mideal\)은 \(x_{i+1},\ldots,x_n\)을 첨가했을 때
\[
N(\Gmod_{i-1}\mid\Nmod_{i-1})
=N(\Gmod_{i-1}\mid\Mmod_{i-1}),
\]
또는
\[
\Norm(\Gmod^*_{i-1}\mid\Nmod^*_{i-1})
=N(\Gmod^*_{i-1}\mid\Mmod^*_{i-1})
\]
이라고 쓸 수 있다. 여기서 (32)에 대응하여
\[
\Nmod_{i-1}=(\Nmod^*_{i-1},\mathfrak C_{i-1})
\]
로 두었다. 따라서 \(\Nmod^*_{i-1}\)도 \(\xi\)의 선형형식 모듈이고 \(\Gmod^*_{i-1}\)은 그 기초모듈이다. \(\Mmod_{i-1}\)이 \(\Nmod_{i-1}\)로 나누어지므로 \(\Mmod^*_{i-1}\)도 \(\Nmod^*_{i-1}\)로 나누어진다. 따라서 정리 IV에 따라 이 첨가 아래 두 모듈 \(\Mmod^*_{i-1},\Nmod^*_{i-1}\), 그러므로 \(\Nmod_{i-1},\Mmod_{i-1}\)도 일치한다. 한편 \(x_{i+1},\ldots,x_n\)의 첨가는 \(\Mmod_{i-1}\), 또는 \(\mideal\)에
\[
b^{(i+1)}(x)G(x)\zero{\mideal},\qquad b^{(i+1)}\ne0
\]
이 되는 모든 다항식 \(G(x)\)를 더하는 것을 뜻한다. 즉 이 첨가로 \(\Mmod_{i-1}\), 또는 \(\mideal\)은 \(\gideal_i\)로 옮겨지고, 마찬가지로 \(\nideal\)은 \(\hideal_i\)로 옮겨진다. 따라서 \(\gideal_i=\hideal_i\)이고, 끝내 \(\gideal_n=\hideal_n\), 즉 \(\mideal=\nideal\)이다.

끝으로 같은 이행원리를 정리 V에 적용하면 다음을 얻는다.

\medskip
\noindent\textbf{정리 X.} \emph{\(R^{(i)}=S^{(i)}T^{(i)}\)가 \(R^{(i)}\)을 다항식의 의미에서 서로소인 두 다항식 \(S^{(i)},T^{(i)}\)으로 분해한 것이라고 하자. 그러면 \(\mideal\)의 인자이고 같은 \((i-1)\)단계 기초아이디얼 \(\gideal_{i-1}\)을 갖는 아이디얼 \(\mathfrak s,\tideal\)이 각각 유일하게 존재하여
\[
S^{(i)}=\Norm(\Gmod_{i-1}\mid\Smod_{i-1}),\qquad
T^{(i)}=\Norm(\Gmod_{i-1}\mid\Tmod_{i-1})
\]
이 된다. \(\mathfrak s,\tideal\)은 각각
\[
s^{(i+1)}(x)T^{(i)}(x)S(x)\zero{\mideal},
\qquad s^{(i+1)}\ne0,
\]
\[
t^{(i+1)}(x)S^{(i)}(x)T(x)\zero{\mideal},
\qquad t^{(i+1)}\ne0
\]
이 되는 다항식 \(S(x)\), 또는 \(T(x)\) 전체로 정의된다. 또한 \(\gideal_i\)는 \(\mathfrak s,\tideal\)의 최소공배아이디얼이다.
\[
\gideal_i=[\mathfrak s,\tideal].
\]}

여기서는 \(\mathfrak s,\tideal\)의 기저원소들에 대응하는 \(s^{(i+1)},t^{(i+1)}\)들의 곱을 취하여, \(s^{(i+1)},t^{(i+1)}\)을 각 아이디얼에 대하여 고정해도 된다는 점만 덧붙이면 된다. 또한 앞에서 보았듯 \(x_{i+1},\ldots,x_n\)을 첨가하면 \(\mideal\)은 \(\gideal_i\)로 바뀐다. 특히 \(R^{(i)}\)의 분해를 기약다항식들의 거듭제곱, 즉 으뜸인자들까지 계속할 수 있으며, 그 결과 \(\gideal_i\)도 이에 대응하는 최소공배아이디얼 표현을 갖는다.

\section*{\S{}6. 본래의 소거론}

정리 VIII, 즉 식 (34)에 따르면 종결식형과 초등인자형은 \(\mideal\)로 나누어지므로 \(\mideal\)의 모든 영점에서 사라진다. 여기서 \glqq{}\(\mideal\)의 영점\grqq{}이란 다음 값의 계를 뜻한다. \(x_{i+1},\ldots,x_n\)을 계수영역에 첨가했다고 생각하고, \(P(u;x_{i+1},\ldots,x_n)\)에서 유도되는 대수적으로 닫힌 체에 속하는 \(x_1,\ldots,x_i\)의 값의 계 가운데 \(\mideal\)의 모든 다항식이 사라지는 것들이다. 이는 \(i=1,\ldots,n\) 각각에 대하여 말한다\srcfn{16)}{모든 체가 본질적으로 유일한 방식으로 대수적으로 닫힌 체까지 확대될 수 있다는 사실은 Steinitz, \emph{Algebraische Theorie der Körper}, \emph{J. f. M.} 137 (1910), 167--309쪽에서 증명되었고, 이로써 대수학의 기본정리에 대한 유리적 등가물이 주어졌다. 실제로 각 \(i=1,\ldots,n\)에 대해서는 \(P(u,x_{i+1},\ldots,x_n)\)의 유한확대체만 다룬다. (E. N.)}. 뒤에서 보이듯 첨가한 \(x_{i+1},\ldots,x_n\)을 \(P(u)\)의 원소들로 바꿀 수도 있다. 거꾸로 종결식형의 영점에서 \(\mideal\)의 영점들을 유도하는 것이 소거론의 내용이다. 이 역방향은 이 절에서 전개할 연속소거에 바탕을 둔다. 다음 절에서 그 과정에 나타나는 형식 \(D^{(i)}(x)\)이 \(E^{(i)}(x)\)으로 나누어짐을 보이면, \(E^{(i)}(x)\)의 어떤 거듭제곱이 \(R^{(i)}(x)\)으로 나누어진다는 사실에서 구하는 역방향이 따라온다.

여기서 제시할 연속소거 이론은 다음 정리에 바탕을 둔다.

\medskip
\noindent\textbf{정리 XI.} \emph{\(\bideal\)이 \(P\)를 계수체로 갖는 한 부정원 \(t\)의 다항식 아이디얼, 따라서 주아이디얼이라고 하자. \(h(t)\)는 \(\bideal\)에 속하는 \(k\)차 다항식이고, \(\Bmod\)은 \(\bideal\)을 \(h(t)\)로 법을 취했을 때 생기는 \(1,t,\ldots,t^{k-1}\)의 선형형식 모듈이라고 하자. \(\bideal\)의 기저다항식 \(f(t)\)의 차수가 \(p\)이면 \(\Bmod\)의 랭크는 \(k-p\)이다.}

가정에 따라 \(h(t)=h_1(t)f(t)\)이고 \(h_1(t)\)의 차수는 \(k-p\)이다. 따라서 다항식
\[
f(t),tf(t),\ldots,t^{k-p-1}f(t)
\]
이 대표하는 \(\bideal\)의 \(h(t)\)에 관한 잉여류들은 선형독립이다. 또한 \(\bideal\)의 각 잉여류는 \(\mathsf{P}\)의 계수를 갖는 선형결합으로 이 \(k-p\)개의 특별한 류들로 표현된다. 한편 \(\bideal\)은 \(h(t)\)를 법으로 하여 \(1,t,\ldots,t^{k-1}\)의 선형형식 모듈로 옮겨지므로 그 랭크는 정확히 \(k-p\)이다.

이제 \(\aideal=\aideal_1\)이 변환된 다항식 아이디얼이고, \(A^{(1)}(x)\)가 \(\aideal_1\)에 속하면서 \(x_1\)에 관하여 정칙인 \(k_1\)차 다항식이라고 하자. \(\mathfrak A_1\)은 \(\aideal_1\)을 \(A^{(1)}(x)\)로 법을 취했을 때 생기는 \(1,x_1,\ldots,x_1^{k_1-1}\)의 선형형식 모듈이고, 그 계수는 다항식 \(a^{(2)}(x)\)이다. 또 \(\aideal_2\)는 \(\mathfrak A_1\)에서 얻는 모든 \(k_1\)차 행렬식이 생성하는 \(x_2,\ldots,x_n\)의 아이디얼이라고 하자. 정리 XI에 따라 \(\aideal_2\)가 영아이디얼인 것은 \(\aideal_1\)의 다항식들이 \(x_1\)에 관하여 차수 \(p_1>0\)인 공통인자, 여기서는 다항식 의미의 인자를 가질 때 그리고 그때뿐이다. \(\aideal_2\)가 영아이디얼이 아니면, \(\aideal\)이 변환된 아이디얼이라는 가정 때문에 \(\aideal_2\)는 \(x_2\)에 관하여 정칙인 \(k_2\)차 다항식 \(A^{(2)}(x)\)을 포함한다. 이는 변환 (12)을 특수한 변환 (25), (26)에서 합성하면 직접 알 수 있다. 다시 \(\mathfrak A_2\)를 \(\aideal_2\)가 \(A^{(2)}(x)\)을 법으로 하여 옮겨지는 \(1,x_2,\ldots,x_2^{k_2-1}\)의 선형형식 모듈이라 하고, \(\aideal_3\)을 \(\mathfrak A_2\)의 모든 \(k_2\)차 행렬식이 생성하는 \(x_3,\ldots,x_n\)의 아이디얼이라 하자. 이 아이디얼은 \(x_3\)에 관하여 정칙인 다항식 \(A^{(3)}(x)\)을 포함해야 한다.

이 절차를 영아이디얼에 이르거나 \(\aideal_{n+1}\)이 단위원아이디얼이 될 때까지 계속한다. \(\aideal_{n+1}\)은 \(x\)와 무관하므로 영아이디얼 또는 단위원아이디얼일 수밖에 없다. 아이디얼 \(\aideal_1,\aideal_2,\ldots\)은 기초로 택한 다항식 \(A^{(\lambda)}(x)\)과 무관하게 \(\aideal\)만으로 유일하게 정해진다. \(\aideal_1=\aideal\)에 대해서는 분명하므로 \(\aideal_1,\ldots,\aideal_i\)까지 성립한다고 가정하자. \(\aideal_i\)가 \(A^{(i)}(x)\)을 법으로 하여 선형형식 모듈 \(\mathfrak A_i\)로 옮겨진다면, \(\mathfrak A_i\)는 \(x_i\)의 차수가 \(k_i\)에 미치지 않는 \(\aideal_i\)의 모든 다항식으로도 특징지을 수 있다. 따라서 \(\mathfrak A_i\)는 \(A^{(i)}(x)\)의 차수 \(k_i\)에만 의존한다. 이제 \(\bar A^{(i)}(x)\)가 \(\aideal_i\)에 속하면서 \(x_i\)에 관하여 정칙인 \(\bar k_i>k_i\)차 다항식이고, \(\overline{\mathfrak A}_i\)가 대응하는 선형형식 모듈, \(\bar\aideal_{i+1}\)이 \(\overline{\mathfrak A}_i\)의 \(\bar k_i\)차 행렬식들이 생성하는 아이디얼이라고 하자. 그러면 모듈등식
\[
\overline{\mathfrak A}_i=
(\mathfrak A_i,A^{(i)},x_iA^{(i)},\ldots,
x_i^{\bar k_i-k_i-1}A^{(i)})
\]
이 성립한다. \(A^{(i)}(x)\)이 \(x_i\)에 관하여 정칙이므로 \(A^{(i)},x_iA^{(i)},\ldots\)를 선형형식의 새 부정원들로 도입할 수 있다. 따라서 \(\mathfrak A_i\)의 \(k_i\)차 행렬식들과 \(\overline{\mathfrak A}_i\)의 \(\bar k_i\)차 행렬식들이 일치하고, \(\bar\aideal_{i+1}=\aideal_{i+1}\)이다\srcfn{17)}{이는 원리상 \(\Gmod^*_{i-1}\mid\Mmod^*_{i-1}\)과 \(\Gmod_{i-1}\mid\Mmod_{i-1}\)의 동형이 바탕을 둔 것과 같은 추론이다.}.

또한 \(\aideal_{i+1}\)은 언제나 \(\aideal_i\)로 나누어진다. \(\aideal_{i+1}\)이 영아이디얼이면 분명하다. 그렇지 않으면 \(\mathfrak A_i\)의 랭크는 \(k_i\)이다. 따라서 정의상 \(\aideal_{i+1}\)은 \(\mathfrak A_i\)의 \(k_i\)차 행렬식들로 이루어진 아이디얼이고, \(\mathfrak A_i\), 따라서 \(\aideal_i\)로 나누어진다. 그러므로 모든 \(\aideal_{i+1}\)은 \(\aideal\)로 나누어진다. 따라서 \(\aideal_{n+1}\)이 단위원아이디얼이면 \(\aideal\)도 단위원아이디얼이다.

이제 \(\aideal\)은 단위원아이디얼과 다르고
\[
\aideal_1\ne0,\ldots,\aideal_i\ne0,\qquad
\aideal_{i+1}=0
\]
이라고 하자. \(D^{(i)}(x)\)는 정리 XI에 따라 존재하는 \(\aideal_i\)의 모든 다항식의 최대공약원, 즉 다항식 의미의 공통인자이다. \(A^{(i)}(x)\)의 인자이므로 \(D^{(i)}(x)\) 자체도 \(x_i\)에 관하여 정칙이고 차수는 \(p_i>0\)이다. \(x_{i+1},\ldots,x_n\)을 \(\mathsf{P}(u)\)에 첨가하자. 그러면 \(D^{(i)}(x)\)은 \(\mathsf{P}(u;x_{i+1}\ldots x_n)\)의 한 대수적 확대체에서 \(p_i\)개의 일차인자로 분해된다. \(x_i-\bar x_i\)가 그 한 인자라 하자. 그러면 \(x_i=\bar x_i\)에서 \(\aideal_i\)의 모든 다항식이 사라지고, 정리 XI에 따라 \(\aideal_{i-1}\)의 다항식들은 최대공약원 \(D^{(i-1)}(x)\)을 얻는다. \(A^{(i-1)}(x)\)의 인자이므로 \(D^{(i-1)}(x)\)도 \(x_{i-1}\)에 관하여 정칙이고 차수는 \(p_{i-1}>0\)이다. 특히 항등적으로 사라질 수 없고, \(\mathsf{P}(u,\bar x_i,x_{i+1},\ldots,x_n)\)의 한 확대체에서 \(p_{i-1}\)개의 일차인자로 분해된다. \(x_{i-1}-\bar x_{i-1}\)가 그 한 인자이면 이에 \(\aideal_{i-2}\)의 최대공약원이 대응한다. 이 과정을 계속하면 \(P(u;x_{i+1},\ldots,x_n)\)의 한 대수적 확대체에 속하는 \(\aideal\)의 공통영점들을 유한개 얻는다. 거꾸로 \(\aideal_\lambda\)는 \(\aideal\)로 나누어지므로 \(\aideal\)의 모든 영점은 \(\aideal_\lambda\)의 영점이기도 하다. 특히
\[
x_1=\bar x_1,\ldots,x_i=\bar x_i,\quad
x_{i+1}=x_{i+1},\ldots,x_n=x_n
\]
이 \(\aideal\)의 영점이면 \(\aideal_{i+1}=0\)이고, \(\aideal_i\)의 다항식들은 일차인자 \(x_i-\bar x_i\)를 갖는 최대공약원 \(D^{(i)}(x)\)을 얻는다. 요약하면 다음과 같다.

\medskip
\noindent\textbf{정리 XII.} \emph{\(\aideal\)이 단위원아이디얼과 다르고
\[
\aideal_1\ne0,\ldots,\aideal_i\ne0,\qquad
\aideal_{i+1}=0
\]
이라고 하자. \(x_{i+1},\ldots,x_n\)을 \(\mathsf{P}(u)\)에 첨가하면 \(\mathsf{P}(u,x_{i+1},\ldots,x_n)\)의 한 대수적 확대체에 속하는 서로 결부된 값의 계 \(\bar x_1,\ldots,\bar x_i\)가 유한개 존재하여, \(x_1=\bar x_1,\ldots,x_i=\bar x_i\)에서 \(\aideal\)의 모든 다항식이 사라진다. 거꾸로 \(\aideal\)의 그러한 영점이 있으면 \(\aideal_{i+1}=0\)이고, \(\aideal_i\)의 모든 다항식은 일차인자 \(x_i-\bar x_i\)를 갖는 공통인자 \(D^{(i)}(x)\)를 갖는다.}

특히 정리 XII는 영점을 갖지 않는 모든 아이디얼 \(\aideal\)이 단위원아이디얼임을 보여 준다.

값의 계가 서로 결부되는 방식까지 명시적으로 보여 주는 분해를 얻기 위해, \(P(u,x_{i+1},\ldots,x_n)\)에 첨가할 새 부정원 \(v\)들을 도입한다. 다음과 같이 두자.
\begin{equation}
x_i=-v_1x_1-\cdots-v_{i-1}x_{i-1}+z_i.
\tag{35}
\end{equation}
(12)와 (35)를 합성하면 다시 (12)형의 치환을 얻는다. 다만 \(u_{i\kappa}\)는 \(w_{i\kappa}=u_{i\kappa}-v_\kappa\)로, 더 일반적으로 \(u_{i+\lambda,\kappa}\)는
\[
w_{i+\lambda,\kappa}
=u_{i+\lambda,\kappa}-u_{i+\lambda,i}v_\kappa
\]
로 바뀐다. 따라서 치환 (35)가 변환된 아이디얼 \(\aideal\)을 \(\bideal\)로 옮기면, \(\bideal\)은 \(\aideal\)에서 \(u_{\mu\nu}\)를 \(w_{\mu\nu}\)로 바꾸고 \(v\)들을 첨가하여 생긴다. 같은 과정으로 \(\bideal\)이 정의하는 아이디얼 \(\bideal_1,\bideal_2,\ldots\)은 \(\aideal_1,\aideal_2,\ldots\)에서 생긴다. 거꾸로 \(v=0\)으로 두면 \(\bideal_i\)에서 \(\aideal_\lambda\)를 얻는다. 따라서 \(\aideal_i,\bideal_i\)는 동시에 영아이디얼이거나 동시에 영아이디얼과 다르다.

다시
\[
\aideal_1\ne0,\ldots,\aideal_i\ne0,\qquad\aideal_{i+1}=0
\]
이라 하자. 그러면 \(\bideal_1\ne0,\ldots,\bideal_i\ne0,\ \bideal_{i+1}=0\)도 성립한다. \(\bar x_1,\ldots,\bar x_i,x_{i+1},\ldots,x_n\)이 서로 결부된 \(\aideal\)의 영점계이고
\[
\bar z_i=\bar x_i+v_1\bar x_1+\cdots+v_{i-1}\bar x_{i-1}
\]
로 정의하면, (35)에 따라 \(\bar x_1,\ldots,\bar x_{i-1},\bar z_i,x_{i+1},\ldots,x_n\)은 \(\bideal\)의 영점이다. 이제 \(H^{(i)}(z,x)\)가 \(\bideal_i\)의 모든 다항식의 최대공약원이라고 하자. 이는 \(v\)에 관하여 정수적이고 원시적이라고 가정할 수 있다. 그러면 정리 XII의 마지막 부분에 따라 \(H^{(i)}(z,x)\)은 인자
\[
z_i-\bar z_i
=z_i-(\bar x_i+v_1\bar x_1+\cdots+v_{i-1}\bar x_{i-1})
\]
를 갖는다. \(x_{i+1},\ldots,x_n\)을 첨가했을 때 나타나는 \(\aideal\)의 각 영점에 그러한 인자가 대응한다.

이 인자들이 \(H^{(i)}(z,x)\)의 분해를 모두 소진함을 보이자. 즉 \(z\)와 \(v\)에 관하여 일차인자들로 분해되지 않는 인자 \(H_1^{(i)}\)를 \(H^{(i)}_1\)에서 떼어 낼 수 없음을 보이면 된다. 그렇다고 가정하자. \(v\)에 관한 원시성 때문에 \(H^{(i)}\)은 적어도 하나의 일차인자 \(z_i-\bar{\bar z}_i\)를 가지며, \(\bar{\bar z}_i\)는 \(\mathsf{P}(u,v,x_{i+1}\ldots x_n)\)의 한 대수적 확대체에 속한다. \(\bar{\bar x}_1\ldots\bar{\bar x}_{i-1},\bar{\bar x}_i,x_{i+1}\ldots x_n\)이 이에 대응하는 \(\aideal\)의 영점이면, 이는 \(x_{i+1},\ldots,x_n\)을 첨가했을 때 나타나는 유한개의 \(\aideal\)의 영점 가운데 하나와 일치해야 하므로 \(v\)와 무관하다. 따라서 \(\bar{\bar z}_i\)는 \(v\)들을 반드시 선형으로 포함하며, 대수적으로나 유리비선형으로 포함하지 않는다. 그러면 가정에 어긋나게 \(H_1^{(i)}\)에서 \(z,v\)에 관한 또 하나의 일차인자
\[
z_i-(\bar x_i+v_1\bar x_1+\cdots+v_{i-1}\bar x_{i-1})
\]
를 떼어 낼 수 있다. 따라서 명시적인 분해
\[
H^{(i)}(z,x)=
\prod_\lambda
\left\{z_i-(\bar x_{i\lambda}+v_1\bar x_{1\lambda}
+\cdots+v_{i-1}\bar x_{i-1,\lambda})\right\}^{\alpha_\lambda}
\]
를 얻는다.

\(H^{(i)}\)의 차수와 각 지수 \(\alpha_\lambda\)는 \(D^{(i)}(x)\)의 대응하는 것들과 같아야 한다. \(H^{(i)}\)은 \(D^{(i)}\)에서 \(u_{\mu\nu}=w_{\mu\nu}\)로 특수화하여 생기고, \(D^{(i)}\)은 \(H^{(i)}\)에서 \(v=0\)으로 특수화하여 생기기 때문이다. 그러므로 이 분해는 또 다음을 보여 준다. 부정원 \(u_{\mu\nu}\)들을 \(\mathsf{P}\)에 첨가했기 때문에, \(\aideal\)에서 시작하는 소거에서도 \(D^{(i)}\)의 각 영점 \(\bar x_i\)에는 \(\aideal_{i-1},\ldots,\aideal_1\)의 매 단계에서 일차인 최대공약원 또는 그 거듭제곱이 대응한다. 따라서 매번 \(\aideal\)의 서로 결부된 영점계 \(\bar x_1,\ldots,\bar x_i,x_{i+1},\ldots,x_n\) 하나만이 생긴다\srcfn{18)}{여기서 제시한 연속소거는 크로네커의 방법과 달리 주어진 아이디얼 \(\aideal\)에만 의존하고 어떤 아이디얼기저에도 의존하지 않는다. 따라서 특히 지수 \(\alpha_\lambda\)들도 기저와 무관하다. 헨첼에 따르면 이 방법으로 소거를 완전히 수행하려면 몫 \(\aideal_i/D^{(i)}\)에 대하여 절차를 계속해야 하지만, 다음 절을 고려하여 여기서는 생략한다. (E. N.)}.

\section*{\S{}7. 종결식형과 소거론}

종결식형과 소거론의 관계를 얻기 위해 제6절의 연속소거를 아이디얼 \(\mideal\)을 그 \((i-1)\)단계 기초아이디얼로 나눈 몫
\[
\aideal=\mideal/\gideal_{i-1}
\]
에 적용한다. 먼저 이 몫이 변환된 아이디얼임을 보여야 한다. \(\mideal\)은 변환되었고, 따라서 제3절의 정리 VI에 따라 \(\gideal_{i-1}\)도 변환되었다. 이제
\[
H(x)\zero{\mideal/\gideal_{i-1}},\qquad
H(x)=\overline H(y)=\sum_i\alpha_i(u)\overline H_i(y)
=\sum_i\alpha_i(u)H_i(x)
\]
이면
\[
\overline H(y)\,\overline\gideal_{i-1,(u)}
\zero{\overline\mideal_{(u)}},
\qquad\text{따라서}\qquad
\overline H(y)\,\overline\gideal_{i-1}
\zero{\overline\mideal_{(u)}}.
\]
그러므로
\[
\overline H_\lambda(y)\,\overline\gideal_{i-1}
\zero{\overline\mideal},
\qquad\text{따라서}\qquad
H_\lambda(x)\gideal_{i-1}\zero{\mideal}.
\]
즉 \(H_\lambda(x)\)는 \(\mideal/\gideal_{i-1}\)로 나누어지고, 이 아이디얼이 변환되었음이 증명되었다. 따라서 제6절의 소거법을 적용할 수 있다.

한편 제4절의 (28)을 고려하면 (30)은 \(C^{(i)}(x)\)가 \(\mideal/\gideal_{i-1}\)로 나누어짐을 보여 준다. 여기서 \(C^{(i)}(x)\)는 \(\Mmod^*_{i-1}\)에서 얻는 항등적으로 영이 아닌 임의의 \(\rho\)차 행렬식이다. \(i>1\)이면 \(C^{(i)}\)은 \(x_1\)에 관하여 0차이므로, 아이디얼 \(\aideal=\aideal_1\)에 대응하는 모듈 \(\mathfrak A_1\)은
\[
C^{(i)},x_1C^{(i)},\ldots,x_1^{k_1-1}C^{(i)}
\]
을 포함한다. 따라서 \((C^{(i)})^{k_1}\)은 \(\aideal_2\)로 나누어진다. 마찬가지로 \((C^{(i)})^{k_1k_2}\)은 \(\aideal_3\)으로 나누어지고, 끝으로
\[
(C^{(i)})^{k_1k_2\cdots k_{i-1}}
\]
은 \(\aideal_i\)로 나누어진다. 그러므로 \(\aideal_1,\ldots,\aideal_i\)는 영아이디얼과 다르며, 이는 \(i=1\)일 때에도 성립한다. 이제 \(D^{(i)}(x)\)가 \(\aideal_i\)의 모든 다항식의 최대공약원이라고 하자. 그 차수가 \(p_i>0\)인 것은 \(\aideal_{i+1}\)이 영아이디얼일 때뿐이다. \(D^{(i)}(x)\)의 정의와 \(\aideal_i\)가 \(\aideal\)로 나누어진다는 사실에서
\[
d^{(i+1)}D^{(i)}(x)\zero{\mideal/\gideal_{i-1}},
\qquad d^{(i+1)}\ne0
\]
을 얻는다. 따라서 \(d^{(i+1)}D^{(i)}\)은 \(x_1\ldots x_{i-1}\)과 무관하고 \(\mideal/\gideal_{i-1}\)에 속하는 다항식이다. 한편 \(E^{(i)}(x)\)은 \(\Mmod_{i-1}\), 또는 \(\Mmod^*_{i-1}\)의 가장 높은 초등인자이며, 정리 II와 제4절에 따라 \(\mideal/\gideal_{i-1}\)에 속하면서 \(x_1,\ldots,x_{i-1}\)과 무관한 모든 다항식의 다항식 의미의 최대공약원으로 정의되었다. 따라서 \(d^{(i+1)}D^{(i)}\)은 \(E^{(i)}(x)\)으로 나누어지고, \(E^{(i)}(x)\)가 \(x_i\)에 관하여 정칙이므로 \(D^{(i)}(x)\) 자체도 \(E^{(i)}(x)\)으로 나누어진다.

처음부터 변환 (12)을 (35)와 합성했다고 해도, 즉 부정원 \(u_{\mu\nu}\)를 \(w_{\mu\nu}\)로 바꾸었다고 해도 같은 논증을 수행할 수 있다. 그러면 \(H^{(i)}(z,x)\)이 대응하는 \(\bar E^{(i)}(z,x)\)으로 나누어진다. \(D^{(i)}(x),H^{(i)}(z,x)\)과 마찬가지로 \(E^{(i)}(x),\bar E^{(i)}(z,x)\), 또 \(R^{(i)}(x),\bar R^{(i)}(z,x)\)도 서로 특수화하여 얻어지므로, 일차인자 분해에서의 중복도들도 서로 일치해야 한다. 여기에 \(E^{(i)}(x)\)의 어떤 거듭제곱이 \(R^{(i)}(x)\)으로 나누어진다는 사실을 고려하면 다음을 얻는다.

\medskip
\noindent\textbf{정리 XIII.} \emph{\(R^{(i)}(x)\)을 \(\mathsf{P}(u,x_{i+1}\ldots x_n)\)의 한 대수적 확대체에서 일차인자 \(x_i-\bar x_i\)들로 분해하면, \(\bar x_i\)를 서로 결부된 값의 계
\[
\bar x_1,\ldots,\bar x_i,x_{i+1},\ldots,x_n
\]
로 보충하는 방법은 유일하며, 이 값의 계는 \(\mideal/\gideal_{i-1}\), 따라서 \(\mideal\)의 영점이 된다. \(R^{(i)}\)의 일차인자들에서 이렇게 생기는 \(\mideal\)의 유한개의 영점, 즉 \(x_{i+1},\ldots,x_n\)을 첨가한 뒤의 유한개의 영점은 명시적인 분해
\begin{equation}
\bar R^{(i)}(z,x)=
\prod_\nu
\left\{z_i-(\bar x_{i\nu}+v_1\bar x_{1\nu}
+\cdots+v_{i-1}\bar x_{i-1,\nu})\right\}^{\lambda_\nu}
\tag{36}
\end{equation}
에 함께 나타난다. \(\bar R^{(i)}(z,x)\)이 \(z_i\)에 관하여 정칙이므로, \(\mathsf{P}(u)\)에 첨가한 \(x_{i+1},\ldots,x_n\)을 \(\mathsf{P}(u)\), 또는 거기서 유도되는 대수적으로 닫힌 체의 임의의 특수한 값의 계 \(\bar x_{i+1},\ldots,\bar x_n\)으로 바꾸어도 이 분해는 유지된다.}

이로써 \(\mideal\)의 영점들이 종결식형의 개별 인자들에 따라 동시에 분리되었다. \(\mathsf{P}(u)\)에 첨가된 부정원 \(x\)의 수는 대응하는 대수적 도형의 차원이라고도 부른다.

\(\bar R^{(i)}(z,x)\), 또는 \(R^{(i)}(x)\)의 대응 분해에서 \(\mathsf{P}(u,x_{i+1},\ldots,x_n)\)에 관하여 켤레인 인자들을 묶자. 일반적인 \(\mathsf{P}\)에서는 이 인자들이 전부 또는 일부 동일할 수 있다. 그러면 \(R^{(i)}(x)\)을 \(P(u,x_{i+1},\ldots,x_n)\)에 관하여 기약인 다항식들의 거듭제곱으로 분해한 식
\[
R^{(i)}(x)=S_1^{(i)}(x)^{\beta_1}\cdots
S_\nu^{(i)}(x)^{\beta_\nu}
\]
을 얻는다. 정리 X에 따라 이 분해에는 표현
\[
\gideal_i=[\mathfrak s_1,\ldots,\mathfrak s_\varkappa]
\]
이 대응한다. 여기서 \(S_\mu^{(i)}(x)^{\beta_\mu}\)은 \(\gideal_{i-1}\)의 아이디얼 \(\mathfrak s_\mu\)에 대한 노름, 또는 모듈 \(\Gmod_{i-1}\)의 \(\mathfrak s_\mu\)에 대응하는 모듈에 대한 노름이다. 따라서 지수 \(\beta_\mu\)의 의미도 분명해진다. 특히 \(\gamma_\mu\)가 \(S_\mu^{(i)}\)의 차수라면 중복도 \(\beta_\mu\gamma_\mu\)는 \(\mathsf{P}(u,x_{i+1}\ldots x_n)\)에 관하여 \(\gideal_{i-1}\)의 \(\mathfrak s_\mu\)에 대한 선형독립인 잉여류의 수와 같다.

끝으로 \(\mathsf{P}\)의 표수가 영인 특수한 경우를 간단히 보자. \(n\)개 행렬식 \(C^{(i)}(x)\), \((i=1,2,\ldots,n)\), 이 각각 \(x_i\)에 관하여 정칙으로 남도록 \(u_{\mu\nu}\)를 \(P\)의 원소 \(\bar u_{\mu\nu}\)로 특수화하면\srcfn{19)}{헨첼은 처음부터 \(u_{\mu\nu}\)에 이 특수화를 가한다. 그래서 \(H^{(i)}(z,x)\), 또는 \(\bar R^{(i)}(z,x)\)이 \(z,v\)에 관한 일차인자들로 분해됨을 보이는 데 훨씬 복잡한 고찰을 써야 한다. 이때 나타나는 지수들은 위의 지수들과 반드시 일치하지는 않는다. (E. N.)}, \(R^{(i)}\)의 정칙성도 유지된다. 이 특수화 아래 \(R^{(i)}\), 그리고 \(\bar R^{(i)}(z,x)\)은 \(\Mmod^*_{i-1}\)의 모든 \(\rho\)차 행렬식의 공통인자로 남지만 반드시 최대공약원일 필요는 없다. 그러나 이 마지막 성질도 유지되도록 언제나 특수화할 수 있다. 실제로 모듈기저의 존재가 보여 주듯 \(\bar R^{(i)}(z,x)\)은 \(\Mmod^*_{i-1}\)에서 얻는 유한개의 \(\rho\)차 행렬식의 최대공약원으로도 특징지을 수 있다. 이렇게 특수화한 \(\bar R^{(i)}(z,x)\)의 분해는 (36)에서 \(u_{\mu\nu}\)를 \(\bar u_{\mu\nu}\)로 바꾸어 단순히 얻어지므로, 소거론 전체가 유지된다.

\begin{center}
(1922년 3월 17일 접수)
\end{center}

\clearpage
\setcounter{footnote}{0}
\phantomsection
\addcontentsline{toc}{section}{23. 대수불변식과 미분불변식}
%
%
%
%
%
%

\section*{23. 대수불변식과 미분불변식}
\begin{center}
Jahresber. d. Deutschen Mathematiker-Vereinigung 32 (1923), 177--184쪽
\end{center}

\begin{center}
{\Large\bfseries 대수불변식과 미분불변식.}\par
\vspace{0.45em}
1922년 9월 18일 라이프치히 강연에 따른 보고.\par
\vspace{1.0em}
괴팅겐의 \textsc{에미 뇌터}.
\end{center}

\begingroup
\renewcommand{\thefootnote}{\arabic{footnote})}
\setcounter{footnote}{0}
\newcommand{\PXXIIIitem}[2]{\par\noindent\hangindent=2.2em\hangafter=1\makebox[2.2em][l]{#1.}#2\par}

대수불변식과 미분불변식의 발전을 보고하면서, 두 분야 모두에서 힐베르트의 표현에 따르면 소박한 시기와 형식적 시기 뒤에 오는 비판적 시기에 논의를 한정하고자 한다. 이 비판적 시기는 대수불변식의 경우 힐베르트 자신의 이름으로, 미분불변식의 경우 리만의 이름으로 특징지어진다. 내용을 기준으로 말하면, 전자에서는 이러한 문제를 다루는 과정에서 본질적으로 발달하고 여기서 그 온전한 위력을 드러낼 수 있었던 대수학의 산술적 방법이, 후자에서는 형식적 변분법의 방법이 그 특징이다. 따라서 나는 특수한 대상 자체를 넘어서는 의의를 가진 산술적 방법을 보고하고, 둘째 부분에서는 변분법의 방법을 따라 이루어지는 미분불변식론의 대수불변식론으로의 환원만을 다루겠다.

\textbf{1.} 잘 알려져 있듯이 \emph{대수불변식}의 바탕에는 $x_1,\ldots,x_n$의 \emph{일반선형군}
\srcnumdisplay{(1)}{x_i=\sum s_{ik}x'_k
\qquad\text{또는 간단히}\qquad x=S(x')}
이 있다. 여기서 $s_{ik}$는 부정원이고, 따라서 $\D=|s_{ik}|$는 영이 아니다. 이제 $f(a,x),g(b,x),\ldots$가 $x$에 관한 차수 $\alpha,\beta,\ldots$의 형식이고 그 계수가 부정원 $a,b,\ldots$라면, (1)에 따라 $a,b,\ldots$의 \emph{유도변환}이 생긴다. 실제로
\[
\begin{aligned}
f(a,x)&=\sum a_{i_1\ldots i_n}x_1^{i_1}\cdots x_n^{i_n}
      =\sum a'_{j_1\ldots j_n}{x'_1}^{j_1}\cdots {x'_n}^{j_n}
      =f(a',x'),\\
g(b,x)&=\sum b_{r_1\ldots r_n}x_1^{r_1}\cdots x_n^{r_n}
      =\sum b'_{s_1\ldots s_n}{x'_1}^{s_1}\cdots {x'_n}^{s_n}
      =g(b',x')
\end{aligned}
\]
에서 $a',b',\ldots$가 $a,b,\ldots$의 선형 동차함수가 되고, 그 계수들은 $s_{ik}$에 관하여 각각 차수 $\alpha,\beta,\ldots$임을 알 수 있다.
\srcnumdisplay{(2)}{a'=A_\alpha(a);\qquad b'=B_\beta(b),\ldots .}
이제 \emph{불변식}이란 변환 (1)과 (2)에 대한 불변식을 말한다. 즉 $a,b,\ldots,x,y,\ldots$의 부정원에 관한 다항함수 --- 여기서 $y=S(y')$이기도 하다 --- 로서, 변환을 적용하면 치환행렬식의 한 거듭제곱이라는 인자만큼 달라진 채 다시 자신을 주는 것이다.
\srcnumdisplay{(3)}{I(a',b',\ldots,x',y',\ldots)
=\D^\rho I(a,b,\ldots,x,y,\ldots).}
임의의 수체 $P$에 계수를 둔 각 불변식은 유리수 계수 또는 정수 계수를 가진 유한개의 특수 불변식의 $P$-선형결합으로 나타낼 수 있으므로, $I$의 계수는 유리수, 나아가 정수라고 가정할 수 있다. 마찬가지로 각 불변식은 각 열에 관하여 따로 동차인 유한개의 특수 불변식의 합이므로, $I$가 각 열에 관하여 따로 동차라고 가정해도 제한이 없다.

역으로 조건 (3)을 요구하면 (1)과 (2)도 다시 결정된다는 점을 지적해 두자. 최근 Ostrowski와 Schur가 보였듯이,\footnote{A. \textsc{Ostrowski}와 I. \textsc{Schur}, Über eine fundamentale Eigenschaft der Invarianten einer binären Form, Math. Zeitschr. 15 (1922), 및 이에 이어지는 Ostrowski의 미출판 연구.} 유도변환은 각 열마다 따로 작용하는 선형변환들 가운데, 명시할 수 있는 몇 가지 예외를 제외하면 $x,y,\ldots$를 포함하지 않는 임의의 불변식 $I$를 허용하는 변환들의 전체로도 특징지을 수 있다.

두 불변식의 곱은 다시 불변식이지만, 합은 가중치, 곧 (3)의 지수 $\rho$가 일치할 때 그리고 그때에만 불변식이다. 그러나 행렬식이 1인 변환들로 제한하면 임의의 합도 다시 불변식이며, 이렇게 제한된 군에 대한 모든 불변식을 빠짐없이 얻는다. 따라서 하나의 정역, 더 정확히는 \emph{동차 정역}을 얻는다. 여기서 동차 정역이란 어떤 다항식이 영역에 속하면 그 동차 성분들도 각각 그 영역에 속하는 영역을 말한다. 이 경우에는 다시 (3)의 각 열에 관하여 따로 동차인 불변식들로 귀착된다.

여기서 불변식론과 산술대수학 전체가 발전하는 계기가 된 근본 문제가 생긴다. \emph{이 정역은 유한한가? 즉 모든 불변식을 $I_1,\ldots,I_k$의 다항식으로 나타낼 수 있게 하는 유한 정수성 기저 $I_1,\ldots,I_k$를 가지는가,
$I=G(I_1,\ldots,I_k)$?} 힐베르트의 해법은 정수성 기저의 문제를 아이디얼 기저의 문제로 환원한다. Gordan의 이원형식에 대한 증명은 조망하기 어려운 기호계산 때문에 확장할 수 없고, Mertens--Hilbert의 증명은 이원형식을 선형형식들로 분해하는 것을 사용하기 때문에 확장할 수 없다. 여기서는 산술적 기본사상을 강조하면서 힐베르트의 논법을 간단히 개괄한 뒤, 그와 이어지는 세 문제군, 곧 유한 단계에 의한 기저의 실제 구성, 부분군에서의 유한성 문제, 정수성을 고려한 유한성 문제를 다루겠다.

\textbf{2.} (유한) 대수적 수체에서의 \emph{아이디얼 정의}를 상기하자. 그 체의 수들로 이루어진 계 $\ma$가 아이디얼이라는 것은 다음을 뜻한다.
\PXXIIIitem{1}{$\ma$는 그 체의 모든 대수적 정수의 영역 $\mo$에 포함된다.}
\PXXIIIitem{2}{$\alpha$와 $\beta$가 $\ma$에 속하면 그 차 $\alpha-\beta$도 $\ma$에 속한다.}
\PXXIIIitem{3}{$\alpha$가 $\ma$에 속하고 $\lambda$가 $\mo$의 임의의 원소이면 $\lambda\alpha$도 $\ma$에 속한다.}
그리고 모든 아이디얼은 \emph{아이디얼 기저}를 가져 $\ma=(\alpha_1,\ldots,\alpha_k)$가 된다는 정리가 알려져 있다. 즉 $\ma$에는 유한개의 원소 $\alpha_1,\ldots,\alpha_k$가 있어서, 조건 2와 3으로 이 원소들에서 $\ma$를 이끌어 낼 수 있다. 다시 말해
$\alpha=\lambda_1\alpha_1+\cdots+\lambda_k\alpha_k$가 $\ma$의 모든 원소를 빠짐없이 나타낸다.\footnote{$k$를 2까지 줄일 수 있다는 사실은 이하의 논의와 무관하다.}

이 아이디얼 정의는 $\mo$가 정역이라는 사실에만 의존한다. 따라서 $\mo$를 임의의 정역으로 바꾸어도 아이디얼의 개념과 아이디얼 기저의 개념은 글자 그대로 유지된다.

이제 힐베르트의 유한성 증명은 다음 세 정리로 나뉜다.
\PXXIIIitem{1}{유리성 영역에 계수를 둔 $n$개 부정원의 모든 다항식의 영역에서는 모든 아이디얼에 대하여 아이디얼 기저 정리가 성립하며, 따라서 임의의 계 $\mS$에 대해서도 성립한다.}
\PXXIIIitem{2}{다항식들로 이루어진 동차 정역 $\mJ$가 유한한 것은 $\mJ$에서 아이디얼 기저 정리가 성립할 때 그리고 그때뿐이다.}
\PXXIIIitem{3}{불변식 영역 $\mJ$에서는 아이디얼 기저 정리가 다음의 강화된 형태로 성립한다. 모든 다항식의 영역에 관하여 만든 아이디얼 기저는 동시에 $\mJ$ 안에서도 아이디얼 기저이다. 즉}
\[
 I=A_1I_1+\cdots+A_kI_k
 \qquad\text{이면 언제나}\qquad
 I=i_1I_1+\cdots+i_kI_k
\]
이다. 여기서 $A$들은 $a,b,\ldots,x,y,\ldots$의 다항식이고, $i$들은 불변식이다.

정리 1의 증명은 원리상 대수적 수체의 경우와 같은 추론에 바탕을 둔다. 두 경우 모두 아이디얼 기저의 존재는 선형형식들의 모듈에 관한 한 일반정리의 직접적인 귀결이다.\footnote{이에 관해서는 대수함수의 산술이론에 관한 나의 비교 보고, Jahresber. 28 (1920), 187쪽 및 188쪽 주 2)를 참조하라.} 정리 1에서 다항식들로 이루어진 모든 유한 정역에 아이디얼 기저가 존재함이 직접 따라오며, 동차 영역에서는 그 역도 성립함을 쉽게 알 수 있다. 정리 3에서만 불변식의 특수한 성질이 사용된다. 힐베르트에게서 정리 3은 $\Omega$-과정에 관한 알려진 정리의 귀결로 얻어진다. 최근 E. Fischer는 다른 일반정리들을 바탕으로 하면, 각 변환과 더불어 그 반변변환 또는 복소켤레변환도 포함한다는 일반선형군의 성질만으로 정리 3을 이끌어 낼 수 있음을 보였다.\footnote{E. \textsc{Fischer}, Über die Endlichkeit der Invarianten, Göttinger Nachrichten 1915 및 라이프치히 강연.}

\textbf{3.} \emph{기저의 실제 구성}은 \emph{함수체} 이론을 도입하여 유한성의 개념을 한 번 더 산술적으로 바꾸어 쓰는 데 바탕을 둔다. 다음이 성립한다.
\PXXIIIitem{4}{다항식들로 이루어진 정역 $\mJ$가 유한한 것은 $\mJ$ 안에 유한 부분정역 $\mJ'$가 있어서 $\mJ$가 $\mJ'$에 대하여 대수적으로 정수일 때 그리고 그때뿐이다. 이 경우 $\mJ'$의 임의의 정수성 기저에 이 대수적 정수 원소들의 한 기본계를 덧붙이면 $\mJ$의 기저를 얻는다.\footnote{Hilbert는 Über die vollen Invariantensysteme, Math. Annalen 42 (1893), § 1과 § 2에서 가정한 유한성으로부터 이 사실들이 따름을 보였다. 역으로 대수적 정수 종속성에서 유한성이 따르는 것은 유리기저의 존재에서 나온다(E. Noether, Körper und Systeme rationaler Funktionen, Math. Annalen 76 (1915), § 12).}}
특히 불변식의 경우처럼 아이디얼 기저 정리가 강화된 형태 3으로 성립하는 동차 영역 $\mJ$를 생각하자. 그러면 그러한 영역 $\mJ'$는 다음 성질을 가진, 언제나 존재하는 어떤 유한개의 다항식 $I_1,\ldots,I_s$에서 정수성 기저를 취하여 만들어 낼 수 있다. 이 다항식들이 모두 영이면 $\mJ$의 모든 다항식도 영이다. 이 사실은 다시 아이디얼론의 한 일반정리에 바탕을 둔다.
\PXXIIIitem{5}{다항식 아이디얼 $\ma$마다 정수 $r$이 존재하여, $\ma$의 모든 영점에서 영이 되는 각 아이디얼 $\mb$에 대하여 $\mb^r$은 $\ma$로 나누어진다.}

불변식의 경우에는 $I_1,\ldots,I_s$의 \emph{차수에 대한 상계}, 따라서 그 개수 $s$에 대한 상계도 정할 수 있으며, 이 상계는 바탕에 놓인 기본형식들의 차수에만 의존한다. 이 대목에서는 불변식의 특수한 성질이 다시 사용된다. 수치적으로 특수화한 기본형식계가 영이 아닌 불변식을 갖는 것은 $\D=|s_{ik}|$가 변환된 계수들에 대하여 대수적으로 정수일 때 그리고 그때뿐이라는 정리가 그것이다.

K. Hentzelt는\footnote{나는 곧 이 유고를 출판할 예정이다.} 이 상계에서 \emph{정수성 기저에 속하는 불변식들의 차수에 대한 상계}를 정했다. 이 상계도 기본형식들의 차수에만 의존한다. 그는 정리 5의 \emph{지수} $r$에 대하여, $\ma$의 한 아이디얼 기저에 속하는 다항식들의 개수와 최대차수에만 의존하고 그 계수들과는 무관하며 주어진 수들에서 유한 단계로 계산할 수 있는 \emph{상계}를 제시함으로써 이를 해냈다. 원리상 이로써 제기한 문제는 완전히 해결된다. 다만 그 상계는 상당히 크다. 예를 들어 삼원 이차형식에서는 판별식, 곧 3차 불변식 하나만으로도 $x$를 포함하지 않는 모든 불변식의 기저가 되지만, Hilbert와 Hentzelt의 방법으로는 상계가 수백만까지 올라간다.

\textbf{4.} 일반선형군의 \emph{부분군}에 대한 불변식의 유한성 문제에서는 아직 부분적인 결과만 얻었다. 그 방법은 거의 전적으로 Study가 직교군에 대하여 사용한 절차, 곧 부분군의 불변식을 일반선형군의 연립불변식으로 환원하는 데 바탕을 둔다. Weitzenböck은 운동불변식과 아핀불변식에 대하여, Deruyts는 반불변식과 평행이동불변식에 대하여 이를 수행하였다.\footnote{문헌은 불변식론에 관한 Weitzenböck의 백과사전 논문 III E 1을 참조하라.} 이 방법이 어디까지 미치는지는 아직 결정되지 않았다. 2절 끝에서 언급한 Fischer의 관찰도 여러 부분군의 유한성 문제를 해결하며, 특히 직교군에 대한 가장 간단한 증명을 준다. 그러나 반불변식의 예는 부분군에서는 유한성이 성립하더라도 아이디얼 기저에 관한 강화된 명제 3이 성립하지 않을 수 있음을 보여 준다. Fischer는 최근 이 점을 지적하였다.\footnote{라이프치히 강연 및 Math. Zeitschrift.}

여기서의 유한성 문제를 체론의 관점에서 정식화하면, 힐베르트가 제기한 \emph{상대적으로 정수인 함수}의 문제에 이른다. 즉 유리함수체에 속하는 모든 다항식으로 이루어진 정역은 언제나 유한한가라는 물음이다. 이 방향의 한 결과로, 유한군의 불변식에 대해서는 초등적으로 증명할 수 있는 특별한 정수성 기저를 직접 제시할 수 있다. 그것은 갈루아 레졸벤트의 계수계이다.\footnote{E. Noether, Der Endlichkeitssatz der Invarianten endlicher Gruppen, Math. Ann. 77 (1916).}

\textbf{5.} 불변식의 유한성 정리를 \emph{정수성}에 관하여 강화할 수 있는가라는 문제도 일반적으로는 아직 해결되지 않았다. 힐베르트가 보였듯이 모든 정수 계수 다항식의 영역에서도 아이디얼 기저 정리는 성립하고, 아이디얼 기저와 정수성 기저 사이의 연관 2도 그대로 유지된다.\footnote{특히 이 연관에서, 내가 Math. Ann. 83 (1921)에서 전개한 일반 아이디얼론의 분해정리가 모든 유한 정역에 대하여 성립함이 따른다.} 그러나 필수적이지는 않은 강화된 형태 3에 대한 증명은 실패한다. \emph{이원 불변식}의 경우에는 \emph{정수 계수 유한성}을 증명할 수 있다.\footnote{E. Noether, Die Endlichkeit des Systems der ganzzahligen Invarianten binärer Formen, Göttinger Nachrichten 1919.} 이 증명은 이원형식에 대한 Mertens--Hilbert 유한성 증명을 정수성에 관하여 강화한 것으로, 정수 계수 아이디얼 기저 정리뿐 아니라 이원형식을 선형형식들로 분해하는 것도 사용한다. 따라서 더 많은 부정원으로 옮길 수는 없다. 반면 비슷한 토대 위에서 \emph{유한군의 불변식}에 대한 \emph{정수 계수 유한성}을 위에서 언급한 증명의 강화로 얻을 수 있다. 다만 여기서도 기저를 실제로 제시하지는 못하고 존재만을 증명한다.\footnote{주 11)에 인용한 논문의 § 4.} 이 문제에서도 함수체 이론과 일반 아이디얼론을 한층 더 발전시켜야만 목표에 이를 수 있을 것이다.

\textbf{6.} 이제 \emph{미분불변식}과, 서론에서 제기한 그것을 대수불변식론으로 환원하는 문제로 넘어가겠다. 여기서 바탕이 되는 군은 잘 알려져 있듯이 $x_i=x_i(y)$에 따라 다음과 같이 정의된다.
\srcnumdisplay{(4)}{d x_i=\sum \frac{\partial x_i}{\partial y_k}d y_k;\qquad
 d\delta x_i=\sum \frac{\partial x_i}{\partial y_k}d\delta y_k+
\sum\frac{\partial^2 x_i}{\partial y_k\partial y_l}d y_k\,\delta y_l;\ \ldots .}
(4)가 유도하는 변환으로 넘어가기 위해 임의의 함수 $f(x,d x)=g(y,d y)$를 바탕으로 삼을 수 있다. $\delta f,\delta^2f,\ldots$가 (4)에 대하여 불변이라고 요구하면, $f$를 $x$와 $d x$로 미분한 도함수들의 변환을 얻는다. 이를테면 $d x$에 관한 동차형식으로 제한하면
\[
 f(x,d x)=\sum a_{i_1\ldots i_n}(x)d x_1^{i_1}\cdots d x_n^{i_n}
 =\sum a'_{j_1\ldots j_n}(y)d y_1^{j_1}\cdots d y_n^{j_n}
 =g(y,d y)
\]
이고, $a'$는 $a$에 관하여 선형이며, $\frac{\partial a'}{\partial y}$는 $a$와 $\frac{\partial a}{\partial x}$에 관하여 선형이고, 그 뒤도 마찬가지이다.
\srcnumdisplay{(5)}{a'=A_0(a),\qquad
\frac{\partial a'}{\partial y}=A_1\!\left(a,\frac{\partial a}{\partial x}\right),\qquad
\frac{\partial^2 a'}{\partial y^2}=A_2\!\left(a,\frac{\partial a}{\partial x},\frac{\partial^2 a}{\partial x^2}\right),\ldots}
따라서 문제는 변환 (4)와 (5)에 대한 불변성이다. 이때 나타나는 모든 양
\[
a,\ \frac{\partial a}{\partial x},\ldots,\ d x,\ \delta x,\ d\delta x,\ldots;\ 
\frac{\partial x}{\partial y},\ \frac{\partial^2 x}{\partial y^2},\ldots
\]
은 부정원으로 보아야 하므로 각 변환의 행렬식은 분명 영이 아니다. 형식적으로 완성된 이론을 특수한 함수에 적용할 때에야 비로소, 어떤 영역에서 변환이 유일하게 역변환될 수 있고 충분한 수의 미분몫이 존재하는 함수들로 제한해야 하는가라는 문제가 생긴다. 이는 대수적 경우 수치적 특수화에서 행렬식 $|s_{ik}|$가 영이 아니라고 가정해야 하는 것과 완전히 같다.

나타나는 모든 인수를 부정원으로 보는 이 관점에서 문제의 군은 보통의 대수적 방법으로 직접 다룰 수 없는 특수한 선형군이다. 이 군에서 $d x$의 변환과 $a'=A_0(a)$의 변환은 각각 (1)과 유도변환 (2)에 일치한다. 여기서 \emph{새 인수들을 도입}하여 변환 (4)와 (5)를 일반선형군 (1)과 그것이 \emph{유도하는 변환} (2)으로 일반적으로 \emph{환원}할 수 있는가라는 문제가 생긴다.

실제로 그렇게 할 수 있다. 고차 미분 $d^2x,\ d\delta x\ldots$는 $f(x,d x)$에 딸린 변분문제에서 생기는 라그랑주 도함수들과, 그것들을 다시 극화하여 얻는 결합들(Weyl과 Schouten의 용어로 접속)로 바꾸어야 한다. 이들은 모두 $d x$에 동변이므로 그 변환은 일반선형군 (1)로 주어진다. 한편 $\frac{\partial a}{\partial x},\frac{\partial^2a}{\partial x^2},\ldots$ 또는 더 일반적으로 동차라고 가정한 $f$를 $x$와 $d x$로 미분한 도함수들은 다음 형식들의 계수로 바꾸어야 한다. 즉 ``고차 변분의 정규형''에서 위에 제시한 방법으로 $d^2x,\ d\delta x\ldots$를 소거하여 생기는 어떤 형식들과 그 ``공변미분''
\[
 [\Omega_i],\quad [\Omega_i^{(1)}],\quad [\Omega_i^{(2)}],\ldots\qquad (i=1,2,\ldots)
\]
의 계수들이다. $[\Omega_i^{(k)}]$는 (4)와 (5)에 대한 불변식이며 $d x$와 $\delta x$에만 의존하므로, 이 계수들은 대수적으로 유도된 변환 (2)을 따른다.\footnote{E. Noether, Invarianten beliebiger Differentialausdrücke, Göttinger Nachrichten 1918.} 이 \emph{환원정리}에 따라 미분불변식론은 대수불변식론에 종속된다.

$f(x,d x)$가 $d x$에 관한 $\alpha$차 동차형식이면 $[\Omega_i]$의 열은 $[\Omega_\alpha]$에서 끝난다. 이차형식의 경우 $[\Omega_2]$는 리만의 곡률형식과 같아지며, 여기서 제시한 구성과정은 리만이 라틴어 현상논문에서 제시한 바로 그 과정이다. 리만의 교수자격취득 강연에 이어 Lipschitz도 이와 독립적으로 같은 과정을 전개하였다.\footnote{주 7)에 언급한 Weitzenböck의 백과사전 논문 제2부 28번에 실린 Riemann과 Lipschitz에 관한 나의 설명도 참조하라.} Christoffel과 Ricci가 이차 미분형식에 대하여 계산으로 증명한 환원정리는 한 점에서 나가는 극값곡선들을 직선으로 바꾸는 ``리만 정규좌표''의 도입에 바탕을 둔다. 이 좌표에서는 $x$의 임의의 변환에 정규좌표의 선형변환이 대응하므로, 일반선형군으로 넘어갈 수 있다.

Weyl과 Schouten의 절차처럼 군을 정의할 때 변분문제 대신 하나의 \emph{접속만을 바탕으로 삼는} 경우에도 그러한 환원정리가 성립하는가라는 문제가 생긴다. 즉 라그랑주 도함수에 대응하여 고차 미분을 소거하는 식 $\psi_i(d\delta x,d x,\delta x)$를 세우고, 이차 미분형식과 비슷하게 그 자체로는 불필요한 조건, 곧 $\psi_i$가 $d x$에 관하여 선형이어야 한다는 조건을 덧붙이는 경우이다. 임의의 동차 $f(x,d x)$에서는 극화된 라그랑주 도함수가 $d x$에 관하여 일차 동차일 뿐이다. $\psi$가 $d x$에 동변이어야 한다고 요구하면 그 계수들의 변환이 유도되고, 따라서 (5)와 비슷하게 군이 정해진다. Weitzenböck은 Weyl의 경우 가장 낮은 차수의 불변식들에 대하여 환원정리를 계산으로 확인했지만, 일반적인 접근법은 아직 없다.\footnote{이에 관해서는 Weyl, Raum, Zeit, Materie; Schouten, Math. Zeitschr. 13 (1922); Weitzenböck, Sitzungsber. d. Wiener Akademie, 1920 및 1921을 참조하라.}

\begin{center}
(1922년 10월 26일 접수.)
\end{center}
\endgroup

\clearpage
\setcounter{footnote}{0}
\phantomsection
\addcontentsline{toc}{section}{24. 소거이론과 일반 아이디얼 이론}
%
%
%
%
%
%

\providecommand{\frakS}{\mathfrak S}
\providecommand{\fraka}{\mathfrak a}
\providecommand{\frakb}{\mathfrak b}
\providecommand{\frakc}{\mathfrak c}
\providecommand{\frakg}{\mathfrak g}
\providecommand{\frakG}{\mathfrak G}
\providecommand{\frakh}{\mathfrak h}
\providecommand{\frakm}{\mathfrak m}
\providecommand{\frakM}{\mathfrak M}
\providecommand{\frakn}{\mathfrak n}
\providecommand{\frako}{\mathfrak o}
\providecommand{\frakp}{\mathfrak p}
\providecommand{\frakq}{\mathfrak q}
\providecommand{\frakr}{\mathfrak r}
\providecommand{\frakt}{\mathfrak t}
\providecommand{\frK}{\mathfrak K}
\providecommand{\barfrakm}{\overline{\mathfrak m}}

\section*{24. 소거이론과 일반 아이디얼 이론}
\begin{center}
Math. Ann. 90 (1923), 229--261쪽\\[0.9em]
{\Large\bfseries 소거이론과 일반 아이디얼 이론}\\[1.0em]
괴팅겐의 에미 뇌터\\[0.6em]
\rule{4em}{0.4pt}
\end{center}

이하에서는 내가 헨첼의 서술에 부여한 산술적 형태의 소거이론\footnote{Kurt Hentzelt, \emph{Zur Theorie der Polynomideale und Resultanten}. E. Noether 편, Math. Ann. 88 (1922), 53--79쪽. 이하 H.-N.으로 인용한다.} --- 이는 다항식영역의 아이디얼 이론이라고 할 수 있다 --- 을 일반 아이디얼 이론\footnote{E. Noether, \emph{Idealtheorie in Ringbereichen}, Math. Ann. 83 (1921), 24--66쪽. 이하 \emph{아이디얼 이론}으로 인용한다.} 안에 배치하고, 동시에 영점에 관한 부분을 새로 정초한다. 두 이론의 기본개념과 체론의 기본개념\footnote{E. Steinitz, \emph{Algebraische Theorie der Körper}, J. f. M. 137 (1910), 167--309쪽. 이하 Steinitz로 인용한다.}은 § 1에 짧게 모아 두었으므로 여기서 참조할 수 있다.

이 배치와 새로운 정초는 네 문제를 중심으로 한다. 차원개념의 산술적 정식화, 영점이론, 노름과 초등인자의 분해정리와 일반 아이디얼 이론의 분해정리 사이의 관계, 그리고 노름과 초등인자형에 의한 소아이디얼과 으뜸아이디얼의 특징짓기가 그것이다.

\emph{차원개념}의 산술적 정식화(§ 4)는 소아이디얼의 사슬로 주어진다. 이는 곡면 위에 곡선이 있고 곡선 위에 점이 있다는 사실의 산술적 등가물이다. 소거이론의 \emph{매개변수 정의}와 일치함이 증명되는 이 산술적 정식화는 조금만 고치면 임의의 환영역으로 옮길 수 있다. 나머지 문제들의 어떤 부분들을 함께 옮기면 거기서 아이디얼의 구조를 정확히 통찰할 수 있으며, 이는 다른 곳에서 보이겠다.

H.-N.에서는 통상적인 방식대로 주어진 아이디얼의 \emph{영점} 문제를 연속소거로 돌아가 직접 다룬다. 이 때문에 검산의 성격을 완전히 피하지는 못한다. 이와 대조되는 것은 한 부정원 다항식에 늘 쓰는 방법이다. 주어진 다항식을 \emph{소함수들의 거듭제곱들의 곱}(으뜸함수들의 곱)으로 나타내고, 각 소함수에 대하여 잉여류체와 \emph{동형인} 영점체를 구성한다. 소함수의 차수는 이 체의 차수를 주고, 그 영점이 소함수의 영점과 일치하는 으뜸함수의 차수는 이 으뜸함수에 따른 잉여류환의 차수를 준다. 각 소함수에 관하여 갈루아체로 넘어가면 일차인수분해를 얻는다. 이로써 원래 다항식의 영점, 일차인수분해, 중복도 문제가 해결된다. 이와 완전히 나란하게 여기서는(§ 3) 한 소아이디얼에 따른 잉여류계를 분수몫을 만들어 잉여류체로 확대하고, 그와 동형인 영점체를 구성한다. 이 영점체는 \(x_{i+1},\ldots,x_n\) 같은 부정원들을 첨가하여 생성한 부분체를 포함한다. 노름의 차수는 이 부분체에 관한 차수를 주며, 동시에 갈루아체로 넘어가면 초등인자형과 따라서 노름의 일차인수분해를 얻는다. 한편 대응하는 으뜸아이디얼의 영점들은 같다. 노름이 소아이디얼의 초등인자형의 한 거듭제곱이 되므로(§ 2) 분해도 함께 주어진다. 노름의 차수는 으뜸아이디얼에 따른 잉여류환이 부정원들에서 유도된 부분체에 관하여 갖는 차수를 준다.

이에 따라 임의의 아이디얼의 영점 문제는 \emph{노름과 초등인자의 분해정리와 일반 아이디얼 이론의 분해정리 사이의 관계}(§ 5)로 환원된다. 아이디얼의 영점들이 그 개별 대응 소아이디얼들의 영점들로 이루어진다는 사실에는 노름을 \emph{극대 으뜸인자}들로 분해하는 것이 대응한다. 이 인자들은 각 대응 소아이디얼의 초등인자형의 거듭제곱과 같다.\footnote{소거이론과 일반 아이디얼 이론의 이 평행성은 크로네커의 소거이론에서는 성립하지 않는다. H.-N., 주 4)를 보라.} 이로써 노름의 일차인수분해가 일반적으로 이루어진다. H.-N.에서는 중복도를 기초아이디얼들과 그 성분들로 해석한다. 노름의 한 극대 으뜸인자의 차수는 \(x_{i+1},\ldots,x_n\)을 첨가했을 때 \((i-1)\)단계 기초아이디얼을 \(i\)단계 기초아이디얼의 한 성분으로 나눈 선형독립 잉여류들의 수와 같다. 이제 \((i-1)\)단계 기초아이디얼은 \((n-i)\)차원보다 높은 모든 소아이디얼들에 의해 유일하게 정의되는 고립성분과 같음이 증명된다. \(i\)단계 기초아이디얼의 한 성분은 노름인자에 대응하는 소아이디얼과 더 높은 차원의 모든 소아이디얼들에 의해 정의되는 고립성분과 같다. 그리고 대응하는 노름인자의 차수는 \(x_{i+1},\ldots,x_n\)을 첨가했을 때, 이 성분의 잉여류들 가운데 \((i-1)\)단계 기초아이디얼의 류들만으로 이루어진 부분환의 차수와 같다.

\emph{소아이디얼과 진으뜸아이디얼의 특징짓기}(§ 6)는 잉여류체를 고찰하면 직접 따라온다. 원래 계수영역이 완전체이면 소아이디얼에는 노름과 초등인자형으로 소함수들이, 진으뜸아이디얼에는 진으뜸함수들이 대응하며, 역도 성립한다. 계수영역이 비완전체이면 필요조건 또는 충분조건만 줄 수 있다. 예들이 보여 주듯 한 소아이디얼의 노름이 진으뜸함수가 될 수도 있고, 한 진으뜸아이디얼의 초등인자형이 소함수가 될 수도 있다. 더 정확히 말하면 표수 \(p\)에서 소아이디얼의 노름은 그 초등인자형의 \(p^g\)제곱(\(g\geqq0\))이 된다. 반면 초등인자형이 소함수인 진으뜸아이디얼의 노름은 임의의 거듭제곱이 될 수 있으며, 이 역시 예들이 보여 준다. 끝으로 \emph{절대소아이디얼}(§ 7), 곧 계수영역을 대수적으로 닫힌 체로 확대해도 소아이디얼로 남는 아이디얼들을 고찰한다. 한 유한 대수적 수체의 계수를 갖는 절대소아이디얼의 경우, 위 특징짓기에서 A. Ostrowski가 절대소함수에 대하여 처음 세운 정리의 일반화가 나온다. 유한개의 예외를 제외하면 수체의 모든 소아이디얼을 법으로 해도 소아이디얼이라는 성질이 유지된다.

이 마지막 문제와 대수적 수체의 아이디얼 이론 사이의 유사성도 지적해 두자. 그러한 아이디얼의 초등인자형은 그 아이디얼로 나누어지는 가장 작은 양의 유리정수라고 보아야 한다(주 10). 따라서 소아이디얼의 경우에는 언제나 소수가 된다. 역으로 한 아이디얼의 노름이 소수인 순간 그 아이디얼은 소아이디얼이다. 증명도 완전히 평행하다. 그러므로 위에서 \emph{말한 예들}은 비완전체에서 고차 소아이디얼과 분기아이디얼의 유사물이 나타나는 반면, 완전체에서는 1차 소아이디얼만 있고 분기아이디얼은 없음을 보인다.

\subsection*{§ 1. 소거이론·일반 아이디얼 이론·체론의 기본개념}

\textbf{1. 영역.} \(\bar{\mathfrak{J}}\)를 추상적으로 정의된 체 \(\mathsf{P}\)의 계수를 갖는 \(y_1,\ldots,y_n\)의 모든 다항식 정역이라 하자. \(y\)들에 부정원 \(u_{\mu\nu}\)를 계수로 갖는 변환을 가한다.\footnote{§ 3 등을 고려하여 이 변환은 H.-N.의 것보다 조금 더 일반적이며, 따라서 그곳의 모든 결과는 더욱 그대로 유지된다.}
\srcnumdisplay{(1)}{y_1=u_{11}x_1+\cdots+u_{1n}x_n;\ \ldots;\ y_n=u_{n1}x_1+\cdots+u_{nn}x_n}
그 역변환은
\srcnumdisplay{(1')}{x_1=t_{11}y_1+\cdots+t_{n1}y_n;\ \ldots;\ x_n=t_{1n}y_1+\cdots+t_{nn}y_n}
이다. 서로 유리적으로 나타낼 수 있으므로 같은 뜻인 \(u_{\mu\nu}\), 또는 \(t_{\mu\nu}\)를 \(\mathsf{P}\)에 첨가하여 계수영역 \(\mathsf{P}(u)=\mathsf{P}(t)\)를 얻는다. \(\mathfrak{J}\)를 \(\mathsf{P}(u)\)의 계수를 갖는 \(x_1,\ldots,x_n\)의 모든 다항식 정역이라 하자. 영역 \(\mathfrak{J}\)와 \(\bar{\mathfrak{J}}\)가 소거이론의 바탕이고, \(\mathfrak{J}\)의 아이디얼 \(\fraka,\frakb,\ldots\)와 \(\bar{\mathfrak{J}}\)의 아이디얼 \(\bar{\fraka},\bar{\frakb},\ldots\)를 고찰한다. 임의의 환에서 아이디얼은 통상적으로 \(\alpha,\beta\)와 함께 \(\alpha-\beta\)를 포함하고, \(\alpha\)와 함께 \(\lambda\alpha\)도 포함하는 것으로 정의한다. 여기서 \(\lambda\)는 임의의 환원소이다. \(\frako\)는 언제나 모든 원소로 이루어진 단위아이디얼을 뜻한다.

\textbf{2. 변환된 아이디얼.} \(\mathfrak{J}\)의 아이디얼 \(\frakm\)이 \(\bar{\mathfrak{J}}\)의 한 아이디얼 \(\barfrakm\)에서 (1)에 따라 \(u_{\mu\nu}\), 또는 \(t_{\mu\nu}\)를 첨가하여 생길 때 이를 변환된 아이디얼이라 한다. \(\bar f(y)\), 또는 \(\bar g(y)\)가 \(\barfrakm\)의 모든 다항식을 돌 때 \(\frakm\)은 모든 선형결합
\[
 f(x)=\sum U_\lambda(u)\bar f_\lambda(y)=\sum U_\lambda(u)f_\lambda(x)
\]
또는
\[
 g(x)=\sum T_\lambda(t)\bar g_\lambda(y)=\sum T_\lambda(t)g_\lambda(x)
\]
로 이루어진다. 특히 모든 \(f_\lambda(x)=\bar f_\lambda(y)\), 또는 같은 뜻으로 모든 \(g_\lambda(x)=\bar g_\lambda(y)\)가 들어 있다. \(f(x)\)와 \(g(x)\)는 \(\mathsf{P}(u)=\mathsf{P}(t)\)의 양만큼을 제외하고 정해지므로, 일반성을 잃지 않고 \(U_\lambda,T_\lambda\)를 각각 \(u,t\)의 거듭제곱곱이라고 가정할 수 있다. \(U,T\)를 임의의 다른 거듭제곱곱으로 바꾸어도 \(f(x),g(x)\)는 다시 \(\frakm\)의 다항식으로 간다. 특히 (1)의 열들, 또는 (1')의 행들을 맞바꾸어도 그러하다. 따라서 \(\frakm\)은 임의의 한 다항식과 함께 \(x\)들을 맞바꾸어 얻는 모든 다항식을 포함한다. 다만 \(u\) 또는 \(t\)에 의존하는 계수들에서도 대응하는 치환을 해야 한다. 이하에 나오는 모든 아이디얼은 반대말을 명시하지 않는 한 변환된 것이라고 가정한다. 변환되지 않은 아이디얼도 (1)과 \(x_i=v_{i1}z_1+\cdots+v_{in}z_n\)을 합성하면 \(\mathsf{P}(u,v)\)의 계수를 갖는 \(z\)의 다항식영역 안의 변환된 아이디얼로 간다. 반면 변환된 아이디얼에서는 이 합성은 \(u_{\mu\nu}\)를 \(u,v\)의 쌍선형 결합들로 바꾸는 것일 뿐이다.

\textbf{3. 표기.} \(f^{(i)},g^{(i)},\ldots,a^{(i)},b^{(i)},\ldots\)는 시종일관 \(\mathfrak{J}\)의 다항식 가운데 \(x_1,\ldots,x_{i-1}\)을 포함하지 않는 것들을 뜻한다.

\textbf{4. 기초아이디얼.} 각 아이디얼 \(\frakm\)에는 다음 규정에 따라 \(n\)개의 기초아이디얼 \(\frakg_0,\frakg_1,\ldots,\frakg_{n-1}\)이 대응한다. \((i-1)\)단계 기초아이디얼 \(\frakg_{i-1}\)은 \(b^{(i)}G(x)\equiv0(\frakm)\)이 되도록 하는 0 아닌 다항식 \(b^{(i)}\)가 존재하는 모든 \(G(x)\), 그리고 오직 그러한 다항식들만을 포함한다. 일반적으로 \(b^{(i)}\)는 \(G(x)\)에 따라 달라진다. 아이디얼 기저의 존재에서 모든 \(G(x)\)에 공통으로 쓸 수 있는 \(B^{(i)}\)가 적어도 하나 존재함도 따라온다. 따라서
\begin{equation}
 B^{(i)}\frakg_{i-1}\equiv0(\frakm),\qquad B^{(i)}\ne0 . \tag{2}
\end{equation}
변환된 아이디얼의 기초아이디얼도 변환된 아이디얼이다(H.-N. 정리 VI). \(i\)단계 기초아이디얼 \(\frakg_i\)는 \(x_{i+1},\ldots,x_n\)을 \(\mathsf{P}(u)\)에 첨가했을 때 \(\frakm\)이 옮겨가는 아이디얼로도 정의된다. 이때 \(x_{i+1},\ldots,x_n\)에 관하여 \emph{다항식인} 것들로 제한해도 일반성을 잃지 않는다. \(\frakg_0\)은 \emph{언제나} 단위아이디얼 \(\frako\)와 같다.

\textbf{5. 아이디얼의 모듈표현, 초등인자와 개별노름.} 각 다항식 \(f(x)\)을 \(x_1,\ldots,x_{i-1}\)의 거듭제곱곱 \(\xi\)에 관한 선형형식으로 보고 다항식 \(a^{(i)}\)를 계수로 삼으면, 아이디얼 \(\frakm\)은 \(a^{(i)}\)들의 영역에 관한 선형형식들의 가군 \(\frakM_{i-1}\)으로 간다. 곧 \(\frakM_{i-1}\)은 임의의 두 선형형식과 함께 그 차를 포함하고, 한 선형형식 \(l(\xi)\)와 함께 임의의 \(a^{(i)}\)에 대한 \(a^{(i)}l(\xi)\)도 포함한다. 그러한 가군 \(\frakM\)의 기초모듈은 \(b^{(i)}g(\xi)\equiv0(\frakM)\), \(b^{(i)}\ne0\)이 되도록 하는 모든 선형형식 \(g(\xi)\)의 집합이다. 따라서 모듈표현에서 기초아이디얼 \(\frakg_{i-1}\)은 \(\frakM_{i-1}\)의 기초모듈 \(\frakG_{i-1}\)으로 간다. \(\frakM_{i-1}\)과 \(\frakG_{i-1}\)은 무한히 많은 부정원 \(\xi\)에 의존하지만, 각 선형형식에는 그 가운데 유한개만 나타난다. \(\frakM_{i-1}\)은 \(x_{i+1},\ldots,x_n\)을 \(\mathsf{P}(u)\)에 첨가했을 때 \emph{단위원과 다른 초등인자를 유한개만 갖는} 특징적인 성질을 가진다. 곧 이 첨가 뒤 \(\frakG_{i-1}\)과 \(\frakM_{i-1}\)은 다음 기저표현을 갖는다. 여기서 \(\zeta\)들은 새 부정원이며, \(\xi\)들과 \(x_i\)에 관하여 다항식인 가역 선형변환으로 이어진다. \(\zeta_i=r_i(\xi)\), \(\xi_k=s_k(\zeta)\)에는 매번 이 부정원들 가운데 유한개만 나타난다.
\begin{align}
\frakG_{i-1}&=(\zeta_0,\zeta_1,\ldots,\zeta_\sigma,\zeta_{\sigma+1},\ldots,\zeta_{\sigma+\nu},\ldots), \notag\\
\frakM_{i-1}&=(E^{(i)}\zeta_0,E_1^{(i)}\zeta_1,\ldots,E_\sigma^{(i)}\zeta_\sigma,
\zeta_{\sigma+1},\ldots,\zeta_{\sigma+\nu},\ldots), \tag{3}
\end{align}
여기서 각 \(E_\mu^{(i)}\)는 그 앞의 것을 나눈다(H.-N. (24), (28)).

\emph{최고 초등인자} \(E^{(i)}\)는 (2)에 나타나는 모든 \(B^{(i)}\)의, 다항식 의미에서의 최대공약인자로도 정의된다. \(x_{i+1},\ldots,x_n\)에 관하여 다항식이고 원시적이라고 잡을 수 있으며, 그러면 \(x_i\)에 관하여 정칙이다. 곧 \(x\)들에 관한 차수가 \(\lambda\)이면 0 아닌 계수를 갖는 항 \(x_i^\lambda\)를 포함한다.

(3)에 나타나는 모든 초등인자의 곱 \(R^{(i)}\)을 \(\frakG_{i-1}\)에서 \(\frakM_{i-1}\)로 가는 노름, 또는 아이디얼 \(\frakm\)의 \emph{개별노름}이라 한다. \(x_{i+1},\ldots,x_n\)을 첨가하면 \(\frakm\)이 \(\frakg_i\)로 간다는 점을 고려하여 기호로 쓰면
\begin{equation*}
 R^{(i)}=E^{(i)}E_1^{(i)}\cdots E_\sigma^{(i)}
 =N(\frakG_{i-1}\mid\frakM_{i-1})=N(\frakg_{i-1}\mid\frakg_i).
\end{equation*}
(3)이 보여 주듯 \(x_{i+1},\ldots,x_n\)을 첨가하면 \(R^{(i)}\)은 \(\frakG_{i-1}\)에서 \(\frakM_{i-1}\)로 가는 전이치환의 행렬식과 같다. \(R^{(i)}\)의 차수는 \(\mathsf{P}(u,x_{i+1},\ldots,x_n)\)에 관하여 선형독립인 \(\frakG_{i-1}\)에 따른 \(\frakM_{i-1}\)의 잉여류들, 따라서 \(\frakg_{i-1}\)에 따른 \(\frakg_i\)의 잉여류들의 수를 준다. \(R^{(i)}\)은 \(x_{i+1},\ldots,x_n\)에 관하여 다항식이고 원시적이라고 잡을 수 있으며, 그러면 \(x_i\)에 관하여 정칙이다. 또한 \(R^{(i)}\)은 다음 성질을 갖고 언제나 존재하는 랭크 \(\varrho\)의 가군 \(\frakM^*_{i-1}\)의 모든 \(\varrho\)행 행렬식들의 다항식 의미의 최대공약인자로 계산할 수 있다. \(\frakM^*_{i-1}\)의 기초모듈을 \(\frakG^*_{i-1}\)이라 할 때, 잉여류계 \(\frakG^*_{i-1}\mid\frakM^*_{i-1}\)은 \(\frakG_{i-1}\mid\frakM_{i-1}\)과 동형이다(H.-N. 정리 VII과 § 5).

\textbf{6. 아이디얼의 초등인자형과 노름(종결식형).} 모든 최고 초등인자의 곱 \(E_\frakm=E^{(1)}E^{(2)}\cdots E^{(n)}\)을 \(\frakm\)의 \emph{초등인자형}, 모든 개별노름의 곱 \(R_\frakm=R^{(1)}R^{(2)}\cdots R^{(n)}\)을 \emph{노름(종결식형)}\footnote{H.-N.에서는 종결식형이라는 명칭만 쓴다. 노름이라는 명칭은 산술적 성질에 대응한다.}이라 한다. 초등인자형과 노름은 아이디얼로 나누어지고, 노름은 초등인자형으로 나누어지며, 초등인자형의 어떤 거듭제곱은 노름으로 나누어진다. 더 일반적으로
{\csname tagsleft@true\endcsname\begin{align}
 E^{(i)}\frakg_{i-1}&\equiv0(\frakg_i),
& E^{(n)}\cdots E^{(i)}\frakg_{i-1}&\equiv0(\frakm)\quad\text{따라서}\notag\\
 R^{(i)}\frakg_{i-1}&\equiv0(\frakg_i),
& R^{(n)}\cdots R^{(i)}\frakg_{i-1}&\equiv0(\frakm),\quad\text{(H.-N. 정리 VIII),} \tag*{\raisebox{8.5pt}[0pt][0pt]{(4)}}
\end{align}}
이 성립한다. 더 나아가 \(\frakm\)이 \(\frakn\)으로 나누어지고 노름 \(R_\frakm,R_\frakn\)이 일치하면 \emph{아이디얼 \(\frakm,\frakn\)도 일치한다}(H.-N. 정리 IX). \(\frakm\)의 한 인자에 대하여 \(R^{(1)},R^{(2)},\ldots\)가 일치하면 기초아이디얼 \(\frakg_1,\frakg_2,\ldots\)도 일치하고 언제나 \(\frakg_0=\frako\)임을 고려하면, 이에 맞추어 다음을 얻는다. \(\frakn\)이 \(\frakm\)의 \emph{진인자}이면 \(\frakn\)의 개별노름 가운데 \(\frakm\)의 대응하는 개별노름과 처음으로 다른 것은 \emph{진인자}이다.\footnote{반면 \(R_\frakn\)의 뒤쪽 인자들은 \(R_\frakm\)의 대응 인자들의 배수가 될 수 있다. 예컨대 \(\frakm\)이 소아이디얼이고 \(\frakn\)이 단위아이디얼이 아닐 때에는 언제나 그렇다.} \emph{\(\frakm\)이 0 아닌 다항식 \(G^{(i)}\ne0\)를 포함하면} 정의상 \(\frakg_0,\frakg_1,\ldots,\frakg_{i-1}\)은 단위아이디얼과 같다. 따라서 개별노름이 선형독립 잉여류들의 수를 정한다는 5의 성질에 의해 \(R^{(1)},\ldots,R^{(i-1)}\)은 \emph{단위원}과 같고, 그러므로 \(E^{(1)},\ldots,E^{(i-1)}\)도 \emph{단위원}과 같다.

\textbf{7. 개별노름과 초등인자의 분해, 기초아이디얼의 성분들.} \(\mathsf{P}(u)\)에 관하여 기약인 다항식을 소함수, 소함수의 거듭제곱을 으뜸함수라 한다. 동시에 소함수가 아닌 으뜸함수, 곧 1보다 높은 거듭제곱을 진으뜸함수라 한다. 한 다항식의 극대 으뜸인자들로의 분해란 각 인자가 으뜸함수이지만 임의의 두 인자의 곱은 더 이상 으뜸이 아닌 분해이다.\footnote{소함수와 으뜸함수의 정의는 11과 정확히 나란히 아이디얼이라는 말만 다항식으로 바꾸어도 된다. 극대 으뜸인자들은 12의 분해에서 극대 으뜸성분들에 대응한다.} 개별노름 \(R^{(i)}=N(\frakg_{i-1}\mid\frakg_i)\)을 극대 으뜸인자 \(Q_\lambda^{(i)}\)들로 분해하는 데에는 \(\frakg_i\)를 최소공배수
\[
\frakg_i=[\frakr_1,\frakr_2,\ldots,\frakr_s]
\]
로 나타내는 것이 대응하며, \(Q_\lambda^{(i)}=N(\frakg_{i-1}\mid\frakr_\lambda)\)이다.
이때 \(\frakr_\lambda\)는 \(\frakg_{i-1}\)을 \((i-1)\)단계 기초아이디얼로 가지며 그 \(i\)단계 기초아이디얼과 일치한다. \(\frakr_\lambda\)들을 \emph{기초아이디얼의 성분들}이라 한다. \(\frakr_\lambda\)의 최고 초등인자는 \(Q_\lambda^{(i)}\)와 \(E^{(i)}\)의 다항식 의미의 최대공약인자, 따라서 \(E^{(i)}\)의 한 극대 으뜸인자와 같다. \(\frakr_\lambda\)들은 위에서 말한 기초아이디얼들에 관한 거동과, 개별노름 또는 최고 초등인자가 각각 \(R^{(i)}\) 또는 \(E^{(i)}\)의 극대 으뜸인자가 되어야 한다는 요구에 의해 유일하게 결정된다(H.-N. 정리 X).

\textbf{8. 차원개념.} \(R^{(i)}\)이 단위원과 다른 첫 개별노름일 때 아이디얼 \(\frakm\)의 최고차원을 \(n-i\)라 한다. 6에 따라 같은 뜻으로 \(\frakg_i\)가 단위아이디얼과 다른 첫 기초아이디얼일 때이다. 단위아이디얼 \(\frako\)의 차원은 \(-1\)이라 한다. 단위원과 다른 개별노름 \(R^{(i)}\)이 하나뿐이어서 \(\frakg_i=\frakg_{i+1}=\cdots=\frakm\)이면 최고차원을 단순히 \(\frakm\)의 차원이라고도 한다. \(\frakm\)이 0 아닌 다항식 \(G^{(i)}\ne0\)을 포함하면 6의 끝에 있는 주에서 알 수 있듯 그 최고차원은 기껏해야 \(n-i\)이다. \(\frakm\)의 최고차원이 \(n-i\)이고 \(\frakn\)이 \(\frakm\)의 인자이면 \(\frakn\)의 최고차원도 기껏해야 \(n-i\)이다.

\textbf{9. 노름(종결식형)과 영점.} 아이디얼 \(\frakm\)의 영점이란 \(\mathsf{P}(u;x_{i+1},\ldots,x_n)\)의 어떤 적당한 대수적 확대체에 속하는 \(x_1,\ldots,x_i\)의 값의 계들(\(i=1,2,\ldots,n\)) 가운데, \(x_{i+1},\ldots,x_n\)을 계수영역에 첨가했을 때 \(\frakm\)의 모든 다항식이 사라지는 것들을 뜻한다. 노름(종결식형)의 각 인자 \(R^{(i)}\)에는 이 첨가 뒤 그 차수만큼의 \(\frakm\)의 영점이 대응한다. \(R^{(i)}\)을 \(u_{\mu\nu}\)와 다른 부정원 \(v\)의 쌍선형 결합 \(w_{\mu\nu}\)로 쓰면 다음 명시적 분해를 갖는다.
\begin{equation}
 R^{(i)}=\prod\{z-(v_1\bar{x}_{1\nu}+\cdots+v_{i-1}\bar{x}_{i-1,\nu}+v_i\bar{x}_{i,\nu})\}^{\lambda_\nu}. \notag
\end{equation}
여기서 \(\bar{x}_{1\nu},\ldots,\bar{x}_{i\nu}\)는 서로 대응하는 영점들의 한 계이다(H.-N. 정리 XIII. § 3에서 새 증명을 줄 것이다). 따라서 최고차원 \(n-i\)인 아이디얼의 영점들 가운데 \((n-i)\)개 매개변수에 의존하는 것들은 있지만, 그보다 더 많은 매개변수에 의존하는 것은 없다. 이로써 8에서 준 차원개념이 통상적인 정식화와 일치함을 보였다.

\textbf{10. 일반 아이디얼 이론의 환영역.} 바탕 영역은 \emph{유한사슬 정리}가 성립하는 가환환이라 하자. 곧 아이디얼의 뜻에서 각 \(\fraka_\lambda\)가 바로 앞의 아이디얼의 진인자인 모든 사슬 \(\fraka,\fraka_1,\ldots,\fraka_\nu,\ldots\)은 유한 단계에서 끝난다. 이 요구는 모든 아이디얼이 아이디얼 기저를 갖는다는 요구와 동등하므로(\emph{아이디얼 이론}, 정리 I), 특히 다항식영역에서 만족한다. 이어지는 11, 12, 13에서는 이 일반 환영역을 바탕으로 삼는다.

\textbf{11. 소아이디얼과 으뜸아이디얼.} 한 곱이 \(\frakp\)로 나누어지면 적어도 한 인자가 \(\frakp\)로 나누어지는 아이디얼 \(\frakp\)를 소아이디얼이라 한다. 한 곱이 \(\frakq\)로 나누어지면 한 인자 또는 각 인자의 어떤 거듭제곱이 \(\frakq\)로 나누어지는 아이디얼 \(\frakq\)를 으뜸아이디얼, 또는 으뜸이라 한다. 기호로 쓰면
\[
 \quad a\not\equiv0(\frakp),\quad b\not\equiv0(\frakp)\quad\text{이면}\quad ab\not\equiv0(\frakp);
\]
\[
 \quad a\not\equiv0(\frakq),\quad b^x\not\equiv0(\frakq)\ \text{가 모든 거듭제곱 \(x\)에 성립하면}\quad ab\not\equiv0(\frakq).
\]
각 으뜸아이디얼 \(\frakq\)에는 유일한 대응 소아이디얼 \(\frakp\)가 있다. 이는 \(\frakq\)의 인자이고 그 어떤 거듭제곱은 \(\frakq\)로 나누어진다. 곧 \(\frakq\equiv0(\frakp)\), \(\frakp^\rho\equiv0(\frakq)\)이다.
이때 \(\frakp\)는 그 어떤 거듭제곱이 \(\frakq\)로 나누어지는 모든 환원소로 이루어진다. 정의가 보여 주듯 모든 소아이디얼은 동시에 으뜸아이디얼이다. 동시에 소아이디얼이 아닌 으뜸아이디얼, 따라서 \(\rho>1\)인 것을 진으뜸아이디얼이라 한다.

\textbf{12. 분해정리.} 아이디얼을 최소공배수로 나타낸 표현 \(\fraka=[\frakq_1,\ldots,\frakq_\alpha]\)에서 어느 \(\frakq\)도 나머지 것들의 최소공배수, 곧 그 보충부분으로 나누어지지 않을 때 이를 최단표현이라 한다. 각 \(\frakq\)가 으뜸이지만 임의의 두 \(\frakq\)의 최소공배수가 으뜸이 아닐 때 극대 으뜸성분들에 의한 표현이라 한다. 모든 아이디얼은 \emph{유한개의 극대 으뜸성분에 의한 최단표현}을 갖는다. 서로 다른 두 그러한 표현에서는 \emph{성분들의 수와, 모두 서로 다른 대응 소아이디얼들이 일치한다}. 이렇게 유일하게 결정되는 소아이디얼들을 그 아이디얼의 소아이디얼 또는 대응 소아이디얼이라 한다(\emph{아이디얼 이론}, 정리 IX).

\textbf{13. 분해의 고립성분.} \(\frakr=[\frakq_{i_1},\ldots,\frakq_{i_\tau}]\)가 \(\fraka\)의 극대 으뜸성분들에 의한 적어도 하나의 최단표현에 나타나는 \(\frakq_{i_\lambda}\)들로 이루어지고, 이들의 대응 소아이디얼 가운데 어느 것도 \(\fraka\)의 나머지 소아이디얼들 가운데 하나로 나누어지지 않는다면 \(\frakr\)을 \(\fraka\)의 분해의 고립성분이라 한다. 분해의 \emph{고립성분들은 그 소아이디얼들에 의해 유일하게 결정된다}. 특히 고립된 극대 으뜸성분들은 유일하게 결정된다(\emph{아이디얼 이론}, 정리 XIII).

\textbf{14. 표수, 완전체와 비완전체, 제1종 확대.} \(\Omega\)에 포함되고 단위원에서 유도되는 소체가 유리수체의 유형이면 \(\Omega\)의 표수는 0이고, 이 소체가 소수 \(p\)에 따른 잉여류계의 유형이면 표수는 \(p\)라고 한다. \(\Omega\)의 계수를 갖는 모든 소함수가 어떤 적당한 확대체에서 서로 다른 일차인자들로 분해되면 \(\Omega\)를 완전체라 하고, 그렇지 않으면 비완전체라 한다. 표수 0인 모든 체는 완전체이다. 표수 \(p\)인 체가 완전체인 것은 각 원소와 함께 그 \(p\)제곱근도 포함할 때 그리고 오직 그때이다. 비완전체의 한 소함수는 \(x^{p^f}\)에 관한 \(r\)차 다항식이고, 어떤 적당한 확대체에서 서로 다른 \(r\)개 일차인자의 \(p^f\)제곱들로 분해된다. \(f\)를 지수라 한다. 어떤 적당한 확대체에서 서로 다른 일차인자들로 분해되어 지수 \(f\)가 0인 소함수를 제1종이라 한다. 모든 원소가 제1종, 곧 제1종 소함수의 영점인 확대를 제1종 확대라 한다. 제1종 원소를 첨가하여 생기는 모든 확대는 제1종이고, 완전체에는 제1종 확대만 있다(Steinitz, § 11과 § 13).

\textbf{15. 유한확대의 축약차수와 지수.} \(\mathfrak{L}\)가 비완전체 \(\Omega\)의 유한확대이면, \(\mathfrak{L}\) 안에는 \(\Omega\)에 관하여 제1종이고 \(\frK\)의 모든 제1종 원소, 그리고 오직 그 원소들만으로 이루어진 차수 \(r\)의 부분체 \(\mathfrak{L}_0\)가 있다. \(\frK\)의 모든 원소의 \(p^f\)제곱은 \(\mathfrak{L}_0\)에 속하고, 정확히 차수 \(rp^f\)인 원소들도 \(\frK\) 안에 있다. \(r\)을 축약차수, \(f\)를 이 유한확대의 지수라 한다(Steinitz, § 14, 1과 § 14, 5).

\subsection*{§ 2. 소아이디얼과 으뜸아이디얼의 초등인자형과 노름. 소아이디얼의 인자}

이제부터는 시종일관 § 1, 1에서 정의한 다항식영역을 다룬다. 먼저 소아이디얼과 으뜸아이디얼도 변환된 아이디얼로 제한하여 다룰 수 있음을 보여야 한다.

\textbf{보조정리 I.} \emph{\(\bar{\mathfrak{J}}\)의 한 소아이디얼 \(\bar{\frakp}\)을 변환하여 얻은 아이디얼 \(\frakp\)는 소아이디얼이고, \(\bar{\mathfrak{J}}\)의 한 으뜸아이디얼 \(\bar{\frakq}\)을 변환하여 얻은 \(\frakq\)는 으뜸아이디얼이다. 달리 말하면 소아이디얼 또는 으뜸아이디얼이라는 성질은 \(u_{\mu\nu}\)를 \(\mathsf{P}\)에 첨가해도 유지된다.}

\(\fraka\not\equiv0(\frakp)\), \(\frakb\not\equiv0(\frakp)\)라고 하자. 여기서 \(\fraka,\frakb\)는 반드시 변환된 아이디얼일 필요는 없다. 예컨대 \(\fraka,\frakb\)의 다항식 \(a(x),b(x)\)를 각각 \(a(x)\not\equiv0(\frakp)\), \(b(x)\not\equiv0(\frakp)\)이도록 택하자. (1)을 (1')로 역변환하면, 일반성을 잃지 않고 \(T\)들을 \(t_{\mu\nu}\)의 거듭제곱곱으로 두어
\[
 a(x)=\sum T_j \bar a_j(y);\qquad b(x)=\sum T_j \bar b_j(y)
\]
를 얻는다. 적어도 하나의 \(\bar a_j(y)\)와 하나의 \(\bar b_j(y)\)는 각각 \(\bar{\frakp}\)로 나누어지지 않아야 한다. \(T\)들을 어떤 사전식 순서로 정렬했을 때 \(\bar a_r(y)\not\equiv0(\bar{\frakp})\), \(\bar b_s(y)\not\equiv0(\bar{\frakp})\)가 각각 최고항이라 하자. 그러면 \(\bar a_r(y)\bar b_s(y)\not\equiv0(\bar{\frakp})\)이고, 따라서 \(a(x)b(x)\not\equiv0(\frakp)\), 곧 \(\fraka\frakb\not\equiv0(\frakp)\)이다. 이로써 \(\frakp\)가 소아이디얼임을 보였다. 이 증명은 소아이디얼에 따른 잉여류계가 영인자 없는 환이라는 사실, 그리고 그 성질이 부정원을 첨가해도 유지된다는 사실에 바탕을 둔다.

이에 대응하여 \(\fraka\not\equiv0(\frakq)\), 그리고 모든 거듭제곱 \(z\)에 대하여 \(\frakb^\varkappa\not\equiv0(\frakq)\)라고 하자. 아이디얼 기저가 존재하므로 \(\fraka\)의 어떤 \(a(x)\)와 \(\frakb\)의 어떤 \(b(x)\)를 \(a(x)\not\equiv0(\frakq)\), 그리고 모든 \(z\)에 대하여 \(b(x)^\varkappa\not\equiv0(\frakq)\)이도록 택할 수 있다. 그렇지 않다고 가정하면 \(\frakb\)의 어떤 거듭제곱이 \(\frakq\)로 나누어질 것이기 때문이다.

\(a(x)b(x)=c(x)\)라 두고, \(\fraka^*,\frakb^*,\frakc^*\)를 각각 \(a(x),b(x),c(x)\)의 계수 \(\bar a_j(y),\bar b_j(y),\bar c_j(y)\)에서 유도되는 아이디얼이라 하자. 모든 거듭제곱 \(z\)에 대하여 \(\frakb^{*\varkappa}\not\equiv0(\bar{\frakq})\)이다. 그렇지 않으면 \(b(x)^\varkappa\)가 \(\frakq\)로 나누어지기 때문이다. \(\fraka^*\not\equiv0(\bar{\frakq})\)이므로 모든 거듭제곱 \(\lambda\)에 대하여 \(\fraka^*\frakb^{*\lambda}\not\equiv0(\bar{\frakq})\)이어야 한다. 그런데 데데킨트--메르텐스 정리\footnote{H.-N., 63쪽과 주 13)을 보라.}에 따라 \(\fraka^*\frakb^{*h}=\frakb^{*h-1}\frakc^*\)이 되게 하는 지수 \(h\)가 존재한다. 따라서 \(\frakc^*\not\equiv0(\bar{\frakq})\)이어야 한다. 이는 \(c(x)\not\equiv0(\frakq)\), 따라서 \(\fraka\frakb\not\equiv0(\frakq)\)을 준다. 이로써 \(\frakq\)가 으뜸임을 보였다.

한 아이디얼을 최소공배수로 나타낼 때에도 변환된 아이디얼들로 제한할 수 있다.

\textbf{보조정리 II.} \emph{\(\bar{\mathfrak{J}}\)에서 \(\barfrakm=[\bar{\frakc}_1,\ldots,\bar{\frakc}_\alpha]\)라는 표현이 있으면 \(\mathfrak{J}\)에서 \(\frakm=[\frakc_1,\ldots,\frakc_\alpha]\)라는 표현을 얻는다. 여기서 \(\frakc_\lambda\)들은 \(\bar{\frakc}_\lambda\)들을 변환하여 얻은 아이디얼들이다. 따라서 극대 으뜸성분들에 의한 모든 최단표현(§ 1, 12)에서도 이 성분들을 변환된 아이디얼들이라고 가정할 수 있다.}

\(\frakm\)의 모든 다항식 \(f(x)\)은 필요하면 \(u_{\mu\nu}\)의 적당한 거듭제곱을 곱한 뒤 \(f(x)=\sum U_i\bar f_i(y)\) 꼴이고, 각 \(\bar f_i(y)\)는 \(\barfrakm\)의 다항식으로서 가정상 모든 \(\bar{\frakc}_\lambda\)로 나누어지므로, \(\frakm\)은 \([\frakc_1,\ldots,\frakc_\alpha]\)로 나누어진다. 역으로 \(f(x)\)가 각 \(\frakc_\lambda\)로 나누어지면 위 꼴이고, 따라서 \(\frakm\)으로 나누어진다. 그러므로 \(\bar{\mathfrak{J}}\)의 극대 으뜸성분들에 의한 분해에서 출발하면 \(\mathfrak{J}\)에서 분해 \(\frakm=[\frakq_1,\ldots,\frakq_\alpha]\)를 얻는다. 여기서 \(\frakq\)들은 보조정리 I에 따라 으뜸아이디얼들의 변환이므로 으뜸이다. 또 \(\bar{\mathfrak{J}}\)의 서로 다른 대응 소아이디얼 \(\bar{\frakp}\)들에는 \(\mathfrak{J}\)의 서로 다른 \(\frakp\)들이 대응하므로 이들은 극대 으뜸성분들이다.

이제 소아이디얼과 으뜸아이디얼을 초등인자형과 노름으로 처음 특징짓는다. 아직 소인 것과 으뜸인 것을 완전히 분리하지는 않는다. 대수적 수체의 아이디얼 이론과 정확히 나란히 다음이 성립한다.\footnote{대수적 수체의 아이디얼에서 최고 초등인자는 그 아이디얼로 나누어지는 가장 작은 양의 유리정수라고 보아야 한다. 따라서 특히 소아이디얼의 경우에는 그 소아이디얼로 나누어지는 소수이고, 으뜸아이디얼(소아이디얼의 거듭제곱)의 경우에는 소수의 거듭제곱이다. 실제로 모든 아이디얼 \(\fraka^*\)는 동시에 체기저 \(\omega_1,\ldots,\omega_n\)의 원소들에 관한 선형형식들의 가군 \(\mathfrak{A}^*\)이고, \((\omega_1,\ldots,\omega_n)\)은 \(\mathfrak{A}^*\)의 기초모듈이다. 따라서 \(\fraka^*\)로 나누어지는 가장 작은 양의 유리정수가 가군 \(\mathfrak{A}^*\)의 최고 초등인자이고, \(\fraka^*\)의 노름은 \(\mathfrak{A}^*\)의 모든 초등인자의 곱이다.}

\textbf{정리 I.} \emph{한 소아이디얼의 초등인자형은 소함수이고, 그 노름은 이 소함수의 한 거듭제곱이다. 으뜸아이디얼의 경우 초등인자형과 노름은 으뜸함수이고, 더 정확히는 같은 소함수의 거듭제곱들이다. 이 소함수 자체는 대응 소아이디얼의 초등인자형이다. 소아이디얼과 으뜸아이디얼에서는 최고차원이 곧 차원이고, 으뜸아이디얼과 그 대응 소아이디얼의 차원은 일치한다.}

단위아이디얼에서는 초등인자형과 노름이 단위원과 같으므로 정리 I이 분명히 성립한다. 이제 소아이디얼 \(\frakp\)의 최고차원이 \((n-i)\)라고 하자. § 1, 6의 (4)에 따라
\[
 E^{(n)}\cdots E^{(i+1)}\frakg_i\equiv0(\frakp),\quad\text{그러나}\quad
 E^{(n)}\cdots E^{(i+1)}\not\equiv0(\frakp)
\]
이다. 그렇지 않으면 § 1, 8에 따라 \(\frakp\)의 최고차원은 기껏해야 \(n-i-1\)일 것이기 때문이다. 따라서 \(\frakg_i\equiv0(\frakp)\)이고, \(\frakp\)는 그 \(i\)단계 기초아이디얼, 그러므로 \(\frakg_{i+1},\frakg_{i+2},\ldots\)와도 같다. 따라서 \(R^{(i+1)}=N(\frakg_i\mid\frakg_{i+1}),R^{(i+2)},\ldots\)은 단위원과 같다. 그 인자인 \(E^{(i+1)},E^{(i+2)},\ldots\)도 단위원과 같다. 그러므로 초등인자형은 \(E^{(i)}\)과 같고, 이것은 한 소함수 \(P^{(i)}\)와 같아야 한다. 실제로 § 1, 5에 따라 \(\frakg_{i-1}\)이 단위아이디얼임을 고려하면 \(E^{(i)}\)은 \(\frakp\)로 나누어지는 모든 \(B^{(i)}\)의 다항식 의미의 최대공약인자이다. \(E^{(i)}=C_1^{(i)}C_2^{(i)}\equiv0(\frakp)\)이면 \(E^{(i)}\)은 그 인자 \(C_1^{(i)}\), 또는 \(C_2^{(i)}\) 가운데 하나로 나누어져야 한다. 또 \(E^{(i)}\)의 어떤 거듭제곱은 \(R^{(i)}\)으로 나누어지므로 \(R^{(i)}\)은 이 소함수 \(P^{(i)}\)의 한 거듭제곱이며, 첫 거듭제곱일 수도 있다. 동시에 \(R^{(i)}\)은 \(\frakp\)의 노름이다. 단위원과 다른 개별노름 \(R^{(i)}\)이 하나뿐이므로 끝으로 \(\frakp\)의 최고차원은 곧 차원이다.

이에 대응하여 으뜸아이디얼 \(\frakq\)의 최고차원이 \(n-i\)라고 하자. 그러면 위와 같이
\[
 E^{(n)}\cdots E^{(i+1)}\frakg_i\equiv0(\frakq),\quad
 (E^{(n)}\cdots E^{(i+1)})^\varkappa\not\equiv0(\frakq)
\]
가 모든 거듭제곱 \(z\)에 대하여 성립한다. 따라서 위와 같이 \(\frakq=\frakg_i\)이고, 그 초등인자형은 \(E^{(i)}\), 노름은 \(R^{(i)}\)이며, 최고차원은 곧 차원이다. 이 차원은 대응 소아이디얼의 차원과 같다. 실제로 \(\frakp^\rho\equiv0(\frakq)\), \(\frakq\equiv0(\frakp)\)에서 \(\frakp\)의 차원 \(n-i\)은 기껏해야 \(\frakq\)의 차원이고, \(\frakq\)의 차원은 기껏해야 \(\frakp^\rho\)의 차원임이 따른다. \(P^{(i)}\)이 \(\frakp\)의 초등인자형일 때 \((P^{(i)})^\rho\equiv0(\frakp^\rho)\)이므로 마지막 최고차원은 기껏해야 \(n-i\)이다. 따라서 \(n-i\)은 \(\frakp\)와 \(\frakq\)의 차원이다. 더 나아가 \((P^{(i)})^\rho\equiv0(\frakq)\)에서 \(E^{(i)}\)은 \(\frakq\)의 모든 \(B^{(i)}\)의 최대공약인자로서 \(P^{(i)}\)의 한 거듭제곱이며, 첫 거듭제곱일 수도 있다. 따라서 \(R^{(i)}\)에도 같은 사실이 성립한다. 이로써 정리 I의 모든 부분을 증명하였다.

차원개념에 의한 소아이디얼의 특징은 다음 정리가 준다.

\textbf{정리 II.} \emph{소아이디얼은 같은 최고차원을 갖는 진인자를 갖지 않는다. 역으로 그러한 진인자를 갖지 않는 모든 아이디얼은 소아이디얼이다.}

증명은 다음 보조정리에 바탕을 둔다.

\textbf{보조정리 III.} \emph{한 체 \(\Omega\)의 계수를 갖는 다항식들의 소아이디얼 \(\frakp\)가 \(\Omega\)에 관하여 선형독립인 잉여류를 유한개만 가지면, \(\frakp\)는 단위아이디얼과 다른 진인자를 갖지 않는다. 달리 말하면 \(\frakp\)에 따른 잉여류계는 체이다.}

가정은 다음을 뜻한다. 임의의 \(k+1\)개 잉여류 사이에는 \(\Omega\)의 계수, 더 정확히는 \(\Omega\)의 모든 원소가 나타내는 잉여류체 \((\Omega)\)의 계수를 갖는 선형종속관계가 있지만, \((\Omega)\)에 관하여 선형독립인 \(k\)개의 잉여류로 이루어진 계가 적어도 하나 있다. 여기서 모든 잉여류를 임의의 선형독립인 \(k\)개 계로 선형표현할 수 있음이 곧바로 따른다. 이제 \(\fraka\)가 \(\frakp\)의 진인자이고, \(a(z)\not\equiv0(\frakp)\)인 다항식 \(a(z)\)가 \(\fraka\)에 속한다고 하자. \(c_1f_1(z)+\cdots+c_kf_k(z)\not\equiv0(\frakp)\)이면
\[
 c_1a(z)f_1(z)+\cdots+c_ka(z)f_k(z)\not\equiv0(\frakp)
\]
도 성립한다. 곧 \(f_1,\ldots,f_k\)의 잉여류들뿐 아니라 \(af_1,\ldots,af_k\)의 잉여류들도 \(k\)개의 선형독립인 계를 이루며, 특히 단위원류를 이들로 선형표현할 수 있다. 따라서 \(\fraka\)는 단위아이디얼과 같다. 달리 말하면 잉여류계는 체이다. 실제로 \(a(z)\not\equiv0(\frakp)\)마다 \(1\equiv a(z)f(z)(\frakp)\)이 되는 \(f(z)\)가 존재한다.

정리 II의 첫 부분을 증명하려면 차원 \(n-i\)인 모든 소아이디얼, 더 일반적으로 최고차원 \(n-i\)인 모든 아이디얼이 \(\mathsf{P}(u;x_{i+1},\ldots,x_n)\)에 관하여 선형독립인 잉여류를 유한개만 가진다는 점만 주목하면 된다. \(\frakg_{i-1}\)이 여기서 단위아이디얼이므로 그 수는 (§ 1, 5) \(R^{(i)}\)의 차수가 주는 수이다. 따라서 \(\frakp\)의 모든 진인자 \(\fraka\)는 \(x_{i+1},\ldots,x_n\)을 \(\mathsf{P}(u)\)에 첨가하면 단위아이디얼로 간다. 달리 말하면 \(\fraka\)의 \(i\)단계 기초아이디얼은 단위아이디얼이고, 따라서 \(\fraka\)의 최고차원은 기껏해야 \(n-i-1\)이다. 변환되지 않은 아이디얼 \(\fraka\)도 인자로 허용하려면 § 1, 2에 따라 새로운 변환된 영역으로 넘어가기만 하면 된다.

역으로 \(\frakp\)가 같은 최고차원 \(n-i\)을 갖는 진인자를 갖지 않는다고 하자. \(\fraka\not\equiv0(\frakp)\), \(\frakb\not\equiv0(\frakp)\)이고 \(\fraka,\frakb\)는 다시 변환된 아이디얼이라고 하자. 아이디얼 의미의 최대공약인자 \((\fraka,\frakp)\)와 \((\frakb,\frakp)\)는 \(\frakp\)의 진인자이므로 가정에 따라
\[
 G^{(i+1)}\equiv0(\fraka,\frakp),\qquad
 H^{(i+1)}\equiv0(\frakb,\frakp),\qquad
 G^{(i+1)}\neq0;\qquad
 H^{(i+1)}\neq0
\]
을 얻는다. 따라서
\[
 G^{(i+1)}H^{(i+1)}\equiv0(\fraka\frakb,\frakp),\qquad
 G^{(i+1)}H^{(i+1)}\neq0.
\]
그러므로 \(\fraka\frakb\equiv0(\frakp)\)이면 \(\frakp\)의 최고차원이 \(n-i\)보다 작아질 것이므로 반드시 \(\fraka\frakb\not\equiv0(\frakp)\)이다. \(\fraka,\frakb\)가 변환된 아이디얼이므로 \(\bar{\frakp}\), 그리고 보조정리 I에 따라 \(\frakp\)도 소아이디얼이다. 이로써 정리 II를 증명하였다.

\subsection*{§ 3. 소아이디얼과 으뜸아이디얼의 영점}

먼저 소아이디얼에 관하여 영점이론(§ 1, 9)을 직접 정초하는 길을 열어 주는 것은 보조정리 III을 확장한 다음 보조정리이다. 이는 임의의 환의 소아이디얼에 대하여 성립한다.

\textbf{보조정리 IV.} \emph{단위아이디얼과 다른 임의의 소아이디얼에 따른 잉여류계는 분수몫을 만들어 한 체, 곧 그 소아이디얼의 잉여류체로 확대할 수 있다.}

임의의 아이디얼에 따른 잉여류계는 환을 이룬다. 합과 차와 곱이 다시 그 계에 속하고, 결합법칙·교환법칙·분배법칙은 잉여류로 넘어가도 유지되기 때문이다. 특별히 소아이디얼이면 이 환에는 영인자가 없다. 소아이디얼의 정의(§ 1, 11)가 사라지지 않는 잉여류들의 곱도 사라지지 않는다고 말하기 때문이다. 소아이디얼이 단위아이디얼과 다르면 이 환에는 영원소와 다른 원소가 적어도 하나 있으므로, 분수몫을 만드는 것, 곧 원소쌍들을 첨가하는 것으로 체까지 확대할 수 있다.\footnote{Steinitz, § 3. 원래 정역에 영과 다른 원소가 \emph{하나} 있다는 가정으로 충분하다. 그러면 분수몫을 만들 때 단위원의 존재가 보장되어 원래 영역이 그 체의 부분영역이 되기 때문이다. Steinitz는 정역에 단위원이 존재한다고 가정한다.} 이로써 잉여류체가 정의된다.

이하에서는 다음 \textbf{표기}를 일관되게 쓴다. 잉여류는 그 류의 임의의 한 원소를 괄호에 넣어 \((x),\ldots,(a(x)),\ldots\)처럼 나타내고, 잉여류들로 이루어진 체도 괄호로 나타낸다. 특히 \((P)\), \((\mathsf{P}(u))\), \((\mathsf{P}(u,x_{i+1},\ldots,x_n))\)는 각각 \(P\), \(\mathsf{P}(u)\), \(\mathsf{P}(u,x_{i+1},\ldots,x_n)\)의 원소들로 나타낼 수 있는 잉여류들 모두와 오직 그 잉여류들로 이루어진 잉여류체를 뜻한다.

또 다음을 미리 말해 둔다. 한 체 \(\mathfrak{K}\)이 바탕체 \(\Omega\)에 관하여 대수적 랭크(초월차수) \(k\)라는 것은, \(\mathfrak{K}\)이 \(\Omega\)에 관하여 대수적으로 독립인 \(k\)개 원소, 곧 초월원소들로 이루어진 계를 적어도 하나 포함하지만 \(\mathfrak{K}\)의 임의의 \(k+1\)개 원소는 \(\Omega\)에 관하여 대수적으로 종속이라는 뜻이다. \(\Omega\)에 관하여 대수적으로 독립인 \(z_1,\ldots,z_s\)에 대해 \(\mathfrak{K}\)을 부분체 \(\Omega(z_1,\ldots,z_s)\)의 대수적 확대체로 나타낼 수 있으면 반드시 \(s=k\)이다.\footnote{임의의 표수에서 성립하는 이 알려진 사실의 증명은 Steinitz, § 22, 9를 보라. \(u\)들을 첨가해도 대수적 랭크는 커질 수 없다. 원래 원소들과 \(\Omega(u)\)의 계수로 이루어진 유리적 결합만을 다루기 때문이다. 또 작아질 수도 없다. \(\mathfrak{K}\)의 원소들 사이의 \(\Omega(u)\)-계수 관계는 유한개의 \(\Omega\)-계수 관계로 분해되기 때문이다.} 마찬가지로 \(u\)들이 \(\mathfrak{K}\)에 들어 있지 않은 부정원들이면, \(\Omega(u)\)에 관한 \(\mathfrak{K}(u)\)의 대수적 랭크도 \(k\)이다.

이제 한 부정원의 소함수에 통상적으로 쓰는 방법\footnote{Steinitz, § 6.}과 완전히 나란하게 다음 정리에 따라 영점체를 구성한다.

\textbf{정리 III.} \emph{차원 \(n-i\)인 소아이디얼 \(\frakp\)의 잉여류체 \((\mathfrak{K})\)에는 그와 동형인 \(\frakp\)의 영점체 \(R\)을 대응시킬 수 있다. 이 체의 대수적 랭크는 \(n-i\), 곧 \(n-i\)개의 매개변수에 의존하며, 특히 \(\mathsf{P}(u,x_{i+1},\ldots,x_n)\)의 유한확대체로 볼 수 있다. \((\mathfrak{K})\)과 \(R\) 사이의 동형은 \((\mathsf{P})\)와 \(\mathsf{P}\), \((\mathsf{P}(u))\)와 \(\mathsf{P}(u)\) 사이의 동형을 포함하고, \(x_{i+1},\ldots,x_n\)을 매개변수로 삼으면 \((\mathsf{P}(u,x_{i+1},\ldots,x_n))\)와 \(\mathsf{P}(u,x_{i+1},\ldots,x_n)\) 사이의 동형도 포함한다. 역으로 대수적 랭크가 \(n-i\), 곧 \(n-i\)개의 매개변수에 의존하는 모든 영점체는 잉여류체 \((\mathfrak{K})\)과 동형이다.}

정리 III의 따름정리로 알려진 다음 정리를 얻는다. \emph{한 아이디얼 \(\fraka\)가 소아이디얼 \(\frakp\)의 모든 영점에서 사라지면, \(\fraka\)는 \(\frakp\)로 나누어진다.}

증명을 위해 먼저 \((\mathsf{P}(u))\)에 관한 잉여류체 \((\mathfrak{K})\)의 대수적 랭크가 \(n-i\)임을 주목한다. \((x_{i+1}),\ldots,(x_n)\)의 류들은 \((\mathsf{P}(u))\)에 관하여 대수적으로 독립이다. 차원이 \(n-i\)인 \(\frakp\)는 \(x_1,\ldots,x_i\)와 무관한 다항식을 포함하지 않기 때문이다. 그러나 다른 모든 류는 이 류들에 종속이다. 초등인자형 \(P^{(i)}\)이 \(\frakp\)에 속하므로 § 1, 2에 따라 \(\lambda=1,2,\ldots,i\) 각각에 대하여 \(x_\lambda,x_{i+1},\ldots,x_n\)에 의존하는 다항식도 \(\frakp\)에 속하기 때문이다. 따라서 \((\mathfrak{K})\)은 \((\mathsf{P}(u,x_{i+1},\ldots,x_n))\)의 유한확대체가 된다. § 1, 5에서 \(\frakg_{i-1}\)은 단위아이디얼이고 정리 I에 따라 \(\frakg_i=\frakp\)임을 고려하면, 이 부분체에 관한 차수는 노름의 차수와 같다.

동형인 영점체로 넘어가기 위해서는 \(\mathsf{P}(u,x_{i+1},\ldots,x_n)\)과 \((\mathsf{P}(u,x_{i+1},\ldots,x_n))\) 사이에 다음 동형이 있음을 더 주목해야 한다. \(\mathsf{P}(u,x_{i+1},\ldots,x_n)\)의 각 원소에 그 원소가 나타내는 잉여류를 대응시키면 된다. 우선 \(x_{i+1},\ldots,x_n\), 또는 \((x_{i+1}),\ldots,(x_n)\)의 다항식들로 이루어진 정역들이 이 대응으로 동형이다. \(\frakp\)에는 \(x_1,\ldots,x_i\)와 무관한 다항식이 없으므로 \(\mathsf{P}(u,x_{i+1},\ldots,x_n)\)의 서로 다른 다항식들에는 서로 다른 류들이 대응하기 때문이다. 동형인 정역들에서 분수몫을 만들면 동형인 체들이 생긴다. \(\mathsf{P}(u,x_{i+1},\ldots,x_n)\)과 \((\mathsf{P}(u,x_{i+1},\ldots,x_n))\) 사이의 이 동형은 \(\mathsf{P}\)와 \((\mathsf{P})\), \(\mathsf{P}(u)\)와 \((\mathsf{P}(u))\) 사이의 동형을 포함한다.

이 동형을 \((\mathfrak{K})\)과 한 영점체 \(R\) 사이의 동형으로 확대하기 위하여, 류 \((x_1),\ldots,(x_i)\)에 어떤 새 원소 \(\bar{x}_1,\ldots,\bar{x}_i\)를 동형으로 대응시키자. 곧 각 다항식 \((a(x))\)에는 다항식 \(a(\bar{x}_1,\ldots,\bar{x}_i,x_{i+1},\ldots,x_n)\)이 대응하고, 합에는 합이, 곱에는 곱이 대응하게 한다. 이제 분수몫을 만들면 --- 분모를 각각 부분체 \((\mathsf{P}(u,x_{i+1},\ldots,x_n))\)와 \(\mathsf{P}(u,x_{i+1},\ldots,x_n)\)의 원소들로 제한해도 된다 --- 잉여류체 \((\mathfrak{K})\)에 체 \(R\)이 대응한다. \(R\)은 \(\mathsf{P}(u,x_{i+1},\ldots,x_n)\)의 노름 차수의 유한확대체이며, 영점체라고 부를 수 있다. 실제로 이 대응은 \(f(x)\equiv0(\frakp)\)이면 \(f(\bar{x}_1,\ldots,\bar{x}_i,x_{i+1},\ldots,x_n)\)이 사라진다고 말한다. 이 전체 논의에서 \((\mathsf{P}(u,x_{i+1},\ldots,x_n))\) 대신 \((\mR)\)의 대수적으로 독립인 \(n-i\)개 원소를 첨가하여 생기는 다른 임의의 부분체를 써도 된다.

역으로 대수적 랭크가 \(n-i\)인 모든 영점체 \(R\)이 잉여류체와 동형임을 쉽게 보일 수 있다. \(R\)은 \(\bar{x}_1,\ldots,\bar{x}_n\)에서 \(\mathsf{P}(u)\)의 계수를 갖는 유리식으로 얻어지므로, 이 원소들 가운데 \(n-i\)개는 대수적으로 독립, 곧 \(\mathsf{P}(u)\)에 관하여 초월적이어야 한다. 특히 \(\bar{x}_{i+1},\ldots,\bar{x}_n\)이 그러해야 한다. \(P^{(i)}\equiv0(\frakp)\)에서 앞에서처럼 \(R\)이 \(\mathsf{P}(u,\bar{x}_{i+1},\ldots,\bar{x}_n)\)의 대수적 확대체임이 따르기 때문이다. 가정상 \(\frakp\)의 모든 다항식 \(f(x)\)에 대하여 \(f(\bar{x})\)가 사라지므로, \((\mR)\) 안의 \((x)\)들의 다항식영역에 속하는 각 류에는 \(\bar{x}\)들의 다항식인 \(R\)의 원소가 하나씩만 대응하고, 합에는 합이, 곱에는 곱이 대응한다.

서로 다른 류들에는 \(R\)의 서로 다른 원소들이 대응한다. 그렇지 않으면 영과 다른 어떤 류, 따라서 적당히 곱한 뒤에는 \((\mathsf{P}(u,x_{i+1},\ldots,x_n))\)의 다항식이 나타내는 어떤 류에 \(R\)의 영원소가 대응하게 되어, 가정에 어긋나게 \(\bar{x}_{i+1},\ldots,\bar{x}_n\) 사이에 대수적 종속관계가 생긴다. 이 대응은 \(R\)에 들어 있는 \(\bar{x}\)들의 모든 다항식을 소진한다. 이제 분수몫을 만들면 --- 다시 분모를 각각 \((\mathsf{P}(u,x_{i+1},\ldots,x_n))\)과 \(\mathsf{P}(u,\bar{x}_{i+1},\ldots,\bar{x}_n)\)의 원소들로 제한해도 된다 --- 이 동형인 정역들에서 동형인 체들이 생긴다.

이제 아이디얼 \(\fraka\)가 \(\frakp\)의 모든 영점, 특히 \(n-i\)개의 매개변수에 의존하는 모든 영점에서 사라진다고 하자. \(a(x)\)가 \(\fraka\)의 임의의 다항식이면 \(a(\bar{x})\)가 사라진다. \(R\)과 \((\mathfrak{K})\) 사이의 동형에 따라 \((a(x))\)은 영류이고, 따라서 \(a(x)\), 나아가 \(\fraka\)가 \(\frakp\)로 나누어진다. 이로써 정리 III과 따름정리를 증명하였다.

정리 III에서 § 1, 9에 준 분해로 넘어가려면, 우선 소아이디얼의 경우에도 지수에 관하여 더 정확한 명제를 얻기 위해 초등인자형을 더 조사해야 한다.

\textbf{보조정리 V.} \emph{표수가 \(p\)일 때 소아이디얼 또는 으뜸아이디얼의 초등인자형 \(E^{(i)}\), 더 일반적으로 최고차원과 단순한 차원이 일치하는 아이디얼 \(\fraka\)의 초등인자형이 \(x_i\)에 관하여 지수 \(f\), 곧 \(x_i^{p^f}\)의 다항식이면, \(x_{i+1},\ldots,x_n\)과 \(E^{(i)}\)에 나타나는 부정원 \(u_{\mu\nu}\), 또는 \(t_{\mu\nu}\)에 관해서도 지수 \(f\)이다. 특히 계수영역이 완전체이면 소아이디얼의 초등인자형 \(P^{(i)}\)은 언제나 \(x_i\)에 관하여 지수 0, 곧 제1종 소함수이다.}

여기서 \(\mathsf{P}\)가 표수 \(p\)의 완전체이면 \(\mathsf{P}(u,x_{i+1},\ldots,x_n)\)은 비완전체가 되므로, § 1, 14만으로는 \(P^{(i)}\)이 제1종이라고 곧바로 결론 내릴 수 없다는 점을 주의해야 한다. 이제 \(E^{(i)}\)이 \(x_i\)에 관하여 지수 \(f\)이지만 다른 부정원 \(x_{i+\lambda}\)를 지수 \(f'<f\)로 포함한다고 하자. \(u_{\mu\nu}\), 또는 \(t_{\mu\nu}\)들을 서로 바꾸면 § 1, 2에 따라 아이디얼로 나누어지는 영과 다른 다항식 \(G^{(i)}\)도 존재한다. 이는 \(x_i^{p^{f'}}\)의 다항식이고, 초등인자형의 정의에 따라 \(E^{(i)}\)으로 나누어진다. 그러나 \(G^{(i)}\)과 \(E^{(i)}\)의 차수가 같으므로, \(G^{(i)}\)은 \(\mathsf{P}(u)\)의 한 원소를 인자로 한 것만 제외하면 \(E^{(i)}\)과 일치해야 하고, 따라서 반드시 \(f'=f\)이다. 그러므로 \(t_{\mu\nu}\)들에 관하여 다항식이고 원시적이라고 가정할 수 있는 \(E^{(i)}\)은 다음 꼴이다.
\[
  E^{(i)}=\sum c_\lambda(t)\,x_i^{\lambda_i p^f}\cdots x_n^{\lambda_n p^f}
       =\sum c_\lambda(t)\,g_\lambda(y^{p^f},t^{p^f})
       =\sum T_\mu(t)\,\bar h_\mu(y^{p^f}),
\]
여기서 \(T_\mu\)는 \(t\)들의 거듭제곱곱이고, \(\bar h\)들은 \(\bar{\fraka}\)로 나누어진다. 따라서 \(T_\mu\)에서 각 \(t\)의 지수를 그 안에 포함된 \(p^f\)의 최대배수로 바꾸어도 \(\fraka\)에 속하는 다항식을 얻는다. 위 표현에서 보듯 이는 \(c_\lambda(t)\)에서 각 \(t\)의 지수를 그 안에 포함된 \(p^f\)의 최대배수로 바꾸는 것과 같다. 이 치환으로 \(E^{(i)}\)은 \(t,x_i,\ldots,x_n\)의 한 다항식으로 가는데, 정의에 따라 이 다항식은 \(E^{(i)}\)으로 나누어지며, 가정한 원시성 때문에 \(t\)들에 관해서도 그러하다. 어느 인수에 관해서도 차수가 늘지 않았으므로 이 다항식은 \(E^{(i)}\)과 동일해야 한다. 특히 \(\mathsf{P}\)가 완전체이면 \(E^{(i)}\)은 \(x^{p^f},t^{p^f}\)의 다항식으로서 \(p^f\)제곱이다. 따라서 소아이디얼을 다룰 때에는 반드시 \(f\)이고, 초등인자형 \(P^{(i)}\)은 제1종 소함수이다.

이제 분해는 다음 정리로 주어진다.

\textbf{정리 IV.} \emph{소아이디얼의 초등인자형 \(P^{(i)}\)은 영점체에서 유도된 갈루아체 안에서 다음과 같이 명시적으로 분해된다.}
\[
 P^{(i)}=\prod (x_i-t_{1i}\bar y_{1\nu}-\cdots-t_{ni}\bar y_{n\nu})^\delta
 =\prod (t_{1i}(y_1-\bar y_{1\nu})+\cdots+t_{ni}(y_n-\bar y_{n\nu}))^\delta,
\]
\emph{여기서 \(\bar y_{1\nu},\ldots,\bar y_{n\nu}\)는 각각 서로 대응하고 \(t_{1i},\ldots,t_{ni}\)와 무관한 원래 부정원 \(y\)들의 영점계이며, 서로 다른 \(\nu\)에는 서로 다른 일차인자가 대응한다. 지수 \(\delta\)는 \(\mathsf{P}\)가 완전체이면 1이고, 비완전체이면 \(p^f\) (\(f\geqq0\))이다. 이로써 새 부정원 \(v\)들을 첨가하고 \(u_{\mu\nu}\) 대신 \(u\)와 \(v\)의 쌍선형 결합을 넣어 쓴 초등인자형의 분해도 함께 정해진다.}
\[
 P^{(i)}(u,v)=\prod (z-v_1\bar x_{1\nu}-\cdots-v_i\bar x_{i\nu})^\delta,
\]
\emph{여기서 \(\delta\)는 위와 같은 뜻이다. 마찬가지로 소아이디얼의 노름, 또는 \(\frakp\)를 대응 소아이디얼로 갖는 으뜸아이디얼의 노름이 \(P^{(i)}\)의 거듭제곱으로 분해되는 것도 함께 주어진다.}

정리 IV의 따름정리는 다음과 같다. \emph{두 소아이디얼의 초등인자형 \(P^{(i)}\)이 일치하면 두 소아이디얼은 동일하다.}

정리 IV의 증명은 초등인자형 \(P^{(i)}\)이 \((\mathsf{P}(u,x_i,x_{i+1},\ldots,x_n))\)을 \((\mathsf{P}(u,x_{i+1},\ldots,x_n))\) 위에서 본 확대체의 기본방정식과 동일함을 보이는 것으로 귀착된다. 먼저 부정원 \(u_{\mu\nu}\), 또는 \(t_{\mu\nu}\)에 대한 의존성을 알아보기 위하여 \(\bar{\frakp}\)의 잉여류체 \((\bar{\mathfrak{K}})\)로 넘어가자. 대수적 랭크의 정의에서 말한 바에 따라 이 체는 \((\bar P)\)에 관하여 대수적 랭크 \(n-i\)이어야 한다. 여기서 \((\bar{\mathsf{P}})\)는 \(\mathsf{P}\)의 원소들이 나타내는 부분체이다. 실제로 \((\mathfrak{K})\)은 \((\bar{\mathfrak{K}})\)과 동형인 한 부분체에 부정원 \(u\), 또는 \(t\)를 첨가하여 생긴다. 이 부분체는 \(\bar{\mathfrak{J}}\)의 다항식들로 나타낼 수 있는 모든 류와 오직 그 류들에서 분수몫을 만들어 얻으며, \((\bar{\mathsf{P}})\)와 \((\mathsf{P})\)가 서로 대응한다. 따라서 \((\mathsf{P}(u))\)에 관한 \((\bar{\mathfrak{K}})\)의 대수적 랭크는 \((\bar{\mathsf{P}})\)에 관한 이 부분체의 대수적 랭크와 같고, 동형에 따라 다시 \((\mathsf{P})\)에 관한 \((\bar{\mathfrak{K}})\)의 대수적 랭크와 같다. \((\mathfrak{K})\)은 류 \((y_1),\ldots,(y_n)\)에서 \((\bar P)\)의 계수를 갖는 유리식으로 얻어지므로 이 류들 가운데 \((n-i)\)개는 대수적으로 독립이어야 하며, \((\bar{\mathfrak{K}})\)은 이 \((n-i)\)개 류를 \((\bar{\mathsf{P}})\)에 첨가하여 생기는 부분체의 유한 대수적 확대체이다.\footnote{이 관찰은 정리 III과 그 따름정리를 \(\bar{\frakp}\)의 영점체 \(\bar R\)에 대하여 직접 진술하고 증명할 수도 있었음을 보여 준다. 그러나 변환된 아이디얼로 넘어가야 임의의 아이디얼에서 같은 차원을 가진 최대 으뜸성분들의 영점에 같은 매개변수들이 나타난다. 그래야 비로소 임의의 아이디얼의 노름과 초등인자형과의 연결이 생긴다.}

\(x_{i+1},\ldots,x_n\)에 나타나는 부정원 \(t_{\sigma\tau}\)들만 첨가하면, 류 \((y_1),\ldots,(y_n)\)과 \((x_{i+1}),\ldots,(x_n)\)을 포함하고 \((\mathsf{P}(t_{\sigma\tau},x_{i+1},\ldots,x_n))\)의 유한 대수적 확대체인 또 하나의 잉여류체를 얻는다. 여기의 체들은 \((\mathfrak{K})\)의 부분체들과 동형이지만 동일하지는 않다. 부정원들이 다르면 잉여류들이 같은 원소들로 나타나더라도 같은 원소들을 포함하지 않아 동일하지 않기 때문이다. 따라서 \((\mathsf{P}(t_{\sigma\tau},x_{i+1},\ldots,x_n))\)에 대한 류 \((y)\)들의 의존성에는 부정원 \(t_{\sigma\tau}\)들만 들어간다.

이제 \(t_{1i},\ldots,t_{ni}\)를 첨가하여 류 \((x_i)=(t_{1i}y_1+\cdots+t_{ni}y_n)\)을 만들고,
\(
 (x_{i\nu})=(t_{1i}y_{1\nu}+\cdots+t_{ni}y_{n\nu})
\)
를 \((\mathsf{P}(t_{\sigma\tau},t_{1i},\ldots,t_{ni},x_{i+1},\ldots,x_n))\)에 관한 갈루아체 안의 서로 다른 \(r\)개 공액값이라 하자. 확대가 제1종인지 아닌지에 따라 인자 \(t-(x_{i\nu})\)들의 곱, 또는 그 \(p^f\)제곱들의 곱은 \((\mathsf{P}(t,x_{i+1},\ldots,x_n))\)에 관한 소함수 \((G)\)이다. 실제로 차수가 \(r\cdot p^f\)인 원소들이 존재하므로(§ 1, 15), 각 원소가 \(t_{1i},\ldots,t_{ni}\)의 특수화로 생기는 이상 \((x_i)\)도 반드시 그런 원소이다. 그리고 \((x_i)\)는 \((G)\)의 영점이다. \(\frakp\)의 동형인 영점체와 거기서 유도된 갈루아체로 넘어가면, 여기서 \(G\)는 일차인자 \(t-t_{1i}\bar y_{1\nu}-\cdots-t_{ni}\bar y_{n\nu}\)들, 또는 그 \(p^f\)제곱들로 분해된다. 한편 \((G)\)는 \(t=(x_i)\)에서 사라지므로 \(G(x_i)\)은 \(\frakp\), 따라서 \(P^{(i)}\)으로 나누어진다. \(G(x_i)\)은 \(\mathsf{P}(u,x_{i+1},\ldots,x_n)\)에 관한 소함수이고 \(x_{i+1},\ldots,x_n\)에 관하여 원시적이므로 \(P^{(i)}\)과 동일하다. 이로써 정리 IV의 첫 분해를 증명하였다.

둘째 분해로 넘어가기 위하여, 부정원을 첨가해도 제1종 확대인지 아닌지라는 성질은 유지됨을 주목하자. \(v_1,\ldots,v_i\)와 모든 \(t_{\mu\nu}\)를 첨가한 뒤 류 \((z)=(v_1x_1+\cdots+v_ix_i)\)와 공액류 \((z_\nu)=(v_1x_{1\nu}+\cdots+v_ix_{i\nu})\)를 만들면, 모든 \(t-(z_\nu)\)의 곱, 또는 그 \(p^f\)제곱들의 곱은 다시 \(t=(z)\)에서 사라지는 소함수 \((H)\)이다. 따라서 \(H\)는 위와 같이 일차인자들의 곱으로 나타낼 수 있고, 변환 (1)에 \(z=v_1x_1+\cdots+v_ix_i\)를 합성하여 \(\frakp\)에서 얻은 소아이디얼의 초등인자형과 동일하다. 상호 특수화 때문에 이 초등인자형은 \(P^{(i)}\)에서 \(u_{\mu\nu}\)들을 \(u\)와 \(v\)의 어떤 쌍선형 결합들로 바꾸어 얻어진다.\footnote{H.-N., § 7을 보라.} 이 분해들로써 노름이 \(P^{(i)}\)의 한 거듭제곱으로 분해되는 것도 주어진다.\footnote{제1종 확대, 특히 완전체에서는 첫째 거듭제곱이고, 비완전체에서는 \(p^g\)제곱이다. 이는 § 6의 정리 XIII과 XIV에서 보일 것이다.} 마찬가지로 \(\frakp\)를 대응 소아이디얼로 갖는 으뜸아이디얼의 노름과 초등인자형의 분해도 주어진다.

끝으로 두 소아이디얼의 초등인자형 \(P^{(i)}\)이 일치하면 위 분해에 따라 그 영점들도 일치한다. 따라서 각 아이디얼은 다른 아이디얼로 나누어지고, 두 아이디얼은 동일하다.

\subsection*{§ 4. 차원개념의 산술적 정식화}

변환된 아이디얼의 도입과 무관한 차원개념의 첫 정식화는 다음 정리가 준다.

\textbf{정리 V.} \emph{한 소아이디얼의 차원은 그 잉여류체의 대수적 랭크로 주어진다. 한 아이디얼의 최고차원은 그 대응 소아이디얼들의 차원수 가운데 가장 큰 것과 같다.}

정리 V의 첫 부분은 정리 IV의 첫째 주에서 이미 증명하였다. 거기서는 \(\bar{\frakp}\)의 잉여류체 \((\bar{\mathfrak{K}})\)의 대수적 랭크가 \(\frakp\)의 잉여류체 \((\mathfrak{K})\)의 대수적 랭크와 같고, 따라서 § 1, 8과 일치하게 \(\frakp\)의 차원과 같음을 보였다. 둘째 부분을 증명하기 위해 \(\frakm=[\frakq_1,\ldots,\frakq_\alpha]\)가 극대 으뜸성분들에 의한 최단표현이라고 하자. 우선 각 \(\frakq\)는 \(\frakm\)의 인자이므로 \(\frakm\)보다 높은 차원을 가질 수 없다. 한편 모든 \(\frakq\)의 곱은 \(\frakm\)으로 나누어지고, 따라서 개별 \(\frakq\)의 모든 초등인자형의 곱도 \(\frakm\)으로 나누어진다. 그러므로 \((n-i)\)가 \(\frakq\)들의 차원수 가운데 가장 크다면 \(\frakm\)은 영과 다른 다항식 \(G^{(i)}\ne0\)을 포함하고, 따라서 그 차원은 \(n-i\)보다 높을 수 없다. 정리 I에 따라 \(\frakq\)들의 차원수는 그 대응 소아이디얼들의 차원수와 일치한다. 변환된 아이디얼을 도입하지 않고 소아이디얼의 차원을 잉여류체의 대수적 랭크로 정의하면, 정리 V의 둘째 부분은 임의의 아이디얼의 최고차원을 정의한다.

한편 소아이디얼들의 사슬을 이용하여 역시 변환된 아이디얼의 도입과 무관하게 차원개념을 정식화하려면, 정리 II를 보충하여 다음을 보여야 한다.

\textbf{정리 VI.} \emph{단위아이디얼과 다른 모든 소아이디얼은 --- 필요하면 부정원들을 첨가한 뒤 --- 차원이 정확히 하나 작은 소아이디얼을 적어도 하나 인자로 갖는다.}

실제로 \(\frakp\)의 차원이 \((n-i)>0\)이고 \(v\)가 새 부정원이면
\[
 \frakc=(\frakp,x_i-v)
\]
는 원하는 성질을 가진 아이디얼이다. 이는 동형대응에서 따른다. 곧 \(\frakc=(\frakp^*,x_i-v)\)이며, 여기서 \(\frakp^*\)는 \(\frakp\)에서 \(x_i\)를 부정원 \(v\)로 바꾸고 \(v\)를 계수영역에 첨가하여 얻는다. 잉여류들도 마찬가지로 류 \((x_i)\)를 \((v)\)로 바꾸어 얻는다. 그런데 \(x_i\)를 \(\mathsf{P}(u)\)에 첨가해도 \(\frakp\)는 자기 자신으로 간다. 실제로
\[
 h(x_i)f(x)\equiv0(\frakp),\qquad h(x_i)\ne0
\]
이면 반드시 \(f(x)\equiv0(\frakp)\)이다. \(\frakp\)가 \(x_i\)에만 의존하는 다항식을 포함한다면 \(x_n\)에만 의존하는 다항식도 포함하여, 가정에 어긋나게 차원이 기껏해야 0이 되기 때문이다. 따라서 대응 \(v\sim x_i\) 아래에서 \(\frakc\)에 따른 영류에는 반드시 \(\frakp\)에 따른 영류가 대응하고, 그 역도 물론 성립한다. 이로써 \(\frakp\)와 \(\frakc\)에 따른 잉여류계들 사이에 동형대응이 생긴다. 뒤의 잉여류계에는 영인자가 없으므로 \(\frakc\)는 소아이디얼이다.

이 동형대응 아래에서 \((x_i)\)와 \((\mathsf{P}(u,x_{i+1},\ldots,x_n))\) 사이의 대수적 종속관계에는 \((\mathsf{P}(u,v))\)에 관한 \((x_{i+1}),\ldots,(x_n)\) 사이의 한 관계가 대응한다. 반면 \(\lambda=1,\ldots,i-1\)에 대하여 \((x_\lambda)\)와 \((\mathsf{P}(u,x_{i+1},\ldots,x_n))\) 사이의 종속관계들은 그대로 유지되어 대수적 랭크를 더 낮추지 않는다. 따라서 대수적 랭크는 정확히 하나 작아졌다. \(\frakp\)의 차원이 0이었다면 \(\frakc\)는 단위아이디얼이 되므로 위 결과는 여전히 성립한다. 따라서 원래 영역 \(\bar{\mathfrak{J}}\)에서는
\[
 \bar{\fraka}=(\bar{\frakp},v+v_1y_1+\cdots+v_ny_n),
\]
가 원하는 성질을 가진 아이디얼이다. 여기서 \(v_\lambda=-t_{\lambda i}\)로 두었다. 특히 \(\mathsf{P}\)가 무한개의 원소를 포함하면 \(v,v_1,\ldots,v_n\)을 언제나 \(\mathsf{P}\)의 원소들로 특수화하여, \(\bar{\fraka}\)의 노름과 초등인자형, 따라서 차원이 특수화된 아이디얼 \(\frakb\)의 것과 같게 할 수 있다.\footnote{H.-N., § 7의 마지막 단락을 보라.} 다만 이때 초등인자형은 더 이상 소함수로 남지 않을 수 있다. 정리 V에 따라 \(\frakb\)는 원하는 차원의 대응 소아이디얼을 적어도 하나 갖는다. 다시 \(\bar{\mathfrak{J}}\)로 돌아가면, 이 경우 정리 VI는 어떠한 부정원도 첨가하지 않고 성립한다.

정리 II와 정리 VI는 곧바로 원하는 정식화를 준다.

\textbf{정리 VII.} \emph{한 소아이디얼 \(\frakp\)의 차원이 \(n-i\)라는 것은 --- 필요하면 부정원들을 첨가한 뒤 --- 다음과 같은 \(1+(n-i)+1\)개의 소아이디얼로 이루어진 사슬이 적어도 하나 존재한다는 뜻이다.}
\[
 \frakp_0=\frakp;\quad \frakp_1,\ldots,\frakp_{n-i};\quad \frakp_{n-i+1},
\]
\emph{여기서 각 아이디얼은 바로 앞 아이디얼의 진인자이고, 항의 수가 더 많은 그러한 사슬은 존재하지 않는다. \(\mathsf{P}\)가 무한개의 원소를 포함하면 이 정의는 부정원을 전혀 도입하지 않고 원래 영역에서도 성립한다.}\footnote{차원의 이 사슬정의는 대수적 수체의 경우 통상적인 소아이디얼들에는 차원 0, 단위아이디얼에는 차원 \(-1\), 영아이디얼에는 차원 1을 준다. 실제로 영아이디얼도 소아이디얼의 정의를 만족한다. 체에서는 사라지지 않는 두 원소의 곱이 언제나 영과 다르기 때문이다. 이 차원정의 아래에서도 정리 II는 대수적 수체에서 성립한다.}

우선 사슬의 마지막 항은 소아이디얼의 정의를 만족하는 단위아이디얼이어야 한다. 그렇지 않으면 \(\frako\)를 덧붙여 사슬을 늘일 수 있기 때문이다. \(\frako\) 자체에는 정리 VII에 따라 § 1, 8과 일치하게 차원 \(-1\)이 주어진다. 이제 \(\frakp\ne\frako\)이고 그 차원이 \(n-i\)라 하자. 정리 II에 따라 같은 차원의 소아이디얼 두 개가 나타날 수 없으므로 사슬의 항 수는 기껏해야 \(1+(n-i)+1\)이다. 그러나 정리 VI에 따라 --- 필요하면 새 부정원들을 첨가하여 --- 각 항마다 차원이 정확히 하나씩 작아지는 이 항 수의 사슬이 적어도 하나 존재한다. 정리 VII를 차원의 정의로 삼으면 --- 이 정의는 원래 영역에서도 그대로 말할 수 있고, \(\mathsf{P}\)가 무한개의 원소를 가지면 정리 VI에 따라 부정원을 도입할 필요도 없다 --- 임의의 아이디얼의 최고차원은 다시 정리 V의 둘째 부분으로 정의된다.

\subsection*{§ 5. 기초아이디얼과 그 성분들을 일반 아이디얼 이론 안에 배치하기. 아이디얼의 영점}

여기서는 기초아이디얼과 그 성분들(§ 1, 7)을 일반 아이디얼 이론의 뜻에서 분해의 고립성분들(§ 1, 13)로 특징짓는다. 이 특징은 노름의 분해정리와 일반 아이디얼 이론의 분해정리가 서로 나란히 진행됨을 보여 주고, § 3에서 소아이디얼과 으뜸아이디얼에 대하여 준 영점이론을 임의의 아이디얼에 대하여 진술할 수 있게 한다. 기초아이디얼에 대해서는 다음이 성립한다.

\textbf{정리 VIII.} \emph{\(i\)단계 기초아이디얼 \(\frakg_i\)는 차원이 \(n-1,n-2,\ldots,n-i\)인 모든 대응 소아이디얼들에 의해 유일하게 결정되는 분해의 고립성분이다. 모든 대응 소아이디얼의 차원이 같으면 아이디얼 자체도 그 차원을 갖는다. 곧 노름에는 인자 \(R^{(i)}\) 하나만 나타난다. 일반적으로 노름에는 대응 소아이디얼들 사이에 서로 다른 차원수가 나타나는 수만큼 단위원과 다른 인자 \(R^{(i)},R^{(i+\sigma)},R^{(i+\tau)},\ldots\)가 있다.}

증명에 앞서 다음을 둔다.

\textbf{보조정리 VI.} \emph{한 아이디얼이 최소공배수 \(\frakm=[\fraka_1,\ldots,\fraka_\alpha]\)로 나타나 있으면, 그 기초아이디얼 \(\frakg_i\)는 \(\fraka\)들의 \(i\)단계 기초아이디얼 \(\frakh_i\)들의 최소공배수이다.}

실제로
\[
 b^{(i+1)}f\equiv0(\frakm),\qquad b^{(i+1)}\ne0
\]
에서
\[
 b^{(i+1)}f\equiv0(\fraka_\lambda),\qquad b^{(i+1)}\ne0
\]
이 따른다. 따라서 \(\frakg_i\)는 \([\frakh_{i1},\ldots,\frakh_{i\alpha}]\)로 나누어진다. 역으로
\[
 c_\lambda^{(i+1)}f\equiv0(\fraka_\lambda),\qquad c_\lambda^{(i+1)}\ne0
\]
에서
\[
 c_1^{(i+1)}\cdots c_\alpha^{(i+1)}f\equiv0(\frakm),\qquad
 c_1^{(i+1)}\cdots c_\alpha^{(i+1)}\ne0
\]
도 따르므로 반대쪽 나눗셈도 성립한다. 이로써 보조정리를 증명하였다.

이제 \(\frakm\)의 대응 소아이디얼들을 차원에 따라 다음과 같이 배열하자. 물론 어떤 차원들은 빠질 수도 있으며, 그때 \(j_\lambda=0\)이다.
\[
 \frakp_{1,1},\ldots,\frakp_{1,j_1};\quad
 \frakp_{2,1},\ldots,\frakp_{2,j_2};\quad\ldots;\quad
 \frakp_{n,1},\ldots,\frakp_{n,j_n}.
\]
일반적으로 \(\frakp_{\lambda,\kappa}\)는 차원 \(n-\lambda\)인 소아이디얼을 뜻하고, \(\frakq_{\lambda,\kappa}\)는 \(\frakm\)의 어떤 고정된 분해에 나타나는 그에 대응하는 으뜸아이디얼이라 하자. 정의에 따라 \(\lambda>i\)인 모든 \(\frakq_{\lambda,\kappa}\)의 \(i\)단계 기초아이디얼은 단위아이디얼이고, \(\lambda\leqq i\)이면 정리 I에 따라 그 기초아이디얼은 \(\frakq_{\lambda,\kappa}\)와 같다. 따라서 보조정리 VI에서
\[
 \tag{5}
 \frakg_i=[\frakq_{1,1},\ldots,\frakq_{1,j_1};\ldots;\frakq_{i,1},\ldots,\frakq_{i,j_i}]
\]
를 얻는다. 이 표현은 \(\frakg_i\)가 분해의 고립성분임을 보여 준다. \(\frakg_i\)에 속하는 어떤 소아이디얼도 모두 더 낮은 차원을 가진 나머지 소아이디얼 가운데 하나에 포함될 수 없기 때문이다.

특별히 대응 소아이디얼들 가운데 차원 \(n-i\)인 것들만 나타나면, (5)에 따라 \(\frakg_0,\ldots,\frakg_{i-1}\)은 단위아이디얼이고 \(\frakg_i=\frakg_{i+1}=\cdots=\frakm\)이다. 따라서 \(R_\frakm=R^{(i)}\)이다. 서로 다른 차원의 소아이디얼들이 있어서, 이를테면 \(j_i,j_{i+\sigma},j_{i+\tau},\ldots\)가 영과 다르면 (5)에 따라
\[
 \frakg_0\ne\frakg_i\ne\frakg_{i+\sigma}\ne\frakg_{i+\tau}\ne\cdots,
\]
이지만
\[
 \frakg_0=\frakg_1=\cdots=\frakg_{i-1}=\frako,
 \qquad
 \frakg_i=\frakg_{i+1}=\cdots=\frakg_{i+\sigma-1},\ldots
\]
이다. 따라서 인자 \(R^{(i)},R^{(i+\sigma)},R^{(i+\tau)}\)와 오직 이 인자들만 단위원과 다르다.\footnote{정리 VIII에서 보조정리 II도 함께 고려하면, 기초아이디얼은 변환된 아이디얼들의 최소공배수로서 다시 변환된 아이디얼임이 따른다. 정리 VIII은 이 사실을 쓰지 않는 H.-N.의 정리들에만 근거하므로, 이로써 내적인 이유까지 드러내는 새 증명을 얻는다. H.-N.에서는 이 정리가 영점이론에서만 쓰이며, 그 이론은 여기서 새로 전개하였다.}

정리 VIII을 이용하면 정리 I을 다음과 같이 강화할 수 있다.

\textbf{정리 IX.} \emph{한 아이디얼이 으뜸일 필요충분조건은 그 초등인자형, 따라서 그 노름이 으뜸함수인 것이다.}

으뜸아이디얼이 이 조건을 만족함은 정리 I에서 보였다. 역으로 \(\frakm\)의 초등인자형이 \(Q^{(i)}=(P^{(i)})^\rho\)이고 \(P^{(i)}\)이 소함수라고 하자. 정리 VIII에 따라 \(\frakm\)의 모든 소아이디얼은 같은 차원 \(n-i\)을 갖는다. \(\frakm=[\frakq_1,\ldots,\frakq_\alpha]\)에서
\[
 Q^{(i)}\equiv0(\frakq_\lambda);\quad\hbox{따라서}\quad
 (P^{(i)})^\rho\equiv0(\frakp_\lambda)\quad\hbox{이고, 그러므로}\quad
 P^{(i)}\equiv0(\frakp_\lambda)
\]
를 얻는다. 이제 모든 \(\frakp_\lambda\)의 차원이 \(n-i\)이고 \(P^{(i)}\)은 소함수이므로, \(P^{(i)}\)은 각 \(\frakp_\lambda\)의 초등인자형과 같다. 따라서 정리 IV의 따름정리에 따라 모든 \(\frakp_\lambda\)가 일치한다. 극대 으뜸성분들을 다루고 있으므로 하나의 \(\frakp\)만 나타날 수 있고, \(\frakm\)은 으뜸이다. \(\rho=1\)일 때 소아이디얼이라는 특수한 경우도 배제되지 않는다. 정리 VIII과 IX에서 기초아이디얼의 성분들을 다음과 같이 배치할 수 있다.

\textbf{정리 X.} \emph{\(R^{(i)}\)의 극대 으뜸인자 \(Q_\lambda^{(i)}=(P_\lambda^{(i)})^{\rho_\lambda}\)에 대응하는 기초아이디얼 \(\frakg_i\)의 성분 \(\frakr_\lambda\)는, 차원이 \(n-1,\ldots,n-i+1\)인 대응 소아이디얼들과 초등인자형이 \(P_\lambda^{(i)}\)인 차원 \(n-i\)의 그 소아이디얼에 의해 유일하게 결정되는 분해의 고립성분이다. 대응 소아이디얼들과 노름의 극대 으뜸인자들은 일대일로 대응한다.}

§ 1, 7에 따라 \(\frakr_\lambda\)는 그 \(i\)단계 기초아이디얼과 일치하고, \(\frakg_{i-1}\)을 \((i-1)\)단계 기초아이디얼로 갖는다. 따라서 정리 VIII, 또는 (5)에 따라 다음 최단표현이 있다.
\[
 \frakr_\lambda=[\frakg_{i-1},\frakq_{\lambda,1},\ldots,\frakq_{\lambda,t_\lambda}],
\]
여기서 모든 \(\frakq\)의 차원은 \(n-i\)이고 서로 다른 소아이디얼들에 속한다. \(Q_\lambda^{(i)}\)의 정의와 \(\frakr_\lambda\)가 그 \(i\)단계 기초아이디얼과 일치한다는 사실에서
\[
 Q_\lambda^{(i)}\frakg_{i-1}\equiv0(\frakq_{\lambda\kappa}),
 \quad(\kappa=1,\ldots,t_\lambda)
\]
를 얻고, 따라서
\[
\begin{aligned}
 (Q_\lambda^{(i)})^{\sigma_\lambda}&\equiv0(\frakq_{\lambda\kappa});
 \quad\hbox{그러므로}\quad
 (P_\lambda^{(i)})^{\rho_\lambda\sigma_\lambda}\equiv0(\frakp_{\lambda\kappa})\\
 &\hbox{따라서}\quad P_\lambda^{(i)}\equiv0(\frakp_{\lambda\kappa})
\end{aligned}
\]
이다. 정리 IX의 증명과 마찬가지로 여기서 \(\frakr_\lambda\)의 표현에는 차원 \(n-i\)인 \(\frakq_\lambda\) 하나와 그 대응 소아이디얼 \(\frakp_\lambda\) 하나만 나타날 수 있고, \(\frakp_\lambda\)의 초등인자형은 \(P_\lambda^{(i)}\)이다.

이제 \(\frakp_\lambda\)가 \(\frakm\)의 대응 소아이디얼임을, 곧 \(\frakq_\lambda\)가 \(\frakm\)을 극대 으뜸성분들로 나타내는 표현 가운데 적어도 하나에 실제로 나타남을 보여야 한다. 그러면 어떤 소아이디얼도 같은 차원이나 더 낮은 차원의 소아이디얼에 포함될 수 없으므로 \(\frakr_\lambda\)가 분해의 고립성분, 더 정확히는 \(\frakp_\lambda\)와 더 높은 차원의 모든 대응 소아이디얼들에 의해 정해지는 성분임도 드러난다. 이를 위해서는 \(\frakq_\lambda\)가 \(\frakg_i\)의 한 최단표현에 나타남만 보이면 된다. 정리 VIII에 따라 그러한 표현은 언제나 \(\frakm\)의 한 최단표현으로 보충할 수 있기 때문이다. 따라서
\[
 \frakg_i=[\frakr_1,\ldots,\frakr_\alpha]
          =[\frakg_{i-1},\frakq_1,\ldots,\frakq_\alpha]
\]
가 최단임만 보이면 된다. 이를테면 \(\frakq_\lambda\)가 그 보충부분에 포함되어
\(
\frakg_{i-1}\frakq_1\cdots\frakq_{\lambda-1}
\frakq_{\lambda+1}\cdots\frakq_\alpha
\)
가 \(\frakq_\lambda\)로 나누어진다고 하자. \(\frakg_{i-1}\not\equiv0(\frakq_\lambda)\)이므로 어떤 \(\mu\ne\lambda\)에 대하여 \(\frakq_\mu\)의 한 거듭제곱이 \(\frakq_\lambda\)로 나누어진다. 대응 소아이디얼 \(\frakp_\mu\)와 \(\frakp_\lambda\)의 차원이 같으므로 정리 II에 따라 이 소아이디얼들이 일치해야 한다. 그러나 정의상 그 초등인자형 \(P_\mu^{(i)}\)과 \(P_\lambda^{(i)}\)은 서로 다르므로 이는 불가능하다. 정리 VIII에 따라 차원 \(n-i\)인 모든 대응 소아이디얼이 \(\frakg_i\)에 대응하며, 위 표현은 각각의 그러한 소아이디얼에 노름의 극대 으뜸인자가 실제로 하나씩 대응함도 보여 준다. 이로써 일대일 대응을 증명하였다.

정리 X에서 정리 IV를 곧바로 다음과 같이 옮길 수 있다.

\textbf{정리 XI.} \emph{노름의 각 극대 으뜸인자 \(Q_\lambda^{(i)}=(P_\lambda^{(i)})^{\rho_\lambda}\)는 다음 곱표현을 갖는다.}
\[
 Q_\lambda^{(i)}
 =\prod (x_i-t_{1i}\bar y_{1\nu}-\cdots-t_{ni}\bar y_{n\nu})^{\delta_\lambda\rho_\lambda}
 =\prod (t_{1i}(y_1-\bar y_{1\nu})+\cdots+t_{ni}(y_n-\bar y_{n\nu}))^{\delta_\lambda\rho_\lambda},
\]
\emph{여기서 \(\bar y_{1\nu},\ldots,\bar y_{n\nu}\)는 각각 원래 부정원 \(y\)에서 대응 소아이디얼 \(\frakp_\lambda\)의 서로 대응하고 \(t_{1i},\ldots,t_{ni}\)와 무관한 영점계이고, \(\delta_\lambda\)는 \(1\), 또는 \(p^{f_\lambda}\)이다. 이로써 노름의 일차인수분해가 완전히 주어진다. \(Q_\lambda^{(i)}\)의 차수는 고립성분 \(\frakr_\lambda\)에 따른 잉여류 가운데 \(\frakg_{i-1}\)의 잉여류들로 나타나는 부분환이 분수몫으로 유도된 잉여류 부분체 \((\mathsf{P}(u,x_{i+1},\ldots,x_n))\)에 관하여 갖는 차수를 준다. 으뜸아이디얼의 경우에는 모든 잉여류로 이루어진 환을 말한다. 새 부정원 \(v\)들을 첨가하면 다음 분해도 얻는다.}
\[
 Q_\lambda^{(i)}(u,v)
 =\prod (z-v_1\bar x_{1\nu}-\cdots-v_i\bar x_{i\nu})^{\delta_\lambda\rho_\lambda}.
\]

정리 X에 따라 \(P_\lambda^{(i)}\)이 \(\frakp_\lambda\)의 초등인자형임을 알았으므로, 실제로는 정리 IV의 분해를 대입했을 뿐이다. 차수에 관한 주도 § 1, 5와 7에서 말한 것을 달리 표현한 것일 뿐이다.\footnote{마찬가지로 H.-N.의 주 2)에 언급된 중복도의 정식화도 정리 X에서 곧바로 따른다. 실제로 \(x_{i+1},\ldots,x_n\)을 첨가하면 잉여류계 \(\frakg_{i-1}\mid\frakr_\lambda\)는 \(\frakt_\lambda\mid\frakr_\lambda\)와 동형이다. 여기서 \(\frakt_\lambda=[\frakg_{i-1},\frakr_1,\ldots,\frakr_{\lambda-1},\frakr_{\lambda+1},\ldots,\frakr_\alpha]\)로 둔다. \(Q_\mu^{(i)}\)들이 서로 소임을 고려하면 선형형식의 가군으로 돌아가는 것(H.-N., 정리 V)에서 이 사실이 곧바로 드러난다. 또 \(\frakt_\lambda\mid\frakr_\lambda\)는 \(\frakt_\lambda\mid\frakq_\lambda\)와 동형이다. 이때 \(\frakt_\lambda\)는 \(\frakq_\lambda\)의 보충부분에 \(x_{i+1},\ldots,x_n\)을 첨가하여 생긴다.}

\subsection*{§ 6. 초등인자형과 노름에 의한 소아이디얼과 진으뜸아이디얼의 특징짓기}

정리 IX에서는 초등인자형과 노름으로 으뜸아이디얼을 특징지었지만, 소아이디얼을 으뜸아이디얼의 특수한 경우로 보는 데 그쳤다. 이제 소아이디얼과 진으뜸아이디얼을 분리해야 하며, 논의는 § 3과 나란히 진행된다. 필요조건이나 충분조건은 다음과 같이 쉽게 줄 수 있다.

\textbf{정리 XII.} \emph{소아이디얼의 필요조건은 초등인자형이 소함수인 것이고, 충분조건은 노름이 소함수인 것이다. 달리 말하면 소아이디얼의 초등인자형은 언제나 소함수이고, 진으뜸아이디얼의 노름은 언제나 진으뜸함수이다.}

소아이디얼의 초등인자형이 소함수임은 정리 I에서 이미 보였다. 이제 노름 \(R_\frakq\)가 소함수와 같다고 하자. 이 가정 아래에서 \(\frakq\)의 모든 진인자 \(\fraka\)가 더 낮은 최고차원을 가짐을 보이면, 정리 II에 따라 \(\frakq\)가 소아이디얼임이 드러난다. 실제로 § 1, 6에 따르면 노름 \(R_\fraka\)의 인자들 가운데 \(R_\frakq\)의 대응 인자와 처음으로 다른 것은 진인자이다. \(R_\frakq=P^{(i)}\)이므로 \(R_\fraka\)의 인자 \(R^{(i)}\)은 \(P^{(i)}\)의 진인자와 같아야 하고, 따라서 단위원이다. 그러므로 \(\fraka\)의 최고차원은 더 낮다.

둘째의 간단한 증명은 다음 정리에도 바탕이 되는 보조정리를 이용한다.

\textbf{보조정리 VII.} \emph{으뜸아이디얼의 초등인자형 \(E^{(i)}\), 더 일반적으로 차원 \(n-i\)의 대응 소아이디얼들만 갖는 한 아이디얼의 초등인자형에 따른 잉여류 가운데 다항식 \(H^{(i)}\)들로 나타나는 잉여류계는, 그 아이디얼에 따른 잉여류의 한 부분계와 동형이다. 초등인자형의 차수는 분수체 \((\mathsf{P}(u,x_{i+1},\ldots,x_n))\)에 관한 이 잉여류계의 차수를 준다.}

실제로 같은 다항식 \(H^{(i)}\)으로 나타나는 류들을 서로 대응시키면 대응을 만들 수 있으며, 합에는 합이, 곱에는 곱이 대응한다. 이 대응은 일대일이기도 하다. 그 아이디얼에 따른 영류에 속하는 모든 다항식 \(H^{(i)}\)은 정의상 \(E^{(i)}\)으로 나누어지기 때문이다. 따라서 영류에는 \(E^{(i)}\)에 따른 영류가 대응하고, 그 역도 물론 성립한다.

이제 한 아이디얼의 노름이 소함수 \(P^{(i)}\)이고, 따라서 초등인자형과 같다고 하자. 그러면 초등인자형 \(P^{(i)}\)과 \(\frakq\)에 따른 잉여류계들의 \((\mathsf{P}(u,x_{i+1},\ldots,x_n))\)에 관한 차수는 일치한다. 따라서 \(P^{(i)}\)에 따른 잉여류계와 동형인 부분계가 \(\frakq\)에 따른 잉여류들을 모두 소진한다. 이 계에는 영인자가 없으므로 \(\frakq\)는 소아이디얼이다.

일반적으로 정리 XII는 뒤의 예들이 보여 주듯 필요충분조건으로 강화할 수 없다. 그러나 계수영역 \(\mathsf{P}\)가 완전체(§ 1, 14)이면 다음 정리가 성립한다.

\textbf{정리 XIII.} \emph{계수영역 \(\mathsf{P}\)가 완전체이면 소아이디얼의 노름과 초등인자형은 소함수이고, 따라서 서로 같다. 진으뜸아이디얼의 노름과 초등인자형은 진으뜸함수이다. 그러므로 완전체에서는 초등인자형이 소함수라는 조건이 소아이디얼의 필요충분조건이다. 이 조건에서 초등인자형 대신 노름을 써도 된다.}

정리 XII를 고려하면, 완전체에서 초등인자형이 소함수인 순간 그 아이디얼이 소아이디얼이고 그 노름도 이 소함수와 같음을 보이면 정리 XIII가 증명된다. 보조정리 V, § 3에서 완전체의 경우 초등인자형이 소함수이기만 하면 제1종 소함수임을 보였다. 따라서 \((\mR)\)을 \((\mathsf{P}(u,x_{i+1},\ldots,x_n))\)에 \((x_i)\)를 첨가하여 얻은 \(P^{(i)}\)에 따른 잉여류체 \((\mathsf{P}(u,x_i,x_{i+1},\ldots,x_n))\)라고 하자. 그러면 \((x_i)\)는 제1종 소함수 \((P^{(i)})\)의 영점이다. 따라서 이를테면 라그랑주 공식이 보여 주듯 \((y_1),\ldots,(y_n)\)은 \((\mR)\)에 들어 있다.

보조정리 VII에 따라 \((\mR)\)과 동형인, 고찰하는 아이디얼 \(\frakq\)에 따른 잉여류의 부분계 --- 여기서는 \((\mathsf{P}(u,x_{i+1},\ldots,x_n))\)의 원소들이 분모에 나타날 뿐이다 --- 는 잉여류 \((y_1),\ldots,(y_n)\)을 포함하고, 따라서 \(\frakq\)에 따른 잉여류계를 모두 소진한다. 이 부분계는 \((\mR)\)과 동형이어서 영인자를 가질 수 없으므로 \(\frakq\)는 소아이디얼이다. \((\mathsf{P}(u,x_{i+1},\ldots,x_n))\)에 관한 이 \(\frakq\)의 잉여류체 \((\mR)\)의 차수는 \(\frakq\)의 노름 차수와 같다. 한편 \((\mR)\)과의 동형에 따라 초등인자형 \(P^{(i)}\)의 차수와도 같으므로, 노름과 초등인자형은 일치해야 한다.

비완전체의 경우에도 정리 XII를 다음과 같이 강화할 수 있다.

\textbf{정리 XIV.} \emph{계수영역 \(\mathsf{P}\)가 표수 \(p\)의 비완전체이면, 한 소아이디얼의 노름은 그 초등인자형의 \(p^g\)제곱이다. 여기서 \(0\leqq g\leqq(i-1)f\)이고, \(f\)는 초등인자형의 지수, \(n-i\)는 소아이디얼의 차원이다.}

정리 XIV는 \(\frakp\)의 잉여류체 \((\mR)\)이 그와 동형인 부분체
\(
(\mR')=(P(u,x_i,x_{i+1},\ldots,x_n))
\)
에 관하여 \(p^g\)차 (\(g\geqq0\))라는 주장으로 귀착된다. 이는 § 1, 15에 제시한 Steinitz의 정리들, 더 정확히는 다음 보조정리에서 곧바로 따른다.

\textbf{보조정리 VIII.} \emph{\(\mathfrak L\)이 비완전체 \(\Omega\)의 축약차수 \(r\), 지수 \(f\)인 유한확대이고, \(\mathfrak L_0\)가 \(\mathfrak L\) 안의 제1종 부분체이며, \(\mathfrak M\)이 \(\mathfrak L_0\)와 \(\mathfrak L\) 사이의 중간체이면, \(\mathfrak L\)의 각 원소가 \(\mathfrak M\)에 관하여 갖는 차수는 \(p\)의 거듭제곱이다. 여기서 \(p\)는 \(\Omega\)의 표수이다.}

실제로 \(\mathfrak L\)의 각 원소 \(z\)의 \(p^{f'}\)제곱 (\(f'\leqq f\))은 \(\mathfrak L_0\)에 속한다. 따라서 \(z\)는 \(\mathfrak L_0\)의 계수를 갖는 소함수 \(G(t)=(t-z)^{p^{f'}}\)의 영점이다. \(H(t)=(t-z)^\delta\)를 \(z\)를 영점으로 갖는 \(\mathfrak M\) 위의 소함수라 하자. 그러면 \(H(t)\)는 \(\mathfrak L_0(z)\)의 계수도 가지므로, 그 계수들은 \(\mathfrak M\)과 \(\mathfrak L_0(z)\)의 교차체, 곧 \(\mathfrak L_0\)와 \(\mathfrak L_0(z)\) 사이의 한 중간체에 속한다. 따라서 \(\mathfrak M\)에 관한 \(z\)의 차수는 이 중간체에 관한 차수와 같고, \(p^{f'}\)의 약수이므로 \(p\)의 거듭제곱이다.

정리 XIV를 증명하려면 다음만 주목하면 된다. \((\mR)\)의 축약차수가 \(r\), 지수가 \(f\)이면 체
\(
(\mR')=(P(u,x_i,x_{i+1},\ldots,x_n))
\)
는 \((\mathsf{P}(u,x_{i+1},\ldots,x_n))\)에 관하여 \(r\cdot p^f\)차이고, 따라서 \((\mathfrak{K})\) 안의 제1종 부분체 \((\mathfrak{K}_0)\)는 \((\mR')\)의 부분체이다. 실제로 \((\mR)\) 안에는 차수 \(r\cdot p^f\)인 원소들이 존재하므로(§ 1, 15), \((\mathfrak{K})\)의 모든 원소가 \((x_i)\)에서 \(t_{\mu\nu}\)들을 특수화하여 생기는 이상 \((x_i)\)도 반드시 이 차수를 갖는다. 따라서 \((\mathfrak{K})\)의 지수는 초등인자형의 지수와 일치하고, \((x_i)^{p^f}\)는 차수 \(r\)이며 \((\mathfrak{K}_0)\)의 원소이므로 \((\mathfrak{K}_0)\)의 원시원소이다. 이제 \((\mathfrak{K})\)은 유한개의 원소, 이를테면 \((x_1),\ldots,(x_{i-1})\)을 \((\mR')\)에 첨가하여 생기므로 보조정리 VIII을 유한번 적용하면 원하는 증명과 동시에 \(g\)의 상계를 얻는다.\footnote{이와 나란히 진행되는 A. Ostrowski, \emph{Zur arithmetischen Theorie der algebraischen Größen}, Gött. Nachr. 1919, 279--298쪽, 특히 대응 정리가 증명 없이 진술된 288쪽을 보라.}

끝으로 정리 XII를 일반적으로 정리 XIV보다 더 강화할 수 없음을 보여 주는 예들을 덧붙인다. 노름이 초등인자형의 제곱인 소아이디얼 (\(g>0\))과, 초등인자형이 소함수인 진으뜸아이디얼들을 제시할 것이다. 뒤의 예는 여기에는 정리 XIV의 유사물이 없고, 노름이 초등인자형의 임의의 거듭제곱이 될 수 있음도 함께 보여 준다.

계수영역 \(\mathsf{P}\)를 법 \(2\)의 잉여류체에 두 부정원 \(\lambda,\mu\)를 첨가하여 얻자. 이는 표수 2의 비완전체이다. 아이디얼
\[
 \bar{\frakp}=(y_1^2+\lambda,\,y_2^2+\mu)
\]
를 고찰하자. 그 잉여류계 \((\bar{\mR})\)이 체 \(\mathsf{P}(\sqrt\lambda,\sqrt\mu)\)와 동형이므로 이는 소아이디얼이다. \((\bar{\mR})\)은 \((\bar{\mathsf{P}})\)에 관하여 4차, 축약차수 \(r=1\), 지수 \(f=2\)이다. 따라서 초등인자형 \(P^{(2)}\)은 2차이고 노름은 4차여서 \(P^{(2)}\)의 제곱이다. § 1, 5의 모듈표현 (3)으로도 이를 쉽게 검산할 수 있으며, 다음을 얻는다.
\[
 P^{(2)}=(u_{11}u_{22}+u_{12}u_{21})^2x_2^2+(u_{11}^2\mu+u_{21}^2\lambda),
\]
또는 적당히 곱하면
\[
 x_2^2+(t_{12}^2\lambda+t_{22}^2\mu).
\]

이제 \(\mathsf{P}\)를 법 \(2\)의 잉여류체에 부정원 \(\lambda\) 하나를 첨가하여 얻고, 아이디얼
\[
 \bar{\frakq}=(y_1^2+\lambda,\,y_2^2+\lambda)
              =(y_1^2+\lambda,(y_1+y_2)^2)
\]
를 바탕으로 삼자. 위에 준 초등인자형에서 \(\mu=\lambda\)로 특수화하면
\[
 P^{(2)}=x_2^2+\lambda(t_{12}^2+t_{22}^2)
\]
를 얻는다. \(t_{12}=(0)\), \(t_{22}=(1)\)로 두면 소함수 \(x_2^2+\lambda\)가 되므로 이것은 소함수이다. 그러나 \(\bar{\frakq}\)는 \((y_1+y_2)^2\)을 포함하지만 \((y_1+y_2)\)는 포함하지 않으므로 진으뜸이다. \(\bar{\frakq}\)에 따라 \(\mathsf{P}\)에 관하여 선형독립인 잉여류가 네 개이므로, 노름은 \(P^{(2)}\)의 제곱, 곧 \(p\)제곱과 같아야 한다.

정리 XIV의 이 마지막 유사물이 언제나 성립하지는 않음을 다음 예가 보여 준다. \(P\)를 법 \(3\)의 잉여류체에 부정원 \(\lambda\)를 첨가하여 얻고,
\[
 \bar{\frakq}=(y_1^3+\lambda,(y_1-y_2)^2)
\]
를 바탕으로 삼자. 그러면 위와 같이 진으뜸아이디얼이다. \(y_2^3+\lambda\)도 \(\bar{\frakq}\)로 나누어짐을 고려하면 초등인자형은
\[
 P^{(2)}=x_2^3+\lambda(t_{12}+t_{22})^3
\]
이므로 소함수이다. 그러나 \(\bar{\frakq}\)에 따른 선형독립인 잉여류가 여섯 개이므로 노름은 \(P^{(2)}\)의 제곱이다. 따라서 여기서는 초등인자형의 \(p^g\)제곱이 나타나지 않는다.

첫 예는 비완전체를 계수영역으로 삼으면 대수적 수체에서와 똑같이 1보다 높은 차수의 소아이디얼들이 나타날 수 있고, 이때 \(p^g\)를 소아이디얼의 차수라고 불러야 함을 보여 준다. 뒤의 예들은 분기아이디얼의 유사물, 곧 한 소수가 한 소아이디얼의 더 높은 거듭제곱인 으뜸아이디얼로 나누어지는 현상으로 보아야 한다. 주 10)에서 보았듯 초등인자형은 대수적 수체에서 한 아이디얼로 나누어지는 가장 작은 유리정수에 대응하기 때문이다.

\subsection*{§ 7. 절대소아이디얼}

\(\mathsf{P}\)의 계수를 갖는 한 소아이디얼이 \(P\)를 확대하여 얻을 수 있는 대수적으로 닫힌 체 \(A\)에서도 소아이디얼로 남으면 이를 \emph{절대소아이디얼}이라 한다.\footnote{잘 알려졌듯 대수적으로 닫힌 체란 한 부정원의 모든 다항식이 일차인자들로 분해되는 체이다. 임의의 체에 속하는 대수적으로 닫힌 체의 존재와 본질적 유일성은 Steinitz를 보라.} 이에 대응하여 \(\mathsf{P}\)의 계수를 갖는 소함수 \(P\)가 \(\mathsf{A}\)에서도 소함수로 남으면 이를 \emph{절대소함수}, 곧 절대기약다항식이라 한다. 절대소아이디얼은 다음 정리로 특징지어진다.

\textbf{정리 XV.} \emph{한 절대소아이디얼의 필요충분조건은 \(u\)들에 관하여 다항식이고 원시적이라고 가정할 수 있는 그 초등인자형이 \(x\)와 \(u\)의 다항식으로서 절대소함수인 것이다. 초등인자형은 노름과 같으므로 이 조건을 노름으로 진술해도 된다.}

증명을 위해 먼저 계수영역을 대수적으로 확대해도 일반적으로 노름과 초등인자형은 그대로 유지됨을 주목한다. 개별 인자 \(R^{(i)}\), 또는 \(E^{(i)}\)은 다항식 의미의 최대공약인자로 정의되며, 이 성질은 계수영역의 대수적 확대 아래에서도 유지되기 때문이다. 따라서 \(\mathsf{P}\)에서 \(\mathsf{A}\)로 넘어갈 때도 \(P^{(i)}\)은 \(\frakp\)의 초등인자형으로 남는다. 이는 계수영역을 완전체로 바꾸는 것이다. 대수적으로 닫힌 체는 정의에서 곧바로 알 수 있듯 완전체이기 때문이다. 따라서 정리 XIII에 따라 \(\frakp\)가 절대소아이디얼일 필요충분조건은 \(P^{(i)}\)이 \(\mathsf{A}(u)\)에 관하여 소함수인 것이다. \(P^{(i)}\)을 \(u\)들에 관하여 다항식이고 원시적이라고 가정하면, 이는 \(x\)와 \(u\)의 다항식으로서 절대소함수라는 뜻이다. 동시에 \(\mathsf{A}\)를 계수영역으로 삼는 순간 정리 XIII에 따라 \(P^{(i)}\)은 \(\frakp\)의 노름과도 같다. 앞서 말한 바에 따라 \(\mathsf{P}\)를 바탕으로 삼아도 똑같이 성립한다. 이로써 정리 XV를 증명하였다.

이제 한 (유한) 대수적 수체 \(\mathfrak K\)의 계수를 갖는 절대소함수 \(P\)에 대해서는 다음 정리가 성립한다.\footnote{A. Ostrowski, 앞의 논문(주 21), 296쪽의 보조정리. 더 간단한 증명은 E. Noether, \emph{Ein algebraisches Kriterium für absolute Irreduzibilität}, Math. Ann. 85 (1922), 26--33쪽, 제7번에 있다.} \(\mathfrak K\)의 유한개 소아이디얼을 제외하면, \(\mathfrak K\)의 모든 소아이디얼을 법으로 해도 \(P\)는 소함수, 더 정확히는 절대소함수로 남는다. 이 정리를 절대소아이디얼로 옮기는 데에는 다음 정리가 쓰인다.

\textbf{정리 XVI.} \emph{계수영역으로 한 체 대신 한 (유한) 대수적 수체 \(\mathfrak K\)의 모든 대수적 정수들의 환 \(\mathfrak o^*\)을 삼으면, 한 다항식아이디얼의 노름과 초등인자형은 \(\mathfrak K\)의 유한개 소아이디얼을 제외한 모든 소아이디얼을 법으로 해도 노름과 초등인자형으로 그대로 유지된다.}

증명을 위해 § 1, 1의 다항식영역을 조금 바꾸어야 한다. \(\bar{\mathfrak{J}}\)는 \(y_1,\ldots,y_n\)의 모든 다항식 가운데 계수가 \(\mathfrak o^*\), 곧 \(\mathfrak K\)의 대수적 정수인 것들로 이루어진다고 하자. \(\mathfrak{J}\)의 계수영역은 \(\mathfrak o^*[u]\), 곧 계수가 \(\mathfrak o^*\)인 \(u\)들의 모든 다항식으로 이루어져 있다고 하자. \(\bar{\frakm}\)이 \(\bar{\mathfrak{J}}\)의 아이디얼이면 변환 (1)에 따라 \(\mathfrak{J}\)의 변환된 아이디얼 \(\frakm\)으로 간다. 노름을 만들 때 \(\mathfrak o^*[u]\)의 유한개 다항식만 분모에 나타남을 보이면 정리 XVI가 곧바로 따른다.

우선 § 1, 5의 모듈표현 (3)으로 개별노름을 만드는 것은 \(x_{i-1}\)에 관하여 정칙인 다항식
\(
C^{(i-1)}=U_{i-1}(u)x_{i-1}^{k}+\hbox{낮은 차수의 항들}
\)
을 법으로 \(x_{i-1}\)의 거듭제곱들을 환원하는 것으로 귀착됨을 주목한다.\footnote{H.-N., § 4를 보라.} 여기서 \(C^{(i-1)}\)의 모든 계수, 특히 \(U_{i-1}\)을 \(\mathfrak o^*[u]\)의 원소라고 가정할 수 있다. \(x_1,\ldots,x_{i-1}\)을 차례로 환원하므로, \(x\)에 관하여 다항식인 모듈표현은 분모에 나타나는 거듭제곱곱 \(U_1^{\lambda_1}\cdots U_{i-1}^{\lambda_{i-1}}\)만 제외하면 \(\mathfrak o^*[u]\)에서도 다항식이다. 그런데 \(u\)의 다항식으로 본 곱 \(U_1U_2\cdots U_n\)의 계수들은 \(\mathfrak K\)의 한 아이디얼 \(\mathfrak r^*\)을 정하고, 이는 \(\mathfrak K\)의 유한개 소아이디얼로만 나누어진다. 이 유한개와 다른 \(\frakp^*\)를 택하고 \(\frakp^*\)에 따른 잉여류체를 계수영역으로 삼으면, \(\mathfrak o^*\)의 각 수를 \(\frakp^*\)에 따른 그 잉여류로 바꾸기만 하여 원래 모듈표현이 유지된다.

또 노름과 초등인자형은 \(\mathfrak o^*[u]\)의 인자를 곱하는 차이까지만 정해지므로, \(\mathfrak o^*[u]\)에 관해서도 \(\frakm\)으로 나누어진다고 가정할 수 있다. 이 가정 아래 이를테면
\(
R^{(i)}=V_i(u)x_i^{l}+\hbox{낮은 차수의 항들}
\)
이라 하자. 그러면 \(C^{(i-1)}\)이 정하는 계수 \(\rho\)의 가군 \(\mathfrak M_{i-1}^*\)의 모든 \(\rho\)차 행렬식이 나누어지는 것을 다룰 때, 분모에는 \(U_1^{\lambda_1}\cdots U_{i-1}^{\lambda_{i-1}}\) 외에 \(V_i\)의 거듭제곱만 나타난다. 따라서 \(\mathfrak K\)에서 \(V_1V_2\cdots V_n\)이 정하는 아이디얼에 포함되는 유한개 소아이디얼과도 다른 \(\frakp^*\)를 택하면, \(\frakp^*\)에 따른 잉여류체를 계수영역으로 삼아도 \(R^{(i)}\)은 이 행렬식들의 공약인자로 남고 \(\mathfrak M_{i-1}^*\)의 계수도 유지된다. \(C^{(i)}\)이 \(\mathfrak M_{i-1}^*\)의 한 \(\rho\)차 행렬식이었기 때문이다.

끝으로 각 \(R^{(i)}\)이 최대공약인자로 남도록 \(\frakp^*\)를 더 골라야 한다. \(D^{(i)}\)가 \(\mathfrak M_{i-1}^*\)의 모든 \(\rho\)차 행렬식을 돌게 하자. 가정에 따라
\(
U_1^{\lambda_1}\cdots U_{i-1}^{\lambda_{i-1}}V_i^\sigma D^{(i)}
=R^{(i)}T^{(i)}
\)
이고, \(T^{(i)}\)들은 다항식 의미에서 \(x\)를 포함하는 공약인자를 갖지 않는다. 따라서 모든 \(T^{(i)}\)에서 유도된 아이디얼은
\(
G^{(i+1)}=W_i(u)x_{i+1}^{m}+\hbox{낮은 차수의 항들},
\qquad W_i\ne0
\)
인 다항식을 포함한다. \(G^{(i+1)}\)을 \(x_{i+1}\)에 관하여 정칙이라고 가정할 수 있음은 변환 (1)에 새 부정원들을 계수로 하는 \(x_{i+1},\ldots,x_n\)의 변환을 합성하면 알 수 있다. 그러므로 \(\mathfrak K\)에서 \(W_1W_2\cdots W_n\)이 정하는 아이디얼에 포함되는 유한개 소아이디얼과도 다른 \(\frakp^*\)를 택하면, \(\frakp^*\)에 따른 잉여류체를 계수영역으로 삼아도 \(R_\frakm\)은 실제로 노름으로 유지된다. 완전히 같은 방식으로 \(\frakp^*\)를 더 골라 \(E_\frakm\)도 초등인자형으로 유지되게 할 수 있다. 이로써 정리 XVI를 증명하였다.

정리 XV와 XVI, 그리고 절대소함수에 관한 위 정리에서 그 일반화인 다음 정리가 따른다.

\textbf{정리 XVII.} \emph{\(\frakp\)가 한 (유한) 대수적 수체 \(\mathfrak K\)의 대수적 정수들을 계수로 갖는 절대소아이디얼이면, \(\mathfrak K\)의 유한개 소아이디얼을 제외한 모든 소아이디얼을 법으로 해도 \(\frakp\)는 소아이디얼, 더 정확히는 절대소아이디얼로 남는다.}

증명을 위해 \(\frakp\)의 노름 \(R_\frakp\)를 정리 XVI에서처럼 \(\mathfrak o^*[u]\)에 관해서도 \(\frakp\)로 나누어지도록 택하자. 그러면
\(
R_\frakp=T(u)P^{(i)}(x)
\)
이고, \(P^{(i)}(x)\)는 \(u\)들에 관하여 다항식이고 원시적이라고 가정할 수 있다. 따라서 정리 XV에 따라 \(P^{(i)}\)은 \(\mathfrak K\)의 계수를 갖는 \(u\)와 \(x\)의 다항식으로 보았을 때 절대소함수이다. 절대소함수에 관한 정리에 따라, \(\mathfrak K\)의 유한개 소아이디얼을 제외하면 모든 소아이디얼을 법으로 해도 \(P^{(i)}\)은 이 성질을 유지한다. 따라서 이 유한개 소아이디얼과 다르고, 정리 XVI에 나오는 또 다른 유한개 소아이디얼과도 다른 \(\frakp^*\)를 택하자. 그러면 \(P^{(i)}\)의 \(\mathfrak o^*\)-계수들을 \(\frakp^*\)에 따른 대응 잉여류들로 바꾸었을 때, \(P^{(i)}\)은 \(\frakp^*\)에 따른 잉여류체를 계수영역 \(\mathsf{P}\)로 삼은 \(\frakp\)의 노름이고, 이 노름은 절대소함수이다. 따라서 \(\mathsf{P}\)를 계수영역으로 삼아도 정리 XV에 따라 \(\frakp\)는 절대소아이디얼로 남는다. 이로써 정리 XVII를 증명하였다.

\begin{center}
괴팅겐, 1923년 5월 8일.\\[1em]
(1923년 3월 9일 접수.)
\end{center}

\clearpage

\clearpage
\setcounter{footnote}{0}
\phantomsection
\addcontentsline{toc}{section}{25. 소거이론과 아이디얼 이론}
%
%
%
%

\section*{25. 소거이론과 아이디얼 이론}

\begin{center}
J. Ber. d. DMV 33 (1924), 116--120쪽
\end{center}

\setcounter{footnote}{0}

\begin{center}
1923년 9월 25일 마르부르크에서 한 강연.\srcfn{1)}{상세한 서술과 참고문헌은 Math. Ann. 90 및 곧 출판될 논문에 있다.}\\[0.5em]
괴팅겐의 \textsc{Emmy Noether}.
\end{center}

이 글에서는 소거이론을 일변수 다항식의 영점이론과 완전히 나란하게 세울 수 있음을 보이고자 한다. 이 구성은 \emph{Hentzelt}의 소거이론을 \emph{Steinitz}가 전개한 체론의 방법 및 내가 전개한 아이디얼 이론의 방법과 결합하는 데 바탕을 둔다.

먼저 일변수 다항식의 경우에 문제를 개괄하겠다. 계수영역으로는 고정된 체 \(P\)를 시종일관 택하며, 소함수 등의 개념은 이 체를 기준으로 한다. \(P\)는 Steinitz의 의미에서 추상적으로 정의된 체여도 되고, 물론 특수한 경우에는 유리수체 또는 예전 소거이론에서 일반적으로 가정하던 모든 복소수의 체와 같아도 된다. 일변수 다항식의 영점이론은 다음 네 단계로 나뉜다.

1. \emph{분해정리.} 이 정리는 문제를 소함수, 곧 \(P\)에서 기약인 다항식에 관한 문제로 환원한다. 실제로
\[
a(x)=p_1(x)^{\varrho_1}\cdots p_r(x)^{\varrho_r}=q_1(x)\cdots q_r(x)
\]
에서 \(p(x)\)들이 서로 다른 소함수이고 \(q(x)\)들이 따라서 ``극대 으뜸함수''라면, \(a(x)\)는 소함수 \(p(x)\)들의 모든 영점에서, 그리고 오직 그 영점들에서만 0이 된다.

2. \emph{소함수 \(p(x)\)의 영점체 구성.} \(\alpha\)를 \(p(x)\)의 한 영점이라고 할 때, Steinitz에 따라 그러한 영점체 \(\mathfrak K=P(\alpha)\)를 \(p(x)\)의 잉여류체 \((\mathfrak K)\)와 동형인 \(P\)의 확대체로 언제나 구성할 수 있다. 거꾸로 \(P\)의 어떤 확대체 안에 있는 영점체도 잉여류체 \((\mathfrak K)\)와 동형이다. 따라서 영점체와 잉여류체는 추상적으로 동일하다.

3. \emph{\(p(x)\)가 정하는 갈루아체에서 \(p(x)\)를 일차인수들로 분해하기.} 2의 구성을 반복하면 본질적으로 유일하게 정해지는 갈루아체 \(P(\alpha_1,\ldots,\alpha_m)\)에 이른다. 이 체는 \(p(x)\)가 일차인수들로 분해되기에 꼭 충분한 체이다.

4. \emph{원래 다항식 \(a(x)\)에서 중복도의 의미.} 1, 2, 3에 의해 임의의 다항식 \(a(x)\)의 영점 문제와 그에 대응하는 일차인수분해가 해결된다. 이때 각 으뜸함수 \(q(x)\)의 차수, 곧 \(q(x)\)에 대한 잉여류들 가운데 \(P\) 위에서 선형독립인 것들의 수를 중복도로 볼 수 있다. 계수체 \(P\)가 대수적으로 닫혀 있으면 \(p(x)\)들이 일차식이 되므로, 이 수는 바로 영점의 중복도와 같아진다.

이제 일반적인 소거이론으로 넘어가자. 즉 \(P\)의 계수를 갖는 \(x_1,\ldots,x_n\)의 모든 다항식환 \(\mathfrak{J}\)를 바탕으로 삼는다. 그러면 소거이론은 이 다항식환의 \emph{한 아이디얼의 영점에 관한 이론}으로 정의할 수 있다. 여기서 아이디얼 \(\mathfrak a\)란 통상적인 뜻대로, \(f,g\)와 함께 \(f-g\)도 포함하고 \(f\)와 함께 \(Af\)도 포함하는 다항식들의 계이다. 여기서 \(A\)는 \(\mathfrak{J}\)의 임의의 다항식이다. \(\mathfrak a\)의 영점이란 \(P\)의 어떤 확대체에 속하는 값의 계 \(\alpha_1,\ldots,\alpha_n\)로서, 아이디얼에 속하는 모든 다항식 \(f\)에 대하여 \(f(\alpha)=0\)인 것을 뜻한다. \(f_1,\ldots,f_t\)가 언제나 존재하는 \(\mathfrak a\)의 아이디얼 기저라면, 즉 \(\mathfrak a\)의 모든 \(f\)에 대하여 \(f=A_1f_1+\cdots+A_tf_t\)이면, \(\mathfrak a\)의 영점은 분명 \(f_1,\ldots,f_t\)의 공통영점과 같다. 거꾸로 유한개의 임의의 다항식계는 그것이 생성하는 아이디얼의 기저로 볼 수 있다. 따라서 위에서 준 소거이론의 영점 정의는 통상적인 정의와 같지만, 임의로 한 기저를 특별히 택하는 데서 생기는 어려움을 피한다. 일변수 다항식에서는 주아이디얼만을 다루므로 이런 어려움이 없다. 기저다항식이 단위원, 곧 \(P\)의 한 원소를 인자로 곱하는 것까지 유일하게 정해지기 때문이다. 그러나 이 경우에도 아이디얼의 관점은 더 정확한 통찰을 준다.

소거이론도 일변수의 경우와 마찬가지로 네 단계로 나뉜다.

I. \emph{분해정리}, 곧 아이디얼을 극대 으뜸성분들로 나타내는 정리. 이 정리는 문제를 소아이디얼에 관한 문제로 환원한다. 곱이 한 아이디얼에 들어갈 때 언제나 적어도 한 인자가 그 아이디얼에 들어가면 그 아이디얼을 소아이디얼이라 한다. 곱이 한 아이디얼에 들어갈 때 적어도 한 인자 또는 다른 인자의 어떤 거듭제곱이 그 아이디얼에 들어가면 으뜸아이디얼이라 한다. 각 으뜸아이디얼 \(\mathfrak q\)에는 유일한 대응 소아이디얼 \(\mathfrak p\)가 있다. \(\mathfrak p\)는 \(\mathfrak q\)의 약수이고, \(\mathfrak p\)의 어떤 거듭제곱은 \(\mathfrak q\)로 나누어진다. 최소공배수로서 다음 표현이 성립한다.
\[
\mathfrak a=[\mathfrak q_1,\cdots,\mathfrak q_r];\qquad
\mathfrak p_1^{\sigma_1}\cdots \mathfrak p_r^{\sigma_r}\equiv 0\,(\mathfrak a).
\]
여기서 \(\mathfrak q\)들은 극대 으뜸성분이다. 즉 각 \(\mathfrak q\)는 으뜸이지만 임의의 두 \(\mathfrak q\)의 최소공배수는 으뜸이 아니다. 또한 일반성을 잃지 않고 어느 \(\mathfrak q\)도 나머지 것들의 최소공배수를 나누지 않는다고, 즉 어느 \(\mathfrak q\)도 단순히 생략할 수 없다고 가정해도 된다. \(\mathfrak p\)들은 \(\mathfrak q\)들의 대응 소아이디얼이다. \(\mathfrak q\)들 자체는 \(\mathfrak a\)에 의해 유일하게 정해지지 않지만, 일변수 경우와의 유사성이 바로 여기에 있는데, 대응 소아이디얼 \(\mathfrak p\)들은 \(\mathfrak a\)에 의해 유일하게 정해진다. 각 \(\mathfrak p\)는 \(\mathfrak a\)의 약수이고 거꾸로 \(\mathfrak p\)들의 어떤 거듭제곱곱은 \(\mathfrak a\)로 나누어지므로, 한 아이디얼의 영점들은 그 대응 소아이디얼들의 영점들과 일치한다.

II. \emph{소아이디얼 \(\mathfrak p\)의 영점체 구성.} 정의에 따라 소아이디얼에 대한 잉여류들은 영인자가 없는 환을 이룬다. \(\mathfrak p\)가 단위아이디얼과 다르면 이 환은 영원소와 다른 원소를 적어도 하나 포함하므로, 분수를 취하여 \(\mathfrak p\)의 잉여류체 \((\mathfrak K)\)로 확대할 수 있다. \((\mathfrak K)\)와 동형인 \(P\)의 확대체 \(\mathfrak K\)를 언제나 구성할 수 있으며, 이것은 \(\mathfrak p\)의 영점체가 된다. 즉 \(\alpha_1,\ldots,\alpha_n\)이 \(\mathfrak p\)의 한 영점일 때 \(\mathfrak K=P(\alpha_1,\ldots,\alpha_n)\)이다. 이때 \(\mathfrak K\)은 \(P\) 위에서 초월차수 \(k\), \(0\leq k<n\)을 갖는다. 곧 \(P\) 위에서 대수적으로 독립인 원소 \(k\)개를 포함하지만, 임의의 \(k+1\)개 원소는 \(P\) 위에서 대수적으로 종속이다. 거꾸로 \(P\)의 어떤 확대체 안에 있는 초월차수 \(k\)의 영점체도 잉여류체 \((\mathfrak K)\)와 동형이다. 따라서 초월차수 \(k\)의 영점체와 잉여류체는 추상적으로 동일하다.

III. \emph{\(\mathfrak p\)가 정하는 갈루아체에서 \(\mathfrak p\)의 노름을 일차인수들로 분해하기.} 원래 변수 \(y_1,\ldots,y_n\)에서 부정원을 계수로 갖는 선형변환으로 변수 \(x_1,\ldots,x_n\)을 얻었다고 하자. 그러면 초월차수 \(k\)의 영점체 \(\mathfrak K\)을 \(P(x_{i+1},\ldots,x_n)\)의 유한 대수적 확대체로 볼 수 있다. 여기서 \(k=n-i\)이다. 따라서 \(P(x_{i+1},\ldots,x_n)\)에 대한 갈루아체 \(\overline{\mathfrak K}\)로 넘어갈 수 있으며, 이 체에서는 영점 \(\alpha_i\)를 갖는 소함수 \(P^{(i)}(x_i,x_{i+1},\ldots,x_n)\)가 일차인수들로 분해된다. 원래 변수 \(y\)로 돌아가면 이 소함수 \(P^{(i)}\)가 기본방정식의 역할을 한다. 실제로 \(\alpha_i\)는 기본형식 \(u_1\beta_1+\cdots+u_n\beta_n\)와 같아지며, 여기서 \(\beta\)는 매개변수 \(x_{i+1},\ldots,x_n\)에 대수적으로 종속인 \(y\)의 한 영점계이다. 따라서 일차인수분해는 모든 영점에 걸친 곱을 준다. 한편 \(\mathfrak p\)를 \(x_1,\ldots,x_{i-1}\)의 거듭제곱 단항곱들에 관한 선형형식들의 모듈로 보고 \(P\) 대신 \(P(x_{i+1},\ldots,x_n)\)를 계수영역으로 택한 뒤, 유한개만 단위원과 다른 모든 초등인자들의 곱으로 노름을 정의하면 \(P^{(i)}\) 또는 그 어떤 거듭제곱을 \(\mathfrak p\)의 노름 \(N(\mathfrak p)\)으로 볼 수도 있다. 계수체 \(P\)가 완전체이면 언제나 \(N(\mathfrak p)=P^{(i)}\)이고, 비완전체이면 노름으로 \(P^{(i)}\)의 어떤 거듭제곱이 나타날 수 있다.

IV. \emph{원래 아이디얼의 노름과 중복도의 의미.} 이에 대응하여 아이디얼 \(\mathfrak a\)를 차례로 \(x_1,\ldots,x_{i-1}\)의 거듭제곱 단항곱들에 관한 선형형식들의 모듈로 보자. 여기서 \((i=1,2,\ldots n)\)이다. 그러면 \(P(x_{i+1},\ldots,x_n)\)를 계수영역으로 삼을 때 이 모듈에는 단위원과 다른 초등인자가 유한개만 있음을 보일 수 있다. 이 유한성 조건은 아이디얼 기저의 존재에서 따라온다. 각 단계의 이 유한개 초등인자들의 곱이 개별노름이고, \(\mathfrak a\)의 노름 \(N(\mathfrak a)\)은 모든 개별노름의 곱으로 정의된다. \(N(\mathfrak a)\)의 극대 으뜸인수 \(Q^{(i)}(x_i,x_{i+1},\ldots,x_n)\)들은 \(\mathfrak a\)의 대응 소아이디얼 \(\mathfrak p\)들과 일대일로 대응한다. 이때 \(Q^{(i)}\)는 \(P^{(i)}\)의 어떤 거듭제곱이며, \(P^{(i)}\), 또는 경우에 따라 그 어떤 거듭제곱은 \(\mathfrak p\)의 노름 \(N(\mathfrak p)\)이다. 따라서
\[
N(\mathfrak a)=N(\mathfrak p_1)^{\varrho_1}\cdots N(\mathfrak p_r)^{\varrho_r}=Q_1\cdots Q_r
\]
라는 분해를 얻는다. 필요한 경우 \(N(\mathfrak p)\)는 최고 초등인자 \(E(\mathfrak p)=P^{(i)}\)로 바꾸어야 한다. 이로써 I, II, III에 따라 \(\mathfrak a\)의 모든 영점에 대하여 노름의 일차인수분해가 이루어진다. \(\mathfrak a\)의 각 소아이디얼 \(\mathfrak p\)에 대응하는 영점들의 중복도는 \(Q^{(i)}\)의 차수로 볼 수 있다. 이 차수는 \(P(x_{i+1},\ldots,x_n)\) 위에서 선형독립인 잉여류들의 수로 해석된다. 더 정확히는 \(\mathfrak p\)와 더 높은 차원의 소아이디얼들이 유일하게 정하는 분해의 고립성분들의 잉여류, 곧 \((i-1)\)단계 기초아이디얼을 \(\mathfrak p\)가 정하는 \(i\)단계 기초아이디얼의 성분으로 나눈 잉여류들이다. 뒤의 성분에는 \(\mathfrak p\)와 모든 더 높은 차원의 소아이디얼이 대응 소아이디얼로 붙지만, 앞의 기초아이디얼에는 더 높은 차원의 소아이디얼들만 붙는다.

일변수의 경우도 주아이디얼로 넘어가면 이 틀에 \emph{완전히 들어온다.} 1의 분해는 \(a(x)\)를 기저다항식으로 보느냐 노름으로 보느냐에 따라 두 분해로 갈라진다. 하나는 I에 따라 주아이디얼을 극대 으뜸성분들의 최소공배수로 나타내는 것이며, 여기서는 소아이디얼들의 거듭제곱곱으로 나타내는 것과 같다. 다른 하나는 IV에 따라 노름을 각 소아이디얼의 노름들의 거듭제곱곱으로 분해하는 것이다. 2에서 영점체를 구성한다는 것은 실제로 \(p(x)\)가 정의하는 소아이디얼의 영점체를 구성한다는 뜻이고, 3에서는 노름을 일차인수들로 분해한다. 4의 중복도 정의는 IV의 정의에 포함된다. 일변수에서는 0단계 기초아이디얼이 단위아이디얼과 같고, 1단계 기초아이디얼이 이미 아이디얼 자체와 같아지기 때문이다.

아이디얼의 노름이 종결식의 의미를 갖는다는 사실로부터 이 이론은 \emph{통상적인 소거이론 안에 자리 잡는다.} 이 관점에서 소거이론의 옛 문제는 다항식환의 아이디얼 이론의 한 부분, 곧 아이디얼의 분해정리뿐 아니라 그 노름의 분해와 두 분해의 관계를 다루는 부분이 된다.

\begin{center}
(1923년 10월 19일 접수.)
\end{center}

\clearpage

\clearpage
\setcounter{footnote}{0}
\phantomsection
\addcontentsline{toc}{section}{26. 대수적 수체에서의 아이디얼 이론의 추상적 구성}
%
%
%

%
%
\section*{26. 대수적 수체에서의 아이디얼 이론의 추상적 구성}

\begin{center}
J. Ber. d. DMV 33 (1924), S. 102
\end{center}

4. E. Noether, 괴팅겐: 대수적 수체에서의 아이디얼 이론의 추상적 구성.

대수적 수체에서의 아이디얼 이론과 완전히 동치이고, 따라서 그러한 아이디얼 이론을 갖는 모든 환을 특징짓는 추상적 성질들이 제시되었다. 그것들은 단위원을 갖고 영인자가 없는 환들로서 이중 사슬 조건(Doppelkettensatz)이 성립하는 것들이다. 즉, 아이디얼들의 모든 약수 사슬(Teilerkette)은 유한 단계에서 멈추고, 영 아이디얼이 아닌 각 아이디얼을 법으로 한 모든 배수 사슬(Vielfachenkette)도 마찬가지로 유한 단계에서 멈춘다. 또한 이 환들은 그 분수체 안에서 정수적으로 닫혀 있다. 마지막 가정을 버리면 유한 오더(order)에서의 아이디얼 이론에 관한 정리들이 얻어지며, 그 일부는 이미 Dedekind(제3판)에 증명 없이 실려 있다.

상세한 논의는 \emph{Mathematische Annalen}에 게재될 예정이다.

\clearpage
\setcounter{footnote}{0}
\phantomsection
\addcontentsline{toc}{section}{27. 아이디얼 이론의 힐베르트 수}
%
%
%

%
%
\section*{27. 아이디얼 이론의 힐베르트 수}

\begin{center}
J. Ber. d. DMV 34 (1925), S. 101
\end{center}

--- \textbf{E. Noether}: 아이디얼 이론의 힐베르트 수. 일반적으로 힐베르트 수(Hilbertsche Anzahlen)란 아이디얼의 잉여류를 세는 데 관계되는 수를 말한다. 힐베르트 자신은 자신의 ``특성 함수''(charakteristische Funktion)에서, 다항식환의 아이디얼에 대해 어떤 경계값 이상에서는 선형 독립인 잉여류의 정확한 개수를 주는 함수를 제시했다. 더 낮은 차수에 대해서는 매콜리가 이 함수를 생성 함수로 보완했고, 오스트롭스키가 그 생성 함수를 명시적으로 계산할 수 있었다. 그러나 매콜리는 또 다른 종류의 수들도 제시했다. 이 수들은 추상적으로 정식화할 수 있기 때문에 일반 아이디얼 이론에 가장 본질적인 것으로 드러났다. 일반 분해 정리에 따라 으뜸 아이디얼 $\mathfrak q$에 대해서만 정의하면 충분하다. 여기서 말하는 수는 대응하는 소 아이디얼 $\mathfrak p$의 잉여류체(residue field; Restklassenkörper)에 관하여 선형 독립인 잉여류의 개수이다. 이 수들은 $\mathfrak q,\mathfrak q/\mathfrak p,\mathfrak q/\mathfrak p^2,\ldots$의 기약 성분(irreduzible Komponenten) 수로 각각 주어지는 부분계들로 나뉜다. $i$번째 부분계의 수는 $\mathfrak q/\mathfrak p^i$에서 $\mathfrak q/\mathfrak p^{i-1}$까지 이어지는 합성열을 이루는 항의 수로도 특징지어진다. 전체 합성열과 그 힐베르트 수들은 일반적인 경우---예컨대 대수적 수체의 유한 오더(order)에서---더 이상 성립하지 않는, 으뜸 아이디얼을 소 아이디얼의 거듭제곱으로 나타내는 표현에 상응하는 대체물을 이룬다.

\clearpage
\setcounter{footnote}{0}
\phantomsection
\addcontentsline{toc}{section}{28. 군 지표와 아이디얼 이론}
%
%
%

%
%
\section*{28. 군 지표와 아이디얼 이론}

\begin{center}
J. Ber. d. DMV 34 (1925), S. 144
\end{center}

\textbf{3. {\latinfont E. Noether}, 괴팅겐: 군 지표와 아이디얼 이론.}

Frobenius의 군 지표 이론, 즉 유한군의 표현론은 완전가약환인 군환의 아이디얼 이론으로 이해된다. 따라서 완전가약환의 덧셈 분해에 관한 동형 정리들을 일반적으로 전개한다. 분해 불가능한 단순 아이디얼들의 직합으로 나타내는 방법은 동형을 제외하면 유일하며, 동일한 양쪽 아이디얼에 속하는 분해 불가능한 한쪽 아이디얼들은 서로 동형이다. 그러므로 서로 동형인 아이디얼들을 하나의 동형류로 묶으면, 서로 다른 분해 불가능 아이디얼류의 수는 양쪽 분해 불가능 아이디얼의 수와 같다.

완전가약인 초복소량의 체계, 특히 군환으로 특수화하면, 분해 불가능 아이디얼류와 기약 표현류가 서로 일대일로 대응한다. 또한 양쪽 분해 불가능 아이디얼과 지표가 서로 대응하며, 한 양쪽 아이디얼의 덧셈적 분해 불가능 성분 수는 그 분해 불가능 아이디얼의 랭크(Rang)와 같다. 이로써 Frobenius의 이론은 이 틀 안에 자리 잡는다.

자세한 논의는 \emph{Math. Ann.}에 게재될 예정이다.

\clearpage
\setcounter{footnote}{0}
\phantomsection
\addcontentsline{toc}{section}{29. 표수 p인 유한 선형군 불변식의 유한성 정리}
%
%
%

%
%
\section*{29. 표수 $p$인 유한 선형군의 불변식에 관한 유한성 정리}

\begin{center}
Nachr. v. d. Ges. d. Wiss. zu Göttingen 1926, S. 28--35
\end{center}

\setcounter{footnote}{0}

\begin{center}
지음\\[0.5em]
\textbf{에미 뇌터(Emmy Noether)}, 괴팅겐.\\[1em]
1926년 7월 30일 회의에서 R. Courant가 제출.
\end{center}

이하에서 나는 선형 치환의 유한군에 대한 불변식의 잘 알려진 유한성 정리\footnote{초등적인 유한성 증명은 E. Noether, Der Endlichkeitssatz der Invarianten endlicher Gruppen. Math. Ann. 77 (1916), S. 89--92를 참조하라.}가 계수 영역으로 보통의 수체 대신 임의의 체, 따라서 특히 표수 $p$인 체\footnote{체론의 개념에 대해서는 E. Steinitz, Algebraische Theorie der Körper. J. f. M. 137 (1910), S. 167--309를 참조하라. 표수는 § 4, 근체(Wurzelkörper)는 § 12, 분수체는 § 3에 있다.}를 택하는 경우에도 성립함을 보인다. 그러나 군의 위수가 $p$로 나누어지는 순간 기존의 유한성 증명들은 작동하지 않는다. 더 깊은 산술적 사실, 곧 지금까지 표수 0에 대해서만 알려져 있던 다음 명제를 이용해야 한다.

\textbf{유한성 판정기준(Endlichkeitskriterium)}\footnote{E. Noether, Algebraische und Differentialinvarianten. Jahresber. d. d. Math.-Ver. 32 (1923), S. 177--184, Nr. 3, Satz 4 및 그곳에 제시된 문헌을 참조하라. 그 판정기준의 마지막 문장에는 ``$\mathfrak R$에 관한 가군 생성계''(정수적 원소들의 기본계)라고 잘못 적혀 있으며, ``$\mathfrak R$의 어떤 부분환에 관한 가군 생성계''가 옳다.}: 체의 계수를 갖는 다항식, 더 일반적으로는 유리함수로 이루어진 정역 $\mathfrak J$가 유한 생성되는 것은, $\mathfrak J$ 안에 유한 생성 부분환 $\mathfrak R$이 포함되어 있고 $\mathfrak J$가 $\mathfrak R$에 대해 정수적일 때 그리고 그때에 한한다. 그러면 $\mathfrak R$의 모든 정수성 기저(Integritätsbasis)에 $\mathfrak R$의 어떤 부분환에 관한 $\mathfrak J$의 가군 생성계(Modulbasis)를 덧붙이면 $\mathfrak J$의 정수성 기저가 된다.

그 자체로 독립적인 관심을 끄는 이 판정기준은 유리 기저(Rationalbasis)의 존재와, 유한 확대의 정수적 원소들로 넘어갈 때 약수사슬정리(Teilerkettensatz)가 보존된다는 사실에 기초한다. 이 두 번째 정리는 바로 앞의 Artin과 v. d. Waerden의 논문에서 임의의 표수에 대해 증명되었지만, 원래의 영역에 관한 어떤 유한성 가정 아래에서만 증명되었다. 그 가정은 유한군의 경우에 충족된다(근환(Wurzelring)이 유한 확대를 이룬다). 나는 임의의 체를 계수 영역으로 하는 경우에 유리 기저가 존재함을 증명한다. 이로써 위의 제한 아래에서 유한성 판정기준을 증명할 수 있다.

유한군의 불변식이 유한성 판정기준을 만족한다는 사실은 대칭함수의 원리에 기초한다. 그러나 표수 0에서는 갈루아 레졸벤트(Galois'sche Resolvente)의 계수들이 이미 정수성 기저를 주는 반면, 일반적인 경우에는 그 계수들이 유한성 판정기준에서 요구하는 유한 생성 부분환만을 생성한다. 이 유한 생성 부분환의 존재로부터 같은 논법에 따라 ``상대'' 불변식(Relativ-Invarianten)과 ``모듈러'' 불변식(Modular-Invarianten)의 유한성이 따라온다.

여기서 다루는 불변식에는 Dickson과 그 학파\footnote{Dickson, On Invariants and the Theory of Numbers. The Madison Colloquium 1913을 참조하라. 그 이후의 문헌은 그때부터 출간된 Transactions of the American Mathematical Society의 여러 권에 실려 있다.}가 다룬 ``법 $p$의 사영 불변식''이 특수한 경우로 포함된다. 이 경우 모듈러 불변식에 대해서는 유한성 정리가 알려져 있었지만, 일반적인 ``형식적'' 불변식에 대해서는 알려져 있지 않았다.

%
%
\subsection*{§ 1. 유한성 판정기준}

1. \textbf{유리 기저의 존재 정리. 제1형식:} 체 $P$의 계수를 갖는 $n$개의 부정원에 관한 유리함수들로 이루어진 임의의 계 $S$에는 \textbf{유한 유리 기저(Rationalbasis)}가 존재한다. 즉 그 계에서 유한 개의 함수를 택하여, 계의 모든 함수를 이 유한 개의 함수로 \textbf{$\mathsf{P}$의 계수를 갖는 유리식으로} 나타낼 수 있다.

\textbf{제2형식:} $\mathfrak K$를 $\mathsf{P}$와 $\mathsf{P}(x_1,\ldots,x_n)$ 사이의 중간체라 하자. 여기서 $x_1,\ldots,x_n$은 부정원이고, $\mathfrak K$의 초월 차수는 $t\leqq n$이다. 또 $y_1,\ldots,y_t$를 $\mathfrak K$ 안의 기약계(irreduzibles System)\footnote{Steinitz의 용어로, 대수적으로 독립인 함수들로 이루어진 계를 기약계라 한다. 적어도 $t$개의 원소로 이루어진 기약계를 하나 포함하지만 그보다 많은 원소로 이루어진 기약계는 포함하지 않는 계를 초월 차수 $t$의 계라 한다. 이하에서 쓰는 기약계에 관한 간단한 정리들은 Steinitz, § 22를 참조하라.}라 하자. 그러면 $\mathfrak K$는 $\mathsf{P}(y_1,\ldots,y_t)$의 \textbf{유한 대수적 확대}가 된다.

두 형식은 실제로 동치다. 모든 중간체에 대해 유리 기저의 존재를 증명하면 충분하다. 계 $S$가 생성하는 체는 그러한 중간체 $\mathfrak K$이고, $\mathfrak K$의 유리 기저는 곧바로 $S$의 유리 기저를 주기 때문이다. 실제로 $\mathfrak K$의 모든 원소는 $S$에서 뽑은 유한 개 함수의 유리식 결합이다. 한편 제2형식에 따르면, $\mathfrak K$ 안의 임의의 기약계 $y_1,\ldots,y_t$에 $\mathfrak K$를 $\mathsf{P}(y_1,\ldots,y_t)$의 유한 확대가 되게 하는 유한 개의 함수를 덧붙이면 그러한 유리 기저를 얻는다. 반대로 제1형식에 따라 $\mathfrak K$의 유리 기저가 주어졌다면, 초월 차수의 정의에 의해 각 기저 원소는 $\mathsf{P}(y_1,\ldots,y_t)$ 위에서 대수적이어야 한다. 따라서 $\mathfrak K$는 $\mathsf{P}(y_1,\ldots,y_t)$의 유한 대수적 확대가 된다.

이 정리는 \textbf{제2형식}으로 증명하겠다. 먼저 $\mathfrak K$의 초월 차수가 $\mathsf{P}(x_1,\ldots,x_n)$과 같은 $n$이라고 하자. $y_1,\ldots,y_n$을 $\mathfrak K$ 안의 기약계라 하면, 따라서 이것은 $\mathsf{P}(x_1,\ldots,x_n)$ 안의 기약계이기도 하다. $y_1,\ldots,y_n$의 초월 차수가 $n$이므로 각 $x_i$는 $\mathsf{P}(y_1,\ldots,y_n)$ 위에서 대수적이다. 그러므로 $\mathsf{P}(x_1,\ldots,x_n)$은 $\mathsf{P}(y_1,\ldots,y_n)$의 유한 대수적 확대이고, 그 중간체인 $\mathfrak K$도 그러하다.

이제 $t<n$이라 하자. $y_1,\ldots,y_t,x_{t+1},\ldots,x_n$이 $\mathsf{P}(x_1,\ldots,x_n)$ 안의 기약계를 이루도록 $x$들의 번호를 정한다. 그러면 $x_{t+1},\ldots,x_n$은 $\mathsf{P}(y_1,\ldots,y_t)$ 위에서 대수적으로 독립이다. $\mathfrak K$가 $\mathsf{P}(y_1,\ldots,y_t)$ 위에서 대수적이라는 사실과 대수적 종속성의 추이법칙에 의해, 이들은 $\mathfrak K$ 위에서도, 그리고 $\mathsf{P}(y_1,\ldots,y_t)$와 $\mathfrak K$ 사이의 모든 중간체 $\mathfrak M$ 위에서도 대수적으로 독립이다. 따라서
\[
\overline{\mathfrak K}=\mathfrak K(x_{t+1},\ldots,x_n),
\qquad
\overline{\mathfrak M}=\mathfrak M(x_{t+1},\ldots,x_n)
\]
은 각각 $\mathfrak K$와 $\mathfrak M$의 순수 초월 확대다. 곧 $\overline{\mathfrak K}$에서 $\mathfrak K$에 속하지 않는 모든 원소는 $\mathfrak K$ 위에서 초월적이고, 마찬가지로 $\overline{\mathfrak M}$에서 $\mathfrak M$에 속하지 않는 모든 원소는 $\mathfrak M$ 위에서 초월적이다. 또한 $\overline{\mathsf{P}}$를 순수 초월 확대 $\mathsf{P}(x_{t+1},\ldots,x_n)$로 두면, $y_1,\ldots,y_t$는 $\overline{\mathsf{P}}$ 위에서 대수적으로 독립이다.

이제 $\overline{\mathsf{P}}$를 계수 영역으로 택함으로써 증명을 위에서 처리한 경우로 환원할 수 있다. 실제로 $\mathsf{P}<\mathfrak K< \mathsf{P}(x_1,\ldots,x_t;\,x_{t+1},\ldots,x_n)$\footnote{$\mathsf{P}<\mathfrak K$는 $\mathsf{P}$가 $\mathfrak K$에 포함된다는 뜻, 곧 $P$가 $\mathfrak K$의 부분체라는 뜻이다.}로부터
\[
\overline P<\overline{\mathfrak K}<\overline P(x_1,\ldots,x_t)
\]
를 얻으며, $\overline{\mathfrak K}$는 $\overline{\mathsf{P}}$ 위에서 $\overline{\mathsf{P}}(x_1,\ldots,x_t)$과 같은 초월 차수 $t$를 갖는다. 따라서 $\overline{\mathfrak K}$는 $\overline{\mathsf{P}}$를 계수 영역으로 하는 \textbf{유한 유리 기저}, 이를테면 $y_1,\ldots,y_t;z_1,\ldots,z_s$를 갖는다. $\overline{\mathfrak K}$는 $\mathfrak K$와 $\overline{\mathsf{P}}$의 합성체(Vereinigungskörper)이므로 $z$들은 $\mathfrak K$에서 택할 수 있다. 이제
\[
\mathfrak K=\mathfrak M=P(y_1,\ldots,y_t;z_1,\ldots,z_s)
\]
임을 보이면 된다. $\mathfrak M$은 $\mathsf{P}(y_1,\ldots,y_t)$와 $\mathfrak K$ 사이의 중간체이므로 $\mathfrak K$의 모든 원소는 $\mathfrak M$ 위에서 대수적이다. 한편 $\overline{\mathfrak M}=\overline{\mathfrak K}$이므로 $\mathfrak K$는 $\overline{\mathfrak M}$에 포함된다. 따라서 $\mathfrak K$는 $\mathfrak M$과 $\overline{\mathfrak M}$ 사이의 중간체다. 그런데 $\overline{\mathfrak M}$에서 $\mathfrak M$에 속하지 않는 모든 원소는 $\mathfrak M$ 위에서 초월적이므로, 반드시 $\mathfrak K=\mathfrak M$이다.

\textbf{따름정리:} $\mathsf{P}<\mathfrak K<\mathfrak L<\mathsf{P}(x_1,\ldots,x_n)$이고 $\mathfrak L$이 $\mathfrak K$ 위에서 대수적이면, $\mathfrak L$은 \textbf{$\mathfrak K$의 유한 대수적 확대}다. 실제로 $\mathfrak L$의 유리 기저의 모든 원소는 $\mathfrak K$ 위에서 대수적이므로, $\mathfrak L$은 $\mathfrak K$의 유한 확대가 된다.

%
%
2. \textbf{유한성 판정기준의 증명.} 이 조건이 필요함은 명백하다. 실제로 $\mathfrak J$가 유한 생성 정역이면, $\mathfrak J$ 자체를 판정기준에 나타나는 유한 생성 부분환 $\mathfrak R$로 잡을 수 있다.

역으로 계수체의 표수가 $p$일 때에는 \textbf{근체(Wurzelkörper) $\mathsf{P}^{1/p}$가 계수체 $P$의 유한 확대라는 가정}\footnote{원논문 28쪽 각주 2 참조.} 아래에서 이 조건이 충분하다. 표수가 0일 때에는 아무 제한 가정도 필요하지 않다. 이를 증명하려면 $\mathfrak J$가 $\mathfrak R$의 어떤 부분환에 관한 \textbf{유한 가군 생성계(Modulbasis)}를 가짐을 보이면 된다.

$\mathfrak K$는 $\mathfrak R$의 분수체\footnotemark[\value{footnote}]를, $\mathfrak L$은 $\mathfrak S$의 분수체를 뜻한다고 하자. $\mathfrak R$과 $\mathfrak J$는 영인자가 없는 환이므로 이 두 분수체가 존재하며, $\mathsf{P}<\mathfrak K<\mathfrak L<\mathsf{P}(x_1,\ldots,x_n)$이다. 따라서 1의 따름정리에 의해 $\mathfrak L$은 $\mathfrak K$의 \textbf{유한 대수적 확대체}다. 또한 가정에 의해 $\mathfrak J$의 모든 원소는 $\mathfrak R$에 대해 정수적이므로, $\mathfrak J$는 $\mathfrak L$에서 $\mathfrak R$에 대해 정수적인 모든 원소로 이루어진 환 $\mathfrak S$의 부분환이다. 그러나 $\mathfrak R$이 $\mathfrak K$ 안에서 정수적으로 닫혀 있다고 가정하지 않았으므로, 유한 확대에서 약수사슬정리(Teilerkettensatz)가 보존된다는 사실을 직접 적용할 수는 없다. 그 대신 먼저, 역시 유한 생성인 $\mathfrak R$의 부분환 $\mathfrak T$로 넘어가야 하며, 이 $\mathfrak T$는 그러한 정수적 닫힘 조건을 만족해야 한다.

%
%
우선 계수체 $P$가 무한히 많은 원소를 가진다고 가정하자. 그러면 $\mathfrak R$, 따라서 $\mathfrak J$도 $\mathfrak T=\mathsf{P}[g_1(x),\ldots,g_t(x)]$에 대해 정수적이 되도록, $\mathfrak R$에서 기약계(irreduzibles System, 곧 $P$ 위에서 대수적으로 독립인 함수계) $g_1(x),\ldots,g_t(x)$를 골라낼 수 있다. 실제로 가정에 따라 존재하는 $\mathfrak R$의 정수성 기저(Integritätsbasis, 곧 $\mathfrak R$를 $P$-대수로 생성하는 유한 집합)를 $f_1(x),\ldots,f_k(x)$라 하자. 이 함수들이 기약계를 이루지 않는다면, 필요할 경우 지표를 다시 매긴 뒤 $f_k$가 만족하는 관계식 $F(f_k;f_1,\ldots,f_{k-1})=0$을 잡자. $f_1,\ldots,f_{k-1}$을
\[
 f_1+t_1f_k,\ldots,f_{k-1}+t_{k-1}f_k;
\]
로 바꾸면, $t_i$들을 부정원으로 택했을 때 $f_k$, 따라서 $f_1,\ldots,f_{k-1}$도 이 새로운 함수들에 대해 정수적이다. 또한 $\mathsf{P}$가 무한히 많은 원소를 가지므로, 이 정수적 종속 관계가 유지되도록 $t_i$들을 $\mathsf{P}$의 원소들로 특수화할 수 있다. 따라서 정수적 종속성의 추이법칙을 적용하면서 이 과정을 유한 번 반복하면 원하는 환 $\mathfrak T=\mathsf{P}[g_1(x),\ldots,g_t(x)]$\footnote{Hilbert, \emph{Über die vollen Invariantensysteme}, § 1, Math. Ann. 42 (1893), S. 313--373 참조.}를 얻는다. 그러면 $\mathfrak L$은 $\mathfrak T$의 분수체 $\mathsf{P}(g_1(x),\ldots,g_t(x))$의 유한 차수의 확대체가 되고, $\mathfrak L$에서 $\mathfrak R$에 대해 정수적인 모든 원소로 이루어진 환 $\mathfrak S$는 $\mathfrak L$에서 $\mathfrak T$에 대해 정수적인 모든 원소로 이루어진 환과 일치한다.

%
%
$\mathfrak T$에서는 모든 아이디얼의 계에 대해 약수사슬정리(Teilerkettensatz, 원문의 역사적 명칭이며 여기서는 아이디얼계에 적용된다)가 성립하고, $\mathfrak T$는 그 분수체 안에서 정수적으로 닫혀 있다. 이는 두 사실 모두 $\mathfrak T$가 $t$개의 부정원에 대한 다항식환과 동형이기 때문이다. $\mathsf{P}^{1/p}$가 $\mathsf{P}$의 유한 차수 확대체라는 가정으로부터 $\mathfrak T^{1/p}$가 $\mathfrak T$-가군으로 유한 생성임이 또한 따른다\footnote{서론에서 인용한 Artin-v. d. Waerden의 논문 참조.}. 따라서 $\mathfrak S$ 안의 모든 $\mathfrak T$-부분가군의 계에 대해 같은 약수사슬정리가 성립하며, 특히 $\mathfrak T$를 포함하는 $\mathfrak S$의 부분환인 $\mathfrak J$는 \textbf{유한 $\mathfrak T$-가군 생성계(Modulbasis, 곧 $\mathfrak S$를 $\mathfrak T$-가군으로 생성하는 유한 집합)}\footnotemark[\value{footnote}]를 가진다. 더 나아가 $\mathfrak L$이 $\mathsf{P}(g_1(x),\ldots,g_t(x))$의 제1종 확대체(Erweiterung erster Art, 곧 분리 확대체)가 되는 경우에는, 계수체 $\mathsf{P}$에 관한 어떠한 제한 가정도 없이 $\mathfrak S$의 유한 $\mathfrak T$-가군 생성계가 존재한다. 따라서 특히 $P$의 표수가 0이면 언제나 그러하다\footnote{예컨대 E. Noether, \emph{Abstrakter Aufbau der Idealtheorie in algebraischen Zahl- und Funktionenk\"orpern}, § 3, Math. Ann. 96 (1926), S. 26--61 참조.}.

$h_1(x),\ldots,h_s(x)$가 그러한 $\mathfrak S$의 유한 $\mathfrak T$-가군 생성계를 이룬다고 하자. 그러면 $f(x)$를 $\mathfrak J$의 임의의 원소라 할 때,
\[
 f(x)=A_1(g)h_1(x)+\cdots+A_s(g)h_s(x)
\]
위와 같이 나타낼 수 있다. 여기서 $A_i(g)$들은 $\mathfrak T$의 원소다. 다시 말해 이들은 $g_1(x),\ldots,g_t(x)$에 관한 다항식이다. 따라서 이 표현은 $g_1(x),\ldots,g_t(x),h_1(x),\ldots,h_s(x)$가 구하고자 한 $\mathfrak S$의 정수성 기저(Integrit\"atsbasis, 곧 $\mathfrak S$를 $P$-대수로 생성하는 유한 집합)를 이룸을 보여 준다.

%
%
체 $\mathsf{P}$가 유한 개의 원소만을 갖더라도, 부정원 $u$ -- 곧 $\mathfrak J$ 위에서 초월적인 원소 -- 를 첨가하여 얻은 체 $\mathsf{P}(u)$는 무한히 많은 원소를 갖는다. 따라서 환 $\mathfrak J(u)$는 $\mathsf{P}(u)$를 계수체로 하는 유한 정수성 기저(Integritätsbasis, 곧 $\mathfrak J(u)$를 $\mathsf{P}(u)$-대수로 생성하는 유한 집합)를 갖는다. 또한 $\mathfrak J(u)$는 $\mathfrak J$와 $\mathsf{P}(u)$가 함께 생성하는 환(Vereinigungsring)이므로, 이 기저는 $u$의 거듭제곱별로 분해함으로써 $\mathfrak J$에서 택할 수 있다. 이때 $\mathfrak S$에서 택한 기저원소 $g_i$들은 이제 더 이상 기약계(irreduzibles System, 곧 $P$ 위에서 대수적으로 독립인 함수계)를 이루지 않는다. 따라서 $\mathfrak T$의 모든 다항식 $f(x)$는 다음과 같이 나타낼 수 있다: $C(u)f(x)=G(u,g_i,h_i)$. 여기서 $C(u)$는 $\mathsf{P}[u]$의 $0$이 아닌 다항식이고, $G$는 $u,g_i,h_i$에 관한 다항식이다. 이 등식에서 $u$의 적당한 거듭제곱의 계수를 비교하면 $g_i(x)$와 $h_i(x)$가 $\mathfrak J$의 정수성 기저를 이룸을 알 수 있다. 이로써 \textbf{유한성 판정기준이 일반적으로 증명되었다}.

%
%
\subsection*{§ 2. 불변식의 유한성}

1. \textbf{표수 $p$인 유한 선형군}이라 함은, 표수 $p$인 체 $\mathsf{P}$의 계수를 갖는, 각각 행렬식이 $0$이 아닌 $h$개의 선형 치환 $A_1,\ldots,A_h$로 이루어진 군 $\mathfrak H$를 말한다. 여기서 $A_k$는 다음 선형 치환
\[
 x_1^{(k)}=a_{11}^{(k)}x_1+\cdots+a_{1n}^{(k)}x_n,\ldots,\quad
 x_n^{(k)}=a_{n1}^{(k)}x_1+\cdots+a_{nn}^{(k)}x_n
\]
을 나타내며, 이를 줄여서 $(x^{(k)})=A_k(x)$라 쓴다. 따라서 군 $\mathfrak H$는 $x_1,\ldots,x_n$으로 이루어진 변수열 $(x)$을 $x_1^{(k)},\ldots,x_n^{(k)}$으로 이루어진 변수열들 $(x^{(k)})$로 보낸다. $A_1,\ldots,A_h$ 가운데 항등치환이 반드시 포함되므로, 변수열들 $(x^{(k)})$ 가운데에는 $(x)$도 포함된다.

%
%
군의 \textbf{다항 (절대)불변식}이라 함은 $x_1,\ldots,x_n$에 관한 다항함수로서 -- 계수는 $\mathsf{P}$\footnote{이는 일반성을 제한하지 않는다. 실제로 표수 $p$인 어떤 체의 계수를 갖는 임의의 불변식은 -- 그 불변식들이 정의되려면 이 체와 $\mathsf{P}$가 공통 확대체(共通擴大體, 곧 두 체를 모두 포함하는 체)를 가져야 한다 -- $\mathsf{P}$의 계수를 갖는 유한 개의 불변식들의 선형결합으로 나타낼 수 있다. 일반적인 경우에는 계수들을 $\mathsf{P}$ 위에서 선형독립인 원소들의 선형결합으로 나타내기만 하면 된다. 그러면 불변성 조건은 그 독립원소들에 따른 각 성분별로 항등적으로 성립한다.}에서 취하고 -- $A_1,\ldots,A_h$를 적용해도 항등적으로 변하지 않는 것을 말한다. 따라서 변수열들 $(x^{(k)})$에 관한 대칭함수 중 $P$의 계수를 갖는 것은 모두 이러한 불변식이다. 역으로 이러한 모든 불변식에 대해서는 다음이 성립한다.
\[
 f(x)=f(x^{(1)})=\cdots=f(x^{(h)})
 \quad\text{또는}\quad
 h\cdot f(x)=f(x^{(1)})+\cdots+f(x^{(h)}).
\]

%
%
따라서 군의 위수 $h$가 표수로 나누어지지 않는 경우 -- 특히 표수가 $0$인 경우 -- 변수열들 $(x^{(k)})$에 관한 대칭함수 중 계수를 $\mathsf{P}$에서 취한 것들이 불변식 전체를 이룬다. 그런데 이 대칭함수들은 유한 정수성 기저를 가지므로, 이 경우 유한성의 증명이 완성된다. 동시에 갈루아 레졸벤트(Galois'sche Resolvente)\footnote{원논문 28쪽 각주 1 참조.}의 계수계 $G_{\alpha\alpha_1\ldots\alpha_n}(x)$가 정수성 기저를 이룬다.
\[
\begin{aligned}
\Phi(z,u)&=\prod\bigl(z-u_1x_1^{(k)}-\cdots-u_nx_n^{(k)}\bigr)\\
&=z^h+\sum G_{\alpha\alpha_1\ldots\alpha_n}(x)
 z^{\alpha}u_1^{\alpha_1}\ldots u_n^{\alpha_n}
 \quad(\alpha\ne h;\ \alpha+\alpha_1+\cdots+\alpha_n=h).
\end{aligned}
\]

%
%
2. $h$가 $p$로 나누어지는 경우에는 이제 $h\cdot f(x)$로부터 $f(x)$를 추론할 수 없으므로, 군의 모든 불변식이 반드시 변수열들 $(x^{(k)})$에 관한 대칭함수일 필요는 없다. 그러나 유한 개의 불변식 $G_{\alpha\alpha_1\ldots\alpha_n}(x)$ -- 곧 갈루아 레졸벤트의 계수들 -- 에서 생성된 환 $\mathfrak R$은 모든 불변식으로 이루어진 계 $\mathfrak J$의 유한 생성 부분환을 이룬다. 그리고 $\mathfrak J$가 $\mathfrak R$에 관하여 유한성 판정기준의 가정들을 만족함을 보여야 한다.

%
%
우선 일반성을 잃지 않고\footnote{원논문 33쪽 각주 참조.} 모든 치환계수 $a_{\mu\nu}^{(k)}$로 생성되는 $\mathsf{P}$의 부분체 $\overline{\mathsf{P}}$를 계수체로 삼을 수 있다. 이 유한 개의 원소로 생성된 체는 소체(Primkörper)에 유한 개의 대수적 또는 초월적 원소를 첨가하여 얻어진다. 그리고 소체가 완전하므로, $\overline{\mathsf{P}}^{1/p}$가 $\overline{\mathsf{P}}$의 유한 차수 확대체임이 따른다.\footnote{원논문 32쪽 각주 2 참조.}

%
%
더 나아가 $\mathfrak J$가 $\mathfrak R$에 대해 정수적임을 보여야 한다. 1에 의하면 $h$개의 변수열 $x^{(k)}$ 가운데에는 본래의 $x_1,\ldots,x_n$도 나타난다. 따라서 $z=u_1x_1+\cdots+u_nx_n$으로 두면 갈루아 레졸벤트 $\Phi(z,u)$는 $u$에 관하여 항등적으로 $0$이 된다. 매번 $u_i$만 $1$로 두고 나머지 모든 $u$를 $0$으로 특수화하면 다음 관계식을 얻는다.
\[
 x_i^h+\sum x_i^\alpha G_{\alpha00\ldots\alpha_i\ldots00}(x)=0
 \quad (\alpha\ne h;\ \alpha+\alpha_i=h),
\]
이 관계식은 각 $x_i$가 $\mathfrak R$에 대해 정수적임을 보여 준다. 그러면 $x$에 관한 임의의 다항식에 대해서도 같은 사실이 성립한다. 특히 $\mathfrak J$는 $\mathfrak R$에 대해 정수적이다. 이로써 유한성 판정기준의 가정들이 만족되며, \textbf{불변식의 유한성이 증명되었다}.

%
%
3. 같은 논법으로 \textbf{상대불변식과 모듈러 불변식의 유한성}도 얻어진다는 점을 덧붙여 둔다. 각 $f(x^{(k)})$가 $f(x)$에 $\overline{\mathsf{P}}$의 0이 아닌 인자를 곱한 것과 같을 때, 다항식 $f(x)$를 이 군의 \textbf{상대불변식}이라 한다. 따라서 절대불변식 전체---특히 $G_{\alpha\alpha_1\ldots\alpha_n}(x)$로부터 얻어지는 환 $\mathfrak R$---는 모든 상대불변식으로부터 얻어지는 환의 유한 생성 부분환을 이룬다. 또한 위와 마찬가지로 $\mathfrak R$에 관하여 유한성 판정기준의 가정들이 충족된다.

%
%
다항식 $f(x)$가 모든 $k$와 $\overline{\mathsf{P}}$에서 취한 모든 구체적 값계 $x_1=\alpha_1,\ldots,x_n=\alpha_n$에 대하여 $f(x)=f(x^{(k)})$를 만족할 때, $f(x)$를 이 군의 \textbf{모듈러 불변식}이라 한다. 모든 절대불변식은 동시에 모듈러 불변식이다. 그러나 $\overline{\mathsf{P}}$가 유한 개의 원소만 가지면 그 역은 성립할 필요가 없다. 여기서도 $G_{\alpha\alpha_1\ldots\alpha_n}(x)$로부터 얻어지는 부분환 $\mathfrak R$에 관하여 유한성 판정기준의 가정들이 충족된다.

\clearpage
\setcounter{footnote}{0}
\phantomsection
\addcontentsline{toc}{section}{30. 대수적 수체와 함수체에서 아이디얼론의 추상적 구성}
%
%
%
%

\providecommand{\mC}{\mathfrak C}
\providecommand{\mE}{\mathfrak E}
\providecommand{\mF}{\mathfrak F}
\providecommand{\mG}{\mathfrak G}
\providecommand{\mH}{\mathfrak H}
\providecommand{\mM}{\mathfrak M}
\providecommand{\mN}{\mathfrak N}
\providecommand{\mO}{\mathfrak O}
\providecommand{\mX}{\mathfrak X}
\providecommand{\mc}{\mathfrak c}
\providecommand{\md}{\mathfrak d}
\providecommand{\me}{\mathfrak e}
\providecommand{\mf}{\mathfrak f}
\providecommand{\mm}{\mathfrak m}
\providecommand{\mn}{\mathfrak n}
\providecommand{\mq}{\mathfrak q}
\providecommand{\mt}{\mathfrak t}
\providecommand{\mv}{\mathfrak v}
\providecommand{\mpideal}{\mathfrak p}

\section*{30. 대수적 수체와 함수체에서 아이디얼론의 추상적 구성.}
\noindent\emph{Math. Ann. 96 (1927), 26--61쪽.}

\setcounter{footnote}{0}

\begin{center}
글쓴이\\[0.5em]
괴팅겐의 에미 뇌터.
\end{center}

이하에서는 아이디얼론이 대수적 수체의 모든 정수적 원소가 이루는 환의 아이디얼론과 일치하는 모든 환, 다시 말해 그 아이디얼들이 소아이디얼들의 거듭제곱곱으로 유일하게 표현되는 모든 환을 추상적으로 특징짓는다. 가장 잘 알려진 예로는 한 부정원의 대수함수체의 모든 정수적 원소가 이루는 환도 이에 속한다. 더 일반적으로 여러 부정원의 경우에도, 크로네커를 따라 최고차원 아이디얼로 제한하여 어떤 몫환인 함수영역으로 넘어가면 마찬가지이다.

이 아이디얼론과 완전히 동치인 공리들은 바탕이 되는 가환환에 대하여 다음과 같다\footnote{아이디얼론의 기본개념 정의는 예컨대 E. Noether, \emph{Idealtheorie in Ringbereichen}, Math. Ann. 83 (1921), 24--66쪽(이하 \emph{아이디얼론}으로 인용)을 보라. 가장 본질적인 개념들은 이 논문에서도, 특히 §§ 1, 4, 5에서 짧게 정식화한다. 완전성을 위하여 이 논문의 직접적 목적에는 필요하지 않은 내용도 얼마간 포함하였다.}.

\begin{description}
\item[I. 약수사슬정리:] 각 아이디얼이 앞선 아이디얼의 진약수인 아이디얼 사슬은 유한 단계에서 멈춘다. 다시 말해, 모든 아이디얼 약수사슬에는 그 뒤의 모든 아이디얼이 서로 같아지는 한 지표가 존재한다.
\item[II. 영아이디얼과 다른 각 아이디얼을 법으로 한 배수사슬정리.] 모두 어떤 고정된 비영아이디얼의 약수인 아이디얼들의 사슬에서, 각 아이디얼이 앞선 아이디얼의 진배수이면 그 사슬은 유한 단계에서 멈춘다.
\item[III.] 곱셈의 단위원이 존재한다.
\item[IV.] 환에는 영인자가 없다.
\item[V. 분수체 안에서의 정수적 닫힘:] 분수체의 원소 가운데 환에 대하여 정수적인 것은 모두 그 환에 속한다.
\end{description}

나는 이 공리들로부터 점차 더 강하게 제한되는 아이디얼론을 전개하여 마침내 찾는 이론에 이르고, 역으로 그 이론에서는 이 공리들이 실제로 모두 성립함을 보인다.

내가 \emph{아이디얼론} § 4에서 보였듯이, 약수사슬정리 I로부터 한 아이디얼을 서로 다른 소아이디얼들에 대응하는 유한개의 으뜸아이디얼의 최소공배수로 나타내는 표현이 나온다. 여기서는 이 증명을 짧게 되풀이한다(§ 6). 아이디얼 몫 개념을 쓰면 증명이 단순해지고, 한 가지 잘못도 바로잡아야 하기 때문이다. 원래 증명은 한 곳에서 가정하지 않은 곱셈 단위원의 존재를 사용한다\footnote{P. Urysohn이 이 잘못을 지적하였다. 내가 고친 증명은 여기 제시하는 E. Artin의 증명보다 조금 번거로웠다. Artin은 이미 더 일찍 스스로 이 틈을 메웠고, 나와 독립적으로 아이디얼 몫을 이용한 단순화도 생각해 냈다. 또한 \glqq{}축약표현\grqq{}을 도입하지 않아도 되는 또 하나의 단순화는 W. Krull의 관련 문제에서 가져왔다. W. Krull, \emph{Algebraische Theorie der Ringe III}, Math. Ann. 92 (1924), 181--213쪽, 정리 I.}.

배수사슬정리를 가정하면(§ 7), 각 비영아이디얼에 따른 잉여류환에서 한 소아이디얼은 모든 원소를 포함하는 단위아이디얼을 제외하고는 진약수를 가질 수 없다. 그러면 단위아이디얼에 대응하는 것들을 제외한 이 아이디얼들의 으뜸성분은 유일하게 결정된다. 이 결과는 조금 덜 날카로운 형태로 이미 Masazo Sono에게서 나타난다\footnote{\emph{On the Reduction of Ideals}, Memoirs of the College of Science, Kyoto University, Series A, 7 (1924), 191--204쪽.}. 그가 잉여류환에 합성열이 존재한다고 둔 가정은 이 환에서의 \glqq{}이중사슬정리\grqq{}와 동치임을 마지막에(§ 10) 짧게 보인다. 여기서 이중사슬정리란 비영아이디얼에 국한하지 않고 약수사슬정리와 배수사슬정리가 모두 성립한다는 뜻이다.

공리 III, 곧 단위원의 존재까지 성립하면 단위아이디얼은 대응 소아이디얼로 나타나지 않는다. 따라서 으뜸성분들은 유일하게 결정되고, 쌍마다 서로소이며, 그 최소공배수는 그 곱과 같다. 공리 I, II, III은 대수적 수체의 모든 유한 오더에 성립한다. 이 경우 데데킨트는 이미 Dirichlet--Dedekind, \emph{Zahlentheorie}, 제3판, 제11보유, § 172에서 완전히 전개된 증명 없이도 한 아이디얼을 쌍마다 서로소인 동종 아이디얼들(으뜸성분들)의 곱으로 유일하게 나타내는 표현을 진술하였다.

가정 IV, 곧 영인자가 없다는 사실로부터 영아이디얼 및 단위아이디얼과 다른 한 아이디얼의 서로 다른 거듭제곱들이 모두 서로 다르고 분수체가 존재한다는 것이 따른다. 끝으로 분수체 안에서 정수적으로 닫혀 있다는 공리 V까지 성립하면 으뜸아이디얼들은 --- IV 때문에 유일하게 결정되는 --- 소아이디얼의 거듭제곱이 되고, 이로써 통상적인 분해정리를 얻는다(§ 8). 이 마지막 증명은 수체의 경우 데데킨트가 이미 제3판 § 172에서 정수적 닫힘으로부터 끌어낸 따름정리에 기초한다. 즉 진분수 원소를 정수적 원소들의 몫으로 나타내되, 분자의 제곱조차 분모로 나누어지지 않게 할 수 있다. 후대의 서술에서는 역시 정수적 닫힘의 결과인 훨씬 덜 투명한 일반화된 가우스 정리나, 그에 대응하면서 조금 더 많은 내용을 말하는 모듈정리들이 이 따름정리를 대신한다. 이 정리들은 모두 적용할 때 기저원소들을 전제하지만 여기서는 아이디얼 자체를 다룬다.

역방향(§ 9)에서는 공리 V 이외의 공리들이 거의 직접 나온다. 공리 V는 진분수 원소를 언제나 그 분자의 어떤 거듭제곱도 분모로 나누어지지 않도록 나타낼 수 있음을 보이면 따른다. 이는 애초 V에서 끌어낸 따름정리를 더 강화한 것이다. 그러나 이 증명은 아이디얼 표현으로부터 통상적인 결론들, 곧 고정된 한 아이디얼을 법으로 한 주아이디얼 표현과 분수아이디얼, 즉 체 안의 아이디얼 몫 이론을 이끌어 낸 뒤에야 가능하다. 소아이디얼 거듭제곱에 의한 표현으로부터 정수적 닫힘이 따른다는 사실도 데데킨트에게 이미 알려져 있었음은 제3판 § 172, 522쪽 아래의 주석이 보여 준다.

공리 I--V에 따르면, 일반적으로는 서로 구별되는 네 분해, 곧 쌍마다 서로소인 최소 아이디얼들로의 분해, 서로 상대소인 아이디얼들로의 분해, 극대 으뜸성분들로의 분해, 기약아이디얼들로의 분해가 서로 일치한다. 역으로 이 일치와 공리 I--IV로부터 정수적 닫힘이 나온다. 영인자를 갖는 환에서는 나머지 공리들이 성립하더라도 --- 이때 V는 몫환에서의 정수적 닫힘으로 정의한다 --- 분해들이 반드시 일치하지는 않는다. 유한 오더의 어떤 아이디얼에 따른 잉여류환이 그 예이다(§ 9).

§ 1에서는 한 환에 대하여 정수적인 원소들의 이론을, 정수적 닫힘에서 나오는 데데킨트의 따름정리까지 개관한다. § 2에서는 계수영역이자 곱셈자영역인 이 환을 갖는 선형형식들의 모듈로 사슬정리들이 어떻게 옮겨지는지를 보인다\footnote{이 형태의 모듈정리는 Prüfer에게서 나왔으며(\emph{Neue Begründung der algebraischen Zahlentheorie}, § 2, Math. Ann. 94 (1924), 198--243쪽), 통상적인 기저정리의 이행보다 더 투명하다.}. 이로부터 § 3에서는 대수적 확대체(제1종)의 정수적 원소들로 넘어갈 때 모든 오더에서 공리 I--IV가 보존되고, 모든 정수적 원소로 이루어진 오더에는 V도 성립한다는 결론을 끌어낸다. 대수적 수체에 관해서는 물론 알려진 이 사실 덕분에 처음에 말한 함수체들도 이 체계 아래 들어온다. 특히 여러 부정원의 정수계수 다항식들에서 유도된 함수영역이 공리들을 만족한다는 간단한 증명으로 크로네커의 아이디얼론도 완전히 포괄된다\footnote{따라서 차이는 판별식정리에서 비로소 나타난다. 한 소수를 법으로 한 함수영역의 잉여류환이 비완전체가 되어 제2종 확대가 생길 수 있기 때문이다.}.

그 뒤에 전개되는 아이디얼론은 분류를 위해서만 쓰인 §§ 2, 3을 전제하지 않는다. §§ 4, 5에서는 공리 I--V와 독립적인 기초를 제시한다. 이로써 이론의 어느 부분이 유한성 가정과 무관한지가 원래의 정초에서보다 더 분명해진다.

\begin{center}
{\bfseries\large § 1.\par}
\vspace{0.35em}
{\bfseries\large 정수적 원소의 이론.\par}
\end{center}

영인자가 없고 곱셈 단위원 \(e\)를 갖는(공리 III, IV) 가환환\footnote{잘 알려져 있듯이, 한 집합에 동치관계와 두 연산인 덧셈과 곱셈이 주어져 통상적인 법칙들을 만족하면 환이 정의된다. 덧셈은 결합적이고 가환적이며 유일하게 역산할 수 있고, 곱셈은 결합적이며 --- 가환환이면 가환적이기도 하고 --- 덧셈에 대하여 분배적이다. 여기서 바탕이 되는 상등의 정의가 집합론적 동일성이 아니라면, 각 부분계에서도 같은 상등의 정의를 유지하기 위하여 그러한 부분계(부분환, 모듈, 아이디얼)는 한 원소와 더불어 그와 같은 모든 원소도 포함해야 한다. 마찬가지로 모든 확대에서는 새 상등의 정의가 옛 정의를 포함한다고 가정해야 한다. \glqq{}유한개의 원소\grqq{}, \glqq{}유일한 대응\grqq{} 등의 모든 개념은 이 상등의 정의에 따라 이해한다. 예컨대 \glqq{}정확히 하나의 원소\grqq{}가 있다는 말은 서로 같은 것들을 제외하면 하나뿐이라는 뜻이다. 같은 원소들을 한 류로 묶고 이 류들을 새 원소로 삼으면 언제나 상등에서 집합론적 동일성으로 넘어갈 수 있다. 상등의 정의에 관한 이 주석들은 R. Hölzer에게서 나왔다.} \(\mT\)를 놓고, \(\mT\) 안에 단위원을 포함하는 고정된 부분환 \(\mR\)을 택한다. 모듈, 오더, 정수적이라는 개념은 줄곧 이 고정된 환 \(\mR\)에 관한 것으로 하며, 뒤에서는 간단히 하려고 표기에서 이를 생략한다\footnote{영인자를 허용하는 대신 \(\mR\)에 약수사슬정리를 가정해도 정수적 원소의 이론은 유지되며, § 2에 따라 보조정리도 증명할 수 있다. 그러나 이 형태의 정수적 원소는 뒤에서 쓰지 않으므로 각주에서만 언급한다. 영인자가 있어도 여기서 다루는 정의들은 모두 유지된다. 잘 알려져 있고 쉽게 알 수 있듯이, 그 대부분은 훨씬 약한 가정만을 필요로 한다. 다만 데데킨트의 정식화 \glqq{}어떤 수가 외피를 가지면 정수적이다\grqq{}는 유지되지 않는다. 영인자가 생기면 두 유한 모듈의 몫이 유한하지 않을 수 있기 때문이다. 그래서 여기서는 오더도 모듈 몫으로가 아니라 직접 정의한다.}.
(\(\mR\)의 원소는 라틴 문자로, 일반적인 \(\mT\)의 원소는 그리스 문자로 나타낸다.)

1. \(\mT\)의 원소들로 이루어진 계 \(\mM\)이 두 원소 \(\alpha,\beta\)와 더불어 그 차도 포함하고, \(\alpha\)와 더불어 모든 \(r\in\mR\)에 대한 \(r\alpha\)도 포함하면, 통상대로 이를 \emph{\(\mR\)-모듈}, 줄여서 \emph{모듈}이라 한다. \(\mM\)의 모든 원소가 \(\mN\)에도 들어 있을 때 \(\mM\)은 \(\mN\)으로 나누어진다,
\[
        \mM\equiv 0(\mN),
\]
고 하고, \(\mM\)을 배수, \(\mN\)을 약수라 한다. 원소가 모두 \(\mR\)에 속하는 \(\mR\)-모듈을 \(\mR\)의 \emph{아이디얼}이라 한다.

분명히 \(\mT\)의 임의의 많은 모듈의 교집합, 곧 최소공배수 \([\ldots\mM_i\ldots]\)도 다시 모듈이고, \(\mT\) 자체도 모듈이다. 따라서 \(\mT\)의 원소들의 임의의 계 \(\Sigma\)에서 \emph{생성되는 모듈}은 \(\Sigma\)를 포함하는 모든 모듈의 교집합으로서 유일하게 존재한다. 임의의 모듈계의 합, 곧 최대공약모듈 \((\ldots\mM_i\ldots)\)은 그 합집합에서 생성되는 모듈이다. 두 모듈의 곱 \(\mM\mN\)은 \(\mu\)가 \(\mM\)의 모든 원소를, \(\nu\)가 \(\mN\)의 모든 원소를 달릴 때 원소 \(\mu\nu\)들에서 생성되는 모듈로 정의한다. 환원소에 대해 결합법칙과 교환법칙이 성립하므로 모듈의 곱도 이를 만족한다. \(\mT\)의 분배법칙에서 모듈에 대한 분배법칙
\[
        (\ldots,\mM_i,\ldots)\mN=(\ldots,\mM_i\mN,\ldots)
\]
도 따른다. 실제로 \(\mT\)의 바로 그 분배법칙에 따라 양변은 모두 모든 원소 \(\mu_i\alpha\)에서 생성되는 모듈이다.

\(\mR\) 자체가 모듈이므로, \(\mR\)이 곱셈자영역이라는 조건은 \(\mR\mM\equiv0(\mM)\)으로 나타낼 수 있다. 가정에 따라 \(\mR\)이 단위원을 포함하므로 \(\mM\equiv0(\mR\mM)\)도 성립하고, 따라서 \(\mM=\mR\mM\)이다.

\(\mM\)이 \emph{유한}하다는 것은 \(\mM\)을 생성하는 원소가 \emph{유한개}인 계 \(\Sigma\)가 적어도 하나 있다는 뜻이다. \(\Sigma\)의 원소 \(\mu_1,\ldots,\mu_r\)를
\[
        \mM=(\mu_1,\ldots,\mu_r)
\]
의 모듈기저라 한다. 두 유한 모듈의 합과 곱, 따라서 유한개의 유한 모듈의 합과 곱도 다시 유한하다. 각각 기저원소들의 합집합과 곱들에서 생성되기 때문이다. 곱에 대한 말은 \(\mM\)의 모든 원소 \(\mu\)가 기저에 따라 \(r_1\mu_1+\ldots+r_k\mu_k\)로 표현된다는 사실\footnote{\(\mR\)이 단위원을 포함하지 않으면 기저표현은 \(r_1\mu_1+\ldots+r_k\mu_k+n_1\mu_1+\ldots+n_k\mu_k\) 꼴이 된다. 여기서 정수 \(n\)은 환원소가 아니라 반복된 덧셈 또는 뺄셈을 나타내는 기호이다.}과 \(\mT\)의 곱셈 교환법칙에서 따른다. 이 두 사실은 여기서 처음으로 쓰인다.

2. \(\mR\)의 원소를 모두 포함하면서 \(\mT\) 안에 놓인 \(\mR\)의 확대환을 \emph{\(\mR\)-오더}, 줄여서 \emph{오더}라 한다. \(\mT\) 안의 임의의 많은 오더의 교집합도 다시 오더이고 \(\mT\)도 오더이므로, \(\mT\)의 원소들의 임의의 계 \(\Sigma\)에서 생성되는 오더가 유일하게 존재한다. 오더의 합과 곱은 모듈의 경우와 마찬가지로 정의할 수 있고, 특히 모든 오더에 대해 \(\mO^2=\mO\)이다.

모든 오더는 동시에 \(\mR\)-모듈이다. 오더가 유한하다는 것은 그것이 유한 모듈이라는 뜻이다. 합과 곱은 모듈의 합과 곱으로 만들어지므로 1에 따라 유한 오더의 합과 곱도 다시 유한 오더이다. 또한 다음 \emph{추이법칙}이 성립한다. \(\mA\)가 유한 \(\mR\)-오더이고 \(\mB\)가 유한 \(\mA\)-오더이면, \(\mB\)는 유한 \(\mR\)-오더이기도 하다. 실제로 \(\mB\)는 동시에 \(\mR\)-오더, 따라서 \(\mR\)-모듈이 된다. 그런데 \(\alpha_1,\ldots,\alpha_r\)가 \(\mR\)에 관한 \(\mA\)의 모듈기저이고 \(\beta_1,\ldots,\beta_s\)가 \(\mA\)에 관한 \(\mB\)의 모듈기저이면, 곱 \(\alpha_i\beta_k\)들이 분명 \(\mR\)에 관한 \(\mB\)의 모듈기저를 이룬다.

3. \(\mT\)의 원소 \(\alpha\)에서 생성되는 오더 \(\mR_\alpha\)가 \emph{유한}하면 \(\alpha\)를 \emph{\(\mR\)에 대하여 정수적}, 줄여서 \emph{정수적}이라 한다. 달리 말하면 모듈의 약수사슬
\[
        \mA_\mu=(\alpha^0,\alpha,\ldots,\alpha^{\mu-1})
\]
이 유한 단계에서 멈추어 \(\mA_n=\mA_{n+1}=\ldots=\mA_{n+\nu}=\ldots\)가 된다는 뜻이다.

이 정의는 4에서 쓸 다음 두 번째 형태로도 쓸 수 있다. \(\mT\)의 원소 \(\alpha\)가 정수적이라는 것은 영모듈과 다른 주모듈 \(\mC\), 곧 한 원소 \(\gamma\ne0\)에서 생성되는 \(\mC\)가 적어도 하나 존재하여 오더 \(\mR_\alpha\)와 \(\mC\)의 곱이 유한 모듈이 된다는 뜻이다. 달리 말하면 모듈의 사슬 \(\mC\mA_\mu\)가 유한 단계에서 멈춘다\footnote{\(\mR\)에 영인자를 허용하면 정칙 주모듈의 존재를 요구해야 한다. 즉 \(\gamma\)가 영인자가 아닌 정칙원소라고 가정해야 하며, 영인자가 없는 환에서는 이는 \(\gamma\ne0\)과 같다.}.

끝으로 이 정의는 통상적인 다음 형태와도 같다. \(\mT\)의 원소 \(\alpha\)가 정수적이라는 것은 \(\mR\)의 원소 \(r_i\)를 계수로 하는 방정식
\[
        \alpha^n+r_1\alpha^{n-1}+\ldots+r_n=0
\]
을 적어도 하나 만족한다는 뜻이다.

이 여러 정의는 실제로 서로 같다. \(\mR_\alpha\)는 \(\alpha\)의 모든 거듭제곱에서 생성되는 모듈과 같으므로, 유한한 \(\mR_\alpha\)에는 언제나 이 거듭제곱들로 이루어진 모듈기저를 택할 수 있다. 예컨대 그 기저가 \(e=\alpha^0,\alpha,\ldots,\alpha^{n-1}\)이면 \(\mA_\mu\)의 사슬이 \(\mA_n\)에서 멈춘다는 뜻이고, 역으로 \(\mA_n=\mA_{n+\nu}\ldots\)에서 이 기저가 나온다. \(\mA_\mu\)의 사슬이 멈춘다는 것은 방정식
\[
        \alpha^n+r_1\alpha^{n-1}+\ldots+r_n=0
\]
이 성립하는 것과도 같다. \(\mR_\alpha\)가 유한하면 \(\mC\mR_\alpha\)도 유한하므로 \(\mC\mA_\mu\)의 사슬이 멈춘다. 역으로 \(\mC\mR_\alpha\)가 유한하여 \(\mC\mA_\mu\)가 멈추면 \(\mA_\mu\)도 멈추고, 따라서 \(\mR_\alpha\)가 유한하다. \(\mC\)를 비영 주모듈로 가정했으므로 \(\mC\mA_n=\mC\mA_{n+\nu}\)에서 언제나 \(\mA_n=\mA_{n+\nu}\)가 나오기 때문이다.

정수적 원소들이 환을 이룬다는 사실과 정수적 종속성의 추이법칙은 다음 보조정리에 기초한다.

\medskip
\noindent\textbf{보조정리.} \(\mS\)가 \(\mT\) 안의 유한 오더이고 \(\alpha\)가 \(\mS\)의 임의의 원소이면, \(\alpha\)에서 생성되는 오더 \(\mR_\alpha\)는 유한하다. 다시 말해, 유한 오더의 모든 원소는 정수적이다.

증명은 통상적인 추론에서 나온다. \(\sigma_1,\ldots,\sigma_n\)이 \(\mS\)의 모듈기저이면, 환의 성질에 따라 \(\alpha\sigma_i\)도 \(\mS\)에 속한다. 따라서
\[
        \alpha\sigma_i=r_{i1}\sigma_1+\ldots+r_{in}\sigma_n,
\]
이고, 여기서
\[
        \left|r_{ik}-\alpha e_{ik}\right|=0,
\]
을 얻는다. \(i\ne k\)이면 \(e_{ik}=0\), \(e_{ii}=e\)이므로 이 식은 \(\alpha\)가 정수적임을 보인다. 행렬식이 사라짐을 보일 때 \(\mT\), 따라서 \(\mS\)에 영인자가 없다는 가정을 쓴다\footnote{\(\mR\)에 약수사슬정리를 가정하되 영인자를 허용하면, 이 보조정리는 § 2의 모듈정리에서 직접 따른다. 그 정리에 따르면 사슬정리가 \(\mS\)의 모듈들로 옮겨진다. 수체에 대한 이 증명 역시 데데킨트에게서 나왔고(제4판, § 173, III), 문헌에서 사슬정리를 처음 적용한 예이다. 아래 4의 따름정리들에서 데데킨트는 그 대신 대수적 수체의 한 아이디얼에 따른 잉여류의 수가 유한하다는 더 특수한 사실을 아직 사용한다.}.

\(\mT\) 안에서 정수적인 모든 원소의 계 \(\mS\)는 오더를 이룬다. \(\mS\)는 \(\mR\)을 포함한다. 실제로 \(\mR\)의 원소들은 정수적이다. \(\mR\)은 단위원을 기저로 하는 유한 \(\mR\)-모듈이기 때문이다. \(\mS\)는 환의 성질도 갖는다. 임의의 두 정수적 원소 \(\alpha,\beta\)에서 생성되는 오더 \(\mR_{\alpha,\beta}\)는 곱 \(\mR_\alpha\cdot\mR_\beta\)와 같아서 유한하다. 따라서 보조정리에 의해 유한 오더의 원소 \(\alpha\pm\beta\)와 \(\alpha\beta\)는 정수적이다.

\medskip
\noindent\textbf{정수적 종속성의 추이법칙.} \(\mS\)가 정수적 원소들로 이루어진 오더이면, \(\mS\)에 대하여 정수적인 \(\mT\)의 모든 원소는 \(\mR\)에 대해서도 정수적이다. 정의에 따라 \(\alpha\)는 방정식
\[
        \alpha^m+\sigma_1\alpha^{m-1}+\ldots+\sigma_m=0
\]
을 만족한다. 그러므로 \(\alpha\)는 \(\sigma_1,\ldots,\sigma_m\)에서 생성되는 오더 \(\overline{\mS}=\mR_{\sigma_1,\ldots,\sigma_m}\)에 대해서도 정수적이다. 이 오더는 곱 \(\mR_{\sigma_1}\cdot\mR_{\sigma_2}\cdot\ldots\cdot\mR_{\sigma_m}\)이므로 유한하다. 따라서 2의 추이법칙에 따라 \(\alpha\)에서 생성되는 유한 \(\overline{\mS}\)-오더 \(\overline{\mS}_\alpha\)는 유한 \(\mR\)-오더이기도 하고, 보조정리에 따라 \(\alpha\)는 유한 오더의 원소로서 정수적이다.

4. \(\mR\)에 대하여 정수적인 \(\mT\)의 모든 원소가 \(\mR\)에 속하면 \(\mR\)을 \emph{\(\mT\) 안에서 정수적으로 닫혀 있다}고 한다. 3의 정수적 원소 정의의 두 번째 형태에 따르면, \(\mT\)의 원소 \(\alpha\)가 \(\mR\)에 속할 필요충분조건은 영모듈과 다른 주모듈 \(\mC\)가 적어도 하나 존재하여 \(\mC\cdot\mR_\alpha\)가 유한 모듈을 이루는 것이다\footnote{\(\mR\)에 영인자를 허용하면 \(\mC\)를 다시 정칙 주모듈로 가정해야 하고, 따름정리 I의 원소 \(c\)도 정칙이라고 가정해야 한다.}.

데데킨트는 정수적으로 닫힌 환이라는 개념에서 아이디얼론에 이 개념을 적용할 수 있게 하는 중요한 따름정리 두 개를 끌어냈다(제3판, § 172).

\medskip
\noindent\textbf{데데킨트의 따름정리 I.} \emph{\(\mR\)이 \(\mT\) 안에서 정수적으로 닫혀 있고, 추가로 \(\mR\)의 아이디얼에 대한 약수사슬정리(공리 I)가 성립한다고 하자. \(\beta\)가 \(\mT\)의 정수적이지 않은 원소이고 \(c\ne0\)이 \(\mR\)의 원소이면, 원소들의 열
\[
        c,\;c\beta,\ldots,c\beta^r,\ldots
\]
에서 \(c,\ldots,c\beta^\sigma\)는 정수적이고 그 뒤의 모든 원소는 정수적이지 않게 하는 지수 \(\sigma\geqq0\)가 존재한다.}

모든 원소 \(c\beta^r\)가 정수적이라면 \(\mR\)의 정수적 닫힘에 따라 모두 \(\mR\)에 속한다. 그러면 모듈의 사슬 \(\mC\mB_1,\mC\mB_2,\ldots,\mC\mB_r,\ldots\), 여기서 \(\mC\)는 \(c\)에서, \(\mB_r\)은 \(\beta^0,\ldots,\beta^{r-1}\)에서 생성되는 모듈인데, 이는 \(\mc_r=(c,c\beta,\ldots,c\beta^{r-1})\)인 아이디얼의 사슬 \(\mc_1,\ldots,\mc_r,\ldots\)로 바뀐다. 이 사슬은 가정에 따라 유한 단계에서 멈춘다. 그러면 \(\mR_\beta\)가 유한하고 \(\beta\)가 정수적이 되어 가정에 모순된다. 따라서 \(c\beta^r\) 가운데 정수적이지 않은 원소가 있다. 더 나아가 하나가 정수적이지 않으면 뒤의 모든 원소도 정수적이지 않다. 실제로 \(c\beta^\tau\)가 정수적, 따라서 \(\mR\)에 속한다면 \(\rho<\tau\)인 \(c\beta^\rho\)는 방정식
\[
        (c\beta^\rho)^\tau-c^{\,\tau-\rho}(c\beta^\tau)^\rho=0
\]
을 만족하므로 역시 정수적이다. 따라서 \(c\beta^{\sigma+1}\)이 최초의 정수적이지 않은 원소라면 \(c\)가 정수적이므로 \(\sigma\geqq0\)이고, 이것이 찾는 지수이다.

확대환 \(\mT\)를 \(\mR\)의 분수체, 곧 모든 원소쌍을 첨가하여 생기는 체\footnote{잘 알려진 Steinitz의 정식화, J. f. M. 137. \(\mR\)에 영인자를 허용하면, 분수체 대신 \(\mR\)의 정칙원소를 분모로 하는 모든 몫쌍을 첨가하여 얻는 몫환을 써야 한다. 이 경우에는 데데킨트의 따름정리 II가 더 이상 성립하지 않는다. 일반적으로 \(n\)이 정칙원소가 되지 않기 때문이다.}로 특수화하면 다음이 따른다.

\medskip
\noindent\textbf{데데킨트의 따름정리 II.} \emph{\(\mR\)이 그 분수체 \(\mK\) 안에서 정수적으로 닫혀 있고 \(\mR\)에서 아이디얼에 대한 약수사슬정리가 성립한다고 하자. \(\beta\)가 \(\mK\)의 정수적이지 않은 원소이면, \(\beta\)를 \(\mR\)의 원소들의 몫
\[
        \beta=m/n
\]
으로 나타내되 \(m^2/n\)도 정수적이지 않게 할 수 있다.}

\(\beta=b/c\), \(c\ne0\)이라 쓰면 열 \(c,c\beta=b,c\beta^2\ldots\)의 첫 두 원소는 확실히 정수적이다. 따라서 따름정리 I의 지수 \(\sigma\)는 \(\sigma\geqq1\)이다. \(m=c\beta^\sigma\), \(n=c\beta^{\sigma-1}\)로 놓으면 --- \(\sigma\geqq1\)이므로 이들은 열의 정수적 원소들이다 --- \(m/n=\beta\)이고 \(m^2/n=c\beta^{\sigma+1}\)은 정수적이지 않다.

\emph{정수적 닫힘의 추이법칙}은 다음 형태로 성립한다. \(\mR\)을 \(\mT\)의 임의의 환이라 하고, \(\mS\)를 \(\mT\) 안에서 \(\mR\)에 대하여 정수적인 모든 원소의 계라 하자. 그러면 \(\mS\)는 \(\mS\)에 대하여 정수적인 \(\mT\)의 모든 원소로도 이루어진다. 실제로 \(\mS\)에 대하여 정수적인 모든 원소는 3의 추이법칙에 따라 \(\mR\)에 대해서도 정수적이므로 \(\mS\)에 속한다.

\subsection*{§ 2. 유한 모듈영역에서의 사슬정리}

사슬정리들이 옮겨진다는 모듈정리는 여기서의 적용에서는 확대환의 모듈에 대해서만 나타난다. 그러나 일반 모듈영역의 경우에도 증명이 같으므로 그러한 일반 영역을 바탕에 둔다.

원소 \(\alpha,\beta\ldots\)의 계 \(M\)에 다음 두 연산, 곧 원소들의 덧셈 연산과 원소에 \(\mR\)의 원소를 곱하는 연산이 주어져 언제나 \(M\)의 원소를 유일하게 낳고, \(\mathsf{M}\)이 덧셈에 관하여 아벨군을 이루며 곱셈에는 결합법칙이 성립하고 분배법칙이 두 형태로 모두 성립하면, \(\mathsf{M}\)을 \emph{환 \(\mR\)에 관한 모듈영역}이라 한다\footnote{\emph{아이디얼론}, § 9를 보라.}.

모듈영역에는 § 1, 1에서 제시한 모듈 정의 가운데 곱셈을 언급하지 않는 것들이 분명 모두 유지된다. 특히 \(\mathsf{M}\)에서 유한개의 원소 \(\xi_1,\ldots,\xi_k\)를 택하여 \(\mathsf{M}\)이 이 원소들에서 생성되는 모듈, 곧 \(\mR\)에 단위원이 있을 때 모든 선형형식 \(r_1\xi_1+\ldots+r_k\xi_k\)의 계와 같으면 \(\mathsf{M}\)을 \emph{유한 모듈영역}이라 한다. 단위원이 없으면 여기에 항 \(n_i\xi_i\)들도 더해진다(각주 8 참조).

\medskip
\noindent\textbf{모듈정리.} \emph{\(\mathsf{M}\)이 단위원을 갖는 가환환 \(\mR\)에 관한 유한 모듈영역이고, \(\mR\)에서 아이디얼의 약수사슬정리 또는 배수사슬정리가 성립하면, \(\mathsf{M}\)에서 모듈의 약수사슬정리 또는 배수사슬정리가 각각 성립한다.}

증명은 \(\mathsf{M}\)의 각 모듈에 \(\mR\)의 유한개의 아이디얼로 이루어진 한 계를 유일하게 대응시키되, 모듈의 진약수 또는 진배수에는 대응 아이디얼 가운데 적어도 하나의 진약수 또는 진배수가 각각 대응하게 하는 데 기초한다.

\(\mathsf{M}\)의 한 원소가 \(\xi_1,\ldots,\xi_i\)의 선형형식으로 적어도 한 가지 방식으로 표현되지만 \(\xi_1,\ldots,\xi_{i-1}\)의 선형형식으로는 표현되지 않으면 그 원소의 \emph{길이가 \(i\)}라고 한다. 이제 \(\mathsf{M}\)의 임의의 모듈 \(\mA\)에 \(\mR\)의 \(k\)개 아이디얼 \(\mathfrak a_1,\ldots,\mathfrak a_k\)를 다음과 같이 대응시킨다. \(\mA_i\)는 길이가 \(\leqq i\)인 \(\mA\)의 모든 원소가 이루는 모듈이고, \(\mathfrak a_i\)는 \(\mA_i\)에서 \(\xi_i\)의 계수들이 이루는 계이다. 우선 다음이 따른다.

\(\mB\)가 \(\mA\)의 \emph{진약수}이고 \(\mathfrak b_1,\ldots,\mathfrak b_k\)가 \(\mB\)에 대응하는 아이디얼들이면, \(\mathfrak b\) 가운데 적어도 하나는 대응하는 \(\mathfrak a\)의 \emph{진약수}이다. 실제로 \(\mA\equiv0(\mB)\)에서 \(\mA_\lambda\equiv0(\mB_\lambda)\), 따라서 \(\lambda=1,2,\ldots,k\)에 대해 \(\mathfrak a_\lambda\equiv0(\mathfrak b_\lambda)\)가 나온다. 가정에 따라 \(\mB\)에는 \(\mA\)에 들어 있지 않은 원소가 있고, 따라서 그런 원소 가운데 최소 길이 \(i\)인 것도 있다. 그러면 \(\mathfrak b_i\)는 \(\mathfrak a_i\)의 진약수이다. 실제로 그러한 최소 길이의 원소
\[
        \beta=b_1\xi_1+\ldots+b_i\xi_i
\]
를 \(\mB\)에서 택하면 \(b_i\equiv0(\mathfrak b_i)\)이지만 \(b_i\not\equiv0(\mathfrak a_i)\)이다. 만일 \(b_i\equiv0(\mathfrak a_i)\)라면 원소
\[
        \alpha=a_1\xi_1+\ldots+b_i\xi_i\equiv0(\mA)
\]
가 존재하고, 따라서 \(\beta\)보다 길이가 짧으면서 \(\beta-\alpha\equiv0(\mB)\), \(\beta-\alpha\not\equiv0(\mA)\)인 원소가 존재하여 최소성에 모순된다.

이제 약수사슬 또는 배수사슬
\[
        \mA_1,\mA_2,\ldots,\mA_r,\ldots
\]
이 주어지고 \(\mR\)에서 각각 약수사슬정리 또는 배수사슬정리가 성립한다고 하자. \(\mathfrak a_{r1},\ldots,\mathfrak a_{rk}\)가 \(\mA_r\)에 대응하는 아이디얼이면, \(\lambda=1,\ldots,k\)마다 열
\[
        \mathfrak a_{1\lambda},\mathfrak a_{2\lambda},\ldots,\mathfrak a_{r\lambda},\ldots
\]
은 약수사슬 또는 배수사슬을 이룬다. 따라서 어떤 지표 \(\mu\)가 존재하여 모든 \(\lambda\)에 대해 아이디얼 \(\mathfrak a_{\mu\lambda}\)가 그 뒤의 모든 아이디얼과 같다. 그러면 위에서 증명한 바에 따라 모듈 \(\mA_\mu\)도 그 뒤의 모든 모듈과 같아진다. 이로써 \emph{사슬정리들이 옮겨졌다}.

곧바로 다음을 얻는다.

\medskip
\noindent\textbf{모듈정리의 따름정리.} \emph{\(\mR\)에서 영아이디얼과 다른 각 아이디얼을 법으로 한 배수사슬정리가 성립하면, \(\mathsf{M}\)에서는 대응 아이디얼 \(\mathfrak c_1,\ldots,\mathfrak c_k\)가 모두 비영아이디얼인 각 모듈 \(\mC\)를 법으로 한 배수사슬정리가 성립한다.}
실제로 이 가정 아래에서는 배수사슬 \(\mathfrak a_{1\lambda},\ldots,\mathfrak a_{r\lambda},\ldots\)가 유한 단계에서 멈춘다.

\medskip
\noindent\textbf{덧붙이는 주석.} \(\mR\)이 단위원 없는 환이어도 약수사슬정리와 비영아이디얼을 법으로 한 배수사슬정리는 여전히 옮겨진다. 실제로 \(\mA\) 대신
\[
        a_1\xi_1+\ldots+a_k\xi_k+n_1\xi_{k+1}+\ldots+n_k\xi_{2k}
\]
꼴의 모든 원소가 이루는 계 \(\overline{\mA}\)를 고찰한다. \(\xi_{k+\lambda}\)를 \(\xi_\lambda\)로 바꾸면 이 계는 다시 \(\mA\)가 된다. 위와 마찬가지로 \(\overline{\mA}\)에 \(2k\)개 아이디얼 \(\mathfrak a_1,\ldots,\mathfrak a_k,\mathfrak n_1,\ldots,\mathfrak n_k\)를 대응시킬 수 있다. 여기서 \(\mathfrak a_\lambda\)는 \(\mR\)의 아이디얼이고 \(\mathfrak n_\lambda\)는 정수환의 아이디얼이다. \(\mathfrak n_\lambda\)에는 약수사슬정리와 각 비영아이디얼을 법으로 한 배수사슬정리가 성립하므로, \(\overline{\mA}\), 따라서 \(\mA\)에 대한 약수사슬정리의 증명은 그대로 유지된다. 배수사슬정리에서는 \(\mC\), 또는 \(\overline{\mC}\)에 대응하는 \(2k\)개 아이디얼이 모두 비영아이디얼이라고 요구해야 한다.

\providecommand{\mK}{\mathfrak{K}}
\providecommand{\mL}{\mathfrak{L}}
\providecommand{\mT}{\mathfrak{T}}
\providecommand{\mC}{\mathfrak{C}}
\providecommand{\mF}{\mathfrak{F}}
\providecommand{\ma}{\mathfrak{a}}
\providecommand{\mb}{\mathfrak{b}}
\providecommand{\mc}{\mathfrak{c}}
\providecommand{\md}{\mathfrak{d}}
\providecommand{\me}{\mathfrak{e}}
\providecommand{\mf}{\mathfrak{f}}
\providecommand{\mm}{\mathfrak{m}}
\providecommand{\mn}{\mathfrak{n}}
\providecommand{\mo}{\mathfrak{o}}
\providecommand{\mq}{\mathfrak{q}}
\providecommand{\mt}{\mathfrak{t}}
\providecommand{\mv}{\mathfrak{v}}
\providecommand{\mpideal}{\mathfrak{p}}

\subsection*{§ 3. 유한 확대체로의 이행}

수체와 함수체를 이 체계 아래 넣으려면, 서론에서 정식화한 공리 I--V가 유한 확대체로 넘어갈 때 어떻게 옮겨지는지를 보여야 한다.

1. \emph{\(\mR\)에서 공리 II, 곧 배수사슬정리를 제외한 공리 I--V가 모두 성립한다고 하자. \(\mK\)는 \(\mR\)의 분수체이고 \(\mL\)은 \(\mK\)의 유한 제1종 확대체라 하자\footnote{Steinitz에 따르면, 모든 원소가 어떤 적당한 확대체에서 서로 다른 일차인자들로 분해되는 소함수의 영점이면 그 대수적 확대체를 \glqq{}제1종\grqq{}이라 한다.}. 그러면 \(\mL\) 안에서 \(\mR\)에 대하여 정수적인 모든 원소의 계 \(\mS\)에는 같은 공리들이 성립하며, \(\mS\) 안의 모든 오더에는 정수적 닫힘을 제외한 이 공리들이 모두 성립한다.}

§ 1, 3에 따라 \(\mS\)는 공리 III, IV, 곧 단위원의 존재와 영인자가 없다는 성질이 성립하는 환을 이룬다. § 1, 4의 끝에 따라 \(\mS\)는 \(\mL\) 안에서 정수적으로 닫혀 있다. 또한 \(\mL\)의 모든 원소는 \(\mR\)의 적당한 원소를 곱하면 \(\mS\)의 원소가 되므로 \(\mL\)은 \(\mS\)의 분수체이다. 따라서 공리 V가 성립한다.

I의 증명은 알려진 추론으로 얻어진다. 제1종 확대인 \(\mL\)은 \(\mK\)에 한 원소 \(\alpha\)를 첨가하여 생기며, 위에 따라 \(\alpha\)가 \(\mR\)에 대하여 정수적이라고 가정할 수 있다. \(\alpha\)와 더불어 \(\mK\)에 대한 \(\mL\)의 갈루아체 안에 있는 켤레원소 \(\alpha',\alpha'',\ldots\)도 정수적이며, 따라서 그 차들의 곱과 이 차곱의 제곱 \(D\)도 정수적이다. \(\alpha,\alpha',\ldots\)가 모두 서로 다르므로 \(D\)는 \(\mK\)의 비영원소이고, \(\mR\)이 \(\mK\) 안에서 정수적으로 닫혀 있으므로 \(\mR\)의 원소이다. 통상적인 추론에 따라 \(\mS\)는 \(\mR\)에 관한 유한 모듈영역 \(M\), 곧 원소 \(\xi_i=\alpha^i/D\)들이 생성하는 영역에 포함된다. 이를 보이려면 \(\mS\)의 한 원소를 \(\xi_i\)들로 나타낸 표현에서 켤레원소들로 넘어간 다음 표현계수에 대한 선형방정식계를 풀기만 하면 된다. § 2의 모듈정리에 따라 \(\mathsf{M}\) 안의 모든 \(\mR\)-모듈, 특히 \(\mS\)의 모든 아이디얼에 약수사슬정리가 성립한다. \(\mS\) 안의 임의의 오더의 아이디얼들도 \(M\) 안의 \(\mR\)-모듈이므로 이들에도 약수사슬정리가 성립하고, 모든 오더에는 공리 III, IV도 성립한다.

2. \emph{\(\mR\)에서 공리 I--V가 모두 성립하면 \(\mS\)에서도 모두 성립한다. \(\mS\) 안의 모든 오더에는 정수적 닫힘을 제외한 모든 공리가 여전히 성립한다.}

1을 고려하면 공리 II만 더 증명하면 된다. § 2의 모듈정리의 따름정리에 따라 다음을 보이면 충분하다. \(\mS\)의 비영아이디얼을 \(M\) 안의 \(\mR\)-모듈 \(\mC\)로 볼 때, \(\mC\)에 대응하는 \(\mR\)의 아이디얼 \(\mc_i\)가 모두 비영아이디얼이어야 한다. 그런데 \(\mC\)는 \(\mK\)에 관하여 \(\mathsf{M}\)과 같은 유한 선형 랭크를 갖는다. 실제로 \(\mC\)는 임의의 한 원소 \(\gamma\)와 더불어 \(\gamma\alpha^i\)도 포함한다. 따라서 각 \(\xi_i\)에 대해 \(c_i\xi_i\)가 \(\mC\)에 속하게 하는 비영원소 \(c_i\in\mR\)가 존재한다. \(\xi_i\)에 의한 표현의 유일성에 따라 --- 지표는 체의 차수보다 작은 값만 달린다 --- \(c_i\)는 \(\mc_i\)에 속하므로 \(\mc_i\)는 비영아이디얼이다. 이 논의는 \(\mathsf{M}\)과 같은 랭크를 갖는 모든 오더에도 유지된다. 더 낮은 랭크의 오더에 대해서는 그 분수체로 넘어간다. 이 체는 \(\mK\)와 \(\mL\)의 중간체로서 역시 제1종이고, \(\mK\)에 대한 그 차수는 오더의 선형 랭크와 같으므로 위 논의가 유지된다. 끝으로 \(\mS\)와 다른 \(\mS\)의 오더는 결코 정수적으로 닫혀 있지 않다. 모든 오더가 \(\mR\)도 포함하므로, 정수적으로 닫혀 있다면 \(\mS\)까지 포함해야 하기 때문이다.

그러므로 수체와 함수체를 이 체계 아래 넣으려면 기초영역, 곧 정수환, 한 부정원의 다항식환, 여러 부정원의 다항식환의 함수영역에서 공리 I--V가 성립함을 보이면 충분하다.

3. \emph{유리정수환 \(\mR\), 또는 한 체를 계수로 하는 한 부정원의 다항식환 \(\mR\)에는 공리 I--V가 성립한다.} I과 II는 각 비영수 또는 한 부정원의 각 비영다항식을 법으로 한 잉여류가 각각 유한개, 또는 선형독립인 것이 유한개라는 사실에서 나온다. 또는 모든 아이디얼이 주아이디얼이라는 사실에서도 나온다. 이 사실로부터 아이디얼이 소아이디얼의 거듭제곱곱으로 유일하게 분해되고, 따라서 한 비영아이디얼을 나누는 아이디얼은 유한개뿐이라는 것이 따른다. 공리 III, IV는 분명 성립하고, IV는 상등의 정의에서 따른다. 정수적 닫힘 V는 다음에서 나온다. \(\mR\)의 원소들은 단위원, 곧 1의 약수들을 제외하면 기약원소들로 유일하게 분해된다. 따라서 \(\mR\)에 들어 있지 않은 분수체의 모든 원소를 \(\mR\)의 서로소인 두 원소의 몫으로 나타낼 수 있고, 그 분자의 어떤 거듭제곱도 분모로 나누어지지 않는다. 그러므로 그러한 원소는 \(\mR\)에 대하여 정수적임을 특징짓는 방정식을 만족할 수 없다.

4. \emph{이제 \(\mR\)을 한 체의 계수 또는 유리정수 계수를 갖는 여러 부정원의 모든 다항식이 이루는 환이라 하자.} 이 경우에도 원소가 단위원을 제외하고 기약원소들로 유일하게 분해되므로, 3에서와 같이 공리 III, IV, V가 성립한다. 약수사슬정리 I은 아이디얼기저의 존재에 관한 힐베르트 정리의 직접적인 결과이지만, 배수사슬정리 II는 여기서 성립하지 않는다.

그러나 최고차원 아이디얼로 제한하는 것과 같은 \emph{함수영역}으로 넘어가면 공리 II까지 성립하게 할 수 있다. 그러면 I도 힐베르트 정리를 쓰지 않고 3에서와 같이 증명할 수 있다. 실제로 \(\mR\)에 한 부정원 \(u\)를 첨가하여, \(\mR\)의 원소를 계수로 하는 \(u\)의 모든 다항식이 이루는 환 \(\mR^*\)을 고찰하고 상등을 계수의 상등으로 정의한다. 통상대로 \(\mR^*\)의 한 다항식의 계수들의 최대공약원 --- 여기서는 아이디얼 약수가 아니라 \(\mR\)의 다항식이라는 뜻의 약수이다 --- 이 \(\mR\)의 단위원이면 그 다항식을 \(\mR\)에 관하여 원시적이라 한다. 그러면 \(\mR^*\)의 모든 다항식은 \(\mR\)의 한 원소와 \(\mR^*\)의 한 원시다항식의 곱이고, \(\mR^*\)의 원시다항식들의 곱은 다시 원시적이다.

\emph{\(\mR^*\)의 분수체 원소 가운데 분모가 \(\mR^*\)의 원시다항식인 것들의 전체는 하나의 환, 곧 \(\mR\)의 함수영역 \(\mF\)를 이룬다.} 함수영역 \(\mF\)에서는 분자와 분모가 모두 \(\mR^*\)의 원시다항식인 원소들이 정확히 단위원이다. 따라서 \(\mF\)의 모든 원소는 \(\mR\)의 한 원소와 \(\mF\)의 한 단위원의 곱으로 유일하게 표현된다. 이로써 \(\mR\) 원소들의 분해정리가 \(\mF\)의 원소들로 옮겨지고, \(\mF\)에는 공리 III, IV와 더불어 3에서처럼 V도 성립한다.

또한 \(\mF\)의 모든 아이디얼은 주아이디얼이다. 실제로 한 아이디얼은 \(\mF\)의 임의의 한 원소와 더불어, 여기에 한 단위원을 곱하여 얻는 \(\mR\)의 원소도 포함한다. 그리고 \(\mR\)의 임의의 두 원소와 더불어 그 최대공약원도 포함한다. 왜냐하면 아이디얼이 \(f(x)=t(x)\overline f(x)\)와 \(g(x)=t(x)\overline g(x)\)를 포함하면 \(t(x)(\overline f(x)+u\overline g(x))\)도 포함하고, \(\overline f,\overline g\)가 서로소여서 그 선형결합이 단위원이 되면 \(t(x)\)도 포함하기 때문이다. 따라서 아이디얼의 임의의 원소 \(f(x)\)에서 시작하여, 아이디얼 안에 \(f(x)\)로 나누어지지 않는 \(g(x)\)가 있으면 최대공약원 \(t(x)\)도 아이디얼에 속하며 \(f(x)\)의 진약수가 된다. 이 과정을 유한번 되풀이하면 아이디얼의 한 기저원소에 이르고, 따라서 그 아이디얼은 주아이디얼이다. 그러므로 \(\mF\)에서 아이디얼은 \(\mR\) 원소들의 분해에 대응하여 소아이디얼들의 거듭제곱곱으로 유일하게 분해되고, 이로부터 3에서와 같이 공리 I, II가 성립한다. 여기서 고려하는 \(\mF\)의 분수체의 확대체들은 \(\mS\)의 한 원소를 첨가하여 생성할 수 있는 것들뿐이다\footnote{이 확대체들에서 \(\mF\)에 대하여 정수적인 원소의 계는 \(\mR\)에서 \(\mF\)로 넘어간 것과 마찬가지로 \(\mS\)에서 대응하는 함수영역으로 넘어가도 얻어진다. 예컨대 J. König, \emph{Algebraische Größen}, Leipzig 1903, 468/69쪽을 보라.}. 1, 2를 함께 고려하면 수체와 함수체에서 공리 I--V가 성립함을 알았다.

\subsection*{§ 4. 동형정리. 직접합}

이하에서는 환 또는 모듈의 성질만을 가정하고 다른 공리는 전제하지 않는다.

1. \(\mathsf{M}\)과 \(\overline{\mathsf{M}}\)이 \(\mR\)에 관한 모듈영역(§ 2)일 때, \(M\)의 각 원소에 \(\overline M\)의 정확히 한 원소가 대응하여 \(\overline M\)을 모두 채우고\footnote{\(\overline M\)에서 통용되는 상등의 정의에 따른 말이다. 각주 9를 보라.}, 이 대응에서 차와 같은 \(\mR\) 원소에 의한 곱셈이 보존되면 \(M\)이 \(\overline M\)에 준동형이라고, 더 정확히는 모듈준동형이라고 하여 \(\mathsf{M}\sim\overline{\overline{\mathsf{M}}}\)으로 쓴다. 즉 \(\beta\sim\overline\beta\), \(\gamma\sim\overline\gamma\)이면 언제나 \((\beta-\gamma)\sim(\overline\beta-\overline\gamma)\), \(r\beta\sim r\overline\beta\)이다.

그러면 \(M\) 안의 \(\mR\)-모듈 \(\mB\)에는 \(\overline M\) 안의 \(\mR\)-모듈 \(\overline{\mB}\)가 준동형으로 대응한다. 또한 \(\overline{\mathsf{M}}\) 안의 한 \(\mR\)-모듈 \(\overline{\mC}\)의 원소들에 대응하는 \(M\)의 모든 원소가 이루는 계 \(\mC^*\)는 \(\mathsf{M}\) 안의 \(\mR\)-모듈이고, \(\overline{\mC}\)가 유일하게 결정하는 대응 모듈로서 \(\mC^*\sim\overline{\mC}\)이다. 특히 \(\mA\)가 \(\overline{\mathsf{M}}\)의 영원소에 대응하는 모듈이면 \(\mC^*\)는 \(\mA\)의 약수이고, \(\mA\)를 법으로 서로 합동인 \(\mathsf{M}\)의 원소들의 류들로 나뉘며, 이 류들은 \(\overline{\mC}\)의 원소들과 일대일로 대응한다. \(M\)의 \(\mB\)에서 \(\overline M\)의 \(\overline{\mB}\)로 갔다가 다시 \(\mB^*\)로 돌아오면 \(\mB^*=(\mB,\mA)\), 곧 \(\mB\)와 \(\mA\)의 최대공약모듈이 된다. 따라서 \(\mB\)가 \(\mA\)의 약수이면 \(\mB^*=\mB\)이다.

\(\mathsf{M}\)과 \(\overline{\mathsf{M}}\)의 원소 대응이 양방향으로 유일하면 두 영역을 동형, 더 정확히는 모듈동형이라 하여 \(\mathsf{M}\simeq\overline{\mathsf{M}}\)으로 쓴다.

\(\mA\)가 \(M\) 안의 임의의 \(\mR\)-모듈일 때, \(\mA\)에 따른 합동을 새 상등관계로 삼으면 \(\mathsf{M}\)과 준동형인 모듈영역 \(\overline{\mathsf{M}}\), 곧 잉여류모듈 \(\mathsf{M}\mid\mA\)가 생긴다. 이때 \(M\)의 각 원소에는 \(\overline M\)에서 그와 같은 모든 원소만이 대응한다. \(\overline M\)에서 상등의 정의로부터 동일성으로 넘어간다는 것은 \(\overline M\)에서 같은 모든 원소를 하나의 류, 곧 잉여류로 묶고 이 잉여류들을 \(\overline M\)의 새 원소로 본다는 뜻이다\footnote{각주 9를 보라.}.

\emph{잉여류모듈로 넘어가면 모든 준동형이 생성된다.} 실제로 \(\mathsf{M}\sim\overline{\overline{\mathsf{M}}}\)이고 \(\mA\)가 \(\overline{\overline{\mathsf{M}}}\)의 영원소에 대응하는 모듈이면, 위에서 보았듯이 \(\overline{\overline{\mathsf{M}}}\)은 잉여류모듈 \(\mathsf{M}\mid\mA\)와 동형이다.

2. \textbf{제1동형정리.} \emph{\(\overline{\mathsf{M}}\)이 잉여류모듈 \(\mathsf{M}\mid\mA\)이고 \(\mC\)가 \(\mA\)의 약수이면
\[
        \overline M\mid\overline{\mC}\simeq M\mid\mC
\]
이다.}
실제로 \(\mA\)에 따른 합동은 \(\mC\)에 따른 합동이기도 하다. 따라서 \(\mA\)에 따라 같은 원소들은 \(\mC\)에 따라서도 같고, 먼저 원소들을 \(\mA\)에 따라 같게 놓아 \(\overline{\mathsf{M}}\)으로 넘어간 뒤 그 가운데 \(\mC\)에 따라 같은 것들을 묶어서 잉여류모듈 \(\mathsf{M}\mid\mC\), 곧 \(\mC\)에 따른 상등을 만들 수 있다. 이는 \(\overline{\mathsf{M}}\)에서 \(\overline{\mC}\)에 따라 같게 놓는 것, 곧 \(\overline{\mathsf{M}}\mid\overline{\mC}\)를 만드는 것과 같다.

\textbf{제2동형정리.} \emph{\(\mB,\mA\)가 \(M\)의 모듈이면
\[
        (\mB,\mA)\mid\mA\simeq \mB\mid[\mB,\mA]
\]
이다.}
실제로 1에 따라 \((\mB,\mA)\mid\mA=\overline{\mB}\)로 놓으면 \(\mB\)는 \(\overline{\mB}\)와 준동형이다. 이때 \([\mB,\mA]\)의 모든 원소와 오직 이 원소들만이 \(\overline{\mB}\)의 영원소에 대응하므로, 다시 1에 따라 위 동형을 얻는다.

3. \(\mR,\overline{\mR}\)이 가환환일 때, \(\mR\)의 각 원소에 \(\overline{\mR}\)의 정확히 한 원소가 대응하여 \(\overline{\mR}\)을 모두 채우고 이 대응에서 차와 곱이 보존되면 \(\mR\)이 \(\overline{\mR}\)에 준동형이라고, 더 정확히는 환준동형이라고 하여 \(\mR\sim\overline{\mR}\)으로 쓴다. 원소 대응이 양방향으로 유일하면 두 환을 동형, 더 정확히는 환동형이라 하여 \(\mR\simeq\overline{\mR}\)으로 쓴다.

1, 2의 모든 논의에서 모듈을 \(\mR\)의 아이디얼로 바꾸고\footnote{비가환환에서는 양쪽아이디얼을 바탕에 두어야 한다.}, 모듈준동형과 모듈동형을 환준동형과 환동형으로 바꾸어도 그대로 성립한다. 특히 \(\mc\)가 \(\ma\)의 약수이면, \(\ma\)에 따른 합동을 새 상등관계로 도입하여 잉여류환 \(\mc\mid\ma\)를 얻는다. 이제 잉여류의 곱도 정의된다는 점에 유의해야 한다.

\(\mc\)는 \(\mc\mid\ma\)와 준동형이고, 모든 준동형은 이 방식으로 생성된다. 2에서와 같은 동형정리들이 성립한다.

\textbf{제1동형정리.} \(\overline{\mR}\)이 잉여류환 \(\mR\mid\ma\)이고 \(\mc\)가 \(\ma\)의 약수이면 환동형의 뜻에서
\[
        \overline{\mR}\mid\overline{\mc}\simeq\mR\mid\mc
\]
이다.

\textbf{제2동형정리.} \(\mb,\ma\)가 \(\mR\)의 아이디얼이면 환동형
\[
        (\mb,\ma)\mid\ma\simeq\mb\mid[\mb,\ma]
\]
이 성립한다\footnote{군론에서 알려진 이 동형정리들은 모듈에 대해서는 조금 더 특수한 형태로 데데킨트에게서 처음 나타난다(제3판, 484쪽). 위의 아이디얼 형태는 Masazo Sono, \emph{On Congruences. I, II, III, IV}, Memoirs of the College of Science, Kyoto Imperial University, 2 (1917), 3 (1918, 1919), 특히 2, 215쪽에 있다. 각각 제2동형정리만 명시적으로 진술되어 있다.}.

4. 이제 서로소인 아이디얼들에 관한 계산법칙들을 모은다\footnote{데데킨트에게서 유래한 형태(제4판, § 178, III--VIII)로 제시한다. 근래의 서술에서 늘 나타나는 단위원을 이용한 계산은 훨씬 번거롭다.}. 가환환 \(\mR\)에 곱셈 단위원 \(e\)가 존재한다고 가정한다. \(\mo\)를 단위아이디얼이라 하면 모든 아이디얼 \(\ma\)에 대해 \(\ma=\mo\ma\)이다. 두 아이디얼 \(\ma,\mb\)의 최대공약아이디얼 \((\ma,\mb)\)가 \(\mo\)이면 둘을 서로소라 한다.

\emph{4\(\alpha\). \((\ma,\mb)=\mo\)이면 \((\ma,\mb\mc)=(\ma,\mc)\)이다.} 실제로
\[
(\ma,\mc)=(\ma,\mb)\cdot(\ma,\mc)=(\ma^2,\ma\mb,\ma\mc,\mb\mc)
\]
는 \((\ma,\mb\mc)\)로 나누어지고, 역으로도 그러하다. 이로부터 다음이 따른다.

\emph{4\(\beta\). \(\mc\mb\equiv0(\ma)\), \((\ma,\mb)=\mo\)이면 \(\mc\equiv0(\ma)\)이다.} 실제로 \(\mc\mb\equiv0(\ma)\)는 \((\ma,\mb\mc)=\ma\)와 같은 말이다. 그러므로 \((\ma,\mb)=\mo\)에 따라 \((\ma,\mc)=\ma\), 곧 \(\mc\equiv0(\ma)\)이기도 하다. 더 나아가 다음이 따른다.

\emph{4\(\gamma\). 아이디얼 \(\ma_1,\ldots,\ma_r\) 각각이 아이디얼 \(\mb_1,\ldots,\mb_s\) 각각과 서로소이면, \(\ma_i\)들의 곱은 \(\mb_i\)들의 곱과 서로소이다.} 실제로 \((\ma,\mb)=\mo\), \((\ma,\mc)=\mo\)에서 \((\ma,\mb\mc)=\mo\)가 나오고, 이를 유한번 되풀이하면 주장에 이른다.

4\(\alpha\), 4\(\gamma\)에서 다음을 얻는다.

\emph{4\(\delta\). 아이디얼 \(\mb_1,\ldots,\mb_r\)이 쌍마다 서로소, 곧 \(i\ne k\)일 때 \((\mb_i,\mb_k)=\mo\)이면, 그 최소공배수는 그 곱과 같다.} \((\ma,\mb)=\mo\), \(\mv=[\ma,\mb]\)라 하자. 그러면 \(\mv=\mo\mv=(\ma\mv,\mb\mv)\equiv0(\ma,\mb)\)이고, \(\ma\mb\equiv0(\mv)\)이므로 \(\mv=\ma\mb\)이다. 4\(\gamma\)에 유의하며 이를 유한번 되풀이하면 주장에 이른다.

\emph{4\(\varepsilon\). 아이디얼 \(\mb_1,\ldots,\mb_r\)이 쌍마다 서로소이고
\[
\ma_i=[\mb_1,\ldots,\mb_{i-1},\mb_{i+1},\ldots,\mb_r]
      =\mb_1\ldots\mb_{i-1}\mb_{i+1}\ldots\mb_r
\]
로 놓으면
\[
(\ma_1,\ldots,\ma_{i-1},\ma_{i+1},\ldots,\ma_r)=\mb_i,
\qquad
(\ma_1,\ldots,\ma_r)=\mo
\]
이다.}
실제로 분배법칙에 따라 \((\ma_1,\ma_2)=\mb_3\ldots\mb_r(\mb_2,\mb_1)=\mb_3\ldots\mb_r\)이다. 따라서 \((\ma_1,\ma_2,\ma_3)=\mb_4\ldots\mb_r(\mb_3,\mb_1\mb_2)=\mb_4\ldots\mb_r\)이고, 번호를 알맞게 매겨 이를 유한번 되풀이하면 주장에 이른다.

\(\mm=[\mb_1,\ldots,\mb_r]\)이라는 표현에서 \(\ma_i\)를 \(\mb_i\)의 보충부분이라 한다.

5. 다시 가환환 \(\mR\)에 곱셈 단위원이 존재한다고 가정한다. \(\mR\)의 모든 원소 \(c\)가 \(a_i\in\ma_i\)인 꼴
\[
        c=a_1+a_2+\ldots+a_r
\]
로 정확히 한 가지 방식으로 표현되면, 환 \(\mR\)을 아이디얼 \(\ma_1,\ldots,\ma_r\)의 \emph{직접합}이라 하고 \(\mR=\ma_1+\ma_2+\ldots+\ma_r\)로 쓴다.

\emph{영아이디얼이 쌍마다 서로소인 아이디얼 \(\mb_1,\ldots,\mb_r\)의 최소공배수이고 \(\ma_i\)가 \(\mb_i\)의 보충부분이면, \(\mR\)은 아이디얼 \(\ma_1,\ldots,\ma_r\)의 직접합이다.} 실제로 \(\mo=(\ma_1,\ldots,\ma_r)\)이므로 \(\mR\)의 모든 원소는 \(\ma_i\)들에 의한 덧셈표현을 적어도 하나 가진다. \([\ma_i,\mb_i]=(0)\)이므로 이 표현은 유일하다. 실제로 \(0=a_1+a_2+\ldots+a_r=a_i+b_i\)이면 \(a_i=0\)이다. 제2동형정리에 따라 \(\ma_i\)는 \emph{잉여류환 \(\mR\mid\mb_i\)와 동형}이다. \glqq{}직교관계\grqq{} \(\ma_i\ma_k=(0)\) (\(i\ne k\))와 \(\ma_i^2=\ma_i\)가 성립한다. 마지막 식은 \(\ma_i=\mo\ma_i\)에서 나온다.

\emph{역으로 \(\mR\)의 직접합 표현
\[
        \mR=\ma_1+\ma_2+\ldots+\ma_r
\]
이 존재하고
\[
\mb_i=(\ma_1,\ldots,\ma_{i-1},\ma_{i+1},\ldots,\ma_r)
     =\ma_1+\ldots+\ma_{i-1}+\ma_{i+1}+\ldots+\ma_r
\]
로 놓으면, \(\mb_i\)들은 쌍마다 서로소이고 그 최소공배수는 영아이디얼이다.}
실제로 \(\mo=(\ma_i,\mb_i)\)에서 \(k\ne i\)일 때 \(\mb_k\)가 \(\ma_i\)의 약수이므로 \(\mo=(\mb_k,\mb_i)\)가 따른다. 또 \(c\)가 모든 \(\mb_i\)로 나누어진다고 하자. 그러면
\[
 c=a_2^{(1)}+\ldots+a_r^{(1)}
  =a_1^{(2)}+a_3^{(2)}+\ldots+a_r^{(2)}
  =\ldots
  =a_1^{(r)}+\ldots+a_{r-1}^{(r)},
\]
이고 각 \(a_i^{(\lambda)}\equiv0(\ma_i)\)이다. 표현의 유일성에 따라 매번 한 성분이 0이 되는 데서 모든 \(\lambda\)에 대해 \(0=a_1^{(\lambda)},\ldots,0=a_r^{(\lambda)}\)이고, 따라서 \(c=0\)이다\footnote{5의 정리들은 직접합과 직접교집합의 관계에 관한 일반정리의 특수한 경우이다. H. Prüfer, \emph{Theorie der Abelschen Gruppen I}, Math. Zeitschr. 20 (1924), 165--187쪽, § 6을 보라. 그곳의 논의는 군 개념을 모듈과 아이디얼이 특수한 경우로 들어갈 만큼 일반적으로 잡아도 그대로 유지된다.}.

\begin{center}
{\bfseries\large § 5.\par}
\vspace{0.35em}
{\bfseries\large 소아이디얼과 으뜸아이디얼.\par}
\end{center}

다시 \(\mR\)을 가환환이라 하되, 단위원의 존재를 포함하여 다른 어떤 공리도 가정하지 않는다. 소아이디얼과 으뜸아이디얼에 관한 기본사실들을 모은다\footnote{\emph{아이디얼론} § 4에서는 약수사슬정리라는 공리 I을 가정하였다. 그래서 소아이디얼과 으뜸아이디얼의 어느 성질이 이 공리와 무관한지가 뚜렷하게 드러나지 않는다.}.

1. \(\mR\)의 한 아이디얼 \(\mpideal\)에 따른 잉여류환 \(\mR\mid\mpideal\)이 영인자가 없는 환이면 \(\mpideal\)을 \emph{약한 소아이디얼}이라 한다. 즉 \(a\not\equiv0(\mpideal)\), \(b\not\equiv0(\mpideal)\)이면 언제나 \(ab\not\equiv0(\mpideal)\)이다. \(\mR\)의 한 아이디얼 \(\mpideal^*\)에 따른 잉여류환 \(\mR\mid\mpideal^*\)에 \emph{영인자아이디얼}이 없으면 \(\mpideal^*\)를 \emph{강한 소아이디얼}이라 한다. 즉 \(\ma\not\equiv0(\mpideal^*)\), \(\mb\not\equiv0(\mpideal^*)\)이면 언제나 \(\ma\mb\not\equiv0(\mpideal^*)\)이다.

\(\ma,\mb\)를 주아이디얼로 특수화하면 \emph{모든 강한 소아이디얼은 약한 소아이디얼이기도 하다}. 역으로 \emph{모든 약한 소아이디얼도 강한 소아이디얼}이므로, 단순히 \emph{소아이디얼}이라 말할 수 있다. 실제로 \(\mpideal\)이 약한 소아이디얼이고 \(\ma\not\equiv0(\mpideal)\), \(\mb\not\equiv0(\mpideal)\)이라 하자. \(\ma,\mb\)에서 각각 \(a\not\equiv0(\mpideal)\), \(b\not\equiv0(\mpideal)\)인 원소 \(a,b\)를 택할 수 있다. 그러면 \(ab\not\equiv0(\mpideal)\), 따라서 \(\ma\mb\not\equiv0(\mpideal)\)이므로 \(\mpideal\)은 강한 소아이디얼이기도 하다.

2. 잉여류환 \(\mR\mid\mq\)에서 모든 영인자의 어떤 거듭제곱이 0이 되면 \(\mR\)의 아이디얼 \(\mq\)를 \emph{약한 으뜸아이디얼}이라 한다. 즉 \(a\not\equiv0(\mq)\)이고 모든 \(\varkappa\)에 대해 \(b^\varkappa\not\equiv0(\mq)\)이면 언제나 \(ab\not\equiv0(\mq)\)이다. \(\mR\mid\mq\)에서 영인자가 되는 \(\mR\)의 모든 원소의 계 \(\mpideal\)은 아이디얼을 이루며, \(\mq\)의 약수인 소아이디얼이다. 이를 \emph{대응 소아이디얼}이라 한다. 실제로 \(a\)와 더불어 \(ra\)도 영인자이고, \(a,b\)와 더불어 \(a-b\)도 영인자이다. 마지막 말은 \(a^\varkappa\), \(b^\lambda\)와 더불어 \((a-b)^{\varkappa+\lambda}\)도 언제나 \(\mq\)로 나누어지기 때문이다. 또한 \(a\)가 영인자가 아니면 정의에 따라 \(a\)의 어떤 거듭제곱도 영인자가 아니다. 따라서 \(a\not\equiv0(\mpideal)\), \(b\not\equiv0(\mpideal)\)이면 \(a^\varkappa b^\varkappa=(ab)^\varkappa\not\equiv0(\mq)\), 따라서 \(ab\not\equiv0(\mpideal)\)이다.

잉여류환 \(\mR\mid\mq^*\)에서 모든 영인자아이디얼의 어떤 거듭제곱이 0이 되면 \(\mR\)의 아이디얼 \(\mq^*\)를 \emph{강한 으뜸아이디얼}이라 한다. 즉 \(\ma\not\equiv0(\mq^*)\)이고 모든 \(\varkappa\)에 대해 \(\mb^\varkappa\not\equiv0(\mq^*)\)이면 언제나 \(\ma\mb\not\equiv0(\mq^*)\)이다. 약한 으뜸아이디얼의 경우와 같이 다음을 보인다. \(\mR\mid\mq^*\)에서 영인자아이디얼이 되는 \(\mR\)의 모든 아이디얼의 최대공약아이디얼은 \(\mq^*\)의 약수인 한 소아이디얼, 곧 \emph{대응 소아이디얼}을 이룬다. 역으로 같은 대응 소아이디얼 \(\mpideal\)을 갖는 모든 으뜸아이디얼을 \glqq{}\(\mpideal\)에 대응한다\grqq{}고 한다.

\(\ma,\mb\)를 주아이디얼로 특수화하면 \emph{모든 강한 으뜸아이디얼은 약한 으뜸아이디얼이기도 하다}. \emph{그러나 일반적으로 그 역은 성립하지 않는다}\footnote{R. Hölzer가 알려 준 다음 예가 이를 보인다. \(\mR\)을 한 체를 계수로 하는 가산개의 부정원 \(x_i\)의 다항식영역이라 하고
\[
\mq=(x_1^2,x_2^3,\ldots,x_\nu^{\nu+1},\ldots,x_i x_k\,[i\ne k],\ldots)
\]
라 하자. 잉여류환의 영인자들은 정확히 \(\mpideal=(x_1,x_2,\ldots,x_\nu,\ldots)\)로 나누어지는 다항식들이며, 이 영인자들의 어떤 거듭제곱은 각각 \(\mq\)로 나누어진다. 따라서 \(\mq\)는 \(\mpideal\)을 대응 소아이디얼로 갖는 약한 으뜸아이디얼이다. 그러나 \(\ma=(x_1,x_3,x_5,\ldots,x_{2\nu+1},\ldots)\), \(\mb=(x_2,x_4,\ldots,x_{2\nu},\ldots)\)로 놓으면 \(\ma\mb\equiv0(\mq)\)이지만 \(\ma,\mb\)의 어떤 거듭제곱도 \(\mq\)로 나누어지지 않으므로 \(\mq\)는 강한 으뜸아이디얼이 아니다.}. \emph{그러나 \(\mR\)에 약수사슬정리라는 공리 I을 가정하면 모든 약한 으뜸아이디얼은 강한 으뜸아이디얼이기도 하다.} 따라서 단순히 \emph{으뜸아이디얼}이라 말할 수 있다. 실제로 \(\mq\)가 약한 으뜸아이디얼이고 \(\ma\not\equiv0(\mq)\), 또 모든 \(\varkappa\)에 대해 \(\mb^\varkappa\not\equiv0(\mq)\)이라 하자. 그러면 \(a\not\equiv0(\mq)\)인 \(a\in\ma\)와 모든 \(\varkappa\)에 대해 \(b^\varkappa\not\equiv0(\mq)\)인 \(b\in\mb\)가 존재한다. 따라서 \(ab\not\equiv0(\mq)\), 곧 \(\ma\mb\not\equiv0(\mq)\)이다. 실제로 \(\mb\)의 각 \(b\)마다 \(b^\lambda\equiv0(\mq)\)이게 하는 지수 \(\lambda\)가 존재한다면, 유한개의 기저원소의 지수를 \(\lambda_1,\ldots,\lambda_r\)라 할 때 \(\mb^{\lambda_1+\ldots+\lambda_r}\equiv0(\mq)\)일 것이다\footnote{공리 I로부터 유한 아이디얼기저를 끌어내려면 \(\mR\)에 고정된 정렬순서를 바탕에 두어야 한다(§ 6 참조).}. 이 마지막 논의는 다음도 보인다. \(\mpideal\)이 대응 소아이디얼이면 \(\mpideal^\varrho\equiv0(\mq)\)이게 하는 최소 지수 \(\varrho\)가 존재하며, 이를 \(\mq\)의 \emph{지수}라 한다.

3. 이하에서는 다시 약수사슬정리를 \emph{가정하지 않는다}. 유한개의 아이디얼의 최소공배수 \([\ma_1,\ldots,\ma_r]\)에서 어떤 \(\ma_i\)도 나머지 아이디얼들의 최소공배수로 나누어지지 않아 어느 \(\ma_i\)도 생략할 수 없으면 이를 \emph{최단표현}이라 한다.

\emph{같은 소아이디얼 \(\mpideal\)에 대응하는 유한개의 약한 또는 강한 으뜸아이디얼의 최소공배수는 다시 \(\mpideal\)을 대응 소아이디얼로 갖는 약한 또는 강한 으뜸아이디얼이다. 서로 다른 소아이디얼들에 대응하는 유한개의 약한 또는 강한 으뜸아이디얼에 의한 최단표현은 약한 또는 강한 으뜸아이디얼이 아니다.}

\(\mf=[\mq_1,\ldots,\mq_r]\)이고 \(\mpideal\)이 약한 으뜸아이디얼 \(\mq_i\)들의 대응 소아이디얼이라 하자. 그러면 \(\mpideal\)은 어떤 거듭제곱이 \(\mf\)로 나누어지는 모든 원소의 전체이기도 하다. \(a\not\equiv0(\mf)\)이고 모든 \(\varkappa\)에 대해 \(b^\varkappa\not\equiv0(\mf)\)이면 \(b\not\equiv0(\mpideal)\)이고, 적어도 한 지표 \(i\)에 대해 \(a\not\equiv0(\mq_i)\)이다. 따라서 \(ab\not\equiv0(\mq_i)\), 곧 \(ab\not\equiv0(\mf)\)이다. 줄곧 원소를 아이디얼로 바꾸면 더는 \(\mb\not\equiv0(\mpideal)\)을 결론낼 수 없다.

이제 \(i\ne k\)일 때 \(\mpideal_i\ne\mpideal_k\)인 약한 으뜸아이디얼들의 최단표현
\[
        \mm=[\mq_1,\ldots,\mq_r],\qquad r\geqq2
\]
를 생각하자. 대응 소아이디얼들 가운데 적어도 하나, 이를테면 \(\mpideal_1\)은 나머지 어느 것에도 들어가지 않는다. \(\mpideal\)들의 모든 진배수사슬은 늦어도 \(r\)단계에서 멈추기 때문이다. 따라서 \(a_i^{\varrho_i}\equiv0(\mq_i)\)이고 \(a_2\not\equiv0(\mpideal_1),\ldots,a_r\not\equiv0(\mpideal_1)\)이게 하는 \(a_i\in\mpideal_i\)를 택할 수 있다. 또 \(q_1\not\equiv0(\mm)\)인 \(q_1\in\mq_1\)을 택하면
\[
q_1a_2^{\varrho_2}\ldots a_r^{\varrho_r}\equiv0(\mm)
\]
이지만 \((a_2^{\varrho_2}\ldots a_r^{\varrho_r})\)의 어떤 거듭제곱도 \(\mm\)으로 나누어지지 않는다. 그렇지 않으면 \(\mpideal_1\)로 나누어져야 하기 때문이다. 따라서 \(\mm\)은 약한 으뜸아이디얼이 아니다. 줄곧 원소를 아이디얼로, 즉 \(q_1\)을 \(\mq_1\)로, \(a_i\)를 \(\ma_i\not\equiv0(\mpideal_1)\)이면서 \(\ma_i^{\varrho_i}\equiv0(\mq_i)\)인 \(\ma_i\)로 바꾸면 강한 으뜸아이디얼에 대한 증명이 나온다\footnote{약수사슬정리를 불필요하게 만든 이 원래 증명의 수정은 B. L. van der Waerden에게서 배웠다.}.

4. \emph{모듈 몫과 아이디얼 몫.} \(\mT\)를 \(\mR\)의 확대환이라 하고, \(\mA,\mB,\mC,\ldots\)를 \(\mT\) 안의 \(\mR\)-모듈이라 하자. \emph{몫} \(\mC=\mA:\mB\)는 다음 조건을 만족하는 모듈로 정의한다. \(\mC\mB\equiv0(\mA)\)이고, \(\mD\mB\equiv0(\mA)\)이면 \(\mD\equiv0(\mC)\)이다. 따라서 \(\mC\)는 \(\mC\mB\equiv0(\mA)\)이게 하는 \glqq{}가장 작은\grqq{} 모듈이다. 이 조건으로 몫은 \(\mD\mB\equiv0(\mA)\)인 모든 모듈 \(\mD\)의 최대공약모듈로서 유일하게 결정된다. 영모듈은 확실히 그러한 모듈이므로 이 \(\mD\)들은 존재한다.

정의로부터 다음이 따른다. \(\mB\)가 \(\overline{\mB}\)의 배수이면 \(\mA:\overline{\mB}\)는 \(\mA:\mB\)의 배수이다. \(\mA\)가 \(\overline{\mA}\)의 배수이면 \(\mA:\mB\)는 \(\overline{\mA}:\mB\)의 배수이다. 또한
\[
        \mA:\mB\mC=(\mA:\mB):\mC=(\mA:\mC):\mB
\]
이다. 확대환 \(\mT\)가 \(\mR\)과 같으면 \emph{모듈 몫}은 \emph{아이디얼 몫} \(\ma:\mb\)이 되고, 위 계산법칙들이 그대로 성립한다. \(\ma\mb\equiv0(\ma)\)이므로 이 경우에는 \(\ma\equiv0(\ma:\mb)\)도 따른다.

5. \emph{\(\ma:\mb=\ma\)이면 \(\mb\)가 \(\ma\)에 대하여 상대소라고 한다.} 따라서 \(\mb\)가 \(\ma\)에 대하여 상대소일 필요충분조건은 \(\md\mb\equiv0(\ma)\)이면 언제나 \(\md\equiv0(\ma)\)인 것이다. 4의 계산법칙에서 다음이 따른다. \emph{\(\mb,\mc\)가 \(\ma\)에 대하여 상대소이면 그 곱과 최소공배수도 \(\ma\)에 대하여 상대소이다.} \(\mb\)가 \(\ma\)에 대하여 상대소이고 \(\ma\)도 \(\mb\)에 대하여 상대소이면 \(\ma,\mb\)를 서로 상대소라고 한다. § 4, 4\(\beta\)에서 서로소인 아이디얼들은 서로 상대소이기도 하다는 것이 따른다.

\setcounter{footnote}{25}

\begin{center}
{\bfseries\large § 6.\par}
\vspace{0.35em}
{\bfseries\large 약수사슬정리를 가정한 아이디얼론.\par}
\end{center}

약수사슬정리라는 \emph{공리 I이 성립}하는 가환환 \(\mR\)을 바탕에 두고, 이하에서는 언제나 \(\mR\)의 \emph{모든 원소에 대한 고정된 정렬순서}도 바탕에 둔다. 그러면 \(\mR\)의 \emph{모든 아이디얼에도 정렬순서}가 주어진다. 실제로 이 두 가정에서 \(\mR\)의 모든 아이디얼이 유한개의 원소로 이루어진 아이디얼기저를 갖는다는 것이 곧바로 나온다. 이제 \(\mR\) 원소의 정렬순서에서 \(\mR\)의 모든 유한 부분집합의 준사전식 순서로 넘어가고\footnote{그 뜻은 다음과 같다. 유한 부분집합 \(\mathfrak A_n\)의 원소 \(a_1,\ldots,a_n\)을 지표 순서가 \(\mR\)의 정렬순서와 일치하도록 나타낸다. 원소가 \(n\)개인 모든 부분집합 \(\mathfrak A_n\)은 \(m>n\)개인 모든 \(\mathfrak A_m\)보다 앞선다. 원소 수가 같은 \(\mathfrak A_n,\mathfrak B_n\)에 대해서는 대응하는 \(b_i\)와 다른 첫 원소 \(a_i\)가 \(\mR\)의 정렬순서에서 \(b_i\)보다 앞서면 \(\mathfrak A_n\)이 \(\mathfrak B_n\)보다 앞선다고 한다. 이로써 유한 부분집합들의 모든 계에는 첫 원소 \(\mathfrak A_n\)이 대응한다. 따라서 \(\mR\)의 정렬순서가 주어지면 다른 선택공리는 필요 없다. 특히 대수적 수체의 경우처럼 \(\mR\)이 가산집합이면 선택공리를 전혀 바탕에 두지 않는다. 아이디얼론에서 선택공리가 하는 역할은 \emph{아이디얼론}, 각주 9에서 자세한 설명 없이 지적하였다.}, 각 아이디얼에는 가능한 여러 기저 가운데 유한 부분집합의 정렬순서에서 첫째인 것을 특수기저로 대응시키면, 유한 부분집합들과 더불어 아이디얼들도 정렬된다.

\textbf{정리 I.} \emph{약수사슬정리를 가정하면 \(\mR\)의 모든 아이디얼은 유한개의 기약아이디얼의 최소공배수로 표현된다. 여기서 기약아이디얼이란 두 진약수의 최소공배수로 표현되지 않는 아이디얼이다.}

증명을 위해 다음을 보이면 된다. 한 아이디얼 \(\mm\)에 정리 I이 성립하지 않으면, \(\mm\)에는 정리 I이 역시 성립하지 않는 진약수가 있다. 그러면 가정한 약수사슬정리에 어긋나게 유한 단계에서 멈추지 않는 약수사슬을 구성할 수 있다. 실제로 \(\mm\)은 가약이어야 한다. 그렇지 않으면 \(\mm=[\mm]\)이 요구되는 표현을 주기 때문이다. \(\mm=[\ma,\mb]\)이면 정리 I은 \(\ma,\mb\)에 동시에 성립할 수 없다. 그렇다면 \(\mm\)에도 대응하는 표현이 생기기 때문이다. 따라서 \(\mm\)에는 정리 I이 성립하지 않는 진약수들이 있고, 그 가운데 아이디얼의 정렬순서에서 첫째를 \(\ma_1\)이라 하자. 같은 방식으로 \(\ma_1\)의 진약수 \(\ma_2\), 일반적으로 \(\ma_i\)의 진약수 \(\ma_{i+1}\)를 구성하면, 잘 정의되면서 유한 단계에서 멈추지 않는 약수사슬을 얻어 가정에 모순된다\footnote{모든 모듈계에 약수사슬정리를 가정하면 정리 I은 모듈영역에도 그대로 성립한다. 정리 I은 순전히 집합론적 성격을 가지며, 아이디얼의 정렬순서를 증명에 쓰는 유일한 정리이기도 하다. 모든 연산과 무관하게 다음 집합론적 개념들만 문제된다. 한 집합 \(\mathfrak M\)에서 그 멱집합의 한 부분집합 \(\Sigma\), 곧 부분집합들의 한 계를 택한다. \(\Sigma\)가 정렬되어 있다고 하고, 그 원소를 \(\Sigma\)-집합이라 하자. \(\Sigma\)에 다음 사슬정리를 가정한다. \(\mathfrak A_\nu\)가 \(\mathfrak A_{\nu-1}\)의 진상위집합인 \(\Sigma\)-집합들의 모든 사슬 \(\mathfrak A_1,\mathfrak A_2,\ldots,\mathfrak A_\nu\ldots\)는 유한 단계에서 멈춘다. 한 \(\Sigma\)-집합 \(\mathfrak A\)가 모두 \(\mathfrak A\)의 진상위집합인 두 \(\Sigma\)-집합의 교집합이면 가약, 그렇지 않으면 기약이라 한다. 그러면 위 논의에 따라 모든 \(\Sigma\)-집합은 유한개의 기약 \(\Sigma\)-집합의 교집합으로 표현된다.}.

가정한 약수사슬정리에 따라 § 5, 2에서 단순히 으뜸아이디얼이라 말할 수 있다. 으뜸성과 기약성의 관계는 다음과 같다.

\textbf{정리 II.} \emph{약수사슬정리를 가정하면 모든 기약아이디얼은 으뜸이다. 다시 말해, 으뜸이 아닌 모든 아이디얼은 가약이다.}

잉여류환 \(\mR\mid\mm\)으로 넘어가도 준동형 때문에 약수사슬정리가 유지되고, \(\mm\)을 최소공배수로 나타내는 표현에는 영아이디얼의 표현이 대응하며, 제1동형정리(§ 4, 3)에 따라 \(\mm\)의 으뜸 약수에는 \(\mR\mid\mm\)의 으뜸아이디얼이 대응하고 그 역도 성립한다. 따라서 잉여류환에서 영아이디얼을 분해하는 문제로 제한할 수 있다. 이 환도 같은 일반성을 가지므로 처음부터 분해할 아이디얼을 \(\mR\)의 영아이디얼이라 해도 된다.

이제 \(\mR\)의 영아이디얼이 으뜸이 아니라고 하자. 그러면 \(\ma\ne(0)\), 모든 \(\varkappa\)에 대해 \(\mb^\varkappa\ne(0)\)이지만 \(\ma\mb=(0)\)인 아이디얼쌍 \(\ma,\mb\)가 적어도 하나 존재하며, \(\mb\)는 주아이디얼이라고 가정해도 된다. 가정에 따라 유한 단계에서 멈추는 아이디얼 몫의 약수사슬
\[
        \ma,\ \ma:\mb,\ldots,\ma:\mb^\nu,\ldots
\]
을 만든다. 예컨대 \(\mt=\ma:\mb^{m-1}=\ma:\mb^m\ldots\)가 이 사슬에서 그 뒤의 모든 아이디얼과 같다고 하자. 그러면 \(\mt=\mt:\mb\), 곧 \(\mb\)가 \(\mt\)에 대하여 상대소이다. 또한 \(\mt\)는 \(\ma\)의 약수이므로 영아이디얼과 다르고, 모든 지수 \(\varkappa\)에 대해 \(\mb^\varkappa\)도 영아이디얼과 다르다. 따라서 정리 II를 증명하려면
\[
        (0)=[\mt,\mb^{m+1}]
\]
을 보이면 충분하다. 정의에 따라
\[
        \mb^{m-1}\mt\equiv0(\ma),
        \qquad\text{따라서}\qquad
        \mb^m\mt=(0).
\]
그러므로
\[
        \mc=[\mt,\mb^{m+1}]\equiv0(\mb^m\mt)
\]
만 보이면 된다. 이는 \(\mb\)가 주아이디얼이고 \(\mt\)에 대하여 상대소라는 가정에서 따른다. \(c\in\mc\)는 \(\mb^{m+1}\)로 나누어지므로, \(n\)을 정수의 기호라 할 때
\[
        c=k b^{m+1}+n b^{m+1}=r b^m
\]
으로 표현된다. 여기서 \(r\)은 다시 \(\mR\)의 원소이다. \(c\equiv0(\mt)\)에서 \(r b^m\equiv0(\mt)\)이고, \(\mt:\mb=\mt\)이므로 \(r\equiv0(\mt)\)이다. 따라서 \(c\equiv0(\mt\mb^m)\)이고 정리 II가 증명되었다.

\textbf{정리 III.} \emph{약수사슬정리를 가정하면 모든 아이디얼은 서로 다른 소아이디얼들에 대응하는 유한개의 으뜸아이디얼, 곧 극대 으뜸성분들의 최소공배수로 된 최단표현을 갖는다\footnote{대응 소아이디얼들과 고립성분들은 유일하게 결정된다. \emph{아이디얼론}, 또는 W. Krull, \emph{Ein neuer Beweis für die Hauptsätze der allgemeinen Idealtheorie}, Math. Ann. 90 (1923), 55--64쪽을 보라. § 7에서는 그곳에 나타나는 단순한 특수한 경우에 유일성을 직접 증명한다. 그곳에서 유일하다고 판정되는 으뜸아이디얼들은 고립성분이다.}.}

그러한 표현을 얻으려면 정리 I에 따라 존재하는 기약아이디얼, 따라서 정리 II에 따라 으뜸아이디얼에 의한 표현을 가능한 대로 최단표현으로 바꾸기만 하면 된다. 같은 소아이디얼에 대응하는 으뜸아이디얼들을 묶으면 § 5, 3에 따라 요구되는 표현을 얻는다. 이 성분들 가운데 어떤 것들의 최소공배수도 더는 으뜸이 아니므로 으뜸성분들을 극대라고 부를 수 있다.

\begin{center}
{\bfseries\large § 7.\par}
\vspace{0.35em}
{\bfseries\large 이중사슬정리를 가정한 아이디얼론.\par}
\end{center}

지금까지 전개한 이론에 비해 생기는 단순화는 1, 2에서 제시할 보조정리들에 기초한다.

1. \emph{영인자가 없는 가환환에서 배수사슬정리가 성립하면 그 환은 체이기도 하다.} \(a\ne0\)일 때 방정식 \(ax=b\)가 언제나 환 안에서 해를 갖는다는 것을 보이면 된다. 해가 둘 이상일 수 없다는 것은 환에 영인자가 없다는 가정에서 따른다.

\(\ma\)를 \(a\)에서 생성되는 주아이디얼이라 하자. 가정에 따라 거듭제곱의 열은 유한 단계에서 멈춘다. 예컨대 \(\ma^m\)이 그 뒤의 모든 것과 같다고 하자. \(\mb\)를 \(b\)에서 생성되는 주아이디얼이라 하면
\[
\ma^m\mb=\ma^{m+1}\mb,
\qquad
\ma^{m+1}\mb=(a^{m+1}b).
\]
특히 원소 \(a^m b\)에 대해, \(n\)을 정수의 기호라 할 때
\[
        a^m b=k a^{m+1}b+n a^{m+1}b=a^{m+1}c
\]
인 표현이 나온다. 여기서 \(c\)는 다시 \(\mR\)의 원소이다. 가정에 따라 영인자가 없으므로 \(b=ac\)이고, 체의 성질이 증명되었다.

2. 1에서 곧바로 다음 보조정리를 얻는다. \emph{가환환에서 배수사슬정리가 성립하면 한 소아이디얼은 \(\mo\)와 다른 진약수를 갖지 않는다.}

마찬가지로 다음 덧붙임이 따른다. \emph{비영아이디얼을 법으로 한 배수사슬정리라는 공리 II만 성립하면, 모든 비영 소아이디얼은 \(\mo\)와 다른 진약수를 갖지 않는다.}

실제로 소아이디얼 \(\mpideal\)에 따른 잉여류환은 정의에 따라 영인자가 없는 환이고, 준동형 때문에 여기서도 배수사슬정리가 성립하므로 1에 따라 체이다. \(\mpideal\)의 약수들과 잉여류환의 아이디얼들이 일대일로 대응하므로 \(\mpideal\)은 \(\mo\)와 다른 진약수를 갖지 않는다\footnote{\(\mR\), 따라서 \(\mo\), 곧 \(\mR\)의 모든 원소로 이루어진 단위아이디얼에는 단위원이 존재하지 않아도 된다. 1에 따라 잉여류환에는 한 단위류가 존재하지만 그렇다. 이 종류의 가장 단순한 예는 --- 덧붙임의 더 약한 경우에 해당한다 --- 모든 짝수의 계이다.}.

\textbf{정리 IV.} \emph{가환환에서 이중사슬정리가 성립하면, 한 아이디얼의 모든 최단표현에서 단위아이디얼 \(\mo\)에 대응하지 않는 으뜸성분들은 유일하게 결정된다.}

\emph{덧붙임.} 공리 I, II를 가정하면 정리 IV는 모든 비영아이디얼에 성립한다.

\[
\mm=[\mq,\mq_1,\ldots,\mq_r]
   =[\overline{\mq},\overline{\mq}_1,\ldots,\overline{\mq}_{\overline r}]
\]
가 \(\mm\)의 극대 으뜸성분들에 의한 최단표현이고, \(\mo,\mpideal_1,\ldots,\mpideal_r\) 및 \(\mo,\overline{\mpideal}_1,\ldots,\overline{\mpideal}_{\overline r}\)가 각각 대응 소아이디얼이라 하자. \(\mo\)에 대응하는 으뜸아이디얼이 나타나지 않으면 \(\mq\), 또는 \(\overline{\mq}\)는 생략한다. 식
\[
        \mo^\varrho\mpideal_1^{\varrho_1}\ldots\mpideal_r^{\varrho_r}
        \equiv0(\mm)
\]
때문에 각 \(\overline{\mpideal}_i\)는 어떤 \(\mpideal_i\)에 들어가야 하고, 보조정리에 따라 그 아이디얼과 같다. 똑같이 각 \(\mpideal_k\)도 어떤 \(\overline{\mpideal}_k\)와 같아야 하므로 \(\mo\)와 다른 소아이디얼들은 일치하여 \(\mpideal_\lambda=\overline{\mpideal}_\lambda\)이다. 또한 \(\mq\mq_1\ldots\mq_r\equiv0(\overline{\mq}_\lambda)\)에서 \(\mq_\lambda\equiv0(\overline{\mq}_\lambda)\)가 나온다. 나머지 성분들은 \(\mpideal_\lambda\)로 나누어지지 않으므로 \(\overline{\mq}_\lambda\)에 대하여 상대소이기 때문이다. 마찬가지로 \(\overline{\mq}_\lambda\equiv0(\mq_\lambda)\)이므로 \(\mo\)에 대응하지 않는 으뜸성분들의 유일성이 증명되었다. 끝으로 \(\mq\)가 실제로 나타나면 최단표현 때문에 \(\overline{\mq}\)도 실제로 나타나야 하므로 대응 소아이디얼의 유일성도 증명된다. 이 유일성 증명은 \(\mo\)에 대응하는 으뜸성분을 포함하지 않는 모든 표현이 동시에 최단임도 보인다.

3. 이중사슬정리의 가정에 단위원의 존재를 더하면 정리 IV는 다음 정리 V로 강화된다.

\textbf{정리 V.} \emph{단위원을 갖는 가환환에서 이중사슬정리가 성립하면, 모든 아이디얼은 쌍마다 서로소인 유한개의 으뜸아이디얼의 곱으로 유일하게 표현된다.}

\emph{덧붙임.} 공리 I, II, III을 가정하면 정리 V는 모든 비영아이디얼에 성립한다.

단위원이 존재하므로 \(\mo^2=\mo\)이고, 따라서 \(\mo\)에 대응하는 모든 으뜸아이디얼은 \(\mo\)와 같다. 그러한 성분은 \(\mo\)와 다른 아이디얼의 최단표현에는 나타날 수 없다. 따라서 정리 IV에 따라 으뜸성분들이 유일하게 결정된다. 동시에 \(\mo\)를 생략하면 모든 표현이 최단표현이 된다. 또한 2의 보조정리에 따라 서로 다른 두 소아이디얼은 언제나 서로소이다. 따라서 § 4, 4의 계산법칙에 따라 으뜸성분들도 서로소이고, 최소공배수는 그 곱이 된다.

가정들로부터 다음 보조정리도 따른다.

\textbf{보조정리.} \emph{단위원과 이중사슬정리를 갖는 가환환에서, 단위아이디얼을 제외하면 한 아이디얼 \(\mm\)에 들어가지 않는 소아이디얼들이 정확히 \(\mm\)에 대하여 상대소인 소아이디얼들이다. 서로 상대소와 서로소라는 개념은 일치한다.}

\emph{덧붙임.} 공리 I, II, III을 가정하는 경우에는 \(\mm\)이 비영아이디얼이라고 해야 한다.

\(\mpideal\)이 \(\mm\)에 들어가지 않으면 \(\mm\)의 모든 대응 소아이디얼과 다르다. 따라서 2의 보조정리에 따라 이들 \(\mo\)와 다른 소아이디얼 가운데 어느 것으로도 나누어지지 않는다. 그러므로 \(\mpideal\)은 모든 으뜸성분에 대하여 상대소이고, 따라서 \(\mm\)에 대하여 상대소이다. 동시에 모든 으뜸성분과 서로소이고, 따라서 \(\mm\)과도 서로소이다.

\(\mpideal\ne\mo\)가 \(\mm\)에 들어가면 어떤 대응 소아이디얼과 같아야 한다. 이를테면 성분 \(\mq=\mq_1\)에 대응한다고 하자. 그러면 \(\mpideal^{\varrho-1}\equiv0(\mq:\mpideal)\)이지만 \(\mpideal^{\varrho-1}\not\equiv0(\mq)\)이므로 \(\mq:\mpideal\)은 \(\mq\)의 진약수이다. 동시에 \(\mq:\mpideal\)은 \(\mpideal^{\varrho-1}\)의 약수로서 \(\ne\mo\)와 다른 소아이디얼 하나로만 나누어지므로 으뜸이다. 유일분해에 따라 \(\mm:\mpideal\)로 나누어지는 곱 \((\mq:\mpideal)\cdot\mq_2\ldots\mq_r\)도 \(\mm\)의 진약수이다. 따라서 \(\mpideal\)은 \(\mm\)에 대하여 상대소가 아니다.

그러므로 아이디얼 \(\mt\ne\mo\)가 \(\mm\)에 대하여 상대소일 필요충분조건은 \(\mt\)의 대응 소아이디얼 가운데 어느 것도 \(\mm\)에 들어가지 않는 것이다. 이 경우 \(\mt\)는 \(\mm\)과 서로소이기도 하다. 역으로 서로소인 두 아이디얼은 언제나 서로 상대소이므로(§ 5, 5), 위 가정 아래에서는 두 개념이 일치한다.

\begin{center}
{\bfseries\large § 8.\par}
\vspace{0.35em}
{\bfseries\large 분수체 안에서 정수적으로 닫힌 경우의 아이디얼론.\par}
\end{center}

지금까지 전개한 아이디얼론에 마지막 두 공리를 더하여 통상적인 더 강한 이론으로 만든다.

1. \emph{보조정리.} \(\mR\)이 단위원을 갖고 영인자가 없는 가환환이며 \(\mR\)에서 약수사슬정리를 가정하면(공리 I, III, IV), \(\ma=\ma\mb\), \(\ma\ne(0)\)에서 언제나 \(\mb=\mo\)가 따른다. 따라서 영아이디얼 및 단위아이디얼과 다른 모든 아이디얼에 대해
\[
        \mc^\varrho\ne\mc^{\varrho+1}
\]
이다\footnote{반대로 같은 가정 아래에서도 \(\ma\mb=\ma\mc\)에서 \(\mb=\mc\)를 결론낼 수는 없다. 다음 예가 이를 보인다. \(\mR\)을 한 체를 계수로 하는 \(x,y\)의 다항식영역이라 하고
\[
\ma=(x,y),\qquad \mb=(x^2,xy,y^2),\qquad \mc=(x^2,y^2)
\]
라 하자. 실제로 \(\ma\mb=\ma\mc=\ma^3\)이지만 \(\mb\ne\mc\)이다.}.

증명은 \(\ma\mb\)를 \(\mb\)에 관한 유한 모듈로 보면 § 1, 3의 보조정리와 같이 나온다. 공리 I에 따라 존재하는 \(\ma\)의 아이디얼기저를 \(\alpha_1,\ldots,\alpha_n\)이라 하면, \(\ma=\ma\mb\)에서 연립방정식
\[
        \alpha_i=b_{i1}\alpha_1+\ldots+b_{in}\alpha_n,
        \qquad b_{ik}\equiv0(\mb),\quad i=1,\ldots,n
\]
을 얻는다. \(\mR\)에 영인자가 없으므로 여기서 \(\lvert b_{ik}-e_{ik}\rvert=0\)이 따른다. \(i\ne k\)이면 \(e_{ik}=0\), \(e_{ii}=e\)이다. 그런데
\[
        e^n+b_1e^{n-1}+\ldots+b_n=0,
        \qquad b_i\equiv0(\mb)
\]
에서 \(e\equiv0(\mb)\), 따라서 \(\mb=\mo\)가 따른다.

2. 이제 공리 I--IV에 분수체 안에서 정수적으로 닫혀 있다는 가정(공리 V)을 더하면 모든 비영 으뜸아이디얼이 그 대응 소아이디얼의 한 거듭제곱이 됨을 보이기만 하면 된다. 그러면 모든 비영아이디얼에 대한 표현
\[
        \mm=\mpideal_1^{\varrho_1}\ldots\mpideal_r^{\varrho_r}
\]
의 유일성은 정리 V와 1의 보조정리로 이미 증명되었다. 동시에 이 소아이디얼들은 단위아이디얼과 다른 진약수를 갖지 않는다는 것도 보였다. 지수가 2인 경우는 데데킨트의 따름정리 II(§ 1, 4)를 이용하여 증명하고, 일반적인 경우는 완전귀납법으로 증명한다.

\emph{보조정리.} \emph{공리 I--V를 가정하면 지수가 2인 으뜸아이디얼은 \(\mq=\mpideal^2\)뿐이다.}

이를 증명하려면
\[
        \mpideal^2\equiv0(\mq),\qquad
        \mq\equiv0(\mpideal),\qquad
        \mq\not\equiv0(\mpideal^2)
\]
에서 반드시 \(\mq=\mpideal\)임을 보이면 된다. 따라서
\[
        c\equiv0(\mq),\qquad
        c\equiv0(\mpideal),\qquad
        c\not\equiv0(\mpideal^2)
\]
라 하자. § 7, 3의 보조정리에서 \(\mo c:\mpideal\)이 \(\mo c\)의 진약수임이 따른다. 따라서
\[
        b\mpideal\equiv0(\mo c)\equiv0(\mq),
        \qquad b\not\equiv0(\mo c)
\]
이게 하는 원소 \(b\)가 존재한다. 그러므로 \(\gamma=b/c\)는 정수적이지 않다. 즉 분수체에는 속하지만 \(\mR\)에는 속하지 않는다.

\(b\not\equiv0(\mpideal)\)임을 보이면 된다. 그러면 \(\mq\)가 \(\mpideal\)에 대응하는 으뜸아이디얼이므로 \(\mpideal\equiv0(\mq)\)가 따른다.

데데킨트의 따름정리 II에 따라 \(\mR\)의 원소 \(m,n\)을 택하여
\[
        \gamma=b/c=m/n
\]
이고 \(m^2/n\)도 정수적이지 않게 할 수 있다. 따라서
\[
        bn=mc,\qquad
        m\not\equiv0(\mo n),\qquad
        m^2\not\equiv0(\mo n).
\]
\(b\mpideal\)에 \(m\)을 곱하면
\[
        mb\mpideal\equiv0(\mo mc),
        \qquad\text{따라서}\qquad
        mb\mpideal\equiv0(\mo nb),
        \qquad\text{그러므로}\qquad
        m\mpideal\equiv0(\mo n).
\]
마지막 결론은 환에 영인자가 없기 때문에 성립한다. 이로부터
\[
        m\not\equiv0(\mpideal)
\]
이 따른다. 실제로 \(m^2\not\equiv0(\mo n)\)이고 \(n\equiv0(\mpideal)\)이다. 마지막 관계는 § 7, 3의 보조정리에서 나온다. \(m\not\equiv0(\mo n)\)이므로 \(\mo n:\mpideal\)은 \(\mo n\)의 진약수이기 때문이다. 이제 네 관계
\[
        m\not\equiv0(\mpideal),\qquad
        n\equiv0(\mpideal),\qquad
        c\equiv0(\mpideal),\qquad
        \quad {}\not\equiv0(\mpideal^2)
        bn=mc
\]
에서 \(b\not\equiv0(\mpideal)\)이 나온다. 실제로 \(\mpideal^2\)가 으뜸이므로 우선 \(mc\not\equiv0(\mpideal^2)\)이고, \(bn=mc\)에서 \(b\not\equiv0(\mpideal)\)이 따른다. 이로써 보조정리가 증명되었다.

3. \emph{보조정리.} \emph{공리 I--V를 가정하면 지수가 \(\varrho\)인 으뜸아이디얼은 \(\mq=\mpideal^\varrho\)뿐이다.}

이 보조정리는 공리를 더 쓰지 않고 2의 보조정리에서 나온다\footnote{보조정리 2가 성립한다고 가정한 보조정리 3의 증명은 Masazo Sono, \emph{On Congruences II}, §§ 9, 10에 있다. 각주 19도 보라.}. 먼저 다음을 보여야 한다.
\[
\text{\(c\equiv0(\mpideal)\), \(c\not\equiv0(\mpideal^2)\)이면 모든 \(\sigma\)에 대해}\qquad
        \mpideal^\sigma=(\mo c^\sigma,\mpideal^{\sigma+1}).
\]
2의 보조정리에 따라 \(\mpideal=(\mo c,\mpideal^2)\)이다. 따라서
\[
        \mpideal^{\sigma-1}=(\mo c^{\sigma-1},\mpideal^\sigma)
\]
를 가정하고 양변에 \(\mpideal\)을 곱하면 분배법칙에 따라
\[
\begin{aligned}
        \mpideal^\sigma
        &=(\mpideal c^{\sigma-1},\mpideal^{\sigma+1})\\
        &=(\mo c^\sigma,\mpideal^2c^{\sigma-1},\mpideal^{\sigma+1})\\
        &=(\mo c^\sigma,\mpideal^{\sigma+1})
\end{aligned}
\]
이다.

증명할 보조정리는 다음 두 번째 형태로 바꿀 수 있다.
\[
\begin{gathered}
        \mq\equiv0(\mpideal^\sigma),\qquad
        \mq\not\equiv0(\mpideal^{\sigma+1}),\\
        \mpideal^{\sigma+\lambda}\equiv0(\mq),
        \qquad\lambda\geqq1
\end{gathered}
\]
이면
\[
        \mpideal^{\sigma+\lambda-1}\equiv0(\mq)
\]
이기도 하다. 실제로 \(\mq\equiv0(\mpideal)\)이고 \(\not\equiv0(\mpideal^{\varrho+1})\)이므로 모든 으뜸아이디얼에 그러한 지수 \(\sigma\geqq1\)가 존재한다. 여기서 \(\mpideal^\varrho\equiv0(\mq)\)이고 1의 보조정리에 따라 \(\mpideal^\varrho\ne\mpideal^{\varrho+1}\)이므로 \(\mq\not\equiv0(\mpideal^{\varrho+1})\)이다. 두 번째 형태를 유한번 적용하면 \(\mpideal^\sigma\equiv0(\mq)\), 따라서 \(\mq=\mpideal^\sigma\)가 된다.

두 번째 형태의 가정과 \(\mpideal^\sigma\)의 위 표현을 함께 쓰면, 모든 \(q\in\mq\)에 대해
\[
        q\equiv bc^\sigma\pmod{\mpideal^{\sigma+1}},
        \qquad b\not\equiv0(\mpideal),
\]
또는
\[
        bc^\sigma\equiv0(\mq,\mpideal^{\sigma+1})
\]
이다. \((\mq,\mpideal^{\sigma+1})\)가 으뜸이므로 여기서
\[
        c^\sigma\equiv0(\mq,\mpideal^{\sigma+1})
\]
또는
\[
        c^{\sigma+\lambda-1}\equiv0(\mq,\mpideal^{\sigma+\lambda})
\]
가 따른다. 따라서 \(\mpideal^{\sigma+\lambda-1}\equiv0(\mq)\)이다.

4. 요약하면 다음과 같다.

\textbf{정리 VI.} \emph{가환환에서 공리 I--V가 성립하면, 영아이디얼 및 단위아이디얼과 다른 모든 아이디얼은 영아이디얼 및 단위아이디얼과 다른 유한개의 소아이디얼의 거듭제곱곱으로 유일하게 표현된다. 이 소아이디얼들은 단위아이디얼과 다른 진약수를 갖지 않는다.}

\begin{center}
{\bfseries\large § 9.\par}
\vspace{0.35em}
{\bfseries\large 가정한 분해에서 공리들을 이끌어 내기.\par}
\end{center}

정리 VI에 따라 공리 I--V에서 나오는 통상적인 아이디얼 분해의 존재로부터 역으로 이 공리들을 이끌어 내기 위해서는 분해를 다음 형태로 가정해야 한다.

\textbf{가정.} \emph{가환환 \(\mR\)에서 영원소 또는 단위원에서 생성되지 않는 모든 아이디얼은 소아이디얼들의 거듭제곱곱으로 유일하게 표현된다. 이 소아이디얼들은 단위원에서 생성되지 않는 단순아이디얼, 곧 \(\mo\)와 다른 진약수를 갖지 않는 아이디얼이다. 역으로 모든 단순아이디얼은 소아이디얼이다. 유일성은 영원소 또는 단위원에서 생성되지 않는 모든 아이디얼에 대해}
\[
        \ma^\varrho\ne\ma^{\varrho+1}
\]
\emph{이라는 강한 형태로 성립한다.}

여기서는 단위원의 존재를 가정하지 않는다. 단위원이 없으면 \glqq{}단위원에서 생성되지 않는다\grqq{}는 조건은 아무 제한도 주지 않는다.

\textbf{1. 공리 III(단위원의 존재)의 증명.} 공리 III이 성립하지 않는다고 하자. 그러면 가정에 어긋나게 \(\mo^2\)이 소아이디얼은 아니지만 단순아이디얼임을 보이면 된다. 강한 유일성 가정에 따라 \(\mo\ne\mo^2\)이다. 따라서 \(\mo\not\equiv0(\mo^2)\)이지만 \(\mo^2\equiv0(\mo^2)\)이므로 \(\mo^2\)은 소아이디얼이 아니다. 그러나 \(\mo^2\)은 단순하다. \(\mo^2\)은 \(\mo\)와 다른 어느 소아이디얼로도 나누어질 수 없으므로, 가정한 곱표현에 따라 \(\mo,\mo^2\)과 다른 약수를 갖지 않기 때문이다.

\textbf{2. 공리 IV(영인자가 없는 환)의 증명.} \(a\)가 영인자, 곧 \(ab=0\)이지만 \(a\ne0\), \(b\ne0\)이라고 하자. \(a,b\)에서 생성되는 주아이디얼의 표현들을 곱하면
\[
        (0)=\mpideal_1^{\varrho_1}\ldots\mpideal_k^{\varrho_k}
\]
를 얻는다. 단위원이 존재하므로 \(\mpideal\)들은 영아이디얼 및 단위아이디얼과 다르다. 어떤 한 인자를 생략하면 영아이디얼의 진약수가 되도록 이 표현이 최단이라고 가정한다. 영아이디얼에는 유일표현에 관한 가정을 두지 않았으므로 이 조건은 저절로 성립하는 것이 아니다. 이 표현에 소아이디얼이 하나만 나타나면
\[
        (0)=\mpideal^\varrho=\mpideal^{\varrho+1}
\]
이어서 강한 유일성 조건에 모순된다. 둘 이상의 소아이디얼이 나타나면
\[
\begin{aligned}
\mpideal_1^{\varrho_1+1}
 &=\bigl(\mpideal_1^{\varrho_1+1},
          \mpideal_1^{\varrho_1}\ldots\mpideal_k^{\varrho_k}\bigr)\\
 &=\mpideal_1^{\varrho_1}
   \bigl(\mpideal_1,\mpideal_2^{\varrho_2}\ldots\mpideal_k^{\varrho_k}\bigr)\\
 &=\mpideal_1^{\varrho_1}.
\end{aligned}
\]
단위원이 존재하므로 \(\mpideal_1\)은 나머지 모든 소아이디얼과 서로소이기 때문이다. 이 경우에도 강한 유일성 조건에 모순된다.

\textbf{3. 공리 I, II(사슬정리)의 증명.} 가정들로부터 모든 비영아이디얼에 대해 \emph{나눗셈에서 곱표현이 따른다}. 즉 \(\mm\equiv0(\ma)\), \(\ma\ne\mo\)이면 \(\mm=\ma\mb\)이면서 \(\mb\not\equiv0(\mm)\)인 \(\mb\)가 존재한다. 실제로
\[
        \mm\equiv0(\ma),\qquad
        \mm=\mpideal_1^{\varrho_1}\ldots\mpideal_r^{\varrho_r},\qquad
        \ma=\overline{\mpideal}_1^{\sigma_1}\ldots
             \overline{\mpideal}_{\bar r}^{\sigma_{\bar r}}
\]
라 하자. 그러면 각 \(\overline{\mpideal}\)은 어떤 \(\mpideal\)과 같고 \(\sigma_i\leqq\varrho_i\)이다. 따라서 일부 \(\sigma_i\)는 0일 수 있게 하여
\[
        \mb=\mpideal_1^{\varrho_1-\sigma_1}\ldots
             \mpideal_r^{\varrho_r-\sigma_r}
\]
로 놓으면 주장에 이른다. 그러므로 \(\mm\)의 약수들은 \(0\leqq\sigma_i\leqq\varrho_i\)의 조합들에 대응하는 유한개뿐이다. 따라서 \(\mm\)에 따른 잉여류환에는 이중사슬정리가 성립한다. 이로써 공리 I, II가 증명되었다. 약수사슬정리는 \(\mR\) 자체에도 성립한다. 영아이디얼의 모든 진약수는 영아이디얼과 다르기 때문이다.

\emph{분수체 안에서 정수적으로 닫혀 있다는 공리 V를 증명하려면 아이디얼 분해에서 나오는 통상적인 결론들이 필요하므로, 이를 먼저 유도해야 한다.}

4. \emph{모든 비영아이디얼에 따른 잉여류환은 주아이디얼환이다.} 단위아이디얼에 따른 잉여류환은 영원소만으로 이루어져 조건이 분명 성립하므로, 여기서 아이디얼은 단위아이디얼과도 다르다고 가정해도 된다.

먼저 아이디얼이 으뜸인 경우 \(\mq=\mpideal^\varrho\)라 하자. \(c\equiv0(\mpideal)\), \(\not\equiv0(\mpideal^2)\)이면 3에 따라 \(\mo c=\mpideal\ma\), \((\ma,\mpideal)=\mo\)인 \(\ma\)가 존재한다. 따라서 모든 \(\sigma\leqq\varrho\)에 대해
\[
        \mpideal^\sigma=(\mo c^\sigma,\mpideal^\varrho)
\]
이고, 으뜸아이디얼에 따른 잉여류환의 모든 아이디얼은 주아이디얼이다. 일반적으로
\[
        \mm=\mq_1\mq_2\ldots\mq_r,
        \qquad \mq_i=\mpideal_i^{\varrho_i},
\]
이고 잉여류환이 이에 따라 직접합
\[
        \mR\mid\mm=\overline{\ma}_1+\ldots+\overline{\ma}_r
\]
으로 표현된다고 하자. \(\overline{\ma}_i\)는 \(\mR\mid\mq_i\)와 동형이므로(§ 4, 5) \(\overline{\ma}_i\)의 모든 아이디얼은 주아이디얼이다. \(\overline{\mc}\)가 잉여류환의 임의의 아이디얼이면 \(\overline{\ma}_i\overline{\mc}\)는 주아이디얼 \(\overline{\ma}_i\overline c_i\)이다. 따라서
\[
\begin{aligned}
        \overline{\mc}
        &=\overline{\mo}\,\overline{\mc}\\
        &=\overline{\ma}_1\overline{\mc}+\ldots+
          \overline{\ma}_r\overline{\mc}\\
        &=\overline{\ma}_1\overline c_1+\ldots+
          \overline{\ma}_r\overline c_r\\
        &=\overline{\ma}_1\overline c+\ldots+
          \overline{\ma}_r\overline c\\
        &=\overline{\mo}\,\overline c,
\end{aligned}
\]
여기서 \(\overline c=\overline c_1+\ldots+\overline c_r\)로 놓았다. 실제로 \(\overline{\ma}_i\overline{\ma}_k=(0)\)이고 \(\overline c_i\equiv0(\overline{\ma}_i)\)이므로, \(\overline{\ma}_i\) 안의 주아이디얼 \(\overline{\ma}_i\overline c_i\)와 \(\overline{\ma}_i\overline c\)는 일치한다.

잘 알려져 있듯이 주아이디얼환에 관한 정리는 다음 형태로도 쓸 수 있다.

\emph{\(\mc\)가 임의의 아이디얼이면, 주어진 한 아이디얼 \(\mb\)와 서로소인 아이디얼을 \(\mc\)에 곱하여 주아이디얼로 만들 수 있다.}

실제로 \(\mm=\mb\mc\)로 놓으면 \(\mc\)는 \(\mm\)을 법으로 주아이디얼이 된다. 즉
\[
        \mc=(\mo c,\mm)=(\mc\ma,\mc\mb)=\mc(\ma,\mb).
\]
따라서 \(\mo c=\mc\ma\)이고 \((\ma,\mb)=\mo\)이다.

\textbf{5. 분수아이디얼의 이론.} 분수체 \(\mK\) 안의 모든 유한 \(\mR\)-모듈을 분수아이디얼이라 한다.

\textbf{정리.} \emph{\(\mK\) 안의 비영 유한 \(\mR\)-모듈들은 곱셈에 관하여 아벨군을 이룬다.}

두 유한 \(\mR\)-모듈의 곱은 다시 유한 \(\mR\)-모듈이고, 곱셈은 결합적이며 가환적이다. \(\mA=\mo\mA\)이므로 단위아이디얼은 이 계의 단위원이다. 따라서 방정식 \(\mathfrak A\mX=\mo\)가 유한 \(\mR\)-모듈의 계 안에서 언제나 해를 갖는다는 것만 보이면 된다.

\textbf{예비주석.} 방정식 \(\mathfrak A\mT=\mathfrak B\)가 정확히 하나의 해를 가지면 \(\mT\)는 모듈 몫 \(\mathfrak B:\mathfrak A\)와 같다(§ 5, 4). 실제로
\[
\mT\equiv0(\mathfrak B:\mathfrak A),\qquad
\mathfrak B=\mathfrak A\mT
\equiv0\bigl(\mathfrak A\cdot(\mathfrak B:\mathfrak A)\bigr)
\equiv0(\mathfrak B),
\]
따라서
\[
        \mathfrak A\cdot(\mathfrak B:\mathfrak A)=\mathfrak B
\]
이다.

먼저 \(\mathfrak A\)가 주모듈 \(\mathfrak A=(\alpha)\)이면 \(\mX=(e/\alpha)\)가 \(\mathfrak A\mX=\mo\)의 해이고, \(\mX\)는 다시 유한 \(\mR\)-모듈이다. 이로부터 다음이 따른다. \emph{모든 유한 \(\mR\)-모듈은 \(\mR\)의 두 아이디얼의 모듈 몫이다.} 이 사실은 분수아이디얼이라는 이름을 정당화한다. 실제로 \(\tau_1,\ldots,\tau_r\)가 \(\mT\)의 모듈기저이고 \(\tau_i=t_i/a\)이며 \(\mt\)가 \(t_i\)들에서 생성되는 \(\mR\)의 아이디얼이면
\[
        \mo a\cdot\mT=\mt.
\]
따라서 예비주석에 따라 \(\mR\) 안에서 몫을 취하여 \(\mT=(\mt:\mo a)\)이다.

이제 \(\mathfrak A\mX=\mo\)의 해를 일반적으로 줄 수 있다. \(\mathfrak A=\ma:\mo c\)로 놓고 \(\ma\mb\)가 주아이디얼 \(\mo a\)가 되도록 \(\mb\)를 택한다. 그러면
\[
        \mathfrak A c=\ma,
        \qquad
        \mathfrak A\mb c=\mo a,
        \qquad
        \mathfrak A\cdot(\mb c:\mo a)=\mo.
\]
여기서 \(\mX\)는 한 아이디얼을 한 주아이디얼로 나눈 몫이므로 다시 유한 \(\mR\)-모듈이다. 이로써 \emph{군의 성질이 증명}되었다. 또 이로부터 \emph{임의의 두 아이디얼의 모듈 몫 \(\ma:\mb\)}는 방정식 \(\mb\mX=\ma\)의 해로서 \emph{유한 \(\mR\)-모듈}이라는 것이 따른다.

6. 군의 성질에서 다음 결론들이 나온다.

6\(\alpha\). \emph{약분할 수 있다.} 즉 \(\mathfrak T=\mathfrak A\mathfrak M:\mathfrak B\mathfrak M\)이면 \(\mathfrak T=\mathfrak A:\mathfrak B\)이고 그 역도 성립한다. 실제로 \(\mathfrak T\mathfrak B\mathfrak M=\mathfrak A\mathfrak M\)에서 \(\mathfrak T\mathfrak B=\mathfrak A\)가 나오며 그 역도 성립한다.

6\(\beta\). \emph{모든 분수아이디얼, 곧 유한 \(\mR\)-모듈은 \(\ma,\mb\)가 서로소인 표현 \(\mathfrak C=\ma:\mb\)를 갖고, 이 조건으로 \(\ma,\mb\)는 유일하게 결정된다.} 표현의 가능성은 5와 6\(\alpha\)에서 나온다. 한편
\[
        \mathfrak C=\ma:\mb=\overline{\ma}:\overline{\mb}
\]
이면 \(\ma\overline{\mb}=\mb\overline{\ma}\)이다. \(\ma,\mb\)도 서로소이고 \(\overline{\ma},\overline{\mb}\)도 서로소라고 가정하면 여기서 유일성이 나온다.

6\(\gamma\). \emph{정수적이지 않은 원소, 곧 \(\mK\)에는 속하지만 \(\mR\)에는 속하지 않는 원소에서 생성되는 모든 주모듈 \(\mo\eta\)는 두 주아이디얼의 몫으로 표현되되 분자의 어떤 거듭제곱도 분모로 나누어지지 않는다.} 기약표현 \(\mo\eta=\mb:\mc\), \((\mb,\mc)=\mo\)를 택하고, 4에 따라 \(\mo b=\mb\ma\), \((\ma,\mc)=\mo\)인 \(\ma\)를 택하자. 그러면
\[
        \mo\eta=\ma\mb:\ma\mc=\mo b:\ma\mc.
\]
또 \(\ma\mc=\mo b:\mo\eta\)이므로 \(\ma\mc\)도 주아이디얼이고, 따라서
\[
        \mo\eta=\mo b:\mo c.
\]
\((\ma,\mc)=\mo\), \((\mb,\mc)=\mo\)이므로 \(\ma\mb\)의 어떤 거듭제곱도 \(\mc\)로 나누어지지 않는다.

7. \emph{분수체 안에서 정수적으로 닫혀 있다는 공리 V의 증명.} 6\(\gamma\)에서 \emph{\(\mK\)의 모든 정수적이지 않은 원소는}
\[
        \eta=a/c
\]
\emph{로 표현되되, \(\mR\)에서 분자의 어떤 거듭제곱도 분모로 나누어지지 않게 할 수 있다.} 실제로 \(\mo\eta=\mo b:\mo c\)이면 \(\eta\)는 몫 \(b/c\)와 \(\mR\)의 한 단위원만큼만 다르다. 따라서 그러한 \(\eta\)는 \(\mR\)에 대하여 정수적임을 특징짓는 어떤 방정식도 만족할 수 없다. 실제로
\[
        \eta^n+r_1\eta^{n-1}+\ldots+r_n=0
\]
이면 \(\mR\)에서 \(a^n\equiv0(\mo c)\)가 되기 때문이다.

8. \emph{따름정리: 가환환 \(\mR\)에서 공리 I--IV가 성립하고 모든 으뜸아이디얼이 기약이면, 분수체 안에서 정수적으로 닫혀 있다는 공리 V도 성립한다.} 공리 I--III에 따라 모든 소아이디얼 거듭제곱 \(\mpideal^\varrho\)는 으뜸아이디얼이다. 특히 \(\mpideal^2\)도 그러하고, 가정에 따라 기약이다. 이제
\[
        \mpideal=(\mo c,\mpideal^2)
\]
임을 보이면 된다. 그러면 § 8, 3의 보조정리에 따라 모든 으뜸아이디얼이 소아이디얼 거듭제곱이 되고, 이는 다른 가정들과 함께 7에 따라 정수적 닫힘을 낳는다.

약수사슬정리에 따라
\[
        \mpideal=(\mo c_1,\ldots,\mo c_k,\mpideal^2)
\]
이고, 여기서 \(c\)의 곱셈자로는 \(\mpideal\)에 따른 잉여류들만 고려하면 되며, \(c\)들은 \(\mpideal\)에 따른 잉여류체에 관하여 선형독립이라고 가정할 수 있다. 이 선형독립성으로부터 \(k>1\)이면
\[
        \mpideal^2=[\mq_1,\mq_2],\qquad
        \mq_1=(\mo c_1,\mpideal^2),\qquad
        \mq_2=(\mo c_2,\ldots,\mo c_k,\mpideal^2)
\]
인 표현이 나온다. 이때 \(\mq_1,\mq_2\)는 \(\mpideal^2\)의 진약수이므로 \(\mpideal^2\)이 가약이 되어 가정에 모순된다. 이로써 \(\mpideal=(\mo c,\mpideal^2)\)와 그에 따르는 결론이 증명되었다.

역으로 공리 I--V의 결과로 모든 으뜸아이디얼이 기약이라는 것은 곧바로 알 수 있다. 소아이디얼 거듭제곱 \(\mpideal^\varrho\)는 \(\mpideal^\sigma,\mpideal^\tau\) 꼴의 진약수들의 최소공배수로 표현될 수 없기 때문이다.

9. \emph{환에 영인자를 허용하면, 공리 I--III과 몫환 안에서의 정수적 닫힘으로부터 모든 으뜸아이디얼이 기약이라는 것은 따르지 않는다}\footnote{각주 19를 보라.}. 실제로 \(\mT\)를 정수적 닫힘을 제외한 공리 I--IV가 모두 성립하는 환이라 하자. 그런 환은 존재한다. 예컨대 \(\mR\)에 공리 I--V가 성립할 때, \(\mR\)의 분수체의 한 유한 확대체 안에서 \(\mR\)에 대하여 정수적인 모든 원소의 계와 다른 유한 오더들이 그러하다(§ 3, 2). 따름정리 8에 따라 \(\mT\)에는 가약 으뜸아이디얼 \(\mq\)가 존재한다. \(\mpideal^\sigma\)를 \(\mq\)의 진배수라 하고 \(\overline{\mR}\)을 잉여류환 \(\mT\mid\mpideal^\sigma\)라 하자. 그러면 \(\mq\)에 대응하는 아이디얼 \(\overline{\mq}\)는 \(\overline{\mR}\)에서 비영 가약 으뜸아이디얼이 된다. \(\overline{\mR}\)에는 공리 I--III이 성립하고 몫환 안에서의 정수적 닫힘도 성립한다. 실제로 \(\mT\)에서 \(\mpideal\)로 나누어지지 않는 모든 원소는 \(\mpideal^\sigma\)와 서로소이므로, \(\overline{\mR}\)의 모든 정칙원소는 단위원이다. 따라서 \(\overline{\mR}\)은 그 몫환과 같다.

\begin{center}
{\bfseries\large § 10.\par}
\vspace{0.35em}
{\bfseries\large 이중사슬정리와 합성열.\par}
\end{center}

이하에서는 정렬되어 있다고 가정한 임의의 모듈영역(§ 2)에 대하여 \emph{이중사슬정리가 성립한다는 가정}, 곧 모든 모듈 약수사슬과 배수사슬이 유한 단계에서 멈춘다는 가정이 \emph{합성열이 존재한다는 가정}과 같음을 보인다\footnote{이 영역의 모듈들은 덧셈에 관하여 아벨군을 이룬다. 곱셈자영역이 \(\mR\)의 원소들로 이루어지므로 정확히는 \emph{일반화된} 아벨군이다. 아벨군이므로 합성열과 주열이라는 개념은 일치한다. 이 절의 정리들은 \glqq{}모듈\grqq{}을 임의의 비가환 일반화군을 포함한 한 군의 모든 정규부분군의 전체로 이해해도 성립한다. 뒤의 표기가 앞의 표기와 조금 다른 것은 군론을 떠올리게 하려는 것이다.}. 모듈영역을 가환환으로 특수화하면 환의 모든 아이디얼계에 대한 대응 사실들이 나온다. 이때 2에서는 줄곧 모듈동형을 환동형으로 바꾸어야 한다.

모듈영역 \(\mG\) 안에 \(\mG\) 자체와 영모듈 \(\mE\) 이외의 \(\mR\)-모듈이 없으면 \(\mG\)를 \emph{단순}하다고 한다. \(\mG\mid\mA\)가 단순하면 모듈 \(\mA\)를 \(\mG\) 안에서 \emph{단순}하다고 한다(극대 정규부분군).

\emph{배수사슬}
\[
        \mG,\mA_1,\ldots,\mA_r,\mE
\]
\emph{의 모든 모듈이 서로 다르고 각 모듈이 앞선 모듈 안에서 단순하면 이를 길이 \(r\)인 \(\mG\)의 합성열이라 한다.} 즉 잉여류모듈, 곧 몫군 \(\mA_{i-1}|\mA_i\)들과 \(\mG|\mA_1\), \(\mA_r\)가 단순 모듈영역이다. 잉여류모듈 \(\mG|\mA\)를 만들면 \(\mG\)에서 \(\mA\ne\mE\)까지의 합성열이 존재하는 경우를 합성열 자체의 경우로 환원할 수 있다.

1. \emph{\(\mG\)에서 이중사슬정리가 성립한다고 가정하면 \(\mG\)에는 합성열이 존재한다.} 여기서 \(\mG\ne\mE\)라고 가정한다. \(\mG\)의 정렬순서와 약수사슬정리로부터 § 6에서처럼 준사전식 순서를 써서 모든 모듈계의 정렬순서를 얻는다. 이 정렬순서와 약수사슬정리의 가정에 따라 \(\mG\) 안에서 단순한 모듈이 적어도 하나 존재한다. 실제로 \(\mE_1\)을 정렬순서에서 \(\mE\)의 첫 진약수라 하고, 일반적으로 \(\mE_i\)를 \(\mE_{i-1}\)의 첫 진약수라 하면 잘 정의된 사슬
\[
        \mG,\mG_1,\ldots,\mG_i,\ldots
\]
은 유한 단계에서 멈춘다. 따라서 반드시 \(\mG\) 안에서 단순한 모듈 \(\mA_1\)에 이른다. \(\mA_1\ne\mE\)이면 같은 방법으로 \(\mA_1\) 안에서 단순한 모듈 \(\mA_2\)를 얻는다. 일반적으로 \(\mA_{i-1}\ne\mE\)이면 \(\mA_{i-1}\) 안에서 단순한 모듈 \(\mA_i\)가 존재한다. 이 방식으로 잘 정의된 배수사슬
\[
        \mG,\mA_1,\ldots,\mA_i,\ldots
\]
을 얻는다. 이 사슬은 배수사슬정리의 가정에 따라 유한 단계에서 멈추므로 \(\mG\)의 합성열을 준다.

2. \emph{\(\mG\)에 합성열이 존재한다고 가정하면 \(\mG\)에서 이중사슬정리가 성립한다.} 증명은 완전귀납법으로 하며, \(\mG\)에 정렬순서를 가정할 필요는 없다. 귀납과정에서는 합성열의 존재만을 가정한 조르당--횔더 정리의 증명도 함께 나온다\footnote{환의 경우는 Masazo Sono(각주 19), \emph{On Congruences I}, §§ 11--13을 보라. 그곳에서는 비가환군에서 합성열에 대응하는 것, 곧 \(\mR,\mA_1,\ldots,\mA_r,\mE\) 꼴의 열도 다룬다. 여기서 각 \(\mA_i\)는 \(\mA_{i-1}\) 안의 단순아이디얼이지만 \(\mR\)의 아이디얼일 필요는 없다. 실제로 비가환군에서도 배수사슬을 각각 \(\mA_i\)가 \(\mA_{i-1}\)의 정규부분군인 \(\mG,\mA_1,\ldots,\mA_i,\ldots\)로 이해하고 약수사슬도 그에 맞게 정의하면 위 논의가 모두 유지된다. 또 Dedekind, \emph{Über die von drei Moduln erzeugte Dualgruppe}, Math. Ann. 53 (1900), 371--403쪽을 보라. 그 논문에서도 훨씬 일반적인 영역(모듈군)에 대해 존재만을 가정하여 조르당--횔더 정리를 증명한다(§ 6, XVI). 여기서는 제2동형정리 대신 조금 더 약한 대응관계가 나타나므로 정리의 진술도 조금 더 약하지만, 증명은 원리상 위의 증명과 같다.}.

먼저 \(\mG\)가 단순, 곧 \(r=0\)이어서 합성열 \(\mG,\mE\)만 존재한다고 하자. 이때 다음 명제들이 성립한다.
\[
\begin{array}{ll}
\alpha) & \text{\(\mG\) 안에서 단순한 모든 모듈을 지나는 합성열을 그을 수 있다.}\\
\beta) & \text{조르당--횔더 정리: 모든 합성열의 길이는 같고, 몫군}\\
       & \text{(잉여류모듈)의 계는 순서를 제외하고 일치한다. 즉 대응하는}\\
       & \text{몫군들은 동형이다.}\\
\gamma) & \text{\(\mG\)와 다른 임의의 모듈을 지나는 합성열을 그을 수 있다.}\\
\delta) & \text{배수사슬정리가 성립한다.}\\
\varepsilon) & \text{약수사슬정리가 성립한다.}
\end{array}
\]
따라서 길이 \(\bar r<r+1\)인 합성열을 갖는 모든 모듈영역에는 \(\alpha)\)--\(\varepsilon)\)이 성립한다고 가정할 수 있다. 이제 \emph{\(\mG\)에 길이 \(r+1\)인 합성열}
\[
        \mG,\mA,\mA_1,\ldots,\mA_r,\mE
\]
\emph{이 존재한다는 가정} 아래에서 이 명제들을 증명한다.

\textbf{2\(\alpha\).} \(\mA\)와 다른, \(\mG\) 안에서 단순한 모듈이 없으면 증명할 것이 없다. \(\mB\ne\mA\)가 \(\mG\) 안에서 단순하다고 하자. \(\mB\)를 지나는 \(\mG\)의 합성열이 존재함을 보여야 한다. \(\mA,\mB\)가 모두 \(\mG\) 안에서 단순하고 서로 다르므로 반드시 \((\mA,\mB)=\mG\)이다. \(\mM=[\mA,\mB]\)로 놓으면 제2동형정리(§ 4, 2)에 따라
\[
        \mG|\mB\simeq\mA|\mM,
        \qquad
        \mG|\mA\simeq\mB|\mM.
\]
따라서 \(\mM\)은 \(\mA,\mB\) 안에서 단순하다. \(\mA\)는 길이 \(r<r+1\)인 합성열을 가지므로 귀납가정 \(\alpha)\), \(\beta)\)에 따라 \(\mA\)에는 \(\mM\)을 지나는 길이 \(r\)인 합성열, 이를테면
\[
        \mA,\mM,\mM_1,\ldots,\mM_{r-1},\mE
\]
이 존재한다. 그러므로
\[
        \mG,\mB,\mM,\mM_1,\ldots,\mM_{r-1},\mE
\]
는 \(\mB\)를 지나는 \(\mG\)의 합성열이다.

\textbf{2\(\beta\).} 조르당--횔더 정리를 증명하기 위하여
\[
\begin{aligned}
(1)\quad&\mG,\mA,\mA_1,\ldots,\mA_r,\mE,\\
(2)\quad&\mG,\mB,\mB_1,\ldots,\mB_s,\mE,
\qquad s\geqq r
\end{aligned}
\]
를 \(\mG\)의 두 합성열이라 하자. \(\mA=\mB\)이면 \(r<r+1\)에 대한 귀납가정으로 모두 끝난다. 그렇지 않으면 2\(\alpha\)에서 구성한 열
\[
\begin{aligned}
(3)\quad&\mG,\mA,\mM,\mM_1,\ldots,\mM_{r-1},\mE,\\
(4)\quad&\mG,\mB,\mM,\mM_1,\ldots,\mM_{r-1},\mE
\end{aligned}
\]
과 비교한다. \(r<r+1\)에 대한 귀납가정에 따라 (1)과 (3)의 몫군계가 일치하고, (4)에서 \(\mB\)를 지나는 합성열의 길이가 \(r\)이므로 (2)와 (4)의 몫군계도 일치한다. 따라서 \(s=r\)이다. 또한 2\(\alpha\)의 동형에 따라 (3)과 (4)가 일치한다. 이로써 (1)과 (2)가 일치하며 조르당--횔더 정리가 증명되었다.

\textbf{2\(\gamma\).} 예비주석. \emph{\(\mA,\mB,\mH\)가 \(\mG\)의 모듈이고 \(\mB\)가 \(\mA\) 안에서 단순하면, \((\mB,\mH)\)와 \((\mA,\mH)\)는 서로 같거나 \((\mB,\mH)\)가 \((\mA,\mH)\) 안에서 단순하다.} 제2동형정리(§ 4, 2)에 따라
\[
        (\mA,\mB,\mH)|(\mB,\mH)
        \simeq \mA|[\mA,(\mB,\mH)]
\]
이다. \(\mB\equiv0(\mA)\)이므로 이는
\[
        (\mA,\mH)|(\mB,\mH)\simeq\mA|\mC,
        \qquad
        \mC=[\mA,(\mB,\mH)]
\]
로 쓸 수 있다. 여기서 \(\mC\equiv0(\mA)\)이다. 한편 \(\mB\equiv0(\mA)\), \(\mB\equiv0((\mB,\mH))\)이므로 \(\mB\equiv0(\mC)\)이기도 하다. 따라서 \(\mC\)는 \(\mA,\mB\) 사이에 놓이고, 가정에 따라 \(\mA\) 또는 \(\mB\)와 같아 예비주석이 따른다.

이제 다시 \(\mG,\mA,\mA_1,\ldots,\mA_r,\mE\)를 \(\mG\)의 한 합성열이라 하고, \(\mH\ne\mG\)를 \(\mG\)의 임의의 모듈이라 하자. \(\mH\)를 지나는 합성열이 존재함을 보이기 위하여 모듈의 열
\[
        \mG,(\mA,\mH),(\mA_1,\mH),\ldots,(\mA_r,\mH),\mH
\]
을 만든다. 예비주석에 따라 이들 가운데 서로 다른 모듈 \(\mG,\mC,\mC_1,\ldots,\mC_s,\mH\)는 \(\mG\)에서 \(\mH\)까지의 합성열을 이룬다. 따라서 \(\mC\)는 \(\mG\) 안에서 단순한 모듈이고, 특수한 경우 2\(\alpha\)에서는 \(\mH\)와 같다. 그러므로 2\(\alpha\)에 따라 \(\mC\)에는 길이 \(r<r+1\)인 합성열이 존재하고, 귀납가정에 따라 \(\mC\)에서 \(\mH\)를 지나는 합성열을 그을 수 있다. 여기에 \(\mG\)를 덧붙이면 \(\mG\)에서 \(\mH\)를 지나는 합성열을 얻는다. 이 논의를 되풀이하면 특히 그러한 합성열을 \(\mG\)에서 \(\mH\)까지 구성한 열의 연장으로 택할 수 있음도 알 수 있다.

\textbf{2\(\delta\).} 배수사슬정리의 증명. 다시 \(\mG\)에 길이 \(r+1\)인 합성열이 존재한다고 가정하고,
\[
        \mG,\mB,\mB_1,\ldots,\mB_\nu,\ldots
\]
를 각 모듈이 앞선 모듈의 진배수인 배수사슬이라 하자. 언제나 \(\mG\)를 첫 항으로 덧붙일 수 있으므로 사슬이 \(\mG\)에서 시작한다고 해도 일반성을 잃지 않는다. 2\(\gamma\)에 따라 \(\mG\) 안에는 \(\mB\)를 지나는 합성열이 있고, 따라서 \(\mB\)는 더 짧은 합성열을 갖는다. 그러므로 귀납가정에 따라 사슬 \(\mB,\mB_1,\ldots,\mB_\nu,\ldots\)은 유한 단계에서 멈추고, \(\mG\)를 덧붙인 사슬도 마찬가지이다.

\textbf{2\(\varepsilon\).} 약수사슬정리의 증명. 다시 \(\mG\)에 길이 \(r+1\)인 합성열이 존재한다고 가정하고,
\[
        \mE,\mC,\mC_1,\ldots,\mC_\nu,\ldots
\]
를 각 모듈이 앞선 모듈의 진약수인 약수사슬이라 하자. 사슬이 \(\mE\)에서 시작한다고 해도 역시 일반성을 잃지 않는다. 2\(\gamma\)에 따라 \(\mG\) 안에는 \(\mC\)를 지나는 합성열이 있고, 따라서 \(\mG\mid\mC\)에는 더 짧은 합성열이 있다. 그러므로 귀납가정에 따라 \(\mG\mid\mC\)의 약수사슬
\[
        \mC\mid\mC,\mC_1\mid\mC,\ldots,\mC_\nu\mid\mC,\ldots
\]
은 유한 단계에서 멈춘다. 제1동형정리(§ 4, 2)에 따른 일대일 대응 때문에 원래의 \(\mC\)에서 시작하는 사슬도 마찬가지이고, \(\mE\)를 덧붙인 사슬도 그러하다.

\emph{이로써 두 가정, 곧 합성열의 존재와 이중사슬정리의 동치가 증명되었다.}

\begin{center}
(1925년 8월 13일 접수.)
\end{center}

\clearpage
\setcounter{footnote}{0}

\clearpage
\setcounter{footnote}{0}
\phantomsection
\addcontentsline{toc}{section}{31. 대수적 수체 또는 함수체의 오더에 대한 판별식 정리}
%
%
%
%
%
%

\providecommand{\mH}{\mathfrak H}
\providecommand{\mRprime}{\mathfrak R'}
\providecommand{\bP}{\overline P}
\providecommand{\bR}{\overline{\mathfrak R}}
\providecommand{\bp}{\overline{\mathfrak p}}
\providecommand{\mc}{\mathfrak c}
\providecommand{\mt}{\mathfrak t}
\providecommand{\mq}{\mathfrak q}
\providecommand{\mpideal}{\mathfrak p}
\providecommand{\mP}{\mathfrak P}

\section*{31. 대수적 수체 또는 함수체의 오더에 대한 판별식 정리}
\begin{center}
\emph{Journal f. d. reine u. angew. Math. 157 (1927), 82--104쪽}

에미 뇌터 (괴팅겐)\\[0.4em]
\rule{9em}{0.4pt}
\end{center}

잘 알려진 \emph{데데킨트 판별식 정리}는 다음과 같이 말한다. 한 소수는 적어도 한 소아이디얼의 제곱으로 나누어질 때 그리고 오직 그때 체의 판별식을 나눈다.

이 글에서는 이 정리가 한 (유한) 대수적 수체의 모든 유한 오더, 곧 유리정수환을 포함하는 모든 \emph{대수적 정수로 이루어진 환}에 대하여 다음 꼴로 유지됨을 보이겠다. \emph{한 소수 \(p\)가 한 오더의 판별식을 나누는 것은, 그 오더에서 성립하는 주아이디얼 \((p)\)의 극대 으뜸성분 분해에 적어도 하나의 진으뜸아이디얼, 곧 소아이디얼이 아닌 으뜸아이디얼이 나타날 때 그리고 오직 그때이다.}\srcfn{1)}{이에 대해서는 E. Noether, \emph{Abstrakter Aufbau der Idealtheorie in algebraischen Zahl- und Funktionenkörpern}, Math. Ann. 96 (1926), 26--61쪽을 보라. 모든 정의는 이하 “아이디얼 이론”이라 부를 이 논문을 참조한다. 이하 “오더”는 언제나 “유한 오더”를, “환”은 가환환을 뜻한다.}

정수환, 곧 그 체의 모든 대수적 정수로 이루어진 계에서는 소아이디얼의 거듭제곱 이외의 으뜸아이디얼이 없으므로, 특수화하면 \emph{데데킨트 정리}가 나온다.

증명에서는 \((p)\)를 법으로 한 오더의 잉여류환으로 옮겨, 환 안에 포함된 한 체에 관하여 유한 랭크인 환, 즉 바로 그 체의 계수를 가지며 단위원을 갖는 가환 초복소량계의 분해이론으로 문제를 환원한다. 여기서 계수체로는 특히 유리정수의 \(p\)를 법으로 한 잉여류체가 나타나고, \((p)\)의 분해에는 잉여류환에서 영아이디얼의 분해가 대응한다.

그러므로 그러한 유한 랭크 환에서 영아이디얼을 유한개의 소아이디얼의 최소공배수로 나타낼 수 있을 때 그 환을 \emph{완전가약}이라 부르면, 문제는 완전가약성의 판정조건을 구하는 것이 된다. 우선 \emph{계수영역이 완전체}이면, 그 환이 완전가약인 것은 대수적으로 닫힌 계수영역을 붙여 얻는 확대환이 완전가약일 때 그리고 오직 그때임을 얻는다. 그런데 이 확대환에서는 판별식이 0이 아니라는 것이 완전가약성의 필요충분조건임을 아주 간단히 알 수 있다\srcfn{1)}{한 특수한 경우에는 Hilbert도 비슷하게 논증하였다. \emph{Zur Theorie der aus \(n\) Haupteinheiten gebildeten komplexen Größen}, Gött. Nachr. 1896.}. 오더의 판별식은 각 소수를 법으로 하여 잉여류환의 판별식으로 옮겨가므로, 이로써 판별식 정리가 증명된다\srcfn{2)}{이와 관련된 연구는 비가환인 경우에도 있으나 \emph{정수환}(극대영역)에 한정되고 유리정수 계수를 전제한다. 막 출판된 A. Speiser, \emph{Allgemeine Zahlentheorie}, Naturf. Gesellschaft Zürich 71 (1926), 8--48쪽 및 거기 적힌 미국 문헌을 보라.}.

수체 대신 함수체를 고찰하면, 이제 오더는 모든 정수 유리함수의 기초영역을, 또는 여러 부정원의 대수함수나 정수 계수 대수함수의 경우에는 정수 계수 다항식 계수를 갖는 모든 정수 함수식의 기초영역을 포함해야 한다. 이때 잉여류환에서 \emph{불완전체인 계수영역}이 나타날 수도 있다. 예컨대 정수 계수 함수식 영역에서 한 소수를 법으로 한 잉여류환을 고찰할 때 그러하다. 계수영역이 불완전체이면 소아이디얼을 제1종과 제2종으로 구분해야 한다. 이는 그 소아이디얼에 따른 잉여류체가 제1종 확대체인지 제2종 확대체인지에 따른다\srcfn{3)}{잘 알려진 Steinitz의 정식화는 J. f. M. 137에 있다. 이 논문의 § 2, 4에 붙인 주도 보라.}. 이에 따라 제1종과 제2종 완전가약성을 구분해야 한다. 위의 판정조건들은 모두 완전체의 경우에 유일하게 나타나는 제1종 완전가약성에 관한 것이다. 따라서 가장 일반적인 경우의 \emph{판별식 정리}는 다음을 말한다. \emph{기초영역의 한 소원소가 한 오더의 판별식을 나누는 것은, \((p)\)의 아이디얼 분해에 적어도 하나의 진으뜸성분 또는 제2종 소아이디얼이 나타날 때 그리고 오직 그때이다.}

이 사실은 정수환에 대하여 \emph{크로네커}가 처음 예들에서 알아차렸다. 그보다 앞선 기념논문에는 수체의 정수환에만 성립하는 정식이 아직 잘못 적혀 있었고, 그나마 완전한 증명은 없었다. 정수환에 대해서는 \emph{오스트롭스키}\srcfn{4)}{A. Ostrowski, \emph{Zur arithmetischen Theorie der algebraischen Größen}, Nachr. Gött. Ges. d. Wissensch. 1919.}가 위 판별식 정리의 증명과 그에 이어지는 더 세밀한 분기정리들을 제시했으며, 상세한 문헌목록도 함께 주었다. 끝으로 각 기초영역의 소아이디얼에 관한 몫환으로 이행하면 \emph{상대체의 모든 오더에 대한 판별식 정리}가 직접 따라온다. 이 이행의 착상은 특수한 경우에는 알려져 있었으나, 그 원리적 중요성은 \emph{W. 크룰}\srcfn{5)}{그의 단치히 강연, Jhrbr. d. D. Math. Ver. 34, 121쪽을 보라.}이 강조하였다. 더 일반적인 연구의 틀에서 몫환의 체계적 이론은 \emph{H. 그렐}\srcfn{6)}{\emph{Beziehungen zwischen den Idealen verschiedener Ringe} (Math. Annalen에 출판 예정).}이 제시한다.

모든 오더와 완전·불완전 계수영역을 한결같이 다루는 이 방법은 아이디얼 이론에 더하여, 유한 랭크 환의 이론과 그에 따른 유한 확대의 일반개념을 처음부터 기초로 삼고 원시원소로 생성되는 확대의 개념을 기초로 삼지 않는 데 있다. 따라서 오더의 모듈기저와 그 원소들 사이의 관계가 체의 원시원소의 거듭제곱과 정의방정식을 대신하며\srcfn{1)}{이 관점은 고차 복소수에 관한 Dedekind의 관점과 가깝다. \emph{Zur Theorie der aus \(n\) Haupteinheiten gebildeten komplexen Größen}, Gött. Nachr. 1885.}, 또는 기본형식의 거듭제곱과 기본방정식을 대신한다.

그러한 모듈기저는 기초영역의 각 소원소를 법으로 하여 잉여류환의 기저로 옮겨간다. 반면 원시원소의 거듭제곱으로 이루어진 기저는 정수환에서조차 그렇게 되지 않을 수 있고, 기본형식의 거듭제곱으로 이루어진 기저도 여러 부정원의 정수 계수 대수함수를 다룰 때에는 그렇다\srcfn{2)}{Ostrowski, 앞의 글, “제1종 비정칙성”을 보라.}. 달리 말하면, 한 부정원 다항식영역의 임의의 정의방정식 또는 기본방정식에 따른 잉여류환은 \(p\)를 법으로 했을 때 정수환의 \((p)\)에 따른 잉여류환과 동형일 필요가 없다. 바로 이 동형의 소멸에 통상적인 표현법의 주된 어려움이 있다. 여기서는 이 보조 다항식영역이 나타나지 않으므로 그런 어려움을 저절로 피한다. 동시에 잉여류환의 덧셈적 분해이론은 방정식의 곱셈적 분해와 동등한 것으로 볼 수 있다. 방정식을 서로소인 인자들로 분해하면 다항식 잉여류환에서도 덧셈적 분해가 이에 대응한다.

판별식 정리는 오더의 일반 분기이론에서 첫 정리에 지나지 않을 수 있음을 덧붙여 둔다. 정수환에서도 판별식 정리 옆에는 분기아이디얼의 더 세밀한 이론이 있어 지수를 더 정확히 결정할 수 있다. 이 지수결정은 방정식의 으뜸인자들의 거듭제곱표현에 바탕을 두며, \(p\)에서 \(q=p^e\)까지의 합성열 길이를 정하는 것으로 해석할 수 있다. 일반 분기이론에서도 이 꼴로 나타날 것으로 예상한다.

\begin{center}
\textbf{§ 1. 체의 확대환}
\end{center}

\textbf{1.} 먼저 임의의 환에 관한 몇 가지 말을 해 둔다. \(\mT\)가 환이고 \(S\)가 \(\mT\)의 원소계일 때, “\(\mT\)에서 \(S\)로부터 유도된 환”이란 \(S\)를 포함하는 \(\mT\)의 모든 환의 교집합을 뜻한다. “\(\mT\)에서 \(S\)로부터 유도된 아이디얼”도 이에 따라 정의하고, \(\mathsf{P}\)가 \(\mT\)의 부분환이면 “\(\mT\)에서 \(S\)로부터 유도된 \(\mathsf{P}\)-모듈”도 마찬가지로 정의한다\srcfn{3)}{이에 대해서는 \emph{아이디얼 이론} § 1을 보라.}.

\(\mR\)이 \(\mT\)의 부분환이고 \(S\)가 \(\mT\)의 원소계이면, \(\mT\)에서 \(\mR\)과 \(S\)로부터 유도된 환을 \(S\)를 \(\mR\)에 환으로 첨가하여 생긴 환 \(\mR[S]\)라 한다. 이 환은 \(S\)의 원소들에 관한, \(\mR\)의 계수를 갖는 모든 다항식으로 이루어지며, 등식관계와 연산관계는 \(\mT\)에서의 등식관계와 연산관계로 정해진다. 이들이 이미 정해져 있으므로 특수한 경우에는 모든 다항식 중 일부 계만, 예컨대 \(\mR\)의 계수를 갖는 \(S\)의 모든 선형형식만 만들면 충분하다. 예를 들어 \(\mR\)이 단위원을 갖는 환이고 \(S\)가 같은 단위원을 갖는 체이면 그러하다(2를 보라). \(c_i\in S\), \(\gamma_i\in\mR\)인 모든 선형결합 \(\sum c_i\gamma_i\)가 환 \(\mR[S]\)를 이룬다.

한편 \(\mR\)만 주어지고, \(\mR\)에 속하지 않으면서 등식관계 또는 서로의 연산관계를 갖는 원소계 \(S\)가 주어졌을 때에도 \(S\)를 \(\mR\)에 첨가한 확대환 \(\mR[S]\)를 구성할 수 있다. 이를 위해 \(\mR\)과 \(S\)를 포함하는 환이 언제나 존재함을, 곧 통상적인 연산법칙을 정하고 \(S\)의 모든 다항식을 \(\mR\)의 계수로 형식적으로 구성한 환이 존재함을 보이면 된다. 등식의 정의에는 연산법칙 외에도 \(\mR\)과 \(S\)에서 성립하는 등식의 정의를 포함해야 한다는 제한만 있을 뿐 완전한 자유가 남아 있다. 따라서 \(\mR\)과 \(S\)에서 같은 원소들로부터 같은 방식으로, 그 등식과 연산에 의해 생기는 다항식들만 반드시 서로 같다고 두면 된다. 이렇게 구성한 \(\mR\)과 \(S\)를 포함하는 환 \(\mT\)는 \(\mT\)에서 \(\mR\)과 \(S\)로부터 유도된 환과 동일하다\srcfn{1)}{Steinitz의 몫체 구성과 비교하라.}.

\textbf{2. 체의 확대환.} 이제부터 단위원을 가지며 한 체 \(P\)를 부분환으로 포함하는, 곧 체의 상환인 환 \(\mathfrak R\)을 가정한다. 그러면 \(\mathfrak R\)의 단위원 \(e\)는 동시에 \(\mathsf{P}\)의 단위원이고, 이 환은 \(\mathsf{P}\)-모듈이 된다.

\(\overline{\mathsf{P}}\)가 \(\mathsf{P}\)의 확대체일 때, \(\overline{\mathsf{P}}\)를 \(\mathfrak R\)에 환으로 첨가하여 생기는 확대환 \(\overline{\mathfrak R}=\mathfrak R[\overline{\mathsf{P}}]\)을 다음과 같이 정의한다.

먼저 \(\mathfrak R\)과 \(\overline{\mathsf{P}}\)의 집합론적 교집합 \([\mathfrak R,\overline{\mathsf{P}}]\)가 \(P\)이고, \(\mathfrak R\)과 \(\overline{\mathsf{P}}\)에서 \(\mathsf{P}\)에 속하지 않는 원소들 사이에는 연산관계가 없다고 가정한다. \(\overline{\mathsf{P}}\)를 \(\mathsf{P}\)의 동치인 확대체로 바꾸거나 \(\mathfrak R\)을 동치인 확대환으로 바꾸면 이 가정은 언제나 달성할 수 있다\srcfn{2)}{Steinitz에 따르면 \(\mathsf{P}\)의 두 확대체가 \(\mathsf{P}\)를 원소별로 자기 자신에 대응시키는 동형사상을 가질 때 “동치”라 한다. \(\mathsf{P}\)의 동치인 확대환도 이에 맞추어 정의한다. 새 기호를 도입하면 언제나 동치인 확대를 만들 수 있다. 동치인 확대를 자유롭게 쓰는 것은 이하에서 본질적이다.}.

\(\overline{\mathfrak R}\)의 원소는 형식적으로 만든 모든 선형결합
\[
   \bar c_{i_1}\gamma_{i_1}+\cdots+\bar c_{i_\sigma}\gamma_{i_\sigma}
\]
으로 정의한다. 여기서 \(\bar c_{i_\lambda}\)는 \(\overline{\mathsf{P}}\)의 모든 원소를, \(\gamma_{i_\lambda}\)는 \(\mathfrak R\)의 모든 원소를 두루 돈다\srcfn{1)}{일반적으로 \(\mathfrak R\) 또는 \(\overline{\mathfrak R}\)의 원소는 그리스 문자로, \(\mathsf{P}\) 또는 \(\overline{\mathsf{P}}\)의 원소는 라틴 문자로 나타낸다.}. 이 정의는 일의적이어야 한다. 곧 \(\gamma\)들을 \(\mathfrak R\)에서 그들과 같은 원소들로, \(\bar c\)들을 \(\overline{\mathsf{P}}\)에서 그들과 같은 원소들로 바꾸어도 \(\overline{\mathfrak R}\)의 원소 역시 같다고 정의한다.

합은 다음과 같이 일의적으로, 즉 한 항을 앞으로 정할 등식의 정의에서 그와 같은 것으로 바꾸어도 결과가 같도록 정의한다.
\[
   \sum \bar c_i\gamma_i+\sum \bar d_i\gamma_i
   =\sum(\bar c_i+\bar d_i)\gamma_i .
\]
(\(\bar c\)와 \(\bar d\) 가운데 0이 나타날 수 있으므로 형식적으로 같은 \(\gamma_i\)들이 나타난다.)

곱도 같은 뜻에서 일의적으로 정의한다.
\[
   \sum \bar c_i\gamma_i\cdot \sum \bar d_k\gamma_k
   =\sum \bar c_i\bar d_k\gamma_i\gamma_k .
\]
등식은 다음과 같이 정의한다. \(\gamma_{i_1},\ldots,\gamma_{i_\sigma}\)가 \(\mathsf{P}\)에 관하여 선형독립이면 두 원소
\[
   \bar c_{i_1}\gamma_{i_1}+\cdots+\bar c_{i_\sigma}\gamma_{i_\sigma}
   \quad\hbox{와}\quad
   \bar d_{i_1}\gamma_{i_1}+\cdots+\bar d_{i_\sigma}\gamma_{i_\sigma}
\]
는 \(\overline{\mathsf{P}}\)에서 각 \(\bar c_{i_\lambda}\)와 \(\bar d_{i_\lambda}\)가 같을 때 그리고 언제나 그때만 같다고 한다. 달리 말하면, \(\mathfrak R\)에서 \(\mathsf{P}\)에 관하여 선형독립인 원소들을 \(\overline{\mathsf{P}}\)에 관해서도 선형독립이라고 정의한다.

합과 곱의 정의로부터 \(\overline{\mathfrak R}\)의 모든 원소에 대한 등식의 정의가 정해진다. 실제로 모든 원소를 \(\mathfrak R\)의 선형독립인 원소들로 표현할 수 있기 때문이다. 관계
\[
   \mu=c_1\gamma_1+\cdots+c_s\gamma_s\quad\hbox{가 }\mathfrak R\hbox{에서 성립하면}
\]
\(\bar b=\bar b e\)를 곱하여 곱의 정의에 따라
\[
   \bar b\mu=\bar b e\cdot(c_1\gamma_1+\cdots+c_s\gamma_s)
      =\bar b c_1\gamma_1+\cdots+\bar b c_s\gamma_s
\]
를 얻고, 합의 정의에 따라 모든 원소를 \(\mathfrak R\)의 선형독립인 원소들로 나타낼 수 있다.

따라서 모든 원소를 \(\mathfrak R\)의 선형독립인 원소들로 나타내면 등식은 \(\overline{\mathsf{P}}\)에서의 계수 일치가 된다. 그러므로 \(\mathfrak R\)과 \(\overline{\mathfrak R}\)에 동시에 속하거나 \(\overline{\mathsf{P}}\)와 \(\overline{\mathfrak R}\)에 동시에 속하는 원소들에 대해서는 그곳에서 유효한 등식과 일치한다. 반면 \(\mathfrak R\)과 \(\overline{\mathsf{P}}\)에서 성립하는 등식은 교집합 가정과 그 부가조건 때문에 \(\overline{\mathfrak R}\)의 나머지 원소들에 대해서는 정의 때 요구한 일의성 이외의 아무것도 말하지 않는다\srcfn{2)}{이 가정이 없으면 \(\overline{\mathsf{P}}\)의 \(P\)에 속하지 않는 원소 \(\bar c\) 가운데 적어도 하나에 대하여 \(\bar c=\bar c e=\gamma\)가 성립한다. 그러면 \(\mathfrak R\)의 등식 정의에 따라 \(e\)와 \(\gamma\)는 \(\mathsf{P}\)에 관해서는 선형독립이지만 \(\overline{\mathsf{P}}\)에 관해서는 종속이다. \(\overline{\mathsf{P}}\)와 \(\mathfrak R\)의 원소 사이의 관계 때문에 위의 등식 정의는 불가능해진다. 예를 들어 기초체를 중간체로 바꿀 때 유한 확대체의 차수가 낮아지는 경우를 생각하면 된다. 반면 위의 확대는 새로운 영인자가 나타날 수 있다는 점으로 특징지어진다. 예컨대 유리계수 다항식영역 \(\mathsf{P}\)의 \(x^2+1\)에 따른 잉여류환은 영인자가 없는 환으로서 체 \(\mathsf{P}(\sqrt{-1})\)와 동형이지만, 계수를 \(\mathsf{P}(\sqrt{-1})\)로 확대하면 \((x+i)(x-i)=(x^2+1)\)이고 어느 인자도 \(x^2+1\)로 나누어지지 않는 데 대응하여 영인자가 나타난다. 모든 유한 수체에서도 이에 대응하는 것이 성립한다. 이것이 고차 복소수에 대하여 Dedekind가 처음 제시한 관점이다. § 6, 4의 끝도 보라.}. 같은 이유로 합과 곱의 일의성을 요구해도 새로운 관계가 생기지 않는다. 같은 원소들을 계수가 같은 것으로 가정할 수 있으므로 형식적으로 같은 합과 곱의 식이 나오기 때문이다.

이제 환 연산의 모든 공리가 충족됨을 곧바로 확인할 수 있다. 따라서 환 첨가로 생성되는 확대환 \(\mathfrak R[\overline{\mathsf{P}}]\)은 언제나 존재한다.

\textbf{3. \(\mathfrak R\)과 \(\overline{\mathfrak R}\)의 아이디얼.} \(\mathfrak a\)가 \(\mathfrak R\)의 아이디얼이면 모듈곱 \(\overline{\mathfrak R}\mathfrak a\)는 \(\overline{\mathfrak R}\)의 아이디얼 \(\overline{\mathfrak a}\)가 되며, 이를 \(\mathfrak a\)의 확대아이디얼이라 한다. 동시에 \(\overline{\mathfrak a}\)는 \(\overline{\mathfrak R}\)에서 \(\mathfrak a\)로부터 유도된 아이디얼이다. 이는 모든 유한합 \(\sum\bar c_i a_i\)로 이루어지며, \(a_i\)는 \(\mathfrak a\)의 모든 원소를, \(\bar c_i\)는 \(\overline{\mathsf{P}}\)의 모든 원소를 두루 돈다.

\(\overline{\mathfrak b}\)가 \(\overline{\mathfrak R}\)의 아이디얼이면 교집합 \([\overline{\mathfrak b},\mathfrak R]\)은 \(\mathfrak R\)의 아이디얼 \(\mathfrak b\)가 되며, 이를 \(\overline{\mathfrak b}\)의 수축아이디얼이라 한다\srcfn{1)}{H. Grell의 명칭이다(주 6에 인용한 논문, 83쪽).}. \(\mathfrak R\)의 각 아이디얼 \(\mathfrak a\)는 어떤 아이디얼의 수축아이디얼이며, 바로 그 확대아이디얼 \(\overline{\mathfrak a}\)의 수축아이디얼이다. 실제로 \([\mathfrak R,\overline{\mathfrak a}]\)의 각 원소 \(a\)는 \(\sum\bar c_i a_i\) 꼴이고, 여기서 \(a_i\)들은 \(\mathsf{P}\) 또는 \(\overline{\mathsf{P}}\)에 관하여 선형독립이라고 가정할 수 있다. 그러면 \(a,a_i\)는 \(\overline{\mathfrak R}\)에서 선형종속이고, 등식의 정의에 따라 \(\mathfrak R\)에서도 선형종속이다. 따라서 \(a\in\mathfrak a\)이다. 역포함도 성립하므로 \(\mathfrak a=[\mathfrak R,\overline{\mathfrak a}]\)이다.

\(\overline{\mathfrak p}\)가 \(\overline{\mathfrak R}\)의 소아이디얼이면 \(\mathfrak p=[\mathfrak R,\overline{\mathfrak p}]\)는 \(\mathfrak R\)의 소아이디얼이다. 실제로 제2동형정리\srcfn{2)}{\emph{아이디얼 이론} § 4, 3. 거기서는 같은 환의 아이디얼에 대해서만 정리를 말했지만, 같은 논증으로 \(\overline{\mathfrak R}\)의 임의의 아이디얼 \(\overline{\mathfrak a}\)에 대해서도 위 꼴로 성립한다. 환 준동형 \(\overline{\mathfrak R}\sim\overline{\mathfrak R}\,|\,\overline{\mathfrak a}\)는 부분환 \(\mathfrak R\)에 대하여 환 준동형 \(\mathfrak R\sim(\mathfrak R,\overline{\mathfrak a})\,|\,\overline{\mathfrak a}\)를 유도한다. 이 준동형에서는 \([\mathfrak R,\overline{\mathfrak a}]\)의 모든 원소와 오직 그 원소들만 영원소에 대응하므로 \((\mathfrak R,\overline{\mathfrak a})\,|\,\overline{\mathfrak a}\simeq \mathfrak R\,|\,[\mathfrak R,\overline{\mathfrak a}]\)이다. 양쪽의 류들은 \(\mathfrak R\)의 같은 원소들로 대표할 수 있다. 증명의 각 단계에서 원소를 그 원소가 대표하는 류에 대응시키기 때문이다.}에 따라 잉여류환 \(\mathfrak R\,|\,\mathfrak p\)는 \(\overline{\mathfrak R}\,|\,\overline{\mathfrak p}\)의 부분환과 동형이고, 따라서 영인자가 없는 환이다. 곧 다음 환 동형이 성립한다.
\[
   (\mathfrak R,\overline{\mathfrak p})/\overline{\mathfrak p}
     \simeq \mathfrak R/[\mathfrak R,\overline{\mathfrak p}],
\]
여기서 \((\mathfrak R,\overline{\mathfrak p})\)는 두 모듈 \(\mathfrak R\)과 \(\overline{\mathfrak p}\)의 합, 곧 최대공약수를 뜻한다.

\begin{center}
\textbf{§ 2. 유한 랭크 환. 직합표현}
\end{center}

\textbf{1. 유한 랭크 환.} \(\mathfrak A\)이 유한 \(\mathsf{P}\)-모듈이면 \(\mathfrak A\)은 \(P\)에 관하여 어떤 유한 랭크 \(k\)를 갖는다. 곧 \(\mathfrak A\)에는 \(P\)에 관하여 선형독립인 원소가 정확히 \(k\)개까지 있고, 그러한 \(k\)개 원소로 이루어진 모든 계는 \(\mathfrak A\)의 모듈기저가 된다. 유한 랭크인 한 \(\mathsf{P}\)-모듈 \(\mathfrak B\)가 같은 유한 랭크인 다른 모듈 \(\mathfrak A\)로 나누어지면, \(\mathfrak B\equiv0\,(\mathfrak A)\), 두 모듈은 동일하다.

이제부터 \(\mathfrak R\)은 \(P\)에 관하여 유한 랭크 \(n\)이라고 가정한다. 그러면 \(\mathfrak R\)의 각 \(\mathsf{P}\)-모듈도 유한 랭크이다. 따라서 \(\mathfrak R\)의 \(\mathsf{P}\)-모듈들의 진약수사슬 또는 진배수사슬에서는 각 모듈의 랭크가 바로 앞 모듈보다 각각 크거나 작아야 하므로, 사슬은 유한하게 끝난다. 특히 \(\mathfrak R\)의 모든 아이디얼계에 대해 “이중사슬정리”가 성립한다.

이로부터\srcfn{1)}{\emph{아이디얼 이론} § 7, 2.} \(\mathfrak R\)의 한 소아이디얼 \(\mathfrak p\)는 단위아이디얼과 다른 진약수를 갖지 않으며, 따라서 단위아이디얼과 다른 소아이디얼에 따른 잉여류환 \(\mathfrak R\,|\,\mathfrak p\)는 체가 된다. \(\mathfrak q\)가 \(\mathfrak p\)에 속하는 으뜸아이디얼, 곧 \(\mathfrak q\equiv0\,(\mathfrak p)\)이고 \(\mathfrak p^{\varrho}\equiv0\,(\mathfrak q)\)이면, 잉여류환 \(\mathfrak R\,|\,\mathfrak q\)는 으뜸환이 된다. 즉 이 환에서는 각 영인자의 어떤 거듭제곱, 여기서는 늦어도 \(\varrho\)제곱이 0이고, \(\mathfrak R\,|\,\mathfrak q\)에는 영인자와 단위만 있다. 마지막 명제는 \(\mathfrak p\)로 나누어지지 않는 각 아이디얼이 \(\mathfrak p\), 따라서 \(\mathfrak q\)와 서로소가 된다는 데서 곧바로 나온다.

이중사슬정리와 단위원의 존재로부터 다시\srcfn{2)}{\emph{아이디얼 이론} § 7, 3, 정리 V.} \(\mathfrak R\)의 각 아이디얼 \(\mathfrak a\)를 유한개의 쌍마다 서로소인 으뜸아이디얼의 곱으로 유일하게 나타낼 수 있다. 서로소성이 있으므로 이 곱은 최소공배수와 같고, 잉여류환 \(\mathfrak R\,|\,\mathfrak a\)를 으뜸환들의 직합으로 나타내는 것에 대응한다\srcfn{3)}{\emph{아이디얼 이론} § 4, 5. 현재의 정리에서 덧붙인 내용은 문헌의 여러 곳에 흩어져 있다.}. 곧 \(\mathfrak R\,|\,\mathfrak a\)는 이 환들의 최대공약수이고, \(\mathfrak R\,|\,\mathfrak a\)의 각 원소의 합표현은 유일하다.

\textbf{2. 유한 랭크 환을 으뜸환들의 직합으로 나타내기.} \(\mathfrak R\)의 영아이디얼 분해가
\[
   (0)=\mathfrak q_1\cdot\mathfrak q_2\cdots\mathfrak q_r
      =[\mathfrak q_1,\mathfrak q_2,\ldots,\mathfrak q_r]
\]
로 주어지고, 이에 대응하는 직합표현이
\[
   \mathfrak R=\mathfrak R_1+\mathfrak R_2+\cdots+\mathfrak R_r,
   \qquad
   \mathfrak R_i\mathfrak R_k=(0)\quad\hbox{for } i\ne k,
   \qquad
   \mathfrak R_i^2=\mathfrak R_i,
   \qquad
   \mathfrak R\mathfrak R_i=\mathfrak R_i
\]
라 하자. 이때 \(\mathfrak R_i\)는 \(\mathfrak R\)의 아이디얼이고,
\[
   \mathfrak R_i
   =\mathfrak q_1\cdots\mathfrak q_{i-1}\mathfrak q_{i+1}\cdots\mathfrak q_r
   =[\mathfrak q_1,\ldots,\mathfrak q_{i-1},\mathfrak q_{i+1},\ldots,\mathfrak q_r],
\]
\[
   \mathfrak q_i=\mathfrak R_1+\cdots+\mathfrak R_{i-1}+\mathfrak R_{i+1}+\cdots+\mathfrak R_r
\]
이며, \(\mathfrak R_i\)는 잉여류환 \(\mathfrak R\,|\,\mathfrak q_i\)와 동형이다. 구체적으로 \(\mathfrak R_i\)의 각 원소 \(\gamma_i\)에는 \(\mathfrak R\,|\,\mathfrak q_i\)에서 \(\gamma_i\)가 대표하는 류가 대응한다\srcfn{4)}{\emph{아이디얼 이론} § 4, 5.}. \(\mathfrak a\)가 \(\mathfrak R\)의 임의의 아이디얼이면 분배법칙에 따라
\[
   \mathfrak a=\mathfrak R\mathfrak a
   =\mathfrak R_1\mathfrak a+\cdots+\mathfrak R_r\mathfrak a
   =\mathfrak a_1+\cdots+\mathfrak a_r.
\]
\(\mathfrak a_i=\mathfrak R_i\mathfrak a_i=\mathfrak R\mathfrak a_i\)이므로 각 성분 \(\mathfrak a_i\)는 \(\mathfrak R_i\)와 \(\mathfrak R\) 양쪽 모두의 아이디얼이다. \(\mathfrak a_i\)는 \(\mathfrak R_i\)와 마찬가지로 유한 \(\mathsf{P}\)-모듈이고, 직합이므로 그 랭크들의 합은 각각 \(\mathfrak a\)와 \(\mathfrak R\)의 랭크이다.

단위원의 합표현이
\[
   e=\varepsilon_1+\cdots+\varepsilon_r,
   \qquad
   \varepsilon_i\varepsilon_k=0\quad\hbox{for } i\ne k,
   \qquad
   \varepsilon_i^2=\varepsilon_i,
   \qquad
   \varepsilon_i\equiv e\,(\mathfrak q_i)
\]
로 주어지면, \(\mathfrak R\)의 임의의 원소 \(\gamma\)는 유일하게
\[
   \gamma=\gamma e=\gamma\varepsilon_1+\cdots+\gamma\varepsilon_r
      =\gamma_1+\cdots+\gamma_r
\]
로 나타난다. 여기서 성분 \(\gamma_i=\gamma_i\varepsilon_i\)는 \(\mathfrak R_i\)의 원소이고, 동시에 \(\mathfrak R_i=\mathfrak R\varepsilon_i\)이다. “직교관계”로부터
\[
   \gamma\delta=(\gamma_1\varepsilon_1+\cdots+\gamma_r\varepsilon_r)
      (\delta_1\varepsilon_1+\cdots+\delta_r\varepsilon_r)
   =\gamma_1\delta_1\varepsilon_1+\cdots+\gamma_r\delta_r\varepsilon_r
\]
가 나온다. 따라서 \(\gamma_i\delta_i\)는 곱의 \(i\)번째 성분이고, \(\gamma_i+\delta_i\)는 합의 \(i\)번째 성분이다. 이는 \(\mathfrak R\)에서 \(\mathfrak R_i\)로 가는 준동형에서도 따라온다.

환 \(\mathfrak R_i\)는 \(\mathsf{P}\)와 동형인 부분체 \(\mathsf{P}_i=\mathsf{P}\varepsilon_i\)를 포함하고 \(\mathsf{P}_i\)에 관하여 유한 랭크이다. 이 랭크는 \(\mathfrak R_i\)가 \(\mathsf{P}\)-모듈로서 갖는 랭크와 같다. 실제로 \(c\in P\)이면 \(c\varepsilon_i\equiv ce\equiv c\,(\mathfrak q_i)\)이므로 \(c\ne0\)일 때 \(c\varepsilon_i\ne0\)이다. 따라서 \(P\)와 \(\mathsf{P}_i\) 사이의 준동형은 동형이 된다. 또한 \(\mathfrak R_i\)의 원소들 사이의 \(P\)에 관한 모든 선형종속관계는 \(\varepsilon_i\)를 곱하면 \(\mathsf{P}_i\)에 관한 관계로 옮겨가고, 역으로 \(\mathsf{P}_i=\mathsf{P}\varepsilon_i\)이므로 \(\mathsf{P}_i\)에 관한 모든 관계는 동시에 \(\mathsf{P}\)에 관한 관계이다. 따라서 두 랭크가 일치한다.

\textbf{3. 확대환의 직합표현.}
\[
   \mathfrak R=\mathfrak R_1+\cdots+\mathfrak R_r
\]
이면 \(\overline{\mathfrak R}\)에는
\[
   \overline{\mathfrak R}=\overline{\mathfrak R}_1+\cdots+\overline{\mathfrak R}_r
\]
라는 표현이 성립한다. 여기서 \(\overline{\mathfrak R}_i\)는 \(\overline{\mathfrak R}\)에서 \(\mathfrak R_i\)의 확대아이디얼이며, 동시에 \(\overline{\mathfrak R}\)에서 \(\mathfrak R_i\)로부터 유도된 환이다. 실제로 \((0)=\mathfrak q_1\cdots\mathfrak q_r\)에 \(\overline{\mathfrak R}\)을 곱하면 곱표현 \((\overline 0)=\overline{\mathfrak q}_1\cdots\overline{\mathfrak q}_r\)를 얻고, \(\overline{\mathfrak q}_i\)들은 다시 쌍마다 서로소이다. 따라서 \(\overline{\mathfrak R}\)의 직합표현이 나오며, 그 \(i\)번째 성분은
\[
   \overline{\mathfrak q}_1\cdots\overline{\mathfrak q}_{i-1}\,
   \overline{\mathfrak q}_{i+1}\cdots\overline{\mathfrak q}_r
\]
이므로 \(\mathfrak R_i\)의 확대아이디얼과 같다. 그런데 § 1, 3에 따라 이 확대아이디얼은 \(\overline{\mathsf{P}}_i=\overline{\mathsf{P}}\varepsilon_i\)에 관하여 \(\mathfrak R_i\)에서 만든 확대환 \(\mathfrak R_i[\overline{\mathsf{P}}_i]\)와 동일하다. 따라서 확대할 때에는 각 성분을 확대하는 것으로 충분하다. 다만 \(\overline{\mathfrak R}_i\)는 일반적으로 더 이상 으뜸환이 아니다.

\textbf{4. 제1종과 제2종 소아이디얼.} 1에 따라 단위아이디얼과 다른 각 소아이디얼 \(\mathfrak p\)에 따른 잉여류환 \(\mathfrak R\,|\,\mathfrak p\)는 체이다. 이 체에는 \(P\)와 동형인 부분체
\[
(P)=(P,\mathfrak p)\,|\,\mathfrak p\simeq P\,|\,[P,\mathfrak p]\simeq P
\]
가 있다. \(\mathfrak p\)는 \(\mathsf{P}\)의 0 아닌 원소를 포함할 수 없기 때문이다. 따라서 \((\mathsf{P})\)는 \(\mathsf{P}\)의 원소로 대표할 수 있는 모든 류와 오직 그 류들로 이루어진다\srcfn{1)}{제2동형정리에 관한 § 1, 3의 주를 보라.}. \(\mathfrak R\,|\,\mathfrak p\)는 \((\mathsf{P})\)에 관하여 유한 랭크이고, 따라서 유한 대수적 확대체이다. 이 확대가 제1종인지 제2종인지에 따라\srcfn{2)}{잘 알려진 Steinitz의 용어로, 대수적 확대는 각 원소가 어떤 확대체에서 서로 다른 일차인자들로 분해되는 소다항식의 영점이면 제1종이고, 그렇지 않으면 제2종이다.} \(\mathfrak p\)를 제1종 또는 제2종 소아이디얼이라 한다.

\begin{center}
\textbf{§ 3. 완전가약환}
\end{center}

\textbf{1.} 환 \(\mR\)은 그 영아이디얼의 곱표현(§ 2, 2)에서 모든 으뜸성분이 소아이디얼이 될 때 \emph{완전가약}이라 한다. § 2, 1과 2에 따라 이는 \(\mR\)을 유한개의 체의 직합으로 나타낼 수 있다는 것과 같다. 모든 소아이디얼이 제1종이면 완전가약환을 제1종 완전가약이라 하고(§ 2, 4), 그렇지 않으면 제2종 완전가약이라 한다.

\(\mathsf{P}\)가 완전체이면 제1종 완전가약성만 존재한다. 특히 \(\mathsf{P}\)가 대수적으로 닫힌 체이면 언제나 그러하다. 이하에서 \(\overline{\mathsf{P}}\)는 특히 \(\mathsf{P}\)로부터 얻은 대수적으로 닫힌 체를 뜻하고, \(\overline{\mathfrak R}=\mathfrak R[\overline{\mathsf{P}}]\)라 둔다.

\textbf{2. 정리.} \emph{환 \(\mR\)이 제1종 완전가약인 것은, 대수적으로 닫힌 \(\overline{\mathsf{P}}\)를 갖는 확대환 \(\bR=\mR[\overline{\mathsf{P}}]\)이 (제1종) 완전가약일 때 그리고 오직 그때이다.}

증명은 세 단계로 이루어진다.

\textbf{2a.} \(\bR\), 더 일반적으로 어떤 확대환이 완전가약이면 \(\mR\)도 제1종 또는 제2종 완전가약이다. 실제로 가정에 따라 \(\bR\)의 영아이디얼은 소아이디얼들의 최소공배수로
\[
        (\overline{0})=[\bp_1,\ldots,\bp_s]
\]
와 같이 나타난다. \(\mR\)과의 교집합을 취하면 영아이디얼의 표현
\[
        (0)=[[\mR,\bp_1],\ldots,[\mR,\bp_s]]
\]
을 얻고, § 1, 3에 따라 \([\mR,\bp_i]\)들은 소아이디얼이다.

\textbf{2b.} \(\mR\)이 제1종 완전가약이면 \(\bR\), 더 일반적으로 \(\mR\)과 \(\bR\) 사이의 모든 중간환도 제1종 완전가약이다.

§ 2, 3에 따라 \(\mR\)의 직합표현에는 \(\bR\)의 직합표현이 대응하고, 각 성분 \(\bR_i\)는 확대환 \(\mR_i[\overline{\mathsf{P}}_i]\)가 된다. 여기서 \(\overline{\mathsf{P}}_i\)는 성분 \(\mR_i\)의 \(P\)와 동형인 부분체 \(\mathsf{P}_i\)의 확대체이고 \(\overline{\mathsf{P}}\)와 동형이다. 따라서 각 성분 \(\mR_i\)만 고찰하면 충분하다. 일반성을 잃지 않고 \(\mR\)이 체이고, 그 부분체 \(\mathsf{P}\)의 유한 제1종 확대체라고 가정할 수 있다. 그러면 \(\mR\)은 제1종 원시원소 하나를 \(\mathsf{P}\)에 첨가하여 생성된다. 달리 말하면 \(\mR\)은 한 부정원 \(x\)의 다항식환 \(\mathsf{P}[x]\)에서 제1종 소다항식 \(f(x)\)에 따른 잉여류환 \(\mathsf{P}[x]\,|\,f(x)\)와 동형이다. \(\Omega\)가 \(\mathsf{P}\)의 확대체이면 \(\Omega[x]\,|\,f(x)\)는 \(\mR[\Omega]\)와 동형이다. 실제로 \(\mathsf{P}[x]\)에서 \(\Omega[x]\)로 옮겨도 잉여류환의 랭크는 \(f(x)\)의 차수와 같게 유지되므로, 이것은 정확히 § 1, 2에서 준 확대환의 구성이다. \(\Omega[x]\)에서 \(f(x)\)는 쌍마다 서로소인 제1종 소다항식들로 분해되고, 각각은 \(\Omega[x]\)에서 소아이디얼을 생성한다. 이는 \(\Omega[x]\,|\,f(x)\), 따라서 동형에 의해 \(\mR[\Omega]\)가 제1종 완전가약임을 뜻한다. 모든 확대체 \(\Omega\), 특히 대수적으로 닫힌 체 \(\overline{\mathsf{P}}\)에 대해 이것이 성립하므로 2b가 증명된다.

\textbf{2c.} \(\mR\)이 제2종 완전가약이면 \(\bR\)은 완전가약이 아니다.

2b에서와 마찬가지로 \(\mR\)을 체, 더 정확히는 \(\mathsf{P}\)의 유한 제2종 확대체라고 가정하면 충분하다. \(\bR\)을 으뜸환들의 직합으로 나타낼 때 적어도 하나의 진으뜸성분이 나타남을 보여야 한다. 이를 위해서는 \(\bR\)의 어떤 0 아닌 원소의 한 거듭제곱이 0임을 보이면 충분하다. 실제로 \(\bar\beta\ne0\)이면 적어도 한 성분 \(\bar\beta_\lambda\ne0\)이고, \(\bar\beta^\sigma=0\)이면 \(\bar\beta_\lambda^\sigma=0\)이다(§ 2, 2, \(\bR\)에서 \(\bR_\lambda\)로 가는 준동형). 따라서 \(\bR_\lambda\)는 체가 아니다.

\(\gamma\)를 \(\mR\)의 지수 \(f\geqq1\)인 제2종 원소라 하고,
\[
        F(\gamma)=\gamma^{kp^f}+c_1\gamma^{(k-1)p^f}+\cdots+c_k=0
\]
을 \(\gamma\)가 \(\mathsf{P}\)에 관하여 만족하는 최소차수 방정식이라 하자. 여기서 \(p\)는 \(\mathsf{P}\)의 표수이다. \(\mR\)의 원소들
\[
        e,\gamma,\gamma^2,\ldots,\gamma^{kp^f-1}
\]
은 \(\mathsf{P}\)에 관하여 선형독립이고, 따라서 등식의 정의에 의해 \(\bR\)에서도 \(\overline{\mathsf{P}}\)에 관하여 선형독립이다. 이제
\[
        G(\gamma)=\gamma^{kp^{f-1}}+\sqrt[p]{c_1}\,\gamma^{(k-1)p^{f-1}}+\cdots+\sqrt[p]{c_k}
\]
로 두면 \(G(\gamma)\)는 \(\bR\)의 0 아닌 원소이다. 실제로 \(p>1\), \(f\geqq1\)이므로 \(kp^{f-1}<kp^f\)이고, 따라서 \(G(\gamma)\)에 나타나는 \(\gamma\)의 거듭제곱들 사이에는 종속관계가 있을 수 없다. 그러나 \((G(\gamma))^p=F(\gamma)=0\)이다. 이로써 2c가 증명된다.

2a와 2c에서 \(\bR\)의 제1종 완전가약성이 \(\mR\)의 제1종 완전가약성을 함의하고, 2b에서 그 역이 나온다. 이로써 2의 정리가 증명되었다.

\textbf{3. 따름정리.} \emph{\(\mR\)이 제1종 완전가약이면 \(\bR\)은 랭크 1인 환들, 곧 \(\overline{\mathsf{P}}\)와 동형인 체들의 직합이 된다. 랭크 1인 환들로의 이러한 분해는 이미 어떤 환 \(\mR[\Omega]\)에서 존재하며, 여기서 \(\Omega\)는 \(\mathsf{P}\)의 유한 확대체이다.}

실제로 2b에 따라 \(\bR\)은 체들의 직합이 되고, 이 체들은 \(\overline{\mathsf{P}}\)와 동형인 부분체들의 유한 확대이므로(§ 2, 2), \(\overline{\mathsf{P}}\)가 대수적으로 닫혀 있다는 데서 그 부분체들과 동일하다. 따라서
\[
        \bR=\overline{\mathsf{P}}\bar\varepsilon_1+\cdots+\overline{\mathsf{P}}\bar\varepsilon_n,
        \qquad e=\bar\varepsilon_1+\cdots+\bar\varepsilon_n
\]
을 얻는다. \(e\)의 성분 \(\bar\varepsilon_\lambda\)들은 동시에 \(\bR\)의 선형독립인 모듈기저를 이룬다.

§ 1, 2에 따라 \(\bar\varepsilon_\lambda\)들을 \(\mR\)의 원소 \(\gamma\)들의 \(\overline{\mathsf{P}}\)-계수 선형결합으로 나타내면, 나타나는 모든 계수의 계는 \(P\)의 어떤 유한 확대체 \(\Omega\)를 정한다. 따라서 \(\bar e_\lambda\in\mR[\Omega]\)이고 모듈기저의 성질로부터
\[
        \mR[\Omega]=\Omega\bar e_1+\cdots+\Omega\bar e_n
\]
을 얻는다. 그러므로 이미 \(\mR[\Omega]\)에서 랭크 1인 환들로의 분해가 성립하고, \(\mR[\Omega]\)와 \(\bR\) 사이의 각 중간환에서는 분해 불가능한 성분들이 \(\Omega\bar\varepsilon_\lambda\)의 확대환으로 생긴다\srcfn{1)}{그러한 최소 확대체 \(\Omega\)는 \(\mR\)의 성분 \(\mR_i=\mR\varepsilon_i\)가 정하는 갈루아체들의 합성체로 특징지을 수 있다. 실제로 2b의 \(f_i(x)\)를 \(\mathsf{P}[x]\,|\,f_i(x)\)가 각각 성분 \(\mR_i\)와 동형이 되도록 택하자. \(\Omega^{(i)}\)를 \(f_i(x)\)를 일차인자들로 분해하기에 꼭 충분한, \(\overline{\mathsf{P}}\) 안의 갈루아 확대체라 하고, \(\Omega\)를 모든 \(\Omega^{(i)}\)의 합성체라 하자. \(f_i(x)\)가 일차인자들로 분해되는 데에는 \(\Omega^{(i)}[x]\,|\,f_i(x)\)를 랭크 1인 체들의 직합으로 나타내는 것이 대응하고, \(\Omega^{(i)}\)는 이것이 가능한 최소 확대체이다. 따라서 동형에 의해 \(\mR_i[\Omega^{(i)}\varepsilon_i]\), 나아가 \(\mR[\Omega]\)에도 같은 사실이 성립한다.}.

\begin{center}
\textbf{§ 4. 행렬표현,}\\
\textbf{유한 랭크 환의 대각합과 판별식}
\end{center}

수체와 함수체의 오더도 포함하기 위하여 이 절과 다음 절에서는 기초환에 대한 가정을 지금까지보다 조금 일반적으로 둔다. 본질적으로 알려진 사실들을 일반적이고 통일된 꼴로 정식화하려는 것이다.

\emph{가정:} \(\mR\)은 영인자가 없는 환 \(\Sigma\)의 확대환이고, \(\mR\)의 단위원은 이미 \(\Sigma\)에 속한다고 하자. \(\mR\)의 각 \(\Sigma\)-모듈, 특히 \(\mR\) 자체에는 \(\Sigma\)에 관하여 선형독립인 원소들로 이루어진 유한 모듈기저가 적어도 하나 존재한다고 하자. 그러면 \(\beta_1,\ldots,\beta_s\)가 \(\alpha_1,\ldots,\alpha_s\)와 같은 \(\Sigma\)-모듈의 선형독립인 모듈기저인 것은 그 사이의 변환행렬식이 \(\Sigma\)의 단위일 때 그리고 오직 그때이다.

이 가정은 지금까지 고찰한 유한 랭크 환에서 분명히 충족되며, 이때 \(\Sigma\)는 체 \(\mathsf{P}\)이다. 수체와 함수체의 오더에서도 충족되며, 이때 \(\Sigma\)는 각각 유리정수환 또는 함수식영역이다\srcfn{2)}{\emph{아이디얼 이론} § 3을 보라.}.

\textbf{1. \(\mR\)과 준동형인 행렬환.} \(\alpha_1,\ldots,\alpha_s\)가 아이디얼 \(\ma\)의 선형독립인 \(\Sigma\)-모듈기저이고 \(\gamma\)가 \(\mR\)의 임의의 원소라 하자. 그러면 \(\gamma\alpha_i\)도 \(\ma\)에 속하므로 \(\Sigma\)의 계수\srcfn{3)}{\(\Sigma\)의 원소는 라틴 문자로 나타낸다.}를 갖는 등식
\[
\gamma\alpha_i=c_{1i}\alpha_1+\cdots+c_{si}\alpha_s,
\qquad (i=1,\ldots,s)
\]
들이 성립한다. \((\tau_1,\ldots,\tau_s)\)가 이 원소들로 이루어진 한 행짜리 행렬을 뜻한다고 하면 행렬꼴로는
\[
(\gamma\alpha_1,\ldots,\gamma\alpha_s)
=(\alpha_1,\ldots,\alpha_s)
\begin{pmatrix}
 c_{11}&\cdots&c_{1s}\\
 \vdots&&\vdots\\
 c_{s1}&\cdots&c_{ss}
\end{pmatrix}
=(\alpha_1,\ldots,\alpha_s)C
\]
이다.

\emph{\(\gamma\)가 \(\mR\)의 모든 원소를 두루 돌면 기저 \(\alpha_i\)에 따라 대응하는 행렬 \(C\)들의 전체는 준동형 행렬환 \(R_\alpha\)를 이룬다. \(R_\alpha\)는 잉여류환 \(\mR\,|\,((0):\ma)\)와 동형이며, 여기서 \((0):\ma\)는 영아이디얼을 \(\ma\)로 나눈 아이디얼 몫이다.} 실제로 차 \(\beta-\gamma\)에는 분명히 \(B-C\)가 대응한다. 곱에 대해서도 같다. 곧
\[
(\beta\gamma\alpha_1,\ldots,\beta\gamma\alpha_s)
=(\beta\alpha_1,\ldots,\beta\alpha_s)C
=(\alpha_1,\ldots,\alpha_s)BC\srcfnmark{1)}
\]
\srcfntext{1)}{표기법은 \(\Sigma\)가 모든 원소와 교환한다고 가정한 비가환환 \(\mR\)에서도 그대로 유지되도록 배열하였다.}
특히 \(\Sigma\)의 원소 \(c\)에는 대각행렬 \(cE\)가 대응한다. 여기서 \(E\)는 단위원 \(e\)에 대응하는 행렬 \((\begin{smallmatrix}e&&\\&\ddots&\\&&e\end{smallmatrix})\)이다. 또한 \(\gamma\ma\)가 영아이디얼인 원소 \(\gamma\)들과 오직 그 원소들에 영행렬이 대응하므로 위의 동형을 얻는다.

\emph{\(\ma\)의 다른 기저 \(\bar\alpha\)는 \(R_\alpha\)와 동형이고 더 정확히 동치인 행렬환 \(R_{\bar\alpha}=P^{-1}R_\alpha P\)를 생성한다.} \(\bar\alpha_1,\ldots,\bar\alpha_s\)가 또 하나의 선형독립인 \(\Sigma\)-모듈기저이면
\[
(\bar a_1,\ldots,\bar a_s)=(a_1,\ldots,a_s)P
\]
이고, 행렬 \(P\)의 행렬식은 \(\Sigma\)의 단위이다. 따라서
\[
(\gamma\bar\alpha_1,\ldots,\gamma\bar\alpha_s)
=(\alpha_1,\ldots,\alpha_s)CP
=(\bar \alpha_1,\ldots,\bar \alpha_s)P^{-1}CP\srcfnmark{1)},
\]
곧 \(R_{\bar \alpha}=P^{-1}R_\alpha P\)이다.

\textbf{2. 아이디얼류와 표현류.} \(\mR\)의 두 아이디얼 \(\ma,\mb\)를 \(\mR\)-모듈로 보아 동형일 때 같은 아이디얼류에 속한다고 한다. 곧 두 아이디얼의 원소들을 일대일 대응시켜 차와 같은 \(\mR\) 원소에 의한 곱셈이 보존되며, \(\alpha\simeq\beta\)이면 임의의 \(\gamma\in\mR\)에 대해 언제나 \(\gamma\alpha\simeq\gamma\beta\)인 경우이다.

1의 행렬표현으로 생성한 두 환 \(R_\ma,R_{\bar a}\)가 동치일 때, 곧 어떤 유니모듈러 행렬 \(P\)가 있어 \(R_{\bar\alpha}=P^{-1}R_\alpha P\)일 때 같은 표현류에 속한다고 한다\srcfn{2)}{특수한 경우에 해결된 다음 문제는 여기서 다루지 않는다. \(\mR\)과 준동형인 모든 행렬환을 1에 따라 \(\mR\)의 한 아이디얼로 생성할 수 있는가, 필요하다면 선형종속인 기저를 허용하면 가능한가? 표현류는 정의상 1에 따라 \(\mR\)의 아이디얼이 생성하는 행렬환들로만 이루어진다.}.

\emph{아이디얼류와 표현류는 일대일로 대응한다.} 1에서 같은 아이디얼의 서로 다른 모듈기저들이 같은 표현류의 환 \(R_\ma,R_{\bar \alpha}\)를 생성함을 보였다. 더 나아가 같은 아이디얼류의 서로 다른 아이디얼 \(\ma,\mb\)도 같은 표현류의 환을 생성한다. 실제로 \(\beta_1,\ldots,\beta_s\)를 동형에 따라 \(\ma\)의 기저 \(\alpha_1,\ldots,\alpha_s\)에 대응하는 \(\mb\)의 기저라 하자. 그러면 \(\gamma\alpha_i\)와 \(\gamma\beta_i\)가 대응하고, \(c_{1i}\alpha_1+\cdots+c_{si}\alpha_s\)와 \(c_{1i}\beta_1+\cdots+c_{si}\beta_s\)도 대응하므로 환 \(R_\alpha\)와 \(R_\beta\)는 동일하다.

역으로 \(R_\alpha\)와 \(R_{\bar\beta}\)가 같은 표현류에 속하여 \(R_{\bar\beta}=P^{-1}R_\alpha P\)라 하자. 여기서 \(\alpha_i\)는 \(\ma\)의 기저이고 \(\bar\beta_i\)는 \(\mb\)의 기저이다. 그러면
\[
(\beta_1,\ldots,\beta_s)=(\bar\beta_1,\ldots,\bar\beta_s)P^{-1}
\]
도 \(\mb\)의 기저이고 \(R_\beta=PR_{\bar\beta}P^{-1}=R_\alpha\)이다. 이제 \(\ma\)의 각 원소 \(\alpha=c_1\alpha_1+\cdots+c_s\alpha_s\)에 \(\mb\)의 원소 \(\beta=c_1\beta_1+\cdots+c_s\beta_s\)를 대응시키면, 이 대응은 일대일이고 합과 차를 보존한다. \(R_\alpha=R_\beta\)에서 \(\gamma\alpha_i\)와 \(\gamma\beta_i\)도 대응한다. 이들은 각각 \(\alpha_i\)와 \(\beta_i\)의 같은 선형형식으로 표현되기 때문이다. 따라서 일반적으로 \(\gamma\alpha\)와 \(\gamma\beta\)가 대응하고, \(\ma,\mb\)는 같은 아이디얼류에 속한다. 단위아이디얼의 류를 \emph{주류}라 한다. 따라서 \(\mR\)의 모든 기저는 주류를 표현류로 낳는다.

\textbf{3. 한 류에 관한 원소의 대각합.} \(C\)가 \(R_\alpha\)에서 \(\mR\)의 원소 \(\gamma\)에 대응하는 행렬이라 하자. \(R_\alpha\)에서 \(R_{\bar\alpha}\)로 옮겨가는 것을 보면 행렬식 \(|C|\), 더 일반적으로 \(t\)가 부정원일 때 \(|tE-C|\)는 표현류의 불변량이다. 2에 따라 이는 동시에 아이디얼류의 불변량, 줄여서 \emph{류불변량}이다. 따라서 \(|tE-C|\)의 \(t\)의 각 계수는 류불변량이다. 특히 \((-t^{s-1})\)의 계수를 \emph{류에 관한 \(\gamma\)의 대각합}이라 하고
\[
S_{(\ma)}(\gamma)=c_{11}+c_{22}+\cdots+c_{ss}
\]
로 쓴다. \(t\)가 없는 계수 \(\pm|C|\)는 \emph{류에 관한 \(\gamma\)의 노름}이라 한다.

\emph{고정한 류 \((\ma)\)에 관한 대각합 \(S_{(\ma)}(\gamma)\)들은 \(\Sigma\)-모듈을 이루며, \(\Sigma\)-모듈로 본 \(\mR\)은 이 모듈에 준동형이다.} 실제로 1에서 증명한 \(\mR\)에서 \(R_\alpha\)로의 환 준동형에 따라 \((\beta-\gamma)\sim(B-C)\), \(c\beta\sim cEB=cB\)이다. 따라서 대각합에 대해서는
\[
S_{(\ma)}(\beta-\gamma)=S_{(\ma)}(\beta)-S_{(\ma)}(\gamma),
\qquad S_{(\ma)}(c\beta)=cS_{(\ma)}(\beta)
\]
이고, 모듈 성질과 준동형성이 나온다.

\emph{주의: \(\mR\)이 영인자가 없는 환이면 원소의 대각합은 류와 무관하게 정의된다. 따라서 단순히 대각합이라고 할 수 있으며, 노름도 마찬가지이다.} 실제로 영인자가 없는 환에서는 영아이디얼이 아닌 모든 아이디얼의 랭크가 환의 랭크와 같다. 따라서 한 아이디얼의 각 기저는 몫체에서 동시에 단위아이디얼의 기저이다. 모든 대각합과 노름은 몫체에서 동시에 단위아이디얼로 생성되므로 \(\mR\)에서 류와 무관하다.

\textbf{4. 한 류에 관한 아이디얼의 판별식.} 판별식 자체가 아니라 그 판별식이 \(\Sigma\)에서 생성하는 주아이디얼이 아이디얼과 류의 불변량이 된다.

\(\varrho_1,\ldots,\varrho_r\)와 \(\bar\varrho_1,\ldots,\bar\varrho_r\)가 \(\mR\)의 한 아이디얼 \(\mathfrak r\)의 두 선형독립인 \(\Sigma\)-모듈기저라 하자. \emph{행렬식 \(|S_{(\ma)}(\varrho_i\varrho_k)|\)와 \(|S_{(\ma)}(\bar\varrho_i\bar\varrho_k)|\)는 \(\Sigma\)의 한 단위의 제곱을 인자로 한 차이만 가진다.} 실제로 이는 두 행렬식 \(|(\varrho_i\varrho_k)|\)와 \(|(\bar\varrho_i\bar\varrho_k)|\)에 대해 성립하며, 두 행렬식은 \(\varrho\)와 \(\bar\varrho\) 사이의 변환행렬식의 제곱만큼 다르다. 3에서 증명한 모듈 준동형성 때문에 \(S_{(\ma)}(\varrho_i\varrho_k)\)와 \(S_{(\ma)}(\bar\varrho_i\bar\varrho_k)\) 사이에는 \(\varrho_i\varrho_k\)와 \(\bar\varrho_i\bar\varrho_k\) 사이와 같은 선형관계가 성립한다. 따라서 동변변환에 의해 행렬식은 같은 인자를 얻는다.

\(\Sigma\)의 단위의 제곱까지 정해지는 각 행렬식 \(|S_{(\ma)}(\varrho_i\varrho_k)|\)을 \emph{류 \((\ma)\)에 관한 \(\mathfrak r\)의 판별식}이라 하고, 그것이 \(\Sigma\)에서 생성하는 주아이디얼을 \emph{판별식아이디얼}이라 한다.

\emph{주의: \(\mR\)이 영인자가 없는 환이면 아이디얼의 판별식은 바탕으로 삼은 류와 무관하다.} 이는 3에 따라 행렬식의 각 원소에 대해 무관성이 성립하기 때문이다.

\begin{center}
\textbf{§ 5. 류의 직합}
\end{center}

\textbf{1. 아이디얼류 또는 표현류의 직합.} 아이디얼 \(\ma\)가 직합 \(\ma=\mb+\mc\)이면 동형성에 따라 그 류의 모든 아이디얼 \(\bar\ma\)도 직합 \(\bar\ma=\bar\mb+\bar\mc\)이 된다. 이때 \(\mR\)-모듈로 본 \(\bar\mb,\bar\mc\)는 각각 \(\mb,\mc\)와 동형이다. 따라서 \emph{아이디얼류의 직합}과 \emph{성분류}를 말할 수 있다.

한 표현류의 어떤 환 \(R_\alpha\)가 \(R_\alpha\)의 아이디얼이기도 한 부분환들의 직합이면, 동형에 의해 그 표현류의 모든 환 \(R_{\bar\alpha}\)에 대해서도 같고, 성분들은 동치인 환들이 된다. 따라서 \emph{표현류의 직합}과 \emph{성분류}를 말할 수 있다.

\emph{아이디얼류의 직합에는 표현류의 직합이 대응한다. 이 뜻에서 류의 직합이라 하겠다.}\srcfn{1)}{표현류의 모든 직합을 아이디얼류의 직합으로 생성할 수 있는지가 어느 정도까지 성립하는지는 다시 다루지 않는다.} \(\beta_1,\ldots,\beta_l\)가 \(\mb\)의 기저이고 \(\gamma_1,\ldots,\gamma_m\)이 \(\mc\)의 기저라 하자. 직합이므로 \(\beta\)들과 \(\gamma\)들을 합친 것은 \(\ma\)의 선형독립인 기저 \(a\)가 된다. 이 기저가 생성하는 행렬환을 \(R_\alpha\)라 하면, \(\mR\)의 임의의 원소 \(\delta\)에 대응하는 \(R_\alpha\)의 행렬은
\[
\begin{pmatrix}D^{(\beta)}&0\\0&D^{(\gamma)}\end{pmatrix}
\]
꼴이다. \(R^{(\beta)}\)를 모든 행렬 \((\begin{smallmatrix}D^{(\beta)}&0\\0&0\end{smallmatrix})\)의 계로 정의하고 \(R^{(\gamma)}\)도 대응하게 정의하자. 그러면 \(R^{(\beta)},R^{(\gamma)}\)가 \(R_\alpha\)의 부분환이고 \(R_\alpha\)가 그들의 직합임을 보여야 한다. 동시에 \(R^{(\beta)},R^{(\gamma)}\)는 각각 \(\mb,\mc\)가 생성하는 환 \(R_\beta,R_\gamma\)와 동형이다.

실제로 준동형에 따라 \(R^{(\beta)}\)에는 \(\delta\mc=0\)인 \(\mR\)의 모든 원소 \(\delta\), 곧 아이디얼 몫 \((0):\mc\)의 원소들이 대응한다. 따라서 \(R^{(\beta)}\)는 준동형에 의해 \(R_\alpha\)의 아이디얼이고, \(R^{(\gamma)}\)도 마찬가지이다. \(R_\alpha\)는 이 아이디얼들의 합이며, 두 환의 곱 \(R^{(\beta)}R^{(\gamma)}\)가 0이므로 직합이다. 끝으로 행렬 \(D^{(\beta)}\)를 \(D^{(\beta)}\)와 영블록을 포함하는 \(R^{(\beta)}\)의 행렬에 대응시키면 \(R^{(\beta)}\)와 \(R_\beta\) 사이의 동형을 얻는다. 이 뜻에서 \(\ma\)의 아이디얼류의 성분류 \((\mb),(\mc)\)에는 \(R_\alpha\)의 성분류가 대응한다.

\textbf{2. 류의 직합에서 대각합과 판별식의 거동.} \(\ma=\mb+\mc\)이고 대각합 \(S_{(\ma)}(\delta)\)를 정할 때 행렬 \((\begin{smallmatrix}D^{(\beta)}&0\\0&D^{(\gamma)}\end{smallmatrix})\)을 바탕으로 삼으면, 1을 고려하여 다음을 읽어낼 수 있다.

\emph{한 류에 관하여 취한 원소의 대각합은 각 성분류에 관하여 취한 대각합들의 합이다.}

곱 \(\mb\mc\)가 0이므로 또한 다음을 읽어낼 수 있다.

\emph{\(\beta\in\mb\)이면 \(S_{(\ma)}(\beta)=S_{(\mb)}(\beta)\)이다. 곧 류 \((\ma)\)에 관하여 취한 \(\beta\)의 대각합은 류 \((\mb)\)에 관하여 취한 대각합과 같다.}

§ 4, 4를 함께 고려하면 다음이 따라온다.

\emph{류 \((\ma)\)에 관하여 취한 성분아이디얼 \(\mb\)의 판별식아이디얼은 성분류 \((\mb)\)에 관하여 취한 \(\mb\)의 판별식아이디얼과 같다.}

끝으로 다음을 얻는다.

\emph{류 \((\ma)\)에 관하여 취한 \(\ma\)의 판별식아이디얼은 각 성분류에 관하여 취한 성분들의 판별식아이디얼의 곱과 같다.} 실제로 판별식과 대각합을 만들 때 모두 직합에 대응하는 \(\beta,\gamma\)의 기저를 바탕으로 삼으면, \((\ma)\)에 관한 \(\ma\)의 판별식아이디얼의 한 기저는
\[
\left|
\begin{array}{cc}
S_{(\ma)}(\beta_i\beta_k)&0\\
0&S_{(\ma)}(\gamma_\mu\gamma_\nu)
\end{array}
\right|
=
\left|
\begin{array}{cc}
S_{(\mb)}(\beta_i\beta_k)&0\\
0&S_{(\mc)}(\gamma_\mu\gamma_\nu)
\end{array}
\right|
\quad
\left(
\begin{matrix}
 i,k=1,\ldots,l\\
\mu,\nu=1,\ldots,m
\end{matrix}
\right).
\]
이다.

\textbf{3. 확대환으로 이행.} \(\overline{\Sigma}\)가 \(\Sigma\)의 영인자 없는 확대환이고 \(\mR\)과 \(\overline{\Sigma}\)의 교집합이 \(\Sigma\)라고 하자. \(\overline{\Sigma}\)를 첨가하여 얻는 확대환 \(\bR=\mR[\overline{\Sigma}]\)는 § 1, 2와 같이 같은 랭크의 환으로 정의한다.

\(\ma,\mathfrak r\)가 \(\mR\)의 아이디얼이고 \(\bar\ma,\bar{\mathfrak r}\)가 \(\bR\)에서 그 확대아이디얼이면, \((\bar\ma)\)에 관하여 취한 \(\bar{\mathfrak r}\)의 판별식아이디얼도 \((\ma)\)에 관하여 취한 \(\mathfrak r\)의 판별식아이디얼을 \(\overline{\Sigma}\)로 확대한 아이디얼이다. 실제로 \(\ma\) 또는 \(\mathfrak r\)의 모든 기저는 \(\bar\ma\) 또는 \(\bar{\mathfrak r}\)의 \(\overline{\Sigma}\)-모듈기저이기도 하므로, \((\bar\ma)\)에 관한 \(\bar{\mathfrak r}\)의 판별식을 만들 때 \(\ma,\mathfrak r\)의 기저를 바탕으로 삼을 수 있다. 특히 주류 \((\bR)\)에 관한 \(\bR\)의 판별식아이디얼은 주류 \((\mR)\)에 관한 \(\mR\)의 판별식아이디얼의 확대아이디얼이다. 이렇게 정의한 \(\Sigma\), 각각 \(\overline{\Sigma}\)의 아이디얼을 줄여서 \emph{환의 판별식아이디얼}이라 하겠다.

\begin{center}
\textbf{§ 6. 제1종 완전가약성의 판별식 판정조건}
\end{center}

이제부터 다시 \(\mR\)을 체 \(\mathsf{P}\)의 확대환이라 하고 \(\overline{\mathsf{P}}\)를 \(\mathsf{P}\)의 확대체라 하자. 제1종 완전가약성의 필요충분조건으로 판별식이 0이 아님을 얻을 것이다. § 5, 3에 따라 \(\mR\)의 판별식아이디얼은 확대환 \(\bR=\mR[\overline{\mathsf{P}}]\)으로 정할 수 있다. § 5, 2에 따라서는 \(\bR\)을 덧셈적으로 이루는 으뜸환들만 고찰하면 충분하다. 따라서 먼저 으뜸환의 판별식아이디얼을 고찰한다.

\textbf{1. 특수한 \(\mathsf{P}\)-모듈기저.} \(\ma,\mt\)가 \(\mR\)의 아이디얼이고 \(\ma\)가 \(\mt\)로 나누어지면, \(\ma\)의 임의의 기저 \(\alpha_1,\ldots,\alpha_s\)에 원소 \(\tau_{s+1},\ldots,\tau_t\)를 덧붙여 \(\mt\)의 선형독립인 모듈기저로 보충할 수 있다\srcfn{1)}{이 기저가 매개하는 임의의 원소 \(\gamma\)의 행렬표현은 \((\begin{smallmatrix}C_{\alpha}&0\\ C_1&C_2\end{smallmatrix})\) 꼴이고, 역으로 \(\mt\)가 매개하는 그러한 행렬표현에서는 \(\mt\)의 배수 \(\ma\)가 존재한다. 완전가약인 경우에는 분해 불가능한 성분들이 영아이디얼과 단위아이디얼 이외의 아이디얼을 갖지 않으므로, 이 성분아이디얼의 표현류에는 그러한 표현이 나타날 수 없다. § 5, 1에서 준 류의 직합 사이의 관계를 함께 고려하면 문헌에서 통상 쓰는 완전가약성의 정의에 이른다.}. 이는 \(\mathsf{P}\)에 관한 유한 랭크라는 가정의 직접적인 결과이다.

이제 \(\mR\)을 \emph{으뜸환}이라 하자. \(\mpideal\)을 영아이디얼에 속하는 소아이디얼, \(\varrho\)를 그 지수라 하자. 곧 \(\mpideal^\varrho=(0)\), \(\mpideal^{\varrho-1}\ne(0)\)이다. \(\mpideal^{\varrho-1}\)의 기저를 \(\mpideal^{\varrho-2}\)의 기저로 보충하고, 일반적으로 \(\mpideal^\lambda\)의 기저를 \(\mpideal^{\lambda-1}\)의 기저로 차례로 보충하면 \(\mR\)의 특수한 \(\mathsf{P}\)-모듈기저를 얻는다. 여기서 \(\mR\)은 \(\mpideal^0\)으로 센다. 이 기저로 다음을 보일 수 있다.

\emph{\(\mR\)이 으뜸환이면, 그에 속하는 소아이디얼 \(\mpideal\)로 나누어지는 모든 원소의 주류(§ 4, 2)에 관한 대각합은 0이다.} 그러한 특수한 \(P\)-모듈기저를
\[
\pi_{0,1},\ldots,\pi_{0,i_0};\quad
\pi_{1,1},\ldots,\pi_{1,i_1};\quad\ldots;\quad
\pi_{\varrho-1,1},\ldots,\pi_{\varrho-1,i_{\varrho-1}}
\]
이라 하자. 여기서
\[
\pi_{\lambda,j}\equiv0\;(\mpideal^\lambda),
\qquad \pi_{\lambda,j}\not\equiv0\;(\mpideal^{\lambda+1}).
\]
\(\pi\equiv0\;(\mpideal)\)이면 \(\pi\pi_{\lambda,j}\equiv0\;(\mpideal^{\lambda+1})\)이다. 따라서 이 기저에 따라 \(\pi\)에 대응하는 행렬의 대각선 성분은 모두 0이고, \(\pi\)의 대각합은 0이다.

\textbf{2. 대수적으로 닫힌 부분체 \(\mathsf{P}\)를 갖는 으뜸환의 판별식아이디얼.} \(\mR\)의 영아이디얼이 소아이디얼이고 \(\mathsf{P}\)가 대수적으로 닫혀 있으면 \(\mR\)은 랭크 1이고, 따라서 \(\mathsf{P}\)와 동일하다(§ 3, 3). 단위원 \(e\)를 모듈기저로 택할 수 있고
\[
S_{(\mR)}(e^2)=S_{(\mR)}(e)=e
\]
이므로 판별식아이디얼은 단위아이디얼과 같다.

대수적으로 닫힌 부분체 \(\mathsf{P}\)를 갖는 진으뜸환의 판별식아이디얼은 영아이디얼과 같음도 보일 수 있다. \(\mR\)에 1에서 준 특수한 모듈기저를 택하거나, 단지 \(\mpideal\)의 임의의 기저를 \(\mR\)의 기저로 보충하자. 잉여류환 \(\mR\,|\,\mpideal\)은 랭크 1인 환이므로 \(i_0=1\)이고, \(\pi_{0,1}\)로 단위원 \(e\)를 택할 수 있다. 그러면 \(e^2\) 이외의 모든 기저원소의 곱은 \(\mpideal\)로 나누어지고, 주류에 관한 그 대각합은 0이다. 따라서 판별식도 적어도 2행인 행렬식으로서 0이다. \(\mR\)이 진으뜸이라고 가정했으므로 \(e\)와 다른 기저원소가 적어도 하나 있어야 하고, 행렬식에서는 \(S_{(\mR)}(e^2)\) 이외의 모든 성분이 0이다.

\textbf{3. 정리.} \emph{부분체 \(P\)에 관하여 유한 랭크인 환 \(\mR\)이 제1종 완전가약인 것은 그 판별식아이디얼이 \(\mathsf{P}\)의 단위아이디얼일 때 그리고 오직 그때이다. 달리 말하면, 단위의 제곱까지 정해지는 판별식이 0이 아닐 때 그리고 오직 그때이다.}

§ 3, 2의 정리에 따라 제1종 완전가약성은 확대환 \(\bR=\mR[\overline{\mathsf{P}}]\)이 완전가약이라는 것, 곧 그 분해 불가능한 성분들이 \(\overline{\mathsf{P}}\)와 동형인 체라는 것과 같다. § 5, 3에 따라 \(\mR\)과 \(\bR\)의 판별식아이디얼은 동시에 영아이디얼이거나 동시에 단위아이디얼이다. § 5, 2에 따라 \(\bR\)의 판별식아이디얼은 각 성분에 관하여 취한 성분들의 판별식아이디얼의 곱이다. 이 개별 판별식아이디얼은 2에서 고찰한 것, 곧 성분들을 \(\bR\)의 아이디얼이 아니라 독립된 환으로 본 것에서, 매번 \(\overline{\mathsf{P}}\)와 동형인 체 \(\overline{\mathsf{P}}_i=\overline{\mathsf{P}}\bar\varepsilon_i\)를 \(\overline{\mathsf{P}}\) 자체로 바꾸어 얻는다. 그러므로 \(\bR\)이 완전가약이면 그 판별식아이디얼은 단위아이디얼들의 곱으로서 단위아이디얼이다. \(\bR\)이 완전가약이 아니면 적어도 한 성분의 판별식아이디얼이 영아이디얼이므로 \(\bR\)의 판별식아이디얼도 영아이디얼이다. 이로써 정리가 증명되었다.

\textbf{4. 제1종 완전가약인 경우 켤레원소에 의한 대각합·노름·판별식의 표현.} § 3, 3에 따라 제1종 완전가약일 때
\[
\mR[\Omega]=\Omega\bar\varepsilon_1+\cdots+\Omega\bar\varepsilon_n
\]
을 랭크 1인 환들, 곧 \(\Omega\)와 동형인 체들의 직합표현이라 하고, \(\Omega\)를 그러한 표현이 가능한 최소 확대체라 하자. 그러면 \(\mR\)의 \(n\)개 성분은 \(K_i\bar\varepsilon_i\)이고, 각 \(K_i\)는 \(\Omega\)의 부분체이며 \(\mR\)은 \(K_i\)에 준동형이다. 주류에 관한 대각합·노름·판별식을 \(K_i\)의 원소들로 표현하고자 한다.

\(\bar\varepsilon_1,\ldots,\bar\varepsilon_n\)을 모듈기저로 택하면 원소
\[
\gamma=c_1\bar\varepsilon_1+\cdots+c_n\bar\varepsilon_n
\]
에는 대각행렬 \(\operatorname{diag}(c_1,\ldots,c_n)\)이 대응한다. \(\gamma\in\mR\)이면 \(c_i\in K_i\)이다. 따라서 주류에 관한 대각합과 노름을 각각 \(S(\gamma),N(\gamma)\)로 쓰면
\[
S(\gamma)=c_1+\cdots+c_n,\qquad N(\gamma)=c_1\cdots c_n
\]
이다.

또 \(\alpha_1,\ldots,\alpha_n\)이 \(\mR\)의 원소들로 이루어진 모듈기저이면 판별식 \(|S(\alpha_i\alpha_k)|\)은 모든 \(K_\lambda\)의 원소들로 표현할 수 있다. 
\[
\alpha_i=\alpha_i^{(1)}\bar\varepsilon_1+\cdots+\alpha_i^{(n)}\bar\varepsilon_n
\]
이라 하자. 그러면
\[
\alpha_i \alpha_k=\alpha_i^{(1)}\alpha_k^{(1)}\bar\varepsilon_1+\cdots+\alpha_i^{(n)}\alpha_k^{(n)}\bar\varepsilon_n
\]
이고 \(S(\alpha_i\alpha_k)=\alpha_i^{(1)}\alpha_k^{(1)}+\cdots+\alpha_i^{(n)}\alpha_k^{(n)}\)이므로
\[
\left|S(\alpha_i\alpha_k)\right|
=
\left|\sum_{\lambda}\alpha_i^{(\lambda)}\alpha_k^{(\lambda)}\right|
=
\left|
\begin{matrix}
 \alpha_1^{(1)}&\cdots&\alpha_n^{(1)}\\
 \vdots&&\vdots\\
 \alpha_1^{(n)}&\cdots&\alpha_n^{(n)}
\end{matrix}
\right|^2.
\]

특히 \(\mR\) 자체가 체, 곧 부분체 \(\mathsf{P}\)의 제1종 확대체이면 \(\mR\)에서 \(K_i\)로 가는 준동형은 \(\mathsf{P}\)의 원소를 자기 자신에 대응시키는 동형이 된다. 실제로 \(\mathsf{P}\)의 성분은 \(\mathsf{P}\bar\varepsilon_i\)이다. 따라서 모든 \(K_i\)가 공통 포괄체 \(\Omega\) 안에 있으므로, 이들은 \(\mathsf{P}\)에 관하여 서로 켤레인 체이다. 또한 0 아닌 판별식 \(|S(\alpha_i \alpha_k)|\)의 위 행렬식 표현에서 보듯 \(n\)개의 동형은 서로 모두 다르다.

\emph{\(\mR\)이 체이면, \(K_1,\ldots,K_n\)으로의 직합표현은 갈루아 확대체 안에 있고 \(\mR\)과 동형이며 \(\mathsf{P}\) 위 \(n\)차인 \(n\)개의 켤레 확대체를 준다. 대각합·노름·판별식의 표현은 켤레원소에 의한 잘 알려진 표현으로 옮겨간다.} 여기서 \(\mR\) 자체는 \(K_i\)들과 동치일 뿐 갈루아 확대체 \(\Omega\) 안에 있는 체라고 가정한 것이 아님을 다시 강조한다\srcfn{1)}{§ 1, 2의 마지막 주를 보라.}.

\begin{center}
\textbf{§ 7. 대수적 수체 또는 함수체의 오더에 대한}\\
\textbf{판별식 정리}
\end{center}

\textbf{1.} 여기서 고찰하는 수체와 함수체는 모두, 정수적 원소를 지정한 다음과 같은 체 유형에 들어간다. 여러 부정원의 대수함수와 정수 계수 대수함수의 경우에는 함수식영역의 확대를 다룬다\srcfn{2)}{\emph{아이디얼 이론} § 3, 3과 4.}.

\(\mH\)를 \emph{영인자가 없고 단위원을 갖는 주아이디얼환}(유리정수환, 한 체의 계수를 갖는 한 부정원 다항식영역, 함수식영역)이라 하고, \(\mathfrak{K}\)를 그 \emph{몫체}라 하자. \(\mL\)을 \(\mathfrak{K}\)의 \emph{유한 제1종 확대체}, \(\mS\)를 \(\mL\)에서 \emph{\(\mH\)에 관하여 정수적인 모든 원소의 환}이라 하자\srcfn{3)}{한 환에 관하여 정수적인 원소의 정의는 \emph{아이디얼 이론} § 1, 3을 보라.}. \emph{\(\mS\)의 부분환 가운데 \(\mH\)를 포함하는 것은 모두 오더라 하고, \(\mS\) 자체는 정수환이라 한다.}

\(\mS\)의 모든 \(\mH\)-모듈, 특히 한 오더의 모든 아이디얼과 오더 자체는 \(\mH\)에 관한 \emph{선형독립인 모듈기저}를 갖는다\srcfn{4)}{\emph{아이디얼 이론} § 3, 1.}. \(\mL\)이 \(\mathfrak{K}\)에 관하여 차수 \(n\)이면 \(\mH\)에 관하여 랭크 \(n\)인 오더들만 고찰하면 충분하다. 더 낮은 랭크의 각 오더는 몫을 취하면 \(\mathfrak{K}\)과 \(\mL\) 사이에 그 랭크와 같은 차수의 중간체를 생성하며, 그 체에서도 모든 가정이 유지되기 때문이다.

\emph{따라서 각 오더 \(\mT\)는 § 4와 § 5의 더 일반적인 가정을 만족하는 환이며 영인자가 없는 환이다. 그러므로 오더의 판별식 \(D_\mT\)는 \(\mH\)의 한 단위의 제곱을 인자로 한 차이까지 유일하게 정의된다.} 이 판별식은 \(\mathfrak{K}\)의 한 단위를 인자로 한 차이까지, \(\mathfrak{K}\)-모듈로 본 체 \(\mL\)의 판별식과도 동일하다. 따라서 \(\mL\)을 제1종 확대라고 가정했으므로 \emph{0이 아니다}(§ 6, 3). 동시에 이 판별식은 § 6, 4에서 준 켤레기저원소들의 행렬식 제곱으로 표현된다.

\(\mT\)의 아이디얼 이론은 다음 정리로 주어진다\srcfn{5)}{\emph{아이디얼 이론} § 7, 3.}. \emph{\(\mT\)에서 영아이디얼이 아닌 모든 아이디얼은 유한개의 쌍마다 서로소인 으뜸아이디얼의 곱으로 유일하게 나타난다. 이 으뜸아이디얼에 속하는 소아이디얼은 단위아이디얼과 다른 진약수를 갖지 않는다.}

\textbf{2. 오더에서 한 체에 관하여 유한 랭크인 환으로 이행.} \(p\)가 \(\mH\)의 소원소라 하자. 곧 \(\mH\)에서 \(p\)로부터 유도된 주아이디얼 \(\mH p\)는 영아이디얼도 단위아이디얼도 아닌 \(\mH\)의 소아이디얼이다. 이에 따라 \(\mT p\)는 \(\mT\)에서 \(p\)로부터 유도된 주아이디얼을 뜻하게 하자\srcfn{1)}{이제부터 주아이디얼을 어느 환에서 고찰하는지를 동시에 나타내기 위해 이 표기법을 쓴다.}. \emph{잉여류환 \(\mR=\mT\,|\,\mT p\)는 체 \(\mH\,|\,\mH p\)와 동형인 한 부분체 \(P\)에 관하여 랭크 \(n\)인 환이 되고, \(\mT\)는 \(\mT\,|\,\mT p\)에 준동형이다.}

\(\mT\sim\mT\,|\,\mT p\)라는 준동형은 모든 잉여류환에 대해 일반적으로 성립한다. \(\mR\)이 \(\mH\,|\,\mH p\)와 동형인 부분체 \(\mathsf{P}\)를 포함한다는 것은 제2동형정리에서 나온다\srcfn{2)}{§ 1의 마지막 주를 보라.}. 실제로 부분환 \((\mH,\mT p)\,|\,\mT p\)는
\[
\mH/[\mH,\mT p]=\mH/\mH p
\]
와 동형이다.

끝으로 \(\mR\)은 \(P\)에 관해서도 랭크 \(n\)이다. \(\alpha_1,\ldots,\alpha_n\)을 \(\mH\)에 관한 \(\mT\)의 선형독립인 모듈기저라 하고, \((\alpha_1),\ldots,(\alpha_n)\)을 \(\mT p\)에 따라 \(a_i\)들이 정하는 잉여류라 하자. 그러면 \((\alpha_i)\)들 사이의 \(P\)에 관한 선형종속은 준동형을 거슬러 합동식
\[
c_1a_1+\cdots+c_na_n\equiv0\;(\mT p)
\]
에서 생기며, 이는 \(\gamma\in\mT\)에 대한 관계
\[
c_1a_1+\cdots+c_na_n=p\gamma
\]
에서 생긴다. \(\gamma\)를 \(\alpha_i\)의 기저로 나타낼 수 있고 그 표현이 유일하므로 각 \(c_i\)는 \(p\)로 나누어지고, 따라서 \(P\)의 원소 \((c_i)\)는 0이다. 이로써 \(\mathsf{P}\)에 관한 \((\alpha_i)\)의 선형독립성이 증명된다.

\emph{이 준동형에 따라 \(\mT\)의 판별식 \(D_\mT\)는 \(\mR\)의 판별식으로 옮겨가고, \(\mT p\)의 아이디얼 분해는 \(\mR\)의 영아이디얼 분해로 옮겨간다.} 방금 증명한 \((\alpha_i)\)의 선형독립성 때문에 \(\mR\)의 판별식은 이 기저에 의해 \(|S((\alpha_i)(\alpha_k))|\)로 만들 수 있다. 따라서 준동형에 의해 \(\alpha_i\)의 기저로 만든 \(\mT\)의 판별식 \(|S(\alpha_i\alpha_k)|\)에서 생긴다. 여기서 대각합의 정의에는 행렬표현을 바탕으로 삼았고, \(\mT\)에는 성립하지만 \(\mR\)에는 반드시 성립하지는 않는 켤레원소에 의한 표현을 바탕으로 삼지 않았다. 그러므로 대각합과 판별식은 환의 원소들로부터 환 연산만으로 생긴다.

\[
\mT p=\mq_1\cdots\mq_r=[\mq_1,\ldots,\mq_r]
\]
가 \(\mT\)에서 \(\mT p\)의 표현이면, 준동형에 따라 \(\mR\)에서는
\[
(0)=[(\mq_1),\ldots,(\mq_r)]
\]
로 옮겨간다. 제1동형정리\srcfn{3)}{\emph{아이디얼 이론} § 4, 3과 § 6(정리 II 다음)을 보라.}에 따라 \(\mR\,|\,(\mq_i)\)는 \(\mT\,|\,\mq_i\)와 동형이다. 따라서 \(\mq_i\)와 \((\mq_i)\)는 동시에 진으뜸아이디얼이거나 동시에 소아이디얼이다. 또한 준동형 때문에 \((\mq_i)\)들은 쌍마다 서로소이다.

\textbf{3. 판별식 정리.} \emph{\(\mH\)의 한 소원소 \(p\)가 한 오더 \(\mT\)의 판별식 \(D_\mT\)를 나누는 것은, \(\mT\)에서 주아이디얼 \(\mT p\)를 쌍마다 서로소인 으뜸성분들로 분해할 때 적어도 하나의 진으뜸성분 또는 적어도 하나의 제2종 소아이디얼이 나타날 때 그리고 오직 그때이다.}\srcfn{4)}{제2종 소아이디얼이 나타나는 예는 다음과 같다. \(\mH\)를 \(x\)의 모든 유리정수 계수 다항식으로부터 유도된 함수식영역, \(\mathfrak{K}\)를 그 몫체, \(\mL\)을 \(\sqrt{x}\)를 첨가한 확대체라 하자. \(\mH\)에 관한 모듈기저 \(1,\sqrt{x}\)를 갖는 오더 \(\mT\)를 고찰한다. 치환 \(x=y^2\)에서 보듯 이는 정수환이다. 판별식은 \(\left|\begin{smallmatrix}1&\sqrt{x}\\1&-\sqrt{x}\end{smallmatrix}\right|^2=2^2x\)이다. 주아이디얼 \(\mT2\)는 \(\mT\)에서 소아이디얼로 남지만 제2종이다. 실제로 모듈기저 \(1,\sqrt{x}\) 때문에 오더 \(\mT\)는 새 부정원 \(t\)에 대하여 잉여류환 \(\mH[t]\,|\,t^2-x\)와 동형이다. 따라서 \(\mT\,|\,\mT2=\mR\)은 \(\mathsf{P}[t]\,|\,t^2-x\)와 동형이며, 여기서 \(\mathsf{P}\)는 \(x\)와 함수식영역에 나타나는 부정원 \(u\)의 모든 유리함수 중 계수를 2로 줄인 것들의 체이다. \(t^2-x\)는 \(\mathsf{P}[t]\)의 소다항식이고 제2종 소다항식이므로 \(\mR\)은 체이며 \(\mathsf{P}\)의 제2종 확대체이다. 따라서 \(\mT2\)는 제2종 소아이디얼이다. 반면 \(\mT x=(\mT\sqrt{x})^2\)는 진으뜸이다. 일반적으로 \(\mR\)이 원소 하나를 \(\mathsf{P}\)에 첨가하여 생성된다는 것은 성립하지 않는다. 부정원 하나 \(x\)를 \(x,y\)로 바꾸기만 하면, \(1,\sqrt{x},\sqrt{y},\sqrt{xy}\)에서 유도된 오더, 다시 정수환인 이 오더에서 주아이디얼 \(\mT2\)는 제2종 소아이디얼이고 \(\mR\)은 원소 둘을 \(\mathsf{P}\)에 첨가해야 생성됨을 알 수 있다. 이 예들에 대해서는 Ostrowski의 앞의 글을 보라.}

\emph{잉여류체 \(\mH\,|\,\mH p\)가 완전체이면, \(D_\mT\)를 나누는 각 소원소 \(p\), 그리고 오직 그런 \(p\)는 적어도 하나의 진으뜸성분을 갖는다. \(\mH\,|\,\mH p\)가 완전체이고 \(\mT\)로 정수환 \(\mS\)를 택하면, \(p\)가 판별식을 나누는 것은 적어도 한 소아이디얼의 제곱으로 나누어질 때 그리고 오직 그때이다.}

실제로 \(\mT\)에서 주아이디얼 \(\mT p\)의 아이디얼 분해와 \(\mR=\mT\,|\,\mT p\)에서 영아이디얼의 분해가 대응하므로(2를 보라), \(\mT p\)에 진으뜸성분 또는 제2종 소아이디얼이 나타난다는 것은 \(\mR\)이 제1종 완전가약이 \emph{아니라}는 뜻이다. § 6, 3의 정리에 따라 이는 \(\mR\)의 판별식이 0일 때 그리고 오직 그때이며, 2의 준동형에 따라 이는 \(p\)가 \(D_\mT\)를 나눈다는 뜻이다.

\begin{center}
\textbf{§ 8. 상대체의 오더에 대한 판별식 정리}
\end{center}

출발환으로 주아이디얼환 \(\mH\) 대신 이미 한 수체나 함수체의 정수환, 더 일반적으로 \emph{곱셈환}\srcfn{1)}{“곱셈환”이라는 짧은 명칭은 W. Krull, \emph{Über Multiplikationsringe}, Heidelberger Berichte 1925, 제5논문 3호에서 빌렸다. 다만 Krull은 여기 나타나는 영인자 없는 환을 “정칙 곱셈환”이라 부른다.}, 곧 정수환의 아이디얼 이론이 성립하는 환을 택하면, 판별식 \(D_\mT\) 대신 \(\mH\)에 관한 \(\mT\)의 판별식아이디얼이 나타난다\srcfn{2)}{문헌에서 “상대수체”의 경우에는 정수환의 판별식아이디얼, 곧 “상대판별식”만 나타난다. Hilbert, \emph{Zahlbericht}, § 14와 Hecke, \emph{Theorie der algebraischen Zahlen}, § 38을 보라.}. 다음 고찰에 의해 판별식 정리는 완전히 옮겨간다.

\textbf{1.} \(\mH\)를 \emph{곱셈환}, 곧 영인자가 없고 단위원을 가지며, 영아이디얼도 단위아이디얼도 아닌 모든 아이디얼을 소아이디얼들의 거듭제곱의 곱으로 유일하게 나타낼 수 있고 그 소아이디얼들이 단위아이디얼과 다른 진약수를 갖지 않는 환이라 하자\srcfn{3)}{곱셈환의 공리계에 대해서는 \emph{아이디얼 이론}, 예컨대 서론을 보라.}. 또 \(\mathfrak{K},\mL,\mS,\mT\)를 § 7, 1과 같이 정의하자. 그러면 \(\mT\)에는 § 7, 1에서 말한 아이디얼 이론이 성립한다\srcfn{4)}{\emph{아이디얼 이론} § 3, 2와 § 7, 3.}.

특히 한 곱셈환 \(\mH\)가 영아이디얼과 단위아이디얼 이외에 소아이디얼을 단 하나만 가지면 \(\mH\)는 주아이디얼환이다. 실제로 \(\mpideal\)이 그 소아이디얼이라 하고
\[
p\equiv0\;(\mpideal),\qquad p\not\equiv0\;(\mpideal^2)
\]
인 \(p\)를 택하자. 가정에 따라 영아이디얼도 단위아이디얼도 아닌 모든 아이디얼은 \(\mpideal\)의 거듭제곱이므로, 주아이디얼 \(\mH p\)도 그러하다. \(p\not\equiv0\;(\mpideal^2)\)이므로 \(\mH p=\mpideal\)이어야 한다. 따라서 \(\mpideal\)의 모든 거듭제곱도 주아이디얼이고, 영아이디얼과 단위아이디얼 역시 주아이디얼이다.

\textbf{2. 대각합과 판별식아이디얼.} \(\mT\)의 각 원소 \(\gamma\)의 대각합은 \emph{\(\mathfrak{K}\)-모듈로 본 \(\mL\)의 원소 \(\gamma\)의 대각합}으로 정의한다. 따라서 \(\mL\)에서 \(\mathfrak{K}\)에 관하여 선형독립인 \(n\)개 원소의 모든 계를 대각합을 정의하는 행렬표현을 생성하는 기저로 택할 수 있다. \(\alpha_1,\ldots,\alpha_n\)이 그러한 계이면 \(\mL\)을 제1종 확대라고 가정했으므로 \(|S(\alpha_i\alpha_k)|\)는 \(\mathfrak{K}\)의 0 아닌 원소이며, 더 정확히는 켤레기저원소들의 행렬식 제곱이다(§ 6, 4). \(\mH\)가 \(\mathfrak{K}\)에서 정수적으로 닫혀 있다는 데서\srcfn{1)}{\emph{아이디얼 이론} § 9, 7.}, 모든 \(a_i\)가 \(\mS\)에 속하면 \(|S(\alpha_i\alpha_k)|\)는 \(\mH\)의 원소임이 따라온다.

\emph{판별식아이디얼:} \(\tau_1,\ldots,\tau_n\)이 \(\mT\)의 원소 \(n\)개로 이루어진 모든 계를 두루 돌게 하자. 각 계에 대해 행렬식 \(|S(\tau_i\tau_k)|\)을 만든다. 이는 \(\mH\)의 원소이고, \(\tau\)들이 선형독립이면 0이 아니다. \emph{이 모든 행렬식이 \(\mH\)에서 유도하는 아이디얼을 \(\mH\)에 관한 오더 \(\mT\)의 판별식아이디얼이라 한다.} 따라서 이는 영아이디얼이 아니다.

또 \(\beta_1,\ldots,\beta_m\), \(m\geqq n\), 이 \(\mT\)의 언제나 존재하는 \(\mH\)-모듈기저라 하자\srcfn{2)}{\emph{아이디얼 이론} § 9와 § 3, 1.}. 그러면 \(\beta\)들 가운데 각각 \(n\)개에 대응하는 유한개의 행렬식이 판별식아이디얼의 아이디얼 기저를 이룬다. 이는 \(\mT\)에서 그 대각합들의 계로 가는 모듈 준동형에서 나온다(§ 4, 3). 실제로 곱 \(\tau_i\tau_k\)들로 만든 행렬식 \(|(\tau_i\tau_k)|\)은 각각 \(n\)개의 \(\beta\)에 대응하는 유한개의 행렬식 \(|(\beta_i\beta_k)|\)의 \(\mH\)-계수 선형결합으로 나타난다. 따라서 \(|S(\tau_i\tau_k)|\)과 \(|S(\beta_i\beta_k)|\)에도 같은 사실이 성립한다. 특히 오더가 선형독립인 모듈기저 \(\alpha_1,\ldots,\alpha_n\)을 가지면, 행렬식 \(|S(\alpha_i\alpha_k)|\)을 판별식아이디얼의 기저로 택할 수 있다. 이 정의는 § 7, 1의 정의와 일치한다.

\textbf{3. 몫환으로 이행.}\srcfn{3)}{이 항에서 인용하는 모든 정리의 증명은 원문 83쪽의 주 6에서 인용한 H. Grell의 논문을 보라.} \(\ma\)가 \(\mT\)의 임의의 아이디얼이면, 몫환 \(\mT_\ma\)는 \(\mT\)의 몫체 \(\mL\)의 원소들 가운데 분모가 \(\mT\)에서 \(\ma\)와 서로소인 원소인 것들의 전체로 정의한다\srcfn{4)}{\(\mT\)의 소아이디얼들은 단위아이디얼과 다른 진약수를 갖지 않으므로, 여기서는 원소가 \(\ma\)에 대해 소라는 것과 아이디얼의 뜻에서 서로소라는 것이 같다.}. 환 \(\mT_\ma\)에서 영아이디얼과 단위아이디얼 이외의 소아이디얼은 \(\ma\)에 속하는 소아이디얼들의 확대아이디얼뿐이다. \(\mT_\ma\)의 영아이디얼·단위아이디얼이 아닌 모든 아이디얼과, 그에 속하는 소아이디얼들이 \(\ma\)에 속하는 소아이디얼들 가운데 포함되는 \(\mT\)의 아이디얼들 사이에는 일대일 대응이 있고 잉여류환들이 동형이다. 영아이디얼과 단위아이디얼도 각각 서로 일대일로 대응한다.

같은 고찰이 곱셈환 \(\mH\)에도 성립한다. 특히 \(\mpideal\)이 소아이디얼이면 \(\mH_\mpideal\)은 영아이디얼과 단위아이디얼 이외에 소아이디얼을 하나만 가진다. 잉여류환의 동형 때문에 \(\mH_\mpideal\)도 \(\mH\)와 마찬가지로 곱셈환이고, 1에 따라 주아이디얼환이다. \(\mpideal\)에서 유도된 소아이디얼의 기저는 \(\mH_\mpideal\)의 한 소원소 \(p\)가 된다. \(\mH\)의 임의의 아이디얼 \(\mc\)의 확대아이디얼은 \(\mc\)의 \(\mpideal\)에 속하는 으뜸성분이 생성한다. 따라서 \(\mpideal\)이 \(\mc\)를 \(\varrho\)제곱으로 나누는지 전혀 나누지 않는지에 따라 확대아이디얼은 \((p)^\varrho\) 또는 단위아이디얼과 같다. 이로부터 다시 다음이 나온다. \(\mH\)의 두 아이디얼 \(\mb,\mc\)의 확대가 \emph{모든} 몫환 \(\mH_\mpideal\)에서 일치하면 \(\mb,\mc\)는 동일하다. 실제로 두 아이디얼은 같은 소아이디얼 \(\mpideal\)들로 같은 거듭제곱만큼 나누어진다.

\textbf{4. 몫환으로 이행할 때의 판별식아이디얼.} \(\mpideal\)이 \(\mH\)의 소아이디얼이고 \(\mP=\mT\mpideal\)이 \(\mT\)에서 그 확대아이디얼이라 하자. \emph{몫환 \(\mT_\mP\)는 주아이디얼환 \(\mH_\mpideal\)에 관하여 선형독립인 원소들로 이루어진 모듈기저를 갖는다. 동시에 \(\mT_\mP\)는 \(\mH_\mpideal\)의 몫체 \(\mathfrak{K}\)의 확대체 \(\mL\) 안의 유한 \(\mH_\mpideal\)-오더이다.}\srcfn{1)}{증명은 H. Grell을 보라.} \emph{따라서 \(\mT_\mP\)에서는 \(\mH_\mpideal\)에 관하여 판별식 정리(§ 7, 3)가 성립한다.}

이때 판별식 \(D_{\mT_\mP}\)는 \(\mH\)에 관한 \(\mT\)의 판별식아이디얼을 \(\mH_\mpideal\)로 확대한 아이디얼의 기저를 이룬다. 실제로 언제나 가능하듯 \(\mT_\mP\)의 \(\mH_\mpideal\)-모듈기저로 \(\mT\)의 원소 \(\alpha_1,\ldots,\alpha_n\)을 택하자. 그러면 \(|S(\alpha_i \alpha_k)|\)는 판별식아이디얼에 속하고, 2에 따라 모든 행렬식 \(|S(\tau_i\tau_k)|\)은 \(\mH_\mpideal\)에서 \(|S(\alpha_i\alpha_k)|\)로 나누어진다\srcfn{2)}{여기서 \(\mH\)가 한 대수적 수체의 정수환이면 Hilbert가 정의한 상대판별식과 일치함을 얻는다. 이는 Hecke가 정의한 것과 일치한다고 알려져 있다. 상대체의 정수환의 \(\mH\)-모듈기저를 \(\omega_1,\ldots,\omega_m\)이라 하자. Hilbert는 특히 유리정수환에 관한 모듈기저를 택한다. Hilbert를 따라 각 \(n\)개 기저원소와 그 켤레원소들로 모든 행렬식을 만들자. “상대판별식”은 이 행렬식들 가운데 같은 것이든 다른 것이든 매번 둘을 곱한 것들이 \(\mH\)에서 생성하는 아이디얼이다. 따라서 상대판별식은 \(\mH_\mpideal\)에서 기저 \(\alpha_1,\ldots,\alpha_n\)과 그 켤레원소들의 행렬식 제곱, 곧 \(|S(\alpha_i\alpha_k)|\)으로 옮겨간다. 그러므로 상대판별식과 (상대 정수환의) 판별식아이디얼의 확대는 모든 몫환 \(\mH_\mpideal\)에서 일치하고, 3의 끝에 따라 두 아이디얼은 동일하다.}.

\textbf{5. 일반 판별식 정리.} \emph{한 곱셈환 \(\mH\)의 소아이디얼 \(\mpideal\)이 \(\mH\)에 관한 한 오더의 판별식아이디얼을 나누는 것은, \(\mT\)에서 \(\mT\mpideal\)을 쌍마다 서로소인 으뜸성분들로 분해할 때 적어도 하나의 진으뜸성분 또는 적어도 하나의 제2종 소아이디얼이 나타날 때 그리고 오직 그때이다.}

3과 4에 따라 \(\mH\)의 소아이디얼 \(\mpideal\)이 \(\mT\)의 판별식아이디얼을 나누는 것은 \(\mH_\mpideal\)의 소원소 \(p\)가 판별식 \(D_{\mT_\mP}\)를 나눌 때 그리고 오직 그때이다. 이는 잉여류환 \(\mT_\mP\,|\,\mT_\mP p\)가 제1종 완전가약이 아니라는 것과 같다. 3에 따라 이 잉여류환은 \(\mT\,|\,\mT\mpideal\)과 동형이므로, 이로써 일반 판별식 정리가 증명된다.

수체의 상대판별식으로 특수화하면 다음을 얻는다.

\emph{한 수체의 정수환의 소아이디얼 \(\mpideal\)이 한 확대체의 상대판별식을 나누는 것은, 확대체의 정수환에서 \(\mpideal\)이 한 소아이디얼의 제곱으로 나누어질 때 그리고 오직 그때이다.}

괴팅겐, 1926년 3월 30일.

\begin{center}
\rule{9em}{0.4pt}
\end{center}

\clearpage
\section*{논문 31에 대한 편집 주석}
\addcontentsline{toc}{section}{논문 31에 대한 편집 주석}
원문 103쪽, \S 8, 제4항 둘째 문단: 인쇄본의 제1항 참조를 여기서는 제2항 참조로 읽는다. $|S(a_i a_k)|$에 관한 필요한 명제는 원문 102쪽의 제2항 끝에 있다. 제3항에서 앞서 나온 제1항 참조는 그대로 옳다.
\smallskip
\noindent{\small AI가 제작한 교정 주석과 정정: OpenAI Codex, GPT-6 Astra, Ultra 추론 강도.}

\clearpage
\setcounter{footnote}{0}
\phantomsection
\addcontentsline{toc}{section}{32. 기약표현의 최소 분해체에 관하여}
%
%
%

%
%
\section*{32. R. 브라우어와 공동 저술: 기약표현의 최소 분해체에 관하여}
\emph{Sitz. Ber. d. Preuß. Akad. d. Wiss. 1927, S. 221--228.}

\begin{center}
{\Large\bfseries 기약표현의 최소 분해체에 관하여.}\\[1.0em]
쾨니히스베르크의 사강사(Privatdozent) Dr.\\
\textsc{리하르트 브라우어(Richard Brauer)}\\
및 괴팅겐의 Prof. Dr.\\
\textsc{에미 뇌터(Emmy Noether)}\\[0.8em]
\emph{(Schur가 제출.)}\\[0.5em]
\rule{4em}{0.4pt}
\end{center}

%
%
기초체 $P$ 위의 수체들 가운데 차수가 가장 작으면서, 그 안에서 유한군의 $P$ 위 기약표현 하나가 절대기약 성분들로 분해되는 수체를 찾는 문제는 I. Schur가 처음 다루었다.\footnote{Arithmetische Untersuchungen über endliche Gruppen, Sitzungsber. d. Berl. Ak. d. Wiss. 1906, S. 164. Beiträge zur Theorie der Gruppen linearer Substitutionen, Transact. of the Am. Math. Soc. Ser. 2, Vol. XV, 1909, S. 159. 둘째 논문에서는 그 결과를 완전가약 초복소수계로 확장한다. 하나의 표현은 그 모든 변환된(상사인) 표현과 함께 하나의 류로 묶는다.}

이제 $P$가 그러한 성분 하나의 지표를 포함한다고 하자. 그러면 I. Schur는 이들 체의 $P$에 대한 차수---곧 지수---가 모두 같은 류에 속하는 그 절대기약 성분들의 수와 일치한다는 것, 따라서 그 체 자체가 절대기약 성분 하나를 분리해 낸다는 요구만으로도 이미 결정된다는 것, 나아가 그러한 분해가 일어나는 임의의 체의 $P$에 대한 차수는 지수의 배수가 된다는 것을 보였다. $P$가 그 지표를 포함하지 않는 경우에도 이러한 체들을 모두 분해체라 부르기로 한다. I. Schur는 한 예를 통하여 최소 차수의 분해체들이 서로 동형일 필요는 없음을 보였다.

$P$가 관련 단순지표를 하나도 포함하지 않는다면, 분해체들은 그러한 단순지표 하나의 지표체와 그 지표체의 $P$에 관한 공액체들을 포함해야 한다. 서로 다른 공액 지표들에 대응하는 절대기약표현들은 서로 공액이기는 하지만, 분명히 서로 다른 류에 속한다. 여기서도 한 류에 속하는 절대기약표현들로 이루어진 한 계를 분리해 내는 데에는 대응하는 지표체와 해당 분해체들 가운데 하나면 충분하다.

%
%
이하에서는 분해체의 문제를 더 살펴본다. \srcspaced{최소 차수의 분해체는 그에 대응하는 비가환체로 특성화할 수 있고, 일반적인 분해체는 이 비가환체들과 불변적으로 결부된 양쪽 단순환으로 특성화할 수 있다.} 이어서 사원수체의 예를 통하여 \srcspaced{극소 분해체의 차수에는 상계가 없음}을 보인다. 여기서 어떤 분해체의 진부분체 가운데 어느 것도 분해체가 아닐 때, 그 분해체를 \srcspaced{극소 분해체}라고 한다. 이 비유계성은 본질적으로 하나의 수론적 존재 정리로부터 따르며, 그 정리의 증명은 바로 이어지는 논고에서 H. Hasse가 제시한다. 그러나 이 예는 또한---다분히 검증하는 성격을 띠는 방식으로---초등적인 계산에 의해서도 다룬다. 곧 범위는 더 좁지만 여기에는 충분한 존재 정리를 초등적으로 증명함으로써, 일반 이론과 그 수론적 존재 정리를 원용하지 않고도 차수들의 비유계성을 보인다.\footnote{이 예는 E. Noether가 고안한 것으로, 차수가 지수를 넘는 극소 분해체가 존재함을 보였던 R. Brauer의 한 예를 변형한 것이다. 이 예를 초등적으로 다룬 것은 R. Brauer이다. 분해체의 특성화는 두 저자의 연구에서 비롯되었지만, 그 밖의 면에서 두 사람은 서로 별개의 목표를 추구한다. E. Noether의 목표는 일반적인 유한성 조건 아래에서 가군론과 아이디얼론을 바탕으로 표현론을 구축하는 것이고, R. Brauer의 목표는 인자계를 이용하여 주어진 가환 완전체인 기초체 위의 유한 차수 비가환체를 구성하는 것이다. 최소 차수 분해체의 특성화에는 두 저자가 서로 독립적으로 도달하였다. R. Brauer는 완전체인 기초체의 경우에, E. Noether는 그러한 제한 없이 이에 도달하였다. 일반적인 분해체의 특성화는 비가환 매장의 가능성에 관한 E. Noether의 질문을 계기로 R. Brauer가 발견하였으며, E. Noether는 이 경우에도 기초 영역이 완전체여야 한다는 제한을 없앨 수 있었다.}

또한 이 예 역시 유한군, 곧 Schur의 예에서도 바탕이 된 사원수군의 분해체에 관한 것임을 덧붙여 둔다. 그곳에 제시된 표현들은 동시에 사원수체의 (동형인) 표현들이기 때문이다.

%
%
\begin{center}
\textbf{1. 분해체의 특성화.}
\end{center}

우선 기본개념을 간단히 정리해 둔다. $\mS$는 체 $\mathrm{P}_0$에 관한 초복소수계라 하자. $P$에서의 $\mS$의 표현 $\Gamma$란 $\mathrm{P}$의 원소를 성분으로 하는 행렬계로서, $\Gamma$가 $\mS$의 준동형상이 되고 $\mathrm{P}_0$의 원소 $\rho_0$에는 대각행렬 $\rho_0E$가 대응하는 것을 말한다. 따라서 $\mathrm{P}$는 $\mathrm{P}_0$를 포함해야 한다. 표현 $\Gamma$는 그 모든 변환표현 $\mathit{P}^{-1}\Gamma \mathit{P}$와 함께 하나의 표현류를 이룬다.

이와 같이 정의한 표현에는 유한군 $\mathfrak G$의 표현도 포함된다. 이에 대응하는 초복소수계 $\mS$, 곧 ``군환''은 모든 ``군수'', 즉 $\mathrm{P}_0$의 계수를 갖는 군 원소들의 모든 선형결합으로 이루어진다. 이때 본래의 군 원소들은 선형독립으로 간주하며, 곱셈은 군의 곱셈과 연산법칙으로 정의한다. ``가약'', ``기약'', ``완전가약'', ``절대기약'' 표현이라는 개념은 통상적인 방식으로 정의하며, 이 명칭들을 표현류에도 확장해 쓴다. $\mS$의 모든 표현류가 $\mathrm{P}$에서 완전가약일 때, $\mS$를 ``$P$에서 완전가약''이라고 한다.

%
%
이제 $\mS$가 $\mathrm{P}_0$에서 완전가약이라고 가정하자. 또한 $\mathrm{P}_0$를 완전체라고 가정하면, $\mathrm{P}_0$의 모든 대수적 확대 $\mathrm{P}$도 마찬가지로 완전체이다. 표현류들은 $\mathrm{P}_0$의 임의의 확대 $\mathrm{P}$에서도 계속 완전가약이다.\footnote{$\mathrm{P}_0$가 완전체가 아닐 때에는 중심의 성분 $\mathfrak{K}$들이 $\mathrm{P}_0$의 ``제1종 확대''가 되는 경우에, 그리고 오직 그 경우에만 완전가약성이 보존된다. 이 가정 아래에서는 본문의 이후 모든 귀결이 성립한다.}

$\mS$의 중심은 유한 개 체의 직합이 된다. 그러한 체 $\mathfrak{K}$와 동형인 $\mathrm{P}_0$의 모든 확대 $P$는 어떤 (단순) 지표의 체가 된다. $\mathrm{P}_0$를 $\mathfrak{K}$와 동형인 $\mathrm{P}$로 확대하면, $\mS$에서 하나의 양쪽 단순환이 분리되어 나오며 이 환에는 그 지표의 절대기약 표현류가 대응한다. 이 환은 베더번 정리에 따라 $\mathrm{P}$ 위의 한 행렬환과, 동형을 제외하면 유일하게 정해지는 비가환체 $\mA$의 직접곱이 된다. $\mA$의 중심은 $P$이고 $P$ 위에서 차수 $m^2$이며, 여기서 $m$은 그 지표에 속하는 슈어 지수를 뜻한다.

\srcspaced{$\mS$의 모든 분해체---그 지표에 속하는 표현에 관하여---는 동시에 $P$에서의 표현들에 관하여 $\mA$의 분해체이며, 그 역도 성립한다.} $\mS$의 공액 표현들은 $\mathrm{P}_0$ 위의 대응하는 양쪽 단순환에서 출발하여 얻는다. 이에 대응하는 체 $\mA_0$는 $\mA$와 동형이 되며 $\mathfrak{K}$를 중심으로 가진다. 따라서 분해체 문제는 대응하는 비가환체 $\mA$와 그 분해 문제로 완전히 환원된다.\footnote{이는 $\mathrm{P}$에서 $\mS$의 기약표현과 가환하는, $\mathrm{P}$의 원소를 성분으로 하는 행렬계로 환원하는 Schur의 방법에 대응한다. 이 행렬들의 전치행렬들은 $P$에서 $\mA$의 한 표현을 이룬다.} 특히 다음 정리가 성립한다.

%
%
\srcspaced{$\mA$의 모든 극대 가환 부분체는 $P$ 위에서 같은 차수 $m$을 가지며, 여기서 $m$은 슈어 지수를 나타낸다. 이 모든 극대 가환 부분체는 최소 차수의 분해체이고, 최소 차수의 모든 분해체는 이들 가운데 하나이다. 모든 분해체의 차수는 $m$의 배수이다. $\mA$에는 절대기약 표현류가 단 하나만 있으며, 마찬가지로 $\mathrm{P}$에 대한 기약 표현류도 단 하나만 있다.}

%
%
일반적인 분해체를 특성화하기 위하여, $\mA_r$를 $\mA$와 $P$에 관한 차수 $r^2$의 행렬환, 곧 $\mathrm{P}$의 원소를 성분으로 하는 모든 $r$차 정사각행렬의 계와의 직접곱이라 하자. 그러면 $\mA_r$는 $\mA$를 대응하는 비가환체로 갖는 양쪽 단순환이 된다. 또한 $\mathrm{P}$에 관하여 초복소수계이고 $\mA$가 대응하는 모든 양쪽 단순환은 이와 같은 방식으로 얻어진다. $\mA_r$는 $\mA$의 원소를 성분으로 하는 모든 $r$차 정사각행렬의 계로도 특성화된다.

%
%
\srcspaced{차수 $mr$인 $\mA$의 모든 분해체---그러한 것이 존재한다면---는 $\mA_r$의 극대 가환 부분체 하나와 동형이다. 역으로 $\mA_r$의 모든 극대 가환 부분체는 분해체이며, 각각 그 차수는 $ms$이다. 여기서 $s$는 $r$의 약수이다. 정칙인 경우에는 $\mA_r$의 모든 극대 가환 부분체의 차수가 $mr$이다. 여기서 정칙인 경우란 다음을 뜻한다. 기초체 $\mathrm{P}$ 위의 임의의 가환체 $\Sigma$가 $\mathrm{P}$ 위에서 유한일 때, $\Sigma$ 위에 임의의 차수를 갖는 가환 확대체가 더 존재한다는 성질을 $P$가 갖는다.}\footnote{따라서 예를 들어 $\mathrm{P}$가 모든 실수의 체이면 정칙인 경우가 아니다.}

%
%
$r>1$일 때 이 분해체들 가운데 극소 분해체도 나타나는지는 $\mA_r$로의 매장만으로는 아직 결정되지 않는다. 실제로 무한히 많은 $r$에 대하여 그런 일이 일어날 수 있음을 다음 예가 보여 준다.

\begin{center}
\textbf{2. 극소 분해체 차수들의 비유계성.}
\end{center}

이 비유계성을 보이기 위하여 이제 사원수체를 예로 든다. 이를 유리수체 $\mathrm{P}$에 관한 초복소수계로 보고, 단위원 $1$ 이외의 기저원소를 $i,j,k$라 하면 잘 알려진 관계식은 다음과 같다.
\[
 i^2=j^2=k^2=-1;\qquad ij=k;\quad jk=i;\quad ki=j;\quad ji=-k;\quad kj=-i;\quad ik=-j.
\]

%
%
\srcspaced{어떤 대수적 수체를 첨가했을 때 $r^2=r$을 만족하면서 $0$과 단위원 $1$과도 다른 ``멱등원'' $r$이 존재하는 경우, 바로 그러한 수체들만이 분해체임을 보일 것이다. 이 수체들은 $(-1)$을 세 제곱의 합으로, 따라서 두 제곱의 합으로 나타낼 수 있는 수체들로도 특성화할 수 있다.}\footnote{$(-1)$을 세 제곱의 합으로 나타낼 수 있으면 언제나 두 제곱의 합으로도 나타낼 수 있다는 것은 항등식
\[
(c^2+d^2)(a^2+b^2+c^2+d^2)=(ac+bd)^2+(ad-bc)^2+(c^2+d^2)^2
\]
에서 알 수 있다. 이 항등식은 예를 들어 사원수의 노름곱 정리로부터 얻어진다. 분해를 위해 멱등원의 존재를 요구하는 대신 노름이 $0$인 원소의 존재를 요구해도 충분했을 것이다.}

실제로
\[
        r=\alpha+\frac12\beta i+\frac12\gamma j+\frac12\delta k
\]
로 두면
\[
        r^2=\left(\alpha^2-\frac14\beta^2-\frac14\gamma^2-\frac14\delta^2\right)
          +\alpha\beta i+\alpha\gamma j+\alpha\delta k
\]
이므로, $r^2=r$은 다음과 동치이다.
\[
4\alpha^2-4\alpha-\beta^2-\gamma^2-\delta^2=0;\quad
2\alpha\beta-\beta=0;\quad 2\alpha\gamma-\gamma=0;\quad 2\alpha\delta-\delta=0.
\]
따라서 $\beta,\gamma,\delta$가 동시에 $0$이어서는 안 된다는 조건 아래, 이는 다시 다음과 동치이다.
\[
        (-1)=\beta^2+\gamma^2+\delta^2.
\]

%
%
그러한 멱등원 각각이 분해를 낳는다는 사실과, 그러한 멱등원의 존재가 분해에 필요하다는 사실은 다음에서 나온다. 완전가약계의 모든 기약표현은 한쪽 단순 아이디얼, 더 정확히는 그에 대응하는 아이디얼류에 의해 생성되며, 그러한 아이디얼류는 각각 그 표현을 생성한다. 이 한쪽 단순 아이디얼들은 모두 한쪽 직접합 분해에 나타나고, 이 분해는 단위원의 성분이 되는 ``원시'' 멱등원들을 결정한다. 역으로 모든 멱등원은
\[
\mS=r\mS+(e-r)\mS
\]
라는 합분해를 주며, 여기서 $e$는 단위원이다. ``원시'' 멱등원들은 단순 아이디얼을 분리해 낸다.\footnote{원시 멱등원과 기약표현의 관계는 이미 M. Herzberger의 학위논문 Über Systeme hyperkomplexer Größen, Kapitel 4, Berlin 1923에 나온다. 단순 아이디얼이 기약표현을 생성한다는 사실은 가환인 경우에는 예를 들어 E. Noether, Der Diskriminantensatz für Ordnungen, Crelle 157 (1926)에, 비가환인 경우에는 W. Krull, Theorie und Anwendung der verallgemeinerten Abelschen Gruppen, Heidelberger Berichte 1926에 증명되어 있다. Speiser의 군론 마지막 글도 참조하라. $a_1,\cdots,a_m$이 분해체에 관한 아이디얼의 기저이고 $c$가 임의의 환 원소여서
\[
a_i c=\gamma^i_1a_1+\cdots+\gamma^i_ma_m
\]
이라면, $(\gamma_{ik})$는 $c$에 대응하는 행렬이다. 모든 기약표현이 아이디얼에 의해 생성된다는 사실은 더 깊은 정리이다.}

비가환체를 그 중심에 관한 초복소수계로 보자. 이를 분해체에 관한 초복소수계로 확대하면 $m$개의 절대단순 한쪽 아이디얼의 직접합이 되며, 여기서 $m$은 슈어 지수이다.

사원수체의 경우에는 $m=2$이다. 따라서 $0$도 $1$도 아닌 모든 멱등원은 이미 절대단순 아이디얼들로의 분해를 낳고, 이에 절대기약 표현류들로의 분해가 대응한다. 이로써 앞에서 제시한 체들이 분해체로 특성화된다.

%
%
더 나아가 다음이 성립한다. \srcspaced{유리수체에 관하여 차수 $2^n$인 순환체 $\Omega_n$ 가운데 $(-1)$을 두 제곱의 합으로 나타낼 수 있는 것들은 모두 사원수체의 극소 분해체이다.} 실제로 $(-1)$을 제곱들의 합으로 나타낼 수 있으므로 분해체는 실수체일 수 없다. 한편 $\Omega_n$의 차수 $2^{n-1}$인 부분체---따라서 $\Omega_n$의 모든 진부분체---는 반드시 실수체이다. 이 부분체는 $\Omega_n$과 모든 실대수적 수의 체의 교집합이기 때문이다. 실제로 $\Omega_n$은 이 교집합에 대한 차수 2의 갈루아체가 되므로, 이 교집합은 반드시 차수 $2^{n-1}$인 그 부분체와 일치한다.

%
%
H. Hasse가 증명한 존재정리\footnote{$\Omega_2$의 존재는 이를테면 Weber가 제시한 4차 순환체의 매개변수 표현에서도 이미 읽어낼 수 있다(Kleines Lehrbuch der Algebra, § 93).}에 따르면 \srcspaced{그러한 체 $\Omega_n$은 모든 $n$에 대해 존재하며, 극소 분해체들의 차수에는 상계가 없다. 여기서는 심지어 지수의 모든 거듭제곱이 극소 분해체의 차수로 나타난다.}\footnote{그 뒤 R. Brauer는 사원수체의 극소 분해체가 모든 짝수 차수에 대해, 곧 분해체에 가능한 모든 차수에 대해 존재한다는 것까지 보였다. 그러한 체 $\mathrm{P}(\vartheta)$는 모든 $n\geqq2$에 대해 다음 방정식으로 정의할 수 있다.
\[
 f(\vartheta)=\vartheta^{2n}+a^2p^2\vartheta^2+b^2p^{2\beta}=0;
\quad \text{따라서}\quad -1=\left(\frac{ap}{\vartheta^{n-1}}\right)^2+
        \left(\frac{bp^\beta}{\vartheta^n}\right)^2,
\]
여기서 $a,b,p,\beta$을 적절히 선택하면 다음 조건들을 성립시킬 수 있다. 1. $f(x)$는 모든 $n$에 대해 기약이다. 2. $\mathrm{P}(\vartheta)$는 $\mathrm{P}$와 다른 부분체를 오직 하나만 가지며, 그것은 $\mathrm{P}(\vartheta^2)$이다. 3. $\mathrm{P}(\vartheta^2)$의 공액체들 가운데에는 실수체인 것도 있으므로 $\mathrm{P}(\vartheta^2)$는 분해체일 수 없다. 가능한 $a,b,p,\beta$의 선택 가운데 한 예는 다음과 같다.
\[
        f(x)=x^{2n}+(25\cdot9)4x^2+9\cdot64.
\]
그 증명은 원시 방정식을 구성하는 Furtwängler의 방법(Math. Ann. 85, S. 34, § 3)을 변형한 것에 바탕을 두지만, 세부 전개는 상당히 번거롭다.}

%
%
\begin{center}
\textbf{3. 사원수체의 분해체에 대한 초등적 고찰.}
\end{center}

$\mathfrak{O}$는 다음 8개의 행렬로 이루어진 사원수군의 표현을 뜻한다.\footnote{Sylvester, Math. Papers Vol. III S. 647, Vol. IV S. 122를 참조하라.}
\[
\pm E=\pm\mat{1&0\\0&1},\qquad
\pm I=\pm\mat{i&0\\0&-i},\qquad
\pm J=\pm\mat{0&1\\-1&0},\qquad
\pm K=\pm\mat{0&i\\ i&0}.
\]
$\Omega$가 생성하는 군환은 바로 초복소수계로 본 사원수체의 한 표현이다. 우리는 $\Omega$를 체 $\mathrm{K}$에서 표현할 수 있는 것은 $-1$을 $\mathrm{K}$에서 두 제곱의 합으로 나타낼 수 있을 때이고 또 그때에만 그러함을 초등적으로 보이고자 한다.

%
%
\emph{a)} $\mathfrak{O}^*=T^{-1}\mathfrak{O} T$가 체 $\mathrm{K}$ 위의 유리 표현이라고 하자. 또 $I^*=T^{-1}IT$, $J^*=T^{-1}JT$라 하자. $J$와 $J^*$는 모두 $\mathrm{K}$ 위에서 유리적인 서로 상사인 행렬이므로, $J^*$를 $J$로 옮기는 $K$ 위의 유리 변환 $R$도 존재한다.
\[
        R^{-1}J^*R=J.
\]
처음부터 $\mathfrak{O}^*$를 역시 $\mathrm{K}$ 위에서 유리적인 표현 $R^{-1}\mathfrak{O}^*R$로 바꾸면, 일반성을 잃지 않고 $J^*=J$이라 가정할 수 있음을 알 수 있다. 이제
\[
        I^*=\mat{a&b\\ c&d}\qquad (a,b,c,d\text{는 }\mathrm{K}\text{의 원소})
\]
라 하자. $IJ=-JI$에서 $I^*J^*=-J^*I^*$가 따르고, $J^*=J$이므로
\[
        a=-d,
        \qquad b=c
\]
를 얻는다. 또 $I^2=-E$에서 $I^{*2}=-E$가 따르므로
\[
        \mat{-1&0\\0&-1}=\mat{a&b\\ b&-a}^2
        =\mat{a^2+b^2&0\\0&a^2+b^2}.
\]
따라서 $a^2+b^2=-1$이고, 실제로 $-1$을 $\mathrm{K}$에서 두 제곱의 합으로 나타낼 수 있다.

%
%
\emph{b)} 체 $\mathrm{K}$에서 $-1$을 두 제곱의 합으로 나타낼 수 있다고 하자. 이를테면 $-1=a^2+b^2$라면 다음과 같이 둔다.
\[
        I^*=\mat{a&b\\ b&-a},\qquad
        J^*=\mat{0&1\\-1&0},\qquad
        \mathit{K}^*=I^*J^*.
\]
쉽게 계산해 보면 $\pm E$, $\pm I^*$, $\pm J^*$, $\pm \mathit{K}^*$는 $\mathrm{K}$ 위에서 유리적이고 $\Omega$와 상사인 군을 이룬다.\footnote{인자계를 이용해도 분해체의 특성화를 쉽게 얻을 수 있다.}

%
%
\srcspaced{더 나아가 무한히 많은 $n$에 대하여, 사원수군의 극소 분해체인 차수 $2^n$의 순환체가 존재함을 보이겠다.}

적당한 $r>0$에 대하여
\[
        (*)\qquad 2^r+1\equiv0\pmod p
\]
를 만족하는 유리소수 $p$를 택하고, $2^n$을 $p-1$을 나누는 2의 가장 큰 거듭제곱이라 하자.

먼저 무한히 많은 $n$에 대하여 그러한 소수 $p$가 존재함을 보인다. 임의의 양의 정수 $k$를 택하고, $p$를 $2^{2^k}+1$의 한 소인수라 하자. 그러면 $r=2^k$로 두었을 때 $p$에 요구한 합동식이 성립한다. 또한 $2\pmod p$의 위수는 $2^{k+1}$이므로, $p-1$을 나누는 2의 가장 큰 거듭제곱 $2^n$에 대하여 $n>k$이다. 따라서 임의로 큰 $n$에 대응하는 소수 $p$가 존재한다.

%
%
$\varepsilon$을 원시 $p$차 단위근이라 하자. 이제 $\mathrm{P}(\varepsilon)$에서 $-1$을 두 제곱의 합
\[
        -1=a^2+b^2=(a+ib)(a-ib)\qquad (a,b\text{는 }\mathrm{P}(\varepsilon)\text{의 원소})
\]
으로 나타낼 수 있음을 보이겠다. 이는 $4p$차 단위근들의 체 $\mathrm{P}(\varepsilon,i)=\mathrm{P}(\varepsilon\cdot i)$ 안에 $\mathrm{P}(\varepsilon)$에 대한 상대노름이 정확히 $-1$인 원소가 존재한다는 것과 동치이다.

$1+i\varepsilon^\nu$ $(\nu=1,2,\ldots,p-1)$의 상대노름은
\[
(1+i\varepsilon^\nu)(1-i\varepsilon^\nu)=1+\varepsilon^{2\nu}
\]
이다. 따라서 모든 $1+\varepsilon^\nu$ $(\nu=1,2,\ldots,p-1)$는 $\mathrm{P}(\varepsilon)$을 기초체로 할 때 $\mathrm{P}(\varepsilon i)$의 원소들의 상대노름이다. 그러므로 다음 $\Pi$도 상대노름이다.
\[
\Pi=\prod_{\nu=0}^{r-1}(1+\varepsilon^{2^\nu})
=(1+\varepsilon)(1+\varepsilon^2)(1+\varepsilon^4)\ldots(1+\varepsilon^{2^{r-1}})
=\sum_{\lambda=0}^{2^r-1}\varepsilon^\lambda.
\]
관계식 $(*)$에 의해 $\varepsilon^{2^r}=\varepsilon^{-1}=\varepsilon^{p-1}$이므로, $\Pi$에 $\varepsilon^{p-1}$을 더하면 $\Pi+\varepsilon^{p-1}$은 각각
\[
1+\varepsilon+\varepsilon^2+\ldots+\varepsilon^{p-1}=0
\]
인 항들의 묶음들의 합이 된다. 따라서
\[
        \Pi=-\varepsilon^{-1}.
\]
또한 $\varepsilon$도 $\mathrm{P}(\varepsilon i)$의 원소, 곧 $\varepsilon^{\frac{p+1}{2}}$의 상대노름이다. 그러므로
\[
        \Pi\cdot\varepsilon=-1
\]
도 상대노름이다. 따라서 $\mathrm{P}(\varepsilon)$에서 $-1$을 두 제곱의 합으로 나타낼 수 있다.\footnote{$n8\pm3$ 꼴의 소수에 대해서는 이 논법이 이미 E. Landau, Über die Darstellung definiter Funktionen durch Quadrate, Math. Ann. Bd. 62, S. 272--285에 나온다.} 그러므로 $\mathrm{P}(\varepsilon)$은 사원수군의 분해체이다.

%
%
$\mathrm{K}$를 $\mathrm{P}(\varepsilon)$에 포함된 차수 $2^n$의 순환체라 하자. $\mathrm{K}$ 역시 분해체임을 주장한다.\footnote{원분체의 부분체로서 차수 $2^n$의 순환 분해체를 택한 것은 하세의 예를 본뜬 것이다.} $K$를 기초체로 할 때 사원수체의 지수를 $m$이라 하면, $m=1$ 또는 $m=2$이다. 그런데 $K$에 대한 상대차수 $\frac{p-1}{2^n}$가 홀수인 분해체 $\mathrm{P}(\varepsilon)$가 존재하므로, 두 번째 경우는 불가능하다. 따라서 $m=1$이다. 이는 곧 $\mathrm{K}$가 분해체라는 뜻이다.

따라서 무한히 많은 $n$에 대하여, 극소 분해체인 차수 $2^n$의 순환체가 존재한다.

\clearpage
\setcounter{footnote}{0}
\phantomsection
\addcontentsline{toc}{section}{33. 산술적 관점에서 본 초복소수와 표현론}
%
%
%

%
%
\section*{33. 산술적 관점에서 본 초복소수와 표현론}
\emph{Atti Congresso Bologna 2 (1928), S. 71--73.}

\begingroup
\renewcommand{\thefootnote}{(\arabic{footnote})}

\begin{center}
\textsc{에미 뇌터(Emmy Noether)} (괴팅겐 -- 독일)\\[-0.1em]
\rule{3em}{0.4pt}\\[0.8em]
\textsc{초복소수와 표현론}\\
\textsc{산술적 관점에서}\footnote{상세한 논의는 Math. Zeitschr.에 실릴 예정이며, 그 속편에서는 비가환 갈루아 이론도 다룰 것이다. 또한 ``분해체(Zerfällungskörper)''에 관해서는 R. Brauer와 E. Noether, Ber. Berl. Ak. 1927을, 일반 구조정리에 관해서는 이 회의록에 실린 G. Köthe의 강연을 참조하라.}
\end{center}

%
%
나는 표현론---특히 군의 표현론---을 가군류와 아이디얼류의 이론으로 이해할 수 있는 방식과, 이 후자의 것들을 다시 일반화된 군 개념, 곧 연산자를 가진 군 아래에 포섭할 수 있는 방식을 보이고자 한다.\footnote{연산자를 가진 군이라는 개념은 W. Krull (Math. Zeitschr. 23)과 O. Schmidt (Math. Zeitschr. 29)에 거슬러 올라간다. 여기서 제시하는, 주어진 표현의 환원 문제에 대한 처리도---표현체가 가환인 경우로 제한하면---W. Krull에게서 비롯한다(Theorie und Anwendung der verallgemeinerten Abelschen Gruppen, Heidelberger Ber. 1926).}

%
%
주어진 체 $P$ 위에서 \emph{군의 표현}이란 잘 알려진 대로 다음을 뜻한다. 군 $\mathfrak{G}$의 원소 $a,b$에 $P$의 원소를 성분으로 하는 행렬 $A,B$를 대응시켜, 곱 $ab$에 곱 $AB$가 대응하도록 한다. 그러면 이 행렬계는 군의 준동형상이 된다. 이 표현은 동시에 ``군환''의 표현도 준다. 군환은 모든 ``군수'', 곧 $a\alpha+b\beta+....+c\gamma$ 꼴의 모든 선형결합으로 이루어진다. 여기서 $a,b,....,c$는 군 원소들 가운데 임의로 유한 개를, $\alpha,\beta,....,\gamma$는 그에 대응하는 $P$의 임의의 유한 개 원소를 나타낸다($a,b,....,c$는 선형독립으로 간주하며, 곱셈은 군의 곱셈과 연산법칙으로 정의한다). 군환의 표현은 덧셈에 관해서도 준동형이 된다. 이로써 군의 표현은 덧셈과 곱셈에 관한 준동형성을 요구하는 \emph{환의 표현}이라는 일반 문제에 포함된다.

%
%
여기서 두 문제가 본질적이다. 첫째는 \emph{주어진 표현의 환원}과 그때 나타나는 기약 성분들의 유일성 문제이다. 둘째는 주어진, 반드시 가환일 필요는 없는 체 위에서 주어진 환이 가질 수 있는 표현 전체---이를테면 기약표현만으로 한정한 전체---에 관한 문제이다.

첫째 문제는 훨씬 간단하다. 이는 표현될 환과---일반적으로 비가환인---표현체에 아무런 특별한 가정도 필요하지 않다는 점에서 이미 드러난다. 표현이 고정된 차수의 행렬들로 이루어진다는 사실이 필요한 모든 유한성 가정을 준다. 이 문제는 연산자를 가진 군의 이론으로 완전히 다룰 수 있으므로, 이제 그 이론을 간단히 살펴보겠다.

%
%
군 $\mathfrak{B}$와 함께 기호들의 계 $\varTheta,H$---곧 연산자들---가 주어져서, $\mathfrak{B}$의 각 $a$에 대해 $\varTheta(a),H(a)$가 유일하게 정의되고 $\mathfrak{B}$의 원소가 되며, 또한 분배법칙
\[
\varTheta(ab)=\varTheta(a)\varTheta(b)
\]
이 성립할 때, 군 $\mathfrak{B}$를 \emph{연산자 영역을 가진 군}이라 한다. 같은 연산자 영역을 가진 두 군이 통상적인 의미에서 동형이고, 또한 $a$와 $\bar a$가 서로 대응하면 $\varTheta(a)$와 $\varTheta(\bar a)$, $H(a)$와 $H(\bar a)$도 각각 서로 대응할 때, 이 두 군을 연산자 동형이라 한다. 이제 부분군으로는 그 자체도 이 연산자 영역을 허용하는 것만을 인정하자. 그러면 군이 단순하다는 것은 항등부분군 이외에는 허용되는 진정한 정규부분군을 갖지 않는다는 뜻이다. 이때 합성열에 관한 요르단--횔더 정리는 다음 형태로 그대로 성립한다. 연산자 영역을 가진 군이 합성열을 하나라도 가지면, 그 인자군들의 계는 연산자 동형의 의미에서 순서의 차이를 제외하고 유일하게 결정된다. 또한 같은 전제 아래, 직접분해 불가능한 인자들로의 분해가 연산자 동형의 의미에서 유일하다는 정리도 성립한다.

%
%
표현 문제와의 연관은 특히 아이디얼과 가군도 연산자를 가진 군에 속한다는 데서 나온다. 이때 아이디얼과 가군은 덧셈에 관한 아벨군으로 보고, 연산자는 환 원소와의 곱셈으로 준다. 한편 모든 표현은 표현가군, 곧 환과 표현체에 관한 가군으로부터 생성된다. 더 정확히 말하면, 서로 연산자 동형인 표현가군들의 류인 각 가군류는 동일한 표현을 생성한다. 이로써 표현의 환원이 유일한가 하는 문제는 합성열 정리와 직접곱 정리로 해결된다. 합성열을 표현가군에 적용하면 각 행렬에 대해 다음 꼴의 환원을 얻는다.
\[
\begin{pmatrix}
B_1&0&0\\
&B_2&0\\
A_{ik}&....&B_r
\end{pmatrix}
\]
여기서 인자군들에 대응하는 대각행렬들은 ``기약'' 표현, 곧 그 자체는 더 이상 이러한 환원을 허용하지 않는 표현을 생성한다.

%
%
그런데 이 인자군들의 류가 유일하게 결정되므로 대각표현들도 마찬가지로 유일하다. 여기서 유일성은 표현의 동치라는 의미에서의 유일성이다. 이와 마찬가지로 직접곱 정리는 ``분해''가 ``직접분해 불가능한'' 성분들로 이루어지는 방식의 유일성을 준다.
\[
\begin{pmatrix}
C_1&&&\\
&C_2&&\\
&&....&\\
&&&C_s
\end{pmatrix},
\]
이때 각 $C_i$는 다시 합성열 정리에 따라 유일하게 환원될 수 있다.

%
%
\emph{둘째 문제}는 초복소수계의 표현, 더 일반적으로는 ``이중 사슬 조건''을 만족하는 환의 표현의 경우를 들어 개략적으로 설명하고자 한다. 그 결과는 모든 기약(단순) 가군류가 아이디얼류라는 것, 따라서 한쪽 아이디얼들로 이루어진 합성열의 인자군들이 이미 모든 기약 표현을 준다는 것이다. 더 정확히 말하면 라디칼로 나눈 잉여류환에서의 합성열이면 충분하며, 이 잉여류환은 ``완전가약'' 환이 된다. 이로써 문제는 완전가약환의 구조를 조사하는 문제로 환원된다. 또한 자기동형체에서의 표현을 출발점으로 삼는 것이 이 문제에 알맞음이 드러난다.

그 뜻은 다음과 같다. 이러한 환의 양쪽 단순 성분에 속하는 모든 단순 오른쪽 아이디얼은 같은 류에 속하므로 동일한 자기동형환을 가진다. 이 자기동형환은 아이디얼의 단순성 때문에 체, 곧 나눗셈환이 되며, 동시에 단순 왼쪽 아이디얼들의 자기동형환이기도 하다. 단순환 자체는 자기동형체 위의 모든 행렬로 이루어진 환과 동형이 된다. 그리고 이 표현으로부터 부분체에 관한 모든 표현, 특히 초복소수계의 계수 영역에 관한 모든 표현이 생겨난다. 이는 자기동형체의 원소 자체를 각각 그에 대응하는 행렬로 바꿈으로써 얻어진다.

이리하여 초복소수계의 경우에 알려진 베더번 구조정리는 일반 군론에서 유래한 자기동형체 개념을 통해 새롭게 해석되고, 동시에 초복소수계 표현론의 원천으로 인식된다. 이로써 비가환체의 갈루아 이론이라는 새로운 문제들이 생겨나며, 이 문제들은 ``분해체'' 문제와 밀접하게 관련된다. 동시에 직접분해 불가능한 아이디얼의 자기동형환을 통해 일반적인 완전가약이 아닌 환의 구조정리를 내다볼 수 있게 된다. 군론은 연산자를 가진 가장 일반적인 군들로까지 확장됨으로써 여기에서도 다시 질서를 부여하는 원리임이 드러난다.

\endgroup

\clearpage
\setcounter{footnote}{0}
\phantomsection
\addcontentsline{toc}{section}{34. 초복소량과 표현론}
%
%
%
%
%

\providecommand{\mE}{\mathfrak E}
\providecommand{\mA}{\mathfrak A}
\providecommand{\mB}{\mathfrak B}
\providecommand{\mC}{\mathfrak C}
\providecommand{\mD}{\mathfrak D}
\providecommand{\mG}{\mathfrak G}
\providecommand{\mH}{\mathfrak H}
\providecommand{\mK}{\mathfrak K}
\providecommand{\mL}{\mathfrak L}
\providecommand{\mM}{\mathfrak M}
\providecommand{\mN}{\mathfrak N}
\providecommand{\mO}{\mathfrak O}
\providecommand{\mP}{\mathfrak P}
\providecommand{\mR}{\mathfrak R}
\providecommand{\mS}{\mathfrak S}
\providecommand{\mT}{\mathfrak T}
\providecommand{\mU}{\mathfrak U}
\providecommand{\mZ}{\mathfrak Z}
\providecommand{\mo}{\mathfrak o}
\providecommand{\ma}{\mathfrak a}
\providecommand{\mb}{\mathfrak b}
\providecommand{\mc}{\mathfrak c}
\providecommand{\md}{\mathfrak d}
\providecommand{\mf}{\mathfrak f}
\providecommand{\mh}{\mathfrak h}
\providecommand{\ml}{\mathfrak l}
\providecommand{\mn}{\mathfrak n}
\providecommand{\mr}{\mathfrak r}
\providecommand{\mt}{\mathfrak t}
\providecommand{\mx}{\mathfrak x}
\providecommand{\mz}{\mathfrak z}
\providecommand{\iso}{\simeq}
\providecommand{\ann}{\operatorname{Ann}}
\providecommand{\tuple}[1]{(#1)}
\providecommand{\mat}[1]{\begin{pmatrix}#1\end{pmatrix}}
\providecommand{\glqq}{“}
\providecommand{\grqq}{”}

\section*{34. 초복소량과 표현론}
\emph{Math. Zs. 30 (1929), 641--692쪽.}

\begin{center}
{\Large\bfseries 초복소량과 표현론\footnote{이하는 B. L. van der Waerden이 나의 1927/28년 겨울학기 강의를 자유롭게 정리한 것이다. 출판을 위한 원고는 우리가 함께 다듬었다. 또한 여러 비판적 논평을 해 준 B. L. van der Waerden에게 감사를 표한다.}.}\\[1em]
글쓴이\\[0.5em]
괴팅겐의 에미 뇌터.\\[0.8em]
\rule{3em}{0.4pt}
\end{center}

\subsection*{서론.}

초복소계에 관한 가장 중요한 일반정리들은 Molien(Math. Annalen 41권 및 Dorpater Ber. 1897)에게서 비롯한다. 그와 본질적으로 독립하여 프로베니우스는 곧이어 초복소계와 그 표현의 이론, 특히 유한군의 표현론을 통일적으로 전개하였다. 그 기초는 데데킨트에게서 유래한 군행렬식, 더 일반적으로는 임의의 초복소계의 행렬식이라는 개념이다. 프로베니우스는 이 행렬식의 서로 다른 기약인자들에 --- 계수영역은 언제나 복소수 전체의 체이다 --- 서로 다른 기약표현류들이 대응하며, 이 방식으로 모든 기약표현류가 빠짐없이 얻어진다는 것을 보였다. 그러한 표현류는 그 \glqq{}지표계\grqq{}로 완전히 특징지어진다. 유한군의 경우 이 지표계는 군의 켤레원소류들에서 생성되는 가환 초복소계의 행렬식을 분해하여 얻어진다. 이 행렬식은 지표들을 계수로 하는 일차인자들로 분해된다. 이는 유한 아벨군의 군행렬식이 그 아벨군의 서로 다른 지표들을 계수로 하는 일차인자들로 분해된다는, 데데킨트가 처음 얻은 결과의 직접적인 일반화이다(프로베니우스와의 서신교환).

그러나 프로베니우스는 개념적으로 단순하고 투명한 이 결과들을 힘겨운 계산으로 얻었다. 후대의 발전은 이 결과들을 더 단순하게 유도하는 동시에 임의의 체를 계수영역으로 삼는 경우로 확장하는 데 향하였다. 이 발전은 초복소계와 표현론에서 완전히 따로 진행되었다. 초복소계는 웨더번에게서 산술적으로 다루어졌다. 모든 초복소계는 양쪽으로 직접분해되지 않는 환들의 직접합으로 유일하게 분해된다. 근기가 없는 계에서는 이 직접분해되지 않는 환들이 본질적으로 유일하게 결정되는 한 대응체 --- 반드시 가환일 필요는 없다 --- 의 원소로 된 모든 \(n\)행 정방행렬의 계와 동형이다\footnote{Dickson, \emph{Algebras and their Arithmetics}(독일어판 \emph{Algebren und ihre Zahlentheorie}, Zürich 1927)의 문헌목록을 보라.}. 표현론은 Burnside와 I. Schur에게서 초복소계와 무관하게 주어진 표현에서 직접 출발하여 행렬정리들을 쓰는 초등적 기초를 얻었다. 특히 Schur는 주어진 한 체에서 기약인 표현이 절대적으로 분해되는 최소차수 수체를 물었다. 절대기약 성분들은 유한개의 켤레 표현류에 속한다. 찾는 수체들의 차수는 이 표현류의 수와 인덱스, 곧 각 켤레류에 속하는 성분의 수의 곱과 같다. Schur 이론의 주된 도구는 한 기약표현과 가환하는 모든 행렬이 이루는 계이다\footnote{Speiser의 군론 제10장의 문헌목록을 보라.}.

이하의 순수하게 산술적인 정초에서는 초복소량과 표현론이 다시 하나의 통일된 전체로 나타난다. 이는 일정한 유한성 조건만을 만족하는 비가환환의 일반이론의 특수한 경우이다. 정확히는 이 환들에 관한 가군류와 아이디얼류의 이론이며, 주된 결과는 기약 가군류들이 이미 대응 아이디얼류들로 모두 얻어진다는 것이다. 특히 근기가 없는 환에서는 모든 가군류가 기약인 것들로 분해되어 완전가약이 된다. 이는 정칙 군행렬식 또는 계행렬식의 기약인자들만으로 이미 모든 기약표현류를 얻고, 근기가 없는 계에서는 표현들이 완전가약이라는 프로베니우스 결과의 산술적 동치이다. 실제로 계행렬식과 그 분해를 고찰하는 것은 노름으로 넘어가는 일로 볼 수 있지만, 여기서는 아이디얼의 직접합분해와 그 합성열을 직접 다룬다. 표현론은 표현가군이라는 개념을 통해 가군류의 이론이 된다.

가군류와 아이디얼류의 도입은 순전히 군론적으로 정초된다. 가군과 아이디얼은 덧셈에 관해서는 아벨군이되, 환원소와의 일정한 곱셈을 허용한다는 제한을 받는 것으로 이해한다. 즉 이들은 \glqq{}연산자를 갖는 군\grqq{}을 이룬다(§ 1). 이러한 군에서는 통상적인 동형 대신 \glqq{}연산자동형\grqq{}이 나타난다(§ 2). 한 고정된 영역 안에서 서로 연산자동형인 군들을 한 류로 묶는다. 이 류 개념은 예컨대 수체의 아이디얼의 경우 통상적인 류 개념과 일치한다. 연산자를 갖는 군의 이론은 W. Krull과 O. Schmidt에게서 유래하며(각주 6 참조), 제I장에서 체계적으로 전개한다. 합성열, 직접곱 또는 직접합, 완전가약군에 관한 정리들도 여기서 여전히 성립한다. 위의 류 구분이라는 의미에서의 이 유일성정리들이 이후 모든 논의의 기초이다.

우선 이 정리들로부터 \emph{임의의 표현}의 기약 대각성분에 관한 일반 \emph{유일성정리}를 얻는다(제III장 § 16). 이는 표현류와 표현가군류, 곧 연산자를 갖는 군의 류가 일대일로 대응한다는 사실에 기초한다. 표현들의 \emph{전체}에 관한 더 나아간 물음은 표현되는 환, 특수하게는 초복소계의 구조를 더 자세히 살필 것을 요구한다. 제II장에서는 한쪽아이디얼의 고찰을 전면에 두는 \emph{군론적} 관점에 기초하여 웨더번의 결과들을 새로 얻고 발전시킨다. 구체적으로 우측아이디얼에 대한 \glqq{}배수사슬정리\grqq{}, 또는 그와 같은 \glqq{}최소조건\grqq{} --- 우측아이디얼들의 모든 집합에는 그 집합 안에서 극소인 것이 적어도 하나 있다 --- 만으로 유한성 조건이 충분하다는 것이 드러난다\footnote{\glqq{}이중사슬정리\grqq{}를 가정하면 웨더번의 추론을 옮길 수 있다. E. Artin, \emph{Zur Theorie der hyperkomplexen Zahlen}, Hamb. Abh. 1927을 보라. 또 A. Suschkewitsch, \emph{Über die endlichen Gruppen ohne das Gesetz der eindeutigen Umkehrbarkeit}, Math. Annalen 99 (1928), 30--50쪽을 보라. 여기서는 역산할 수 없는 단 하나의 결합적 연산만을 갖는 유한 영역에 대해 나란한 구조정리들을 전개한다. Suschkewitsch가 \glqq{}핵\grqq{}을 우측군과 좌측군으로 분해하는 것은 \emph{단위원을 갖는 양쪽단순환}을 우측아이디얼과 좌측아이디얼들의 합으로 나타내는 것에 대응한다.}. 근기가 없고 최소조건을 만족하는 환들은 단위원을 갖는 우측 완전가약환들과 정확히 같다. 그러한 환에서 양쪽으로 직접분해되지 않는 한 환의 모든 단순 우측아이디얼은 같은 류에 속한다. 따라서 연산자를 갖는 군으로서 같은 자기동형환을 가지며, 아이디얼이 단순하므로 이 환은 체가 된다. 바로 이 아이디얼류의 자기동형체가 웨더번의 행렬표현을 매개한다.

자기동형체에서의 이 행렬표현은 완전가약인 경우에 자기동형체 또는 그 부분체들에 관한 표현을 요구하면 표현문제를 동시에 해결한다. 특히 초복소계의 계수영역에 관한 경우가 그렇다. 가군류를 아이디얼류로 환원하는 정리의 직접적인 결과로, 웨더번의 표현과 그 자기동형체의 원소들을 자연스러운 방식으로 한 부분체의 행렬들로 바꾸어 얻는 표현 이외에는 다른 기약표현이 없다.

따라서 완전가약인 경우 서로 다른 기약표현류의 수는 아이디얼류의 수와 같고, 이는 서로 다른 양쪽 직접분해불능환, 곧 단순환의 수와 같다. 마지막 수는 다시 중심의 직접분해불능 성분들의 수와 같은데, 이 성분들은 가환체가 된다. 대수적으로 닫힌 계수영역 위의 초복소계에서는 이 가환체들이 중심의 서로 다른 환준동형, 곧 기약표현들로 완전히 결정된다. 이것들은 한 \emph{수인자}를 제외하면 지표들이다.

얻은 결과들은 근기를 갖는 계의 기약표현도 개관하게 해 준다. 이러한 표현은 근기에 따른 잉여류환의 표현으로 환원되며, 그 잉여류환은 근기가 없는 계가 되기 때문이다.

\emph{뒤에서 보이듯이}, 근기가 없는 초복소계에서는 자기동형체로부터 \glqq{}분해체\grqq{}의 이론도 나온다. 분해체란 계수영역의 가환확대 가운데 표현들이 절대기약인 것들로 분해되는 것, 달리 말해 절대단순 우측아이디얼들의 직접합분해가 일어나는 것이다. \emph{분해체}들은 자기동형체의 분해체들과 같고, \emph{최소차수의 모든 분해체}는 자기동형체의 \emph{극대 가환부분체}와 동형이다. 위에서 말한 Schur의 연구와의 관계는 주어진 표현과 가환하는 행렬들의 \emph{전치행렬}들이 바로 자기동형체의 한 표현을 준다는 사실로 세워진다\footnote{R. Brauer--E. Noether, \emph{Über minimale Zerfällungskörper irreduzibler Darstellungen}, Sitz.-Ber. Preuß. Akad. Wiss. 1927, 221쪽을 보라.}. 이 정리들은 비가환체의 갈루아 이론이라는 틀 안에서 체계적으로 전개할 예정이다.

바탕에 둔 개념들, 곧 연산자준동형과 자기동형환은 최소조건만을 만족하는 근기 있는 일반환의 구조도 낳는다. 이는 다른 곳에서 전개될 것이다.

\section*{제I장. 군론적 기초.}

\subsection*{§ 1. 연산자를 갖는 군\footnotemark.}
\footnotetext{여기서 설명하는 개념들은 W. Krull, \emph{Über verallgemeinerte endliche Abelsche Gruppen}, Math. Zeitschr. 23 (1925), 161--196쪽과 O. Schmidt, \emph{Über unendliche Gruppen mit endlicher Kette}, Math. Zeitschr. 29 (1928), 34--41쪽에서 유래한다.}

\(\mG\)를 유한 또는 무한 군이라 하고 원소를 \(a,b,\ldots\)라 쓰자. 군 \(\mG\)의 연산자영역 \(\Omega\)란 새 기호 \(\mathsf{H},\Theta,\ldots\)의 집합으로서, 모든 \(a\in\mG\), \(\Theta\in\Omega\)에 대해 \(\Theta a\in\mG\)가 유일하게 정의되고 분배법칙
\[
        \Theta(ab)=\Theta a\cdot\Theta b
\]
이 성립하는 것이다. 따라서 각 연산자는 군에서 자기 자신으로 가는 준동형을 정의하며, 군은 자기 자신 또는 한 부분군으로 사상된다. 단위원은 단위원으로, 역원은 역원으로 간다.

서로 다른 연산자들이 서로 다른 준동형을 정의하면 그 연산자영역을 \emph{절대적}이라 한다. 임의의 연산자영역에서 같은 준동형을 생성하는 모든 원소를 같게 놓으면 절대 연산자영역을 얻는다. 절대 연산자영역은 군에서 자기 자신으로 가는 모든 준동형의 집합의 한 부분집합과 일대일로 대응한다.

고정된 연산자영역 \(\Omega\)에 관하여, 모든 \(a\in\mH\), \(\Theta\in\Omega\)에 대해 \(\Theta a\in\mH\)인 부분군 \(\mH\)를 \(\mG\)의 \emph{허용부분군}이라 한다.

\emph{예.} 1. 연산자들이 내부자기동형 \(\Theta a=c^{-1}ac\)라 하자. 허용부분군은 정규부분군들이다.

2. 연산자들이 모든 자기동형이라 하자. 허용부분군은 모든 자기동형 아래에서 자기 자신으로 가는 \glqq{}특성부분군\grqq{}들이다.

3. \(\mG\)가 환, 곧 덧셈에 관해서는 아벨군이고 곱셈도 정의되어
\[
\begin{aligned}
r(a+b)&=ra+rb,\\
(a+b)r&=ar+br,\\
ab\cdot c&=a\cdot bc
\end{aligned}
\]
를 만족한다고 하자. 각 원소 \(r\)은 두 연산자 \(rx\), \(xr\)를 동시에 정의한다. 허용부분군은 \glqq{}아이디얼\grqq{} \(\mathfrak a\)이다. 연산 \(rx\)를 허용하여 \(r\mathfrak a\subseteqq\mathfrak a\)인 것은 \emph{좌측아이디얼}, 연산 \(xr\)를 허용하여 \(\mathfrak a r\subseteqq\mathfrak a\)인 것은 \emph{우측아이디얼}, 두 연산을 모두 허용하는 것은 \emph{양쪽아이디얼}이다. 환 \(\mG\)가 \emph{체}, 곧 \(\mG-\{0\}\)이 곱셈에 관하여 군이면 모든 아이디얼은 자명하다, 즉 \(\{0\}\) 또는 \(\mG\)이다\footnote{따라서 환에도 체에도 곱셈의 교환법칙을 가정하지 않는다.}.

4. \emph{환 \(\mo\)에 관한 가군.} \(\mo\)를 환, \(\mM\)을 덧셈으로 쓴 아벨군이라 하자. \(\mo\)의 원소와 \(\mM\)의 원소 사이에 곱 \(r\cdot a\)가 정의되어 다시 \(\mM\)에 들어 있고,
\[
\left.
\begin{aligned}
 r(a+b)&=ra+rb,\\
 rs\cdot a&=r\cdot sa,\\
 (r+s)a&=ra+sa
\end{aligned}
\right\}\qquad r,s\in\mo,\quad a,b\in\mM
\]
를 만족한다고 하자. 그러면 \(\mM\)을 \emph{\(\mo\)-가군}이라 한다. 더 정확히는 \(\mo\)의 곱셈자들을 왼쪽에 쓰므로 \emph{\(\mo\)-좌측가군}이다. 환 \(\mo\)는 동시에 연산자영역이다. 이 영역이 절대적, 곧 서로 다른 환원소들이 서로 다른 연산을 낳으면 \(\mo\)를 가군 \(\mM\)의 \emph{절대 곱셈자영역}이라 한다.

위와 같이 같은 연산을 생성하는 원소들을 같게 놓으면 임의의 곱셈자영역에서 절대 곱셈자영역으로 넘어갈 수 있다. 절대 곱셈자영역은 다시 환이다. 가군 \(\mM\)의 허용부분군을 \emph{부분가군}이라 한다. 특별히 \(\mo=\mM\)이면 좌측아이디얼로 돌아온다\footnote{같은 방식으로 우측가군에는
\[
\begin{aligned}
(a+b)r&=ar+br,\\
a\cdot rs&=ar\cdot s,\\
a(r+s)&=ar+as
\end{aligned}
\]
를 요구한다.}.

5. \emph{쌍가군.} \(\mM\)이 동시에 \(\mo\)-좌측가군과 \(\mo'\)-우측가군이고, 모든 \(a\in\mo\), \(\alpha\in\mM\), \(a'\in\mo'\)에 대해
\[
        a\cdot\alpha a'=a\alpha\cdot a'
\]
이면 \(\mM\)을 \emph{쌍가군}이라 한다.

6. \emph{아벨군의 자기동형환.}\footnote{A. Châtelet, \emph{Les Groupes Abéliens finis}, Paris, Gauthier-Villars, 1925, 99쪽을 보라.} 덧셈으로 쓴 아벨군에서 자기 자신으로 가는 준동형들에는
\[
        (H+\Theta)a=Ha+\Theta a,\qquad
        (H\Theta)a=H(\Theta a)
\]
로 덧셈과 곱셈을 정의할 수 있다. 이렇게 정의한 연산들도 분명 다시 준동형이고, 따라서 그 계에 속한다. 이 계는 환을 이루며, 4에 따라 그 아벨군을 이 환에 관한 가군으로 볼 수 있다.

이제부터 \glqq{}부분군\grqq{}은 언제나 허용부분군을 뜻한다. 두 허용부분군의 교집합 \(\mA\cap\mB\)은 다시 허용부분군이다. 두 부분군 가운데 적어도 하나가 정규부분군이면 그 곱 \(\mA\mB\)도 허용부분군이다.

\subsection*{§ 2. 동형정리.}

같은 연산자영역을 갖는 군 \(\mG,\overline{\mG}\) 사이의 사상이 첫째 통상적인 의미의 준동형 \(ab\rightarrow \bar a\bar b\)이고, 둘째 \(a\)가 \(\bar a\)로 가면 \(\Theta a\)가 \(\Theta\bar a\)로 갈 때 이를 \emph{연산자준동형}이라 한다\footnote{연산자준동형의 개념은 \(\mG,\overline{\mG}\)가 서로 다른 연산자영역을 갖는 경우로 확장할 수 있다. 이때 준동형은 \(\mG\)를 \(\overline{\mG}\)로 사상할 뿐 아니라 연산자영역들도 서로 사상하며, \(a\)가 \(\bar a\)로, \(\Theta\)가 \(\overline\Theta\)로 가면 \(\overline\Theta\bar a\)가 \(\Theta a\)의 상이 되게 한다. 이 일반적 형태에서 연산자준동형은 환준동형도 포함한다. § 8을 보라.}. 기호는 \(\mG\sim\overline{\mG}\)이다. 대응이 일대일이면 \emph{연산자동형}이라 하고 \(\mG\simeq\overline{\mG}\)로 쓴다. 통상적인 군과 마찬가지로 다음 \emph{준동형정리}를 증명한다.

\emph{\(\overline{\mG}\)가 \(\mG\)의 연산자준동형상이면, \(\overline{\mG}\)는 한 몫군 \(\mG/\mN\)과 연산자동형이다. 여기서 \(\mN\)은 \(\overline{\mG}\)의 단위원에 대응하는 \(\mG\)의 모든 원소로 이루어진 허용 정규부분군이다. 역으로 모든 허용 정규부분군 \(\mN\)은 \(\mG\)의 연산자영역을 허용하고 \(\mG\)의 연산자준동형상인 몫군 \(\mG/\mN\)을 정의한다.}

자기 자신과 단위군 이외의 정규부분군을 갖지 않는 군을 \emph{단순}하다고 한다. 단순군의 모든 준동형은 동형이거나 모든 원소를 단위원으로 보낸다.

\emph{제1동형정리.} \emph{\(\overline{\mG}\)가 \(\mG\)의 준동형상이고 \(\overline{\mA}\)가 \(\overline{\mG}\)의 정규부분군이며, \(\mA\)가 \(\overline{\mA}\)의 원소들에 대응하는 \(\mG\)의 모든 원소의 전체이면 \(\mA\)도 정규부분군이고}
\[
        \overline{\mG}/\overline{\mA}\simeq\mG/\mA. \tag{1}
\]
\emph{증명.} \(\mG\sim\overline{\mG}\), \(\overline{\mG}\sim\overline{\mG}/\overline{\mA}\)이므로 \(\mG\sim\overline{\mG}/\overline{\mA}\)이다. 따라서 \(\overline{\mG}/\overline{\mA}\)는 \(\mG\)의 한 몫군과 동형이다. 대응 정규부분군은 \(\overline{\mG}/\overline{\mA}\)의 단위원에 대응하는 원소들의 전체이며, 이것이 바로 \(\mA\)이다.

\emph{덧붙임.} 준동형정리에 따라 \(\overline{\mG}=\mG/\mN\)으로 놓으면 \(\mA\)는 반드시 \(\mN\)을 포함한다. \(\mA\)에서 \(\overline{\mA}\)를 \(\overline{\mA}=\mA/\mN\)으로 되찾을 수 있고, 대응 \(\mA\leftrightarrow\overline{\mA}\)는 일대일이다. 식 (1)은
\[
        (\mG/\mN)/(\mA/\mN)\simeq\mG/\mA
\]
로도 쓸 수 있다.

\emph{제2동형정리.} \emph{\(\mA\)가 \(\mG\)의 부분군이고 \(\mB\)가 정규부분군이면 \(\mA\cap\mB\)는 \(\mA\)의 정규부분군이고}
\[
        \mA\mB/\mB\simeq \mA/\mA\cap\mB.
\]
\emph{증명.} 준동형 \(\mG\sim\mG/\mB\)에서는 특히 \(\mA\)의 원소 \(a\)에 잉여류 \(a\mB\)가 대응하고, 이 잉여류들은 함께 군 \(\mA\mB/\mB\)를 이룬다. 따라서 \(\mA\sim\mA\mB/\mB\)이고, 준동형정리에서 주장이 따른다.

군 \(\mG\)에서 자기 자신 또는 한 부분군으로 가는 연산자준동형상을 \(\mG\)의 \emph{자기준동형}이라 한다. 특별히 \(\mG\)가 연산자를 갖는 덧셈 아벨군이고 § 1의 예 6에서처럼 준동형의 합과 곱을 정의하면, 군에서 자기 자신으로 가는 연산자준동형들은 \emph{자기동형환}이라는 한 환을 이룬다.

\emph{연산자를 갖는 단순 아벨군의 자기동형환은 체이다.}

\emph{증명.} 다른 허용부분군이 없으므로 모든 준동형은 군을 자기 자신으로 보내거나 영군으로 보낸다. 모든 것을 0으로 보내지 않는 준동형들은 위에서 단순군에 관해 말한 바에 따라 동형이고, 한 군에서 자기 자신으로 가는 동형사상들은 군을 이룬다. 따라서 환에서 영연산자를 빼면 군이 되므로 환 자체는 체이다.

특별히 \(\mG\)가 \(\mo\)-좌측가군이고 연산자준동형들을 우측연산자로 쓰면, \(\mG\)는 \emph{왼쪽의 \(\mo\)와 오른쪽의 자기동형환에 관한 쌍가군}이 된다. 실제로 새 연산자 \(\Gamma\)들을 준동형으로 택했다는 사실은
\[
        (a+b)\Gamma=a\Gamma+b\Gamma,
        \qquad
        ra\cdot\Gamma=r\cdot a\Gamma
\]
로 나타난다. 이 식들은 이미 앞서 해석한 나머지 식들과 함께 쌍가군을 특징짓는다.

역으로 한 쌍가군이 주어지면 바로 이 \emph{식들}은 우측 곱셈자영역의 각 원소가 좌측연산자들에 관한 자기준동형을 유도한다는 것과 그 역을 표현한다.

\subsection*{§ 3. 합성열.}

\(\mG\)에 허용부분군들의 유한열
\[
        \mG\supset \mA_1\supset \mA_2\supset \ldots\supset \mA_r=\mE\tag{2}
\]
이 존재하고 --- \(\mE\)는 단위원 \(e\)만으로 이루어진 군이다 --- 각 \(\mA_i\)가 앞선 군의 정규부분군이며, 열의 연속한 두 항 사이에 같은 성질을 가진 다른 항을 끼워 넣을 수 없으면 이 열을 \emph{합성열}이라 한다.

몫군
\[
        \mG/\mA_1,\quad \mA_1/\mA_2,\ldots,\quad \mA_{r-1}/\mE
\]
를 \emph{합성인자}라 한다. 이들은 \emph{단순군}, 곧 자기 자신과 단위군 이외의 허용 정규부분군을 갖지 않는 군이다.

연산자들 가운데 특히 모든 내부자기동형을 포함하면 열 (2)는 모두 정규부분군으로 이루어지고 이를 \emph{주열}이라 한다. 모든 자기동형을 포함하면 \emph{특성열}이라 한다. § 1의 뒤쪽 예들에서는 \(\mo\)의 아이디얼, 가군 또는 쌍가군의 합성열이 된다.

\emph{조르당--횔더 정리.} \emph{한 군 \(\mG\)에 서로 다른 두 합성열}
\[
\mG\supset \mA_1\supset \mA_2\supset \ldots\supset \mA_r=\mE,
\qquad
\mG\supset \mB_1\supset \mB_2\supset \ldots\supset \mB_s=\mE
\]
\emph{이 존재하면 길이는 같아 \(r=s\)이고, 합성인자들}
\[
        \mG/\mA_1,\quad \mA_1/\mA_2,\ldots,\quad\mA_{r-1}/\mE
\]
\emph{은 어떤 순서로}
\[
        \mG/\mB_1,\quad \mB_1/\mB_2,\ldots,\quad\mB_{s-1}/\mE
\]
\emph{와 동형이다.}

증명은 예컨대 E. Noether, \emph{Abstrakter Aufbau usw.}, Math. Annalen 96 (1926), § 10, 57쪽을 보라. 그곳에서는 동시에 다음도 증명한다.

\emph{\(\mG\)에 합성열이 하나 존재하면 \(\mG\)의 모든 정규부분군 \(\mH\)를 지나는 합성열을 그을 수 있다.}

합성열의 존재는 유한군에 대해 분명하며, 더 일반적으로 극대조건과 극소조건이 성립하는 군들에 대해서도 분명하다.

\emph{극대조건}은 부분군들의 모든 집합에 그 집합의 다른 부분군에 더 이상 포함되지 않는 극대부분군이 존재한다는 뜻이다. 이는 모든 상승사슬 \(\mA_1\subset\mA_2\subset\mA_3\ldots\)이 유한개의 항 뒤에 멈춘다는 것과 동치이다.

\emph{극소조건}은 부분군들의 모든 집합에 그 집합의 다른 부분군을 더 이상 포함하지 않는 극소부분군이 존재한다는 뜻이다. 또는 모든 하강사슬 \(\mA_1\supset\mA_2\supset\mA_3\ldots\)이 유한개의 항 뒤에 멈춘다는 뜻이다.

\emph{부분군으로 정규부분군만 허용하면}, 예컨대 아벨군의 경우에는 \emph{합성열의 존재에서 역으로 극대조건과 극소조건이 따른다.}

\emph{증명.} 모든 정규부분군은 합성열을 가지며, 그 길이를 그 정규부분군의 길이라고 한다. \(\mA\subset\mB\)이면 \(\mB\)의 합성열을 \(\mA\)를 지나도록 그을 수 있으므로 \(\mA\)의 길이는 \(\mB\)의 길이보다 작다. 따라서 최소길이의 모든 부분군은 극소이고 최대길이의 모든 부분군은 극대이다.

\subsection*{§ 4. 직접곱과 직접교집합.}

다음 조건들이 성립하면 군 \(\mG\)를 두 인자의 직접곱 \(\mG=\mA\times\mB\)라 한다.
\[
\begin{array}{ll}
1.& \mA,\mB\text{는 }\mG\text{의 정규부분군이다},\\
2.& \mA\mB=\mG,\\
3.& \mA\cap\mB=\mE.
\end{array}
\]
\emph{이 정의는 다음과 분명히 동치이다.} \(\mG\)의 모든 원소 \(g\)는 \(a\in\mA\), \(b\in\mB\)인 \(g=ab\)로 유일하게 표현되고, \(\mA\)의 원소들은 \(\mB\)의 원소들과 가환한다, 즉 \(ab=ba\)이다.

\(\mB_i=\mA_1\ldots\mA_{i-1}\mA_{i+1}\ldots\mA_n\)으로 놓을 때 모든 \(i\)에 대해 \(\mG=\mA_i\times\mB_i\)가 직접곱이면, 군 \(\mG\)를 \(n\)개 인자의 직접곱 \(\mG=\mA_1\times\ldots\times\mA_n\)이라 한다.

\emph{이 정의는 다시 다음과 동치이다.} \(\mG\)의 모든 \(g\)는
\[
        g=a_1a_2\ldots a_n,\qquad a_i\in\mA_i
\]
로 유일하게 표현되고, \(\mA_i\)의 원소들은 \(\mA_k\)의 원소들과 가환한다.

다음 사실들을 자주 쓴다.
\begin{enumerate}
\item \(\mA_1\times\ldots\times\mA_n=\mK\), \(\mK\times\mH=\mG\)이면 \(\mA_1\times\ldots\times\mA_n\times\mH=\mG\)이다.
\item \(\mG=\mA_1\times\ldots\times\mA_n=\mC_1\times\ldots\times\mC_n\)이고 \(\mC_i\subseteqq\mA_i\)이면 \(\mC_i=\mA_i\)이다. 이는 \(\mA_i\)의 원소들을 \(\mC_j\)들로 나타내는 표현에서 따른다.
\item \(\mG=\mA\times\mB\), \(\mK\supseteqq\mA\)이면 \(\mK=\mA\times(\mK\cap\mB)\)이다. 이는 \(\mK\)의 원소들을 \(ab\) 꼴로 나타낼 때 둘째 인자가 \(\mK,\mB\)에 모두 속한다는 데서 따른다.
\item \(\mG=\mA\times\mB\)이면 \(\mA\simeq\mG/\mB\), \(\mB\simeq\mG/\mA\)이다(제2동형정리). 여기서 다시 다음이 따른다.
\item \(\mG=\mA\times\mB\)이고 \(\mA,\mB\)가 각각 길이 \(m,n\)의 합성열을 가지면 \(\mG\)는 길이 \(m+n\)의 합성열을 가진다.
\item \(\mA\simeq\overline{\mA}\), \(\mB\simeq\overline{\mB}\)이면 \(\mA\times\mB\simeq\overline{\mA}\times\overline{\mB}\)이다.
\end{enumerate}

단위군 \(\ne\mE\)와 다른 인자들의 직접곱으로 표현되지 않는 군을 직접분해불능이라 한다.

분명히 \emph{극소조건을 갖는 모든 군은 유한개의 직접분해불능군의 직접곱이다}\footnote{가환인 경우 W. Krull이, 일반적인 경우 O. Schmidt가 보였듯이(각주 6 참조), 극대조건과 극소조건 아래에서 이 표현은 동형을 제외하면 유일하다. 직접분해불능이라는 개념을 바꾸지 않고 모든 내부자기동형을 연산자영역에 더할 수 있으므로, 정규부분군에 대한 유한성 조건 또는 주열의 존재만을 요구해도 충분하다.}.

문제의 군들을 덧셈으로 쓰면 곱과 직접곱이라는 개념은 합과 직접합으로 바뀐다. 군들의 합은 \((\mA_1,\mA_2,\ldots)\), 직접합은 \(\mA_1+\mA_2+\ldots\)로 쓴다.

\glqq{}직접교집합\grqq{}이라는 개념은 뒤에서 쓰지 않지만, 직접교집합과 직접곱의 관계에 관한 다음 정리들은 다음 장의 아이디얼론을 합 대신 교집합으로도 정식화할 수 있음을 보여 준다. 이로써 알려진 가환이론들과의 관계가 세워진다.

다음 조건들이 성립하면 군 \(\mD\)를 \(\mG\)에 관한 두 군 \(\mR,\mS\)의 직접교집합이라 한다.
\[
\begin{array}{ll}
1.& \mR,\mS\text{는 }\mG\text{의 정규부분군이다},\\
2.& \mR\cap\mS=\mD,\\
3.& \mR\mS=\mG.
\end{array}
\]
마찬가지로 \(\mR_i=\mS_1\cap\ldots\cap\mS_{i-1}\cap\mS_{i+1}\cap\ldots\cap\mS_n\)으로 놓을 때 모든 \(i\)에 대해 \(\mD=\mR_i\cap\mS_i\)가 직접교집합이면, \(\mD\)를 \(n\)개 군 \(\mS_1,\ldots,\mS_n\)의 직접교집합이라 한다.

정의에 따라 \(\mD\)는 \(\mG\)의 정규부분군이어야 한다. \(\mD=\mS_1\cap\ldots\cap\mS_n\)이 직접교집합이고 준동형 \(\mG\sim\mG/\mD\)에서 군 \(\mG,\mS,\mD\)가 \(\overline{\mG},\overline{\mS},\overline{\mE}\)로 가면 \(\overline{\mE}=\overline{\mS}_1\cap\ldots\cap\overline{\mS}_n\)은 직접교집합이다. 역으로 그러한 모든 표현 \(\overline{\mE}=\overline{\mS}_1\cap\ldots\cap\overline{\mS}_n\)에서 직접표현 \(\mD=\mS_1\cap\ldots\cap\mS_n\)으로 돌아갈 수 있다. 이 대응으로 임의의 \(\mD\)에 대한 직접교집합 정리들은 특수한 경우 \(\mD=\mE\)로 환원된다.

\(\mD=\mE\)이고 인자가 둘이면 직접곱과 직접교집합의 조건은 일치한다. \(n\)개 인자에 대해서는 직접곱과 직접교집합 사이에 다음 일대일 관계가 있다.

\emph{I. \(\mG=\mA_1\times\ldots\times\mA_n=\mA_i\times\mB_i\)이면 \(\mE=\mB_1\cap\ldots\cap\mB_n=\mB_i\cap\mC_i\)는 직접교집합이고 \(\mC_i=\mA_i\)이다.}

\emph{II. \(\mE=\mS_1\cap\ldots\cap\mS_n=\mR_i\cap\mS_i\)가 직접교집합이면 \(\mG=\mR_1\times\ldots\times\mR_n=\mR_i\times\mT_i\)이고 \(\mT_i=\mS_i\)이다.}

\emph{I의 증명.} \(\mB=\mB_1\cap\ldots\cap\mB_n=\mB_i\cap\mC_i\)라 하자. 그러면 \(\mC_i\supseteqq\mA_i\)이고, \(\mC_i\subseteqq\mA_i\)임을 보이면 된다. 예컨대 \(c\in\mC_1\)이면 \(c\in\mB_2,\ldots,\mB_n\)이므로 \(\mA\)들에 의한 성분표현은
\[
        c=a_1a_2\ldots a_n
         =a_1ea_3\ldots a_n
         =a_1a_2e\ldots a_n
         =\ldots
         =a_1\ldots a_{n-1}e
\]
이다. 모든 \(\mA_i\)의 곱이 직접곱이므로 이 표현들은 일치한다. 첫 자리를 제외한 각 자리에는 한 번씩 \(e\)가 있으므로 \(c=a_1\), 곧 \(\mC_1\subseteqq\mA_1\)이다. 따라서 \(\mC_1=\mA_1\)이고 마찬가지로 \(\mC_i=\mA_i\)이다. 이제 \(\mB=\mB_i\cap\mC_i=\mB_i\cap\mA_i=\mE\), \(\mB_i\mC_i=\mB_i\mA_i=\mG\)이므로 이는 직접교집합이다.

\emph{II의 증명.} \(\mH=\mR_1\ldots\mR_n=\mR_i\mT_i\)라 하자. 그러면 \(\mT_i\subseteqq\mS_i\)이고, \(\mS_i\subseteqq\mT_i\)임을 보이면 된다. 예컨대 \(s\in\mS_1\)이라 하자. \(\mG=\mS_i\times\mR_i\)에 따라
\[
        s=s_1e=s_2r_2=\cdots=s_nr_n.
\]
\(t_1=er_2\ldots r_n=t_ir_i\)로 놓으면, \(\mR_i,\mS_i\)의 원소별 가환성에 따라
\[
        st_1^{-1}=s_it_i^{-1}
\]
이다. 이는 모든 \(\mS_i\)의 원소이므로 \(e\)와 같다. 따라서 \(\mS_1\subseteqq\mT_1\), 곧 \(\mT_1=\mS_1\)이고, 마찬가지로 \(\mT_i=\mS_i\)이다. 그러면 \(\mH=\mR_i\mT_i=\mR_i\mS_i=\mG\), \(\mR_i\cap\mT_i=\mR_i\cap\mS_i=\mE\)이므로 \(\mG\)는 \(\mR_i\)들의 직접곱이다.

\clearpage

\providecommand{\mC}{\mathfrak{C}}
\providecommand{\mD}{\mathfrak{D}}
\providecommand{\mK}{\mathfrak{K}}
\providecommand{\mR}{\mathfrak{R}}
\providecommand{\mS}{\mathfrak{S}}
\providecommand{\mT}{\mathfrak{T}}
\providecommand{\mO}{\mathfrak{O}}
\providecommand{\ma}{\mathfrak{a}}
\providecommand{\mb}{\mathfrak{b}}
\providecommand{\md}{\mathfrak{d}}
\providecommand{\mf}{\mathfrak{f}}
\providecommand{\mh}{\mathfrak{h}}
\providecommand{\mn}{\mathfrak{n}}
\providecommand{\mt}{\mathfrak{t}}
\providecommand{\mx}{\mathfrak{x}}
\providecommand{\mz}{\mathfrak{z}}
\providecommand{\mL}{\mathfrak{L}}
\providecommand{\iso}{\simeq}
\providecommand{\ann}{\operatorname{Ann}}
\providecommand{\tuple}[1]{(#1)}

\subsection*{§ 5. 완전가약군.}

한 군이 유한개의 단순군의 직접곱
\[
        \mG=\mA_1\times\ldots\times\mA_n
\]
이면 그 군을 \emph{완전가약}이라 한다. 이 경우 군들
\[
        \mA_1\times\ldots\times\mA_n
        \supset \mA_1\times\ldots\times\mA_{n-1}
        \supset\ldots\supset \mA_1\supset \mE
\]
은 합성열을 이룬다. 그 길이는 \(n\)이고 합성인자들은
\[
        \simeq\mA_n,\ \mA_{n-1},\ldots,\mA_1
\]
이다.

\subsection*{§ 6. 체에 관한 가군. 초복소계.}

한 체, 반드시 가환일 필요는 없는 체에 관한 유한기저 가군은 연산자를 갖는 완전가약군의 거의 자명한 예이다.

\(\mG\)를 \(\mathsf{K}\)-우측가군이라 하고 \(\mathsf{K}\)의 단위원 \(e\)가 동시에 단위연산자, 곧 모든 \(a\in\mG\)에 대해 \(ae=a\)라고 하자. 한 원소 \(a\)에서 생성되는 모든 부분가군 \(a\mathsf{K}\)는 단순하다. 실제로 그 안의 임의의 비영원소에서 생성되는 가군과 같다. 따라서 \(\mH\)가 어떤 원소들에서 생성되는 부분가군이고 \(a\notin\mH\)이면 \(\mH\cap a\mathsf{K}=\mE\), 따라서 \(\mH+a\mathsf{K}\)는 직접합이다. 이처럼 한 기저원소 \(a_1\)에서 시작하여 새 기저원소 \(a_2,a_3,\ldots\)를 계속 더하면 \(\mG\)를 직접합
\[
        \mG=a_1K+a_2K+\cdots+a_nK
\]
으로 얻는다. \emph{따라서 \(\mG\)는 완전가약이다.}

합성열의 길이 \(n\)을 \(\mathsf{K}\)에 관한 \(\mG\)의 \emph{랭크}라 한다. 기저 \(a_1,\ldots,a_n\)은 원소들의 표현이 유일하고 \(e\)를 단위연산자로 가정했으므로 선형독립이다.

모든 부분가군 \(\mH\)는 직접가수이다.
\[
        \mG=\mH+\mR
        =(h_1K+\ldots+h_sK)+(r_1K+\ldots+r_{n-s}K).
\]
달리 말해, \emph{\(\mH\)의 모든 선형독립기저를 \(\mG\)의 선형독립기저로 보충할 수 있다.} \(r_i\)들은 원래의 기저원소 \(a_i\)들 가운데서 택할 수도 있다.

\((a_1,\ldots,a_n)\), \((c_1,\ldots,c_n)\)을 선형독립기저라 하자. 그러면
\[
        c_k=\sum a_i\pi_{ik},
\]
또는 행렬로
\[
        (c_1,\ldots,c_n)=(a_1,\ldots,a_n)P;
        \qquad
        P=\begin{pmatrix}
        \pi_{11}&\ldots&\pi_{1n}\\
        \vdots&&\vdots\\
        \pi_{n1}&\ldots&\pi_{nn}
        \end{pmatrix}.
\]
역으로
\[
        (a_1,\ldots,a_n)=(c_1,\ldots,c_n)Q.
\]
따라서
\[
        (a_1,\ldots,a_n)=(a_1,\ldots,a_n)PQ,
        \qquad
        (c_1,\ldots,c_n)=(c_1,\ldots,c_n)QP.
\]
\(a_i,c_i\)들이 선형독립이므로
\[
        PQ=QP=E.
\]

동시에 환이기도 한 특수한 \(K\)-가군이 \glqq{}초복소계\grqq{}이다.

가환체 \(K\)에 관한 우측가군이기도 한 환 \(\mo\)가 다음 조건들을 만족하면 \emph{\(\mathsf{K}\)에 관한 초복소계}, 문헌에서 말하는 \glqq{\(\mathsf{K}\) 위의 대수\grqq{}}라 한다.
\begin{enumerate}
\item 랭크가 유한하다. 선형독립기저를 \(u_1,\ldots,u_n\)이라 하자.
\item \(ab\cdot\varkappa=a\cdot b\varkappa=a\varkappa\cdot b\)이다. 이를 \(\mathsf{K}\)가 \(\mo\)와 가환적으로 결합되어 있다고 표현한다.
\item \(\mathsf{K}\)의 단위원 \(e\)가 동시에 단위연산자이다. 즉 모든 \(a\in\mo\)에 대해 \(a\varepsilon=a\)이다.\footnotemark
\end{enumerate}
\footnotetext{조건 2는 \(\mo\)를 \(\mo\)-좌측가군으로 보거나 \(\mo\)-우측가군으로 볼 때 \(\mathsf{K}\)의 한 원소에 의한 곱셈이 \(\mo\)의 준동형이라는 뜻이다. 조건 3에 따라 이 준동형들은 실제로 동형이다.}

\[
        \Bigl(\sum_i u_i\alpha_i\Bigr)
        \Bigl(\sum_k u_k\beta_k\Bigr)
        =\sum_{i,k}(u_i u_k)(\alpha_i\beta_k)
\]
이므로, \(\mathsf{K}\) 이외에 각 곱 \(u_i u_k\)를 \(u_j\)들로 나타내는 방법
\[
        u_i u_k=\sum_j u_j\gamma^j_{ik}
\]
을 주는 \emph{곱셈표}가 주어지면 계는 유일하게 결정된다. \(\gamma^j_{ik}\)들은 결합법칙에서 나오는 알려진 관계들을 만족해야 한다.

\(\mo\)가 단위원 \(e\), 곧 모든 \(a\)에 대해 \(ea=ae=a\)인 원소를 갖는다고 가정하는 경우가 많다. 그러면 \(K\)의 원소 \(x\)를 \(e\varkappa\)와 동일시하여 \(\mathsf{K}\)를 \(\mo\)의 부분체로 볼 수 있다. 이때 초복소계는 \emph{중심에 놓인 한 체에 관하여 유한랭크인 환}이라고도 기술할 수 있다.

유한군의 원소들을 기저원소 \(u_i\)로 삼고 군의 곱셈을 곱셈으로 삼아 얻는 \emph{군환}은 초복소계의 한 예이다. 이때 체 \(\mathsf{K}\)는 임의이다.

\subsection*{§ 7. 행렬.}

한 환 \(\mo\)의 원소 \(\pi_{ik}\)로 된 정방행렬들은 통상적인 행렬곱
\[
        \sum_j\pi_{ij}\varrho_{jk}=\sigma_{ik}
\]
과 덧셈 \(\pi_{ik}+\varrho_{ik}=\sigma_{ik}\)에 관하여 환을 이룬다.

체 \(\mathsf{K}\)의 원소로 된 행렬이 우측역행렬과 좌측역행렬을 모두 가지면 \emph{정칙}이라 한다.

§ 6에서 보았다. \emph{\(\mathsf{K}\)-가군의 두 선형독립기저 사이의 변환행렬은 정칙이다.}

\emph{정칙이기 위해서는 우측역행렬만으로 충분하다.} 실제로 선형독립기저 \((a_1,\ldots,a_n)\)을 갖는 \(\mathsf{K}\)-우측가군, 곧 부정원 \(a_1,\ldots,a_n\)의 선형형식가군을 만들고
\[
        (c_1,\ldots,c_n)=(a_1,\ldots,a_n)P
\]
로 놓자. 그러면
\[
        (c_1,\ldots,c_n)P^{-1}
        =(a_1,\ldots,a_n)PP^{-1}
        =(a_1,\ldots,a_n).
\]
따라서 \((c_1,\ldots,c_n)\)은 다시 기저이고 선형독립이므로 행렬 \(P\)는 정칙이다. 또한 우측역행렬은 좌측역행렬과 같다.

마찬가지로 좌측역행렬만으로도 충분하다.

\emph{\(P\)가 좌측역행렬을 가지면 \(P\)는 좌측 영인자가 아니다.}
\emph{증명.} \(PA=0\)에서 \(A=P^{-1}PA=0\)이 따른다.

따라서 \(P\)는 우측역행렬도 하나만 가지며, § 6에 따라 그것은 좌측역행렬과 일치한다.

통상대로 \(P_\varkappa\)가 행렬 \((\pi_{ik}\varkappa)\)를 뜻하고 \(C_{ik}\)가 \(ik\) 자리에는 단위원, 나머지 자리에는 0을 갖는 행렬이라 하면
\[
        (\pi_{ik})=\sum C_{ik}\pi_{ik}.
\]
따라서 \(\mathsf{K}\)의 행렬들은 랭크 \(n^2\)인 \(K\)-가군을 이룬다. \(C_{ik}\)들은
\[
        C_{ij}C_{jk}=C_{ik},
        \qquad
        C_{ij}C_{\bar j k}=0,\quad j\ne\bar j
\]
를 만족한다. 그러므로 \(\mathsf{K}\)의 행렬들은 유한 \(\mathsf{K}\)-가군이고, \(\mathsf{K}\)가 가환이면 초복소계를 이룬다.

\section*{제II장. 비가환 아이디얼론.}

\subsection*{§ 8. 환의 준동형정리.}

\(\ma\)가 \(\mo\)의 양쪽아이디얼이면 잉여류영역 \(\mo/\ma\)는 \(\mo\)-가군일 뿐 아니라 환이기도 하다는 것은 쉽게 알 수 있다. \(\mo\)의 모든 준동형상은 다시 환이며, \(\ma\)가 양쪽아이디얼인 잉여류환 \(\mo/\ma\)와 동형이다. 증명은 군의 경우와 같다.

특히 \(\mo\)가 체이면 \((0)\)과 \(\mo\) 이외의 아이디얼은 없다. 따라서 이 경우 준동형은 동형이거나 모든 원소를 0에 대응시킨다.

특히 환 \(\mo\)가 \(\mathsf{K}\)-우측가군이고, \(\mathsf{K}\)가 원소별로 \(\mo\)와 가환하는 환이라고 하자.
\[
        ab\cdot x=a\cdot bx=ax\cdot b
\]
(예컨대 가환체 \(\mathsf{K}\)에 관한 초복소계를 다루는 경우가 그렇다.) 그러면 허용되는 우측ㆍ좌측ㆍ양쪽아이디얼은 동시에 \(\mathsf{K}\)-가군인 것들뿐이며, 환준동형에 더하여 \(\mathsf{K}\)와의 곱셈에 관한 연산자준동형도 요구한다. 허용아이디얼은 \(\mo\)와 \(\mathsf{K}\)에 관한 쌍가군이 된다. (초복소계에서 단순히 \glqq{}아이디얼\grqq{}이라 할 때에는 언제나 바탕 계수체 \(\mathsf{K}\)에 관하여 허용되는 것만을 뜻한다.) 이 확장에서도 위의 준동형정리는 성립한다.

환준동형의 한 예는 승수영역 \(\mo\)에서 절대승수영역으로 넘어가는 것이다(§ 1, 예 4). 따라서 절대승수영역은 잉여류환 \(\mo/\md\)와 동형이며, 여기서 \(\md\)는 \(\mM\)을 소멸시키는 \(\mo\)의 원소들이 이루는 양쪽아이디얼이다.

준동형정리가 성립하므로 § 2에서와 같이 환동형과 연산자동형에 관한 제1ㆍ제2 동형정리가 따른다.

\(\mo\)의 부분집합들의 모든 곱 \(\mA\mB\)는 가군곱으로 이해한다. 곧 \(\mA\mB\)는 \(a\in\mA\), \(b\in\mB\)인 모든 합 \(\sum ab\)의 전체이고, \(a\mB\)는 \(b\in\mB\)인 모든 합 \(\sum ab\)의 전체이다. \(\mB\)가 어떤 환을 연산자영역으로 하는 우측가군이면 곱 \(\mA\mB\), 또는 \(a\mB\)도 우측가군이다. (특히 우측아이디얼, 계수영역 \(\mathsf{K}\)에 관한 허용아이디얼 등이다.) \((\mA,\mB)\)가 가군 \(\mA,\mB\)의 합을 나타낼 때 다음 계산법칙이 성립한다.
\[
        \mA\mB\cdot\mC=\mA\cdot\mB\mC,
\]
\[
        \mA(\mB,\mC)=(\mA\mB,\mA\mC),
        \qquad
        (\mA,\mB)\mC=(\mA\mC,\mB\mC).
\]
직접합은 \(\mA+\mB\)로 나타낸다(§ 4).

이하에서는 우측아이디얼에 관하여 극대조건 또는 극소조건(§ 3)을 만족하는 환을 자주 고찰한다. 특히 초복소계가 여기에 속한다. 위의 약속에 따라 \(\mathsf{K}\)-가군만 아이디얼로 허용되며, 따라서 그 랭크가 유계이기 때문이다(§ 6).

\subsection*{§ 9. 멱등원. 우측아이디얼로의 직접합분해.}

\(\mo\)를 단위원을 갖는 환이라 하자.

\(\mr\)가 우측아이디얼이면 \(\mr\mo\subseteqq\mr\)이고, 단위원 때문에 사실 \(\mr\mo=\mr\)이다. 좌측아이디얼에 대해서도 마찬가지로 \(\mo\ml=\ml\)이다.

원소 \(c\)가 \(c^2=c\)를 만족하면 \emph{멱등원}이라 한다.

\emph{\(\mo=\mr_1+\ldots+\mr_n\)이 우측아이디얼들로의 분해이고}
\[
        e=e_1+\ldots+e_n,
        \qquad e_i\in\mr_i
\]
\emph{이면}
\[
        e_i^2=e_i,
        \qquad e_i e_k=0\quad(i\ne k)
        \quad\hbox{(\glqq{}직교관계\grqq{})},
\]
\[
        \mr_i=e_i\mo
\]
\emph{이다.}

\emph{증명.} \(r\)을 \(\mr_1\)의 한 원소라 하자. 그러면
\[
        r=er=e_1r+\ldots+e_nr,
        \qquad e_ir\in\mr_i
\]
이지만, 한편
\[
        r=r+0+\ldots+0
\]
이기도 하므로
\[
        e_1r=r,
        \qquad e_ir=0\quad(i\ne1)
\]
이다. 첫째 식에서 \(\mr_1\subseteqq e_1\mo\)가 따르고, 다른 한편 \(e_1\mo\subseteqq\mr_1\)이므로 \(e_1\mo=\mr_1\)이다. \(r=e_1\)로 특수화하면
\[
        e_1^2=e_1,
        \qquad e_i e_1=0\quad(i\ne1)
\]
을 얻는다. 지수 1을 \(k\)로 바꾸어도 같은 결과가 성립한다.

\emph{역으로 \(e=\sum e_i\), \(e_i e_k=0\) \((i\ne k)\), \(e_i^2=e_i\)이고 \(\mr_i=e_i\mo\)로 놓으면 \(\mo=\mr_1+\ldots+\mr_n\)이다.}

\emph{증명.} \(\mo\)의 모든 원소 \(r\)은
\[
        r=er=e_1r+\ldots+e_nr
\]
이므로
\[
        \mo=(e_1\mo,\ldots,e_n\mo)
\]
이다. 그런데 이 표시는 유일하다. 실제로
\[
        0=e_1a_1+\ldots+e_na_n
\]
에서 \(e_i\)를 곱하면
\[
        e_i^2a_i=e_i a_i=0
\]
이 따르기 때문이다.\footnote{또 다른 역도 성립한다. \(e_i e_k=0\), \(e_i^2=e_i\), \(c=\sum e_i\)이면 \(c\)는 환 \(e_1\mo+\ldots+e_n\mo\)의 좌측단위원이다. 합이 직접합이라는 것은 위와 아주 똑같이 알 수 있다.}

따라서 우측표현
\[
        \mo=\mr_1+\ldots+\mr_n=e_1\mo+\ldots+e_n\mo
\]
에서 좌측표현
\[
        \mo=\ml_1+\ldots+\ml_n=\mo e_1+\ldots+\mo e_n
\]
이 따른다.

\(\mr_i\)들이 직접분해불능이면 \(\ml_i\)들도 그러하다. 실제로 예컨대 \(\ml_1\)의 분해가 있다면 이는
\[
        \ml_1=\mo e_1=\mo e_1'+\mo e_1'',
        \qquad
        \mo=\mo e_1'+\mo e_1''+\mo e_2+\ldots+\mo e_n
\]
을 뜻한다. 여기서 우측분해
\[
        \mo=e_1'\mo+e_1''\mo+e_2\mo+\ldots+e_n\mo
        =\mr_1'+\mr_1''+\mr_2+\ldots+\mr_n,
\]
\[
        \mr_1\cong \mo/(\mr_2+\ldots+\mr_n)=\mr_1'+\mr_1''
\]
을 얻으므로 \(\mr_1\)은 분해가능하게 된다.

\subsection*{§ 10. 양쪽아이디얼로의 분해.}

\(\mo\)가 다시 단위원을 갖는 환이고
\begin{equation}
        \mo=\ma_1+\ldots+\ma_n\tag{1}
\end{equation}
이 \(\mo\)를 양쪽아이디얼들로 분해한 것이면, 아이디얼들에 대해서도 직교관계
\begin{equation}
        \ma_i\ma_k=0\quad(i\ne k),
        \qquad
        \ma_i^2=\ma_i                                             \tag{2}
\end{equation}
가 성립한다.

\emph{증명.} \(\mb_1=\ma_2+\ldots+\ma_n\), \(\mo=\ma_1+\mb_1\)이라 하자. 그러면 \(\ma_1\mb_1\subseteqq \ma_1\cap\mb_1=0\), 따라서 \(\ma_1\mb_1=0\)이고, 더구나
\[
        \ma_1\ma_i=0\quad(i\ne1)
\]
이다. 또한
\[
        \ma_1=\ma_1\mo=\ma_1(\ma_1+\mb_1)
        =(\ma_1^2,\ma_1\mb_1)=\ma_1^2.
\]

역으로, \emph{(1)과 (2)에서 가군 \(\ma_i\)들이 양쪽아이디얼이라는 것이 따른다.}
\emph{증명.}
\[
        \mo\ma_i=(\ma_1+\cdots+\ma_n)\ma_i=\ma_i^2=\ma_i,
        \qquad
        \ma_i\mo=\ma_i.
\]

\emph{환 \(\ma_i\)의 우측아이디얼은 \(\mo\)의 우측아이디얼이다.}
\emph{증명.}
\[
        \mr\mo=\mr(\ma_i+\mb_i)=(\mr\ma_i,\mr\mb_i)
        \subseteq(\mr,\ma_i\mb_i)=(\mr,0)=\mr.
\]

\emph{\(\mr\)가 \(\mo\)의 우측아이디얼이면 \(\mr\)는 \(\mr\ma_i\subseteqq\ma_i\)인 우측아이디얼들의 직접합이다.}
\emph{증명.} \(\mr=\mr\mo=(\mr\ma_1,\ldots,\mr\ma_n)\)이고 \(\mr\ma_i\mo=\mr\ma_i\)이므로 \(\mr\ma_i\)는 \(\mr\)에 포함되는 우측아이디얼이다. 또한 \(\mr\ma_i\subseteqq\ma_i\)이므로 합은 직접합이다.
\[
        \mr=\mr\ma_1+\ldots+\mr\ma_n.
\]

\emph{특히 \(\mr\)가 직접분해불능이면 \(\mr\)는 성분을 하나만 가질 수 있으므로 어떤 \(\ma_i\) 안에 들어가야 한다.}

이 정리들에 따라 각각의 \(\ma_i\)의 아이디얼론을 알면 \(\mo\)의 아이디얼론을 다룰 수 있다. 그러므로 대개 양쪽으로 분해불능인 환에 국한한다.

\emph{서로 다른 \(\ma_i\)에 들어 있는 두 우측아이디얼은 결코 연산자동형일 수 없다.}
\emph{증명.} \(\ma_1\) 안에 있는 아이디얼은 다른 모든 \(\ma_k\)에 의해서는 소멸되지만 \(\ma_1\)에 의해서는 소멸되지 않는다. 반면 \(\ma_2\) 안에 있는 아이디얼은 \(\ma_1\)에 의해 소멸된다.

따라서 환 \(\mo\)를 양쪽으로 분해불능인 아이디얼들의 합 \(\ma_1+\ldots+\ma_n\)으로 분해할 수 있고, 각 \(\ma_i\)를 다시 한쪽으로 분해불능인 우측아이디얼들의 합 \(\mr_i'+\mr_i''+\ldots\)로 분해할 수 있다면, 서로 연산자동형인 \(\mr\)들은 같은 \(\ma_i\)에 속하는 것들 사이에서 찾아야 한다. 그러나 다음 예가 보이듯 같은 \(\ma_i\) 안에도 서로 다른 \(\mr\)의 류, 곧 연산자동형이 아닌 류들이 있을 수 있다.

\(\mathsf{K}\)를 유리수체라 하고
\[
        \mo=e_1\mathsf{K}+e_2\mathsf{K}+u\mathsf{K}
\]
를 초복소계라 하자. \(e_i\)와 \(u\)는 \(\mathsf{K}\)의 수들과 가환하며, 곱셈표는
\[
\begin{array}{c|ccc}
     & e_1&e_2&u\\ \hline
 e_1 & e_1&0&u\\
 e_2 & 0&e_2&0\\
 u   & 0&u&0
\end{array}
\]
이다. 단위원은 \(e=e_1+e_2\)이다. \(e_i\)들은 직교관계를 만족한다. 따라서 우측분해는
\[
        \mo=e_1\mo+e_2\mo=(e_1,u)+(e_2)
\]
이고, 좌측분해는
\[
        \mo=\mo e_1+\mo e_2=(e_1)+(e_2,u)
\]
이다. \(e_2\mo\)와 \(\mo e_1\)은 분해불능임이 명백하므로 \(\mo e_2\)와 \(e_1\mo\)도 분해불능이어야 한다.

양쪽분해는 불가능하다. 그러한 분해가 있다면 한 성분은 \(e_1\mo\)를, 다른 성분은 \(e_2\mo\)를 포함해야 한다. 그런데 첫째 성분은 \(u\)를 포함하고, 둘째 성분도 \(ue_2=u\)이므로 \(u\)를 포함한다.

그러나 \(e_1\mo\)와 \(e_2\mo\)는 \(\mathsf{K}\)에 관한 랭크가 서로 다르므로 연산자동형이 아니다.

\emph{\(\mo=\ma_1+\ldots+\ma_n\)이고 \(\ma_i\)들이 양쪽으로 분해불능이면, 이들은 유일하게 결정된다.}

\emph{증명.} 또한 \(\mo=\mc_1+\ldots+\mc_n\)이라 하자. 그러면
\[
        \ma_i=\ma_i\mo=(\ma_i\mc_1,\ldots,\ma_i\mc_n).
\]
마지막 합은 \(\ma_i\mc_k\subseteqq\mc_k\)이고 동시에 \(\ma_i\mc_k\subseteqq\ma_i\)이므로 직접합이다. 그런데 \(\ma_i\)들이 분해불능이므로 \(\ma_i\mc_k\)는 하나, 곧 \(\ma_i\mc_{j_i}\)를 제외하면 모두 0이어야 한다. 따라서
\[
        \ma_i=\ma_i\mc_{j_i}.
\]
마찬가지로
\[
        \mc_{j_i}=\mo\mc_{j_i}
        =\ma_1\mc_{j_i}+\cdots+\ma_n\mc_{j_i},
\]
이고 \(\ma_i\mc_{j_i}\ne0\)이므로 나머지는 모두 0이며,
\[
        \mc_{j_i}=\ma_i\mc_{j_i}=\ma_i
\]
이다. 증명 끝.

\subsection*{§ 11. 중심.}

환 \(\mo\)의 중심 \(\mZ\)는 모든 환원소와 가환하는 \(z\)들의 전체이다(모든 \(a\)에 대해 \(za=az\)). \(\mZ\)는 가환환이다.

\(\mo\)의 모든 한쪽아이디얼 \(\ma\)는 \(\mZ\) 안에 \glqq{}수축아이디얼\grqq{} \(\ma\cap\mZ\)을 갖는다. 실제로 \(\ma\)와 \(\mZ\)가 모두 \(\mZ\)-가군이므로 \(\ma\cap\mZ\)도 그러하다.

\(\mZ\)의 모든 아이디얼 \(\mA\)는 확장아이디얼, 곧 \(\mA\)가 \(\mo\)에서 생성하는 아이디얼 \(\ma=(\mA,\mo\mA)\)를 갖는다. 이것은
\[
        \ma\mo=(\mo\mA,\mA)\mo=(\mo\mA\mo,\mA\mo)
        = (\mo^2\mA,\mo\mA)\subseteqq \mo\mA\subseteqq\ma
\]
이므로 양쪽아이디얼이다. \(\mo\)에 단위원이 있다고 가정하면 \(\ma=\mo\mA\)로 쓸 수 있다. 이 경우 다음 정리가 성립한다.

\emph{\(\mo\)의 모든 양쪽분해는 중심의 한 분해와 일대일로 대응한다. 즉}
\begin{equation}
        \mo=\ma_1+\ldots+\ma_n,
        \qquad \mA_i=\ma_i\cap\mZ                         \tag{1}
\end{equation}
\emph{에서}
\begin{equation}
        \mZ=\mA_1+\ldots+\mA_n,
        \qquad \mo\mA_i=\ma_i                              \tag{2}
\end{equation}
\emph{가 따르며, 역도 성립한다.}

\emph{증명.} \(z\)를 중심원소라 하고
\begin{equation}
        z=z_1+\ldots+z_n\tag{3}
\end{equation}
을 (1)에 따른 그 분해라 하자. 그러면 \(z_i\)들은 중심원소이다. 실제로
\[
        za=z_1a+\ldots+z_na=az=az_1+\ldots+az_n
\]
이고, 합이 직접합이며 \(\ma_i\)들이 양쪽아이디얼이므로
\[
        z_i a=az_i
\]
이다. 따라서 \(z_i\in\ma_i\cap\mZ=\mA_i\)이고, (3)은 주장한 합분해를 뜻한다. 특히 \(e=\sum e_i\)가 되므로 \(e_i\)들은 중심원소이다. 끝으로
\[
        z=ze_1+\ldots+ze_n
\]
이므로
\[
        \mA_i=\mZ e_i,
        \qquad
        \mo\mA_i=\mo\mZ e_i=\mo e_i=\ma_i.
\]

역으로 \(\mZ=\mA_1+\ldots+\mA_n\)이 중심의 한 분해이면 § 9에 따라
\[
        \mA_i=\mZ e_i,
        \qquad e_i^2=e_i,
        \qquad e_i e_k=0\quad(i\ne k).
\]
따라서 직교관계와 § 9에 의해
\[
\begin{aligned}
        \mo&=\mo e=\mo e_1+\ldots+\mo e_n
        =\mo\mZ e_1+\ldots+\mo\mZ e_n\\
        &=\mo\mA_1+\ldots+\mo\mA_n
        =\ma_1+\ldots+\ma_n .
\end{aligned}
\]
이다. 이제 \(\ma_i\cap\mZ=\mA_i'\)라 놓으면 주장의 첫째 부분에서
\[
        \mZ=\mA_1'+\ldots+\mA_n'\quad\hbox{(직접합)}
\]
임을 안다. 그런데 \(\mA_i'\supseteqq\mA_i\)이므로 \(\mA_i'=\mA_i\)이다(§ 4, 2).

\emph{따름정리.} \(\ma_i\)와 \(\mA_i\)는 동시에 양쪽으로 직접분해가능하거나 동시에 그렇지 않다.

\subsection*{§ 12. 멱영아이디얼.}

아이디얼 \(\mc\)가 \(\mc^\varrho=0\)을 만족하면 \emph{멱영}이라 한다.

\emph{예.} § 10의 아이디얼 \((u)\).

\emph{\(\mo\)에 멱영 우측아이디얼 \(\mr\)가 있으면 멱영 좌측아이디얼, 나아가 양쪽아이디얼도 있다.}

\emph{증명.} \(\mr^\varrho=0\)이라 하자. 그러면 \(\mo\mr\)는 좌측아이디얼이다. \(\mo\)가 단위원을 포함하지 않아 \(\mo\mr\)가 0일 수도 있는 경우에는 \((\mo\mr,\mr)\)가 좌측아이디얼이며,
\[
        (\mo\mr)^\varrho=\mo\mr\mo\mr\ldots\mo\mr
        \subseteqq\mo\mr\mr\ldots\mr=\mo\mr^\varrho=0.
\]

\emph{두 멱영 우측아이디얼의 합은 다시 멱영 우측아이디얼이다.}

\emph{증명.} \(\mc^\varrho=0\), \(\md^\sigma=0\)이라 하자.
\[
        (\mc,\md)^{\varrho+\sigma-1}
        =(\ldots,\md\mc\ldots\md\ldots\mc,\ldots)
\]
의 각 항에는 \(\mc\) 인자가 적어도 \(\varrho\)개 있거나 \(\md\) 인자가 적어도 \(\sigma\)개 있다. 첫째 경우에는
\[
        \md\mc\cdots\md\cdots\mc\cdots
        \subseteq\md\mc\mc\cdots\mc=\md\mc^\varrho=0
\]
이고, 둘째 경우에는 \(\md\) 인자들에 관하여 마찬가지이다. 따라서
\[
        (\mc,\md)^{\varrho+\sigma-1}=0.
\]

이제 \(\mo\)의 우측아이디얼들에 대한 극대조건이 성립하면 극대한 멱영 우측아이디얼이 있다. 이것은 다른 모든 멱영 우측아이디얼을 포함한다. 그렇지 않으면 그러한 아이디얼 하나와 합을 만들어 더 큰 것을 얻을 수 있기 때문이다. 이를 \(\mc\)라 하자. \(\mo\mc\)도 멱영 우측아이디얼이므로 \(\mo\mc\subseteqq\mc\)이고, 따라서 \(\mc\)는 좌측아이디얼이다. 모든 멱영 좌측아이디얼 \(\md\)도 \(\mc\)에 포함된다. 실제로 \((\md,\md\mo)\)는 멱영 우측아이디얼이므로 \(\md\subseteqq\mc\)이다. \emph{따라서 다른 모든 우측 및 좌측 멱영아이디얼을 포함하는 극대 멱영 양쪽아이디얼 \(\mc\), 곧 \glqq{}근기\grqq{}가 있다.}

§ 10의 예에서는 \((u)\)가 극대 멱영아이디얼이다.

근기가 영아이디얼이면 \glqq{}근기 없는 환\grqq{}(\glqq{}반단순환\grqq{}, \glqq{}데데킨트계\grqq{})이라 한다. 근기에 따른 잉여류환은 언제나 근기 없는 환이다.

\emph{\(\mZ\)가 \(\mo\)의 중심이고 \(\mc\)가 \(\mo\)의 근기이면, \(\mC=\mc\cap\mZ\)는 \(\mZ\)의 근기이다.}

\emph{증명.} \(\mC\)가 멱영이라는 것은 명백하다. \(\mZ\)에 이를 진정하게 포함하는 멱영아이디얼 \(\overline{\mC}\)가 있다면, 이것을 \(\mc\)에 완전히 포함되지 않는 멱영아이디얼 \(\overline{\mc}\)로 확장할 수 있다.
\[
        (\overline{\mC}\mo)^\varrho
        =\overline{\mC}^{\varrho}\mo^\varrho=0.
\]

\emph{따름정리.} \(\mo\)가 근기 없는 환이면 \(\mZ\)도 그러하다. 그러나 § 10의 예가 보이듯 그 역은 성립하지 않는다. \(\mo/\mc\)의 중심은 \(\mZ/\mC\)와 동형인 환을 포함하지만, 이것은 진부분환일 수 있다(§ 10의 예).

\subsection*{§ 13. 완전가약환.}

§ 5에 따라, 한 환이 유한개의 단순 우측아이디얼의 직접합이면 \emph{우측 완전가약}이라 한다.

\emph{단위원을 갖는 우측 완전가약환은 \(\ne(0)\)이 아닌 멱영아이디얼을 갖지 않으며, 따라서 근기를 갖지 않는다.}

\emph{증명.} 모든 우측아이디얼 \(\mr\)는 직접가수이다(§ 5).
\[
        \mo=\mt+\mr=e_1\mo+e_2\mo.
\]
\(e_2^2=e_2\)이므로 \(e_2^\varrho=e_2\)이기도 하다. \(\mr\)가 멱영이라면 \(e_2^\varrho=0\), 따라서 \(e_2=0\), 따라서 \(\mr=0\)이다. 증명 끝.

이제 두 가지 역을 보이자. 첫째, \emph{근기가 없고 우측아이디얼들에 대한 극소조건을 만족하는 환은 우측아이디얼에 관하여 완전가약이다.}

\emph{증명.} \(\mt\ne0\)을 극소 우측아이디얼이라 하자. \(\mt\)가 멱등원을 포함함을 보이겠다.

\(\mt^2\subseteqq\mt\)이고 \(\mt^2\ne0\)이므로 \(\mt^2=\mt\)이다. \(\mt^2\)도 우측아이디얼이기 때문이다. 따라서 \(a\mt\ne0\)인 \(a\in\mt\)가 있다. 그런데 \(a\mt\)도 우측아이디얼이므로 \(a\mt=\mt\)이어야 한다.

\(a\)에 의해 소멸되는 \(\mt\)의 모든 \(b\), 곧 \(ab=0\)인 \(b\)들의 전체는 우측아이디얼이고 \(\mt\)와 같지 않으므로 0이다. 따라서 \(ab=0\)이면 \(b=0\)이다.

\(\mt=a\mt\)이므로 \(a\)는 \(ac\)의 꼴로 나타난다. 곧 \(a=ac\), \(c\ne0\)이다. 따라서
\[
        ac=ac^2,
        \qquad a(c-c^2)=0,
        \qquad c-c^2=0,
        \qquad c^2=c.
\]
이제 \(\mt\)가 직접가수임을 보이자. \(c\mo\)는 \(\mt\)에 포함되는 우측아이디얼이고, \(c^2=c\in c\mo\)이므로 0이 아니다. 따라서 \(c\mo=\mt\)이다. \(\mo\)의 모든 원소 \(r\)은
\[
        r=cr+(r-cr)
        \qquad\hbox{(한쪽 페어스 분해)}
\]
로 나타난다. \(cr\)인 원소들은 아이디얼 \(\mt\)를 이루고, \(r-cr\)인 원소들은 \(c\)에 의해 소멸되는 다른 우측아이디얼 \(\mr\)를 이룬다. 따라서
\[
        \mo=(\mt,\mr).
\]
\(\mt\)의 원소들은 \(c\)에 의해 소멸되지 않으므로 합은 직접합이다.
\begin{equation}
        \mo=\mt+\mr.                                             \tag{1}
\end{equation}

이제 극소조건에 따라 \(\mo\)를 직접분해불능인 아이디얼들의 직접합
\[
        \mo=\mr_1+\ldots+\mr_n
\]
으로 나타내면 \(\mr_i\)들은 단순, 곧 극소이어야 한다. 실제로 예컨대 \(\mr_1\)이 단순하지 않고 \(\mt\)가 \(\mr_1\) 안의 극소아이디얼이면 (1)에서
\[
        \mr_1=\mt+\mr\cap\mr_1\qquad\hbox{(§ 4, 3)}
\]
을 얻는다. 그러면 \(\mr_1\)이 분해가능하게 되어 모순이다. 따라서 \(\mo\)는 완전가약이다.

둘째, \emph{근기가 없고 극소조건을 만족하는 환은 단위원을 갖는다.}

먼저 좌측단위원의 존재를 보이기 위해서는 다음을 보이면 충분하다. \emph{우측아이디얼 \(\ma=c_1\mo\ne\mo\)가 좌측단위원 \(c_1\)을 가지면, 더 큰 길이를 가지면서 역시 좌측단위원 \(c\)을 갖는 우측아이디얼 \(c\mo\)가 존재한다.} 실제로 완전가약성, 따라서 모든 길이의 유계성은 이미 보였으므로 이 과정을 유한번 거치면 \(\mo\) 자체의 좌측단위원에 이른다.

그러므로 \(\ma=c_1\mo\), \(c_1^2=c_1\)이라 하자. 식
\[
        r=c_1r+(r-c_1r)
\]
에 의해 위에서와 같이 페어스 분해 \(\mo=c_1\mo+\mr\)가 주어진다. 위에서와 같이 \(c_2\)를 \(\mr\) 안의 멱등원이라 하자. 그러면 \(c_2^2=c_2\), \(c_1c_2=0\)이다. \(c_1\)이 \(\mr\)의 모든 원소를 소멸시키기 때문이다. 이제 \(e_1=c_1\), \(e_2=c_2-c_2c_1\)로 놓으면
\[
        e_1^2=e_1,
        \quad e_1e_2=0,
        \quad e_2e_1=0,
        \quad e_2c_2=c_2,
        \quad e_2\ne0,
        \quad e_2^2=e_2.
\]
따라서 각주 13의 논평에 따라 \(c=e_1+e_2\)는 환
\[
        c\mo=e_1\mo+e_2\mo
\]
의 좌측단위원이다. 이 환은 \(e_2^2=e_2\ne0\)이므로 우측아이디얼로서 \(\ma\)보다 더 긴 길이를 갖는다.

구성한 좌측단위원 \(e\)가 우측단위원이기도 함을 보이기 위해 좌측아이디얼들에 대한 페어스 분해
\[
        \mo=\mo e+\ml
\]
를 만들자. \(\ml e=0\), \(\ml=e\ml\)이므로 \(\ml^2=\ml e\ml=0\)이다. 가정에 따라 멱영 좌측아이디얼은 존재할 수 없으므로 \(\ml=0\)이다. 따라서 \(e\)는 우측단위원이기도 하며, 결국 단위원이다.

\emph{정리. 우측 완전가약성과 단위원의 존재에서 양쪽 완전가약성이 따른다.}
\[
        \mo=\md_1+\cdots+\md_s,
        \qquad \md_i\ \hbox{는 양쪽단순이다.}
\]

앞의 증명에서와 같이, 모든 극소 양쪽아이디얼 \(\ma\)가 직접가수임을 보이면 충분하다. \(\ma\)는 우측아이디얼로서 직접가수이다.
\[
        \mo=\ma+\mr=e_1\mo+e_2\mo.
\]
이에 대응하는 좌측분해는
\[
        \mo=\mc+\ml=\mo e_1+\mo e_2
\]
이다. 따라서
\[
        \mr\mc\subseteqq \mr\ma\subseteqq \mr\cap\ma=0,
        \qquad
        \ml\ma=\mo e_2e_1\mo=0,
\]
\[
        (\ma\ml)^2=\ma\ml\ma\ml=0,
        \qquad\hbox{따라서 }\ma\ml=0,
\]
\[
        \mr=\mr\mo=(\mr\mc,\mr\ml)=\mr\ml,
        \qquad
        \ml=\mo\ml=(\ma\ml,\mr\ml)=\mr\ml.
\]
그러므로 \(\mr=\ml\)이고, 따라서 \(\mr\)는 양쪽아이디얼이다. 이에 따라 \(\ma\)는 양쪽 직접가수이다. 그 뒤는 한쪽 완전가약성의 증명과 같다.

\(\md_i\)들은 환으로서도 양쪽단순이다. \(\md_i\)의 모든 아이디얼은 \(\mo\)의 아이디얼이기 때문이다. 이제 그 구조를 조사한다.

\subsection*{§ 14. 단위원을 갖는 양쪽단순 완전가약환.}

\(\mo\)를 그러한 환이라 하자.
\[
        \mo=\mr_1+\ldots+\mr_n=e_1\mo+\ldots+e_n\mo,
        \qquad \mr_i\ \hbox{는 단순},
\]
이면
\[
        \mo=\ml_1+\ldots+\ml_n=\mo e_1+\ldots+\mo e_n,
        \qquad \ml_i\ \hbox{는 직접분해불능}\quad\hbox{(§ 9)}
\]
이 따른다. \(\ml_i\mr_i\)는 0이 아닌 양쪽아이디얼이다(\(e_i^2\)가 \(\ml_i\mr_i\) 안에 있기 때문이다). 따라서 \(\ml_i\mr_i=\mo\)이다.

\(\mA_{ik}=\mr_i\ml_k\)로 놓으면
\[
        \mo=(\mr_1,\ldots,\mr_n)\cdot(\ml_1,\ldots,\ml_n)
        =(\ldots,\mA_{ij},\ldots)
\]
은 직접합이다. 실제로 \(0=\sum_{i,j}a_{ij}\)이면 \(a_{ij}\in\mr_i\)이고 \(\sum\mr_i\)가 직접합이므로
\[
        0=\sum_j a_{ij},
\]
이며, 다시 \(a_{ij}\in\ml_j\)이고 \(\sum\ml_j\)가 직접합이므로
\[
        0=a_{ij}
\]
이다.

또한 \(\mA_{ij}\mr_j=\mr_i\ml_j\mr_j=\mr_i\mo=\mr_i\)이므로 \(\mA_{ij}\ne0\)이다. \(\mA_{ij}\)에서 \(a_{ij}\ne0\)을 택하자. 그러면
\[
        a_{ij}\mr_j\subseteqq\mr_i,
        \qquad
        a_{ij}\mr_j\ne0\quad\hbox{(\(a_{ij}e_j=a_{ij}\)이므로)}
\]
이고, 따라서
\[
        a_{ij}\mr_j=\mr_i.
\]
각 \(r_j\in\mr_j\)에 대해
\[
        a_{ij}r_j=r_i
\]
로 놓으면 사상 \(r_j\to r_i\)는 \(\mr_j\)에서 \(\mr_i\)로 가는 연산자준동형이다. 0에 대응되는 원소들은 \(\mr_j\)에 포함되는 아이디얼을 이루므로 영아이디얼이다. 따라서 이 사상은 \emph{동형}이다. 그러므로 다음을 얻는다.

\emph{\(\mo\)의 한 분해에 나타나는 임의의 두 단순 우측아이디얼 \(\mr_i,\mr_j\)는 연산자동형이며, 동형들은 \(\mA_{ij}\)의 원소들에 의해 주어진다.} 서로 다른 임의의 두 단순 우측아이디얼 \(\mr,\mr'\)은 언제나 \(\mo\)의 적어도 한 분해에 함께 나타나므로, 그러한 임의의 두 아이디얼도 동형이다.

\emph{\(\mr_j\)에서 \(\mr_i\)로 가는 모든 준동형은 \(\mA_{ij}\)의 원소들에 의해 주어진다.}

\emph{증명.} \(e_j\to a_{ij}\)이면 \(e_j^2\to a_{ij}e_j\)이므로 \(a_{ij}=a_{ij}e_j\)는 \(\mr_i\ml_j=\mA_{ij}\)에 속한다. 또한
\[
        \mr_j=e_j\mr_j\to a_{ij}\mr_j=\mr_i.
\]

\emph{특히 \(\mr_i\)에서 자기 자신으로 가는 준동형들은 \(\mA_{ii}\)의 원소들에 의해 주어진다.}

\(\mA_{ii}\)의 두 원소의 곱에는 준동형들의 곱이, 합에는 그 합이 대응한다(정의는 § 1, 6 참조). 서로 다른 \(a_{ii}\)는 서로 다른 자기동형을 준다. 실제로 \(e_i\)는 \(a_{ii}e_i=a_{ii}\)로 간다. 따라서 환 \(\mA_{ii}\)는 \(\mr_i\)의 자기동형환과 동형이다.

그런데 단순아이디얼의 자기동형환은 체이다(§ 2). 따라서 \(\mA_{ii}\)는 체이다. 또한 모든 \(\mr_i\)가 연산자동형이므로 그 자기동형환들도 환동형이다. 곧 \(\mA_{ii}\)는 환동형을 제외하면 \(\mo\)에 의해 유일하게 결정된다. 가능한 서로 다른 \(\mA_{ii}\)들은 단순 우측아이디얼의 추상적으로 정의된 자기동형환을 구체적으로 실현하는 서로 다른 방식일 뿐이다.

\emph{주정리. 단위원을 갖는 모든 양쪽단순 완전가약환 \(\mo\)는 어떤 체 \(K\)의 원소로 된 \(n\)행 정방행렬환과 동형이다. (체 \(\mathsf{K}\)는 \(\mo\)의 우측아이디얼들의 자기동형체와 동형이다.)}

\emph{증명\textsuperscript{13a)}. 행렬단위 \(c_{ik}\)의 구성.} \(\Gamma_{11}\)을 \(\mr_1\)의 항등자기동형이라 하고, \(\Gamma_{i1}\mr_1=\mr_i\)인 임의의 동형 \(\Gamma_{i1}\)들을 택한 뒤
\[
        \Gamma_{ik}=\Gamma_{i1}\Gamma_{k1}^{-1}
\]
로 놓자. 그러면 일반적으로
\[
        \Gamma_{ik}\Gamma_{kl}=\Gamma_{il}
\]
이다. \(\Gamma_{ik}\)가 \(\mA_{ik}\)의 원소 \(c_{ik}\)에 의해 주어지면
\[
        c_{ik}=e_i c_{ik}=c_{ik}e_k,
        \qquad
        c_{ik}c_{kl}e_l=c_{il}e_l,
\]
따라서
\[
        c_{ik}c_{kl}=c_{il},
        \qquad
        c_{ij}c_{kl}=c_{ij}e_j e_k c_{kl}=0\quad(j\ne k)
\]
이다. 그러므로 \(c_{ik}\)들은 행렬단위이다(§ 7).

\footnotetext{최근 B. L. van der Waerden은 추상적 자기동형체를 직접 써서 더 간단히 증명하는 방법을 찾아냈다. 이 증명은 그의 대수학 저서(\emph{Grundlehren der mathematischen Wissenschaften}, Berlin: Julius Springer)에 실릴 예정이다. [1929. 6. 11.]}

\emph{\(\mathsf{K}\)의 구성.} 대응
\[
        a_{ii}=c_{i1}a_{11}c_{1i}
\]
는 \(\mA_{11}\)의 각 \(a_{11}\)에 \(\mA_{ii}\)의 \glqq{}켤레원소\grqq{} \(a_{ii}\)를 대응시킨다.

이 대응은 환동형이다. \(a\)와 \(c\)에 대응하는 동형 \(\Delta,\Gamma\)가
\[
        \Delta_{ii}=\Gamma_{i1}\Delta_{11}\Gamma_{i1}^{-1}
\]
로 연결되기 때문이다. 이 관계가 환동형을 나타낸다는 것은 잘 알려져 있다. 이제 원소
\[
        \alpha=a_{11}+\ldots+a_{nn}
        =a_{11}+c_{21}a_{11}c_{12}+\ldots+c_{n1}a_{11}c_{1n}
\]
를 만들자. 각 \(a_{11}\)에는 하나의 \(\alpha\)가 대응하며, 직접합이므로 이 대응은 일대일이다. 모든 \(\alpha\)의 전체는 체 \(\mathsf{K}\simeq\mA\)이다. 실제로
\[
        \alpha=a_{11}+\ldots+a_{nn},
        \qquad
        \beta=b_{11}+\ldots+b_{nn}
\]
이면
\[
        \alpha+\beta=(a_{11}+b_{11})+\ldots+(a_{nn}+b_{nn}),
        \qquad
        \alpha\cdot\beta=a_{11}b_{11}+\ldots+a_{nn}b_{nn}
\]
이다. 따라서 대응 \(a_{11}\to\alpha\)는 동형이다.

또
\[
        c_{ik}\alpha=c_{ik}a_{kk}
        =c_{ik}c_{k1}a_{11}c_{1k}=c_{i1}a_{11}c_{1k},
\]
\[
        \alpha c_{ik}=a_{ii}c_{ik}
        =c_{i1}a_{11}c_{1i}c_{ik}=c_{i1}a_{11}c_{1k},
\]
이므로 각 \(\alpha\)는 모든 \(c_{ik}\)와 가환한다.

끝으로 같은 식들에서
\[
        c_{ik}K=c_{ik}\mA_{kk}
        =c_{ik}\mr_k\ml_k=\mr_i\ml_k=\mA_{ik}
\]
가 따르며, 따라서
\[
        \mo=\sum c_{ik}\mathsf{K}.
\]

\emph{대응.} \(\mo\)의 각 원소를 \(\sum c_{ik}\alpha_{ik}\) 꼴로 나타내어 그 원소에 행렬 \((\alpha_{ik})\)를 대응시키면, 행렬곱의 정의에서 곧바로 알 수 있듯 이 대응은 동형이다. 이로써 주정리가 증명되었다.

\emph{역. 체 \(\Lambda\)의 원소로 된 \(n\)행 정방행렬환은 양쪽단순이고 완전가약이다. \(\Lambda\)는 우측아이디얼들의 자기동형환과 동형이고, \(\Lambda e\)는 앞 정리에서 구성한 \(\mathsf{K}\)이다.}

\emph{증명.} 행렬단위(§ 7)를 \(c_{ik}\)라 하고
\[
        \mr_i=\sum_k c_{ik}\Lambda
\]
로 놓자. 그러면 \(\mo=\mr_1+\ldots+\mr_n\)임은 명백하다. \(\mr_i\)들은 단순 우측아이디얼이다. 실제로 0이 아닌 모든 원소 \(\sum c_{ik}\alpha_k\ne0\)는 \(\mr_i\) 전체를 생성한다. \(\alpha_j\ne0\)이면
\[
        \Bigl(\sum c_{ik}\alpha_k\Bigr)
        (c_{jl}\alpha_j^{-1})=c_{il}
\]
이고, \(\sum c_{il}\beta_l\)은 \(\mr_i\)의 모든 원소를 지난다.

\(\mr_i\)들은 연산자동형이다. 곧 \(\mr_i=c_{ik}\mr_k\)이다.

여기서 \(\mo\)가 양쪽으로 분해불능임이 따른다. 따라서 이미 증명한 완전가약성과 단위원의 존재에 의해 § 13의 정리로부터 \(\mo\)는 양쪽단순이다. 이는 0이 아닌 \(\sum\sum c_{ik}\alpha_{ik}\) 하나가 양쪽아이디얼 \(\mo\) 전체를 생성한다는 사실에서도 알 수 있다.

또한
\[
        \mr_i=c_{ii}\mo,
        \qquad
        \ml_i=\mo c_{ii},
\]
\[
        e=\sum c_{ii},
        \qquad c_{ii}=e_i,
\]
\[
        \mA_{ii}=\mr_i\cap\ml_i=c_{ii}\Lambda\simeq\Lambda.
\]
끝으로 \(a_{11}=c_{11}\lambda\)로 놓으면
\[
        \alpha=a_{11}+c_{21}a_{11}c_{12}+\ldots
        =c_{11}\lambda+c_{22}\lambda+\ldots=e\lambda
\]
이므로 모든 것이 증명되었다.

\emph{\(\mo\)의 중심은 \(\mathsf{K}\)의 중심이며, 따라서 가환체이다.}

\emph{증명.} \(\mZ(\mathsf{K})\)는 모든 \(x\)와 가환하는 \(\zeta\)들의 전체이다. 그런데 이들은 모든 \(c_{ik}\)와도 가환하며, 따라서 모든 합 \(\sum c_{ik}a_{ik}\)와도 가환한다. 그러므로
\[
        \mZ(\mathsf{K})\subseteqq\mZ(\mo).
\]
역으로 \(z\in\mZ(\mo)\)이고
\[
        z=\sum_\mu\sum_\nu c_{\mu\nu}\gamma_{\mu\nu}
\]
라 하자.
\[
        zc_{ik}=c_{ik}z,
        \qquad
        \sum_\mu c_{\mu k}\gamma_{\mu i}
        =\sum_\nu c_{i\nu}\gamma_{k\nu}.
\]
따라서
\[
        \gamma_{\mu i}=0\quad(\mu\ne i),
        \qquad
        c_{ik}\gamma_{ii}=c_{ik}\gamma_{kk},
        \qquad
        \gamma_{ii}=\gamma_{kk}=\gamma,
\]
\[
        z=\sum c_{ii}\gamma=e\gamma=\gamma.
\]
따라서 \(z\in K\)이고 모든 \(x\)와 가환하므로 \(z\in\mZ(K)\)이다.

행렬환은 물론 좌측 완전가약이기도 하고, 단위원을 갖는 모든 우측 완전가약환은 행렬환들의 직접합이므로 다음이 따른다.

\emph{단위원을 갖는 모든 우측 완전가약환은 좌측으로도 완전가약이며, 그 역도 성립한다.}

또 다음이 성립한다. \emph{단위원을 갖는 완전가약환의 중심은 양쪽단순 행렬환들에 대응하는 가환체들의 직접합이다.}

다음 정리도 성립한다.\textsuperscript{14)}

\emph{서로 다른 두 분해 \(\mo=\sum c_{ik}\mathsf{K}\), \(\mo=\sum c'_{ik}\mathsf{K}'\)는 내부자기동형 \(\alpha'\to x^{-1}\alpha x\), \(c'_{ik}=x^{-1}c_{ik}x\)에 의해 서로 옮겨진다.}
\footnotetext{E. Artin, 앞의 논문, 258쪽 참조.}

\emph{증명.} \(\mo=\mr_1+\ldots+\mr_n=\mr_1'+\ldots+\mr_n'\)이라 하고, \(a,b\)를 아이디얼 \(\mr_1,\mr_1'\) 사이의 서로 역인 두 동형을 주는 원소들이라 하자.
\[
        \mr_1=a\mr_1',
        \qquad
        \mr_1'=b\mr_1,
        \qquad
        ab=c_{11},
        \qquad
        ba=c'_{11}.
\]
\[
        x=\sum_i c_{i1}ac'_{1i},
        \qquad
        y=\sum_i c'_{i1}bc_{1i}
\]
로 놓으면
\[
        xy=e,
        \qquad
        y=x^{-1},
        \qquad
        x^{-1}c_{ik}x=c'_{ik},
        \qquad
        x^{-1}\alpha x=\alpha'
\]
이다.

\section*{제III장. 가군론과 표현론.}

\subsection*{§ 15. 표현과 표현가군\textsuperscript{15)}.}

\(\mo\)를 환, \(\mathsf{K}\)를 단위원을 갖는 환이라 하자. (뒤의 응용에서는 \(\mathsf{K}\)가 언제나 체이다.)

\emph{\(\mo\)의 \(K\) 안에서의 \(n\)차 표현이란 준동형}
\[
        \mo\sim\mathfrak D
\]
\emph{을 말한다. 여기서 \(\mathfrak D\)는 \(K\)의 원소로 된 \(n\)행 정방행렬들이 이루는 환이다.}

\(\mo\)의 \(\mathsf{K}\)에 관한 \emph{표현가군}이란 좌측 \(\mo\)-가군이자 우측 \(\mathsf{K}\)-가군인 쌍가군 \(\mM\)(§ 1, 5)을 말한다.
\[
        \mo\mM\subseteqq\mM,
        \qquad
        \mM \mathsf{K}\subseteqq\mM.
\]
또한 이는 유한개의 한 원소 \(\mathsf{K}\)-가군의 직접합
\[
        \mM=x_1\mathsf{K}+\ldots+x_n\mathsf{K}
\]
이어야 하며, \(\mathsf{K}\)의 단위원은 항등연산자이어야 한다.

\emph{모든 표현가군은 하나의 표현을 낳는다.} \(c\in\mo\)이고
\[
        cx_k=\sum x_i\gamma_{ik},
\]
또는
\[
        c(x_1,\ldots,x_n)=(x_1,\ldots,x_n)C,
        \qquad C=(\gamma_{ik})
\]
라 하자. 그러면 행렬 \(C\)들은 \(c\)의 표현을 이룬다. 실제로
\[
        (b+c)x_k=\sum x_i(\beta_{ik}+\gamma_{ik}),
\]
\[
%
        bcx_k
        =b\sum x_j\gamma_{jk}
          =\sum b x_j\gamma_{jk}
          =\sum\sum x_i\beta_{ij}\gamma_{jk}
        =\sum x_i(\sum\beta_{ij}\gamma_{jk})
%
\]
또는
\[
        bc(x_1,\ldots,x_n)=(x_1,\ldots,x_n)BC
\]
이기 때문이다.

\emph{역으로, \(\mo\)의 모든 표현 \(\mathfrak D\)에는 하나의 표현가군, 더 정확히는 그 가군의 한 특정 기저가 속한다.} 실제로 \(\mM\)을 \(x_1,\ldots,x_n\)의 형식적 선형형식
\[
        y=\sum x_i\alpha_i
\]
전체로 이해하자. 그러면 \(\mM\)은 우측 \(\mathsf{K}\)-가군이다. 원소 \(c\)에 행렬 \(C=(\gamma_{ik})\)가 대응할 때, 더 나아가
\srcnumdisplay[3.2em]{(1)}{%
\begin{gathered}
        cx_k=\sum x_i\gamma_{ik},\\
        c\biggl(\sum_k x_k\alpha_k\biggr)
        =\sum_i\sum_k x_i\gamma_{ik}\alpha_k
\end{gathered}}
로 정의하자. 준동형관계
\[
        c+d\to C+D,
        \qquad cd\to CD
\]
에서
\[
        (c+d)x_i=cx_i+dx_i,
        \qquad cdx_i=c\cdot dx_i
\]
가 따르고, 합 \(\sum x_i\alpha_i\)에 대해서도 마찬가지이다. 나머지 쌍가군 성질
\[
        c(y+z)=cy+cz,
        \qquad c(y\alpha)=(cy)\alpha
\]
은 자명하다. 따라서 \(\mM\)은 표현가군이고, (1)에 의해 정확히 표현 \(\mo\sim\mathfrak D\)에 속한다.

이제 \((y_1,\ldots,y_n)\)을 같은 표현가군의 다른 기저라 하자.
\[
        (y_1,\ldots,y_n)=(x_1,\ldots,x_n)\cdot P,
        \qquad
        (x_1,\ldots,x_n)=(y_1,\ldots,y_n)\cdot P^{-1}.
\]
\(x\)라는 기저에 관하여 \(b\sim B\)이면
\[
        b(y_1,\ldots,y_n)=b(x_1,\ldots,x_n)P
        =(x_1,\ldots,x_n)BP
        =(y_1,\ldots,y_n)P^{-1}BP.
\]
표현 \(b\sim B\)와 \(b\sim P^{-1}BP\)를 \emph{동치}라 하고 같은 \emph{표현류}에 넣는다. 모든 정칙행렬 \(P\)는 \(\mM\)의 기저를 다시 기저로 옮기므로 다음이 증명되었다.

\emph{모든 표현가군은 하나의 특정 표현류를 유일하게 결정한다.}\textsuperscript{15a)}

{\footnotesize\noindent
15) 가환체에 관한 표현가군은 W. Krull, \emph{Theorie und Anwendung der verallgemeinerten Abelschen Gruppen}, Sitzungsber. Heidelberger Ak. 1926, 제1논문에서 처음 나타난다.\par
15a) 교정 때의 추가(1929년 4월 14일). B. L. van der Waerden이 내게 알려 준 대로, \emph{선형변환과 행렬}이라는 개념을 분리하면 특정한 기저 선택과 무관한 불변적 관계를 직접 얻을 수 있다. 선형변환은 두 선형형식가군 사이의 준동형이고, 행렬은 특정한 기저 선택에서 이 준동형을 나타낸 표현이다. \emph{I. 모든 표현가군은 승수영역 \(\mo\)에 그 가군에서 자기 자신으로 가는 선형변환들의 한 준동형계를 유일하게 대응시킨다.} 실제로 § 2 끝부분에 따라 쌍가군 \(\mM\)의 좌측승수들은 우측연산자(여기서는 \textbf{\(\mathsf{K}\)와의 곱셈})에 관하여 \(\mM\)에서 자기 자신으로 가는 연산자준동형을 생성하며, 그 대응은 준동형이다. \emph{II. \(\mo\)와 준동형인, 한 선형형식가군에서 자기 자신으로 가는 모든 선형변환계는 표현가군을 낳는다.} 다시 § 2 끝부분에 따라 이 준동형들은 좌측승수로서 쌍가군을 낳는다. 따라서 모든 \(m\in\mM\), \(c\in\mo\)에 대해 \(cm=\Gamma m\)로 놓으면 승수영역 \(\mo\)에 관해서도 마찬가지이다. 여기서 \(c\mapsto\Gamma\)는 처음 주어진 환준동형 대응이다. 우측ㆍ좌측 곱셈에 관하여 연산자동형인 표현가군들에는 같은 변환들이, 연산자동형이 아닌 표현가군들에는 서로 다른 변환들이 대응한다. 특정한 기저를 골라 각 변환을 행렬로 나타내면 I은 표현 \(\mo\sim\mathfrak D\)를 주고, II는 그러한 모든 표현이 하나의 표현가군, 더 정확히는 그 가군의 한 특정 기저를 낳는다고 말한다. 여기서 표현가군과 표현류의 대응이 다시 따르며, 계산은 전혀 필요하지 않다.\par}

다음은 명백하다.

\emph{연산자동형인 두 표현가군에는 같은 표현류가 대응한다.}

그러나 역도 성립한다.

\emph{두 표현가군 \((x_1,\ldots,x_n)\)과 \((\bar x_1,\ldots,\bar x_n)\)이 같은 표현을 생성하면, 대응 \(x_i\to\bar x_i\), \(\sum x_i\alpha_i\to\sum\bar x_i\alpha_i\)는 동형을 준다.} 곱셈법칙이 행렬 \(C\)들에 의해 완전히 결정되기 때문이다.

특수한 경우로, \(\mM\)의 두 기저가 같은 표현을 생성하면 이들은 \(\mM\)에서 자기 자신으로 가는 한 연산자동형에 의해 서로 옮겨진다.

달리 말해, 고정된 \(P\)에 대해 모든 \(C\)가 \(C=PCP^{-1}\)을 만족하면 \(y_i=\sum_kx_k\pi_{ik}\)일 때 대응 \(y_i\mapsto x_i\)는 \(\mM\)에서 자기 자신으로 가는 동형이다. 역으로 쌍가군 \(\mM\)에서 자기 자신으로 가는 모든 동형 \(y_i\mapsto x_i\)는 모든 \(C\)와 가환하는 행렬 \(P\)에 의해 주어진다.

표현 개념은 가환 승수영역 \(P\)와 가환적으로 결합된 환들로도 확장할 수 있다(§ 8). 이 경우에는 중심에 \(P\)를 포함하는 표현환 \(\mathsf{K}\)만 취하고, 표현이 환준동형 \(\mo\sim\mathfrak D\)일 뿐 아니라 연산자준동형일 것도 요구한다. 즉 \(a\mapsto A\)이면 모든 \(\varrho\in P\)에 대해 \(a\varrho\mapsto A\varrho\)이어야 한다.
표현가군에 관해서 이 요구는 계산법칙
\[
        am\cdot\varrho\ [=a\cdot m\varrho]=a\varrho\cdot m
\]
을 추가한다. 이때 가군과 환은 흔히 \glqq{}\(\mathsf{P}\)와 가환적으로 결합되어 있다\grqq{}고 한다.

\subsection*{§ 16. 가약표현.}

\emph{\(\mA\)가 표현가군 \(\mM\)의 부분가군이고, \(\mA\)의 기저 \(z_1,\ldots,z_t\)에 \(y_1,\ldots,y_r\)을 덧붙여 \(\mM\)의 기저를 고를 수 있어서}
\[
        \mM=y_1\mathsf{K}+\ldots+y_r\mathsf{K}+z_1\mathsf{K}+\ldots+z_t\mathsf{K}\textsuperscript{16)}
\]
\emph{이면, 표현들은}
\[
        C=\begin{pmatrix}R&0\\ S&T\end{pmatrix}
\]
\emph{의 꼴이다. 여기서 행렬 \(T\)들만으로 \(\mA\)가 생성하는 \(t\)차 표현을 이루고, 행렬 \(R\)들은 \(\mM/\mA\)가 생성하는 \(r\)차 표현을 이룬다.}

\emph{증명.} 
\[
        (cy_1,\ldots,cy_r,cz_1,\ldots,cz_t)
        =(y_1,\ldots,y_r,z_1,\ldots,z_t)\cdot C.
\]
\(cz_i\)들은 \(\mA\)의 원소이므로 \(z_i\)들만으로 나타난다. 따라서 \(C\)의 오른쪽 위 블록은 0이다. \(C\)의 나머지 블록들을 위 배열과 같이 \(R,S,T\)라 하면, 특히
\[
        (cz_1,\ldots,cz_t)=(z_1,\ldots,z_t)\cdot T
\]
이다. 그러므로 \(T\)들은 \(\mA\)가 매개하는 표현을 이룬다. 또한
\[
        (cy_1,\ldots,cy_r)\equiv
        (y_1,\ldots,y_r)\cdot R\pmod{\mA},
\]
이고 \(y_i\)들은 \(\bmod\,\mA\) 선형독립기저를 이룬다. 따라서 \(R\)들은 \(\mM/\mA\)가 생성하는 표현을 이룬다.

역으로
\[
        C=\begin{pmatrix}R&0\\ S&T\end{pmatrix}
\]
꼴의 \glqq{}가약\grqq{} 표현이 주어지고 \(R,T\)가 정방행렬이면, 이에 속하는 표현가군에서 각 \(c\)와 마지막 \(t\)개 기저원소 \(z_1,\ldots,z_t\)의 곱은 이 \(z_i\)들만으로 나타난다. 곧 \(\mA=(z_1,\ldots,z_t)\)는 부분가군이다.

{\footnotesize\noindent 16) 달리 말하면 \(\mA\)가 \(\mathsf{K}\)-가군으로서 직접가수라는 뜻이다. \(\mathsf{K}\)가 체이면 이 가정은 언제나 성립한다.\par}

\emph{따름정리.}
\[
        \mM\supset\mA_1\supset\mA_2\supset\ldots\supset\mA_e=0
\]
을 \(\mM\)의 한 합성열이라 하고, 각 \(\mA_i\)가 앞 가군의 \(\mathsf{K}\)-가군으로서 직접가수라 하자. 그러면
\[
        z_{11},\ldots,z_{1r_1},\ z_{21},\ldots,z_{2r_2},\ldots,
        z_{e1},\ldots,z_{er_e}
\]
라는 기저를 골라 \(i>\nu\)인 \(z_{ik}\)들이 \(\mA_\nu\)의 기저를 이루게 할 수 있다. 이때 \(\mM\)이 매개하는 표현은
\srcnumdisplay[3.2em]{(2)}{%
 C=\begin{pmatrix}
 R_{11}&0&\ldots&0\\
 R_{12}&R_{22}&\ldots&0\\
 \vdots&\vdots&\ddots&\vdots\\
 R_{1e}&R_{2e}&\ldots&R_{ee}
 \end{pmatrix}}
의 꼴이고, \(R_{\nu\nu}\)들은 합성인자 \(\mA_{\nu-1}/\mA_\nu\)가 매개하는 표현을 이룬다.

각 표현 \(R_{\nu\nu}\)은 기약이다. 합성인자 \(\mA_{\nu-1}/\mA_\nu\)가 단순하기 때문이다. \(\mathsf{K}\)가 체라서 모든 부분가군이 직접가수라고 가정하면, 역으로 (2)의 꼴로 가능한 데까지 가약화된 모든 표현은 합성열을 낳는다. 조르당--횔더 정리에 따라 합성인자들은 연산자동형과 순서를 제외하면 유일하므로, \emph{\(R_{\nu\nu}\)들은 동치표현과 순서를 제외하면 유일하게 결정된다.}

\(\mM=\mA+\mB\)가 두 표현가군 \(\mA=(y_1,\ldots,y_r)\), \(\mB=(z_1,\ldots,z_s)\)의 직접합이면, 기저 \((y_1,\ldots,y_r,z_1,\ldots,z_s)\)가 매개하는 표현은 명백히
\[
        C=\begin{pmatrix}R&0\\0&S\end{pmatrix}
\]
의 꼴이다. 여기서 \(R\)은 \(\mA\)가, \(S\)는 \(\mB\)가 매개하는 원소 \(c\)의 표현이며, 역도 성립한다. 먼저 가군 \(\mM\)을 직접분해불능 성분들의 직접합으로 나타내고, 각 성분에서 위와 같은 합성열을 찾으면 알맞은 기저를 골랐을 때 표현은
\[
\begin{array}{@{}r@{\quad}c@{}}
(3)&
C=\left(
\begin{array}{ccc}
\boxed{\begin{matrix}
R_{11}&0\\
\vdots&\ddots\\
R_{ik}&R_{rr}
\end{matrix}} & & 0\\[1.4em]
& \boxed{\begin{matrix}
S_{11}&0\\
\vdots&\ddots\\
S_{ik}&S_{ss}
\end{matrix}} & \\ & & \ddots \\
0 & & \boxed{\begin{matrix}
T_{11}&0\\
\vdots&\ddots\\
T_{ik}&T_{tt}
\end{matrix}}
\end{array}
\right)
\end{array}
\]
의 꼴이다.

다시 \(\mathsf{K}\)가 체이면, 역으로 (3)의 꼴로 가능한 데까지 가약화된 모든 표현은 가군을 직접분해불능 가수들로 나타내고 그 안에서 합성인자들을 택하는 분해를 낳는다. 크룰과 오토 슈미트의 정리(각주 11)에 따라 직접분해불능 표현성분들의 류는 순서를 제외하면 유일하게 결정된다.

가군 \(\mM\)이 완전가약이면 모든 상자는 하나의 기약표현만으로 이루어지며, 표현 (3)을 완전가약이라 한다.

\subsection*{§ 17. 단위원을 갖는 환에서 표현가군의 직접합분해.}

\emph{\(\mM\)을 \(\mo\)-가군, \(\mo\)를 단위원을 갖는 환이라 하자. 그러면}
\[
        \mM=\mM_e+\mM_0,
\]
\emph{여기서 단위원은 \(\mM_e\)에서는 항등연산자이고 \(\mM_0\)에서는 영연산자이다. (\(\mM\)이 우측 \(K\)-가군이면 \(\mM_e,\mM_0\)도 그러하다.)}

\emph{증명.} \(\mM\)의 모든 \(m\)은
\[
        m=em+(m-em)
\]
으로 나타난다. \(em\) 전체를 \(\mM_e\), \(m-em\) 전체를 \(\mM_0\)라 하자. 그러면 \(\mM_e,\mM_0\)가 가법군이며, \(\mM\)이 우측 \(K\)-가군일 때 이들도 우측 \(K\)-가군이라는 것은 자명하다. 또한
\[
        r em=erm
\]
이므로 \(\mM_e\)도 \(\mo\)-가군이다. 마찬가지로
\[
        r(m-em)=rm-rem=rm-rm
\]
이므로 \(\mM_0\)도 \(\mo\)-가군이다. \(em\)들에는 \(e\)가 항등연산자이고 \((m-em)\)들에는 영연산자이다. 따라서 \(\mM_e\)와 \(\mM_0\) 양쪽에 속하는 것은 0뿐이며 합 \(\mM_e+\mM_0\)는 직접합이다.

표현가군의 경우 \(\mM_0\)는 모든 원소에 0을 대응시키는 자명표현을 생성한다. 가수 \(\mM_0\)를 떼어 내면 \(\mM_e\)가 매개하는 표현이 남고, \emph{이 표현에서는 단위원에 단위행렬이 대응한다.} 그러므로 이제부터는 그러한 표현에만 국한할 수 있고 또 그렇게 하겠다.

\emph{\(\mo=\ma_1+\ldots+\ma_s\)가 양쪽아이디얼들의 직접합이고, \(e=e_1+\ldots+e_s\)가 \(e\)의 분해이며, \(\mM\)이 \(e\)를 항등연산자로 갖는 \(\mo\)-가군이면}
\[
        \mM=\ma_1\mM+\ldots+\ma_s\mM
\]
\emph{은 직접합이다.} 실제로 \(\mM=\mo\mM=(\ma_1\mM,\ldots,\ma_s\mM)\)임은 명백하다. \(\mb_i=\ma_1+\cdots+\widehat{\ma_i}+\cdots+\ma_s\)로 놓으면
\[
        \mM=(\ma_i\mM,\mb_i\mM)
\]
이고, \(e_i\)가 \(\ma_i\mM\)에서는 항등연산자, \(\mb_i\mM\)에서는 영연산자이므로 이 합은 직접합이다.

이 정리에 따라 우리는 대개 양쪽으로 분해불능인 환의 표현에 국한한다.

\subsection*{§ 18. 완전가약환의 가군론과 표현론.}

\emph{\(\mo\)가 완전가약이고 양쪽단순이며, \(\mM\)이 \(\mo\)의 단위원을 항등연산자로 갖는 유한 \(\mo\)-가군이라고 하자. 그러면 \(\mM\)은 완전가약이고, 그 단순성분들은 단순 좌측아이디얼 \(\ml_i\)들과 연산자동형이다.}

\emph{증명.} 
\[
        \mo=\ml_1+\ldots+\ml_n,
        \qquad
        \mM=(\mo m_1,\ldots,\mo m_k)
\]
라 하자. 그러면
\[
\makebox[0pt][r]{(4)\qquad}\mM=(\ldots,\ml_i m_k,\ldots).
\]
0이 아닌 \(\ml_i m_k\)들은 \(\ml_i\)와 동형이다. 대응 \(a\mapsto am_k\)가 연산자동형이기 때문이다. 따라서 이 \(\ml_i m_k\)들은 단순하다. 표시 (4)에서 앞선 항들의 합에 이미 포함되는 \(\ml_i m_k\)들을 버리면 합은 직접합이 된다.

\emph{따름정리.} 또한 \(\mM\)이 직접분해불능이면 \(\mM\)은 단순이고 어떤 \(\ml_i\)와 동형이다.

이 단순 \(\mM\)의 자기동형체를 \(\Delta\), \(\ml_i\)의 자기동형체를 \(\Lambda\)라 하자. \(\mM\)과 \(\ml_i\) 사이의 연산자동형에 의해 환동형 \(\Delta\simeq\Lambda\)가 성립한다. § 1에 따라 \(\mM\)과 \(\ml_i\)를 각각 \(\Delta\), \(\Lambda\)를 우측영역으로 하는 쌍가군으로 볼 수 있으며, 이 쌍가군들은 표현가군이기도 하다. 실제로 \(\ml_i\), 따라서 \(\mM\)은 자기동형체에 관하여 유한랭크이고(§ 14), 그 체의 단위원은 항등연산자이다. 이들이 매개하는 표현들은 환동형 \(\Delta\simeq\Lambda\)에 의해 서로 옮겨진다. \emph{따라서 이 표현류를 연구할 때에는 \(\ml_1\)가 그 자기동형체에서 매개하는 표현에 국한할 수 있다.}

특히 \(\ml_1\)에 대해 행렬단위들로 된 기저
\[
        \ml_1=(c_{11},\ldots,c_{n1})
\]
를 택하고, 자기동형체 \(\Lambda\)의 실현으로 앞에서 \(\mathsf{K}\)라 부른 \(\mo\)의 부분체(§ 14)를 택하자. 이때 앞의 논의에서 우측아이디얼을 좌측아이디얼로 바꾸고, 그에 따라 자기동형들을 오른쪽에 써야 한다. 
\[
        a=\sum c_{ik}\alpha_{ik}
\]
로 나타내면
\[
        a\longmapsto(\alpha_{ik})
\]
는 \(\ml_1\)이 매개하는 \(\mathsf{K}\) 안에서의 표현이다. 실제로
\[
\begin{aligned}
        a\cdot c_{k1}
        &=\left(\sum\sum c_{ij}\alpha_{ij}\right)\cdot c_{k1}\\
        &=\sum c_{ik}\alpha_{ik}c_{k1}
         =\sum c_{i1}\alpha_{ik},
\end{aligned}
\]
또는
\[
        a(c_{11},\ldots,c_{n1})
        =(c_{11},\ldots,c_{n1})(\alpha_{ik})
\]
이다.

\emph{\(\Gamma\)를 \(K\)의 부분체라 하고, \(K\)가 \(\Gamma\)에 관하여 유한랭크 \(t\)를 갖는다고 하자.}
\[
        K=z_1\Gamma+\cdots+z_t\Gamma.
\]
\emph{그러면 \(K\)는 \(\Gamma\)에 관한 자기 자신의 표현가군이 된다. \(K\)의 원소 \(\beta\)는}
\[
        \beta z_j=\sum_i z_i\beta_{ij},
        \qquad B=(\beta_{ij})
\]
\emph{로 얻는 \(\Gamma\)의 행렬 \(B\)에 의해 표현된다.}

이제 \(\ml_1\)을 부분체 \(\Gamma\)에 관한 표현가군으로 보면 다음을 얻는다.

\emph{\(\ml_1\)이 매개하는 \(\mo\)의 \(\Gamma\) 안에서의 표현은, 위에서 준 \(\mo\)의 \(\mathsf{K}\) 안에서의 표현}
\[
        a\to(\alpha_{ik})
\]
\emph{에서 원소 \(\alpha_{ik}\)를 \(\mathsf{K}\)의 \(\Gamma\) 안에서의 표현에서 이에 대응하는 행렬 \(A_{ik}\)로 바꾸어 얻는다.}

\emph{증명.}
\[
        \ml_1=\sum c_{i1}\mathsf{K}=\sum\sum c_{i1}\varkappa_j\Gamma.
\]
기저 \((\ldots,c_{i1}\varkappa_j,\ldots)\)에 관하여
\[
c_{ik}\longmapsto
\begin{array}{c@{\quad}c}
 & \begin{array}{ccc}
 \boxed{\scriptstyle 1}&\boxed{\scriptstyle k}&\boxed{\scriptstyle n}
 \end{array}\\[-.2em]
\begin{array}{c}
\boxed{\scriptstyle 1}\\ \boxed{\scriptstyle i}\\ \boxed{\scriptstyle n}
\end{array}
&
\begin{pmatrix}
0&\ldots&0&\ldots&0\\
\vdots&&\vdots&&\vdots\\
0&\ldots&E_t&\ldots&0\\
\vdots&&\vdots&&\vdots\\
0&\ldots&0&\ldots&0
\end{pmatrix}
\end{array},
\]
이다. 여기서 \(E_t\)는 단위행렬이고 0은 \(t\)차 영행렬이다. 또한
\[
        \alpha\longmapsto
        \begin{pmatrix}
        A& &0\\
         &\ddots&\\
        0& &A
        \end{pmatrix}.
\]
따라서
\[
        a=\sum c_{ik}\alpha_{ik}\to
        \begin{pmatrix}
        A_{11}&\ldots&A_{1n}\\
        \vdots&&\vdots\\
        A_{n1}&\ldots&A_{nn}
        \end{pmatrix}.
\]

임의의 \(\mM\)으로 넘어가는 것은 다시 표현환에 환동형을 가하는 것을 뜻한다. \(\mo\)가 양쪽단순환들의 직접합이면 이 환들은 차례로 동형으로 표현되고, 나머지 환들은 0으로 표현된다.

\textbf{논평.} 여기서 연구한 좌측아이디얼의 자기동형체와 그 유한부분체 안에서의 표현들은, 단순 \(\mo\)-가군(또는 좌측아이디얼)이 매개할 수 있는 표현들 가운데 동형을 제외하면 유일한 것들이다. 실제로 앞에서 한 논평(§ 2)에 따라, \(\mo\)-가군 \(\mM\)을 임의의 체 \(\mathsf{K}\)에 관한 쌍가군으로 보면 \(\mathsf{K}\)의 원소들은 가군준동형들을 생성한다. 따라서 체 \(\mathsf{K}\)는 이 \(\mo\)-가군의 자기동형체의 한 부분체에 준동형으로 대응해야 한다. 단위원이 항등연산자이므로 이 준동형은 동형이다. 또한 전체 자기동형체는 해당 부분체에 관하여 유한차수이어야 한다. 그렇지 않으면 표현가군도 무한랭크가 되어 불가능하기 때문이다.

끝으로 여기서 조사한 기약표현들이 전부는 아니고, 단순 \(\mo\)-가군들이 매개하는 것들뿐이라는 점을 덧붙인다. 일반적으로는 표현가군들, 예컨대 부분체 \(\Sigma\)를 좌측영역으로 하고 자기 자신을 우측영역으로 하는 쌍가군으로 본 체 \(\Omega\)가 더 있다. 이들은 \(\mo\)-가군으로서는 가약이지만 쌍가군으로서는 기약이고, 그러므로 기약표현을 낳는다.

\subsection*{§ 19. 가군과 표현가군의 단순 합성인자.}

\(\mo\)를 좌측아이디얼에 대한 극대조건과 극소조건을 만족하는 환, \(\mc\)를 극대 멱영아이디얼이라 하자. 그러면 다음이 성립한다.

\emph{모든 단순 \(\mo\)-가군 \(\mM\)은 \(\mo\)에 의해 소멸되거나, 근기 없는 환 \(\mo/\mc\)의 한 단순 좌측아이디얼 \(\ml\)과 동형이다. 뒤의 경우 그 가군은 \(\mo/\mc\)의 양쪽아이디얼 가운데 \(\ml\)을 포함하지 않는 모든 것에 의해 소멸되고, 절대승수영역(§ 1)은 \(\ml\)을 포함하는 양쪽단순환과 동형이다.}

\textbf{증명.} \(\mc^{\varrho}=0\)이라 하자. \(\mc\mM=0\)이어야 한다. 그렇지 않으면 \(\mc\mM=\mM\)이고, 따라서 \(\mM=\mc\mM=\mc^2\mM=\ldots=0\)이 되어 버리기 때문이다. 그러므로 \(\mM\)을 \(\mo/\mc\)-가군으로도 볼 수 있다. \(\mc\)에 따른 같은 잉여류의 모든 원소는 \(\mM\)의 원소를 곱할 때 같은 결과를 준다. 근기 없는 환인 \(\mo/\mc\)는 단위원을 가지므로, \(\mM\)은 \(\mo/\mc\)에 의해 소멸되는 가군과 그 단위원이 항등연산자인 가군의 직접합이다(§ 17). \(\mM\)이 단순이므로 이 두 가수 가운데 하나만 나타날 수 있다. 단위원이 항등연산자이면 § 18에서와 같이 단순 좌측아이디얼과의 동형이 따른다. 실제로 \(m\ne0\)이 \(\mM\)의 원소이면 \((\mo/\mc)m\ne0\)이고, 따라서 \(\mo/\mc\)의 적어도 한 단순 좌측아이디얼 \(\ml\)에 대해 \(\ml m\ne0\)이다. 그러면 \(\ml m\)은 \(\mM\) 전체이어야 한다. 나머지 주장들은 그러한 좌측아이디얼에 대해서는 명백하며, 그와 동형인 모든 가군으로 곧바로 옮겨진다.

\textbf{따름정리.} \emph{\(\mM\)이 합성열을 갖는 \(\mo\)-가군이면, 그 단순 합성인자들은 \(\mo\)에 의해 소멸되거나 \(\mo/\mc\)의 단순 좌측아이디얼들과 동형이다.}

\(\mathsf{P}\)가 \(\mo\)와 가환적으로 결합된 환이면 모든 결과가 그대로 성립한다. 다만 가군 대신 우측의 \(\mathsf{P}\)와 좌측의 \(\mo\)에 관한 쌍가군으로서
\[
        am\cdot\rho=a\cdot m\rho=a\rho\cdot m
        \qquad\hbox{(§ 15)}
\]
을 만족하는 것들을 고찰해야 한다. 또한 \(\mathsf{P}\)의 단위원도 항등연산자이어야 한다. 부분가군으로는 \(\mathsf{P}\)와의 곱셈을 허용하는 것들만을 언제나 고찰한다(아이디얼의 경우와 같다). 우리의 증명에서 구성한 부분가군 \(\mc\mM,\mc^2\mM,\ldots\)는 이 조건을 만족한다.

특히 \(\mathsf{P}\)가 체이고 \(\mM\)이 유한랭크 \(\mathsf{P}\)-가군이면 \(\mM\)은 동시에 표현가군이다. \(\mM\)이 매개하는 표현은 환준동형일 뿐 아니라 \(\mathsf{P}\)에 관한 연산자준동형이다(§ 15). 이 성질을 갖는 \(\mathsf{P}\) 안에서의 모든 표현은 그러한 가군들로 얻어진다. 기약표현에는 단순가군이 대응하며, 역도 성립한다. 따라서 이 절의 정리들은 곧바로 \(\mo\)의 \(\mathsf{P}\) 안에서의 표현에 관한 정리로 해석할 수 있다. \emph{모든 기약표현은 영표현을 제외하면 \(\mo/\mc\)의 단순 좌측아이디얼들이 매개한다. 따라서 임의의 표현을 § 16에 따라 합성열로 가약화하면, 주대각선에는 0 이외에 \(\mo/\mc\)의 단순 좌측아이디얼들이 매개하는 표현들과 동치인 것들만 나타난다.}

\textbf{논평.} \(\mo\)와 \(\mathsf{P}\)에 관하여 위의 가정들을 두면 \(\mo\) 자체에는 유한성 조건을 가정할 필요가 없다. 모든 표현은 동시에 절대승수영역의 동형표현이기 때문이다. 이 절대승수영역은 유한차 행렬환의 동형원상으로서 스스로 유한랭크 \(\mathsf{P}\)-가군이 된다. 또 정의상 \(\mathsf{P}\)-가군인 것들만 허용아이디얼로 고찰하므로, 이 절대승수영역에서는 극대조건과 극소조건이 성립한다.

\(\mathsf{P}\)를 가환체라고 가정하면 이 절대승수영역은 \(\mathsf{P}\)에 관한 초복소계가 된다. 주대각선에 0 이외의 것이 하나라도 나타나면 이 가정은 언제나 충족된다. 실제로 § 18의 논평에 따라 \(\mathsf{P}\)는 \(\ml\)의 자기동형체의 한 부분체와 동형이고, 이를 \(\mo/\mc\)의 부분체 \(K\)로 실현할 때(§ 14) 가환적 결합 때문에 \(P\)는 \(K\)의 중심에 놓여야 한다.

\section*{제IV장. 군과 초복소계의 표현.}

\subsection*{§ 20. 초복소계의 분류.}

가환체 \(P\)에 관한 초복소계 \(\mo\)를 고찰한다. 약속에 따라 \(\mo\)의 아이디얼로는 \(\mathsf{P}\)-가군만을 고찰하므로 좌측ㆍ우측아이디얼에 대한 극대조건과 극소조건이 성립한다(§ 8 참조). 따라서 제II장에서 전개한 환론 전체가 적용된다.

이하에서 초복소계 \(\mo\)의 \emph{표현}이라 하면 언제나 체 \(\mathsf{P}\) 안에서의 표현으로서, \(a\mapsto A\)이면 모든 \(\varrho\in P\)에 대해 \(a\varrho\mapsto A\varrho\)인 것, 곧 \(\mathsf{P}\)에 관한 가군준동형을 뜻한다. 표현가군에 대해서는 이것이 \(a\varrho\cdot m=a\cdot m\varrho\)를 뜻한다(§ 15 끝 참조). 따라서 표현가군에는 § 19의 정리들이 적용되며, 그 결과 \emph{모든 기약표현은 \(\bar{\mo}=\mo/\mc\)의 단순 좌측아이디얼들이 매개하고, 따라서 서로 동치가 아닌 기약표현의 수는 분해 \(\bar{\mo}=\bar{\mo}e^{(1)}+\ldots+\bar{\mo}e^{(s)}\)의 양쪽단순 가수 \(\ma\)의 수와 같다.}

그러한 좌측아이디얼 \(\ml_\nu\)가 정의하는 \(\mo\)의 표현을 실제로 구하려면, 먼저 \(\mo\)의 원소들에서 \(\bar{\mo}\)의 대응 잉여류들로 넘어간다. \(\mathsf{P}\)의 원소와의 곱셈은 \(\bar{\mo}\)의 모든 아이디얼에서 동형을 이루므로 체 \(\mathsf{P}\)는 각 단순 좌측아이디얼의 자기동형체 \(\mathsf{K}_\nu\)의 부분체 \(e^{(\nu)}\mathsf{P}\)에 준동형이고, 단위원이 항등연산자이므로 사실 동형이다. 단순 좌측아이디얼 \(\ml_\nu\)가 매개하는 \(P\) 안에서의 표현은 먼저 이 부분체 \(e^{(\nu)}\mathsf{P}\) 안에서 \(\ml_\nu\)가 매개하는 표현을 구한 뒤, 동형 \(e^{(\nu)}\mathsf{P}\simeq \mathsf{P}\)로 \(\mathsf{P}\)에 옮겨 얻는다. 자기동형체 \(\mathsf{K}_\nu\)는 \(\mathsf{P}\)에 관하여 유한랭크이다. 따라서 이는 \(\ml_\nu\)의 자기동형체의 유한차수 부분체 안에서의 표현을 구하는 문제이다.

그러므로 구하는 표현을 얻기 위해 먼저 \(\bar{\mo}\)의 각 \(c\)를 양쪽성분들로
\[
        c=c_1+\ldots+c_s
\]
나타내자. 다른 \(c_i\)들은 아이디얼 \(\ml_\nu\)를 소멸시켜 영행렬로 표현되므로 \(c_\nu\)를 나타내는 행렬만 찾으면 된다. § 18에 따라 다음과 같이 구한다. \(c_\nu\)를 \(\ma_\nu\)의 행렬단위 \(c_{ik}^{(\nu)}\)로
\[
        c_\nu=\sum c_{ik}^{(\nu)}\alpha_{ik}^{(\nu)}
\]
나타내고, 행렬 \((\alpha_{ik}^{(\nu)})\)의 각 원소 \(\alpha_{ik}^{(\nu)}\)를 \(\mathsf{K}_\nu\)가 매개하는 \(\mathsf{K}_\nu\)의 \(\mathsf{P}\) 안에서의 표현에서 이에 대응하는 행렬 \(A_{ik}^{(\nu)}\)로 바꾼다.

\textbf{논평.} 체 \(\mathsf{K}_\nu\)는 \(\mathsf{P}\)에 관하여 유한랭크이다. \(x\)의 거듭제곱들 사이에는 반드시 선형종속관계가 있으므로 각 원소 \(x\)는 \(e^{(\nu)}\mathsf{P}\)의 계수를 갖는 방정식 \(f(\varkappa)=0\)을 만족한다. 특히 \(\mathsf{P}\)가 대수적으로 닫혀 있으면 이 방정식은 일차인자들로 분해된다. \(\mathsf{K}\)가 체이므로 \(x\)는 이미 한 일차인자의 영점이고, 따라서 \(x\in e^{(\nu)}P\)이다. 곧 \(\mathsf{K}_\nu=e^{(\nu)}\mathsf{P}\)이다. 이 경우에는 행렬 \((\alpha_{ik}^{(\nu)})\) 자체가 기약표현을 이룬다.

이 논평과 § 19 끝부분에서 \glqq{}번사이드 정리\grqq{}가 나온다.

\emph{\(\mathsf{P}\)가 대수적으로 닫혀 있고 환 \(\mo\)와 가환적으로 결합되어 있다고 하자. 그러면 \(P\) 안에서의 모든 \(n\)차 기약표현 \(\mo\sim\mD\)는 정확히 \(n^2\)개의 선형독립 행렬을 포함한다.}\textsuperscript{17)}

이 정식화에서 통상적인 정식화가 곧바로 따른다. \glqq{}\(\mathsf{P}\)의 원소로 된 \(n\)행 정방행렬들의 한 계가 임의의 두 원소와 함께 그 곱도 포함한다고 하자. 이 계는, 그 계의 모든 행렬의 원소 \(\alpha_{ik}\)들이 만족하는 \(\mathsf{P}\) 계수의 선형동차관계
\[
        \sum \alpha_{ik}\varrho_{ik}=0
\]
가 존재하지 않을 때 그리고 그때에만 기약이다.\grqq{} 실제로 그러한 행렬계에 모든 선형결합을 덧붙이면 환이자 \(\mathsf{P}\)-가군을 얻으며, 이 환을 자기 자신의 표현으로 볼 수 있다.

\textbf{증명.} § 19에 따라 \(\mD\)는 단위원을 갖는 양쪽단순 완전가약환 \(\bar{\mo}\)와 동형이다. 그런데 그러한 환은 좌측아이디얼의 자기동형체에 관한 완전행렬환, 이를테면 모든 \(m\)행 정방행렬들의 환이고, 대수적 폐성 때문에 이 자기동형체는 \(\mathsf{P}\)와 같다. \(\bar{\mo}\)의 모든 기약표현, 따라서 주어진 표현도 단순 좌측아이디얼이 생성한다. 그 좌측아이디얼의 랭크는 \(m\)이고 표현의 차수는 \(n\)이므로 \(m=n\)이다. 따라서 \(\bar{\mo}\), 그리고 \(\mD\)에는 정확히 \(n^2\)개의 선형독립 원소가 있다.

마찬가지로 \glqq{}일반화된 번사이드 정리\grqq{}\textsuperscript{18)}도 쉽게 따른다.

\emph{같은 가정 아래 \(\mD\)가 \(\mo\)의 완전가약표현이고, 차수가 \(n_1,n_2,\ldots,n_s\)인 서로 동치가 아닌 성분들로만 분해된다고 하자. 그러면 \(\mD\)는 \(\mathsf{P}\)에 관하여 랭크 \(n_1^2+\ldots+n_s^2\)를 갖는다.}

\textbf{증명.} 다시 \(\bar{\mo}\)를 절대승수영역이라 하자. 이는 반드시 초복소계이다. 서로 동치가 아닌 표현들은 \(\bar{\mo}/\bar{\mc}\)의 서로 다른 양쪽단순 성분 \(\ma_i\)에 속하는 좌측아이디얼 \(\ml_1,\ldots,\ml_s\)이 생성한다. 여기서 \(\bar{\mc}\)는 근기이다. 따라서 주어진 표현은 \(\ml_1+\ldots+\ml_s\)가 생성하며, \(\ma_1+\ldots+\ma_s\)를 절대승수영역으로 갖는다. 위와 같은 랭크 고찰로부터 주장이 따른다.

이제 \(\mo\)를 임의의 체에 관한 근기 없는 초복소계라 하자. § 13에 따라 \(\mo\)는 완전가약이고 단위원을 갖는다. § 17과 § 18의 정리들에서 다음이 따른다. \emph{\(\mo\)의 모든 표현은 완전가약이다.}

역으로 다음이 성립한다. \emph{\(\mathsf{P}\)와 가환적으로 결합된 환 \(\mo\)의 모든 완전가약표현은 근기 없는 초복소계를 절대승수영역으로 갖는다.}

\textbf{증명.} 절대승수영역 \(\bar{\mo}\)는 \(\mathsf{P}\)의 행렬들로 된 환이자 \(\mathsf{P}\)-가군과 동형이므로 초복소계이다. 근기 \(\bar{\mc}\)의 원소들은 § 19에 따라 모든 기약표현에서, 따라서 전체 표현에서 0으로 표현된다. 그러므로 \(\bar{\mc}=0\)이다.

(표현에 서로 동치인 성분들만 나타나고 이들이 모두 한 좌측아이디얼 \(\ml\)에 의해 매개되면 절대승수영역은 양쪽단순이다.)

초복소계 \(\mo\)의 \emph{정칙표현}이란 단위아이디얼 \(\mo\) 자체가 매개하는 표현을 말한다. (군환에서는 특히 군원소 \(u_1,\ldots,u_h\)로 된 기저를 택한다. § 6 참조.) \(\mo\)의 합성인자들 가운데 \(\mo/\mc\)의 모든 좌측아이디얼이 나타나므로, 정칙표현을 § 16의 식 (2)에 따라 합성열을 써서 알맞은 기저로 변환하면 모든 기약표현이 이미 그 대각행렬 \(R_{ii}\)들 안에 들어 있다.

기저원소 \(u_1,\ldots,u_h\)를 나타내는 행렬들을 알면 표현은 알려진다. 특히 군환에서 \(u_1,\ldots,u_h\)가 한 군의 원소들이면, 그 군의 모든 행렬 준동형표현은 군환의 표현도 생성한다. 따라서 군의 표현문제는 초복소계 표현문제의 특수한 경우이다.

{\footnotesize\noindent
17) 따라서 \(\mo\)를 초복소계라고 가정하지 않는다. 예컨대 무한군의 원소들을 매번 유한개씩 취한 모든 선형결합으로 이루어진 군환일 수 있다.\par
18) G. Frobenius 및 I. Schur, \emph{Über die Äquivalenz der Gruppen linearer Substitutionen}, Sitzungsber. Berlin 1906.\par}

\subsection*{§ 21. 바탕체의 확대. 중심의 표현.}

\(\mo=a_1\mathsf{P}+\ldots+a_h\mathsf{P}\)가 초복소계이고 \(\Omega\)가 바탕체 \(\mathsf{P}\)의 확대체이면, 기저원소 \(a_i\)에 대한 예전 곱셈법칙을 그대로 써서
\[
        \mo\Omega=a_1\Omega+\ldots+a_h\Omega
\]
를 만들 수 있다.\textsuperscript{18a)} \(\mo\Omega\)는 다시 \(\Omega\)에 관한 초복소계이다.

\(\mo\)의 모든 표현은 \(\mo\Omega\)의 표현을 낳는다. 기저원소 \(a_i\)들의 표현을 알면 모든 원소의 표현을 알기 때문이다.\textsuperscript{19)}

{\footnotesize\noindent
18a) 더 정확히 말하면, 예전 곱셈법칙을 갖는 새 기호 \(\bar a_i\)를 도입한다. 그러면 \(\bar a_1\Omega+\ldots+\bar a_h\Omega\)는 \(\mo\)와 동형인 부분환을 포함하므로, 나중에 \(\bar a\)를 \(a\)와 동일시할 수 있다.\par
19) 물론 다시 \(\mathsf{P}\)-가군준동형인 표현만을 고찰한다. 곧 \(a\mapsto A\)이면 모든 \(\varrho\in P\)에 대해 \(a\varrho\mapsto A\varrho\)이어야 하며, \(\Omega\)에 대해서도 이에 대응하는 조건을 둔다.\par}

아이디얼이나 표현가군, 그리고 이에 대응하는 표현이 대수적으로 닫힌 체로 넘어간 뒤에도 기약으로 남으면 \emph{절대기약}이라 한다.

\(\mo\Omega\)가 근기 없으면 \(\mo\)도 근기 없다. \(\mo\)의 멱영아이디얼 \(\mc\)는 \(\mo\Omega\)의 확장아이디얼 \(\mc\Omega\)를 낳기 때문이다. 그 역은 아래에서 보듯 성립하지 않는다.

\(\mo\Omega\)의 중심은 확대된 중심 \(\mZ\Omega\)와 같다. \(\mo\)가 완전가약이면 \(\mZ\)는 체들의 직접합
\[
        \mZ=\mZ_1+\cdots+\mZ_s
\]
이다. \(\mo\Omega\)로 넘어갈 때 \(\mo\)가 완전가약으로 남으면, 곧 새 멱영아이디얼이 생기지 않으면, 새 중심
\[
        \mZ\Omega=\mZ_1\Omega+\cdots+\mZ_s\Omega
\]
은 \(\Omega\) 안에서 다시 체들의 직접합으로 분해된다. \(\Omega\)가 대수적으로 닫혀 있을 때 이 체들은 § 20의 논평에 따라 \(\Omega\)에 관한 랭크가 1이고 \(\Omega\)와 동형이다.

중심 \(\mZ\)의 표현에는 다음 정리들이 성립한다.

\emph{가환 초복소계 \(\mZ\)의 대수적으로 닫힌 체 \(\Omega\) 안에서의 모든 기약표현은 1차이다. 달리 말하면 기약표현들은 \(\mZ\)에서 \(\Omega\)로 가는 준동형들과 같다.}

\textbf{증명.} \(\mZ\)의 모든 기약표현은 \(\mZ_\Omega\)의 기약표현을 낳고, 따라서 \(\mC\)가 \(\mZ_\Omega\)의 근기일 때 \(\mZ_\Omega/\mC\)의 기약표현도 낳는다. \(\mZ_\Omega/\mC\)는 근기 없는 가환계이므로 § 14 끝부분에 따라 체들의 직접합이다. \(\Omega\)가 대수적으로 닫혀 있으므로 이 체들은 1차이고, 모든 기약표현을 생성한다(§ 20).

동시에 1차 동치표현은 반드시 같으므로(\(\lambda^{-1}\alpha\lambda=\alpha\)) 다음이 따른다. \emph{\(\mZ\)에서 \(\Omega\)로 가는 서로 다른 준동형의 수는 \(\mZ_\Omega/\mC\)의 랭크와 같다.}

마지막 논평은 \(\mZ\)가 \(\mathsf{P}\) 위의 체인 특수한 경우에 다음 정리를 준다.

\emph{\(\mZ\)가 가환체이면, \(\mZ\)가 \(\mathsf{P}\)의 제1종 확대일 때 그리고 그때에만 \(\mZ\)는 완전가약, 곧 근기 없다.}

실제로 체에 대한 준동형은 동형이다. 그 수는 \(\mZ_\Omega/\mC\)의 랭크와 같으므로, \(\mC=0\)일 때 그리고 그때에만 체차수, 곧 \(\mZ\)가 제1종 확대일 때의 수와 같다.

\(\mZ\)가 완전가약이고
\[
        \mZ=\mZ_1+\cdots+\mZ_t
\]
이면 모든 표현에서 각 체 \(\mZ_\nu\)도 준동형으로 사상된다. 기약표현에서는 이 체들 가운데 하나가 동형으로 사상되고 나머지는 0으로 표현된다(§ 19).

이 정리들을 초복소계에 적용하면 우선 근기 없는 계에 대해 다음을 얻는다.

\emph{\(\mo_\Omega\), 따라서 \(\mo\)가 근기 없는 계이면, \(\mo\)의 모든 절대기약표현에서 중심원소 \(z\)는 대각행렬 \(E\zeta^{(\nu)}\)로 표현되고, \(z\to\zeta^{(\nu)}\)는 \(\mZ\)의 \(\Omega\) 안에서의 1차 표현이다. 이 방식으로 매개되는 \(\mo\)의 절대기약 표현류와 \(\mZ\)의 절대기약 표현류 사이의 관계는 일대일이다. 따라서 이 표현류들의 수는 중심의 랭크와 같다.}\textsuperscript{20)}

{\footnotesize\noindent
20) 대응하는 정리들은 대수적으로 닫힌 체 \(\Omega\) 안에서의 표현뿐 아니라 각각의 자기동형체 안에서의 표현에도 성립하며, 이는 다른 곳에서 전개할 예정이다.\par}

\textbf{증명.} 양쪽 직접분해불능 아이디얼들로의 분해
\[
        \mo_\Omega=\ma_1+\ldots+\ma_t
\]
에는 중심의 체들로의 분해
\[
        \mZ_\Omega=\mZ_1+\ldots+\mZ_t
\]
가 일대일로 대응한다. 여기서 \(\mZ_\Omega\)는 \(\mo_\Omega\)의 중심이다. \(\mo\)의 기약표현에서는 한 \(\ma_\nu\)가 동형으로 사상되고 나머지는 0으로 표현되며, 표현류는 \(\ma_\nu\)에 의해 유일하게 결정된다. \(\ma_\nu\)는 완전행렬환으로 표현되고, 그러한 환의 중심은 대각행렬 \(E\zeta^{(\nu)}\)들로 이루어진다. 대응 \(z\to\zeta^{(\nu)}\)는 \(\mZ\)의 준동형이고, 이때 한 \(\mZ_\nu\)는 동형으로 사상되며 나머지는 0으로 표현된다. 이로써 모두 증명되었다.

근기 없는 \(\mo\)에 언제 역시 근기 없는 \(\mo_\Omega\)가 대응하는가라는 문제에 대해서는 다음도 따른다.

\emph{\(\mZ\)가 \(\mathsf{P}\)의 제2종 확대인 성분 \(\mZ_\nu\)를 가지면, \(\mo_\Omega\)는 반드시 근기를 갖는다.}\textsuperscript{21)} \(\mZ_\Omega\)가 이미 근기를 갖기 때문이다.

{\footnotesize\noindent
21) \(\mZ\)가 제1종 성분들만 가지면 \(\mo_\Omega\)는 근기 없게 되며, 이 역시 뒤의 다른 곳에서 증명할 예정이다. 표수가 0인 경우 I. Schur는 이에 동치인 정리, 곧 \(\mathsf{P}\) 안에서의 기약표현은 계수영역을 어떻게 확대해도 완전가약으로 남는다는 정리를 이미 증명했다(\emph{Beiträge zur Theorie der Gruppen linearer Substitutionen}, Transact. Am. Math. Soc. 15 (1909), 159쪽).\par}

\subsection*{§ 22. 아벨군에의 적용.}

\(\mG=\mathfrak{C}_1\times\mathfrak{C}_2\times\ldots\times\mathfrak{C}_r\)을 위수가 \(h_i\)인 순환군 \(\mathfrak{C}_i\)들로 분해된 유한 아벨군이라 하자. 그 원소들은
\[
        a=c_1^{\lambda_1}\ldots c_r^{\lambda_r}
\]
이다. \(\mathsf{P}\)를 그 표수가 군의 위수 \(h=h_1h_2\ldots h_r\)를 나누지 않는 체라 하자.

군환 \(\mZ\)는 모든 합
\[
        \sum a_\lambda\varrho_\lambda
        =
        \sum c_1^{\lambda_1}\ldots c_r^{\lambda_r}\varrho_\lambda
        \qquad(\varrho_\lambda\in P)
\]
으로 이루어지며, 대응
\[
        \sum \varrho_\lambda z_1^{\lambda_1}\cdots z_r^{\lambda_r}
        \longmapsto
        c_1^{\lambda_1}\cdots c_r^{\lambda_r}\varrho_\lambda
\]
에 의해 다항식영역 \(\mathsf{P}[z_1,\ldots,z_r]\)의 준동형상이다. 따라서 \(\mZ\simeq \mathsf{P}[z_1,\ldots,z_r]/\mathfrak{m}\)이고, 쉽게
\[
        \mathfrak{m}=(z_1^{h_1}-e,\ldots,z_r^{h_r}-e)
\]
를 얻는다.

확대체 \(\Omega\) 안에서 \(z_i^{h_i}-e\)는 서로 다른 일차인자 \(z_i-\varepsilon_i\)들로 분해된다. 따라서 \(\mathfrak{m}\)은 각각 하나의 영점 \((\varepsilon_1^{(\nu)},\ldots,\varepsilon_r^{(\nu)})\)만 갖는 서로 다른 소아이디얼
\[
        \mP^{(\nu)}
        =(z_1-\varepsilon_1^{(\nu)},\ldots,z_r-\varepsilon_r^{(\nu)})
\]
들의 곱 또는 교집합이다. 교집합을 취하는 데에는 직접합분해가 대응한다(§ 4).
\[
        \mZ_\Omega=\mZ_1+\cdots+\mZ_h.
\]
여기서 각각 \(\mZ_\nu\simeq\mZ_\Omega/\mP^{(\nu)}\)이고, 따라서 표현
\[
        c_1^{\lambda_1}\ldots c_r^{\lambda_r}
        \longmapsto
        \varepsilon_1^{(\nu)\lambda_1}\cdots
        \varepsilon_r^{(\nu)\lambda_r},
        \qquad\hbox{또는}\qquad
        a\to\chi^{(\nu)}(a)
\]
를 매개한다. 이 표현들이 바로 \emph{지표}이다. 그 수는 일반이론이 요구하는 대로 \(h_1h_2\ldots h_r=h\)이다.

\subsection*{§ 23. 초복소계의 행렬식.}

\(\mo=a_1\mathsf{P}+\ldots+a_h\mathsf{P}\)를 초복소계라 하자. \(\mathsf{P}\)에 \(h\)개의 부정원 \(x_1,\ldots,x_h\)를 첨가하여
\[
        \mo^*=a_1\mathsf{P}(x)+\ldots+a_h\mathsf{P}(x)
\]
를 만든다. \(x\)들은 \(a_i\)들과 가환한다고 하며, 이로써 \(\mo^*\)의 계산법칙이 정해진다.

\(\mo^*\) 안에는 \glqq{}\(\mo\)의 일반원소\grqq{}
\[
        w=a_1x_1+\ldots+a_hx_h
\]
가 있다. 한 표현에서 \(a_i\to A_i\)이면 \(w\)에는
\[
        W=A_1x_1+\ldots+A_hx_h
\]
가 대응한다. \(W\)를 그 표현에 속하는 \emph{계행렬}이라 한다. 특히 \(a_i\)들이 군을 이루어 \(\mo\)가 군환이면 \emph{군행렬}이라 한다. 정칙표현의 경우에는 \emph{정칙 계행렬}을 얻는다.

\(W\)의 원소 \(w_{ik}\)들은 \(x\)의 선형형식이다. 따라서 \(n\)차 표현의 경우 \glqq{}\emph{계행렬식}\grqq{} \(|W|\)는 \(n\)차이다. 특히 \emph{정칙 계행렬식}은 \(h\)차이다.

동치표현으로 넘어가도 \(|\mathsf{P}W\mathsf{P}^{-1}|=|\mathsf{P}|\,|W|\,|\mathsf{P}^{-1}|=|W|\)이므로 계행렬식은 변하지 않는다. \((a_1,\ldots,a_h)\)에서 새 기저 \((b_1,\ldots,b_h)\)로, \(w=\sum a_ix_i\)에서 \(w=\sum b_iy_i\)로 넘어갈 때에는 예전 원소들에 정칙 치환행렬을 갖는 치환
\[
        x_i=\sum \varrho_{ik}y_k
\]
을 시행하여 \(W\)의 새 원소들과 새 행렬식을 얻는다.

표현가군의 합성열이 주어지면 알맞은 기저에서 행렬 \(W\)는
\[
        W=\begin{pmatrix}
        W_1&&0\\
        &W_2&\\
        &\ddots&\\
        W_{ik}&&W_r
        \end{pmatrix}
\]
의 꼴이고,
\[
        |W|=|W_1|\cdot|W_2|\cdot\ldots\cdot|W_r|.
\]
임의의 표현에 나타나는 \(|W_i|\)들은 \(\mo\), 더 나아가 \(\mc\)가 근기일 때 \(\mo/\mc\)의 정칙표현에 나타나는 것 이외에는 없다.

\emph{\(\mathsf{P}\)가 대수적으로 닫혀 있으면, 기약표현에 속하는 행렬식 \(|W_i|\)는 \(x\)들에 관한 기약다항식이고, 서로 동치가 아닌 표현들에는 서로 다른 기약인자들이 속한다.}

\textbf{증명.} \(\mo\)의 모든 기약표현은 \(\mo/\mc\)의 표현이기도 하므로 근기 없는 환 \(\mo/\mc\)에 국한할 수 있다. 그 안에서 행렬단위 \(c_{ik}\)들을 기저로 도입하면 일반원소는
\[
        w=\sum c_{ik}^{(\nu)}x_{ik}^{(\nu)}
\]
이다. 기약표현들의 행렬은 \(W_\nu=(x_{ik}^{(\nu)})\)이다. 함수 \(|W_\nu|=|x_{ik}^{(\nu)}|\)들은 잘 알려진 대로 기약이고 서로 다름이 명백하다.

\(|W_\nu|\)들을 계산하려면 \(\mo/\mc\)의 정칙 계행렬을 그 기약인자들로 분해하면 된다. 각 기약인자는 기약표현의 차수가 나타내는 횟수만큼 나타난다. 대응 아이디얼 \(\ml_\nu\)가 합성열에서 그만큼 나타나기 때문이다.

\(\mo\)의 정칙 계행렬을 써도 되지만, 이 경우에는 각 기약인자가 더 자주 나타난다. 곧 대응하는 \(\ml_\nu\)가 좌측아이디얼들의 합성인자로 나타나는 횟수만큼이다. 이 정칙표현에서는 \(\mo\)를 좌측아이디얼로 보았다. \(\mo\)를 우측아이디얼로 보면 두 번째 정칙 계행렬, 곧 프로베니우스의 \glqq{}안티스트로프 행렬\grqq{}이 나온다. 이는 같은 기약인자들, 곧 모든 기약표현의 계행렬식들을 포함하지만 그 지수들은 다를 수 있다(§ 10의 예 참조).

가환계의 계행렬식은 모든 기약표현이 1차이므로 일차인자들로 분해된다. 이 일차인자들은 그 자체가 기약표현이며, 따라서 아벨군의 경우 지표들을 준다. 이 사실이 데데킨트가 비아벨군의 군행렬식을 연구할 때의 출발점이었다.

\subsection*{§ 24. 대각합과 지표.}

초복소계 \(\mo\)의 한 표현 \(\mo\sim\mathfrak D\)에서 원소 \(a\)에 행렬 \(A\)가 대응하면
\[
        \mathit{Sp}_{\mathfrak D}(a)=\mathit{Sp} A
\]
로 놓는다.

\emph{대각합은 선형함수이다.}
\[
        \operatorname{Sp}(c+d)
        =\operatorname{Sp}(c)+\operatorname{Sp}(d),
        \qquad
        \operatorname{Sp}(c\alpha)
        =\alpha\operatorname{Sp}(c).
\]

\emph{동치표현들은 같은 대각합을 갖는다.}

\emph{가약표현에서의 대각합은 합성인자들이 매개하는 표현들에서의 대각합의 합이다.}

\textbf{증명.}
\[
        \mathit{Sp}\begin{pmatrix}
        A_{11}&&0\\
        &A_{22}&\\
        &\ddots&\\
        A_{ik}&&A_{rr}
        \end{pmatrix}
        =\operatorname{Sp}(A_{11})+\operatorname{Sp}(A_{22})
        +\ldots+\operatorname{Sp}(A_{rr}).
\]

\emph{주대각합}은 정칙표현에서의 대각합이다.

\emph{축약대각합}은 서로 다른 기약표현들에서의 대각합의 합이다.

\emph{\(c\)가 극대 멱영아이디얼 \(\mc\)의 원소이면 모든 표현에서 \(\mathit{Sp}(c)=0\)이다.}

\textbf{증명.} \(\mo/\mc\)의 단순 좌측아이디얼들이 매개하는 표현들만 조사하면 된다. 다른 모든 표현에는 이 합성인자들만 나타나고, 대각합은 가법적으로 합성되기 때문이다.

그런데 \(\mathfrak{c}\)는 \(\mo/\mc\)의 모든 원소를 소멸시킨다. 따라서 \(\mc\)의 모든 \(c\)에는 영행렬이 대응하고 \(\mathit{Sp}(c)=0\)이다. 증명 끝.

\(\mo\)가 양쪽단순이고 \(\mathsf{P}\)가 대수적으로 닫혀 있어서 \(\mo\)가 행렬환
\[
        \mo=\sum c_{ik}P
\]
이라고 하자. \(\mo\)의 기약표현에서 원소 \(a=\sum c_{ik}\alpha_{ik}\)에는 행렬 \((\alpha_{ik})\)가 대응한다. 따라서
\[
\begin{aligned}
        \hbox{축약대각합 }\operatorname{Sp}_{\ml}(a)
          &=\sum_i\alpha_{ii},\\
        \hbox{주대각합 }\operatorname{Sp}_{\mo}(a)
          &=n\cdot\hbox{축약대각합}
           =n\cdot\sum_i\alpha_{ii}.
\end{aligned}
\]
특히
\[
\begin{aligned}
        \operatorname{Sp}(c_{ik})&=0\quad(i\ne k),\\
        \mathit{Sp}_{\ml}(c_{ii})&=e,\\
        \mathit{Sp}_{\mo}(c_{ii})&=n\cdot e.
\end{aligned}
\]

\emph{\(\mathsf{P}\)에서 대수적으로 닫힌 체 \(\Omega\)로 넘어가도 주대각합은 변하지 않는다.} 같은 기저에서 계산할 수 있기 때문이다. 축약대각합에는 이 명제가 성립하지 않는다.

근기 없는 계의 원소 \(a\)가 절대기약표현들에서 갖는 대각합을 \emph{지표}라 하고 \(\chi(a)\), 또는 어느 표현을 뜻하는지 나타내려면 \(\chi^{(\nu)}(a)\)로 쓴다.

\(n_\nu\)차 기약표현에서 중심원소들은 § 21에 따라 대각행렬 \(E\zeta^{(\nu)}\)로 표현된다. 여기서 \(z\to\zeta^{(\nu)}\)는 중심의 1차 표현, 곧 중심에서 \(\Omega\)로 가는 준동형이다. 따라서 그 대각합은 \(n_\nu\zeta^{(\nu)}\)이다. \emph{그러므로 중심의 준동형 \(\Theta\)와 지표 \(\chi\)는 관계}
\srcnumdisplay{(1)}{%
        \chi(z)=n_\nu\cdot\Theta(z)}
\emph{로 연결된다.} 가환인 경우에는 \(n_\nu=1\)이고 지표 자체가 준동형을 준다(§ 22 참조). 이하에서 언제나 가정하듯 체 \(\Omega\)의 표수가 0이면 (1)을 \(n_\nu\)로 나눌 수 있다.
\[
        \Theta(z)=\frac{\chi(z)}{n_\nu}.
\]
\(\Theta(z)\)의 준동형 성질은 식
\[
        \frac{\chi(z)}{n_\nu}
        +\frac{\chi(z')}{n_\nu}
        \frac{\chi(z+z')}{n_\nu},
\]
\[
        \frac{\chi(z)}{n_\nu}
        \cdot\frac{\chi(z')}{n_\nu}
        =\frac{\chi(zz')}{n_\nu}
\]
으로 표현된다.

\emph{한 표현류는 행렬들의 대각합만으로 이미 유일하게 결정된다.} (모든 행렬의 대각합을 알려면 주어진 표현에서 기저원소들의 대각합만 알면 충분함은 물론이다.)

\textbf{증명.} 표현가군 \(\mR\)은 각 기약 표현가군 \(\mM_\nu\)가 직접가수로 몇 번 나타나는지를 알면 결정된다. 그 횟수가 \(p_\nu\)이면 표현에서 \(e^{(\nu)}\)의 대각합은 \(p_\nu n_\nu\)이다. 여기서 \(n_\nu\)는 기약표현의 차수이다. 따라서
\[
        p_\nu=\frac{\mathit{Sp} e^{(\nu)}}{n_\nu}.
\]

\subsection*{§ 25. 판별식.}

\(\mo=a_1\mathsf{P}+\ldots+a_h\mathsf{P}\)를 초복소계라 하자.

\noindent\emph{판별행렬}은 그 원소들이 주대각합 \(\mathit{Sp}(a_i a_k)\)인 행렬이다.

\noindent\emph{축약판별행렬}은 대수적으로 닫힌 체에서 만든 축약대각합에 관하여 같은 방식으로 얻는 행렬이다.

이 행렬의 행렬식을 \emph{판별식}, 각각 \emph{축약판별식}이라 한다.

\noindent\emph{바탕체를 확대해도 판별식은 변하지 않는다.}

\noindent\emph{다른 기저로 넘어가면 판별식에는 변환행렬의 행렬식의 제곱이 곱해진다.}

\textbf{증명.}
\begin{align*}
(a_1,\ldots,a_h)&=(b_1,\ldots,b_h)P,\\
(a_i a_1,\ldots,a_i a_h)&=(a_i b_1,\ldots,a_i b_h)P,\\
\bigl(\operatorname{Sp}(a_i a_1),\ldots,
      \operatorname{Sp}(a_i a_h)\bigr)
&=\bigl(\operatorname{Sp}(a_i b_1),\ldots,
        \operatorname{Sp}(a_i b_h)\bigr)P.
\end{align*}
따라서 행렬로 쓰면
\srcnumdisplay{(1)}{%
\bigl(\mathit{Sp}(a_i a_k)\bigr)
=\bigl(\operatorname{Sp}(a_i b_k)\bigr)P.}
마찬가지로 \(\widetilde P\)가 전치행렬이면
\[
\begin{pmatrix}a_1\\[-1mm]\vdots\\[-1mm]a_h\end{pmatrix}
=\widetilde P
\begin{pmatrix}b_1\\[-1mm]\vdots\\[-1mm]b_h\end{pmatrix}
\]
이고, 앞에서와 같이
\srcnumdisplay{(1a)}{%
\bigl(\operatorname{Sp}(a_i b_k)\bigr)
=\widetilde P\bigl(\operatorname{Sp}(b_i b_k)\bigr).}
(1)과 (1a)에서
\[
\bigl(\mathit{Sp}(a_i a_k)\bigr)
=\widetilde P\cdot
 \bigl(\operatorname{Sp}(b_i b_k)\bigr)\cdot P
\]
를 얻는다. 행렬식을 취하면 주장이 따른다.

따라서 판별식은 \(P\)의 제곱수 인자를 제외하면 정해질 뿐이다. 그러나 판별식이 0인지 아닌지는 불변성질이다.

(1)에서는 다음도 따른다. \emph{서로 다른 두 기저를 써서 \(\left|\bigl(\mathit{Sp}(a_i b_k)\bigr)\right|\)로 판별식을 0이 아닌 한 인자를 제외하고 구할 수도 있다.}

\noindent\emph{\(\mo\)가 멱영아이디얼 \(\mc\)를 가지면 판별식은 0이다.}

\textbf{증명.} \(\mo\)의 기저로 \((c_1,\ldots,c_t,d_{t+1},\ldots,d_h)\)를 택하되, \((c_1,\ldots,c_t)\)가 \(\mc\)의 기저가 되게 하자. 판별식은
\[
\left|\begin{array}{cc}
\mathit{Sp}(c_i c_k)&\mathit{Sp}(c_i d_k)\\
\mathit{Sp}(d_i c_k)&\mathit{Sp}(d_i d_k)
\end{array}\right|
=
\left|\begin{array}{cc}
0&0\\
0&\mathit{Sp}(d_i d_k)
\end{array}\right|=0.
\]
(축약판별식에도 이 명제가 성립한다.)

\textbf{따름정리.} \emph{대수적으로 닫힌 체로 넘어간 뒤 멱영아이디얼이 나타나는 경우에도 판별식은 0이다.}

\(\mo=\ma_1+\ldots+\ma_s\)가 양쪽 직접합이고, \(M,M_1,\ldots,M_s\)가 각각 환 \(\mo,\ma_1,\ldots,\ma_s\)의 판별행렬이라 하자. 그러면
\[
M=\begin{pmatrix}
M_1& & &0\\
 &M_2& & \\
 & &\ddots& \\
0& & &M_s
\end{pmatrix},
\]
따라서 \(|M|=|M_1|\,|M_2|\cdot\ldots\cdot|M_s|\)이다.

\textbf{증명.} \(\ma_1,\ldots,\ma_s\)의 기저들을 합쳐 \(\mo\)의 기저를 택하자. \(a\in\ma_i\)이면
\[
        \operatorname{Sp}_{\mo}(a)
        =\operatorname{Sp}_{\ma_i}(a)
\]
이다. 양쪽 모두 각각 환 \(\mo\), \(\ma_i\)에서의 주대각합을 뜻한다. 또한 \(a_i\in\ma_i\), \(a_k\in\ma_k\)이고 \(i\ne k\)이면 \(a_i a_k=0\), 따라서 \(\mathit{Sp}(a_i a_k)=0\)이다. 여기서 주장이 따르며, 축약판별식에 대해서도 똑같이 따른다.

행렬환 \(\sum c_{ik}\mathsf{P}\)의 \emph{판별식}을 만들기 위해 두 서로 다른 기저를 택하자. 한 번은 \(c_{ik}\)들을 첫째 지수에 따라, 다른 한 번은 둘째 지수에 따라 배열한다. 곱셈표는 다음과 같다.
\[
\begin{array}{c|c|c|c}
& \overbrace{c_{11}\;\ldots\;c_{1n}}^{\mathfrak r_1}
& \overbrace{c_{21}\;\ldots\;c_{2n}}^{\mathfrak r_2}
& \ldots \\ \hline
\mathfrak l_1\left\{\begin{array}{c}c_{11}\\[-1mm]\vdots\\[-1mm]c_{n1}\end{array}\right.
& (c_{ik})&0&0\\[1.3em]\hline
\mathfrak l_2\left\{\begin{array}{c}c_{12}\\[-1mm]\vdots\\[-1mm]c_{n2}\end{array}\right.
&0&(c_{ik})&0\\[1.3em]\hline
\vdots&0&0&(c_{ik})
\end{array}
\]
판별식은
\[
D=\left|\begin{array}{cccc}
\mathit{Sp}(c_{11})& & &0\\
&\mathit{Sp}(c_{22})& & \\
& &\ddots& \\
0& & &\mathit{Sp}(c_{nn})
\end{array}\right|^{n}
=\begin{cases}
 e,&\hbox{축약대각합의 경우},\\
 n^{n^2}\cdot e&\hbox{주대각합의 경우}.
\end{cases}
\]

\textbf{따름정리.} \emph{\(\mo\)가 대수적으로 닫힌 확대체에서 차수가 \(n_1,\ldots,n_s\)인 행렬환들로 분해되면 판별식은}
\[
        D=n_1^{n_1^2}n_2^{n_2^2}\ldots n_s^{n_s^2}\cdot e
\]
\emph{이고, 축약판별식은}
\[
        D_{\rm red}=e
\]
\emph{이다.}

따라서 근기 없는 계에서는 축약판별식이 언제나 0이 아니다. 다른 판별식은 \(n_i\) 가운데 어느 것도 체의 표수로 나누어지지 않을 때에만 0이 아니다. 그러므로 다음이 성립한다.

\noindent\emph{축약판별식이 0이 아니라는 것은 체 \(\mathsf{P}\)를 대수적으로 닫은 뒤 근기가 없는 계가 되기 위한 필요충분조건이다.}

\noindent\emph{표수가 0인 경우 판별식이 0이 아니라는 것은 체 \(\mathsf{P}\)를 대수적으로 닫은 뒤 근기가 나타나지 않기 위한 필요충분조건이다.}\textsuperscript{22)}

{\footnotesize\noindent
22) 가환계에 대해서는 E. Noether, \emph{Diskriminantensatz für Ordnungen\ldots}, J. f. M. 157 (1927), 82--104쪽, §§ 4--6 참조. 그곳의 증명법은 같지만 세부적으로 더 복잡하다. 한 기저에서 다른 기저로 넘어가는 과정(§ 4, 4)은 여기, 곧 이 절 첫머리에서 준 과정으로 바꾸어야 한다. 행렬식 \(\left|(\varrho_i\varrho_k)\right|\)는 0이기 때문이다.\par}

\subsection*{§ 26. 군환의 분류.}

\noindent\emph{\(a_1,\ldots,a_h\)가 유한군의 원소들이고 표수가 \(h\)의 약수가 아닌 체에서 군환을 만들면, 판별식 \(D\ne0\)이고 따라서 군환은 근기 없는 환이다.}

\textbf{증명.} 먼저 다음을 보이자. 이제부터 \(\mathit{Sp}\)는 주대각합을 뜻한다.
\[
        \operatorname{Sp}(e)=he,
        \qquad
        \operatorname{Sp}(a_i)=0\quad(a_i\ne e).
\]
정칙표현에서
\[
e\longmapsto
\begin{pmatrix}
e&&&0\\
0&e&\cdots&0\\
&&\ddots&\\
0&&&e
\end{pmatrix},
\qquad
\operatorname{Sp}(e)=he
\]
이다. \(a_i\ne e\)이면
\[
        a_i a_k=a_\lambda,
        \qquad \lambda\ne k.
\]
따라서 곱 \(a_i a_1,\ldots,a_i a_h\)를 \(a_1,\ldots,a_h\)로 나타내는 행렬은 주대각선의 모든 원소가 0이고, 그러므로 \(\mathit{Sp}(a_i)=0\)이다.

이제 두 기저 \(a_1,\ldots,a_h\)와 \(a_1^{-1},\ldots,a_h^{-1}\)를 택한다. 행렬 \(\bigl(\mathit{Sp}(a_i a_k^{-1})\bigr)\)은
\[
\begin{pmatrix}
he&&&0\\
&he&&\\
&&\ddots&\\
0&&&he
\end{pmatrix}.
\]
이다. 따라서
\[
        D=h^h e\ne0.
\]
증명 끝.

군원소의 \emph{류}란 군환에서 만든 합
\srcnumdisplay{(1)}{K_a=\sum s^{-1}as}
을 말한다. 여기서는 서로 다른 \(s^{-1}as\)들에 대해서만 합한다. \emph{류 \(K_a\)들은 중심원소이다.} 모든 군원소와 가환하기 때문이다. \emph{류 \(K_a\)들은 중심을 생성한다.} 실제로 군환의 원소 \(\sum a_i\varrho_i\)가 모든 \(s\)와 가환하면
\[
\begin{aligned}
\sum_i a_i\varrho_i
&=s^{-1}\biggl(\sum_i a_i\varrho_i\biggr)s
 =\sum_i(s^{-1}a_i s)\varrho_i\\
&=\frac1h\sum_i\biggl(\sum_s s^{-1}a_i s\biggr)\varrho_i
 =\frac1h\sum_i\biggl(\frac h{h_i}K_i\biggr)\varrho_i.
\end{aligned}
\]
여기서 \(h\)는 군원소의 수이고 \(h_i\)는 류 \(K_i\)의 원소 수이다.

따라서 중심의 랭크는 켤레군원소류의 수와 같고, 그러므로 \emph{절대기약표현의 수 역시 이 류의 수와 같다.}

행렬 \(A\)와 \(S^{-1}AS\)는 같은 대각합을 갖는다. 따라서 군원소 \(a\)와 \(s^{-1}as\)는 같은 대각합을 갖는다.

이제 (1)의 양변에서 대각합을 취하면
\[
        \chi(K_i)=h_i\chi(a_i)
\]
를 얻는다. 말로 쓰면 다음과 같다.

\noindent\emph{군원소의 지표는 그 류의 지표를 그 류의 원소 수로 나눈 것과 같다.}

\begin{center}
(1928년 8월 12일 접수.)
\end{center}

\emph{\(\mG\)가 완전가약이면 모든 정규부분군은 직접인자이고, 다른 인자는 주어진 \(\mG\)의 곱분해에 나타나는 단순군들의 곱으로 택할 수 있다. 정규부분군 \(\mH\)도 완전가약이다.}

\emph{증명.} 
\begin{equation}
        \mG=\mH\mA_1\mA_2\ldots\mA_n.\tag{3}
\end{equation}
\(\mH_i=\mH\mA_1\mA_2\ldots\mA_i\)로 놓으면 \(\mH_{i+1}=\mH_i\mA_{i+1}\)이다. 이제 \(\mA_{i+1}\subseteqq\mH_i\), 따라서 \(\mH_{i+1}=\mH_i\)이면 (3)에서 인자 \(\mA_{i+1}\)을 생략한다. 그렇지 않으면 \(\mH_i\cap\mA_{i+1}\)은 \(\mA_{i+1}\)의 진정규부분군이므로 \(\mE\)와 같고, \(\mH_i\times\mA_{i+1}\)은 직접곱이다. 불필요한 인자들을 생략하면 (3)은
\[
        \mG=\mH\times\mA_{i_1}\times\ldots\times\mA_{i_r}
\]
이 되어 \(\mH\)가 직접인자임을 보인다. 또한
\[
        \mH\simeq
        \mG/\mA_{i_1}\times\cdots\times\mA_{i_r}
        \simeq
        \mA_{j_1}\times\cdots\times\mA_{j_\mu},
\]
여기서 \(\mA_{j_1},\ldots,\mA_{j_\mu}\)는 나머지 \(\mA_i\)들이다. 따라서 \(\mH\)도 완전가약이다.

\emph{따름정리.} \emph{완전가약군을 직접분해불능 인자들로 나타내는 모든 분해는 단순인자들로의 분해이며, 따라서 합성열을 낳는다. 이로부터 인자들이 동형을 제외하면 유일하게 결정된다는 것이 따른다.}

\clearpage
\setcounter{footnote}{0}
\phantomsection
\addcontentsline{toc}{section}{35. 정수적 함수로 이루어진 극대영역에 관하여}
%
%
%
%
%
%

\section*{35. 정수적 함수로 이루어진 극대영역에 관하여}

\begin{center}
에미 뇌터 (괴팅겐)\par
\emph{Rec. Soc. Math. Moscou 36 (1929), 65--72쪽}
\end{center}

다음 고찰은 상대적으로 정수인 함수에 관한 힐베르트의 문제와 밀접하게 맞닿아 있다\textsuperscript{1)}. 이 고찰은 E. 헤케가 우연히 던진 물음에서 나왔다. 그는 보조정리로 필요했던 한 특수한 경우를 직접 계산하여 해결할 수 있었다\textsuperscript{2)}.

문제는 다음과 같다. 주어진 영역 \(\mathfrak Z\) 안에서 \(x_1,\ldots,x_n\)의 정수적 함수 --- 다항식, 거듭제곱급수 또는 그러한 것들의 몫 --- 로 이루어진 \emph{극대영역}이란 정수 분모를 첨가하여 더는 확대할 수 없는 영역을 뜻한다고 하자. 즉 \(g(x)\in\mathfrak Z\)이고 \(a\)가 0이 아닌 유리정수일 때, \(a\cdot g(x)\)가 영역에 속하면 언제나 \(g(x)\)도 영역에 속하는 것이다. \(\mathfrak Z\) 안의 각 정역 \(\mathfrak I\)은 하나의 극대영역 \(\mathfrak M(\mathfrak I)\)을 유일하게 생성한다. 이것은 적어도 하나의 \(a\ne0\)에 대하여 \(a\cdot g(x)\in\mathfrak I\)가 되는 \(\mathfrak Z\)의 모든 함수 \(g(x)\)로 이루어진다.

원래 정역 \(\mathfrak I\)이 유한 생성된 경우, 곧 \(\mathfrak Z\)의 유한개 함수 \(f_1(x),\ldots,f_r(x)\)에 관한 모든 정수 계수 다항식으로 이루어진 경우에 이 극대영역 \(\mathfrak M(\mathfrak I)\)의 분모 \(a\)에 관하여 무엇을 말할 수 있는지를 묻는다. 여기서 \(\mathfrak Z\)로는 \(x_1,\ldots,x_n\)의 모든 정수 계수 다항식 또는 거듭제곱급수를 택하고, 몫을 다룰 때에는 \(f(x)\)들에 나타나는 분모와 그 거듭제곱 이외의 분모가 없는 것들만 택한다.

거듭제곱급수나 그 몫을 다룰 때에는 \(f(x)\) 사이의 종속성이 오직 대수적이라는 부가가정 아래, \(a\)들을 언제나 다음과 같이 택할 수 있음을 보이겠다. 모든 \(a\)를 통틀어 서로 다른 소수는 유한개만 그 인수로 나타난다. 다만 이 소수들은 일반적으로 임의의 높은 거듭제곱으로 나타난다.

마지막 사실은 즉시 분명하다. \(a\cdot g(x)\in\mathfrak I\)이지만 어떠한 지수 \(\varrho\)에 대해서도 \(a^{\varrho-1}g(x)^\varrho\)가 \(\mathfrak I\)에 속하지 않는다면 모든 거듭제곱 \(a^\varrho\)이 나타난다. 예를 들어 \(\mathfrak I=[2x]\)이면 \(x^\varrho=(2x)^\varrho/2^\varrho\)이다.

유계성의 물음은 다음처럼만 정식화할 수 있다. 일반적으로 무한한 \(\mathfrak M(\mathfrak I)\)의 생성계를, 대응하는 분모에서 소수들의 지수가 유계가 되도록 제시할 수 있는가? 관련된 소수가 유한개이므로, 알려진 논법\textsuperscript{3)}에 따라 이 물음은 \(\mathfrak M(\mathfrak I)\)의 유한 생성성과 같다. 따라서 이는 정수적 형식으로 표현한 상대적으로 정수인 함수에 관한 힐베르트의 문제와 같다. 여기서 사용하는 방법으로는 이 문제를 결정할 수 없다\textsuperscript{4)}.

분모 \(a\)에 서로 다른 소수 \(p\)가 유한개만 나타남을 보이는 증명은 다음 사실에 근거한다. 그러한 각 \(p\)에는 \(\mathfrak I\)을 생성하는 \(f(x)\)들 사이의 적어도 하나의 \(p\)를 법으로 하는 관계가 대응하며, 이 관계는 정수 계수로는 \(x\)에 관한 항등식이 아니다. 달리 말해, \(\mathfrak o\)가 부정원 \(u_1,\ldots,u_r\)의 정수 계수 다항식환이면 대응 \(u_i\mapsto f_i(x)\)에 따라 \(\mathfrak I\)은 \(\mathfrak o\)의 준동형상이다. 이 준동형은 \(\mathfrak o\)의 한 소아이디얼 \(\mathfrak a\)를 통하여 주어진다. 즉 \(\mathfrak I\)은 잉여류환 \(\mathfrak o/\mathfrak a\)와 동형이다. 이에 따라 \(\mathfrak I_p\)는 \(\mathfrak o_p\)의 준동형상이다. 아래첨자는 \(p\)를 법으로 하는 잉여류로 넘어감을 뜻하며, \(\mathfrak I\)의 분모에 나타나는 유한개의 소수를 제외하면 이 이행이 가능하다. 이 준동형은 다시 \(\mathfrak o_p\)의 한 소아이디얼 \(\mathfrak c_p\)를 통하여 주어지며, \(\mathfrak c_p\)는 \(\mathfrak a_p\)를 포함한다. \(p\)가 \(\mathfrak M(\mathfrak I)\)의 분모에 나타나는 것은 \(\mathfrak c_p\)가 \(\mathfrak a_p\)의 진상위아이디얼, 곧 아이디얼 의미의 진약수가 될 때 그리고 그때뿐이다. 이는 \(\mathfrak I_p\)의 초월 차수가 \(\mathfrak I\)의 초월 차수보다 낮아지거나 \(\mathfrak a_p\)가 소아이디얼의 성질을 잃을 때에만 일어날 수 있다\textsuperscript{5)}. 이런 일이 유한개의 \(p\)에 대해서만 일어난다는 사실은 다음 두 결과에서 따라온다. \(p\)를 법으로 해도 대수적 종속성은 함수행렬식의 소멸을 수반하며 (§ 1), 더 나아가 \(\mathfrak a\)는 절대소아이디얼이므로 유한개의 \(p\)를 법으로 할 때에만 소아이디얼의 성질을 잃는다\textsuperscript{6)}.

유리정수 대신 한 유한 대수적 수체의 대수적 정수들을 놓고, 소수 대신 그 수체의 소아이디얼들을 놓아도 모든 고찰은 조금만 고치면 그대로 성립한다는 점도 덧붙인다(§ 4).

\bigskip
\paragraph{§ 1. 표수 \(p\)인 계수영역에서의 함수행렬식.}

\textbf{1.}
\(f_1(x_1,\ldots,x_n),\ldots,f_t(x_1,\ldots,x_n)\)을 표수 \(p\)인 체 \(P\)를 계수체로 하는 유리함수, 또는 더 일반적으로 형식적으로 정의된 거듭제곱급수나 그 몫이라 하자. 도함수 \(\partial f_i/\partial x_k\)도 형식적으로 정의되고 통상의 계산법칙이 모두 유지된다고 하자. \(\varOmega\)는 \(\mathrm{P}\)의 대수적으로 닫힌 확대체라 하자.

그러면 표수 0에서와 마찬가지로 다음이 성립한다.

\textbf{정리.}
\emph{\(f(x)\)들 사이에 \(\varOmega\)의 계수를 갖는 대수적 종속관계가 있으면, 함수행렬 \((\partial f_i/\partial x_k)\) \((i=1,\ldots,t;\ k=1,\ldots,n)\)의 랭크는 \(t\)보다 작다.}\textsuperscript{7)}

\emph{증명.}
다항식환 \(\varOmega[u_1,\ldots,u_t]\)에서, \(G(f)=0\)이 \(x\)에 관한 항등식이 되는 모든 다항식 \(G(u)\)의 소아이디얼을 \(\mathfrak m\)이라 하자. 가정에 따라 \(\mathfrak m\)은 영아이디얼과 다르다. \(\mathfrak m\)에서 차수가 가장 낮은 다항식 \(G(u)\ne0\), 따라서 상수도 아닌 것을 택하자. 특히 \(G(u)\)는 기약이다. 그러면 통상적인 방법으로
\[
\begin{aligned}
\left(\frac{\partial G}{\partial u_1}\right)_{f}\frac{\partial f_1}{\partial x_1}
 +\cdots+
\left(\partial G/\partial u_t\right)_f\partial f_t/\partial x_1&=0,\\[-1mm]
&\vdots\\[-1mm]
\left(\frac{\partial G}{\partial u_1}\right)_{f}\frac{\partial f_1}{\partial x_n}
 +\cdots+
\left(\partial G/\partial u_t\right)_f\partial f_t/\partial x_n&=0
\end{aligned}
\]
을 얻는다. 따라서 함수행렬의 랭크가 \(t\)보다 작거나, 아니면 모든 \(\left(\partial G/\partial u_i\right)_f\)가 사라진다. 그런데 \(G(u)\)를 차수가 가장 낮은 다항식으로 택했으므로, 후자의 경우에는 모든 \(\partial G/\partial u_i\) 자체가 사라진다.

그러므로 \(G(u)=H(u_1^p,\ldots,u_t^p)\)이다. 계수체 \(\varOmega\)는 대수적으로 닫혀 있으므로 완전체이고, 따라서 더 나아가
\[
G(u)=H(u_i^p)=(T(u))^p
\]
가 되어 \(G(u)\)의 선택에 모순된다. 물론 이 증명은 표수 0에서도 성립하며, 그 경우에는 마지막 추론이 필요 없다.

\textbf{2. 따름정리.}
\(F_1(x_1,\ldots,x_n),\ldots,F_t(x_1,\ldots,x_n)\)을 함수행렬의 랭크가 정확히 \(t\)인 정수적 함수들, 곧 다항식, 거듭제곱급수 또는 그 몫이라 하자. 소수 \(p\)는 \(F_i(x)\)들의 분모에도, 함수행렬의 모든 \(t\)행렬식에도 공통으로 나누어 들어가지 않도록 택하자.

그러면 \(F_1(x),\ldots,F_t(x)\)는 \(p\)를 법으로 해도 대수적으로 독립이다.

\emph{증명.}
\(\mathrm{P}\)를 \(p\)에 대한 잉여류체라 하고, \(f_1(x),\ldots,f_t(x)\)를 \(F_1(x),\ldots,F_t(x)\)의 정수 계수들을 \(p\)에 대한 잉여류로 바꾸어 얻는 \(P\) 위의 함수들이라 하자. \(p\)가 \(F_i(x)\)의 분모에 나타나지 않으므로 \(f_i(x)\)들은 정의된다. 함수행렬 \((\partial f_i/\partial x_k)\)도 \((\partial F_i/\partial x_k)\)의 정수 계수들을 \(p\)에 대한 잉여류로 바꾸어 얻는다. \((\partial F_i/\partial x_k)\)에는 \(F_i(x)\)들에 없는 다른 분모가 나타나지 않는다. \(p\)의 선택에 따라 \((\partial f_i/\partial x_k)\)의 랭크는 \(t\)로 유지된다. 따라서 \(f_i(x)\)들은 \(\varOmega\)에서 대수적으로 독립이고, 더욱이 \(\mathrm{P}\)에서도 대수적으로 독립이다.

\bigskip
\paragraph{§ 2. 정리의 정식화.}

\textbf{1.}
서론에서와 같이 \(\mathfrak I\)은 \(x_1,\ldots,x_n\)의 정수적 함수, 곧 유리정수 계수의 다항식이나 거듭제곱급수 또는 그 몫으로 이루어진 정역이라 하자. 거듭제곱급수는 형식적으로 정의된 것이어도 되고 어떤 영역에서 수렴하는 것이어도 된다. 사용되는 것은 해당 함수들이 정역을 이룬다는 사실뿐이다.

\(\mathfrak Q\)를 \(\mathfrak I\)의 분수체라 하고, \(\mathfrak Z\)를 \(\mathfrak I\)과 \(\mathfrak Q\) 사이의 정역이라 하자. 서론에서와 같이, \(a\ne0\)이 정수이고 \(g(x)\in\mathfrak Z\)일 때 \(a\cdot g(x)\)가 \(\mathfrak M\)에 속하면 언제나 \(g(x)\)도 \(\mathfrak M\)에 속하는 경우 \(\mathfrak M\)을 \(\mathfrak Z\) 안에서 \emph{극대}라고 한다. \(\mathfrak I\)에 대해서는 \(\mathfrak I\)이 \(\mathfrak Z\) 안에서 생성하는 극대영역 \(\mathfrak M(\mathfrak I)\)이 유일하게 존재한다. 이는 \(\mathfrak I\)을 포함하는 \(\mathfrak Z\) 안의 모든 극대영역의 교집합이다. 실제로 \(\mathfrak Z\) 자체가 그러한 영역이고, 임의의 개수의 극대영역들의 교집합도 극대이며 \(\mathfrak I\)을 포함한다. \(\mathfrak M(\mathfrak I)\)은 적어도 하나의 \(a\ne0\)에 대하여 \(a\cdot g(x)\in\mathfrak I\)가 되는 \(\mathfrak Z\)의 모든 함수 \(g(x)\)로 이루어진다.

실제로 이 \(g(x)\)들은 \(\mathfrak I\)을 포함하는 \(\mathfrak Z\) 안의 극대영역 \(\mathfrak N\)을 이룬다. \(b\ne0\), \(h(x)\in\mathfrak Z\), \(b\cdot h(x)\in\mathfrak I\)이면 \(ab\cdot h(x)\in\mathfrak I\)이고 \(ab\ne0\)이므로 \(h(x)\in\mathfrak N\)이다. 한편 \(\mathfrak Z\) 안에서 \(\mathfrak I\)을 포함하는 모든 극대영역 \(\mathfrak M\)은 \(\mathfrak N\)을 포함한다. 실제로 \(a\cdot g(x)\in\mathfrak I\)이면 \(\mathfrak I\subseteq\mathfrak M\)이므로 \(a\cdot g(x)\in\mathfrak M\)이고, 따라서 \(g(x)\in\mathfrak M\)이다.

\textbf{2.}
이제 \(\mathfrak I\)을 단위원을 가진 유한 생성 정역이라 하고, \(f_1(x),\ldots,f_r(x)\)를 그 정수성 기저라 하자. 즉 \(\mathfrak I\)은 \(f_i(x)\)들에 관한 모든 정수 계수 다항식으로 이루어진다. 더 나아가 \(\mathfrak Z\)는 \(x\)에 관한 모든 정수 계수 다항식 또는 거듭제곱급수로 이루어지거나, 유한개의 \(f_i(x)\)에 나타나는 분모 이외의 분모가 없는 그러한 것들의 몫으로 이루어진다고 하자. 이 분모들은 물론 임의의 거듭제곱으로 허용한다.

또 적어도 하나의 \(f_i(x)\), 일반적으로는 모든 \(f_i(x)\)에 \(x\)가 실제로 나타난다고 가정하자. 그렇지 않으면 \(\mathfrak I\)과 \(\mathfrak Z\)가 단지 한 분수 \(1/c\)의 모든 정수 계수 다항식으로 이루어져 증명할 것이 없기 때문이다.

다항식 또는 유리함수를 다루는 경우에는 분수체 \(\mathfrak Q\)가 유리수체 \(\mathbf{K}\)에 관하여 초월 차수 \(t\), \(1 \leqq t \leqq r\)를 갖는다. 예를 들어 \(f_1(x),\ldots,f_t(x)\)가 기약계(주 \textsuperscript{5)})이면, \(\mathfrak Q\)는 순수초월체 \(\mathbf{K}(f_1(x),\ldots,f_t(x))\)의 유한 대수적 확대이다. 동시에 \(f_1(x),\ldots,f_r(x)\)의 함수행렬과 \(f_1(x),\ldots,f_t(x)\)의 함수행렬은 모두 정확히 랭크 \(t\)를 갖는다\textsuperscript{8)}.

거듭제곱급수나 그 몫의 경우에도 \(\mathfrak Q\)가 이 구조를 유지하도록 다음 부가가정을 둔다. \(t\)가 \(f_i(x)\)들의 함수행렬의 랭크이고, 동시에 예컨대 \(f_1(x),\ldots,f_t(x)\)의 함수행렬의 랭크도 \(t\)라고 하자. 따라서 이 함수들은 § 1의 1에 따라 기약계를 이룬다. 이때 다시 \(\mathfrak Q\)가 순수초월체 \(\mathbf{K}(f_1(x),\ldots,f_t(x))\)의 유한 대수적 확대라고 가정한다. 즉 \(f_{t+1}(x),\ldots,f_r(x)\)가 \(\mathbf{K}(f_1(x),\ldots,f_t(x))\)에 관하여 대수적이라고 가정한다\textsuperscript{9)}. 또한 표수가 0이므로 모든 \(\partial f_i/\partial x_k\)가 사라질 수 없어서 다시 \(t \geqq 1\)이다.

이 조건은 정수성 기저에서 택한 함수 \(t\)개 가운데 함수행렬의 랭크가 정확히 \(t\)인 모든 계에 대해서도 성립한다. 실제로 그러한 함수들은 기약계를 이루므로 나머지 함수들이 그 계에 대하여 대수적으로 종속된다.

\textbf{3.}
이제 \(\mathfrak o\)를 부정원 \(u_1,\ldots,u_r\)의 정수 계수 다항식환이라 하자. 대응 \(u_i\mapsto f_i(x)\)가 직접 보여 주듯이 \(\mathfrak I\)은 \(\mathfrak o\)의 환준동형상이다. 따라서 \(\mathfrak I\)은 \(\mathfrak o/\mathfrak a\)와 환으로서 동형이다. 여기서 \(\mathfrak I\)은 영인수가 없는 환이므로 \(\mathfrak a\)는 소아이디얼이고, \(G(f)\)가 \(x\)에 관하여 항등적으로 사라지는 모든 다항식 \(G(u)\)로 이루어진다.

서론에서 이미 언급했듯이 증명할 정리는 다음과 같다.

\textbf{정리.}
\emph{2의 부가가정 아래 \(\mathfrak o_p,\mathfrak a_p,\mathfrak I_p\)를 \(\mathfrak o,\mathfrak a,\mathfrak I\)의 정수 계수들을 \(p\)를 법으로 하는 잉여류로 바꾸어 얻는 영역들이라 하자. 그러면 많아야 유한개의 소수 --- \(\mathfrak I_p\)가 정의되려면 \(f_i(x)\)들의 분모에 나타나는 소수들도 여기에 포함된다 --- 를 제외하면, \(\mathfrak o_p\)에서 \(\mathfrak I_p\)로 가는 준동형은 \(\mathfrak a_p\)를 통하여 주어진다. 따라서 \(\mathfrak I_p\)는 \(\mathfrak o_p/\mathfrak a_p\)와 환으로서 동형이고, 다시 \(\mathfrak o/(\mathfrak a,p\mathfrak o)\)와 동형이다.}

이 정리로부터 \(\mathfrak M(\mathfrak I)\)의 분모 \(a\)에 소수가 유한개만 나타난다는 사실은 다음 따름정리를 통하여 얻어진다.

\textbf{따름정리.}
\emph{\(p\)를 \(\mathfrak I_p\cong\mathfrak o_p/\mathfrak a_p\)가 되도록 택했다고 하자. \(p\cdot g(x)\in\mathfrak I\)이고 \(g(x)\in\mathfrak Z\)이면 언제나 \(g(x)\in\mathfrak I\)이다.}

달리 말하면, \(\mathfrak I\)은 유한개의 예외소수와 다른 모든 소수에 관하여 \(\mathfrak Z\) 안에서 극대이다.

\(\mathfrak o_p/\mathfrak a_p\)는 \(\mathfrak o/(\mathfrak a,p\mathfrak o)\)와 동형이다. 실제로 정의상 \(\mathfrak o_p=\mathfrak o/p\mathfrak o\)이고, 이에 따라 \(\mathfrak a_p=(\mathfrak a,p\mathfrak o)/p\mathfrak o\)이다. 그러므로 제1동형정리에 따라 \(\mathfrak o_p/\mathfrak a_p\), 그리고 가정에 의해 \(\mathfrak I_p\)도 \(\mathfrak o/(\mathfrak a,p\mathfrak o)\)와 동형이다.

이는 다음을 뜻한다.
\[
H(f_i(x))\equiv0\pmod p\quad\text{\(\mathfrak I\) 안에서}
\quad\Longrightarrow\quad
H(u)\equiv0\pmod{(\mathfrak a,p\mathfrak o)}
\quad\text{\(\mathfrak o\) 안에서}.
\]

이를 다음과 같이 바꿔 쓸 수 있다. \(H(f_i(x)) = pg(x)\), 여기서 \(H(f_i(x))\in\mathfrak I\), \(g(x)\in\mathfrak Z\)이면 \(H(u_i)-pF(u_i)\equiv0\pmod{\mathfrak a}\)이다. 따라서 \(H(f_i(x))=pF(f_i(x))\)이고 \(g(x)=F(f_i(x))\), 즉 \(g(x)\in\mathfrak I\)이다.

이를 유한번 반복하면 다음을 얻는다. \(a\)가 어느 예외소수로도 나누어지지 않을 때 \(a\cdot g(x)\in\mathfrak I\), \(g(x)\in\mathfrak Z\)이면 언제나 \(g(x)\in\mathfrak I\)이다. 이로써 따름정리가 증명되었다.

\bigskip
\paragraph{§ 3. 정리의 증명.}

\textbf{1. 보조정리.}
\emph{\(\mathfrak I\)과 \(\mathfrak o\) 사이의 준동형을 규정하는 소아이디얼 \(\mathfrak a\)(§ 2, 3)는 절대소아이디얼이다.}

\(\mathfrak o^*\)를 \(u_1,\ldots,u_r\)에 관한 다항식환으로서 임의의 대수적 수, 분수형 대수적 수도 계수로 허용한 것이라 하자. \(\mathfrak a^*\)는 \(\mathfrak o^*\)에서 \(\mathfrak a\)가 생성하는 아이디얼이라 하자. 따라서 \(\mathfrak a^*\)는
\[
a_1G_1(u)+\cdots+a_sG_s(u)
\]
꼴의 모든 선형결합으로 이루어진다. 여기서 \(a_i\)는 임의의 대수적 수이고 \(G_i(u)\)는 \(\mathfrak a\)의 다항식이다. \(\mathfrak a^*\)가 소아이디얼임을 보여야 한다.

\(\mathfrak a^*\)가 \(\mathfrak o^*\)에서 \(\mathfrak I^*\)로 가는 준동형을 규정함을 보이면 증명이 끝난다. 여기서 \(\mathfrak I^*\)는 \(f_i(x)\)에 관한 다항식 가운데 계수로 임의의 대수적 수를 허용한 것들로 이루어진 정역이다. 즉 \(H(u)\in\mathfrak o^*\)이고 \(H(f_i(x))=0\)이면 \(H(u)\in\mathfrak a^*\)임을 보이면 된다.

\(\omega_1,\ldots,\omega_s\)를 \(H(u)\)의 유한개의 계수들이 생성하는 유한 수체의 선형독립 기저라 하자. 그러면
\[
a\cdot H(u)=\omega_1H_1(u)+\cdots+\omega_sH_s(u)
\]
로 쓸 수 있다. 여기서 \(H_i(u)\in\mathfrak o\)이고 \(a\ne0\)은 유리정수이다. \(H(f)=0\)에서
\[
\omega_1H_1(f)+\cdots+\omega_sH_s(f)=0
\]
을 얻고, 따라서 \(H_1(f)=0,\ldots,H_s(f)=0\)이다. 그러므로 \(H_1(u),\ldots,H_s(u)\in\mathfrak a\)이고 결국 \(H(u)\in\mathfrak a^*\)이다.

\textbf{2.}
이제 정리는 다음 형태로 증명할 수 있다.

\(p\)를 다음 조건들이 성립하도록 택하자.
\begin{itemize}
\item[I.] \(\mathfrak I_p\)가 정의된다.
\item[II.] \(\mathfrak a_p\)가 소아이디얼, 더구나 절대소아이디얼로 남는다.
\item[III.] \(\mathfrak I_p\)의 초월 차수가 \(\mathfrak I\)의 초월 차수보다 낮아지지 않는다.
\end{itemize}
그러면 \(\mathfrak o_p\)에서 \(\mathfrak I_p\)로 가는 준동형은 \(\mathfrak a_p\)를 통하여 주어진다. 이 조건들은 많아야 유한개의 예외소수를 빼면 모든 소수에 대하여 충족된다.

조건 I, II, III 각각에 관하여 예외소수가 유한개뿐이라는 사실은 앞의 결과에서 곧바로 따라온다.

I은 분명하다. 예외소수는 \(f_i(x)\)들의 분모에 나타나는 유한개의 소수뿐이다. II는 보조정리에 따라 \(\mathfrak a\)가 절대소아이디얼이라는 사실에서 나온다. 따라서 유한개를 제외한 모든 소수를 법으로 해도 이 성질을 유지한다(주 \textsuperscript{6)} 참조). 끝으로 III은 § 1, 2의 따름정리에서 증명하였다. 실제로 § 2, 2에 따라 \(f_1(x),\ldots,f_t(x)\)의 함수행렬은 정확히 랭크 \(t\)이다.

이제 이 예외소수들과 다른 \(p\)를 택하자. \(\mathfrak o_p\)에서 \(\mathfrak I_p\)로 가는 준동형은 한 아이디얼 \(\mathfrak c_p\)를 통하여 주어진다. \(\mathfrak I_p\)도 영인수가 없는 환이므로 \(\mathfrak c_p\)는 소아이디얼이고, 그 초월 차수는 \(\mathfrak I_p\)의 초월 차수로 주어지므로 적어도 \(t\)이다. 실제로 \(\mathfrak I_p\cong\mathfrak o_p/\mathfrak c_p\)이므로 \(\mathfrak c_p\)의 잉여류체는 \(\mathfrak I_p\)의 분수체와 동형이다.

\(\mathfrak a_p\)도 소아이디얼이고 \(\mathfrak c_p\)의 배수, 곧 부분아이디얼이므로, \(\mathfrak a_p\)와 \(\mathfrak c_p\)는 초월 차수가 같을 때 그리고 그때에만 일치한다\textsuperscript{10)}. \(\mathfrak c_p\)의 배수인 \(\mathfrak a_p\)의 초월 차수도 적어도 \(t\)이고, 실제로 \(u_1,\ldots,u_t\)의 \(\mathfrak a_p\)에 대한 잉여류들이 기약계를 이룬다. 이 초월 차수가 정확히 \(t\)임을 보이면 모든 것이 증명된다.

여기서 거듭제곱급수나 그 몫의 경우에는 \(\mathfrak I\)의 분수체 \(\mathfrak Q\)의 구조에 관한 부가가정(§ 2, 2)을 사용한다. 그 가정에 따라 소아이디얼 \(\mathfrak a\)의 초월 차수는 언제나 \(t\)였다. 번호는 \(u_1,\ldots,u_t\)의 \(\mathfrak a\)에 대한 잉여류들이 기약계를 이루고 \(u_{t+1},\ldots,u_r\)의 잉여류들이 그 기약계에 대하여 대수적으로 종속되도록 정하였다. 따라서 \(\mathfrak a\)에는 정수 계수이고 원시적인 다음 다항식들이 들어 있다.
\[
G_1(u_{t+1};u_1,\ldots,u_t),\ldots,
G_{r-t}(u_r;u_1,\ldots,u_t). \tag{1}
\]
이들은 \(\mathfrak a_p\)의 다항식
\[
G^*_1(u_{t+1};u_1,\ldots,u_t),\ldots,
G^*_{r-t}(u_r;u_1,\ldots,u_t) \tag{2}
\]
로 넘어가며, 원시성이 있으므로 사라지지 않는다. 그런데 각 \(G^*_\lambda\)에는 \(u_{t+\lambda}\)가 실제로 나타나야 한다. 그렇지 않다면 \(p\)의 선택과 모순되게 \(u_1,\ldots,u_t\)의 \(\mathfrak a_p\)에 대한 잉여류들 사이에 대수적 종속관계가 생겼을 것이다. 이로써 \(\mathfrak a_p\)의 초월 차수가 \(t\)임을 보였다. 따라서 \(\mathfrak a_p=\mathfrak c_p\)이고 \(\mathfrak a_p\)가 준동형을 규정하므로 모든 것이 증명되었다.

\bigskip
\paragraph{§ 4. 대수적 정수 계수로의 확장.}

서론 끝에서 이미 말했듯이, 유리정수 대신 한 유한 대수적 수체 \(\mathbf{K}\)의 대수적 정수들을 기초로 삼고 소수 \(p\) 대신 이 수체의 소아이디얼 \(\mathfrak p\)를 기초로 삼아도 모든 고찰은 조금만 고치면 그대로 성립한다. § 1은 완전히 그대로 유지되고, § 2에서 정리를 정식화하기까지의 고찰도 마찬가지다. 다만 증명(§ 3, 2)과 § 2, 3의 따름정리에는 몇 가지를 덧붙여야 한다.

\textbf{1.}
§ 3, 2에서는 다항식 \(G_1(u),\ldots,G_{r-t}(u)\)을 나누는 유한개의 소아이디얼을 예외소아이디얼에 더해야 한다[§ 3, 2의 (1)]. 이제는 이 다항식들을 원시적이라고 가정할 수 없기 때문이다. 그러면 다항식 \(G^*_1(u),\ldots,G^*_{r-t}(u)\)이 사라지지 않음이 보장된다[§ 3, 2의 (2)]. 이는 조건 III를 강화한 것이다. 절대기약성에 관한 정리는 바뀌지 않으므로 조건 II도 그대로이고, 물론 조건 I도 그대로이다.

\textbf{2.}
최종결과인 § 2, 3의 따름정리는 그곳의 마지막 서술에 맞추어 다음처럼 정식화한다.

\textbf{따름정리.}
\emph{\(K\)의 대수적 정수 \(\alpha\)가 유한개의 예외소아이디얼 가운데 어느 것으로도 나누어지지 않는다면, \(\alpha\cdot g(x)\in\mathfrak I\), \(g(x)\in\mathfrak Z\)에서 언제나 \(g(x)\in\mathfrak I\)가 따라온다.}

\(K\)의 대수적 정수 \(\lambda_1,\ldots,\lambda_s\)를 각각 하나의 예외소아이디얼로는 나누어지지만 그 제곱이나 다른 예외소아이디얼들로는 나누어지지 않도록 택하자. 그러면 분모 \(\alpha\)로는 이 \(\lambda\)들의 거듭제곱곱만으로 충분하다.

\emph{증명.}
\[
(\alpha)=\mathfrak p_1^{\varrho_1}\ldots\mathfrak p_g^{\varrho_g}
\]
를 \(\alpha\)의 소아이디얼분해라 하고, \(\mathfrak p_1,\ldots,\mathfrak p_g\)가 예외소아이디얼들과 다르다고 하자. 따라서 각 \(\mathfrak p_i\)에 대하여
\[
\mathfrak I_{\mathfrak p_i}\cong
\mathfrak o_{\mathfrak p_i}/\mathfrak a_{\mathfrak p_i}.
\]
\(K\)의 모든 대수적 정수의 환에서 출발하여, \((\alpha)\)와 서로소인 대수적 정수들, 곧 \(\mathfrak p_1,\ldots,\mathfrak p_g\) 어느 것으로도 나누어지지 않는 대수적 정수들만을 분모로 모두 첨가해 얻는 국소화환 \(\mathfrak R\)로 넘어가자. \(\mathfrak R\)에서도 \(\mathfrak p_1,\ldots,\mathfrak p_g\)는 소아이디얼로 남지만 이제 주아이디얼이 되고, 다른 모든 소아이디얼은 단위아이디얼과 같아진다. \((\alpha)\)의 분해는 그대로 유지된다\textsuperscript{11)}. 따라서 \(\mathfrak R\)을 계수영역으로 삼아도 따름정리와 그 증명은 § 2, 3의 형태 그대로 성립한다.

이는 다음을 뜻한다. \(\alpha\cdot g(x)\in\mathfrak I\)이면 \(g(x)=F(f_i(x))\)이고, 여기서 다항식 \(F(u)\)의 계수들은 \(\mathfrak R\)에 속한다. 이 계수들의 공통분모를 \(\delta\)라 하자. 그러면 \(\delta\)는 \(\alpha\)와 서로소이다. 따라서 같은 사실을 다음처럼 쓸 수 있다. \(\alpha\cdot g(x)\in\mathfrak I\)이면 \(\delta\cdot g(x)\in\mathfrak I\)인 \(\alpha\)와 서로소인 \(\delta\)가 존재한다. 그러므로 \((\delta,\alpha)\)는 단위아이디얼이고, 결국 \(\alpha\cdot g(x)\in\mathfrak I\)에서 \(g(x)\in\mathfrak I\)가 따라온다. 이로써 위 따름정리의 첫째 주장을 증명하였다.

둘째 주장을 증명하기 위하여 \(\mathfrak R^*\)를 유한개의 예외소아이디얼로 정해지는 \(\mathbf{K}\) 안의 국소화환이라 하고, \(\lambda_1,\ldots,\lambda_s\)를 \(\mathfrak R^*\)에서 이 소아이디얼들의 생성원이라 하자. 이 생성원들은 대수적 정수로 택할 수 있다. \(\mathfrak R^*\)에서 \(\alpha\)는 \(\mathfrak R^*\)의 한 단위원을 제외하면 \(\lambda\)들의 거듭제곱곱이다. 다시 대수적 정수들로 돌아가면
\[
\gamma\alpha=\beta\lambda_1^{\sigma_1}\ldots\lambda_s^{\sigma_s}
\]
를 얻는다. 여기서 \(\gamma,\beta\)는 어느 예외소아이디얼로도 나누어지지 않는다. \(\alpha\cdot g(x)\in\mathfrak I\)이면
\[
\beta\cdot\lambda_1^{\sigma_1}\ldots\lambda_s^{\sigma_s}g(x)\in\mathfrak I
\]
이고, 따라서 따름정리의 첫째 주장에 의하여
\[
\lambda_1^{\sigma_1}\ldots\lambda_s^{\sigma_s}g(x)\in\mathfrak I
\]
이다.

\begin{center}
\rule{4cm}{0.4pt}\\[0.6em]
\textbf{주}
\end{center}

{\small
\setlength{\parindent}{0pt}

\noindent\textsuperscript{1)} D. Hilbert, \emph{Mathematische Probleme}, Gött. Nachr. 1900, 253--297쪽, 문제 14.

\noindent\textsuperscript{2)} E. Hecke, \emph{Über die Konstruktion relativ-Abelscher Zahlkörper durch Modulfunktionen von zwei Variabeln}, Math. Ann. 74 (1913), 465--510쪽, § 2. 헤케의 경우에는 유한 생성 정역 \(\mathfrak I\)이 두 변수에 관한 거듭제곱급수의 몫 세 개로 생성되며, 이들 사이에는 하나의 대수적 관계가 있다.

\noindent\textsuperscript{3)} 힐베르트의 앞의 논문 또는 E. Noether, \emph{Die Endlichkeit des Systems der ganzzahligen Invarianten binärer Formen}, Gött. Nachr. 1919, 138--156쪽, § 1, IV의 더 상세한 서술을 보라.

\noindent\textsuperscript{4)} 따라서 예를 들어 한 기본형식계의 정수 계수 (사영)불변식의 경우, 통상적 의미의 한 완전불변식계로 그것들을 나타낼 수 있으며 그 분모에는 서로 다른 소수가 \emph{유한개}만 나타남을 보인 것이다. 이는 통상적인 \(\varOmega\)-과정의 방법에서는 따라오지 않는 결과이다. 그러나 정수 계수 유한성 자체에 대해서는 아무것도 말하지 않는다. 이 유한성은 지금까지 이원 기본형식의 경우에만 알려져 있다. 3)에 인용한 글을 보라.

\noindent\textsuperscript{5)} 한 계의 \glqq{}초월 차수\grqq{}란 그 계가 대수적으로 종속되는 대수적으로 독립인 함수들의 불변적인 개수를 뜻한다. 대수적으로 독립인 함수들로 이루어진 계를 기약계라 한다. 한 소아이디얼의 초월 차수(차원)는 그 잉여류체의 초월 차수로 이해한다. 소아이디얼은 같은 초월 차수를 가진 진소약수를 갖지 않는다. 예컨대 B. L. van der Waerden, \emph{Zur Nullstellentheorie der Polynomideale}, Math. Ann. 96 (1926), 183--208쪽을 보라.

\noindent\textsuperscript{6)} E. Noether, \emph{Eliminationstheorie und allgemeine Idealtheorie}, Math. Ann. 90 (1923), 229--261쪽, § 7. 이 정리는 소거론\textsuperscript{*)}을 거쳐, 절대기약다항식은 유한개의 소수를 법으로 할 때에만 그 성질을 잃는다는 더 간단한 정리에서 따라온다. A. Ostrowski, \emph{Zur arithmetischen Theorie der algebraischen Größen}, Gött. Nachr. 1919, 279--298쪽, 특히 296쪽; 더 간단한 증명은 E. Noether, \emph{Ein algebraisches Kriterium für absolute Irreduzibilität}, Math. Ann. 85 (1922), 26--33쪽, 7번을 보라. \(\mathfrak I\)의 정수성 기저가 그 초월 차수보다 정확히 함수 하나만 더 포함하여 \(\mathfrak a\)가 주아이디얼이 되는 경우, 예컨대 헤케가 다룬 경우(주 2)에는 소거론을 피하고 이 더 간단한 정리만 사용해도 충분하다. 곳곳에서 소수 대신 한 유한 대수적 수체의 소아이디얼을 놓을 수 있다. 이 정리는 \(\mathfrak a\)가 영아이디얼인 경우에도 분명히 성립한다.

\noindent\textsuperscript{7)} 원래 서술에서는 계수체 \(\mathrm{P}\)를 완전체라고 가정하였다. \(\mathrm{P}\)에서 \(\varOmega\)로 넘어감으로써 임의의 체, 완전체가 아닌 체까지 확장하는 것은 B. L. van der Waerden에게서 배웠다. 완전체 \(P\)에서도 표수 \(p\)에서는 역이 성립하지 않음을 \(f=x^p,\ df/dx=0\)이라는 예가 보여 준다.

\medskip
\hrule width 2.7cm
\smallskip
\noindent\textsuperscript{*)} \(\mathfrak a\)의 계수 표수가 0이므로, 정확히 말하면 이 논문의 보조정리 V와 무관하게 그러하다. 실제로 그 보조정리의 증명은 옳지 않다. 거기서 제시한 치환을 하면 모든 것이 항등적으로 사라질 수 있다. 올바른 증명은 B. L. van der Waerden, \emph{Eine Verallgemeinerung des Bézoutschen Theorems}, Math. Ann. 99 (1928), 497--541쪽, 주 37에 비로소 나온다.

같은 논문의 257쪽에도 쉽게 알 수 있는 수정이 하나 더 필요하다. 정리 XIV 역시 § 7에서는 사용되지 않는다. 거기에는 \((\mathfrak R)\)의 모든 원소가 \((x_i)\) 안의 \(t_{\mu,\nu}\)를 \((\mathfrak R')\)의 원소들로 특수화하여 생긴다고 적혀 있다. 실제로는 더 일반적인 선형형식
\[
\sum s_{\lambda_1\ldots\lambda_n}
(y_1)^{\lambda_1}\cdots(y_n)^{\lambda_n}
\]
에서 특수화하여 생기지만, 그 지수는 \((x_i)\)의 지수와 같다.

\noindent\textsuperscript{8)} 표수 0에 관한 이 함수행렬식 정리의 순수대수적 증명은 B. L. van der Waerden, \emph{Algebraische Theorie der Differentiation}, Nieuw Archief v. Wiskunde (2) 15 (1926), 111--120쪽을 보라.

\noindent\textsuperscript{9)} 이 부가가정은 헤케가 다룬 경우(주 2)처럼 해석함수들을 다루고 \(f_1(x),\ldots,f_t(x)\)가 해석적으로 독립이며 \(f_{t+1}(x),\ldots,f_r(x)\)가 이들에 대하여 대수적으로 종속될 때 충족된다. 그러나 이 가정은 형식적으로 정의된 거듭제곱급수의 경우에도 명확한 뜻을 갖도록 의도적으로 함수행렬의 랭크라는 형식으로 서술하였다. 이로써 모든 고찰은 순수하게 형식대수적인 범위에 머물며, 해석함수에 관한 정리는 특수한 경우에 그 가정이 충족됨을 보이는 데에만 필요하다.

\noindent\textsuperscript{10)} 예컨대 B. L. van der Waerden, \emph{Zur Nullstellentheorie der Polynomideale}, Math. Ann. 96 (1926), 183--208쪽, § 3을 보라.

\noindent\textsuperscript{11)} 예컨대 H. Grell, \emph{Zur Theorie der Ordnungen in algebraischen Zahl- und Funktionenkörpern}, Math. Ann. 97 (1927), 524--558쪽, § 2, 1을 보라.
}

\bigskip
\begin{center}
\rule{4cm}{0.4pt}\\[0.8em]
{\Large\bfseries 정수적 함수의 극대영역에 관하여}\par
\vspace{0.7em}
{\bfseries 에미 뇌터 (괴팅겐)}\par
\vspace{0.5em}
\emph{러시아어 요약의 한국어 번역}
\end{center}

영인수가 없는 각 환(정역) \(\mathfrak I\)은 같은 유형의 한 극대영역 \(\mathfrak M(\mathfrak I)\)을 정한다. 이 극대영역은 적어도 하나의 \(a\ne0\)에 대하여 \(a\cdot g(x)\in\mathfrak I\)가 되는 모든 함수 \(g(x)\)로 이루어진다.

이 논문에서는 \(\mathfrak I\)이 유한 생성되었다는 가정 아래 이 극대영역 \(\mathfrak M(\mathfrak I)\)의 \glqq{}분모\grqq{} \(a\)를 연구한다. 모든 \(a\)를 통틀어 인수로 나타나는 소수가 유한개뿐이도록 \(a\)들을 언제나 택할 수 있음을 증명한다. 다만 \(g(x)\)들이 거듭제곱급수 또는 그러한 급수들의 몫인 경우에는, \(\mathfrak I\)을 정하는 함수 \(f_1(x),\ldots,f_r(x)\) 사이에 오직 대수적 관계만 있을 수 있다고 가정한다.

증명은 문제를 일반 아이디얼론의 몇 가지 물음으로 환원하는 데에 바탕을 둔다. 이 환원에서는 인수 \(p\)가 분모 \(a\)에 나타나는 것을 함수 \(f_i\)들 사이의 \(p\)를 법으로 하는 관계로 해석한다.

\begin{center}
(\emph{Rec. Math.} XXXVI:1, 1929.)
\end{center}

\clearpage
\setcounter{footnote}{0}
\phantomsection
\addcontentsline{toc}{section}{36. 아이디얼 미분과 디퍼런트}
%
%
%

%
%
\section*{36. 아이디얼 미분과 디퍼런트}

\begin{center}
\emph{J. Ber. d. DMV 39 (1930), S. 17}
\end{center}

2. E. Noether, 괴팅겐: 아이디얼 미분과 디퍼런트.

대수적 수체의 디퍼런트(different)가 그것을 정의하는 \emph{아이디얼}의 \emph{미분몫}으로 어떻게 해석될 수 있는지를 보인다. 이 정의가 통상적인 정의와 일치한다는 사실은 그 자체로 흥미로운 구조 정리들로부터 따른다. 이 정리들은 유사한 의미에서 라그랑주 보간 공식의 일반화를 이룬다.

상세한 논의는 \emph{Mathematische Annalen}에 게재될 예정이다.

\clearpage
\setcounter{footnote}{0}
\phantomsection
\addcontentsline{toc}{section}{37. 고차 분기가 없는 체의 정규기저}
%
%
%
%
%
%

\section*{37. 고차 분기가 없는 체의 정규기저}

\begin{center}
\emph{Journal f. d. reine u. angew. Math. 167 (1932), 147--152쪽}
\end{center}

\begin{center}
괴팅겐의 \emph{에미 뇌터}.
\end{center}

갈루아 수체 \(K/k\)의 정규기저란 \(k\)의 정수환 \(\frako\)에 대한 \(K\)의 정수환 \(\frakO\)의 기저로서 켤레원소들로 이루어진 것을 말한다. 그러한 기저가 존재하기 위한 필요조건은 \(K/k\)에 고차 분기체가 나타나지 않는다는 것임이 알려져 있다.\footnote{A. Speiser, \emph{Gruppendeterminante und Körperdiskriminante}, § 6, Math. Ann. 77 (1916), 546--562쪽을 보라.} 이는 한 소에 국한하여, 곧 \(k\)의 \(\frakp\)-진 확대 \(k_\frakp\)와 이에 대응하는 \(K\)의 확대 \(K_\frakp\)로 넘어가도 성립한다. 이하에서는 그 역이 성립함을 보인다.

\emph{\(\frakp\)가 \(K/k\)의 차수를 나누지 않는 모든 소 \(\frakp\)에서 정규기저가 존재한다. 특히 \(K/k\)에 고차 분기가 없으면 모든 분기소에서 정규기저가 존재하며, 그곳에서 판별식은 군행렬식의 제곱이 된다.}

이 마지막 명제에는 다음 설명을 덧붙여야 한다. E. Artin의 추측과 달리 그의 도체들\footnote{E. Artin, \emph{Über die gruppentheoretische Struktur der Diskriminante}, J. f. Math. 164 (1931), 1--11쪽.}로 넘어가려면 더 많은 고찰이 필요하다. 군행렬식을 기약지표에 따라 분해하면 일반화된 근수들이 생기고, 이를 알맞게 묶으면 지표들을 첨가한 기초체에서의 분해를 얻는다. 가장 간단한 경우들에는 근수들의 아이디얼론적 분해를 이용하여 이 새 인자들을 Artin 도체와 동일시할 수 있다.\footnote{[1931. 10. 6.] 일반적인 경우에도 아이디얼론적 고찰이 더 필요함은 M. Deuring이 알려 준 다음 예에서 드러난다. 이 예는 임의로 변형할 수 있다. 어떤 비극대 오더의 판별식은 군행렬식의 제곱으로 나타나지만 켤레 지표들에는 서로 다른 아이디얼인자들이 대응한다. \(k\)를 5차 단위근체, \(K=k(\sqrt[5]{2})\)라 하자. 오더 \([1,2\sqrt[5]{2},\sqrt[5]{2^2},\sqrt[5]{2^3},\sqrt[5]{2^4}]\)를 생각하되, 이 글의 고찰에 따라 이 오더가 정규기저를 갖는 소 \((2)\)에서 생각한다. 지표들에 따른 대응 군행렬식의 분해에서는 위의 기저원소들이 바로 인자들로 나온다. 따라서 \glqq{}도체\grqq{}들은 \(1\)을 제외하면
\[
2\sqrt[5]{2\cdot}\sqrt[5]{2^4},\quad
\sqrt[5]{2^2}\sqrt[5]{2^3},\quad
\sqrt[5]{2^3}\sqrt[5]{2^2},\quad
\sqrt[5]{2^4\cdot}\,2\sqrt[5]{2},
\]
곧 \(4,2,2,4\)가 되어 켤레 지표들에 대해 서로 다르다.}
더 나아가 정규기저가 존재하지 않는 일반적인 경우에도 \glqq{}갈루아 가군에 따른 합성열\grqq{}을 이용하면 일반화된 근수들로 언제나 분해할 수 있음을 말해 두고자 한다. 여기에서도 위와 마찬가지로 확대된 기초체에서의 분해가 생기며, 다시 아이디얼론적 문제가 시작된다.

위 가정 아래에서 정규기저의 존재를 증명하는 근거는 다음과 같다. \(\frakO/\frako\)를 갈루아 가군, 곧 \(\frako\)-가군으로 보되 갈루아군의 치환들을 연산자영역으로 두면, 이는 \(\frako\)로 확대된 정수군환의 한쪽 아이디얼과 연산자 동형사상이 된다. 이 정수군환은 모든 보통 분기소에서 극대 오더이다. 그런데 \(\frakp\)-진 반단순계의 극대 오더에 속하는 아이디얼들은 주아이디얼이다.\footnote{H. Hasse, \emph{Über \(\frakp\)-adische Schiefkörper und ihre Bedeutung für die Arithmetik hyperkomplexer Zahlsysteme}, 정리 46, 47, Math. Ann. 104 (1931), 495--534쪽을 보라. 가환계에서는 \(k\)의 \(\frakp\)에 대한 국소화환으로 넘어가기만 해도 주아이디얼 성질이 성립한다. 그러므로 아벨군과 아벨체를 다룰 때에는 이 더 약한 확대를 써서 소를 정의해도 된다.} 따라서 여기에서는 한 원소에 군 또는 군환의 치환들을 적용하여 아이디얼을 얻는다. 갈루아 가군으로 되돌리면 이것이 바로 정규기저의 존재를 뜻한다. 고차 분기소에서는 정수군환이 비극대 오더가 되고 \(\frakO/\frako\)에 대응하는 아이디얼은 비주아이디얼이 된다.

상수체의 표수가 \(K/k\)의 차수를 나누지 않는 한, 이 모든 고찰은 한 변수의 함수체에도 그대로 성립한다.

\subsection*{§1. \(\frakp\)-진으로 확대한 정수군환}

\(\frako\)와 \(\frako_\frakp\)는 각각 대수적 수체 \(k\)와 그 \(\frakp\)-진 확대 \(k_\frakp\)의 정수환이며, \(\frakp\)는 \(k\)의 소아이디얼이라고 하자. 원소 수가 \(n\)인 군 \(\Gg\)의 유리수계수 군환을 \((\Gg)\), 정수군환을 \([\Gg]\)로 쓴다. \((\Gg)_k\)와 \((\Gg)_{k_\frakp}\)는 각각 \(k\)와 \(k_\frakp\)에 계수를 둔 군환을 뜻하고, \([\Gg]_\frako\)와 \([\Gg]_{\frako_\frakp}\)는 정수군환을 각각 \(\frako\)와 \(\frako_\frakp\)에 계수를 둔 군환으로 확대한 것을 뜻한다.

따라서 \((\Gg)_k\)와 \((\Gg)_{k_\frakp}\)는 라디칼이 없는 계, 곧 반단순계이고, \([\Gg]_\frako\)와 \([\Gg]_{\frako_\frakp}\)는 각각 그 안의 오더이다.

\paragraph{정리 1.}
\emph{소아이디얼 \(\frakp\)가 \(\Gg\)의 원소 수 \(n\)을 나누지 않으면 \([\Gg]_{\frako_\frakp}\)는 \((\Gg)_{k_\frakp}\)의 극대 오더이다. \(\frakp\)가 \(n\)을 나누면 \([\Gg]_{\frako_\frakp}\)는 극대가 아니다.}

정수군환의 판별식, 따라서 \([\Gg]_\frako\)와 \([\Gg]_{\frako_\frakp}\)의 판별식은 원소 수 \(n\)의 한 거듭제곱과 같다.\footnote{예를 들어 E. Noether, \emph{Hyperkomplexe Größen und Darstellungstheorie}, § 26, Math. Ztschr. 30 (1929), 641--692쪽을 보라.} 그러므로 \(\frakp\)가 \(n\)을 나누지 않으면 \([\Gg]_{\frako_\frakp}\)의 판별식은 단위원이다. 이는 \(\frako_\frakp\)를 고정한 채 \([\Gg]_{\frako_\frakp}\)를 더 큰 오더로 확대할 수 없다는 뜻이다. 그러한 확대가 존재하면 판별식에서 비단위원인 제곱인자가 떨어져 나오기 때문이다. 그런데 \(\frako_\frakp\)는 이미 \(k_\frakp\)의 정수환, 따라서 그 안의 극대 오더라고 가정하였다. 이로써 이하에서 유일하게 쓰일 정리 1의 첫 부분이 증명되었다.

반대로 \(\frakp\)가 \(n\)을 나누면 자명표현에 대응하는 직합
\[
\mathfrak C=
E^{(1)}[\Gg]_{\frako_\frakp}
+(1-E^{(1)})[\Gg]_{\frako_\frakp},
\qquad
E^{(1)}=\frac1n\sum S
\]
은 \([\Gg]_{\frako_\frakp}\)의 진확대이다. 실제로 \(E^{(1)}=\frac1n\sum S\)이므로 \(\mathfrak C\)는 진확대이며, \(E^{(1)}\)과 \(S\in\Gg\)에 대한 \((1-E^{(1)})S\)를 종속인 생성원들로 갖는 오더이다. 여기서 \(E^{(1)}S=E^{(1)}\)이다.

\paragraph{정리 2.}
\emph{\(\frakp\)가 \(n\)을 나누지 않으면 \([\Gg]_{\frako_\frakp}\)에 관한 모든 정수 또는 분수 아이디얼은 주아이디얼이다.}

정리 1에 따라 \([\Gg]_{\frako_\frakp}\)는 \(\frakp\)-진 반단순계의 극대 오더이므로 그 아이디얼들은 주아이디얼이다.\footnote{Hasse, 앞의 논문. 그곳에서는 단순계에 대해서만 주아이디얼 성질을 증명하지만 알려진 추론에 따라 반단순계에도 성립한다. 극대 오더의 성분들은 다시 극대이고, 따라서 생각하는 아이디얼의 성분들은 이를테면 기저 \(a_i\)를 가진 주아이디얼이다. 그러면 \(a=\sum a_i\)가 본래 아이디얼의 기저가 된다. 아벨군에 관해서는 주 3)도 보라.}

\subsection*{§2. 갈루아 가군, 연산자 동형사상, 정규기저}

\(K/k\)가 갈루아이고 그 군이 \(\Gg\)라고 하자. \(z^S\)는 \(S\in\Gg\)라는 치환으로 \(z\)에서 얻는 \(K\)의 원소를 뜻한다. \(\Gg\)의 치환들은 \(K/k\)를 \(k\)-가군, 곧 \(k\)를 연산자영역으로 갖는 가법 아벨군으로 볼 때 연산자 자기동형사상들을 만든다. 따라서 다음과 같이 정하면 \(K/k\)를 \((\Gg)_k\)에 관한 가군으로 볼 수 있다.
\[
z\Bigl(\sum_i S_i c_i\Bigr)=\sum_i z^{S_i}c_i,
\qquad z\in K,\quad S_i\in\Gg,\quad c_i\in k.
\]

\paragraph{정의.}
\emph{\(K/k\) 안의 모든 \((\Gg)_k\)-가군을 유리 갈루아 가군이라 한다. 마찬가지로 \(K/k\) 안의 \([\Gg]_\frako\)-가군들을 정수 갈루아 가군이라 한다. 특히 \(\frakO/\frako\)는 정수 갈루아 가군이다.}

\paragraph{정리 3.}
\emph{유리 갈루아 가군으로서 \(K/k\)는 \((\Gg)_k\)와 연산자 동형사상이다.}

이는 \(K/k\)에 언제나 체의 정규기저, 곧 켤레원소 \(z^S\)들로 이루어진 \(k\)-기저가 존재한다는 사실에서 따라온다.\footnote{\(a_1,\ldots,a_n\)이 \(K/k\)의 한 기저이고 \(u_i\)가 부정원들이면, \(S\)가 \(\Gg\)의 원소들을 돌 때 \(\sum u_i a_i^S\)들은 선형독립이다. 따라서 선형독립성을 보존하도록 \(u_i\)들을 \(k\)의 원소들로 특수화할 수도 있다.} 대응
\[
S\to z^S
\sum S_i c_i\to\sum z^{S_i}c_i
\quad(c_i\in k),\qquad
ST\to z^{ST}
\]
은 연산자 준동형사상을 주고, 양변의 \(k\)-랭크가 같으므로 동형사상이 된다.

\paragraph{정리 4.}
\emph{정수 갈루아 가군으로서 \(\frakO/\frako\)는 \((\Gg)_k\) 안의 최고랭크 \(n\)인 한 \([\Gg]_\frako\)-가군과 연산자 동형사상이다. 마찬가지로 \(\frakO_\frakp/\frako_\frakp\)는 \((\Gg)_{k_\frakp}\) 안의 랭크 \(n\)인 한 \([\Gg]_{\frako_\frakp}\)-가군과 연산자 동형사상이다.}

정리 3의 \((\Gg)_k\)-동형사상 \(K/k\to(\Gg)_k\)는 \([\Gg]_\frako\)-가군 \(\frakO/\frako\)에 \((\Gg)_k\) 안의 랭크 \(n\)인 한 \([\Gg]_\frako\)-가군을 대응시킨다. 또 정리 3의 대응에서 \(c_i\)들이 \(k_\frakp\)의 모든 원소를 돌게 하면, \((\Gg)_k\)-동형사상을 \(K_\frakp/k_\frakp\)와 \((\Gg)_{k_\frakp}\) 사이의 \((\Gg)_{k_\frakp}\)-동형사상으로 확대할 수 있다. 이 동형사상은 \(\frakO_\frakp/\frako_\frakp\)에 대한 \([\Gg]_{\frako_\frakp}\)-동형사상을 유도한다. 여기서 \(K_\frakp/k_\frakp\)는 \(k_\frakp\) 위의 초복소계로 보아야 한다. 이는 \(K/k\)에서 계수영역을 확대하여 생기며, 일반적으로 영인자를 가진 계, 더 정확히는 서로 동형인 체들의 합인 반단순계가 된다.

\paragraph{정리 5.}
\emph{\(K/k\)의 차수 \(n\)을 나누지 않는 모든 소 \(\frakp\)에서 \(\frakO_\frakp/\frako_\frakp\)는 정규기저를 갖는다.}\footnote{아벨체의 경우에는 주 3)과 5)에 따라 국소화환으로도 소를 정의할 수 있다.}

정리 1에 따라 \([\Gg]_{\frako_\frakp}\)는 그곳에서 극대 오더이다. 따라서 \((\Gg)_{k_\frakp}\) 안에서 랭크 \(n\)인 \([\Gg]_{\frako_\frakp}\)-가군들은 정수 또는 분수 아이디얼이고, 정리 2에 따라 주아이디얼이다. 그러므로 \(\frakO_\frakp/\frako_\frakp\)에 대응하는 아이디얼이
\[
W[\Gg]_{\frako_\frakp}
\]
라면, 이는 \(\frako_\frakp\)-기저 \(WS_i\)를 갖는다. \(w\)가 \(\frakO_\frakp/\frako_\frakp\)에서 \(W\)에 대응하는 원소라 할 때, 연산자 동형사상은 \(w^{S_i}\)들이 \(\frakO_\frakp/\frako_\frakp\)의 \(\frako_\frakp\)-기저임을 말한다. 따라서 \(n\)개의 켤레원소 \(w^{S_i}\)가 \(\frakO_\frakp/\frako_\frakp\)의 정규기저를 이룬다.

\(\frakp\)가 \(n\)을 나누면 정규기저가 존재할 수 없으므로 \(\frakO_\frakp/\frako_\frakp\)에 대응하는 가군도 주가군일 수 없다.

\paragraph{정리 3에 대한 추가 설명.}
정리 3에서 증명한 \((\Gg)_k\)와 \(K/k\)의 연산자 동형사상\footnote{그 증명은 \(K\)가 \(k\) 위의 제1종 갈루아체이고 \(k\)가 무한히 많은 원소를 가진다는 것만을 가정한다. 이 마지막 제한은 불필요하다. M. Deuring은 유한체에도 유효한 정리 3의 증명을 고안하였고, 그 증명에서 거꾸로 체의 정규기저 존재가 따라온다. 동시에 유한체와 무한체에 한결같이 적용되는 원시원소 존재의 증명도 얻어진다.}에서 클라인의 형식문제에 관한 Speiser의 결과들\footnote{앞의 논문. 갈루아 가군의 개념은 그곳에서 추상화한 것이다. \(k\)의 표수가 \(n\)을 나누면 갈루아 가군에 따른 합성열과 그 기약인자들을 다루게 된다.}이 따라오며, 다음과 같이 정식화할 수 있다.

\emph{\(K/k\)의 갈루아군 \(\Gg\)의 \(k\) 위 기약표현은 모두 \(K/k\) 안의 한 기약 유리 갈루아 가군이 생성한다. 절대기약표현은 모두 \(K_{\mathsf{Z}}/\mathsf{Z}\) 안의 한 갈루아 가군이 생성한다.}

여기서 \(\mathsf{Z}\)는 대응하는 군환 성분의 분해체로서 \(k\)를 포함한다. \(K_{\mathsf{Z}}/\mathsf{Z}\)는 \(\mathsf{Z}\) 위의 초복소계로 보며 \(K/k\)의 계수영역 확대에서 생긴다. 특히 \(Z\)를 \(k\) 위에서 \(K\)와 선형적으로 서로소가 되게 택할 수 있으면 \(K_{\mathsf{Z}}\)는 체로 남고 \(K_{\mathsf{Z}}/\mathsf{Z}\)는 부가확대가 된다. 이것이 Speiser가 다룬 경우이다.

명제 자체는 연산자 동형사상에 따라 \((\Gg)_k\) 또는 \((\Gg)_{\mathsf{Z}}\)에 대한 대응 명제에서 따라온다. 그곳에서는 \((\Gg)_k\) 또는 \((\Gg)_{\mathsf{Z}}\)에 관한 기약 아이디얼들이 표현을 생성한다.

\(v_1,\ldots,v_t\)가 그러한 갈루아 가군의 기저이고 \(S\to \bar S\)가 대응하는 표현이면
\[
\begin{pmatrix}\vdots\\ v_i^S\\ \vdots\end{pmatrix}
=\bar S\begin{pmatrix}\vdots\\ v_i\\ \vdots\end{pmatrix}.
\]
이것이 Klein과 Speiser의 정식화이다.

\subsection*{§3. 고차 분기가 없는 체에서 군행렬식으로서의 판별식}

판별식을 분해하려면 각 분기소에서의 분해를 살펴보는 것으로 충분함이 알려져 있다. 따라서 고차 분기가 없는 체에서는 정규기저를 갖는 소들만을 다루면 된다.

\paragraph{정리 6.}
\emph{고차 분기가 없는 갈루아체 \(K/k\)에서는 각 분기소에서 판별식이 갈루아군의 군행렬식의 제곱과 같아지며, 이때 부정원들은 정규기저로 바뀐다. 절대기약표현들에 따라 군행렬식은 일반화된 근수들로 분해되고, 판별식은 지표들을 첨가한 기초체의 아이디얼들로 분해된다. 이때 복소켤레 지표들에는 같은 아이디얼들이 대응한다.}

\(w^S\)를 정규기저로 택하면 판별식 \(\Delta\)는
\[
D=\bigl|w^{ST^{-1}}\bigr|,\qquad S,T\in\Gg
\]
의 제곱과 같다. 그런데 \(D\)는 부정원 \(x_S\)들을 \(w^S\)로 바꾼 군행렬식이다.

근수들로의 분해는 앞에서 인용한 Speiser의 추론에서 얻어진다. 절대기약표현들에 따라 군행렬식과 군행렬을 각각
\[
D=D_1^{f_1}\cdots D_t^{f_t},
\qquad
M=f_1M_1+\cdots+f_tM_t
\]
로 분해하자. \(S\to\bar S\)가 \(M_\lambda\)에 대응하는 표현이고 \(\bar S\)들이 확대된 기초체에 속하면
\[
M_\lambda=\sum_S w^S\bar S,
\qquad
M_\lambda^{T^{-1}}=\sum_S w^{ST^{-1}}\bar S
=\sum_R w^R\bar R\bar T=M_\lambda\bar T
\]
이고, 행렬식으로 넘어가면
\[
D_\lambda^{T^{-1}}=D_\lambda|\bar T|
=D_\lambda\varepsilon_T
\]
를 얻는다. 이로써 \(D_\lambda\)들이 일반화된 근수임을 알 수 있다. \(\varepsilon_T\)는 \(\bar T\)의 행렬식이므로 단위근이다. 이는 \(\Gg\)를 교환자부분군으로 나눈 몫군의 기약표현이다. 순환군의 경우 \(D_\lambda\)들은 단순히 라그랑주 근수
\[
\sum w^S\chi_\lambda(S)
\]
가 된다.

일반적으로 \(D_\lambda\)는 \(K_\frakp/k_\frakp\)에서 계수영역 \(k_\frakp\)에 대응하는 지표를 첨가하여 얻는 초복소계에 속하며, 더 정확히는 이 계의 \emph{정수적} 원소이다. 부정원 \(x_S\)들로 만든 행렬식의 계수들이 지표체의 대수적 정수이고 \(w^S\)들이 \(K_\frakp\)의 대수적 정수이기 때문이다.

군행렬식의 분해에서 판별식의 분해로 넘어가기 위해 매번 한 표현을 그 수반표현과 묶는다. \(\lambda,\bar\lambda\)가 서로 수반인 표현이면 동시에
\[
D_\lambda^{T^{-1}}=D_\lambda\varepsilon_T,
\qquad
D_{\bar\lambda}^{T^{-1}}=D_{\bar\lambda}\varepsilon_T^{-1}
\]
이다. 또 부정원 \(x_S\)들에 대하여 만든 \(\lambda,\bar\lambda\)의 행렬식들은 서로 복소켤레이고, 일반적으로 켤레 지표들에는 부정원 \(x_S\)들에 관한 켤레 행렬식들이 대응한다.

\(\Delta_\lambda=D_\lambda D_{\bar\lambda}\)로 두면 \(\lambda\)번째 지표체의 실부분체만 첨가하면 되고
\[
\Delta_\lambda^T=\Delta_\lambda
\qquad(\Gg\text{의 모든 }T)
\]
가 성립한다. 따라서\footnote{여기서는 한 체의 계수영역을 초복소적으로 확대하므로 통상적인 갈루아 이론의 추론을 수정해야 한다. \(\mathsf{P}\)를 생각하는 \(k_\frakp\)의 확대체라 하고
\[
(K_\frakp)_{\mathsf{P}}=(K_\frakp)_{\mathsf{P}} e^{S_1}+\cdots+(K_\frakp)_{\mathsf{P}} e^{S_r}
\]
를 체들로의 직합분해라 하자. 그러면 \( (K_\frakp)_{\mathsf{P}} e^{S_i}/\mathsf{P}e^{S_i}\)는 갈루아가 되고, 그 군은 \(\frakp\)의 한 소인자의 분해군의 부분군이다. 또한 \(e^{S_i}\)들은 잉여류들에 관하여 켤레이다. \(\Delta_\lambda^T=\Delta_\lambda\)에서 먼저 \(\Delta_\lambda\)의 \(i\)번째 성분이 \(\mathsf{P}e^{S_i}\)에 속함이 따라오므로, 이는 \(P\)의 \(\gamma_i\)를 써서 \(\gamma_i e^{S_i}\)가 된다. 이어 \(e^{S_i}\)들의 켤레성에서 모든 \(\gamma_i\)가 같음이 따라온다. 따라서 실제로 \(\Delta_\lambda=\gamma\sum e^{S_i}=\gamma\)는 \(\mathsf{P}\)에 속한다.} \(\Delta_\lambda\)는 \(\lambda\)번째 지표체의 실부분체를 첨가한 기초체에 속한다.

\emph{따라서 판별식 분해}
\[
\Delta=\Delta_1^{f_1}\cdots\Delta_t^{f_t}
\]
\emph{는 지표체들의 실부분체를 첨가한 기초체에서의 분해이며, 복소켤레 지표들에는 같은 인자들이 대응한다. 나머지 켤레 지표들에 대해서는 더 이상의 결론을 곧바로 내릴 수 없다[2a) 참조].}

이로써 정리 6의 모든 부분이 증명되었다. \(\Delta_\lambda\)들이 생성하는 아이디얼들이 실제로 이미 기초체 \(k\)에 속하고 켤레 지표들에 대하여 같아짐을 보일 수 있다면, 합성지표들에 대응하는 \(\Delta_\lambda\)들의 거듭제곱곱을 대응시키는 것만으로 이 아이디얼들은 Artin 도체들과 일치한다. 실제로 함수방정식이 성립하고, 정규기저에서 직접 따라오는 사실, 곧 부분군의 자명표현에는 대응하는 부분체의 판별식이 대응한다는 사실도 성립한다(앞서 인용한 Speiser 참조). 각 소에서 일치하면 전체 아이디얼도 일치한다. 유리수체 위에서 기약인 지표들로 제한하면 위 사실들만으로 일치가 따라온다. 절대기약지표에 대해서는 가장 간단한 경우에 그 일치를 직접 확인하겠다.

\paragraph{정리 7.}
\emph{\(k\)가 유리수체이고 \(K/k\)가 판별식과 서로소인 소수차수 \(l\)의 순환확대, 따라서 고차 분기가 없는 확대이면, 비자명표현들에 대하여 \(\Delta_\lambda\)들이 생성하는 아이디얼들은 서로 같고 \(K/k\)의 도체와 일치한다.}

이는 근수들의 알려진 분해에서 따라온다.\footnote{Hilbert, Zahlbericht, § 108. 정리 5에 따라 정규기저의 존재가 확정되어 있으므로 --- 주 7)에 따라 여기서는 소를 국소화환으로 정의해도 된다 --- 절대순환체가 원분체라는 Zahlbericht의 사실을 더는 끌어올 필요가 없다.} 실제로 \(K/k\)의 한 분기소 \(p\)에서 \(l\)차 단위근체 \(k_\varepsilon\)와 \(K_\varepsilon\)에 대하여
\[
(p)=\frakp_1\cdots\frakp_{l-1},
\qquad
\frakp_i=\mathfrak P_i^{\,l}
\]
가 성립한다. 여기서 \(\frakp\)는 \(k_\varepsilon\)에, \(\mathfrak P\)는 \(K_\varepsilon\)에 속한다. \(k_\varepsilon\)는 \(K\)와 선형적으로 서로소이므로 \(K_\varepsilon\)는 체로 남고, 주 7)에 따라 여기서는 \(\frakp\)-진 확대가 필요 없다. 또한 근수들에 대하여
\[
(\Omega_\lambda)=
\mathfrak P_1^{r_1}\cdots\mathfrak P_{l-1}^{r_{l-1}},
\qquad
(\Omega_{\bar\lambda})=
\mathfrak P_1^{l-r_1}\cdots\mathfrak P_{l-1}^{l-r_{l-1}}
\]
가 성립하며 \(r_1,\ldots,r_{l-1}\)은 \(1,\ldots,l-1\)의 한 순열이다. 여기서는 근수 \(\Omega_\lambda\)가 군행렬식의 인자 \(D_\lambda\)와 동일하므로
\[
(\Delta_\lambda)=(D_\lambda D_{\bar\lambda})
=(\Omega_\lambda\Omega_{\bar\lambda})
=\mathfrak P_1^l\cdots\mathfrak P_{l-1}^l
=\frakp_1\cdots\frakp_{l-1}=(p)
\]
를 얻는다. 따라서 \(K/k\)의 도체와 일치한다.

\begin{center}
1931년 8월 24일 접수.
\end{center}

\clearpage
\setcounter{footnote}{0}
\phantomsection
\addcontentsline{toc}{section}{38. 대수론의 한 주정리의 증명}
%
%
%
%
%
%

\section*{38. R. Brauer, H. Hasse와 공동: 대수론의 한 주정리의 증명}

\begin{center}
\emph{Journal f. d. reine u. angew. Math. 167 (1932), 399--404쪽}
\end{center}

\begin{center}
\textbf{대수론의 한 주정리의 증명.}\\[0.5em]
쾨니히스베르크의 \emph{R. Brauer}, 마르부르크의 \emph{H. Hasse},
괴팅겐의 \emph{E. Noether}\footnote{이 글의 집필은 H. Hasse가 맡았다.}.
\end{center}

마침내 우리의 공동 노력으로 다음 정리가 옳음을 증명할 수 있었다. 이 정리는 대수적 수체 위의 대수에 대한 구조론뿐 아니라 그 너머에도 근본적인 의의를 갖는다.

\paragraph{주정리.}
\emph{대수적 수체 위의 모든 정규 나눗셈대수는 순환대수이다. 달리 말하면 Dickson형이다.}

이 성과는 본질적으로 \(p\)-진 방법에 힘입은 것이므로, 이 방법의 창시자인 Kurt Hensel 선생의 70번째 생일에 바칠 수 있어 특히 기쁘다.

우리의 증명은 세 환원으로 이루어지며 우리 셋이 각각 하나씩 기여하였다.\footnote{이하에서는 이 환원들을 생겨난 순서대로 제시하는데, 이는 체계적인 순서와 반대이다.}

\paragraph{1.}
첫 번째 환원은 H. Hasse가 최근 전개한 대수적 수체 위의 순환대수 이론을 바탕으로 제시하였다.\footnote{H. Hasse, \emph{Theorie der zyklischen Algebren über einem algebraischen Zahlkörper}, Gött. Nachr. 1931. 특히 E. Noether가 한 강의에서 전개한 분해체와 교차곱의 이론 --- 그 논문과 여기에서 모두 근본적이다 --- 을 포함하여 증명을 상세히 서술한 논문이 \emph{Theory of cyclic algebras over an algebraic number field}라는 제목으로 곧 Trans. Amer. Math. Soc.에 실릴 예정이다. 이하에서는 이 논문을 H로 인용한다. H, 1--6은 용어의 차이만 제외하면 앞의 짧은 글과 같다.}

\paragraph{환원 1.}
다음 명제를 증명하면 주정리가 증명된다.
\[
\text{I. \(\Omega\) 위에서 모든 곳에서 분할되는 대수는 모두 \( \sim\Omega\)이다.}
\]
이하에서 표현을 줄이기 위하여 \glqq{}\(\Omega\) 위의 대수 \(A\)\grqq{}는 언제나 \glqq{}대수적 수체 \(\Omega\) 위의 정규 단순대수 \(A\)\grqq{}, 곧 중심이 \(\Omega\)인 단순 초복소계를 뜻한다. 또 \(A\sim\Omega\)는 \(A\)가 \(\Omega\) 위의 완전행렬대수라는 뜻이다. 다시 말해 상사성, 곧 \(\Omega\) 위에서 대응하는 나눗셈대수들이 같다는 의미에서 \(A\)가 \(\Omega\) 자체를 \(\Omega\) 위의 나눗셈대수로 하여 정해지는 특별한 류에 속한다는 뜻이다[H, 5]. 마지막으로 \(\Omega\) 위에서 모든 곳에서 분할되는 대수란, \(\Omega\)의 각 소 \(\mathfrak p\)에서 그 \(\mathfrak p\)-진 확대, 곧 계수체 \(\Omega\)를 대응하는 \(\mathfrak p\)-진 체 \(\Omega_{\mathfrak p}\)로 확대하여 얻는 대수가 분할되는 대수를 뜻한다.
\[
A_{\mathfrak p}\sim\Omega_{\mathfrak p}.
\]

\paragraph{증명.}
\(D\)를 \(\Omega\) 위의 나눗셈대수라 하자. \(\mathsf{Z}\)가 \(\Omega\) 위의 순환체이고, \(\Omega\)의 각 소 \(\mathfrak p\)에서 \(\mathsf{Z}\)의 \(\mathfrak p\)-차수가 \(D\)의 \(\mathfrak p\)-인덱스의 배수라고 하자[H, 17 Bb의 주 참조]. 그러면 \(Z\)에서 대응하는 소인자 \(\mathfrak P\)마다 \(\mathfrak P\)-진 체 \(\mathsf{Z}_{\mathfrak P}\)가 \(D_{\mathfrak p}\)의 분해체이다[H, 18.1]. 따라서 \(D\)의 계수체를 \(\mathsf{Z}\)로 확대하여 얻는 \(Z\) 위의 대수 \(D_{\mathsf{Z}}\)는 모든 곳에서 분할된다. 이제 모든 \(\Omega\)에 대하여 I이 이미 성립한다고 하면 \(D_{\mathsf{Z}}\sim \mathsf{Z}\)이다. 이는 \(D\)가 순환 분해체 \(\mathsf{Z}\)를 가지므로 순환적으로 표시될 수 있음을 뜻한다[H, 5]. 그러면 \(D\)는 그 차수가 \(D\) 자신의 차수와 같은 순환 분해체도 가진다. 곧 \(D\)는 순환대수이다[H, 6, 정리 6]. 따라서 I을 가정하면 주정리가 증명된다.

\paragraph{2.}
두 번째 환원에는 R. Brauer가 결정적인 계기를 주었다. 그는 나눗셈대수의 지수(Exponent)의 정확한 값에 관한 관련 문제 --- 주정리로 이 문제도 함께 해결된다 --- 를 Sylow 군정리를 이용하여 가해 분해체의 경우로 환원하는 방법을 H. Hasse에게 편지로 알려 주었다. 이 생각은 순환성 문제를 한 단계 더 환원하는 데에도 쓸 수 있었다.

\paragraph{환원 2.}
다음 명제를 증명하면 명제 I이 증명된다.
\[
\text{II. \(\Omega\) 위에서 가해 분해체를 가지며 모든 곳에서 분할되는 대수는 모두 \(\sim\Omega\)이다.}
\]

\paragraph{증명.}
\(A\)가 \(\Omega\) 위에서 모든 곳에서 분할되는 대수라고 하자. \(\mathsf{K}\)는 \(A\)의 갈루아 분해체, \(p\)는 임의의 소수, \(\Sigma\)는 \(K/\Omega\)의 \(p\)에 관한 Sylow체 가운데 하나라고 하자. 즉 \(\Sigma\)는 갈루아군의 \(p\)-Sylow부분군의 불변체이다. 그러면 \(A_\Sigma\)도 \(\Sigma\) 위에서 모든 곳에서 분할되는 대수이며 가해 분해체 \(\mathsf{K}\)를 갖는다. 모든 \(\Omega\)에 대하여 II가 이미 성립한다고 하면 \(A_\Sigma\sim\Sigma\)이다. 이는 \(A\)가 \(p\)와 서로소인 차수의 분해체 \(\Sigma\)를 갖는다는 뜻이다. 따라서 \(A\)의 인덱스는 이 차수의 약수이므로 \(p\)와 서로소이다. 이것이 모든 소수 \(p\)에 대하여 성립하므로 인덱스는 1이고 \(A\sim\Omega\)이다. 따라서 II를 가정하면 I이 증명된다.\footnote{Sylow 군정리를 이용하여 가해 분해체로 환원하는 생각은 R. Brauer가 앞서도 사용한 바 있다. 곧 인덱스의 모든 소인자가 지수에도 나타남을 보이는 데 사용하였다(\emph{Über den Zusammenhang von arithmetischen und invariantentheoretischen Eigenschaften von Gruppen linearer Substitutionen}, Berl. Akad.-Ber. 1926). 최근 A. A. Albert는 이 생각뿐 아니라 R. Brauer와 E. Noether의 이론에 속하는 여러 일반정리에 대하여 표현론에 의존하지 않는 간단한 증명들을 전개하였다. 1. \emph{On direct products, cyclic division algebras, and pure Riemann matrices}; 2. \emph{On direct products}; 둘 다 Trans. Amer. Math. Soc. 33 (1931). 여기서 말하는 환원은 특히 둘째 논문의 정리 23을 보라.

\emph{인쇄 중 추가.} A. A. Albert는 H. Hasse에게서 아벨대수에 대한 주정리를 증명했다는 편지 내용을 전해 듣고(아래 본문 참조), 우리와 독립적으로 즉시 다음을 이끌어 냈다.

a) 차수가 \(2^e\)인 경우의 주정리,

b) 아래의 정리 1(지수 = 인덱스),

c) 환원 2의 기본사상뿐 아니라 이어지는 환원 3의 사상. 물론 환원 1과의 관계는 없었으므로 결과는 다음과 같았다. \(\Omega\) 위에서 소수 거듭제곱 차수 \(p^e\)를 가진 나눗셈대수 \(D\)에는, \(\Omega\) 위에서 \(p\)와 서로소인 차수를 가진 확대체 \(\Omega'\)가 있어서 \(D_{\Omega'}\)가 순환대수가 된다.

이 세 결과는 모두 그동안 완성된 우리의 주정리 증명으로 포괄되었다. 그러나 이들은 A. A. Albert도 주정리의 증명에 독립적인 몫을 가졌음을 보여 준다. 마지막으로 A. A. Albert는 우리의 주정리 증명을 안 뒤, 우리의 중심명제 I이 출판 중인 그의 논문(Bull. Amer. Math. Soc. 37 (1931))의 정리 13, 10, 9에서 몇 줄 만에 따라온다는 사실도 지적하였다. 그 정리들의 증명은 본질적으로 우리의 환원 2와 3과 같은 추론에 바탕을 둔다.}

\paragraph{3.}
세 번째 환원은 H. Hasse의 통지에서 계기를 얻은 E. Noether가 제시하였다. Hasse의 통지는 아벨 분해체라는 특수 경우에 주정리를 증명한 것이었으며, 일반적인 환원 1과 대응하는 인자계의 공식적 환원을 이용하였다. E. Noether는 그 핵심을 분리하여 바로 세 번째 환원을 얻었다. R. Brauer도 E. Noether와 독립적으로 이 세 번째 환원을 이미 생각해 두었다. 처음 두 환원을 알고 있던 H. Hasse는 이 환원이 최종 증명으로 이어짐을 알아보았다.

\paragraph{환원 3.}
다음 명제를 증명하면 명제 II가 증명된다.
\[
\begin{gathered}
\text{III. \(\Omega\) 위에서 소수차수의 순환 분해체를 가지며}\\
\text{모든 곳에서 분할되는 대수는 모두 \(\sim\Omega\)이다.}
\end{gathered}
\]

\paragraph{증명.}
\(A\)가 \(\Omega\) 위에서 모든 곳에서 분할되고 가해 분해체 \(K\)를 갖는 대수라고 하자.
\[
K=\Lambda_0>\Lambda_1>\cdots>\Lambda_r=\Omega
\]
를 \(\mathsf{K}\)와 \(\Omega\) 사이의 체들의 사슬로 택하되, \(\Lambda_i/\Lambda_{i+1}\)가 언제나 소수차수의 순환확대가 되게 하자. 그러면 \(A_{\Lambda_1}\)도 \(\Lambda_1\) 위에서 모든 곳에서 분할되며 소수차수의 순환 분해체 \(\mathsf{K}=\Lambda_0\)를 갖는다. 모든 \(\Omega\)에 대하여 III이 이미 성립한다고 하면 \(A_{\Lambda_1}\sim\Lambda_1\)이다. 이는 \(\Lambda_1\)이 \(A\)의 분해체라는 뜻이다. 이제 \(\Lambda_0\) 대신 \(\Lambda_1\)에서 시작하여 같은 방식으로 \(A_{\Lambda_2}\sim\Lambda_2\)를 증명하고, 이를 계속하면 마침내 \(A_{\Lambda_r}\sim\Lambda_r\), 곧 \(A\sim\Omega\)를 얻는다. 따라서 III을 가정하면 II가 증명된다.

\paragraph{4.}
그런데 III은 H. Hasse의 결과[H, 24]로 이미 확립되어 있다. 따라서 환원 3, 2, 1을 거꾸로 적용하면 주정리가 옳음이 따라온다.

E. Noether는 이와 함께 다음을 지적한다. III의 타당성은, 더 일반적으로 임의 차수의 순환 분해체에 대해서도, Hasse의 노름정리\footnote{\label{fn:p38-hasse-normensatz}H, 3.11에 인용되어 있으며 H. Hasse, \emph{Beweis eines Satzes und Widerlegung einer Vermutung über das allgemeine Normenrestsymbol}, Gött. Nachr. 1931에서 증명되었다.}로 귀착된다. 그러나 환원 3에는 소수차수라는 특수 경우, 곧 본래의 Hilbert--Furtwängler 노름정리\footnote{H. Hasse, \emph{Bericht über neuere Untersuchungen und Probleme in der Theorie der algebraischen Zahlkörper II}, Jahresber. der D. M.-V., Erg.-Bd. 6 (1930), § 8 및 그곳의 문헌을 보라.}만 필요하다. 그러므로 환원 III은 Hasse 노름정리의 새롭고 간단한 증명을 준다. 그 증명은 다음과 같다.

\(Z\)가 \(\Omega\) 위의 순환체이고 \(\alpha\in\Omega\)가 \(\Omega\)의 모든 소 \(\mathfrak p\)에 대하여 노름잉여기호
\[
\left(\frac{\alpha,Z}{\mathfrak p}\right)=1
\]
을 만족한다고 하자. 그러면 순환대수
\[
A=(\alpha,Z)
\]
\footnote{\label{fn:p38-hasse-notation}H, 1의 표기. \(\alpha\)와 결부된 \(Z\)의 자기동형사상 \(S\)는 여기에서 중요하지 않아 생략하였다.}
는 모든 곳에서 분할된다[H, 17.7]. 환원 3에 따라 --- Hilbert--Furtwängler 노름정리만을 사용하여 --- \(A\sim\Omega\)이다. 따라서 \(\alpha\)는 \(Z\)의 한 원소의 노름이다[H, 15.4].

노름정리와 관련하여, 명제 I은 그 정리를 고차의 경우, 비아벨의 경우까지 올바르게 일반화한 것으로 보아야 한다. 반면 H. Hasse가 보였듯이\textsuperscript{\ref{fn:p38-hasse-normensatz}} 글자 그대로의 일반화는 일반적으로 옳지 않다.

\begin{center}
\textbf{귀결들(H. Hasse)}
\end{center}

\paragraph{5.}
주정리의 가장 가까운 귀결로, H. Hasse가 전개한 대수적 수체 위의 순환대수 이론이 이제 대수적 수체 위의 정규 단순대수, 특히 정규 나눗셈대수에 대한 일반적인 구조론과 불변식론이 되었다는 점을 들 수 있다. 특히 이제 다음 명제가 일반적으로 증명되었다.

\paragraph{정리 1.}
대수적 수체 위의 정규 단순대수의 지수(Exponent)는 그 인덱스(Index)와 같다.

또 명제 I에서 다음 주목할 사실이 따라온다.

\paragraph{정리 2.}
\(\Omega\) 위의 진정한 정규 나눗셈대수의 기초아이디얼은 \(\Omega\)의 적어도 한 소, 사실 적어도 두 소로 나누어진다.

여기서 기초아이디얼은 예를 들어 축약 디퍼런트의 축약 노름으로 정의한다.\footnote{유리 사원수대수라는 특수 경우에는 H. Brandt가 도입한 \glqq{}기초수\grqq{} 개념으로 귀착된다(\emph{Idealtheorie in Quaternionenalgebren}, Math. Ann. 99 (1928)). 여기서는 수가 아니라 아이디얼을 다루므로 \glqq{}기초아이디얼\grqq{}이라 해야 했다. 이를 Dedekind가 디퍼런트와 판별식에 사용한 기초아이디얼 및 기초수라는 명칭과 혼동해서는 안 된다. 그 명칭은 바로 이 이유 때문에 상대체에는 쓸 수 없다. 축약 디퍼런트의 정의는 다음 각주를 보라. E. Noether는 아이디얼의 축약 노름도 \glqq{}소별로\grqq{} 정의한다. 각 소에서 모든 아이디얼이 주아이디얼이라는 사실에 맞추어, 각 소에서 주아이디얼 기저원소의 축약 수노름을 취하면 된다.} 또 무한소는 그곳에서 대수가 실수체나 복소수체가 아니라 사원수대수로 환원될 때 그리고 그때에만 기초아이디얼에 포함한다.

\paragraph{증명.}
H. Hasse의 결과\footnote{\label{fn:p38-hasse-padisch}H. Hasse, \emph{Über \(\mathfrak p\)-adische Schiefkörper und ihre Bedeutung für die Arithmetik hyperkomplexer Zahlsysteme}, Math. Ann. 104 (1931), 정리 42, 59.}에 따라 \(\Omega\)의 소 \(\mathfrak p\)가 기초아이디얼을 나누는 것은 \(\mathfrak p\)-인덱스가 1과 다를 때 그리고 그때뿐이다. 모든 \(\mathfrak p\)-인덱스가 1이라면 대수는 모든 곳에서 분할되며, 명제 I에 따라 진정한 나눗셈대수가 아니다. 한 \(\mathfrak p\)-인덱스만 1과 다를 수도 없다는 사실은, 그 위수가 \(\mathfrak p\)-인덱스인 노름잉여기호들[H, 17.7]에 곱정리, 곧 상호법칙[H, 3.8]이 성립한다는 데서 따라온다.

\paragraph{6.}
더 나아가 명제 I을 이용하면 \(\Omega\) 위의 한 대수 \(A\)에 속하는 모든 분해체 \(K\)를 결정하고 산술적으로 특징지을 수 있다. 순환적인 \(A,\mathsf{K}\)의 특수 경우에 H. Hasse가 이미 확립한 사실[H, 6, 정리 2]의 정확한 일반화로 다음을 얻는다.

\paragraph{정리 3.}
\(\Omega\) 위의 대수 \(A\)에 대하여, \(\Omega\) 위의 대수적 수체 \(\mathsf{K}\)가 분해체인 것은 \(\Omega\)의 모든 소 \(\mathfrak p\)와 \(K\)에서 그 모든 소인자 \(\mathfrak P_i\)에 대하여 \(K\)의 \(\mathfrak P_i\)-차수 \(n_{\mathfrak P_i}\)가 \(A\)의 \(\mathfrak p\)-인덱스 \(m_{\mathfrak p}\)의 배수일 때 그리고 그때뿐이다.

\(\mathsf{K}\)의 \(\mathfrak P_i\)-차수는 \(\mathfrak P_i\)-진 체 \(\mathsf{K}_{\mathfrak P_i}\)의 \(\mathfrak p\)-진 체 \(\Omega_{\mathfrak p}\) 위 차수이다. 즉
\[
\mathfrak p=\prod_i\mathfrak P_i^{e_{\mathfrak P_i}},
\qquad
\qquad \mathsf{N}_{\mathsf{K}\Omega}(\mathfrak P_i)=\mathfrak p^{f_{\mathfrak P_i}}
\]
가 \(K\)에서 \(\mathfrak p\)의 분해라면, 이는 잉여차수 \(f_{\mathfrak P_i}\)와 \(\mathfrak p\)에 대한 \(\mathfrak P_i\)의 분기지수 \(e_{\mathfrak P_i}\)의 곱
\[
n_{\mathfrak P_i}=f_{\mathfrak P_i}e_{\mathfrak P_i}
\]
이다[H, 4]. 또 \(A\)의 \(\mathfrak p\)-인덱스는 \(\mathfrak p\)-진 확대 \(A_{\mathfrak p}\)의 인덱스 \(m_{\mathfrak p}\)이다[H, 6].

\paragraph{증명.}
\(\mathsf{K}\)가 \(A\)의 분해체, 곧 \(A_{\mathsf{K}}\sim \mathsf{K}\)라는 것은 명제 I에 따라 모든 \(\mathfrak P_i\)에 대하여 \(A_{\mathsf{K}_{\mathfrak P_i}}\sim \mathsf{K}_{\mathfrak P_i}\)라는 것과 동치이다.

\(\mathfrak p\)가 유한소일 때
\[
A_{\mathfrak p}\sim(\pi,W_{\mathfrak p})
\]
\textsuperscript{\ref{fn:p38-hasse-notation}}
를 \(A_{\mathfrak p}\)의 산술적으로 뛰어난 순환표현이라 하자[H, 16 및 앞서 인용한 논문\textsuperscript{\ref{fn:p38-hasse-padisch}}, 정리 38]. 여기서 \(W_{\mathfrak p}\)는 \(\Omega_{\mathfrak p}\) 위의 차수 \(m_{\mathfrak p}\)인 비분기체이고, \(\pi\)는 \(\Omega_{\mathfrak p}\)에서 정확히 \(\mathfrak p^1\)로 나누어지는 수이다. 또 \(\mathit{K}_{\mathfrak P_i}\)를 합성체라 하고
\[
\Delta_{\mathfrak P_i}=K_{\mathfrak P_i}\cap W_{\mathfrak p}
\]
를 \(W_{\mathfrak p}\)와 \(\mathsf{K}_{\mathfrak P_i}\)의 교집합이라 하자. \(\mathsf{K}_{\mathfrak P_i}\) 안의 가장 큰 비분기 부분체는 \(\Omega_{\mathfrak p}\) 위에서 차수 \(f_{\mathfrak P_i}\)이고, \(\Omega_{\mathfrak p}\) 위에는 각 차수의 비분기체가 하나뿐이므로 \(\Delta_{\mathfrak P_i}\)는
\[
\Omega_{\mathfrak p}\text{ 위에서 차수 }
d_{\mathfrak P_i}=(m_{\mathfrak p},f_{\mathfrak P_i})
\]
이다. 따라서 \(\mathit{K}_{\mathfrak P_i}\)는 \(\mathsf{K}_{\mathfrak P_i}\) 위에서 차수 \(\frac{m_{\mathfrak p}}{d_{\mathfrak P_i}}\)이고 비분기이다.

이제
\[
A_{K_{\mathfrak P_i}}
=(A_{\mathfrak p})_{K_{\mathfrak P_i}}
\sim(\pi,K_{\mathfrak P_i}W_{\mathfrak p})
\]
이다[H, 15.5]. 따라서 \(A_{\mathsf{K}_{\mathfrak P_i}}\sim \mathsf{K}_{\mathfrak P_i}\)인 것은 \(\pi\)가 \(\mathit{K}_{\mathfrak P_i}\)에서 \(\mathsf{K}_{\mathfrak P_i}\)로의 노름일 때 그리고 그때뿐이다. \(K_{\mathfrak P_i}W_{\mathfrak p}/K_{\mathfrak P_i}\)가 비분기이므로, 이는 \(\mathfrak P_i\)에서 \(\pi\)의 위수 \(e_{\mathfrak P_i}\)가 이 확대의 차수 \(\frac{m_{\mathfrak p}}{d_{\mathfrak P_i}}\)의 배수일 때 그리고 그때뿐이다. 그런데
\[
\left(
\frac{m_{\mathfrak p}}{d_{\mathfrak P_i}},
\frac{f_{\mathfrak P_i}}{d_{\mathfrak P_i}}
\right)=1
\]
이므로 이는
\[
\frac{f_{\mathfrak P_i}}{d_{\mathfrak P_i}}e_{\mathfrak P_i}
=\frac{n_{\mathfrak P_i}}{d_{\mathfrak P_i}}
\]
가 \(\frac{m_{\mathfrak p}}{d_{\mathfrak P_i}}\)의 배수라는 것, 곧 \(n_{\mathfrak P_i}\)가 \(m_{\mathfrak p}\)의 배수라는 것과 같다.

\(\mathfrak p\)가 무한소이고 자명한 경우 \(m_{\mathfrak p}=1\)이 아니라면 \(m_{\mathfrak p}=2\)이고 \(A_{\mathfrak p}\)는 실수체 \(\Omega_{\mathfrak p}\) 위의 사원수대수이다. 이때
\[
A_{K_{\mathfrak P_i}}
=(A_{\mathfrak p})_{K_{\mathfrak P_i}}\sim K_{\mathfrak P_i}
\]
인 것은 \(\mathsf{K}_{\mathfrak P_i}\)가 복소수체인 것, 곧
\[
n_{\mathfrak P_i}=f_{\mathfrak P_i}=2=m_{\mathfrak p}
\]
인 것과 동치이다. 여기에는 \(e_{\mathfrak P_i}\)가 나타나지 않는다. 따라서 이 경우에도 명제가 얻어진다. \(n_{\mathfrak P_i}\)와 \(m_{\mathfrak p}\)는 여기서 1 또는 2만을 값으로 가질 수 있다.

\paragraph{7.}
정리 3에서 답한 문제를 뒤집으면 중요한 결과들을 얻는다. 곧 고정된 대수 \(A\)의 모든 분해체 \(\mathsf{K}\)를 찾는 대신, \(\Omega\) 위의 대수적 수체 \(\mathsf{K}\)를 고정하고 \(K\)가 분할하는 \(\Omega\) 위의 모든 대수 \(A\)를 생각한다. 이 대수들은 \(\mathsf{K}\)가 정하는 부분군 \(\mathfrak K\)를 이루며, 이는 \(\Omega\) 위의 모든 대수, 더 정확히는 서로 상사인 대수류 전체로 이루어진 R. Brauer 군 \(\mathfrak A\)의 부분군이다.\footnote{R. Brauer, \emph{Über Systeme hyperkomplexer Größen}, Jahresber. d. D. M.-V. 38 (1929), 47/48쪽. H, 13.1도 보라.} \(K/\Omega\)가 갈루아라고 가정하면, 상대 아벨 수체에 관한 이론인 유체론의 주정리들을 일반 상대 갈루아 수체로 확장한 것으로 볼 수 있는 다음 정리들에 이른다.

\paragraph{정리 4.}
\emph{(분해정리.)} \(\Omega\)의 소아이디얼 \(\mathfrak p\)가 \(\mathsf{K}\)의 상대판별식을 나누지 않을 때, \(\mathsf{K}\)에서 \(\mathfrak p\)의 소인자들이 갖는 상대차수 \(f\)는 \(K\)에 대응하는 군 \(\mathfrak K\)의 모든 대수 \(A\)에 대하여
\[
A_{\mathfrak p}^{\,f}\sim\Omega_{\mathfrak p}
\]
가 되게 하는 가장 작은 양의 지수이다.

\paragraph{증명.}
\(f\)는 \(K\)에서 \(\mathfrak p\)의 모든 소인자에 공통인 \(\mathfrak p\)-차수이다. 한편 \(A\)의 \(\mathfrak p\)-인덱스는 \(A_{\mathfrak p}\)의 인덱스이고, 정리 1에 따라 그 지수와 같다. 그러므로 정리 3에 따라 \(f\)는 위 최소지수의 배수이다. 이제 \(\mathfrak K\) 안에 정확한 \(\mathfrak p\)-인덱스 \(f\)를 가진 대수 \(A\)가 존재함만 보이면 된다.

Frobenius 밀도정리에 따라 \(K\)의 상대판별식을 나누지 않으면서 그 소인자들이 상대차수 \(f\)를 갖는 다른 소아이디얼 \(\mathfrak p'\)가 \(\Omega\)에 존재한다. \(Z\)를 \(\Omega\) 위의 순환체로 택하되 \(\mathfrak p\)-차수와 \(\mathfrak p'\)-차수가 모두 \(f\)가 되게 하자. 또 \(\alpha\in\Omega\)를 택하여 노름잉여기호
\[
\left(\frac{\alpha,Z}{\mathfrak p}\right)
\quad\text{와}\quad
\left(\frac{\alpha,Z}{\mathfrak p'}\right)
\]
가 위수 \(f\)인 서로 역인 값을 갖게 하고, \(Z\)의 도체를 나누는 나머지 모든 \(\mathfrak q\)에 대하여
\[
\left(\frac{\alpha,Z}{\mathfrak q}\right)=1
\]
이 되게 하자. 일반화된 산술진행 정리에 따라, \(\alpha\)가 \(\mathfrak p,\mathfrak p'\), \(\mathfrak q\)들을 제외하면 많아야 하나의 소아이디얼 \(\mathfrak r\)만을 정확히 일차로 포함하도록 택할 수도 있다. 그러면 노름잉여기호의 곱정리, 곧 상호법칙에 따라
\[
\left(\frac{\alpha,Z}{\mathfrak r}\right)=1
\]
도 성립한다[H, 3.8--10]. 따라서 대수
\[
(\alpha,Z)=A
\]
는 \(\mathfrak p\)-인덱스와 \(\mathfrak p'\)-인덱스가 \(f\)이고, \(\Omega\)의 다른 모든 소에서는 인덱스가 1이다[H, 17.7]. \(\mathfrak p\)에 대한 가정과 \(\mathfrak p'\)의 선택을 고려하면 정리 3에서 \(\mathsf{K}\)가 \(A\)의 분해체임이 따라온다. 이로써 실제로 군 \(\mathfrak K\) 안에 \(\mathfrak p\)-인덱스 \(f\)인 대수 \(A\)가 존재함을 보였다.

\paragraph{정리 5.}
\emph{(유일성 및 순서정리.)} \(\mathsf{K}\leqq \mathsf{K}'\)이면 \(\mathfrak K\leqq\mathfrak K'\)이고, 역도 성립한다.

특히 갈루아체 \(K\)와 그 대수군 \(\mathfrak K\)의 대응은 일대일이며 역으로도 유일하다.

\paragraph{증명.}
a) \(\mathsf{K}\leqq \mathsf{K}'\)이면 자명하게 \(\mathfrak K\leqq\mathfrak K'\)이다. \(\mathsf{K}\)가 분할하는 모든 대수 \(A\)는 당연히 \(\mathsf{K}'\)도 분할하기 때문이다.

b) 역으로 \(\mathfrak K\leqq\mathfrak K'\)라 하자. 분해정리에 따라 상대판별식을 나누지 않는 각 \(\mathfrak p\)에 대하여 \(f_{\mathfrak p}\mid f'_{\mathfrak p}\)이다. 따라서 특히 \(f'_{\mathfrak p}=1\)인 거의 모든 \(\mathfrak p\)에 대하여 \(f_{\mathfrak p}=1\)이다. 알려진 해석적 추론\footnote{예를 들어 H. Hasse[앞에서 인용한 주 6의 논문], § 25, III을 보라.}에 따라 여기서 \(\mathsf{K}\leqq \mathsf{K}'\)가 따라온다.

\paragraph{8.}
마지막으로 주정리가 I. Schur가 다룬 문제, 곧 유한군의 절대기약표현이 가능한 수체들을 찾는 문제\footnote{I. Schur, \emph{Arithmetische Untersuchungen über endliche Gruppen linearer Substitutionen}, Berl. Akad.-Ber. 1906.}에 본질적인 진전을 가져옴을 설명하겠다.

\paragraph{정리 6.}
유한군 \(\mathfrak G\)의 모든 절대기약표현은 원분체에서 가능하다. 예를 들어 \(\mathfrak G\)의 위수가 \(n\)이면, 충분히 큰 \(h\)에 대하여 \(n^h\)차 단위근체에서는 언제나 가능하다.

\paragraph{증명.}
\(\mathfrak G\)의 유리수계수 군환 \(G\)로 넘어가면, \(\mathfrak G\)의 절대기약표현 \(\Gamma_i\)들은 반단순대수 \(G\)의 단순 성분 \(G_i\)들의 절대기약표현이 된다. 그 중심은 각각 대응하는 지표체 \(\Omega_i\)이다.\footnote{a) R. Brauer와 E. Noether, \emph{Über minimale Zerfällungskörper irreduzibler Darstellungen}, Berl. Akad.-Ber. 1927, § 1. b) R. Brauer, \emph{Über Systeme hyperkomplexer Zahlen}, Math. Zeitschr. 30 (1929), 정리 3. c) E. Noether, \emph{Hyperkomplexe Größen und Darstellungstheorie}, Math. Zeitschr. 30 (1929), §§ 21, 24, 26.} \(\mathfrak G\)가 유한하므로 이 중심들은 원분체이고, 특히 \(n\)차 단위근체의 부분체이다.

정리 3에 따라 \(\Omega_i\) 위의 순환체 \(\mathsf{Z}_i\)가 \(G_i\)의 분해체인 것은, \(\Omega_i\)의 모든 소 \(\mathfrak p\)에서 그 \(\mathfrak p\)-차수 \(n_{i\mathfrak p}\)가 \(G_i\)의 \(\mathfrak p\)-인덱스 \(m_{i\mathfrak p}\)의 배수일 때이다. \(m_{i\mathfrak p}\)가 1과 다른 것은 \(G_i/\Omega_i\)의 기초아이디얼을 나누는 \(\mathfrak p\)들뿐이고, 따라서 \(G_i\)의 절대판별식을 나누는 \(\mathfrak p\)들뿐이다. 이 판별식은 \(n^n\)을 나눈다. 실제로 \(n^n\)은 \(G\)의 한 비극대 오더의 판별식이다.\footnote{E. Noether[앞의 주 7c], § 26을 보라.} 따라서 문제의 \(\mathfrak p\)들은 많아야 \(n\)을 나누는 소인자들이다. 이 \(\mathfrak p\)들에 대하여 \(n_{i\mathfrak p}\)를 \(m_{i\mathfrak p}\)의 배수로 만들기 위해서는, 각 \(\mathfrak p\)가 \(\mathsf{Z}_i\)에서 \(m_{i\mathfrak p}\)로 나누어지는 분기지수를 갖도록 요구하면 충분하다. 충분히 큰 \(h\)에 대한 \(n^h\)차 단위근체가 이를 만족함은 쉽게 알 수 있다.

인용한 논문의 마지막 정리에서 I. Schur는 지금까지 알려진 모든 경우에 이미 \(n\)차 단위근체만으로 충분하다고 밝혔다. 이것이 언제나 성립하는지, 또 여기서 전개한 방법들로 이 문제를 결정할 수 있는지는 앞으로의 연구에 맡긴다.

\begin{center}
1931년 11월 11일 접수.
\end{center}

\clearpage
\setcounter{footnote}{0}
\phantomsection
\addcontentsline{toc}{section}{39. 초복소계와 가환대수 및 수론의 관계}
%
%
%
%
%
%

\section*{39. 초복소계와 가환대수 및 수론의 관계}

\begin{center}
\emph{Verhandl. Intern. Math.-Kongreß Zürich 1 (1932), 189--194쪽}
\end{center}

\begin{center}
\textsc{괴팅겐의 에미 뇌터}
\end{center}

\begingroup
\renewcommand{\thefootnote}{\arabic{footnote})}
\setcounter{footnote}{0}

\paragraph{1.}
초복소계, 곧 대수의 이론은 최근 몇 년 동안 크게 발전하였다. 그러나 이 이론이 가환적인 문제들에 갖는 의의가 분명해진 것은 극히 최근의 일이다. 오늘 나는 비가환적인 것이 가환적인 것에 갖는 이 의의를 보고하고자 한다. 구체적으로는 가우스까지 거슬러 올라가는 두 고전적 문제, 곧 주종정리와 그에 밀접하게 연결된 노름정리를 따라 이를 살펴보겠다. 이 문제들의 정식화는 시대에 따라 거듭 달라졌다. 가우스에게서는 이차형식론의 결말로 나타났고, 뒤에는 유체론을 통하여 상대 순환 수체와 상대 아벨 수체를 특성화하는 데 본질적인 구실을 하였다. 마지막으로는 자기동형사상과 대수의 분할에 관한 정리로 표현할 수 있으며, 이 마지막 정식화는 동시에 그 정리들을 임의의 상대 갈루아 수체로 옮겨 준다.

뒤에서 자세히 전개할 이 개요를 통하여, 비가환적인 것을 가환적인 것에 적용하는 \emph{원리}도 설명하고자 한다. \emph{대수론을 이용하여 이차형식이나 순환체에 관한 알려진 사실들을 불변적이고 단순한 형태, 곧 대수의 구조적 성질에만 의존하는 형태로 정식화하려 한다. 이러한 불변적 정식화를 일단 증명하면 --- 위에서 든 예들에서는 실제로 증명되어 있다 --- 그 사실들은 저절로 임의의 갈루아체로 옮겨진다.}

\paragraph{2.}
세부를 전개하기 전에 여러 방법과 그 밖의 결과들을 전체적으로 개관하겠다. 우선 가장 큰 어려움은 일반 갈루아체에 맞는 정식화를 얻는 데 있으며, 초복소적 방법이 없었다면 이를 시작할 실마리조차 없었음을 말해 두어야 한다. 앞서 든 예에서 비가환적인 것으로 넘어가는 바탕은 \glqq{}교차곱\grqq{}과 그 곱셈상수인 \glqq{}인자계\grqq{}를 이용하여 \emph{체와 군을 동시에 고찰}하는 것이다(3절 참조). 이렇게 하면 기초체 위의 단순 정규대수 하나를 얻으며, 본질적으로 모든 그러한 대수를 이 방식으로 생성할 수 있다. 이러한 교차곱은 Dickson이 처음 고찰하였다.\footnote{그의 저서 \emph{Algebren und ihre Zahlentheorie}, Zürich 1927, § 34를 보라.} 한편 인자계의 이론은 Speiser, Schur, R. Brauer가 절대기약표현의 문제라는 전혀 다른 출발점에서 발전시켰다.\footnote{예를 들어 R. Brauer, \emph{Untersuchungen über die arithmetischen Eigenschaften von Gruppen linearer Substitutionen}, Math. Z. 28 (1928) 및 그 논문의 주 2)에 제시된 문헌을 보라.} 두 이론을 합친 뒤에야 가환적인 문제를 다룰 만큼 단순하면서도 넓게 적용되는 구성을 얻을 수 있었다.\footnote{나는 이 구성을 1929/30년 겨울학기 강의에서 처음 전개하였다. 그 내용은 H. Hasse, \emph{Theory of cyclic algebras}, Transac. 134 (1932), 제2장에 실려 있다. 이 강연이 다루는 전 분야에 관해서는 M. Deuring의 초복소수와 수론적 응용에 관한 보고가 길잡이가 되며, 이는 \glqq{}Ergebnisse der Mathematik\grqq{} 총서에 출판될 예정이다.}

동시에 알려진 사실들에도 새롭고 투명한 증명이 생긴다. 여기서는 H. Hasse가 제시하여 곧 Math. Ann.에 실릴 순환체의 \emph{상호법칙}에 대한 초복소적 증명을 들고자 한다.\footnote{H. Hasse, \emph{Die Struktur der R. Brauerschen Algebrenklassengruppe über einem algebraischen Zahlkörper (insbesondere Normenrestsymbol und Reziprozitätsgesetz)}, Math. Ann. 107 (1932/33).} 이 증명은 교차곱 이론에 바탕을 둔 그의 노름잉여기호의 불변적 정식화를 이용한다. 또 C. Chevalley가 최근 같은 토대 위에서 제시한 국소 유체론의 초복소적 기초도 있다. 다만 이를 위해서는 인자계에 관한 새로운 대수적 정리들을 더 발전시켜야 했다.\footnote{C. Chevalley, \emph{Sur la théorie du symbole de restes normiques}; J. f. Math. 169에 실릴 예정이다.} 그러나 교차곱의 방법만으로는 아무래도 \emph{갈루아 수체의 이론 전체}를 얻지 못한다는 제한도 덧붙여야 한다. 위에서 말한 Hasse의 증명을 앞의 원리에 따라 잇는 Artin의 아직 출판되지 않은 새 결과들은 완전한 동형정리 대신 개수의 등식만을 주기 때문이다.

완전한 동형, 다만 연산자 동형사상을 주는 방법은 대수적인 범위 안에 이미 존재한다. 이는 A. Speiser의 착상을 이어 간 것이다.\footnote{A. Speiser, \emph{Gruppendeterminante und Körperdiskriminante}, Math. Ann. 77 (1916).} 갈루아체를 \glqq{}갈루아 가군\grqq{}, 곧 기초체 위의 가군으로서 갈루아군의 치환들을 연산자로 허용하는 가군으로 보는 관점이다. 체와 군환(군대수) 사이에는 다음 의미의 연산자 동형사상이 존재한다. 원소들이 일대일로 대응하고 기초체 위의 선형형식들이 서로 대응하며, 체에서의 갈루아군 치환에는 군환에서의 곱셈이 대응한다. 내가 세운 이 정리\footnote{E. Noether, \emph{Normalbasis bei Körpern ohne höhere Verzweigung}, 정리 3, J. f. Math. 167 (1932). 그 증명에는 틈이 있다.}는 M. Deuring이 증명하였다.\footnote{M. Deuring, \emph{Galoissche Theorie und Darstellungstheorie}, Math. Ann. 107 (1932).} 그는 이를 바탕으로 갈루아 이론을 증명하였고, 이때 연산자 동형사상이 군과 체의 대응을 실현한다. 역시 Deuring이 얻은 더 깊은 구조정리들은 Artin \(L\)-급수의 형식적 사실들과 나란히 놓이며 Artin 도체에 구조적으로 접근하는 길을 준다. 일반 군지표로 만드는 이 Artin \(L\)-급수와 도체\footnote{E. Artin, \emph{Über eine neue Art von \(L\)-Reihen}, Math. Sem. Hamburg 3 (1924); \emph{Zur Theorie der \(L\)-Reihen mit allgemeinen Gruppencharakteren}, 같은 곳 8 (1931); \emph{Die gruppentheoretische Struktur der Diskriminanten algebraischer Zahlkörper}, J. f. M. 164 (1931).}는 Speiser의 연구\textsuperscript{6)}와 더불어 수론과 표현론 사이의 첫 연결이며, 아벨체를 넘어선 첫 진전이다. 이 연구들은 전체 발전에 강한 자극을 주었고, 특히 갈루아 가군론은 여기에서 방향을 얻었다.

\paragraph{3.}
이제 처음에 내세운 문제들, 곧 노름정리와 주종정리를 자세히 살펴보겠다. 먼저 \emph{교차곱의 정의}를 제시한다. \(K/k\)를 \(n\)차 갈루아체라 하고 그 군을 \(\Gg\)라 하자. 교차곱이란 \(K\)와 \(\Gg\)를 하나의 대수 \(A\) 안에 동시에 매장하여 \(K\)의 자기동형사상들이 내부자기동형사상이 되게 하는 것이다. \(n\)개의 군 원소에 대응하는 기호들을 \(u_{S_1},\ldots,u_{S_n}\)이라 하자. 그러면 먼저 \(A\)를 \(K\) 위에서 랭크 \(n\)인 선형형식 가군으로 둔다.
\srcnumdisplay[2.8em]{(1)}{\begin{gathered}
A=u_{S_1}K+\cdots+u_{S_n}K\\
\text{\normalfont(즉 \(A\)는 \(a_i\)가 \(K\)의 임의의 원소인}\\
\text{\normalfont 모든 선형형식 \(u_{S_1}a_1+\cdots+u_{S_n}a_n\)으로 이루어진다.)}
\end{gathered}}
내부자기동형사상 --- \(u_S\), 더 일반적으로는 \(u_SK^\ast\)가 생성한다\footnote{\(K^\ast\)는 \(K\)에서 0을 뺀 집합이며, 이하 이 표기를 계속 사용한다.} --- 을 요구하면 \(A\)는 환, 따라서 \(k\) 위에서 랭크 \(n^2\)인 대수가 된다. 그 요구는 모든 \(z\in K\)에 대하여 다음과 같이 표현된다.
\srcnumdisplay[2.8em]{(2)}{u_S^{-1}zu_S=z^S\footnotemark,\qquad \text{또는}\qquad zu_S=u_Sz^S}
\footnotetext{\(z^S\)는 통상적인 표기대로 치환 \(S\)로 \(z\)에서 생기는 원소를 뜻한다.}
\srcnumdisplay[2.8em]{(3)}{u_Su_T=u_{ST}a_{S,T}\quad \text{\(a_{S,T}\in K^\ast\)}}
\srcnumdisplay[2.8em]{(4)}{\begin{gathered}
a_{ST,R}\,a_{S,T}^{\,R}=a_{S,TR}\,a_{T,R}\\
\text{\normalfont(\([u_Su_T]u_R=u_S[u_Tu_R]\)에서 오는 결합법칙)}
\end{gathered}}
\(A\)를 인자계 \(a_{S,T}\)에 대한 \(K\)와 \(\Gg\)의 교차곱이라 한다. \(A\)는 \(k\) 위의 단순 정규대수, 곧 대응하는 나눗셈대수 \(D\) 위의 \(r\)차 행렬환 \(D_r\)이 되고, \(K\)는 극대 가환 부분체, 따라서 분해체가 됨을 증명할 수 있다. 여기서 분해체라는 말은 계수영역 \(k\)를 \(K\)와 동형인 체로 확대하면 대수가 \emph{분할}, 곧 그 중심 위의 완전행렬환이 된다는 뜻이다. 역으로 주어진 나눗셈대수 \(D\)에 대해서도 위와 같은 방식으로 교차곱으로 생성할 수 있는 행렬환 \(D_r\)이 언제나 존재한다.

\(c_S\in K^\ast\)를 택하여 \(u_S\)에서 같은 자기동형사상을 생성하는 \(v_S=u_Sc_S\)로 옮기면 \glqq{}연관된\grqq{} 인자계가 생긴다.
\srcnumdisplay[2.8em]{(5)}{\bar a_{S,T}=a_{S,T}\frac{c_S^{\,T}c_T}{c_{ST}}}
연관된 인자계들은 한 류 \((a)\)로 묶는다. 마찬가지로 \(A\)와 상사인 모든 대수, 곧 \(r=1,2,\ldots\)인 모든 \(D_r\)를 한 류 \(\mathfrak A\)로 묶는다. \emph{류 \(\mathfrak A\)와 \((a)\)는 일대일로 대응한다. 분해체 \(K\)를 고정한 류들은 직접곱 구성에 관하여 아벨군을 이루며, 이는 인자계의 류들을 성분별로 곱하여 얻는 군과 동형이다. 단위원은 각각 분할대수류와 모든 변환량 \(c_S^{\,T}c_T/c_{ST}\)의 계이다.}
이것이 R. Brauer가 발견한 대수류군이다.

\paragraph{4.}
이제 \emph{순환 분해체}로 \emph{특수화}하여 \emph{노름개념}과의 관계, 따라서 처음에 말한 원리에 따른 일반화된 \emph{노름정리}의 정식화에 이르겠다. \(Z\)가 순환체이고 \(S\)가 그 군의 생성치환이면, 대응하는 대수를 순환대수라 한다. \(S\)의 거듭제곱들에 \(u\)의 거듭제곱들을 대응시켜 다음과 같이 둘 수 있다.
\[
\begin{aligned}
(1')\qquad &A=Z+uZ+\cdots+u^{n-1}Z,\\
(2')\qquad &zu=uz^S,\\
(3')\quad &u^n=a,\\
(4')\quad &a\text{는 기초체 }k^\ast\text{에 속한다},\\
(5')\quad &v=uc\text{로 두면 }\bar a=a\cdot N(c).
\end{aligned}
\]
따라서 여기서 각 인자계는 기초체에 속하는 단 하나의 원소 \(a\)로 이루어지며 --- \(A=(\alpha,Z)\)로 표기한다 --- 인자계의 단위원류는 \(Z^\ast\)에서 오는 \emph{노름}들로 주어진다. 대수류군은 몫군 \(k^\ast/N(Z^\ast)\)와 동형이다.
그러므로 순환대수 \((\alpha,Z)\)가 분할되는 것은 \(a\)가 \(Z\)의 한 원소의 노름일 때 그리고 그때뿐이다.

노름과 분할 사이의 이 관계는 \glqq{}노름정리\grqq{}를 다음과 같이 정식화한다. \emph{분할대수에 관한 정리: 한 대수가 모든 소에서 분할되면 전역적으로도 분할된다.} 여기서 \glqq{}소\grqq{}는 수론에서 통상적인 방식으로, 기초체 \(k\)를 그 \(\mathfrak p\)-진 확대 \(k_{\mathfrak p}\)로 바꾸어 정의한다. \(\mathfrak p\)는 \(k\)의 소아이디얼이며, 유한 개의 무한소에 대해서는 \(k\)와 그 켤레체들을 실수체로 확대한다.

실제로 이 명제에는 순환체의 노름정리가 들어 있다. 앞서 본 바에 따라 순환대수 \((\alpha,Z)\)에 대해서는 이 정리가 다음 명제와 동치이기 때문이다. \(a\)가 모든 유한소와 무한소에서 \(\mathfrak p\)-진 노름이면, \(a\)는 \(Z\)의 한 수의 노름이다. 또는 \(\mathfrak p\)-진 확대로 넘어가지 않고 말하면 다음과 같다. \emph{\(a\)가 \(k\)의 모든 소아이디얼 \(\mathfrak p\)에 대하여 노름잉여이고 몇 가지 부호조건도 만족하면, \(a\)는 \(Z\)의 한 수의 노름이다.} 이 마지막 형태가 바로 알려진 해석적 보조수단을 이용하여 유체론에서 증명되는 노름정리이다. 분할대수에 관한 일반정리는 이 순환적 특수 경우로부터 순전히 대수적·산술적 고찰로 증명할 수 있다.\footnote{R. Brauer, H. Hasse, E. Noether, \emph{Beweis eines Hauptsatzes in der Theorie der Algebren}, J. f. M. 167 (1932).} Hasse는 여기서 중요한 첫 귀결을 얻었다. \emph{대수적 수체 위의 모든 단순 정규대수는 순환대수이다.} 오래 추측되어 온 이 사실의 증명을 찾는 과정에서 위의 일반적 정식화가 생겨났다.

\paragraph{5.}
분할대수에 관한 정리의 두 번째 귀결은, 다시 순전히 대수적·산술적인 추론으로 얻어지는 처음의 \emph{주종정리}이다.\footnote{그 증명은 Math. Ann.에 실릴 예정이다.} 그 불변적 정식화는 교차곱을 정의하는 관계 (2)부터 (5)가 순전히 곱셈적이어서, \(K^\ast\)를 다음 조건만을 만족하는 아벨군 \(\mathfrak J\)로 바꾸어도 뜻을 유지한다는 사실에 바탕을 둔다. 즉 \(\mathfrak J\)의 자기동형군이 \(\Gg\)와 동형인 부분군 하나를 포함하면 된다. 이때 (1)은 군론에서 말하는 \glqq{}\(\mathfrak J\)에 의한 \(\Gg\)의 확장\grqq{}으로 바뀐다. \(\mathfrak J\)를 \(K\)의 모든 아이디얼로 이루어진 군으로 택하면 인자계도 아이디얼들의 계가 된다. \(\mathfrak J\)의 류분할은 인자계에 대한 류분할을 유도하는데, 일반적으로 이는 본래의 아이디얼류분할보다 더 세밀하다. \(u_S\)들의 곱셈이 절대 아이디얼류에 관하여 일의적이라는 요구는 변환량 \(c_S^{\,T}c_T/c_{ST}\), 곧 원소 인자계의 단위원류가 아이디얼 인자계의 단위원류에 속한다는 것만을 말하기 때문이다. 그러나 알려진 경우로 특수화해 보면 실제로는 이보다 조금 덜 세밀한 분할만으로 충분하다. 다음과 같이 정의한다. \emph{모든 유한 및 무한 분기소에서 분할되는 대수를 생성하는 원소 \(a_{S,T}\)들을 인자계의 단위원류에 넣는다.}

이렇게 생기는 \(\Gg\)의 확장을 \(\Gg^\ast\)로 쓰고, \(\mathfrak c\)의 절대 아이디얼류를 \((\mathfrak c)\)로 쓰자. 그러면 다음을 얻는다.

\emph{주종정리의 불변적 정식화: 위에서 정의한 류분할 아래에서 치환 \(v_S=u_S(\mathfrak c_S)\)가 \(\Gg^\ast\)의 자기동형사상을 만든다고 하자. 이러한 모든 \((\mathfrak c_S)\)가 주종을 이룬다. 그러면 그 자기동형사상은 내부자기동형사상이며, 하나의 아이디얼류 \((\mathfrak b)\)에 의해 생성된다.}

알려진 특수 경우들은 다음의 동치인, 조금 더 명시적인 정식화에서 따라 나온다. \emph{아이디얼류 \((\mathfrak c_S)\)로부터 만든 변환량 \((\mathfrak c_S^{\,T})(\mathfrak c_T)/(\mathfrak c_{ST})\)이 인자계의 단위원류에 속한다고 하자. 그러면 모든 \(S\in\Gg\)에 대하여 \((\mathfrak c_S)=(\mathfrak b)/(\mathfrak b^{S})\)가 되도록 하는 아이디얼류 \((\mathfrak b)\)가 존재한다. 즉 \((\mathfrak c_S)\)들은 기호적 \((1-S)\)제곱이다.} 실제로 가정은 바로 자기동형사상 성질을 표현하고, 그 자기동형사상이 내부적이라는 사실은
\[
v_S=(\mathfrak b)^{-1}u_S(\mathfrak b)
   =u_S(\mathfrak b)^{1-S}
   =u_S(c_S)
\]
로 표현된다.

순환적인 경우로 특수화하면, 정규화를 고려하여 다음을 얻는다. \(N([\mathfrak c])\)가 인자계의 단위원류에 속하면 \((\mathfrak c)\)는 기호적 \((1-S)\)제곱, 곧 \((\mathfrak c)=(\mathfrak b)^{1-S}\)가 된다.

\paragraph{6.}
여기서 순환체와 이차형식에 관한 알려진 형태로 나아가기 위해, 먼저 이 정리가 모든 아이디얼류에 대하여 진술되었지만 통상적인 방식대로 분기소와 서로소인 아이디얼들로 제한해도 그 내용은 같음을 말해 둔다.

그러면 여기서 정의한 주종은 이차체의 경우 가우스의 주종으로 바뀐다. 아이디얼류에는 이차형식이, 류의 노름에는 그 형식들이 나타내는 수가 대응하기 때문이다. 단위원류가 \(K\)의 분기소에서 분할되는 대수들을 생성한다는 것은, 이 표현되는 수들이 분기소에서 이차잉여가 된다는 뜻이다. 따라서 대응하는 형식들은 주형식의 총특성(Totalcharakter)을 가지며 가우스의 주종을 이룬다. 기호적 \((1-S)\)제곱이 제곱연산으로 바뀐다는 사실은 잘 알려져 있다.

순환체의 경우에는 다음 형태가 된다. 주종은 그 노름이 모든 유한 및 무한 분기소에서 노름잉여가 되는 아이디얼류 전체로 이루어진다. 이는 잘 알려져 있듯이 \glqq{}도체에 대한 노름잉여\grqq{}와 동치이므로 통상적인 정리가 특수 경우로 얻어진다.

그 밖에도 임의의 갈루아체라는 일반적인 경우에 분기소들로만 이루어진 도체를 도입할 수 있다. 이때 각 분기소에서 인자계를 정규화하면 그 레이는 단위원류의 원소들만을 포함한다.

여기서 2절의 개관에서 말한 Artin 도체들과의 관계가 문제로 생긴다. Artin 도체들도 같은 소아이디얼들로 이루어져 있기 때문이다. 더 나아가 두 번째 초복소적 방법인 갈루아 가군론과의 관계도 문제가 된다. 이 두 방법이 어디까지 미칠지는 앞으로 밝혀져야 한다.

\endgroup

\clearpage
\setcounter{footnote}{0}
\phantomsection
\addcontentsline{toc}{section}{40. 비가환 대수}
%
%
%
%
%
%

\providecommand{\frA}{\mathfrak A}
\providecommand{\frako}{\mathfrak o}
\providecommand{\frB}{\mathfrak B}
\providecommand{\frC}{\mathfrak C}
\providecommand{\frG}{\mathfrak G}
\providecommand{\frL}{\mathfrak L}
\providecommand{\frM}{\mathfrak M}
\providecommand{\frN}{\mathfrak N}
\providecommand{\frR}{\mathfrak R}
\providecommand{\frS}{\mathfrak S}
\providecommand{\frT}{\mathfrak T}
\providecommand{\Gcal}{\mathfrak G}
\providecommand{\mbarA}{\overline{\mathfrak A}}

\section*{40. 비가환 대수}

\begin{center}
\emph{Math. Zs. 37 (1933), 514--541쪽}
\end{center}

\begin{center}
\textbf{비가환 대수학}\\[0.5em]
괴팅겐의 에미 뇌터
\end{center}

가환대수학의 주요 정리들은 갈루아 이론에 들어 있음이 알려져 있다. 그에 앞서는 것은 분리체와 분해체, 곧 주어진 다항식이 일차인자를 하나 분리해 내거나 완전히 일차인자들로 분해되기에 충분한 체들의 이론이다.

여기서는 이에 대응하는 대수학의 부분들을 비가환적으로, 특히 초복소적인 경우에 전개한다. 원리적으로 비가환 방법, 곧 비가환체에서의 표현을 사용한다. 끝에서는 앞에서 말한 가환대수학의 정리들을 가환체에서의 표현을 사용하여 완전히 평행하게 정초할 수 있음을 보인다.

바탕이 되는 표현론 앞에는 짧은 자기동형사상 이론이 놓인다(§ 1). 이 표현론은 표현가군 이론에 바탕을 둔 기존 이론을 발전시킨 것이다\footnote{E. Noether, \emph{Hyperkomplexe Größen und Darstellungstheorie}, Math. Zeitschr. 30 (1929), 641--692쪽. 이하 \emph{표현론}으로 인용한다. van der Waerden, \emph{Moderne Algebra II}의 서술도 보라.}. 특히 직접 표현과 함께 반대 표현을 고찰한다. 한쪽은 다른 쪽으로 환원할 수 있으며, 반대 표현은 이제 반대 표현가군에 의해 생성된다. 장점은 이 반대 표현가군, 곧 쌍가군을 한 확대환 위의 한쪽 가군으로도 볼 수 있다는 데 있다(§ 2). 이로써 비가환체에서의 표현은 확대환의 아이디얼 이론으로 환원된다. 기약 표현류는 확대환의 기약 아이디얼류에 대응하며, 이는 초복소계를 그 가환 계수영역에서 표현할 때 그 계 자체에 성립하는 사실들과 정확히 유사하다(§§ 3, 4).

여기서부터 비가환체 위 행렬환의 구조정리들이 나온다. 핵심은 \emph{모든 부분환은 이 체에 의한 표현, 또는 반대동형인 체에 의한 반대 표현을 뜻한다}는 관찰이다(§ 5). 이 반대동형인 체는 가환인 경우의 최소 분리체·분해체에 대한 첫 유사물이다. 실제로 모든 1차 반대 표현을 매개하고, 중심에 관하여 유한 랭크이면 --- 지금까지는 이를 가정하지 않았다 --- 랭크 1인 직합성분들로의 완전한 분해도 매개한다. 이것이 체의 갈루아 이론을 주며(§ 6), 동시에 반대동형인 나눗셈대수들, 곧 중심에 관하여 유한 랭크인 체들이 R. 브라우어의 대수류군에서 서로 역인 류를 생성한다는 사실을 나타낸다(§ 7).

더 일반적으로 단순계에 성립하며 본문에서는 먼저 제시할 갈루아 이론의 둘째 정초는, 서로 원소별로 교환하는 부분환에 관한 한 정리의 거의 직접적인 귀결이다. 그 정리 자체도 앞서 둔 자기동형사상 고찰에 거의 바로 이어진다.

여기까지 고찰은 순전히 비가환적이다. 그러나 \emph{가환} 분해체의 문제도 § 5에서 먼저 말한 관찰에 따라 분해체를 나눗셈대수 안에 표현함으로써 비가환적으로 다룬다(§ 7). 마지막에는 앞에서 말한 가환적인 경우로의 이행(§ 8)과 임의의 계에 대한 분해체 이론(§ 9)을 제시한다. 이로써 가환체에서의 통상적인 표현론과의 관계도 얻는다.

R. 브라우어는 가환체에서의 이 표현론과 “무리” 인자계에 분해체 이론을 세웠다\footnote{우리의 공동 논문 \emph{Über minimale Zerfällungskörper irreduzibler Darstellungen}, Sitz. Ber. d. Preuß. Ak. d. Wiss. 1927, 221--228쪽을 보라. 222쪽의 주에는 “경위 보고”가 들어 있다. 주요 정리들은 서로 독립적으로 거의 동시에 생겼다. 또 R. Brauer, \emph{Über Systeme hyperkomplexer Zahlen}, § 3, Math. Zeitschr. 30 (1929), 79--107쪽을 보라. A. A. Albert는 뒤에 이 정리들을 독립적으로 다시 발견하였다.}. 이 가환적 정초 때문에 그는 --- 뒤의 Albert도 마찬가지로 --- 중심을 완전체라고 가정해야 했지만, 비가환적 정초에서는 이 제한이 필요 없다. 같은 제한 아래, 역시 가환체에서의 표현론을 사용하여, 나의 비가환체 갈루아 이론을 알게 된 뒤 R. Brauer와 K. Shoda는 이 이론을 더 발전시켰다. R. Brauer는 앞서 말한 교환 부분환 정리를 주었고, K. Shoda는 독립적으로 반단순계까지 포함하는 완전한 이론을 주었다\footnote{R. Brauer, \emph{Über die algebraische Struktur von Schiefkörpern}, § 2, Journ. f. Math. 166 (1932), 241--252쪽. K. Shoda, \emph{Über die galoissche Theorie der halbeinfachen hyperkomplexen Systeme}, Math. Ann. 107 (1932), 252--258쪽. J. Levitzki도 이 결과들을 전개하였다.}.

끝으로 체의 경우에 대해서는 van der Waerden II, § 128에 짧은 서술이 있음을 말해 둔다. 이는 내가 이 문제들을 처음 강의한 1928년 여름학기 강의를 따른 것이다. van der Waerden의 서술은 그 강의보다 여러 곳을 단순화하였고, 나는 그 일부는 두 번째 강의(1929/30년 겨울학기)\footnote{첫 강의는 G. Köthe가, 둘째 강의는 M. Deuring이 정리하였다.}에서, 일부는 비로소 여기에서 받아들여 더 전개하였다. 특히 갈루아 이론의 초복소적 증명법을 가환적인 경우로 옮긴 § 8은 둘째 강의에서 처음 나왔다. 그때 비가환적 증명(§ 6, 3)은 첫 강의보다 본질적으로 단순화되어 있었다. 교환 부분환 정리(§ 5)와 이에 바탕을 둔 비가환 갈루아 이론의 둘째 증명법(§ 6, 1과 2)은 둘째 강의 뒤 R. Brauer의 주 3 논문과 van der Waerden § 128을 알게 된 다음에야 더해졌다. 다만 § 5, 3에서는 거기보다 훨씬 약한 유한성 가정만 쓴다.

첫 강의의 어떤 논법들, 곧 차아이디얼의 교집합을 만드는 방법은 더 복잡하지만 무한 랭크인 계와 정수 계수 계에도 옮길 수 있다는 장점이 있다\footnote{이 방법은 본질적으로 G. Köthe, \emph{Schiefkörper unendlichen Ranges über dem Zentrum}, § 5, Math. Ann. 105 (1931), 15--39쪽에 재현되어 있다. 주 10도 보라. 이 교집합 방법은 이제 양쪽 아이디얼들이 불변가군에 속한다는 거의 자명한 사실(§ 4, 1)로 대체되었다. 이 사실을 처음 알아차린 이는 van der Waerden이다. 이로써 모든 것을 더 단순한 직합분해로 정초할 수 있지만, 무한 랭크인 경우와 정수 계수인 경우에는 이 방법이 성립하지 않는다.}. 가환 정수 계수 계에서는 이로부터 아이디얼 미분과 디퍼런트의 관계에 이른다\footnote{프라하 강연, Jahresber. d. Deutsch. Math. Ver. 39 (1930), 17쪽(비스듬한 쪽매김)을 보라.}. 이에 대해서는 다른 기회에 더 상세히 다루겠다.

분해체 이론의 주된 응용은 교차곱과 그 인자계의 이론에 있으며, 이 이론은 다시 수론적 응용의 기초가 된다. 여기서는 더 다루지 않는다\footnote{교차곱의 이론은 둘째 강의에서 전개하였다. 여기서 준 이론을 전제하지 않도록 조금 고친 꼴이 H. Hasse, \emph{Theory of cyclic algebras over an algebraic numberfield}, 제II장, Transactions of the Amer. Math. Soc. 34 (1932), 171--214쪽에 재현되어 있다. 강의에 더 밀접한 서술은 M. Deuring의 \emph{Ergebnisse der Mathematik} 보고서에 실릴 예정이다.}.

\begin{center}
\textbf{§ 1}\\[0.4em]
\textbf{자기동형사상, 가군, 쌍가군}
\end{center}

표현가군에 바탕을 둔 표현론은 연산자를 갖거나 갖지 않는 아벨군의 자기동형환 이론에 기초한다(§ 2도 보라). 여기서 암묵적으로 쓰는 논법, 곧 사상과 연산법칙 사이의 관계를 반복하지 않도록 몇 가지 간단한 정리로 정식화한다.

\noindent\textbf{1. 곱셈적 사상. 결합법칙.} 먼저 \(\frG\)를 연산자가 없는 군이라 하고, \(\frA\)를 그 절대 자기동형영역, 곧 \(\frG\)에서 자기 자신으로 가는 모든 준동형의 계라 하자. 곱 \(\sigma\tau\)를
\[
g(\sigma\tau)=(g\sigma)\tau
\quad\text{(\(g\in\frG,\ \sigma,\tau\in\frA\))} \tag{1}
\]
로 정의하면 \(\frA\)는 곱셈에 닫혀 있고 결합법칙을 만족한다.

\(\frG\)에 기호 \(\Theta,H,\ldots\)의 집합 \(\frB\)를 주어 연산자 있는 군으로 만든다는 것은, 연산 \(g\Theta,gH,\ldots\)가 \(\frG\)의 일의적으로 정의된 원소를 주고 \(\frG\)의 자기준동형을 생성한다는 뜻이다. 즉 \((g\cdot h)\Theta=g\Theta\cdot h\Theta\)이다. 따라서 \(\frB\)에서 절대 자기동형영역 \(\frA\)의 한 부분집합 \(\overline{\frB}\)로 가는 일의적 사상, 일반적으로는 일대일일 필요가 없는 사상이 존재한다. 앞서 말한 논법은 여기서 다음과 같다.

\emph{연산자영역 \(\frB\)가 곱셈에 닫혀 있으면, \(\frB\)에서 \(\overline{\frB}\)로 가는 일의적 사상이 곱셈적 준동형인 것은 (1)에 대응하는 결합관계}
\[
g(\Theta H)=(g\Theta)H
\quad\text{(\(g\in\frG,\ \Theta,H\in\frB\))} \tag{1a}
\]
\emph{가 성립할 때 그리고 오직 그때이다. 왼쪽 연산자를 다루는 경우에는 이에 대응하여 반대 준동형을 정의한다.}

실제로 \(\frB\)에서 \(\overline{\frB}\)로 가는 사상은 모든 \(g\in\frG\)에 대한 \(g\Theta=g\sigma\)로 주어지므로
\[
(g\Theta)H=(g\sigma)\tau=g(\sigma\tau)
\]
이다. 따라서 (1a)는 곱셈적 준동형의 필요충분조건이다. 왼쪽 연산자가 반대 준동형을 만드는 까닭은 왼쪽에 쓴 자기동형사상 \(\sigma^\ast g,\tau^\ast\sigma^\ast g\)는 오른쪽에서 왼쪽으로 읽어야 하기 때문이다. 곧 \(\tau^\ast\sigma^\ast g=g\sigma\tau\)인 반면, 연산자는 언제나 왼쪽에서 오른쪽으로 읽는다.

\noindent\textbf{2. 연산자준동형 사상.} \(\frG\)가 연산자 있는 군이면 연산자준동형은 잘 알려진 대로
\[
\begin{array}{r@{\quad}c@{\quad}c@{\quad}c@{\qquad}c@{\qquad}r@{\quad}c@{\quad}c@{\quad}c}
(2)&(g\Theta)\sigma&=&(g\sigma)\Theta
&\text{또는}&
(2*)&(\Theta g)\sigma&=&\Theta(g\sigma)
\end{array}
\]
로, 곧 교환적 결합 또는 이어지는 결합법칙으로 정의한다. 두 연산자영역에 1의 논법, 곧 자기동형영역으로의 사상을 적용하면 다음을 얻는다.

\emph{\(\Theta,H,\ldots\)를 원소로 갖는 연산자영역 \(\frB\)와 \(\bar\Theta,\bar H,\ldots\)를 원소로 갖는 연산자영역 \(\frC\)가 주어졌다고 하자. \(\frC\)의 원소들이 \(\frB\)-자기동형사상을 생성하고 동시에 \(\frB\)의 원소들이 \(\frC\)-자기동형사상을 생성하는 것은, 두 영역이 \(\frG\)와 교환적으로 결합되거나 이어지는 결합법칙}
\[
\begin{array}{r@{\quad}c@{\quad}c@{\quad}c@{\qquad}r@{\quad}c@{\quad}c@{\quad}c}
(2a)&(g\Theta)\bar H&=&(g\bar H)\Theta,&
(2a*)&(\Theta g)\bar H&=&\Theta(g\bar H)
\end{array}
\]
\emph{이 성립할 때 그리고 오직 그때이다.}

\noindent\textbf{3. 환 위의 가군과 쌍가군.} 환 \(\frR\) 위의 \emph{오른쪽 가군} \(\frM\)은 \(\frR\)의 원소들을 오른쪽 연산자로 갖는 덧셈 아벨군으로 정의한다. 이때 결합관계 (1a), 연산자를 정의하는 관계 \((g+h)\Theta=g\Theta+h\Theta\)와 함께 분배관계
\[
\begin{array}{r@{\quad}c@{\quad}c@{\quad}c}
(3a)&g(\Theta+H)&=&g\Theta+gH
\end{array}
\]
가 성립해야 한다. 왼쪽 가군도 이에 대응하여 정의한다.

\emph{쌍가군}에는 두 종류가 있다. 두 환 \(\frR,\frS\) 위의 오른쪽 가군을 쌍가군이라 하는 것은, \(\frR,\frS\) 각각에 대한 관계 (1a), (3a)에 더하여 \(\frR,\frS\)가 (2a)에 따라 \(\frM\)과 교환적으로 결합될 때이다. \(\frR\)-왼쪽·\(\frS\)-오른쪽 가군은 다른 관계들에 (2a*), 곧 이어지는 결합법칙이 더해질 때 쌍가군이라 한다.

이 관계들의 의미는 \emph{아벨}군의 경우 절대 자기동형영역이 하나의 환, 곧 절대 자기동형환을 이룬다는 데서 나온다\footnote{비아벨군에서는 “일반화된” 환을 다룬다. 이에 대해서는 H. Fitting의 Math. Annalen 게재 예정 논문을 보라[뒤에 Math. Annalen 107, 514--542쪽에 실림].}. 구체적으로 사상에 관한 위 논법은 다음을 준다.

\emph{덧셈으로 쓴 아벨군 \(\frM\)의 연산자영역 \(\frR\)이 환이면, \(\frR\)에서 절대 자기동형환의 부분집합 \(\overline{\frR}\)으로 가는 일의적 사상이 환 준동형인 것은 \(\frM\)이 \(\frR\) 위 오른쪽 가군, 곧 (1a), (3a)를 만족할 때 그리고 오직 그때이다. 왼쪽 가군이면 반대 환 준동형이다.}

\emph{\(\frM\)이 환 \(\frR,\frS\) 위 오른쪽 가군이면, \(\frM\)이 쌍가군, 곧 (2a)를 만족하는 것은 \(\frR\)을 \(\frS\)-자기동형환의 부분환에, 동시에 \(\frS\)를 \(\frR\)-자기동형환의 부분환에 환 준동형으로 보낼 수 있을 때 그리고 오직 그때이다. \(\frM\)이 \(\frS\)-오른쪽 가군이자 \(\frR\)-왼쪽 가군이면, 쌍가군 조건 (2a*)는 \(\frS\)에서 \(\frR\)-자기동형환의 부분환으로 가는 준동형과 \(\frR\)에서 \(\frS\)-자기동형환의 부분환으로 가는 반대 준동형을 뜻한다.}

마지막 명제는 알려진 귀결들을 다음처럼 정밀화한다.

\emph{1) \(\frR\)이 단위원을 갖는 환이면, \(\frR\)은 \(\frR\)-왼쪽 가군으로서의 자기동형환과 직접 동형이고, \(\frR\)-오른쪽 가군으로서의 자기동형환과 반대동형이다.} 실제로 이어지는 결합법칙 (2a*)가 성립하므로 준동형 명제가 적용된다. 단위원의 존재 때문에 모든 자기동형사상은 \(\frR\)의 어떤 \(a\)에 대한 \(e\to a\) 꼴이다. 따라서 동형이 성립하고 \(\frR\)이 \(\frR\)-자기동형환 전체를 채운다.

\emph{2) \(\frR\)이 일반적으로 비가환인 체 \(\mathsf{A}\) 위의 완전행렬환}
\[
\frR=\sum c_{ik}\mathsf{A}=\sum \mathsf{A} c_{ik}
\]
\emph{이면, \(\mathsf{A}\)는 단순 오른쪽 아이디얼들의 자기동형체와 반대동형이고 단순 왼쪽 아이디얼들의 자기동형체와 직접 동형이다.} 실제로 각 단순 오른쪽 아이디얼은 연산자동형인 \(\mathsf{A}\)-왼쪽·\(\frR\)-오른쪽 쌍가군이 되고, \(\mathsf{A}\)가 \(\frR\)-자기동형사상 전체를 채운다. 왼쪽 아이디얼도 이에 대응한다.

\noindent\textbf{4. 오른쪽 쌍가군에서 곱환 위 한쪽 가군으로 이행.} 오른쪽 쌍가군의 중요성은 본질적으로 이 이행에 바탕을 둔다.

\noindent\textbf{정리.} \emph{\(\frM\)이 두 부분환 \(\frR,\frS\)를 포함하고 이들이 원소별로 서로 교환하는 한 환 \(\frT\) 위 오른쪽 가군이면, \(\frM\)을 \(\frR,\frS\) 위 쌍가군으로도 볼 수 있다. 역으로 \(\frR,\frS\) 위 쌍가군이 주어지고, \(\frR,\frS\)가 원소별로 교환하는 곱환 \(\frT\)가 존재하면, \(\frM\)을 \(\frT\) 위 오른쪽 가군으로도 볼 수 있다. 이때 \(\frT\)에 의한 연산은 주어진 \(\frR,\frS\)에 의한 연산의 연장이다.}

첫 명제는 \(\frT\)가 정의에 따라 절대 자기동형환의 한 부분환 \(\overline{\frT}\)에 준동형으로 사상되고, \(\frR,\frS\)가 원소별로 교환하는 부분환 \(\overline{\frR},\overline{\frS}\)로 옮겨간다는 데서 나온다. 이는 바로 쌍가군을 특징짓는 관계 (2a)이다.

역으로 \(\frM\)이 \(\frR,\frS\)-쌍가군, 더 정확히는 양쪽 모두에 대한 오른쪽 가군으로 주어지면, \(\frR,\frS\)는 절대 자기동형환의 원소별로 교환하는 부분환 \(\overline{\frR},\overline{\frS}\)에 준동형으로 사상된다. 절대 자기동형환 안에는 \(\overline{\frR},\overline{\frS}\)의 곱환 \(\overline{\frT}\)가 존재한다\footnote{여기서 곱환은 공통 상환을 갖는 두 환이 생성하는 환일 뿐이며, 곱이 직합일 필요는 없다. 주어진 교집합에 대해 곱환이 언제나 존재하는 것은 아님을 다음 예가 보인다. \(\frako\)를 정수환이라 하고 \(\frR=\frako[x]\), \(\frS=\frako[y]\), \(2x-1=0\), \(2y-1=0\)이라 하자. \(x,y\) 사이에 연산관계를 정하지 않으면 \(\frako\)가 \(\frR,\frS\)의 교집합이다. 그러나 \(\frR,\frS\)가 한 공통 상환에 들어 있으면 \((2x-1)y-x(2y-1)=0\), 따라서 \(x=y\)이다. 그러므로 교집합이 \(\frako\)인 곱환은 없다.}. 원소별 교환성 때문에 이 환은
\[
\sum \bar r_i\bar s_i+\bar r+\bar s
\]
꼴의 모든 원소로 이루어진다. \(\overline{\frR},\overline{\frS}\)가 단위원을 가지면 덧붙인 항 \(\bar r,\bar s\)는 필요 없다. 이제 \(\frR,\frS\)가 원소별로 교환하는 곱환 \(\frT\)도 존재하면, 이에 따라 \(\frT\)는 \(\sum r_i s_i+r+s\) 꼴의 모든 원소로 이루어진다. 준동형
\[
\frR\to\overline{\frR},\qquad
\frS\to\overline{\frS}
\]
는 대응
\[
\sum r_i s_i+r+s
\to
\sum\bar r_i\bar s_i+\bar r+\bar s
\]
에 의해 준동형 \(\frT\to\overline{\frT}\)로 연장된다. 따라서 연산
\[
m\Bigl(\sum r_i s_i+r+s\Bigr)
=\sum(mr_i)s_i+mr+ms
\]
는 일의적이며 § 1, 3에 따라 주어진 \(\frR,\frS\)-가군을 포함하는 \(\frT\)-가군을 정의한다.

\begin{center}
\textbf{§ 2}\\[0.4em]
\textbf{반대 표현과 직접 표현}
\end{center}

\noindent\textbf{1. 선형형식 가군.} § 1에서 전개한 이론에서 표현론으로 넘어가는 과정은 거기의 가군들을 선형형식 가군으로 특수화하고 그 자기동형환을 고찰하는 데 바탕을 둔다.

\(\frS\)가 단위원을 갖는 환일 때 \(\frS\)-오른쪽 가군 \(\frM\)을 \(\frS\) 안의 선형형식 가군이라 하는 것은, \(\frM\)이 \(n\)개의 한 항짜리 \(\frS\)-가군의 직합
\[
\frM=m_1\frS+\ldots+m_n\frS
\]
이고 각 \(m_i\frS\)가 \(\frS\)와 연산자동형, 곧 \(m_i s=0\)이면 언제나 \(s=0\)인 경우이다. 이때 \(\frS\)는 절대 자기동형환의 한 부분환에 단지 준동형일 뿐 아니라 동형으로 사상된다.

\noindent\textbf{정리 1.} \emph{\(\frM\)이 \(\frS\) 안의 선형형식 가군이면, \(\frM\)의 \(\frS\)-자기동형환 \(\mA\)는 \(\frS\)의 원소를 성분으로 하는 모든 \(n\)행 행렬의 환 \(\mbarA\)와 단지 준동형일 뿐 아니라 반대동형이다.}

\(\frM\)의 임의의 \(\frS\)-자기동형사상 \(\alpha\)는 대응 \(m_i\to\overline m_i=m_i\alpha\)로 완전히 기술된다. 정의에 따라 이로부터
\[
\sum m_i s_i\to
\sum\overline m_i s_i
=\sum(m_i\alpha)s_i
=\sum(m_i s_i)\alpha
\]
가 나오기 때문이다. 행렬표기로
\[
(m_1\alpha,\ldots,m_n\alpha)=(m_1,\ldots,m_n)A
\]
라 쓰면, \(\alpha\)가 모든 자기동형사상을 두루 돌 때 대응 \(\alpha\to A\)는 \(\mA\)와 \(\mbarA\) 사이의 일대일 대응을 준다. 실제로 \(\frM\)을 선형형식 가군이라고 가정했으므로 \(A\)는 \(\alpha\)로 유일하게 정해지고, \(\frS\)의 원소를 성분으로 하는 모든 \(n\)차 행렬 \(A\)는 자기동형사상을 생성한다. 또한
\[
\alpha+\beta\to A+B,\qquad
\alpha\beta\to BA
\]
이다. 마지막 것은
\[
(m_1\alpha\beta,\ldots,m_n\alpha\beta)
=(m_1,\ldots,m_n)A\beta
=(m_1,\ldots,m_n)\beta A
\quad[\text{\(ms\beta=m\beta s\)이므로}]
=(m_1,\ldots,m_n)BA
\]
에서 나온다.

\noindent\textbf{2. 표현과 표현가군.} 환 \(\frR\)의 \(\frS\) 안 \(n\)차 반대 표현 또는 직접 표현을, \(\frR\)에서 \(\frS\)의 원소를 성분으로 하는 모든 \(n\)행 행렬의 환의 한 부분환으로 가는 반대 환 준동형 또는 직접 환 준동형으로 이해하면, 정리 1은 다음과 같이도 쓸 수 있다.

\noindent\textbf{정리 1'.} \emph{\(\frS\) 안 \(n\)차 선형형식 가군의 자기동형환 \(\mA\)는 \(\frS\)의 모든 \(n\)행 행렬로 이루어진 완전행렬환 \(\mbarA\)에 의한 동형인, 곧 충실한 반대 표현을 갖는다.}

여기서부터 직접 표현에 관한 통상적인 정리들이 반대 표현에도 성립한다.

\noindent\textbf{정의.} \(\frS\) 안의 선형형식 가군 \(\frM\)이 \(\frR,\frS\) 위 쌍가군이면서 동시에 \(\frR\)-오른쪽 가군이자 \(\frS\)-오른쪽 가군이면 \(\frR\)의 \(\frS\) 안 반대 표현가군이라 한다. \(\frM\)이 쌍가군이면서 \(\frR\)-왼쪽·\(\frS\)-오른쪽 가군이면 직접 표현가군이라 한다.

\noindent\textbf{정리 2.} \emph{각 반대 표현가군 또는 직접 표현가군은 \(\frR\)의 \(\frS\) 안 동치인 반대 표현 또는 직접 표현들의 한 류를 생성하며, 모든 반대 표현과 직접 표현은 이렇게 생성된다.}

먼저 \(\frM\)이 반대 표현가군이라 하자. § 1, 3에 따라 \(\frR\)은 \(\frM\)의 \(\frS\)-자기동형환의 한 부분환 \(\overline{\frR}\)에 직접 준동형으로 사상된다. § 2, 정리 1'에 따라 \(\overline{\frR}\)은 충실한 반대 표현을 가지며, 이는 \(\frR\)의 반대 표현이기도 하다. 역으로 반대 표현 \(\frR\to\frR^*\)가 주어지면, 이를 \(\frR^*\)에서 \(\frS\)-자기동형환의 부분환 \(\overline{\frR}\)으로 가는 반대동형과 합성하여 \(\frR\)에서 \(\overline{\frR}\)으로 가는 직접 준동형을 얻는다. 따라서 \(\frM\)은 \(\frR\)-가군이 되고, § 1, 3에 따라 쌍가군, 곧 표현가군이 된다. 다른 \(\frS\)-기저로 바꾸면 동치인 표현들의 류가 나온다.

직접 표현가군의 경우에는 \(\frR\)이 \(\overline{\frR}\)에 반대 준동형으로 사상된다. 이를 \(\overline{\frR}\)의 반대 표현과 합성하면 \(\frR\)의 직접 표현을 얻는다. 역으로 직접 표현에서 \(\overline{\frR}\)으로 가는 반대 준동형 사상이 나오며, § 1, 3에 따라 \(\frM\)은 직접 표현가군이 된다.

\noindent\textbf{주의.} 물론 두 종류의 표현은 서로 환원할 수 있다. \(\frR\)의 직접 표현은 \(\frR\)의 반대환의 반대 표현이고 그 역도 성립한다. 다른 환원은 \(\frS\)에서 반대환으로 옮기고 행렬들을 전치행렬로 바꾸는 것이다. 이는 \(\frR\)-왼쪽·\(\frS\)-오른쪽 가군에서 \(\frR,\frS\)-왼쪽 가군으로 옮기는 데 대응한다.

\noindent\textbf{3. 반대 표현가군에서 확대환 위 가군으로 이행.} 반대 표현가군의 장점은 무엇보다 § 1, 4에 따라 확대환으로 이행할 수 있다는 데 있다. 지금의 특수한 경우, 곧 \(\frM\)이 \(\frS\) 위 선형형식 가군인 경우에는 표현가군에 대한 다음 이행정리가 된다.

\emph{\(\frM\)이 두 부분환 \(\frR,\frS\)를 포함하고 이들이 원소별로 서로 교환하는 환 \(\frT\) 위 오른쪽 가군이며, \(\frM\)이 \(\frS\) 위 선형형식 가군이면, \(\frM\)을 \(\frR\)의 \(\frS\) 안 반대 표현가군으로도 볼 수 있다.}

\emph{역으로 \(\frR\)의 \(\frS\) 안 반대 표현가군이 주어지고, \(\frR,\frS\)가 원소별로 교환하는 곱환 \(\frT\)가 존재하면, \(\frM\)을 \(\frT\)-가군으로도 볼 수 있다. 이때 \(\frT\)-연산은 주어진 \(\frR,\frS\)-연산의 연장이다.}

표현론에서 본질적으로 쓰는 것은 다음 귀결이다.

\noindent\textbf{표현가군 이행정리의 귀결.}\\*
\(\frM\)이 \(\frR\)의 \(\frS\) 안 반대 표현가군이면서 곱환 \(\frT\) 위 가군이기도 하면, \(\frM\)에서 한 가군 \(\frN\)으로 가는 \(\frR,\frS\)-동형은 \(\frT\)-동형과 같다. \emph{따라서 동형인 \(\frT\)-가군들의 각 류에는 \(\frR\)의 \(\frS\) 안 반대 표현류 하나가 대응하고 그 역도 성립한다.}

이 정리들을 § 1, 2와 3에 결합하면 한 표현과 교환하는 행렬들과 \(\frR,\frS\)-자기동형사상들 사이의, 이하에서 기초가 되는 관계를 얻는다.

\noindent\textbf{교환행렬에 관한 정리.}\\*
\emph{\(\frR\to\frR^*\)가 \(\frR\)의 \(\frS\) 안 \(n\)차 반대 표현이고, \(\frB^*\)가 \(\frR^*\)의 모든 원소와 교환하는 \(\frS\) 안 \(n\)차 행렬들의 환이면, \(\frB^*\)는 이 표현을 생성하는 표현가군의 \(\frR,\frS\)-자기동형사상들, 곧 동시에 \(\frR\)-자기동형사상인 \(\frS\)-자기동형사상들의 반대동형 표현을 준다. 곱환 \(\frT\)가 존재하면 이는 \(\frT\)-자기동형사상들의 표현이기도 하다.}

\(\frB\)를 \(\frM\)의 모든 \(\frR,\frS\)-자기동형사상들의 환이라 하자. \(\frS\)-자기동형환 \(\mA\)의 부분환으로서 \(\frB\)는 \(\frS\) 안 반대동형 표현을 갖는다. § 1, 2와 3에 따라 \(\frB\)는 \(\mA\)의 자기동형사상 가운데 \(\frR\)과 교환적으로 결합된 것, 곧 관계 (2a)를 만족하는 것들로 이루어진다. 따라서 \(\frB\)는 \(\overline{\frR}\)의 각 원소와 교환하는 자기동형사상 전체이며, 여기서 \(\overline{\frR}\)은 \(\mA\) 안의 \(\frR\)의 상이다. \(\frB,\overline{\frR}\)의 표현은 단순한 반대 준동형이 아니라 반대동형이므로 이것이 증명할 사실이다.

\noindent\textbf{주의.} \(\frR,\frS\)-자기동형사상은 대응 \(m_i\to m_i'\)를 뜻하며, \(m_i'\)들에 의해 같은 표현 \(\frR^*\)이 생긴다. 다만 \(\frB^*\)의 행렬들이 반드시 가역은 아니라는 사실에 대응하여, \(m_i'\)들이 반드시 \(\frM\)의 \(\frS\)-기저를 이루지는 않는다. 합 \(\sum m_i'\frS\)가 더 이상 직합이 아닐 때 “표현”은 확장된 뜻으로 이해해야 한다.
\[
(m_1',m_2',\ldots m_n')a=(m_1',\ldots m_n')A
\]
는 여전히 맞지만 \(A\)는 더 이상 \(a\)로 유일하게 정해지지 않는다.

이행정리에 관해 한 가지를 더 지적한다. 대각행렬 \(E\cdot s\)는 \(\frS\)의 \emph{직접 동형} 표현, 곧 \(\frS\)의 “항등” 표현을 준다. 이것이 직접 동형인 까닭은 \(\frM\)의 \(\frS\)-자기동형환이 \(\frS\)와 반대동형인 부분환을 가지며, 지금의 표현은 그 부분환의 반대 표현이기 때문이다. 실제로 \(\frM\)의 각 한 항짜리 성분 \(m_i\frS\)는 오른쪽 가군으로서 \(\frS\)와 연산자동형이고, 따라서 그 \(\frS\)-자기동형환은 \(\frS\)와 반대동형이다(§ 1, 3, 귀결 1).

\(\frR,\frS\)로 일의적으로 정의되는, \(\frT\)에서 \(\frS\) 안의 행렬들로 가는 대응은 \emph{환 \(\frT\)의 표현이 아니라, \(\frR\)의 반대 표현을 \(\frS\)에 관하여 연산자준동형인 사상으로 확대한 것}이다.
\[
\sum r_i s_i\to\sum R_i s_i.
\]
특히 \(\frR\)의 반대 표현은 동시에 \emph{\(\frR\)의 중심에 놓이는 교집합 \([\frR,\frS]\)에 관하여 연산자준동형}이다.

\begin{center}
\textbf{§ 3}\\[0.4em]
\textbf{한 체에 관한 가군}
\end{center}

이하에서는 일반적으로 비가환인 체 안의 표현만 다룰 것이므로, 한 체 위 선형형식 가군에 관한 몇 가지 간단한 정리를 모아 둔다.

\noindent\textbf{1. 전체 가군의 주어진 기저에 관한 부분가군의 정규기저.} \(\mathsf{A}\)를 체, \(\frN=x_1\mathsf{A}+\ldots+x_n\mathsf{A}\)를 랭크 \(n\)인 선형형식 가군이라 하고, \(\frL=z_1\mathsf{A}+\ldots+z_l \mathsf{A}\)를 랭크 \(l\leqq n\)인 부분가군이라 하자. \emph{\(z_i\)들이 다음 꼴이면 \(x\)에 관한 정규기저라 한다.}
\[
z_i=x_i-(x_{l+1}\alpha_{i,l+1}+\ldots+x_n\alpha_{i,n}),
\qquad (i=1,\ldots,l).
\]
\emph{모든 가군 \(\frL\)은 \(x\)들의 번호를 알맞게 바꾸면 정규기저를 갖는다.}

실제로 번호를 알맞게 바꾸면
\[
\frN=\frL+x_{l+1}\mathsf{A}+\ldots+x_n\mathsf{A}
\]
이고, 따라서
\[
x_i\equiv x_{l+1}\alpha_{i,l+1}+\ldots+x_n\alpha_{i,n}
\quad(\frL),\qquad(i=1,\ldots,l). \tag{2}
\]
그러므로 \(l\)개 원소 \(z_i=x_i-(x_{l+1}\alpha_{i,l+1}+\ldots)\)는 \(\frL\)에 속한다. 이들은 \(x_{l+1},\ldots,x_n\)과 함께 \(\frN\)의 기저를 이루므로 \(\frL\) 전체를 생성한다.

\noindent\textbf{2. 확대가군.} 다시 \(A\)를 체, \(P\)를 부분체라 하자. 랭크 \(n\)인 \(A\)-가군 \(\frN\)이 한 \(\mathsf{P}\)-가군 \(\frM\)의 \emph{확대가군}, \(\frN=\frM_{\mathsf{A}}\), 이라는 것은 \(\frM\)이 \(\frN\)의 \emph{같은 랭크}를 갖는 부분가군인 경우이다. \(\frM=y_1\mathsf{P}+\ldots+y_n\mathsf{P}\)이면
\[
\frN=\frM_A=y_1A+\ldots+y_nA.
\]
역으로 각 \(\mathsf{P}\)-가군 \(\frM=y_1\mathsf{P}+\ldots+y_n\mathsf{P}\)에는 가군 동형까지 유일한 확대가군 \(\frM_{\mathsf{A}}\)가 존재한다. 실제로 가군 \(\bar y_1A+\ldots+\bar y_nA\)는 \(\frM\)과 동형인 부분가군 \(\overline{\frM}\)을 포함하며, 이를 \(\frM\)과 동일시할 수 있다.

\noindent\textbf{귀결 a).} \(z_1,\ldots,z_t\)가 \(\frM\)의 \(P\)에 관하여 선형독립인 원소들이면, \(\frN=\frM_{\mathsf{A}}\)에서도 \(A\)에 관하여 선형독립이다. 실제로 이들을 \(\frM\), 따라서 \(\frN\)의 기저로 보충할 수 있다. 그러므로 \(\frM\)의 각 \(P\)-가군 \(\frT=z_1\mathsf{P}+\ldots+z_t\mathsf{P}\)는 \(\frM_{\mathsf{A}}\) 안에서 확대가군 \(\frT_{\mathsf{A}}=z_1\mathsf{A}+\ldots+z_t\mathsf{A}\)를 생성한다.

\noindent\textbf{귀결 b).} \(\frM\)의 각 \(P\)-가군 \(\frT=z_1\mathsf{P}+\ldots+z_t\mathsf{P}\)는 수축가군, 곧 그 확대가군 \(\frT_{\mathsf{A}}\)와 \(\frM\)의 교집합이다.
\[
\frT=\frT_A\cap\frM.
\]
실제로 \(\frT\leqq \frT_{\mathsf{A}}\cap\frM\)이다. 역으로 \(a\in\frT_A\cap\frM\)이면 \(a=\sum z_i\alpha_i\)이므로 \(\frM\)의 원소 \(a,z_1,\ldots,z_t\)는 \(\mathsf{A}\)에 관하여 종속이고, 귀결 a에 따라 \(\mathsf{P}\)에 관해서도 종속이다. 따라서 \(a\in\frT\)이다.

\noindent\textbf{3. 불변가군에 관한 정리.} 다시 \(\frN=\frM_{\mathsf{A}}\)를 랭크 \(n\)인 \(P\)-가군 \(\frM\)의 확대가군이라 하자. 또 \(P\)가 \(A\)의 환 자기동형사상들로 이루어진 한 군 \(\frG\)에 대한 완전 불변체라 하자. \(\frG\)는 다음 규정으로 \(\frN=\frM_{\mathsf{A}}\)의 연산자영역이 된다.
\[
G\cdot m=m\quad(m\in\frM,\ G\in\frG),
\]
\[
G\cdot\sum m_i\alpha_i=\sum m_i\,G(\alpha_i)
\quad(m_i\in\frM,\ \alpha_i\in A).
\]

\noindent\textbf{보조정리.} \emph{이 규정 아래 \(\frM\)은 \(\frN\)의 원소 가운데 \(\frG\)에 의해 고정되는 모든 원소로 이루어진다.} 실제로 \(x_1,\ldots,x_n\)이 \(\frM\)의 \(P\)-기저이고 \(w=x_1\alpha_1+\ldots+x_n\alpha_n\)이면, 모든 \(G\in\frG\)에 대한 \(G(w)=w\)에서 \(G(\alpha_i)=\alpha_i\)가 나오고, 가정에 따라 \(\alpha_i\in P\)이다.

\noindent\textbf{정리.} \emph{\(\frG\)를 연산자영역으로 볼 때 허용되는 부분가군은 \(\frM\)의 부분가군 \(\frL\)들의 확대가군 \(\frL_{\mathsf{A}}\)와 오직 그것들이다.}

한쪽 방향은 분명하다. 역으로 \(\frL\)이 \(\frG\)에 관하여 허용된다고 하자. \(\frL\)의 \(x\)에 관한 정규기저를 \(z_1,\ldots,z_l\)이라 하자(1을 보라). 여기서 \(x\)는 \(\frM\)의 \(P\)-기저로 택하여 \(\frG\)가 원소별로 고정한다. 가정에 따라
\[
G(z_i)=x_i-\bigl(x_{l+1}G(\alpha_{i,l+1})
              +\ldots+x_nG(\alpha_{i,n})\bigr)
\]
는 모든 \(G\in\frG\)에 대해 \(\frL\)의 원소이다. 따라서 고정한 \(G\)에 대하여 \(z\)들의 선형결합으로 나타난다. \(x_1,\ldots,x_l\)의 계수를 비교하면 모든 \(G\in\frG\)에 대해 \(G(z_i)=z_i\)이다. 보조정리에 따라 \(z_i\in\frM\)이고, \(\frL\)은 \(\frM\)의 \(P\)-가군 \(\frT=z_1\mathsf{P}+\ldots+z_l \mathsf{P}\)의 \emph{확대}이다. 2에 따라 \(\frT=\frL\cap\frM\)이다\footnote{이 절의 정리들은 간단한 정렬정리 논법에 의해 \(\frM\)의 \(\mathsf{P}\)-랭크가 무한인 경우에도 유지된다. G. Köthe, \emph{Ein Beitrag zur Theorie der kommutativen Ringe ohne Endlichkeitsvoraussetzung}, § 1, Gött. Nachr. 1931, 195--207쪽을 보라.}.

\begin{center}
\textbf{§ 4}\\[0.5em]
\textbf{초복소계와 그 표현류}
\end{center}

\noindent\textbf{1. 확대정리.} 이제 § 2의 표현정리들을 § 3의 가군정리들과 결합한다. 일반 환 \(\frR\) 대신 가환체 \(P\) 위에서 단위원을 갖는 \emph{초복소계} \(S,T\ldots\), 곧
\[
S=x_1\mathsf{P}+\ldots+x_n\mathsf{P}
\]
를 고찰한다. 임의의 표현환 \(\frS\) 대신에는 체 \(\mathsf{A},\mathsf{B}\ldots\)만 고찰하며, 달리 말하지 않으면 그 중심은 \(\mathsf{P}\)라고 한다. 이 경우 \(S,A\)가 원소별로 교환하고 교집합이 \(P\)인 곱환이 언제나 존재한다. 이는 직접곱 \(S\times \mathsf{A}\)이며 동시에 § 3의 뜻에서 \(S\)의 확대가군
\[
S_{\mathsf{A}}=x_1\mathsf{A}+\ldots+x_n\mathsf{A}
\]
이므로 확대환이라고도 한다.

이 특수화 아래 불변가군에 관한 정리(§ 3, 3)는 다음 정리가 된다.

\noindent\textbf{확대정리.} \emph{\(S_{\mathsf{A}}\)의 각 양쪽 아이디얼은 \(S\)의 한 아이디얼의 확대아이디얼이며, 더 정확히는 그 아이디얼과 \(S\)의 교집합을 확대한 것이다.}

실제로 바탕에 둔 군 \(\frG\)가 \(A\)의 내부자기동형사상군이면, \(A\)의 중심인 \(P\)는 완전 불변체이다. \(A\)와 \(S\)의 원소별 교환성은 § 3, 3에 따라 \(\frG\)가 \(S_{\mathsf{A}}\)로 확장되고 \(S\)가 완전 불변영역이 된다는 뜻이다. \(S_{\mathsf{A}}\)의 각 양쪽 아이디얼 \(\mathfrak a\)는 \(\alpha^{-1}\mathfrak a\alpha\leqq\mathfrak a\)이므로 허용되는 부분군이다. 따라서 § 3, 3에 따라 교집합가군 \(\mathfrak a\cap S\)의 확대이며, 이 교집합은 \(S\)의 아이디얼이다.

\noindent\textbf{보충.} \(S\)가 가환이면 \(S\)는 \(S_{\mathsf{A}}\)의 중심이 된다. 실제로 \(S\)는 중심에 있고 \(\mathsf{A}\)가 생성하는 내부자기동형사상들에 대한 완전 불변영역이다. \emph{더 일반적으로 \(S\)의 중심은 \(S_{\mathsf{A}}\)의 중심이기도 하다.}

\noindent\textbf{표기.} \(\mathsf{A}_r,\mathsf{B}_n\ldots\)는 언제나 각각 \(\mathsf{A},\mathsf{B}\ldots\) 위의 표시한 차수의 행렬환을 뜻한다. 곧
\[
A_r=\sum A c_{ik}=\sum c_{ik}A.
\]

이하에서는 단순계에 대한 확대정리의 귀결을 본질적으로 쓴다. \(S\)가 단순계이면\footnote{단순계는 통상대로 “양쪽 단순”을 뜻한다.} \(S_{\mathsf{A}}\)도 단순이다.
\[
S_A=\sum Bc_{ik}=\sum c_{ik}B=B_t,
\]
여기서 대응하는 체 \(\mathsf{B}\)의 중심은 \(\mathsf{P}\)를 포함한다. \(\mathsf{A}\), 따라서 \(\mathsf{B}\)도 \(\mathsf{P}\) 위 무한 랭크일 수 있으며, \(S\)만 초복소적이라고 가정한다. \(S,A\)의 공통 중심이 \(\mathsf{P}\)이면 \(\mathsf{P}\)는 \(S_{\mathsf{A}}\)의 중심이기도 하다.

더 일반적으로 \(S\)가 단순 초복소계이면 직접곱 \(S\times \mathsf{A}_r\)도 단순이다. 실제로 \(S\times \mathsf{A}_r\)는 \(S_{\mathsf{A}}\times \mathsf{P}_r\)와 일치하므로 \(S_{\mathsf{A}}=\mathsf{B}_t\)이면 \(\mathsf{B}_{tr}\)와 같다.

\noindent\textbf{2. 표현류.} 초복소계의 반대 표현 또는 직접 표현은 통상대로 계수영역 \(\mathsf{P}\)에 관하여 연산자준동형이고(§ 2의 끝) 영표현과 다른 것을 뜻한다. 이행정리 및 그 귀결(§ 2, 3)과 1의 확대정리를 결합하면 다음을 얻는다.

\noindent\textbf{표현류에 관한 정리.} \emph{\(S\)가 단위원을 갖고 \(\mathsf{P}\)를 계수영역으로 하는 초복소계이며 \(\mathsf{A}\)가 중심 \(\mathsf{P}\)인 체이면, \(S\)가 갖는 \(\mathsf{A}\) 안의 서로 다른 기약 반대 표현류의 수는 \(S_{\mathsf{A}}\)를 그 라디칼로 나눈 잉여류환에서 단순 오른쪽 아이디얼들의 류 수와 같다. \(S\)가 단순계이면 \(\mathsf{A}\) 안에 기약 반대 표현류 하나와 기약 직접 표현류 하나를 정확히 갖는다.}\footnote{\normalfont \(\mathsf{A}\)가 \(S\)에 딸린 자기동형체인 경우에는 \emph{표현론}, 주 20을 보라.}

실제로 이행정리의 귀결(§ 2, 3)에 따라 단순 반대 표현류와 \(S_{\mathsf{A}}\) 위 단순 가군류는 일대일로 대응한다. \(S_{\mathsf{A}}\)는 체 \(\mathsf{A}\)에 관하여 유한 랭크이므로 한쪽 아이디얼에 대한 최대·최소조건을 만족하며, 따라서 명제의 첫 부분은 \emph{표현론} § 19에서 바로 나온다. \(S\)가 단순이면 확대정리에 따라 \(S_{\mathsf{A}}\)도 단순이다. 최대·최소조건 때문에 동시에 완전가약이고, 따라서 단순 한쪽 아이디얼류가 하나뿐이다. 곧 \(S\)는 \(\mathsf{A}\) 안에 기약 반대 표현류가 하나뿐이다. \(S\)와 반대동형인 계도 단순이고 기약 반대 표현류가 하나뿐이므로, \(S\)에는 단순 직접 표현류도 하나뿐이다.

\noindent\textbf{3. 랭크 관계.} \(S\)가 단순일 때 \(S_{\mathsf{A}}=\sum Bc_{ik}\)가 \(\mathsf{A}\)와 \(\mathsf{B}\) 양쪽 모두에 관하여 유한 랭크라는 사실은 다음 랭크 관계를 준다.
\[
\begin{array}{rl}
(1)&n=rt,\quad
n=(S:P),\quad t^2=(S_A:B),\quad
r=\text{기약 반대 표현의 차수}.
\end{array}
\]

실제로 \(S_{\mathsf{A}}\) 자체가 \(S\)의 \(\mathsf{A}\) 안 반대 표현가군이다. 그 표현은 이미 \(\mathsf{P}\)에 들어 있다. \(\mathsf{B}\) 위 행렬환으로서 \(S_{\mathsf{A}}\)는 \(t\)개의 단순 오른쪽 아이디얼, 곧 \(S_{\mathsf{A}}\)-가군으로 분해되며, 각 아이디얼의 \(\mathsf{A}\)-랭크는 기약 반대 표현의 차수 \(r\)과 같다. \(n\)은 동시에 \(S_{\mathsf{A}}\)의 \(\mathsf{A}\)-랭크이므로 랭크 관계가 따라온다.

\paragraph{4. \(\mathsf{A}\)가 \(\mathsf{P}\)에 관하여 유한 랭크일 때 랭크 관계의 강화.}
이 경우 세 관계
\[
\begin{array}{rl}
(1)& (S:\mathsf{P})=n=rt,\\
(2)& (\mathsf{B}:\mathsf{P})\cdot t=(\mathsf{A}:\mathsf{P})\cdot r,\\
(3)& (S:\mathsf{P})\cdot(\mathsf{B}:\mathsf{P})=(\mathsf{A}:\mathsf{P})\cdot r^2.
\end{array}
\]
가 성립한다. 여기서 (2)는 \(S_{\mathsf{A}}\)의 단순 오른쪽 아이디얼의 \(P\)-랭크를 3에 따라 \(\mathsf{B}\)-랭크와 \(\mathsf{A}\)-랭크로 각각 나타낸 것이다. (3)은 행렬환 \(\mathsf{B}_n\) 안의 단순 오른쪽 아이디얼의 랭크를 나타내며, (1)을 사용하여 (2)에 \(r\)을 곱하면 나온다.

\begin{center}
§ 5

\textbf{행렬환에 표현정리 적용}
\end{center}

§ 4에서 고찰한 \(S\)의 \(\mathsf{A}\) 안 반대 표현은 \(S\)에서 \(\overline A_r\)의 한 부분환으로 가는 직접 준동형으로도 볼 수 있다. 여기서 \(A\)는 \(A\)와 반대동형인 체이다. \emph{이하 여러 절의 원리는 역으로 행렬환 \(\overline A_r\)의 연구를 반대동형인 체 \(A\) 안의 표현론으로 환원하는 것이다.}

\noindent\textbf{1. \(\overline A_r\) 안의 가약·기약 매장.} \(\mathsf{A},A\)를 중심 \(P\) 위에서 유한 또는 무한 랭크인 반대동형체라 하자\footnote{일반적으로 반대인 체 또는 계는 서로 대응하는 그리스 문자와 라틴 문자로 나타낸다.}. \emph{\(P\) 위 단순계 \(S\)가 \(\overline A_f\) 안에 기약 또는 가약으로 매장 가능하다는 것은, \(S\)가 \(A\) 안에 \(f\)차 기약 또는 가약 반대 표현을 갖는다는 뜻이다.} \(S\)가 \(\overline A_r\) 안에 기약으로 매장 가능하면, \(s>1\)인 모든 행렬환 \(\overline A_{rs}\) 안에는 가약으로 매장 가능하다. 실제로 이 차수들과 오직 이 차수들의 가약 표현을 갖는다(§ 4, 2). 따라서 \(S\)는 \(\overline A_r,\overline A_{rs}\)의 어떤 부분환들과 직접 동형이다. 이 동형은 \(\mathsf{P}\)의 항등사상의 연장이다\footnote{표현의 경우와 마찬가지로 이하의 동형은 언제나 \(\mathsf{P}\)의 항등사상의 연장, 곧 \(\mathsf{P}\) 위의 동형이다. 매번 따로 밝히지는 않는다.}. § 4, 3에 따라 \(r\)은 \(S\)의 \(P\)-랭크 \(n\)의 약수이다.

이 방법으로 \(P\)를 포함하고 \(P\) 위 유한 랭크인, 여러 \(\overline A_f\), \(f=1,2\ldots\), 의 모든 단순 부분환을 얻는다. 실제로 그러한 각 부분환은 \(A_f\)의 한 부분환과 반대동형이므로 \(A\) 안에 가약 또는 기약 반대 표현을 갖는다.

\noindent\textbf{2. 내부자기동형사상에 관한 정리.} \emph{\(S^{(1)},S^{(2)}\)가 \(\overline A_f\)의 두 단순 부분환으로 \(P\)를 포함하고 \(P\) 위 유한 랭크이며, \(S^{(1)},S^{(2)}\)가 \(P\) 위에서 동형이면, 이 동형은 \(\overline A_f\)의 한 내부자기동형사상으로 생성된다.}

실제로 \(S^{(1)},S^{(2)}\)는 한 단순계 \(S\)를 \(\mathsf{A}_f\)에 일반적으로 가약으로 매장한 것, 따라서 \(A\) 안의 \(f\)차 직접 표현들이다. 이들은 \(A\) 안의 같은 가약 직접 표현류에 속하므로, 둘을 옮기는 행렬이 자기동형사상을 생성한다.

\(A\) 자체가 \(\mathsf{P}\) 위 유한 랭크여서 초복소적이면 이 정리는 다음의 알려진 정리가 된다.

\emph{\(A_f\)에서 중심 \(\mathsf{P}\)를 포함하는 두 단순 부분계가 \(\mathsf{P}\) 위에서 동형이면, \(A_f\)의 한 내부자기동형사상에 의해 서로 옮겨간다. 특히 \(A_f\)의 모든 자기동형사상은 내부자기동형사상이다.}\footnote{\normalfont \(A\)가 \(\mathsf{P}\) 위 무한 랭크이면 이는 더 이상 성립하지 않는다. 주 5에서 인용한 Köthe의 논문을 보라.}

\noindent\textbf{3. 원소별로 교환하는 부분환에 관한 정리.}
이 정리도 § 2, 3의 교환행렬 정리에서, 이 절의 처음에 말한 원리에 따라 나온다.

\emph{\(S\)가 \(\overline A_f\)의 단순 부분계로 \(\mathsf{P}\)를 포함하고, \(R\)이 \(\overline A_f\)에서 \(S\)의 모든 원소와 교환하는 원소들의 전체이면, \(R\)도 행렬환이다. 곧 \(R=B_s\)이다. \(S_{\mathsf{A}}\)에 딸린 체 \(B\)와 \(R\)에 딸린 체 \(B\)는 반대동형이다. \(R,S\)의 교집합은 \(S\)의 중심이다. \(R\)이 체인 것은 \(S\)가 \(\overline A_f\) 안에 기약으로 매장되었을 때 그리고 오직 그때이다.}

실제로 § 2, 3에 따라 \(R\)은 \(S\)의 \(\mathsf{A}\) 안 반대 표현가군 \(\frM\)의 자기동형환과 직접 동형이다. 이는 \(\frM\)을 \(S_{\mathsf{A}}\)-가군으로 본 자기동형환이다. \(\frM\)이 \(s\)개의 단순 \(S_{\mathsf{A}}\)-가군으로 분해되고, § 4, 1처럼 \(S_{\mathsf{A}}=\sum \mathsf{B}c_{ik}\)라 두면, \(R\)은 \(\sum Bc_{ik}\)와 동형이다. 여기서 \(B,\overline B\)는 반대동형이다. 실제로 \(\mathsf{B}\)는 \(S_{\mathsf{A}}\)의 단순 오른쪽 아이디얼, 곧 단순 \(S_{\mathsf{A}}\)-가군의 자기동형체와 반대동형이다(§ 1, 3의 끝). \(R\)이 체 \(R=B\)인 것은 \(s\), 곧 \(S\)가 기약으로 매장된 경우와 정확히 같다. \(R,S\)의 교집합이 \(S\)의 중심이라는 사실은 \(R\)의 정의에서 곧바로 나온다.

\noindent\textbf{4. 행렬환이 초복소적인 경우의 강화된 교환정리.} 이 경우 § 4, 4의 랭크 관계에서 다음 강화가 나온다.

\emph{중심 \(\mathsf{P}\) 위 유한 랭크인 체 \(A\)에 대하여, \(\overline A_f\)에서 \(\mathsf{P}\)를 포함하는 단순 부분계들은 쌍 \(S,\overline S\)로 나뉜다. \(\overline S\)는 \(S\)의 모든 원소와 교환하는 원소 전체이고 그 역도 성립한다. \(S_{\mathsf{A}}\)와 \(\overline S\)에 딸린 체들은 반대동형이며, \(S\)와 \(\overline S_{\mathsf{A}}\)에 딸린 체들도 반대동형이다. \(S,\overline S\)의 교집합은 공통 중심과 같다. 한 부분환이 체인 것은 다른 부분환이 기약으로 매장되었을 때 그리고 오직 그때이다. \(S,\overline S\)의 \(P\)-랭크의 곱은 \(\overline A_f\)의 랭크와 같다. \(f=rs=\bar r\,\bar s\)이고 \(S\)가 \(\overline A_r\)에 기약으로, \(\overline S\)가 \(\overline A_{\bar r}\)에 기약으로 매장 가능하면, \(S,\overline S\)에 딸린 체들은 \(f=g s\bar s\)인 \(A_g\) 안의 서로 교환하는 부분환 한 쌍과 동형이다.}

본질적으로는 랭크들의 곱에 관한 명제만 증명하면 된다. 먼저 \(S\)가 기약으로 매장되어 \(\overline S\)가 체이고 \(\mathsf{B}\)와 반대동형이라고 하자. 여기서 \(S_{\mathsf{A}}=\mathsf{B}_t\)이다. 이제 \(f=r\)이고 \(\overline S\)의 랭크는 \(\mathsf{B}\)의 랭크와 같으므로 § 4, 4의 랭크 관계 (3)이 곱에 관한 명제를 준다. \(S\)가 가약으로 매장되었으면 \(\overline S=B_s\), \(f=rs\)이고, (3)에 \(s^2\)을 곱하면 명제를 얻는다. 이로부터 \(S\)가 \(\overline S\)의 모든 원소와 교환하는 원소 전체임이 곧바로 나오고, 마지막 명제를 제외한 나머지는 3에서 모두 따라온다.

마지막 명제는 기약 매장으로 두 번 이행하여 얻는다. \(f=rs\)인 \(\overline A_r\) 안에서는 \(S\)가 기약으로 매장되고 \(S,R\)은 서로 교환하는 계이며, \(R\)은 체 \(B\)와 동형이다. 한편 \(S,\overline S\)의 반대성 때문에 \(S=\overline B_s\)이다. 여기서 \(\overline B\)는 \(B\)가 \(S\)에 대해 갖는 것과 같은 의미를 \(\overline S\)에 대해 갖는다. 따라서 \(R\)은 \(\overline A_r\)에 가약으로 매장되고, \(r=g\bar s\)인 \(\overline A_g\)에 기약으로 매장된다. 여기서는 서로 교환하는 두 계가 각각 \(B,\overline B\)와 동형이다.

\begin{center}
§ 6

\textbf{단순계의 갈루아 이론}
\end{center}

§ 5, 4의 강화된 교환정리는 중심을 기초영역으로 한 단순계의 갈루아 이론을 나타낸다. 먼저 체를, 그 다음 일반 단순계를 고찰한다.

\noindent\textbf{1. 중심에 관하여 유한 랭크인 체의 갈루아 이론.} \(A\)의 \emph{갈루아군} \(\Gcal\)은 중심 \(P\)의 항등사상을 연장하는 모든 자기동형사상의 군, 따라서 § 5, 2에 따라 모든 내부자기동형사상의 군으로 정의한다. 그러므로
\[
\Gcal\simeq A^*/P^*,
\]
여기서 \(A^*,P^*\)는 각각 \(A,P\)의 0 아닌 원소들의 곱셈군이다. 따라서 \(\Gcal\)의 부분군 \(\mathfrak H\)에는 \(\mathsf{P}^*\)를 포함하는 \(A^*\)의 부분군 \(H^*\)가
\[
\mathfrak H\simeq H^*/P^*
\]
에 따라 일대일로 대응한다. \(\Gcal\)의 부분군 \(\mathfrak H\)는 \(H^*\)에 0을 덧붙인 것이 한 체 \(H\)가 될 때 닫혀 있다고 한다. \(C\)가 군 \(\mathfrak H\)를 허용한다는 것은 \(C\)와 \(H\)가 원소별로 교환한다는 것과 같다. 따라서 교환정리(§ 5, 4)는 다음 정리가 된다.

\begin{center}
\textbf{\srcspaced{비가환체의 갈루아 이론의 기본정리}}
\end{center}

\emph{\(\mathsf{P}\)와 \(A\) 사이의 체 \(C\)들과 \(\Gcal\)의 닫힌 부분군 \(\mathfrak H\)들을 일대일로 대응시켜, \(C\)가 \(\mathfrak H\)의 완전 불변체이고 \(\mathfrak H\)가 \(C\)의 완전 불변군이 되게 할 수 있다.}

\noindent\textbf{\srcspaced{2. 단순계의 갈루아 이론.}}
\(A_f\)가 중심 \(\mathsf{P}\)인 단순계이면 갈루아군 \(\Gcal\)은 다시 내부자기동형사상들의 군이고
\[
\Gcal\simeq A_f^*/P^*
\]
이다. 이제 \(A_f^*\)는 \(A_f\)의 정칙원들의 곱셈군을 뜻한다. \(\Gcal\)의 부분군 \(\mathfrak H\simeq H^*/P^*\)는 \(H^*\)가 생성하는 부분환 \(H\)가 단순계이고 \(H^*\)가 \(H\)의 모든 정칙원으로 이루어질 때 단순-닫힌 부분군이라 한다. 교환정리에서 갈루아 이론으로 넘어가는 것은 다음 보조정리로 주어진다.

\noindent\textbf{\srcspaced{보조정리.}} \emph{\(A_f\)에서 \(P\)를 포함하는 각 단순계는 그 정칙원들의 군 \(H^*\)에 의해 생성된다.} \(\mathsf{P}\)가 무한히 많은 원소를 가지면 이는 분명하다. 부정원들로 만든 “일반원소”가 정칙이고, 알맞게 특수화하면 \(\mathsf{P}\) 위의 정칙 기저원소들을 만들 수 있기 때문이다. Shoda에 따르면 보조정리는 일반적으로도 성립한다\footnote{주 3에서 인용한 Shoda의 논문. 거기서는 부정원의 도입을 피한다.}.

보조정리에 따라 \(S\)와 한 단순계 \(\overline S\)가 교환한다는 것은 다시 \(S\)가 \(\overline S\)의 정칙원들이 유도하는 자기동형사상들을 허용한다는 것과 같다. 따라서 여기서도 교환정리(§ 5, 4)는 다음 정리가 된다.

\noindent\textbf{\srcspaced{단순계의 갈루아 이론의 기본정리.}} \emph{\(A_f\)에서 \(\mathsf{P}\)를 포함하는 단순계 \(S\)들과 \(\Gcal\)의 단순-닫힌 부분군 \(\mathfrak H\)들을 일대일로 대응시켜, \(S\)가 \(\mathfrak H\)의 완전 불변영역이고 \(\mathfrak H\)가 \(S\)의 완전 불변군이 되게 할 수 있다.}

\noindent\textbf{\srcspaced{3. 연장원리에 의한 증명.}}
체의 경우에 기본정리의 둘째 증명을 제시한다. 이 증명은 동형, 곧 1차 표현의 연장원리에 바탕을 두며 랭크를 세지 않으므로 더 약한 유한성 가정만 필요하다. 특히 무한 랭크인 체 \(S\)로 어느 방향에서 옮길 수 있는지가 드러난다. 또한 기약 반대 표현류가 하나뿐이라는 사실을 쓰지 않으므로 가환적인 경우에도 그대로 성립한다(§ 8, 2). 먼저 몇 가지 보조정리를 둔다.

\noindent\textbf{\srcspaced{가정.}} \(P\) 위 단순계 \(S\)가 \(A\) 안에 1차 반대 표현을 갖는다고 하자. 따라서 \(S\)는 체이다. 여기서 \(\mathsf{A}\)는 중심 \(\mathsf{P}\) 위 유한 또는 무한 랭크일 수 있다. \(S_{\mathsf{A}}\)를 \(n\)개의 단순하고 서로 연산자동형인 오른쪽 아이디얼로 분해한 것이(§ 4, 3)
\[
S_{\mathsf{A}}=\mathfrak r_1+\ldots+\mathfrak r_n
=e_1S_{\mathsf{A}}+\ldots+e_nS_{\mathsf{A}}
=e_1\mathsf{A}+\ldots+e_n\mathsf{A} 
\]
로 주어졌다고 하자. \(e_i\)가 생성하는 반대 표현은
\[
e_i s=e_i\sigma
\]
로 정의된다.

\noindent\textbf{\srcspaced{보조정리 1.}} \emph{\(e_1,\ldots,e_n\)이 생성하는 \(n\)개의 반대 표현은 모두 서로 다르다. 따라서 같은 표현류에 속하는 서로 다른 표현들이다. 그러므로 \(S\)는 적어도 그 랭크만큼 서로 다른 표현을 갖는다.}

증명을 위해 § 2의 끝에 따라 \(S\)의 표현들을 \(S_{\mathsf{A}}\)의 대응들로 확장하여 \(\mathsf{A}\)에 관하여 연산자준동형이 되게 한다. 여기서는 \(S_{\mathsf{A}}\)의 각 \(w\)에 대하여 \(\omega\in A\)인
\[
e_i w=e_i\omega
\]
로 이 대응을 정의한다. 서로 다른 \(i,j\)에 대하여 \(e_i,e_j\)가 같은 표현을 생성한다면 모든 \(w\in S_A\)에 대해 \(e_iw=e_jw\)이어야 한다. 그러나 \(e_ie_i=e_i\), \(e_je_i=0\)이므로 불가능하다.

\noindent\textbf{\srcspaced{보조정리 2.}} \emph{\(T\)가 \(P,S\) 사이의 체이고 \(S\)의 \(T\)-랭크가 \(s\)이면, \(T\)에서 \(\mathsf{A}\)로 가는 각 반대동형은 서로 다른 연장을 적어도 \(s\)개 갖는다.}

\(T\)에서 \(\mathsf{A}\)로 가는 반대동형, 곧 \(T\)의 \(A\) 안 반대 표현이
\[
T_{\mathsf{A}}=E_1T_{\mathsf{A}}+\ldots+E_hT_{\mathsf{A}}
=E_1\mathsf{A}+\ldots+E_h\mathsf{A};\qquad E_it=E_i\tau
\]
라는 분해로 매개된다고 하자. 여기서 \(E_i\)들은 멱등원이다. 이는 일반성을 제한하지 않는다. 기저원소 \(E_i\alpha\)가 생성하는 표현은 \(\alpha^{-1}\cdot E_i\alpha\), 곧 한 멱등원으로도 생성되기 때문이다. 오른쪽에서 \(S\)를 곱하면
\[
S_A=E_1S_A+\ldots+E_hS_A
\]
를 얻는다. 이때 모든 \(E_iS_{\mathsf{A}}\)의 \(A\)-랭크는 \(s\)이다. 실제로 \(S\)에 \(T\)-기저를 대입하고 \(S,A\)의 교환성을 쓰면 랭크가 많아야 \(s\)임을 알 수 있다.
\[
S=Tw_1+\ldots+Tw_s,
\]
\[
E_iS_{\mathsf{A}}=E_iT_{\mathsf{A}}w_1+\ldots+E_iT_{\mathsf{A}}w_s
=E_iw_1\mathsf{A}+\ldots+E_iw_s\mathsf{A} .
\]
합 \(S_{\mathsf{A}}\)의 랭크가 \(n=hs\)이므로 각 \(E_iS_{\mathsf{A}}\)의 랭크는 정확히 \(s\)이다. 따라서
\[
E_iS_A=r_{i1}+\ldots+r_{is}
=e_{i1}A+\ldots+e_{is}A,
\]
이고 보조정리 1에 따라 \(e_{i1},\ldots,e_{is}\)가 생성하는 \(s\)개의 표현은 모두 서로 다르다. 이들은 모두 \(E_i\)가 생성하는 \(T\)의 표현의 연장이다. 실제로 \(E_it=E_i\tau\)에서
\[
(e_{i1}+\ldots+e_{is})t=(e_{i1}+\ldots+e_{is})\tau
\]
가 나오고, 직합분해 때문에 \(e_{il}t=e_{il}\tau\)이다.

\paragraph{정의.}
\(T\)가 유도하는, \(S\)에서 \(\mathsf{A}\)로 가는 반대동형들의 계 \(\mathfrak I\)의 류분할 \(\mathfrak I_T\)를 다음처럼 정의한다. \(T\) 위에서 같은 동형을 유도하는 \(\mathfrak I\)의 동형들과 오직 그것들만 서로 동치라고 한다.

\noindent\textbf{\srcspaced{기본정리.}} \emph{\(\mathfrak I_T\)가 \(T\)가 유도한 류분할이면, \(T\)는 이 류분할에 관한 극대 부분체이다.}

실제로 \(L\)이 \(T,S\) 사이의 중간체이고 \(T\) 위 차수가 \(l\)이면, 보조정리 2에 따라 \(\mathfrak I_T\)의 각 류는 \(\mathfrak I_L\)의 류 적어도 \(l\)개로 갈라진다. 이 정식에서 통상적인 기본정리로 넘어가려면, 동형들의 집합 \(\mathfrak I\)의 원소들을 그 가운데 고정한 하나의 역원으로 곱하여 자기동형사상군으로 옮기면 된다. 이때 \(\mathfrak I_T\)의 한 류는 \(T\)의 불변군으로 옮겨간다. 닫힌 부분군들이 완전 불변군이라는 명제도 여기에 포함되어 있다. \(\mathfrak H\simeq H^*/\mathsf{P}^*\)에 대해서는 체 \(H\)를 중간체 \(T\)로 대입하기만 하면 된다.

\begin{center}
§ 7

\textbf{\srcspaced{유사한 대수들의 류. 분해체}}
\end{center}

이제부터는 중심 \(P\)에 관하여 \emph{유한 랭크}인 체 \(\mathsf{A},A\)와 행렬환 \(A_r\)만 다룬다. 곧 \(\mathsf{P}\) 위의 단순 정규대수이며, 이를 줄여서 \emph{\(\mathsf{P}\) 위의 대수} 또는 단순히 대수라 하겠다. \(\mathsf{A},A\)는 이때 \emph{나눗셈대수}이다. 한 \emph{유사한 대수들의 류}, \(A_r\sim A_s\), 에는 같은 딸린 \(A\)를 갖는 모든 \(A_r\)를 묶는다. 동형인 \(A\)들은 동일시한다.

\noindent\textbf{\srcspaced{1. 대수류들의 군.}} \(A_r,B_s\)가 \(\mathsf{P}\) 위의 대수이면 그 곱, 곧 \(\mathsf{P}\) 위의 직접곱도 그러하다. § 3의 끝에 따라
\[
A_r\times B_s=C_t
\]
이고, 여기서 \(C\)는 동형까지 유일하게 정해지는 중심 \(\mathsf{P}\)인 나눗셈대수이다. \(A_r,B_s\)에 \(\mathsf{P}\) 위 행렬환들을 곱해 보면 이 곱셈이 류들에 대해 일의적으로 정해짐을 알 수 있다. 따라서 류들은 곱셈에 닫힌 가환계를 이룬다. 더 나아가 다음이 성립한다.

\noindent\textbf{\srcspaced{R. 브라우어의 정리.}} \emph{유사한 대수들의 류는 직접곱에 관하여 아벨군을 이룬다.}

실제로 이 계는 \(\mathsf{P}\) 위의 행렬환들, 곧 \(\mathsf{P}\)와 유사한 대수들로 이루어진 항등류를 갖는다. 또한 각 류 \((A)\)에 역류 \((A)^{-1}\)가 있으며, 이는 \(A\)와 반대동형인 \(\mathsf{A}\)에 유사한 대수들로 이루어진다. 이는 § 5, 3의 교환정리에서 \(S=A\)로 특수화하면 나온다. 실제로 \(A\)는 자기 자신 안에 기약으로 매장되므로 \(B=\mathsf{P}\), \(A_{\mathsf{A}}=\mathsf{P}_t\)이다\footnote{대수류군에 관한 더 세밀한 정리들은 인자계 이론을 쓴다. 여기서는 다루지 않는다. 여기까지는 가환체에 관한 고찰이 전혀 나타나지 않았음을 지적한다. 예컨대 랭크 \((A:\mathsf{P})\)가 제곱수라는 사실은 사용하지도 증명하지도 않았다.}.

\noindent\textbf{\srcspaced{2. 한 대수류의 분해체.}}\quad 분해체 이론도 전적으로 § 5의 처음에 말한 원리에 바탕을 둔다. 가환 확대체들을 \(\mathsf{A}\) 또는 \(A\) 안에 표현한다. 먼저 다음 보조정리를 둔다.

\medskip
\noindent\textbf{\srcspaced{보조정리.}} \emph{\(A\)가 중심 \(\mathsf{P}\)인 체이고 \(\mathsf{Z}\)가 \(\mathsf{P}\)의 유한 가환 확대체이면 \(A_{\mathsf{Z}}\)는 중심 \(\mathsf{Z}\)인 단순계이다.} 실제로 § 4, 1에 따라 \(A_{\mathsf{Z}}\)와 반대동형인 계 \(Z_{\mathsf{A}}\)에 이 사실이 성립한다. 따라서 \(P\) 위에서 \(A\)와 유사한 대수들의 각 류 \((A)\)에 대하여, \(Z\)는 \(Z\) 위에서 \(A_{\mathsf{Z}}\)와 유사한 단순대수들의 \emph{확대류} \((A)_{\mathsf{Z}}\)를 생성한다.

확대류에 딸린 나눗셈대수 \(D\)는 § 5, 3의 교환정리에서 곧바로 나온다.

\noindent\textbf{\srcspaced{확대류에 딸린 나눗셈대수에 관한 정리.}} \emph{\((A)_{\mathsf{Z}}\)가 확대류이고 \(\mathsf{Z}\)가 \(A_r\) 안에 기약으로 매장, 곧 표현되었다고 하자. 그러면 \((A)_{\mathsf{Z}}=(D)\)이며, \(D\)는 \(A_r\)에서 \(Z\)의 모든 원소와 교환하는 원소들의 전체이다.} § 5, 3의 단순계 \(S\)를 \(Z\)로 바꾸고 \(Z_A,A_Z\)가 반대동형임을 고려하면 된다. \(\mathsf{Z}\)가 가약으로 매장되었으면 \(Z\)의 모든 원소와 교환하는 원소 전체는 \(D\) 위의 행렬환이 된다.

\(P\)의 가환 확대체 \(Z\)는 확대류 \((A)_{\mathsf{Z}}\)가 완전히 분할될 때, 곧 \(\mathsf{Z}\) 위의 항등류, 즉 \(\mathsf{Z}\) 위 모든 행렬환의 류와 같을 때 류 \((A)\)의 \emph{분해체}라 한다.

\noindent\textbf{\srcspaced{분해체의 특성화 정리.}}
\emph{\(P\)의 유한 가환 확대체 \(Z\)가 류 \((A)\)의 분해체인 것은, \(Z\)를 \(A_r\) 안에 기약으로 매장한 것이 \(A_r\)의 극대 가환 부분체가 될 때 그리고 오직 그때이다.}

\(\mathsf{Z}\)가 \((A)\)의 분해체라는 성질은 \((D)=(\mathsf{Z})\)와 같다. 따라서 딸린 나눗셈대수에 관한 정리에 따라 기약으로 매장된 \(Z\)는 극대 가환 부분체여야 한다. 역으로 이 조건이 성립하면 반드시 \(D=Z\)이다. 실제로 \(D\)의 각 원소 \(a\)를 \(Z\)에 첨가하면 가환 확대체를 생성한다.

\noindent\textbf{\srcspaced{주의.}}\quad
1. 동시에 \emph{극대 가환 부분체} \(Z\)를 \emph{기약으로} 매장하면 \emph{극대 가환 부분환}을 생성함이 따라온다. \(Z\)의 모든 원소와 교환하는 원소 전체가 체를 이루기 때문이다.

2. 이 정리는 분해체의 \emph{존재}도 준다. 딸린 나눗셈대수 \(A\)는 반드시 극대 가환 부분체를 갖기 때문이다.

\medskip
\noindent\textbf{\srcspaced{3. 랭크 관계, 인덱스, 무한 가환 확대.}}\quad § 5, 4의 랭크 명제는 먼저 단순대수의 중심에 관한 랭크가 제곱수라는 알려진 사실을 준다. \(Z\)가 \(A\)의 극대 가환 부분체이면 \(A\) 안의 모든 매장은 기약이므로
\[
(Z:P)\cdot(Z:P)=(A:P).
\]
따라서
\[
(A:P)=m^2,\qquad(Z:P)=m,
\]
이고 \(m\)은 \emph{슈어 인덱스}이다. 이는 동시에 나눗셈대수 \(A\) 자체에 매장 가능한 모든 분해체의 차수이다. 일반 분해체의 차수 \(n\)에 대해서는
\[
n=m\cdot r
\]
을 얻는다\footnote{일반적으로는 진부분체가 분해체가 아닌 “최소” 분해체가 모든 차수에 걸쳐 존재하기도 한다. 주 2에서 인용한 Brauer--Noether 논문을 보라.}. 예를 들어 \((A_r:\mathsf{P})=m^2\cdot r^2\)에서 나오며, § 4, 3의 랭크 관계에서도 나온다. 후자의 이유는 인덱스 \(m\)이 \(A_{\mathsf{Z}}\)의 단순성분 수와도 같고, \(A_{\mathsf{Z}}\)가 \(\mathsf{Z}\) 위 랭크 \(m^2\)인 행렬환이 되기 때문이다. 더 일반적으로 \(A_\Lambda\sim D\)이고 \(L\)이 \(\Lambda\)의 \(A_r\) 안 기약 매장이며 나눗셈대수 \(D\)의 중심 \(L\) 위 랭크 \((D:L)\)가 \(d^2\)이면
\[
l\cdot d=m\cdot r.
\]
실제로 § 5, 4와 딸린 나눗셈대수에 관한 정리에 따라
\[
(A_r:P)=(L:P)\cdot(D:P),
\quad\text{즉}\quad m^2r^2=l^2d^2
\]
이다.

분해체의 존재로부터 \(P\)의 무한 가환 확대체 \(\Omega\), 대수적이든 초월적이든, 에 대하여 다음이 따라온다.

\noindent\emph{확대류 \((A)_\Omega\)는 중심 \(\Omega\)인 단순대수들의 류이다. 특히 \(\Omega\)가 대수적으로 닫혀 있으면 \(A_\Omega\)는 \(m\)차 행렬환이다. 따라서 인덱스는 “절대 성분수”와 같다.}

실제로 \(\mathsf{Z}\)가 \(A\)의 분해체이고 \(\bar\Omega\)가 \(\Omega,Z\)의 합성체이면 \(A_{\bar\Omega}\)는 \(\bar\Omega\) 위 행렬환이다. 따라서 라디칼이 없고 중심이 \(\bar\Omega\)이다. 그러므로 \(A_\Omega\)도 라디칼이 없으며 중심은 \(\Omega\)이다. 실제로 \(A_\Omega\)의 \(\Omega\) 밖 원소가 중심에 있으면 \(A_{\bar\Omega}\)의 중심에도 있게 된다. 따라서 \(A_\Omega\)는 \(\Omega\) 위 단순대수이고, \(\Omega\)가 대수적으로 닫혀 있으면 반드시 \(\Omega\) 위 행렬환이다.

\medskip
\noindent\textbf{\srcspaced{4. 분리가능 분해체의 존재.}}\quad
\emph{모든 류 \((A)\)는 분리가능 분해체를 가지며, 심지어 \(A\) 자체에 매장 가능한 것을 갖는다.}\footnote{G. Köthe, \emph{Über Schiefkörper mit Unterkörpern zweiter Art über dem Zentrum}, Journ. f. Math. 166 (1932), 182--184쪽을 보라. 여기의 훨씬 간단한 증명은 M. Zorn의 지적에서 나왔다.}

기초체 \(\mathsf{P}\)의 표수를 \(p\), 인덱스를 \(m=s\cdot p^g\)라 하자. 여기서 \(s\)는 \(p\)와 서로소이다. \(\mathsf{Z}\)가 \(A\)의 차수 \(m\)인 분해체이고 \(\mathsf{Z}_0\)가 그 안에 든 제1종 확대, 따라서 \((\mathsf{Z}:\mathsf{Z}_0)=p^f\), 이면 3에 따라 \(A_{\mathsf{Z}_0}\sim D\)의 인덱스는 \(p^f\)이다. \(\mathsf{Z}_0\)의 분리가능 확대체 \(\mathsf{Z}_1\) 가운데 \(D_{\mathsf{Z}_1}\)의 인덱스가 \(p^f\)의 진약수가 되는 것이 존재함을 보이겠다. 대수적으로 닫힌 \(\Omega\)에 대한 \(D_\Omega\)는 3에 따라 라디칼이 없으므로 \(D\)의 축약판별식은 0이 아니다. 따라서 \(D\)에는 축약대각합이 0이 아닌 원소 \(d\)가 적어도 하나 있다. 이 \(d\)는 기초체 \(\mathsf{Z}_0\)에 속할 수 없다. \(\mathsf{Z}_0\)의 각 원소 \(\alpha\)의 대각합은 \(p^f\alpha=0\)이기 때문이다. 그러므로 \(\mathsf{Z}_1=\mathsf{Z}_0(d)\)는 진확대이고 분리가능 확대체이며, \(D_{\mathsf{Z}_1}\)의 인덱스는 \(p^f\)의 진약수이다. 이를 유한 번 반복하면 인덱스와 같은 차수 \(m\)을 갖는 \((A)\)의 분리가능 분해체에 이른다.

\begin{center}
\textbf{§ 8}\\[0.35em]
\textbf{가환적인 경우의 분해체·분리체·갈루아 이론}
\end{center}

중심 위 단순계의 분해체에서 임의의 단순계의 분해체로 넘어가려면 중심의 분해이론을 더 전개해야 한다. 여기에 § 6, 3과 평행한 갈루아 이론도 덧붙인다. 이 절의 모든 체와 계는 가환이다.

\noindent\textbf{\srcspaced{1. 가환 단순계의 분해체와 분리체.}}\quad 가환계 \(Z\)가 \(\mathsf{P}\) 위에서 단순, 따라서 체이고 \(\Omega\)가 \(\mathsf{P}\)의 대수적으로 닫힌 확대체라 하자. 그러면 다음은 알려져 있다(\emph{표현론} § 21).

\noindent\emph{\(Z_\Omega\)가 라디칼 없는 계로 남는 것은 \(Z\)가 \(\mathsf{P}\) 위 분리가능, 곧 제1종 확대일 때 그리고 오직 그때이다.}

실제로 이때 그리고 오직 이때 \(Z_\Omega\)는 \(\Omega\) 위 랭크와 같은 수의 1차 동형사상을 \(\Omega\) 안에 가지며, 따라서 같은 수의 서로 다른 표현을 갖는다. 여기서는 표현과 표현류가 일치한다.

\(Z_\Omega\)에 라디칼이 나타날 수 있고 따라서 절대단순 성분들로의 직합분해가 반드시 존재하지는 않는다는 사실은 다음 정의들을 낳는다.

\noindent\textbf{\srcspaced{정의.}}\quad \(\mathsf{P}\)의 확대체 \(\Lambda\)는 \(Z_\Lambda\)의 합성인자들이 모두 1차일 때, 달리 말해 모든 절대기약 표현이 이미 \(\Lambda\) 안에 있을 때 분해체라 한다. \(\mathsf{P}\)의 확대체 \(T\)는 \(Z_{\mathsf{T}}\)가 적어도 하나의 1차 합성인자를 분리해 낼 때, 달리 말해 \(Z\)의 절대기약 표현이 적어도 하나 \(\mathsf{T}\) 안에 존재할 때 분리체라 한다.

분해체와 분리체는 다음 정리로 특성화된다.

\noindent\textbf{\srcspaced{정리.}} \emph{한 체 \(\mathsf{T}\)가 분리체인 것은 \(Z\)와 동형인 부분체를 포함할 때 그리고 오직 그때이다. 한 체 \(\Lambda\)가 분해체인 것은 \(Z\)에 딸린 갈루아체와 동형인 부분체 \(\Gamma\)를 포함할 때 그리고 오직 그때이다.}

이는 절대기약 표현을 이용한 분해체와 분리체의 정의에서 바로 나온다. 특히 비가환적인 경우와 유사하게 다음 귀결을 얻는다.

\noindent\textbf{\srcspaced{귀결.}} \emph{한 체가 \(\mathsf{Z}\)와 동형일 때 그리고 오직 그때 \(\mathsf{Z}\)의 최소 분리체, 곧 진부분체가 분리체가 아닌 체이다. 이는 그 체를 \(Z\)에 기약 가환적으로 매장한 것이 \(Z\)의 극대 가환 부분체를 이룬다는 뜻이다. 최소 분리체가 동시에 분해체인 것은 \(Z\)가 \(\mathsf{P}\) 위 정규체, 곧 갈루아체일 때 그리고 오직 그때이다.}

단순대수를 그 중심 위에서 정규라 부르는 이유가 여기에 있다. 고정한 \(P\) 위 대수적으로 닫힌 \(\Omega\) 안에서는 최소 분해체가 \(\mathsf{Z}\)에 딸린 갈루아체 하나만 존재한다. 이는 일반적으로 동형이 아닌 최소 분해체가 무한히 많은 비가환적인 경우와 대조된다\footnote{R. Brauer--E. Noether, 앞의 글.}. 그 이유는 가환적인 경우 직합분해가 유일한 반면 비가환적인 경우에는 연산자동형까지 유일할 뿐이라는 데 있다. 같은 말로 하면, 가환적인 경우에는 서로 다른 절대기약 표현이 유한개뿐인 반면 비가환적인 경우에는 같은 류에 속하면서도 서로 다른 것이 무한히 많기 때문이다(2를 보라).

\medskip
\noindent\textbf{\srcspaced{2. \(Z/\mathsf{P}\)가 분리가능체일 때 갈루아 이론의 초복소적 정초.}}\quad § 6, 3과 정확히 유사하게, \(\mathsf{P}\) 위 분리가능인 임의의 \(Z\)에 대해 동형 이론이 성립한다. \(Z\)가 갈루아체인 경우에는 통상적인 자기동형사상군으로 옮길 수 있다.

\medskip
\noindent\textbf{\srcspaced{가정.}}\quad \(P\) 위 단순계 \(Z\)를 유한 분리가능 확대체라 하고, \(\Omega\)를 \(\mathsf{P}\) 위 유한 또는 무한 분해체라 하자. \(\mathsf{Z}=\mathsf{Z}^{(1)}\)은 \(Z\)와 동형인 부분체, \(\Gamma\)는 그에 딸린 갈루아체라 하자. 1에 따라 \(Z\)가 분리가능이므로 직합분해
\[
Z=\mathfrak r^{(1)}+\ldots+\mathfrak r^{(n)}
=e^{(1)}Z_\Omega+\ldots+e^{(n)}Z_\Omega
=e^{(1)}\Omega+\ldots+e^{(n)}\Omega 
\]
가 성립한다. \(e^{(i)}\)가 생성하는 표현 \(Z\to\mathsf{Z}^{(i)}\), \(\mathsf{Z}^{(i)}\leqq\Gamma\), 는
\[
e^{(i)}z=e^{(i)}\zeta^{(i)}
\]
로 정의된다.

\noindent\textbf{\srcspaced{보조정리 1.}} \emph{\(e^{(1)},\ldots,e^{(n)}\)이 생성하는 \(n\)개 표현은 모두 서로 다르다.} 증명은 § 6, 3과 같고, 1에서도 나온다. 이들은 서로 다른 표현류에 속하기 때문이다.

\noindent\textbf{\srcspaced{보조정리 2.}} \emph{\(T\)가 \(P,Z\) 사이의 중간체이고 \(Z\)의 \(T\)-랭크가 \(s\)이면, \(T\)에서 \(\Omega\)로 가는 각 동형은 적어도 \(s\)개, 따라서 정확히 \(s\)개의 연장을 갖는다.}

증명은 § 6, 3과 같다. 여기서 “적어도”를 “정확히”로 강화할 수 있는 것은 \(Z_\Omega\)의 분해가 유일하기 때문이다. 이에 따라 \(Z\)는 \(\Omega\) 안에 적어도 \(n\)개가 아니라 정확히 \(n\)개의 동형을 갖는다(1 끝의 주의를 보라).

§ 6, 3과 마찬가지로 이제 \(Z\)의 동형들로 이루어진 유한집합 \(\mathfrak I\)에 \(T\)가 \emph{유도하는 류분할} \(\mathfrak I_T\)를 정의한다. 곧 \(T\) 위에서 같은 동형을 유도하는 \(\mathfrak I\)의 동형들과 오직 그것들만 \(\mathfrak I_T\)에서 서로 동치라고 한다. 그리고 다음을 증명한다.

\noindent\textbf{\srcspaced{기본정리.}} \emph{\(\mathfrak I_T\)가 \(T\)가 유도하는 류분할이면 \(T\)는 이 류분할에 관한 극대 부분체이다.}

\(Z\)가 갈루아체인 경우에는 § 6, 3과 똑같이 한 동형의 역과 합성하여 자기동형사상군 및 통상적인 정식으로 옮길 수 있다. 부분군에 관한 대응 명제는 그 자체로 흥미로운 멱등원을 고찰하면 나온다.

\medskip
\noindent\textbf{\srcspaced{3. \(Z_\Omega\)의 멱등원.}}\quad \(S\)가 \(Z\)를 각각의 \(\mathsf{Z}^{(i)}\)로 옮기는 \(n\)개 치환을 두루 돌게 하자. 이는 \(\mathsf{Z}\to Z\)와 \(Z\to\mathsf{Z}^{(i)}\)라는 동형들의 합성이며, 첫 것은 \(Z\to\mathsf{Z}\)의 역동형을 뜻한다. \(Z\)는 갈루아체일 필요가 없다. \(Z_{\mathsf{Z}}\)의 원소 \(a\)에 대하여 치환 \(S\)를 연산자로 정의하되, \(Z\) 위에서는 항등사상을 유도한다고 정한다. 그러면 각 \(a\)에 대해 \(Z_{\mathsf{Z}^{(i)}}\)에 놓이는 켤레원소 \(a^S\)가 정의된다. 이 규정 아래 다음이 성립한다.

\noindent\textbf{\srcspaced{정리.}} \emph{\(Z_\Omega\)의 \(n\)개 멱등원 \(e^{(i)}\)는 서로 켤레이다. \(e=e^{(1)}\)라 하면 \(e^{(i)}=e^S\)이다. 이들은 \(n\)개의 켤레 확대계 \(Z_{\mathsf{Z}^{(i)}}\)에 각각 놓인다.}

실제로 \(S\)를 적용하면 정의관계 \(ez=e\zeta^{(1)}\), \(e^2=e\)는
\[
e^Sz=e^S\zeta^{(i)},\qquad(e^S)^2=e^S
\]
로 옮겨간다. 분해의 유일성 때문에 멱등원들도 유일하게 정해지므로 \(e^{(i)}=e^S\)이다. 또한 \(\mathsf{Z}\)는 분리체여서 표현 \(ez=e\zeta\)를 이미 매개하므로 \(e\in Z_Z\), 따라서 \(e^{(i)}\in Z_{Z^{(i)}}\)이다.

특히 \(Z\)가 갈루아체이면 \emph{갈루아 이론 기본정리에서 부분군에 관한 부분}이 곧바로 따라온다. 실제로 \(\mathfrak H\)가 갈루아군 \(\Gcal\)의 부분군이면, 직합성분이 유일하게 정해진다는 데서 \(H\in\mathfrak H\)에 대한 원소 \(\sum e^H\)는 \(\mathfrak H\)의 모든 치환과 오직 그것들만 허용한다. 군 \(\Gcal\)은 \(\mathsf{Z}\)에 대해 정의했으므로 이는 \(\mathsf{Z}\)의 갈루아 이론이며, 동형에 의해 \(Z\)에 관한 대응 결과도 함께 얻는다.

\paragraph{4. 멱등원과 여차기저의 관계.}
다시 \(Z\)가 반드시 갈루아체라고 가정하지 않는다.

\noindent\textbf{\srcspaced{정리.}} \emph{\(a_1,\ldots,a_n\)이 \(Z/P\)의 기저이고 멱등원을 \(e=a_1\beta_1+\ldots+a_n\beta_n\), 따라서 \(e^S=a_1\beta_1^S+\ldots+a_n\beta_n^S\), 로 두면, \(\beta_1^S,\ldots,\beta_n^S\)는 \(e^S\)가 매개하는 사상에서 \(a_1,\ldots,a_n\)의 여차기저 \(b_1,\ldots,b_n\)의 상이다.}

실제로 \(e^S\)가 매개하는 사상
\[
e^Sz=e^S\zeta^S,
\]
특히 \(e^Sa_i=e^S\alpha_i^S\)에서 \(e^S\cdot e^S=e^S\), \(e^S\cdot e^R=0\) (\((S\ne R)\))이므로 대응 \(e^S\to1\), \(e^R\to0\)이 성립한다. 따라서
\[
1=\alpha_1^S\beta_1^S+\ldots+\alpha_n^S\beta_n^S;\qquad
0=\alpha_1^S\beta_1^R+\ldots+\alpha_n^S\beta_n^R\quad(S\ne R),
\]
곧 행렬로 쓰면 여차기저의 정의방정식
\[
\begin{pmatrix}
\alpha_1 & \ldots & \alpha_n\\
\ldots & \ldots & \ldots\\
\alpha_1^S & \ldots & \alpha_n^S\\
\ldots & \ldots & \ldots
\end{pmatrix}
\begin{pmatrix}
\beta_1 & \ldots & \beta_1^R & \ldots\\
\ldots&\ldots&\ldots&\ldots\\
\beta_n & \ldots & \beta_n^R & \ldots
\end{pmatrix}
=E
\]
을 얻는다. 대각합을 통한 정의로 옮기는 것은 다음에서 나온다. \(a_i=\sum_S e^S\alpha_i^S\), \(b_i=\sum_S e^S\beta_i^S\)에 대하여 행렬관계의 두 행렬을 바꾸면
\[
\Sp(a_i b_i)=1,\qquad
\Sp(a_i b_k)=0\quad(i\ne k)\srcfnmark{21)}
\]
가 된다.
\srcfntext{21)}{\emph{표현론} § 25를 보라.}
이때 \(a_i,b_i\)는 모든 치환 \(S\)를 허용하므로 \(Z\)에 놓인다(§ 3, 3의 보조정리)\srcfn{22)}{여차기저와 멱등원 사이의 관계는 Dedekind, \emph{Zur Theorie der aus \(n\) Haupteinheiten gebildeten komplexen Größen}에 처음 나타난다. \emph{Gesammelte Werke} II, 4--5쪽을 보라.}.

여차기저들이 반대변환한다는 사실은 멱등원 \(e^S\)가 \(Z\)의 불변량이고 바탕으로 삼은 기저와 무관하다는 데서 나온다. 곧 \(\mathsf{Z}/\mathsf{P}\)의 다른 기저 \(\bar a_1,\ldots,\bar a_n\)으로 옮겨갈 때 모든 \(a_1\beta_1^S+\ldots+a_n\beta_n^S\)는 자기 자신으로 옮겨간다.

\begin{center}
\textbf{§ 9}\\[0.8ex]
\textbf{임의의 계의 분해체와 분리체}
\end{center}

\noindent\textbf{\srcspaced{1. 정의와 단순한 경우로의 환원.}}
§ 8의 정의들에 대응하여 일반적인 경우에는 다음과 같이 정의한다.

\noindent\textbf{\srcspaced{정의.}} \(S\)가 \(\mathsf{P}\) 위 초복소계라 하자. \(\mathsf{P}\)의 가환 확대체 \(\Lambda\)는 \(S_\Lambda\)가 한쪽 아이디얼의 합성열에서 절대단순 합성인자들로 분해될 때, 달리 말해 \(S\)의 \(\Lambda\) 안 모든 기약 표현이 이미 절대기약일 때 \(S\)의 분해체라 한다. \(\mathsf{P}\)의 확대체 \(T\)는 \(S_{\mathsf{T}}\)가 적어도 하나의 절대단순 합성인자를 분리해 낼 때, 곧 \(\mathsf{T}\) 안의 기약 표현 적어도 하나가 이미 절대기약일 때 분리체라 한다.

이 정의들은 § 8에서 가환적인 경우에, § 7에서 단순대수에 대해 준 정의들을 분명히 포함한다. 후자에 대해서는 중심을 가환적으로 확대할 때마다 \(A_\Lambda\)가 완전가약이므로 합성인자와 직합분해의 성분이 일치하기 때문이다. 여기서 중심 위 완전행렬환과 오직 그것들만 절대단순 성분, 따라서 절대기약 표현을 준다. 또 \(A_\Lambda\)는 양쪽 단순이어서 서로 연산자동형인 성분들로 분해되므로 절대단순 성분 하나를 분리해 내기만 해도 완전히 분할된다. 따라서 분리체와 분해체가 일치한다.

정의에서 다음 사실들이 바로 나온다. \emph{\(S\)의 분해체는 \(\mathsf{P}\)에서 기약인 각 표현의 분해체 하나씩을 합성한 체들이다. \(S\)의 분리체는 \(\mathsf{P}\)에서 기약인 어떤 표현의 분리체이면 된다.} 단순계, 곧 라디칼로 나눈 잉여류환의 성분들과 \(\mathsf{P}\) 안의 기약 표현들이 일대일로 대응하므로 \emph{단순계만 고찰하면} 충분하다. 이 경우 결과들은 § 7과 § 8을 결합하여 얻는다.

\noindent\textbf{\srcspaced{2. 단순계의 분해체와 분리체.}} 먼저 중심 \(Z\)가 \(P\) 위 \emph{분리가능}인 단순계 \(A\)를 고찰한다. § 8, 1과 § 8, 3에서 다음이 나온다.

\emph{\(\Omega\)가 \(Z\)의 분해체일 때 \(Z_\Omega\)의 직합분해에 대응하는 \(A_\Omega\)의 양쪽 분해는 \(A\)와 동형인 \(n\)개의 켤레 성분 \(e^{(i)}A\)를 준다. 이 성분들은 각자의 중심 \(e^{(i)}\mathsf{Z}^{(i)}=e^{(i)}Z\) 위 단순대수가 된다. 여기서 치환 \(S\)는 \(A\) 위 항등사상을 유도하도록 규정하여 \(A_{\mathsf{Z}}\)의 연산자로 정의한다.} (§ 8, 3을 보라.)

실제로 § 8, 3에 따라 멱등원들이 켤레이므로, \(A\)에 대한 치환의 정의상 성분 \(e^{(i)}A\)들도 켤레이다. \(A\)가 양쪽 단순이므로 \(e^{(i)}A\)는 \(A\)와 환 동형이다. 또한 \(e^{(i)}Z=e^{(i)}\mathsf{Z}^{(i)}=e^{(i)}Z_{\mathsf{Z}^{(i)}}\)이므로
\[
e^{(i)}A=e^{(i)}A_{Z^{(i)}}.
\]
따라서 중심은 계수영역과 일치하며, 성분들은 정규대수이다.

이제 § 7, 2의 결과들은 다시 다음을 준다.

\emph{\(A\)가 \(\mathsf{P}\) 위 단순계이고 중심 \(Z\)가 \(\mathsf{P}\) 위 분리가능이면 \(A_\Omega\)도 라디칼이 없고 따라서 완전가약이다. 여기서 \(\Omega\)는 임의의 유한 또는 무한 확대체일 수 있다.}

실제로 § 7, 2에 따라 \(e^{(i)}A_{\mathsf{Z}^{(i)}}\)는 계수를 어떻게 확대해도 단순이고 따라서 라디칼이 없다. 그러므로 직합도 라디칼이 없다. 이 직합은 \(A_\Omega\)이거나, \(\Omega\)가 \(Z\)의 분해체가 아니면 \(A_\Omega\)에서 계수를 확대하여 생긴다. § 7, 2로부터 다음도 나온다.

\begin{center}
\textbf{\srcspaced{분리체와 분해체의 특성화}}
\end{center}

\emph{\(A\)가 분리가능 중심 \(Z\)를 갖는 나눗셈대수이면, 분리체의 기약 매장은 \(A_r\)에서 \(Z\)를 포함하는 극대 가환 부분체들과 정확히 일치한다. 분해체는 분리체 하나와 중심의 분해체 하나, 특히 중심에 딸린 갈루아체를 합성한 체들과 정확히 일치한다. 중심이 갈루아체가 아니어도 최소 분리체가 분해체일 수 있다.}

실제로 분리체는 \(Z\)와 동형인 체를 포함해야 한다. 분리체의 매장에 관한 명제는 이제 앞의 사실들과 § 7, 2에서 나온다. 분해체는 더 나아가 \(Z\)의 한 분해체 \(\Omega\)를 포함해야 한다. 그런데 \(\Omega\) 위의 모든 분리체는 동시에 분해체이다. 성분 \(e^{(i)}A_\Omega\)가 단순이고 \(\Omega\)와 동형인 중심을 갖기 때문이다. 따라서 \(Z\)가 갈루아체이면 최소 분리체와 최소 분해체가 일치한다. 그러나 가환적인 경우와 달리 비갈루아인 경우에도 최소 분리체가 동시에 분해체일 수 있다. 바로 그러한 최소 분리체가 \(Z\)에 딸린 갈루아체를 포함하는 경우이다. 예를 들어 \(P\)가 유리수체일 때 \(\mathsf{P}(\sqrt[3]{2})\)와 동형인 중심을 갖는 사원수체에서는 \(\mathsf{P}(\sqrt[3]{2})\)에 딸린 갈루아체가 최소 분리체이자 분해체이다.

끝으로 중심이 \emph{비분리가능이면 \(Z_\Omega\), 따라서 \(A_\Omega\)는 라디칼을 갖는 계이다.} 예를 들어 \(\mathfrak C\)를 \(Z_\Omega\)의 라디칼, \(\mathfrak c\)를 \(A_\Omega\)에서 그 확대아이디얼이라 하자. 그러면 \(Z_\Omega/\mathfrak C\)에는 § 8, 2와 정확히 유사하게 켤레 성분들로의 직합분해가 성립한다. 위와 똑같이 \(A_\Omega/\mathfrak c\)의 직합분해를 얻으며, 그 성분 \(\bar e^{(i)}A\)들은 \(A\)와 동형이고 중심은 \(\bar e^{(i)}\mathsf{Z}^{(i)}\)이다. 따라서 분리체와 분해체에 관한 정리들도 유지된다. 더 나아가 랭크를 세면 \(\mathfrak C\)의 합성열에 대응하는 합성인자들이 \(A_\Omega/\mathfrak c\)와 단순히 준동형일 뿐 아니라 연산자동형임을 알 수 있다.

기약 표현의 말로 쓰면 다음과 같이 요약된다.

\noindent\textbf{\srcspaced{요약.}} \emph{한 초복소계의 \(\mathsf{P}\) 안 기약 표현이 주어졌을 때, 이 표현이 절대 완전가약인 것은 그에 딸린 중심이 \(\mathsf{P}\) 위 분리가능일 때 그리고 오직 그때이다. 어느 경우든 표현의 분해체와 분리체는 위와 같이 딸린 단순계 안의 매장으로 특성화할 수 있다. 특히 분리체의 최소차수는 중심의 차수와 중심 위 대응 대수류의 인덱스의 곱과 같고, 모든 분리체의 차수는 그 배수이다. 이 표현은 합성열의 뜻에서, 중심에 포함된 가장 큰 분리가능체의 차수만큼 서로 다르지만 서로 켤레인 절대기약 표현류들로 분해된다. 각 류는 인덱스와 그 분리가능 확대에 대한 중심의 차수의 곱만큼 반복된다.}

모든 확대체가 분리가능인 완전체의 경우에는 이렇게 I. Schur의 알려진 결과들을 다시 얻으며, R. Brauer에게 이미 있는 매장에 의한 특성화를 덧붙이게 된다.

\begin{center}
(1932년 6월 8일 접수)
\end{center}

\clearpage
\setcounter{footnote}{0}
\phantomsection
\addcontentsline{toc}{section}{41. 상대 갈루아 수체에 대한 주종정리}
%
%

%
%
\section*{41. 상대 갈루아 수체에 대한 주종정리.}

\begin{center}
지음\\[0.4ex]
\textbf{괴팅겐의 에미 뇌터.}\\[0.8ex]
\emph{Math. Ann. 108 (1933), 411--419}
\end{center}

최근에는 비가환적 방법, 특히 대수 이론을 이용하면 상대 순환 수체와 상대 아벨 수체에 관한 알려진 정리들을 임의의 상대 갈루아 체에 대해 정식화하고 증명할 수 있음이 밝혀졌다\srcfn{1)}{이에 대해서는 나의 취리히 강연, Verhandl. des internat. Math. Kongr. Zürich, 1932, Bd. I에 실린 \emph{Hyperkomplexe Systeme in ihren Beziehungen zur kommutativen Algebra und Zahlentheorie}를 보라. 거기에서는 특히 노름정리와 주종정리의 초복소적 정식화도 상세히 논의하였고 문헌도 제시하였다. 분할대수에 관한 정리에 대해서는 또한 H. Hasse, \emph{Die Struktur der R. Brauerschen Algebrenklassengruppe}, (6.1), Math. Annalen 107 (1932)의 서술을 보라.}. 무엇보다도 나는 일반적으로 분할대수에 관한 정리로 표현되는 ``노름정리''를 상기시키고자 한다. 이하에서는 이에 대응하여 주종정리도 내부자기동형사상에 관한 정리로, 또는 ``교차표현''에 관한 정리로 일반적인 초복소적 정식화를 할 수 있음을 보인다. 그 증명은 앞서 말한 분할대수에 관한 정리로부터 순전히 대수적·산술적 고찰을 통해 얻어진다. 반면 앞의 정리에는 순환적인 경우의 노름정리---소수 차수만으로 충분하다---가 초월적 핵심으로 들어 있다. 정식화를 더 잘 이해할 수 있도록, 이하에서 사용하지는 않을 ``최소형 주종정리''\srcfn{2)}{임의의 체에서 주종정리에 대응하는 대수적 유사물이기 때문에 이를 ``최소형''이라고 부른다. 이에 반해 ``소형''에서는 이미 \(p\)-진 기초체로 옮겨 가는 것이 의미를 갖는다.}를 먼저 제시한다. 이는 초등대수적 정리\srcfn{3)}{이 정리는 A. Speiser, \emph{Zahlentheoretische Sätze aus der Gruppentheorie}, S. 3, Math. Zeitschr. 5 (1919), S. 1--6에 나온다. 그곳에서도 이미 교차표현에 관한 정리로 진술되어 있다.}로서, 순환적인 특수 경우에는 힐베르트의 수론 보고서에 나오는 잘 알려진 정리 90, 즉 \(N(a)\)가 1인 것과 \(a=b^{1-S}\)로 쓸 수 있는 것이 동치라는 정리로 귀착된다.

또한 나는 이 정리도 내부자기동형사상에 관한 정리로, 또는 교차표현에 관한 정리로 진술하고, 내부자기동형사상과 대수의 교차곱표현에 관한 일반 정리들에 의하여 증명한다. 본래의 주종정리에서는 뒤의 교차곱 대신 아이디얼류군과 갈루아군의 교차곱이 나타나지만, 인자계에는 더 세밀하게 유도된 아이디얼류 분할을 사용한다. 이로써 유체론의 레이류 분할에 대응하는 유사물이 포착된다. 여기에는 자기동형사상과 표현에 관한 일반 정리들이 아직 마련되어 있지 않으므로---주종정리는 이 방향의 첫걸음으로 볼 수 있다---앞서 말한 대로 분할대수에 관한 정리를 이용한 환원을 제시한다. 이로써 이 정리는 아이디얼류 대신 아이디얼에 대한 대응 정리로 환원된다. 여기에서는 한 대수의 극대 오더들 사이의 관계에 관한 브란트의 정리들\srcfn{4)}{극대 오더에 관한 브란트 이론은 H. Hasse, \emph{Über \(\mathfrak p\)-adische Schiefkörper und ihre Bedeutung für die Arithmetik hyperkomplexer Zahlsysteme}, Math. Ann. 104 (1931)에 서술되어 있으며 상세한 새 증명도 제시되어 있다.}을 자기동형사상 정리들과 동치인 것으로 볼 수 있다. 실제로 이 정리들을 교차곱의 극대 오더에 관한 고찰\srcfn{5)}{이에 관해서는 C. Chevalley (Hamb. Ber.), H. Hasse, E. Noether의 논문이 각각 한 편씩 출판될 예정이며, 뒤의 두 편은 Herbrand 기념 논문집에 실릴 예정이다.}과 결합하면, 개별 소로 되돌아가야 하므로 더 복잡하기는 하지만 본질적으로 최소형의 증명과 나란히 진행되는 증명이 얻어진다. 따라서 나는 E. Artin이 제시한 거의 자명한 증명을 택한다. 이는 여기에서 슈파이저의 증명\textsuperscript{3)}에 대응하는 한층 더 단순한 유사물이다.

%
%
\subsection*{§ 1.}
\begin{center}
\textbf{최소형 주종정리.}
\end{center}

이 절에서 \(k\)는 임의의 (가환) 체---반드시 수체일 필요는 없다---를, \(K/k\)는 \(n\)차 분리 갈루아 확대를, \(\mathfrak G\)는 이에 대응하는 갈루아군을 뜻한다.

\noindent\textbf{1.} 먼저 교차곱에 관한 알려진 사실들을 간단히 모아 둔다\srcfn{6)}{이 이론은 나의 강의(1929/30)에 이어 H. Hasse, \emph{Theory of cyclic algebras}, Kap. 2, Am. Transact. 34 (1932)에 서술되어 있다. 또한 \emph{Ergebnisse der Mathematik}에 출판될 예정인 초복소수에 관한 M. Deuring의 보고서를 보라. 여기의 요약은 주 1에서 인용한 나의 취리히 강연에서 가져왔다.}. 교차곱은 \(K\)의 자기동형사상들이 내부자기동형사상이 되도록 \(K\)와 \(\mathfrak G\)를 하나의 대수 \(A\)에 동시에 매입하는 것을 뜻한다. \(n\)개의 군 원소에 대응하는 기호들을 \(u_{S_1},\ldots,u_{S_n}\)이라 하자. 그러면 먼저 \(A\)를 \(K\) 위에서 랭크 \(n\)인 선형형식 모듈로 둔다:

%
%
\srcdisplay{(1)}{A=u_{S_1}K+\ldots+u_{S_n}K.}

\(u_S\), 더 일반적으로는 \(u_SK^*\)가 만들어 내는 내부자기동형사상이라는 조건\srcfn{7)}{\(K^*\)는 \(K\)에서 0을 제외하여 얻으며, 이 표기를 이하 일반적으로 사용한다.}에 따라 \(A\)는 환, 따라서 \(k\) 위에서 차원 \(n^2\)인 대수가 된다. 실제로 이 조건은 다음과 같이 표현된다.

\srcdisplay{(2)}{u_S^{-1}zu_S=z^S\quad\text{또는}\quad zu_S=u_Sz^S}

여기서 \(z\)는 \(K\)의 임의의 원소이고,

\srcdisplay{(3)}{u_Su_T=u_{ST}a_{S,T}\quad\text{단, }a_{S,T}\text{는 }K^*\text{에 속한다},}

\srcdisplay{(4)}{a_{ST,R}a_{S,T}^{R}=a_{S,TR}a_{T,R}}

(결합법칙은 \([u_Su_T]u_R=u_S[u_Tu_R]\)와 동치이다).

이제 임의의 두 원소의 곱을 또한
\[
\sum u_Sb_S\cdot\sum u_Tc_T=\sum u_Su_Tb_S^Tc_T
\]
로 정의하면, \(A\)는 인자계 \(a_{S,T}\)에 대한 \(K\)와 \(\mathfrak G\)의 교차곱이 된다. 이를 \(A=(a_{S,T},K,\mathfrak G)\)로 표기한다. \(A\)가 \(k\) 위의 단순 정규대수(중심단순대수), 곧 이에 대응하는 나눗셈대수 \(D\) 위의 \(r\)차 행렬환 \(D_r\)이 되고, \(K\)가 극대 가환 부분체, 따라서 분해체가 됨을 증명할 수 있다. 역으로 주어진 나눗셈대수 \(D\)에 대해서도, 위와 같은 방식으로 교차곱으로 생성할 수 있는 행렬환 \(D_r\)가 항상 존재한다.

\(u_S\)에서 \(v_S=u_Sc_S\)로 옮겨 가되 \(c_S\in K^*\)라 하면, 이는 같은 자기동형사상을 만들어 내며 다음의 ``연관된'' 인자계가 생긴다.

\srcdisplay{(5)}{\bar a_{S,T}=a_{S,T}c_S^Tc_T/c_{ST}.}

특히 (5)에 따라 인자계 \(\bar a_{S,T}=c_S^Tc_T/c_{ST}\)는 1에 연관된다. 이 \(c_S^Tc_T/c_{ST}\)들을 변환량이라고 한다. 연관된 인자계들은 하나의 류 \((a)\)로 묶고, 특히 변환량 \(c_S^Tc_T/c_{ST}\)들은 류 \((1)\)로 묶는다. 마찬가지로 \(A\)와 상사인 모든 대수, 즉 \(r=1,2,\ldots\)인 모든 \(D_r\)를 하나의 류 \(\mathfrak A\)로 묶는다. 류 \(\mathfrak A\)와 \((a)\)는 일대일로 대응한다. \emph{분해체 \(K\)를 고정한 류들은 직접곱 구성에 관하여 아벨군을 이루며, 이는 인자계의 류들을 성분별로 곱하여 얻는 군과 동형이다. 단위원은 각각 분할대수류와 모든 변환량 \(c_S^Tc_T/c_{ST}\)의 계이다.} 이것이 R. 브라우어의 대수류군이다.

%
%
이제 순환대수인 특수 경우를 간단히 제시한다. \(Z/k\)가 순환이고 \(S\)가 그 군의 생성 치환이라 하자. 그러면 \(S\)의 거듭제곱들에 \(u\)의 거듭제곱들을 대응시킬 수 있으며, 다음을 얻는다.

\srcdisplay{(1')}{A=Z+uZ+\ldots+u^{n-1}Z.}
\srcdisplay{(2')}{zu=uz^S.}
\srcdisplay{(3')}{u^n=\alpha.}
\srcdisplay{(4')}{\alpha\text{는 기초체 }k^*\text{에 속한다}.}
\srcdisplay{(5')}{\bar\alpha=\alpha\cdot N(c),\quad\text{단 }v=uc\text{로 둔다}.}

따라서 여기에서 각 인자계는 기초체에 속하는 단 하나의 원소 \(\alpha\)로 이루어지며---\(A=(\alpha,Z,S)\)로 표기한다---인자계의 단위원류는 \(Z^*\)에서 오는 노름들로 주어지고, 대수류군은 몫군 \(k^*/N(Z^*)\)와 동형이 된다.

%
%
\noindent\textbf{2.} \emph{최소형 주종정리}의 정식화에서는 관계 (2)부터 (5)가 순전히 곱셈적이라는 사실을 사용한다. 따라서 이 관계들은 동시에 \(n\)개의 잉여류 \(u_SK^*\)의 원소들로 이루어진 \(\mathfrak G\)의 확장군 \(\mathfrak G^*\)을 정의한다:

\srcdisplay{(6)}{\mathfrak G^*=\{u_{S_1}K^*,\ldots,u_{S_n}K^*\}.}

\(\mathfrak G^*\)은 \(K^*\)을 아벨 정규부분군으로 가지며, \(\mathfrak G^*/K^*\)은 \(\mathfrak G\)와 동형이다\srcfn{8)}{따라서 관계 (2)부터 (5)는 단순히 군의 확장에 관한 정의관계이다. 이러한 확장은, \(K^*\)을 오직 곱셈군으로 보았을 때 그 전체 자기동형군 안에 \(\mathfrak G\)의 준동형상인 부분군이 존재하기만 하면 성립한다는 것이 알려져 있다. 교차곱의 특수성은 이 준동형이 동형이 되고, 더 나아가 바로 그 부분군으로 곱셈적이면서 동시에 덧셈적인 자기동형사상 전체, 즉 \(K\)의 갈루아군을 택한다는 데 있다. 이로써 군의 확장과 더불어 환의 확장도 얻어진다.}. \(\mathfrak G^*\)은 \(K\)를 자기 자신으로 변환하는 모든 (가역) 원소 \(g^*\), 즉 \(g^{*-1}Kg^*=K\)인 원소들로 이루어진다. 그러므로 \(K\)의 환 자기동형사상은 정확히 \(\mathfrak G^*\)의 원소들에 의한 변환으로 유도된다. 실제로 그러한 \(g^*\)는 모두 \(K\)의 환 자기동형사상을 유도한다. 역으로 그러한 자기동형사상은 모두 어떤 원소 \(v=u_Sw\)에 의한 변환으로 유도되는데, 여기서 \(w\)는 \(K\) 위에 항등사상을 유도한다. 따라서 \(w\)는 \(K\)의 모든 원소와 가환하므로 \(K\)에 속한다. \(K\)가 최대 가환 부분환이므로 \(w\)는 \(K^*\)에 속한다.

\emph{최소형 주종정리. 제1정식: \(K^*\)의 항등자기동형사상의 연장인 \(\mathfrak G^*\)의 모든 군 자기동형사상은 내부자기동형사상이며, \(K^*\)의 한 원소에 의해 유도된다. 제2정식: 변환량 \(c_S^Tc_T/c_{ST}\)가 모두 1이면, \(c_S\)들은 기호적 \((1-S)\)제곱이다. 즉, 모든 \(S\in\mathfrak G\)에 대하여 \(c_S=b^{1-S}=b/b^S\)가 되게 하는 \(b\in K^*\)가 존재한다. 제3정식: 군 \(\mathfrak G\)은 \(K^*\) 안에서 인자계 1에 속하는 1차 교차표현류를 단 하나만 가진다.}

이때 표현 \(u_S\to C_S\)가 인자계 \(a_{S,T}\)에 대한 교차표현이라는 것은 \(C_S^TC_T=C_{ST}a_{S,T}\)가 성립한다는 뜻이다. 두 표현은 \(C_S=B^{-S}D_SB\)일 때 같은 류에 속한다.

%
%
먼저 제1정식을 증명한 다음 그것이 제2정식 및 제3정식과 일치함을 보이겠다. 이 증명은 주어진 종류의 군 자기동형사상이 모두 \(A\)의 환 자기동형사상으로 연장될 수 있고, 이 환 자기동형사상이 알려진 바와 같이 내부자기동형사상이라는 사실에 근거한다. 실제로 이 자기동형사상 아래에서 \(u_S\)의 상을 \(v_S\)라 하자. 그러면 가정에 따라 관계 (2)부터 (4)가 모두 보존되므로, \(\sum v_SK\)도 같은 인자계에 대한 \(\mathfrak G\)와 \(K\)의 교차곱을 이룬다. 대응 \(\sum u_Sb_S\to\sum v_Sb_S\)은 \(K\)의 원소들을 자기 자신에 대응시키는 환 자기동형사상을 준다. 따라서 앞에서 \(\mathfrak G^*\)에 관하여 말한 바에 의해, 이 자기동형사상은 \(K^*\)의 한 원소에 의해 유도된다.

제2정식으로 넘어가는 근거는 (2)에 의해 \(v_S\)가 \(v_S=u_Sc_S\) 꼴이어야 한다는 데 있다. 또한 인자계도 보존되므로 (5)에 따라 그에 딸린 변환량들은 모두 1이어야 한다. 역으로 이것이 성립하면, (2)부터 (5)는 다시 치환 \(v_S=u_Sc_S\)가 앞에서 지정한 종류의 \(\mathfrak G^*\)의 자기동형사상을 준다는 것을 보인다. 따라서 제1정식에 의해
\[
v_S=b^{-1}u_Sb=u_Sb/b^S\qquad\text{이고 따라서}\qquad c_S=b^{1-S}.
\]
특히 순환인 경우에는 제2정식의 가정이 (5')에 따라 \(N(c)=1\)로 표현되며, \(c=b^{1-S}\)는 잘 알려진 정리를 준다\srcfn{9)}{여기서 제시한 증명은 특히 \(k\)가 유한개의 원소만을 가질 때에도 성립한다. 이 경우 힐베르트의 증명은 성립하지 않지만 슈파이저의 증명은 여전히 유효하다.}.

제3정식은 제2정식과 곧바로 동치이다. 실제로 제2정식의 가정은 \(u_S\to c_S\)가 인자계 1에 대한 1차 교차표현을 이룬다는 뜻이다. 명제 \(c_S=b^{1-S}\)는 이 표현이 일표현과 동치임, 즉 \(c_S\sim1\)임을 뜻한다.

다만 이 제3정식은 훨씬 더 일반적인 다음 정리의 특수한 경우일 뿐임을 지적해 둔다. 주어진 인자계 \(a_{S,T}\)에 대하여 불가약 교차표현류는 단 하나만 존재한다. 그 차수는 \(m\)이며, 여기서 \(m\)은 대응하는 대수의 슈어 지수\srcfn{10)}{I. Schur, Einige Bemerkungen zu der vorstehenden Arbeit des Herrn A. Speiser\textsuperscript{3)}. Math. Zeitschr. 5 (1919), S. 7--10. 6)에서 인용한 나의 강의에 나오는 교차표현 가군에 근거한 증명에 관해서는, 6)에서 인용한 M. Deuring의 보고서를 보라.}, 곧 대응하는 나눗셈대수의 차수이다.

%
%
\subsection*{§ 2.}
\begin{center}
\textbf{주종정리.}
\end{center}

\noindent\textbf{1. 류분할에 관한 예비 고찰.} 유체론에서 순환 상대수체에 관한 주종정리에는 두 가지 서로 다른 류분할, 즉 상위체의 절대 아이디얼류와 하위체의 레이류가 나타난다는 것이 알려져 있다. 일반적인 경우에는 이에 대응하여 상위체의 아이디얼에 대한 하나의 류분할이 인자계에 대하여 더 세밀한 아이디얼류분할을 유도한다.

우선 \(\mathfrak I\)가 \(K\)의 모든 아이디얼로 이루어진 군을 뜻한다고 하자. \(\mathfrak G\)가 \(\mathfrak I\) 위에 자기동형사상을 유도하므로, (2)부터 (5)에 따라 \(\mathfrak G\)에 의한 확장이 정의되며 \(n\)개의 잉여류 \(\{\ldots u_S\mathfrak I\ldots\}\)의 원소들로 이루어진다. \(\mathfrak I\)에서 동치관계를 부여하여 절대 아이디얼류의 군 \(\overline{\mathfrak I}\)로 넘어가도, \(n\)개의 잉여류 \(u_S\overline{\mathfrak I}\)는 여전히 정의된다. 주류가 \(\mathfrak G\)에 의해 자기 자신으로 변환되므로, \(\mathfrak G\)는 \(\overline{\mathfrak I}\) 위에도 자기동형사상을 유도하기 때문이다\srcfn{11)}{일반적으로 \(\mathfrak I\)의 부분군 중 \(\mathfrak G\)의 작용을 허용하는 것을 주류로 택하기만 하면 \(\{\ldots u_S\overline{\mathfrak I}\ldots\}\)가 정의된다. 이 경우 그 부분군의 아이디얼로부터 얻는 변환량이 주 아이디얼의 변환량을 대신한다.}. 이는 \(\mathfrak I\)에서 \(\overline{\mathfrak I}\)로 이끄는 동치관계를 \(\{\ldots u_S\mathfrak I\ldots\}\)에서 \(\{\ldots u_S\overline{\mathfrak I}\ldots\}\)로 옮길 수 있다는 뜻이다. 즉, 동치인 아이디얼 \(\mathfrak a\)와 \(\mathfrak b\)에는 서로 동치인 곱 \(u_S\mathfrak a\)와 \(u_S\mathfrak b\)를 대응시킬 수 있다. 특히 \(u_S\)는 \(u_S(c_S)\) 꼴의 모든 곱과 동치이며, 여기서 \((c_S)\)는 \(\mathfrak I\)의 모든 주 아이디얼을 두루 돈다. 이제 곱셈관계 (3)가 이 동치관계의 의미에서 일의적으로 정해진다고 요구하는 것은, 주 아이디얼로부터 얻은 모든 변환량 \((c_S^T)(c_T)/(c_{ST})\)이 단위원 인자계와 동치라는 뜻이다. 다시 말하면 다음과 같다.

\emph{\(n\)개의 잉여류로 이루어진 체계 \(\{\ldots u_S\overline{\mathfrak I}\ldots\}\)에서, \(\overline{\mathfrak I}\)가 \(K\)의 절대 아이디얼류의 군을 뜻하고 특히 \(H\)가 주류를 뜻한다고 하자. 그러면 \(u_SH\)에 대한 곱셈관계 (3)는, 유도된 인자계의 아이디얼류분할에서 주류가 주 아이디얼로부터 얻는 모든 변환량, 즉 \(H\)의 생성원 \((c_S)\)로 이루어진 변환량을 포함할 때 그리고 오직 그때 일의적으로 정의된다.}

이 사실에 이어 레이류분할의 일반화로서 다음 정의를 제시한다.

\emph{정의: 인자계의 유도된 아이디얼류분할을 다음과 같이 정한다. 주류는 모든 주 아이디얼 \((a_{S,T})\) 가운데, 각각 적어도 하나의 생성원 \(a_{S,T}\)가 있어서 대수 \((a_{S,T},K,\mathfrak G)\)가 \(K\)의 모든 분기소에서 분할되는 것들로 이루어진다.}

이 주 아이디얼들 \((a_{S,T})\)은 인자계의 결합 연산에 관하여 하나의 군을 이룬다. 이 군은 변환량 \((c_S^T)(c_T)/(c_{ST})\)을 포함한다. 실제로 인자계 \(c_S^Tc_T/c_{ST}\)는 모든 곳에서 분할되는 대수를 유도한다.

%
%
\noindent\textbf{2. 주종정리의 정식화.} 최소형 정리와 마찬가지로 이 정리를 서로 동치인 세 가지 정식으로 진술하겠다.

\emph{제1정식: 인자계의 유도된 아이디얼류분할을 바탕으로 할 때, 치환 \(v_S=u_S\bar{\mathfrak c}_S\)\srcfn{12)}{\(\bar{\mathfrak c}_S\)는 \(\mathfrak c_S\)의 아이디얼류를 뜻한다.}가 \(n\)개의 잉여류 \(\{\ldots u_S\overline{\mathfrak I}\ldots\}\)의 자기동형사상을 유도한다고 하자. 즉, 이 류분할의 의미에서 \(v_S\)에 대하여 동일한 관계 (3)이 성립한다고 하자. 그러면 이 자기동형사상은 내부자기동형사상이며 하나의 아이디얼류 \(\bar{\mathfrak b}\)에 의해 유도된다.}

\emph{제2정식: 아이디얼류 \(\bar{\mathfrak c}_S\)로부터 얻은 변환량 \(\bar{\mathfrak c}_S^T\bar{\mathfrak c}_T/\bar{\mathfrak c}_{ST}\)이 모두 인자계의 유도된 류분할의 주류에 속한다고 하자. 아이디얼류로 이루어진 이 ``벡터''들의 전체 \(\{\ldots\bar{\mathfrak c}_S\ldots\}\)가 주종을 이룬다. 그러면 류 \(\bar{\mathfrak c}_S\)들은 기호적 \((1-S)\)제곱이다. 즉, 모든 \(S\in\mathfrak G\)에 대하여 \(\bar{\mathfrak c}_S=\bar{\mathfrak b}^{1-S}=\bar{\mathfrak b}/\bar{\mathfrak b}^S\)가 되게 하는 아이디얼류 \(\bar{\mathfrak b}\)가 존재한다.}

\emph{제3정식: 인자계에 대한 유도된 아이디얼류분할을 바탕으로 할 때, 군 \(\mathfrak G\)은 \(\overline{\mathfrak I}\) 안에서 인자계의 일류에 속하는 1차 교차표현류를 단 하나만 가진다.}

서로 다른 정식들이 동치라는 것은 §~1에서와 같이 따라온다. 제1정식에서 제2정식으로 넘어가는 과정은 더 간단한데, 이제 자기동형사상이 이미 \(v_S=u_S\bar{\mathfrak c}_S\)에 의해 유도된다고 가정했기 때문이다. 이 가정은 필요하다. 이제는 \(\overline{\mathfrak I}\)가 더 이상 최대 가환부분군일 필요가 없기 때문이다(\(\overline{\mathfrak I}\)의 모든 류가 모호(ambig)할 수도 있다). 제3정식에 관해서는 다음을 지적해야 한다. \(\overline{\mathfrak I}\)에서는 곱셈 연산만 정의되어 있으므로, 표현행렬로는 각 행과 각 열에 0이 아닌 원소가 하나씩만 있는 것들만 나타난다. 1차 표현에서는 이것이 아무런 제한도 되지 않는다.

%
%
\noindent\textbf{3. 주종정리의 증명.} 여기서는 두 번째 정식화의 증명을 제시한다. 증명은 두 보조정리에 기초한다.

\emph{보조정리 1. (아이디얼의 주종정리): 아이디얼 \(\mathfrak c_S\)로부터 만든 변환량 \(\mathfrak c_S^T\mathfrak c_T/\mathfrak c_{ST}\)가 단위아이디얼을 이루면, \(\mathfrak c_S\)는 기호적 \((1-S)\)제곱이다. 즉 \(\mathfrak G\)의 모든 \(S\)에 대하여 \(\mathfrak c_S=\mathfrak b^{1-S}\)이다. 다시 말해 이는 첫 번째 및 세 번째 정식화와 동치이며, 첫 번째 정식화에서는 아이디얼류의 경우와 마찬가지로 자기동형사상이 \(v_S=u_S\mathfrak c_S\)에 의해 생성된다고 가정해야 한다.}

\emph{증명}\srcfn{13)}{이에 관해서는 서론의 끝부분을 보라.}. \(\mathfrak b=\sum\mathfrak c_S\)로 두면 원하는 바가 성립한다. 여기서 합은 가군합, 즉 최대공약아이디얼을 뜻한다. 실제로 다음을 얻는다.
\[
\mathfrak b=\sum\mathfrak c_S;\quad
\mathfrak b^T=\sum\mathfrak c_S^T;\quad
\mathfrak b^T\mathfrak c_T=\sum\mathfrak c_S^T\mathfrak c_T
=\sum\mathfrak c_{ST}=\mathfrak b;\quad
\text{따라서}\quad\mathfrak c_T=\mathfrak b^{1-T}.
\]

\emph{보조정리 2. (분할대수에 관한 정리의 다른 정식화): 대수 \(A=(a_{S,T},K,\mathfrak G)\)가 \(K\)의 모든 (유한 및 무한) 분기소에서 분할되고, 주아이디얼 \((a_{S,T})\)가 아이디얼의 변환량, 즉 \((a_{S,T})=\mathfrak c_S^T\mathfrak c_T/\mathfrak c_{ST}\)를 이루면, \(A\)는 모든 곳에서, 곧 전역적으로 분할된다. 따라서 \((a_{S,T})\)는 주아이디얼의 변환량 \((a_{S,T})=(d_S^T)(d_T)/(d_{ST})\)가 된다\srcfn{14)}{더 정확히는 \(a_{S,T}=d_S^Td_T/d_{ST}\)이다. 그러나 주종정리의 증명에서는 주아이디얼에 관한 더 약한 명제만을 사용한다.}. 역으로, 모든 곳에서 분할되는 각 대수는 이 조건을 만족한다.}

%
%
증명은 \(\mathfrak p\)-진 확대에서 나타나는 \(A\)의 잘 알려진 거동에 기초한다. \(k_{\mathfrak p}\)는 \(k\)의 \(\mathfrak p\)-진 확대를 뜻하고, 이에 대응하여 \(K_{\mathfrak P}\)는 \(K\)의 \(\mathfrak P\)-진 확대를 뜻한다. 여기서 \(\mathfrak P\)는 \(\mathfrak p\) 위에 놓인 소아이디얼이다. \(A_{\mathfrak p}\)는 \(k\)에서 \(k_{\mathfrak p}\)로 이행하여 얻는 \(A\)의 확대라 하자. 또한 \(\mathfrak Z\)를 \(\mathfrak P\)에 대한 분해군이라 하고, \(a_{\mathfrak Z}\)를 인자계의 대응하는 부분이라 하자(즉 \(a_{\mathfrak Z}\)는 \(S\)와 \(T\)가 \(\mathfrak Z\)에 속하는 \(a_{S,T}\)들로 이루어진다). 그러면
\[
A_{\mathfrak p}\sim(a_{\mathfrak Z},K_{\mathfrak P}/k_{\mathfrak p},\mathfrak Z)\srcfn{15)}{이에 관해서는 Hasse\textsuperscript{6)}, 14를 보라. 조금 더 간단한 증명은 다음과 같다. \(Z\)를 \(\mathfrak P\)에 대한 분해체라 하자. 그러면 잘 알려진 바와 같이 \(k_{\mathfrak p}\)는 이에 동형인 \(Z\)를 포함한다. 이때 \(A_Z\sim(a_{\mathfrak Z},K/Z,\mathfrak Z)\)가 성립한다. 실제로 \(A_Z\)는 \(A\)에서 \(Z\)의 모든 원소와 가환하는 원소 전체와 상사이다(예컨대 v. d. Waerden, Bd. II, S. 210 또는 Math. Zeitschr.에 게재될 나의 비가환대수 논문을 보라). \(A_Z\)를 \(A_{\mathfrak p}\)로 확대하면 위 결과를 얻는다. 서로 켤레인 \(\mathfrak P\) 가운데 어느 것을 택하는지는 무관하다. \(\mathfrak G=\sum\mathfrak ZT\)일 때 생기는 대수들은 서로 동형이며 \(u_T\)에 의한 변환으로 서로 옮겨진다.}.
\]

특히 \(\mathfrak p\)가 \emph{비분기}이면, 즉 \(\mathfrak Z\)가 순환군이면, 이는 비분기 관성체 \(K_{\mathfrak P}/k_{\mathfrak p}\)에 의한 \(A_{\mathfrak p}\)의 표준적 표현이다\srcfn{16)}{예컨대 4)에서 인용한 H. Hasse의 논문 § 3을 보라.}.

보조정리 2의 가정 아래에서 이 비분기 경우에도 \(A_{\mathfrak p}\)가 분할됨을 보여야 한다. 무한소 \(\mathfrak p\)에 대해서는 \(K_{\mathfrak P}=k_{\mathfrak p}\)이므로 \(A_{\mathfrak p}\)가 분할된다. 유한소 \(\mathfrak p\)에 대해서는 우선 인자계 \(a_{\mathfrak Z}\)가 \(K_{\mathfrak P}\)의 단위원들로 이루어진 인자계와 연관됨을 알 수 있다. 실제로 \(K_{\mathfrak P}\)에서는 모든 아이디얼 \(\mathfrak c_S\)가 주아이디얼 \((c_S)\)가 된다. 따라서 \((a_{S,T})\)에 관한 가정은 단위원 \(e_{S,T}\)를 써서 \(a_{S,T}=e_{S,T}c_S^Tc_T/c_{ST}\)라는 뜻이며, 인자계 \(a_{\mathfrak Z}\)는 단위원 \(e_{S,T}\)들로 이루어진 인자계와 연관된다.

%
%
더 나아가 \emph{정규화된 순환표현} (1')에서 (5')로 이행해도 이 성질은 보존된다. 따라서 \(R\)이 \(\mathfrak Z\)의 생성원일 때 \(A_{\mathfrak p}=(\alpha,K_{\mathfrak P}/k_{\mathfrak p},R)\)에서 \(\alpha\)는 단위원이 된다. 실제로 \(e_{S,T}\)가 인자계라는 사실은 다음 관계식으로 표현된다.
\[
u_{R^i}u_{R^j}=u_{R^{i+j}}e_{R^i,R^j}.
\]
그런데 \(u_R=u\)로 두면 이로부터 다음 관계식을 얻는다.
\[
u^i=u_{R^i}e_{R,R^{i-1}}e_{R,R^{i-2}}\ldots e_{R,R};
\]
따라서 특히 \(u_E=e_{E,E}\)이고 \(f\)가 \(\mathfrak Z\)의 위수일 때
\[
u^f=\alpha=e_{E,E}e_{R,R^{f-1}}\ldots e_{R,R}.
\]
그러나 \(K_{\mathfrak P}/k_{\mathfrak p}\)가 비분기이면 모든 단위원이 노름이므로\textsuperscript{16)}, \(A_{\mathfrak p}\)가 분할된다. 따라서 \(A\)는 모든 곳에서 분할되고, 그러므로 분할대수에 관한 정리에 따라 전역적으로 분할된다. 달리 말하면, 가정으로부터 \(a_{S,T}\)가 원소들로부터 얻은 변환량, 즉 \(K^*\)의 원소 \(d_S\)를 써서 \(a_{S,T}=d_S^Td_T/d_{ST}\)임이 따른다. 따라서 더 약하게는 \((a_{S,T})\)가 주아이디얼의 변환량이 된다. 역으로 \(A\)가 모든 곳에서 분할되면, 각 장소에서 그러하므로 주아이디얼 \((a_{S,T})\) 역시 아이디얼의 변환량이 되며, 특히 \(A\)는 분기소에서 분할된다. 이는 실제로 분할대수에 관한 정리의 가정을 달리 정식화한 것이다.

%
%
두 보조정리로부터 \emph{주종정리의 증명}이 곧바로 따라 나온다. 각 \(\mathfrak c_S\)가 류 \(\bar{\mathfrak c}_S\)에 속하는 하나의 아이디얼을 뜻한다고 하자. 그러면 류들에 관한 가정은---인자계의 유도된 류분할의 정의를 고려하면---\(\mathfrak c_S\)의 변환량이 보조정리 2의 가정을 만족하는 주아이디얼 \((a_{S,T})\)이 된다는 바로 그 사실을 말한다. 따라서 주아이디얼 \((d_S)\)를 적절히 택하여 \(\mathfrak d_S=\mathfrak c_S/(d_S)\)로 둔 아이디얼들이 여전히 류 \(\bar{\mathfrak c}_S\)에 속하면서 보조정리 1의 가정을 만족하게 할 수 있다. 즉 \(\mathfrak d_S\)의 변환량은 단위아이디얼과 같아진다. 그러므로 보조정리 1에 따라 \(\mathfrak d_S=\mathfrak b^{1-S}\)를 만족하는 아이디얼 \(\mathfrak b\)가 존재한다. 이는 류 \(\bar{\mathfrak c}_S\)가 류 \(\bar{\mathfrak b}\)의 기호적 \((1-S)\)제곱이 된다는 뜻이다.

순환 경우의 주종정리가 이미 알려진 정리로 귀착된다는 사실은 다음으로부터 따른다. 즉 정규화를 기초로 삼으면, 유도된 류분할의 주류는 모든 분기소에서 \(\alpha\)가 노름잉여가 되는 주아이디얼 \((\alpha)\)들의 군으로 귀착된다.

\begin{center}
(1932년 10월 27일 접수.)
\end{center}

\clearpage
\setcounter{footnote}{0}
\phantomsection
\addcontentsline{toc}{section}{42. 분할 교차곱과 그 극대 오더}
%
%

%
%
\section*{42. 분할 교차곱과 그 극대 오더}

\begin{center}
\emph{Actualités scientifiques et industrielles 148 (1934), 5--15쪽}
\end{center}

분할 교차곱의 가장 단순한 경우\srcfn{(1)}{교차곱의 이론은 나의 한 강의에 이어 H. Hasse, \emph{Theory of cyclic algebras over an algebraic numberfield}, Transact. Amer. Math. Soc. 34 (1932), 제II절에 서술되어 있다. 여기서는 최초의 기본개념들만을 사용한다. 브란트의 극대 오더 이론은 더 광범한 결과들과 함께 H. Hasse, \emph{Ueber $p$-adische Schiefkörper und ihre Bedeutung für die Arithmetik hyperkomplexer Zahlsysteme}, Math. Ann. 104 (1931) (이하 H.로 인용)에 서술되고 상세히 새로 정초되어 있다.}---곧 분할대수를 이루며 따라서 $1$과 동치인 인자계를 갖는 경우---에는 그 사정을 상당한 범위까지 명시적으로 추적할 수 있다. 이는 비갈루아 분해체(극대 가환 부분체)의 경우에도 여전히 성립한다. 여기에 속하는 단순한 대수적 사실들은 § 1에서 전개한다. § 2에서는 기초체가 대수적 수체일 때, 바탕이 되는 분해체 $k$의 가군과 대각합-쌍대 가군을 이용하여 모든 극대 오더와 그 아이디얼의 명시적 표현을 제시한다\srcfn{(2)}{이 명시적 표현들은 오래전부터 알고 있었다. 이어지는 고찰의 계기는 아래 (4)에서 인용한 H. Hasse의 정리를 전해 들은 일이었다.}. 또한 모든 극대 오더 가운데 $k$와의 교집합이 같은 것들을 하나의 ``구역''으로 묶어 구분한다. 이로써 구역들은 $k$의 오더들과 일대일로 대응하며, $k$의 정수환에는 ``주구역''이 대응한다. 한 구역 안의 극대 오더들은 해당 오더의 가군들로 만든 아이디얼에 의한 변환으로 서로 옮겨진다. 특히 주구역에서는 $k$의 아이디얼들(정수환에 관한 가군들)을 확장한 아이디얼들이 쓰인다. § 3에서는 임의의 단순대수로 넘어간다. 구역의 구분을 세분하여---교집합뿐 아니라 분기소에서 극대 오더들의 $p$-진 성분도 일치해야 한다---분할의 경우에 얻은 정리들을 그대로 옮길 수 있다. 특히 이때에도 한 주구역의 극대 오더들은 $k$의 아이디얼을 확장한 아이디얼들에 의한 변환으로 서로 옮겨진다.

주구역에 관한 정리들은 Chevalley\srcfn{(3)}{Sur certains idéaux d'une algèbre simple, Abh. math. Sem. Hambg., 1934.}와 Hasse\srcfn{(4)}{Ueber gewisse Ideale in einer einfachen Algebra, 이 논문 바로 앞에 실린 글; 이 모음집의 I.}에게서도 볼 수 있다. Chevalley는 분기소에 관한 위의 세분된 구분을 두지 않으면 이 정리들이 일반적으로 더 이상 성립하지 않음도 보였다.

\begin{center}
\textbf{§ I. --- \emph{분할 교차곱의 행렬단위와\\
분해체의 쌍대기저.}}
\end{center}

%
%
\noindent\textbf{1. ---} 이 절에서 $\Omega$는 임의의 기초체를 뜻한다. $k/\Omega$를 차수 $n$의 갈루아 분리확대라 하고, 그 갈루아군을 원소 $\mathrm{S},\mathrm{T},\ldots$를 갖는 $\Gg$라 하며, 이에 대응하는 연산자를 $u_s,u_t,\ldots$라 하자. 여기서는 \emph{인자계가 $1$인 교차곱}, 곧 분할대수
\[
\mathrm{K}=\Gg\ast k=u_{s_1}k+\ldots+u_{s_n}k,\qquad u_su_t=u_{st},
\]
\[
zu_s=u_sz^s\quad\text{($k$의 모든 $z$에 대하여)}
\]
를 고찰한다. 인자계가 $1$이면 다음 관계들도 얻어진다. 즉
\[
E_\Gg=\sum_Su_S
\]
로 놓으면, $\Gg$의 모든 $S$에 대하여
\srcnumdisplay[3.2em]{(1)}{\mathrm{E}_{\mathfrak g}u_s=u_s\mathrm{E}_{\mathfrak g}=\mathrm{E}_{\mathfrak g}}
가 성립한다(곧 $\mathrm{E}_{\mathfrak g}$가 $\Gg$의 자명표현을 생성한다\srcfn{(5)}{$k$의 표수가 $n$을 나누지 않으면 $\frac1n\mathrm{E}_{\mathfrak g}$가 자명표현의 멱등원이 된다. 그러나 $\mathrm{K}$의 표준형은 표수가 $n$을 나누는 경우에도 성립한다. 유일한 가정은 $k/\Omega$가 분리확대라는 것이다.}). 또한
\srcnumdisplay[3.2em]{(2a)}{z\mathrm{E}_{\mathfrak g}=\sum_su_sz^s,}
\srcnumdisplay[3.2em]{(2b)}{\mathrm{E}_{\mathfrak g}z\mathrm{E}_{\mathfrak g}=\mathrm{E}_{\mathfrak g}\cdot\sum_su_sz^s=\mathrm{E}_{\mathfrak g}\cdot\Sp(z),}
이며, 여기서 $\Sp(z)$는 $k/\Omega$의 원소 $z$의 대각합(트레이스)을 뜻한다. 이 관계들로부터 다음 정리가 나온다.

\medskip
\noindent\textsc{정리. ---} \emph{$a_1,\ldots,a_n$이 $k/\Omega$의 한 기저이고 $\widetilde a_1,\ldots,\widetilde a_n$이 이에 대한 쌍대기저라 하자. 그러면 $c_{ik}=\widetilde a_i\mathrm{E}_{\mathfrak g} a_k$는 $\mathrm{K}$의 행렬단위계를 이룬다. 따라서 $\mathrm{K}$는 다음 꼴의 가군곱으로 나타낼 수 있다.}
\[
\mathrm{K}=k\mathrm{E}_{\mathfrak g}k=k\cdot\sum_su_s\cdot k
 =\sum_{i,k}(\widetilde a_i\mathrm{E}_{\mathfrak g}a_k)\,\Omega.
\]
실제로 (2b)에 의해
\[
c_{ij}c_{lk}
 =\widetilde a_i\mathrm{E}_{\mathfrak g}a_j\cdot\widetilde a_l\mathrm{E}_{\mathfrak g}a_k
 =\widetilde a_i\mathrm{E}_{\mathfrak g}a_k\cdot\Sp(a_j\widetilde a_l)
 =c_{ik}\cdot\Sp(a_j\widetilde a_l)
\]
이다. 그런데 쌍대기저의 정의식 $\Sp(a_j\widetilde a_l)=\delta_{jl}$ 때문에 이것이 바로 행렬단위의 성질이다.

\emph{비고 1.} $\mathrm{Z}$가 $\Omega$의 임의의 (가환) 확대체라면, $a_i$를 $k_{\mathrm Z}/\mathrm{Z}$의 기저로 이해하고 $k_{\mathrm Z}$가 체들의 직합이 되는 경우에도 위의 모든 사실은 그대로 성립한다. 이때 군의 치환들은 $\mathrm{Z}$ 위에서 항등사상을 유도하도록 정의된다. 이는 $K_Z$ 안의 $u_s$가 $\mathrm{Z}$의 모든 원소와 가환한다는 사실과 일치한다.

\emph{비고 2.} $\mathrm{K}=k\mathrm{E}_{\mathfrak g}k$라는 사실은 우변이 단순대수 $\mathrm{K}$의 영이 아닌 아이디얼이라는 데서도 곧바로 쉽게 따른다.

%
%
\medskip
\noindent\textbf{2. ---} 이제 기초체 $\Omega$의 \emph{표수가 $0$}이라고 가정하자. \emph{$k$가 $\mathrm{K}/\Omega$의 비갈루아 분해체인 경우에도 분할대수 $\mathrm{K}$에 대한 1의 표현 $\mathrm{K}=k\mathrm{E}_{\mathfrak g}k$는 그대로 성립한다.} 이는 갈루아 분해체로 옮겨 가면 알 수 있다. $\bar k$를 $k$의 갈루아 폐포, $\bar\Gg$를 그 군, $\bar u_s$를 대응하는 연산자라 하고, $\bar{\mathrm K}=\bar k\mathrm{E}_{\bar{\mathfrak g}}\bar k$를 분할 교차곱이라 하자. 부분체 $k$에 위수 $\mathrm{h}$인 부분군 $\Hh$가 대응한다고 하자. 다음과 같이 놓는다.
\[
\bar\Gg=\Hh\mathrm{S}_1+\ldots+\Hh\mathrm{S}_n
 =\mathrm{S}_1^{-1}\Hh+\ldots+\mathrm{S}_n^{-1}\Hh;
\qquad
\mathrm{E}_{\mathfrak h}=\frac1h\sum\bar u_{\mathfrak h},
\]
\[
\mathrm{E}_{\mathfrak g}=\frac1h\mathrm{E}_{\bar{\mathfrak g}}=\sum_i u_{s_i},
\qquad\text{여기서}\qquad
u_{s_i}=\mathrm{E}_{\mathfrak h}\bar u_{s_i}.
\]
특히 멱등원 $\mathrm{E}_{\mathfrak h}$는 $\Hh$의 자명표현을 생성하고, $\mathrm{E}_{\mathfrak g}$는 $\mathrm{E}_{\bar{\mathfrak g}}$와 마찬가지로 $\bar\Gg$의 자명표현을 생성한다. \emph{이렇게 정의한 $\mathrm{E}_{\mathfrak g}$와 $u_s$에 대해서는 관계 (1)이 한쪽으로 성립하고 관계 (2)는 완전히 성립한다.} 여기서 $z$는 $k$의 모든 원소를 지나며, 따라서 $z^s$는 공액체 $k^s$의 원소들을 지난다. 이 공액체는 $\bar k$의 부분체이므로 $\bar{\mathrm K}$의 부분체이기도 하다. 먼저 다음이 성립한다.
\srcnumdisplay[3.2em]{(3)}{\mathrm{E}_{\mathfrak g}=\mathrm{E}_{\mathfrak g}\mathrm{E}_{\mathfrak h}=\mathrm{E}_{\mathfrak h}\mathrm{E}_{\mathfrak g}=\mathrm{E}_{\mathfrak h}\mathrm{E}_{\mathfrak g}\mathrm{E}_{\mathfrak h},}
\srcnumdisplay[3.2em]{(4)}{z\mathrm{E}_{\mathfrak h}=\mathrm{E}_{\mathfrak h}z,\quad\text{따라서}}
\srcnumdisplay[3.2em]{(5)}{zu_s=u_sz^s.}
실제로 (3)은 $\mathrm{E}_{\bar{\mathfrak g}}$에 대하여 $\mathrm{E}_{\mathfrak h}$의 정의와 (1)로부터 성립하므로 $\mathrm{E}_{\mathfrak g}$에 대해서도 성립한다. 한편 (4)는 $k$의 원소인 $z$가 $\Hh$의 치환들을 허용한다는 것, 곧 $z\bar u_h=\bar u_hz$라는 뜻이다. (4)의 양변에 오른쪽에서 $\bar u_{s_i}$를 곱하면 (5)를 얻는다. 이제 (1)은 한쪽으로 성립한다.
\srcnumdisplay[3.2em]{(1$'$)}{\mathrm{E}_{\mathfrak g}u_s=\mathrm{E}_{\mathfrak g},\quad\text{왜냐하면 }\mathrm{E}_{\mathfrak g}u_s=\mathrm{E}_{\mathfrak g}\mathrm{E}_{\mathfrak h}\bar u_s=\mathrm{E}_{\mathfrak g}\bar u_s=\mathrm{E}_{\mathfrak g};}
(5)를 합하면 (2a)가 나오고, 여기에 (1$'$)을 쓰면 (2b)가 나온다.

$\mathrm{K}=k\mathrm{E}_{\mathfrak g}k$라는 증명은 이제 1에서와 똑같이 얻어진다. 실제로 $\widetilde a_i\mathrm{E}_{\mathfrak g}a_k$가 행렬단위임을 보일 때에는 (2b)만을 사용하였다. $\mathrm{K}$는 $\Omega$ 위의 전행렬환이므로 차수 $n$인 모든 체를 부분체로 포함하며, 특히 $k$도 포함한다.

%
%
\medskip
\noindent\textbf{3. ---} § 3을 염두에 두고, $\mathrm{K}$가 분할되지 않는 경우에도 $K$에서 $\bar{\mathrm K}$로 옮겨 갈 수 있음을 덧붙이고자 한다\srcfn{(6)}{이로부터 비갈루아의 경우에 대한 교차곱 표현의 유사물을 얻을 수 있다. 이는 아마도 한 학위논문에서 더 전개될 것이다.}. $k$뿐 아니라 $\bar k$도 분해체이므로 $\mathrm{K}$는 $\bar{\mathrm K}$, 곧 $\bar\Gg$와 $\bar k$로 이루어진 교차곱과 상사이다. 이때 $k$가 분해체이므로 $k$에 대응하는 부분군 $\Hh$에는 인자계 $1$이 대응한다. 따라서 멱등원 $\mathrm{E}_{\mathfrak h}$는 2에서와 같이 정의된다. 이를 거꾸로 적용하면 $\mathrm{K}$가 $\mathrm{E}_{\mathfrak h}\bar{\mathrm K}\mathrm{E}_{\mathfrak h}$와 동형임을 얻는다. 실제로 이 대수는 왼쪽 아이디얼 $\bar{\mathrm K}\mathrm{E}_{\mathfrak h}$의 자기준동형환과 동형이므로 $\bar{\mathrm K}$와 상사이다. 그 랭크는 $\mathrm{K}$의 랭크와 일치한다. 왜냐하면 $\bar{\mathrm K}\mathrm{E}_{\mathfrak h}$의 $\bar k$ 위 랭크는 $n$이고($\bar u_{s_i}^{-1}\mathrm{E}_{\mathfrak h}$들은 $\bar k$ 위에서 독립이다), $\bar{\mathrm K}$의 $\bar k$ 위 랭크는 $nh$이기 때문이다. 이는 $\bar{\mathrm K}$가 $\bar{\mathrm K}\mathrm{E}_{\mathfrak h}$와 동형인 $h$개의 아이디얼로 분해된다는 뜻이며, 자기준동형환에서는 차수 $h$의 전행렬대수를 곱하는 데 해당한다.

$\mathrm{K}$, 따라서 $\bar{\mathrm K}$가 분할되는 경우에는 이렇게 하여 2의 결과를 다시 얻는다. 실제로
\[
\begin{aligned}
\mathrm{E}_{\mathfrak h}\bar{\mathrm K}\mathrm{E}_{\mathfrak h}
 =\mathrm{E}_{\mathfrak h}\bar k\mathrm{E}_{\bar{\mathfrak g}}\bar k\mathrm{E}_{\mathfrak h}
 =(\mathrm{E}_{\mathfrak h}\bar k\mathrm{E}_{\mathfrak h})\mathrm{E}_{\bar{\mathfrak g}}(\mathrm{E}_{\mathfrak h}\bar k\mathrm{E}_{\mathfrak h})
 &=\Sp(\bar k/k)\mathrm{E}_{\mathfrak h}\mathrm{E}_{\bar{\mathfrak g}}\mathrm{E}_{\mathfrak h}\Sp(\bar k/k)\\
 &=k\mathrm{E}_{\mathfrak g}k.
\end{aligned}
\]

%
%
\begin{center}
\textbf{§ II. --- \emph{분할 교차곱의 극대 오더와\\
그 구역별 구분.}}
\end{center}

이제부터 $\Omega$를 (유한) 대수적 수체\srcfn{(7)}{이 절의 고찰은 $\Omega$가 한 부정원의 함수체인 경우에도 그대로 성립한다.}, $\mathrm{K}$를 $\Omega$ 위의 차수 $n$인 분할대수, $k$를 갈루아 또는 비갈루아 극대 가환 부분체라 하자. 그러면 § 1에 의해 $\mathrm{K}=k\mathrm{E}_{\mathfrak g}k$이다. 또한 $\mathfrak a,\widetilde{\mathfrak a}$를 $k$ 안의 서로 대각합-쌍대인 유한 $\oo$-가군의 한 쌍이라 하자($o$는 $\Omega$의 정수환이고 $\oo$는 $k$의 정수환이다). 문제는 $k$에 상대적인 $K$의 극대 오더들의 구조이며, 이는 다음 정리들로 기술된다.

\medskip
\noindent\textsc{정리 1. ---} \emph{가군곱 $\mathfrak O_{\mathfrak a}=\widetilde{\mathfrak a}\mathrm{E}_{\mathfrak g}\mathfrak a$는 $\mathrm{K}$의 극대 오더이다. 특히 $\mathfrak O=\mathfrak O_{\mathfrak o}=\widetilde{\mathfrak o}\mathrm{E}_{\mathfrak g}\mathfrak o$도 극대 오더이다.}

\medskip
\noindent\textsc{정리 2. ---} \emph{$\mathfrak O$에서 $\mathfrak O_{\mathfrak a}$로의 변환은 $\mathfrak O$에 관한 왼쪽 아이디얼 $\mathfrak L=\widetilde{\mathfrak o}\mathrm{E}_{\mathfrak g}\mathfrak a$와 이에 역인 오른쪽 아이디얼 $\mathfrak R=\widetilde{\mathfrak a}\mathrm{E}_{\mathfrak g}\mathfrak o$에 의해 이루어진다.}

\emph{비고.} 서로 쌍대인 두 가군 $\mathfrak a,\widetilde{\mathfrak a}$가 곱 $\mathfrak L\mathfrak R$ 안에서 만날 때 실제로 역관계가 생기는 것은 (단순한 곱 $\mathfrak a\widetilde{\mathfrak a}$의 경우와 달리) $\mathrm{E}_{\mathfrak g}$가 일으키는 대각합 연산 때문이다.

\medskip
\noindent\textsc{정리 3. ---} \emph{$\mathfrak O$에 관한 모든 왼쪽 및 오른쪽 아이디얼은 정리 2에서 제시한 것들로 모두 주어진다. 따라서 극대 오더 $\mathfrak O_{\mathfrak a}$들은 $\mathrm{K}$의 모든 극대 오더를 빠짐없이 이룬다. 임의의 두 극대 오더 $\mathfrak O_{\mathfrak a}$와 $\mathfrak O_{\mathfrak b}$는 $\mathfrak L_{\mathfrak a}=\widetilde{\mathfrak a}\mathrm{E}_{\mathfrak g}\mathfrak b$, $\mathfrak R_{\mathfrak a}=\widetilde{\mathfrak b}\mathrm{E}_{\mathfrak g}\mathfrak a$에 의한 변환으로 서로 옮겨지며, $\mathfrak L_{\mathfrak a}$와 $\mathfrak R_{\mathfrak a}$는 $\mathfrak O_{\mathfrak a}$에 관한 모든 왼쪽 및 오른쪽 아이디얼을 빠짐없이 이룬다.}

\medskip
\noindent\textsc{정리 4. ---} \emph{$\mathfrak O_{\mathfrak a}$가 $\mathrm{K}$의 모든 극대 오더를 지날 때, 교집합 $[\mathfrak O_{\mathfrak a},k]$는 $k$의 모든 오더를 지나며 각 경우 이 교집합은 $\mathfrak a$의 오더와 같다. $k$ 안의 같은 오더에 속하는 극대 오더들을 하나의 구역으로 묶으면, 한 구역 안의 변환은 $\mathfrak O_{\mathfrak a}$에 관한 왼쪽 아이디얼 $\widetilde{\mathfrak a}\mathrm{E}_{\mathfrak g}\mathfrak b$와 이에 역인 오른쪽 아이디얼들에 의해 이루어진다. 여기서 $\mathfrak a$와 $\mathfrak b$는 $k$ 안의 해당 오더에 관한 가군들이다.}

\medskip
\noindent\textsc{정리 4a. ---} \emph{특히 주구역---곧 $k$의 정수환 $o$를 포함하는 모든 극대 오더의 구역---에 속한 극대 오더들의 변환은 $\widetilde{\mathfrak a}\mathrm{E}_{\mathfrak g}\mathfrak b$ 꼴의 아이디얼들에 의해 이루어지며, 여기서 $\mathfrak a$와 $\mathfrak b$는 $k$의 아이디얼들이다. 다시 말해 $\mathfrak L$이 주구역의 한 극대 오더의 왼쪽 아이디얼이고 $\mathfrak L$의 오른쪽 오더가 $o$를 포함한다면, $\mathfrak L$은 한 $\mathfrak o$-아이디얼의 확장이다.}

%
%
증명들은 각 소로 옮겨 가는 데 바탕을 둔다. 이를 위해서는 매번 $p$에 대한 국소화환으로 옮겨 가는 것으로 충분하며, 여기서 $p$는 $\Omega$의 모든 소아이디얼을 지난다. $o$를 $\Omega$의 정수환이라 하고 $o_p$를 $p$에 대한 국소화환, 곧 $\Omega$의 수 가운데 분모가 $p$와 서로소인 것들로 나타낼 수 있는 모든 $p$-정수적 원소의 주아이디얼환이라 하자. $\mathfrak C$가 $\mathrm{K}$의 임의의 $o$-가군이면, 확장가군 $\mathfrak C_p=\mathfrak C\oo_p$를 그 $p$-성분이라 부르기로 한다. 그러면 다음 보조정리가 성립한다.

\medskip
\noindent\emph{보조정리. $\mathfrak C$는 모든 확장가군 $\mathfrak C_p$의 교집합이다.} 실제로 한 원소 $w$가 $\mathfrak C_p$에 들어 있으면, $p$로 나누어지지 않는 $o$의 원소 $\sigma$가 존재하여 $\sigma w$가 $\mathfrak C$에 들어간다. 이제 $w$가 모든 $\mathfrak C_p$의 교집합에 들어 있다고 하고, $\sigma w$가 $\mathfrak C$에 들어가게 하는 $o$의 모든 $\sigma$로 이루어진 아이디얼을 $g$라 하자. 방금 보인 바에 따르면 $g$는 어느 $p$로도 나누어질 수 없으므로 $g=o$이다\srcfn{(8)}{소를 $p$-진 확대에 의해 정의하면 대응하는 보조정리는 다음과 같다. 모든 $p$-진 성분 $\mathfrak C^{(p)}$의 교집합을 $\mathrm{K}$와 취하면 $\mathfrak C$가 된다. 이는 교집합 $[\mathfrak C^{(p)},\mathrm{K}]=\mathfrak C_p$라는 사실에서 따른다. 그 사실은 $\mathfrak C_p$의 한 $o_p$-기저가 동시에 $\mathfrak C^{(p)}$의 기저가 됨을 이용하면 알 수 있다. [H.의 정리 66에서는 이 보조정리가 (최대 랭크인) 아이디얼에 대해서만 진술되어 있다.]}. 이 보조정리는 가군 $\mathfrak C$의 랭크에 관하여 아무 가정도 하지 않는다는 점을 덧붙여 둔다.

\emph{따름정리들.} 따라서 모든 가군은 그 성분들에 의해 결정된다. $\mathfrak D$가 $\mathfrak C$의 진상위가군이면, 적어도 하나의 $\mathfrak D_p$가 $\mathfrak C_p$의 진상위가군이 된다. 특히 $\mathfrak M$이 $\mathrm{K}$의 오더이고 모든 성분 $\mathfrak M_p$가 $o_p$에 관한 극대 오더라면, $\mathfrak M$도 극대 오더이다.

%
%
\medskip
\noindent\emph{정리 1의 증명.} $\mathfrak O_{\mathfrak a}$는 오더이다. 실제로 $a$와 $\widetilde{\mathfrak a}$가 유한 $\oo$-가군이므로 $\mathfrak O_{\mathfrak a}$도 유한 $\oo$-가군이고, 더욱이 최대 랭크이다. 또한
\[
\mathfrak O_a^2<\mathfrak O_a;
\quad\text{실제로}\quad
\mathfrak O_a^2
 =\widetilde aE_\Gg a\cdot\widetilde aE_\Gg a
 =\widetilde aE_\Gg a\cdot\Sp(\widetilde a a)\text{\textsuperscript{(9)}}
\]
\srcfntext{(9)}{$\Sp(\mathfrak c)$는 일반적으로 $\mathfrak c$의 모든 원소의 대각합들로 이루어진 $o$의 아이디얼을 뜻한다.}
이다(§ 1에 의함). $\Sp(\widetilde{\mathfrak a}\mathfrak a)<o$이므로 $\mathfrak O_{\mathfrak a}^2<\mathfrak O_{\mathfrak a}$가 따른다.

$\mathfrak O_{\mathfrak a}$가 극대임을(따라서 $o$를 포함함을) 보이려면, 위의 따름정리에 따라 각 성분이 극대임을 보이는 것으로 충분하다. $o_p$가 주아이디얼환이므로 $\mathfrak a_p$와 $\widetilde{\mathfrak a}_p$에는 서로 독립인 쌍대기저 $a_1,\ldots,a_n$과 $\widetilde a_1,\ldots,\widetilde a_n$이 존재한다. § 1에 따라 곱 $\widetilde a_i\mathrm{E}_{\mathfrak g}a_k$들은 행렬단위계를 이룬다. 따라서
\[
\widetilde a_pE_\Gg a_p=\sum_{i,k}c_{ik}o_p
\]
는 극대이다.

\emph{비고.} $\mathfrak O_{\mathfrak a}$가 극대이므로 $\mathfrak O_{\mathfrak a}^2=\mathfrak O_{\mathfrak a}$이고, 이에 따라 $\Sp(\widetilde{\mathfrak a}\mathfrak a)=o$도 얻어진다. 이 관계가 모든 소에서 성립한다는 것은 쌍대기저로부터 알 수 있으므로, 보조정리에 의해 이 결론을 직접 얻을 수도 있다.

%
%
\medskip
\noindent\emph{정리 2의 증명.} 증명은 방금 확립한 대각합 규칙에서 따른다. 그 규칙에 의해 다음을 얻는다.
\[
\begin{aligned}
\mathfrak O\mathfrak L
 &=\widetilde{\mathfrak o}\mathrm{E}_{\mathfrak g}\mathfrak o\,\widetilde{\mathfrak o}\mathrm{E}_{\mathfrak g}\mathfrak a
 =\widetilde{\mathfrak o}\mathrm{E}_{\mathfrak g}\mathfrak a=\mathfrak L,
&\qquad
\mathfrak R\mathfrak O
 &=\widetilde{\mathfrak a}\mathrm{E}_{\mathfrak g}\mathfrak o\,\widetilde{\mathfrak o}\mathrm{E}_{\mathfrak g}\mathfrak o
 =\widetilde{\mathfrak a}\mathrm{E}_{\mathfrak g}\mathfrak o=\mathfrak R,\\
\mathfrak L\mathfrak R
 &=\widetilde{\mathfrak o}\mathrm{E}_{\mathfrak g}\mathfrak a\,\widetilde{\mathfrak a}\mathrm{E}_{\mathfrak g}\mathfrak o
 =\widetilde{\mathfrak o}\mathrm{E}_{\mathfrak g}\mathfrak o=\mathfrak O,
&
\mathfrak R\mathfrak L
 &=\widetilde{\mathfrak a}\mathrm{E}_{\mathfrak g}\mathfrak o\,\widetilde{\mathfrak o}\mathrm{E}_{\mathfrak g}\mathfrak a
 =\widetilde{\mathfrak a}\mathrm{E}_{\mathfrak g}\mathfrak a=\mathfrak O_{\mathfrak a}.
\end{aligned}
\]

%
%
\medskip
\noindent\emph{정리 3의 증명.} $\mathfrak L$을 $\mathfrak O$에 관한 왼쪽 아이디얼이라 하자. 특히
\[
\mathfrak L=\widetilde oE_\Gg\mathfrak L
 =(\widetilde oE_\Gg l_1,\ldots,\widetilde oE_\Gg l_t)
\]
이며, 여기서 예를 들어 $l_1,\ldots,l_t$는 $\mathfrak L$의 한 $o$-기저이다. $l_\mu=\sum_i u_{s_i}z_{\mu,i}$로 놓고 $z_{\mu,i}$를 $k$의 원소라 하면\srcfn{(10)}{비갈루아 $k$의 경우에는 $z$를 $\bar z\mathrm{E}_{\mathfrak h}$로, 따라서 $a$를 $\bar a\mathrm{E}_{\mathfrak h}$로 바꾸어야 한다. 그러면
\[
\mathrm{E}_{\mathfrak g}l=\mathrm{E}_{\mathfrak g}\bar a\mathrm{E}_{\mathfrak h}=\mathrm{E}_{\mathfrak g}a
\]
가 되며, 여기서 $a$는 $k$의 원소이다. 실제로 (3), (4), (5)에 의해
\[
k\mathrm{E}_{\mathfrak g}k=k\mathrm{E}_{\mathfrak g}k\mathrm{E}_{\mathfrak h}=\sum u_s\{k^s,k\}\mathrm{E}_{\mathfrak h}
\]
이다.},
\[
E_\Gg l_\mu=E_\Gg\sum_i z_{\mu,i}=E_\Gg a_\mu
\]
를 얻는다. 따라서 $\mathfrak L=\widetilde{\mathfrak o}\mathrm{E}_{\mathfrak g}\mathfrak a$이다. $\mathfrak L$이 정칙(최대 랭크)인 경우에는 $a$의 랭크가 $n$, 곧 최대 랭크여야 하므로 $\widetilde aE_\Gg o$가 이에 역인 오른쪽 아이디얼이 된다. 이들을 합성하면 $\widetilde{\mathfrak a}\mathrm{E}_{\mathfrak g}\mathfrak b$가 $\mathfrak O_{\mathfrak a}$의 모든 왼쪽 아이디얼을 지나고, $\widetilde bE_\Gg a$가 이에 역인 모든 오른쪽 아이디얼을 지남을 알 수 있다.

%
%
\medskip
\noindent\emph{정리 4의 증명.} 교집합 $t=[k,\mathfrak O_a]$는 오더가 된다. 실제로 $\mathfrak O_{\mathfrak a}$는 정수적인 원소들로만 이루어져 있다. $t$는 $\mathfrak O_{\mathfrak a}$의 부분집합으로서 오른쪽 승수환이 되며, 이 조건을 만족하는 $k$ 안의 가장 큰 오더이다. 따라서 적어도 $a$의 오더 $o_a$를 포함한다. $t=o_a$임을 보이기 위해서는 다시 각 성분이 같음을 보이면 충분하다. 한 원소 $t$가 $t_p$에 속한다고 하자. 그러면 모든 $c_{ik}=\widetilde a_i\mathrm{E}_{\mathfrak g} a_k$에 대하여 $c_{ik}t$는 정수적 계수(곧 $o_p$의 계수)를 갖는 $c_{ik}$들의 선형결합이다. 다른 한편 $a_kt$는 $\Omega_p$의 계수를 갖는 $a_k$들의 선형결합이다. $c_{ik}$에 의한 표현이 유일하므로 이 계수들은 $o_p$에 들어가야 한다. 곧 $t$는 $\mathfrak a_p$의 오더에 속한다. 따라서 $t$는 $a$의 오더와 같다. $k$의 모든 오더가 교집합으로 나타날 수 있다는 것은 $k$의 가군 $a$가 특히 모든 오더들을 지나기 때문에 따른다. 이제 변환에 관한 주장은 정리 3에서 곧바로 나온다.

\medskip
\noindent\emph{정리 4a의 증명.} 이 정리는 정리 4를 주구역에 특수화한 것일 뿐이다. 먼저 출발 극대 오더 $\mathfrak O=\widetilde oE_\Gg o$의 경우를 보자. $o$가 $\mathfrak L=\widetilde{\mathfrak o}\mathrm{E}_{\mathfrak g}\mathfrak a$의 오른쪽 오더에도 포함되어 있다면 $\widetilde aE_\Gg a$도 주구역에 속한다. 따라서 $a$는 $k$의 아이디얼이고 $\mathfrak L=\mathfrak Oa$이다. $\mathfrak O_{\mathfrak a}$에서 출발하면 마찬가지로 $\mathfrak L=\widetilde aE_\Gg b$가 되고, $b$, 따라서 $a^{-1}b$도 $k$의 아이디얼이다. 그러므로 $\mathfrak L=\widetilde aE_\Gg aa^{-1}b=\mathfrak O_a a^{-1}b$를 얻는다.

%
%
\begin{center}
\textbf{§ III. --- \emph{임의의 교차곱의 극대 오더.}}
\end{center}

분할 교차곱에 대하여 이끌어 낸 정리들은 구역별 구분을 세분하기만 하면 일반적인 경우에도 성립한다. 주구역에 관한 사실들은 그대로 성립하고, 나머지는 각 소에 관한 명제로 정식화해야 한다. 여기서 소는 이제 $p$-진 확대라는 뜻으로 이해한다.

\emph{정의.} $K$를 $\Omega$ 위의 단순 정규대수(중심단순대수)라 하고 $k$를 극대 가환 부분체라 하자. $K$의 극대 오더들 가운데 $k$와의 교집합이 같고, 또한 $K$의 분기소에서 $p$-진 성분들이 같은 것들을 $k$에 관한 하나의 구역으로 묶는다. $k$와의 교집합이 정수환이면 이를 주구역이라 한다.

주구역에 대해서는 다음이 성립한다.

\medskip
\noindent\textsc{정리 1. ---} \emph{한 주구역의 극대 오더들은 $k$의 아이디얼들을 확장한 아이디얼에 의한 변환으로 서로 옮겨진다. 다시 말해 한 주구역의 극대 오더들의 왼쪽 아이디얼 가운데 디퍼런트(different)와 서로소이고 그 오른쪽 오더가 $o$를 포함하는 것들은 $k$의 아이디얼을 확장한 것들이다.}

\medskip
\noindent\textsc{정리 2. ---} \emph{$K$의 차수가 소수이면 주구역은 하나뿐이며, 이 주구역은 $k$-아이디얼의 확장아이디얼에 의해 자기 자신으로 옮겨진다. 하나의 고정된 극대 오더 $\mathfrak O$에 대응하는 아이디얼들은 $k$의 아이디얼을 $\mathfrak O$의 양쪽 아이디얼들과 함께 확장하여 얻는 것들이다.}

\emph{증명.} 한 주구역의 두 극대 오더는 유한 개의 소에서만 서로 다르며, 정의에 따라 이들은 $K$가 분할되는 소들이다. 이런 소에서는 $K$가 § 1과 § 2에서 다룬 꼴이 된다. $k$가 갈루아일 때에는 인자계가 $1$과 동치이므로 $1$인 인자계로 바꿀 수 있어 이것이 명백하다. 비갈루아의 경우에도 § 1의 2와 3에 의해 같은 결론이 따른다. 따라서 각각의 그런 소에서, 그러므로 일반적으로도\srcfn{(11)}{이에 대해서는 H. § 8 또는 앞에 실린 Hasse의 논문 (4)를 참조하라.} 확장아이디얼이 쓰인다.

$K$의 차수가 소수이면 $K$는 분할대수이거나(§ 2) 나눗셈대수이다. 후자의 경우 모든 분기소에서도 나눗셈대수로 남으므로 주구역은 하나뿐이고, 따라서 다시 확장아이디얼에 의한 변환이 성립한다. 가장 일반적인 변환은 확장아이디얼에, 극대 오더를 자기 자신으로 변환하는 아이디얼들을 합성하여 얻는다. 곧 양쪽 아이디얼들을 합성하는데, 이들은 잘 알려진 대로 분기소에서만 중심의 아이디얼, 따라서 확장아이디얼과 다르다. 이로써 나머지 주장들도 따른다.

또한 다음이 성립한다.

\medskip
\noindent\textsc{정리 3. ---} \emph{임의의 한 구역의 극대 오더들은 디퍼런트와 서로소인 아이디얼에 의해 서로 옮겨진다. 이 아이디얼들은 각 소에서 § 2의 정리 3에 제시된 바와 같이 해당 오더의 가군들로 구성된다. 하나의 고정된 아이디얼에 대하여 이 가군들의 성분은 유한 개의 소에서만 그 오더와 다르다.}

이는 § 2에서 곧바로 따른다. 다만 각 소에서 연산자 $u_s$를 그에 대응하는 동치인 연산자로 바꾸어야 함을 덧붙여 둔다.

$k$의 각 오더마다 그 오더를 교집합으로 갖는 극대 오더가 존재하는지는 분기소에서 더 정밀하게 조사해야 한다. 주구역은 언제나 존재한다. 이는 교차곱 표현, 경우에 따라서는 일반화된 교차곱 표현으로부터 알 수 있다. 인자계를 정수적인 것으로 잡으면 이 표현은 $o$를 포함하는 한 오더를 주기 때문이다\srcfn{(12)}{다른 증명은 Hasse와 Chevalley의 (4)와 (3)에 있다.}.

%
%
또 하나의 문제는 구역에 따른 산술적 구분을 통하여 대수적 분해체들을 어느 정도까지 특성화할 수 있는가 하는 것이다. 다시 말해 서로 다른 분해체들은 언제나 서로 다른 구역 구분을 낳는가, 또는 이미 서로 다른 주구역을 낳는가? 그리고 모든 분해체의 주구역들만으로 이미 모든 극대 오더를 빠짐없이 얻을 수 있는가?

하나의 극대 오더만으로는 그러한 특성화가 분명히 불가능하다는 것을 가장 단순한 예들이 보여 준다. 예를 들어 사원수체의 극대 오더
\[
\left[1,i,j,\frac{1+i+j+k}{2}\right]
\]
는 정수환 $[1,i]$뿐 아니라 3차 단위근체의 정수환
\[
\left[1,\frac{1+i+j+k}{2}\right]
\]
도 포함한다.

\begin{flushleft}
괴팅겐, 1932년 8월.
\end{flushleft}

\clearpage
\setcounter{footnote}{0}
\phantomsection
\addcontentsline{toc}{section}{43. 아이디얼 미분과 디퍼런트}
%
%
%
%
%
%

\section*{43. 아이디얼 미분과 디퍼런트}

\begin{center}
\emph{Journal f. d. reine u. angew. Math. 188 (1950), 1--21쪽}\\[0.75ex]
\emph{에미 뇌터} \(\dagger\)\srcfnmark{1)}.
\end{center}

\srcfntext{1)}{E. Noether가 남긴 이 논문은 1927/28년 겨울에 집필되었으며, H. Grell이 친절하게 편집자에게 제공하였다. 이 논문은 1929년 프라하에서 열린 D. M. V. 학술대회에서 발표한 같은 제목의 강연을 상세히 서술한 것이다. Jahresbericht D. M. V. 39 (1930), 비스듬한 쪽매김 17쪽을 보라. § 6, 3부터는 원고를 더 정확히 다듬으려 했던 것으로 보인다. 그러나 전문가는 사소한 마무리만 하고 그대로 둔 이 개략적 서술을 어렵지 않게 이해할 것이다.}

\subsection*{서론}

대수적 수체의 분기이론의 주요 정리는 잘 알려진 대로 다음과 같다. 어떤 소아이디얼이 그 아래의 유리소수 \(p\)에서 \(\varrho\)제곱으로 나타나면, 디퍼런트(분기아이디얼)는 적어도 그 소아이디얼의 \((\varrho-1)\)제곱으로 나누어진다. 더욱이 \(\varrho\)가 \(p\)로 나누어지지 않으면 정확히 \((\varrho-1)\)제곱으로, \(\varrho\)가 \(p\)로 나누어지면 더 높은 거듭제곱으로 나누어진다.

이 정리는 다항식 \(f(x)\)의 미분몫 \(f'(x)\)에 관한 다음 사실의 유사물이다. \(f(x)\)에서 어떤 일차인자가 \(\varrho\)제곱으로 나타나면 \(f'(x)\)는 적어도 그 인자의 \((\varrho-1)\)제곱으로 나누어진다. 더욱이 \(\varrho\)가 계수영역의 표수로 나누어지지 않으면 정확히 \((\varrho-1)\)제곱으로, 그 표수로 나누어지면 더 높은 거듭제곱으로 나누어진다.

이하에서는 이것이 형식적인 유사 이상임을 보인다. \emph{디퍼런트는 수체 \(\mathsf{K}\)의 정의아이디얼을, \(x_1,\ldots,x_n\)의 어떤 대응하는 정수 계수 다항식영역에서 미분한 몫으로 볼 수 있다. 이 미분몫은 \(x=\omega\), 곧 \(x_1=\omega_1,\ldots,x_n=\omega_n\)에서 취하며, 여기서 \(\omega_1,\ldots,\omega_n\)은 \(\mathsf{K}\)의 정수환 \(\mathfrak o\)의 가군기저이다.} 이때 정의아이디얼 \(\mathfrak M\)은 \(\omega\)들 사이의 관계 전체에 대응한다. 즉 \(f(\omega)\)가 사라지는 모든 정수 계수 다항식 \(f(x)\)로 이루어진다. 한편 미분몫 \(\mathfrak M'[x\to\omega]\)은 \(x=\omega\)에서 취한 차분몫으로 정의된다.

실제로 \(f(\omega)=0\)이므로 \(\mathfrak M\)을 모든 차 \(f(x)-f(\omega)\)로 이루어진 아이디얼로 본다. 이어 차분아이디얼 \(\mathfrak P=(x_1-\omega_1,\ldots,x_n-\omega_n)\)과 차분몫 \(\mathfrak A=\mathfrak M:\mathfrak P\)를 만든다. 여기서 몫은 통상적인 아이디얼 몫이며, \(\mathsf{K}\)의 정수, 곧 \(\mathfrak o\)의 원소를 계수로 갖는 \(x\)의 다항식영역에서 취한다. 그러면 디퍼런트 \(\mathfrak D\)는
\[
\mathfrak D=\mathfrak A[x\to\omega]
\]
로 정의된다. 차분아이디얼과 차분몫 자체를 정의하려면 계수영역을 확대해야 한다. 따라서 이것은 한 부정원 \(x\)의 다항식의 미분몫
\[
f'(\xi)=(f(x)-f(\xi)):(x-\xi)
\]
을 \(x=\xi\)에서 취하는 정의의 직접적인 일반화이다. 여기서도 차들을 정의하려면 \(\xi\)를 첨가하여 다항식영역을 확대해야 한다. 실제로 \(\omega\)의 기저가 하나의 원소의 거듭제곱들로 이루어지는 순간 아이디얼 미분몫도 다항식 미분몫으로 환원된다.

이처럼 여기서 주는 디퍼런트의 정의는 알려진 것에 직접 이어지지만, 아이디얼 \(\mathfrak M\)이 선택한 기저에 의존하므로 불변이 아닌 중간 대상들을 도입한다. 정수 계수 다항식영역 대신 정수환 \(\mathfrak o\)와 동형인 환 \(\mathfrak O\) --- 이는 \(\mathfrak M\)에 따른 잉여류환과 동형이 된다 --- 을 도입하고 그 계수영역을 \(\mathfrak o\)로 확대하여 \(\mathfrak O_{\mathfrak o}\)를 만들면, 이와 동등한 \emph{완전히 불변인 정의}를 얻는다.\srcfn{2)}{계수영역의 확대, 곧 직접곱 구성은 § 1에서 이 논문의 목적에 필요한 것보다 더 일반적으로 비가환환에 대해서도 다룬다. 여기에는 가환환과 유한기저로 특수화한 § 2만으로 충분하다.} 차분아이디얼 \(\mathfrak P\)는 여기서 \(\mathfrak O\)와 \(\mathfrak o\) 사이의 동형에 따라 서로 대응하는 \(x\)와 \(\xi\)의 모든 차 \(x-\xi\)에서 유도된다. 차분몫 \(\mathfrak A\)는 영아이디얼을 \(\mathfrak P\)로 나눈 몫과 같다. \(\mathfrak A\)에서 각 \(x\)를 그와 동형인 \(\xi\)로 바꾸면 디퍼런트를 얻는다.

이하에서는 이 불변적 정의를 맨 앞에 둔다. 곧 계수확대의 가능성에 관한 몇 가지 조건만 만족한다는 가정 아래, 임의의 가환환이 한 부분환에 관하여 갖는 디퍼런트를 직접 정의한다. 이 조건들은 예컨대 가군기저가 존재하면, 필요할 경우 어떤 몫환들로 넘어가서 언제나 만족한다. 이에 따라 특히 수체의 임의의 오더의 디퍼런트와, 예컨대 한 소수 거듭제곱에 따른 오더의 잉여류환의 디퍼런트도 정의된다. 상대디퍼런트와 여러 부정원의 대수함수인 경우의 디퍼런트도 마찬가지이다. 디퍼런트의 미분적 정의가 통상적인 정의와 일치한다는 것은 직접합분해에 바탕을 둔 정확한 구조 연구에서 나온다. 이 구조 연구는 직접곱 구성으로 수체에 대응시킨 갈루아 확대환을 다룬다\srcfn{3)}{이는 나의 판별식 논문의 바탕이 된 고찰을 더 발전시킨 것이지만, 그 논문을 알 필요는 없다. E. Noether, \emph{Der Diskriminantensatz für die Ordnungen eines algebraischen Zahl- und Funktionenkörpers}, J. Math. 157 (1927), 82--104쪽을 보라.}. 한 부정원의 다항식에 따른 잉여류환이라는 특수한 경우에는 이 연구가 라그랑주 보간공식의 분석으로 귀착된다. 더 정확히 말하면 다음과 같다.

수체 \(\mathsf{K}\)와 동형이고 \(\mathfrak O\)를 포함하는 환 \(\mathfrak K\)를 도입하고, 그 계수영역인 유리수체를 \(K\)의 갈루아체 \(F\)로 확대한다. 이 확대계 \(\mathfrak K_\Gamma\)는 \(K\)의 \(n\)개 동형사상(켤레체들로의 이행)에 대응하는 \(n\)개의 단순성분의 직접합이 된다. 이 성분들은 차분몫 \(\mathfrak A\)와 그 켤레들의 확대이며, 동시에 \(\mathfrak K_\Gamma\)의 영아이디얼은 차분아이디얼과 그 켤레들에서 유도되는 \(n\)개 아이디얼의 교집합이 된다. 단순성분들은 \(\Gamma e^{(i)}\) 꼴이고, 여기서 \(e^{(i)}\)는 단위원의 성분이다. \(x=\xi\)일 때 \(e^{(i)}\)가 단위원으로 가므로 차분몫은 \(\mathfrak d e^{(1)}\)와 같아진다. 여기서 \(\mathfrak d\)는 \(\mathfrak o\)의 디퍼런트이다. 디퍼런트 \(\mathfrak d\)는 체의 원소들 가운데 \(e^{(1)}\)와 곱했을 때 \(\mathfrak O_{\mathfrak o}\), 곧 오더와 자기 자신의 직접곱에 속하는 모든 원소로 이루어진다. 켤레들에도 이에 대응하는 것이 성립한다.

여기서 여모듈에 의한 정의를 얻으려면 \(e^{(i)}\)가 \(\mathfrak O\), 또는 \(\mathfrak K\)의 기저 \(x_1,\ldots,x_n\) --- 이 기저는 \(\omega\)들에 대응한다 --- 에 관한 선형형식이 된다는 점에 주목해야 한다. 그런데 \(e^{(i)}\)에서 \(x\)들의 계수들은 바로 \(\mathfrak o\)와 그 켤레들의 여모듈의 기저를 이룬다. 차분몫 \(\mathfrak A\)가 \(\mathfrak o\)의 계수를 갖는다는 사실에서 디퍼런트가 \(\mathfrak o\)를 그 여모듈로 나눈 몫과 같다는 것이 곧바로 따른다. 이로써 데데킨트의 정의를 얻는다. 이 고찰들은 정수환에 국한되지 않으며, 모든 오더의 디퍼런트를 그 오더를 여모듈로 나눈 몫으로 특징짓는다.

정수환의 경우, 미분적 정의가 기본방정식에 의한 정의와 일치한다는 것은 위에서 말한 사실에서 나온다. 즉 기저가 하나의 원소, 여기서는 기본형의 거듭제곱들로 이루어지면 아이디얼 미분몫이 다항식 미분몫으로 환원된다. 따라서 이 마지막 정의와 데데킨트의 정의 사이의 개념적 관계도 주어진다.

정수환의 경우에 기본방정식을 바탕으로 삼으면 처음에 말한 주요 정리가 바로 따라온다는 것은 잘 알려져 있다. 지수를 정확히 결정하려면 충분히 큰 \(t\)에 대하여 \(p^t\)에 따른 정수환의 잉여류환의 디퍼런트 \(\mathfrak d[p^t]\)를 고찰한다. \(t=n\)이면 충분하며, 이차체에서는 이것이 필요하기도 하다. 이때 \(\mathfrak d\)를 \(p^t\)로 보아 얻는 것이 정확히 \(\mathfrak d[p^t]\)라는 사실이 본질적이다. 여기서 환의 직접합에 디퍼런트들의 직접합이 대응하므로, 다시 \(\mathfrak p^{t\varrho}\)에 따른 잉여류환으로 국한할 수 있다. 그런데 이 마지막 환은 \(p^t\)를 법으로 한 계수를 갖는 한 부정원 다항식에 따른 잉여류환과 동형이다. 그리고 이 다항식을 미분하면 --- 다른 길로 그런 다항식에 이른 Ö. Ore가 보였듯이\srcfn{4)}{출판 시 주: 이 인용은 원고에서 채워지지 않았다. 개략만 제시된 이 논문의 § 7에 있는 문헌 안내를 보라.} --- 보충수들을 이용하여 가능한 모든 지수를 정확히 파악할 수 있다.

끝으로 구조 연구에서는 잘 알려진 정리, 곧 \(\mathsf{K}\)의 상위체의 디퍼런트는 그 상대디퍼런트와 \(\mathsf{K}\)의 디퍼런트의 곱과 같다는 정리도 나온다. 이는 직접합분해의 단위원들에 같은 곱셈 성질이 성립하고 따라서 여모듈들에도 성립한다는 사실의 귀결이다.

위에서 말한 사실들은 어떤 몫환들로 넘어감으로써 모두 상대디퍼런트에도 성립한다. 이는 내가 판별식 논문에서 상대판별식에 대하여 전개한 것과 정확히 같다. 함수영역으로 넘어가면 표수 0인 계수체를 갖는 대수함수체에도 모든 것이 그대로 유지된다. 반면 표수 \(p\)와 정수적 대수함수의 경우에는 구조적 성질들은 유지되지만(본질적으로는 함수영역으로 넘어가지 않아도 된다), 비분리 확대체가 나타나므로 분기이론은 달라진다. 여기서는 이를 더 전개하지 않는다. 가장 중요한 결과는 위에서 개략한, 여러 부정원의 대수함수에도 성립하는 완전히 불변인 구조정리들(§ 5와 § 6)이라고 하겠다. 여모듈로 넘어가는 것은 이미 계수항등식들로 풀어 쓰는 일이며, 알려진 것을 이 틀에 넣기 위해서만 수행하였다.

\subsection*{§ 1. 환의 계수영역 확대 또는 직접곱 구성. 정의아이디얼\srcfnmark{5)}}
\srcfntext{5)}{이 절에 관한 비판적인 지적에 대하여 B. L. van der Waerden에게 감사한다.}

이 절에서는 곱셈의 교환법칙을 가정하지 않는 일반적인 환을 바탕으로 삼는다. 다만 간단하게 하기 위하여 단위원의 존재를 가정하며, 반대말을 명시하지 않는 한 논문 전체에서 그렇게 한다.

\paragraph{1. 직접곱의 정의.}
환 \(\mathfrak O\times\mathfrak o\), 또는 \(\mathfrak O_{\mathfrak o}\)가 다음 조건들을 만족할 때 이를 \(\mathfrak h\)에 관한 환 \(\mathfrak O\)와 \(\mathfrak o\)의 \emph{직접곱}, 또는 \emph{\(\mathfrak O\)의 계수영역 \(\mathfrak h\)를 \(\mathfrak o\)로 확대하여 얻은 환}\srcfn{6)}{이하에서 주로 쓸 두 번째 표기는 계수영역의 확대를 나타내기 위한 것이다. 첫 번째 것은 직접곱의 통상적인 명칭이다.}이라 한다.

\begin{enumerate}
\item[(1)] \(\mathfrak O_{\mathfrak o}\)는 \(\mathfrak O\)와 \(\mathfrak o\)를 부분환으로 포함하고 이 환들의 곱과 같다. 곧 \(\mathfrak O_{\mathfrak o}\) 안에서 \(\mathfrak O\)와 \(\mathfrak o\)를 포함하는 모든 환의 교집합과 같다.

\item[(2)] \(\mathfrak O\)와 \(\mathfrak o\)의 교집합은 \(\mathfrak h\)이고, \(\mathfrak h\)는 \(\mathfrak O_{\mathfrak o}\)의 단위원을 포함한다.

\item[(3)] \(\mathfrak O\)와 \(\mathfrak o\)의 원소들은 서로 교환한다. 따라서 \(\mathfrak O_{\mathfrak o}\)는 모든 쌍선형 결합
\[
x_{i_1}y_{i_1}+\cdots+x_{i_r}y_{i_r}
\]
전체로 이루어진다. 여기서 \(x\)들과 \(y\)들은 각각 \(\mathfrak O\)와 \(\mathfrak o\)에서 취한 임의의 유한한 원소계들을 돈다. 또한 \(\mathfrak h\)는 \(\mathfrak O\)와 \(\mathfrak o\)의 중심, 따라서 \(\mathfrak O_{\mathfrak o}\)의 중심에 놓인다.

\item[(4)] \(\mathfrak O_{\mathfrak o}\)의 두 원소가 같은 것은, \(\mathfrak O\)에서 한편으로, \(\mathfrak o\)에서 다른 한편으로 성립하는 관계들에 의해 적어도 한 방식으로 형식상 같아질 때 그리고 오직 그때이다. 이 말은 다음 뜻이다. \(\mathfrak O_{\mathfrak o}\)에서
\[
x_1y_1+\cdots+x_ny_n=0
\quad\hbox{이라 하자.}
\]
그러면 \(\mathfrak h\)의 \(h\), \(\mathfrak O\)의 \(\gamma\), \(\mathfrak o\)의 \(\delta\)에 대하여 형식상 사라지는 결합 \((\gamma h)\delta-\gamma(h\delta)\)들의 합을 더함으로써 이 식을
\[
(z_1\alpha_1+\cdots+z_r\alpha_r)+(t_1\beta_1+\cdots+t_s\beta_s)
\]
꼴로 바꿀 수 있다. 여기서 \(z_1=0,\ldots,z_r=0\), \(\beta_1=0,\ldots,\beta_s=0\)이고, \(\alpha,\beta\)는 \(\mathfrak o\)에, \(z,t\)는 \(\mathfrak O\)에 속한다. 따라서 \(z=0\)은 \(\mathfrak O\) 안의 \(x,\gamma,h\) 사이의 관계를, \(\beta=0\)은 \(\mathfrak o\) 안의 \(y,\delta,h\) 사이의 관계를 뜻한다.\srcfn{7)}{단위원의 존재를 가정하지 않으면 (3)과 (4)에 선형인 항들을 더 넣어야 한다.}
\end{enumerate}

\(\mathfrak O_{\mathfrak o}\)의 두 원소가 같은 것은 그 차가 위 방식으로 사라질 때 그리고 오직 그때이다.

따라서 직접곱이라는 요구에 의하여 등식관계와 연산관계(덧셈과 곱셈)는 인자들의 그것으로 유일하게 결정된다. 여기서 다음이 따른다. \(\overline{\mathfrak O},\overline{\mathfrak o}\)와 그 교집합 \(\overline{\mathfrak h}\)가 각각 \(\mathfrak O,\mathfrak o\), \(\mathfrak h\)와 동형이고 직접곱 \(\mathfrak O\times\mathfrak o\)가 존재하면, \(\overline{\mathfrak O}\times\overline{\mathfrak o}\)도 존재하고 \(\mathfrak O\times\mathfrak o\)와 동형이 된다.

\(\overline{\mathfrak O}\)와 \(\overline{\mathfrak o}\)의 곱과 같고 이 두 환의 원소들이 서로 교환하는 모든 환 \(\mathfrak S\)는 \(\mathfrak O\times\mathfrak o\)의 준동형상이다. 이때 \(\overline{\mathfrak O}\)에서 \(\mathfrak O\)로, 또는 \(\overline{\mathfrak o}\)에서 \(\mathfrak o\)로 가는 사상은 동형으로 남는다. 실제로 \(\mathfrak O\times\mathfrak o\)에서의 등식은 \(\mathfrak S\)에서의 등식에 포함되고, \(\mathfrak O\) 또는 \(\mathfrak o\)에서는 각각 \(\mathfrak O\times\mathfrak o\) 안의 등식과 일치한다.\srcfn{8)}{§ 1, 4도 비교하라.}

\paragraph{2. 부분환에 관한 직접곱이 존재하기 위한 필요충분조건.}
\(\mathfrak O\)와 \(\mathfrak o\)가 환이고 그 교집합이 양쪽 중심에 놓인 환 \(\mathfrak h\)라고 하자. 그러면 \(\mathfrak h\)에 관한 직접곱이 반드시 존재하는 것은 아니다. 다음 예가 이를 보인다. \(\mathfrak h\)를 유리정수환, \(\mathfrak O=\mathfrak h[x]\)이고 \(1-2x=0\), \(\mathfrak o=\mathfrak h[y]\)이고 \(1-2y=0\)이라 하자. 여기서 \(y\)는 새 기호이다. 따라서 \(\mathfrak O\)와 \(\mathfrak o\)는 \(\mathfrak h\)에 관한 동치인 환들이고, 그 집합론적 교집합은 \(\mathfrak h\)이지만 기호 \(x\)와 \(y\) 사이에는 아무 연산관계도 주어져 있지 않다. \(\mathfrak O\)와 \(\mathfrak o\)를 포함하고 그 교집합이 \(\mathfrak h\)인 환 \(R\)은 존재하지 않는다. 실제로 그런 \(R\)에서는 \(x\)와 \(y\) 사이에도 연산이 생기므로
\[
0=(1-2x)y-x(1-2y)=y-x
\]
가 따른다. 여기서는 \(x\)와 \(y\)의 교환가능성조차 쓰지 않았다.\srcfn{9)}{한 고정된 환 안에 놓인 가환환의 몫환 --- 여기서는 \(\frac12\)의 첨가 --- 은 바로 유일하다. 이는 서로 동치이지만 서로 다른 두 갈루아체 확대를 하나의 공통 포괄체에 매장할 수 없다는 사실의 유사물이다. 다만 여기서는 직접곱 구성을 통하여 영인자를 갖는 환에 매장할 수 있다.}

\emph{직접곱이 존재하기 위한 필요충분조건}은 다음과 같다.

\emph{\(\mathfrak O\)와 \(\mathfrak o\)를 포함하는 환 \(R\)이 적어도 하나 존재하여, \(\mathfrak O\)와 \(\mathfrak o\)의 교집합이 \(\mathfrak h\)이고 두 환의 원소들이 서로 교환해야 한다.}

직접곱이 존재하면 그 자체가 그러한 환 \(R\)이므로 이 조건은 필요하다. 역으로 직접곱을 구성하기 위해서는 \(\mathfrak O\)와 \(\mathfrak o\)의 \(\mathfrak h\)에 속하지 않는 원소들 사이에 아무 연산관계도 없다고 가정해도 된다. 이는 \(\mathfrak h\)의 동치인 확대, 곧 다른 기호들을 도입한 확대들로 넘어가면 언제나 이룰 수 있고, 이하에서는 늘 그렇게 가정한다.\srcfn{9a)}{§ 2, 1과 각주 12)를 보라.} 또 \(\mathfrak h\)는 \(\mathfrak O\)와 \(\mathfrak o\)의 단위원 \(e\)를 포함한다고 하자(다만 주 7도 보라). 이제 \(\mathfrak T\)를 모든 쌍선형 결합
\[
x_{i_1}y_{i_1}+\cdots+x_{i_r}y_{i_r}
 =y_{i_1}x_{i_1}+\cdots+y_{i_r}x_{i_r}
\]
의 집합으로 정의한다. 여기서 \(x\)와 \(y\)들은 각각 \(\mathfrak O\)와 \(\mathfrak o\)의 임의의 유한한 원소계들을 돌며, 따라서 \(\mathfrak T\)에서 \(x\)와 \(y\)는 서로 교환한다. 그러한 두 결합은 1, (4)에서 정확히 말한 의미로 \(\mathfrak O\)와 \(\mathfrak o\) 안의 관계들에 의해 형식상 같아질 때 그리고 오직 그때 같다고 정한다. 또한 \(\alpha,\beta\) 가운데 0이 나타날 수도 있을 때
\[
\alpha_{i_1}x_{i_1}+\cdots+\alpha_{i_r}x_{i_r}
\quad\hbox{와}\quad
\beta_{i_1}x_{i_1}+\cdots+\beta_{i_r}x_{i_r}
\]
의 합과 곱을 각각
\[
(\alpha_{i_1}+\beta_{i_1})x_{i_1}+\cdots+(\alpha_{i_r}+\beta_{i_r})x_{i_r}
\quad\hbox{와}\quad
\sum\alpha_{i_\lambda}\beta_{i_\mu}x_{i_\lambda}x_{i_\mu}
\]
로 정의한다. 이 정의는 위의 등식이라는 뜻에서 유일하다. 실제로 모든 영관계(1, (4)의 의미)는 \(\mathfrak o\) 또는 \(\mathfrak O\)의 원소를 오른쪽이나 왼쪽에서 곱해도 다시 그러한 영관계가 되고, 영관계들의 합과 차도 마찬가지이다. 계 \(\mathfrak T\)는 \(\mathfrak O\)와 \(\mathfrak o\)를 포함하는 환을 이루며, 등식·덧셈·곱셈에 관하여 직접곱의 조건들을 만족한다. 또한 일반적인 매장 가정의 귀결로 \(\mathfrak T\)는 교집합 조건도 만족한다. 실제로 \(\mathfrak S\)를 \(R\) 안에서의 \(\mathfrak O\)와 \(\mathfrak o\)의 곱이라 하자. 그러면 1에 따라 \(\mathfrak S\)는 \(\mathfrak T\)의 준동형상이고, \(\mathfrak O\)와 \(\mathfrak o\) 각각에 관해서는 동형이다. \(\mathfrak O\)의 \(x\)와 \(\mathfrak o\)의 \(y\)에 대하여 \(\mathfrak T\)에서 \(x=y\)이면 \(\mathfrak S\)에서도 같은 관계가 따른다. 따라서 가정상 \(\mathfrak S\)에서 두 원소 모두 \(\mathfrak h\)에 속하며, \(\mathfrak O\)와 \(\mathfrak o\) 각각에 관한 동형성 때문에 \(\mathfrak T\)에서도 마찬가지이다. 그러므로 매장 가정은 필요충분하다.

\paragraph{3. 부분환에 관한 환의 정의아이디얼.}
\(\mathfrak O\)의 원소들 \(\ldots,z,\ldots\)로 이루어진 계 \(S\)가 \(\mathfrak h\)와 \(S\)에 의해 \(\mathfrak O\) 안에서 생성되는 환이 \(\mathfrak O\)와 같을 때 --- 곧 \(\mathfrak O\) 안에서 \(\mathfrak h\)와 \(S\)를 포함하는 모든 환의 교집합과 같을 때 --- \(S\)를 \(\mathfrak h\)에 관한 \(\mathfrak O\)의 \emph{생성계}라 한다. 기호로는 \(\mathfrak O=\mathfrak h[S]\)이다. 다시 \(\mathfrak h\)가 \(\mathfrak O\)의 중심의 부분환이라고 가정한다. 그러면 \(\mathfrak O\)는 비가환 다항식영역 \(\mathfrak h[..Z..]\)의 준동형상이다. 여기서 \(Z\)들은 부정원이고 그 계는 \(S\)와 같은 기수를 갖는다. 비가환 다항식영역이란 모든 유한합 \(\sum h_iP_i=\sum P_i h_i\)의 계를 뜻한다. 여기서 \(P_i\)는 지수 0에 대응하는 단위원까지 포함한 거듭제곱곱
\[
P_i=Z_{i_1}^{\varrho_1}Z_{i_2}^{\varrho_2}\cdots Z_{i_n}^{\varrho_n}
\]
이고(물론 같은 첨자가 되풀이될 수 있다), 등식은 서로 다른 거듭제곱곱들에 관한 계수들의 등식으로 정의된다. 이 준동형에서 \(\mathfrak O\)는 잉여류환 \(\mathfrak h[..Z..]/\mathfrak M\)과 동형이 된다. 여기서 양쪽 아이디얼 \(\mathfrak M\)은 \(Z\)와 \(z\)가 준동형에 따라 대응할 때 \(\mathfrak O\)에서 \(F(z)\)가 사라지는 모든 다항식 \(F(Z)\), 곧 \(\mathfrak h[..Z..]\)의 원소로 이루어진다. \emph{아이디얼 \(\mathfrak M\)을 \(\mathfrak h\)에 관한 정의아이디얼, 더 정확히는 \(z\)들의 생성계에 대응하는 정의아이디얼이라 한다.} 이 정의에서 다음이 드러난다. 동형인 환들이 서로 동형으로 대응하는 부분환 \(\mathfrak h,\overline{\mathfrak h}\)와 생성계 \(S,\overline S\)를 가지면 그 정의아이디얼들도 서로 동형이다. 특히 \(\mathfrak h\)가 \(\overline{\mathfrak h}\)와 일치하면 정의아이디얼들도 같다고 잡을 수 있다.

\(\mathfrak h\)가 \(\mathfrak O\)의 중심에 있다고 가정했으므로, 하나의 원소 \(z\)만으로 이루어진 생성계 \(S\)가 존재하면 \(\mathfrak O\)는 반드시 가환이다. 또 \(\mathfrak M\)이 주아이디얼이면 \(\mathfrak M\)의 기저다항식 \(G(Z)\)에 대응하는 등식 \(G(z)=0\)을 \(\mathfrak h\)에 관한 \(\mathfrak O\)의 정의방정식이라 한다.

\paragraph{4. 직접곱의 정의아이디얼.}
(3)에서와 같이 \(\mathfrak O\)가 \(\mathfrak h[..Z..]/\mathfrak M\)과 동형이고 \(\mathfrak M\)이 정의아이디얼이라고 하자. \(\mathfrak h[..Z..]\)와 함께 다항식영역 \(\mathfrak o[..Z..]\), 곧 모든 유한합
\[
\sum \gamma_iP_i=\sum P_i\gamma_i
\qquad\hbox{(\(\gamma_i\)는 \(\mathfrak o\)의 원소)}
\]
의 계를 고찰하며, 등식은 계수들의 등식으로 정의한다. \(\mathfrak o[Z]\)는 \(\mathfrak h[Z]\)를 포함하고 \(\mathfrak h\)에 관한 \(\mathfrak h[..Z..]\)와 \(\mathfrak o\)의 직접곱이다. 달리 말하면 \(\mathfrak h[..Z..]\)의 계수영역을 확대하여 얻는다.\srcfn{10)}{실제로 등식조건과 연산조건이 만족되고, \(\gamma=h+h_1P_1+\cdots+h_rP_r\)이면 \(\mathfrak o[..Z..]\)에서의 등식에 따라 \(\gamma=h\)이다. § 2, 2도 보라. \(\mathfrak o[..Z..]\)가 직접곱이라는 사실은 이 항에서는 쓰지 않는다.} \(\mathfrak M_{\mathfrak o}\)를 \(\mathfrak o[..Z..]\)에 관한 \(\mathfrak M\)의 확대아이디얼, 곧 \(\mathfrak o[..Z..]\) 안에서 \(\mathfrak M\)을 포함하는 모든 양쪽 아이디얼의 교집합이라 하자. 이는 \(\mathfrak M\)의 원소들을 \(\mathfrak o\)의 계수로 선형결합한 모든 원소의 계이고, \(\mathfrak o\)가 \(Z\)들과 교환하므로 이미 양쪽 아이디얼이다. \emph{\(\mathfrak h\)에 관한 직접곱 \(\mathfrak O_{\mathfrak o}\)가 존재하면 \(\mathfrak O_{\mathfrak o}\)는 \(\mathfrak o[..Z..]/\mathfrak M_{\mathfrak o}\)와 동형이다.} 실제로 \(\mathfrak O_{\mathfrak o}\)는 \(\mathfrak o[..Z..]\)의 준동형상이므로, \(z\)들의 생성계에 대응하는 \(\mathfrak o\)에 관한 \(\mathfrak O_{\mathfrak o}\)의 정의아이디얼에 따른 잉여류환과 동형이다. 그런데 1, (4)의 등식 정의에 따라 \(\mathfrak O_{\mathfrak o}\)에서 \(z\)들 사이에 성립하는 관계들은 \(\mathfrak O\)에서 성립하는 관계들의 \(\mathfrak o\) 계수 선형결합으로 쓸 수 있다. 따라서 \(\mathfrak M_{\mathfrak o}\)가 정의아이디얼이다.

\paragraph{주.}
같은 고찰은 직접곱과 그 정의아이디얼에 대한 대칭적이지만 더 번거로운 표현도 준다. \(\mathfrak o\)가 \(\mathfrak h[..Y..]/\mathfrak N\)과 동형이라고 하자. 여기서 \(Y\)들은 \(Z\)들과 다른 기호이므로 다항식영역 \(\mathfrak h[..Y..Z..]\)가 정의된다. 그러면 \(\mathfrak O_{\mathfrak o}\)는
\[
\mathfrak h[..Y..Z..]/(\mathfrak M,\mathfrak N,..YZ-ZY..)
\]
와 동형이다. 오른쪽의 정의아이디얼은 \(\mathfrak M,\mathfrak N\)과 모든 결합 \(YZ-ZY\)에 의해 생성된다. 마지막 결합들은 \(\mathfrak o\)와 \(\mathfrak O\) 원소들의 교환가능성을 나타내고, 앞의 두 아이디얼은 \(\mathfrak O_{\mathfrak o}\)의 관계들이 \(\mathfrak O\)와 \(\mathfrak o\) 각각의 관계들의 귀결임을 나타낸다.

\subsection*{§ 2. 독립 가군기저를 갖는 환. 직접곱의 구성}

\paragraph{1. 직접곱의 구성.}
\(\mathfrak O\)의 원소들 \(\ldots,t,\ldots\)로 이루어진 계 \(T\)에서 \(\mathfrak O\) 안에 유도되는 \(\mathfrak h\)-가군이 \(\mathfrak O\)와 같을 때 \(T\)를 \(\mathfrak O\)의 \(\mathfrak h\)-가군기저라 한다. 그러면 \(\mathfrak h\)에 단위원이 존재하므로 \(\mathfrak O\)는 모든 선형결합 \(h_{i_1}t_{i_1}+\cdots+h_{i_r}t_{i_r}\)로 이루어진다. 가군기저는
\[
h_{i_1}t_{i_1}+\cdots+h_{i_r}t_{i_r}=0
\quad\hbox{이면 언제나}\quad
h_{i_1}=0,\ldots,h_{i_r}=0
\]
이 따를 때 독립이라 한다. \emph{\(\mathfrak O\)가 단위원을 포함하는 독립 가군기저를 가지면, \(\mathfrak h\)의 임의의 확대환 \(\mathfrak o\)에 대하여 직접곱 \(\mathfrak O_{\mathfrak o}\)가 존재한다. 여기서 필요한 조건은 \(\mathfrak h\)가 \(\mathfrak o\)의 중심의 부분환이고 집합론적 교집합 \([\mathfrak O,\mathfrak o]\)가 \(\mathfrak h\)라는 것뿐이다.}\srcfn{11)}{대칭적으로 \(\mathfrak o\)에 기저 조건을 부과해도 된다. § 2, 2에 따라 일반성을 잃지 않고 --- E. Noether, \emph{Der Diskriminantensatz für die Ordnungen eines algebraischen Zahl- oder Funktionenkörpers}, 같은 저널 157 (1927), 82--104쪽, §§ 1, 2를 보라 --- \(\mathfrak O\)와 \(\mathfrak o\)의 \(\mathfrak h\)에 속하지 않는 원소들 사이에는 아무 연산관계도 없다고 가정한다.} § 1, 2에 따라 대응하는 매장환 \(\mathfrak R\)을 구성하면 충분하다. 여기서는 이 환 자체가 이미 직접곱임이 드러날 것이다. \(\mathfrak R\)을 모든 선형결합
\[
\gamma_{i_1}t_{i_1}+\cdots+\gamma_{i_r}t_{i_r}
=t_{i_1}\gamma_{i_1}+\cdots+t_{i_r}\gamma_{i_r}
\]
전체로 정의한다. 여기서 \(t_{i_\nu}\)는 독립기저 \(T\)의 임의의 유한한 원소계들을 돌고 \(\gamma\)들은 이에 대응하여 \(\mathfrak o\)에서 취한다. 등식은 \(t_{i_\nu}\)들의 계수들이 같다는 뜻으로 정의하고, 합과 곱은 § 1, 2에서처럼 통상적인 계산법칙으로 정의한다. 기저의 성질 때문에 임의의 두 \(t_{i_\nu}\)의 곱은 유한개의 기저원소를 \(\mathfrak h\)의 계수로 선형결합한 것이므로 \(\mathfrak R\)은 환이다. \(\mathfrak h\)가 \(\mathfrak o\)의 부분환이므로 \(\mathfrak R\)은 \(\mathfrak O\)를 포함하고, \(T\)가 단위원을 포함하므로 \(\mathfrak o\)도 포함한다. 또 \(\gamma\)들은 \(h_{i_\nu}\) 및 \(t_{i_\nu}\)와 각각 교환하므로, \(\mathfrak R\) 안에서 \(\mathfrak O\)와 \(\mathfrak o\)의 원소들은 서로 교환한다. 끝으로 \(\mathfrak R\)에서 \(\mathfrak O\)와 \(\mathfrak o\)의 교집합은 \(\mathfrak h\)이다. 실제로 관계
\[
\gamma=\gamma e=h_1t_1+\cdots+h_rt_r
\]
에서 등식의 정의와 \(e\)가 기저원소라는 가정에 따라 하나의 \(t\)는 \(e\)와 같고 나머지 계수들은 사라지며, 따라서 \(\gamma=h\)이다. 그러므로 직접곱이 존재한다. 이는 \(\mathfrak O\)와 \(\mathfrak o\)의 곱이고 교환조건과 등식조건을 만족하는 \(\mathfrak R\)과 같다. 기저원소들이 독립이므로 이 마지막 조건은 계수들이 같다는 조건으로 환원된다. 동시에 \(\mathfrak O\)의 이 기저 \(T\)는 \(\mathfrak O_{\mathfrak o}\)의 \(\mathfrak o\)-기저가 된다.

\paragraph{2. 확대가군과 수축가군.}
\(\mathfrak B\)가 \(\mathfrak O\)의 \(\mathfrak h\)-가군이면 그 확대 \(\mathfrak B_{\mathfrak o}\)란 \(\mathfrak O_{\mathfrak o}\) 안에서 \(\mathfrak B\)를 포함하는 모든 \(\mathfrak o\)-가군의 교집합인 \(\mathfrak O_{\mathfrak o}\)의 \(\mathfrak o\)-가군을 뜻한다. \(\mathfrak C\)가 \(\mathfrak O_{\mathfrak o}\)의 \(\mathfrak o\)-가군이면 그 수축가군 \(\mathfrak C_{\mathfrak h}\)는 \(\mathfrak O\) 안에서 교집합 \([\mathfrak C,\mathfrak O]\)으로 정의된다. \(\mathfrak B\)가 동시에 \(\mathfrak O\)의 아이디얼이면 \(\mathfrak B_{\mathfrak o}\)는 \(\mathfrak O_{\mathfrak o}\)의 아이디얼이고, \(\mathfrak C\)가 \(\mathfrak O_{\mathfrak o}\)의 아이디얼이면 \([\mathfrak C,\mathfrak O]\)는 \(\mathfrak O\)의 아이디얼이다.

\paragraph{정리.}
\emph{\(\mathfrak O\)의 한 가군 \(\mathfrak B\)가 독립 \(\mathfrak h\)-가군기저를 가지며 이 기저를 \(\mathfrak O\)의 독립 \(\mathfrak h\)-가군기저로 보충할 수 있으면, \(\mathfrak B\)는 \(\mathfrak O_{\mathfrak o}\) 안에서 그 확대가군 \(\mathfrak B_{\mathfrak o}\)를 수축한 가군이다.}

\paragraph{증명.}
\(T\)가 \(\mathfrak B\)의 독립 \(\mathfrak h\)-가군기저이고, 가정에 따라 \(\mathfrak O\)의 원소 \(\overline t\)들로 이루어진 계 \(\overline T\)를 더하여 \(\mathfrak O\)의 독립 \(\mathfrak h\)-가군기저로 만들 수 있다고 하자. 그러면 \(T\)는 동시에 \(\mathfrak B_{\mathfrak o}\)의 독립 \(\mathfrak o\)-기저이고, \(T\)와 \(\overline T\)를 합한 것은 \(\mathfrak O_{\mathfrak o}\)의 그러한 기저이다.

\(\mathfrak B\)는 교집합 \([\mathfrak B_{\mathfrak o},\mathfrak O]\)에 포함된다. 역으로 \(c\)가 이 수축가군의 원소라고 하자. \(\mathfrak B_{\mathfrak o}\)의 원소로서 \(c\)는 \(T\)의 원소들에 의한 표현을 가지며, 이는 동시에 \(\mathfrak O_{\mathfrak o}\)의 \(\mathfrak o\)-기저 \(T,\overline T\)에 의한 표현이다.
\[
c=\gamma_1t_1+\cdots+\gamma_rt_r
\qquad\hbox{(\(\gamma\)들은 \(\mathfrak o\)의 원소).}
\]
\(\mathfrak O\)의 원소로서 \(c\)는 \(\mathfrak O\)의 \(\mathfrak h\)-기저 \(T,\overline T\)에 의한 표현
\[
c=h_{i_1}t_{i_1}+\cdots+h_{i_r}t_{i_r}
 +h_{j_1}\overline t_{j_1}+\cdots+h_{j_s}\overline t_{j_s}
\]
을 갖는다. \(\mathfrak O\)가 \(\mathfrak O_{\mathfrak o}\)의 부분환이고 \(\mathfrak h\)가 \(\mathfrak o\)의 부분환이므로 이 역시 동시에 \(\mathfrak O_{\mathfrak o}\)의 \(\mathfrak o\)-기저 \(T,\overline T\)에 의한 표현이다. 따라서 두 표현은 계수별로 같아야 한다. \(\overline t\)들의 계수는 사라지고
\[
c=h_{i_1}t_{i_1}+\cdots+h_{i_r}t_{i_r}
\]
이므로 \(c\)는 \(\mathfrak B\)의 원소이다. 이로써 \(\mathfrak B\)가 수축가군임을 알았다.

\paragraph{보충 주.}
이 증명은 정리에서 말한 것보다 약간 약한 가정만 사용하고 나머지 가정을 증명한다. 따라서 다음 강화된 형태가 성립한다.

\emph{\(\mathfrak O\) 안에 독립 \(\mathfrak h\)-가군기저 \(T,\overline T\)가 있고 \(T\)가 \(\mathfrak B\)에 속하며, 동시에 \(T\)가 \(\mathfrak B_{\mathfrak o}\)의 독립 \(\mathfrak o\)-가군기저이면 \(\mathfrak B\)는 그 확대 \(\mathfrak B_{\mathfrak o}\)의 수축가군이다. 그러면 \(T\)는 또한 \(\mathfrak B\)의 독립 \(\mathfrak h\)-가군기저임이 드러난다.}

\paragraph{3. 직접곱이 존재하기 위한 충분조건.}
(예컨대 상대디퍼런트의 경우처럼) 계수영역 \(\mathfrak h\)를 확대하고 나서야 독립 가군기저가 존재할 수도 있다. 이 경우에는 다음 충분조건에서 직접곱의 존재를 얻을 수 있다.

\emph{\(\mathfrak h\)의 확대환 \(\mathfrak f\)가 존재하여 \(\mathfrak h\)에 관한 직접곱 \(\mathfrak O_{\mathfrak f}\)와 \(\mathfrak o_{\mathfrak f}\), 그리고 \(\mathfrak f\)에 관한 직접곱 \(\mathfrak O_{\mathfrak f}\times\mathfrak o_{\mathfrak f}\)가 존재하면, \(\mathfrak h\)에 관한 \(\mathfrak O\times\mathfrak o\)가 존재한다.}

\(\mathfrak R=\mathfrak O_{\mathfrak f}\times\mathfrak o_{\mathfrak f}\)가 § 1, 2의 의미에서 매장환임을 보이면 된다. 실제로 \(\mathfrak O\)와 \(\mathfrak o\)는 각각 \(\mathfrak O_{\mathfrak f}\)와 \(\mathfrak o_{\mathfrak f}\)의 부분환으로서 원소별로 서로 교환한다. 또
\[
[\mathfrak O,\mathfrak o]\subseteqq[\mathfrak O_{\mathfrak f},\mathfrak o_{\mathfrak f}]=\mathfrak f
\]
이고 \([\mathfrak O,\mathfrak o]\subseteqq\mathfrak O\)도 성립한다. 따라서 \([\mathfrak O,\mathfrak o]\subseteqq[\mathfrak O,\mathfrak f]=\mathfrak h\)이고, 결국 \([\mathfrak O,\mathfrak o]=\mathfrak h\)이다. 이로써 \(\mathfrak R\)이 매장환임을 보였다.

\paragraph{주.}
이 증명은 교집합에 관한 가정을 \(\mathfrak O_{\mathfrak f}\)에 대해서만 사용하고 \(\mathfrak o_{\mathfrak f}\)에는 사용하지 않는다. 따라서 마지막 환은 단순히 \(\mathfrak o\)와 \(\mathfrak f\)의 곱이어도 된다.

\subsection*{§ 3. 부분환에 관한 환의 디퍼런트. 정의아이디얼의 미분몫}

이제부터 나타나는 모든 환은 가환이라고 가정한다.\srcfn{12)}{이 가정이 없으면 오른쪽 디퍼런트와 왼쪽 디퍼런트를 정의하고, 이 둘이 어느 범위까지 서로 일치하는지, 그리고 반단순계의 오더인 경우에 알려진 구성들과 어느 범위까지 일치하는지를 연구해야 한다(§ 4의 연구를 일반화한 것이 된다).} 또한 환 \(\mathfrak O\)와 \(\mathfrak o\)는 부분환 \(\mathfrak h\)의 동형인, 더 정확히는 동치인 확대들이라고 늘 가정한다. 곧 \(\mathfrak O\)에서 \(\mathfrak o\)로 가는 사상은 \(\mathfrak h\)의 항등사상을 연장한다. 그리고 \(\mathfrak h\)에 관한 직접곱 \(\mathfrak O_{\mathfrak o}\)가 존재한다고 가정한다.

\paragraph{1. 디퍼런트의 정의.}
\(\mathfrak h\)에 관한 직접곱 \(\mathfrak O_{\mathfrak o}\)가 존재한다고 하자. \(x\)는 \(\mathfrak O\)의 모든 원소를 돌고 \(\xi\)는 매번 동형에 따라 \(\mathfrak o\)에서 대응하는 원소라 하자. 모든 차 \(x-\xi\)가 생성하는 \(\mathfrak O_{\mathfrak o}\)의 양쪽 \emph{차분아이디얼}을 \(\mathfrak P\)라 쓰자.
\[
\mathfrak P=\{..(x-\xi)..\}\srcfnmark{13)}
\]
\srcfntext{13)}{이 표기는 언제나 원소 전체에서 유도되는 아이디얼에 사용한다.}
\emph{차분몫} \(\mathfrak A\)는 영아이디얼을 \(\mathfrak P\)로 나눈 아이디얼 몫으로 정의한다. 따라서 \(a\mathfrak P\)가 사라지는 \(\mathfrak O_{\mathfrak o}\)의 모든 원소 \(a\)로 이루어진다.
\[
\mathfrak A=(0):\mathfrak P.
\]
\(\mathfrak h\)에 관한 \(\mathfrak o\)의 \emph{디퍼런트} \(\mathfrak d\)는 \(\mathfrak A[x\to\xi]\)로 정의한다. 곧 \(\mathfrak A\)의 각 원소 \(a=\alpha_{i_1}x_{i_1}+\cdots+\alpha_{i_r}x_{i_r}\)에서 \(x\)들을 그와 동형으로 대응하는 \(\xi\)들로 바꾸어 얻는다. 이에 대응하여 \(\mathfrak h\)에 관한 \(\mathfrak O\)의 디퍼런트 \(\mathfrak D\)는 \(\mathfrak A[\xi\to x]\)로 정의한다. \(\mathfrak A\)는 \(x\)와 \(\xi\)에 관하여 대칭적으로 정의되므로, \(\mathfrak d\)와 \(\mathfrak D\)는 \(\xi\)와 \(x\)를 맞바꾸어 얻은 것으로서 \(\mathfrak o\)와 \(\mathfrak O\)의 동형에 따라 서로 대응한다. \(\mathfrak A\)가 \(\mathfrak O_{\mathfrak o}\)의 아이디얼, 따라서 특히 \(\mathfrak o\)-가군이며 \(\mathfrak O\)-가군이므로 \(\mathfrak d\)와 \(\mathfrak D\)는 각각 \(\mathfrak o\)와 \(\mathfrak O\)의 아이디얼이다.

\paragraph{2. 정의아이디얼의 미분몫.}
\(\mathfrak M\)이 (§ 1, 3) \(z\)들의 생성계에 대응하는, \(\mathfrak h\)에 관한 \(\mathfrak O\)의 정의아이디얼이라고 하자. 그러면 (§ 1, 3) \(\mathfrak M\)은 동형인 생성계 \(\xi\)들에 대응하는, \(\mathfrak h\)에 관한 \(\mathfrak o\)의 정의아이디얼이기도 하다. 또 (§ 1, 4) \(\mathfrak O_{\mathfrak o}\)는 \(\mathfrak o[..Z..]/\mathfrak M_{\mathfrak o}\)와 동형이다. 이제 다항식영역 \(\mathfrak o[..Z..]\)에서 차분아이디얼과 차분몫은
\[
\mathfrak Q=\{..(Z-\xi)..\}
\]
와
\[
\mathfrak C=\mathfrak M_{\mathfrak o}:\mathfrak Q
=\{..(F(Z)-F(\xi))..\}:\{..(Z-\xi)..\}
\]
로 정의된다. \emph{미분몫} \(\mathfrak M'\)은 \(\mathfrak C[\xi\to Z]\)로 정의한다. 여기서 이 정의가 \(\xi\)라는 생성원들에 의한 \(\mathfrak o\) 원소들의 \emph{표현}과 무관하게 유일하다는 데 주의해야 한다. 실제로 \(\mathfrak C\)는 아이디얼 몫이므로 \(\mathfrak M_{\mathfrak o}\)를 포함한다. \(\mathfrak o\)의 한 원소를 나타내는 두 표현의 차는 어떤 다항식 \(F(\xi)\)이고, 따라서 \(Z\)를 대입하면 \(\mathfrak M\), 그러므로 이미 \(\mathfrak C\)에 속하는 \(F(Z)\)를 얻는다. 아이디얼 \(\mathfrak M'\)은 \(\mathfrak M\)을 포함하는 \(\mathfrak h[..Z..]\)의 아이디얼이다. \emph{미분몫 \(\mathfrak M'\)에 \(Z\to\xi\)를 대입하면 \(\mathfrak o\)의 디퍼런트 \(\mathfrak d\)가 된다.} 준동형사상 아래에서는 아이디얼뿐 아니라 그 몫들도 서로 대응하기 때문이다. 실제로 \(\mathfrak o[..Z..]\)와 \(\mathfrak O_{\mathfrak o}\) 사이의 준동형, 곧 \(Z\)와 \(z\)의 대응 아래에서 \(\mathfrak A\)는 \(\mathfrak C\)의 준동형상이다. 이는 \(\mathfrak Q\)와 \(\mathfrak M_{\mathfrak o}\)에 관하여 \(\mathfrak P\)와 영아이디얼에 같은 사실이 성립하기 때문이다.\srcfn{13a)}{출판 시 주: 몫들이 대응한다고 결론내리기 위해서는 \(\mathfrak M_{\mathfrak o}\)가 \((0)\)의 원상 전체라는 사실이 중요하다.} 따라서 \(\mathfrak C\)에서 \(\xi\to Z\)를 대입하는 것에는 이 준동형 아래 \(\mathfrak A\)에서 \(\xi\to z\)를 대입하는 것이 대응한다. 그러므로 \(\mathfrak O\)의 디퍼런트 \(\mathfrak D\)는 \(Z\)와 \(z\)의 대응 아래 \(\mathfrak M'\)의 준동형상이다. 따라서
\[
\mathfrak M'[Z\to z]=\mathfrak D
\qquad\hbox{이고}\qquad
\mathfrak M'[Z\to\xi]=\mathfrak D[z\to\xi]=\mathfrak d.
\]

\paragraph{3. 정의방정식이 있는 경우의 디퍼런트.}
\(\mathfrak O\)가 정의방정식 \(G(z)=0\)을 갖는다고 하자(§ 1, 3의 끝). 또 최고차항의 계수는 단위원이라고 \emph{가정}하자. 곧
\[
G(Z)=Z^n+h_1Z^{n-1}+\cdots+h_n
\]
이고 계수 \(h_i\)들은 \(\mathfrak h\)에 속한다. 그러면 다항식 미분몫과의 관계는 다음과 같다.

\paragraph{정리.}
\emph{\(\mathfrak h\)에 관한 \(\mathfrak o\)와 \(\mathfrak O\)의 디퍼런트는 정의되어 있고(곧 \(\mathfrak h\)에 관한 직접곱 \(\mathfrak O_{\mathfrak o}\)가 존재하고), 각각 기저원소 \(G'(\xi)\)와 \(G'(z)\)를 갖는 주아이디얼이다. 여기서 \(G'(Z)\)는 다항식 미분몫이다. 아이디얼 미분몫 \(\mathfrak M'\)은 \((G'(Z),G(Z))\)와 같다.}

\paragraph{증명.}
먼저 가정에서 원소 \(e,z,\ldots,z^{n-1}\)이 \(\mathfrak O\)의 독립 가군기저를 이룬다는 것이 따른다. \(G(Z)\)의 최고차계수가 단위원이므로 이 거듭제곱들은 실제로 \(\mathfrak h\)-가군기저이다. 또한 \(\mathfrak M\)은 \(n\)차보다 낮은 다항식을 포함할 수 없으므로 이 기저는 독립이다. 실제로 다항식 \(K(Z)G(Z)\)의 차수는 역시 최고차계수에 관한 가정 때문에 두 인자의 차수의 합과 같다. 따라서 가군기저의 존재에 의해 \(\mathfrak h\)에 관한 직접곱 \(\mathfrak O_{\mathfrak o}\)가 존재하고, 디퍼런트가 정의된다. 나머지는 \(\mathfrak o[..Z..]\)에서 차분몫 \(\mathfrak C\)를 고찰하면 따른다(2를 보라). 이 차분몫은 다음 기저다항식을 갖는 주아이디얼이다.
\[
%
C(Z,\xi)=(Z^{n-1}+Z^{n-2}\xi+\cdots+\xi^{n-1})
+h_1(Z^{n-2}+\cdots+\xi^{n-2})+\cdots+h_{n-1}.
%
\]
실제로 \(C\)는 \(G(Z)-G(\xi)\)를 \(Z-\xi\)로 형식적으로 나누어 얻으므로 \(\mathfrak C\)에 속한다. \(Z-\xi\)가 \(\mathfrak Q\)의 기저이기 때문이다. 임의의 다항식은 \(C\)를 법으로 하여 \(Z\)에 관하여 \(n-1\)차에 이르지 않는 다항식 \(D\)로 줄일 수 있다. 따라서 곱 \(D(Z-\xi)\)는 \(n\)차에 이르지 않고, 사라지지 않는 한 \(\mathfrak M\)에 속할 수 없다. 그런데 \(D(Z-\xi)=0\)이면 \(Z-\xi\)에서 \(Z\)의 계수가 단위원이므로 \(D=0\)이다. 따라서 \(C\)가 \(\mathfrak C\)의 기저이다. 이제 \(C\)와 그 배수들에 \(\xi\to Z\)를 대입하면(\(G(Z)=C(Z-\xi)\)와 그 배수들도 함께 대입하면)
\[
\mathfrak M'=(G'(Z),G(Z))
\]
을 얻고, 따라서
\[
\mathfrak d=(G'(\xi))
\qquad\hbox{이고}\qquad
\mathfrak D=(G'(z)).
\]

\subsection*{§ 4. 직접합의 디퍼런트}

\(\mathfrak O\)에 다음 가정들을 둔다.

\begin{enumerate}
\item[(1)] \(\mathfrak O\)는 아이디얼들의 직접합
\[
\mathfrak O=\mathfrak R_1+\cdots+\mathfrak R_r
=\mathfrak R_i+\mathfrak S_i
=\mathfrak Oe_1+\cdots+\mathfrak Oe_r
\]
이다. 여기서 \(e_i\)들은 단위원의 성분들이다.

\item[(2)] 각 \(\mathfrak R_i\)는 독립 \(\mathfrak h_i\)-가군기저 \(T_i\)를 갖는다. 여기서 \(\mathfrak h_i\)는 \(\mathfrak R_i\) 안의 \(\mathfrak h\)의 성분, 곧 \(\mathfrak h_i=\mathfrak h e_i\)이다. 이 기저 \(T_i\)는 \(\mathfrak R_i\)의 단위원 \(e_i\)를 원소로 포함한다.

\item[(3)] 모든 교집합 \([\mathfrak h,\mathfrak S_i]\)는 영원소와 같다. 즉 \(\mathfrak S_i\)는 상수를 포함하지 않는다.
\end{enumerate}

동형 때문에 \(\mathfrak o\)에도 같은 가정들이 성립한다.

\paragraph{정리.}
\emph{위 가정들 아래에서 \(\mathfrak h\)에 관한 \(\mathfrak O\)의 디퍼런트는 \(\mathfrak h_i\)에 관한 \(\mathfrak R_i\)들의 디퍼런트들의 직접합이다.}

증명은 여러 개별 사실을 차례로 보이는 데 바탕을 둔다.

\paragraph{I.}
먼저 \(\mathfrak O\)도 \(e\)를 포함하는 독립 \(\mathfrak h\)-가군기저를 가지며, 따라서 \(\mathfrak h\)에 관한 \(\mathfrak O\)와 \(\mathfrak o\)의 디퍼런트가 정의된다는 것이 따른다. 이 기저는 \(\mathfrak R_i\)의 기저 \(T_i\)들의 합집합 \(T\)이다. 실제로
\[
x=xe_1+\cdots+xe_r
\]
와
\[
xe_i=(h_1e_i)t_{1i}e_i+\cdots+(h_\lambda e_i)t_{\lambda i}e_i
=h_1(t_{1i}e_i)+\cdots+h_\lambda(t_{\lambda i}e_i)
\]
에서 \(T\)가 실제로 \(\mathfrak O\)의 \(\mathfrak h\)-기저임이 나온다. 이 기저는 또한 독립이다. \(T\)의 원소들 사이의 한 관계에서는 직접합 때문에 각 \(T_i\)의 원소들 사이의 관계들이 따르고, 따라서 모든 계수에 대하여 \(he_i=0\)이다. 그런데 가정 (3)에 따라 이것은 \(h=0\)을 준다. 실제로
\[
h\equiv h e_i\ (\mathfrak S_i)
\]
이므로 \(he_i=0\)이면 \(h\)는 \([\mathfrak h,\mathfrak S_i]\)에 속한다. 또 \(T\)는 \(\mathfrak R_i\)의 단위원 \(e_i\)들을 원소로 포함하고 그 합은 \(e\)이다. 따라서 \(T\)에서 한 \(e_i\)를 \(e\)로 바꾸면 모든 조건을 만족하는 \(\mathfrak h\)-기저 \(\overline T\)를 얻는다. 여기서 \(\mathfrak O_{\mathfrak o}\)와 \(\mathfrak h\)에 관한 \(\mathfrak O\)의 디퍼런트가 존재한다.

\paragraph{II.}
이제 교집합 조건 (3)은 더 강한 조건
\[
[\mathfrak o,\mathfrak S_{i(\mathfrak o)}]=0
\]
을 함의한다. 여기서 \(\mathfrak S_{i(\mathfrak o)}\)는 \(\mathfrak S_i\)가 \(\mathfrak O_{\mathfrak o}\) 안에서 유도하는 아이디얼이다. 다음과 같이 쓰자.
\[
\alpha=\alpha e=\beta_1t_1+\cdots+\beta_\lambda t_\lambda.
\]
여기서 \(t_1,\ldots,t_\lambda\)는 \(T\)에서 \(T_i\)를 빼서 얻는 기저의 원소들이고, \(\alpha,\beta\)는 \(\mathfrak o\)에 속한다. \(e_i\)를 \(e\)로 바꾸어 \(\overline T\)를 만들면 \(e,t_1,\ldots,t_\lambda\)는 \(\mathfrak O\)의 \(\mathfrak h\)-기저 \(\overline T\)의 서로 다른 원소들이다. 그런데 \(\mathfrak O_{\mathfrak o}\)에서의 등식은 그러한 모든 기저에서 계수들이 같다는 등식으로 환원되므로, 관계
\[
\alpha=\alpha e=\beta_1t_1+\cdots+\beta_\lambda t_\lambda
\]
에서 \(\alpha=0\)이 따른다. 이로써 \emph{강화된 교집합 성질을 증명하였다}. 특히 여기서 (위와 같이)
\[
\alpha\text{가 }\mathfrak o\text{의 원소이고}\quad\alpha\ne0
\]
이면 \(\alpha e_i\ne0\)임이 따른다.

\paragraph{III.}
이 교집합 조건의 직접적인 귀결은 다음과 같다. \(\mathfrak o\)는 \(\mathfrak O\)의 분해에 대응하는 그 \emph{성분들} \(\mathfrak o e_i\)와 \emph{환동형}이다. 이 동형은
\[
\mathfrak h\simeq\mathfrak h_i=\mathfrak h e_i
\]
라는 동형을 포함한다. 실제로 \(e_i^2=e_i\)이고 \(e_i\)가 \(\mathfrak o\)의 모든 원소와 교환하므로 \(\mathfrak o e_i\)는 \(\mathfrak o\)의 환동형상이다. 또 \(\alpha\ne0\)이면 \(\alpha e_i\ne0\)이므로 이 대응은 일대일이다. 이때 부분환 \(\mathfrak h\)에는 성분 \(\mathfrak h_i=\mathfrak h e_i\)가 대응한다.

\paragraph{IV.}
\(\mathfrak o\)의 대응하는 분해
\[
\mathfrak o=\mathfrak r_1+\cdots+\mathfrak r_r
=\mathfrak r_i+\mathfrak s_i
=\mathfrak o\varepsilon_1+\cdots+\mathfrak o\varepsilon_r
\]
도 이용하면 \emph{더 많은 동형}이 따른다. 곧
\[
\mathfrak o\simeq\mathfrak o e_i,
\quad \mathfrak O\simeq\mathfrak O\varepsilon_i,
\quad \mathfrak r_\lambda\simeq\mathfrak r_\lambda e_i,
\quad \mathfrak R_\lambda\simeq\mathfrak R_\lambda\varepsilon_i;
\quad \mathfrak h\simeq\mathfrak h e_i\simeq\mathfrak h\varepsilon_i\simeq\mathfrak h e_i\varepsilon_i.
\]
관계 \(\mathfrak O\simeq\mathfrak O\varepsilon_i\)는 위 고찰에서 \(\mathfrak O\) 대신 \(\mathfrak o\)에서 출발할 때 \(\mathfrak o\simeq\mathfrak o e_i\)에 정확히 대응하는 것이다. 이어지는 두 관계도 서로 대응한다. 그 첫째 것은 \(\mathfrak o\simeq\mathfrak o e_i\) 아래 부분환 \(\mathfrak r_\lambda\)에도 \(\mathfrak r_\lambda e_i\)가 대응한다는 데서 나온다. 마찬가지로
\[
\mathfrak h\simeq\mathfrak h e_i\simeq\mathfrak h\varepsilon_i
\]
를 얻는다. 그런데 \(\mathfrak h\varepsilon_i\)는 \(\mathfrak o\), 그리고 \(\mathfrak r_i\)의 부분환이므로 끝으로
\[
\mathfrak h\varepsilon_i\simeq\mathfrak h\varepsilon_i e_i
\]
도 얻는다.

\paragraph{V.}
\(\mathfrak O\)의 직접합 표현에서 다음이 따른다. \emph{\(\mathfrak C=(0):\mathfrak B\)가 \(\mathfrak O\) 안에서 영아이디얼을 아이디얼 \(\mathfrak B\)로 나눈 몫이면, 각 성분 \(\mathfrak C_i\)는 \(\mathfrak R_i\)의 영아이디얼을 성분 \(\mathfrak B_i\)로 나눈 몫과 같다.} 실제로
\[
\mathfrak B\mathfrak C=(\mathfrak B e_1)(\mathfrak C e_1)+\cdots+(\mathfrak B e_r)(\mathfrak C e_r)=0
\]
에서 \emph{직접합}이므로 각각
\[
(\mathfrak B e_i)(\mathfrak C e_i)=0
\]
이 따른다. 역으로 \(\mathfrak R_i\)의 원소 \(r e_i\)가 조건 \((\mathfrak B e_i)(r e_i)=0\)을 만족하면 \(i\ne k\)인 경우에도 \((\mathfrak B e_k)(r e_i)=0\)이다.

따라서 \(r e_i\)는 \(\mathfrak C\)와 \(\mathfrak R_i\)의 원소이고, 그러므로 성분 \(\mathfrak C_i=\mathfrak C e_i\)에 속한다. 이로써 \(\mathfrak C_i=\mathfrak C e_i\)가 영아이디얼을 \(\mathfrak B_i\)로 나눈 몫임을 알았다. 또한
\[
r e_i\mathfrak B_i\equiv0\quad(\mathfrak S_i)
\]
이면 언제나 \(r e_i\mathfrak B_i=0\)이므로 이는 \(\mathfrak R_i\)의 영아이디얼을 \(\mathfrak B_i\)로 나눈 몫이기도 하다.\srcfn{14)}{이 정리는 2에서 쓴 정리, 곧 준동형사상 아래에서 아이디얼뿐 아니라 그 몫들도 서로 대응한다는 정리에서도 바로 따른다.} \(\mathfrak O\)뿐 아니라
\[
\mathfrak O_{\mathfrak o}=\mathfrak O_{\mathfrak o}e_1+\cdots+\mathfrak O_{\mathfrak o}e_r
\]
도 직접합이므로, \emph{이 정리는 \(\mathfrak O_{\mathfrak o}\)의 아이디얼들에도 성립한다.}

\paragraph{VI. 디퍼런트의 성분분해.}
정의에 따라(§ 3, 1을 보라) \(\mathfrak D=\mathfrak A[\xi\to x]\)였으므로
\[
\mathfrak D=\mathfrak D e_1+\cdots+\mathfrak D e_r
=\mathfrak A[\xi\to x]e_1+\cdots+\mathfrak A[\xi\to x]e_r.
\]
그런데 \(e_i\)는 \(\mathfrak O\)에 속하므로 자기 자신에 대응하고, \(\varepsilon_i\)는 동형 아래 \(e_i\)로 가며, \(e_i^n=e_i\)이므로
\[
\mathfrak A[\xi\to x]e_i=\mathfrak A e_i[\xi\to x]
=\mathfrak A e_i\varepsilon_i[\xi\to x].
\]
한편 V에 따라 \(\mathfrak A e_i\varepsilon_i\)는 몫 \((0):\mathfrak P e_i\varepsilon_i\), 그리고 \(\mathfrak R_i\mathfrak r_i\)의 영아이디얼을 \(\mathfrak P e_i\varepsilon_i\)로 나눈 몫과도 같다. 실제로 \(\mathfrak O\)와 \(\mathfrak o\)의 직접합분해를 곱하면
\[
\mathfrak O_{\mathfrak o}=\sum\mathfrak O_{\mathfrak o}e_i\varepsilon_k
=\sum\mathfrak R_i\mathfrak r_k
\]
라는 분해를 얻는다. \(\sum e_i\varepsilon_k=e\)이고 단위원들에 직교관계가 성립하므로 이 합은 직접합이다.

\paragraph{VII.}
\emph{\(\mathfrak D e_i=\mathfrak A e_i\varepsilon_i[\xi\to x]=(0):\mathfrak P e_i\varepsilon_i\)가 \(\mathfrak h_i\)에 관한 \(\mathfrak R_i\)의 디퍼런트임을 보이자.} 이는 다음 사실들에 바탕을 둔다. \(\mathfrak R_i\mathfrak r_i\)는 \(\mathfrak h e_i\varepsilon_i\)에 관한 \(\mathfrak R_i\varepsilon_i\times\mathfrak r_i e_i\)와 같다. 그 차분아이디얼은 \(\mathfrak P e_i\varepsilon_i\), 차분몫은 \(\mathfrak A e_i\varepsilon_i\)이다. 한편 \(\mathfrak D e_i\varepsilon_i\)는 \(\mathfrak R_i\varepsilon_i\)의 디퍼런트이고, 이에 대응하여 \(\mathfrak d\varepsilon_i e_i\)는 \(\mathfrak r_i e_i\)의 디퍼런트이다. 이제 동형
\[
\mathfrak R_i\varepsilon_i\simeq\mathfrak R_i
\qquad\text{와}\qquad
\mathfrak r_i\simeq\mathfrak r_i e_i
\]
에서 \(\mathfrak D e_i\)는 \(\mathfrak h e_i\)에 관한 \(\mathfrak R_i\)의 디퍼런트이고, \(\mathfrak d\varepsilon_i\)는 \(\mathfrak h\varepsilon_i\)에 관한 \(\mathfrak r_i\)의 디퍼런트임이 따른다.

\paragraph{VIII.}
\emph{\(\mathfrak R_i\mathfrak r_i\)는 \(\mathfrak h e_i\varepsilon_i\)에 관한 직접곱 \(\mathfrak R_i\varepsilon_i\times\mathfrak r_i e_i\)와 같다.} \(\mathfrak R_i\)가 \(\mathfrak h_i\)-가군기저 \(T_i\)를 가지므로 동형 (IV) \(\mathfrak R_i\simeq\mathfrak R_i\varepsilon_i\)에서 \(T_i\varepsilon_i\)가 \(\mathfrak R_i\varepsilon_i\)의 \(\mathfrak h e_i\varepsilon_i\)-가군기저임이 따른다. 이제 IV에 따라 \(\mathfrak r_i e_i\)는 \(\mathfrak R_i\varepsilon_i\)와 동형이고 --- \(\mathfrak o\simeq\mathfrak O\)에 따라 \(\mathfrak r_i\simeq\mathfrak R_i\)이므로 --- \(\mathfrak h e_i\varepsilon_i\)의 확대환이다. 따라서 직접곱은 \(T_i\varepsilon_i\)의 원소들을 \(\mathfrak r_i e_i\)의 계수로 선형결합한 모든 선형형식으로 이루어지며, 등식은 계수들의 등식으로 정의된다. 그런데 \(\mathfrak O\times\mathfrak o\) 안에서 이것은 정확히 \(\mathfrak R_i\mathfrak r_i\)이다. 실제로 \(T_i=T_i e_i\), \(\mathfrak r_i=\mathfrak r_i\varepsilon_i\)이고 \(T_i\)는 \(\mathfrak O\)의 한 가군기저에 속한다.

\paragraph{IX.}
\emph{\(\mathfrak R_i\varepsilon_i\times\mathfrak r_i e_i\)의 차분아이디얼은 \(\mathfrak P e_i\varepsilon_i\)이고, 차분몫은 \(\mathfrak A e_i\varepsilon_i\)이다.} 이를 보려면 \(\mathfrak P\)를
\[
\mathfrak P=\{\ldots,xe_1-\xi\varepsilon_1,\ldots,xe_r-\xi\varepsilon_r,\ldots\}
\]
로 쓰면 된다. 그러면
\[
\mathfrak P e_i\varepsilon_i
=\{\ldots,xe_i\varepsilon_i-\xi\varepsilon_i e_i,\ldots\}
\]
이므로 이는 실제로 \(\mathfrak R_i\varepsilon_i\times\mathfrak r_i e_i\)의 차분아이디얼과 같다. 그런데 (VI에 따라) \(\mathfrak A e_i\varepsilon_i\)는 이 환의 영아이디얼을 \(\mathfrak P e_i\varepsilon_i\)로 나눈 몫이므로 차분몫이다.

\paragraph{X.}
\emph{\(\mathfrak R_i\varepsilon_i\)의 디퍼런트는 \(\mathfrak D e_i\varepsilon_i\)와 같다.} \(\mathfrak R_i\varepsilon_i\)의 디퍼런트는 \(\mathfrak A e_i\varepsilon_i\)에서 \(\mathfrak r_i e_i\)의 각 원소를 그와 동형인 원소로 바꾸어 얻는다. 곧 매번 \((\xi\varepsilon_i)e_i\)를 \((xe_i)\varepsilon_i\)로 바꾼다. 다시 말해 \(\xi\varepsilon_i\)도 \(\mathfrak o\)의 어떤 \(\xi\)이고 이는 \(xe_i\)로 가므로, \(\mathfrak A e_i\varepsilon_i\)에서 \(\xi\to x\)를 대입하고 --- 이 대입은 \((\xi\varepsilon_i)e_i\)를 \((xe_i)e_i=xe_i\)로 보낸다 --- 이어 \(\varepsilon_i\)를 곱한다. 그런데 (VI에 따라) \(\mathfrak A e_i\varepsilon_i[\xi\to x]=\mathfrak D e_i\)였으므로, \(\mathfrak h e_i\varepsilon_i\)에 관한 \(\mathfrak R_i\varepsilon_i\)의 디퍼런트는 \(\mathfrak D e_i\varepsilon_i\)와 같다. 대칭적으로 \(\mathfrak r_i e_i\)의 디퍼런트는 \(\mathfrak d e_i\varepsilon_i\)와 같다.

\paragraph{XI.}
\emph{\(\mathfrak R_i\)의 디퍼런트는 \(\mathfrak D e_i\)와 같다.} 실제로 동형 아래 \(\mathfrak R_i\)의 디퍼런트는 \(\varepsilon_i\)를 곱하여 \(\mathfrak R_i\varepsilon_i\)의 디퍼런트로 간다. 마지막 것이 \(\mathfrak D e_i\varepsilon_i\)와 같으므로 전자는 \(\mathfrak D e_i\)와 같다. 따라서
\[
\mathfrak D=\mathfrak D e_1+\cdots+\mathfrak D e_r,
\]
\emph{그리고 \(\mathfrak D e_i\)는 \(\mathfrak R_i\)의 디퍼런트이다.} 대칭적으로
\[
\mathfrak d=\mathfrak d\varepsilon_1+\cdots+\mathfrak d\varepsilon_r,
\]
\emph{그리고 \(\mathfrak d\varepsilon_i\)는 \(\mathfrak r_i\)의 디퍼런트이다.} 이로써 정리가 증명되었다.

\subsection*{§ 5. 체의 갈루아 확대환. 구조정리}

\paragraph{1. 바탕이 되는 환들.}
\(\mathsf{K}\)가 \emph{기초체 \(P\) 위의 \(n\)차 체}이고 분리가능 확대라고 하자. \(\Gamma\)를 이에 속하는, 동형을 제외하고 유일하게 결정되는 갈루아체라 하자. 따라서 \(\Gamma\)는 \(n\)개의 켤레체 \(\mathsf{K}^{(i)}\)를 포함하고 그 곱이며, \(\mathsf{K}=\mathsf{K}^{(1)}\)로 둔다. \(\mathfrak K\)는 \(\mathsf{K}\), 따라서 모든 \(\mathsf{K}^{(i)}\)와 \(\mathsf{P}\)에 관하여 동형인 동치 확대이며, 교집합 \([\mathfrak K,\Gamma]\)가 \(\mathsf{P}\)라고 하자. 곧 \(\Gamma\)와 \(\mathfrak K\)의 공통 포괄체는 없다고 하자.\srcfn{15)}{물론 \(K\)와 \(\mathfrak K\)의 \(P\)에 속하지 않는 원소들 사이에는 아무 연산관계도 없다고 가정한다. 이는 일반성을 제한하지 않는다(§ 1, 2, § 2, 1과 각주 11을 보라).} \(\mathfrak K\)는 단위원을 포함하는 유한 독립 \(P\)-가군기저를 가지므로 § 2, 1에 따라 직접곱 \(\mathfrak K_\Gamma,\mathfrak K_K,\mathfrak K_{K^{(i)}}\)가 존재한다. \(\mathfrak K_\Gamma\)는 그 \(n\)개 부분환 \(\mathfrak K_{\mathsf{K}^{(i)}}\)의 곱이다. \(\mathfrak K_\Gamma\)를 \(\mathfrak K\)의 \emph{갈루아 확대환}이라 한다. \(\mathfrak K\)의 모든 독립 \(P\)-기저가 \(\mathfrak K_{\mathsf{K}}\)의 독립 \(K\)-기저가 되고 켤레들에 대해서도 마찬가지이므로, \(\mathfrak K_\Gamma\)는 동시에 \(K\)에 관한 \(\mathfrak K_{\mathsf{K}}\)와 \(\Gamma\)의 직접곱이다.\srcfn{16)}{모든 켤레에 똑같이 성립하는 정리들은 대개 첨자를 생략하고 \(\mathfrak K_{\mathsf{K}}\)에 대해서만 말한다.} 이 환들의 아이디얼 사이에는 다음 관계가 성립한다.

\paragraph{정리.}
\emph{\(\mathfrak K\)의 모든 아이디얼은 \(\mathfrak K_{\mathsf{K}}\) 또는 \(\mathfrak K_\Gamma\) 안의 그 확대아이디얼을 수축한 아이디얼이고, \(\mathfrak K_{\mathsf{K}}\)의 모든 아이디얼은 \(\mathfrak K_\Gamma\) 안의 그 확대아이디얼을 수축한 아이디얼이다.}

\paragraph{증명.}
\(\mathsf{P}\)와 \(\mathsf{K}\)가 체이므로 \(\mathfrak K\)의 모든 \(P\)-가군과 \(\mathfrak K_{\mathsf{K}}\)의 모든 \(K\)-가군은 각각 유한 독립 \(\mathsf{P}\)-가군기저와 \(\mathsf{K}\)-가군기저를 가지며, 이를 각각 \(\mathfrak K\)와 \(\mathfrak K_{\mathsf{K}}\)의 그러한 기저로 보충할 수 있다. 이제 정리는 § 2, 2의 정리에서 따른다.

\(\mathfrak P_{\mathsf{K}^{(i)}}^{(i)}\)를 \(\mathfrak K_{\mathsf{K}^{(i)}}\)의 차분아이디얼이라 쓰자. 곧
\[
\mathfrak P_{\mathsf{K}^{(i)}}^{(i)}=\{\ldots,x-\xi^{(i)},\ldots\},
\]
특히
\[
\mathfrak P_{\mathsf{K}}=\{\ldots,x-\xi,\ldots\}.
\]
\(\xi\)와 \(\xi^{(i)}\)는 각각 동형에 따라 \(x\)에 대응하는 \(\mathsf{K}\)와 \(\mathsf{K}^{(i)}\)의 원소이므로, 이 \(n\)개 차분아이디얼은 서로 켤레이다.

\(\mathfrak A_{\mathsf{K}}\)와 \(\mathfrak A_{\mathsf{K}^{(i)}}\)를 각각 \(\mathfrak K_{\mathsf{K}}\)에서의 차분몫 \((0):\mathfrak P_{\mathsf{K}}\)와 그 켤레들이라 하자. \(\mathfrak K_\Gamma\) 안의 확대아이디얼들은 \(\mathfrak P_\Gamma\)와 \(\mathfrak A_\Gamma\)로 쓴다. 이 항의 확대아이디얼과 수축아이디얼에 관한 정리에서 특히 다음이 따른다.

\emph{차분아이디얼 \(\mathfrak P_{\mathsf{K}}\)와 차분몫 \(\mathfrak A_{\mathsf{K}}\)는 그 확대 \(\mathfrak P_\Gamma\)와 \(\mathfrak A_\Gamma\)를 수축한 아이디얼들이다.}

\paragraph{2. 갈루아 확대환의 직접합분해.}
차분아이디얼·차분몫과 갈루아 확대환의 구조 사이의 관계는 다음 정리가 준다.

\paragraph{구조정리.}
\emph{갈루아 확대환 \(\mathfrak K_\Gamma\)의 영아이디얼은 \(n\)개 켤레 차분아이디얼들의 확대들의 교집합이고, \(\mathfrak K_\Gamma\) 자체는 \(n\)개 켤레 차분몫들의 확대들의 직접합이다. \(\mathfrak K_{\mathsf{K}^{(i)}}\)는 그에 속하는 차분아이디얼과 차분몫의 직접합이고, 그 영아이디얼은 이 둘의 교집합이다.}

\paragraph{증명.}
각 \(\mathfrak P_\Gamma^{(i)}\)에 따른 잉여류환은 \(\Gamma\)에 관하여 랭크 1이다. 실제로 모든 원소는 \(\mathfrak P_\Gamma^{(i)}\)를 법으로 하여 \(\Gamma\)의 한 원소와 합동이지만, \(\mathfrak P_\Gamma^{(i)}\)는 \(\Gamma\)의 0 아닌 원소를 포함할 수 없다. \(x\to\xi^{(i)}\)일 때 \(\mathfrak P_\Gamma^{(i)}\)가 사라지기 때문이다. 분리가능성에 따라 \(\mathfrak P_\Gamma^{(i)}\)들은 모두 서로 다르므로 쌍마다 서로소이다. 곧 각 두 아이디얼의 합은 \(\mathfrak K_\Gamma\)와 같다. 따라서 아이디얼들
\[
\mathfrak P_\Gamma^{(1)},
[\mathfrak P_\Gamma^{(1)},\mathfrak P_\Gamma^{(2)}],\ldots,
[\mathfrak P_\Gamma^{(1)},\ldots,\mathfrak P_\Gamma^{(n)}]
\]
의 랭크는 매번 적어도 하나씩 감소한다. 모두의 교집합은 영아이디얼이고, 한편 쌍마다 서로소이므로 모든 켤레보다 적은 수의 교집합은 영아이디얼과 다르다. 따라서 랭크는 각 단계에서 정확히 하나씩 감소하였다. 알려진 정리들에서 이제 다음과 같은 \(\mathfrak K_\Gamma\)의 유일하게 결정되는 직접합표현이 따른다.\srcfn{17)}{예컨대 E. Noether, \emph{Abstrakter Aufbau der Idealtheorie in algebraischen Zahl- und Funktionenkörpern}, Math. Ann. 96 (1927), 26--61쪽, § 4, 5를 보라.}
\begin{align}
\mathfrak K_\Gamma
&=\mathfrak C_\Gamma^{(1)}+\cdots+\mathfrak C_\Gamma^{(n)}
=\Gamma e^{(1)}+\cdots+\Gamma e^{(n)}
=\mathfrak C_\Gamma^{(i)}+\mathfrak P_\Gamma^{(i)}, \tag{1}\\
e&=e^{(1)}+\cdots+e^{(n)},\qquad
\mathfrak C_\Gamma^{(i)}=
[\mathfrak P_\Gamma^{(1)},\ldots,\mathfrak P_\Gamma^{(i-1)},
\mathfrak P_\Gamma^{(i+1)},\ldots].\notag
\end{align}
\(\mathfrak C_\Gamma=\mathfrak A_\Gamma\)임을 보여야 한다. 이는 \(\mathfrak K_{\mathsf{K}}\)에서 대응하는 일치를 보이면 따른다. 여기서는 분해
\begin{equation}
\mathfrak K_{\mathsf{K}}=\mathfrak C_{\mathsf{K}}+\mathfrak P_{\mathsf{K}}
=\mathfrak K_{\mathsf{K}} e^{(1)}+\mathfrak K_{\mathsf{K}}(e-e^{(1)}).\tag{2}
\end{equation}
가 성립한다. 실제로 \(e^{(1)}\)는 \(\mathfrak K_{\mathsf{K}}\)에 속한다. \(\mathfrak C_\Gamma^{(1)}\)가 \(\mathfrak P_\Gamma^{(2)},\ldots,\mathfrak P_\Gamma^{(n)}\)의 모든 치환을 허용하므로, \(\mathfrak C_\Gamma^{(1)}\) 안의 단위원 성분인 \(e^{(1)}\)에도 같은 것이 성립하기 때문이다. 따라서 \(e^{(1)}\)을 \(\mathfrak K_{\mathsf{K}}\)의 한 독립 \(K\)-기저로 나타내면 --- 이 기저는 동시에 \(\mathfrak K_\Gamma\)의 독립 \(\Gamma\)-기저이므로 가능하다 --- 계수들은 각각 \(\mathsf{K}^{(2)},\ldots,\mathsf{K}^{(n)}\)의 모든 치환을 허용하며, 따라서 \(\mathsf{K}\)에 속한다. \((e^{(1)})^2=e^{(1)}\)이므로 합
\[
\mathfrak K_{\mathsf{K}}=\mathfrak K_{\mathsf{K}} e^{(1)}+\mathfrak K_{\mathsf{K}}(e-e^{(1)})
\]
는 직접합이다. 이를 \(\mathfrak K_\Gamma\)로 확대해도 다시 직접합이며, (1)에 따라 \(\mathfrak K_\Gamma=\mathfrak C_\Gamma^{(1)}+\mathfrak P_\Gamma^{(1)}\)이 된다. 따라서 § 2, 2의 끝에 따라
\[
\mathfrak K_{\mathsf{K}} e^{(1)}=[\mathfrak C_\Gamma,\mathfrak K_{\mathsf{K}}]
=[\Gamma e^{(1)},\mathfrak K_{\mathsf{K}}]=\mathsf{K}e^{(1)}=\mathfrak C_{\mathsf{K}};
\]
\[
\mathfrak K_{\mathsf{K}}(e-e^{(1)})=[\mathfrak P_\Gamma,\mathfrak K_{\mathsf{K}}]=\mathfrak P_{\mathsf{K}}.
\]
이로써 (2)를 증명하였다.

그런데 (2)에서 \(\mathfrak C_{\mathsf{K}}=\mathfrak A_{\mathsf{K}}\), 따라서 \(\mathfrak C_\Gamma=\mathfrak A_\Gamma\)가 따른다. 실제로 \(e^{(1)}(e-e^{(1)})=0\)이므로 \(\mathfrak C_{\mathsf{K}}\mathfrak P_{\mathsf{K}}=0\)이다. 역으로 \(b\mathfrak P_{\mathsf{K}}=0\)이면 \(b(e-e^{(1)})=0\)도 되어 \(b=be=be^{(1)}\)이다. 따라서
\[
\mathfrak C_{\mathsf{K}}=(0):\mathfrak P_{\mathsf{K}}=\mathfrak A_{\mathsf{K}}
\]
이고, 확대하면 \(\mathfrak C_\Gamma=\mathfrak A_\Gamma\)이다. 그러므로
\begin{align}
\mathfrak K_\Gamma
&=\mathfrak A_\Gamma^{(1)}+\cdots+\mathfrak A_\Gamma^{(n)}
=\Gamma e^{(1)}+\cdots+\Gamma e^{(n)};
&0&=[\mathfrak P_\Gamma^{(1)},\ldots,\mathfrak P_\Gamma^{(n)}]; \tag{3}\\
\mathfrak K_{\mathsf{K}^{(i)}}
&=\mathfrak A_{\mathsf{K}^{(i)}}+\mathfrak P_{\mathsf{K}^{(i)}}^{(i)}
=\mathsf{K}^{(i)}e^{(i)}+\mathfrak P_{\mathsf{K}^{(i)}}^{(i)};
&0&=[\mathfrak A_{K^{(i)}}^{(i)},\mathfrak B_{K^{(i)}}^{(i)}].\notag
\end{align}
이로써 정리의 모든 부분을 증명하였다. 끝으로 이것은 \emph{완전히 불변인 구조정리}임을 강조해 둔다. 나타나는 모든 구성 \(\mathfrak K_K,\mathfrak K_\Gamma,\mathfrak B_K,\mathfrak A_K\)가 기저 없이 정의되기 때문이다.\srcfn{18)}{여기서 가능한, 정의방정식과 대응하는 다항식의 일차인수분해를 이용한 정식화는 불변적이지 않다.}

\paragraph{3. 켤레원소들에 의한 \(\mathfrak K\)의 성분표현.}
켤레 차분아이디얼들과 차분몫들에는 \(\mathfrak K\)의 원소들을 켤레성분들로 나타내는 표현이 대응한다.

\paragraph{정리.}
\emph{
\[
x=\xi^{(1)}e^{(1)}+\cdots+\xi^{(n)}e^{(n)}
\]
가 \(\mathfrak K\)의 원소 \(x\)의 \(\mathfrak K\)-성분표현이면, 이 \(n\)개 성분은 서로 켤레이다. 특히 \(\xi^{(1)},\ldots,\xi^{(n)}\)은 \(\mathfrak K\)에서 \(\mathsf{K}^{(1)},\ldots,\mathsf{K}^{(n)}\)으로 가는 동형들에 따라 \(x\)에 대응하는 \(n\)개 켤레값이다.}

\(e^{(1)},\ldots,e^{(n)}\)이 서로 켤레라는 것은 2의 관계 (3)에서 알 수 있다. 실제로 \(\mathfrak P_{\mathsf{K}}\)를 \(\mathfrak P_{\mathsf{K}^{(i)}}^{(i)}\)로 보내는 \(\Gamma\)의 자기동형사상들은 동시에 \(e^{(1)}\)을 \(e^{(i)}\)로 보낸다. 마지막 것은 실제로 \(\mathfrak K_{\mathsf{K}^{(i)}}\)의 원소이다. 같은 관계들은 \(\mathfrak P_{\mathsf{K}^{(i)}}^{(i)}e^{(i)}=0\)이므로 \(xe^{(i)}=\xi^{(i)}e^{(i)}\), 따라서
\[
(x-\xi^{(i)})e^{(i)}=0
\]
도 준다. 또는 \(\mathfrak K\)는 체로서 그 성분과 동형이고, 그 성분은 \(x\)와 \(\xi^{(i)}e^{(i)}\)가 동형으로 대응하는 \(\mathsf{K}^{(i)}e^{(i)}\)이므로 같은 사실이 따른다.

정리에서 다음이 나온다.

\paragraph{따름정리.}
\emph{\(\mathfrak B\)가 \(\mathfrak K\)의 절대가군(곧 임의의 두 원소와 함께 그 차도 포함하는 가군)이고 \(\mathfrak B^{(1)},\ldots,\mathfrak B^{(n)}\)이 \(\mathfrak K_\Gamma\) 안의 그 성분들이면, \(\mathfrak B\)는 교집합 \([\mathfrak K,\mathfrak B^{(1)}+\cdots+\mathfrak B^{(n)}]\)과 같다.}

\(\mathfrak B\)는 이 교집합에 포함된다. 역으로 \(\mathfrak B^{(1)}+\cdots+\mathfrak B^{(n)}\)의 한 원소가 \(\mathfrak K\)에 속하려면 정리에 따라 켤레성분들의 합이어야 하며, 따라서 \(\mathfrak B\)의 원소이다.

\paragraph{4. \(\mathfrak K\)의 여차기저들과 단위원 성분들의 관계.}
불변인 구조에서 기저표현으로 넘어가면, 이 구조는 기저들을 여차기저의 쌍들로 묶어 준다.

\paragraph{정리.}
\emph{\(\mathfrak K\)의 모든 \(P\)-기저 \(t_1,\ldots,t_n\)에는 여차기저 \(T_1,\ldots,T_n\)이 대응하고, 그 여차기저는 다시 원래의 \(t\)들의 기저이다. 여차기저 \(T_1,\ldots,T_n\)과 동형인 \(\mathsf{K}^{(i)}\)의 기저 \(\mathsf{A}_1^{(i)},\ldots,\mathsf{A}_n^{(i)}\)는 \(t\)들에 의한 기저표현
\[
e^{(i)}=\mathsf{A}_1^{(i)}t_1+\cdots+\mathsf{A}_n^{(i)}t_n
\]
에서 단위원 성분 \(e^{(i)}\)의 계수들로 정의된다. \(\alpha_1^{(i)},\ldots,\alpha_n^{(i)}\)가 \(t_i\)들에 대응하는 기저들이면 \(\alpha\)와 \(\mathsf{A}\)의 행렬들은 서로 역행렬이다. \(\mathfrak K\)의 두 기저 \((s)\)와 \((t)\)가 변환
\[
(s)=(t)P
\]
로 이어지면, 그 여차기저들 사이에는 반변관계
\[
(S)=P^{-1}(T)
\]
가 성립한다.}

\paragraph{증명.}
\(\mathfrak K\)의 한 \(P\)-기저 \(t_1,\ldots,t_n\)을 바탕으로 단위원의 \(n\)개 성분의 기저표현을 \(e^{(i)}=A_1^{(i)}t_1+\cdots+A_n^{(i)}t_n\)이라 하자. 이 기저는 \(\mathfrak K_\Gamma\)의 \(\Gamma\)-기저이기도 하다. 2와 3에 따라 \(\mathsf{A}_1^{(i)},\ldots,\mathsf{A}_n^{(i)}\)는 \(\mathsf{K}^{(i)}\)에 속하고, 서로 다른 \(e^{(i)}\)들에는 켤레인 계수계 \(A\)들이 대응한다. \(\alpha_1^{(i)},\ldots,\alpha_n^{(i)}\)가 \(\mathsf{K}^{(i)}\)에서 \(t_1,\ldots,t_n\)에 동형으로 대응하는 기저이면 다음 관계들도 성립한다.
\begin{equation}
e^{(i)}[t\to\alpha^{(i)}]=e;
\qquad e^{(i)}[t\to\alpha^{(j)}]=0\quad\text{(\(i\ne j\)).} \tag{1}
\end{equation}
실제로 \(\mathfrak P_{\mathsf{K}^{(i)}}^{(i)}\)는 원소 \(t_1-\alpha_1^{(i)},\ldots,t_n-\alpha_n^{(i)}\) --- 이들은 \(\mathfrak P_{\mathsf{K}^{(i)}}^{(i)}\)의 종속 \(P\)-가군기저를 이룬다 --- 를 포함하므로 \(t=\alpha^{(i)}\)를 대입하면 사라진다. 따라서
\[
e\equiv e^{(i)}\;(\mathfrak P_{\mathsf{K}^{(i)}}^{(i)})
\]
에서 (1)의 첫 관계가 나오며, 둘째 관계는 \(i\ne j\)일 때 \(e^{(i)}\)가 모든 \(\mathfrak P_{\mathsf{K}^{(j)}}^{(j)}\)의 원소라는 데서 나온다. 관계 (1)을 계수별로 풀면 행렬식으로 다음과 같이 쓸 수 있다.
\begin{equation}
\begin{pmatrix}
\alpha_1^{(1)}&\ldots&\alpha_n^{(1)}\\
\vdots&&\vdots\\
\alpha_1^{(n)}&\ldots&\alpha_n^{(n)}
\end{pmatrix}
\begin{pmatrix}
A_1^{(1)}&\cdots&A_1^{(n)}\\
\vdots&&\vdots\\
A_n^{(1)}&\cdots&A_n^{(n)}
\end{pmatrix}=E_n. \tag{2}
\end{equation}
여기서 \(E_n\)은 단위행렬이다. 따라서 \(\mathsf{A}\)들은 이미 \(\alpha\)들에 의해 유일하게 결정된다. 관계 (2)는 두 행렬의 행렬식이 사라지지 않음을 보인다. 더 나아가 \(\mathsf{A}_1^{(i)},\ldots,\mathsf{A}_n^{(i)}\)는 \(\mathsf{P}\)에 관하여 선형독립이다. 어떤 선형종속관계가 있으면 켤레계들에도 성립하여 \(\mathsf{A}\)-행렬의 랭크를 줄이는데, 이는 \(E_n\)의 랭크가 \(n\)인 것과 모순이다. 따라서 매번 \(\mathsf{A}_1^{(i)},\ldots,\mathsf{A}_n^{(i)}\)는 \(\mathsf{K}^{(i)}\)의 기저이다. 이 기저들은 서로 켤레이므로 모두 \(\mathfrak K\)의 같은 기저 \(T_1,\ldots,T_n\)과 동형이며, 3에 따라 \(T_\lambda\)는
\[
T_\lambda=\mathsf{A}_\lambda^{(1)}e^{(1)}+\cdots+\mathsf{A}_\lambda^{(n)}e^{(n)}
\]
로 결정된다. 기저 \(T_1,\ldots,T_n\)을 \(t_1,\ldots,t_n\)의 여차기저라 하고, 이에 대응하여 \(\mathsf{A}_1^{(i)},\ldots,\mathsf{A}_n^{(i)}\)를 \(\alpha_1^{(i)},\ldots,\alpha_n^{(i)}\)의 여차기저라 한다. \(\mathsf{A}\)들은 (2)에 따라 이미 유일하게 결정되고, 이 관계는 전치행렬로 넘어가면 \(T\), 곧 \(\mathsf{A}\)의 기저에서 \(\alpha\)들을 결정하므로 원래의 \(t\)들의 기저는 \(T\)의 여차기저이다. 따라서 단위원 성분들이 주는 대응 아래 모든 기저는 여차기저들의 쌍들로 나뉜다. \(s,S\)가 여차기저들의 또 다른 쌍이고
\[
e^{(i)}=\mathsf{B}_1^{(i)}s_1+\cdots+\mathsf{B}_n^{(i)}s_n
\qquad\text{이며}\qquad (s)=(t)P
\]
라고 하자. 그러면
\[
\mathsf{A}_1t_1+\cdots+\mathsf{A}_nt_n=\mathsf{B}_1s_1+\cdots+\mathsf{B}_ns_n
\]
에 \((s)=(t)P\)를 대입하여 \(t\)에 관한 항등식으로 보면 \(\mathsf{B},\mathsf{A}\)가 \(s,t\)에 반변적으로 변환함이 따른다. 따라서 동형에 의하여 \(S,T\)에도 같은 사실이 성립한다. 이로써 정리의 모든 부분을 증명하였다.

\paragraph{주.}
\(c=C_1t_1+\cdots+C_nt_n\)이 한 \(\mathsf{P}\)-기저에 의한 \(\mathfrak K\) 원소의 표현이고 \(c=c_1T_1+\cdots+c_nT_n\)이 여차기저에 의한 같은 원소의 표현이면
\[
C_\lambda=\operatorname{Sp}(cT_\lambda),
\qquad c_\lambda=\operatorname{Sp}(ct_\lambda)
\]
이다. 여기서 통상적으로 \(\operatorname{Sp}(x)\)는 켤레들의 합
\[
\operatorname{Sp}(x)=\xi^{(1)}+\xi^{(2)}+\cdots+\xi^{(n)}
\]
을 뜻하며 \(x\)는 \(\mathfrak K\)의 임의의 원소이다. 실제로
\[
c=\gamma^{(1)}e^{(1)}+\cdots+\gamma^{(n)}e^{(n)}
=\sum_\lambda\gamma^{(1)}A_\lambda^{(1)}t_\lambda+
\cdots+\sum_\lambda\gamma^{(n)}A_\lambda^{(n)}t_\lambda
\]
이므로 계수들을 비교하면 \(C_\lambda=\operatorname{Sp}(cT_\lambda)\)를 얻는다. 둘째 관계는 \(t\)와 \(T\)를 맞바꾸면 마찬가지로 따른다.\srcfn{19)}{이 항의 고찰은 불변인 구조정리들이 불변이 아닌 기저정리들보다 본질적으로 더 단순함을 보인다. 마지막 정리들은 알려진 것과의 관계를 잇기 위해서만 필요했다. 실제로 갈루아 확대환의 분해와 단위원 성분들의 도입은 여차기저들(다음 절의 여모듈)이 하는 모든 일을 해낸다. 그곳에서도 디퍼런트와 여모듈의 관계는 분류를 위해서만 제시한다. 5의 라그랑주 보간공식에 관한 고찰도 마찬가지이다.}

\paragraph{5. 정의방정식을 바탕으로 하는 경우. 라그랑주 보간공식.}
[출판 시 주: 이 항은 원고에서 채워지지 않은 채로 남았다.]

\subsection*{§ 6. 오더의 디퍼런트}

지금까지 체와 그 갈루아 확대환에만 관한 구조정리들을 이제 정수성의 뜻에서 강화한다. 이를 위해 \(\mathsf{P}\)의 어떤 부분환들과, 그에 관하여 \(\mathfrak K\)와 \(\mathsf{K}\)에서 정수적인 원소들로 이루어진 오더들을 지정한다. § 5의 모든 표기를 그대로 쓴다.\srcfn{20)}{한 환에 관하여 정수적인 것과 정수적으로 닫힌 것의 개념은 각주 17에서 인용한 논문의 § 1을 보라.}

\paragraph{1. 바탕이 되는 환들.}
\(\mathfrak h\)가 \(\mathsf{P}\)의 부분환이고 \(\mathsf{P}\)가 \(\mathfrak h\)의 몫체라고 하자. 또 \(\mathfrak h\)는 \emph{\(\mathsf{P}\)에서 정수적으로 닫혀 있다}고 하자. 곧 \(\mathsf{P}\)의 원소 가운데 \(\mathfrak h\)에 관하여 정수적인 것은 모두 \(\mathfrak h\)에 속한다. \(\mathfrak O\)를 \(\mathfrak K\)의 한 \(\mathfrak h\)-오더라 하자. 즉 \(\mathfrak h\)를 포함하는 \(\mathfrak K\)의 부분환이고 그 몫체는 \(\mathfrak K\)이며, 반드시 독립일 필요는 없는 유한 \(\mathfrak h\)-가군기저를 갖는다. 그러면 \(\mathfrak O\)의 모든 원소는 \(\mathfrak h\)에 관하여 정수적이고, 그 \(\mathfrak h\)-가군기저는 동시에 \(\mathfrak K\)의, 일반적으로는 종속인 \(P\)-가군기저이다. \(\mathfrak o\)를 \(\mathsf{K}\) 안에서 \(\mathfrak O\)와 동형인 오더, \(\mathfrak o^{(i)}\)를 \(\mathsf{K}^{(i)}\) 안의 그 켤레들이라 하자. \emph{\(n\)개의 켤레 오더가 \(\Gamma\) 안에서 생성하는 환을 \(\mathfrak g\)라 하자.} 그러면 \(\mathfrak g\)도 오더이다. 실제로 \(\mathfrak g\)는 \(\mathfrak h\)를 포함한다. \(\mathfrak o^{(i)}\)들의 \(\mathfrak h\)-기저 원소들을 모두 곱하여 얻는 원소들은 \(\mathfrak g\)의 \(\mathfrak h\)-기저이자 \(\Gamma\)의 \(P\)-기저이다.

\emph{직접곱 \(\mathfrak O_{\mathfrak o}\)와 \(\mathfrak O_{\mathfrak g}\)가 존재한다.} 이는 § 2, 3의 충분조건에서 나온다. 거기의 환 \(\mathfrak f\)를 \(\mathsf{P}\)로 바꾸면 된다. 우선 \(\mathfrak K_\Gamma\) 안에서 \(\mathfrak h\)에 관한 직접곱 \(\mathfrak O_{\mathsf{P}}=\mathfrak K\)가 존재한다. 실제로 \(\mathfrak h\)는 교집합 \([\mathfrak O,\mathsf{P}]\)에 들어 있다. 역으로 그 교집합의 모든 원소는 \(\mathfrak O\)의 원소로서 정수적이고, 동시에 \(\mathsf{P}\)의 원소이므로 정수적으로 닫혔다는 가정에 따라 \(\mathfrak h\)에 속한다. 이제 \(\mathfrak O\)에서 취한 임의의 \(n\)개 선형독립 원소를 \(P\)의 계수로 선형결합하면 직접곱으로서 \(\mathfrak K\)를 얻는다. 마찬가지로
\[
\mathfrak o_{\mathsf{P}}=\mathsf{K}
\qquad\text{이고}\qquad
\mathfrak g_{\mathsf{P}}=\Gamma.
\]
\(\mathfrak h\)에 관한 직접곱 \(\mathfrak K_{\mathsf{K}}=\mathfrak O_{\mathsf{P}}\times\mathfrak o_{\mathsf{P}}\)에서 § 2, 3의 조건에 따라 \(\mathfrak h\)에 관한 \(\mathfrak O_{\mathfrak o}\)의 존재가 따른다. 마찬가지로 \(\mathfrak K_\Gamma=\mathfrak O_{\mathsf{P}}\times\mathfrak g_{\mathsf{P}}\)에서 \(\mathfrak O_{\mathfrak g}\)의 존재가 따른다. \(\mathfrak O_{\mathfrak o}\)와 함께 켤레환 \(\mathfrak O_{\mathfrak o^{(i)}}\)들도 존재한다. \(\mathfrak O_{\mathfrak g}\)를 \emph{오더 \(\mathfrak O\)의 갈루아 확대환}이라 한다.

따라서 \(\mathfrak h\)에 관한 \(\mathfrak O\)와 \(\mathfrak o\)의 디퍼런트가 정의되고, 켤레 디퍼런트들도 마찬가지이다. 
\[
\mathfrak P=\{\ldots,y-\eta,\ldots\}
\]
를 \(\mathfrak O_{\mathfrak o}\)의 차분아이디얼,
\[
\mathfrak P^{(i)}=\{\ldots,y-\eta^{(i)},\ldots\}
\]
를 \(\mathfrak O_{\mathfrak o^{(i)}}\)의 그 켤레들이라 하자. 또
\[
\mathfrak A^{(i)}=(0):\mathfrak P^{(i)}
\quad\text{--- 몫은 \(\mathfrak O_{\mathfrak o^{(i)}}\)에서 취한다 ---}
\]
를 켤레 \emph{차분몫}들이라 하고, 특히
\[
\mathfrak A=(0):\mathfrak P\quad\text{ 안에서}\mathfrak O_{\mathfrak o}
\]
라 하자. 나아가
\[
\mathfrak D=\mathfrak A[\eta\to y]
\]
를 \(\mathfrak O\)의 디퍼런트,
\[
\mathfrak d^{(i)}=\mathfrak A^{(i)}[y\to\eta^{(i)}]
\]
를 \(n\)개의 켤레 오더 \(\mathfrak o^{(i)}\)의 디퍼런트들이라 하고, 특히
\[
\mathfrak d=\mathfrak A[y\to\eta]
\]
를 \(\mathfrak o\)의 디퍼런트라 하자.\srcfn{21)}{\(\mathsf{P}\)에서 정수적으로 닫힌 \(\mathfrak h\)라는 가정은 예컨대 \(\mathfrak h\)가 한 수체의 정수환인 경우, 특히 \(\mathfrak h\)가 유리정수들의 계인 경우에 만족한다. 따라서 상대체에서도 임의의 오더들의 디퍼런트와 상대디퍼런트가 정의된다. 또 \(\mathfrak h\)는 \(n\)개 부정원의 다항식영역일 수도 있으며, 이는 \(n\)개 부정원의 대수함수체에서 디퍼런트를 준다.}

\paragraph{2. 구조정리들.}
체의 갈루아 확대환에 관한 구조정리(§ 5, 2)에는 다음이 대응한다.

\paragraph{오더의 갈루아 확대환에 관한 구조정리.}
\emph{오더와 체에 관한 차분아이디얼과 차분몫은 다음 관계에 있다. \(\mathfrak P_{\mathsf{K}}\)는 \(\mathfrak O_{\mathfrak o}\)의 차분아이디얼 \(\mathfrak P\)의 확대이고, 마찬가지로 \(\mathfrak A_{\mathsf{K}}\)는 \(\mathfrak O_{\mathfrak o}\)의 차분몫 \(\mathfrak A\)의 확대이며 \(\mathfrak A\)는 \(\mathfrak A_{\mathsf{K}}\)의 수축이다. 차분몫, 디퍼런트, 단위원 성분 사이에는}
\[
\mathfrak A=\mathfrak d e^{(1)},
\qquad
\mathfrak A^{(i)}=\mathfrak d^{(i)}e^{(i)}
\]
\emph{라는 관계가 있고, 따라서}
\[
\mathfrak d=[\mathfrak o,\mathfrak A^{(1)}+\cdots+\mathfrak A^{(n)}].
\]
\emph{그러므로 \(\mathfrak o\)의 디퍼런트 \(\mathfrak d\)는 \(\varkappa e^{(1)}\)이 \(\mathfrak O_{\mathfrak o}\)에 속하도록 하는 \(K\)의 모든 원소 \(x\)로 이루어지며, 영아이디얼과 다르다.}

\(\mathfrak P_{\mathsf{K}}\)가 \(\mathfrak P\)의 확대라는 것은 \(\mathfrak K=\mathfrak O_P\), \(K=\mathfrak o_P\)인 데서 따른다. 따라서 \(\mathfrak P_{\mathsf{K}}\)의 모든 \(x-\xi\)는 \(\mathfrak P\)의 \(y-\eta\)들을 \(K\)의 계수로 선형결합하여 나타낼 수 있다. 여기서 더 나아가 \(\mathfrak A_{\mathsf{K}}=(0):\mathfrak P_{\mathsf{K}}\)는 \(\mathfrak K_{\mathsf{K}}\)에서 취한 몫 \((0):\mathfrak P\)로도 정의된다. 그러므로 \(\mathfrak O_{\mathfrak o}\) 안의 \(\mathfrak A=(0):\mathfrak P\)는 \(\mathfrak A_{\mathsf{K}}\)에 속하는 \(\mathfrak O_{\mathfrak o}\)의 모든 원소로 이루어진다. 곧 \(\mathfrak A=[\mathfrak O_{\mathfrak o},\mathfrak A_{\mathsf{K}}]\)이고 \(\mathfrak A_{\mathsf{K}}\)의 수축이다. \(\mathfrak A_{\mathsf{K}}=\mathsf{K}e^{(1)}\)이므로 \(\mathfrak A=\mathfrak z e^{(1)}\)이고, \(\mathfrak A\)가 \(\mathfrak o\)-가군이므로 \(\mathfrak z\)는 \(K\) 안의 \(\mathfrak o\)-가군이다. 그런데 디퍼런트의 정의에 따라
\[
\mathfrak d=\mathfrak A[x\to\xi]
=(\mathfrak z e^{(1)})[x\to\xi]
=\mathfrak z\bigl(e^{(1)}[x\to\xi]\bigr)=\mathfrak z e=\mathfrak z
\]
이다(§ 5, 4, (1)을 보라). 따라서 \(\mathfrak z=\mathfrak d\)이다. 차분몫 \(\mathfrak A\)는 \(\mathfrak o\)의 디퍼런트 \(\mathfrak d\)의 \(K\)-성분이고, 켤레들에 대해서도 마찬가지이다. 여기서 § 5, 3의 따름정리에 따라 위의 \(\mathfrak d\)에 관한 관계가 나온다. 그리고
\[
\mathfrak A=\mathfrak d e^{(1)}=[\mathsf{K}e^{(1)},\mathfrak O_{\mathfrak o}]
\]
는 \(\mathfrak d\)가 \(\varkappa e^{(1)}\)이 \(\mathfrak O_{\mathfrak o}\)에 속하도록 하는 \(K\)의 모든 원소 \(x\)로 이루어짐을 보인다. 따라서 \(\mathfrak d\)는 영아이디얼과 다르다. 실제로 \(\mathfrak K_{\mathsf{K}}=\mathfrak O_{\mathsf{K}}\)의 원소 \(e^{(1)}\)은 \(\mathfrak O\)의 유한개 원소들을 \(\mathsf{K}\)의 계수로 선형결합한 것이고, 따라서 \(\mathfrak O_{\mathfrak o}\)의 한 원소를 \(\mathfrak o\), 더 나아가 \(\mathfrak h\)의 한 원소로 나눈 몫이다. 이 분모는 \(\mathfrak d\)의 0 아닌 원소를 준다. 따라서 \(\mathfrak d\)는 \(\mathsf{P}\)에 관하여 랭크 \(n\), 곧 \(\mathfrak d_{\mathsf{P}}=\mathsf{K}\)이고, 그러므로 \(\mathfrak A_{\mathsf{K}}\)는 \(\mathfrak A\)의 확대이다.

\begin{center}
\textbf{여기부터는 개요이다.}
\end{center}

\paragraph{3. 독립기저가 있는 경우 여모듈과의 관계.}
\(t_1,\ldots,t_n\)이 \(\mathfrak O\)의 독립 \(\mathfrak h\)-기저이면 \(\mathfrak K\)의 모든 기저 \((s)=(t)P\)도 \(P\)가 \(\mathfrak h\)에 속하고 단일모듈러이면 그러하다. 그런데 \((S)=P^{-1}(T)\)는 \(T\)뿐 아니라 \(S\)도 같은 하나의 \(\mathfrak h\)-가군의 기저임을 보인다. 이 가군을 \(\mathfrak O\)의 여모듈이라 한다.

이제 2의 구조정리 마지막 부분은 \(\mathfrak d\)가 \(K\) 안의 몫
\[
\mathfrak o:\mathfrak c
\]
과 같다고 말한다. 여기서 \(\mathfrak e\)는 여모듈이다. 실제로
\[
e^{(1)}=\mathsf{A}_1t_1+\cdots+\mathsf{A}_nt_n
\]
이고 \(\mathsf{A}_\lambda\)들이 \(\mathfrak e\)의 한 기저라고 하자. 그러면 \(\mathfrak d\)는 \(\delta e^{(1)}\)이 \(\mathfrak O_{\mathfrak o}\)에 속하도록 하는 모든 원소 \(\delta\), 곧 \(\delta \mathsf{A}_\lambda\)가 \(\mathfrak o\)에 속하도록 하는 모든 \(\delta\)로 이루어진다. 따라서 \(\mathfrak d=\mathfrak o:\mathfrak e\)이다.

\paragraph{특수한 경우.}
\(e,z,\ldots,z^{n-1}\)을 가군기저로 갖는 정칙 오더:
\[
\mathfrak d=\{f'(z)\}.
\]

\paragraph{주.}
§ 5, 4의 주에서 \(\mathfrak O\)의 여모듈 \(\mathfrak E\)에 대한 또 하나의 알려진 정의가 나온다. \emph{\(\mathfrak E\)는 \(\mathfrak O\)의 임의의 \(o\)에 대하여 대각합 \(\operatorname{Sp}(co)\), 또는 같은 뜻으로 대각합 \(\operatorname{Sp}(ct)\)가 정수적, 곧 \(\mathfrak h\)에 속하도록 하는 \(\mathfrak K\)의 모든 원소 \(c\)로 이루어진다.} \(\mathfrak o\)와 \(\mathsf{K}\)에 관한 \(\mathfrak e\)에도 이에 대응하는 것이 성립한다.

실제로 \(c=c_1T_1+\cdots+c_nT_n\)이고 \(c_i\)들이 \(\mathfrak h\)에 속하는 것이, \(t\)의 여차기저로 \(\mathfrak E\)의 한 원소 \(c\)를 나타낸 표현이라 하자. 그러면 \(c_\lambda=\operatorname{Sp}(ct_\lambda)\)이므로 \(\mathfrak h\)에 속한다. 역으로 \(\operatorname{Sp}(ct_\lambda)\)가 \(\mathfrak h\)에 속하면 \(c_\lambda\)가 \(\mathfrak h\)에 속하고, 따라서 \(c\)가 \(\mathfrak E\)에 속한다.

\paragraph{보충.}
\emph{오더 안에 독립 \(\mathfrak h\)-가군기저가 존재하면, 차분몫 \(\mathfrak A^{(i)}\)는 \(\mathfrak P^{(i)}\)와 다른 모든 차분아이디얼의 교집합이다. 더 정확히는}
\[
\mathfrak A^{(i)}=
[\mathfrak O_{\mathfrak o^{(i)}},
\mathfrak P_{\mathfrak g}^{(1)},\ldots,
\mathfrak P_{\mathfrak g}^{(i-1)},
\mathfrak P_{\mathfrak g}^{(i+1)},\ldots,
\mathfrak P_{\mathfrak g}^{(n)}].
\]
실제로 § 2, 2의 강화된 형태에 따라 이때 \(\mathfrak P\)도 \(\mathfrak P_{\mathsf{K}}\)의 수축아이디얼이다. \(e,t_2,\ldots,t_n\)이 \(\mathfrak O\), 따라서 \(\mathfrak O_{\mathfrak o}\)의 가군기저이면
\[
e,t_2-\alpha_2,\ldots,t_n-\alpha_n
\]
도 그러하다. 그런데 \(t_2-\alpha_2,\ldots,t_n-\alpha_n\)은 동시에 \(\mathfrak P_{\mathsf{K}}\)의 \(\mathsf{K}\)-가군기저이다. 따라서 § 2, 2의 보충에 있는 가정들이 만족된다. 한편
\[
\mathfrak A_\Gamma=[\mathfrak P_\Gamma^{(2)},\ldots,\mathfrak P_\Gamma^{(n)}];
\qquad
\mathfrak A_\mathfrak g=[\mathfrak A_\Gamma,\mathfrak O_\mathfrak g]
=[\ldots,[\mathfrak P_\Gamma^{(\lambda)},\mathfrak O_\mathfrak g],\ldots]
=[\mathfrak P_\mathfrak g^{(2)},\ldots,\mathfrak P_\mathfrak g^{(n)}]
\]
였다. 그런데 \(\mathfrak A=[\mathfrak A_\mathfrak g,\mathfrak O_{\mathfrak o}]\)이므로
\[
\mathfrak A=[\mathfrak O_{\mathfrak o},\mathfrak P_\mathfrak g^{(2)},\ldots,\mathfrak P_\mathfrak g^{(n)}].
\]

\paragraph{4. 기본방정식이 존재하는 경우 그 미분몫과의 관계.}
오더 \(\mathfrak O\)에 \(n\)개 부정원 \(u_1,\ldots,u_n\)을 첨가하여 확대한다. 곧 \(\mathfrak U\)가 다항식영역 \(\mathfrak h[u_1,\ldots,u_n]\)을 뜻할 때 직접곱 \(\mathfrak O_{\mathfrak U}=\mathfrak O[u_1,\ldots,u_n]\)을 만든다. 그러면 \(\mathfrak O\)의 모든 아이디얼 \(\mathfrak C\)의 확대 \(\mathfrak C_{\mathfrak U}\)는 \(\mathfrak C\)의 계수들을 갖는 \(u\)의 모든 다항식으로 이루어지고, 역으로 \(\mathfrak C\)는 \(\mathfrak C_{\mathfrak U}\)의 모든 계수의 계로 특징지어진다.

이제 기본방정식이 존재한다고 하자. 즉 \(\mathfrak h[u]\)의 하나 이상의 원시다항식 \(H(u)\)들에 따른 몫환으로 넘어갈 때 \(e,U,\ldots,U^{n-1}\)이 \(\mathfrak h\)-가군기저를 이룬다고 하자. 그러면 아이디얼 \(\mathfrak C_{\mathfrak U}\)들은 \(\mathfrak C_{\mathfrak U}\)에 단위원들, 곧 \(H(u)\)와 그 음·양의 거듭제곱들을 곱하여 얻는 계들로 넘어간다. 따라서 \(\mathfrak C\)를 되찾으려면 먼저 이 단위원들을 곱하여 \(u\)에 관하여 정수적인 \(\mathfrak C_{\mathfrak U}\)에 이른 다음 그 계수들에서 \(\mathfrak C\)로 갈 수 있다. 특히 \(\mathfrak O\)의 디퍼런트 \(\mathfrak D\)가 \(\mathfrak D_{\mathfrak U}\), 곧 \(\mathfrak O_{\mathfrak U}\)의 디퍼런트로 간다고 하자. 그런데 \(\mathfrak O_{\mathfrak U}\)에는 원시다항식 \(H(u)\)들로 몫을 취함으로써 정의방정식 \(G(U)=0\)이 존재한다. 따라서 \(G'(U)\)는 디퍼런트의 기저다항식이고, 곧 디퍼런트는 \(G'(U)\)의 계수들에서 유도된다. (\((G'(U))\)가 실제로 \(\mathfrak D_{\mathfrak U}\)만을 포함하고 더 큰 것을 포함하지 않음을 더 정확히 보여야 한다. 아마도 \(\mathfrak o\)가 \(\mathsf{K}\)에서 정수적으로 닫힌 경우에만 성립할 것이다. 실제로 \(H(u)\)의 인자들을 곱하면 \(\mathfrak O_{\mathfrak U}\)에 놓인 더 많은 원소들이 \(\mathfrak D_{\mathfrak U}\)로 생길 수 있다. 정수적으로 닫힌 경우(다섯 공리)에는 계수아이디얼들, 곧 내용들이 서로 곱해진다는 정리에서 이 사실이 따른다. 이에 따라 \(p\)를 법으로 한 기본방정식을 고찰하여 분기이론을 세울 수 있다.)

\subsection*{§ 7. \(p^t\)를 법으로 한 정수환의 분기이론}

한 수체를 바탕으로 삼고 \(\mathfrak h\)를 그 정수환이라 하자. 또는 정수환을 한 소아이디얼로 나눈 몫환에 관한 상대체를 바탕으로 삼아도 된다. 이때 \(p\)는 이 몫환의 소아이디얼의 기저원소를 뜻한다. 부정원 \(u_1,\ldots,u_n\)을 첨가했다고 하자. 그러면 정수환과 \(p^t\)에 따른 잉여류환은 모두 \(\mathfrak h\)-가군기저 \(e,U,\ldots,U^{n-1}\)을 갖는다. 따라서 두 경우 모두 \(G'(u)\)가 디퍼런트를 준다. \emph{정수환의 디퍼런트는 \(p^t\)를 법으로 하면 \(p^t\)에 따른 잉여류환의 디퍼런트로 간다.} (임의의 오더에도 이것이 성립하는가?) 이제
\[
p=\mathfrak p_1^{\varrho_1}\cdots\mathfrak p_r^{\varrho_r}
\]
이면 \(p^t\)에 따른 잉여류환 \(\mathfrak R\)은 직접합이고, 그 성분들은 각각 \(\mathfrak p_i^{t\varrho_i}\)에 따른 잉여류환과 동형이다. 곧
\[
\mathfrak R=\mathfrak R e_1+\cdots+\mathfrak R e_r.
\]
이때 \(i\ne k\)이면 \(e_k\)는 \(\mathfrak p_i^{t\varrho_i}\)로 나누어진다. 따라서
\[
\mathfrak d[p^t]=\mathfrak d[p^t]e_1+\cdots+\mathfrak d[p^t]e_r
=\mathfrak d_1+\cdots+\mathfrak d_r
\]
이고 \(t\)가 충분히 크면, \(\mathfrak p_i\)는 \(\mathfrak d_i\)에서 \(\mathfrak d[p^t]\), 따라서 \(\mathfrak d\)에서와 같은 거듭제곱으로 나타난다. 그런데 § 4에 따라 \(\mathfrak d_i\)는 \(\mathfrak R_i\)의 디퍼런트이다. 그러므로 \(\mathfrak p^{\varrho t}\)에 따른 잉여류환의 디퍼런트만 고찰하면 충분하다. \(\xi\)를 (Zahlbericht, 정리 29와 30) 정수이도록 택하여
\[
\varphi(\xi)\equiv0(\mathfrak p),
\quad \not\equiv0(\mathfrak p^2)
\]
이라 하자. 여기서 \(\varphi(x)\)는 \(p\)를 법으로 한 차수 \(f\)의 소다항식이다. \(t\geqq2\)일 때 이 \(\xi\)를 \(\xi^*=\xi e_1\)로 바꾸자. 여기서 \(e_1\)은 \(p^t\)를 법으로 한 잉여류환 \(\mathfrak R\)의 분해에 나타나는 첫 번째 멱등원이다. 그러면 \(\varphi(\xi^*)\equiv0(\mathfrak p)\), \(\not\equiv0(\mathfrak p^2)\)일 뿐 아니라 \(\mathfrak p_1\equiv(p,\varphi(\xi^*))\)이다. \(p^t\)를 법으로 한 정수들에 관한 한 기저는 \(\mu=0,1,\ldots,f-1\), \(\nu=0,1,\ldots,t-1\)인 \(\xi^\mu\varphi(\xi)^\nu\)들로 이루어진다. 독립 \(\mathfrak h\)-기저(\(p^t\)에 따른 잉여류들)는
\begin{equation}
\xi^{\mu_0}\varphi(\xi)^{\nu_0}
\quad\text{이며}\quad
\mu_0=0,1,\ldots,f-1;
\quad \nu_0=0,1,\ldots,\varrho-1. \tag{1}
\end{equation}
\((\varphi(\xi),\mathfrak p^{\varrho t})=\mathfrak p\)이기 때문이다. 따라서 잉여류환에서 \((p)=\mathfrak p^\varrho\)이므로 주아이디얼에서 원소로 넘어가면 \(\varphi(\xi)=pM(\xi)\)이고, 여기서 \(M(\xi)\)는 잉여류환의 단위원, 곧 \(M(\xi)\not\equiv0(\mathfrak p)\)이다. 그러므로 \(p^t=0\)을 이용하여 \(\xi^{\varrho t}\)을 차례로 더 낮은 거듭제곱들로 나타낼 수 있다. 먼저 \(M(\xi)\)에서, 이어 잉여류환에서 그렇게 할 수 있으므로 (1)은 실제로 기저이다. 이 기저에 관한 정의방정식은
\[
F(x)=\varphi(x)^\varrho+pM(x)
\]
이다. 여기서 \(M(x)\)는 \(p\)를 법으로 하여 \(\varphi(x)\)로 나누어지지 않는다(Ö. Ore, Math. Ann. 96, 348쪽). 따라서
\[
M(\xi)\not\equiv0(\mathfrak p).
\]
이 기저는 또한 독립이다. 실제로
\[
a_0(\xi)+a_1(\xi)\varphi(\xi)+\cdots+a_{\varrho-1}(\xi)\varphi(\xi)^{\varrho-1}
\equiv0(\mathfrak p^{t\varrho})
\]
이면 먼저 \(\equiv0(\mathfrak p)\)이고, 따라서 \(a_0(\xi)\equiv0(\mathfrak p)\), 곧
\[
a_0(x)\equiv0(p)
\]
이다. \(\mathfrak p\)를 차례로 약분하면
\[
a_0(x)\equiv0(p),\ldots,a_{\varrho-1}(x)\equiv0(p)
\]
을 얻고, 다시 \(p\)를 약분하면
\[
a_0(x)\equiv0(p^2),\ldots,a_{\varrho-1}(x)\equiv0(p^2)
\]
를 얻으며, 같은 과정을 계속하여 마침내
\[
a_0(x)\equiv0(p^t),\ldots,a_{\varrho-1}(x)\equiv0(p^t)
\]
를 얻는다. 따라서 잉여류환의 디퍼런트는 다항식 미분몫 \(F'(\xi)\)로 주어진다. 모든 것이 Ore, 345쪽, § 3으로 환원되었다. Ore의 정리 10(Dedekind--Hensel도 비교하라)에서 모든 \(\mathfrak p\)에 \(p^t=p^n\)이면 충분하다는 것이 따르고, 더 나아가 Ore의 결과들에서 보충수들이 나온다.

몫환으로 넘어가면 이 모든 것은 상대디퍼런트에도 동시에 성립한다. 이에 따라 Ore, Math. Ann. 97, 594쪽, § 4는 Ore, Math. Ann. 96에서 곧바로 나온다. 상위체의 디퍼런트를 하위체의 디퍼런트와 상대디퍼런트의 곱으로 나타낼 수 있다는 보충만 더해진다(이 논문의 서론을 보라).

\begin{center}
1949년 10월 25일 접수.
\end{center}

%
%

\clearpage
\setcounter{footnote}{0}
\phantomsection
\addcontentsline{toc}{section}{초복소량의 대수학}
%
%
%
%
%

\begin{center}
\vspace*{3.5em}
{\LARGE\bfseries 초복소량의 대수학\par}
\vspace{3.5em}
{\large 에미 뇌터 교수의 강의\par}
{\large 1929/30년 겨울학기\par}
\vspace{1.5em}
{\large M. 도이링 교수 정리\par}
\vspace{4em}
{\Large\bfseries 목차\par}
\end{center}
\vspace{1em}
\begingroup
\small
\setlength{\parindent}{0pt}
\newcommand{\tocline}[2]{\noindent #1\dotfill #2\par}
\newcommand{\tocsec}[3]{\noindent\hspace*{1.2em}\S\,#1.\quad #2\dotfill #3\par}
\tocline{\textbf{서론}}{3}
\medskip
\tocline{\textbf{제I장: 표현}}{3}
\tocsec{1}{직접 표현의 정의 --- 반대 표현의 정의}{3}
\tocsec{2}{표현류}{4}
\tocsec{3}{표현가군}{4}
\tocsec{4}{표현가군과 표현의 관계}{5}
\medskip
\tocline{\textbf{제II장: 가환체의 갈루아 이론}}{7}
\tocsec{5}{초복소계의 계수 확대}{7}
\tocsec{6}{가환계의 기약표현}{8}
\tocsec{7}{한 체의 동형들}{8}
\tocsec{8}{분해체}{9}
\tocsec{9}{\(\mathsf{Z}_1\)에서 \(\mathsf{Z}_i\)로 가는 동형들}{9}
\tocsec{10}{갈루아군}{9}
\tocsec{11}{갈루아 이론의 주정리}{9}
\tocsec{12}{성분 \(e_i\)의 형식적 의미}{11}
\tocsec{13}{일반 연장정리}{12}
\medskip
\tocline{\textbf{제III장: 아벨군}}{13}
\tocsec{14}{군환}{13}
\tocsec{15}{아벨군의 군환}{14}
\tocsec{16}{지표관계}{15}
\tocsec{17}{아벨군의 갈루아 이론}{15}
\medskip
\tocline{\textbf{제IV장: 양쪽단순환}}{18}
\tocsec{18}{보조정리}{18}
\tocsec{19}{양쪽단순 초복소계의 계수체의 비가환 확대체 안에서의 표현}{19}
\tocsec{20}{비가환체}{22}
\tocsec{21}{비가환체의 갈루아 이론}{26}
\clearpage
\tocline{\textbf{제V장: 인자계}}{29}
\tocsec{22}{중심이 주어진 체들의 군}{29}
\tocsec{23}{인자계}{30}
\tocsec{24}{인자계의 곱셈}{33}
\tocsec{25}{갈루아 극대 가환부분체를 갖는 \(\mathfrak K_r\)의 정규표현}{39}
\tocsec{26}{교차표현의 곱셈}{43}
\medskip
\tocline{\textbf{제VI장: 교차곱 이론}}{44}
\tocsec{27}{\(\mathfrak K_r\)의 교차곱 표현}{44}
\tocsec{28}{인자계의 곱정리}{47}
\tocsec{29}{최소에서의 주종정리}{50}
\tocsec{30}{순환 분해체에의 특수화}{50}
\tocsec{31}{순환 특수 경우의 응용}{52}
\endgroup
\clearpage

\section*{서론}

이 강의는 한편으로 유한군의 표현론에서, 다른 한편으로 초복소수계의 대수적ㆍ산술적 이론에서 성장해 나온 일반 환표현론을 다룬다. 표현론을 이처럼 일반적으로 파악하면 종래의 이론(프로베니우스, 슈어)에 비해 다음과 같은 장점이 있다.

군 자체가 아니라 군환, 더 일반적으로는 초복소계나 임의의 환을 고찰하므로 환론과 아이디얼론을 쓸 수 있고, 이에 따라 복잡하고 불투명한 많은 계산에서 이론을 해방할 수 있다. 일반 환론과 아이디얼론을 쓰면 이론이 단순해질 뿐 아니라, 예컨대 한 군과 그 부분군들의 관계라는 문제에서 더 멀리 나아가고 새로운 응용영역도 열린다. 일반 표현론의 도움으로 갈루아 이론에 새롭고 매우 우아한 기초를 줄 수 있으며, 어쩌면 더 중요하게도 비가환체의 갈루아 이론을 세울 수 있다.

군론ㆍ환론ㆍ아이디얼론의 기본개념과 표현론의 일반정리는 뇌터 양의 논문 \emph{Hyperkomplexe Größen und Darstellungstheorie}, Mathematische Zeitschrift 30, 641쪽에 있다. 이하 M. Z.로 인용하는 이 논문은 이 강의록을 보완하는 것으로 보아야 한다. 뇌터 양의 1927/28년 겨울학기 강의 \emph{Hyperkomplexe Größen und Gruppencharaktere}를 정리한 글에서도 M. Z.와 같은 내용을 찾을 수 있다.

여기서는 직접 표현가군에 비해 형식적으로만 새로운 반대 표현가군과 반대 표현을 도입하기 위해 표현과 표현가군의 개념을 짧게 설명한다.

\section*{제I장. 표현}

\(\mathfrak{o}\)를 표현될 환, \(\mathsf T\)를 표현이 이루어지는 다른 환이라 하자.

\subsection*{§ 1. 직접 표현의 정의}

\(\mathfrak{o}\)의 \(\mathsf T\) 안에서의 \(n\)차 \emph{직접 표현}이란, \(\mathsf T\)의 원소로 된 \(n\)행 정방행렬환 \(\mathsf T_n\)의 한 부분환 \(\mathfrak D\)로 \(\mathfrak{o}\)가 환준동형으로 사상되는 것을 말한다. 기호로는
\[
        \mathfrak{o}\longrightarrow\mathfrak D\subseteqq\mathsf T_n
\]
이다.

\subsection*{반대 표현의 정의}

\(\mathfrak{o}\)의 \(\mathsf T^*\) 안에서의 \(n\)차 \emph{반대 표현}이란, \(\mathsf T^*\)의 원소로 된 \(n\)행 정방행렬환 \(\mathsf T_n^*\)의 한 부분환 \(\mathfrak D^*\)로 \(\mathfrak{o}\)가 반준동형으로 사상되는 것을 말한다. 반준동형이라는 말은 \(\mathfrak{o}\)의 원소 \(c,d\)에 \(\mathfrak D^*\)의 원소 \(c^*,d^*\)가 대응할 때, \(\mathfrak{o}\)의 원소 \(cd\)에는 \(\mathfrak D^*\)의 원소 \(d^*\cdot c^*\)가, \(c+d\)에는 \(c^*+d^*\)가 대응한다는 뜻이다. 기호로는
\[
        \mathfrak{o}\nrightarrow\mathfrak D^*
        \subseteqq\mathsf T_n^*
\]
이다.

반준동형의 예를 들자. 가환체의 \(n\)행 정방행렬에 그 전치행렬을 대응시키면 이 대응은 반대 환동형이다.

일반적으로 임의의 환 \(\mathsf T\)에 대해 그와 반대동형인 환 \(\mathsf T^*\)를 다음과 같이 구성할 수 있다.

\(\mathsf T\)의 각 원소 \(c\)에 기호 \(c^*\)를 대응시키고, 이 기호들의 덧셈과 곱셈을
\[
        c^*+d^*=(c+d)^*,
        \qquad
        c^*\cdot d^*=(dc)^*
\]
로 정의한다. 기호 \(c^*\)들의 집합은 \(\mathsf T\)와 반대동형인 환 \(\mathsf T^*\)이다.

\subsection*{§ 2. 표현류}

\(\mathfrak{o}\)의 \(\mathsf T\) 안에서의 직접 표현 \(\mathfrak D\)를 한 정칙행렬 \(P\)로 원소별 변환하면 동형인 표현 \(\mathfrak D'\)로 바꿀 수 있다. 곧 \(\mathfrak{o}\)의 원소 \(c\)에 \(C\subseteqq\mathfrak D\) 대신 \(P^{-1}CP\subseteqq P^{-1}\mathfrak D P\)를 대응시킨다.

이 방식으로 서로 변환될 수 있는 두 표현을 \emph{동치}라 한다.

한 표현과 동치인 모든 표현은 하나의 \emph{표현류}(직접 표현류)를 이룬다.

반대 표현의 동치도 똑같이 정의한다. 한 반대 표현을 정칙행렬로 변환하여 다른 하나를 얻을 수 있으면 두 반대 표현은 동치이다. 고정된 한 반대 표현과 동치인 모든 표현은 하나의 \emph{반대 표현류}를 이룬다.

\subsection*{§ 3. 표현가군}

\emph{직접 표현가군}의 정의.

\(\mathfrak{o}\)의 \(\mathsf T\)에 관한 직접 표현가군이란 다음 조건을 만족하는 모든 \(\mathfrak{o},\mathsf T\)-쌍가군 \(\mathfrak M\)을 말한다. 곧 \(\mathfrak M\)은 덧셈으로 쓴 아벨군이며, \(\mathfrak{o}\)에 관해서는 좌측가군이고 \(\mathsf T\)에 관해서는 우측가군이며, 모든 \(c\in\mathfrak{o}\), \(m\in\mathfrak M\), \(\tau\in\mathsf T\)에 대해 결합법칙
\[
        (cm)\tau=c(m\tau)
\]
을 만족한다. 또한 \(\mathsf T\)-가군으로서는 유한개의 한 원소 가군의 직접합으로
\[
        \mathfrak M=x_1\mathsf T+x_2\mathsf T+\cdots+x_n\mathsf T
\]
라 쓸 수 있고, \(x_i\tau=0\)이면 \(\tau=0\)이어야 한다.

\emph{반대 표현가군}의 정의.

\(\mathfrak{o}\)의 \(\mathsf T^*\)에 관한 반대 표현가군이란 \(\mathfrak{o}\)-좌측가군이자 \(\mathsf T^*\)-좌측가군인 모든 \(\mathfrak M^*\)로서, 모든 \(c\in\mathfrak{o}\), \(m\in\mathfrak M^*\), \(\tau\in\mathsf T^*\)에 대해 교환법칙
\[
        c(\tau^*m^*)=\tau^*(cm)
\]
을 만족하고 유한개의 단순 \(\mathsf T^*\)-가군의 직접합
\[
        \mathfrak M^*
        =\mathsf T^*x_1+\cdots+\mathsf T^*x_n
\]
으로 쓸 수 있는 것을 말한다. 여기서 \(\tau^*x_i^*=0\)이면 \(\tau^*=0\)이다.

\subsection*{§ 4. 표현가군과 표현의 관계}

\textbf{1.} \emph{모든 표현가군은 하나의 표현류를 유일한 방식으로 낳는다. 직접 표현가군은 직접 표현류를, 반대 표현가군은 반대 표현류를 낳는다.}

그 방식은 다음과 같다.

직접 표현가군 \(\mathfrak M\)에 \(\mathfrak{o}\)의 원소 \(c\)를 곱하는 것은 \(\mathfrak M\)에서 자기 자신으로 가는 준동형이다. 결합법칙 \((cm)\tau=c(m\tau)\) 때문에 이 준동형은 \(\mathsf T\)의 원소들을 연산자로 허용한다. 따라서 이는 \(\mathfrak M\)에서 자기 자신으로 가는 \emph{선형변환}이다. 곧 대응 \(m\mapsto cm=m'\)은 다음 성질을 갖는다.
\[
\text{\(y_i\mapsto y_i'(=cy_i)\)이면 }
\sum y_i\tau_i\mapsto
\sum y_i'\tau_i\left(=c\sum y_i\tau_i\right).
\]

두 선형변환
\[
\mathfrak T_1:\ m\mapsto m'(=c_1m),
\qquad
\mathfrak T_2:\ m\mapsto m''(=c_2m)
\]
의 합 \(\mathfrak T_1+\mathfrak T_2\)을 변환 \(m\mapsto m'+m''=(c_1+c_2)m\)으로 정의하고, 그 곱 \(\mathfrak T_1\cdot\mathfrak T_2\)를 합성변환 \(m\mapsto(m'')'=c_1c_2m\)으로 정의하자. 그러면 \(\mathfrak{o}\)의 원소들이 생성하는 서로 다른 선형변환들의 집합은 \(\mathfrak{o}\)의 환준동형상 \(\overline{\mathfrak D}\)이다.

이제 \(\mathfrak M\)의 임의의 기저 \(x_1,\ldots,x_n\)을 택하자. \(c\)가 생성하는 \(\mathfrak M\)의 선형변환은 기저원소 \(x_i\)들이 어떻게 변환되는지를 알면 정해진다. 곧 연립방정식
\[
        x_j'=cx_j=\sum_i x_i\gamma_{ij}
\]
의 행렬 \((\gamma_{ij})=C\)를 알면 충분하다. 다른 원소 \(x=\sum x_j\tau_j\)에 대해서는 단순히 \(x'=\sum x'_j\tau_j\)이기 때문이다. 이 행렬 \(C\)들은 서로 다른 선형변환들과 일대일로 대응하고 서로를 결정하며, 곧바로 알 수 있듯 \(\overline{\mathfrak D}\)와 동형인 환을 이룬다. 따라서 이는 \(\mathfrak{o}\)의 환준동형상, 곧 \(\mathfrak{o}\)의 \(n\)차 표현이다.

\(\overline{\mathfrak D}\)와 표현 사이의 동형은 다음 식들에 들어 있다. \(c\in\mathfrak{o}\)가 생성하는 선형변환에 행렬 \(C\)가 대응하여
\[
        c(x_1,\ldots,x_n)=(x_1,\ldots,x_n)C
\]
이고, 같은 방식으로 \(d\in\mathfrak{o}\)와 \(D\in\mathsf T_n\)가 대응하여
\[
        d(x_1,\ldots,x_n)=(x_1,\ldots,x_n)D
\]
이면
\[
\begin{aligned}
(c+d)(x_1,\ldots,x_n)
&=(cx_1+dx_1,\ldots,cx_n+dx_n)\\
&=(x_1,\ldots,x_n)(C+D)
\end{aligned}
\]
이고
\[
\begin{aligned}
cd(x_1,\ldots,x_n)
&=c(dx_1,\ldots,dx_n)
 =c\bigl((x_1,\ldots,x_n)D\bigr)\\
&=(cx_1,\ldots,cx_n)D
 =(x_1,\ldots,x_n)CD.
\end{aligned}
\]

그러므로 \(\mathfrak{o}\)가 생성하는 \(\mathfrak M\)의 자기 선형변환들을 \(\mathfrak M\)의 특정한 기저 \(x_\nu\)로 실현하면 \(\mathfrak{o}\)의 \(\mathsf T\) 안에서의 \(n\)차 표현을 얻는다.

\((y_1,\ldots,y_n)\)이 \(\mathfrak M\)의 다른 기저이면, 이는 정칙행렬 \(P\)를 곱하여 \(x_1,\ldots,x_n\)에서 얻어진다.
\[
        (y_1,\ldots,y_n)=(x_1,\ldots,x_n)P.
\]
한 선형변환이 기저 \(x_\nu\)에서 행렬 \(C\)로 실현되면 기저 \(y_\nu\)에서는 \(P^{-1}CP\)로 실현된다. 실제로
\[
        c(x_1,\ldots,x_n)=(x_1,\ldots,x_n)C
\]
이면
\[
\begin{aligned}
c(y_1,\ldots,y_n)
&=c(x_1,\ldots,x_n)P
 =\bigl((x_1,\ldots,x_n)C\bigr)P\\
&=(x_1,\ldots,x_n)CP
 =\bigl((y_1,\ldots,y_n)P^{-1}\bigr)CP\\
&=(y_1,\ldots,y_n)P^{-1}CP.
\end{aligned}
\]

여기서 곧바로 다음이 따른다.

\emph{\(\mathfrak M\)의 모든 기저를 써서 \(\mathfrak{o}\)가 유도하는 \(\mathfrak M\)의 선형변환들을 실현하면, \(\mathfrak M\)은 한 표현류의 모든 표현을 준다.}

\textbf{2.} \emph{모든 표현류는 1에서 설명한 방식으로 한 표현가군에 의해 생성된다.} 다시 직접 표현에 국한하자.

\emph{증명.} 그 류의 한 특정한 표현을 택하고 준동형 대응 \(c\mapsto C\)로 주어졌다고 하자. 표현의 차수가 \(n\)이면 \(n\)개의 부정원 \(x_1,\ldots,x_n\)으로 \(\mathsf T\)-우측가군
\[
        \mathfrak M=x_1\mathsf T+\cdots+x_n\mathsf T
\]
를 만든다. 통상적인 계산법칙
\[
\sum x_i\tau_i+\sum x_i\bar\tau_i
 =\sum x_i(\tau_i+\bar\tau_i),
\]
\[
\sum x_i\tau_i=\sum x_i\bar\tau_i
\quad\Longleftrightarrow\quad
\tau_i=\bar\tau_i\ \text{(모든 \(i\))},
\qquad
\left(\sum x_i\tau_i\right)\varrho
=\sum x_i\tau_i\varrho
\]
을 쓴다. 이제 \(C\)가 \(\mathfrak{o}\)의 원소 \(c\)에 대응하는 표현행렬이고
\[
        (x_1',\ldots,x_n')=(x_1,\ldots,x_n)C
\]
일 때
\[
        c\sum x_i\tau_i=\sum x_i'\tau_i
\]
로 정하여 \(\mathfrak M\)을 \(\mathfrak{o}\)-좌측가군으로 만든다.

이 규정으로 \(\mathfrak M\)은 \(\mathfrak{o},\mathsf T\)-쌍가군이 된다. 실제로 첫째, 준동형관계
\[
        c+d\longmapsto C+D,
        \qquad
        cd\longmapsto CD
\]
에서
\[
        (c+d)x_i=cx_i+dx_i
\]
와
\[
\begin{aligned}
cd\cdot x_i&=c\cdot dx_i,\\
c\left(\sum x_i\tau_i+\sum x_i\bar\tau_i\right)
&=c\sum x_i\tau_i+c\sum x_i\bar\tau_i
\end{aligned}
\]
가 따르며, 이것들이 좌측가군 성질이다. 둘째,
\[
        c\left(\sum x_i\tau_i\varrho\right)
        =\left(c\sum x_i\tau_i\right)\cdot\varrho
\]
임은 명백하며, 이것이 양쪽을 관통하는 결합법칙이다.

따라서 \(\mathfrak M\)은 표현가군이고, 기저 \(x_1,\ldots,x_n\)을 써서 그 선형변환들을 실현하면 1에서 설명한 방식으로 처음의 표현을 생성한다. 다른 기저들은 그 류의 나머지 표현들을 준다.

반대 표현과 반대 표현가군의 경우에도 모든 것이 똑같다. 가군
\[
        \mathfrak M^*
        =\mathsf T^*x_1^*+\dots+\mathsf T^*x_n^*
\]
은 \(\mathfrak{o}\)의 원소 \(c\)에 식
\[
c\begin{pmatrix}x_1\\ \vdots\\ x_n\end{pmatrix}
=\begin{pmatrix}cx_1\\ \vdots\\ cx_n\end{pmatrix}
=C\begin{pmatrix}x_1\\ \vdots\\ x_n\end{pmatrix}
\]
의 행렬 \(C\)를 대응시켜 \(\mathfrak{o}\)의 \(\mathsf T^*\) 안에서의 반대 표현을 생성한다. 서로 다른 기저들은 다시 한 류의 서로 다른 표현들을 준다.

\section*{제II장. 가환체의 갈루아 이론}

지금까지 전개한 표현론만으로도 갈루아 이론에 기초를 줄 수 있다. 이 방법은 통상적인 방법과 달리 특별한 생성원계나 부정원의 다항식에 관한 정리에 의존할 필요가 없다는 장점이 있다.\srcfn{1}{이하에서는 M. Z. §§ 15--20과 비교하라. M. Z. § 21에는 뒤따르는 갈루아 이론 부분의 첫 정리들이 들어 있다. 초복소계의 정의는 M. Z. § 6을 보라.}

\subsection*{§ 5. 초복소계의 계수 확대}

거의 자명하므로 증명은 생략하고 다음 일반정리를 정식화한다. \(\mathfrak{o}=c_1\mathsf P+\cdots+c_n\mathsf P\)를 가환체 \(\mathsf P\) 위의 초복소계라 하자. \(\Omega\)를 중심에 \(\mathsf P\)를 포함하는 \(\mathsf P\)의 확대환이라 하자. \emph{확대계} \(\mathfrak{o}_\Omega\)를 \(\omega_i\in\Omega\)인 모든 형식적 합 \(\sum c_i\omega_i\)의 집합으로 정의한다. 먼저
\[
\begin{array}{l}
\displaystyle
\sum c_i\omega_i=\sum c_i\omega_i'
\quad\Longleftrightarrow\quad
\omega_i=\omega_i'\ \text{(모든 \(i\))},\\[0.4em]
\displaystyle
\sum c_i\omega_i+\sum c_i\omega_i'
=\sum c_i(\omega_i+\omega_i'),\\[0.4em]
\displaystyle
\omega\sum c_i\omega_i=\sum c_i\omega\omega_i,
\qquad
\sum c_i\omega_i\cdot\omega=\sum c_i\omega_i\omega,\\[0.4em]
\displaystyle
\sum c_i\omega_i=\sum\omega_i c_i
\end{array}
\]
로 정하여 이를 \(\Omega\)와 가환적으로 결합된 \(\Omega\)-가군으로 만든다. 더 나아가
\[
        \sum c_i\omega_i\cdot\sum c_i\omega_i
        =\sum_i\sum_jc_ic_j\omega_i\omega_j
\]
로 정하여 \(\mathfrak{o}\)의 확대환으로 만든다. 여기서 \(c_ic_j\)는 \(\mathfrak{o}\)의 원소로서 이미 정의되어 있다. \(\mathfrak{o}_\Omega\)는 \(\Omega\) 위에서 랭크 \(n\)이고
\[
        \mathfrak{o}_\Omega
        =c_1\Omega+\cdots+c_n\Omega
\]
로 쓸 수 있다.

\emph{\(\mathfrak{o}_\Omega\)의 정의는 기저 \(c_1,\ldots,c_n\)에 의존하지 않는다. \(\Omega\)가 가환체이면 \(\mathfrak{o}_\Omega\)는 \(\Omega\) 위의 초복소계이다.}

\subsection*{§ 6. 가환계의 기약표현}

\(\mathfrak Z=z_1\mathsf P+\cdots+z_n\mathsf P\)를 \(\mathsf P\) 위의 가환계라 하자. \(\mathsf P\)의 대수적으로 닫힌 확대체 \(\Omega\) 안에서 \(\mathfrak Z\)의 모든 기약표현을 구하고자 한다. \(\mathfrak Z\)의 \(\Omega\) 안에서의 모든 표현은 \(\mathfrak Z_\Omega=z_1\Omega+\cdots+z_n\Omega\)의 표현을 낳는다. 이미 \(\mathsf P\)에 들어 있는 기저원소들의 표현을 알면 전체 표현을 세우기에 충분하기 때문이다. 역으로 \(\mathfrak Z_\Omega\)의 모든 표현은 \(\mathfrak Z\)의 표현을 포함하므로 \(\mathfrak Z_\Omega\)의 표현을 물어도 된다. M. Z. § 20에 따라 이 표현들은 \(\mathfrak Z_\Omega\)의 정칙표현, 곧 \(\mathfrak Z_\Omega\) 자체를 표현가군으로 써서 얻는 표현에 들어 있다. 더 나아가 \(\mathfrak C\)를 \(\mathfrak Z_\Omega\)의 근기라 할 때 표현가군 \(\mathfrak Z_\Omega/\mathfrak C\)에 국한해도 충분하다(M. Z. § 19).

그런데 \(\mathfrak Z_\Omega/\mathfrak C\)는 완전가약계이므로 \(\Omega\) 위의 유한차수 체들의 직접합이다. \(\Omega\)가 대수적으로 닫혀 있으므로 이 가수들은 \(\Omega\)와 동형이다.
\[
        \mathfrak Z_\Omega/\mathfrak C
        =e_1\Omega+\cdots+e_t\Omega.
\]
\(e_\nu\)들은 단위원의 성분들이다. 그러므로 다음이 성립한다.

\emph{\(\mathfrak Z_\Omega\)는 \(\Omega\) 안에서 정확히 \(t\)개의 서로 다른 기약 표현류를 가지며, 이들은 모두 1차이다. 여기서 \(t\)는 \(\mathfrak Z_\Omega/\mathfrak C\)의 랭크이므로 \(n\) 이하이다.}

각 표현류는 표현을 하나만 포함한다. 가환체의 1차 표현은 변환해도 \(\pi^{-1}\alpha\pi=\alpha\)이므로 바뀌지 않기 때문이다.

따라서 \(\mathfrak Z_\Omega\)에서 \(\Omega\)의 부분환들로 가는 서로 다른 준동형상 \(\Theta_1,\ldots,\Theta_t\)가 정확히 \(t\)개 있다.

\subsection*{§ 7. 한 체의 동형들}

이제 \(\mathfrak Z\)를 체라 하자. \(\mathfrak Z_\Omega\)에서 \(\Omega\)의 한 부분환으로 가는 각 준동형 \(\Theta_i\)는 \(\mathfrak Z\)를 \(\Omega\)의 부분체 \(\mathsf{Z}_i\)로 동형사상해야 한다. \(Z_i=0\)이거나 \(\mathsf{Z}_i\)가 \(\mathfrak Z\)의 동형상이기 때문이다. 여기서는 \(\mathfrak Z\)의 단위원에 \(\Omega\)의 단위원이 대응하므로 \(Z_i=0\)일 수 없다. 따라서 다음을 얻는다.

\emph{\(\mathfrak Z\)가 \(\mathsf P\)의 \(n\)차 유한체확대이면, \(\Omega\) 안에서 \(\mathfrak Z\)와 동형인 서로 다른 체 \(\mathsf{Z}_i\)의 수는 \(\mathfrak Z_\Omega/\mathfrak C\)의 랭크 \(t\)와 같다. 따라서 \(\mathsf{Z}_i\)의 수는 \(n\) 이하이다.}

\(\mathsf{Z}_i\)의 수가 \(n\)과 같으면 \(\mathfrak Z\)를 \(\mathsf P\) 위의 \emph{제1종}이라 하고, 그 수가 \(n\)보다 작으면 \emph{제2종}이라 한다.

\emph{\(\mathfrak Z\)는 \(\mathfrak Z_\Omega\)가 완전가약일 때 그리고 그때에만 \(\mathsf P\) 위의 제1종이다.}

여기서 이 갈루아 이론의 기초가 통상적인 것보다 갖는 장점을 볼 수 있다. 통상적인 갈루아 이론에서는 \(\mathsf{Z}_i\) 가운데 하나에서 다른 것들로 가는 동형들을 고찰하므로 비대칭이 생긴다. 여기서는 이를 피한다. 한 체의 켤레체 수가 그 차수보다 많지 않음을 보이기 위해 보통은 원시원소를 쓰지만, 여기서는 \(\mathfrak Z_\Omega\)의 모든 기약표현이 정칙표현에 들어 있다는 사실에서 단순히 따른다.

\subsection*{§ 8. 분해체}

\(\mathsf{Z}_i\)들은 \(\mathfrak Z\)의 \(\Omega\) 안에서의 기약표현들이다. 따라서 \(\mathsf P\subseteqq\Gamma\subseteqq\Omega\)인 체 \(\Gamma\)가 모든 \(\mathsf{Z}_i\)를 포함하면 \(\mathfrak Z_\Gamma/\mathfrak C\)는 이미 \(t\)개의 1차 체 \(\Gamma\)로 분해되어야 한다. 역으로 \(\mathfrak Z_\Gamma/\mathfrak C\)가 \(t\)개의 체로 분해되고 이들이 \(\mathfrak Z_\Gamma\)의 1차 기약표현들이 되려면 \(\mathsf{Z}_i\)들이 \(\Gamma\)에 들어 있어야 한다. \(\mathfrak Z_\Gamma/\mathfrak C\)가 \(t\)개 체의 직접합이 되는 성질을 갖는 체 \(\Gamma\)를 \(\mathfrak Z\)의 분해체라 부르면 다음을 얻는다.

\emph{\(\mathfrak Z\)와 동형인 \(\mathsf{Z}_i\)들의 갈루아체, 곧 그 합성체는 \(\mathfrak Z\)의 최소분해체이다.}

\subsection*{§ 9. \(\mathsf{Z}_1\)에서 \(\mathsf{Z}_i\)로 가는 동형들}

\(\mathfrak Z\)에 국한하면 \(\Theta_i\)는 \(\mathfrak Z\)에서 \(\mathsf{Z}_i\)로 가는 동형사상이다. 따라서 \(S_i=\Theta_i\Theta_1^{-1}\)은 \(\mathsf{Z}_1\)에서 \(\mathsf{Z}_i\)로 가는 동형사상이다. \(\Theta_i\)들이 서로 다르므로 \(S_i\)들도 서로 다르다. 이 \(S_i\)들은 통상적인 갈루아 이론에서 고찰하는, 한 체 \(\mathsf{Z}_1\)에서 그 켤레체들로 가는 동형들이다. 다른 것은 없다. 다른 \(\Theta_i\)가 없기 때문이다.

\subsection*{§ 10. 갈루아군}

\(\mathfrak Z\)가 갈루아체여서 \(\mathsf{Z}_i\)들이 원소집합으로 서로 같으면, \(S_i\)들은 \(\mathsf P\)를 고정하는 \(\mathsf{Z}\)의 모든 자기동형의 군을 이룬다. 실제로 모든 곱 \(S_iS_j\)는 \(\mathsf{Z}\)의 자기동형이고, \(S_i\) 이외의 자기동형은 없으므로 \(S_iS_j=S_k\)인 \(S_k\)가 있다. 따라서 \(S_i\)들은 \(Z\)에서 자기 자신으로 가며 \(\mathsf P\)의 각 원소를 고정하는 모든 자기동형의 군을 이룬다. 이것이 \emph{\(\mathsf P\)에 관한 \(Z\)의 갈루아군}이다.

\subsection*{§ 11. 갈루아 이론의 주정리}

\(\mathfrak Z\)를 \(\mathsf P\)의 제1종 갈루아확대라 하자. 갈루아 이론의 주정리를 증명한다.

\emph{\(\mathsf P\)와 \(\mathfrak Z\) 사이의 체 \(\Gamma\)들과 \(\mathsf P\)에 관한 \(Z\)의 갈루아군 \(\mathfrak G\)의 부분군 \(\mathfrak H\)들을 다음과 같이 일대일로 대응시킬 수 있다. 체 \(\Gamma\)에 대응하는 군 \(\mathfrak H\)는 \(\Gamma\)의 모든 원소를 고정하는 \(Z\)의 모든 자기동형으로 이루어지고, \(\Gamma\)는 \(\mathfrak H\)에 속하는 모든 자기동형 아래에서 고정되는 \(\mathsf{Z}\)의 모든 원소의 전체이다.}

이를 짧게 쓰면 다음과 같다.

\begin{enumerate}
\item \emph{\(\mathfrak H\)가 \(\Gamma\)의 불변군이면 \(\Gamma\)는 \(\mathfrak H\)의 불변체이다.}
\item \emph{\(\Gamma\)가 \(\mathfrak H\)의 불변체이면 \(\mathfrak H\)는 \(\Gamma\)의 불변군이다.}
\end{enumerate}

\paragraph{1은 다음 보조정리로 환원된다.}

\emph{\(\mathsf P \subseteqq \Sigma \subseteqq \mathsf T \subseteqq \mathsf Z\)이고 \(\mathsf T\)가 \(\Sigma\) 위에서 차수 \(t\)이면, \(\Sigma\)에서 한 켤레체로 가는 모든 동형은 \(\mathsf T\)에서 켤레체들로 가는 서로 다른 동형 \(t\)개로 연장될 수 있다. 그 켤레체들은 반드시 \(\mathsf Z\) 안에 들어 있다.}

실제로 \(\mathfrak H\)가 \(\Sigma\)의 불변군이고 \(\mathsf T\)가 \(\mathfrak H\)의 불변체라고 가정하면 다음과 같이 추론할 수 있다. \(\mathfrak H\)의 원소들은 \(\Sigma\)의 항등동형을 \(\mathsf Z\) 전체로 연장한 것들이다. 이 연장들을 \(\mathsf T\)에만 적용하면 보조정리에 따라 \(t\)개의 류로 나뉜다. 같은 류의 동형들은 \(\mathsf T\) 위에서 같고, 서로 다른 류의 동형들은 \(\mathsf T\) 위에서 다르다. 그런데 \(\mathsf T\)가 \(\mathfrak H\)의 불변체이므로 \(t=1\), \(\mathsf T=\Sigma\)이어야 한다. 증명 끝.

보조정리를 증명하기 위해 \(\mathsf P\)와 \(\mathsf Z\) 사이의 체 \(\mathfrak S,\mathfrak T\)로 돌아가자. \(\Theta_1\)은 이들에 각각 \(\Sigma,\mathsf T\)를 대응시킨다. 이제 다음을 증명해야 한다.

\emph{\(\mathfrak S\)에서 \(\Sigma_i\)로 가는 모든 동형 \(\mathsf{H}_i\)는 \(\mathfrak T\)에서 \(\mathsf Z\)의 부분체 \(\mathsf T_i\)로 가는 동형으로 정확히 \(t\)가지 서로 다른 방식으로 연장된다.}

\(\mathfrak Z\)가 \(\mathsf P\) 위의 제1종이라는 가정 때문에 \(\mathfrak Z_\Omega/\mathfrak C\)는 \(\mathfrak Z_\Omega\) 자체와 같고, 따라서 \(\mathfrak S_\Omega\)는 그 차수 \(s\)만큼의 체로 분해된다.
\[
        \mathfrak S_\Omega
        =E_1\Omega+\cdots+E_s\Omega.
\]
\(E_\nu\Omega\)들은 표현가군으로서 \(\mathfrak S\)에서 \(s\)개의 부분체 \(\Sigma_1,\ldots,\Sigma_s\)로 가는 \(s\)개의 동형을 생성한다. \(\mathfrak T\)의 동형들을 찾으려면 \(\mathfrak T_\Omega\)가 필요하다. \(E_\nu\Omega\)들이 \(\mathfrak S_\Omega\)의 아이디얼이므로
\[
        E_\nu\Omega=\mathfrak S_\Omega E_\nu.
\]
또한 명백히
\[
        \mathfrak T_\Omega
        =\mathfrak T\Omega
        =\mathfrak T\mathfrak S_\Omega.
\]
따라서
\[
\begin{aligned}
\mathfrak T_\Omega
&=E_1\mathfrak S_\Omega\mathfrak T+\cdots
  +E_s\mathfrak S_\Omega\mathfrak T\\
&=E_1\mathfrak T_\Omega+\cdots+E_s\mathfrak T_\Omega.
\end{aligned}
\]
이는 \(\mathfrak T_\Omega\)의 아이디얼분해이다. \(E_\nu\cdot\mathfrak T_\Omega\)는 \(t\)개의 단순아이디얼로 분해된다. 실제로 차수관계
\[
[\mathfrak T_\Omega\cdot E_\nu:\mathfrak S_\Omega E_\nu]
\leqq[\mathfrak T_\Omega:\mathfrak S_\Omega]
=[\mathfrak T:\mathfrak S]=t
\]
가 성립하고, \(\mathfrak S_\Omega E_\nu=\Omega E_\nu\cong\Omega\)이므로
\[
[\mathfrak T_\Omega\cdot E_\nu:\Omega]
=[\mathfrak T_\Omega E_\nu:\mathfrak S_\Omega E_\nu]
\leqq t.
\]
모든 차수 \([\mathfrak T_\Omega E_\nu:\Omega]\)의 합은 \(s\cdot t\)이어야 하므로
\[
[\mathfrak T_\Omega\cdot E_\nu:\Omega]=t,
\qquad
E_\nu\mathfrak T_\Omega
=e_\nu^{(1)\cdot}\Omega+\cdots+e_\nu^{(t)}\Omega.
\]
각 \(e_\nu^{(i)}\Omega\)는 \(\mathfrak T_\Omega\)의 기약표현을 준다. 이 표현들을 \(\mathfrak S_\Omega\)에만 적용하면 \(E_\nu\cdot\Omega\)가 생성하는 표현과 일치함을 보이면 증명이 끝난다. 이는 명백하다. 실제로
\[
s\cdot E_\nu=E_\nu\sigma_\nu\quad
\qquad
\bigl(s\in\mathfrak S_\Omega\text{는 }\sigma_\nu\text{로 표현된다}\bigr)
\]
에서
\[
        \sum s\cdot e_\nu^{(i)}
        =\sum e_\nu^{(i)}\sigma_\nu
\]
가 따르고, 합표현의 유일성 때문에
\[
        s\cdot e_\nu^{(i)}=e_\nu^{(i)}\sigma_\nu
\]
이다. 따라서 새 표현들도 원소 \(s\)에 원소 \(\sigma_\nu\)를 대응시킨다.

\paragraph{2. \(\mathsf T\)를 \(\mathfrak H\)의 불변체라 하자.}

\(\mathfrak H\)가 \(\mathsf T\)의 불변군임을 확인하려면, \(\mathfrak H\)에 속하지 않는 한 자기동형 아래에서 고정되지 않는 \(\mathsf T\)-원소 하나를 제시하면 충분하다.

\(\mathsf Z\)의 갈루아군 \(\mathfrak G\)의 원소 \(S_i\)들은 \(\mathsf P\)를 고정하는 \(\mathsf Z\)의 자기동형들이다. 모든 \(z\in\mathfrak Z\)에 대해
\[
        S_i z=z
\]
로 정하여 이들을 \(\mathfrak Z_{\mathsf Z}\)의 자기동형으로 만든다. \(S_i\)를
\[
        \mathfrak Z_{\mathsf Z}
        =e_1\mathsf Z+\cdots+e_n\mathsf Z
\]
에 적용하면 유일하게 정해지는 직접가수 \(e_\nu\mathsf Z\)들을 서로 바꾸어야 한다. 특히
\[
        S_i e_\nu=e_{\nu'},
        \qquad
        S_i e_\nu\ne S_k e_\nu\quad(i\ne k)
\]
이다. \(S_1,\ldots,S_h\)가 \(\mathfrak H\)의 원소들이면 임의의 \(e_\nu\), 이를테면 \(e'\)를 택하여 \(i\leqq h\)인 \(S_i\)들이 여기서 생성하는 모든 \(e_\nu\)를 고찰하자.
\[
        e_1=S_1e',\ldots,e_h=S_he'.
\]
\(i\leqq h\), \(j\leqq h\)이면
\[
        S_i e_j=S_iS_j e'=S_k e'=e_k
        \qquad(k\leqq h)
\]
이므로 \(\mathfrak H\)의 각 \(S_i\)는 \(e_1,\ldots,e_h\)를 서로 치환한다. 따라서
\[
        E_1=e_1+\cdots+e_h
\]
은 모든 \(S_i\in\mathfrak H\) 아래에서 불변인 \(\mathfrak Z\)-원소이다. \(\mathfrak Z\)의 한 기저로 \(E_1\)을 나타내면 그 계수들은 \(\mathfrak H\) 아래에서 불변인 \(\mathsf Z\)-원소이어야 하므로 \(\mathsf T\)에 속한다.

\(E_1\)은 \(\mathfrak H\)에 속하지 않는 모든 \(S_i\) 아래에서 불변이 아니다. 그러한 \(S_i\)에 대해
\[
        S_i E_1=e_i+\cdots
\]
는 \(e_1,\ldots,e_h\) 사이에 없는 성분 \(e_i\)를 포함하지만, \(E_1\)의 합표현은 유일하기 때문이다.

\subsection*{§ 12. 성분 \(e_i\)의 형식적 의미}

\(\mathfrak Z\)의 고정된 기저 \(z_1,\ldots,z_n\)을 택하자.
\[
        \mathfrak Z=z_1\mathsf P+\cdots+z_n\mathsf P.
\]
이제는 \(\mathfrak Z\)가 제1종이라는 것만을 가정한다. 동형 \(\Theta_i\)는 \(z_1,\ldots,z_n\)을 \(\mathsf Z_i\)의 기저 \(\zeta_1^{(i)},\ldots,\zeta_n^{(i)}\)로 사상하며, 그 관계는
\[
        z_\nu e_i=e_i\zeta_\nu^{(i)}
\]
이다.

역으로 단위원의 성분 \(e_j\)는 \(\mathsf Z\)의 계수를 써서 \(z_\nu\)들로
\[
        e_j=\sum_\nu\varrho_\nu^{(j)}z_\nu
\]
나타나야 한다. 따라서
\[
\begin{aligned}
e_j=e_j^2
&=\sum_\nu\varrho_\nu^{(j)}z_\nu e_j\\
&=\sum_\nu\varrho_\nu^{(j)}
  \zeta_\nu^{(j)}e_j
\end{aligned}
\]
이고
\[
\begin{aligned}
0=e_je_k
&=\sum_\nu\varrho_\nu^{(j)}z_\nu e_k\\
&=\sum_\nu\varrho_\nu^{(j)}
  \zeta_\nu^{(k)}e_k.
\end{aligned}
\]
그러므로
\[
        \sum_\nu\varrho_\nu^{(j)}\zeta_\nu^{(j)}=1,
        \qquad
        \sum_\nu\varrho_\nu^{(j)}\zeta_\nu^{(k)}=0
        \quad(j\ne k).
\]
이는 행렬
\[
\begin{pmatrix}
\zeta_1^{(1)}&\cdots&\zeta_n^{(1)}\\
\cdots&\cdots&\cdots\\
\cdots&\cdots&\cdots\\
\cdots&\cdots&\cdots\\
\zeta_1^{(n)}&\cdots&\zeta_n^{(n)}
\end{pmatrix}
\quad\text{와}\quad
\begin{pmatrix}
\varrho_1^{(1)}&\cdots&\varrho_1^{(n)}\\
\cdots&\cdots&\cdots\\
\cdots&\cdots&\cdots\\
\cdots&\cdots&\cdots\\
\varrho_n^{(1)}&\cdots&\varrho_n^{(n)}
\end{pmatrix}
\]
가 서로 역행렬이라는 뜻이다. 달리 말하면 \(\varrho_1^{(i)},\ldots,\varrho_n^{(i)}\)는 \(\zeta_1^{(i)},\ldots,\zeta_n^{(i)}\)에 대한 여기저이다. 이 기저는 수체의 디퍼런트 이론에서 역할을 한다.

\subsection*{§ 13. 일반 연장정리}

10쪽에서 증명한 \(\mathsf{H}_i\)의 연장에 관한 정리는 일반화할 수 있다.

\(\mathfrak Z=z_1\mathsf P+\cdots+z_n\mathsf P\)를 가환 완전가약 초복소계라 하자. \(\mathfrak S,\mathfrak T\)를 다음과 같이 \(\mathfrak Z\)의 부분계라 하자. 이 상황에서는 이들도 반드시 완전가약이다.
\[
        \mathsf P\subseteqq\mathfrak S
        \subseteqq\mathfrak T\subseteqq\mathfrak Z.
\]
M. Z. § 13에 따라 세 계는 각각 단위원을 갖는다. 그러나 이 세 단위원은 같을 필요가 없다. \(\mathfrak S,\mathfrak T\)가 \(\mathfrak Z\)의 서로 다른 아이디얼이면 확실히 같지 않다. \(E\)를 \(\mathfrak S\)의 단위원이라 하자. 계 \(\mathfrak T E\)는 유한랭크 \(\mathfrak S\)-가군이다.
\[
        E\mathfrak T
        =a_1\mathfrak S+\cdots+a_t\mathfrak S.
\]
일반적으로 \(\mathfrak T\) 자체는 그렇지 않다.

\emph{증명.} \(\mathfrak S\)는 체들의 직접합
\[
        \mathfrak S=\mathfrak K_1+\cdots+\mathfrak K_r
\]
이다. 이에 \(\mathfrak T E\)의 직접합분해가 대응한다. \(\mathfrak T E\)의 각 가수는 \(\mathfrak K_i\) 가운데 하나를 포함해야 한다. 그렇지 않으면 \(E\)에 의해 소멸되지만 \(E\)는 \(\mathfrak T E\)의 단위원이기 때문이다. 따라서 \(\mathfrak T E\)는 \(\mathfrak K_i\)들의 확대환들의 직접합, 곧
\[
        \mathfrak T E
        =a_{11}\mathfrak K_1+\cdots+a_{1p_1}\mathfrak K_1
        +\cdots+a_{rp_r}\mathfrak K_r
\]
꼴의 직접합이다. 따라서
\[
        E\mathfrak T
        =a_{11}\mathfrak S+\cdots+a_{rp_r}\mathfrak S.
\]
증명 끝.

\begin{center}
\textbf{일반 연장정리.}
\end{center}

\emph{\(\mathfrak S\)의 \(\Omega\) 안에서의 모든 기약표현은 한 준동형 \(\mathsf{H}_i\)로 주어지며, 이를 \(\mathfrak T\)의 표현으로 정확히 \(t\)가지 서로 다른 방식으로 확대할 수 있다.}

증명은 10--11쪽의 것과 같다. 다음 식들이 성립한다.
\[
        \mathfrak S_\Omega
        =E_1\Omega+\cdots+E_s\Omega,
        \qquad
        E_\nu\Omega=\mathfrak S_\Omega E_\nu,
\]
\[
\begin{aligned}
\mathfrak T_\Omega\cdot E
&=\mathfrak T\cdot E\Omega
 =\mathfrak T\mathfrak S\Omega
 =\mathfrak T\mathfrak S_\Omega,
\end{aligned}
\]
\[
\begin{aligned}
\mathfrak T_\Omega\cdot E
&=\mathfrak T(E_1\mathfrak S_\Omega+\cdots
  +E_s\mathfrak S_\Omega)\\
&=E_1\mathfrak T_\Omega+\cdots
  +E_s\mathfrak T_\Omega,
\end{aligned}
\]
\[
\begin{aligned}
[\mathfrak T_\Omega E_\nu:\Omega]
&=[\mathfrak T_\Omega E_\nu:\Omega E_\nu]\\
&=[\mathfrak T_\Omega E_\nu:E_\nu\mathfrak S_\Omega]\\
&\leqq[\mathfrak T_\Omega\cdot E:\mathfrak S_\Omega]
 =[\mathfrak T\cdot E:\mathfrak S]=t,
\end{aligned}
\]
\[
        \sum[\mathfrak T_\Omega\cdot E_\nu:\Omega]=st.
\]
따라서
\[
        [\mathfrak T_\Omega E_\nu:\Omega]=t.
\]
그러므로 \(\mathfrak T_\Omega E_\nu\)는 \(t\)개의 아이디얼로 분해된다. 각 아이디얼은 \(\mathfrak T_\Omega E\), 따라서 \(\mathfrak T_\Omega\)의 한 가수이므로 \(\mathfrak T_\Omega\)의 표현을 주며, 이 표현은 10--11쪽에서와 같이 \(\mathsf{H}_\nu\)가 주는 표현의 연장임이 드러난다. \(\mathfrak T_\Omega\)의 나머지 기약표현들, 곧 \(\mathfrak T_\Omega\cdot E_\nu\)에 포함되지 않는 \(\mathfrak T_\Omega\)의 가수들이 생성하는 표현들은 \(\mathsf{H}_\nu\)가 생성하는 \(\mathfrak S_\Omega\)의 표현을 연장한 것일 수 없다. 이 아이디얼들이 \(\mathfrak T_\Omega\cdot E_\nu\), 따라서 더구나 \(\mathfrak S_\Omega\)에 의해 소멸되기 때문이다.

\section*{제III장. 아벨군}

\subsection*{§ 14. 군환}

\(\mathfrak G\)를 원소 \(a_1,\ldots,a_h\)를 갖는 유한군, \(\mathsf P\)를 가환체라 하자. 군의 곱셈을 곱셈으로 써서 초복소계
\[
        \mathfrak o[\mathfrak G]
        =a_1\mathsf P+\cdots+a_h\mathsf P
\]
를 만들고, 이를 \emph{\(\mathsf P\) 안에서의 \(\mathfrak G\)의 군환}이라 한다(M. Z. § 6; Speiser, \emph{Theorie der Gruppen endlicher Ordnung}, § 48).

\(\mathfrak o[\mathfrak G]\)의 \(\mathsf P\) 안에서의 표현, 또는 \(\mathsf P\) 위의 대수적으로 닫힌 체 \(\Omega\) 안에서의 표현을 묻자. 앞서 말했듯 기약표현들은 \(\mathfrak o[\mathfrak G]/\mathfrak C\)의 단순 좌측아이디얼들이 생성한다. 따라서 갈루아 이론에서와 마찬가지로 \(\mathfrak o[\mathfrak G]\)가 언제 근기를 갖는가가 중요하다. 이에 관하여 다음 일반정리가 성립한다.

\begin{center}
\emph{\(\mathfrak G\)의 위수 \(h\)가 \(\mathsf P\)의 표수 \(p\)로 나누어지지 않을 때 그리고 그때에만 \(\mathfrak o[\mathfrak G]\)는 완전가약이다.}
\end{center}

먼저
\[
        h\equiv0\pmod p\quad\Longrightarrow\quad
        \mathfrak C\ne0
\]
을 증명하자.
\[
        \mathfrak a
        =\mathsf P(a_1+\cdots+a_h)
\]
는 \(\mathfrak o\)의 양쪽아이디얼이다. 실제로 군원소 \(a_i\)들의 성질 때문에
\[
\begin{aligned}
\mathsf P\cdot(a_1+\cdots+a_h)a
&=\mathsf P\cdot(a_1a+\cdots+a_ha)\\
&=\mathsf P(a_1+\cdots+a_h)=\mathfrak a
\end{aligned}
\]
이고
\[
\begin{aligned}
a\cdot\mathsf P(a_1+\cdots+a_h)
&=\mathsf P(aa_1+\cdots+aa_h)\\
&=\mathsf P(a_1+\cdots+a_h)=\mathfrak a.
\end{aligned}
\]
\(\mathfrak a\)는 물론 0이 아니지만 멱영이다.
\[
        \left(\sum_i a_i\right)^2
        =\sum_i\sum_j a_i a_j
        =h\sum_j a_j=0.
\]

그 역
\[
        h\not\equiv0\pmod p\quad\Longrightarrow\quad
        \mathfrak C=0
\]
은 여기서는 아벨군에 대해서만 증명한다. 이 부분은 M. Z. § 26에서 일반적으로 증명되어 있다.

\subsection*{§ 15. 아벨군의 군환}

\(\mathfrak A=\{a_1,\ldots,a_h\}\)를 아벨군이라 하자. 그러면 \(\mathfrak o\)가 가환이므로 가환 초복소계에 관한 일반정리들을 적용할 수 있다.

\emph{\(h\)가 \(\mathsf P\)의 표수 \(p\)로 나누어지지 않으면 \(\mathfrak o[\mathfrak A]\)는 \(\Omega\) 안에서 서로 다른 기약표현을 정확히 \(h\)개 갖는다. \(\mathfrak o\)가 가환이므로 이들은 1차이다.}

그러면 \(\mathfrak o[\mathfrak A]/\mathfrak C\)는 \(\mathfrak o\)와 같은 랭크를 가지며, 따라서 \(\mathfrak C=0\)이다. 이는 바로 위에서 남겨 둔 정리의 둘째 부분을 아벨군에 대해 증명한다.

\emph{증명.} \(\mathfrak A\)는 위수가 \(h_i\)인 순환군 \(\mathfrak Z_i\)들의 직접곱이다.
\[
        \mathfrak A
        =\mathfrak Z_1\times\cdots\times\mathfrak Z_t,
        \qquad
        h=h_1h_2\ldots h_t .
\]
\(z_i\)가 \(\mathfrak Z_i\)의 생성원이면, 한 1차 기약표현에서 이에 대응하는 원소는 \(h_i\)차 단위근 \(\varepsilon_\nu^{(h_i)}\)이어야 한다. 역으로 각 \(z_i\)에 하나의 \(h_i\)차 단위근 \(\varepsilon_\nu^{(h_i)}\)를 대응시키면 \(\mathfrak o\)의 \(\Omega\) 안에서의 기약표현을 얻는다. \(h\not\equiv0\pmod p\)이면 더구나 각 \(h_i\not\equiv0\pmod p\)이므로 \(\mathfrak o\)의 서로 다른 표현은 정확히 \(h_1h_2\ldots h_t=h\)개이다. 증명 끝.

(\(h\equiv0\pmod p\)이면 적어도 한 \(h_i\equiv0\pmod p\)이고, 그 \(h_i\)에 대해서는 서로 다른 \(\varepsilon_\nu^{(h_i)}\)가 많아야 \(h_i/p\)개뿐이므로 표현은 \(h\)개보다 적다. 따라서 \(\mathfrak o_\Omega\)는 0이 아닌 근기를 가져야 한다. 그러나 이 사실만으로는 \(\mathfrak o\)의 근기에 관해 곧바로 아무것도 따르지 않는다.)

\(\mathfrak o_\Omega\)에서 \(\Omega\)의 부분환들로 가는 서로 다른 \(h\)개의 준동형상을 \(\Theta_1,\ldots,\Theta_h\)라 하자. 각 \(\Theta_i\)에는 표현가군으로서 \(\mathfrak o_\Omega\)의 한 가수 \(e_i\Omega\)가 속한다. 때로는 \(\Theta_i\)를 \(\mathfrak A\)에만 적용하고 이를 \(\mathfrak A\)의 표현이라 부른다. \(\mathfrak A\)를 단위근들의 군들로 사상하는 이 준동형들이 \emph{\(\mathfrak A\)의 \(h\)개 지표}이다. 특별한 지표로 모든 \(\mathfrak A\)-원소에 1을 대응시키는 \emph{주지표}가 있다. 이에 속하는 \(\Theta_i\)를 \(\Theta_1\)이라 하자.

주지표에 의한 표현은 아벨군만 아니라 모든 군에서 가능하다. 이 표현은 이미 \(\mathsf P\) 안에 있으므로 \(\mathfrak o\)가 완전가약이면 \(\mathfrak o\)는 이 표현을 주는 단순가수 \(E\cdot\mathfrak o\)를 떼어 내야 한다. \(E\)는 쉽게 계산된다. \(E=\sum_i x_i a_i\)라 하자. 첫째,
\[
        aE=E\cdot1,
        \qquad\text{곧}\qquad
        \sum_i x_i aa_i=\sum_i x_i a_i
\]
이다. \(a=a_j a_i^{-1}\)로 놓으면
\[
\begin{aligned}
&x_1a_ja_i^{-1}a_1+\cdots+x_i a_j+\cdots+x_ha_ja_i^{-1}a_h\\
&\qquad=x_1a_1+\cdots+x_ha_h
\end{aligned}
\]
이므로 계수들을 비교하여
\[
        x_i=x_j
\]
를 얻는다.

둘째,
\[
        E=E^2
        =x^2\left(\sum_i a_i\right)^2
        =x^2h\sum_i a_i=xhE.
\]
\(h\equiv0\pmod p\)이면 \(E=0\)이다. 곧 \(\mathfrak o\)는 주지표의 표현가군을 포함하지 않으므로 완전가약일 수 없었다. 이로써 13쪽 정리의 한 방향을 다시 증명했다.

그러나 \(h\not\equiv0(p)\), 따라서 \(\mathfrak o\)가 완전가약이면 곧바로
\[
        x=\frac1h,
        \qquad
        E=\frac1h(a_1+\cdots+a_h)
\]
가 따른다.

\subsection*{§ 16. 지표관계}

\(E=(1/h)\cdot\sum a_i\)가 주지표를 준다는 사실에서 지표관계가 곧바로 나온다.
\[
        a\cdot e_i=e_i\cdot\Theta_i a
\]
이어야 하므로
\[
        e_i\cdot\Theta_i e_j
        =\begin{cases}
          e_i,&i=j,\\
          0,&i\ne j,
        \end{cases}
\]
또는
\[
        \Theta_i e_j
        =\begin{cases}
          1,&i=j,\\
          0,&i\ne j
        \end{cases}
\]
이다. 특히
\[
        \Theta_i e_1
        =\begin{cases}
          1,&i=1,\\
          0,&i\ne1
        \end{cases}
\]
이고 \(e_1=(1/h)\sum a_i\)이므로
\[
        \sum_\nu\Theta_i a_\nu
        =\begin{cases}
          h,&i=1,\\
          0,&i\ne1.
        \end{cases}
\]
이것이 잘 알려진 지표관계이다.

\emph{한 지표를 모든 군원소에 걸쳐 합하면 주지표의 경우에는 \(h\), 나머지 경우에는 0이다.}

\subsection*{§ 17. 아벨군의 갈루아 이론}

아벨군의 군환에 대해서도 가환체의 갈루아 이론에서 증명한 것과 비슷한 정리들을 세울 수 있다. 갈루아 이론에서는 \(\Theta_i\Theta_1^{-1}\)들이 한 군 \(\mathfrak G\)를 이루었다. 주정리는 \(\mathfrak G\)의 부분군들과 \(\mathsf{Z}\)의 부분체들을 일정한 불변성질에 따라 일대일로 대응시킨다.

\(\mathfrak o[\mathfrak A]\)의 준동형 \(\Theta_i\)들도 한 군 \(\mathsf A\)로 묶을 수 있으며, 그 부분군들은 \(\mathfrak A\)의 부분군들과 비슷한 관계에 있다.

\emph{곱 \(\Theta_i\Theta_k=\Theta\)를}
\[
        \Theta a=\Theta_i a\cdot\Theta_k a
\]
\emph{로 정의하면, 모든 \(\Theta_i\)의 집합은 \(\mathfrak A\)와 동형인 군 \(\mathsf A\)이다.}

\emph{증명.} \(\Theta=\Theta_i\cdot\Theta_k\)는
\[
\begin{aligned}
\Theta a\cdot\Theta b
&=\Theta_i a\cdot\Theta_k a\cdot
  \Theta_i b\cdot\Theta_k b\\
&=\Theta_i a\cdot\Theta_i b\cdot
  \Theta_k a\cdot\Theta_k b\\
&=\Theta_i(ab)\cdot\Theta_k(ab)
 =\Theta(ab)
\end{aligned}
\]
이므로 \(\mathfrak A\)의 준동형이다. \(\mathfrak A\)에는 \(\Theta_i\) 이외의 준동형이 없으므로 \(\Theta=\Theta_j\)인 \(j\)가 있다. 따라서 \(\Theta_i\)들은 실제로 \(h\)개 원소를 갖는 군을 이룬다.

\(\varepsilon_{h_\nu}\)를 원시 \(h_\nu\)차 단위근이라 하자. \(\Theta_i\)는 \(\mathfrak A\)의 순환인자 \(\mathfrak Z_\nu\)의 생성원 \(z_\nu\)를 \(\varepsilon_{h_\nu}\)의 한 거듭제곱으로 사상한다.
\[
        \Theta_i z_\nu
        =\varepsilon_{h_\nu}^{\varrho_\nu^{(i)}}.
\]
준동형 \(\Theta_i\)에 \(\mathfrak A\)의 원소
\[
        \prod_\nu z_\nu^{\varrho_\nu^{(i)}}
\]
를 대응시키자. 이 대응은 \(\mathsf A\)에서 \(\mathfrak A\)로 가는 일대일사상이다. \(\Theta_i\)는 \(\varrho_\nu^{(i)}\)들로 결정되므로 서로 다른 \(\Theta_i\)에는 서로 다른 곱이 대응하기 때문이다. 또한 이 사상은 동형이다. 한편
\[
\begin{aligned}
\Theta_i\Theta_j z_\nu
&=\Theta_i z_\nu\cdot\Theta_j z_\nu\\
&=\varepsilon_{h_\nu}^{
  \varrho_\nu^{(i)}+\varrho_\nu^{(j)}},
\end{aligned}
\]
다른 한편
\[
\prod_\nu z_\nu^{\varrho_\nu^{(i)}}
\cdot
\prod_\nu z_\nu^{\varrho_\nu^{(j)}}
=\prod_\nu z_\nu^{
\varrho_\nu^{(i)}+\varrho_\nu^{(j)}}.
\]
증명 끝.

이제 다음을 정의한다.

\emph{\(\mathsf A\)의 부분군 \(\mathsf B\)의 불변영역}은 모든 \(\Theta\in\mathsf B\)에 의해 1로 사상되는 \(\mathfrak A\)의 원소들이 이루는 부분군 \(\mathfrak B\)이다.

\emph{\(\mathfrak A\)의 부분군 \(\mathfrak B\)의 불변군}은 모든 \(a\in\mathfrak B\)를 1로 사상하는 \(\mathsf A\)의 원소들이 이루는 부분군 \(\mathsf B\)이다.

주정리는 다음과 같다.

\emph{\(\mathfrak A\)의 부분군 \(\mathfrak B\)들과 \(\mathsf A\)의 부분군 \(\mathsf B\)들을 일대일로 대응시킬 수 있다. 서로 대응하는 \(\mathfrak B,\mathsf B\)에 대해 \(\mathsf B\)는 \(\mathfrak B\)의 불변군이고 \(\mathfrak B\)는 \(\mathsf B\)의 불변영역이다.}

따라서 다음 두 명제를 보여야 한다.

\begin{enumerate}
\item \emph{\(\mathsf B\)가 \(\mathfrak B\)의 불변군이면 \(\mathfrak B\)는 \(\mathsf B\)의 불변영역이다.}
\item \emph{\(\mathfrak B\)가 \(\mathsf B\)의 불변영역이면 \(\mathsf B\)는 \(\mathfrak B\)의 불변군이다.}
\end{enumerate}

\paragraph{1.}
이는 10쪽에서와 같이 일반 연장정리에서 따른다. \(\mathsf B\)가 \(\mathfrak B\)의 불변군이고, \(\mathfrak B'\)가 \(\mathsf B\)의 불변영역이라고 하자. 그러면 \(\mathfrak B'\)는 \(\mathfrak B\)를 포함한다. \(\mathsf B\)의 원소들은 \(\mathfrak B\)를 1로 사상하는 하나의 준동형을 \(\mathfrak A\) 전체로 연장한 것들이다. 이 연장들을 \(\mathfrak B'\)에만 적용하면 일반 연장정리에 따라 \(t\)개의 류로 나뉜다. 같은 류의 준동형들은 \(\mathfrak B'\)에서 같고, 서로 다른 류의 준동형들은 \(\mathfrak B'\)에서 다르다. 여기서 \(t\)는 \(\mathfrak o[\mathfrak B']\)의 \(\mathfrak o[\mathfrak B]\) 위에서의 차수, 곧 \(\mathfrak B'\) 안에서 \(\mathfrak B\)의 지수이다. \(\mathfrak o[\mathfrak B]\)와 \(\mathfrak o[\mathfrak B']\)가 \(\mathfrak A\) 전체와 같은 단위원을 갖기 때문이다. 그런데 \(\mathfrak B'\)는 \(\mathsf B\)의 불변영역이므로 \(t=1\), 따라서 \(\mathfrak B'=\mathfrak B\)이어야 한다. 증명 끝.

\textbf{2.}는 1로 환원한다.

\(\mathfrak A\)의 원소 \(a\)를 \(\mathsf A\)의 준동형으로 보자. 곧 \(a\Theta\)를 단위근 \(\Theta a\)와 같게 놓아
\[
        a\Theta=\Theta a
\]
로 정한다. 그러면 \(a\)는 실제로 \(\mathsf A\)의 준동형이다. 왜냐하면
\[
        a\Theta\cdot a\Theta'
        =\Theta a\cdot\Theta'a
        =\Theta\Theta'a
        =a\Theta\Theta'
\]
이기 때문이다.

두 준동형 \(a,b\)의 곱은 통상적인 군곱이다.
\[
        ab\Theta
        =a\Theta\cdot b\Theta
        =\Theta a\cdot\Theta b
        =\Theta(ab).
\]

군원소로서 서로 다른 \(a,b\)는 준동형으로서도 서로 다르다. 모든 \(\Theta\)에 대해 \(a\Theta=b\Theta\)이면 모든 \(\Theta\)에 대해 \(\Theta a=\Theta b\), 곧 \(\Theta(ab^{-1})=1\)이다. 따라서 \(\Theta\)들은 \(ab^{-1}\)이 생성하는 \(\mathfrak A\)의 순환부분군 \(\mathfrak Z\)의 항등준동형을 \(\mathfrak A\) 전체로 연장한 것들이다. 연장정리에 따라 \(\mathfrak A\)에서 \(\mathfrak Z\)의 지수는 \(h\)이다. 곧 \(\mathfrak Z=\{e\}\)이고 \(a=b\)이다.

이제 새 관점에서 한 명제
\[
\begin{array}{ll}
\text{a)}&\mathfrak B\text{는 }\mathsf B\text{의 불변군이다},\\
\text{b)}&\mathsf B\text{는 }\mathfrak B\text{의 불변영역이다}
\end{array}
\]
를 옛 관점으로 번역하자.

a)에서 \(\mathfrak B\)는 모든 \(\Theta\in\mathsf B\)에 대해 \(a\Theta=1\)인 모든 \(a\in\mathfrak A\)로 이루어진다. 이는 모든 \(\Theta\in\mathsf B\)에 대해 \(\Theta a=1\)인 모든 \(a\in\mathfrak A\)로 이루어진다는 뜻이다. 따라서 a)의 번역은
\[
        \text{a') }\quad
        \mathfrak B\text{는 }\mathsf B\text{의 불변영역이다}
\]
이다.

마찬가지로 b)의 번역은
\[
        \text{b') }\quad
        \mathsf B\text{는 }\mathfrak B\text{의 불변군이다}
\]
이다.

그러므로 1에서 이미 증명한 명제
\[
\mathfrak B\text{가 }\mathsf B\text{의 불변군이면 }
\mathsf B\text{는 }\mathfrak B\text{의 불변영역이다}
\]
를 번역하면 바로 증명하려는 명제
\[
\mathfrak B\text{가 }\mathsf B\text{의 불변영역이면 }
\mathsf B\text{는 }\mathfrak B\text{의 불변군이다}
\]
가 된다.

15쪽의 지표관계를 \(\mathfrak A\) 대신 \(\mathsf A\)에 대해 쓰면
\[
\sum_{i=1}^{h} \Theta_i a_\nu =
\begin{cases}
h & \nu=1\\
0 & \nu\ne 1.
\end{cases}
\]
를 얻는다.

\emph{모든 지표를 한 군원소에서 합하면 단위원의 경우에는 \(h\), 나머지 군원소의 경우에는 0이다.}

\section*{제IV장. 양쪽단순환}

\subsection*{§ 18. 보조정리}

\(\mathsf K\)를 체, \(\mathfrak G\)를 \(\mathsf K\)의 자기동형들, 곧 자기 자신으로 가는 동형사상들의 한 군이라 하자. \(\mathfrak G\)에 속하는 모든 자기동형 아래에서 변하지 않는 \(\mathsf K\)-원소들의 집합은 \(\mathsf K\)의 부분체 \(\mathsf P\)이고, 이를 \(\mathfrak G\)의 \emph{불변체}라 한다.

0이 아닌 임의의 \(\mathsf K\)-원소 \(\tau\)에 대해 \(\mathsf K\)의 원소 \(\alpha\)에 \(\tau^{-1}\alpha\tau\)를 대응시키면 \(\mathsf K\)의 자기동형을 얻는다. 군론의 대응하는 명칭을 따라 이를 \emph{내부자기동형}이라 한다. 한 체 \(\mathsf K\)의 모든 내부자기동형이 이루는 군의 불변체는 \(\mathsf K\)의 중심이다.

이제 체 \(\mathsf K\), 그 자기동형군 \(\mathfrak G\), 그리고 그 불변체 \(\mathsf P\)가 주어졌다고 하자.
\[
        \mathfrak M
        =x_1\mathsf P+\cdots+x_n\mathsf P
\]
를 랭크 \(n\)인 \(\mathsf P\)-가군이라 하자. \(\mathfrak G\)의 원소 \(\gamma\)가 \(\mathfrak M\)의 원소들에 작용하는 방식을
\[
        \gamma\cdot x_i=x_i,
        \qquad
        \gamma\sum\varrho_i x_i=\sum\varrho_i x_i
\]
로 정의한다. 그러면 \(\mathfrak G\)는 확대가군
\[
        \mathfrak M_{\mathsf K}
        =x_1\mathsf K+\cdots+x_n\mathsf K
\]
의 자기동형군이 되고, \(\mathfrak M_{\mathsf K}\) 안에서 \(\mathfrak G\)의 불변영역은 \(\mathfrak M\)이다. 이 가정 아래 다음 정리가 성립한다.

\begingroup\itshape
\[
        \mathfrak C
        =z_1\mathsf K+\cdots+z_r\mathsf K
\]
가 \(\mathfrak M_{\mathsf K}\)의 부분가군이고, \(\mathfrak G\)의 모든 원소 \(\gamma\)가 \(\mathfrak C\)를 자기 자신으로 옮긴다고 하자. 각 원소를 고정할 필요는 없다. 그러면 \(\mathfrak C\)는 \(\mathfrak M\)의 한 부분가군 \(\mathfrak c\)를 \(\mathsf K\)로 확대한 가군이다.
\[
        \mathfrak C=\mathfrak c_{\mathsf K}.
\]
\endgroup

\emph{증명.} \textbf{1. 예비논평.} 반드시
\[
        \mathfrak c=\mathfrak C\cap\mathfrak M
\]
이어야 한다. 일반적으로
\[
        \mathfrak c
        =\mathfrak c_{\mathsf K}\cap\mathfrak M
\]
이 성립하기 때문이다.

\(\mathfrak c\subseteqq\mathfrak c_{\mathsf K}\cap\mathfrak M\)은 명백하다. 역포함을 보이기 위해 \(\mathfrak c\)의 기저 \(y_1,\ldots,y_r\)을 택하자.
\[
        \mathfrak c
        =y_1\mathsf P+\cdots+y_r\mathsf P,
\]
그러면
\[
        \mathfrak c_{\mathsf K}
        =y_1\mathsf K+\cdots+y_r\mathsf K.
\]
\(\mathfrak c_{\mathsf K}\cap\mathfrak M\)의 모든 원소 \(c\)는 \(\mathfrak c_{\mathsf K}\)-원소로서 기저표현
\[
        c=\sum_{i=1}^{r}y_i\varkappa_i
\]
를 갖고, \(\mathfrak M\)-원소로서는 기저표현
\[
        c=\sum_{i=1}^{r}y_i\varrho_i
          +\sum_{j=r+1}^{n}y_j\varrho_j
\]
를 갖는다. 여기서 \(y_{r+1},\ldots,y_n\)은 \(y_1,\ldots,y_r\)을 \(\mathfrak M\)의 한 기저로 보충하는 \(\mathfrak M\)-원소들이다(M. Z. § 6).

두 표현 모두 \(\mathfrak M_{\mathsf K}\) 안의 기저표현으로 볼 수 있으므로 서로 같고,
\[
        \varkappa_i=\varrho_i\quad(i=1,\ldots,r),
        \qquad
        \varrho_j=0\quad(j=r+1,\ldots,n).
\]
따라서
\[
        c=\sum_{i=1}^{r}y_i\varrho_i
\]
이고, \(c\in\mathfrak c\)이다. 증명 끝.

\textbf{2.} 정리를 증명하기 위해 원소 \(z_i\)들이 \(\mathfrak M\)에 속하는 \(\mathfrak C\)의 기저 \(z_1,\ldots,z_r\)을 구성하자. 그러면
\[
        \mathfrak c
        =z_1\mathsf P+\cdots+z_r\mathsf P
\]
로 놓았을 때 실제로
\[
        \mathfrak C=\mathfrak c_{\mathsf K}
\]
이다.

\(\mathfrak C\)의 임의의 기저 \(y_1,\ldots,y_r\)을 택하고, \(\mathfrak M\)의 기저 \(x_1,\ldots,x_n\), 따라서 \(\mathfrak M_{\mathsf K}\)의 기저 가운데 알맞은 원소 \(x_{r+1},\ldots,x_n\)을 덧붙여 \(\mathfrak M_{\mathsf K}\)의 기저로 보충하자. M. Z. § 5에 따라 이것이 가능하다.
\[
        \mathfrak M_{\mathsf K}
        =\mathfrak C+x_{r+1}\mathsf K+\cdots+x_n\mathsf K.
\]
여기서 \(\mathfrak M\)의 기저의 앞 \(r\)개 원소 \(x_1,\ldots,x_r\)에 대해
\[
        x_i=z_i-\sum_{j=r+1}^{n}x_j\varkappa_j^{(i)},\qquad,
        \qquad
        z_i\in\mathfrak C,\quad
        \varkappa_j^{(i)}\in\mathsf K
\]
인 식들을 얻는다.

\(z_i\)들은 \(x_{r+1},\ldots,x_n\)과 함께 \(\mathfrak M_{\mathsf K}\)의 기저를 이루므로 선형독립이고, 따라서 \(z_1,\ldots,z_r\)은 \(\mathfrak C\)의 기저이다. 이제 이것이 구하는 형태의 기저이고 \(z_i\in\mathfrak M\)임을 보이자.

\(\gamma\in\mathfrak G\)라 하자. 한편
\[
\begin{array}{rcl}
(1)\qquad \gamma z_i
&=&\displaystyle
\gamma x_i+\sum_{j=r+1}^{n}
\gamma x_j\cdot\gamma\varkappa_j^{(i)}\\
&=&\displaystyle
x_i+\sum_{j=r+1}^{n}x_j\gamma\varkappa_j^{(i)}.
\end{array}
\]
다른 한편 \(\mathfrak C\)가 모든 \(\gamma\in\mathfrak G\)를 허용하므로
\[
\begin{array}{rcl}
(2)\qquad
\gamma z_i
&=&\displaystyle
\sum_{\mu=1}^{r}z_\mu\sigma_\mu^{(i)},
\qquad
\sigma_\mu^{(i)}\in\mathsf K.
\end{array}
\]
(2)는 다시
\[
\begin{array}{rcl}
(3)\qquad
\gamma z_i
&=&\displaystyle
\sum_{\mu=1}^{r}x_\mu\sigma_\mu^{(i)}
+\sum_{j=r+1}^{n}x_j
 \sum_{\mu=1}^{r}
 \varkappa_j^{(\mu)}\sigma_\mu^{(i)}
\end{array}
\]
로 쓸 수 있다. (1)과 (3)의 계수를 비교하면
\[
        \sigma_j^{(i)}
        =\begin{cases}
          1,&i=j,\\
          0,&i\ne j.
        \end{cases}
\]
따라서
\[
        \gamma\cdot z_i=z_i.
\]
그러므로 \(z_i\)들은 \(\mathfrak G\)의 불변영역, 곧 \(\mathfrak M\)에 속한다. 증명 끝.

\subsection*{§ 19. 양쪽단순 초복소계의 계수체의 비가환 확대체 안에서의 표현}

\[
        \mathfrak S
        =x_1\mathsf P+\cdots+x_n\mathsf P
\]
를 양쪽단순, 따라서 완전가약인 초복소계라 하고, \(\mathsf K\)를 중심이 \(\mathsf P\)인 체라 하자.

\textbf{정리 1.} \emph{\(\mathfrak S_{\mathsf K}\)는 양쪽단순이고, 따라서 완전가약이기도 하다.}

\emph{증명.} \textbf{1.} \(\mathfrak S_{\mathsf K}\)의 모든 양쪽아이디얼 \(\mathfrak a\)는 \(\mathsf K\)의 내부자기동형군 \(\mathfrak G\) 아래에서 불변이다. 실제로 0이 아닌 \(\varkappa\in\mathsf K\)에 대해
\[
        \varkappa^{-1}\mathfrak a\varkappa=\bar{\mathfrak a}
\]
로 놓으면 아이디얼 성질 때문에 \(\bar{\mathfrak a}\subseteqq\mathfrak a\)이다. 그런데 \(\bar{\mathfrak a}\)와 \(\mathfrak a\)는 \(\mathsf K\)-가군으로서 같은 랭크를 가지므로 \(\bar{\mathfrak a}=\mathfrak a\)이어야 한다.

\textbf{2.} 바로 앞의 보조정리에 따라 \(\mathfrak S_{\mathsf K}\)의 양쪽아이디얼 \(\mathfrak a\)는 \(\mathfrak S\)의 한 부분가군 \(\mathfrak A\)의 확장아이디얼이어야 한다.
\[
        \mathfrak a=\mathfrak A_{\mathsf K}.
\]
\(\mathfrak A=\mathfrak S\cap\mathfrak a\)이므로 \(\mathfrak A\)는 \(\mathfrak S\)의 양쪽아이디얼이다. 실제로
\[
\mathfrak S\mathfrak A
\subseteqq\mathfrak S\mathfrak S\cap
\mathfrak S\mathfrak a
\subseteqq\mathfrak S\cap\mathfrak a
=\mathfrak A,
\]
\[
\mathfrak A\mathfrak S
\subseteqq\mathfrak S\mathfrak S\cap
\mathfrak a\mathfrak S
\subseteqq\mathfrak S\cap\mathfrak a
=\mathfrak A.
\]
\(\mathfrak S\)의 양쪽아이디얼은 \(0\)과 \(\mathfrak S\)뿐이므로 \(\mathfrak S_{\mathsf K}\)의 양쪽아이디얼도 \(0\)과 \(\mathfrak S_{\mathsf K}\)뿐이다. 증명 끝.

이제 \(\mathfrak S\)의 \(\mathsf K\) 안에서의 기약표현을 구하되, 형식적인 장점이 있는 반대 표현을 세우겠다. 직접 표현과 반대 표현은 일대일로 대응하므로 어느 쪽을 택해도 결국 같다.

\(\mathfrak M\)을 \(\mathfrak S\)의 \(\mathsf K\) 안에서의 반대 표현가군이라 하자. 곧 \(\mathfrak M\)은 유한랭크 \(\mathfrak S\)-좌측가군이자 \(\mathsf K\)-좌측가군이고
\[
        \mathfrak M
        =\mathsf K\cdot y_1+\cdots+\mathsf K\cdot y_r
\]
이며 교환법칙
\[
        s(\varkappa m)=\varkappa (sm)
\]
을 만족한다.

\(\mathfrak M\)을 \(\mathfrak S_{\mathsf K}\)-좌측가군으로도 볼 수 있음을 증명하자. 먼저 \(\mathfrak S_{\mathsf K}\)의 원소와 \(\mathfrak M\)의 원소의 곱을 정의해야 한다. \(s\in\mathfrak S\), \(x\in\mathsf K\), \(m\in\mathfrak M\)일 때
\[
        sx\cdot m=s\cdot(xm)
\]
로 놓고, 이에 따라 일반 원소 \(\sum s_ix_i\in\mathfrak S_{\mathsf K}\)에 대해
\[
        \sum s_ix_i\cdot m
        =\sum s_i(x_i m)
\]
로 놓는다.

이 정의는 \(\mathfrak S_{\mathsf K}\)-원소를 곱 \(s_i\varkappa _i\)들의 합으로 나타내는 방식에 의존하지 않는다. 이를 보이려면
\[
        \sum_{i=1}^{p}s_i\varkappa_i=0                                      \tag{1}
\]
에서
\[
        \sum_{i=1}^{p}s_i(\varkappa _i m)=0
\]
이 따름을 보이면 된다. \(s_1,\ldots,s_q\)가 선형독립이고 \(s_{q+1},\ldots,s_p\)가
\[
        s_j=\sum_{i=1}^{q}s_i\varrho_i^{(j)},
        \qquad j=q+1,\ldots,p
\]
로 \(s_1,\ldots,s_q\)에 의해 나타난다고 하자. 이 표현은 이미 \(\mathfrak S\) 안에서 가능하고 유일하므로 계수들은 \(\mathsf P\)에 속해야 한다. (1)에서
\[
        x_i+\sum_{j=q+1}^{p}
        \varrho_i^{(j)}\varkappa_j=0,
        \qquad i=1,\ldots,q
\]
을 얻는다. 따라서 가군 성질과 결합법칙을 차례로 쓰면
\[
\begin{aligned}
\sum_{i=1}^{p}s_i(\varkappa _i m)
&=\sum_{i=1}^{q}s_i(x_i m)
  +\sum_{j=q+1}^{p}
    \left(\sum_{i=1}^{q}s_i\varrho_i^{(j)}\right)(x_jm)\\
&=\sum_{i=1}^{q}s_i(x_i m)
  +\sum_{j=q+1}^{p}\sum_{i=1}^{q}
    s_i\bigl(\varrho_i^{(j)}x_jm\bigr)\\
&=\sum_{i=1}^{q}
  s_i\left(
  \left(x_i+\sum_{j=q+1}^{p}
  \varrho_i^{(j)}x_j\right)m\right)\\
&=0.
\end{aligned}
\]
여기서는 \(s'(sm)=s's\cdot m\), \(x'(xm)=x'x\cdot m\)을 썼다.

\(\mathfrak M\)의 \(\mathfrak S_{\mathsf K}\)-가군 성질 가운데 분배법칙
\[
        (S+S')m=Sm+S'm,
        \qquad
        S(m+m')=Sm+Sm'
\]
은 명백하다. 결합법칙 \(S(S'm)=SS'\cdot m\)만 더 살펴보면 된다. 분배법칙이 있으므로
\[
        sx\cdot(s'x'm)=sxs'x'\cdot m
\]
인 경우에 국한할 수 있다. 정의에 따라 한편
\[
\begin{aligned}
sx\cdot(s'x'\cdot m)
&=s\bigl(x(s'(x'm))\bigr),
\end{aligned}
\]
다른 한편
\[
\begin{aligned}
sxs'x'\cdot m
&=ss'xx'\cdot m
 =ss'\bigl(xx'm\bigr)\\
&=s\bigl(s'(\varkappa(\varkappa'm))\bigr)
 =s\bigl(x(s'(x'm))\bigr).
\end{aligned}
\]
마지막 등식에는 두 교환관계
\[
        xs'=s'x,
        \qquad
        s'(xm)=x(s'm)
\]
을 썼다.

따라서 \(\mathfrak M\)은 실제로 유한 \(\mathfrak S_{\mathsf K}\)-좌측가군이다. 역으로 모든 유한 \(\mathfrak S_{\mathsf K}\)-좌측가군은 \(\mathfrak S\)의 \(\mathsf K\) 안에서의 반대 표현가군이다. 유한 \(\mathsf K\)-가군임은 명백하고,
\[
        s(xm)=sx\cdot m=xs\cdot m=x(sm)
\]
이므로 교환조건도 만족하기 때문이다(§ 24의 더 간단한 증명도 비교하라).

M. Z. § 18에 따라 모든 유한 \(\mathfrak S_{\mathsf K}\)-가군은 \(\mathfrak S_{\mathsf K}\)의 단순 좌측아이디얼들과 \(\mathfrak S_{\mathsf K}\)-연산자동형인 단순가군들의 합이다. 역으로 \(\mathfrak S_{\mathsf K}\)의 모든 좌측아이디얼은 \(\mathfrak S\)의 \(\mathsf K\) 안에서의 표현가군이고, 모든 \(\mathfrak S_{\mathsf K}\)-좌측아이디얼은 서로 연산자동형이다. 따라서 다음 정리가 성립한다.

\textbf{정리 2.} \emph{\(\mathfrak S\)의 \(\mathsf K\) 안에서의 유일한 기약 반대 표현류는 \(\mathfrak S_{\mathsf K}\)의 한 단순 좌측아이디얼이 생성하는 것이다.}

이 기약표현의 차수가 \(r\)이면 임의의 표현의 차수는 \(r\)의 배수이다. 반대로 기약표현들을 나란히 붙이면 \(r\)의 임의의 배수 \(rs\)에 대해 \(rs\)차 표현을 만들 수 있다. 이는 \(s\)개의 단순가군의 직접합인 표현가군을 만드는 데 대응한다.

우리는 이미 \(\mathfrak S\) 자체를 표현가군으로 써서 얻는 \(\mathfrak S\)의 \(\mathsf K\) 안에서의 한 표현, 곧 주표현(M. Z. § 20)을 알고 있으며, 이는 이미 \(\mathsf P\) 안에 있다. 그 차수는 \(n\)이므로 \(r\)은 \(n\)을 나눈다.
\[
        n=rt.
\]
몫 \(t\)는 쉽게 구할 수 있다. \(r\)은 한 좌측아이디얼의 \(\mathsf K\) 위에서의 랭크이고 \(n\)은 \(\mathfrak S_{\mathsf K}\) 전체의 랭크이므로, \(t\)는 \(\mathfrak S_{\mathsf K}\)의 좌측아이디얼 가수의 수이다. 따라서 \(t^2\)은 \(\mathfrak S_{\mathsf K}\)가 그 단순아이디얼들의 자기동형체 위에서 갖는 랭크이다.

이제 \(\mathfrak S\)가 가환체인 경우로 고찰을 특수화한다.

\textbf{정리 3.} \emph{\(\mathfrak S\)가 가환체이면 \(\mathfrak S\)는 \(\mathfrak S_{\mathsf K}\)의 중심이다.}

\emph{증명.} \(\mathfrak S\)의 각 원소는 모든 \(\mathfrak S\)-원소 및 모든 \(\mathsf K\)-원소와 가환하므로 모든 \(\mathfrak S_{\mathsf K}\)-원소와 가환한다. 따라서 \(\mathfrak S\)는 \(\mathfrak S_{\mathsf K}\)의 중심에 포함된다. 역으로 한 \(\mathfrak S_{\mathsf K}\)-원소가 모든 \(\mathfrak S_{\mathsf K}\)-원소와 가환하려면 적어도 모든 \(\mathsf K\)-원소와 가환해야 한다. 그러면 그 계수들은 \(\mathsf K\)의 중심, 곧 \(\mathsf P\)에 속하므로 그 원소는 \(\mathfrak S\)에 속한다. 증명 끝.

가환환에서는 반대동형과 직접동형 사이에 차이가 없으므로 모든 반대 표현을 직접 표현으로도 볼 수 있다. 따라서 \(\mathfrak S\)의 \(\mathsf K\) 안에서의 기약표현은 직접 표현가군으로도 얻어져야 한다. 이 직접 표현가군은 \(\mathfrak S_{\mathsf K}\)의 단순 우측아이디얼이다. 실제로 \(\mathfrak r\)은 랭크 \(r\)인 유한 \(\mathsf K\)-우측가군이고, \(\mathfrak r\mathfrak S=\mathfrak S\mathfrak r\)이므로 \(\mathfrak S\)-좌측가군이기도 하며, 관통 결합법칙도 당연히 만족한다. 따라서 이는 랭크 \(r\)인 직접 표현가군이다. 기약 표현류가 하나뿐이므로 주어진 표현을 생성해야 한다.

\subsection*{§ 20. 비가환체}

이제
\[
        \mathfrak K
        =y_1\mathsf P+\cdots+y_m\mathsf P
\]
를 중심이 \(\mathsf P\)인 비가환체라 하자. \(\mathfrak K\)에 관하여 유한 대수적 가환체확대의 이론과 비슷한 일련의 정리들을 전개하겠다.

\(\mathsf{Z}=\xi_1\mathsf P+\cdots+\xi_n\mathsf P\)가 \(\mathsf P\) 위의 초복소계를 뜻하면 두 확대계 \(\mathfrak K_{\mathsf{Z}}\)와 \(\mathsf{Z}_{\mathfrak K}\)는 같다. 실제로
\[
        \mathfrak K_Z
        =y_1\xi_1\mathsf P+\cdots+y_m\xi_n\mathsf P
\]
로 쓸 수 있고
\[
        Z_{\mathfrak K}
        =\xi_1y_1\mathsf P+\cdots+\xi_ny_m\mathsf P
\]
로 쓸 수 있는데, \(y_i\xi_j=\xi_jy_i\)이므로 두 표현은 같다.

이하에서 \(\mathsf{Z}\)는 언제나 가환체이므로 § 19의 정리들을 적용할 수 있다. § 19의 정리 1에 따라 \(\mathfrak K_{\mathsf{Z}}\)는 언제나 양쪽단순이다.

\textbf{정리 1.} \emph{\(\Omega\)를 \(\mathsf P\)의 대수적으로 닫힌 확대체라 하자. 그러면 \(\mathfrak K_\Omega\)의 중심은 \(\Omega\)이다.}

\emph{증명.} § 19의 대응하는 정리 3을 쓴다. \(\Omega\)가 \(\mathfrak K_\Omega\)의 중심에 놓인다는 것은 명백하다. \(\alpha\)가 \(\mathfrak K_\Omega\)의 중심의 임의의 원소이면 그 계수들은 \(\mathsf P\)의 한 유한 대수적 확대체 \(Z\), 곧 이 계수들이 생성하는 체에 속한다. 그러면 \(\alpha\)는 \(\mathfrak K_{\mathsf{Z}}\)의 중심에도 속해야 하는데, 그 중심은 \(\mathsf{Z}\)이다. 따라서 \(\alpha\in Z\subseteqq\Omega\)이다. 증명 끝.

아주 비슷한 방식으로 다음을 증명한다.

\textbf{정리 2.} \emph{\(\mathfrak K_\Omega\)는 양쪽단순이다.}

\emph{증명.} \(\mathfrak a=z_1\Omega+\cdots+z_r\Omega\)를 양쪽아이디얼이라 하자. \(z_1,\ldots,z_r\)의 계수들은 \(\mathsf P\)의 한 유한확대체 \(Z\)에 속한다. \(\mathfrak K_{\mathsf{Z}}\) 안에서
\[
        \mathfrak a'=z_1Z+\cdots+z_rZ
\]
는 양쪽아이디얼이다. \(\mathfrak K_{\mathsf{Z}}\)가 양쪽단순이므로 \(\mathfrak a'\)는 0 또는 \(\mathfrak K_{\mathsf{Z}}\)이고, 따라서 \(\mathfrak a\)는 0 또는 \(\mathfrak K_\Omega\)이다. 증명 끝.

이 정리에 따라 \(\mathfrak K_\Omega\)는 그 중심 \(\Omega\)의 한 확대체 위의 행렬환이다. 그 체는 \(\mathfrak K_\Omega\)의 부분환으로서 \(\Omega\) 위에서 유한차수이어야 하므로 \(\Omega\) 자체와 같다.

\textbf{정리 3.} \emph{\(\mathfrak K_\Omega\)는 \(\Omega\) 위의 행렬환이다.}
\[
        \mathfrak K_\Omega=\sum\Omega c_{ik}.
\]
\emph{따라서 \(\mathfrak K\)가 그 중심 \(\mathsf P\) 위에서 갖는 차수는 제곱수 \(m=t^2\)이다.}

\(\mathfrak K_\Omega\)가 분해되는 우측아이디얼, 또는 좌측아이디얼의 수 \(t\)를 \(\mathfrak K\)의 \emph{절대성분수}라 부르겠다.

\(\mathfrak K_\Omega\)가 \(t\)개의 단순아이디얼로 분해된다는 \(\Omega\)의 성질은 이미 \(\mathsf P\)의 유한확대 \(\mathsf{Z}\)들도 가질 수 있다. 예컨대 \(c_{ik}\)의 계수들을 \(\mathsf P\)에 첨가하여 얻는 \(\mathsf{Z}\)가 그렇다. 이제 이 분해체들을 조사한다.

\textbf{정의.} \(\mathsf P\)의 유한 대수적 확대체 \(Z\)에 대하여 \(\mathfrak K_{\mathsf{Z}}\)가 분해되는 단순 우측아이디얼의 수가 이미 절대성분수 \(t\)와 같으면 \(Z\)를 \(\mathfrak K\)의 \emph{분해체}라 한다.\srcfn{1}{9쪽의 분해체 개념과 비교하라.}

\textbf{정리 4.} \emph{\(\mathsf{Z}\)가 \(\mathfrak K\)의 분해체이면 \(\mathsf{Z}\)는 \(\mathfrak K_{\mathsf{Z}}\)의 우측아이디얼들의 자기동형체이다. 역으로 \(\mathfrak K_{\mathsf{Z}}\)의 우측아이디얼들의 자기동형체가 \(\mathsf{Z}\)이면 \(\mathsf{Z}\)는 분해체이다.}

\emph{증명.} \textbf{1.} \(\mathsf{Z}\)가 분해체이고 \(\mathfrak T\)가 \(\mathfrak K_{\mathsf{Z}}\)의 우측아이디얼들의 자기동형체라고 하자. 그러면
\[
        \mathfrak K_{\mathsf{Z}}=\sum\mathfrak T c_{ik}
\]
이고 따라서
\[
        \mathfrak K_\Omega
        =\sum\mathfrak T_\Omega c_{ik}.
\]
\(\mathfrak K_\Omega\)가 \(\Omega\) 위에서 갖는 랭크는 \(t^2\)과 \(\mathfrak T_\Omega\)의 랭크의 곱인데, 다른 한편 \(t^2\)이다. 따라서 \(\mathfrak T_\Omega\), 그리고 \(\mathfrak T\)의 랭크는 1이고 \(\mathfrak T=\mathsf{Z}\)이다.

\textbf{2.} \(\mathsf{Z}\)가 \(\mathfrak K_{\mathsf{Z}}\)의 단순아이디얼들의 자기동형체이면
\[
        \mathfrak K_{\mathsf{Z}}=\sum\mathsf Zc_{ik},
\]
따라서
\[
        \mathfrak K_\Omega=\sum\Omega c_{ik}.
\]
그러므로 \(c_{ik}\)의 수는 \(t^2\)과 같아야 한다. 증명 끝.

\textbf{정리 5.} \emph{\(\mathfrak K\)의 분해체 \(Z\)의 \(\mathfrak K\) 안에서의 기약표현 \(\mathfrak Z\)는 행렬환 \(\mathfrak K_r\)의 극대 가환부분체이다. 여기서 \(r\)은 \(\mathsf{Z}\)의 기약표현의 차수이다. 22쪽 정리 2에 따라 기약 표현류는 하나뿐이다.}

\emph{역으로 \(\mathsf P\)의 가환확대체 \(Z\)의 기약표현 \(\mathfrak Z\)가 \(\mathfrak K_r\)의 극대 가환부분체이면 \(Z\)는 분해체이다.}

\emph{증명.} \textbf{1.} 분해체 \(Z\)가 \(\mathsf P\) 위에서 차수 \(n\)이면 \(Z\)의 기약표현 \(\mathfrak Z\)의 차수 \(r\)은 \(n/t\)이다(22쪽). 이제
\[
        \mathfrak Z\subseteqq\mathfrak Z^*
        \subseteqq\mathfrak K_r
\]
이고 \(\mathfrak Z^*\)가 \(\mathsf P\) 위에서 차수 \(n^*\geqq n\)인 극대 가환체라고 하자. \(\mathfrak Z^*=\mathfrak Z\)를 보여야 하지만 \(n^*=n\)을 보이면 충분하다. \(\mathfrak Z^*\)와 동형인 \(Z\)의 확대체 \(\mathsf{Z}^*\)의 기약표현 차수를 \(r^*\)라 하자. 그러면 다시 \(n^*=r^*t\)이다. 다른 한편 \(\mathfrak Z^*\) 안에는 \(\mathsf{Z}^*\)의 \(r\)차 표현이 있으므로 \(r^*\)은 \(r\)을 나누고, 따라서 \(r^*\leqq r\)이다. 그러므로
\[
        n\leqq n^*=r^*t\leqq rt=n
\]
에서 \(n^*=n\)이 따른다.

\textbf{2.} \(\mathsf P\)의 가환확대체 \(\mathsf{Z}\)의 \(\mathfrak K\) 안에서의 \(r\)차 기약표현 \(\mathfrak Z\)가 \(\mathfrak K_r\)의 극대 가환부분체라고 하자. \(\mathfrak K_{\mathsf{Z}}\)의 단순아이디얼들의 자기동형체를 \(\mathfrak T\)라 하면
\[
        \mathfrak K_Z
        =\sum_{1}^{t'}\mathfrak T c_{ik}.
\]
\(\mathfrak T\)는 \(\mathsf{Z}\)를 포함한다. \(\mathfrak T=\mathsf{Z}\)임을 보여야 한다. \(\mathfrak K_r\)의 한 단순 우측아이디얼 \(\mathfrak r\)은 \(\mathsf{Z}\)의 \(\mathfrak K\) 안에서의 표현가군일 뿐 아니라 \(\mathfrak T\)의 표현가군이기도 하다. 실제로
\[
        \mathfrak r
        =\mathfrak T c_{i1}+\cdots+\mathfrak T c_{ik},
\]
따라서 \(\mathfrak T\mathfrak r=\mathfrak r\)이고, \(\mathfrak r\)은 \(\mathfrak T\)-좌측가군이자 물론 \(\mathfrak K\)-우측가군이다.

그러므로 \(\mathfrak K_r\)에는 \(\mathfrak T\)와 동형이고 \(\mathfrak Z\)를 포함하는 부분체 \(\overline{\mathfrak T}\)가 있다. \(\mathfrak T\ne Z\), 따라서 \(\overline{\mathfrak T}\ne\mathfrak Z\)이면 \(\mathfrak Z\)에 속하지 않는 \(\overline{\mathfrak T}\)-원소 \(\alpha\)를 \(\mathfrak Z\)에 첨가하여 \(\mathfrak K_r\) 안에서 \(\mathfrak Z\)의 진정한 가환확대체를 얻는다. 이는 \(\mathfrak Z\)의 극대성과 모순이다. 증명 끝.

정리 5에서 곧바로 다음이 따른다.

\textbf{정리 6.} \emph{\(\mathfrak K\)의 극대 가환부분체들은 \(\mathsf P\) 위에서 차수 \(t\)이고 분해체이다.}\srcfn{1}{이제 모든 분해체가 \(\mathfrak K\)의 한 극대 가환부분체를 포함할 것이라는 추측은 옳지 않다. 이 성질을 갖지 않는 분해체가 있을 수 있으며, 이에 속하는 \(r\)은 반드시 1보다 크다. R. Brauer 및 E. Noether, \emph{Über minimale Zerfällungskörper irreduzibler Darstellungen}, Sitzungsberichte der Preußischen Akademie der Wissenschaften 1927 XXXII, 221쪽을 보라.}

이제 \(\mathfrak K_r\)의 부분환들의 동형에 관한 정리를 전개한다. 비가환체 \(\mathfrak K\)의 갈루아군의 구조와 그에 관련된 사실들을 더 정확히 아는 일은 이 정리에 기초한다.

\textbf{정리 7.} \emph{\(\mathfrak S_1,\mathfrak S_2\)를 \(\mathfrak K\) 위의 행렬환 \(\mathfrak K_r\) 안에서 \(\mathsf P\)를 포함하는 양쪽단순 부분환들이라 하고, 둘 사이에 \(\mathsf P\)의 모든 원소를 고정하는 동형 \(\mathfrak S_1\cong\mathfrak S_2\)가 있다고 하자. 그러면 이 동형은 한 정칙 \(\mathfrak K_r\)-원소 \(\tau\)에 의한 변환이다.}

\emph{\(\sigma_1\in\mathfrak S_1\)과 \(\sigma_2\in\mathfrak S_2\)가 서로 대응하면}
\[
        \tau^{-1}\sigma_1\tau=\sigma_2.
\]

\emph{증명.} \(\mathfrak S_1,\mathfrak S_2\)와 반대동형인 한 환 \(\Sigma\)를 구성하고, 그 안에서 \(\mathsf P\)와 동형인 부분체를 \(\mathsf P\)와 동일시한다. 그러면 \(\mathfrak S_1,\mathfrak S_2\)는 같은 환 \(\Sigma\)의 \(\mathfrak K\) 안에서의 같은 차수 \(r\)인 두 반대 표현이다. 22쪽 정리 2에 따라 기약 반대 표현류가 하나뿐이고, 따라서 주어진 각 차수의 표현류도 하나뿐이다. 그 차수는 기약류의 배수이다. 그러므로 \(\mathfrak S_1,\mathfrak S_2\)는 같은 류에 속하고 변환으로 서로 옮겨진다. 증명 끝.

정리 7에서 곧바로 다음이 따른다.

\textbf{정리 8.} \emph{\(\mathsf P\)의 모든 원소를 고정하는 \(\mathfrak K\)의 자기동형군, 곧 그 중심에 관한 \(\mathfrak K\)의 갈루아군 \(\mathfrak G\)는 \(\mathfrak K\)의 내부자기동형군이다.}

정리 7에서 \(r=1\), \(\mathfrak S_1=\mathfrak S_2=\mathfrak K\)로 놓으면 된다.

다음 두 정리는 지금까지 전개한 일반이론이 특수한 문제를 다루는 데 유용함을 보여 준다.

\textbf{정리 9.} \emph{중심이 실수체 \(\mathsf{R}\)인 유한차수 비가환체는 사원수체뿐이다.}

\emph{증명.} \(\mathfrak K\)를 중심이 \(\mathsf{R}\)인 비가환체로서 \(\mathsf{R}\) 위에서 유한랭크인 것이라 하자. 정리 6에 따라 \(\mathfrak K\)가 \(\mathsf{R}\) 위에서 갖는 차수는 \(\mathfrak K\)의 한 극대 가환부분체의 차수 \(t\)의 제곱이다. 그런데 그러한 부분체는 복소수체와 동형인 체일 수밖에 없으므로 \(t=2\)이고, \(\mathfrak K\)는 \(\mathsf{R}\) 위에서 차수 4이다. \(\mathfrak Z_1\)을 \(\mathfrak K\)의 한 극대 가환부분체라 하자. 그러면 \(i^2=-1\)인 \(i\)를 써서 \(\mathfrak Z_1=\mathsf{R}(i)\)로 놓을 수 있다. \(\mathfrak Z_2=\mathsf{R}(-i)\)도 \(\mathfrak Z_1\)과 원소집합으로 같고 그와 동형인 체이며, 동형은 \(i\mapsto-i\)로 주어진다. 정리 7에 따라
\[
        j^{*-1}ij^*=-i
\]
인 \(j^*\in\mathfrak K\)가 있다.

\(j^*\notin\mathfrak Z_1\)이므로 \(\mathfrak Z_1(j^*)\)가 \(\mathfrak Z_1\) 위에서 갖는 랭크는 적어도 2이다. 따라서 \(\mathfrak Z_1(j^*)\)는 \(\mathsf{R}\) 위에서 랭크 4이고 \(\mathfrak K\)와 같다. 곧
\[
        \mathfrak K
        =\mathsf R+\mathsf Ri+\mathsf Rj^*+\mathsf Rij^*.
\]
또
\[
\begin{aligned}
j^{*-2}1j^{*2}&=1\\
j^{*-2}ij^{*2}&=i\\
j^{*-2}j^*j^{*2}&=j^*\\
j^{*-2}ij^*j^{*2}&=ij^*
\end{aligned}
\]
이므로 \(j^{*2}\)는 \(\mathfrak K\)의 모든 원소와 가환하고 중심 \(\mathsf{R}\)에 속한다. 따라서 \(j^{*2}\)는 실수 \(r\)이다. \(r\)은 양수일 수 없다. 그렇다면 가환체 \(\mathsf{R}(j^*)\) 안에서 다항식 \(x^2-r\)이 네 근 \(+\sqrt r,-\sqrt r,+j^*,-j^*\)을 갖기 때문이다. 따라서 \(j^{*2}=-m^2\), \(m\in R\)이다. \(j=j^*m^{-1}\)로 놓으면 \(j^2=-1\)이고 \(1,i,j,k=ij\)는 사원수 단위들이다. 증명 끝.

이 증명은 \(\mathsf{R}\) 대신 임의의 실폐체에 대해서도 그대로 성립한다.

\textbf{정리 10(웨더번 정리).} \emph{유한개의 원소로 이루어진 모든 체는 가환이다.}

\emph{증명.} \(\mathfrak K\)를 원소가 유한개인 체, \(\mathsf P\)를 그 중심이라 하자. \(\mathfrak K\)가 \(\mathsf P\) 위에서 갖는 차수를 \(t^2\)이라 하자. \(\mathfrak K\)의 모든 극대 가환부분체는 \(\mathsf P\) 위에서 차수 \(t\)이므로 원소 수가 같다. 유한체 이론의 잘 알려진 정리인 무어 정리에 따라 이 체들은 서로 동형이다. 모든 유한체가 그 소체의 갈루아확대이므로 이 동형들은 \(\mathsf P\)의 각 원소가 자기 자신에 대응하도록 고를 수 있다. 정리 7에 따라 \(\mathfrak K\)의 모든 극대 가환부분체는 한 \(\mathfrak Z\)에서 변환으로 얻어진다.

그런데 \(\mathfrak K\)는 그 모든 극대 가환부분체들의 합집합이다. \(\mathsf P\)에 원소 하나를 첨가하면 가환체가 생성되기 때문이다. 이제 0을 뺀 \(\mathfrak K\)의 곱셈군 \(\mathfrak K^*\)만을 보면, \(\mathfrak K^*\)는 부분군 \(\mathfrak Z^*\)와 켤레인 부분군들의 합집합이다. \(t>1\)이면 이는 모순이다. 실제로 \(\mathfrak K^*\)의 원소 수를 \(M\), \(\mathfrak Z^*\)의 원소 수를 \(N\), \(\mathfrak K^*\) 안에서 \(\mathfrak Z^*\)의 지수를 \(L\)이라 하면 \(M=NL\)이고 서로 다른 켤레부분군의 수는 \(L\) 이하이다. 그 수가 \(L\)과 같더라도 이 켤레부분군들이 군 전체를 모두 채우려면 임의의 두 부분군이 원소를 공유해서는 안 된다. 그러나 이들은 모두 단위원을 공유한다.

\subsection*{§ 21. 비가환체의 갈루아 이론}

\(\mathsf{G}\)를 그 중심 \(\mathsf P\)에 관한 \(\mathfrak K\)의 갈루아군이라 하자. 정리 8에 따라 \(G\)는
\[
        \overline G=\mathfrak K^*/\mathsf P^*
\]
와 동형이며, 이 동형을 고정해 둔다. \(G\)의 부분군 \(H\)에 동형 \(\mathsf{G}\cong\mathfrak K^*/\mathsf P^*\)로 대응하는 \(\mathfrak K^*\)의 부분군을 \(\mathfrak H^*\)라 하자. \(\mathfrak H^*\)에 \emph{0을 덧붙였을 때} 체가 되면 \(H\)를 \emph{닫힌} 부분군이라 한다.

\textbf{정리 11. 갈루아 이론의 주정리.} \emph{\(G\)의 닫힌 부분군 \(H\)들과 \(\mathsf P\)와 \(\mathfrak K\) 사이의 체 \(\mathfrak S\)들을 일대일로 대응시킬 수 있다. 서로 대응하는 \(H,\mathfrak S\)에 대해 \(\mathsf{H}\)는 \(\mathfrak S\)의 불변군이고 \(\mathfrak S\)는 \(\mathsf{H}\)의 불변체이다.}

\emph{증명.} 정리를 세 부분으로 나눈다.

\begin{enumerate}
\item 체 \(\mathfrak S\)의 불변군 \(H\)는 닫혀 있다.
\item \(\mathfrak S\)의 불변군이 \(\mathsf{H}\)이면 \(\mathfrak S\)는 \(\mathsf{H}\)의 불변체이다.
\item \(\mathsf{H}\)의 불변체가 \(\mathfrak S\)이면 \(\mathsf{H}\)는 \(\mathfrak S\)의 불변군이다.
\end{enumerate}

\textbf{1.} 체 \(\mathfrak S\)의 불변군 \(H\)에 동형 \(G\simeq\mathfrak K^*/\mathsf P^*\)로 대응하는 \(\mathfrak H^*\)는 \(\mathfrak S\)의 각 원소와 가환하는 0 아닌 모든 \(\mathfrak K\)-원소의 집합이다. 여기에 0을 덧붙이면 당연히 체가 되므로 \(H\)는 닫혀 있다.

\textbf{2.} 셋째 부분정리를 둘째 부분정리로 환원하자. 체 \(\mathfrak S\)에는 \(G\)의 부분군 \(S\), 곧 \(\mathfrak S^*/\mathsf P^*\)에 대응하는 \(\mathsf{G}\)의 부분군이 속한다. 마찬가지로 군 \(\mathsf{H}\)에는 \(\mathfrak K\)의 부분체 \(\mathfrak H\)가 속하여
\[
        H\cong\mathfrak H^*/\mathsf P^*
\]
이다. \(\mathfrak S\)가 \(\mathsf{H}\)의 불변체이면 \(\mathfrak S\)는 \(\mathfrak H\)의 모든 원소와 가환하는 모든 \(\mathfrak K\)-원소로 이루어진다. 달리 말하면 \(\mathsf{S}\)는 \(\mathfrak H\)의 불변군이다. 이에 대응하여 \(\mathsf{H}\)가 \(\mathfrak S\)의 불변군이면 \(\mathfrak H\)는 \(\mathsf{S}\)의 불변체이다. 따라서 셋째 부분정리는 다음처럼 정식화할 수도 있다.

\[
\text{3'. }\quad
\mathfrak H\text{의 불변군이 }S\text{이면 }
\mathfrak H\text{는 }S\text{의 불변체이다.}
\]
이는 바로 아직 증명해야 할 둘째 부분정리와 같은 형식이다.

\textbf{3.} 이 증명은 몇 가지 작은 차이를 제외하면 가환인 경우와 같다.

\(\mathfrak K\)와 반대동형인 체 \(\mathsf{K}\)를 만들고 그 중심을 \(\mathsf P\)와 동일시한다. \(\mathfrak K\)의 자기동형 대신 \(\mathfrak S\)에서 \(\mathsf{K}\)로 가는 반대동형들을 고찰할 수 있다.

가환인 경우에 쓴 연장정리와 비슷한 보조정리를 다시 증명하자.

\(\mathfrak S,\mathfrak T\)를 \(\mathsf P\)와 \(\mathfrak K\) 사이의 체들이라 하자.
\[
        \mathsf P\subseteqq\mathfrak S
        \subseteqq\mathfrak T\subseteqq\mathfrak K.
\]
\(\mathfrak S\)에서 \(\mathsf{K}\)의 한 부분체로 가는 모든 반대동형, 곧 \(\mathfrak S\)의 \(\mathsf{K}\) 안에서의 1차 반대 표현은 \(\mathfrak T\)의 반대동형으로 적어도 \(s=[\mathfrak T:\mathfrak S]\)가지 방식으로 연장될 수 있다고 주장한다.

체 \(\mathsf P,\mathfrak S,\mathfrak T,\mathfrak K\)를 \(K\)로 확대하면
\[
        K\subseteqq\mathfrak S_K
        \subseteqq\mathfrak T_K
        \subseteqq\mathfrak K_K.
\]
20쪽 정리 1에 따라 각 \(\mathfrak S_{\mathsf{K}}\)는 양쪽단순이고 완전가약이며, 서로 연산자동형인 단순 좌측아이디얼들은 \(\mathfrak S\)의 \(\mathsf{K}\) 안에서의 유일한 기약 표현류의 표현가군들이다. \(\mathsf{K}\)의 구성에 의해 \(\mathsf{K}\) 안에는 \(\mathfrak S\)와 반대동형인 체가 있으므로 이 표현류는 1차이다. 따라서 단순 좌측아이디얼들은 모두 랭크 1이고, \(\mathfrak S_K\)는 \(\mathsf P\) 위에서의 차수만큼의 좌측아이디얼들의 합이다.
\[
        \mathfrak S_K
        =K E_1+\cdots+K E_p,
        \qquad
        K E_\nu=\mathfrak S_K E_\nu.
\]

그러면
\[
\begin{aligned}
\mathfrak T_K&=\mathfrak T K
 =\mathfrak T\mathfrak S_K,\\
\mathfrak T_K
&=\mathfrak T\mathfrak S_K E_1+\cdots
  +\mathfrak T\mathfrak S_K E_p\\
&=\mathfrak T_{\mathsf{K}} E_1+\cdots+\mathfrak T_{\mathsf{K}} E_p.
\end{aligned}
\]
각 아이디얼 \(\mathfrak T_{\mathsf{K}}\cdot E_\nu\)는 \(s\)개의 단순아이디얼로 분해된다. 실제로
\[
\begin{aligned}
[\mathfrak T_K E_\nu:K]
&=[\mathfrak T_K E_\nu:K E_\nu]\\
&=[\mathfrak T_K E_\nu:\mathfrak S_K E_\nu]\\
&\leqq[\mathfrak T_K:\mathfrak S_K]
 =[\mathfrak T:\mathfrak S]=s
\end{aligned}
\]
이고
\[
        \sum [\mathfrak T_{\mathsf{K}} E_\nu:\mathsf{K}]=sp
\]
이므로
\[
        [\mathfrak T_{\mathsf{K}} E_\nu:\mathsf{K}]=s,
        \qquad
        \mathfrak T_K E_\nu
        =K e_1^{(\nu)}+\cdots+K e_s^{(\nu)}.
\]

\(\mathsf K e_i^{(\nu)}\)가 주는 \(\mathfrak T\)의 표현은 \(\mathsf K E_\nu\)가 주는 \(\mathfrak S\)의 표현을 연장한다. 실제로 \(\alpha\in\mathfrak S\)이면
\[
        \alpha E_\nu=\sigma E_\nu,
\]
\[
        \sum_i\alpha e_i^{(\nu)}
        =\sigma\sum_i e_i^{(\nu)},
        \qquad
        \alpha e_i^{(\nu)}
        =\sigma e_i^{(\nu)}.
\]

이제 \(\mathsf K e_i^{(\nu)}\)들이 주는 표현들이 모두 서로 다름만 보이면 된다. 가환인 경우에는 서로 다른 표현류에 속했으므로 명백했다. 여기서는 반대로 모두 같은 표현류에 속하므로 표현 자체를 살펴야 한다. \(\mathsf K e_i^{(\nu)}\)가 주는 \(\mathfrak T_{\mathsf{K}}\) 전체의 \(\mathsf{K}\)-준동형상들을 고찰하자. 기저원소들이 어떻게 표현되는지만 알면 전체 사상을 구성할 수 있으므로 이 사상들은 \(\mathfrak T\)의 표현들로 이미 완전히 정해진다. 따라서 \(\mathsf K e_i^{(\nu)}\)가 주는 \(\mathfrak T_{\mathsf{K}}\)의 사상들이 모두 다름만 보이면 충분하다. \(e_i^{(\nu)}\)를 한 번은 표현가군 \(\mathsf K e_i^{(\nu)}\)로, 다른 한 번은 \(\mathsf K e_j^{(\nu)}\)로 표현하면
\[
        e_i^{(\nu)2}=e_i^{(\nu)}
\]
이므로 첫째 경우에는 1이, 
\[
        e_i^{(\nu)}e_j^{(\nu)}=0
\]
이므로 둘째 경우에는 0이 \(e_i^{(\nu)}\)를 나타내는 원소로 나온다. 따라서 두 표현은 실제로 서로 다르다. 가환인 경우에도 똑같이 추론할 수 있지만 그곳에서는 필요하지 않다.

이제부터는 10쪽과 16쪽에서와 똑같이 추론한다. \(\mathsf{H}\)가 \(\mathfrak S\)의 불변군이고 \(\mathfrak T\)가 \(\mathsf{H}\)의 불변체라고 하자. \(\mathsf{H}\)의 원소들은 \(\mathfrak S\)의 항등자기동형을 \(\mathfrak K\) 전체로 연장한 모든 것들이다. 방금 증명한 보조정리에 따라 이들 가운데 \(\mathfrak T\) 위에서 서로 다른 것이 적어도 \(s\)개 있으므로 \(s=1\), \(\mathfrak T=\mathfrak S\)이어야 한다. 증명 끝.

\section*{제V장. 인자계}

\subsection*{§ 22. 중심이 주어진 체들의 군}

고정된 가환체 \(\mathsf P\)를 모든 논의의 바탕으로 삼는다. 자기동형체 \(\mathfrak K\)가 중심이 \(\mathsf P\)이고 \(\mathsf P\) 위에서 유한랭크인 체인 모든 행렬환 \(\mathfrak K_r\)의 전체는 \glqq{}직접곱\grqq{}에 관하여 곱셈적으로 닫혀 있다. 그러한 두 환 \(\mathfrak K_r,\mathfrak L_s\)의 직접곱은 \(\mathfrak K_r\)의 바탕체를 \(\mathfrak L_s\)로 확대하여 얻는 환 \(\mathfrak K_r\mathfrak L_s\)이다.
\[
        \mathfrak K_r\times\mathfrak L_s
        =\mathfrak K_r\cdot\mathfrak L_s.
\]
22쪽의 논평에 따라 이 직접곱은 가환이다.
\[
        \mathfrak K_r\times\mathfrak L_s
        =\mathfrak L_s\times\mathfrak K_r.
\]

모든 \(\mathfrak K_r\)의 계가 곱셈적으로 닫혀 있다는 주장은 \(\mathfrak K_r\times\mathfrak L_s\)가 다시 한 행렬환이고 그 자기동형체 \(\mathfrak M\)의 중심이 \(\mathsf P\)라는 뜻이다. \(\mathsf P\)에 관한 유한랭크임은 자명하다. 20쪽 정리 1에 따라 \(\mathfrak K_r\times\mathfrak L_s\)는 양쪽단순이고 완전가약이므로 \(\mathsf P\)의 한 확대체 \(\mathfrak M\) 위의 행렬환이다. \(\mathfrak M\)의 중심이 \(\mathsf P\)라는 것은 첫째 \(\mathsf P\)가 반드시 그 중심에 포함되고, 둘째 그 중심이 \(\mathfrak K\), 또는 \(\mathfrak L\)의 중심 \(\mathsf P\)에 포함되어야 한다는 데서 곧바로 따른다.

곱셈 \(\mathfrak K_r\times\mathfrak L_s\)는 명백히 결합적이다. \(\mathfrak K_r\)들은 곱셈적으로 닫힌 계이지만 군은 아니다. \(\mathfrak K_r\times\mathfrak L_s=\mathfrak M_m\)이면 반드시 \(m\geqq rs\)이므로, \(r>m\)일 때 등식 \(\mathfrak K_r\times\mathfrak L_s=\mathfrak M_m\)을 만족하는 \(\mathfrak L_s\)는 찾을 수 없기 때문이다.

이제 같은 \(\mathfrak K\)를 갖는 모든 \(\mathfrak K_r\)을 한 류 \(\{\mathfrak K\}\)로 묶자. \(\mathfrak K_r,\mathfrak K_{r'}\)이 같은 류 \(\{\mathfrak K\}\)에 속하고 \(\mathfrak L_s,\mathfrak L_{s'}\)이 같은 류 \(\{\mathfrak L\}\)에 속하면, \(\mathfrak K_r\times\mathfrak L_s\)와 \(\mathfrak K_{r'}\times\mathfrak L_{s'}\)도 같은 류에 속한다.

\(\mathfrak K_r\times\mathfrak L_s\)가 \(\mathfrak K\times\mathfrak L\)과 같은 류에 속함을 보이면 충분하다. 곧 \(\mathfrak K\times\mathfrak L=\mathfrak M_m\)이면 \(\mathfrak K_r\times\mathfrak L_s\)도 \(\mathfrak M\) 위의 행렬환임을 보이면 된다. 이는 쉽게 알 수 있다.
\[
\begin{aligned}
\mathfrak K_r\times\mathfrak L
&=\left(\sum_1^r\mathfrak Kc_{ik}\right)\mathfrak L\\
&=\sum_1^r\mathfrak K_{\mathfrak L}c_{ik}
 =\sum_1^r\mathfrak M_m c_{ik}.
\end{aligned}
\]
이는 \(\mathfrak M_m\) 위의 \(r\)차 행렬환이고, 따라서 \(\mathfrak M\) 위의 \(mr\)차 행렬환이다. 똑같이 계속하면
\[
\begin{aligned}
\mathfrak K_r\times\mathfrak L_s
&=\left(\sum_1^s\mathfrak Ld_{ik}\right)_{\mathfrak K_r}\\
&=\sum_1^s\mathfrak L_{\mathfrak K_r}d_{ik}
 =\sum_1^s\mathfrak M_{mr}d_{ik}
 =\mathfrak M_{mrs}.
\end{aligned}
\]
이로써 주장이 증명되었다. 특히
\[
        \mathfrak L\times\mathfrak K=\mathfrak M_m
        \quad\Longrightarrow\quad
        \mathfrak K_r\times\mathfrak L_s
        =\mathfrak M_{mrs}
\]
임을 기억해 두자.

따라서 두 류 \(\{\mathfrak K\}\), \(\{\mathfrak L\}\)의 곱을 각 류의 대표들의 곱이 속하는 류로 정의하면, 이 정의는 대표의 선택에 의존하지 않는다.

그러므로 사상
\[
        \mathfrak K_r\longmapsto\{\mathfrak K\}
\]
은 모든 \(\mathfrak K_r\)의 계에서 모든 류 \(\{\mathfrak K\}\)의 집합으로 가는 준동형이다. 따라서 류들의 집합 \(\mathscr K\)는 곱셈적으로 닫힌 계이다. 이제 \(\mathscr K\)가 군임을 증명하자. 단위원은 명백히 류 \(\{\mathsf P\}\)이다.
\[
        \mathfrak K_r\times\mathsf P=\mathfrak K_r.
\]

따라서 각 류 \(\{\mathfrak K\}\)에 역류 \(\{\mathfrak K\}^{-1}\)가 존재함만 보이면 된다. \(\overline{\mathfrak K}\)가 \(\mathfrak K\)와 반대동형인 체를 뜻할 때
\[
        \{\mathfrak K\}^{-1}
        =\{\overline{\mathfrak K}\}
\]
로 놓아야 한다. 달리 말하면 \(\mathfrak K\times\overline{\mathfrak K}\)가 \(\mathsf P\) 위의 행렬환이어야 한다. \(\mathfrak K\)가 \(\mathsf P\) 위에서 랭크 \(t^2\), 따라서 \(\mathfrak K\times\overline{\mathfrak K}\)가 랭크 \(t^4\)이면 이 곱은 \(t^2\)개의 단순 좌측아이디얼로 분해되어야 한다. 실제로 \(\mathfrak K\times\overline{\mathfrak K}\)의 한 단순 좌측아이디얼은 \(\overline{\mathfrak K}\)의 \(\mathfrak K\) 안에서의 1차 기약 반대 표현의 표현가군으로서 \(\mathfrak K\) 위에서 랭크 1, 따라서 \(\mathsf P\) 위에서 랭크 \(t^2\)이다.

\(\mathsf P\)와 \(\mathfrak K\) 사이에 중심이 역시 \(\mathsf P\)인 체가 없으면 \(\mathfrak K\)를 \emph{단순체}라 부르겠다.

\textbf{정리.} \emph{모든 체 \(\mathfrak K\)는 단순 부분체들의 직접곱이다.}
\[
        \mathfrak K
        =\mathfrak K^{(1)}\times\cdots
        \times\mathfrak K^{(p)}.
\]

\emph{증명.} 중심이 \(\mathsf P\)인 \(\mathfrak K\)의 부분체 \(\mathfrak K^{(1)}\)이 주어지면 중심이 \(\mathsf P\)인 \(\mathfrak K\)의 부분체 \(\mathfrak S\)로서
\[
        \mathfrak K^{(1)}\times\mathfrak S
        =\mathfrak K
\]
을 만족하는 것이 있음을 보이면 충분하다. 적어도
\[
        \{\mathfrak K^{(1)}\}
        \times\{\mathfrak S\}
        =\{\mathfrak K\}
\]
을 만족하는 류 \(\{\mathfrak S\}\)는 존재한다. 곧
\[
        \{\mathfrak S\}
        =\{\overline{\mathfrak K^{(1)}}\}
         \times\{\mathfrak K\}.
\]
\[
        \overline{\mathfrak K^{(1)}}\times\mathfrak K
        =\mathfrak S_\lambda
\]
라 하자. 그러면
\[
\begin{aligned}
\mathfrak K^{(1)}\times\mathfrak S_\lambda
&=\mathfrak K^{(1)}
  \times\overline{\mathfrak K^{(1)}}\times\mathfrak K\\
&=\mathsf P_{t^2}\times\mathfrak K
 =\mathfrak K_{t^2},
\qquad
t^2=\operatorname{rank}_{\mathsf P}\mathfrak K^{(1)}.
\end{aligned}
\]
\(\mathfrak K^{(1)}\times\mathfrak S_\lambda\)는 차수가 적어도 \(\lambda\)인 행렬환이므로 \(\lambda\leqq t^2\)이다. 한편
\[
        \mathfrak S_\lambda
        =\overline{\mathfrak K^{(1)}}\times\mathfrak K
\]
는 \(\overline{\mathfrak K^{(1)}}\)가 \(\mathfrak K\) 안에서 1차 반대 표현을 허용하므로 \(\mathfrak K\) 위에서 랭크 1인 단순아이디얼들로 분해되어야 한다. 따라서 \(\lambda\geqq t^2\), 곧 \(\lambda=t^2\)이고
\[
        \mathfrak K_{t^2}
        =\mathfrak K^{(1)}\times\mathfrak S_{t^2}.
\]
그러므로 \(\mathfrak K^{(1)}\times\mathfrak S\) 자체가 체이어야 하고
\[
\begin{aligned}
\mathfrak K^{(1)}\times\mathfrak S_{t^2}
&=\bigl(\mathfrak K^{(1)}
  \times\mathfrak S\bigr)_{t^2}\\
&=\mathfrak K_{t^2}.
\end{aligned}
\]

\(\mathfrak K^{(1)}\times\mathfrak S\)는 \(\mathfrak K_{t^2}\)의 한 분해에서 단순 좌측아이디얼들의 자기동형체이고, \(\mathfrak K\)는 다른 한 분해에서 단순 좌측아이디얼들의 자기동형체이다. \(\mathfrak K_{t^2}\)의 서로 다른 임의의 두 좌측아이디얼을 한 분해에 함께 넣을 수 있으므로 \(\mathfrak K\)와 \(\mathfrak K^{(1)}\times\mathfrak S\)는 동형이며, 이 동형은 \(\mathsf P\)의 각 원소를 고정하도록 고를 수 있다. 따라서 25쪽 정리 7에 따라 이 동형은 \(\mathfrak K_{t^2}\)의 한 내부자기동형이 생성한다.
\[
\begin{aligned}
\mathfrak K
&=\lambda^{-1}
  \bigl(\mathfrak K^{(1)}\times\mathfrak S\bigr)\lambda\\
&=\lambda^{-1}\mathfrak K^{(1)}\lambda
  \times\lambda^{-1}\mathfrak S\lambda,
\qquad
\lambda\in\mathfrak K_{t^2}.
\end{aligned}
\]
\(\lambda^{-1}\mathfrak K^{(1)}\lambda\)는 \(\mathfrak K^{(1)}\)과 동형인 \(\mathfrak K\)의 부분체이다. 그러므로 다시 25쪽 정리 7에 따라 알맞은 \(\mu\in\mathfrak K\)가 있어
\[
        \mu^{-1}\lambda^{-1}
        \mathfrak K^{(1)}\lambda\mu
        =\mathfrak K^{(1)}
\]
이고, 따라서
\[
\begin{aligned}
\mu^{-1}\mathfrak K\mu
=\mathfrak K
&=\mathfrak K^{(1)}
  \times\mu^{-1}\lambda^{-1}
  \mathfrak S\lambda\mu.
\end{aligned}
\]
이로써 주장이 증명되었다.

\subsection*{§ 23. 인자계}

이하의 고찰에서는 바탕체 \(\mathsf P\)가 완전체라고 가정한다. 그 대신 다음과 같은 더 일반적인 가정을 둘 수도 있다. 곧 중심이 \(\mathsf P\)인 비가환체 \(\mathfrak K\) 가운데, 모든 \(\mathfrak K_r\)가 \(\mathsf P\) 위에서 제1종인 극대 가환부분체를 포함하는 것만을 고찰한다. (그 뒤 쾨테가 보인 바에 따르면, 이는 모든 \(\mathfrak K\)에 대하여 성립한다.)

중심이 \(\mathsf P\)인 비가환체 \(\mathfrak K\)의 한 분해체를 \(Z\)라 하자. \(Z\)의 \(\mathfrak K\) 안에서의 기약표현 \(\mathfrak Z\)가 차수 \(r\)이라고 하자. 그러면 \(\mathfrak Z\)는 \(\mathfrak K_r\)의 극대 가환부분체이다. \(\Gamma\)를 \(\mathsf{Z}\)의 갈루아체라 하면
\[
        \mathfrak Z_\Gamma
        =e_1\Gamma+\cdots+e_n\Gamma
        \qquad\text{(9쪽)}
\]
이다. 여기서 \(n\)은 \(\mathsf{Z}\)의 \(\mathsf P\) 위에서의 차수이다. 따라서 \(\Gamma\)로 확대된 환 \(\mathfrak K_r\)에도 분해
\[
        \mathfrak K_{r\Gamma}
        =\mathfrak K_{r\Gamma}e_1
         +\cdots+
         \mathfrak K_{r\Gamma}e_n
\]
가 있고, 곧
\[
        \mathfrak K_{r\Gamma}
        =\mathfrak l_1+\cdots+\mathfrak l_n,
        \qquad
        \mathfrak l_i=\mathfrak K_{r\Gamma}e_i
\]
이다. \(\mathfrak l_i\)들은 명백히 좌측아이디얼이며, 사실 단순 좌측아이디얼이다. 실제로 \(\mathsf{Z}\)가 분해체이므로 \(\mathfrak K_{r\Gamma}\)는 \(\Gamma\) 위의 차수 \(r\cdot t=n\)인 행렬환이다. (23쪽 정리 4 및 22쪽의 \(n=rt\).) 이렇게 얻은 \(\mathfrak K_{r\Gamma}\)의 단순 좌측아이디얼들로의 분해는 \(\mathfrak K\)와 \(\mathsf{Z}\)의 대수적 성질에 특별히 맞추어져 있다. 실제로 18쪽에서처럼 \(\Gamma/\mathsf P\)의 갈루아군 \(\mathfrak G\)의 치환 \(S\)가 모든 \(\mathfrak K_{r\Gamma}\)-원소에 작용하는 방식을, 각 \(\mathfrak K_r\)-원소 \(z\)에 대하여
\[
        S(z)=z
\]
로 정의하면 아이디얼 \(\mathfrak l_i\)들은 서로 켤레이다. 즉 \(S\)를 어떤 \(\mathfrak l_j\)에 적용하면 다른 하나가 되어
\[
        S(\mathfrak l_j)=\mathfrak l_i
\]
이다. 단위원의 각 비분해 성분 \(e_j\)가 다른 하나 \(e_i\)로 옮겨가야 하기 때문이다. (11쪽도 보라.)

표현가군 \(e_i\Gamma\)에 의해 \(\mathfrak Z\)는 \(\Gamma\)의 한 부분체 \(\mathsf{Z}_i\) 위로 사상된다. \(\mathsf{Z}_i\)에 의한 \(\mathfrak Z\)의 표현이 이미 \(\mathsf{Z}_i\) 안에 들어 있으므로 \(\mathfrak Z_{\mathsf{Z}_i}\)는 표현가군 \(e_i\mathsf{Z}_i\)를 분리해 내야 한다. 따라서 \(e_i\)는 이미 \(\mathfrak Z_{\mathsf{Z}_i}\) 안에 있다. \(\mathsf{Z}_i\)와 \(\mathsf{Z}_k\)의 합성체를 \(\{\mathsf{Z}_i,\mathsf{Z}_k\}\)라 하자. \(\mathfrak K_{r\{\mathsf{Z}_i,\mathsf{Z}_k\}}\)는 두 좌측아이디얼
\[
        \bar{\mathfrak l}_i
        =\mathfrak K_{r\{\mathsf{Z}_i,\mathsf{Z}_k\}}e_i,
        \qquad
        \bar{\mathfrak l}_k
        =\mathfrak K_{r\{\mathsf{Z}_i,\mathsf{Z}_k\}}e_k
\]
를 분리해 내며, 이들은 단순하다. 왜냐하면
\(
 \mathfrak l_i
 =\mathfrak K_{r\Gamma}\bar{\mathfrak l}_i
\)
이기 때문이다. 그러므로
\[
\begin{aligned}
\bar{\mathfrak l}_i
 &=\sum_{\lambda=1}^{m}
   \{Z_i,Z_k\}c_{\lambda i}^{(ik)},\\
\bar{\mathfrak l}_k
 &=\sum_{\lambda}
   \{\mathsf{Z}_i,\mathsf{Z}_k\}c_{\lambda k}^{(ik)}.
\end{aligned}
\]
여기서 \(c_{\lambda\varrho}^{(ik)}\)는 \(\mathfrak K_{r\{\mathsf{Z}_i,\mathsf{Z}_k\}}\)의 한 행렬단위계이다.

\(\mathfrak K_{r\{\mathsf{Z}_i,\mathsf{Z}_k\}}\)를 연산자 영역으로 하는 \(\bar{\mathfrak l}_i\)에서 \(\bar{\mathfrak l}_k\) 위로의 각 동형사상 \(a_i\mapsto a_k\)는 \(e_i\mapsto p_{ik}\)인 원소 \(p_{ik}\)로 특징지어진다. 실제로 이때
\[
        a_i=a_ie_i\mapsto a_ip_{ik},
        \qquad
        \bar{\mathfrak l}_i\cdot p_{ik}=\bar{\mathfrak l}_k
\]
이다. 가능한 모든 그러한 \(p_{ik}\)는 명백히
\[
        p_{ik}=c_{ik}^{(ik)}\gamma_{ik},
        \qquad
        0\ne\gamma_{ik}\in\{Z_i,Z_k\}
\]
이다. 따라서 \(a_i\mapsto a_ip_{ik}\)는 \(\mathfrak l_i\)에서 \(\mathfrak l_k\) 위로의 연산자동형도 준다. 이제 그러한 \(p_{ik}\)들의 특정한 계를 켤레조건을 부과하여 선택한다.

치환 \(S\)가 첨자 \(i\)에 작용하는 방식을
\[
        S(i)=k,
        \qquad\text{만일}\qquad
        S(Z_i)=Z_k
\]
로 정의하자. 각 첨자쌍 \((\nu,\mu)\)에는 \(S(\nu,\mu)=(1,k)\) 꼴인 켤레쌍이 적어도 하나 있다. 켤레 첨자쌍의 각 류에서 고정된 한 쌍 \((1,k)\)를 고르고, 이 \((1,k)\)에 대하여
\[
        p_{1k}=c_{1k}^{(1k)}\gamma_{1k},
        \qquad
        0\ne\gamma_{1k}\in\{Z_1,Z_k\}
\]
를 그 밖에는 임의로 정한다. 이어서 \(S(1,k)=(\nu,\mu)\)이면
\[
        p_{\nu\mu}=S(p_{1k})
\]
로 둔다. 그러면 실제로 \(p_{\nu\mu}\)는
\[
        p_{\nu\mu}
        =c_{\nu\mu}^{(\nu\mu)}\gamma_{\nu\mu},
        \qquad
        0\ne\gamma_{\nu\mu}
        \in\{Z_\nu,Z_\mu\}
\]
꼴이다. 왜냐하면 \(S\)를 동형사상
\[
        a_1\mapsto a_1p_{1k},
        \qquad\text{특히}\qquad
        e_1\mapsto p_{1k},
        \qquad
        \bar{\mathfrak l}_1\longrightarrow
        \bar{\mathfrak l}_k
\]
에 작용시키면 동형사상
\[
        a_\nu\mapsto a_\nu p_{\nu\mu},
        \qquad\text{특히}\qquad
        e_\nu\mapsto p_{\nu\mu},
        \qquad
        \bar{\mathfrak l}_\nu\longrightarrow
        \bar{\mathfrak l}_\mu
\]
를 얻기 때문이다.

\(c_{\lambda\varrho}\)를 \(\mathfrak K_{r\Gamma}\)의 한 행렬단위계라 하면 \(p_{\nu\mu}\)를
\[
        p_{\nu\mu}=c_{\nu\mu}\bar\gamma_{\nu\mu}
\]
꼴로도 쓸 수 있음은 분명하다. 따라서 관계식
\[
        c_{ij}c_{1k}=0\quad(j\ne1),
        \qquad
        c_{ij}c_{jk}=c_{ik}
\]
에서 다음 관계식들이 따라온다.
\[
\begin{aligned}
p_{ij}p_{1k}&=0,
   &&\text{만일 }j\ne1,\\
p_{ij}p_{jk}&=\alpha_{ik}^{(j)}p_{ik},
   &&\alpha_{ik}^{(j)}\in\Gamma
     \quad\text{(더 정확히는 }\{Z_i,Z_j,Z_k\}\text{ 안에 있다).}
\end{aligned}
\]
이 \(n^3\)개의 양 \(\alpha\)는 유사행렬단위계 \(p_{ik}\)에 속하는, \emph{\(\mathsf{Z}\)를 바탕으로 한 \(\mathfrak K\)의 인자계}를 이룬다.

이 \(\alpha\)들은 서로 켤레이다. 곧
\[
        S(i,k,j)=(\nu,\mu,\tau)
\]
이면
\[
        S\bigl(\alpha_{ik}^{(j)}\bigr)
        =\alpha_{\nu\mu}^{\tau}
\]
이다.

유사행렬단위계 \(p_{ik}^*\)는 모두, 처음의 \(p_{1k}\)들에 \(\{Z_1,Z_k\}\)의 임의의 \(\delta_{1k}\ne0\)를 곱하여
\[
        p_{1k}^*=p_{1k}\delta_{1k}
\]
로 둔 다음, 다시 \(S(1,k)=(\nu,\mu)\)일 때
\[
        p_{\nu\mu}^*=S(p_{1k}^*)
\]
로 놓음으로써 한 유사행렬단위계 \(p_{ik}\)에서 얻어진다. 따라서
\[
        p_{\nu\mu}^*=p_{\nu\mu}\delta_{\nu\mu},
        \qquad
        \delta_{\nu\mu}\in\{Z_\nu,Z_\mu\}
\]
이다. \(\delta_{\nu\mu}\)들은 서로 켤레이며,
\[
        S(i,k)=(\nu,\mu)
        \quad\Longrightarrow\quad
        S(\delta_{ik})=\delta_{\nu\mu}.
\]
\(p_{ik}^*\)에 속하는 인자계를 \(\alpha^*\)라 하면
\[
        \alpha_{ik}^{(j)*}
        =\frac{\delta_{ij}\delta_{jk}}{\delta_{ik}}
         \alpha_{ik}^{(j)}
\]
이다.

인자계 \(\alpha^*\)와 \(\alpha\)를 \emph{연관되었다}고 한다. \(p_{ik}^*\)와 \(p_{ik}\)도 마찬가지이다. \(\mathsf{Z}\)를 바탕으로 한 \(\mathfrak K\)의 모든 인자계는 연관된 인자계들의 한 류 \(\{\alpha\}\)를 이룬다.

이 류 \(\{\alpha\}\)는 그 정의의 바탕으로 삼은 \(\mathfrak K_r\) 안에서의 \(\mathsf{Z}\)의 표현 \(\mathfrak Z\)에 의존하지 않는다. 실제로 한 \(\mathfrak Z\)에서 다른 \(\mathfrak Z'\)로의 이행은 \(\mathfrak K_r\)의 한 원소, 따라서 \(\mathfrak G\)-불변인 원소 \(\lambda\)에 의한 변환으로 할 수 있다. 그러므로 \(\mathfrak Z\)에 속하는 계 \(p_{ik}\)에서 \(\mathfrak Z'\)에 속하는 계
\[
        p'_{ik}=\lambda^{-1}p_{ik}\lambda
\]
가 생긴다. 그러나 \(p_{ik}\)와 \(p'_{ik}\)의 인자계는 서로 동일하다.

인자계란 \(\mathfrak K_{n\Gamma}\)의 곱셈표, 다만 \(\mathfrak K\)의 특별한 대수적 성질에 맞추어진 곱셈표에 지나지 않는다.

\subsection*{§ 24. 인자계의 곱셈}

비가환체 \(\mathfrak K\)의 분해체라는 말 대신 \(\{\mathfrak K\}\)의 분해체라고도 한다. 이것은 뜻이 있다. \(\mathfrak K_{\mathsf{Z}}\)와 더불어 모든 환 \(\mathfrak K_{r\mathsf{Z}}\)도 절대기약 성분들로 분해되고, 그 역도 성립하기 때문이다.

중심이 \(\mathsf P\)인 임의의 두 체 \(\mathfrak K,\mathfrak S\)에는 공통 분해체가 있다. 예를 들어 \(\mathfrak K\)의 한 분해체와 \(\mathfrak S\)의 한 분해체의 합성체가 그러하다.

\textbf{정리 1.} \emph{체 \(Z\)를 분해체로 갖는 류 \(\{\mathfrak K\}\)들은 모든 \(\{\mathfrak K\}\)의 군 \(\mathscr K\)의 한 부분군 \(\mathscr K_{\mathsf{Z}}\)를 이룬다.}

\emph{증명.} \(\mathfrak K\)와 \(\mathfrak S\)가 \(\mathsf{Z}\)를 분해체로 가지면 \(\mathfrak K_{\mathsf{Z}}\)와 \(\mathfrak S_{\mathsf{Z}}\)가 분해되므로 \((\mathfrak K\times\mathfrak S)_{\mathsf{Z}}\)도 절대기약 성분들로 분해된다. 즉 \(\mathsf{Z}\)는 \(\{\mathfrak K\times\mathfrak S\}\)의 분해체이다. (23쪽 정리 4.)

서로 반대동형인 체들은 모든 분해체를 공유한다.

이제 보조 고찰을 하자. 단순 좌측아이디얼들로의 분해
\[
        \mathfrak K_r
        =\mathfrak L_1+\cdots+\mathfrak L_r,
        \qquad
        \mathsf P_m
        =\mathfrak z_1+\cdots+\mathfrak z_m
\]
가 주어졌다고 하자. 그러면
\[
\begin{aligned}
\mathfrak K_{rm}
 &=\mathfrak K_r\times\mathsf P_m\\
 &=\mathfrak L_1\mathfrak z_1+
   \cdots+
   \mathfrak L_i\mathfrak z_j+
   \cdots+
   \mathfrak L_r\mathfrak z_m
\end{aligned}
\]
은 \(\mathfrak K_{rm}\)을 \(0\)이 아닌 \(rm\)개의, 따라서 단순한 좌측아이디얼들로 분해한 것이다. 그러므로 \(\mathfrak K_{rm}\)의 길이가 \(r\)인 각 좌측아이디얼 \(\mathsf{A}\)는 길이가 \(r\)인 좌측아이디얼
\[
        \mathfrak L_1\mathfrak z_1+
        \cdots+
        \mathfrak L_r\mathfrak z_1
        =\mathfrak K_r\mathfrak z_1
\]
과 한 내부자기동형으로 연결된다. 곧
\[
\begin{aligned}
A
 &=\tau^{-1}\mathfrak K_r\tau\,
   \tau^{-1}\mathfrak z_1\tau\\
 &=\mathfrak K'_r\cdot\mathfrak z'_1.
\end{aligned}
\]

\(\mathfrak Z\subset\mathfrak K_r\)에 속하는, 켤레인 좌측아이디얼들로의 분해
\[
        \mathfrak K_{r\Gamma}
        =\mathfrak l_1+\cdots+\mathfrak l_n,
        \qquad
        \mathfrak l_i=\mathfrak K_{r\Gamma}e_i
\]
가 주어졌다고 하자. 그러면
\[
\begin{aligned}
S(\tau^{-1}\mathfrak l_i\tau\cdot\,\mathfrak z'_1)
 &=S(\tau^{-1}\mathfrak l_i\tau)\,\mathfrak z'_1\\
 &=\tau^{-1}(S\mathfrak l_i)\tau\cdot\,\mathfrak z'_1\\
 &=\tau^{-1}\mathfrak l_k\tau\cdot\,\mathfrak z'_1
\end{aligned}
\]
이므로
\[
\tag{1}
        A_\Gamma
        =\mathfrak l'_1\mathfrak z'_1+
         \cdots+
         \mathfrak l'_n\mathfrak z'_1
\]
은 \(\mathsf{A}_\Gamma\)를 켤레인 좌측아이디얼들로 분해한 것이다. \(\tau^{-1}e_i\tau=e'_i\)라 두고 \(E'\)가 \(\mathfrak z'_1\)의 단위원을 뜻한다고 하면
\[
        \mathfrak l'_i\mathfrak z'_1
        =\mathfrak K_{rm\Gamma}e'_iE'
\]
이다.

켤레인 좌측아이디얼들로의 다른 임의의 분해
\[
\tag{2}
        \mathsf{A}_\Gamma=\mathfrak q_1+\cdots+\mathfrak q_n,\qquad,
        \qquad
        S\mathfrak q_i=\mathfrak q_k
        \quad\text{만일}\quad
        S(\mathfrak l'_i\mathfrak z'_1)
        =\mathfrak l'_k\mathfrak z'_1
\]
는 (1)에서 한 \(\mathfrak K_{rm}\)-원소 \(\lambda\)에 의한 변환으로 생긴다.
\[
        \mathfrak q_i
        =\lambda^{-1}\mathfrak l'_i\mathfrak z'_1\lambda.
\]
따라서 \(e_i^*\)가 \(\mathfrak q_i\)의 단위원이면
\[
        e_i^*=\lambda^{-1}e'_iE'\lambda
\]
이다.

이 주장은 다음 보조정리에서 곧바로 따라온다.

\textbf{보조정리.} \emph{\(\mathsf{A}\)를 \(\mathfrak K_n\)의 한 좌측아이디얼이라 하고, 가정상 존재하는 분해}
\[
        A
        =\mathfrak l_1+\cdots+\mathfrak l_n
        =\mathfrak K_{n\Gamma}e_1+
         \cdots+
         \mathfrak K_{n\Gamma}e_n
\]
\emph{이 켤레인 좌측아이디얼들로의 분해라고 하자. 여기서 일반성을 잃지 않고 \(e_i\)들을 서로 켤레라고 가정할 수 있다.}

\emph{\(\mathfrak z\)가 \(\mathsf{A}\)와 연산자동형이고 같은 가정들이 성립한다고 하자.}
\[
        \mathfrak z
        =\mathfrak q_1+\cdots+\mathfrak q_n
        =\mathfrak K_{n\Gamma}e_1^*+
         \cdots+
         \mathfrak K_{n\Gamma}e_n^*.
\]
\emph{따라서 \(\mathfrak q_i\)들과 \(e_i^*\)들은 서로 켤레이다. 그러면 \(\mathfrak K_n\)의 한 가역원 \(\lambda\)가 존재하여}
\[
        \mathfrak q_i=\lambda^{-1}\mathfrak l_i\lambda
\]
\emph{이다.}

\emph{증명.} 1. \emph{\(e_i\)들을 서로 켤레라고 가정할 수 있다.} 실제로 \(e_1\)이 \(\mathfrak l_1\)의 임의의 생성 멱등원, 즉 \(\mathfrak l_1=\mathfrak K_{n\Gamma}e_1\)이라고 하자. \(\Gamma\)의 한 자기동형을 적용하면 \(\mathfrak l_i=\mathfrak K_{n\Gamma}e_i\)를 얻는데, 이때 \(e_i\)는 \(e_1\)과 켤레이며 다시 생성 멱등원이다.

2. \emph{각 \(\mathfrak l\)은 각 \(\mathfrak q\)와 동형이다.} 실제로 서로 켤레인 \(\mathfrak l\)들은 모두 같은 절대길이를 갖고, \(\mathfrak z\)와 \(\mathsf{A}\)가 동형이라는 가정 아래 \(\mathfrak l\)들의 수와 \(\mathfrak q\)들의 수가 일치한다.

3. 대응
\[
\tag{3}
\begin{gathered}
e_1\mapsto s_1=e_1s_1,
\qquad
e_1^*\mapsto t_1=e_1^*t_1,\\
s_1t_1=e_1,
\qquad
t_1s_1=e_1^*
\end{gathered}
\]
으로 \(\mathfrak l_1\)과 \(\mathfrak q_1\) 사이의 한 동형사상과 그 역을 고정하자. 모든 자기동형을 적용하면 (3)은
\[
\begin{gathered}
e_i\mapsto s_i=e_is_i,
\qquad
e_i^*\mapsto t_i=e_i^*t_i,\\
s_it_i=e_i,
\qquad
t_is_i=e_i^*
\end{gathered}
\]
로 옮겨간다. 이로써 동형사상
\[
        \mathfrak q_i=\mathfrak l_is_i,
        \qquad
        \mathfrak l_i=\mathfrak q_it_i
\]
가 고정된다.

4. 이제
\[
        s=s_1+\cdots+s_n,
        \qquad
        t=t_1+\cdots+t_n
\]
로 두면
\[
\begin{gathered}
\mathfrak q_i=\mathfrak l_is,
\qquad
\mathfrak l_i=\mathfrak q_it,\\
st=e_1+\cdots+e_n=\xi,
\qquad
ts=e_1^*+\cdots+e_n^*=\xi^*
\end{gathered}
\]
를 얻는다. 여기서 \(\xi\)와 \(\xi^*\)는 \(\mathfrak K_n\)에 속하며
\[
        A\xi=A,
        \qquad
        \mathfrak z\xi^*=\mathfrak z
        \quad
        (\mathfrak l_i\xi
          =\mathfrak l_i\mathfrak l_i
          =\mathfrak l_i\text{이므로})
\]
이고
\[
        \xi^2=\xi,
        \qquad
        \xi^{*2}=\xi^*
\]
이다. 이제
\[
        e=\overline E+E=\overline E^*+E^*
\]
로 두어
\[
\begin{aligned}
\mathfrak K_n
 &=\overline A+A
  =\mathfrak K_n\overline\xi+\mathfrak K_n\xi\\
 &=\overline{\mathfrak z}+\mathfrak z
  =\mathfrak K_n\overline\xi^*+\mathfrak K_n\xi^*
\end{aligned}
\]
라고 하면 \(\overline{\mathsf A}\)와 \(\overline{\mathfrak z}\)도 동형이 된다. 따라서
\[
        \overline E\mapsto
        \overline s=\overline E\,\overline s,
        \qquad
        \overline E^*\mapsto
        \overline t=\overline E^*\overline t
\]
가 서로 역인 동형사상이고
\[
        \lambda=s+\overline s,
        \qquad
        \mu=t+\overline t
\]
로 두면
\[
\begin{gathered}
\lambda\mu=\mu\lambda=e,
\qquad
\mathfrak l_i\lambda=\mathfrak q_i,
\qquad
\lambda^{-1}\mathfrak q_i=\mathfrak q_i,\\
\lambda^{-1}\mathfrak q_i\subseteqq\mathfrak q_i
=\lambda\cdot\lambda^{-1}\mathfrak q_i
\subseteqq\lambda^{-1}\mathfrak q_i.
\end{gathered}
\]
따라서
\[
        \lambda^{-1}\mathfrak l_i\lambda=\mathfrak q_i,
\]
증명 끝.

분해 (1)에 대하여 다시 원소
\[
        r_{ik}=c_{ik}^{\prime(ik)}\gamma_{ik},
        \qquad
        \gamma_{ik}\in\{Z_i,Z_k\}
\]
를 정할 수 있다. 이들은 \(e'_iE'\mapsto r_{ik}\)에 의해 아이디얼 \(\mathfrak l'_i\mathfrak z'_1\)과 \(\mathfrak l'_k\mathfrak z'_1\) 사이의 연산자동형을 주고, \(S(i,k)=(\nu,\mu)\)이면
\[
        S(r_{ik})=r_{\nu\mu}
\]
인 켤레조건도 만족한다. 이 원소들은
\[
        r_{ij}r_{jk}=\alpha_{ik}^{(j)}r_{ik}
\]
에 의해 \(\mathsf{A}_\Gamma\)의 한 인자계 \(\alpha\)를 정한다. 이제 중요한 점은 \(\mathsf{A}_\Gamma\)의 인자계를 정하는 데 임의의 분해 (2)를 똑같이 쓸 수 있다는 것이다. 실제로 (1)에 속하는 계 \(r_{ik}\)를
\[
        r'_{ik}=\lambda^{-1}r_{ik}\lambda
\]
로 바꾸면 \(e_i^*\mapsto r'_{ik}\)는 \(\mathfrak q_i\)에서 \(\mathfrak q_k\) 위로의 연산자동형을 준다. \(\lambda\)가 갈루아군에 불변이므로 \(r'_{ik}\)들은 켤레조건을 만족하며, \(r'_{ik}\)들은 \(r_{ik}\)들과 같은 인자계를 갖는다.

\(\mathfrak K\)의 유사행렬단위계 \(p_{ik}\)를 택하면
\[
        r_{ik}=\tau^{-1}p_{ik}\tau\cdot\,E'
\]
은 같은 인자계를 갖는 \(\mathsf{A}_\Gamma\)의 한 \(r_{ik}\)-계이다. 따라서 다음 정리를 얻는다.

\textbf{정리.} \emph{\(\mathsf{A}_\Gamma\)를 켤레인 좌측아이디얼들로 분해한 임의의 분해}
\[
        \mathsf{A}_\Gamma=\mathfrak q_1+\cdots+\mathfrak q_n
\]
\emph{을 사용하여 만든 \(\mathsf{A}_\Gamma\)의 인자계들은 \(\mathfrak K_r\)의 인자계들과 같다.}

다음 정리는 \(Z\)를 바탕으로 한 \(\{\mathfrak K\}\), \(\{\mathfrak S\}\)의 인자계가 각각 \(\alpha_{ik}^{(j)}\), \(\beta_{ik}^{(j)}\)이면
\[
        \varepsilon_{ik}^{(j)}
        =\alpha_{ik}^{(j)\cdot}\beta_{ik}^{(j)}
\]
가 \(Z\)를 바탕으로 한 \(\{\mathfrak K\times\mathfrak S\}\)의 인자계임을 보인다. 따라서 연관된 인자계들의 두 류 \(\{\alpha\}\), \(\{\beta\}\)의 \emph{곱 \(\{\alpha\}\times\{\beta\}\)}을 류
\[
        \bigl\{\alpha_{ik}^{(j)}\beta_{ik}^{(j)}\bigr\}
\]
로 정의하는 것은 뜻이 있다.

\textbf{정리 2.} \emph{\(\{\mathfrak K\}\)가 \(\mathsf{Z}\)를 바탕으로 하여 연관된 인자계들의 류 \(\{\alpha\}\)를 갖는다고 하자. 분해체 \(Z\)를 갖는 모든 \(\{\mathfrak K\}\)의 류 \(\{\alpha\}\)들은 하나의 군을 이루며, 이 군은 대응}
\[
        \{\mathfrak K\}\longmapsto\{\alpha\}
\]
\emph{에 의해 \(\mathscr K_{\mathsf{Z}}\)와 동형으로 대응한다.}

\textbf{증명을 위해} 다음 두 가지를 보여야 한다.

1. \(\{\mathfrak K\}\)와 \(\{\mathfrak S\}\)가 각각 인자계 \(\alpha_{ik}^{(j)}\)와 \(\beta_{ik}^{(j)}\)를 가지면
\[
        \varepsilon_{ik}^{(j)}
        =\alpha_{ik}^{(j)}\beta_{ik}^{(j)}
\]
는 \(\{\mathfrak K\times\mathfrak S\}\)의 인자계이다.

2. \(\{\mathfrak K\}\)가 인자계 \(\alpha_{ik}^{(j)}=1\)을 가지면 \(\mathfrak K=\mathsf P\)이다.

1. \(\mathfrak K\)와 \(\mathfrak S\) 안에서의 \(\mathsf{Z}\)의 기약표현의 차수가 각각 \(r,s\)라고 하고
\[
        \mathfrak K_{r\Gamma}
        =\mathfrak l_1+\cdots+\mathfrak l_n,
        \qquad
        \mathfrak S_{s\Gamma}
        =\mathfrak m_1+\cdots+\mathfrak m_n
\]
을 켤레인 아이디얼들로의 분해라 하자. 여기서 \(\mathfrak l_i\)와 \(\mathfrak m_i\)는 서로 대응하는 아이디얼들이라고 한다. 곧 이들은 각각 \(\mathfrak Z_\Gamma\)와 \(\mathfrak Z'_\Gamma\)의 대응하는 성분들에서 얻어진다. 그러면
\[
 (\mathfrak K_r\times\mathfrak S_s)_\Gamma
 =\mathfrak l_1\mathfrak m_1+
  \cdots+
  \mathfrak l_i\mathfrak m_j+
  \cdots+
  \mathfrak l_n\mathfrak m_n
\]
은 \(0\)이 아닌 \(n^2\)개의, 따라서 단순한 아이디얼들로의 분해이다. 그러므로
\[
        A_1
        =\mathfrak l_1\mathfrak m_1+
         \mathfrak l_2\mathfrak m_2+
         \cdots+
         \mathfrak l_n\mathfrak m_n
\]
은 절대길이가 \(n\)인 아이디얼이다. \(\mathfrak m_i\)들이 \(\mathfrak l_i\)들과 똑같이 변환되므로 \(A_1\)은 모든 \(S\)에 대하여 불변이다. 따라서
\[
        \mathsf{A}_1=\mathsf{A}_\Gamma
\]
인 \(\mathfrak K_r\times\mathfrak S_s\)-아이디얼 \(A\)가 존재한다. (18쪽 보조정리.)

\(\mathfrak K\times\mathfrak S\)의 단순아이디얼들의 자기동형체 \(\mathfrak M\) 안에서의 \(\mathfrak Z\)의 기약표현의 차수를 \(u\)라 하자. 그러면 \(\mathsf{A}\)는 \(u\)개의 단순아이디얼로 분해된다. 그 절대길이가 \(n\)이기 때문이다. 따라서 우리의 보조정리에 의해 \(\mathsf{A}\)의 인자계들은 \(\mathfrak M\)의 인자계들이기도 하다.

그런데 \(\mathsf{A}_\Gamma\)를 켤레인 아이디얼들로 분해한 한 분해는 정의에 사용한
\[
\tag{3}
        A_\Gamma
        =\mathfrak l_1\mathfrak m_1+
         \mathfrak l_2\mathfrak m_2+
         \cdots+
         \mathfrak l_n\mathfrak m_n
\]
이다. \(\alpha_{ik}^{(j)}\), \(\beta_{ik}^{(j)}\)가 각각 \(\mathfrak K\), \(\mathfrak S\)의 유사행렬단위 \(p_{ik}\), \(q_{ik}\)에 속한다고 하자. 그러면
\[
        r_{ik}=p_{ik}q_{ik}
\]
가 분해 (3)에 관한 \(\mathsf{A}_\Gamma\)의 유사행렬단위계임은 명백하다. 여기에 속하는 인자계
\[
        \varepsilon_{ik}^{(j)}
        =\alpha_{ik}^{(j)}\beta_{ik}^{(j)}
\]
는 따라서 \(\mathfrak M\)의 인자계이기도 하다.

2. \(\{\mathfrak K\}\)가 인자계 \(\alpha_{ik}^{(j)}=1\), 이에 속하는 유사행렬단위 \(p_{ik}\), 그리고 켤레인 아이디얼들로의 분해
\[
        \mathfrak K_{r\Gamma}
        =\mathfrak l_1+\cdots+\mathfrak l_n
\]
을 갖는다고 하자. 우리는 \(\mathfrak K=\mathsf P\), 또는 \(\mathfrak K_r=\mathsf P_r\), 또는 \(\mathfrak K_r\)에 랭크가 \(n\)인 좌측아이디얼 \(\mathfrak l\)이 존재함을 보여야 한다. 실제로 \(\mathfrak K_r\)의 단순 좌측아이디얼은 랭크가 \(t^2r\)이므로, 정수 \(h\)에 대하여 \(n=t^2rh\)여야 하며, 따라서 \(t=1\)이어야 하기 때문이다.

\(\mathfrak K_{r\Gamma}\)-아이디얼
\[
        \Omega
        =\mathfrak l_1(p_{11}+\cdots+p_{1n})
\]
은 \(0\)이 아니다. 실제로 \(e_1\sum_i p_{1i}\ne0\)이다. 따라서 \(\Omega\)는 단순하며 랭크가 \(n\)이다. 또한 \(\Omega\)는 \(\mathfrak G\)-불변이다. 곧
\[
\begin{aligned}
S\mathfrak L
 &=\mathfrak l_1\sum_\nu S p_{1\nu}\\
 &=\mathfrak l_i\sum_\mu p_{i\mu}\\
 &=\mathfrak l_i\sum_\mu p_{i1}p_{1\mu},
\end{aligned}
\]
여기서는 \(\alpha_{ik}^{(j)}=1\)을 썼다. 따라서
\[
\begin{aligned}
S\mathfrak L
 &=\mathfrak l_ip_{i1}\sum_\mu p_{1\mu}\\
 &=\mathfrak l_1\sum_\mu p_{1\mu}
  =\mathfrak L.
\end{aligned}
\]
그러므로
\[
        \Omega=\mathfrak l_\Gamma,
\]
\emph{여기서 \(\mathfrak l\)은 랭크가 \(n\)인 \(\mathfrak K_r\)-아이디얼이다.}

\textbf{정리 3.} \emph{군 \(\mathscr K\)의 각 원소 \(\{\mathfrak K\}\)는 유한 위수 \(\lambda\)를 가지며, \(\lambda\)는 \(\mathfrak K\)의 절대성분수 \(t\)의 약수이다.}

\emph{첫째 증명(브라우어).} \(\mathfrak K_{r\Gamma}\)를 켤레인 좌측아이디얼들로 분해한 두 분해
\[
\begin{gathered}
\mathfrak K_{r\Gamma}
=\mathfrak l_1+\cdots+\mathfrak l_n
=\mathfrak z_1+\cdots+\mathfrak z_n,\\
\mathfrak l_i=\mathfrak K_{r\Gamma}e_i,
\qquad
\mathfrak z_i=\mathfrak K_{r\Gamma}e_i^*
\end{gathered}
\]
를 택한다. \(a_i\mapsto a_ir_{ik}\), 특히 \(e_i\mapsto r_{ik}\)가 \(\mathfrak l_i\)에서 \(\mathfrak z_k\) 위로의 한 연산자동형이라고 하자. 그 역에서 \(e_k^*\)는 \(\mathfrak l_i\)의 한 원소 \(s_{ki}\)로 옮겨가며
\[
        s_{ki\cdot}r_{ik}=e_k^*
\]
이다. \(r_{ik}\)들을 서로 켤레가 되도록, 곧
\[
        S(r_{ik})=r_{\nu\mu}
        \quad\text{만일}\quad
        S(i,k)=(\nu,\mu)
\]
로 고르면 \(s_{ik}\)들도 서로 켤레이다. 실제로
\[
\begin{aligned}
S(s_{ki})S(r_{ik})
 &=S(s_{ki}r_{ik})
  =S(e_k^*)
  =e_\mu^*\\
 &=s_{\mu\nu}r_{\nu\mu}
  =s_{\mu\nu}S(r_{ik}),
\end{aligned}
\]
따라서 \(S(s_{ki})=s_{\mu\nu}\)이다. \(e_i\mapsto e_ir_{ij}s_{jk}\)가 \(\mathfrak l_i\)에서 \(\mathfrak l_k\) 위로의 한 연산자동형을 주므로
\[
        r_{ij}s_{jk}
        =p_{ik}\xi_{ik}^{(j)},
        \qquad
        0\ne\xi_{ik}^{(j)}\in\Gamma
\]
이어야 한다.

\(r\)들, \(s\)들, \(p\)들이 서로 켤레이므로 \(\xi\)들도 서로 켤레이다. 곧
\[
        S(i,j,k)=(\nu,\tau,\mu)
        \quad\Longrightarrow\quad
        S(\xi_{ik}^{(j)})=\xi_{\nu\mu}^{\tau}.
\]

등식
\[
\begin{aligned}
r_{i\lambda}s_{\lambda j}r_{j\lambda}s_{\lambda k}
 &=r_{i\lambda}e_\lambda^*s_{\lambda k}\\
 &=r_{i\lambda}s_{\lambda k}
\end{aligned}
\]
에서
\[
        p_{ij}\xi_{ij}^{(\lambda)}
        p_{jk}\xi_{jk}^{(\lambda)}
        =p_{ik}\xi_{ik}^{(\lambda)}
\]
가 따라온다. 그런데 \(p_{ij}p_{jk}=p_{ik}\alpha_{ik}^{(j)}\)이므로
\[
        \alpha_{ik}^{(j)}
        =\frac{\xi_{ik}^{(\lambda)}}
        {\xi_{ij}^{(\lambda)}\xi_{jk}^{(\lambda)}}.
\]
이제
\[
        \delta_{\nu\mu}
        =\prod_\lambda\xi_{\nu\mu}^{(\lambda)}
\]
로 두자. \(\delta_{\nu\mu}\)는 \(\nu\)와 \(\mu\)를 고정하는 모든 \(S\)에 불변이므로 \(\{\mathsf{Z}_\nu,\mathsf{Z}_\mu\}\) 안에 있다. 또 \(\xi_{\nu\mu}^{(j)}\)들이 서로 켤레이므로 \(\delta_{\nu\mu}\)들도 서로 켤레이다. 이제
\[
\begin{aligned}
\alpha_{ik}^{(j)n}
 &=\frac{\prod_\lambda\xi_{ik}^{(\lambda)}}
        {\prod_\lambda\xi_{ij}^{(\lambda)}
         \prod_\lambda\xi_{jk}^{(\lambda)}}\\
 &=\frac{\delta_{ik}}{\delta_{ij}\delta_{jk}}
\end{aligned}
\]
이다. 따라서 연관된 인자계에 관한 공식(32쪽)에 의해 \(\alpha_{ik}^{(j)n}\)은 계 \(\beta_{ik}^{(j)}=1\)과 연관되어 있다. \(\{\mathfrak K\}\)의 위수는 \(\{\alpha\}\)의 위수와 같으므로 \(\{\mathfrak K\}\)의 위수는 \(n\)의 약수이다. 특히 \(\mathfrak Z\)를 \(\mathfrak K\)의 한 극대 가환부분체로 택하면 \(n=t\)이다. 따라서 \(\lambda\)는 \(t\)의 약수이다.

\emph{둘째 증명(슈어).} 이 증명은 \(p_{ik}\)를 \(\Gamma\) 위의 가역인 \(n\)차 행렬 \(P_{ik}\)들로 표현하는 데 바탕을 둔다. 이 행렬들은 \(p_{ik}\)들과 같은 켤레조건을 만족한다. \(p_{ik}\)가 \(P_{ik}\)에 대응하면 \(P_{ik}\)는 \(\{\mathsf{Z}_i,\mathsf{Z}_k\}\)에 속하고, \(S(p_{ik})\)는 \(S(P_{ik})\)에 대응한다. 이 표현은 준동형이다. 즉 \(P_{ik}\)와 \(P_{kj}\)가 각각 \(p_{ik}\)와 \(p_{kj}\)에 대응하면 \(P_{ik}P_{kj}\)는 \(p_{ik}p_{kj}\)에 대응하므로
\[
        P_{ik}P_{kj}
        =\alpha_{ij}^{(k)}P_{ij}
\]
이다. (\(k_1\ne k_2\)일 때 \(P_{ik_1}P_{k_2j}=0\)인 경우는 생길 수 없다.) 이러한 \(p_{ik}\)의 표현이 있으면 다음과 같이 결론을 얻는다. 행렬식을 취하면
\[
        |P_{ik}|\cdot\,|P_{kj}|
        =\alpha_{ij}^{(k)n\cdot}|P_{ij}|
\]
이다. 행렬 \(P_{ik}\)들이 가역이므로 \(|P_{\nu\mu}|\ne0\)이다. 따라서 \(|P_{\nu\mu}|=\delta_{\nu\mu}\)로 두면
\[
        \alpha_{ij}^{(k)n}
        =\frac{\delta_{ik}\delta_{kj}}{\delta_{ik}}
\]
이고, 그 뒤는 브라우어의 증명과 같다.

\emph{행렬 \(P_{ik}\)들에 의한 \(p_{ik}\)의 표현.}

분해
\[
        \mathfrak K_{r\Gamma}
        =\mathfrak l_1+\cdots+\mathfrak l_n
\]
의 아이디얼 \(\mathfrak l_i\)들은 \(\Gamma\) 안에서의 \(\mathfrak K_r\)의 반대 표현가군들이다. (32쪽 정리 2.) \(\mathfrak Z\)에 관한 \(\mathfrak K_r\)의 한 고정된 기저 \(k_1,\ldots,k_n\)을 택하면
\[
        \mathfrak K_{r\Gamma}
        =k_1\mathfrak Z_\Gamma+\cdots+k_n\mathfrak Z_\Gamma,
\]
따라서
\[
\begin{aligned}
\mathfrak l_i
 &=\mathfrak K_{r\Gamma}e_i\\
 &=k_1\mathfrak Z_re_i+
   \cdots+
   k_n\mathfrak Z_re_i\\
 &=k_1e_i\Gamma+\cdots+k_ne_i\Gamma.
\end{aligned}
\]
그러므로
\[
\begin{aligned}
\mathfrak l_k
 &=\mathfrak l_ip_{ik}\\
 &=k_1p_{ik}\Gamma+\cdots+k_np_{ik}\Gamma.
\end{aligned}
\]
\(\mathfrak l_k\)의 두 기저
\[
        (k_1e_i,\ldots,k_ne_i),
        \qquad
        (k_1p_{ik},\ldots,k_np_{ik})
\]
는 한 가역행렬 \(P_{ik}\)를 곱하여 서로 옮겨갈 수 있다.
\[
        P_{ik}\cdot
        \begin{pmatrix}
          k_1e_i\\ \vdots\\ k_ne_i
        \end{pmatrix}
        =
        \begin{pmatrix}
          k_1p_{ik}\\ \vdots\\ k_np_{ik}
        \end{pmatrix},
        \qquad
        \text{줄여서}\quad
        P_{ik}(k_\varrho e_i)=(k_\varrho p_{ik}).
\]
이로부터 \(P_{ik}\)들의 켤레성이 따라온다.
\[
\begin{aligned}
S(P_{ik})(k_\varrho S(e_k))
 &=(k_\varrho S(p_{ik})),
        \quad\text{즉}\\
S(P_{ik})(k_\varrho e_\mu)
 &=(k_\varrho p_{\nu\mu})
  =P_{\nu\mu}(k_\varrho e_\mu).
\end{aligned}
\]
따라서
\[
        S(P_{ik})=P_{\nu\mu}
        \quad\text{만일}\quad
        S(i,k)=(\nu,\mu).
\]
다시 이로부터 \(P_{ik}\)가 \(\{\mathsf{Z}_i,\mathsf{Z}_k\}\)에 속함이 따라온다. 실제로 이 합성체는 \(i\)와 \(k\)를 고정하는 \(S\)들의 불변체이다. 준동형성도 쉽게 알 수 있다.
\[
\begin{split}
P_{ik}P_{kj}(k_\varrho e_j)
 &=P_{ik}(k_\varrho p_{kj})
  =(k_\varrho p_{ik}p_{kj})\\
 &=\alpha_{ij}^{(k)}(k_\varrho p_{ij})
  =\alpha_{ij}P_{ij}(k_\varrho e_j),
\quad\text{따라서}\\
&\hspace{-2em}\cdot P_{ik}P_{kj}
  =\alpha_{ij}^{(k)}P_{ij},
\quad\text{증명 끝.}
\end{split}
\]

정리 3은 다음과 같이 강화할 수 있다.

\textbf{정리 4.} \emph{소수 \(p\)가 \(\mathfrak K\)의 절대성분수 \(t\)의 소인수이면 \(p\)는 \(\{\mathfrak K\}\)의 위수 \(\lambda\)의 약수이기도 하다.} (브라우어.)

\emph{증명.} \(p^\sigma\)를 분해체 \(\mathsf{Z}\)의 차수 \(n\)을 나누는 \(p\)의 최고 거듭제곱이라 하자. \(\mathfrak H\)를 \(\mathfrak G\)의 위수가 \(p^\sigma\)인 한 부분군이라 하고, \(\mathsf{T}\)를 \(\mathfrak H\)의 불변체라 하자. 그러한 부분군은 \emph{쉴로 정리}에 의해 존재한다. (슈파이저, 제1판, §19, 42쪽, 정리 43.)

그러면
\[
        [\Gamma:T]=p^\sigma,
        \qquad
        [T:\mathsf P]=\frac{n}{p^\sigma}
\]
이고, 뒤의 수는 \(p\)로 나누어지지 않으며 더구나 \(t\)로도 나누어지지 않는다. 따라서 \(\mathsf{T}\)는 분해체가 아니다. \(\mathfrak K_{\mathsf{T}}=\mathfrak S_m\)의 자기동형체 \(\mathfrak S\)는 \(\mathsf P\)와 다르다. 그런데 \(\mathfrak S\)는 분해체 \(\Gamma\)를 가지며, \(\mathfrak S\)의 중심 \(T\) 위에서의 그 차수가 \(p^\sigma\)이다. 그러므로 \(\mathfrak S\)의 절대성분수는 \(p\)의 거듭제곱 \(p^\beta>1\)이고, 따라서 정리 3에 의해 \(\{\mathfrak S\}\)의 위수도 \(p\)의 거듭제곱 \(p^\alpha>1\)이다. 이제
\[
        \{\mathfrak K\}^{\lambda}=\{\mathsf P\},
        \qquad
        \{\mathfrak K_{\mathsf{T}}\}^{\lambda}=\{\mathsf{T}\}
\]
이므로
\[
        \lambda\equiv0\pmod {p^\alpha}.
\]
증명 끝. (\(\lambda=t\)일 필요가 없음을 보이는 예들은 브라우어가 제시하였다\srcfn{1}{Brauer: Untersuchungen über die arithmetischen Eigenschaften von Gruppen linearer Substitutionen II. Math. Zeitschr. 31, S. 733 § 5, (1930).}.)

\textbf{정리 5.} \emph{\(t\)의 소인수분해가}
\[
        t=\prod_i p_i^{\varrho_i},
        \qquad
        \varrho_i\geqq1,
        \qquad
        i\ne j\Longrightarrow p_i\ne p_j
\]
\emph{라고 하자. 그러면}
\[
        \mathfrak K
        =\mathfrak K^{(1)}\times
         \mathfrak K^{(2)}\times\cdots,
\]
\emph{여기서 \(\mathfrak K^{(i)}\)는 절대성분수가 \(p_i^{\varrho_i}\)인 체이다.}

\emph{증명.} 정리 3과 4에 따라 \(\{\mathfrak K\}\)의 위수 \(\lambda\)는
\[
        \lambda=\prod_i p_i^{\alpha_i},
        \qquad
        1\leqq\alpha_i\leqq\varrho_i
\]
이다. \(\{\mathfrak K\}\)가 생성하는 \(\mathscr K\)의 순환부분군 \(\mathscr \mathsf{Z}\)는 위수가 \(p_i^{\alpha_i}\)인 부분군 \(\mathscr \mathsf{Z}_i\)들의 직접곱이다. 따라서
\[
        \{\mathfrak K\}
        =\prod_i\{\mathfrak K^{(i)}\}.
\]
이로부터 \(\mathfrak K\)의 모든 분해체 \(\mathsf{Z}\)가 각 \(\mathfrak K^{(i)}\)의 분해체이기도 함이 따라온다. 따라서 \(\mathfrak K^{(i)}\)의 절대성분수 \(p_i^{\sigma_i}\)는 \(t\)를 나누며 \(\sigma_i\leqq\varrho_i\)이다. 이제
\[
        \prod_i\mathfrak K^{(i)}=\mathfrak K_m
\]
이라 하면
\[
        \prod_i p_i^{\varrho_i}
        =\prod_i p_i^{\sigma_i\cdot}\,m,
        \qquad
        \sigma_i\geqq\varrho_i.
\]
따라서
\[
        \sigma_i=\varrho_i,
        \qquad
        m=1,
        \qquad
        \mathfrak K=\prod \mathfrak K^{(i)},
\]
증명 끝.

\subsection*{§ 25. 갈루아 극대 가환부분체를 갖는 \(\mathfrak K_r\)의 정규표현\protect\footnotemark}
\footnotetext{§27부터 더 간단한 증명이 주어진다.}

\(\mathfrak K_r\)의 극대 가환부분체 \(\mathfrak Z\)가 갈루아라고 하자. \(\mathfrak Z/\mathsf P\)의 \(n\)개 자기동형 \(\mathsf{H}_i\)는 25쪽 정리 7에 따라 \(\mathfrak K_r\)의 내부자기동형이다. 즉 원소 \(u_i\)들이 존재하여 \(\mathfrak Z\)의 모든 \(z\)에 대하여
\[
        \mathsf{H}_i z=u_i^{-1}zu_i
\]
이다. 각 \(u_i\)는 \(\mathfrak Z\)의 \(0\)이 아닌 인자 \(y\)를 곱하는 것까지 정해진다. 실제로 \(\mathfrak Z\)의 모든 \(z\)에 대하여
\[
        u^{-1}zu=u'^{-1}zu'
\]
이면
\[
        (u^{-1}u')^{-1}z(u^{-1}u')=z
\]
이다. 따라서 \(u^{-1}u'\)를 \(\mathfrak Z\)에 첨가하면 \(\mathfrak K_r\)의 한 가환부분체를 얻는다. \(\mathfrak Z\)의 극대성 때문에 이 부분체는 \(\mathfrak Z\)와 같으므로 \(u^{-1}u'\in\mathfrak Z\)이다.

이제 다음 명제를 증명한다.
\[
        \mathfrak K_r
        =u_1\mathfrak Z+\cdots+u_n\mathfrak Z.
\]
이를 위해서는 \(u_i\)들이 \(\mathfrak Z\) 위에서 선형독립임을 보이면 충분하다. 우리는 \(u_i\)들을 직접 제시하여 증명한다. 다음 분해가 주어졌다고 하자.
\[
\begin{gathered}
\mathfrak Z_{\mathsf{Z}}=e_1\mathsf{Z}+\cdots+e_n\mathsf{Z},\\
\mathfrak K_{r\mathsf{Z}}=\mathfrak l_1+\cdots+\mathfrak l_n,\quad,
\qquad
\mathfrak l_i=\mathfrak K_{r\mathsf{Z}}e_i,
\end{gathered}
\]
그리고 \(p_{ik}\)를 이 분해에 속하는 한 유사행렬단위계라 하자. 8쪽에서 표현가군 \(e_i\mathsf{Z}\)가 주는 \(\mathfrak Z\)에서 \(\mathsf{Z}\) 위로의 동형사상을 \(\Theta_i\)라 썼다. 그러면
\[
        S_i=\Theta_i\Theta_1^{-1}
\]
는 \(\mathsf{Z}\)의 \(n\)개 자기동형이다. 다만 이는 \(\mathsf{Z}\) 자체에 적용할 때에만 그러하며, 여기서 계속 바탕에 두는 \(\mathfrak K_{r\mathsf{Z}}\) 전체에 대한 \(S_i\)의 작용과는 구별해야 한다. 편의를 위해 \(S_i=\Theta_i\Theta_1^{-1}\)에서 \(\Theta_i\) 대신 \(\Theta_{S_i}\)라 쓰고, \(p_{ik}\)의 첨자 \(i,k\)도 해당 자기동형 \(S\)들로 바꾸어
\[
        p_{ik}=p_{S_iS_k}
\]
라 쓰기로 한다.

\(\mathfrak Z\)의 \(n\)개 자기동형은 \(E\)를 항등치환이라 할 때
\[
        \mathsf{H}=\Theta_E^{-1}\Theta
\]
이다. 이제 \(\mathsf{H}\)들의 번호를
\[
        \mathsf{H}_S=\Theta_E^{-1}\Theta_{S^{-1}}
\]
로 정한다.

이 조건들 아래에서
\[
        u_T=\sum_Sp_{S,ST}
            =\sum_SS(p_{E,T})
\]
로 둘 수 있다. 실제로 다음이 성립한다.

1. \(u_T^{-1}\)가 존재하며
\[
        u_T^{-1}
        =\sum_Sp_{ST,S}
         \frac{1}
         {\alpha_{S,S}^{ST}
          \overline\gamma_{S,ST}
          \overline\gamma_{ST,S}}
\]
이다. (\(\bar\gamma\)의 뜻은 32쪽을 보라.)

2. \(u_T\)는 \(\mathfrak G\)-불변이므로 \(\mathfrak K_r\)의 원소이다.
\[
        R(u_T)
        =\sum_SRS(p_{E,T})
        =\sum_{S'}S'(p_{E,T})
        =u_T.
\]

3. \(u_T\)는 \(\mathsf{H}_T\)를 생성한다. 즉 \(z\in\mathfrak Z\)이면
\[
        u_T^{-1}zu_T=H_T(z),
        \qquad\text{또는}\qquad
        z u_T=u_T \mathsf{H}_T(z).
\]
실제로
\[
        z=\sum_S\Theta_S(z)e_S,
        \qquad
        u_T=\sum_SS(e_Ep_{E,T})
            =\sum_Se_SS(p_{E,T})
\]
이므로
\[
\begin{aligned}
zu_T
 &=\sum_{S,S'}\Theta_S(z)e_Se_{S'}S(p_{E,T})\\
 &=\sum_S\Theta_S(z)S(p_{E,T}).
\end{aligned}
\]
한편
\[
        u_T=\sum_SS(p_{E,T}e_T)
            =\sum_SS(p_{E,T})e_{ST}
\]
이고, \(\mathfrak Z\)-원소들에 적용할 때 다음 등식이 성립한다.
\[
\begin{aligned}
\Theta_SH_T
 &=\Theta_S\Theta_E^{-1}\Theta_{T^{-1}}\\
 &=S\Theta_{T^{-1}}
  =ST^{-1}\Theta_E
  =\Theta_{ST^{-1}}.
\end{aligned}
\]
따라서
\[
        H_T(z)
        =\sum_Se_S\Theta_SH_T(z)
        =\sum_Se_S\Theta_{ST^{-1}}(z),
\]
그러므로
\[
\begin{aligned}
u_T\mathsf{H}_T(z)
 &=\sum_{S,S'}S(p_{E,T})e_{ST}\,
   e_{S'}\Theta_{S'T^{-1}}(z)\\
 &=\sum_SS(p_{E,T})\Theta_S(z).
\end{aligned}
\]
결국
\[
        z\cdot u_T=u_T\mathsf{H}_T(z).
\]

4. \(u_S\)들은 \(\mathfrak Z\) 위에서 선형독립이다. 실제로
\[
        \sum_Tz_Tu_T=0
\]
이라고 하자. \(z=\sum_S\Theta_S(z)e_S\)이므로
\[
\begin{aligned}
\sum_{T,S,L}\Theta_S(z_T)e_Sp_{L,LT}
 &=\sum_{S,T}\Theta_S(z_T)p_{S,ST}=0.
\end{aligned}
\]
따라서 \(\Theta_S(z_T)=0\)이고, 결국 \(z_T=0\)이다.

이제부터 \(\mathsf{H}_S(z)\) 대신 기호적 거듭제곱 \(z^S\)를 쓴다. \(u_S\)들 사이에는 다음 꼴의 관계식이 성립해야 한다.
\[
        u_Ru_T=u_{RT}a_{R,T},
        \qquad
        0\ne a_{R,T}\in\mathfrak Z.
\]
실제로
\[
\begin{aligned}
u_Ru_T
 &=\sum_{S,S'}S(p_{E,R})S'(p_{E,T})\\
 &=\sum_Sp_{S,SR}p_{SR,SRT}\\
 &=\sum_Sp_{S,SRT}\alpha_{S,SRT}^{(SR)}\\
 &=\sum_Sp_{S,SRT}\alpha_{S,SRT}^{(SR)}e_{SRT}
\end{aligned}
\]
이고
\[
        a_{R,T}
        =\sum_Se_{SRT}\Theta_{SRT}(a_{R,T}),
\]
따라서
\[
        u_{RT}a_{R,T}
        =\sum_Sp_{S,SRT}\Theta_{SRT}(a_{R,T})
\]
이다. 그러므로
\[
\begin{aligned}
\Theta_{SRT}(a_{R,T})
 &=\alpha_{S,SRT}^{(SR)},\\
a_{R,T}
 &=\sum_Se_{SRT}\alpha_{S,SRT}^{(SR)}\\
 &=\sum_Le_L
   \alpha_{LT^{-1}R^{-1},L}^{(LT^{-1})}.
\end{aligned}
\]

\(a_{R,T}\)들의 계를 \(\mathfrak K_r\)의 \emph{소인자계}라 부르기로 하고, \(\alpha_{BC}^{(A)}\)를 대인자계라 부르기로 한다. 대인자계 \(\alpha_{BC}^{(A)}\)와 소인자계 \(a_{R,T}\)는 서로를 결정한다. 두 대인자계의 곱
\[
        \varepsilon_{BC}^{(A)}
        =\alpha_{BC}^{(A)}\beta_{BC}^{(A)}
\]
에 대응하는 소인자계는 \(\alpha_{BC}^{(A)}\), \(\beta_{BC}^{(A)}\)에 대응하는 두 소인자계 \(a_{R,T}\), \(b_{R,T}\)의 곱
\[
        f_{R,T}=a_{R,T}b_{R,T}
\]
이다. 이는 위 공식들에서 곧바로 따라온다. 대단위원자계 \(\alpha_{BC}^{(A)}=1\)에는 소단위원자계 \(a_{R,T}=1\)이 대응하고, 그 역도 성립한다.

\(u_S\)를
\[
        u'_S=u_Sr_S,
        \qquad
        0\ne r_S\in\mathfrak Z
\]
로 바꿀 수 있다. 여기서 \(r_S\)는 임의로 고른다. \(u'_S\)에 속하는 소인자계를 \(a'_{R,T}\)라 하면
\[
\tag{1}
        u'_Ru'_T=u'_{RT}a'_{R,T}.
\]
\(a'_{R,T}\)를 \(a_{R,T}\)와 연관된 소인자계라 부른다. 그런데
\[
\begin{aligned}
u_Rr_Ru_Tr_T
 &=u_{RT}r_{RT}a'_{R,T}\\
 &=u_Ru_Tr_R^Tr_T\\
 &=u_{RT}a_{R,T}r_R^Tr_T,
\end{aligned}
\]
따라서
\[
\tag{2}
        a'_{R,T}
        =a_{R,T}\frac{r_R^Tr_T}{r_{RT}}.
\]

\(u'_S=u_Sr_S\)에는 다음 유사행렬단위계가 대응한다.
\[
        p'_{S,ST}=p_{S,ST}r_T.
\]
따라서 소인자계 \(a'_{R,T}\)는, \(a_{R,T}\)에 속하는 대인자계 \(\alpha_{BC}^{(A)}\)와 연관된 한 대인자계 \(\alpha_{BC}^{(A)\prime}\)에 속한다. 연관된 소인자계들은 연관된 대인자계들에 대응하고, 그 역도 성립한다. 연관된 대인자계들의 류 \(\{\alpha\}\)와 이에 대응하는 연관된 소인자계들의 류 \(\{a\}\) 사이의 대응은 군 동형사상이다.

이제 \(u_S\)들을 \(\mathfrak Z\) 위의 \(n\)차 행렬들로 표현한다. 이 표현은 지금까지 고찰한 표현들과 본질적으로 다르다.

\(\mathfrak K_r\)의 \(\mathfrak Z\)-기저
\[
\tag{I}
        k_{S_1},\ldots,k_{S_n}
\]
에 오른쪽에서 \(u_S\)를 곱하면
\[
\tag{II}
        k_{S_1}u_S,\ldots,k_{S_n}u_S
\]
를 얻는다. (I)을 (II)로 옮기는 가역행렬 \(A_S\), 곧
\[
\tag{III}
        (k_{S_1}u_S,\ldots,k_{S_n}u_S)
        =(k_{S_1},\ldots,k_{S_n})A_S
\]
인 \(A_S\)를 \(u_S\)를 표현하는 행렬이라 하자. 이 표현은 \(u_S\)처럼 \(\mathfrak Z/\mathsf P\)의 한 자기동형을 생성하는 모든 \(\mathfrak K_r\)-원소에 적용된다. 즉 \(\mathfrak Z\)에서 임의로 고른 \(r_S\ne0\)에 대한 모든 원소 \(u_Sr_S\)에 적용된다. 이 원소들은 하나의 군 \(\mathfrak G^*\)를 이루며, \(\mathfrak Z\)의 곱셈군 \(\mathfrak Z^*\)를 정규부분군으로 포함한다. 몫군 \(\mathfrak G^*/\mathfrak Z^*\)는 \(\mathfrak Z/\mathsf P\)의 갈루아군 \(\mathfrak G\)와 동형이다. 실제로 \(\mathfrak Z^*\)에 대한 \(\mathfrak G^*\)의 각 잉여류의 모든 원소가 그 잉여류에 대응하는 \(\mathfrak Z\)의 자기동형을 생성한다.

이 군 \(\mathfrak G^*\)를 방금 설명한 방식으로
\[
        u_S\longmapsto A_S
\]
라 표현한다. 그러나 이 표현은 준동형도 아니고 반대준동형도 아니다. 곧
\[
        u_S\longmapsto A_S,
        \qquad
        u_T\longmapsto A_T
\]
에서 \(u_Su_T\mapsto A_SA_T\)가 따라오지 않는다. 오히려
\[
\begin{aligned}
(k_{S_1},\ldots,k_{S_n})u_Su_T
 &=(k_{S_1},\ldots,k_{S_n})A_Su_T\\
 &=(k_{S_1},\ldots,k_{S_n})u_TA_S^T\\
 &=(k_{S_1},\ldots,k_{S_n})A_TA_S^T,
\end{aligned}
\]
따라서
\[
        u_Su_T\longmapsto A_TA_S^T.
\]

이 \(\mathfrak G^*\)의 표현을 \emph{교차표현}이라 부른다. \(h_{S_1},\ldots,h_{S_n}\)이 \(\mathfrak K_r\)의 \(\mathfrak Z\)에 관한 다른 한 기저이고
\[
\tag{IV}
        (h_{S_1},\ldots,h_{S_n})
        =(k_{S_1},\ldots,k_{S_n})M
\]
이라고 하자. 그러면 (III)과 (IV)에서
\[
        (h_{S_1},\ldots,h_{S_n})u_S
        =(h_{S_1},\ldots,h_{S_n})M^{-1}A_SM^S
\]
를 얻는다. 즉 \(u_S\)가 \(k\)-기저에 관하여 \(A_S\)로 표현되고 \(M\)이 \(k\)-기저를 \(h\)-기저로 바꾸는 행렬이면, \(u_S\)는 \(h\)-기저에 관하여 \(M^{-1}A_SM^S\)로 표현된다.

교차표현을 정의하는 공식 (III)은 교차표현이 환 \(\mathfrak K_r\)에 대하여 갖는 관계가 보통의 반대 표현이 그 표현가군에 대하여 갖는 관계와 비슷함을 보여 준다. 다만 여기서는 보통의 반대 표현을 왼쪽이 아니라 오른쪽에서 쓴다. 차이는 반대 표현가군에서의 교환관계
\[
        (m^*\tau^*)c=(m^*c)\tau^*
        \qquad\text{(5쪽 참조)}
\]
대신 여기서는
\[
        (kz)u_S=(ku_S)z^S,
        \qquad
        k\in\mathfrak K_r,
        \quad
        z\in\mathfrak Z
\]
가 나타난다는 데서 생긴다.

\(\mathfrak K_r\)의 한 우측아이디얼 \(\mathfrak r\)의 기저를 사용해서도 \(\mathfrak G^*\)의 교차표현을 만들 수 있다. \(\mathfrak K_r\)의 각 우측아이디얼은 이런 방식으로 서로 동치인 표현들의 한 류를 생성한다. 여기서 표현
\[
        u_S\longmapsto A_S,
        \qquad
        u_S\longmapsto B_S
\]
는 어떤 행렬 \(M\)에 대하여
\[
        B_S=M^{-1}A_SM^S
\]
이면 동치라고 한다.

이제 \(\mathfrak K_r\)의 우측아이디얼들이 생성하는 \(\mathfrak G^*\)의 \(\mathfrak Z\) 안에서의 표현들에 국한하지 않고, \(\mathfrak G^*\)의 \(\mathfrak Z\) 안에서의 \(m\)차 교차표현을 일반적으로 다음과 같이 정의할 수 있다.
\begin{enumerate}
\item 각 \(u_S\)에 \(\mathfrak Z\) 위의 \(m\)차 행렬 \(A_S\)가 대응한다.
\item \(u_S\mapsto A_S\), \(u_T\mapsto A_T\)이면 \(u_Su_T\mapsto A_TA_S^T\)이다.
\end{enumerate}

랭크가 \(m\)인 \(\mathfrak Z\)-우측가군
\[
        \mathfrak M
        =x_1\mathfrak Z+\cdots+x_m\mathfrak Z
\]
을 택하고
\[
        (x_1u_S,\ldots,x_mu_S)
        =(x_1,\ldots,x_m)A_S
\]
로 정하여 이를 교환관계
\[
        (mz)u_S=(mu_S)z^S
\]
를 갖는 \(\mathfrak G^*\)-우측가군으로 만든다. 이는 5쪽에서 주어진 표현으로부터 표현가군을 만든 것과 완전히 같다. 이제 20쪽에서와 같이
\[
        m\mathbin{\cdot}zu_S
        =mz\mathbin{\cdot}u_S
        =mu_S\mathbin{\cdot}z^S
\]
로 정하고, 일반적으로
\[
\begin{aligned}
m\mathbin{\cdot}\sum_Sz_Su_S
 &=\sum_Smz_S\mathbin{\cdot}u_S\\
 &=\sum_Smu_S\mathbin{\cdot}z_S^S
\end{aligned}
\]
로 정하여 \(\mathfrak M\)을 유한 \(\mathfrak K_r\)-우측가군으로 만든다. 여기서는 이 정의에 특정한 기저를 사용했으므로 20쪽의 계산에 해당하는 것은 전혀 필요하지 않다.

그런데 M. Z. §18에 따라 모든 유한 단순 \(\mathfrak K_r\)-가군 \(\mathfrak M\)은 \(\mathfrak K_r\)의 한 단순 우측아이디얼과 연산자동형이다. 따라서 22쪽 정리 2와 같이 \(\mathfrak G^*\)의 \(\mathfrak Z\) 안에서의 기약 교차표현류는 \emph{오직 하나}, 곧 단순 우측아이디얼들에 속하는 류뿐임을 알 수 있다.

이 기약표현의 차수는 \(\mathfrak K\)의 지수(절대 지수) \(t\)이다. 이는 간단히 랭크를 세면 알 수 있다. 실제로 그 차수는 한 단순 우측아이디얼 \(\mathfrak r\)의 \(\mathfrak Z\)-랭크 \((\mathfrak r:\mathfrak Z)\)와 같다. 22쪽에 따라
\[
        (\mathfrak Z:\mathsf P)=n=rt,
        \qquad
        (\mathfrak K_r:\mathsf P)=n^2,
        \qquad
        (\mathfrak r:\mathsf P)=n\cdot\,t
\]
이다. 마지막 등식은
\[
        \mathfrak K_r
        =\mathfrak r_1+\cdots+\mathfrak r_r
\]
이기 때문이다. 이제
\[
        (\mathfrak r:\mathsf P)
        =(\mathfrak r:\mathfrak Z)
         (\mathfrak Z:\mathsf P),
\]
즉
\[
        n\cdot t=(\mathfrak r:\mathfrak Z)\cdot n
\]
이므로 \((\mathfrak r:\mathfrak Z)=t\)이다.

\subsection*{§ 26. 교차표현의 곱셈}

지금까지 교차표현을 언제나 \(\mathfrak G^*\)의 표현으로 이해하였다. 그러나 이를 \(\mathfrak Z/\mathsf P\)의 갈루아군 \(\mathfrak G\)의 표현으로도 생각할 수 있다. 곧 \(\mathfrak G\)의 각 원소 \(S\)에 행렬 \(A_S\)를 대응시키는데, 이 행렬은 변환
\[
        u_S^{-1}zu_S=z^S
\]
에 의해 자기동형 \(S\)를 생성하는 \(\mathfrak K_r\)의 한 원소 \(u_S\)를 표현한다.

그러한 표현 \(S\mapsto A_S\)을 특징짓기 위해서는\srcfn{1}{초복소계와 독립적으로 주어진 그러한 표현
\[
        S\mapsto A_S,
        \qquad
        T\mapsto A_T,
        \qquad
        ST\mapsto A_TA_S^T
        =A_{ST\cdot}a_{S,T}^{-1}
\]
(또는 그 전치표현)은 슈파이저가 인자계 \(a_{S,T}\)를 정의할 때의 출발점이었다(Math. Zschr. 5). 따라서 이 인자계들은 여기서 다른 바탕 위에 도입한 인자계들과 동일한 것으로 드러난다.} 그 표현을 어느 \(\mathfrak K_r\) 안에서 취했는지, 어떤 \(u_S\)들을 썼는지, 그리고 \(\mathfrak K\)의 \(\mathfrak Z\) 위의 어떤 기저 \(k_{S_1},\ldots,k_{S_n}\)을 썼는지를 말해야 한다. 기저를 바꾸면 동치인 표현이 생기고, \(u\)들을 바꾸면 연관된 인자계 \(a_{S,T}\)들이 생긴다. 따라서 본질적인 것은 표현을 취한 \(\mathfrak K_r\)들, 즉 연관된 인자계들의 류 \(\{a_{S,T}\}\)이다.

이제 \(\mathfrak G\)의 그러한 두 표현이 주어졌다고 하자.
\[
\begin{array}{lll}
S\mapsto u_S\mapsto A_S
 &\text{이에 속하는 소인자계}&a_{R,T},\\
S\mapsto v_S\mapsto B_S
 &\text{이에 속하는 소인자계}&b_{R,T}.
\end{array}
\]

두 행렬 \(A=(a_{ik})\), \(B=(b_{\mu\nu})\)의 크로네커곱 \(A\times B\)이란 행렬
\[
        A\times B=(a_{ik}b_{\mu\nu})
\]
을 뜻한다.

\noindent\textbf{명제.} \emph{\(S\to A_S\times B_S\)도 \(\mathfrak G\)의 한 교차표현이다. 더 정확히 \(A_S\)가 \(\mathfrak K_r\)에서, \(B_S\)가 \(\mathfrak L_s\)에서 생기면 이 표현은 곱 \(\mathfrak K_r\times\mathfrak L_s\)에 속한다.}

\noindent\emph{증명.} \(S\mapsto A_S\times B_S\)가 교차표현임을 보이자. 다음 대응들이 주어져 있다.
\[
        S\mapsto A_S\times B_S,
        \qquad
        T\mapsto A_T\times B_T,
\]
그리고
\[
\begin{aligned}
ST&\mapsto A_{ST}\times B_{ST}\\
  &=(A_TA_S^T\times B_TB_S^T)
    a_{S,T}^{-1}b_{S,T}^{-1}.
\end{aligned}
\]
따라서 다음을 증명해야 한다.
\[
\begin{aligned}
(A_T\times B_T)(A_S\times B_S)^T
 a_{S,T}^{-1}b_{S,T}^{-1}
 &=(A_TA_S^T\times B_TB_S^T)
   a_{S,T}^{-1}b_{S,T}^{-1}.
\end{aligned}
\]
이는 다음 고찰에서 따라온다.

행렬 \(A\)와 \(B\)를 행렬단위 \(c_{ik}\)와 \(d_{\mu\nu}\)로 나타내자. 이 행렬단위들은 \(\mathfrak K_r\) 또는 \(\mathfrak L_s\)의 행렬단위와 아무 관계가 없다. 곧
\[
        A=\sum c_{ik}\alpha_{ik},
        \qquad
        B=\sum d_{\mu\nu}\beta_{\mu\nu}.
\]
그러면
\[
\begin{aligned}
A\times B=B\times A
 &=\sum c_{ik}d_{\mu\nu}
       \alpha_{ik\cdot}\beta_{\mu\nu}.
\end{aligned}
\]
즉 이는 새로운 행렬단위
\[
        c_{ik}d_{\mu\nu}=d_{\mu\nu}c_{ik}
\]
에 관한 한 행렬과 같다. 따라서 \(A\)-행렬과 \(B\)-행렬의 곱셈은 두 행렬환의 직접곱
\[
        \sum c_{ik}\mathfrak Z
        \times
        \sum d_{\mu\nu}\mathfrak Z
\]
안에서의 원소 곱셈과 일대일로 대응한다. 그러므로 각 \(A\)-행렬은 각 \(B\)-행렬과 교환한다. 다만 \(A\)-행렬들끼리 또는 \(B\)-행렬들끼리 교환하는 것은 아니다. 바로 이 교환가능성에서 \(a_{S,T}^{-1}b_{S,T}^{-1}\)를 곱한 뒤 증명할 관계식이 따라온다.

\section*{제VI장. 교차곱 이론}

\subsection*{§ 27. \(\mathfrak K_r\)의 교차곱 표현}

§24에서 인자계로부터 전개한 교차곱 이론을 이제 독립적으로 다시 정초한다. 이 교차곱의 의의는 갈루아체와 그 군들을 \(\mathfrak K_r\) 안에서 \emph{함께} 기술한다는 데 있다.

\noindent\textbf{정리 1.} \emph{\(\mathfrak Z\)가 \(\mathfrak K_r\)의 갈루아 분해체, 따라서 극대 가환부분체이고 \(\mathfrak G\)가 그 갈루아군이면, \(\mathfrak K_r\)는 교차곱 \(\mathfrak Z\cdot\mathfrak G\)와 같다. 역으로 각 교차곱 \(\mathfrak Z\cdot\mathfrak G\)는 이러한 종류의 \(\mathfrak K_r\)를 생성한다.}

\noindent\textbf{증명과 정의. 정의 1.} \(\mathfrak G=\{E,S,\ldots,T\}\)를 \(\mathfrak Z\)의 갈루아군이라 하자. \(\mathfrak K_r\)가 교차곱과 같다는 것은
\[
        \mathfrak K_r
        =\mathfrak Zu_E+\cdots+\mathfrak Zu_T
\]
이고
\[
        zu_S=u_Sz^S,
        \qquad
        u_Su_T=a_{S,T}u_{ST}
\]
임을 뜻한다. 여기서 \(a_{S,T}\in\mathfrak Z\)는 소인자계를 이룬다.

\noindent\textbf{정의 2.} \emph{생성군} \(\mathfrak G^*\)는 변환에 의해 \(\mathfrak Z\) 전체를 자기 자신으로 옮기는 \(\mathfrak K_r\)의 모든 가역원 \(g\)의 집합으로 정의한다. 곧
\[
        g^{-1}\mathfrak Zg=\mathfrak Z,
        \qquad
        g^{-1}zg=z^S
        \quad(S\in\mathfrak G).
\]
\(\mathfrak Z\)의 각 자기동형은 어떤 \(g\)에 의해 생성되므로 \(\mathfrak G\)는 \(\mathfrak G^*\)의 군준동형상이다.
\[
        \mathfrak G^*\longrightarrow\mathfrak G,
        \qquad
        \mathfrak G\cong\mathfrak G^*/\mathfrak T^*,
        \qquad
        \mathfrak T^*\supseteqq\mathfrak Z^*.
\]
나중에 \(\mathfrak T^*=\mathfrak Z^*\)임을 보일 것이다. 아직 \(\mathfrak Z\)가 극대 가환부분환이라는 것은 증명하지 않았다. 이 사실은 독립적으로도 증명할 수 있다.

\noindent\textbf{1.} \emph{교차곱 표현의 증명.} 자기동형 \(E,\ldots,T\)를 생성하는 \(\mathfrak K_r\)의 원소들을 \(u_E,u_S,\ldots,u_T\)라 하자. 따라서 \(zu_S=u_Sz^S\)이고, 특히 2에 의해 \(u_E=e\)이다. 합
\[
        \mathfrak M
        =\{\mathfrak Zu_E,\ldots,\mathfrak Zu_T\}
\]
이 직접합
\[
        \mathfrak Zu_E+\cdots+\mathfrak Zu_T
\]
과 같음을 보이면 된다. 그러면 랭크가 일치하므로 이 합은 \(\mathfrak K_r\)와 같다. 이는 다음과 같이 알 수 있다(슈바르츠). \(\mathfrak M\)은 \(\mathfrak Z\)-좌측가군이자 \(\mathfrak Z\)-우측가군, 따라서 \(\mathfrak Z\)의 \(\mathfrak Z\) 안에서의 표현가군이다. \(\mathfrak M\)은 서로 다른 \(n\)개 표현을 생성한다. \(\mathfrak Z\)는 제1종 가환체이므로 \(\mathfrak M\)의 랭크는 적어도 \(n\)이다. 나머지 증명도 꼭 같다.

\noindent\textbf{2.} 다음이 성립한다.
\[
        \mathfrak G\cong\mathfrak G^*/\mathfrak Z^*,
        \qquad\text{즉}\qquad
        \mathfrak T^*=\mathfrak Z^*.
\]
다시 말해 변환으로 \(\mathfrak Z\) 위에 항등사상을 유도하는 것은 \(\mathfrak Z\)의 원소들뿐이다.

실제로 \(\mathfrak Z\)의 모든 \(z\)에 대하여 \(tz=zt\)라고 하고
\[
        t=z_0+z_1u_S+\cdots+z_{n-1}u_T
\]
로 쓰자. 그러면
\[
        zt=zz_0+\cdots+zz_{n-1}u_T
\]
이고
\[
\begin{aligned}
tz
 &=z_0z+\cdots+z_{n-1}u_Tz\\
 &=z_0z+z^{S^{-1}}z_1u_S+
   \cdots+z^{T^{-1}}z_{n-1}u_T.
\end{aligned}
\]
합이 직접합이므로
\[
        (z-z^{S^{-1}})z_1u_S=0,
        \ldots,
        (z-z^{T^{-1}})z_{n-1}u_T=0.
\]
원소 \(u_S,\ldots\)가 가역이므로
\[
        (z-z^{S^{-1}})z_1=0,
        \ldots,
        (z-z^{T^{-1}})z_{n-1}=0
        \qquad\text{모든 }z\text{에 대하여}
\]
이다. 이는 랭크가 \(n\)이라는 데서도 따라온다. 그런데 정의에 따라 \(E\)와 다른 각 \(S,\ldots\)에는 \(z-z^{S^{-1}}\ne0\)인 \(z\)가 적어도 하나 있고, 마찬가지로 \(\bar z-\bar z^{S^{-1}}\ne0\)인 \(\bar z\)가 있다. 따라서 \(z_i=0\) \((i=1,\ldots,n-1)\)이고 \(t=z_0\in\mathfrak Z^*\)이다. 증명 끝.

\noindent\textbf{3.} \(u\)들은 곱셈법칙
\[
        u_Su_T=a_{S,T}u_{ST}
\]
을 만족하며, \(a_{S,T}\)들은 교차 결합법칙
\[
        a_{R,S}a_{RS,T}
        =a_{S,T}^{R^{-1}}a_{R,ST}
\]
을 만족한다. 이 \(a_{S,T}\)들이 소인자계이다.

실제로 2에 의해
\[
        \mathfrak Z^*u_S\mathfrak Z^*u_T
        =\mathfrak Z^*u_{ST},
\]
따라서
\[
        u_Su_T=a_{S,T}u_{ST}
\]
이고, 여기서 \(a_{S,T}\in\mathfrak Z^*\), 즉 \(a_{S,T}\in\mathfrak Z\)이고 \(a_{S,T}\ne0\)이다. 다시
\[
        u_R(u_Su_T)=(u_Ru_S)u_T
\]
이므로
\[
\begin{aligned}
u_Ra_{S,T}u_{ST}
 &=a_{S,T}^{R^{-1\cdot}}u_Ru_{ST}\\
 &=a_{S,T}^{R^{-1}}a_{R,ST}u_{RST}\\
 &=a_{R,S}u_{RS\cdot}u_T\\
 &=a_{R,S}a_{RS,T}u_{RST}.
\end{aligned}
\]
이것이 앞서 적은 곱셈법칙이다.

\noindent\textbf{4. 역.} 교차곱은 어떤 \(\mathfrak K_r\)와 같다. 직접합
\[
        \mathfrak M
        =\mathfrak Zu_E+\cdots+\mathfrak Zu_T
\]
에 관계식
\[
        zu_S=u_Sz^S,
        \qquad
        u_Su_T=a_{S,T}u_{ST}
\]
으로 환 구조를 준다. 여기서 \(a_{S,T}\)들은 교차 결합관계들을 만족하므로 \(u_Ru_Su_T\)가 일의적으로 정의된다. 곱셈은
\[
\begin{array}{ll}
1.&z_1\cdot z_2u_S=z_1z_2\cdot u_S,\\[2pt]
2.&zu_S\cdot u_T=z\cdot u_Su_T
\end{array}
\]
로 정한다. 이로써 \(\mathfrak M\)은 결합환이 된다.

\noindent\textbf{역정리 2.} \emph{교차곱 \(\mathfrak Z\cdot\mathfrak G\)는 중심이 \(\mathsf P\)인 양쪽단순환이다.} (R. 브라우어의 착상에 따름.) 증명은 다음 보조정리들에 바탕을 둔다.

\noindent\textbf{보조정리 1.} \emph{\(\mathfrak Z\)의 모든 원소와 교환하는 원소는 \(\mathfrak Z\)에 속한다.}

\noindent\emph{증명.} \(c_{\lambda_S}\in\mathfrak Z\)인
\[
        w=\sum_Sc_{\lambda_S}u_S
\]
가 \(\mathfrak Z\)의 모든 \(z\)에 대하여 \(wz=zw\)를 만족한다고 하자. 그러면
\[
\begin{aligned}
\sum_Szc_{\lambda_S}u_S
 &=\sum_Sc_{\lambda_S}u_Sz\\
 &=\sum_Sc_{\lambda_S}z^{S^{-1}}u_S.
\end{aligned}
\]
따라서 모든 \(z\)와 \(S\)에 대하여
\[
        zc_{\lambda_S}=c_{\lambda_S}z^{S^{-1}},
        \qquad\text{즉}\qquad
        c_{\lambda_S}(z-z^{S^{-1}})=0.
\]
여기서는 \(u\)들의 선형독립성을 썼다. \(S\ne E\)이면 \(z-z^{S^{-1}}\ne0\)인 \(z\)가 적어도 하나 있으므로 \(c_{\lambda_S}=0\)이다. 따라서 \(w\in\mathfrak Z\)이다.

\noindent\textbf{보조정리 2.} \emph{\(\mathfrak Z\)를 좌우 연산자 영역으로 할 때 \(\mathfrak Zu_S=u_S\mathfrak Z\)와 양쪽동형인 각 \(\mathfrak Z\)-쌍가군 \(\mathfrak A\)는 \(\mathfrak Zu_S\)와 동일하다.}

\noindent\emph{2의 증명.} \(\mathfrak A\)는 좌측가군으로서 \(\mathfrak Zu_S\)와 동형이다. 이 동형에서 \(u_S\mapsto a\in\mathfrak A\)라 하면 \(\mathfrak Zu_S\mapsto\mathfrak Za\)이고, 따라서 \(\mathfrak A=\mathfrak Za\)이다. 양쪽동형성에서
\[
        za=az^S
        \tag{1}
\]
이 따라온다. 이는 \(u_Sz\mapsto az\)이기 때문이다. 이로부터 \(\mathfrak A=\mathfrak Zu_S\)가 따라온다. 실제로
\[
        a=u_Su_S^{-1}a=u_Sw
\]
로 두면
\[
        za=zu_Sw=u_Sz^Sw,
\]
한편 (1)에 의해
\[
        za=u_Swz^S.
\]
따라서 \(u_S\)를 곱하여 모든 \(z\in\mathfrak Z\)에 대하여 \(z^Sw=wz^S\)를 얻는다. \(z^S\)도 \(\mathfrak Z\) 전체를 달리므로 \(w\)는 \(\mathfrak Z\)의 모든 원소와 교환하고, 1에 따라 \(w\in\mathfrak Z\)이다.

\noindent\textbf{보조정리 3.} \emph{\(\mathfrak Z\mathbin{\cdot}\mathfrak G\)의 각 \(\mathfrak Z\)-쌍가군은}
\[
        \mathfrak Zu_{S_1}+\cdots+\mathfrak Zu_{S_k}
\]
\emph{꼴이다. 여기서 \(u_{S_1},\ldots,u_{S_k}\)는 \(u_S,\ldots,u_T\) 가운데 서로 다른 임의의 \(k\)개이다.}

\noindent\emph{증명.} 교차곱 \(\mathfrak Z\mathbin{\cdot}\mathfrak G\)는 \(\mathfrak Z\)-쌍가군으로서 완전가약이다. 각 가수 \(\mathfrak Zu_S\)의 랭크가 1이기 때문이다. 또한 모든 가수는 쌍가군으로서 서로 다른 류에 속하므로 서로 양쪽동형이 아니다. 따라서 각 \(\mathfrak Z\)-쌍가군 \(\mathfrak M\)은 완전가약성에 의해 직접가수이고 그 자체가 완전가약이다. 곧
\[
        \mathfrak M=\sum_i\mathfrak A_{S_i},
\]
여기서 \(\mathfrak A_{S_i}\)는 \(\mathfrak Zu_{S_i}\)와 동형이다. 따라서 2에 의해 \(\mathfrak A_{S_i}=\mathfrak Zu_{S_i}\)이다. 모든 \(\mathfrak A_{S_i}\)는 서로 다른 류에 속한다.

\noindent\textbf{보조정리 4.} \emph{\(\mathfrak Z\mathbin{\cdot}\mathfrak G\) 안에서 \(\mathfrak Z\)를 포함하는 각 환, 특히 \(\mathfrak Z\mathbin{\cdot}\mathfrak G\) 자체는 양쪽단순이다. \(\mathfrak Z\)를 포함하는 부분환들은 \(\mathfrak G\)의 부분군들과 일대일로 대응한다.}

\noindent\emph{증명.} \(\mathfrak Z\)를 포함하는 각 환 \(\mathfrak o\)는 \(\mathfrak Z\)-쌍가군이므로 3에 따라
\[
        \mathfrak o=\sum_i\mathfrak Zu_{S_i}
\]
꼴이다. \(\mathfrak o\)가 환이므로 두 \(u_{S_i}\)의 곱도 다시 \(\mathfrak o\)에 속하고, 따라서 이에 대응하는 \(\mathfrak Z\)-가군도 \(\mathfrak o\)에 속한다. 원소가 유한개이므로 \(u_{S_i}\)들은 인자계를 무시하면 하나의 군을 이룬다. 즉
\[
        \mathfrak o=\mathfrak Z\mathbin{\cdot}\mathfrak H
\]
이고 \(\mathfrak H\)는 \(\mathfrak G\)의 부분군이다. 더 정확히는 \(\mathfrak Zu_{S_i}\)들이 곱셈계를 이룬다. 이제 \(\mathfrak o\)의 \(0\)이 아닌 임의의 양쪽아이디얼 \(\mathfrak a\)를 택하자. 이는 \(\mathfrak o\)-쌍가군, 따라서 \(\mathfrak Z\)-쌍가군이자 \(\mathfrak o\)의 부분가군이다. 그러므로 \(\mathfrak Zu_{S_i}\) 하나와 함께 군 \(\mathfrak H\) 전체를 포함하며, 따라서 \(\mathfrak o=\mathfrak Z\cdot\mathfrak H\)이다.

\noindent\textbf{보조정리 5.} \emph{\(\mathsf P\)는 \(\mathfrak Z\cdot\mathfrak G\)의 중심이다.}

\noindent\emph{증명.} \(w\)가 \(\mathfrak Z\cdot\mathfrak G\)의 모든 원소와 교환하면 특히 \(\mathfrak Z\)의 모든 원소와 교환하므로 1에 따라 \(w\in\mathfrak Z\)이다. 이제 \(wu_S=u_Sw^S\)이므로, \(w\)가 교환하면 \(w=w^S\)이다. 따라서 \(w\in\mathsf P\)이다.

\noindent\emph{논평 1.} 이 증명은 \(\mathfrak Z\)가 \(\mathsf P\) 위의 무한 대수적 갈루아체인 경우에도 정리가 성립함을 보여 준다. 완전가약성에 관한 정리들이 그대로 성립하기 때문이다. (크룰, Zahlringe, Math. Ann. 99.) 이때 \(\mathfrak Z\mathbin{\cdot}\mathfrak G\)는 모든 유한합
\[
        \sum_{S_\lambda}u_{S_\lambda\cdot}c_{S_\lambda}
\]
의 계이다.

다만 4에서는 \(\mathfrak o=\mathfrak Z\mathbin{\cdot}\mathfrak H\)라고 가정해야 한다. 이는 전체 곱 \(\mathfrak Z\mathbin{\cdot}\mathfrak G\)에 대해서는 성립한다. 그러나 \(\mathfrak o\)가 \(\mathfrak Z\) 위에서 무한이면 \(u_{S_i}u_{S_j}\)가 속한다는 사실만으로 그 계가 군을 이룬다고 더는 결론지을 수 없기 때문이다.

\noindent\emph{논평 2.} 여기서 직접 정초할 때는 §§25/26에서 전개한 교차 행렬표현 이론을 이어 붙여야 한다. 뒤의 §28 끝에서 사용하는 사실은, \(t\)가 \(\mathfrak K\)의 절대성분수(지수)일 때 기약 교차표현의 차수가 \(t\)라는 것이다. 여기에서 곱셈상수로 도입한 인자계와 문헌에서 통용되는 개념이 동일함을 이 표현론이 보여 준다는 점을 다시 지적해 둔다.

\subsection*{§ 28. 인자계의 곱정리}

다음이 주어졌다고 하자.
\[
\begin{gathered}
\mathfrak K_r
 =\mathfrak Z\mathbin{\cdot}\mathfrak G
 =\mathfrak Zu_E+\cdots+\mathfrak Zu_T,
\qquad
u_Su_T=u_{ST}a_{S,T},\\
\overline{\mathfrak K}_r
 =\overline{\mathfrak Z}\mathbin{\cdot}\overline{\mathfrak G}
 =\overline{\mathfrak Z}\bar u_E+
  \cdots+
  \overline{\mathfrak Z}\bar u_T,
\qquad
\bar u_S\bar u_T=\bar u_{ST}\bar b_{S,T}.
\end{gathered}
\]
여기서 \(\mathfrak Z\)와 \(\overline{\mathfrak Z}\)는 서로 동형인 갈루아체이고, \(\mathfrak G\)와 \(\overline{\mathfrak G}\)는 각각 그 군이다. 그러면
\[
\begin{aligned}
\mathfrak K_r\times\overline{\mathfrak K}_r
 &=\widetilde{\mathfrak Z}
   \mathbin{\cdot}\widetilde{\mathfrak G}
   \times\mathsf P_n\\
 &=\mathfrak A_f\times\mathsf P_n
  =\mathfrak A_j,
\end{aligned}
\]
여기서
\[
        \tilde u_S\tilde u_T
        =\tilde u_{ST}\tilde c_{S,T}
\]
이고 \(\widetilde{\mathfrak Z}\)는 \(\mathfrak Z\) 및 \(\overline{\mathfrak Z}\)와 동형이다. 이는 랭크를 세는 것과 \(\mathfrak Z\)도 분해체라는 사실에서 따라온다.

\noindent\textbf{명제.} \emph{다음이 성립한다.}
\[
        \tilde c_{S,T}
        =\tilde a_{S,T}\tilde b_{S,T},
\]
\emph{여기서 \(\tilde a,\tilde b\)는 \(\widetilde{\mathfrak Z}\) 안에서 각각 \(a,\bar b\)와 동형인 원소이다.}

증명은 \(\mathfrak A_j\)의 절대길이가 \(n\)인 임의의 좌측아이디얼의 자기동형환이 \(\mathfrak A_f=\widetilde{\mathfrak Z}\widetilde{\mathfrak G}\)와 동형이라는 사실, 그리고 알맞은 \(\mathfrak A_j\)의 좌측아이디얼에 대해서는 이 자기동형환을 직접 결정할 수 있다는 사실에 바탕을 둔다. 더 일반적으로 \(n\) 대신 \(s\)를 써서, 편의상 \(\mathfrak K_r\)에 관하여 진술한 다음 정리가 성립한다.

\noindent\textbf{1. 정리 1.} \emph{\(\mathfrak K_{rs}=\mathfrak K_r\times\mathsf P_s\)의 절대길이가 \(n\)인 한 좌측아이디얼 \(\mathfrak l\)의 자기동형환은 \(\mathfrak K_r\)와 동형이다.}

증명은 다음 보조정리들에서 따라온다.

\noindent\textbf{보조정리 1.} \emph{\(\mathfrak K_{rs}\)의 절대길이가 \(n\)인 모든 좌측아이디얼은 서로 연산자동형이다. 따라서 그 자기동형환들은 환으로서 동형이다. 특히 \(\mathfrak b\)가 \(\mathsf P_s\)의 한 단순아이디얼이면}
\[
        \mathfrak l=\mathfrak K_r\times\mathfrak b
\]
\emph{의 절대길이는 \(n\)이다.}

\noindent\emph{증명.} 이 좌측아이디얼들의 절대길이가 \(n\)이라는 것, 즉 이들이 절대비분해 단순아이디얼 \(n\)개로 분해된다는 것은 \(\mathfrak K_{rs}\) 안에서의 길이도 같다는 뜻이다. 곧 \(\mathfrak K_{rs}\) 안에서 같은 수의 단순아이디얼 가수들이 나타난다. 따라서 이들은 서로 연산자동형이고 그 자기동형환들은 환으로서 동형이다. 또 \(\mathfrak K_r=\mathfrak Z\cdot\mathfrak G\)의 절대길이는 \(n\), 즉 랭크가 \(n^2\)이다. 따라서 \(\mathfrak K_r\times\mathsf P_s\)의 절대길이는 \(n\cdot s\)이고, \(\mathfrak K_r\times\mathfrak b\)의 절대길이는 \(n\)이다.

\noindent\textbf{보조정리 2.} \emph{\(\mathfrak o\)가 단위원을 갖는 환이고 \(\mathfrak l=\mathfrak o e_1\)이 직접가수인 좌측아이디얼, 따라서 \(e_1\)이 멱등원이면 \(\mathfrak l\)의 자기동형환은 \(e_1\mathfrak o e_1\)과 동형이다.}

\noindent\emph{증명.} 원소 \(e_1re_1\)을 곱하면 \(\mathfrak l\)에서 자기 자신으로 가는 연산자준동형을 얻는다. 서로 다른 원소 \(e_1re_1\)은 \(e_1\)에 서로 다른 원소를 대응시키므로 서로 다른 준동형을 준다. 모든 준동형은 이런 식으로 생긴다. 실제로
\[
        e_1\mapsto re_1,
        \qquad
        re_1=e_1re_1
\]
인데, 뒤의 등식은 \(e_1e_1=e_1\)이기 때문이다. 또한 \(e_1\mathfrak o e_1\) 안의 합과 곱은 대응하는 준동형들의 합과 곱에 대응한다.

\noindent\textbf{보조정리 3.} \emph{행렬단위 \(c_{ik}\)를 써서}
\[
        \mathsf P_s=\sum_1^s c_{ik}\mathsf P,
        \qquad
        \mathfrak b=c_{11}\mathsf P+
        \cdots+c_{s1}\mathsf P
\]
\emph{로 두면 \(\mathfrak K_r\times\mathfrak b\)의 자기동형환은 \(\mathfrak K_rc_{11}\)이고, 이는 \(\mathfrak K_r\)와 동형이다.}

\noindent\emph{증명.} \(c_{11}\)은 \(\mathfrak K_r\)의 모든 원소와 교환하고 어떤 \(\mathfrak K_r\)-원소도 소멸시키지 않는다. 실제로 \(c_{11}\)은 \(\mathfrak K_r\times\mathsf P_s\)의 생성 멱등원이고
\[
\begin{aligned}
c_{11}(\mathfrak K_r\times\mathsf P_s)c_{11}
 &=\mathfrak K_r\times c_{11}\mathsf P_sc_{11}\\
 &=\mathfrak K_rc_{11}.
\end{aligned}
\]
이 보조정리들에서 정리 1이 곧바로 따라온다.

\noindent\textbf{2. \(\mathfrak K_r\times\overline{\mathfrak K}_r\)의 결정과 절대길이가 \(n\)인 좌측아이디얼들로의 분해.} 정의에 따라
\[
\begin{aligned}
\mathfrak A_j
 &=\mathfrak K_r\times\overline{\mathfrak K}_r\\
 &=\mathfrak Z\mathbin{\cdot}\mathfrak G
   \times
   \overline{\mathfrak Z}\mathbin{\cdot}\overline{\mathfrak G}.
\end{aligned}
\]
여기서 막대를 붙인 원소와 붙이지 않은 원소는 원소별로 교환한다. 따라서
\[
\begin{aligned}
\mathfrak A_j
 &=(\mathfrak G\times\overline{\mathfrak G})
   (\mathfrak Z\times\overline{\mathfrak Z})\\
 &=\mathfrak G\times\overline{\mathfrak G}
   (\mathfrak Ze_1+\cdots+\mathfrak Ze_n),
\end{aligned}
\]
여기서
\[
        \mathfrak Ze_i=\overline{\mathfrak Z}e_i.
\]
이는 한 갈루아체에 그 자신을 곱하여 얻는 분해이다. 모든 가수는 \(\mathfrak Z\)에 관하여 같은 랭크, 따라서 같은 절대길이를 갖고, \(\mathfrak A_j\)의 절대길이는 \(=n^2\)이므로 각 가수의 절대길이는 \(n\)이다. 이하에서는
\[
        \mathfrak l
        =(\mathfrak G\times\overline{\mathfrak G})
         \mathfrak Ze_1
\]
을 바탕에 둔다. 여기서
\[
        ze_1=e_1\bar z
\]
에 의해 \(\mathfrak Z\)와 \(\overline{\mathfrak Z}\) 사이의 한 특정한 동형사상을 고정한다.

\noindent\textbf{3. \(e_1\mathfrak A_je_1\)의 교차표현과 곱정리의 증명.} 1과 2에 따라 곱 \(\mathfrak K_r\times\overline{\mathfrak K}_r\)의 인자계들은
\[
        e_1\mathfrak A_je_1\equiv\mathfrak A
\]
의 한 교차표현의 인자계들과 동형이다.

\noindent\textbf{정리 2.} \emph{다음이 성립한다.}
\[
\begin{gathered}
\mathfrak A
 =\widetilde{\mathfrak Z}
  \mathbin{\cdot}\widetilde{\mathfrak G},\\
\widetilde{\mathfrak Z}
 =\mathfrak Ze_1
 =\overline{\mathfrak Z}e_1,\\
\widetilde{\mathfrak G}
 =\{\ldots,e_1u_S\bar u_Se_1,\ldots\},
\end{gathered}
\]
\emph{그리고}
\[
\begin{aligned}
e_1u_S\bar u_Se_1\,
e_1u_T\bar u_Te_1
 &=e_1u_{ST}\bar u_{ST}e_1\,
   a_{S,T}e_1\bar b_{S,T}e_1.
\end{aligned}
\]
\emph{따라서 \(\tilde a_{S,T}=a_{S,T}e_1\), \(\tilde b_{S,T}=\bar b_{S,T}e_1\)로 두면 곱정리가 증명된다. \(\tilde a\)와 \(\tilde b\)는 같은 체 \(\widetilde{\mathfrak Z}\) 안에 있다.}

\noindent\emph{증명.} \(e_1\)은 이 환의 단위원이 된다. \(\tilde z=ze_1\)의 켤레들은 \(\tilde z^Se_1\)로 주어진다. 그러므로 원소 \(u_S\bar u_Se_1\)들이 이 치환들을 유도함을 보이면 되고, 그러면 교차표현이 증명된다.

\noindent\textbf{보조정리 1.} 다음이 성립한다.
\[
        e_1u_S\bar u_S
        =u_S\bar u_Se_1
        =e_1u_S\bar u_Se_1.
\]

\noindent\emph{증명.} 다음과 같이 두자.
\[
        u_S^{-1}e_1u_S=e^S,
        \qquad\text{즉}\qquad
        e_1u_S=u_Se^S.
\]\srcfn{1}{이 정의에 따라 \(e_1=e^E\)이다.}
그러면 \(e^S\)들은 \(e_1,e_2,\ldots,e_n\)의 켤레들을 모두 달린다. 이제 다음이 성립한다.\srcfn{2}{\(\sum a_i\bar A_i\bar u_R=\sum a_i\bar u_R\bar A_i^R=\bar u_R\sum(a_i^R)^{R^{-1}}\bar A_i^{-R}\)이기 때문이다.}
\[
        e_1\bar u_R=\bar u_Re^{R^{-1}}.
        \tag{1}
\]
실제로 \(a_i\)가 \(\mathfrak Z\)의 임의의 기저를 달리고 \(\bar A_i\)가 이에 대한 \(\overline{\mathfrak Z}\)의 보완기저를 달리면
\[
        e_1=\sum_i a_i\bar A_i
            =\sum_i a_i^R\bar A_i^R.
\]
따라서 \(a_i^R\)와 \(\bar A_i^R\)에 대해서도 같은 사실이 성립하고, 이로부터 \(e_1\bar u_R=\bar u_Re^{R^{-1}}\), 즉 (1)을 얻는다. (1)에서
\[
\begin{aligned}
e_1u_S\bar u_S
 &=u_Se^S\bar u_S\\
 &=u_S\bar u_Se^{SS^{-1}}\\
 &=u_S\bar u_Se_1
\end{aligned}
\]
이 따라온다. 이로써 보조정리 1이 증명되었다.

\noindent\textbf{보조정리 2.} 다음이 성립한다.
\[
        ze_1(e_1u_S\bar u_Se_1)
        =e_1u_S\bar u_Se_1(z^Se_1).
\]

\noindent\emph{증명.} 실제로
\[
\begin{aligned}
ze_1(e_1u_S\bar u_Se_1)
 &=zu_S\bar u_Se_1
      &&\text{(보조정리 1)}\\
 &=u_Sz^S\bar u_Se_1\\
 &=u_S\bar u_Se_1(z^Se_1)\\
 &=b=e_1b.
\end{aligned}
\]
여기서는 \(z\)가 \(\bar u\) 및 \(e_1\)과 교환함을 썼다. 이제 \(u'=u_S\bar u_Se_1\)로 두자. 보조정리 2는 \(u'\)가 \(\mathfrak Ze_1\)의 치환들을 유도한다는 뜻이므로 정리에서 적은 교차표현이 성립한다.

\noindent\textbf{보조정리 3.} \emph{그 인자계는}
\[
        a_{S,T}e_1\cdot\,\bar b_{S,T}e_1
\]
\emph{과 같다.}

\noindent\emph{증명.}
\[
\begin{aligned}
e_1u_S\bar u_Se_1\,
e_1u_T\bar u_Te_1
 &=u_S\bar u_S\cdot u_T\bar u_Te_1\\
 &=u_{ST}a_{S,T\cdot}\bar u_{ST}\bar b_{S,T}e_1\\
 &=u_{ST}\bar u_{ST\cdot}a_{S,T}\bar b_{S,T}e_1\\
 &=e_1u_{ST}\bar u_{ST}e_1\,
   a_{S,T}e_1\cdot\,\bar b_{S,T}e_1.
\end{aligned}
\]
이로써 2의 정리, 따라서 인자계의 곱정리가 증명되었다.

\noindent\textbf{4. 연관된 인자계류들에 대한 곱정리.}
\[
        \{a_{S,T}\}\cdot\{b_{S,T}\}
        =\{a_{S,T\cdot}b_{S,T}\}.
\]
이는 3에서 곧바로 따라온다. 실제로 새 생성원 \(v,\bar v\)로 옮겨가는 것은 \(v\bar ve_1\)으로 옮겨가는 것이기도 하다. 왜냐하면 \(u=vc\), \(\bar u=\bar v\bar d\)이면
\[
\begin{aligned}
u\bar ue_1
 &=vc\cdot\bar v\bar de_1\\
 &=v\bar v\cdot c\bar de_1\\
 &=v\bar ve_1\cdot c\bar de_1
\end{aligned}
\]
이기 때문이다.

\noindent\textbf{5.} \(\{a\}\)가 \(\{\mathfrak K\}\)의 인자계이고 \(t\)가 \(\{\mathfrak K\}\)의 지수이면 \(\{a\}^t\)는 단위원류와 같다. 이는 \(\mathfrak G\)의 기약 교차표현의 차수가 \(t\)라는 사실에서 따라오며, §24 37쪽에서 \(p_{ik}\)의 표현 뒤에 주어진 슈어의 증명과 꼭 같다.

\noindent\textbf{6.} \(\mathfrak Z\)를 분해체로 갖는 \(\{\mathfrak K\}\)들의 군 \(\mathscr K_{\mathfrak Z}\)와 연관된 인자계류 \(\{a_{S,T}\}\)들의 군은 서로 동형이다. §27의 구성에 따라 \(\{\mathfrak K\}\)와 \(\{a_{S,T}\}\) 사이에는 일대일 대응이 있고, 4에 따라 이 대응은 준동형이므로 동형이다. 특히 단위원류들은 서로 대응한다.

\noindent\emph{교차곱이 \(\mathsf P\), 즉 중심 위의 행렬환과 같을 필요충분조건은 그 인자계가 단위원자계와 연관되는 것이다.}

\subsection*{§ 29. 최소형 주종정리}

\(\mathfrak K_r=\mathfrak Z\cdot\mathfrak G\)를 교차곱이라 하고, \(\mathfrak G^*\)를 내부자기동형에 의해 \(\mathfrak Z\) 전체를 자기 자신으로 옮기는 모든 원소의 군이라 하자. 그 잉여류들을
\[
        \mathfrak Z^*u_E,
        \quad
        \mathfrak Z^*u_S,
        \quad\ldots
\]
와 같이 쓴다.

\noindent\textbf{정의.} \emph{최소형 주종} \(\mathscr S\)은 \(\mathfrak Z^*\) 위의 항등사상을 연장하는 \(\mathfrak G^*\)의 모든 군 자기동형으로 이루어진다. 이 명칭의 까닭은 순환체에 특수화하면 드러난다. §30을 보라.

\noindent\textbf{정리.} \emph{주종의 각 자기동형은}
\[
        u_S\longmapsto
        u_Sb^Sb^{-1}
        =u_Sb^{S-1},
        \qquad b\in\mathfrak Z
\]
\emph{꼴이고, 역으로 이러한 각 대응은 주종의 한 자기동형을 생성한다.}

\noindent\emph{증명.} 군 자기동형이 \(\mathfrak Z\)의 모든 \(z\)에 대하여
\[
        u_S\longmapsto v_S,
        \qquad
        z\longmapsto z
\]
로 주어지고 곱들이 서로 대응한다고 하자. 그러면 이 자기동형은 대응
\[
        \sum_Su_Sc_{\lambda_S}
        \longmapsto
        \sum_Sv_Sc_{\lambda_S},
        \qquad c_{\lambda_S}\in\mathfrak Z
\]
에 의해 \(\mathfrak K_r\)의 한 환 자기동형으로 확장된다. 그런데 그러한 모든 자기동형은 알려진 바와 같이 내부자기동형이다.

\(b\)를 그 생성원이라 하자. 그러면
\[
        v_S=bu_Sb^{-1},
        \qquad
        z=bzb^{-1}.
\]
\(b\)는 모든 \(z\)와 교환하므로 \(\mathfrak Z\)에 속한다. 따라서
\[
        v_S=bu_Sb^{-1}=u_Sb^Sb^{-1}.
\]
역으로 \(b\in\mathfrak Z\)에 대하여 \(v_S=u_Sb^Sb^{-1}\)로 두면 \(v_S=bu_Sb^{-1}\)이다. 그러므로 이는 \(\mathfrak Z\)의 각 원소를 고정하는 \(\mathfrak K_r\)의 한 자기동형이며, 정의에 따라 \(\mathfrak G^*\) 전체도 자기 자신으로 옮긴다. 따라서 이는 주종의 자기동형이다.

\subsection*{§ 30. 순환 분해체에의 특수화}

특히 순환체 \(\mathfrak Z\)로 교차곱을 만들면 알려진 노름정리들을 해석할 수 있다. 먼저 다음이 성립한다.

\noindent\textbf{정리 1.} \emph{\(\mathsf P\) 위의 고정된 순환 분해체 \(\mathfrak Z\)를 갖는 \(\{\mathfrak K\}\)들의 군 \(\mathscr K_{\mathfrak Z}\)는 군}
\[
        \mathsf P^*/N\mathfrak Z^*
\]
\emph{와 동형이다. 여기서 \(N\mathfrak Z^*\)는 \(\mathfrak Z\)의 모든 \(z\ne0\)의 노름 \(Nz\)들로 이루어진 군이다.}

\noindent\emph{증명.} 증명은 교차표현을 정규화하고 그에 관한 몇 가지 보조정리를 쓰는 데 바탕을 둔다.

\noindent\textbf{정의.} 순환 교차표현이
\[
        \mathfrak K_r
        =\mathfrak Z+\mathfrak Zu+
         \cdots+
         \mathfrak Zu^{n-1}
\]
꼴이고
\[
        zu=uz^S,
        \qquad
        u^n=a
\]
이면 이를 \emph{정규화되었다}고 한다. 여기서 \(S\)는 \(\mathfrak G\)의 한 생성원이다. 즉 \(r=0,\ldots,n-1\)에 대하여 더 일반적인 \(u_{S^r}=b_ru^r\), \(b_r\in\mathfrak Z\) 대신 \(u_{S^r}=u^r\)로 둔다.

\noindent\textbf{보조정리 1.} \emph{순환 교차표현은 언제나 정규화되었다고 가정할 수 있다. 이때 \(a\in\mathsf P\)이고, \(\mathsf P\)의 각 \(a\)는 그러한 표현을 준다.}

\noindent\emph{증명.} 첫째 주장은 명백하다. 실제로 \(u^{-1}zu=z^S\)에서
\[
        u^{-r}zu^r=z^{S^r}
\]
이 따라오므로 \(u^r\)가 치환 \(z\mapsto z^{S^r}\)을 생성한다. 또 \(u^n\)은 항등사상 \(z\mapsto z^{S^n}\)을 유도하므로
\[
        u^n=a,
        \qquad
        a\in\mathfrak Z.
\]
그런데 \(u^n\)은 모든 \(u^r\)과 교환하므로 \(a\)도 그러하고, 따라서 \(a^{S^r}=a\)이다. 그러므로 \(a\in\mathsf P\)이다.

따라서 \(a\)가 유일하게 본질적인 곱셈상수이다. 결합조건들은 아무 내용도 갖지 않게 되며, \(\mathsf P\)의 각 \(a\)는 한 결합계를, 따라서 한 \(\mathfrak K_r\)를 준다.

\noindent\textbf{보조정리 2.} \emph{연관된 인자계들의 전체 류 \(\{a_{S,T}\}\) 대신 그 류에 들어 있는 정규화된 인자계들, 즉 정규화된 표현에서 생기는 인자계들만을 고찰해도 충분하다. 이 부분집합을 정규화된 류라 하자. 전체 류들의 군은 정규화된 류들의 군과 동형이다.}

\noindent\emph{증명.} 보조정리 1에 따라 어느 정규화된 류도 비어 있지 않으므로 대응은 일대일이다. 연산은 임의의 대표원으로 정의되므로 이전과 같다.

\noindent\textbf{보조정리 3.} \emph{정규화된 인자계들의 군은 곱셈군 \(\mathsf P^*\)와 동형이고, 정규화된 류들의 군은}
\[
        \mathsf P^*/N\mathfrak Z^*
\]
\emph{와 동형이다.}

\noindent\emph{증명.} 보조정리 1에 따라 각 정규화된 인자계는
\[
        [1,1,\ldots,a,a,\ldots]
\]
꼴이다. 실제로
\[
        u^\nu\cdot u^\mu=u^{\nu+\mu}
        \qquad(\nu+\mu<n),
\]
그리고
\[
        u^\nu\cdot u^\mu=au^\varrho
        \qquad
        (\nu+\mu=n+\varrho,
         \ 0\leqq\varrho\leqq n-2).
\]
대응
\[
\begin{gathered}
[1,\ldots,a,a,\ldots]\mapsto a,\\
[1,1,\ldots,b,b,\ldots]\mapsto b,\\
[1,\ldots,ab,\ldots]\mapsto ab
\end{gathered}
\]
은 \(\mathsf P^*\) 위로의 군 동형사상이다. 다시 보조정리 1에 따라 \(a\)가 \(\mathsf P^*\)의 모든 원소를 달리기 때문이다.

서로 연관된 두 정규화된 인자계는 \(c\in\mathfrak Z\)에 대한 변환
\[
        v=uc
\]
으로 서로 옮겨간다. 그런데
\[
        v=uc,
        \qquad
        v^2=u^2c^Sc,
\]
그리고
\[
        v^r=u^r c^{S^{r-1}}c^{S^{r-2}}\ldots c,
        \qquad
        v^n=u^nN(c).
\]
따라서 \(u^n=a\)이면 \(v^n=aN(c)\)이다. 그러므로 한 정규화된 류는
\[
        [1,\ldots,aN(c),aN(c),\ldots]
\]
꼴의 인자계들로 이루어지며, 여기서 \(c\)는 \(\mathfrak Z\)의 모든 \(0\)이 아닌 원소를 달린다. 이로써 보조정리의 둘째 주장이 증명되었다.

보조정리 2와 3, 그리고 그곳의 \(\{\mathfrak K\}\)와 인자계류 \(\{a_{S,T}\}\) 사이의 동형을 함께 쓰면 위 정리 1이 따라온다.

\noindent\textbf{정리 2.} \emph{최소형 주종의 군은}
\[
        N(c)=1
\]
\emph{을 만족하는 \(\mathfrak Z\)-원소 \(c\)들의 군과 동형이다. 따라서 \(N(c)=1\)이면 \(\mathfrak Z\)의 한 \(b\ne0\)가 존재하여}
\[
        c=b^S\cdot b^{-1}=b^{S-1}
\]
\emph{이다. 오른쪽은 기호적 \((S-1)\)제곱이다.}

\noindent\emph{증명.} 보조정리 3의 증명에서 \(N(c)=1\)이라는 것은 변환 \(v=uc\)가 주종의 한 자기동형을 나타낸다는 것과 같은 조건임을 얻었다. 서로 다른 \(c\)들은 서로 다른 자기동형들에 대응하고 그 역도 성립하며, 곱에는 곱이 대응한다.

\noindent\textbf{논평.} 정리 2는 주종이라는 명칭을 정당화한다. 그 뜻에 맞추어 몫군
\[
        \mathfrak Z^*/\mathscr H
        \quad(\mathscr H=\text{주종})
        \cong
        \mathfrak Z^*/\mathfrak Z^{*S-1}
\]
을 최소형 종들의 군이라 부를 수 있다. 이 군은 \(N(\mathfrak Z^*)\)와 동형이다. 따라서 정리 1과 함께 쓰면
\[
\begin{aligned}
\text{최소형 종군}
 &=\mathfrak Z^*/\mathscr H
  \cong\mathfrak Z^*/\mathfrak Z^{*S-1}
  \cong N(\mathfrak Z^*)\\
 &=\mathfrak Z\text{에 속하는 }\{\mathfrak K\}\text{들의 군}
  =\mathscr K_{\mathfrak Z}
  \cong\mathsf P^*/N(\mathfrak Z^*)
\end{aligned}
\]
이다.

\subsection*{§ 31. 순환 특수 경우의 응용}

순환 교차표현은 §20에서 증명한 두 정리, 곧 모든 유한체가 가환이라는 정리와 실수체의 진정한 비가환 체 확대는 사원수체뿐이라는 정리에 대한 새 증명을 준다. 이 증명들은 최소형 유체론의 증명이라고 부를 수도 있을 것이다.

\noindent\textbf{1. 유한개의 원소로 이루어진 모든 체는 가환이다.} (§20 정리 10.)

증명은 가환 유한체, 즉 갈루아체에 관한 다음 두 가지 잘 알려진 사실에 바탕을 둔다.

\noindent\textbf{1.} \emph{갈루아체는 그 각 부분체 위에서 순환체이다.}

\noindent\textbf{2.} \emph{\(\mathfrak Z\)가 갈루아체, 따라서 \(\mathsf P\) 위의 순환체이면 \(\mathsf P\)의 각 원소는 어떤 \(\mathfrak Z\)-원소의 노름이다. 즉}
\[
        \mathsf P^*=N(\mathfrak Z^*).
\]

이제 \(\mathfrak K\)가 유한하고 \(\mathsf P\)가 그 중심이며 \(\mathfrak Z\)가 한 순환 분해체라고 하자. 그러면 §30 끝과 위 1에 따라 \(\mathfrak Z\)에 속하는 \(\{\mathfrak K\}\)들의 군 \(\mathscr K_{\mathfrak Z}\)는
\[
        \mathsf P^*/N\mathfrak Z^*
\]
와 동형이다. 따라서 2에 의해 이는 단위원 \(\{\mathsf P\}\)와 같다. 이는 중심이 \(\mathsf P\)이면서 \(\mathsf P\)와 다른 체는 없다는 뜻이다. 곧 \(\mathfrak K=\mathsf P\)이다.

\noindent\textbf{2. 실수체 \(\mathsf P\) 위의 진정한 비가환 체 확대는 사원수체뿐이다.} (§20 정리 9.)

\noindent\emph{증명.} \(\mathsf P\)는 중심이어야 한다. 실제로 \(\mathsf P\)의 가환 확대체로는 복소수체 \(\mathfrak Z\)만 있고, \(\mathfrak Z\)는 대수적으로 닫혀 있으므로 이를 중심으로 갖는 비가환 체 확대는 없다. 동시에 분해체로도 \(\mathfrak Z\)만을 생각하면 되고, \(\mathfrak Z\)는 \(\mathsf P\) 위의 순환체이다. 또한 모든 양수는 노름이므로 군
\[
        \mathsf P^*/N\mathfrak Z^*
\]
의 위수는 2이다. 따라서 \(\{\mathsf P\}\) 말고 오직 하나의 류 \(\{\mathfrak K\}\)만 있으며, 이것이 사원수체이다. 동시에 연관된 인자계의 류에서 \(-1\)을 정규화된 인자계로 택하면 정규형을 얻는다. 실제로
\[
        \mathfrak K=\mathfrak Z+\mathfrak Zu,
        \qquad
        u^2=-1.
\]
따라서
\[
        \mathfrak K
        =\mathsf P+\mathsf Pi+(\mathsf P+\mathsf Pi)u,
        \qquad
        i^2=-1,
        \qquad
        u^2=-1.
\]

\noindent\textbf{1에 대한 논평.} 잘 알려진 사실 2도 §30 끝을 바탕으로 다음과 같이 증명할 수 있다. \(\mathsf P^*\)의 모든 원소, 그리고 그 원소들만이 주종의 항등자기동형을 생성하며 \(S-1\)을 적용하면 단위원으로 옮겨가므로
\[
        \mathscr H=\mathfrak Z^{*S-1}
        \cong\mathfrak Z^*/\mathsf P^*.
\]
따라서 \(\mathfrak Z\)가 유한이면 \(\mathscr S\)의 원소 수는 \(\mathfrak Z^*\)에서 \(\mathsf P^*\)의 지수와 같다. 그러므로 \(\mathfrak Z^*/\mathscr S\)의 원소 수는 \(\mathsf P^*\)의 원소 수와 같고, 따라서
\[
        N(\mathfrak Z^*)
        =\mathfrak Z^*/\mathscr S
\]
에 대해서도 마찬가지이다. 결국
\[
        \mathsf P^*=N(\mathfrak Z^*).
\]
(종의 수가 애매류의 수와 같다는 명제와 비교하라. 다만 이 수론적 표현에는 어디에서나 ``최소형''이라는 말을 덧붙여야 한다.)

\noindent\textbf{3.} \emph{\(\mathsf P\)가 \(p\)-진 바탕체이고 \(\mathfrak Z\)가 \(\mathsf P\) 위의 \(n\)차 순환체이면, \(\mathfrak Z\)에 속하는 \(\{\mathfrak K\}\)들의 군 \(\mathscr K_{\mathfrak Z}\)는 \(\mathfrak Z\)의 갈루아군과 동형이다.}

이는 §30 끝과
\[
        \mathsf P^*/N\mathfrak Z^*
        \cong\mathfrak G
\]
라고 말하는 국소 유체론에서 곧바로 따라온다.

\noindent\textbf{3에 대한 논평.} 이로부터 \(\mathscr K_{\mathfrak Z}\)의 원소들의 군위수, 즉 \(\{\mathfrak K\}\)의 지수가 그 인덱스와 같아짐이 따라온다. 실제로 \(\mathscr K_{\mathfrak Z}\)의 한 생성원의 지수는 \(n\)이다. 따라서 그 인덱스도 \(n\)의 약수이면서 지수의 배수이므로 \(n\)이다. 그러므로 군 \(\mathscr K_{\mathfrak Z}\)는 인덱스가 \(n\)의 약수이고 \(\mathfrak Z\)가 분해체인 류들만을 포함한다. 더 나아가 H. 하세는 Math. Ann. 104에서 \(\mathscr K_{\mathfrak Z}\)가 인덱스가 \(n\)의 약수인 모든 류 \(\{\mathfrak K\}\)로 이루어짐을 보였다. 따라서 \(p\)-진 \(\mathsf P\) 위의 \(n\)차 순환 분해체 \(\mathfrak Z\)는 어느 것이나 같은 체류군 \(\mathscr K_{\mathfrak Z}\)를 준다.

다음 두 사실,

\noindent\textbf{1.} \(\mathsf P\) 위의 모든 \(\{\mathfrak K\}\) 가운데 인덱스가 \(n\)의 약수인 것들의 군은 위수 \(n\)인 순환군이라는 사실,

\noindent\textbf{2.} \(\mathsf P\) 위의 각 \(n\)차 순환 분해체가 \(\{\mathfrak K\}\)들의 전체 군을 생성한다는 사실은

국소 유체론의 역정리의 본질적인 내용으로 볼 수 있다.

\clearpage

\clearpage
\setcounter{footnote}{0}
\phantomsection
\addcontentsline{toc}{section}{카프페러 논문과 E. 뇌터 공동 부록}
%
%
%
%
%
%

\section*{대수함수에 관한 뇌터 기본정리의 필요충분 중복도 조건: H. 카프페러}

\begin{center}
(E. 뇌터와 공동 집필한 부록 포함)

\emph{Math. Ann. 97 (1927), 559--567쪽}
\end{center}

다음 내용은 뇌터 기본정리를 필요충분 중복도 조건의 형태로 강화한 것이다\srcfn{1)}{상세한 서술은 같은 제목으로 쓴 필자의 논문을 보라. 이 논문은 A. 뢰비(프라이부르크)가 기획한 하이델베르크 학술원 보고 1927년 특별호 \glqq{}대수학 기여\grqq{}에 여러 저자의 글과 함께 실렸다.}. 여기와 이하에서는 아이디얼론의 개념이나 정리를 사용하지 않지만, 이 정리의 아이디얼론적 이해\srcfn{2)}{예컨대 B. L. van der Waerden, \emph{Zur Nullstellentheorie der Polynomideale}, Math. Annalen 96 (1926), 203쪽을 보라.}에 기대어 뇌터 기본정리가 포괄하는 사실들을 다음의 서로 분리된 두 명제로 묶는다\srcfn{3)}{인용한 논문의 § 1에서는 두 명제를 각각 직접, 곧 아이디얼론의 개념 없이 증명하였다.}.

\textbf{정리 Ia.} \(x,y\)에 관한 서로소인 두 다항식 \(\varphi,\psi\) --- 계수는 임의의 복소수여도 된다 --- 의 각 공통영점 \(x=\alpha,\ y=\beta\)에 대하여 다음 성질을 지닌 최소의 양의 정수 \(\varrho\)가 존재한다. \(\mathfrak{p}=(x-\alpha,y-\beta)\)이고 \(\sigma\)가 \(\varrho\)보다 큰 임의의 정수일 때, 모듈 \((\varphi,\psi,\mathfrak{p}^\varrho)\)이 나타내는 다항식 전체는 모듈 \((\varphi,\psi,\mathfrak{p}^\sigma)\)이 나타내는 다항식 전체와 정확히 같다\srcfn{4)}{문자 \(\mathfrak{p}\)와 \(\mathfrak{q}\)는 여기서 나타내는 것들이 최근 문헌에서 쓰이는 소아이디얼 \(\mathfrak{p}\) 및 으뜸아이디얼 \(\mathfrak{q}\)의 개념과 내용상 일치함을 상기시키려고 택하였다.}.

\textbf{정리 Ib.} 이렇게 개념상 유일하게 정해지는 다항식 전체를 \(\mathfrak{q}\)라 쓰고, 서로 다른 \(s\)개의 영점 가운데 \(i\)번째 영점 \(x=\alpha_i,\ y=\beta_i\)에 관한 경우에는 \(\mathfrak{q}_i\)라 쓰면 다음이 성립한다.

\(x,y\)에 관한 주어진 다항식 \(K\)가 가분성 \(K\equiv0(\varphi,\psi)\)을 갖기 위한 필요충분조건은 다음 \(s\)개 합동식이 동시에 성립하는 것이다.
\[
K\equiv0(\mathfrak{q}_1),\quad K\equiv0(\mathfrak{q}_2),\quad\ldots,\quad K\equiv0(\mathfrak{q}_s).
\]

역사적으로 보면 --- 이 문제에 관한 막스 뇌터의 기초 논문 서론을 참조하라\srcfn{5)}{Max Noether, \emph{Über einen Satz aus der Theorie der algebraischen Funktionen}, Math. Annalen 6 (1872).} --- 기본정리는 다음의 근본 문제를 풀려는 노력에서 나왔다. \emph{\(x,y\)에 관한 주어진 각 다항식 \(K\)가 \((\varphi,\psi)\)로 나누어지는지 아닌지를 판정할 수 있어야 한다.} 즉 \(x,y\)에 관한 두 다항식 \(\lambda,\mu\)가 더 존재하여
\[
K=\lambda\cdot\varphi+\mu\cdot\psi
\]
가 \(x,y\)에 관한 항등식이 되는지도 판정해야 한다.

이 문제는 뇌터의 정리만으로 해결되지 않는다. 그 정리는 제기된 문제를 더 간단한 문제, 곧 개별 합동식 \(K\equiv0(\mathfrak{q})\)이 성립하기 위한 필요충분조건을 묻는 문제로 \emph{환원}하는, 물론 근본적인 환원으로 보아야 한다\srcfn{6)}{이 환원은 초등정수론에서 잘 알려진 다음 환원과 대략 비교할 수 있다. 자연수 \(A\)가 자연수 \(B\)로 나누어진다는 성질은 \(A\)가 각각 \(p_1^{a_1},p_2^{a_2},\ldots,p_s^{a_s}\)로 나누어진다는 성질과 동치이다. 여기서 \(p_1,p_2,\ldots,p_s\)는 \(B=p_1^{a_1}p_2^{a_2}\ldots p_s^{a_s}\)의 서로 다른 소인수이다. 이제 소수 거듭제곱에 대응하는 것으로 바로 본문에서 \(\mathfrak{q}\)로 나타낸 \glqq{}으뜸아이디얼\grqq{}의 개념을 이해해야 한다. E. Noether, \emph{Idealtheorie in Ringbereichen}, Math. Annalen 83 (1921), 서론을 보라.}. 그렇다면 언제 \(K\equiv0(\mathfrak{q})\)인가? 이 마지막 물음에 답해야 비로소 M. 뇌터가 처음 제기한 문제도 해결된다. 이것이 이 논문의 목표이며, 실제로 다음과 같이 달성된다. \(K\equiv0(\mathfrak{q})\)이 성립하기 위한 필요충분조건인 유한개의 중복도 조건을 제시하며, 이 조건들은 실제로도 어렵지 않게 확인할 수 있다. \(K\equiv0(\mathfrak{q})\)의 충분조건이지만 필요조건은 아닌 다음 조건은 잘 알려져 있고 문헌에서 흔히 사용되고 인용된다.

합동식 \(K\equiv0(\varphi,\psi)\)이 성립하려면, \(\varphi=0,\ \psi=0\)의 각 교점이 다항식 \(K\)에서 \glqq{}충분히\grqq{} 높은 중복도로 나타나는 것으로 \emph{충분하다}.

이 명제는 정리 I에 특수한 경우로 포함된다. 실제로 \(K\)에서 중복도가 \(\varrho\)인 점은 \(K\equiv0(\mathfrak{p}^\varrho)\)과 같은 뜻이고, 그러면 더욱이 \(K\equiv0(\varphi,\psi,\mathfrak{p}^\varrho)\), 곧 \(K\equiv0(\mathfrak{q})\)이다. 그러나 그 역은 성립하지 않는다.

예고한 문제의 해법은 간단히 다음과 같이 특징지을 수 있다. 각 \(\mathfrak{q}\)에는 \(t\)개의 \emph{양의} 정수 \((g_1),(g_2),\ldots,(g_t)\)의 계를 대응시킬 수 있고, 모듈 \((K,\mathfrak{q})\)에는 같은 수의 \emph{양수 또는 음수인} 정수 \((K_1),(K_2),\ldots,(K_t)\)를 대응시킬 수 있어서 다음 \emph{주정리}가 성립한다.

\textbf{정리 II.} \emph{\(K\equiv0(\mathfrak{q})\)이 성립하기 위한 필요충분조건은 다음 \(t\)개 조건이 동시에 성립하는 것이다.}
\[
(K_1)\geqq(g_1),\quad (K_2)\geqq(g_2),\quad \ldots,\quad (K_t)\geqq(g_t).
\]
\emph{이때}
\[
(g_1)+(g_2)+\ldots+(g_t)=\mu
\]
\emph{이며, \(\mu\)는 \(\mathfrak{q}\)에 대응하고 종결식으로 정의되는 교점중복도를 뜻한다.}

증명을 위하여 일반성을 잃지 않고 \(x=0,\ y=0\)을 \((\varphi,\psi)\)에서 \(q\)에 대응하는 영점이라 하자. 또 일반적으로 \(A(x,y)\)가 \(x,y\)에 관한 임의의 다항식일 때\srcfn{7)}{\(A(x,y)\)로 분수유리함수 \(C(x,y):D(x,y)\)도 허용하면 \((A)=(C)-(D)\)가 된다. 따라서 이 경우에는 \((A)\)가 음의 값을 가질 수도 있다.}, 기호 \((A)\)는 \(A(0,y)\)에서 \(y\)가 인수로 나타나는 횟수를 뜻한다고 하자. 다만 \(A(0,y)\)가 항등적으로 사라지면 \((A)=\infty\)로 둔다. \(\mathfrak{q}\)에 대응하는 수 \((g_1),(g_2),\ldots,(g_t)\)는 다음 방정식계로 정의되며, 이 방정식계가 절차의 방법론적 기초이다.

\begin{equation}
\tag{1}
\begin{aligned}
f_1\cdot g_1(0,y):y^{(g_1)} - g_1\cdot f_1(0,y):y^{(g_1)} &= x\cdot h_1\\
f_2\cdot g_2(0,y):y^{(g_2)} - g_2\cdot f_2(0,y):y^{(g_2)} &= x\cdot h_2\\
&\vdots \qquad\qquad\vdots\qquad\qquad\vdots\\
f_r\cdot g_r(0,y):y^{(g_r)} - g_r\cdot f_r(0,y):y^{(g_r)} &= x\cdot h_r
\end{aligned}
\end{equation}

이 계는 오직 주어진 다항식쌍 \(\varphi,\psi\)만으로 구성된다. 실제로 다항식 \(f_1,g_1\)은 순서를 바꿀 수 있다는 점을 제외하면 각각 \(\varphi,\psi\)와 같다고 정의한다. 그 순서는 \((f_1)\geqq(g_1)\)이 되도록 정한다. 이렇게 정할 때, 그리고 그때에만 \(h_1\)이 다항식이 된다.

다항식 \(f_2,g_2\)는 부조건 \((f_2)\geqq(g_2)\) 아래 각각 \(h_1,g_1\)과 같다고 정의한다. 일반적으로 다항식쌍 \(f_r,g_r\)은 부조건 \((f_r)\geqq(g_r)\) 아래 각각 \(h_{r-1},g_{r-1}\)과 같다고 정의한다.

방정식계 (1)은 설명한 방식으로 끝없이 계속할 수 있다. 그러나 주정리에서는 처음 \(t\)개의 방정식만 만들면 되며, \(t\)는 다음 사실로 정해진다.

\textbf{보조정리 I.} 자연수 \(t\)가 존재하여 [더구나 \(t\leqq\mu\)이며, 여기서 \(\mu\)는 점 \(x=0,\ y=0\)이 \(\varphi=0,\ \psi=0\)에 나타나는 교점중복도이다], (1)로 정의되는 처음 \(t\)개의 수 \((g_1),(g_2),\ldots,(g_t)\)는 모두 \emph{0보다 크고}, 동시에 \((g_{t+1})\)과 그 뒤의 모든 수는 \emph{0과 같다}.

실제로 종결식의 곱정리를 고려하면, 점 \(x=0,\ y=0\)에서 \(f_r=0,\ g_r=0\)의 교점중복도는
\[
\mu-(g_1)-(g_2)-\ldots-(g_r)
\]
와 같다. 여기서 \emph{늦어도 \(\mu\)단계 뒤에는 중단됨}이 따라온다. 동시에
\begin{equation}
\tag{2}
\mu=(g_1)+(g_2)+\ldots+(g_t)
\end{equation}
도 얻는다. 문제의 영점이 알려져 있기만 하면 교점중복도를 정하는 유리적이고 실제로도 쉽게 실행할 수 있는 일반 방법을 주기 때문에 이 공식은 주목할 만하다.

이렇게 수 \((g_i)\)를 알게 된 뒤에는 수 \((K_i)\)를 (1)에 대응하는 다음 방정식계로 정한다.
\begin{equation}
\tag{3}
\begin{aligned}
K_1\cdot g_1(0,y):y^{(g_1)} - g_1\cdot K_1(0,y):y^{(g_1)} &= x\cdot K_2\\
K_2\cdot g_2(0,y):y^{(g_2)} - g_2\cdot K_2(0,y):y^{(g_2)} &= x\cdot K_3\\
&\vdots \qquad\qquad\vdots\qquad\qquad\vdots\\
K_r\cdot g_r(0,y):y^{(g_r)} - g_r\cdot K_r(0,y):y^{(g_r)} &= x\cdot K_{r+1}
\end{aligned}
\end{equation}

\(K_1\)은 \(K\)와 같다고 정의한다. \(K_2\)는 명백히 \((K_1)\geqq(g_1)\)일 때 그리고 그때에만 다항식이다. 이 조건이 충족되면 계속하여 \(K_3\)은 \((K_2)\geqq(g_2)\)일 때 그리고 그때에만 다항식이다. 일반적으로 \(K_r\)이 이미 다항식이면 \(K_{r+1}\)은 \((K_r)\geqq(g_r)\)일 때 그리고 그때에만 다항식이다.

이제 주정리(정리 II)의 증명은 다음 보조정리에 근거한다.

\textbf{보조정리 II.} \((f_i,g_i,p^\varrho)=q^{(i)}\)로 두면, 무한 정수열 \(1,2,\ldots\)의 각 \(i\)에 대하여 다음이 성립한다.

합동식
\[
K_i\equiv0(\mathfrak{q}^{(i)})
\]
이 성립하기 위한 \emph{필요충분조건은 다음 두 조건이 동시에 성립하는 것이다.}
\[
(K_i)\geqq(g_i)\quad \text{및}\quad K_{i+1}\equiv0(\mathfrak{q}^{(i+1)}).
\]

이 보조정리의 증명\srcfn{8)}{\glqq{}부록\grqq{}의 주해 I에 있는 증명도 보라.}은 방정식계 (1)과 (3)을 결합하면 얻어진다. 보조정리 I에 따라 \(g_{t+1}\)은 더 이상 원점에서 사라지지 않는다는 것, 따라서 모듈 \(\mathfrak{q}^{(t+1)}\)은 \emph{모든} 다항식으로 이루어져 단위아이디얼이 되고 결국 \(K_{t+1}\equiv0(\mathfrak{q}^{(t+1)})\)이 자명해진다는 것을 고려하면, 보조정리 II를 \(t\)번 적용하여 주정리의 증명을 얻는다.

\begin{center}
\rule{4cm}{0.4pt}
\end{center}

증명은 위에서 인용한 필자의 논문에 있으므로 생략하고, 주정리에 두 가지 추가 명제만 덧붙인다.

\textbf{추가 I.} 주정리의 수 \(t\)는
\[
x^v\equiv0(\mathfrak{q})
\]
가 되게 하는 최소의\srcfn{9)}{따라서 여기서는 정확한 지수가 얻어진다. 훨씬 일반적인 문제를 다룬 대신 상계만 제시한 그레테 헤르만의 추정과 대비된다. Frl. Grete Hermann, \emph{Die Frage der endlich vielen Schritte in der Theorie der Polynomideale} (K. Hentzelt가 남긴 정리들을 사용), Math. Annalen 95 (1926).} 양의 수 \(\nu\)와 같다.

\textbf{추가 II.} 방정식계 (1), (3)과 거기에서 \(x,y\)를 서로 바꾸어 얻는 방정식계들을 적절히 결합하면, \(\mathfrak{q}\)에 \glqq{}속하는\grqq{} 지수 \(\varrho\), 곧 정리 Ia의 최소지수 \(\varrho\)의 정확한 수치를 정하는 유리적 방법을 얻는다.

E. 베르티니\srcfn{10)}{E. Bertini, Math. Annalen 34 (1889).}는 \(\varrho\)의 일반공식을 제시하였다. 그러나 그 공식이 언제나 \(\varrho\)의 최솟값을 주는 것은 아님이 \(\varphi=x^3+y^4;\ \psi=x^2+y^3\)이라는 예를 통하여 인용한 논문에서 밝혀진다. 끝으로 베르티니의 공식 자체도 우리의 주정리로 새롭게 근거 지을 수 있음을 덧붙인다.

\begin{center}
\rule{4cm}{0.4pt}
\end{center}

해당 \(\mathfrak{q}=(\varphi,\psi,\mathfrak{p}^\varrho)\)의 영점이 두 곡선 \(\varphi=0,\ \psi=0\) 가운데 적어도 하나에서 \emph{단순점}이라는 제한을 두면, 주정리의 흥미로운 \emph{특수화}가 따라온다. 예컨대 그 영점이 \(\psi\)에서 단순점이면, 주정리의 수 \(t\)가 \emph{정확히} \(\mu\)이고 주정리의 \(t\)개 부등식을 \emph{하나의} 중복도 조건으로 바꿀 수 있음을 보일 수 있다. 이로부터 다음 일반정리를 얻는다.

\textbf{정리 III.} \(\varphi=0,\ \psi=0\)의 \emph{모든} 교점이 \(\psi\)에서 \emph{단순점}이라고 하자. \(K\)가 \(\psi\)와 서로소라는 조건 아래, 합동식 \(K\equiv0(\varphi,\psi)\)이 성립하기 위한 \emph{필요충분조건}은 \(K\)가 이 모든 점에서 사라지되, 각 점에서 곡선쌍 \(K=0,\ \psi=0\)의 교점중복도가 곡선쌍 \(\varphi=0,\ \psi=0\)의 교점중복도보다 \emph{적어도} 큰 것이다.

정리 III는 또 다른 특별한 이유로 주목할 만하다. 거기서 요구되는 중복도 조건의 충족 여부는 (1)에 근거한 공식 (2)과 무관하게 유한번의 \emph{미분}으로 판정할 수 있다. 이는 필자가 1923년에 제시한 다음 \glqq{}정리\grqq{}를 인용하면 즉시 분명해진다\srcfn{11)}{H. Kapferer, \emph{Über die Multiplizität der Schnittpunkte von zwei algebraischen Kurven}, Jahresber. der D. M.-V. 32 (1923).}.

\(A,B\)가 \(x,y\)에 관한 임의의 서로소 다항식이고 \(P\)가 두 다항식의 공통점으로서 \emph{\(B\)의 단순점}이라 하자. 이 점이 교점으로서 갖는 중복도가 수 \(\mu\)와 같기 위한 필요충분조건은, 같은 점이 곡선쌍 \(A'=0,\ B=0\)에서 갖는 중복도가 \emph{정확히} \(\mu-1\)인 것이다. 여기서
\[
A'=\frac{\partial A}{\partial x}\cdot\frac{\partial B}{\partial y}
   -\frac{\partial A}{\partial y}\cdot\frac{\partial B}{\partial x}.
\]

\begin{center}
\rule{4cm}{0.4pt}
\end{center}

끝으로 주정리(정리 II)는 조금만 고치면 세 변수 \(x,y,z\)에 관한 동차다항식에도 옮길 수 있다. 위의 특수한 경우에는 이 확장이 특히 우아하다. 실제로 다음을 보일 수 있다.

\textbf{정리 III의 추가 명제.} \(\varphi,\psi,K\)를 \emph{세} 변수에 관한 \emph{동차}다항식으로 이해해도, 정리 III는 문구까지 포함하여 전혀 바뀌지 않고 성립한다\srcfn{12)}{세 변수의 동차 합동식 \(K(x,y,z)\equiv0(\varphi(x,y,z),\psi(x,y,z))\)이 성립하기 위한 조건을 묻는 문제는, 강하게 특수화된 가정 아래이기는 하지만 A. Ostrowski의 \emph{Über die Maxwellsche Erzeugung der Kugelfunktionen}, Jahresb. der D. M.-V. 33 (1925)의 핵심 문제이기도 하다. 그 글에서는 \(\varphi\)가 모두 일차인수로 분해되고, \(\varphi\)와 \(K\)가 같은 차수이며, \(\psi=x^2+y^2+z^2\)라고 가정한다. 따라서 여기서 \(\psi\)에는 특이점이 전혀 없으므로, 오스트롭스키 논문의 주정리는 우리의 정리 III의 특수한 경우로 볼 수 있다.}.

\section*{E. 뇌터와 공동 집필한 부록\srcfn{13)}{아이디얼론적 해석은 E. 뇌터에게서, (5)와 (6)의 바탕이 되는 증명은 H. 카프페러에게서 나왔다.}}

방정식계 (1)은 동시에 \emph{\(\mathfrak{q}\)에 대한 완전잉여계}를 정하는 유리적 방법을 준다. 여기서 완전잉여계란 계수영역 \(\mathsf{P}\)에 관하여 \(\mathfrak{q}\)를 법으로 하는 \emph{유한개의 선형독립 잉여류}에 대한 완전대표계를 뜻한다. 이로써 주정리 증명의 각 단계가 매우 투명해진다. 동시에 \emph{종결식중복도 \(\mu\)가 아이디얼 \(\mathfrak{q}\)의 길이와 일치한다}는 사실이 부수적 결과로 얻어진다.

즉 이는 잉여류환의 랭크, 다시 말해 \(\mathsf{P}\)에 관하여 선형독립인 잉여류의 수와 같다\srcfn{14)}{일반성을 잃지 않고 체 \(\mathsf{P}\)가 대수적으로 닫혀 있다고 가정한다. 미분 조건을 제외하면 앞의 모든 정리는 실제로 이 범위에서 성립한다. \(\mathfrak{q}\)의 길이는 \(\mathfrak{p}\)에서 \(\mathfrak{q}\)로 이어지는 아이디얼들의 합성열의 길이로도 정의된다. E. Noether, Jahresber. d. d. Math.-Ver. 34 (1925), 101쪽(사선 쪽매김)을 보라. 더 자세한 것은 H. Grell, Math. Annalen 97 (1927), 529쪽을 보라. 두 중복도 수가 일치함을 \(n\)변수에서 보이는 기존 문헌의 증명들은 매우 복잡하다. 가장 간단한 것은 아직 출판되지 않은 헨첼트의 유고 증명이다.}.

(1)과 보조정리 II의 기호를 쓰면, 완전잉여계는 다음 정리로 특징지어진다.

\textbf{정리 IV.} \emph{\(\mathfrak{q}\)에 대한 한 완전잉여계는 다음 다항식들로 주어진다.}
\[
h_t,\ h_ty,\ldots,h_ty^{(g_t)-1};\ \ldots;\ h_i,\ h_iy,\ldots,h_iy^{(g_i)-1};\ \ldots;\ h_1,\ h_1y,\ldots,h_1y^{(g_1)-1}.
\]
\emph{이때 \(h_i,h_iy,\ldots,h_iy^{(g_i)-1}\)는 각각 \(\mathfrak{q}^{(i)}\)에 대한 \(\mathfrak{q}^{(i+1)}\)의 완전잉여계를 이룬다.}

정의에서 곧바로 \(\mathfrak{q}^{(i+1)}=(\mathfrak{q}^{(i)},h_i)\), 따라서 \(\mathfrak{q}^{(i)}\equiv0(\mathfrak{q}^{(i+1)})\)을 얻는다. 그러므로 다음을 보이면 정리 IV가 증명된다.
\begin{equation}
\tag{4}
h_i x\equiv0(\mathfrak{q}^{(i)});\qquad h_i y^{(g_i)}\equiv0(\mathfrak{q}^{(i)})
\end{equation}
그리고
\begin{equation}
\tag{5}
\text{\(h_iF(y)\equiv0(q^{(i)})\)이면 \(F(y)\equiv0(y^{(g_i)})\)이다.}
\end{equation}

실제로 관계 (4)는 \(h_i,h_iy,\ldots,h_iy^{(g_i)-1}\)의 선형결합들이 \(\mathfrak{q}^{(i)}\)에 대한 \(\mathfrak{q}^{(i+1)}\)의 잉여류를 모두 준다는 뜻이고, (5)는 선형독립성을 나타낸다. \(q\)의 길이가 유한하므로 이제 \(\mathfrak{q}^{(t+1)}\)은 단위아이디얼과 같아진다. 따라서 여기서 정리 IV의 증명과 함께 \(\mathfrak{q}\)의 길이가 \((g_1)+\ldots+(g_t)\)와 같고, (2)에 따라 종결식중복도와 일치한다는 사실도 얻는다.

(4)의 첫째 관계는 (1)에서 바로 읽을 수 있다. 둘째 관계는 다음과 같이 얻는다.

\(h_ig_i\equiv0(\mathfrak{q}^{(i)})\)이고, 따라서 (4)의 첫째 관계에 따라 \(h_ig_i(0,y)\equiv0(\mathfrak{q}^{(i)})\)이다. 다시 말하면
\[
h_iy^{(g_i)}\cdot(g_i(0,y):y^{(g_i)})\equiv0(q^{(i)}).
\]
\(g_i(0,y):y^{(g_i)}\not\equiv0(y)\)이므로 여기서 \(h_iy^{(g_i)}\equiv0(\mathfrak{q}^{(i)})\)이 따라온다.

(5)의 증명에 앞서 (1)에서 바로 읽을 수 있는 다음 \emph{주해}를 둔다. \(f_i,g_i\)의 공약수는 \(y\)로 나누어지지 않는 \(y\)의 다항식 \(d(y)\)뿐일 수 있다.
\[
f_i=d_i(y)\bar f_i,\qquad g_i=d_i(y)\bar g_i,
\]
그리고 \(\bar f_i,\bar g_i\)는 서로소이다. \(d(y)\not\equiv0(\mathfrak{p})\)이므로 \(q^{(i)}=(\bar f_i,\bar g_i,p^\varrho)\)이기도 하다. 이것은 모든 \(\sigma\geqq\varrho\)에 대하여 \((\bar f_i,\bar g_i,\mathfrak{p}^\sigma)\)와 같고, 그러한 각 \(\sigma\)에 대하여 \(\mathfrak{q}\)의 약수이다. 따라서 \(\mathfrak{q}^{(i)}\)는 \((\bar f_i,\bar g_i)\)에서 원점에 속하는 으뜸성분이다.

\(\gamma_1,\delta_1;\ldots;\gamma_r,\delta_r\)를 \(x=0,\ y=0\)과 다른 \((\bar f_i,\bar g_i)\)의 영점들이라 하고, \(\tau_1,\ldots,\tau_r\)를 그에 대응하는 \((\bar f_i,\bar g_i)\)의 으뜸성분들의 지수라 하자. 곱
\[
P=((x-\gamma_1)+u(y-\delta_1))^{\tau_1}\ldots
((x-\gamma_r)+u(y-\delta_r))^{\tau_r}
\]
에서 부정원 \(u\)를 \(\mathsf{P}\)의 한 원소로 특수화하되 \(P\not\equiv0(\mathfrak{p})\)이 유지되게 하자. 특수화한 \(P\)와 \(d(y)\)의 곱을 \(G(x,y)\)라 두면
\[
G(x,y)\not\equiv0(p),\qquad q^{(i)}\cdot G(x,y)\equiv0(f_i,g_i)
\]
가 된다.

가정 \(h_iF(y)\equiv0(\mathfrak{q}^{(i)})\)에서
\[
h_iF(y)G(x,y)\equiv0(f_i,g_i)
\]
가 따라온다. 이것을 (1)과 함께 항등식의 형태로 쓰면
\[
h_iF(y)G=A f_i+B g_i;\qquad
h_ix=(g_i(0,y):y^{(g_i)})\cdot f_i-(f_i(0,y):y^{(g_i)})g_i
\]
이다. 따라서
\[
\bar f_i(xA-F(y)H)\equiv0(\bar g_i);\qquad
H=G(x,y)\cdot g_i(0,y):y^{(g_i)}\not\equiv0(\mathfrak{p})
\]
를 얻는다.

\(\bar f_i,\bar g_i\)가 서로소이므로 계속하여
\[
xA-F(y)H\equiv0(\bar g_i)
\]
이다. 따라서
\[
F(y)H\equiv0(\bar g_i,x)\equiv0(\bar g_i,x,p^\sigma)=(y^{(g_i)},x),
\]
결국 \(F(y)\equiv0(y^{(g_i)})\)이다. 이로써 (5), 따라서 정리 IV가 증명되었다.

\textbf{주해 I.} 관계 (4), (5)는 보조정리 II와 함께 아이디얼 몫 사이의 서로 역인 다음 관계들로 쓸 수 있으며\srcfn{15)}{몫 \(\mathfrak a:\mathfrak b\)는 \(\mathfrak b c\equiv0(\mathfrak a)\)이 되게 하는 다항식 \(c(x,y)\) 전체를 뜻한다.}, 둘째 관계의 증명은 동시에 보조정리 II의 증명이 된다.
\begin{equation}
\tag{6}
\mathfrak{q}^{(i)}:\mathfrak{q}^{(i+1)}=(\mathfrak{q}^{(i)},x)=(y^{(g_i)},x);\qquad
\mathfrak{q}^{(i)}:(\mathfrak{q}^{(i)},x)=\mathfrak{q}^{(i+1)}.\srcfnmark{16)}
\end{equation}
\srcfntext{16)}{이 서로 역인 두 관계 가운데 하나는 기약아이디얼의 일반론에 따라 다른 하나를 함의한다. Macaulay, \emph{Modular Systems}, 72, 73번, Cambridge Tracts 19 (1916)을 보라.}

첫째 관계는 (4), (5)를 한데 묶은 것이다. 둘째 관계에 관해서는 다음을 덧붙인다. (4)에서 \(x\cdot \mathfrak{q}^{(i+1)}\equiv0(\mathfrak{q}^{(i)})\)을 증명하였다. 그 역은 (5)와 비슷하게 얻어진다. \(xM\equiv0(\mathfrak{q}^{(i)})\)에서 위와 같이
\[
xMG\equiv0(f_i,g_i),\quad\text{따라서}\quad
xMG\cdot g_i(0,y):y^{(g_i)}\equiv0(xh_i,g_i)
\]
를 얻는다.

정의상 \(g_i\not\equiv0(x)\)이다. 실제로 \((f_i)\geqq(g_i)\)이다. 따라서
\[
M\cdot(G\cdot g_i(0,y):y^{(g_i)})\equiv0(h_i,g_i)\equiv0(\mathfrak{q}^{(i+1)}),
\quad\text{따라서}\quad M\equiv0(\mathfrak{q}^{(i+1)})
\]
을 얻는다.

이제 보조정리 II의 조건이 필요함은 다음과 같이 따라온다. \(K_i\equiv0(\mathfrak{q}_i)\)에서 \(K_i(0,y)\equiv0(y^{(g_i)},x)\)를 얻으므로 \((K_i)\geqq(g_i)\)이다. 따라서 (3)에 의하여
\[
xK_{i+1}\equiv0(\mathfrak{q}^{(i)}),\quad
\text{따라서 --- (6)의 둘째 관계에 의하여 ---}\quad
K_{i+1}\equiv0(\mathfrak{q}^{(i+1)}).
\]

조건이 충분함은 즉시 읽어 낼 수 있다. 보조정리 II를 유한번 적용하면 주정리를 얻는다. 마찬가지로 (6)의 둘째 관계를 반복하면 추가 I, 곧 \(\mathfrak{q}:x^t=\mathfrak{q}^{(t+1)}\)이 단위아이디얼과 같음을 얻는다. 더구나 \(\mathfrak{q}:x^{t-1}=\mathfrak{q}^{(t)}\)이므로 \(t\)는 그러한 거듭제곱 가운데 가장 낮은 것이다.

\textbf{주해 II.} \(\mathfrak{q}\)로 나누어지지 않는 다항식 \(K\)에 대해서는 다음도 말할 수 있다. \(K\equiv0(\mathfrak{q}^{(i+1)})\)이지만 \(K\not\equiv0(\mathfrak{q}^{(i)})\)이면
\[
(K,q^{(i)})=(h_iy^{(K)},q^{(i)}).
\]
실제로 정리 IV에 따라 \(c(y)\not\equiv0(y)\)인 \(c(y)\)에 대하여 \(K\equiv h_iy^\lambda c(y)\pmod{q^{(i)}}\)이므로
\[
(K,\mathfrak{q}^{(i)})=(h_iy^\lambda,\mathfrak{q}^{(i)})
\]
이다.

이때 \(y^{(g_i)-\lambda}K\equiv0(\mathfrak{q}^{(i)})\)이고, (5)에 따라 이것은 그러한 거듭제곱 가운데 가장 낮은 것이다. 한편 보조정리 II에 따르면 \((g_i)-(K)\)가 이 가장 낮은 거듭제곱이다. 따라서 \(\lambda=(K)\)이다.

\begin{center}
(1926년 10월 24일 접수.)
\end{center}

\clearpage
\setcounter{footnote}{0}
\phantomsection
\addcontentsline{toc}{section}{참고문헌 및 판본 부록}
%
%
%
%

\clearpage
\begin{center}
{\Large\bfseries 참고문헌}
\end{center}

\begin{enumerate}
\setlength{\itemsep}{0.35em}
\item 삼원 사차형식의 형식계 구성에 관하여.\\
Sitz. Ber. d. Physikal.-mediz. Sozietät in Erlangen 39 (1907), 176--179쪽
\item 삼원 사차형식의 형식계 구성에 관하여.\\
J. Reine Angew. Math. 134 (1908), 23--90쪽 및 표 2개
\item \(n\)개 변수 형식의 불변식론에 관하여.\\
J. Ber. d. DMV 19 (1910), 101--104쪽
\item \(n\)개 변수 형식의 불변식론에 관하여.\\
J. Reine Angew. Math. 139 (1911), 118--154쪽
\item 유리함수체.\\
J. Ber. d. DMV 22 (1913), 316--319쪽
\item 체와 유리함수계.\\
Math. Ann. 76 (1915), 161--196쪽
\item 유한군 불변식의 유한성 정리.\\
Math. Ann. 77 (1916), 89--92쪽
\item 임의로 많은 기본형식들의 계의 불변식에 대한 다항식적 표현.\\
Math. Ann. 77 (1916), 93--102쪽. (16번도 참조)
\item 정수적 초월수로 이루어진 가장 일반적인 영역.\\
Math. Ann. 77 (1916), 103--128쪽. (16번도 참조)
\item 동형사상의 함수방정식.\\
Math. Ann. 77 (1916), 536--545쪽
\item 지정된 군을 갖는 방정식.\\
Math. Ann. 78 (1918), 221--229쪽. (16번도 참조)
\item 임의의 미분식의 불변식.\\
Nachr. v. d. Ges. d. Wiss. zu Göttingen 1918, 37--44쪽
\item 불변 변분문제.\\
Nachr. v. d. Ges. d. Wiss. zu Göttingen 1918, 235--257쪽
\item 한 변수 대수함수의 산술이론과 다른 이론들 및 수체론과의 관계.\\
J. Ber. d. DMV 28 (1919), 182--203쪽
\item 이원형식의 정수 계수 불변식계의 유한성.\\
Nachr. v. d. Ges. d. Wiss. zu Göttingen 1919, 138--156쪽
\item 형식론의 급수전개에 관하여.\\
Math. Ann. 81 (1920), 25--30쪽
\item W. 슈마이들러와 공동: 비가환 영역의 가군, 특히 미분식 및 차분식으로 이루어진 가군.\\
Math. Z. 8 (1920), 1--35쪽
\item 전쟁 중 전사한 K. 헨첼트의 소거이론 논문에 관하여.\\
J. Ber. d. DMV 30 (1921), 101쪽
\item 환영역의 아이디얼 이론.\\
Math. Ann. 83 (1921), 24--66쪽
\item 절대기약성의 한 대수적 판정기준.\\
Math. Ann. 85 (1922), 26--33쪽
\item 형식적 변분법과 미분불변식.\\
Encyklopädie d. math. Wiss. III, 3 (1922), 68--71쪽(R. Weitzenböck, Differentialinvarianten 안에 수록)
\item K. 헨첼트 \(\dagger\)의 원고 편집: 다항식 아이디얼과 종결식 이론.\\
Math. Ann. 88 (1923), 53--79쪽
\item 대수불변식과 미분불변식.\\
J. Ber. d. DMV 32 (1923), 177--184쪽
\item 소거이론과 일반 아이디얼 이론.\\
Math. Ann. 90 (1923), 229--261쪽
\item 소거이론과 아이디얼 이론.\\
J. Ber. d. DMV 33 (1924), 116--120쪽
\item 대수적 수체에서의 아이디얼 이론의 추상적 구성.\\
J. Ber. d. DMV 33 (1924), 102쪽
\item 아이디얼 이론의 힐베르트 수.\\
J. Ber. d. DMV 34 (1925), 101쪽
\item 군 지표와 아이디얼 이론.\\
J. Ber. d. DMV 34 (1925), 144쪽
\item 표수 \(p\)인 유한 선형군의 불변식에 관한 유한성 정리.\\
Nachr. v. d. Ges. d. Wiss. zu Göttingen 1926, 28--35쪽
\item 대수적 수체와 함수체에서 아이디얼 이론의 추상적 구성.\\
Math. Ann. 96 (1927), 26--61쪽
\item 대수적 수체 또는 함수체의 오더에 대한 판별식 정리.\\
J. Reine Angew. Math. 157 (1927), 82--104쪽
\item R. 브라우어와 공동: 기약표현의 최소 분해체에 관하여.\\
Sitz. Ber. d. Preuß. Akad. d. Wiss. 1927, 221--228쪽
\item 산술적 관점에서 본 초복소량과 표현론.\\
Atti Congresso Bologna 2 (1928), 71--73쪽
\item[34.] 초복소량과 표현론.\\
Math. Z. 30 (1929), 641--692쪽
\item[35.] 정수적 함수로 이루어진 극대영역에 관하여.\\
Rec. Soc. Math. Moscou 36 (1929), 65--72쪽
\item[36.] 아이디얼 미분과 디퍼런트.\\
J. Ber. d. DMV 39 (1930), 17쪽
\item[37.] 고차 분기가 없는 체의 정규기저.\\
J. Reine Angew. Math. 167 (1932), 147--152쪽
\item[38.] R. 브라우어 및 H. 하세와 공동: 대수론의 한 주정리의 증명.\\
J. Reine Angew. Math. 167 (1932), 399--404쪽
\item[39.] 초복소계와 가환대수 및 수론의 관계.\\
Verhandl. Intern. Math.-Kongreß Zürich 1 (1932), 189--194쪽
\item[40.] 비가환 대수.\\
Math. Z. 37 (1933), 514--541쪽
\item[41.] 상대 갈루아 수체에 대한 주종정리.\\
Math. Ann. 108 (1933), 411--419쪽
\item[42.] 분할 교차곱과 그 극대 오더.\\
Actualités scientifiques et industrielles 148 (1934), 15쪽
\item[43.] 아이디얼 미분과 디퍼런트.\\
J. Reine Angew. Math. 188 (1950), 1--21쪽
\end{enumerate}

\noindent H. 카프페러, 뇌터의 대수함수 기본정리에 대한 필요충분 중복도 조건\\
(E. 뇌터와 공동으로 쓴 부록 포함).\\
Math. Ann. 97 (1927), 559--567쪽

\clearpage
\begin{center}
{\Large\bfseries 에미 뇌터의 단신 및 서평 목록}
\end{center}

\begin{center}
\emph{독일수학회 연보에 실린 단신}
\end{center}

\begingroup
\setlength{\parindent}{0pt}
\setlength{\parskip}{0.15em}
임의의 미분식의 불변식에 관하여, 27, 제II부 (1918), 28쪽\par
불변 변분문제, 27, 제II부 (1918), 47쪽\par
이원형식의 정수 계수 불변식의 유한성, 28, 제II부 (1919), 29쪽\par
정수 계수 다항식과 거듭제곱급수에 관하여, 28, 제II부 (1919), 29--30쪽\par
가군론의 분해정리, 29, 제II부 (1920), 54쪽\par
기본약수와 일반 아이디얼 이론, 30, 제II부 (1921), 32쪽\par
립시츠의 한 공식에 대한 지적, 30, 제II부 (1921), 48쪽\par
일반 아이디얼 이론에서 힐베르트 정리의 유사물, 32, 제II부 (1923), 20--21쪽\par
공리화 문제에 대한 군론적 논평, 32, 제II부 (1923), 21쪽\par
임의의 체에서의 갈루아 이론, 33, 제II부 (1925), 120쪽\par
아이디얼 이론의 힐베르트 수, 34, 제II부 (1926), 101쪽\par
아이디얼의 미분몫과 분기이론, 38, 제II부 (1929), 81쪽\par
\endgroup

\begin{center}
\emph{독일수학회 연보에 실린 서평}
\end{center}

\begingroup
\setlength{\parindent}{0pt}
\setlength{\parskip}{0.35em}
E. Study, \emph{Einleitung in die Theorie der Invarianten linearer Transformationen auf Grund der Vektorenrechnung}. 제1부. (Die Wissenschaft, Bd. 71) Braunschweig 1923, Vieweg.\\
36, 제II부 (1927), 71쪽\par
\emph{Leopold Kroneckers Werke}, 제4권. K. Hensel 편집. VI 및 509쪽. Leipzig 및 Berlin 1929, B. G. Teubner.\\
40, 제II부 (1931), 11--12쪽\par
\emph{Leopold Kroneckers Werke}, 제5권. K. Hensel 편집. 527쪽. Leipzig 및 Berlin 1930, B. G. Teubner.\\
41, 제II부 (1932), 17쪽\par
\emph{Leopold Kroneckers Werke}, 제III.2권. K. Hensel 편집. 215쪽. Leipzig 및 Berlin 1931, B. G. Teubner.\\
42, 제II부 (1933), 38--39쪽\par
D. E. Rutherford, \emph{Modular Invariants}. Cambridge Tracts No. 27. Cambridge 1932, VII 및 84쪽.\\
43, 제II부 (1934), 94--95쪽\par
\endgroup

\clearpage
\begin{center}
\emph{에미 뇌터가 편집에 참여한 책}
\end{center}

\begingroup
\setlength{\parindent}{0pt}
\setlength{\parskip}{0.5em}
R. Dedekind, \emph{Gesammelte mathematische Werke} I--III. R. Fricke, Emmy Noether 및 Øystein Ore 편집. Vieweg, Braunschweig 1930--1932\\
(제I권: 1930, 제II권: 1931, 제III권: 1932)\par
G. Cantor, \emph{Briefwechsel mit R. Dedekind}. Emmy Noether 및 J. Cavaillès 편집, 60쪽. Actualités scient. et ind. 581, Paris 1937\par
\endgroup

\end{document}
